% Auto-generated from data/segments.json + layout.json (+ transforms) — do not edit by hand.
% volume: 26Khu09
% mode: printing
% segments: 2165
% transform_rules: 8
% toc_marks: 268

% Volume layout defaults (layout.json ``layout``).
\csromanlayoutapply{1}{1.5}{21.6pt}{8pt}{8pt}{65pt}{2.5em}%

\setcsromanpitaka{สุตฺตนฺตปิฏก}
\setcsromannikaya{ขุทฺทกนิกาย}
\csromanheader{2}{\textbf{ขุทฺทกนิกาย}}
\csromanmatikaanchor{mtk.0}
\csromantocmarknopage{chapter}{ปฏิสมฺภิทามคฺคปาฬิ}{mtk.0}
\bookmark[dest={mtk.0},level=0]{ปฏิสมฺภิทามคฺคปาฬิ}
\gambhira{ปฏิสมฺภิทามคฺค\-ปาฬิ}
\namakkaram{นโม ตสฺส ภควโต อรหโต สมฺมา\-สมฺพุทฺธสฺส.}
\csromanmatikaanchor{mtk.1}
\csromantocmarknopage{chapter}{1. มหาวคฺค}{mtk.1}
\bookmark[dest={mtk.1},level=0]{1. มหาวคฺค}
\chapterhead{1. \textbf{มหาวคฺค}}
\csromanheader{2}{\textbf{มาติกา}}
\proseitem{1}{\csromanfolio{1}โสตาวธาเน ปญฺญา สุตมเย ญาณํ.}

\proseitem{2}{สุตฺวาน สํวเร ปญฺญา สีลมเย ญาณํ.}

\proseitem{3}{สํวริตฺวา สมาทหเน ปญฺญา สมาธิ\-ภาวนามเย ญาณํ.}

\proseitem{4}{ปจฺจยปริคฺคเห ปญฺญา ธมฺมฏฺฐิติ\-ญาณํ. 5. อตีตา\-นาคต\-ปจฺจุปฺปนฺนานํ ธมฺมานํ สงฺขิปิตฺวา ววตฺถาเน ปญฺญา สมฺมสเน ญาณํ. 6. ปจฺจุปฺปนฺนานํ ธมฺมานํ วิปริณามา\-นุปสฺสเน ปญฺญา อุทยพฺพยา\-นุปสฺสเน ญาณํ.}

\proseitem{7}{อารมฺมณํ ปฏิสงฺขา ภงฺคานุปสฺสเน ปญฺญา วิปสฺสเน ญาณํ.}

\proseitem{8}{ภยตุปฏฺฐาเน ปญฺญา อาทีนเว ญาณํ. 9. มุญฺจิตุกมฺยตา\-ปฏิสงฺขา\-สนฺติฏฺฐนา ปญฺญา สงฺขารุ\-เปกฺขาสุ ญาณํ. 10. พหิทฺธา วุฏฺฐาน\-วิวฏฺฏเน ปญฺญา โคตฺรภุญาณํ. 11. ทุภโต วุฏฺฐาน\-วิวฏฺฏเน ปญฺญา มคฺเค ญาณํ. 12. ปโยค\-ปฺปฏิปฺปสฺสทฺธิ\-ปญฺญา ผเล ญาณํ.}

\proseitem{13}{\csromanfolio{2}ฉินฺนวฏุมา\-นุปสฺสเน\csromansharedfootnote{b8d26f9c5bc8b96e}{ฉินฺนมนุปสฺสเน (สฺยา), ฉินฺนวฏฺฏมนุปสฺสเน (สีฏฺฐ)} ปญฺญา วิมุตฺติญาณํ. 14. ตทา สมุทาคเต\csromansharedfootnote{64f994fce2bd7fd1}{สมุปาคเต (สฺยา)} ธมฺเม ปสฺสเน ปญฺญา ปจฺจเวกฺขเณ ญาณํ. 15. อชฺฌตฺต\-ววตฺถาเน ปญฺญา วตฺถุนานตฺเต ญาณํ. 16. พหิทฺธา\-ววตฺถาเน ปญฺญา โคจรนานตฺเต ญาณํ. 17. จริยา\-ววตฺถาเน ปญฺญา จริยานานตฺเต ญาณํ. 18. จตุธมฺม\-ววตฺถาเน ปญฺญา ภูมินานตฺเต ญาณํ. 19. นวธมฺม\-ววตฺถาเน ปญฺญา ธมฺมนานตฺเต ญาณํ. 20. อภิญฺญาปญฺญา ญาตฏฺเฐ ญาณํ. 21. ปริญฺญาปญฺญา ตีรณฏฺเฐ ญาณํ. 22. ปหาเน ปญฺญา ปริจฺจาคฏฺเฐ ญาณํ. 23. ภาวนาปญฺญา เอกรสฏฺเฐ ญาณํ. 24. สจฺฉิกิริยา\-ปญฺญา ผสฺสนฏฺเฐ\csromansharedfootnote{5db1d456bd87b57c}{ผุสนฏฺเฐ (ก)} ญาณํ. 25. อตฺถนานตฺเต ปญฺญา อตฺถ\-ปฏิสมฺภิเท ญาณํ. 26. ธมฺมนานตฺเต ปญฺญา ธมฺม\-ปฏิสมฺภิเท ญาณํ. 27. นิรุตฺตินานตฺเต ปญฺญา นิรุตฺติ\-ปฏิสมฺภิเท ญาณํ. 28. ปฏิภาน\-นานตฺเต\csromansharedfootnote{e02fadba60fb0261}{ปฏิภาณนานตฺเต (สฺยา)} ปญฺญา ปฏิภาน\-ปฏิสมฺภิเท ญาณํ. 29. วิหารนานตฺเต ปญฺญา วิหารฏฺเฐ ญาณํ. 30. สมาปตฺติ\-นานตฺเต ปญฺญา สมาปตฺตฏฺเฐ ญาณํ. 31. วิหารสมาปตฺติ\-นานตฺเต ปญฺญา วิหารสมาปตฺต\-ฏฺเฐ ญาณํ. 32. อวิกฺเขป\-ปริสุทฺธตฺตา อาสว\-สมุจฺเฉเท ปญฺญา อานนฺตริก\-สมาธิมฺหิ ญาณํ. 33. ทสฺสนา\-ธิปเตยฺยํ สนฺโต จ วิหาราธิคโม ปณีตา\-ธิมุตฺตตา ปญฺญา อรณวิหาเร ญาณํ.}

\csromanheader{2}{มาติกา}
\proseitem{34}{\csromanfolio{3}ทฺวีหิ พเลหิ สมนฺนาคตตฺตา ตโย จ สงฺขารานํ ปฏิปฺปสฺสทฺธิยา โสฬสหิ ญาณจริยาหิ นวหิ สมาธิจริยาหิ วสิภาวตา ปญฺญา นิโรธ\-สมาปตฺติยา ญาณํ. 35. สมฺปชานสฺส ปวตฺต\-ปริยาทาเน ปญฺญา ปรินิพฺพาเน ญาณํ. 36. สพฺพธมฺมานํ สมฺมา สมุจฺเฉเท นิโรเธ จ อนุปฏฺฐานตา ปญฺญา สมสีสฏฺเฐ ญาณํ. 37. ปุถุนานตฺต\-เตช\-ปริยาทาเน ปญฺญา สลฺเลขฏฺเฐ ญาณํ. 38. อสลฺลีน\-ตฺต\-ปหิตตฺต\-ปคฺคหฏฺเฐ ปญฺญา วีริยารมฺเภ ญาณํ. 39. นานา\-ธมฺมปฺปกาสนตา ปญฺญา อตฺถ\-สนฺทสฺสเน ญาณํ. 40. สพฺพธมฺมานํ เอก\-สงฺค\-หตา\-นานตฺเต\-กตฺต\-ปฏิเวเธ ปญฺญา ทสฺสน\-วิสุทฺธิญาณํ. 41. วิทิตตฺตา ปญฺญา ขนฺติญาณํ. 42. ผุฏฺฐตฺตา ปญฺญา ปริโยคาหเณ\csromansharedfootnote{33a9563229fa2701}{ปริโยคาหเน (สฺยา)} ญาณํ. 43. สโมทหเน ปญฺญา ปเทสวิหาเร ญาณํ. 44. อธิปตตฺตา ปญฺญา สญฺญาวิวฏฺเฏ ญาณํ. 45. นานตฺเต ปญฺญา เจโตวิวฏฺเฏ ญาณํ. 46. อธิฏฺฐาเน ปญฺญา จิตฺตวิวฏฺเฏ ญาณํ. 47. สุญฺญเต ปญฺญา ญาณวิวฏฺเฏ ญาณํ. 48. โวสคฺเค\csromansharedfootnote{972be7100b1dc6d6}{โวสฺสคฺเค (พหูสุ)} ปญฺญา วิโมกฺข\-วิวฏฺเฏ ญาณํ. 49. ตถฏฺเฐ ปญฺญา สจฺจวิวฏฺเฏ ญาณํ. 50. กายมฺปิ จิตฺตมฺปิ เอก\-ววตฺถานตา สุขสญฺญญฺจ ลหุสญฺญญฺจ อธิฏฺฐาน\-วเสน อิชฺฌนฏฺเฐ ปญฺญา อิทฺธิวิเธ ญาณํ. 51. วิตกฺก\-วิปฺผาร\-วเสน นานตฺเต\-กตฺต\-สทฺทนิมิตฺตานํ ปริโยคาหเณ ปญฺญา โสตธาตุวิสุทฺธิญฺญาณํ.}

\proseitem{52}{\csromanfolio{4}ติณฺณนฺนํ จิตฺตานํ วิปฺผารตฺตา อินฺทฺริยานํ ปสาทวเสน นานตฺเต\-กตฺต\-วิญฺญาณจริยา ปริโยคาหเณ ปญฺญา เจโตปริย\-ญาณํ. 53. ปจฺจย\-ปฺปวตฺตานํ ธมฺมานํ นานตฺเต\-กตฺต\-กมฺม\-วิปฺผาร\-วเสน ปริโยคาหเณ ปญฺญา ปุพฺเพ\-นิวาสา\-นุสฺสติ\-ญาณํ. 54. โอภาสวเสน นานตฺเต\-กตฺต\-รูปนิมิตฺตานํ ทสฺสนฏฺเฐ ปญฺญา ทิพฺพจกฺขุ\-ญาณํ. 55. จตุสฏฺฐิยา อากาเรหิ ติณฺณนฺนํ อินฺทฺริยานํ วสิภาวตา ปญฺญา อาสวานํ ขเย ญาณํ. 56. ปริญฺญฏฺเฐ ปญฺญา ทุกฺเข ญาณํ. 57. ปหานฏฺเฐ ปญฺญา สมุทเย ญาณํ. 58. สจฺฉิกิริยฏฺเฐ ปญฺญา นิโรเธ ญาณํ. 59. ภาวนฏฺเฐ ปญฺญา มคฺเค ญาณํ. 60. ทุกฺเข ญาณํ. 61. ทุกฺขสมุทเย ญาณํ. 62. ทุกฺขนิโรเธ ญาณํ. 63. ทุกฺขนิโรธ\-คามินิยา ปฏิปทาย ญาณํ. 64. อตฺถ\-ปฏิสมฺภิเท ญาณํ. 65. ธมฺม\-ปฏิสมฺภิเท ญาณํ. 66. นิรุตฺติ\-ปฏิสมฺภิเท ญาณํ. 67. ปฏิภาน\-ปฏิสมฺภิเท ญาณํ. 68. อินฺทฺริยปโรปริยตฺต\-ญาณํ. 69. สตฺตานํ อาสยานุสเย ญาณํ. 70. ยมกปาฏิหีเร ญาณํ. 71. มหากรุณา\-สมาปตฺติยา ญาณํ. 72. สพฺพญฺญุต\-ญฺญาณํ.}

\proseitem{73}{\csromanfolio{5}อนาวรณญาณํ.}

\prosewithcloser{อิมานิ เตสตฺตติ ญาณานิ, อิเมสํ เตสตฺตติยา ญาณานํ สตฺตสฏฺฐิ ญาณานิ สาวก\-สาธารณานิ, ฉ ญาณานิ อสาธารณานิ สาวเกหิ.}{มาติกา นิฏฺฐิตา.}
\csromanmatikaanchor{mtk.2}
\csromanmatikaanchor{mtk.3}
\csromantocmarknopage{section}{1. ญาณกถา}{mtk.2}
\bookmark[dest={mtk.2},level=1]{1. ญาณกถา}
\csromantocmarkref{subsection}{1. สุตมยญาณนิทฺเทส}{mtk.3}
\bookmark[dest={mtk.3},level=2]{1. สุตมยญาณนิทฺเทส}
\csromanheader{2}{1. สุตมย\-ญาณ\-นิทฺเทส}
\prose{1. กถํ โสตาวธาเน ปญฺญา สุตมเย ญาณํ\csromandash{}}

\prose{“อิเม ธมฺมา อภิญฺเญยฺยา”ติ โสตาวธานํ, ตํปชานนา ปญฺญา สุตมเย ญาณํ. (1)}

\prose{“อิเม ธมฺมา ปริญฺเญยฺยา”ติ โสตาวธานํ, ตํปชานนา ปญฺญา สุตมเย ญาณํ. (2)}

\prose{“อิเม ธมฺมา ปหาตพฺพา”ติ โสตาวธานํ, ตํปชานนา ปญฺญา สุตมเย ญาณํ. (3)}

\prose{“อิเม ธมฺมา ภาเวตพฺพา”ติ โสตาวธานํ, ตํปชานนา ปญฺญา สุตมเย ญาณํ. (4)}

\prose{“อิเม ธมฺมา สจฺฉิกาตพฺพา”ติ โสตาวธานํ, ตํปชานนา ปญฺญา สุตมเย ญาณํ. (5)}

\prose{“อิเม ธมฺมา หานภาคิยา”ติ โสตาวธานํ, ตํปชานนา ปญฺญา สุตมเย ญาณํ. (6)}

\prose{“อิเม ธมฺมา ฐิติภาคิยา”ติ โสตาวธานํ, ตํปชานนา ปญฺญา สุตมเย ญาณํ. (7)}

\prose{“อิเม ธมฺมา วิเสสภาคิยา”ติ โสตาวธานํ, ตํปชานนา ปญฺญา สุตมเย ญาณํ. (8)}

\prose{“อิเม ธมฺมา นิพฺเพธ\-ภาคิยา”ติ โสตาวธานํ, ตํปชานนา ปญฺญา สุตมเย ญาณํ. (9)}

\prose{\csromanfolio{6}“สพฺเพ สงฺขารา อนิจฺจา”ติ โสตาวธานํ, ตํปชานนา ปญฺญา สุตมเย ญาณํ. (10)}

\prose{“สพฺเพ สงฺขารา ทุกฺขา”ติ โสตาวธานํ, ตํปชานนา ปญฺญา สุตมเย ญาณํ. (11)}

\prose{“สพฺเพ ธมฺมา อนตฺตา”ติ โสตาวธานํ, ตํปชานนา ปญฺญา สุตมเย ญาณํ. (12)}

\prose{“อิทํ ทุกฺขํ อริยสจฺจนฺ”ติ โสตาวธานํ, ตํปชานนา ปญฺญา สุตมเย ญาณํ. (13)}

\prose{“อิทํ ทุกฺข\-สมุทยํ\csromansharedfootnote{7268e5b3db3909f2}{ทุกฺขสมุทโย (สฺยา)} อริยสจฺจนฺ”ติ โสตาวธานํ, ตํปชานนา ปญฺญา สุตมเย ญาณํ. (14)}

\prose{“อิทํ ทุกฺขนิโรธํ\csromansharedfootnote{2d36811a5c085b7c}{ทุกฺขนิโรโธ (สฺยา) อฏฺฐกถา โอโลเกตพฺพา.} อริยสจฺจนฺ”ติ โสตาวธานํ, ตํปชานนา ปญฺญา สุตมเย ญาณํ. (15)}

\prose{“อิทํ ทุกฺขนิโรธ\-คามินี ปฏิปทา อริยสจฺจนฺ”ติ โสตาวธานํ, ตํปชานนา ปญฺญา สุตมเย ญาณํ. (16)}

\proseitem{2}{กถํ “อิเม ธมฺมา อภิญฺเญยฺยา”ติ โสตาวธานํ, ตํปชานนา ปญฺญา สุตมเย ญาณํ\csromandash{}}

\prose{เอโก ธมฺโม อภิญฺเญยฺโย, สพฺเพ สตฺตา อาหารฏฺฐิติกา. เทฺว ธมฺมา อภิญฺเญยฺยา, เทฺว ธาตุโย. ตโย ธมฺมา อภิญฺเญยฺยา, ติสฺโส ธาตุโย. จตฺตาโร ธมฺมา อภิญฺเญยฺยา, จตฺตาริ อริยสจฺจานิ. ปญฺจ ธมฺมา อภิญฺเญยฺยา, ปญฺจ วิมุตฺตา\-ยตนานิ. ฉ ธมฺมา อภิญฺเญยฺยา, ฉ อนุตฺตริยานิ. สตฺต ธมฺมา อภิญฺเญยฺยา, สตฺต นิทฺทส\-วตฺถูนิ. อฏฺฐ ธมฺมา อภิญฺเญยฺยา, อฏฺฐ อภิภา\-ยตนานิ. นว ธมฺมา อภิญฺเญยฺยา, นว อนุปุพฺพวิหารา. ทส ธมฺมา อภิญฺเญยฺยา, ทส นิชฺชร\-วตฺถูนิ. (10)}

\proseitem{3}{สพฺพํ ภิกฺขเว อภิญฺเญยฺยํ. กิญฺจ ภิกฺขเว สพฺพํ อภิญฺเญยฺยํ? จกฺขุ\csromansharedfootnote{b15203dbe0312afb}{จกฺขุํ (สฺยา, ก)} ภิกฺขเว อภิญฺเญยฺยํ, รูปา อภิญฺเญยฺยา, จกฺขุวิญฺญาณํ อภิญฺเญยฺยํ, จกฺขุ\-สมฺผสฺโส อภิญฺเญยฺโย, ยมฺปิทํ\csromansharedfootnote{8f8ac9f5ceb5d6e1}{ยมิทํ (ก) สํ 2. 258 ปิฏฺเฐ ปสฺสิตพฺพา.} จกฺขุ\-สมฺผสฺส\-ปจฺจยา อุปฺปชฺชติ เวทยิตํ สุขํ วา ทุกฺขํ วา อทุกฺขมสุขํ วา, ตมฺปิ อภิญฺเญยฺยํ. โสตํ อภิญฺเญยฺยํ, \csromanfolio{7}สทฺทา อภิญฺเญยฺยา. ฆานํ อภิญฺเญยฺยํ, คนฺธา อภิญฺเญยฺยา. ชิวฺหา อภิญฺเญยฺยา, รสา อภิญฺเญยฺยา. กาโย อภิญฺเญยฺโย, โผฏฺฐพฺพา อภิญฺเญยฺยา. มโน อภิญฺเญยฺโย, ธมฺมา อภิญฺเญยฺยา, มโนวิญฺญาณํ อภิญฺเญยฺยํ, มโนสมฺผสฺโส อภิญฺเญยฺโย, ยมฺปิทํ มโน\-สมฺผสฺส\-ปจฺจยา อุปฺปชฺชติ เวทยิตํ สุขํ วา ทุกฺขํ วา อทุกฺขมสุขํ วา, ตมฺปิ อภิญฺเญยฺยํ. (30)}

\prose{รูปํ อภิญฺเญยฺยํ, เวทนา อภิญฺเญยฺยา, สญฺญา อภิญฺเญยฺยา, สงฺขารา อภิญฺเญยฺยา, วิญฺญาณํ อภิญฺเญยฺยํ. (5)}

\prose{จกฺขุ อภิญฺเญยฺยํ, โสตํ อภิญฺเญยฺยํ, ฆานํ อภิญฺเญยฺยํ, ชิวฺหา อภิญฺเญยฺยา, กาโย อภิญฺเญยฺโย, มโน อภิญฺเญยฺโย.  รูปา อภิญฺเญยฺยา, สทฺทา อภิญฺเญยฺยา, คนฺธา อภิญฺเญยฺยา, รสา อภิญฺเญยฺยา, โผฏฺฐพฺพา อภิญฺเญยฺยา, ธมฺมา อภิญฺเญยฺยา.  จกฺขุวิญฺญาณํ อภิญฺเญยฺยํ, โสตวิญฺญาณํ อภิญฺเญยฺยํ, ฆานวิญฺญาณํ อภิญฺเญยฺยํ, ชิวฺหาวิญฺญาณํ อภิญฺเญยฺยํ, กายวิญฺญาณํ อภิญฺเญยฺยํ, มโนวิญฺญาณํ อภิญฺเญยฺยํ.  จกฺขุ\-สมฺผสฺโส อภิญฺเญยฺโย, โสตสมฺผสฺโส อภิญฺเญยฺโย, ฆานสมฺผสฺโส อภิญฺเญยฺโย, ชิวฺหาสมฺผสฺโส อภิญฺเญยฺโย, กายสมฺผสฺโส อภิญฺเญยฺโย, มโนสมฺผสฺโส อภิญฺเญยฺโย.  จกฺขุ\-สมฺผสฺสชา เวทนา อภิญฺเญยฺยา, โสต\-สมฺผสฺสชา เวทนา อภิญฺเญยฺยา, ฆาน\-สมฺผสฺสชา เวทนา อภิญฺเญยฺยา, ชิวฺหา\-สมฺผสฺสชา เวทนา อภิญฺเญยฺยา, กาย\-สมฺผสฺสชา เวทนา อภิญฺเญยฺยา, มโน\-สมฺผสฺสชา เวทนา อภิญฺเญยฺยา.  รูปสญฺญา อภิญฺเญยฺยา, สทฺทสญฺญา อภิญฺเญยฺยา, คนฺธสญฺญา อภิญฺเญยฺยา, รสสญฺญา อภิญฺเญยฺยา, โผฏฺฐพฺพ\-สญฺญา อภิญฺเญยฺยา, ธมฺมสญฺญา อภิญฺเญยฺยา.  รูปสญฺเจตนา อภิญฺเญยฺยา, สทฺทสญฺเจตนา อภิญฺเญยฺยา, คนฺธ\-สญฺเจตนา อภิญฺเญยฺยา, รสสญฺเจตนา อภิญฺเญยฺยา, โผฏฺฐพฺพ\-สญฺเจตนา อภิญฺเญยฺยา, ธมฺม\-สญฺเจตนา อภิญฺเญยฺยา.  รูปตณฺหา อภิญฺเญยฺยา, สทฺทตณฺหา อภิญฺเญยฺยา, คนฺธตณฺหา อภิญฺเญยฺยา, รสตณฺหา อภิญฺเญยฺยา, โผฏฺฐพฺพ\-ตณฺหา อภิญฺเญยฺยา, ธมฺมตณฺหา อภิญฺเญยฺยา.  รูปวิตกฺโก อภิญฺเญยฺโย, สทฺทวิตกฺโก อภิญฺเญยฺโย, คนฺธวิตกฺโก อภิญฺเญยฺโย, รสวิตกฺโก อภิญฺเญยฺโย, โผฏฺฐพฺพ\-วิตกฺโก อภิญฺเญยฺโย, ธมฺมวิตกฺโก อภิญฺเญยฺโย.  รูปวิจาโร อภิญฺเญยฺโย, สทฺทวิจาโร \csromanfolio{8}อภิญฺเญยฺโย, คนฺธวิจาโร อภิญฺเญยฺโย, รสวิจาโร อภิญฺเญยฺโย, โผฏฺฐพฺพ\-วิจาโร อภิญฺเญยฺโย, ธมฺมวิจาโร อภิญฺเญยฺโย. (60)}

\prose{ปถวีธาตุ อภิญฺเญยฺยา, อาโปธาตุ อภิญฺเญยฺยา, เตโชธาตุ อภิญฺเญยฺยา, วาโยธาตุ อภิญฺเญยฺยา, อากาสธาตุ อภิญฺเญยฺยา, วิญฺญาณธาตุ อภิญฺเญยฺยา. (6)}

\prose{ปถวีกสิณํ อภิญฺเญยฺยํ, อาโปกสิณํ อภิญฺเญยฺยํ, เตโชกสิณํ อภิญฺเญยฺยํ, วาโยกสิณํ อภิญฺเญยฺยํ, นีลกสิณํ อภิญฺเญยฺยํ, ปีตกสิณํ อภิญฺเญยฺยํ, โลหิตกสิณํ อภิญฺเญยฺยํ, โอทาตกสิณํ อภิญฺเญยฺยํ, อากาสกสิณํ อภิญฺเญยฺยํ, วิญฺญาณกสิณํ อภิญฺเญยฺยํ. (10)}

\prose{เกสา อภิญฺเญยฺยา, โลมา อภิญฺเญยฺยา, นขา อภิญฺเญยฺยา, ทนฺตา อภิญฺเญยฺยา, ตโจ อภิญฺเญยฺโย, มํสํ อภิญฺเญยฺยํ, นฺหารู\csromansharedfootnote{62ab890a51680ea4}{นหารู (สฺยา)} อภิญฺเญยฺยา, อฏฺฐี อภิญฺเญยฺยา, อฏฺฐิมิญฺชา อภิญฺเญยฺยา\csromansharedfootnote{2aa030006f6f6f0e}{อฏฺฐิมญฺชํ อภิญฺเญยฺยํ (สฺยา, ก)}, วกฺกํ อภิญฺเญยฺยํ, หทยํ อภิญฺเญยฺยํ, ยกนํ อภิญฺเญยฺยํ, กิโลมกํ อภิญฺเญยฺยํ, ปิหกํ อภิญฺเญยฺยํ, ปปฺผาสํ อภิญฺเญยฺยํ, อนฺตํ อภิญฺเญยฺยํ, อนฺตคุณํ อภิญฺเญยฺยํ, อุทริยํ อภิญฺเญยฺยํ, กรีสํ อภิญฺเญยฺยํ, ปิตฺตํ อภิญฺเญยฺยํ, เสมฺหํ อภิญฺเญยฺยํ, ปุพฺโพ อภิญฺเญยฺโย, โลหิตํ อภิญฺเญยฺยํ, เสโท อภิญฺเญยฺโย, เมโท อภิญฺเญยฺโย, อสฺสุ อภิญฺเญยฺยํ, วสา อภิญฺเญยฺยา, เขโฬ อภิญฺเญยฺโย, สิงฺฆาณิกา อภิญฺเญยฺยา, ลสิกา อภิญฺเญยฺยา, มุตฺตํ อภิญฺเญยฺยํ, มตฺถลุงฺคํ อภิญฺเญยฺยํ. (32)}

\prose{จกฺขายตนํ อภิญฺเญยฺยํ, รูปายตนํ อภิญฺเญยฺยํ. โสตายตนํ อภิญฺเญยฺยํ, สทฺทายตนํ อภิญฺเญยฺยํ. ฆานายตนํ อภิญฺเญยฺยํ, คนฺธายตนํ อภิญฺเญยฺยํ. ชิวฺหายตนํ อภิญฺเญยฺยํ, รสายตนํ อภิญฺเญยฺยํ.กายายตนํ อภิญฺเญยฺยํ, โผฏฺฐพฺพา\-ยตนํ อภิญฺเญยฺยํ. มนายตนํ อภิญฺเญยฺยํ, ธมฺมายตนํ อภิญฺเญยฺยํ. (12)}

\prose{จกฺขุธาตุ อภิญฺเญยฺยา, รูปธาตุ อภิญฺเญยฺยา, จกฺขุวิญฺญาณ\-ธาตุ อภิญฺเญยฺยา. โสตธาตุ อภิญฺเญยฺยา, สทฺทธาตุ อภิญฺเญยฺยา, โสตวิญฺญาณ\-ธาตุ อภิญฺเญยฺยา. ฆานธาตุ อภิญฺเญยฺยา, คนฺธธาตุ อภิญฺเญยฺยา, ฆานวิญฺญาณ\-ธาตุ อภิญฺเญยฺยา. ชิวฺหาธาตุ อภิญฺเญยฺยา, \csromanfolio{9}รสธาตุ อภิญฺเญยฺยา, ชิวฺหาวิญฺญาณ\-ธาตุ อภิญฺเญยฺยา. กายธาตุ อภิญฺเญยฺยา, โผฏฺฐพฺพ\-ธาตุ อภิญฺเญยฺยา, กายวิญฺญาณ\-ธาตุ อภิญฺเญยฺยา. มโนธาตุ อภิญฺเญยฺยา, ธมฺมธาตุ อภิญฺเญยฺยา, มโนวิญฺญาณ\-ธาตุ อภิญฺเญยฺยา. (18)}

\proseitem{4}{จกฺขุนฺทฺริยํ อภิญฺเญยฺยํ, โสตินฺทฺริยํ อภิญฺเญยฺยํ, ฆานินฺทฺริยํ อภิญฺเญยฺยํ, ชิวฺหินฺทฺริยํ อภิญฺเญยฺยํ, กายินฺทฺริยํ อภิญฺเญยฺยํ, มนินฺทฺริยํ อภิญฺเญยฺยํ, ชีวิตินฺทฺริยํ อภิญฺเญยฺยํ, อิตฺถินฺทฺริยํ อภิญฺเญยฺยํ, ปุริสินฺทฺริยํ อภิญฺเญยฺยํ, สุขินฺทฺริยํ อภิญฺเญยฺยํ, ทุกฺขินฺทฺริยํ อภิญฺเญยฺยํ, โสมนสฺสิ\-นฺทฺริยํ อภิญฺเญยฺยํ, โทมนสฺสิ\-นฺทฺริยํ อภิญฺเญยฺยํ, อุเปกฺขินฺทฺริยํ อภิญฺเญยฺยํ, สทฺธินฺทฺริยํ อภิญฺเญยฺยํ, วีริยินฺทฺริยํ\csromansharedfootnote{e93cc4052a230ead}{วิริยินฺทฺริยํ (สฺยา)} อภิญฺเญยฺยํ, สตินฺทฺริยํ อภิญฺเญยฺยํ, สมาธินฺทฺริยํ อภิญฺเญยฺยํ, ปญฺญินฺทฺริยํ อภิญฺเญยฺยํ, อนญฺญาต\-ญฺญสฺสามี\-ตินฺทฺริยํ อภิญฺเญยฺยํ, อญฺญินฺทฺริยํ อภิญฺเญยฺยํ, อญฺญาตาวิ\-นฺทฺริยํ อภิญฺเญยฺยํ. (22)}

\prose{กามธาตุ อภิญฺเญยฺยา, รูปธาตุ อภิญฺเญยฺยา, อรูปธาตุ อภิญฺเญยฺยา.  กามภโว อภิญฺเญยฺโย, รูปภโว อภิญฺเญยฺโย, อรูปภโว อภิญฺเญยฺโย.  สญฺญาภโว อภิญฺเญยฺโย, อสญฺญาภโว อภิญฺเญยฺโย, เนว\-สญฺญา\-นา\-สญฺญาภโว อภิญฺเญยฺโย.  เอกโวการภโว อภิญฺเญยฺโย, จตุโวการภโว อภิญฺเญยฺโย, ปญฺจโวการภโว อภิญฺเญยฺโย. (12)}

\prose{ปฐมํ ฌานํ\csromansharedfootnote{39a9e9feb3c33fc1}{ปฐมชฺฌานํ (สฺยา) เอวมีทิเสสุ ฐาเนสุ.} อภิญฺเญยฺยํ, ทุติยํ ฌานํ อภิญฺเญยฺยํ, ตติยํ ฌานํ อภิญฺเญยฺยํ, จตุตฺถํ ฌานํ อภิญฺเญยฺยํ.  เมตฺตา\-เจโตวิมุตฺติ อภิญฺเญยฺยา, กรุณา\-เจโตวิมุตฺติ อภิญฺเญยฺยา, มุทิตา\-เจโตวิมุตฺติ อภิญฺเญยฺยา, อุเปกฺขา\-เจโตวิมุตฺติ อภิญฺเญยฺยา.  อากาสา\-นญฺจา\-ยตนสมาปตฺติ อภิญฺเญยฺยา, วิญฺญาณญฺจา\-ยตนสมาปตฺติ อภิญฺเญยฺยา, อากิญฺจญฺญายตน\-สมาปตฺติ อภิญฺเญยฺยา, เนว\-สญฺญา\-นา\-สญฺญา\-ยตนสมาปตฺติ อภิญฺเญยฺยา. (12)}

\prose{อวิชฺชา อภิญฺเญยฺยา, สงฺขารา อภิญฺเญยฺยา, วิญฺญาณํ อภิญฺเญยฺยํ, นามรูปํ อภิญฺเญยฺยํ, สฬายตนํ อภิญฺเญยฺยํ, ผสฺโส อภิญฺเญยฺโย, เวทนา อภิญฺเญยฺยา, ตณฺหา อภิญฺเญยฺยา, อุปาทานํ อภิญฺเญยฺยํ, ภโว อภิญฺเญยฺโย, ชาติ อภิญฺเญยฺยา, ชรามรณํ อภิญฺเญยฺยํ. (12)}

\proseitem{7}{\csromanfolio{10}ทุกฺขํ อภิญฺเญยฺยํ, ทุกฺขสมุทโย อภิญฺเญยฺโย, ทุกฺขนิโรโธ อภิญฺเญยฺโย, ทุกฺขนิโรธ\-คามินี ปฏิปทา อภิญฺเญยฺยา. รุปํ อภิญฺเญยฺยํ, รูปสมุทโย อภิญฺเญยฺโย, รูปนิโรโธ อภิญฺเญยฺโย, รูปนิโรธ\-คามินี ปฏิปทา อภิญฺเญยฺยา. เวทนา อภิญฺเญยฺยา ฯเปฯ. สญฺญา อภิญฺเญยฺยา ฯเปฯ. สงฺขารา อภิญฺเญยฺยา ฯเปฯ. วิญฺญาณํ อภิญฺเญยฺยํ. จกฺขุ อภิญฺเญยฺยํ ฯเปฯ. ชรามรณํ อภิญฺเญยฺยํ, ชรามรณ\-สมุทโย อภิญฺเญยฺโย, ชรามรณ\-นิโรโธ อภิญฺเญยฺโย, ชรามรณนิโรธ\-คามินี ปฏิปทา อภิญฺเญยฺยา. (808)}

\prose{ทุกฺขสฺส ปริญฺญฏฺโฐ อภิญฺเญยฺโย, ทุกฺข\-สมุทยสฺส ปหานฏฺโฐ อภิญฺเญยฺโย, ทุกฺข\-นิโรธสฺส สจฺฉิกิริยฏฺโฐ อภิญฺเญยฺโย, ทุกฺขนิโรธ\-คามินิยา ปฏิปทาย ภาวนฏฺโฐ อภิญฺเญยฺโย. รูปสฺส ปริญฺญฏฺโฐ อภิญฺเญยฺโย, รูปสมุทยสฺส ปหานฏฺโฐ อภิญฺเญยฺโย, รูปนิโรธสฺส สจฺฉิกิริยฏฺโฐ อภิญฺเญยฺโย, รูปนิโรธ\-คามินิยา ปฏิปทาย ภาวนฏฺโฐ อภิญฺเญยฺโย. เวทนาย ฯเปฯ. สญฺญาย. สงฺขารานํ. วิญฺญาณสฺส. จกฺขุสฺส ฯเปฯ. ชรามรณสฺส ปริญฺญฏฺโฐ อภิญฺเญยฺโย, ชรามรณ\-สมุทยสฺส ปหานฏฺโฐ อภิญฺเญยฺโย, ชรามรณ\-นิโรธสฺส สจฺฉิกิริยฏฺโฐ อภิญฺเญยฺโย, ชรามรณนิโรธ\-คามินิยา ปฏิปทาย ภาวนฏฺโฐ อภิญฺเญยฺโย. (808)}

\prose{ทุกฺขสฺส ปริญฺญา\-ปฏิเวธฏฺโฐ อภิญฺเญยฺโย, ทุกฺข\-สมุทยสฺส ปหาน\-ปฏิเวธฏฺโฐ อภิญฺเญยฺโย, ทุกฺข\-นิโรธสฺส สจฺฉิกิริยา\-ปฏิเวธฏฺโฐ อภิญฺเญยฺโย, ทุกฺขนิโรธ\-คามินิยา ปฏิปทาย ภาวนา\-ปฏิเวธฏฺโฐ อภิญฺเญยฺโย. รูปสฺส ปริญฺญา\-ปฏิเวธฏฺโฐ อภิญฺเญยฺโย, รูปสมุทยสฺส ปหาน\-ปฏิเวธฏฺโฐ อภิญฺเญยฺโย, รูปนิโรธสฺส สจฺฉิกิริยา\-ปฏิเวธฏฺโฐ อภิญฺเญยฺโย, รูปนิโรธ\-คามินิยา ปฏิปทาย ภาวนา\-ปฏิเวธฏฺโฐ อภิญฺเญยฺโย. เวทนาย ฯเปฯ. สญฺญาย. สงฺขารานํ. วิญฺญาณสฺส. จกฺขุสฺส ฯเปฯ. ชรามรณสฺส ปริญฺญา\-ปฏิเวธฏฺโฐ อภิญฺเญยฺโย, ชรามรณ\-สมุทยสฺส ปหาน\-ปฏิเวธฏฺโฐ อภิญฺเญยฺโย, ชรามรณ\-นิโรธสฺส สจฺฉิกิริยา\-ปฏิเวธฏฺโฐ อภิญฺเญยฺโย, ชรามรณนิโรธ\-คามินิยา ปฏิปทาย ภาวนา\-ปฏิเวธฏฺโฐ อภิญฺเญยฺโย. (808)}

\proseitem{8}{ทุกฺขํ อภิญฺเญยฺยํ, ทุกฺขสมุทโย อภิญฺเญยฺโย, ทุกฺขนิโรโธ อภิญฺเญยฺโย, ทุกฺขสฺส สมุทย\-นิโรโธ อภิญฺเญยฺโย, ทุกฺขสฺส ฉนฺทราค\-นิโรโธ อภิญฺเญยฺโย, ทุกฺขสฺส อสฺสาโท อภิญฺเญยฺโย, ทุกฺขสฺส \csromanfolio{11}อาทีนโว อภิญฺเญยฺโย, ทุกฺขสฺส นิสฺสรณํ อภิญฺเญยฺยํ.  รูปํ อภิญฺเญยฺยํ, รูปสมุทโย อภิญฺเญยฺโย, รูปนิโรโธ อภิญฺเญยฺโย, รูปสฺส สมุทย\-นิโรโธ อภิญฺเญยฺโย, รูปสฺส ฉนฺทราค\-นิโรโธ อภิญฺเญยฺโย, รูปสฺส อสฺสาโท อภิญฺเญยฺโย, รูปสฺส อาทีนโว อภิญฺเญยฺโย, รูปสฺส นิสฺสรณํ อภิญฺเญยฺยํ. เวทนา อภิญฺเญยฺยา ฯเปฯ. สญฺญา อภิญฺเญยฺยา. สงฺขารา อภิญฺเญยฺยา. วิญฺญาณํ อภิญฺเญยฺยํ. จกฺขุ อภิญฺเญยฺยํ ฯเปฯ. ชรามรณํ อภิญฺเญยฺยํ, ชรามรณ\-สมุทโย อภิญฺเญยฺโย, ชรามรณ\-นิโรโธ อภิญฺเญยฺโย, ชรามรณสฺส สมุทย\-นิโรโธ อภิญฺเญยฺโย, ชรามรณสฺส ฉนฺทราค\-นิโรโธ อภิญฺเญยฺโย, ชรามรณสฺส อสฺสาโท อภิญฺเญยฺโย, ชรามรณสฺส อาทีนโว อภิญฺเญยฺโย, ชรามรณสฺส นิสฺสรณํ อภิญฺเญยฺยํ. (1616)}

\prose{ทุกฺขํ อภิญฺเญยฺยํ, ทุกฺขสมุทโย อภิญฺเญยฺโย, ทุกฺขนิโรโธ อภิญฺเญยฺโย, ทุกฺขนิโรธ\-คามินี ปฏิปทา อภิญฺเญยฺยา, ทุกฺขสฺส อสฺสาโท อภิญฺเญยฺโย, ทุกฺขสฺส อาทีนโว อภิญฺเญยฺโย, ทุกฺขสฺส นิสฺสรณํ อภิญฺเญยฺยํ.  รูปํ อภิญฺเญยฺยํ, รูปสมุทโย อภิญฺเญยฺโย, รูปนิโรโธ อภิญฺเญยฺโย, รูปนิโรธ\-คามินี ปฏิปทา อภิญฺเญยฺยา, รูปสฺส อสฺสาโท อภิญฺเญยฺโย, รูปสฺส อาทีนโว อภิญฺเญยฺโย, รูปสฺส นิสฺสรณํ อภิญฺเญยฺยํ. เวทนา อภิญฺเญยฺยา ฯเปฯ. สญฺญา อภิญฺเญยฺยา. สงฺขารา อภิญฺเญยฺยา. วิญฺญาณํ อภิญฺเญยฺยํ. จกฺขุ อภิญฺเญยฺยํ ฯเปฯ. ชรามรณํ อภิญฺเญยฺยํ. ชรามรณ\-สมุทโย อภิญฺเญยฺโย, ชรามรณ\-นิโรโธ อภิญฺเญยฺโย, ชรามรณนิโรธ\-คามินี ปฏิปทา อภิญฺเญยฺยา, ชรามรณสฺส อสฺสาโท อภิญฺเญยฺโย, ชรามรณสฺส อาทีนโว อภิญฺเญยฺโย, ชรามรณสฺส นิสฺสรณํ อภิญฺเญยฺยํ. (1414)}

\prose{อนิจฺจา\-นุปสฺสนา อภิญฺเญยฺยา, ทุกฺขา\-นุปสฺสนา อภิญฺเญยฺยา, อนตฺตา\-นุปสฺสนา อภิญฺเญยฺยา, นิพฺพิทา\-นุปสฺสนา อภิญฺเญยฺยา, วิราคา\-นุปสฺสนา อภิญฺเญยฺยา, นิโรธา\-นุปสฺสนา อภิญฺเญยฺยา, ปฏินิสฺสคฺคา\-นุปสฺสนา อภิญฺเญยฺยา.  รูเป อนิจฺจา\-นุปสฺสนา อภิญฺเญยฺยา, รูเป ทุกฺขา\-นุปสฺสนา อภิญฺเญยฺยา, รูเป อนตฺตา\-นุปสฺสนา อภิญฺเญยฺยา, รูเป นิพฺพิทา\-นุปสฺสนา อภิญฺเญยฺยา, รูเป วิราคา\-นุปสฺสนา อภิญฺเญยฺยา, รูเป นิโรธา\-นุปสฺสนา อภิญฺเญยฺยา, รูเป ปฏินิสฺสคฺคา\-นุปสฺสนา อภิญฺเญยฺยา. เวทนาย ฯเปฯ. สญฺญาย. สงฺขาเรสุ. วิญฺญาเณ. จกฺขุสฺมิํ ฯเปฯ. ชรามรเณ \csromanfolio{12}อนิจฺจา\-นุปสฺสนา อภิญฺเญยฺยา, ชรามรเณ ทุกฺขา\-นุปสฺสนา อภิญฺเญยฺยา, ชรามรเณ อนตฺตา\-นุปสฺสนา อภิญฺเญยฺยา, ชรามรเณ นิพฺพิทา\-นุปสฺสนา อภิญฺเญยฺยา, ชรามรเณ วิราคา\-นุปสฺสนา อภิญฺเญยฺยา, ชรามรเณ นิโรธา\-นุปสฺสนา อภิญฺเญยฺยา, ชรามรเณ ปฏินิสฺสคฺคา\-นุปสฺสนา อภิญฺเญยฺยา. (1414)}

\proseitem{9}{อุปฺปาโท อภิญฺเญยฺโย, ปวตฺตํ อภิญฺเญยฺยํ, นิมิตฺตํ อภิญฺเญยฺยํ, อายูหนา\csromansharedfootnote{193ed89303ee7982}{อายุหนา (สฺยา) เอวมุปริปิ.} อภิญฺเญยฺยา, ปฏิสนฺธิ อภิญฺเญยฺยา, คติ อภิญฺเญยฺยา, นิพฺพตฺติ อภิญฺเญยฺยา, อุปปตฺติ อภิญฺเญยฺยา, ชาติ อภิญฺเญยฺยา, ชรา อภิญฺเญยฺยา, พฺยาธิ อภิญฺเญยฺโย, มรณํ อภิญฺเญยฺยํ, โสโก อภิญฺเญยฺโย, ปริเทโว อภิญฺเญยฺโย, อุปายาโส อภิญฺเญยฺโย. (15)}

\prose{อนุปฺปาโท อภิญฺเญยฺโย, อปฺปวตฺตํ อภิญฺเญยฺยํ, อนิมิตฺตํ อภิญฺเญยฺยํ, อนายูหนา\csromansharedfootnote{d94558479fd81312}{อนายุหนา (สฺยา) เอวมุปริปิ.} อภิญฺเญยฺยา, อปฺปฏิสนฺธิ อภิญฺเญยฺยา, อคติ อภิญฺเญยฺยา, อนิพฺพตฺติ อภิญฺเญยฺยา, อนุปปตฺติ อภิญฺเญยฺยา, อชาติ อภิญฺเญยฺยา, อชรา อภิญฺเญยฺยา, อพฺยาธิ อภิญฺเญยฺโย, อมตํ อภิญฺเญยฺยํ, อโสโก อภิญฺเญยฺโย, อปริเทโว อภิญฺเญยฺโย, อนุปายาโส อภิญฺเญยฺโย. (15)}

\prose{อุปฺปาโท อภิญฺเญยฺโย, อนุปฺปาโท อภิญฺเญยฺโย. ปวตฺตํ อภิญฺเญยฺยํ, อปฺปวตฺตํ อภิญฺเญยฺยํ. นิมิตฺตํ อภิญฺเญยฺยํ, อนิมิตฺตํ อภิญฺเญยฺยํ. อายูหนา อภิญฺเญยฺยา, อนายูหนา อภิญฺเญยฺยา. ปฏิสนฺธิ อภิญฺเญยฺยา, อปฺปฏิสนฺธิ อภิญฺเญยฺยา. คติ อภิญฺเญยฺยา, อคติ อภิญฺเญยฺยา. นิพฺพตฺติ อภิญฺเญยฺยา, อนิพฺพตฺติ อภิญฺเญยฺยา. อุปปตฺติ อภิญฺเญยฺยา, อนุปปตฺติ อภิญฺเญยฺยา. ชาติ อภิญฺเญยฺยา, อชาติ อภิญฺเญยฺยา. ชรา อภิญฺเญยฺยา, อชรา อภิญฺเญยฺยา. พฺยาธิ อภิญฺเญยฺโย, อพฺยาธิ อภิญฺเญยฺโย. มรณํ อภิญฺเญยฺยํ, อมตํ อภิญฺเญยฺยํ. โสโก อภิญฺเญยฺโย, อโสโก อภิญฺเญยฺโย. ปริเทโว อภิญฺเญยฺโย, อปริเทโว อภิญฺเญยฺโย. อุปายาโส อภิญฺเญยฺโย, อนุปายาโส อภิญฺเญยฺโย. (30)}

\prose{อุปฺปาโท ทุกฺขนฺติ อภิญฺเญยฺยํ, ปวตฺตํ ทุกฺขนฺติ อภิญฺเญยฺยํ, นิมิตฺตํ ทุกฺขนฺติ อภิญฺเญยฺยํ, อายูหนา ทุกฺขนฺติ อภิญฺเญยฺยํ, ปฏิสนฺธิ ทุกฺขนฺติ อภิญฺเญยฺยํ \csromanfolio{13}คติ ทุกฺขนฺติ อภิญฺเญยฺยํ, นิพฺพตฺติ ทุกฺขนฺติ อภิญฺเญยฺยํ, อุปปตฺติ ทุกฺขนฺติ อภิญฺเญยฺยํ, ชาติ ทุกฺขนฺติ อภิญฺเญยฺยํ, ชรา ทุกฺขนฺติ อภิญฺเญยฺยํ, พฺยาธิ ทุกฺขนฺติ อภิญฺเญยฺยํ, มรณํ ทุกฺขนฺติ อภิญฺเญยฺยํ, โสโก ทุกฺขนฺติ อภิญฺเญยฺยํ, ปริเทโว ทุกฺขนฺติ อภิญฺเญยฺยํ, อุปายาโส ทุกฺขนฺติ อภิญฺเญยฺยํ. (15)}

\proseitem{10}{อนุปฺปาโท สุขนฺติ อภิญฺเญยฺยํ, อปฺปวตฺตํ สุขนฺติ อภิญฺเญยฺยํ, อนิมิตฺตํ สุขนฺติ อภิญฺเญยฺยํ, อนายูหนา สุขนฺติ อภิญฺเญยฺยํ, อปฺปฏิสนฺธิ สุขนฺติ อภิญฺเญยฺยํ, อคติ สุขนฺติ อภิญฺเญยฺยํ, อนิพฺพตฺติ สุขนฺติ อภิญฺเญยฺยํ, อนุปปตฺติ สุขนฺติ อภิญฺเญยฺยํ, อชาติ สุขนฺติ อภิญฺเญยฺยํ, อชรา สุขนฺติ อภิญฺเญยฺยํ, อพฺยาธิ สุขนฺติ อภิญฺเญยฺยํ, อมตํ สุขนฺติ อภิญฺเญยฺยํ, อโสโก สุขนฺติ อภิญฺเญยฺยํ, อปริเทโว สุขนฺติ อภิญฺเญยฺยํ, อนุปายาโส สุขนฺติ อภิญฺเญยฺยํ. (15)}

\prose{อุปฺปาโท ทุกฺขํ, อนุปฺปาโท สุขนฺติ อภิญฺเญยฺยํ, ปวตฺตํ ทุกฺขํ, อปฺปวตฺตํ สุขนฺติ อภิญฺเญยฺยํ, นิมิตฺตํ ทุกฺขํ, อนิมิตฺตํ สุขนฺติ อภิญฺเญยฺยํ, อายูหนา ทุกฺขํ, อนายูหนา สุขนฺติ อภิญฺเญยฺยํ, ปฏิสนฺธิ ทุกฺขํ, อปฺปฏิสนฺธิ สุขนฺติ อภิญฺเญยฺยํ, คติ ทุกฺขํ, อคติ สุขนฺติ อภิญฺเญยฺยํ, นิพฺพตฺติ ทุกฺขํ, อนิพฺพตฺติ สุขนฺติ อภิญฺเญยฺยํ, อุปปตฺติ ทุกฺขํ, อนุปปตฺติ สุขนฺติ อภิญฺเญยฺยํ, ชาติ ทุกฺขํ, อชาติ สุขนฺติ อภิญฺเญยฺยํ, ชรา ทุกฺขํ, อชรา สุขนฺติ อภิญฺเญยฺยํ, พฺยาธิ ทุกฺขํ, อพฺยาธิ สุขนฺติ อภิญฺเญยฺยํ, มรณํ ทุกฺขํ, อมตํ สุขนฺติ อภิญฺเญยฺยํ, โสโก ทุกฺขํ, อโสโก สุขนฺติ อภิญฺเญยฺยํ, ปริเทโว ทุกฺขํ, อปริเทโว สุขนฺติ อภิญฺเญยฺยํ, อุปายาโส ทุกฺขํ, อนุปายาโส สุขนฺติ อภิญฺเญยฺยํ. (30)}

\prose{อุปฺปาโท ภยนฺติ อภิญฺเญยฺยํ, ปวตฺตํ ภยนฺติ อภิญฺเญยฺยํ, นิมิตฺตํ ภยนฺติ อภิญฺเญยฺยํ, อายูหนา ภยนฺติ อภิญฺเญยฺยํ, ปฏิสนฺธิ ภยนฺติ อภิญฺเญยฺยํ, คติ ภยนฺติ อภิญฺเญยฺยํ, นิพฺพตฺติ ภยนฺติ อภิญฺเญยฺยํ, อุปปตฺติ ภยนฺติ อภิญฺเญยฺยํ, ชาติ ภยนฺติ อภิญฺเญยฺยํ. ชรา ภยนฺติ อภิญฺเญยฺยํ, พฺยาธิ ภยนฺติ อภิญฺเญยฺยํ, มรณํ ภยนฺติ อภิญฺเญยฺยํ, โสโก ภยนฺติ อภิญฺเญยฺยํ, ปริเทโว ภยนฺติ อภิญฺเญยฺยํ, อุปายาโส ภยนฺติ อภิญฺเญยฺยํ. (15)}

\prose{อนุปฺปาโท เขมนฺติ อภิญฺเญยฺยํ, อปฺปวตฺตํ เขมนฺติ อภิญฺเญยฺยํ, อนิมิตฺตํ เขมนฺติ อภิญฺเญยฺยํ, อนายูหนา เขมนฺติ อภิญฺเญยฺยํ, อปฺปฏิสนฺธิ \csromanfolio{14}เขมนฺติ อภิญฺเญยฺยํ, อคติ เขมนฺติ อภิญฺเญยฺยํ, อนิพฺพตฺติ เขมนฺติ อภิญฺเญยฺยํ, อนุปปตฺติ เขมนฺติ อภิญฺเญยฺยํ, อชาติ เขมนฺติ อภิญฺเญยฺยํ, อชรา เขมนฺติ อภิญฺเญยฺยํ, อพฺยาธิ เขมนฺติ อภิญฺเญยฺยํ, อมตํ เขมนฺติ อภิญฺเญยฺยํ, อโสโก เขมนฺติ อภิญฺเญยฺยํ, อปริเทโว เขมนฺติ อภิญฺเญยฺยํ, อนุปายาโส เขมนฺติ อภิญฺเญยฺยํ. (15)}

\prose{อุปฺปาโท ภยํ, อนุปฺปาโท เขมนฺติ อภิญฺเญยฺยํ, ปวตฺตํ ภยํ, อปฺปวตฺตํ เขมนฺติ อภิญฺเญยฺยํ, นิมิตฺตํ ภยํ, อนิมิตฺตํ เขมนฺติ อภิญฺเญยฺยํ, อายูหนา ภยํ, อนายูหนา เขมนฺติ อภิญฺเญยฺยํ, ปฏิสนฺธิ ภยํ, อปฺปฏิสนฺธิ เขมนฺติ อภิญฺเญยฺยํ, คติ ภยํ, อคติ เขมนฺติ อภิญฺเญยฺยํ, นิพฺพตฺติ ภยํ, อนิพฺพตฺติ เขมนฺติ อภิญฺเญยฺยํ, อุปปตฺติ ภยํ, อนุปปตฺติ เขมนฺติ อภิญฺเญยฺยํ, ชาติ ภยํ, อชาติ เขมนฺติ อภิญฺเญยฺยํ, ชรา ภยํ, อชรา เขมนฺติ อภิญฺเญยฺยํ, พฺยาธิ ภยํ, อพฺยาธิ เขมนฺติ อภิญฺเญยฺยํ, มรณํ ภยํ, อมตํ เขมนฺติ อภิญฺเญยฺยํ, โสโก ภยํ, อโสโก เขมนฺติ อภิญฺเญยฺยํ, ปริเทโว ภยํ อปริเทโว เขมนฺติ อภิญฺเญยฺยํ, อุปายาโส ภยํ, อนุปายาโส เขมนฺติ อภิญฺเญยฺยํ. (30)}

\prose{อุปฺปาโท สามิสนฺติ อภิญฺเญยฺยํ, ปวตฺตํ สามิสนฺติ อภิญฺเญยฺยํ, นิมิตฺตํ สามิสนฺติ อภิญฺเญยฺยํ, อายูหนา สามิสนฺติ อภิญฺเญยฺยํ, ปฏิสนฺธิ สามิสนฺติ อภิญฺเญยฺยํ, คติ สามิสนฺติ อภิญฺเญยฺยํ, นิพฺพตฺติ สามิสนฺติ อภิญฺเญยฺยํ, อุปปตฺติ สามิสนฺติ อภิญฺเญยฺยํ, ชาติ สามิสนฺติ อภิญฺเญยฺยํ, ชรา สามิสนฺติ อภิญฺเญยฺยํ, พฺยาธิ สามิสนฺติ อภิญฺเญยฺยํ, มรณํ สามิสนฺติ อภิญฺเญยฺยํ, โสโก สามิสนฺติ อภิญฺเญยฺยํ, ปริเทโว สามิสนฺติ อภิญฺเญยฺยํ, อุปายาโส สามิสนฺติ อภิญฺเญยฺยํ. (15)}

\prose{อนุปฺปาโท นิรามิสนฺติ อภิญฺเญยฺยํ, อปฺปวตฺตํ นิรามิสนฺติ อภิญฺเญยฺยํ, อนิมิตฺตํ นิรามิสนฺติ อภิญฺเญยฺยํ, อนายูหนา นิรามิสนฺติ อภิญฺเญยฺยํ, อปฺปฏิสนฺธิ นิรามิสนฺติ อภิญฺเญยฺยํ, อคติ นิรามิสนฺติ อภิญฺเญยฺยํ, อนิพฺพตฺติ นิรามิสนฺติ อภิญฺเญยฺยํ, อนุปปตฺติ นิรามิสนฺติ อภิญฺเญยฺยํ, อชาติ นิรามิมิสนฺติ อภิญฺเญยฺยํ, อชรา นิรามิสนฺติ อภิญฺเญยฺยํ, อพฺยาธิ นิรามิสนฺติ อภิญฺเญยฺยํ, อมตํ นิรามิสนฺติ อภิญฺเญยฺยํ, อโสโก นิรามิสนฺติ อภิญฺเญยฺยํ, อปริเทโว นิรามิสนฺติ อภิญฺเญยฺยํ, อนุปายาโส นิรามิสนฺติ อภิญฺเญยฺยํ. (15)}

\prose{\csromanfolio{15}อุปฺปาโท สามิสํ, อนุปฺปาโท นิรามิสนฺติ อภิญฺเญยฺยํ, ปวตฺตํ สามิสํ, อปฺปวตฺตํ นิรามิสนฺติ อภิญฺเญยฺยํ, นิมิตฺตํ สามิสํ, อนิมิตฺตํ นิรามิสนฺติ อภิญฺเญยฺยํ, อายูหนา สามิสํ, อนายูหนา นิรามิสนฺติ อภิญฺเญยฺยํ, ปฏิสนฺธิ สามิสํ, อปฺปฏิสนฺธิ นิรามิสนฺติ อภิญฺเญยฺยํ, คติ สามิสํ, อคติ นิรามิสนฺติ อภิญฺเญยฺยํ, นิพฺพตฺติ สามิสํ, อนิพฺพตฺติ นิรามิสนฺติ อภิญฺเญยฺยํ, อุปปตฺติ สามิสํ, อนุปปตฺติ นิรามิสนฺติ อภิญฺเญยฺยํ, ชาติ สามิสํ, อชาติ นิรามิสนฺติ อภิญฺเญยฺยํ, ชรา สามิสํ, อชรา นิรามิสนฺติ อภิญฺเญยฺยํ, พฺยาธิ สามิสํ, อพฺยาธิ นิรามิสนฺติ อภิญฺเญยฺยํ, มรณํ สามิสํ, อมตํ นิรามิสนฺติ อภิญฺเญยฺยํ, โสโก สามิสํ, อโสโก นิรามิสนฺติ อภิญฺเญยฺยํ, ปริเทโว สามิสํ, อปริเทโว นิรามิสนฺติ อภิญฺเญยฺยํ, อุปายาโส สามิสํ, อนุปายาโส นิรามิสนฺติ อภิญฺเญยฺยํ. (30)}

\prose{อุปฺปาโท สงฺขาราติ อภิญฺเญยฺยํ, ปวตฺตํ สงฺขาราติ อภิญฺเญยฺยํ, นิมิตฺตํ สงฺขาราติ อภิญฺเญยฺยํ, อายูหนา สงฺขาราติ อภิญฺเญยฺยํ, ปฏิสนฺธิ สงฺขาราติ อภิญฺเญยฺยํ, คติ สงฺขาราติ อภิญฺเญยฺยํ, นิพฺพตฺติ สงฺขาราติ อภิญฺเญยฺยํ, อุปปตฺติ สงฺขาราติ อภิญฺเญยฺยํ, ชาติ สงฺขาราติ อภิญฺเญยฺยํ, ชรา สงฺขาราติ อภิญฺเญยฺยํ, พฺยาธิ สงฺขาราติ อภิญฺเญยฺยํ, มรณํ สงฺขาราติ อภิญฺเญยฺยํ, โสโก สงฺขาราติ อภิญฺเญยฺยํ, ปริเทโว สงฺขาราติ อภิญฺเญยฺยํ, อุปายาโส สงฺขาราติ อภิญฺเญยฺยํ. (15)}

\prose{อนุปฺปาโท นิพฺพานนฺติ อภิญฺเญยฺยํ, อปฺปวตฺตํ นิพฺพานนฺติ อภิญฺเญยฺยํ, อนิมิตฺตํ นิพฺพานนฺติ อภิญฺเญยฺยํ, อนายูหนา นิพฺพานนฺติ อภิญฺเญยฺยํ, อปฺปฏิสนฺธิ นิพฺพานนฺติ อภิญฺเญยฺยํ, อคติ นิพฺพานนฺติ อภิญฺเญยฺยํ, อนิพฺพตฺติ นิพฺพานนฺติ อภิญฺเญยฺยํ, อนุปปตฺติ นิพฺพานนฺติ อภิญฺเญยฺยํ, อชาติ นิพฺพานนฺติ อภิญฺเญยฺยํ, อชรา นิพฺพานนฺติ อภิญฺเญยฺยํ, อพฺยาธิ นิพฺพานนฺติ อภิญฺเญยฺยํ, อมตํ นิพฺพานนฺติ อภิญฺเญยฺยํ, อโสโก นิพฺพานนฺติ อภิญฺเญยฺยํ. อปริเทโว นิพฺพานนฺติ อภิญฺเญยฺยํ, อนุปายาโส นิพฺพานนฺติ อภิญฺเญยฺยํ. (15)}

\prose{อุปฺปาโท สงฺขารา, อนุปฺปาโท นิพฺพานนฺติ อภิญฺเญยฺยํ, ปวตฺตํ สงฺขารา, อปฺปวตฺตํ นิพฺพานนฺติ อภิญฺเญยฺยํ, นิมิตฺตํ สงฺขารา, อนิมิตฺตํ นิพฺพานนฺติ อภิญฺเญยฺยํ, อายูหนา สงฺขารา, อนายูหนา นิพฺพานนฺติ อภิญฺเญยฺยํ, ปฏิสนฺธิ สงฺขารา, อปฺปฏิสนฺธิ นิพฺพานนฺติ อภิญฺเญยฺยํ, คติ สงฺขารา, อคติ นิพฺพานนฺติ อภิญฺเญยฺยํ, นิพฺพตฺติ สงฺขารา, อนิพฺพตฺติ นิพฺพานนฺติ อภิญฺเญยฺยํ, อุปปตฺติ สงฺขารา, อนุปปตฺติ นิพฺพานนฺติ อภิญฺเญยฺยํ, ชาติ สงฺขารา, อชาติ นิพฺพานนฺติ อภิญฺเญยฺยํ, ชรา สงฺขารา, อชรา \csromanfolio{16}นิพฺพานนฺติ อภิญฺเญยฺยํ, พฺยาธิ สงฺขารา, อพฺยาธิ นิพฺพานนฺติ อภิญฺเญยฺยํ, มรณํ สงฺขารา, อมตํ นิพฺพานนฺติ อภิญฺเญยฺยํ, โสโก สงฺขารา, อโสโก นิพฺพานนฺติ อภิญฺเญยฺยํ, ปริเทโว สงฺขารา, อปริเทโว นิพฺพานนฺติ อภิญฺเญยฺยํ, อุปายาโส สงฺขารา, อนุปายาโส นิพฺพานนฺติ อภิญฺเญยฺยํ. (30) ปฐม\-ภาณวาโร.}

\proseitem{11}{ปริคฺคหฏฺโฐ อภิญฺเญยฺโย, ปริวารฏฺโฐ อภิญฺเญยฺโย, ปริปูรฏฺโฐ\csromansharedfootnote{71239795947609cb}{ปริปูรณฏฺโฐ (ก), ปริปูริฏฺโฐ (สีฏฺฐ)} อภิญฺเญยฺโย, เอกคฺคฏฺโฐ อภิญฺเญยฺโย, อวิกฺเขปฏฺโฐ อภิญฺเญยฺโย, ปคฺคหฏฺโฐ อภิญฺเญยฺโย, อวิสารฏฺโฐ อภิญฺเญยฺโย, อนาวิลฏฺโฐ อภิญฺเญยฺโย, อนิญฺชนฏฺโฐ อภิญฺเญยฺโย, เอกตฺตุ\-ปฏฺฐาน\-วเสน จิตฺตสฺส ฐิตฏฺโฐ อภิญฺเญยฺโย, อารมฺมณฏฺโฐ อภิญฺเญยฺโย, โคจรฏฺโฐ อภิญฺเญยฺโย, ปหานฏฺโฐ อภิญฺเญยฺโย, ปริจฺจาคฏฺโฐ อภิญฺเญยฺโย, วุฏฺฐานฏฺโฐ อภิญฺเญยฺโย, วิวฏฺฏนฏฺโฐ\csromansharedfootnote{c019b3dfce93b988}{นิวตฺตนฏฺโฐ (ก)} อภิญฺเญยฺโย, สนฺตฏฺโฐ อภิญฺเญยฺโย, ปณีตฏฺโฐ อภิญฺเญยฺโย, วิมุตฺตฏฺโฐ อภิญฺเญยฺโย, อนาสวฏฺโฐ อภิญฺเญยฺโย, ตรณฏฺโฐ อภิญฺเญยฺโย, อนิมิตฺตฏฺโฐ อภิญฺเญยฺโย, อปฺปณิหิตฏฺโฐ อภิญฺเญยฺโย, สุญฺญตฏฺโฐ อภิญฺเญยฺโย, เอกรสฏฺโฐ อภิญฺเญยฺโย, อนติวตฺตน\-ฏฺโฐ อภิญฺเญยฺโย, ยุคนทฺธฏฺโฐ\csromansharedfootnote{f3ef268cc634cd17}{ยุคนนฺธนฏฺโฐ (ก)} อภิญฺเญยฺโย, นิยฺยานฏฺโฐ อภิญฺเญยฺโย, เหตุฏฺโฐ\csromansharedfootnote{6fc8d664db30d8bc}{เหตฏฺโฐ (สฺยา)} อภิญฺเญยฺโย, ทสฺสนฏฺโฐ อภิญฺเญยฺโย, อาธิปเตยฺย\-ฏฺโฐ อภิญฺเญยฺโย. (31)}

\proseitem{12}{สมถสฺส อวิกฺเขปฏฺโฐ อภิญฺเญยฺโย, วิปสฺสนาย อนุปสฺสนฏฺโฐ อภิญฺเญยฺโย, สมถ\-วิปสฺสนานํ เอกรสฏฺโฐ อภิญฺเญยฺโย, ยุคนทฺธสฺส อนติวตฺตน\-ฏฺโฐ อภิญฺเญยฺโย. (4)}

\prose{สิกฺขาย สมาทานฏฺโฐ อภิญฺเญยฺโย, อารมฺมณสฺส โคจรฏฺโฐ อภิญฺเญยฺโย, ลีนสฺส จิตฺตสฺส ปคฺคหฏฺโฐ อภิญฺเญยฺโย, อุทฺธตสฺส จิตฺตสฺส นิคฺคหฏฺโฐ อภิญฺเญยฺโย, อุโภ\-วิสุทฺธานํ อชฺฌุเปกฺข\-นฏฺโฐ อภิญฺเญยฺโย, วิเสสาธิคม\-ฏฺโฐ อภิญฺเญยฺโย, อุตฺตริ\-ปฏิเวธฏฺโฐ อภิญฺเญยฺโย, สจฺจาภิสมย\-ฏฺโฐ อภิญฺเญยฺโย, นิโรเธ ปติฏฺฐา\-ปกฏฺโฐ อภิญฺเญยฺโย. (9)}

\prose{\csromanfolio{17}สทฺธิ\-นฺทฺริยสฺส อธิโมกฺขฏฺโฐ อภิญฺเญยฺโย, วีริยิ\-นฺทฺริยสฺส ปคฺคหฏฺโฐ อภิญฺเญยฺโย, สตินฺทฺริยสฺส อุปฏฺฐานฏฺโฐ อภิญฺเญยฺโย, สมาธิ\-นฺทฺริยสฺส อวิกฺเขปฏฺโฐ อภิญฺเญยฺโย, ปญฺญินฺทฺริยสฺส ทสฺสนฏฺโฐ อภิญฺเญยฺโย. (5)}

\prose{สทฺธาพลสฺส อสฺสทฺธิเย อกมฺปิยฏฺโฐ อภิญฺเญยฺโย, วีริยพลสฺส โกสชฺเช อกมฺปิยฏฺโฐ อภิญฺเญยฺโย, สติพลสฺส ปมาเท อกมฺปิยฏฺโฐ อภิญฺเญยฺโย, สมาธิพลสฺส อุทฺธจฺเจ อกมฺปิยฏฺโฐ อภิญฺเญยฺโย, ปญฺญาพลสฺส อวิชฺชาย อกมฺปิยฏฺโฐ อภิญฺเญยฺโย. (5)}

\prose{สติสมฺโพชฺฌงฺคสฺส อุปฏฺฐานฏฺโฐ อภิญฺเญยฺโย, ธมฺมวิจย\-สมฺโพชฺฌงฺคสฺส ปวิจยฏฺโฐ อภิญฺเญยฺโย, วีริย\-สมฺโพชฺฌงฺคสฺส ปคฺคหฏฺโฐ อภิญฺเญยฺโย, ปีติสมฺโพชฺฌงฺคสฺส ผรณฏฺโฐ อภิญฺเญยฺโย, ปสฺสทฺธิ\-สมฺโพชฺฌงฺคสฺส อุปสมฏฺโฐ อภิญฺเญยฺโย, สมาธิ\-สมฺโพชฺฌงฺคสฺส อวิกฺเขปฏฺโฐ อภิญฺเญยฺโย, อุเปกฺขา\-สมฺโพชฺฌงฺคสฺส ปฏิสงฺขา\-นฏฺโฐ อภิญฺเญยฺโย. (7)}

\prose{สมฺมาทิฏฺฐิยา ทสฺสนฏฺโฐ อภิญฺเญยฺโย, สมฺมา\-สงฺกปฺปสฺส อภินิโรปน\-ฏฺโฐ อภิญฺเญยฺโย, สมฺมาวาจาย ปริคฺคหฏฺโฐ อภิญฺเญยฺโย, สมฺมา\-กมฺมนฺตสฺส สมุฏฺฐานฏฺโฐ อภิญฺเญยฺโย, สมฺมาอาชีวสฺส โวทานฏฺโฐ อภิญฺเญยฺโย, สมฺมาวายามสฺส ปคฺคหฏฺโฐ อภิญฺเญยฺโย, สมฺมาสติยา อุปฏฺฐานฏฺโฐ อภิญฺเญยฺโย, สมฺมา\-สมาธิสฺส อวิกฺเขปฏฺโฐ อภิญฺเญยฺโย. (8)}

\proseitem{13}{อินฺทฺริยานํ อาธิปเตยฺย\-ฏฺโฐ อภิญฺเญยฺโย, พลานํ อกมฺปิยฏฺโฐ อภิญฺเญยฺโย, โพชฺฌงฺคานํ นิยฺยานฏฺโฐ อภิญฺเญยฺโย, มคฺคสฺส เหตุฏฺโฐ อภิญฺเญยฺโย, สติปฏฺฐานานํ อุปฏฺฐานฏฺโฐ อภิญฺเญยฺโย, สมฺม\-ปฺปธานานํ ปทหนฏฺโฐ อภิญฺเญยฺโย, อิทฺธิปาทานํ อิชฺฌนฏฺโฐ อภิญฺเญยฺโย, สจฺจานํ ตถฏฺโฐ อภิญฺเญยฺโย, ปโยคานํ\csromansharedfootnote{5ac181bf9ae76b4b}{มคฺคานํ (สพฺพตฺถ) อฏฺฐกถา ปสฺสิตพฺพา.} ปฏิปฺปสฺสทฺธ\-ฏฺโฐ อภิญฺเญยฺโย, ผลานํ สจฺฉิกิริยฏฺโฐ อภิญฺเญยฺโย. (10)}

\prose{วิตกฺกสฺส อภินิโรปน\-ฏฺโฐ อภิญฺเญยฺโย, วิจารสฺส อุปวิจารฏฺโฐ อภิญฺเญยฺโย, ปีติยา ผรณฏฺโฐ อภิญฺเญยฺโย, สุขสฺส อภิสนฺท\-นฏฺโฐ อภิญฺเญยฺโย, จิตฺตสฺส เอกคฺคฏฺโฐ อภิญฺเญยฺโย. (5)}

\prose{อาวชฺชนฏฺโฐ อภิญฺเญยฺโย, วิชานนฏฺโฐ อภิญฺเญยฺโย, ปชานนฏฺโฐ อภิญฺเญยฺโย, สญฺชานนฏฺโฐ อภิญฺเญยฺโย, เอโกทฏฺโฐ \csromanfolio{18}อภิญฺเญยฺโย, อภิญฺญาย ญาตฏฺโฐ อภิญฺเญยฺโย, ปริญฺญาย ตีรณฏฺโฐ อภิญฺเญยฺโย, ปหานสฺส ปริจฺจาคฏฺโฐ อภิญฺเญยฺโย, ภาวนาย เอกรสฏฺโฐ อภิญฺเญยฺโย, สจฺฉิกิริยาย ผสฺสนฏฺโฐ อภิญฺเญยฺโย, ขนฺธานํ ขนฺธฏฺโฐ อภิญฺเญยฺโย, ธาตูนํ ธาตุฏฺโฐ\csromansharedfootnote{ff9745a9c36280dd}{ธาตฏฺโฐ (สฺยา)} อภิญฺเญยฺโย, อายตนานํ อายตนฏฺโฐ อภิญฺเญยฺโย, สงฺขตานํ สงฺขตฏฺโฐ อภิญฺเญยฺโย, อสงฺขตสฺส อสงฺขตฏฺโฐ อภิญฺเญยฺโย. (15)}

\proseitem{14}{จิตฺตฏฺโฐ อภิญฺเญยฺโย, จิตฺตา\-นนฺตริย\-ฏฺโฐ อภิญฺเญยฺโย, จิตฺตสฺส วุฏฺฐานฏฺโฐ อภิญฺเญยฺโย, จิตฺตสฺส วิวฏฺฏนฏฺโฐ อภิญฺเญยฺโย, จิตฺตสฺส เหตุฏฺโฐ อภิญฺเญยฺโย, จิตฺตสฺส ปจฺจยฏฺโฐ อภิญฺเญยฺโย, จิตฺตสฺส วตฺถุฏฺโฐ\csromansharedfootnote{106083fce763f253}{วตฺถฏฺโฐ (สฺยา)} อภิญฺเญยฺโย, จิตฺตสฺส ภูมฏฺโฐ\csromansharedfootnote{0be13744a32a9c64}{ภุมฺมฏฺโฐ (สฺยา, สีฏฺฐ) (เอกตฺเต อุปนิพนฺธนฏฺโฐ อภิญฺเญยฺโย) (ก) อฏฺฐกถา ปสฺสิตพฺพา.} อภิญฺเญยฺโย, จิตฺตสฺส อารมฺมณฏฺโฐ อภิญฺเญยฺโย, จิตฺตสฺส โคจรฏฺโฐ อภิญฺเญยฺโย, จิตฺตสฺส จริยฏฺโฐ อภิญฺเญยฺโย, จิตฺตสฺส คตฏฺโฐ อภิญฺเญยฺโย, จิตฺตสฺส อภินีหารฏฺโฐ อภิญฺเญยฺโย, จิตฺตสฺส นิยฺยานฏฺโฐ อภิญฺเญยฺโย, จิตฺตสฺส นิสฺสรณฏฺโฐ อภิญฺเญยฺโย. (15)}

\proseitem{15}{เอกตฺเต อาวชฺชนฏฺโฐ อภิญฺเญยฺโย, เอกตฺเต วิชานนฏฺโฐ อภิญฺเญยฺโย, เอกตฺเต ปชานนฏฺโฐ อภิญฺเญยฺโย, เอกตฺเต สญฺชานนฏฺโฐ อภิญฺเญยฺโย, เอกตฺเต เอโกทฏฺโฐ อภิญฺเญยฺโย, ( ) เอกตฺเต ปกฺขนฺท\-นฏฺโฐ อภิญฺเญยฺโย, เอกตฺเต ปสีทนฏฺโฐ อภิญฺเญยฺโย, เอกตฺเต สนฺติฏฺฐ\-นฏฺโฐ อภิญฺเญยฺโย, เอกตฺเต วิมุจฺจนฏฺโฐ อภิญฺเญยฺโย, เอกตฺเต “เอตํ สนฺตนฺ”ติ ปสฺสนฏฺโฐ อภิญฺเญยฺโย, เอกตฺเต ยานีกตฏฺโฐ อภิญฺเญยฺโย, เอกตฺเต วตฺถุกตฏฺโฐ อภิญฺเญยฺโย, เอกตฺเต อนุฏฺฐิตฏฺโฐ อภิญฺเญยฺโย, เอกตฺเต ปริจิตฏฺโฐ อภิญฺเญยฺโย, เอกตฺเต สุสมารทฺธ\-ฏฺโฐ อภิญฺเญยฺโย, เอกตฺเต ปริคฺคหฏฺโฐ อภิญฺเญยฺโย, เอกตฺเต ปริวารฏฺโฐ อภิญฺเญยฺโย, เอกตฺเต ปริปูรฏฺโฐ อภิญฺเญยฺโย, เอกตฺเต สโมธานฏฺโฐ อภิญฺเญยฺโย, เอกตฺเต อธิฏฺฐานฏฺโฐ อภิญฺเญยฺโย, เอกตฺเต อาเสวนฏฺโฐ อภิญฺเญยฺโย, เอกตฺเต ภาวนฏฺโฐ อภิญฺเญยฺโย, เอกตฺเต พหุลีกมฺม\-ฏฺโฐ อภิญฺเญยฺโย, เอกตฺเต สุสมุคฺคต\-ฏฺโฐ อภิญฺเญยฺโย, เอกตฺเต สุวิมุตฺตฏฺโฐ \csromanfolio{19}อภิญฺเญยฺโย, เอกตฺเต พุชฺฌนฏฺโฐ อภิญฺเญยฺโย, เอกตฺเต อนุพุชฺฌ\-นฏฺโฐ อภิญฺเญยฺโย, เอกตฺเต ปฏิพุชฺฌ\-นฏฺโฐ อภิญฺเญยฺโย, เอกตฺเต สมฺพุชฺฌ\-นฏฺโฐ อภิญฺเญยฺโย, เอกตฺเต โพธนฏฺโฐ อภิญฺเญยฺโย, เอกตฺเต อนุโพธนฏฺโฐ อภิญฺเญยฺโย, เอกตฺเต ปฏิโพธ\-นฏฺโฐ อภิญฺเญยฺโย, เอกตฺเต สมฺโพธนฏฺโฐ อภิญฺเญยฺโย, เอกตฺเต โพธิปกฺขิย\-ฏฺโฐ อภิญฺเญยฺโย, เอกตฺเต อนุโพธิปกฺขิย\-ฏฺโฐ อภิญฺเญยฺโย, เอกตฺเต ปฏิโพธิปกฺขิย\-ฏฺโฐ อภิญฺเญยฺโย, เอกตฺเต สมฺโพธิ\-ปกฺขิย\-ฏฺโฐ อภิญฺเญยฺโย, เอกตฺเต โชตนฏฺโฐ อภิญฺเญยฺโย, เอกตฺเต อุชฺโชตนฏฺโฐ อภิญฺเญยฺโย, เอกตฺเต อนุโชตนฏฺโฐ อภิญฺเญยฺโย, เอกตฺเต ปฏิโชตนฏฺโฐ อภิญฺเญยฺโย, เอกตฺเต สญฺโชตนฏฺโฐ อภิญฺเญยฺโย. (42)}

\proseitem{16}{ปตาปนฏฺโฐ\csromansharedfootnote{dfde6eef21f6f96f}{ปกาสนฏฺโฐ (ก)} อภิญฺเญยฺโย, วิโรจนฏฺโฐ อภิญฺเญยฺโย, กิเลสานํ สนฺตาปนฏฺโฐ อภิญฺเญยฺโย, อมลฏฺโฐ อภิญฺเญยฺโย, วิมลฏฺโฐ อภิญฺเญยฺโย, นิมฺมลฏฺโฐ อภิญฺเญยฺโย, สมฏฺโฐ อภิญฺเญยฺโย, สมยฏฺโฐ อภิญฺเญยฺโย, วิเวกฏฺโฐ อภิญฺเญยฺโย, วิเวก\-จริย\-ฏฺโฐ อภิญฺเญยฺโย, วิราคฏฺโฐ อภิญฺเญยฺโย, วิราค\-จริย\-ฏฺโฐ อภิญฺเญยฺโย, นิโรธฏฺโฐ อภิญฺเญยฺโย, นิโรธ\-จริย\-ฏฺโฐ อภิญฺเญยฺโย, โวสคฺคฏฺโฐ\csromansharedfootnote{d893c1b62b599789}{โปสฺสคฺคฏฺโฐ (สฺยา, ก)} อภิญฺเญยฺโย, โวสคฺค\-จริย\-ฏฺโฐ อภิญฺเญยฺโย วิมุตฺตฏฺโฐ อภิญฺเญยฺโย, วิมุตฺติ\-จริย\-ฏฺโฐ อภิญฺเญยฺโย. (18)}

\prose{ฉนฺทฏฺโฐ อภิญฺเญยฺโย, ฉนฺทสฺส มูลฏฺโฐ อภิญฺเญยฺโย, ฉนฺทสฺส ปาทฏฺโฐ อภิญฺเญยฺโย, ฉนฺทสฺส ปธานฏฺโฐ อภิญฺเญยฺโย, ฉนฺทสฺส อิชฺฌนฏฺโฐ อภิญฺเญยฺโย, ฉนฺทสฺส อธิโมกฺขฏฺโฐ อภิญฺเญยฺโย, ฉนฺทสฺส ปคฺคหฏฺโฐ อภิญฺเญยฺโย, ฉนฺทสฺส อุปฏฺฐานฏฺโฐ อภิญฺเญยฺโย, ฉนฺทสฺส อวิกฺเขปฏฺโฐ อภิญฺเญยฺโย, ฉนฺทสฺส ทสฺสนฏฺโฐ อภิญฺเญยฺโย. (10)}

\prose{วีริยฏฺโฐ อภิญฺเญยฺโย, วีริยสฺส มูลฏฺโฐ อภิญฺเญยฺโย, วีริยสฺส ปาทฏฺโฐ อภิญฺเญยฺโย, วีริยสฺส ปธานฏฺโฐ อภิญฺเญยฺโย, วีริยสฺส อิชฺฌนฏฺโฐ อภิญฺเญยฺโย, วีริยสฺส อธิโมกฺขฏฺโฐ อภิญฺเญยฺโย, วีริยสฺส ปคฺคหฏฺโฐ อภิญฺเญยฺโย, วีริยสฺส อุปฏฺฐานฏฺโฐ อภิญฺเญยฺโย, วีริยสฺส อวิกฺเขปฏฺโฐ อภิญฺเญยฺโย, วีริยสฺส ทสฺสนฏฺโฐ อภิญฺเญยฺโย. (10)}

\prose{\csromanfolio{20}จิตฺตฏฺโฐ อภิญฺเญยฺโย, จิตฺตสฺส มูลฏฺโฐ อภิญฺเญยฺโย, จิตฺตสฺส ปาทฏฺโฐ อภิญฺเญยฺโย, จิตฺตสฺส ปธานฏฺโฐ อภิญฺเญยฺโย, จิตฺตสฺส อิชฺฌนฏฺโฐ อภิญฺเญยฺโย, จิตฺตสฺส อธิโมกฺขฏฺโฐ อภิญฺเญยฺโย, จิตฺตสฺส ปคฺคหฏฺโฐ อภิญฺเญยฺโย, จิตฺตสฺส อุปฏฺฐานฏฺโฐ อภิญฺเญยฺโย, จิตฺตสฺส อวิกฺเขปฏฺโฐ อภิญฺเญยฺโย, จิตฺตสฺส ทสฺสนฏฺโฐ อภิญฺเญยฺโย. (10)}

\prose{วีมํสฏฺโฐ อภิญฺเญยฺโย, วีมํสาย มูลฏฺโฐ อภิญฺเญยฺโย, วีมํสาย ปาทฏฺโฐ อภิญฺเญยฺโย, วีมํสาย ปธานฏฺโฐ อภิญฺเญยฺโย, วีมํสาย อิชฺฌนฏฺโฐ อภิญฺเญยฺโย, วีมํสาย อธิโมกฺขฏฺโฐ อภิญฺเญยฺโย, วีมํสาย ปคฺคหฏฺโฐ อภิญฺเญยฺโย, วีมํสาย อุปฏฺฐานฏฺโฐ อภิญฺเญยฺโย, วีมํสาย อวิกฺเขปฏฺโฐ อภิญฺเญยฺโย, วีมํสาย ทสฺสนฏฺโฐ อภิญฺเญยฺโย. (10)}

\proseitem{17}{ทุกฺขฏฺโฐ อภิญฺเญยฺโย\csromandash{}ทุกฺขสฺส ปีฬนฏฺโฐ อภิญฺเญยฺโย, ทุกฺขสฺส สงฺขตฏฺโฐ อภิญฺเญยฺโย, ทุกฺขสฺส สนฺตาปฏฺโฐ อภิญฺเญยฺโย, ทุกฺขสฺส วิปริณามฏฺโฐ อภิญฺเญยฺโย.  สมุทยฏฺโฐ อภิญฺเญยฺโย\csromandash{} สมุทยสฺส อายูหนฏฺโฐ\csromansharedfootnote{07a40c5dd7402cea}{อายุหนฏฺโฐ (สฺยา)} อภิญฺเญยฺโย, สมุทยสฺส นิทานฏฺโฐ อภิญฺเญยฺโย, สมุทยสฺส สญฺโญคฏฺโฐ อภิญฺเญยฺโย, สมุทยสฺส ปลิโพธฏฺโฐ อภิญฺเญยฺโย.  นิโรธฏฺโฐ อภิญฺเญยฺโย\csromandash{}นิโรธสฺส นิสฺสรณฏฺโฐ อภิญฺเญยฺโย, นิโรธสฺส วิเวกฏฺโฐ อภิญฺเญยฺโย, นิโรธสฺส อสงฺขตฏฺโฐ อภิญฺเญยฺโย, นิโรธสฺส อมตฏฺโฐ อภิญฺเญยฺโย.  มคฺคฏฺโฐ อภิญฺเญยฺโย\csromandash{}มคฺคสฺส นิยฺยานฏฺโฐ อภิญฺเญยฺโย, มคฺคสฺส เหตุฏฺโฐ อภิญฺเญยฺโย, มคฺคสฺส ทสฺสนฏฺโฐ อภิญฺเญยฺโย, มคฺคสฺส อาธิปเตยฺย\-ฏฺโฐ อภิญฺเญยฺโย. (16)}

\prose{ตถฏฺโฐ อภิญฺเญยฺโย, อนตฺตฏฺโฐ อภิญฺเญยฺโย, สจฺจฏฺโฐ อภิญฺเญยฺโย, ปฏิเวธฏฺโฐ อภิญฺเญยฺโย, อภิชานนฏฺโฐ อภิญฺเญยฺโย, ปริชานนฏฺโฐ อภิญฺเญยฺโย, ธมฺมฏฺโฐ อภิญฺเญยฺโย, ธาตุฏฺโฐ อภิญฺเญยฺโย, ญาตฏฺโฐ อภิญฺเญยฺโย, สจฺฉิกิริยฏฺโฐ อภิญฺเญยฺโย, ผสฺสนฏฺโฐ อภิญฺเญยฺโย, อภิสมยฏฺโฐ อภิญฺเญยฺโย. (12)}

\proseitem{18}{เนกฺขมฺมํ อภิญฺเญยฺยํ, อพฺยาปาโท อภิญฺเญยฺโย, อาโลกสญฺญา อภิญฺเญยฺยา, อวิกฺเขโป อภิญฺเญยฺโย, ธมฺม\-ววตฺถานํ อภิญฺเญยฺยํ, ญาณํ อภิญฺเญยฺยํ, ปาโมชฺชํ\csromansharedfootnote{898b8ed305ba8a74}{ปามุชฺชํ (สฺยา)} อภิญฺเญยฺยํ. (7)}

\prose{\csromanfolio{21}ปฐมํ ฌานํ อภิญฺเญยฺยํ, ทุติยํ ฌานํ อภิญฺเญยฺยํ, ตติยํ ฌานํ อภิญฺเญยฺยํ, จตุตฺถํ ฌานํ อภิญฺเญยฺยํ, อากาสา\-นญฺจา\-ยตนสมาปตฺติ อภิญฺเญยฺยา, วิญฺญาณญฺจา\-ยตนสมาปตฺติ อภิญฺเญยฺยา, อากิญฺจญฺญายตน\-สมาปตฺติ อภิญฺเญยฺยา, เนว\-สญฺญา\-นา\-สญฺญา\-ยตนสมาปตฺติ อภิญฺเญยฺยา. (8)}

\prose{อนิจฺจา\-นุปสฺสนา อภิญฺเญยฺยา, ทุกฺขา\-นุปสฺสนา อภิญฺเญยฺยา, อนตฺตา\-นุปสฺสนา อภิญฺเญยฺยา, นิพฺพิทา\-นุปสฺสนา อภิญฺเญยฺยา, วิราคา\-นุปสฺสนา อภิญฺเญยฺยา, นิโรธา\-นุปสฺสนา อภิญฺเญยฺยา, ปฏินิสฺสคฺคา\-นุปสฺสนา อภิญฺเญยฺยา, ขยานุปสฺสนา อภิญฺเญยฺยา, วยานุปสฺสนา อภิญฺเญยฺยา, วิปริณามา\-นุปสฺสนา อภิญฺเญยฺยา, อนิมิตฺตา\-นุปสฺสนา อภิญฺเญยฺยา, อปฺปณิหิตา\-นุปสฺสนา อภิญฺเญยฺยา, สุญฺญตา\-นุปสฺสนา อภิญฺเญยฺยา, อธิปญฺญา\-ธมฺมวิปสฺสนา อภิญฺเญยฺยา, ยถา\-ภูต\-ญาณ\-ทสฺสนํ อภิญฺเญยฺยํ, อาทีนวา\-นุปสฺสนา อภิญฺเญยฺยา, ปฏิสงฺขา\-นุปสฺสนา อภิญฺเญยฺยา, วิวฏฺฏนา\-นุปสฺสนา อภิญฺเญยฺยา. (18)}

\proseitem{19}{โสตาปตฺติมคฺโค อภิญฺเญยฺโย, โสตาปตฺติ\-ผลสมาปตฺติ อภิญฺเญยฺยา, สกทาคามิ\-มคฺโค อภิญฺเญยฺโย, สกทาคามิ\-ผลสมาปตฺติ อภิญฺเญยฺยา, อนาคามิมคฺโค อภิญฺเญยฺโย, อนาคามิ\-ผลสมาปตฺติ อภิญฺเญยฺยา อรหตฺตมคฺโค อภิญฺเญยฺโย, อรหตฺตผล\-สมาปตฺติ อภิญฺเญยฺยา. (8)}

\prose{อธิโมกฺข\-ฏฺเฐน สทฺธินฺทฺริยํ อภิญฺเญยฺยํ, ปคฺคหฏฺเฐน วีริยินฺทฺริยํ อภิญฺเญยฺยํ, อุปฏฺฐาน\-ฏฺเฐน สตินฺทฺริยํ อภิญฺเญยฺยํ, อวิกฺเขป\-ฏฺเฐน สมาธินฺทฺริยํ อภิญฺเญยฺยํ, ทสฺสนฏฺเฐน ปญฺญินฺทฺริยํ อภิญฺเญยฺยํ.}

\prose{อสฺสทฺธิเย อกมฺปิยฏฺเฐน สทฺธาพลํ อภิญฺเญยฺยํ, โกสชฺเช อกมฺปิยฏฺเฐน วีริยพลํ อภิญฺเญยฺยํ, ปมาเท อกมฺปิยฏฺเฐน สติพลํ อภิญฺเญยฺยํ, อุทฺธจฺเจ อกมฺปิยฏฺเฐน สมาธิพลํ อภิญฺเญยฺยํ, อวิชฺชาย อกมฺปิยฏฺเฐน ปญฺญาพลํ อภิญฺเญยฺยํ.}

\prose{อุปฏฺฐาน\-ฏฺเฐน สติสมฺโพชฺฌงฺโค อภิญฺเญยฺโย, ปวิจยฏฺเฐน ธมฺมวิจย\-สมฺโพชฺฌงฺโค อภิญฺเญยฺโย, ปคฺคหฏฺเฐน วีริย\-สมฺโพชฺฌงฺโค อภิญฺเญยฺโย, ผรณฏฺเฐน ปีติสมฺโพชฺฌงฺโค อภิญฺเญยฺโย, อุปสมฏฺเฐน ปสฺสทฺธิ\-สมฺโพชฺฌงฺโค อภิญฺเญยฺโย, อวิกฺเขป\-ฏฺเฐน สมาธิ\-สมฺโพชฺฌงฺโค อภิญฺเญยฺโย, ปฏิสงฺขา\-นฏฺเฐน อุเปกฺขา\-สมฺโพชฺฌงฺโค อภิญฺเญยฺโย.}

\prose{\csromanfolio{22}ทสฺสนฏฺเฐน สมฺมาทิฏฺฐิ อภิญฺเญยฺยา, อภินิโรปน\-ฏฺเฐน สมฺมาสงฺกปฺโป อภิญฺเญยฺโย, ปริคฺคห\-ฏฺเฐน สมฺมาวาจา อภิญฺเญยฺยา, สมุฏฺฐาน\-ฏฺเฐน สมฺมากมฺมนฺโต อภิญฺเญยฺโย, โวทานฏฺเฐน สมฺมาอาชีโว อภิญฺเญยฺโย, ปคฺคหฏฺเฐน สมฺมาวายาโม อภิญฺเญยฺโย, อุปฏฺฐาน\-ฏฺเฐน สมฺมาสติ อภิญฺเญยฺยา, อวิกฺเขป\-ฏฺเฐน สมฺมาสมาธิ อภิญฺเญยฺโย.}

\prose{อาธิปเตยฺย\-ฏฺเฐน อินฺทฺริยา อภิญฺเญยฺยา, อกมฺปิยฏฺเฐน พลา อภิญฺเญยฺยา, นิยฺยานฏฺเฐน โพชฺฌงฺคา อภิญฺเญยฺยา, เหตุฏฺเฐน มคฺโค อภิญฺเญยฺโย, อุปฏฺฐาน\-ฏฺเฐน สติปฏฺฐานา อภิญฺเญยฺยา, ปทหนฏฺเฐน สมฺมปฺปธานา อภิญฺเญยฺยา, อิชฺฌนฏฺเฐน อิทฺธิปาทา อภิญฺเญยฺยา, ตถฏฺเฐน สจฺจา อภิญฺเญยฺยา. (33)}

\prose{อวิกฺเขป\-ฏฺเฐน สมโถ อภิญฺเญยฺโย, อนุปสฺสนฏฺเฐน วิปสฺสนา อภิญฺเญยฺยา, เอกรสฏฺเฐน สมถ\-วิปสฺสนา อภิญฺเญยฺยา, อนติวตฺตน\-ฏฺเฐน ยุคนทฺธํ อภิญฺเญยฺยํ. (4)}

\prose{สํวรฏฺเฐน สีลวิสุทฺธิ อภิญฺเญยฺยา, อวิกฺเขป\-ฏฺเฐน จิตฺตวิสุทฺธิ อภิญฺเญยฺยา, ทสฺสนฏฺเฐน ทิฏฺฐิวิสุทฺธิ อภิญฺเญยฺยา, มุตฺตฏฺเฐน วิโมกฺโข อภิญฺเญยฺโย, ปฏิเวธ\-ฏฺเฐน วิชฺชา อภิญฺเญยฺยา, ปริจฺจาค\-ฏฺเฐน วิมุตฺติ อภิญฺเญยฺยา, สมุจฺเฉท\-ฏฺเฐน ขเย ญาณํ อภิญฺเญยฺยํ, ปฏิปฺปสฺสทฺธ\-ฏฺเฐน อนุปฺปาเท ญาณํ อภิญฺเญยฺยํ. (8)}

\proseitem{20}{ฉนฺโท มูลฏฺเฐน อภิญฺเญยฺโย, มนสิกาโร สมุฏฺฐาน\-ฏฺเฐน อภิญฺเญยฺโย, ผสฺโส สโมธาน\-ฏฺเฐน อภิญฺเญยฺโย, เวทนา สโมสรณ\-ฏฺเฐน อภิญฺเญยฺยา, สมาธิ ปมุขฏฺเฐน อภิญฺเญยฺโย, สติ อาธิปเตยฺย\-ฏฺเฐน อภิญฺเญยฺยา, ปญฺญา ตทุตฺต\-รฏฺเฐน อภิญฺเญยฺยา, วิมุตฺติ สารฏฺเฐน อภิญฺเญยฺยา, อมโตคธํ นิพฺพานํ ปริโยสาน\-ฏฺเฐน อภิญฺเญยฺยํ. (9)}

\prose{เย เย ธมฺมา อภิญฺญาตา โหนฺติ, เต เต ธมฺมา ญาตา โหนฺติ, ตํญาตฏฺเฐน ญาณํ, ปชานนฏฺเฐน ปญฺญา, เตน วุจฺจติ “อิเม ธมฺมา อภิญฺเญยฺยาติ โสตาวธานํ, ตํปชานนา ปญฺญา สุตมเย ญาณนฺ”ติ. (1)}

\vspace{\tipitakaparskip}
\csromancenter{ทุติย\-ภาณวาโร.}
\proseitem{21}{\csromanfolio{23}กถํ “อิเม ธมฺมา ปริญฺเญยฺยา”ติ โสตาวธานํ, ตํปชานนา ปญฺญา สุตมเย ญาณํ\csromandash{}}

\prose{เอโก ธมฺโม ปริญฺเญยฺโย, ผสฺโส สาสโว อุปาทานิโย. เทฺว ธมฺมา ปริญฺเญยฺยา, นามญฺจ รูปญฺจ. ตโย ธมฺมา ปริญฺเญยฺยา, ติสฺโส เวทนา. จตฺตาโร ธมฺมา ปริญฺเญยฺยา, จตฺตาโร อาหารา. ปญฺจ ธมฺมา ปริญฺเญยฺยา, ปญฺจุ\-ปาทาน\-กฺขนฺธา. ฉ ธมฺมา ปริญฺเญยฺยา, ฉ อชฺฌตฺติกานิ อายตนานิ. สตฺต ธมฺมา ปริญฺเญยฺยา, สตฺต วิญฺญาณ\-ฏฺฐิติโย. อฏฺฐ ธมฺมา ปริญฺเญยฺยา, อฏฺฐ โลกธมฺมา. นว ธมฺมา ปริญฺเญยฺยา, นว สตฺตาวาสา. ทส ธมฺมา ปริญฺเญยฺยา, ทสายตนานิ. (10)}

\prose{สพฺพํ ภิกฺขเว ปริญฺเญยฺยํ. กิญฺจ ภิกฺขเว สพฺพํ ปริญฺเญยฺยํ? จกฺขุ ภิกฺขเว ปริญฺเญยฺยํ, รูปา ปริญฺเญยฺยา, จกฺขุวิญฺญาณํ ปริญฺเญยฺยํ, จกฺขุ\-สมฺผสฺโส ปริญฺเญยฺโย, ยมฺปิทํ จกฺขุ\-สมฺผสฺส\-ปจฺจยา อุปฺปชฺชติ เวทยิตํ สุขํ วา ทุกฺขํ วา อทุกฺขมสุขํ วา, ตมฺปิ ปริญฺเญยฺยํ. โสตํ ปริญฺเญยฺยํ, สทฺทา ปริญฺเญยฺยา ฯเปฯ. ฆานํ ปริญฺเญยฺยํ, คนฺธา ปริญฺเญยฺยา. ชิวฺหา ปริญฺเญยฺยา, รสา ปริญฺเญยฺยา. กาโย ปริญฺเญยฺโย, โผฏฺฐพฺพา ปริญฺเญยฺยา ฯเปฯ. มโน ปริญฺเญยฺโย, ธมฺมา ปริญฺเญยฺยา. มโนวิญฺญาณํ ปริญฺเญยฺยํ, มโนสมฺผสฺโส ปริญฺเญยฺโย. ยมฺปิทํ มโน\-สมฺผสฺส\-ปจฺจยา อุปฺปชฺชติ เวทยิตํ สุขํ วา ทุกฺขํ วา อทุกฺขมสุขํ วา, ตมฺปิ ปริญฺเญยฺยํ.  รูปํ ปริญฺเญยฺยํ. เวทนา ปริญฺเญยฺยา. สญฺญา ปริญฺเญยฺยา. สงฺขารา ปริญฺเญยฺยา. วิญฺญาณํ ปริญฺเญยฺยํ. จกฺขุ ปริญฺเญยฺยํ ฯเปฯ. ชรามรณํ ปริญฺเญยฺยํ ฯเปฯ อมโตคธํ นิพฺพานํ ปริโยสาน\-ฏฺเฐน ปริญฺเญยฺยํ. เยสํ เยสํ ธมฺมานํ ปฏิลาภ\-ตฺถาย วายมนฺตสฺส, เต เต ธมฺมา ปฏิลทฺธา โหนฺติ. เอวํ เต ธมฺมา ปริญฺญาตา เจว โหนฺติ ตีริตา จ.}

\proseitem{22}{เนกฺขมฺมํ ปฏิลาภ\-ตฺถาย วายมนฺตสฺส เนกฺขมฺมํ ปฏิลทฺธํ โหติ, เอวํ โส ธมฺโม ปริญฺญาโต เจว โหติ ตีริโต จ. อพฺยาปาทํ ปฏิลาภ\-ตฺถาย วายมนฺตสฺส อพฺยาปาโท ปฏิลทฺโธ โหติ, เอวํ โส ธมฺโม ปริญฺญาโต เจว โหติ ตีริโต จ. อาโลกสญฺญํ ปฏิลาภ\-ตฺถาย วายมนฺตสฺส อาโลกสญฺญา ปฏิลทฺธา โหติ, เอวํ โส ธมฺโม ปริญฺญาโต เจว โหติ ตีริโต จ. อวิกฺเขปํ ปฏิลาภ\-ตฺถาย วายมนฺตสฺส อวิกฺเขโป ปฏิลทฺโธ โหติ, เอวํ โส \csromanfolio{24}ธมฺโม ปริญฺญาโต เจว โหติ ตีริโต จ. ธมฺม\-ววตฺถานํ ปฏิลาภ\-ตฺถาย วายมนฺตสฺส ธมฺม\-ววตฺถานํ ปฏิลทฺธํ โหติ, เอวํ โส ธมฺโม ปริญฺญาโต เจว โหติ ตีริโต จ. ญาณํ ปฏิลาภ\-ตฺถาย วายมนฺตสฺส ญาณํ ปฏิลทฺธํ โหติ, เอวํ โส ธมฺโม ปริญฺญาโต เจว โหติ ตีริโต จ. ปาโมชฺชํ ปฏิลาภ\-ตฺถาย วายมนฺตสฺส ปาโมชฺชํ ปฏิลทฺธํ โหติ, เอวํ โส ธมฺโม ปริญฺญาโต เจว โหติ ตีริโต จ. (7)}

\prose{ปฐมํ ฌานํ ปฏิลาภ\-ตฺถาย วายมนฺตสฺส ปฐมํ ฌานํ ปฏิลทฺธํ โหติ, เอวํ โส ธมฺโม ปริญฺญาโต เจว โหติ ตีริโต จ. ทุติยํ ฌานํ ฯเปฯ. ตติยํ ฌานํ. จตุตฺถํ ฌานํ ปฏิลาภ\-ตฺถาย วายมนฺตสฺส จตุตฺถํ ฌานํ ปฏิลทฺธํ โหติ, เอวํ โส ธมฺโม ปริญฺญาโต เจว โหติ ตีริโต จ. อากาสา\-นญฺจา\-ยตนสมาปตฺติํ ปฏิลาภ\-ตฺถาย วายมนฺตสฺส อากาสา\-นญฺจา\-ยตนสมาปตฺติ ปฏิลทฺธา โหติ, เอวํ โส ธมฺโม ปริญฺญาตฺโต เจว โหติ ตีริโต จ. วิญฺญาณญฺจา\-ยตนสมาปตฺติํ ปฏิลาภ\-ตฺถาย วายมนฺตสฺส วิญฺญาณญฺจา\-ยตนสมาปตฺติ ปฏิลทฺธา โหติ, เอวํ โส ธมฺโม ปริญฺญาโต เจว โหติ ตีริโต จ. อากิญฺจญฺญายตน\-สมาปตฺติํ ปฏิลาภ\-ตฺถาย วายมนฺตสฺส อากิญฺจญฺญายตน\-สมาปตฺติ ปฏิลทฺธา โหติ, เอวํ โส ธมฺโม ปริญฺญาโต เจว โหติ ตีริโต จ. เนว\-สญฺญา\-นา\-สญฺญา\-ยตนสมาปตฺติํ ปฏิลาภ\-ตฺถาย วายมนฺตสฺส เนว\-สญฺญา\-นา\-สญฺญา\-ยตนสมาปตฺติ ปฏิลทฺธา โหติ, เอวํ โส ธมฺโม ปริญฺญาโต เจว โหติ ตีริโต จ. (8)}

\prose{อนิจฺจา\-นุปสฺสนํ ปฏิลาภ\-ตฺถาย วายมนฺตสฺส อนิจฺจา\-นุปสฺสนา ปฏิลทฺธา โหติ, เอวํ โส ธมฺโม ปริญฺญาโต เจว โหติ ตีริโต จ. ทุกฺขา\-นุปสฺสนํ ปฏิลาภ\-ตฺถาย วายมนฺตสฺส ทุกฺขา\-นุปสฺสนา ปฏิลทฺธา โหติ, เอวํ โส ธมฺโม ปริญฺญาโต เจว โหติ ตีริโต จ. อนตฺตา\-นุปสฺสนํ ปฏิลาภ\-ตฺถาย วายมนฺตสฺส อนตฺตา\-นุปสฺสนา ปฏิลทฺธา โหติ, เอวํ โส ธมฺโม ปริญฺญาโต เจว โหติ ตีริโต จ.  นิพฺพิทา\-นุปสฺสนํ ปฏิลาภ\-ตฺถาย วายมนฺตสฺส นิพฺพิทา\-นุปสฺสนา ปฏิลทฺธา โหติ, เอวํ โส ธมฺโม ปริญฺญาโต เจว โหติ ตีริโต จ. วิราคา\-นุปสฺสนํ ปฏิลาภ\-ตฺถาย วายมนฺตสฺส วิราคา\-นุปสฺสนา ปฏิลทฺธา โหติ, เอวํ โส ธมฺโม ปริญฺญาโต เจว โหติ ตีริโต จ. \csromanfolio{25}นิโรธา\-นุปสฺสนํ ปฏิลาภ\-ตฺถาย วายมนฺตสฺส นิโรธา\-นุปสฺสนา ปฏิลทฺธา โหติ, เอวํ โส ธมฺโม ปริญฺญาโต เจว โหติ ตีริโต จ. ปฏินิสฺสคฺคา\-นุปสฺสนํ ปฏิลาภ\-ตฺถาย วายมนฺตสฺส ปฏินิสฺสคฺคา\-นุปสฺสนา ปฏิลทฺธา โหติ, เอวํ โส ธมฺโม ปริญฺญาโต เจว โหติ ตีริโต จ.  ขยานุปสฺสนํ ปฏิลาภ\-ตฺถาย วายมนฺตสฺส ขยานุปสฺสนา ปฏิลทฺธา โหติ, เอวํ โส ธมฺโม ปริญฺญาโต เจว โหติ ตีริโต จ. วยานุปสฺสนํ ปฏิลาภ\-ตฺถาย วายมนฺตสฺส วยานุปสฺสนา ปฏิลทฺธา โหติ, เอวํ โส ธมฺโม ปริญฺญาโต เจว โหติ ตีริโต จ. วิปริณามา\-นุปสฺสนํ ปฏิลาภ\-ตฺถาย วายมนฺตสฺส วิปริณามา\-นุปสฺสนา ปฏิลทฺธา โหติ, เอวํ โส ธมฺโม ปริญฺญาโต เจว โหติ ตีริโต จ.  อนิมิตฺตา\-นุปสฺสนํ ปฏิลาภ\-ตฺถาย วายมนฺตสฺส อนิมิตฺตา\-นุปสฺสนา ปฏิลทฺธา โหติ, เอวํ โส ธมฺโม ปริญฺญาโต เจว โหติ ตีริโต จ. อปฺปณิหิตา\-นุปสฺสนํ ปฏิลาภ\-ตฺถาย วายมนฺตสฺส อปฺปณิหิตา\-นุปสฺสนา ปฏิลทฺธา โหติ, เอวํ โส ธมฺโม ปริญฺญาโต เจว โหติ ตีริโต จ. สุญฺญตา\-นุปสฺสนํ ปฏิลาภ\-ตฺถาย วายมนฺตสฺส สุญฺญตา\-นุปสฺสนา ปฏิลทฺธา โหติ, เอวํ โส ธมฺโม ปริญฺญาโต เจว โหติ ตีริโต จ.}

\prose{อธิปญฺญา\-ธมฺมวิปสฺสนํ ปฏิลาภ\-ตฺถาย วายมนฺตสฺส อธิปญฺญา\-ธมฺมวิปสฺสนา ปฏิลทฺธา โหติ, เอวํ โส ธมฺโม ปริญฺญาโต เจว โหติ ตีริโต จ. ยถา\-ภูต\-ญาณ\-ทสฺสนํ ปฏิลาภ\-ตฺถาย วายมนฺตสฺส ยถา\-ภูต\-ญาณ\-ทสฺสนํ ปฏิลทฺธํ โหติ, เอวํ โส ธมฺโม ปริญฺญาโต เจว โหติ ตีริโต จ. อาทีนวา\-นุปสฺสนํ ปฏิลาภ\-ตฺถาย วายมนฺตสฺส อาทีนวา\-นุปสฺสนา ปฏิลทฺธา โหติ, เอวํ โส ธมฺโม ปริญฺญาโต เจว โหติ ตีริโต จ. ปฏิสงฺขา\-นุปสฺสนํ ปฏิลาภ\-ตฺถาย วายมนฺตสฺส ปฏิสงฺขา\-นุปสฺสนา ปฏิลทฺธา โหติ, เอวํ โส ธมฺโม ปริญฺญาโต เจว โหติ ตีริโต จ. วิวฏฺฏนา\-นุปสฺสนํ ปฏิลาภ\-ตฺถาย วายมนฺตสฺส วิวฏฺฏนา\-นุปสฺสนา ปฏิลทฺธา โหติ, เอวํ โส ธมฺโม ปริญฺญาโต เจว โหติ ตีริโต จ. (18)}

\prose{โสตาปตฺติ\-มคฺคํ ปฏิลาภ\-ตฺถาย วายมนฺตสฺส โสตาปตฺติมคฺโค ปฏิลทฺโธ โหติ. เอวํ โส ธมฺโม ปริญฺญาโต เจว โหติ ตีริโต จ. สกทาคามิ\-มคฺคํ ปฏิลาภ\-ตฺถาย วายมนฺตสฺส สกทาคามิ\-มคฺโค ปฏิลทฺโธ โหติ, เอวํ โส ธมฺโม ปริญฺญาโต เจว โหติ ตีริโต จ. \csromanfolio{26}อนาคามิมคฺคํ ปฏิลาภ\-ตฺถาย วายมนฺตสฺส อนาคามิมคฺโค ปฏิลทฺโธ โหติ, เอวํ โส ธมฺโม ปริญฺญาโต เจว โหติ ตีริโต จ. อรหตฺตมคฺคํ ปฏิลาภ\-ตฺถาย วายมนฺตสฺส อรหตฺตมคฺโค ปฏิลทฺโธ โหติ, เอวํ โส ธมฺโม ปริญฺญาโต เจว โหติ ตีริโต จ. (4)}

\prose{เยสํ เยสํ ธมฺมานํ ปฏิลาภ\-ตฺถาย วายมนฺตสฺส, เต เต ธมฺมา ปฏิลทฺธา โหนฺติ. เอวํ เต ธมฺมา ปริญฺญาตา เจว โหนฺติ ตีริตา จ. ตํญาตฏฺเฐน ญาณํ, ปชานนฏฺเฐน ปญฺญา, เตน วุจฺจติ “อิเม ธมฺมา ปริญฺเญยฺยาติ โสตาวธานํ, ตํปชานนา ปญฺญา สุตมเย ญาณนฺ”ติ. (2)}

\proseitem{23}{กถํ “อิเม ธมฺมา ปหาตพฺพา”ติ โสตาวธานํ, ตํปชานนา ปญฺญา สุตมเย ญาณํ\csromandash{}}

\prose{เอโก ธมฺโม ปหาตพฺโพ, อสฺมิมาโน.  เทฺว ธมฺมา ปหาตพฺพา, อวิชฺชา จ ภวตณฺหา จ.  ตโย ธมฺมา ปหาตพฺพา, ติสฺโส ตณฺหา.  จตฺตาโร ธมฺมา ปหาตพฺพา, จตฺตาโร โอฆา.  ปญฺจ ธมฺมา ปหาตพฺพา, ปญฺจ นีวรณานิ.  ฉ ธมฺมา ปหาตพฺพา, ฉ ตณฺหากายา.  สตฺต ธมฺมา ปหาตพฺพา, สตฺตานุสยา.  อฏฺฐ ธมฺมา ปหาตพฺพา, อฏฺฐ มิจฺฉตฺตา.  นว ธมฺมา ปหาตพฺพา, นว ตณฺหามูลกา.  ทส ธมฺมา ปหาตพฺพา, ทส มิจฺฉตฺตา.}

\proseitem{24}{เทฺว ปหานานิ สมุจฺเฉท\-ปฺปหานํ, ปฏิปฺปสฺสทฺธิ\-ปฺปหานํ. สมุจฺเฉทปฺปหานญฺจ โลกุตฺตรํ ขยคามิ\-มคฺคํ ภาวยโต, ปฏิปฺปสฺสทฺธิ\-ปฺปหานญฺจ ผลกฺขเณ.  ตีณิ ปหานานิ, กามานเมตํ นิสฺสรณํ ยทิทํ เนกฺขมฺมํ, รูปานเมตํ นิสฺสรณํ ยทิทํ อารุปฺปํ. ยํ โข ปน กิญฺจิ ภูตํ สงฺขตํ ปฏิจฺจ\-สมุปฺปนฺนํ นิโรโธ, ตสฺส นิสฺสรณํ. เนกฺขมฺมํ ปฏิลทฺธสฺส กามา ปหีนา เจว โหนฺติ ปริจฺจตฺตา จ, อารุปฺปํ ปฏิลทฺธสฺส รูปา ปหีนา เจว โหนฺติ ปริจฺจตฺตา จ, นิโรธํ ปฏิลทฺธสฺส สงฺขารา ปหีนา เจว โหนฺติ ปริจฺจตฺตา จ.  จตฺตาริ ปหานานิ, ทุกฺขสจฺจํ ปริญฺญา\-ปฏิเวธํ ปฏิวิชฺฌนฺโต ปชหาติ, สมุทยสจฺจํ ปหาน\-ปฏิเวธํ ปฏิวิชฺฌนฺโต ปชหาติ, นิโรธสจฺจํ สจฺฉิกิริยา\-ปฏิเวธํ ปฏิวิชฺฌนฺโต ปชหาติ, มคฺคสจฺจํ ภาวนา\-ปฏิเวธํ ปฏิวิชฺฌนฺโต ปชหาติ.  ปญฺจ ปหานานิ วิกฺขมฺภน\-ปฺปหานํ ตทงฺค\-ปฺปหานํ สมุจฺเฉท\-ปฺปหานํ ปฏิปฺปสฺสทฺธิ\-ปฺปหานํ นิสฺสรณ\-ปฺปหานํ. วิกฺขมฺภนปฺปหานญฺจ นีวรณานํ ปฐมํ ฌานํ ภาวยโต, ตทงฺคปฺปหานญฺจ ทิฏฺฐิคตานํ \csromanfolio{27}นิพฺเพธ\-ภาคิยํ สมาธิํ ภาวยโต, สมุจฺเฉทปฺปหานญฺจ โลกุตฺตรํ ขยคามิ\-มคฺคํ ภาวยโต, ปฏิปฺปสฺสทฺธิ\-ปฺปหานญฺจ ผลกฺขเณ, นิสฺสรณ\-ปฺปหานญฺจ นิโรโธ นิพฺพานํ.}

\prose{สพฺพํ ภิกฺขเว ปหาตพฺพํ. กิญฺจ ภิกฺขเว สพฺพํ ปหาตพฺพํ? จกฺขุ ภิกฺขเว ปหาตพฺพํ, รูปา ปหาตพฺพา, จกฺขุวิญฺญาณํ ปหาตพฺพํ, จกฺขุ\-สมฺผสฺโส ปหาตพฺโพ. ยมฺปิทํ จกฺขุ\-สมฺผสฺส\-ปจฺจยา อุปฺปชฺชติ เวทยิตํ สุขํ วา ทุกฺขํ วา อทุกฺขมสุขํ วา, ตมฺปิ ปหาตพฺพํ. โสตํ ปหาตพฺพํ, สทฺทา ปหาตพฺพา ฯเปฯ. ฆานํ ปหาตพฺพํ, คนฺธา ปหาตพฺพา. ชิวฺหา ปหาตพฺพา, รสา ปหาตพฺพา. กาโย ปหาตพฺโพ, โผฏฺฐพฺพา ปหาตพฺพา. มโน ปหาตพฺโพ, ธมฺมา ปหาตพฺพา. มโนวิญฺญาณํ ปหาตพฺพํ, มโนสมฺผสฺโส ปหาตพฺโพ. ยมฺปิทํ มโน\-สมฺผสฺส\-ปจฺจยา อุปฺปชฺชติ เวทยิตํ สุขํ วา ทุกฺขํ วา อทุกฺขมสุขํ วา, ตมฺปิ ปหาตพฺพํ. รูปํ ปสฺสนฺโต ปชหาติ, เวทนํ ปสฺสนฺโต ปชหาติ, สญฺญํ ปสฺสนฺโต ปชหาติ, สงฺขาเร ปสฺสนฺโต ปชหาติ, วิญฺญาณํ ปสฺสนฺโต ปชหาติ. จกฺขุํ ฯเปฯ. ชรามรณํ ฯเปฯ อมโตคธํ นิพฺพานํ ปริโยสาน\-ฏฺเฐน ปสฺสนฺโต ปชหาติ. เย เย ธมฺมา ปหีนา โหนฺติ, เต เต ธมฺมา ปริจฺจตฺตา โหนฺติ, ตํญาตฏฺเฐน ญาณํ, ปชานนฏฺเฐน ปญฺญา, เตน วุจฺจติ “อิเม ธมฺมา ปหาตพฺพาติ โสตาวธานํ, ตํปชานนา ปญฺญา สุตมเย ญาณนฺ”ติ. (3) ตติย\-ภาณวาโร.}

\proseitem{25}{กถํ “อิเม ธมฺมา ภาเวตพฺพา”ติ โสตาวธานํ, ตํปชานนา ปญฺญา สุตมเย ญาณํ\csromandash{}}

\prose{เอโก ธมฺโม ภาเวตพฺโพ, กายคตาสติ สาตสหคตา.  เทฺว ธมฺมา ภาเวตพฺพา, สมโถ จ วิปสฺสนา จ.  ตโย ธมฺมา ภาเวตพฺพา, ตโย สมาธี.  จตฺตาโร ธมฺมา ภาเวตพฺพา, จตฺตาโร สติปฏฺฐานา.  ปญฺจ ธมฺมา ภาเวตพฺพา, ปญฺจงฺคิโก สมาธิ\csromansharedfootnote{c51205a2095d6cff}{สมฺมาสมาธิ (สฺยา)}.  ฉ ธมฺมา ภาเวตพฺพา, ฉ อนุสฺสติ\-ฏฺฐานานิ.  สตฺต ธมฺมา ภาเวตพฺพา, สตฺต โพชฺฌงฺคา.  อฏฺฐ ธมฺมา ภาเวตพฺพา, อริโย อฏฺฐงฺคิโก มคฺโค.  นว ธมฺมา ภาเวตพฺพา, นว ปาริสุทฺธิ\-ปธานิย\-งฺคานิ.  ทส ธมฺมา ภาเวตพฺพา, ทส กสิณายตนานิ.}

\proseitem{26}{\csromanfolio{28}เทฺว ภาวนา โลกิยา จ ภาวนา, โลกุตฺตรา จ ภาวนา. ติสฺโส ภาวนา, รูปาวจร\-กุสลานํ ธมฺมานํ ภาวนา อรูปาวจร\-กุสลานํ ธมฺมานํ ภาวนา อปริยาปนฺน\-กุสลานํ ธมฺมานํ ภาวนา. รูปาวจร\-กุสลานํ ธมฺมานํ ภาวนา อตฺถิ หีนา, อตฺถิ มชฺฌา อตฺถิ ปณีตา, อรูปาวจร\-กุสลานํ ธมฺมานํ ภาวนา อตฺถิ หีนา, อตฺถิ มชฺฌา, อตฺถิ ปณีตา, อปริยาปนฺน\-กุสลานํ ธมฺมานํ ภาวนา ปณีตา.}

\proseitem{27}{จตสฺโส ภาวนา, ทุกฺขสจฺจํ ปริญฺญา\-ปฏิเวธํ ปฏิวิชฺฌนฺโต ภาเวติ, สมุทยสจฺจํ ปหาน\-ปฺปฏิเวธํ ปฏิวิชฺฌนฺโต ภาเวติ, นิโรธสจฺจํ สจฺฉิกิริยา\-ปฏิเวธํ ปฏิวิชฺฌนฺโต ภาเวติ, มคฺคสจฺจํ ภาวนา\-ปฏิเวธํ ปฏิวิชฺฌนฺโต ภาเวติ, อิมา จตสฺโส ภาวนา. (4)}

\prose{อปราปิ จตสฺโส ภาวนา, เอสนาภาวนา, ปฏิลาภ\-ภาวนา, เอกรสภาวนา, อาเสวนาภาวนา. กตมา เอสนาภาวนา? สพฺเพสํ สมาธิํ สมา\-ป\-ชฺช\-นฺตานํ ตตฺถ ชาตา ธมฺมา เอกรสา โหนฺตีติ อยํ เอสนาภาวนา.  กตมา ปฏิลาภ\-ภาวนา? สพฺเพสํ สมาธิํ สมาปนฺนานํ ตตฺถ ชาตา ธมฺมา อญฺญมญฺญํ นาติวตฺตนฺตีติ อยํ ปฏิลาภ\-ภาวนา.  กตมา เอกรสภาวนา? อธิโมกฺข\-ฏฺเฐน สทฺธินฺทฺริยํ ภาวยโต สทฺธิ\-นฺทฺริยสฺส วเสน จตฺตาริ อินฺทฺริยานิ เอกรสา โหนฺตีติ อินฺทฺริยานํ เอกรสฏฺเฐน ภาวนา, ปคฺคหฏฺเฐน วีริยินฺทฺริยํ ภาวยโต วีริยิ\-นฺทฺริยสฺส วเสน จตฺตาริ อินฺทฺริยานิ เอกรสา โหนฺตีติ อินฺทฺริยานํ เอกรสฏฺเฐน ภาวนา. อุปฏฺฐาน\-ฏฺเฐน สตินฺทฺริยํ ภาวยโต สตินฺทฺริยสฺส วเสน จตฺตาริ อินฺทฺริยานิ เอกรสา โหนฺตีติ อินฺทฺริยานํ เอกรสฏฺเฐน ภาวนา. อวิกฺเขป\-ฏฺเฐน สมาธินฺทฺริยํ ภาวยโต สมาธิ\-นฺทฺริยสฺส วเสน จตฺตาริ อินฺทฺริยานิ เอกรสา โหนฺตีติ อินฺทฺริยานํ เอกรสฏฺเฐน ภาวนา, ทสฺสนฏฺเฐน ปญฺญินฺทฺริยํ ภาวยโต ปญฺญินฺทฺริยสฺส วเสน จตฺตาริ อินฺทฺริยานิ เอกรสา โหนฺตีติ อินฺทฺริยานํ เอกรสฏฺเฐน ภาวนา.}

\prose{อสฺสทฺธิเย อกมฺปิยฏฺเฐน สทฺธาพลํ ภาวยโต สทฺธาพลสฺส วเสน จตฺตาริ พลานิ เอกรสา โหนฺตีติ พลานํ เอกรสฏฺเฐน ภาวนา, โกสชฺเช อกมฺปิยฏฺเฐน วีริยพลํ ภาวยโต วีริยพลสฺส วเสน \csromanfolio{29}จตฺตาริ พลานิ เอกรสา โหนฺตีติ พลานํ เอกรสฏฺเฐน ภาวนา, ปมาเท อกมฺปิยฏฺเฐน สติพลํ ภาวยโต สติพลสฺส วเสน จตฺตาริ พลานิ เอกรสา โหนฺตีติ พลานํ เอกรสฏฺเฐน ภาวนา, อุทฺธจฺเจ อกมฺปิยฏฺเฐน สมาธิพลํ ภาวยโต สมาธิพลสฺส วเสน จตฺตาริ พลานิ เอกรสา โหนฺตีติ พลานํ เอกรสฏฺเฐน ภาวนา, อวิชฺชาย อกมฺปิยฏฺเฐน ปญฺญาพลํ ภาวยโต ปญฺญาพลสฺส วเสน จตฺตาริ พลานิ เอกรสา โหนฺตีติ พลานํ เอกรสฏฺเฐน ภาวนา.}

\prose{อุปฏฺฐาน\-ฏฺเฐน สติสมฺโพชฺฌงฺคํ ภาวยโต สติสมฺโพชฺฌงฺคสฺส วเสน ฉ โพชฺฌงฺคา เอกรสา โหนฺตีติ โพชฺฌงฺคานํ เอกรสฏฺเฐน ภาวนา, ปวิจยฏฺเฐน ธมฺมวิจย\-สมฺโพชฺฌงฺคํ ภาวยโต ธมฺมวิจย\-สมฺโพชฺฌงฺคสฺส วเสน ฉ โพชฺฌงฺคา เอกรสา โหนฺตีติ โพชฺฌงฺคานํ เอกรสฏฺเฐน ภาวนา, ปคฺคหฏฺเฐน วีริย\-สมฺโพชฺฌงฺคํ ภาวยโต วีริย\-สมฺโพชฺฌงฺคสฺส วเสน ฉ โพชฺฌงฺคา เอกรสา โหนฺตีติ โพชฺฌงฺคานํ เอกรสฏฺเฐน ภาวนา, ผรณฏฺเฐน ปีติสมฺโพชฺฌงฺคํ ภาวยโต ปีติสมฺโพชฺฌงฺคสฺส วเสน ฉ โพชฺฌงฺคา เอกรสา โหนฺตีติ โพชฺฌงฺคานํ เอกรสฏฺเฐน ภาวนา, อุปสมฏฺเฐน ปสฺสทฺธิ\-สมฺโพชฺฌงฺคํ ภาวยโต ปสฺสทฺธิ\-สมฺโพชฺฌงฺคสฺส วเสน ฉ โพชฺฌงฺคา เอกรสา โหนฺตีติ โพชฺฌงฺคานํ เอกรสฏฺเฐน ภาวนา, อวิกฺเขป\-ฏฺเฐน สมาธิ\-สมฺโพชฺฌงฺคํ ภาวยโต สมาธิ\-สมฺโพชฺฌงฺคสฺส วเสน ฉ โพชฺฌงฺคา เอกรสา โหนฺตีติ โพชฺฌงฺคานํ เอกรสฏฺเฐน ภาวนา, ปฏิสงฺขา\-นฏฺเฐน อุเปกฺขา\-สมฺโพชฺฌงฺคํ ภาวยโต อุเปกฺขา\-สมฺโพชฺฌงฺคสฺส วเสน ฉ โพชฺฌงฺคา เอกรสา โหนฺตีติ โพชฺฌงฺคานํ เอกรสฏฺเฐน ภาวนา.}

\prose{ทสฺสนฏฺเฐน สมฺมาทิฏฺฐิํ ภาวยโต สมฺมาทิฏฺฐิยา วเสน สตฺต มคฺคงฺคา เอกรสา โหนฺตีติ มคฺคงฺคานํ เอกรสฏฺเฐน ภาวนา, อภินิโรปน\-ฏฺเฐน สมฺมาสงฺกปฺปํ ภาวยโต สมฺมา\-สงฺกปฺปสฺส วเสน สตฺต มคฺคงฺคา เอกรสา โหนฺตีติ มคฺคงฺคานํ เอกรสฏฺเฐน ภาวนา, ปริคฺคห\-ฏฺเฐน สมฺมาวาจํ ภาวยโต สมฺมาวาจาย วเสน สตฺต มคฺคงฺคา เอกรสา โหนฺตีติ มคฺคงฺคานํ เอกรสฏฺเฐน ภาวนา, สมุฏฺฐาน\-ฏฺเฐน สมฺมากมฺมนฺตํ ภาวยโต สมฺมา\-กมฺมนฺตสฺส วเสน สตฺต มคฺคงฺคา เอกรสา โหนฺตีติ มคฺคงฺคานํ เอกรสฏฺเฐน ภาวนา, โวทานฏฺเฐน สมฺมาอาชีวํ ภาวยโต สมฺมาอาชีวสฺส วเสน สตฺต มคฺคงฺคา เอกรสา โหนฺตีติ มคฺคงฺคานํ เอกรสฏฺเฐน \csromanfolio{30}ภาวนา, ปคฺคหฏฺเฐน สมฺมาวายามํ ภาวยโต สมฺมาวายามสฺส วเสน สตฺต มคฺคงฺคา เอกรสา โหนฺตีติ มคฺคงฺคานํ เอกรสฏฺเฐน ภาวนา, อุปฏฺฐาน\-ฏฺเฐน สมฺมาสติํ ภาวยโต สมฺมาสติยา วเสน สตฺต มคฺคงฺคา เอกรสา โหนฺตีติ มคฺคงฺคานํ เอกรสฏฺเฐน ภาวนา, อวิกฺเขป\-ฏฺเฐน สมฺมาสมาธิํ ภาวยโต สมฺมา\-สมาธิสฺส วเสน สตฺต มคฺคงฺคา เอกรสา โหนฺตีติ มคฺคงฺคานํ เอกรสฏฺเฐน ภาวนา. อยํ เอกรสาภาวนา.}

\prose{กตมา อาเสวนาภาวนา? อิธ ภิกฺขุ ปุพฺพณฺหสมยํ อาเสวติ, มชฺฌนฺหิก\-สมยมฺปิ\csromansharedfootnote{79d2af287168d6a8}{มชฺฌนฺติกสมยมฺปิ (สฺยา, ก)} อาเสวติ, สายนฺห\-สมยมฺปิ อาเสวติ, ปุเรภตฺตมฺปิ อาเสวติ, ปจฺฉา\-ภตฺตมฺปิ อาเสวติ, ปุริเมปิ ยาเม อาเสวติ, มชฺฌิเมปิ ยาเม อาเสวติ, ปจฺฉิเมปิ ยาเม อาเสวติ, รตฺติมฺปิ อาเสวติ, ทิวาปิ อาเสวติ, รตฺตินฺทิวาปิ\csromansharedfootnote{246dcdda4a1e6de4}{รตฺติทิวาปิ (ก)} อาเสวติ, กาเฬปิ อาเสวติ, ชุเณฺหปิ อาเสวติ, วสฺเสปิ อาเสวติ, เหมนฺเตปิ อาเสวติ, คิเมฺหปิ อาเสวติ, ปุริเมปิ วโยขนฺเธ อาเสวติ, มชฺฌิเมปิ วโยขนฺเธ อาเสวติ, ปจฺฉิเมปิ วโยขนฺเธ อาเสวติ. อยํ อาเสวนาภาวนา. อิมา จตสฺโส ภาวนา.}

\proseitem{28}{อปราปิ จตสฺโส ภาวนา, ตตฺถ ชาตานํ ธมฺมานํ อนติวตฺตน\-ฏฺเฐน ภาวนา, อินฺทฺริยานํ เอกรสฏฺเฐน ภาวนา, ตทุปค\-วีริย\-วาห\-นฏฺเฐน ภาวนา, อาเสวนฏฺเฐน ภาวนา. กถํ ตตฺถ ชาตานํ ธมฺมานํ อนติวตฺตน\-ฏฺเฐน ภาวนา? กามจฺฉนฺทํ ปชหโต เนกฺขมฺม\-วเสน ชาตา ธมฺมา อญฺญมญฺญํ นาติวตฺตนฺตีติ ตตฺถ ชาตานํ ธมฺมานํ อนติวตฺตน\-ฏฺเฐน ภาวนา, พฺยาปาทํ ปชหโต อพฺยาปาทวเสน ชาตา ธมฺมา อญฺญมญฺญํ นาติวตฺตนฺตีติ ตตฺถ ชาตานํ ธมฺมานํ อนติวตฺตน\-ฏฺเฐน ภาวนา, ถินมิทฺธํ\csromansharedfootnote{371f7b46f67e77bd}{ถีนมิทฺธํ (สฺยา, สีฏฺฐ)} ปชหโต อาโลกสญฺญา\-วเสน ชาตา ธมฺมา อญฺญมญฺญํ นาติวตฺตนฺตีติ ตตฺถ ชาตานํ ธมฺมานํ อนติวตฺตน\-ฏฺเฐน ภาวนา, อุทฺธจฺจํ ปชหโต อวิกฺเขป\-วเสน ชาตา ธมฺมา อญฺญมญฺญํ นาติวตฺตนฺตีติ ตตฺถ ชาตานํ ธมฺมานํ อนติวตฺตน\-ฏฺเฐน ภาวนา, วิจิกิจฺฉํ ปชหโต ธมฺมววตฺถาน\-วเสน ชาตา ธมฺมา อญฺญมญฺญํ นาติวตฺตนฺตีติ ตตฺถ ชาตานํ ธมฺมานํ อนติวตฺตน\-ฏฺเฐน ภาวนา, อวิชฺชํ ปชหโต ญาณวเสน ชาตา ธมฺมา อญฺญมญฺญํ นาติวตฺตนฺตีติ ตตฺถ ชาตานํ ธมฺมานํ อนติวตฺตน\-ฏฺเฐน \csromanfolio{31}ภาวนา, อรติํ ปชหโต ปาโมชฺชวเสน ชาตา ธมฺมา อญฺญมญฺญํ นาติวตฺตนฺตีติ ตตฺถ ชาตานํ ธมฺมานํ อนติวตฺตน\-ฏฺเฐน ภาวนา, นีวรเณ ปชหโต ปฐมชฺฌาน\-วเสน ชาตา ธมฺมา อญฺญมญฺญํ นาติวตฺตนฺตีติ ตตฺถ ชาตานํ ธมฺมานํ อนติวตฺตน\-ฏฺเฐน ภาวนา, วิตกฺกวิจาเร ปชหโต ทุติยชฺฌาน\-วเสน ชาตา ธมฺมา อญฺญมญฺญํ นาติวตฺตนฺตีติ ตตฺถ ชาตานํ ธมฺมานํ อนติวตฺตน\-ฏฺเฐน ภาวนา, ปีติํ ปชหโต ตติยชฺฌาน\-วเสน ชาตา ธมฺมา อญฺญมญฺญํ นาติวตฺตนฺตีติ ตตฺถ ชาตานํ ธมฺมานํ อนติวตฺตน\-ฏฺเฐน ภาวนา, สุขทุกฺเข ปชหโต จตุตฺถชฺฌาน\-วเสน ชาตา ธมฺมา อญฺญมญฺญํ นาติวตฺตนฺตีติ ตตฺถ ชาตานํ ธมฺมานํ อนติวตฺตน\-ฏฺเฐน ภาวนา.}

\prose{รูปสญฺญํ ปฏิฆสญฺญํ นานตฺตสญฺญํ ปชหโต อากาสา\-นญฺจา\-ยตนสมาปตฺติวเสน ชาตา ธมฺมา อญฺญมญฺญํ นาติวตฺตนฺตีติ ตตฺถ ชาตานํ ธมฺมานํ อนติวตฺตน\-ฏฺเฐน ภาวนา, อากาสา\-นญฺจา\-ยตนสญฺญํ ปชหโต วิญฺญาณญฺจายตน\-สมาปตฺติวเสน ชาตา ธมฺมา อญฺญมญฺญํ นาติวตฺตนฺตีติ ตตฺถ ชาตานํ ธมฺมานํ อนติวตฺตน\-ฏฺเฐน ภาวนา, วิญฺญาณญฺจา\-ยตนสญฺญํ ปชหโต อากิญฺจญฺญายตน\-สมาปตฺติวเสน ชาตา ธมฺมา อญฺญมญฺญํ นาติวตฺตนฺตีติ ตตฺถ ชาตานํ ธมฺมานํ อนติวตฺตน\-ฏฺเฐน ภาวนา, อากิญฺจญฺญายตน\-สญฺญํ ปชหโต เนว\-สญฺญา\-นา\-สญฺญา\-ยตนสมาปตฺติวเสน ชาตา ธมฺมา อญฺญมญฺญํ นาติวตฺตนฺตีติ ตตฺถ ชาตานํ ธมฺมานํ อนติวตฺตน\-ฏฺเฐน ภาวนา.}

\prose{นิจฺจสญฺญํ ปชหโต อนิจฺจา\-นุปสฺสนาวเสน ชาตา ธมฺมา อญฺญมญฺญํ นาติวตฺตนฺตีติ ตตฺถ ชาตานํ ธมฺมานํ อนติวตฺตน\-ฏฺเฐน ภาวนา, สุขสญฺญํ ปชหโต ทุกฺขานุปสฺสนา\-วเสน ชาตา ธมฺมา อญฺญมญฺญํ นาติวตฺตนฺตีติ ตตฺถ ชาตานํ ธมฺมานํ อนติวตฺตน\-ฏฺเฐน ภาวนา, อตฺตสญฺญํ ปชหโต อนตฺตา\-นุปสฺสนาวเสน ชาตา ธมฺมา อญฺญมญฺญํ นาติวตฺตนฺตีติ ตตฺถ ชาตานํ ธมฺมานํ อนติวตฺตน\-ฏฺเฐน ภาวนา, นนฺทิํ ปชหโต นิพฺพิทา\-นุปสฺสนาวเสน ชาตา ธมฺมา อญฺญมญฺญํ นาติวตฺตนฺตีติ ตตฺถ ชาตานํ ธมฺมานํ อนติวตฺตน\-ฏฺเฐน ภาวนา, ราคํ ปชหโต วิราคา\-นุปสฺสนาวเสน ชาตา ธมฺมา อญฺญมญฺญํ นาติวตฺตนฺตีติ ตตฺถ ชาตานํ ธมฺมานํ อนติวตฺตน\-ฏฺเฐน ภาวนา, สมุทยํ ปชหโต นิโรธา\-นุปสฺสนาวเสน ชาตา ธมฺมา อญฺญมญฺญํ นาติวตฺตนฺตีติ ตตฺถ ชาตานํ ธมฺมานํ อนติวตฺตน\-ฏฺเฐน \csromanfolio{32}ภาวนา, อาทานํ ปชหโต ปฏินิสฺสคฺคา\-นุปสฺสนาวเสน ชาตา ธมฺมา อญฺญมญฺญํ นาติวตฺตนฺตีติ ตตฺถ ชาตานํ ธมฺมานํ อนติวตฺตน\-ฏฺเฐน ภาวนา, ฆนสญฺญํ ปชหโต ขยานุปสฺสนา\-วเสน ชาตา ธมฺมา อญฺญมญฺญํ นาติวตฺตนฺตีติ ตตฺถ ชาตานํ ธมฺมานํ อนติวตฺตน\-ฏฺเฐน ภาวนา, อายูหนํ ปชหโต วยา\-นุปสฺสนา\-วเสน ชาตา ธมฺมา อญฺญมญฺญํ นาติวตฺตนฺตีติ ตตฺถ ชาตานํ ธมฺมานํ อนติวตฺตน\-ฏฺเฐน ภาวนา. ธุวสญฺญํ ปชหโต วิปริณามา\-นุปสฺสนา\-วเสน ชาตา ธมฺมา อญฺญมญฺญํ นาติวตฺตนฺตีติ ตตฺถ ชาตานํ ธมฺมานํ อนติวตฺตน\-ฏฺเฐน ภาวนา, นิมิตฺตํ ปชหโต อนิมิตฺตา\-นุปสฺสนา\-วเสน ชาตา ธมฺมา อญฺญมญฺญํ นาติวตฺตนฺตีติ ตตฺถ ชาตานํ ธมฺมานํ อนติวตฺตน\-ฏฺเฐน ภาวนา, ปณิธิํ ปชหโต อปฺปณิหิตา\-นุปสฺสนา\-วเสน ชาตา ธมฺมา อญฺญมญฺญํ นาติวตฺตนฺตีติ ตตฺถ ชาตานํ ธมฺมานํ อนติวตฺตน\-ฏฺเฐน ภาวนา, อภินิเวสํ ปชหโต สุญฺญตา\-นุปสฺสนาวเสน ชาตา ธมฺมา อญฺญมญฺญํ นาติวตฺตนฺตีติ ตตฺถ ชาตานํ ธมฺมานํ อนติวตฺตน\-ฏฺเฐน ภาวนา, สาราทานา\-ภินิเวสํ ปชหโต อธิปญฺญา\-ธมฺมวิปสฺสนาวเสน ชาตา ธมฺมา อญฺญมญฺญํ นาติวตฺตนฺตีติ ตตฺถ ชาตานํ ธมฺมานํ อนติวตฺตน\-ฏฺเฐน ภาวนา, สมฺโมหา\-ภินิเวสํ ปชหโต ยถาภูตญาณ\-ทสฺสนวเสน ชาตา ธมฺมา อญฺญมญฺญํ นาติวตฺตนฺตีติ ตตฺถ ชาตานํ ธมฺมานํ อนติวตฺตน\-ฏฺเฐน ภาวนา, อาลยา\-ภินิเวสํ ปชหโต อาทีนวา\-นุปสฺสนาวเสน ชาตา ธมฺมา อญฺญมญฺญํ นาติวตฺตนฺตีติ ตตฺถ ชาตานํ ธมฺมานํ อนติวตฺตน\-ฏฺเฐน ภาวนา, อปฺปฏิสงฺขํ ปชหโต ปฏิสงฺขา\-นุปสฺสนาวเสน ชาตา ธมฺมา อญฺญมญฺญํ นาติวตฺตนฺตีติ ตตฺถ ชาตานํ ธมฺมานํ อนติวตฺตน\-ฏฺเฐน ภาวนา, สญฺโญคา\-ภินิเวสํ ปชหโต วิวฏฺฏนา\-นุปสฺสนา\-วเสน ชาตา ธมฺมา อญฺญมญฺญํ นาติวตฺตนฺตีติ ตตฺถ ชาตานํ ธมฺมานํ อนติวตฺตน\-ฏฺเฐน ภาวนา.}

\prose{ทิฏฺเฐกฏฺเฐ กิเลเส ปชหโต โสตาปตฺติมคฺค\-วเสน ชาตา ธมฺมา อญฺญมญฺญํ นาติวตฺตนฺตีติ ตตฺถ ชาตานํ ธมฺมานํ อนติวตฺตน\-ฏฺเฐน ภาวนา, โอฬาริเก กิเลเส ปชหโต สกทาคามิมคฺค\-วเสน ชาตา ธมฺมา อญฺญมญฺญํ นาติวตฺตนฺตีติ ตตฺถ ชาตานํ ธมฺมานํ อนติวตฺตน\-ฏฺเฐน ภาวนา, อณุสหคเต กิเลเส ปชหโต อนาคามิมคฺค\-วเสน ชาตา ธมฺมา อญฺญมญฺญํ นาติวตฺตนฺตีติ ตตฺถ ชาตานํ ธมฺมานํ อนติวตฺตน\-ฏฺเฐน ภาวนา, สพฺพกิเลเส ปชหโต อรหตฺตมคฺค\-วเสน ชาตา ธมฺมา อญฺญมญฺญํ \csromanfolio{33}นาติวตฺตนฺตีติ ตตฺถ ชาตานํ ธมฺมานํ อนติวตฺตน\-ฏฺเฐน ภาวนา, เอวํ ตตฺถ ชาตานํ ธมฺมานํ อนติวตฺตน\-ฏฺเฐน ภาวนา.}

\prose{กถํ อินฺทฺริยานํ เอกรสฏฺเฐน ภาวนา? กามจฺฉนฺทํ ปชหโต เนกฺขมฺม\-วเสน ปญฺจินฺทฺริยานิ เอกรสา โหนฺตีติ อินฺทฺริยานํ เอกรสฏฺเฐน ภาวนา, พฺยาปาทํ ปชหโต อพฺยาปาทวเสน ปญฺจินฺทฺริยานิ เอกรสา โหนฺตีติ อินฺทฺริยานํ เอกรสฏฺเฐน ภาวนา ฯเปฯ สพฺพกิเลเส ปชหโต อรหตฺตมคฺค\-วเสน ปญฺจินฺทฺริยานิ เอกรสา โหนฺตีติ อินฺทฺริยานํ เอกรสฏฺเฐน ภาวนา, เอวํ อินฺทฺริยานํ เอกรสฏฺเฐน ภาวนา.}

\prose{กถํ ตทุปค\-วีริย\-วาห\-นฏฺเฐน ภาวนา? กามจฺฉนฺทํ ปชหโต เนกฺขมฺม\-วเสน วีริยํ วาเหตีติ ตทุปค\-วีริย\-วาห\-นฏฺเฐน ภาวนา, พฺยาปาทํ ปชหโต อพฺยาปาทวเสน วีริยํ วาเหตีติ ตทุปค\-วีริย\-วาห\-นฏฺเฐน ภาวนา ฯเปฯ สพฺพกิเลเส ปชหโต อรหตฺตมคฺค\-วเสน วีริยํ วาเหตีติ ตทุปค\-วีริย\-วาห\-นฏฺเฐน ภาวนา, เอวํ ตทุปค\-วีริย\-วาห\-นฏฺเฐน ภาวนา.}

\prose{กถํ อาเสวนฏฺเฐน ภาวนา? กามจฺฉนฺทํ ปชหนฺโต เนกฺขมฺมํ อาเสวตีติ อาเสวนฏฺเฐน ภาวนา, พฺยาปาทํ ปชหนฺโต อพฺยาปาทํ อาเสวตีติ อาเสวนฏฺเฐน ภาวนา ฯเปฯ สพฺพกิเลเส ปชหนฺโต อรหตฺตมคฺคํ อาเสวตีติ อาเสวนฏฺเฐน ภาวนา, เอวํ อาเสวนฏฺเฐน ภาวนา. อิมา จตสฺโส ภาวนา.}

\prose{รูปํ ปสฺสนฺโต ภาเวติ, เวทนํ ปสฺสนฺโต ภาเวติ, สญฺญํ ปสฺสนฺโต ภาเวติ, สงฺขาเร ปสฺสนฺโต ภาเวติ, วิญฺญาณํ ปสฺสนฺโต ภาเวติ, จกฺขุํ ฯเปฯ ชรามรณํ ฯเปฯ อมโตคธํ นิพฺพานํ ปริโยสาน\-ฏฺเฐน ปสฺสนฺโต ภาเวติ. เย เย ธมฺมา ภาวิตา โหนฺติ, เต เต ธมฺมา เอกรสา โหนฺติ, ตํญาตฏฺเฐน ญาณํ, ปชานนฏฺเฐน ปญฺญา, เตน วุจฺจติ “อิเม ธมฺมา ภาเวตพฺพาติ โสตาวธานํ, ตํปชานนา ปญฺญา สุตมเย ญาณนฺ”ติ. (4)}

\vspace{\tipitakaparskip}
\csromancenter{จตุตฺถ\-ภาณวาโร.}
\proseitem{29}{\csromanfolio{34}กถํ “อิเม ธมฺมา สจฺฉิกาตพฺพา”ติ โสตาวธานํ, ตํปชานนา ปญฺญา สุตมเย ญาณํ\csromandash{}}

\prose{เอโก ธมฺโม สจฺฉิกาตพฺโพ, อกุปฺปา เจโตวิมุตฺติ.  เทฺว ธมฺมา สจฺฉิกาตพฺพา, วิชฺชา จ วิมุตฺติ จ.  ตโย ธมฺมา สจฺฉิกาตพฺพา, ติสฺโส วิชฺชา.  จตฺตาโร ธมฺมา สจฺฉิกาตพฺพา, จตฺตาริ สามญฺญผลานิ.  ปญฺจ ธมฺมา สจฺฉิกาตพฺพา, ปญฺจ ธมฺมกฺขนฺธา.  ฉ ธมฺมา สจฺฉิกาตพฺพา, ฉ อภิญฺญา.  สตฺต ธมฺมา สจฺฉิกาตพฺพา, สตฺต ขีณาสว\-พลานิ.  อฏฺฐ ธมฺมา สจฺฉิกาตพฺพา, อฏฺฐ วิโมกฺขา.  นว ธมฺมา สจฺฉิกาตพฺพา, นว อนุปุพฺพ\-นิโรธา.  ทส ธมฺมา สจฺฉิกาตพฺพา, ทส อเสกฺขา ธมฺมา.}

\prose{สพฺพํ ภิกฺขเว สจฺฉิกาตพฺพํ. กิญฺจ ภิกฺขเว สพฺพํ สจฺฉิกาตพฺพํ? จกฺขุ ภิกฺขเว สจฺฉิกาตพฺพํ, รูปา สจฺฉิกาตพฺพา, จกฺขุวิญฺญาณํ สจฺฉิกาตพฺพํ, จกฺขุ\-สมฺผสฺโส สจฺฉิกาตพฺโพ, ยมฺปิทํ จกฺขุ\-สมฺผสฺส\-ปจฺจยา อุปฺปชฺชติ เวทยิตํ สุขํ วา ทุกฺขํ วา อทุกฺขมสุขํ วา, ตมฺปิ สจฺฉิกาตพฺพํ. โสตํ สจฺฉิกาตพฺพํ, สทฺทา สจฺฉิกาตพฺพา ฯเปฯ. ฆานํ สจฺฉิกาตพฺพํ, คนฺธา สจฺฉิกาตพฺพา. ชิวฺหา สจฺฉิกาตพฺพา, รสา สจฺฉิกาตพฺพา. กาโย สจฺฉิกาตพฺโพ, โผฏฺฐพฺพา สจฺฉิกาตพฺพา. มโน สจฺฉิกาตพฺโพ, ธมฺมา สจฺฉิกาตพฺพา, มโนวิญฺญาณํ สจฺฉิกาตพฺพํ, มโนสมฺผสฺโส สจฺฉิกาตพฺโพ, ยมฺปิทํ มโน\-สมฺผสฺส\-ปจฺจยา อุปฺปชฺชติ เวทยิตํ สุขํ วา ทุกฺขํ วา อทุกฺขมสุขํ วา, ตมฺปิ สจฺฉิกาตพฺพํ. รูปํ ปสฺสนฺโต สจฺฉิกโรติ, เวทนํ ปสฺสนฺโต สจฺฉิกโรติ, สญฺญํ ปสฺสนฺโต สจฺฉิกโรติ, สงฺขาเร ปสฺสนฺโต สจฺฉิกโรติ, วิญฺญาณํ ปสฺสนฺโต สจฺฉิกโรติ. จกฺขุํ ฯเปฯ. ชรามรณํ ฯเปฯ อมโตคธํ นิพฺพานํ ปริโยสาน\-ฏฺเฐน ปสฺสนฺโต สจฺฉิกโรติ. เย เย ธมฺมา สจฺฉิกตา โหนฺติ, เต เต ธมฺมา ผสฺสิตา\csromansharedfootnote{ae8b5f2e28d949ba}{ผุสิตา (สฺยา)} โหนฺติ. ตํญาตฏฺเฐน ญาณํ, ปชานนฏฺเฐน ปญฺญา, เตน วุจฺจติ “อิเม ธมฺมา สจฺฉิกาตพฺพาติ โสตาวธานํ, ตํปชานนา ปญฺญา สุตมเย ญาณนฺ”ติ. (5)}

\proseitem{30}{กถํ “อิเม ธมฺมา หานภาคิยา”, “อิเม ธมฺมา ฐิติภาคิยา”, “อิเม ธมฺมา วิเสสภาคิยา”, “อิเม ธมฺมา นิพฺเพธ\-ภาคิยา”ติ โสตาวธานํ, ตํปชานนา ปญฺญา สุตมเย ญาณํ\csromandash{}}

\prose{\csromanfolio{35}ปฐมสฺส ฌานสฺส ลาภิํ กามสหคตา สญฺญา\-มนสิการา สมุทาจรนฺติ, หานภาคิโย ธมฺโม, ตทนุธมฺมตา สติ สนฺติฏฺฐติ, ฐิติภาคิโย ธมฺโม, อวิตกฺก\-สหคตา สญฺญา\-มนสิการา สมุทาจรนฺติ, วิเสสภาคิโย ธมฺโม, นิพฺพิทา\-สหคตา สญฺญา\-มนสิการา สมุทาจรนฺติ วิราคู\-ปสํหิตา, นิพฺเพธ\-ภาคิโย ธมฺโม.}

\prose{ทุติยสฺส ฌานสฺส ลาภิํ วิตกฺก\-สหคตา สญฺญา\-มนสิการา สมุทาจรนฺติ, หานภาคิโย ธมฺโม, ตทนุธมฺมตา สติ สนฺติฏฺฐติ, ฐิติภาคิโย ธมฺโม, อุเปกฺขา\-สุขสหคตา สญฺญา\-มนสิการา สมุทาจรนฺติ, วิเสสภาคิโย ธมฺโม, นิพฺพิทา\-สหคตา สญฺญา\-มนสิการา สมุทาจรนฺติ วิราคู\-ปสํหิตา, นิพฺเพธ\-ภาคิโย ธมฺโม.}

\prose{ตติยสฺส ฌานสฺส ลาภิํ ปีติสุขสหคตา สญฺญา\-มนสิการา สมุทาจรนฺติ, หานภาคิโย ธมฺโม, ตทนุธมฺมตา สติ สนฺติฏฺฐติ, ฐิติภาคิโย ธมฺโม, อทุกฺขมสุข\-สหคตา สญฺญา\-มนสิการา สมุทาจรนฺติ, วิเสสภาคิโย ธมฺโม, นิพฺพิทา\-สหคตา สญฺญา\-มนสิการา สมุทาจรนฺติ วิราคู\-ปสํหิตา, นิพฺเพธ\-ภาคิโย ธมฺโม.}

\prose{จตุตฺถสฺส ฌานสฺส ลาภิํ อุเปกฺขา\-สหคตา สญฺญา\-มนสิการา สมุทาจรนฺติ, หานภาคิโย ธมฺโม, ตทนุธมฺมตา สติ สนฺติฏฺฐติ, ฐิติภาคิโย ธมฺโม, อากาสา\-นญฺจา\-ยตนสหคตา สญฺญา\-มนสิการา สมุทาจรนฺติ, วิเสสภาคิโย ธมฺโม, นิพฺพิทา\-สหคตา สญฺญา\-มนสิการา สมุทาจรนฺติ วิราคู\-ปสํหิตา, นิพฺเพธ\-ภาคิโย ธมฺโม.}

\prose{อากาสา\-นญฺจา\-ยตนสฺส ลาภิํ รูปสหคตา สญฺญา\-มนสิการา สมุทาจรนฺติ, หานภาคิโย ธมฺโม, ตทนุธมฺมตา สติ สนฺติฏฺฐติ, ฐิติภาคิโย ธมฺโม, วิญฺญาณญฺจา\-ยตนสหคตา สญฺญา\-มนสิการา สมุทาจรนฺติ, วิเสสภาคิโย ธมฺโม, นิพฺพิทา\-สหคตา สญฺญา\-มนสิการา สมุทาจรนฺติ วิราคู\-ปสํหิตา, นิพฺเพธ\-ภาคิโย ธมฺโม.}

\prose{วิญฺญาณญฺจา\-ยตนสฺส ลาภิํ อากาสา\-นญฺจา\-ยตนสหคตา สญฺญา\-มนสิการา สมุทาจรนฺติ, หานภาคิโย ธมฺโม, ตทนุธมฺมตา สติ สนฺติฏฺฐติ, ฐิติภาคิโย ธมฺโม, อากิญฺจญฺญายตน\-สหคตา สญฺญา\-มนสิการา สมุทาจรนฺติ, วิเสสภาคิโย ธมฺโม, นิพฺพิทา\-สหคตา สญฺญา\-มนสิการา สมุทาจรนฺติ วิราคู\-ปสํหิตา, นิพฺเพธ\-ภาคิโย ธมฺโม.}

\prose{\csromanfolio{36}อากิญฺจญฺญา\-ยตนสฺส ลาภิํ วิญฺญาณญฺจา\-ยตนสหคตา สญฺญา\-มนสิการา สมุทาจรนฺติ, หานภาคิโย ธมฺโม, ตทนุธมฺมตา สติ สนฺติฏฺฐติ, ฐิติภาคิโย ธมฺโม. เนว\-สญฺญา\-นา\-สญฺญา\-ยตนสหคตา สญฺญา\-มนสิการา สมุทาจรนฺติ, วิเสสภาคิโย ธมฺโม, นิพฺพิทา\-สหคตา สญฺญา\-มนสิการา สมุทาจรนฺติ วิราคู\-ปสํหิตา, นิพฺเพธ\-ภาคิโย ธมฺโม. ตํญาตฏฺเฐน ญาณํ ปชานนฏฺเฐน ปญฺญา, เตน วุจฺจติ “อิเม ธมฺมา หานภาคิยา”, “อิเม ธมฺมา ฐิติภาคิยา”, “อิเม ธมฺมา วิเสสภาคิยา”, “อิเม ธมฺมา นิพฺเพธ\-ภาคิยา”ติ โสตาวธานํ, ตํปชานนา ปญฺญา สุตมเย ญาณํ. \mbox{(4-9)}}

\proseitem{31}{กถํ “สพฺเพ สงฺขารา อนิจฺจา”, “สพฺเพ สงฺขารา ทุกฺขา”, “สพฺเพ ธมฺมา อนตฺตา”ติ โสตาวธานํ, ตํปชานนา ปญฺญา สุตมเย ญาณํ\csromandash{} รูปํ อนิจฺจํ ขยฏฺเฐน, ทุกฺขํ ภยฏฺเฐน, อนตฺตา อสารกฏฺเฐนาติ โสตาวธานํ, ตํปชานนา ปญฺญา สุตมเย ญาณํ. เวทนา. สญฺญา. สงฺขารา. วิญฺญาณํ. จกฺขุ ฯเปฯ. ชรามรณํ อนิจฺจํ ขยฏฺเฐน, ทุกฺขํ ภยฏฺเฐน, อนตฺตา อสารกฏฺเฐนาติ โสตาวธนํ, ตํปชานนา ปญฺญา สุตมเย ญาณํ. ตํญาตฏฺเฐน ญาณํ, ปชานนฏฺเฐน ปญฺญา, เตน วุจฺจติ “สพฺเพ สงฺขารา อนิจฺจา”, “สพฺเพ สงฺขารา ทุกฺขา”, “สพฺเพ ธมฺมา อนตฺตา”ติ โสตาวธานํ, ตํปชานนา ปญฺญา สุตมเย ญาณํ. \mbox{(3-12)}}

\proseitem{32}{กถํ “อิทํ ทุกฺขํ อริยสจฺจํ”, “อิทํ ทุกฺข\-สมุทยํ อริยสจฺจํ”, “อิทํ ทุกฺขนิโรธํ อริยสจฺจํ”, “อิทํ ทุกฺขนิโรธ\-คามินี ปฏิปทา อริยสจฺจนฺ”ติ โสตาวธานํ, ตํปชานนา ปญฺญา สุตมเย ญาณํ\csromandash{}}

\proseitem{33}{ตตฺถ กตมํ ทุกฺขํ อริยสจฺจํ? ชาติปิ ทุกฺขา, ชาราปิ ทุกฺขา, มรณมฺปิ ทุกฺขํ, โสก\-ปริเทว\-ทุกฺข\-โทมนสฺสุ\-ปายาสาปิ ทุกฺขา, อปฺปิเยหิ สมฺปโยโค ทุกฺโข, ปิเยหิ วิปฺปโยโค ทุกฺโข, ยมฺปิจฺฉํ น ลภติ, ตมฺปิ ทุกฺขํ, สํขิตฺเตน ปญฺจุ\-ปาทาน\-กฺขนฺธา\csromansharedfootnote{b485e781af8d7fac}{ปญฺจุปาทานกฺขนฺธาปิ (ก) สติปฏฺฐานสุตฺเต สจฺจวิภงฺเค จ ปสฺสิตพฺพํ.} ทุกฺขา.}

\prose{ตตฺถ กตมา ชาติ? ยา เตสํ เตสํ สตฺตานํ ตมฺหิ ตมฺหิ สตฺตนิกาเย ชาติ สญฺชาติ โอกฺกนฺติ อภินิพฺพตฺติ ขนฺธานํ ปาตุภาโว อายตนานํ ปฏิลาโภ, อยํ วุจฺจติ ชาติ.}

\prose{\csromanfolio{37}ตตฺถ กตมา ชรา? ยา เตสํ เตสํ สตฺตานํ ตมฺหิ ตมฺหิ สตฺตนิกาเย ชรา ชีรณตา ขณฺฑิจฺจํ ปาลิจฺจํ วลิตฺตจตา อายุโน สํหานิ อินฺทฺริยานํ ปริปาโก, อยํ วุจฺจติ ชรา.}

\prose{ตตฺถ กตมํ มรณํ? ยา เตสํ เตสํ สตฺตานํ ตมฺหา ตมฺหา สตฺตนิกายา จุติ จวนตา เภโท อนฺตรธานํ มจฺจุ มรณํ กาลกิริยา\csromansharedfootnote{cd699c4c929c7902}{กาลํกิริยา (ก)} ขนฺธานํ เภโท กเฬวรสฺส นิกฺเขโป ชีวิตินฺทฺริยสฺสุ\-ปจฺเฉโท, อิทํ วุจฺจติ มรณํ.}

\prose{ตตฺถ กตโม โสโก? ญาติพฺยสเนน\csromansharedfootnote{1164a00620b0261f}{โย ญาติพฺยสเนน (สฺยา) ปสฺส สจฺจวิภงฺเค 104 ปิฏฺเฐ.} วา ผุฏฺฐสฺส, โภคพฺยสเนน วา ผุฏฺฐสฺส, โรคพฺยสเนน วา ผุฏฺฐสฺส, สีลพฺยสเนน วา ผุฏฺฐสฺส, ทิฏฺฐิ\-พฺยสเนน วา ผุฏฺฐสฺส, อญฺญตร\-ญฺญตเรน พฺยสเนน สมนฺนาคตสฺส, อญฺญตร\-ญฺญตเรน ทุกฺขธมฺเมน ผุฏฺฐสฺส โสโก โสจนา โสจิตตฺตํ อนฺโตโสโก อนฺโตปริโสโก เจตโส ปริชฺฌายนา โทมนสฺสํ โสกสลฺลํ, อยํ วุจฺจติ โสโก.}

\prose{ตตฺถ กตโม ปริเทโว? ญาติพฺยสเนน\csromansharedfootnote{1164a00620b0261f}{โย ญาติพฺยสเนน (สฺยา) ปสฺส สจฺจวิภงฺเค 104 ปิฏฺเฐ.} วา ผุฏฺฐสฺส, โภคพฺยสเนน วา ผุฏฺฐสฺส, โรคพฺยสเนน วา ผุฏฺฐสฺส, สีลพฺยสเนน วา ผุฏฺฐสฺส, ทิฏฺฐิ\-พฺยสเนน วา ผุฏฺฐสฺส, อญฺญตร\-ญฺญตเรน พฺยสเนน สมนฺนาคตสฺส, อญฺญตร\-ญฺญตเรน ทุกฺขธมฺเมน ผุฏฺฐสฺส อาเทโว ปริเทโว อาเทวนา ปริเทวนา อาเทวิตตฺตํ ปริเทวิตตฺตํ วาจา ปลาโป วิปฺปลาโป ลาลปฺโป ลาลปฺปนา ลาลปฺปิตตฺตํ, อยํ วุจฺจติ ปริเทโว.}

\prose{ตตฺถ กตมํ ทุกฺขํ? ยํ กายิกํ อสาตํ กายิกํ ทุกฺขํ, กาย\-สมฺผสฺสชํ อสาตํ ทุกฺขํ เวทยิตํ, กาย\-สมฺผสฺสชา อสาตา ทุกฺขา เวทนา, อิทํ วุจฺจติ ทุกฺขํ.}

\prose{ตตฺถ กตมํ โทมนสฺสํ? ยํ เจตสิกํ อสาตํ เจตสิกํ ทุกฺขํ, เจโต\-สมฺผสฺสชํ อสาตํ ทุกฺขํ เวทยิตํ, เจโต\-สมฺผสฺสชา อสาตา ทุกฺขา เวทนา, อิทํ วุจฺจติ โทมนสฺสํ.}

\prose{ตตฺถ กตโม อุปายาโส? ญาติพฺยสเนน วา ผุฏฺฐสฺส, โภคพฺยสเนน วา ผุฏฺฐสฺส, โรคพฺยสเนน วา ผุฏฺฐสฺส, สีลพฺยสเนน วา ผุฏฺฐสฺส, ทิฏฺฐิ\-พฺยสเนน วา ผุฏฺฐสฺส, อญฺญตร\-ญฺญตเรน พฺยสเนน สมนฺนาคตสฺส, \csromanfolio{38}อญฺญตร\-ญฺญตเรน ทุกฺขธมฺเมน ผุฏฺฐสฺส อายาโส อุปายาโส อายาสนา อุปายาสนา อายาสิตตฺตํ อุปายาสิตตฺตํ, อยํ วุจฺจติ อุปายาโส.}

\prose{ตตฺถ กตโม อปฺปิเยหิ สมฺปโยโค ทุกฺโข? อิธ ย’สฺส เต โหนฺติ อนิฏฺฐา อกนฺตา อมนาปา รูปา สทฺทา คนฺธา รสา โผฏฺฐพฺพา, เย วา ปนสฺส เต โหนฺติ อนตฺถกามา อหิตกามา อผาสุกามา อโยคกฺเขม\-กามา, ยา เตหิ สทฺธิํ สงฺคติ สมาคโม สโมธานํ มิสฺสีภาโว, อยํ วุจฺจติ อปฺปิเยหิ สมฺปโยโค ทุกฺโข.}

\prose{ตตฺถ กตโม ปิเยหิ วิปฺปโยโค ทุกฺโข? อิธ ย’สฺส เต โหนฺติ อิฏฺฐา กนฺตา มนาปา รูปา สทฺทา คนฺธา รสา โผฏฺฐพฺพา, เย วา ปนสฺส เต โหนฺติ อตฺถกามา หิตกามา ผาสุกามา โยคกฺเขมกามา มาตา วา ปิตา วา ภาตา วา ภคินี วา มิตฺตา วา อมจฺจา วา ญาตี วา สาโลหิตา วา, ยา เตหิ สทฺธิํ อสงฺคติ อสมาคโม อสโมธานํ อมิสฺสีภาโว, อยํ วุจฺจติ ปิเยหิ วิปฺปโยโค ทุกฺโข.}

\prose{ตตฺถ กตมํ ยมฺปิจฺฉํ น ลภติ, ตมฺปิ ทุกฺขํ? ชาติธมฺมานํ สตฺตานํ เอวํ อิจฺฉา อุปฺปชฺชติ “อโห วต มยํ น ชาติธมฺมา อสฺสาม, น จ วต โน ชาติ อาคจฺเฉยฺยา”ติ. น โข ปเนตํ อิจฺฉาย ปตฺตพฺพํ, อิทมฺปิ ยมฺปิจฺฉํ น ลภติ, ตมฺปิ ทุกฺขํ. ชราธมฺมานํ สตฺตานํ ฯเปฯ. พฺยาธิ\-ธมฺมานํ สตฺตานํ. มรณ\-ธมฺมานํ สตฺตานํ. โสกปริเทวทุกฺขโทมนสฺสุปายาส- ธมฺมานํ สตฺตานํ เอวํ อิจฺฉา อุปฺปชฺชติ “อโห วต มยํ น โสก\-ปริเทว\-ทุกฺข\-โทมนสฺสุ\-ปายาส\-ธมฺมา อสฺสาม, น จ วต โน โสก\-ปริเทว\-ทุกฺข\-โทมนสฺสุ\-ปายาสา อาคจฺเฉยฺยุนฺ”ติ, น โข ปเนตํ อิจฺฉาย ปตฺตพฺพํ, อิทมฺปิ ยมฺปิจฺฉํ น ลภติ, ตมฺปิ ทุกฺขํ.}

\prose{ตตฺถ กตเม สํขิตฺเตน ปญฺจุ\-ปาทาน\-กฺขนฺธา ทุกฺขา? เสยฺยถิทํ\csromansharedfootnote{b9679a7da29d1751}{เสยฺยถีทํ (สฺยา, สีฏฺฐ)}, รูปุปาทาน\-กฺขนฺโธ\csromansharedfootnote{f3c3ed8c97f38229}{รูปูปาทานกฺขนฺโธ (สฺยา) เอวมุปริปิ.} เวทนุ\-ปาทาน\-กฺขนฺโธ สญฺญุ\-ปาทาน\-กฺขนฺโธ สงฺขารุ\-ปาทาน\-กฺขนฺโธ วิญฺญาณุ\-ปาทาน\-กฺขนฺโธ, อิเม วุจฺจนฺติ สํขิตฺเตน ปญฺจุ\-ปาทาน\-กฺขนฺธา ทุกฺขา. อิทํ วุจฺจติ ทุกฺขํ อริยสจฺจํ.}

\proseitem{34}{\csromanfolio{39}ตตฺถ กตมํ ทุกฺข\-สมุทยํ อริยสจฺจํ? ยายํ ตณฺหา โปโนภวิกา\csromansharedfootnote{041ccc09cd1816ea}{โปโนพฺภคิกา (สฺยา)} นนฺทิ\-ราค\-สหคตา ตตฺร ตตฺรา\-ภินนฺทินี. เสยฺยถิทํ, กามตณฺหา ภวตณฺหา วิภวตณฺหา, สา โข ปเนสา ตณฺหา กตฺถ อุปฺปชฺชมานา อุปฺปชฺชติ, กตฺถ นิวิสมานา นิวิสติ? ยํ โลเก ปิยรูปํ สาตรูปํ, เอตฺเถสา ตณฺหา อุปฺปชฺชมานา อุปฺปชฺชติ, เอตฺถ นิวิสมานา นิวิสติ. กิญฺจ โลเก ปิยรูปํ สาตรูปํ? จกฺขุ โลเก ปิยรูปํ สาตรูปํ, เอตฺเถสา ตณฺหา อุปฺปชฺชมานา อุปฺปชฺชติ, เอตฺถ นิวิสมานา นิวิสติ. โสตํ โลเก. ฆานํ โลเก. ชิวฺหา โลเก. กาโย โลเก. มโน โลเก ปิยรูปํ สาตรูปํ, เอตฺเถสา ตณฺหา อุปฺปชฺชมานา อุปฺปชฺชติ, เอตฺถ นิวิสมานา นิวิสติ.  รูปา โลเก ปิยรูปํ สาตรูปํ, เอตฺเถสา ตณฺหา อุปฺปชฺชมานา อุปฺปชฺชติ, เอตฺถ นิวิสมานา นิวิสติ. สทฺทา โลเก ปิยรูปํ สาตรูปํ ฯเปฯ. ธมฺมา โลเก.  จกฺขุวิญฺญาณํ โลเก ฯเปฯ. มโนวิญฺญาณํ โลเก.  จกฺขุ\-สมฺผสฺโส โลเก ฯเปฯ. มโนสมฺผสฺโส โลเก.  จกฺขุ\-สมฺผสฺสชา เวทนา โลเก ฯเปฯ. มโน\-สมฺผสฺสชา เวทนา โลเก.  รูปสญฺญา โลเก ฯเปฯ. ธมฺมสญฺญา โลเก.  รูปสญฺเจตนา โลเก ฯเปฯ. ธมฺม\-สญฺเจตนา โลเก.  รูปตณฺหา โลเก ฯเปฯ. ธมฺมตณฺหา โลเก.  รูปวิตกฺโก โลเก ฯเปฯ. ธมฺมวิตกฺโก โลเก.  รูปวิจาโร โลเก ฯเปฯ. ธมฺมวิจาโร โลเก ปิยรูปํ สาตรูปํ, เอตฺเถสา ตณฺหา อุปฺปชฺชมานา อุปฺปชฺชติ, เอตฺถ นิวิสมานา นิวิสติ. อิทํ วุจฺจติ ทุกฺข\-สมุทยํ อริยสจฺจํ.}

\proseitem{35}{ตตฺถ กตมํ ทุกฺขนิโรธํ อริยสจฺจํ? โย ตสฺสาเยว ตณฺหาย อเสส\-วิราค\-นิโรโธ จาโค ปฏินิสฺสคฺโค มุตฺติ อนาลโย, สา โข ปเนสา ตณฺหา กตฺถ ปหียมานา ปหียติ, กตฺถ นิรุชฺฌมานา นิรุชฺฌติ? ยํ โลเก ปิยรูปํ สาตรูปํ, เอตฺเถสา ตณฺหา ปหียมานา ปหียติ, เอตฺถ นิรุชฺฌมานา นิรุชฺฌติ. กิญฺจ โลเก ปิยรูปํ สาตรูปํ? จกฺขุ โลเก ปิยรูปํ สาตรูปํ, เอตฺเถสา ตณฺหา ปหียมานา ปหียติ, เอตฺถ นิรุชฺฌมานา นิรุชฺฌติ ฯเปฯ. ธมฺมวิจาโร โลเก ปิยรูปํ สาตรูปํ, เอตฺเถสา ตณฺหา ปหียมานา ปหียติ, เอตฺถ นิรุชฺฌมานา นิรุชฺฌติ. อิทํ วุจฺจติ ทุกฺขนิโรธํ อริยสจฺจํ.}

\proseitem{36}{\csromanfolio{40}ตตฺถ กตมํ ทุกฺขนิโรธ\-คามินี ปฏิปทา อริยสจฺจ? อยเมว อริโย อฏฺฐงฺคิโก มคฺโค. เสยฺยถิทํ, สมฺมาทิฏฺฐิ สมฺมาสงฺกปฺโป สมฺมาวาจา สมฺมากมฺมนฺโต สมฺมาอาชีโว สมฺมาวายาโม สมฺมาสติ สมฺมาสมาธิ.}

\prose{ตตฺถ กตมา สมฺมาทิฏฺฐิ? ทุกฺเข ญาณํ, ทุกฺขสมุทเย ญาณํ, ทุกฺขนิโรเธ ญาณํ, ทุกฺขนิโรธ\-คามินิยา ปฏิปทาย ญาณํ, อยํ วุจฺจติ สมฺมาทิฏฺฐิ.}

\prose{ตตฺถ กตโม สมฺมาสงฺกปฺโป? เนกฺขมฺม\-สงฺกปฺโป อพฺยาปาทสงฺกปฺโป อวิหิํสาสงฺกปฺโป, อยํ วุจฺจติ สมฺมาสงฺกปฺโป.}

\prose{ตตฺถ กตมา สมฺมาวาจา? มุสาวาทา เวรมณี\csromansharedfootnote{f96664fc1fbeb341}{เวรมณิ (ก)}, ปิสุณาย วาจาย เวรมณี, ผรุสาย วาจาย เวรมณี, สมฺผปฺปลาปา เวรมณี, อยํ วุจฺจติ สมฺมาวาจา.}

\prose{ตตฺถ กตโม สมฺมากมฺมนฺโต? ปาณาติปาตา เวรมณี, อทินฺนาทานา เวรมณี, กาเมสุ\-มิจฺฉาจารา เวรมณี, อยํ วุจฺจติ สมฺมากมฺมนฺโต.}

\prose{ตตฺถ กตโม สมฺมาอาชีโว? อิธ อริยภาวโก มิจฺฉาอาชีวํ ปหาย สมาอาชีเวน ชีวิกํ\csromansharedfootnote{f8e93fdd53c8c38b}{ชีวิตํ (ก)} กปฺเปติ, อยํ วุจฺจติ สมฺมาอาชีโว.}

\prose{ตตฺถ กตโม สมฺมาวายาโม? อิธ ภิกฺขุ อนุปฺปนฺนานํ ปาปกานํ อกุสลานํ ธมฺมานํ อนุปฺปาทาย ฉนฺทํ ชเนติ วายมติ, วีริยํ อารภติ, จิตฺตํ ปคฺคณฺหาติ ปทหติ. อุปฺปนฺนานํ ปาปกานํ อกุสลานํ ธมฺมานํ ปหานาย ฯเปฯ. อนุปฺปนฺนานํ กุสลานํ ธมฺมานํ อุปฺปาทาย ฯเปฯ. อุปฺปนฺนานํ กุสลานํ ธมฺมานํ ฐิติยา อสมฺโมสาย ภิยฺโยภาวาย เวปุลฺลาย ภาวนาย ปาริปูริยา ฉนฺทํ ชเนติ วายมติ, วีริยํ อารภติ, จิตฺตํ ปคฺคณฺหาติ ปทหติ, อยํ วุจฺจติ สมฺมาวายาโม.}

\prose{ตตฺถ กตมา สมฺมาสติ? อิธ ภิกฺขุ กาเย กายานุปสฺสี วิหรติ อาตาปี สมฺปชาโน สติมา วิเนยฺย โลเก อภิชฺฌา\-โทมนสฺสํ. เวทนาสุ ฯเปฯ. จิตฺเต ฯเปฯ. ธมฺเมสุ ธมฺมานุปสฺสี วิหรติ อาตาปี สมฺปชาโน สติมา วิเนยฺย โลเก อภิชฺฌา\-โทมนสฺสํ, อยํ วุจฺจติ สมฺมาสติ.}

\prosewithcloserruled{ตตฺถ กตโม สมฺมาสมาธิ? อิธ ภิกฺขุ วิวิจฺเจว กาเมหิ วิวิจฺจ อกุสเลหิ ธมฺเมหิ สวิตกฺกํ สวิจารํ วิเวกชํ ปีติสุขํ ปฐมํ ฌานํ \csromanfolio{41}อุปสมฺปชฺช วิหรติ. วิตกฺก\-วิจารานํ วูปสมา อชฺฌตฺตํ สมฺปสาทนํ เจตโส เอโกทิภาวํ อวิตกฺกํ อวิจารํ สมาธิชํ ปีติสุขํ ทุติยํ ฌานํ อุปสมฺปชฺช วิหรติ. ปีติยา จ วิราคา อุเปกฺขโก จ วิหรติ สโต จ สมฺปชาโน สุขญฺจ กาเยน ปฏิสํเวเทติ, ยํ ตํ อริยา อาจิกฺขนฺติ “อุเปกฺขโก สติมา สุขวิหารี”ติ, ตติยํ ฌานํ อุปสมฺปชฺช วิหรติ. สุขสฺส จ ปหานา ทุกฺขสฺส จ ปหานา ปุพฺเพว โสมนสฺส\-โทมนสฺสานํ อตฺถงฺคมา อทุกฺขมสุขํ อุเปกฺขา\-สติ\-ปาริสุทฺธิํ จตุตฺถํ ฌานํ อุปสมฺปชฺช วิหรติ. อยํ วุจฺจติ สมฺมาสมาธิ. อิทํ วุจฺจติ ทุกฺขนิโรธ\-คามินี ปฏิปทา อริยสจฺจํ.  ตํญาตฏฺเฐน ญาณํ, ปชานนฏฺเฐน ปญฺญา, เตน วุจฺจติ “อิทํ ทุกฺขํ อริยสจฺจํ”, “อิทํ ทุกฺข\-สมุทยํ อริยสจฺจํ”, “อิทํ ทุกฺขนิโรธํ อริยสจฺจํ”, “อิทํ ทุกฺขนิโรธ\-คามินี ปฏิปทา อริยสจฺจนฺ”ติ โสตาวธานํ, ตํปชานนา ปญฺญา สุตมเย ญาณํ. เอวํ โสตาวธาเน ปญฺญา สุตมเย ญาณํ. \mbox{(4-16)}}{สุตมย\-ญาณ\-นิทฺเทโส ปฐโม.}
\csromanmatikaanchor{mtk.4}
\csromantocmarkref{subsection}{2. สีลมยญาณนิทฺเทส}{mtk.4}
\bookmark[dest={mtk.4},level=2]{2. สีลมยญาณนิทฺเทส}
\csromanheader{2}{2. สีลมย\-ญาณ\-นิทฺเทส}
\proseitem{37}{กถํ สุตฺวาน สํวเร ปญฺญา สีลมเย ญาณํ\csromandash{}ปญฺจ สีลานิ ปริยนฺต\-ปาริสุทฺธิสีลํ อปริยนฺต\-ปาริสุทฺธิสีลํ ปริปุณฺณ\-ปาริสุทฺธิสีลํ อปรามฏฺฐ\-ปาริสุทฺธิสีลํ ปฏิปฺปสฺสทฺธิปาริสุทฺธิสีลนฺติ.}

\prose{ตตฺถ กตมํ ปริยนฺต\-ปาริสุทฺธิสีลํ? อนุปสมฺปนฺนานํ ปริยนฺต\-สิกฺขาปทานํ, อิทํ ปริยนฺต\-ปาริสุทฺธิสีลํ.  กตมํ อปริยนฺต\-ปาริสุทฺธิสีลํ? อุปสมฺปนฺนานํ อปริยนฺต\-สิกฺขาปทานํ, อิทํ อปริยนฺต\-ปาริสุทฺธิสีลํ.  กตมํ ปริปุณฺณ\-ปาริสุทฺธิสีลํ? ปุถุชฺชน\-กลฺยาณกานํ กุสลธมฺเม ยุตฺตานํ เสกฺข\-ปริยนฺเต\csromansharedfootnote{9a1fee1699d04151}{เสขปริยนฺเต (ก)} ปริปูร\-การีนํ กาเย จ ชีวิเต จ อนเปกฺขานํ ปริจฺจตฺต\-ชีวิตานํ, อิทํ ปริปุณฺณ\-ปาริสุทฺธิสีลํ.  กตมํ อปรามฏฺฐ\-ปาริสุทฺธิสีลํ? สตฺตนฺนํ เสกฺขานํ, อิทํ อปรามฏฺฐ\-ปาริสุทฺธิสีลํ.  กตมํ ปฏิปฺปสฺสทฺธิ\-ปาริสุทฺธิสีลํ? ตถาคต\-สาวกานํ ขีณาสวานํ ปจฺเจก\-พุทฺธานํ ตถาคตานํ อรหนฺตานํ สมฺมา\-สมฺพุทฺธานํ, อิทํ ปฏิปฺปสฺสทฺธิ\-ปาริสุทฺธิสีลํ.}

\proseitem{38}{\csromanfolio{42}อตฺถิ สีลํ ปริยนฺตํ, อตฺถิ สีลํ อปริยนฺตํ. ตตฺถ กตมํ ตํ สีลํ ปริยนฺตํ? อตฺถิ สีลํ ลาภ\-ปริยนฺตํ, อตฺถิ สีลํ ยสปริยนฺตํ, อตฺถิ สีลํ ญาติปริยนฺตํ, อตฺถิ สีลํ องฺคปริยนฺตํ, อตฺถิ สีลํ ชีวิต\-ปริยนฺตํ.}

\prose{กตมํ ตํ สีลํ ลาภ\-ปริยนฺตํ? อิเธกจฺโจ ลาภเหตุ ลาภปจฺจยา ลาภการณา ยถา\-สมาทินฺนํ\csromansharedfootnote{daf90c8c144c1907}{ยถาสมาทิณฺณํ (ก)} สิกฺขาปทํ วีติกฺกมติ, อิทํ ตํ สีลํ ลาภ\-ปริยนฺตํ.  กตมํ ตํ สีลํ ยสปริยนฺตํ? อิเธกจฺโจ ยสเหตุ ยสปจฺจยา ยสการณา ยถา\-สมาทินฺนํ สิกฺขาปทํ วีติกฺกมติ, อิทํ ตํ สีลํ ยสปริยนฺตํ.  กตมํ ตํ สีลํ ญาติปริยนฺตํ? อิเธกจฺโจ ญาติเหตุ ญาติปจฺจยา ญาติการณา ยถา\-สมาทินฺนํ สิกฺขาปทํ วีติกฺกมติ, อิทํ ตํ สีลํ ญาติปริยนฺตํ.  กตมํ ตํ สีลํ องฺคปริยนฺตํ? อิเธกจฺโจ องฺคเหตุ องฺคปจฺจยา องฺคการณา ยถา\-สมาทินฺนํ สิกฺขาปทํ วีติกฺกมติ, อิทํ ตํ สีลํ องฺคปริยนฺตํ.  กตมํ ตํ สีลํ ชีวิต\-ปริยนฺตํ? อิเธกจฺโจ ชีวิตเหตุ ชีวิตปจฺจยา ชีวิตการณา ยถา\-สมาทินฺนํ สิกฺขาปทํ วีติกฺกมติ, อิทํ ตํ สีลํ ชีวิต\-ปริยนฺตํ. เอวรูปานิ สีลานิ ขณฺฑานิ ฉิทฺทานิ สพลานิ กมฺมาสานิ น ภุชิสฺสานิ น วิญฺญุ\-ปฺปสตฺถานิ\csromansharedfootnote{62c8c52a2e6a59ff}{น วิญฺญูปสตฺถานิ (ก)} ปรามฏฺฐานิ อสมาธิสํวตฺตนิกานิ\csromansharedfootnote{b9acf57b222deb8a}{น สมาธิสํวตฺตนิกานิ (ก)} น อวิปฺปฏิสาร\-วตฺถุกานิ น ปาโมชฺช\-วตฺถุกานิ น ปีติวตฺถุกานิ น ปสฺสทฺธิ\-วตฺถุกานิ น สุขวตฺถุกานิ น สมาธิ\-วตฺถุกานิ น ยถา\-ภูต\-ญาณ\-ทสฺสนวตฺถุกานิ น เอกนฺต\-นิพฺพิทาย น วิราคาย น นิโรธาย น อุปสมาย น อภิญฺญาย น สมฺโพธาย น นิพฺพานาย สํวตฺตนฺติ, อิทํ ตํ สีลํ ปริยนฺตํ.}

\prose{กตมํ ตํ สีลํ อปริยนฺตํ? อตฺถิ สีลํ น ลาภ\-ปริยนฺตํ, อตฺถิ สีลํ น ยสปริยนฺตํ, อตฺถิ สีลํ น ญาติปริยนฺตํ, อตฺถิ สีลํ น องฺคปริยนฺตํ, อตฺถิ สีลํ น ชีวิต\-ปริยนฺตํ.}

\prose{กตมํ ตํ สีลํ น ลาภ\-ปริยนฺตํ? อิเธกจฺโจ ลาภเหตุ ลาภปจฺจยา ลาภการณา ยถา\-สมาทินฺนํ สิกฺขาปทํ วีติกฺกมาย จิตฺตมฺปิ น อุปฺปาเทติ, กิํ โส วีติกฺกมิสฺสติ, อิทํ ตํ สีลํ น ลาภ\-ปริยนฺตํ.  กตมํ ตํ สีลํ น ยสปริยนฺตํ? อิเธกจฺโจ ยสเหตุ ยสปจฺจยา ยสการณา ยถา\-สมาทินฺนํ สิกฺขาปทํ วีติกฺกมาย จิตฺตมฺปิ น \csromanfolio{43}อุปฺปาเทติ, กิํ โส วีติกฺกมิสฺสติ, อิทํ ตํ สีลํ น ยสปริยนฺตํ.  กตมํ ตํ สีลํ น ญาติปริยนฺตํ? อิเธกจฺโจ ญาติเหตุ ญาติปจฺจยา ญาติการณา ยถา\-สมาทินฺนํ สิกฺขาปทํ วีติกฺกมาย จิตฺตมฺปิ น อุปฺปาเทติ, กิํ โส วีติกฺกมิสฺสติ, อิทํ ตํ สีลํ น ญาติปริยนฺตํ.  กตมํ ตํ สีลํ น องฺคปริยนฺตํ? อิเธกจฺโจ องฺคเหตุ องฺคปจฺจยา องฺคการณา ยถา\-สมาทินฺนํ สิกฺขาปทํ วีติกฺกมาย จิตฺตมฺปิ น อุปฺปาเทติ, กิํ โส วีติกฺกมิสฺสติ, อิทํ ตํ สีลํ น องฺคปริยนฺตํ.  กตมํ ตํ สีลํ น ชีวิต\-ปริยนฺตํ? อิเธกจฺโจ ชีวิตเหตุ ชีวิตปจฺจยา ชีวิตการณา ยถา\-สมาทินฺนํ สิกฺขาปทํ วีติกฺกมาย จิตฺตมฺปิ น อุปฺปาเทติ, กิํ โส วีติกฺกมิสฺสติ, อิทํ ตํ สีลํ น ชีวิต\-ปริยนฺตํ. เอวรูปานิ สีลานิ อขณฺฑานิ อจฺฉิทฺทานิ อสพลานิ อกมฺมาสานิ ภุชิสฺสานิ วิญฺญุ\-ปฺปสตฺถานิ อปรามฏฺฐานิ สมาธิ\-สํวตฺตนิกานิ อวิปฺปฏิสาร\-วตฺถุกานิ ปาโมชฺช\-วตฺถุกานิ ปีติวตฺถุกานิ ปสฺสทฺธิ\-วตฺถุกานิ สฺ อุขวตฺถุกานิ สมาธิ\-วตฺถุกานิ ยถา\-ภูต\-ญาณ\-ทสฺสนวตฺถุกานิ เอกนฺต\-นิพฺพิทาย วิราคาย นิโรธาย อุปสมาย อภิญฺญาย สมฺโพธาย นิพฺพานาย สํวตฺตนฺติ, อิทํ ตํ สีลํ อปริยนฺตํ.}

\proseitem{39}{กิํ สีลํ, กติ สีลานิ, กิํสมุฏฺฐานํ สีลํ, กติธมฺม\-สโมธานํ สีลํ.}

\prose{กิํ สีลนฺติ เจตนา สีลํ, เจตสิกํ สีลํ, สํวโร สีลํ, อวีติกฺกโม สีลํ.  กติ สีลานีติ ตีณิ สีลานิ, กุสลสีลํ, อกุสลสีลํ, อพฺยากตสีลํ.  กิํสมุฏฺฐานํ สีลนฺติ กุสล\-จิตฺต\-สมุฏฺฐานํ กุสลสีลํ, อกุสล\-จิตฺต\-สมุฏฺฐานํ อกุสลสีลํ, อพฺยากต\-จิตฺต\-สมุฏฺฐานํ อพฺยากตสีลํ.  กติธมฺม\-สโมธานํ สีลนฺติ สํวร\-สโมธานํ สีลํ, อวีติกฺกม\-สโมธานํ สีลํ, ตถาภาเว ชาต\-เจตนา\-สโมธานํ สีลํ.}

\proseitem{40}{ปาณาติปาตํ สํวรฏฺเฐน สีลํ, อวีติกฺกม\-ฏฺเฐน สีลํ. อทินฺนาทานํ สํวรฏฺเฐน สีลํ, อวีติกฺกม\-ฏฺเฐน สีลํ. กาเมสุ\-มิจฺฉาจารํ สํวรฏฺเฐน สีลํ, อวีติกฺกม\-ฏฺเฐน สีลํ. มุสาวาทํ สํวรฏฺเฐน สีลํ, อวีติกฺกม\-ฏฺเฐน สีลํ. ปิสุณํ วาจํ\csromansharedfootnote{e8b7016534dc393d}{ปิสุณาวาจํ (สฺยา, ก) ที 1. 4 ปิฏฺเฐ ปสฺสิตพฺพา.} สํวรฏฺเฐน สีลํ, อวีติกฺกม\-ฏฺเฐน สีลํ. ผรุสํ วาจํ\csromansharedfootnote{630b0a3d51a601f3}{ผรุสวาจํ (สฺยา, ก)} สํวรฏฺเฐน สีลํ, อวีติกฺกม\-ฏฺเฐน สีลํ. สมฺผปฺปลาปํ สํวรฏฺเฐน สีลํ, อวีติกฺกม\-ฏฺเฐน \csromanfolio{44}สีลํ. อภิชฺฌํ สํวรฏฺเฐน สีลํ, อวีติกฺกม\-ฏฺเฐน สีลํ. พฺยาปาทํ สํวรฏฺเฐน สีลํ, อวีติกฺกม\-ฏฺเฐน สีลํ. มิจฺฉาทิฏฺฐิํ สํวรฏฺเฐน สีลํ, อวีติกฺกม\-ฏฺเฐน สีลํ.}

\proseitem{41}{เนกฺขมฺเมน กามจฺฉนฺทํ สํวรฏฺเฐน สีลํ, อวีติกฺกม\-ฏฺเฐน สีลํ. อพฺยาปาเทน พฺยาปาทํ สํวรฏฺเฐน สีลํ, อวีติกฺกม\-ฏฺเฐน สีลํ. อาโลกสญฺญาย ถินมิทฺธํ. อวิกฺเขป\-ฏฺเฐน อุทฺธจฺจํ. ธมฺม\-ววตฺถาเนน วิจิกิจฺฉํ. ญาเณน อวิชฺชํ. ปาโมชฺเชน อรติํ.}

\prose{ปฐเมน ฌาเนน\csromansharedfootnote{6fa30d05cbad9a7d}{ปฐมชฺฌาเนน (สฺยา) เอวมีทิเสสุ ฐาเนสุ.} นีวรเณ. ทุติเยน ฌาเนน วิตกฺกวิจาเร. ตติเยน ฌาเนน ปีติํ. จตุตฺเถน ฌาเนน สุขทุกฺขํ. อากาสา\-นญฺจา\-ยตนสมาปตฺติยา รูปสญฺญํ ปฏิฆสญฺญํ นานตฺตสญฺญํ. วิญฺญาณญฺจายตน\-สมาปตฺติยา อากาสา\-นญฺจา\-ยตนสญฺญํ. อากิญฺจญฺญายตน\-สมาปตฺติยา วิญฺญาณญฺจา\-ยตนสญฺญํ. เนว\-สญฺญา\-นา\-สญฺญา\-ยตนสมาปตฺติยา อากิญฺจญฺญายตน\-สญฺญํ.}

\prose{อนิจฺจา\-นุปสฺสนาย นิจฺจสญฺญํ. ทุกฺขา\-นุปสฺสนาย สุขสญฺญํ. อนตฺตา\-นุปสฺสนาย อตฺตสญฺญํ. นิพฺพิทา\-นุปสฺสนาย นนฺทิํ. วิราคา\-นุปสฺสนาย ราคํ. นิโรธา\-นุปสฺสนาย สมุทยํ. ปฏินิสฺสคฺคา\-นุปสฺสนาย อาทานํ. ขยานุปสฺสนาย ฆนสญฺญํ. วยา\-นุปสฺสนาย อายูหนํ. วิปริณามา\-นุปสฺสนาย ธุวสญฺญํ. อนิมิตฺตา\-นุปสฺสนาย นิมิตฺตํ. อปฺปณิหิตา\-นุปสฺสนาย ปณิธิํ. สุญฺญตา\-นุปสฺสนาย อภินิเวสํ. อธิปญฺญา\-ธมฺม\-วิปสฺสนาย สาราทานา\-ภินิเวสํ. ยถา\-ภูต\-ญาณ\-ทสฺสเนน สมฺโมหา\-ภินิเวสํ. อาทีนวา\-นุปสฺสนาย อาลยา\-ภินิเวสํ. ปฏิสงฺขานุปสฺสานาย อปฺปฏิสงฺขํ. วิวฏฺฏนา\-นุปสฺสนาย สญฺโญคา\-ภินิเวสํ.}

\prose{โสตาปตฺติ\-มคฺเคน ทิฏฺเฐกฏฺเฐ กิเลเส. สกทาคามิ\-มคฺเคน โอฬาริเก กิเลเส. อนาคามิมคฺเคน อณุสหคเต กิเลเส. อรหตฺต\-มคฺเคน สพฺพกิเลเส สํวรฏฺเฐน สีลํ, อวีติกฺกม\-ฏฺเฐน สีลํ.}

\prose{ปญฺจ สีลานิ ปาณาติปาตสฺส ปหานํ สีลํ, เวรมณี สีลํ, เจตนา สีลํ, สํวโร สีลํ, อวีติกฺกโม สีลํ. เอวรูปานิ สีลานิ จิตฺตสฺส อวิปฺปฏิสาราย สํวตฺตนฺติ, ปาโมชฺชาย สํวตฺตนฺติ, ปีติยา สํวตฺตนฺติ, \csromanfolio{45}ปสฺสทฺธิยา สํวตฺตนฺติ, โสมนสฺสาย สํวตฺตนฺติ, อาเสวนาย สํวตฺตนฺติ, ภาวนาย สํวตฺตนฺติ, พหุลีกมฺมาย สํวตฺตนฺติ, อลงฺการาย สํวตฺตนฺติ, ปริกฺขาราย สํวตฺตนฺติ, ปริวาราย สํวตฺตนฺติ, ปาริปูริยา สํวตฺตนฺติ, เอกนฺต\-นิพฺพิทาย วิราคาย นิโรธาย อุปสมาย อภิญฺญาย สมฺโพธาย นิพฺพานาย สํวตฺตนฺติ.}

\prose{เอวรูปานํ สีลานํ สํวร\-ปาริสุทฺธิ อธิสีลํ, สํวร\-ปาริสุทฺธิยา ฐิตํ จิตฺตํ น วิกฺเขปํ คจฺฉติ, อวิกฺเขป\-ปาริสุทฺธิ อธิจิตฺตํ, สํวร\-ปาริสุทฺธิํ สมฺมา ปสฺสติ, อวิกฺเขป\-ปาริสุทฺธิํ สมฺมา ปสฺสติ. ทสฺสน\-ปาริสุทฺธิ อธิปญฺญา. โย ตตฺถ สํวรฏฺโฐ, อยํ อธิสีลสิกฺขา. โย ตตฺถ อวิกฺเขปฏฺโฐ, อยํ อธิจิตฺต\-สิกฺขา. โย ตตฺถ ทสฺสนฏฺโฐ, อยํ อธิปญฺญา\-สิกฺขา.}

\prose{อิมา ติสฺโส สิกฺขาโย อาวชฺชนฺโต สิกฺขติ, ชานนฺโต สิกฺขติ, ปสฺสนฺโต สิกฺขติ, ปจฺจเวกฺขนฺโต สิกฺขติ, จิตฺตํ อธิฏฺฐหนฺโต สิกฺขติ, สทฺธาย อธิมุจฺจนฺโต สิกฺขติ, วีริยํ ปคฺคณฺหนฺโต สิกฺขติ, สติํ อุปฏฺฐเปนฺโต สิกฺขติ, จิตฺตํ สมาทหนฺโต สิกฺขติ, ปญฺญาย ปชานนฺโต สิกฺขติ, อภิญฺเญยฺยํ อภิชานนฺโต สิกฺขติ, ปริญฺเญยฺยํ ปริชานนฺโต สิกฺขติ, ปหาตพฺพํ ปชหนฺโต สิกฺขติ, สจฺฉิกาตพฺพํ สจฺฉิกโรนฺโต สิกฺขติ, ภาเวตพฺพํ ภาเวนฺโต สิกฺขติ.}

\prose{ปญฺจ สีลานิ อทินฺนาทานสฺส. กาเมสุ\-มิจฺฉาจารสฺส. มุสาวาทสฺส. ปิสุณาย วาจาย. ผรุสาย วาจาย. สมฺผ\-ปฺปลาปสฺส. อภิชฺฌาย. พฺยาปาทสฺส. มิจฺฉาทิฏฺฐิยา.}

\prose{เนกฺขมฺเมน กามจฺฉนฺทสฺส. อพฺยาปาเทน พฺยาปาทสฺส. อาโลกสญฺญาย ถินมิทฺธสฺส. อวิกฺเขเปน อุทฺธจฺจสฺส. ธมฺม\-ววตฺถาเนน วิจิกิจฺฉาย. ญาเณน อวิชฺชาย. ปาโมชฺเชน อรติยา.}

\prose{ปฐเมน ฌาเนน นีวรณานํ. ทุติเยน ฌาเนน วิตกฺก\-วิจารานํ. ตติเยน ฌาเนน ปีติยา. จตุตฺเถน ฌาเนน สุขทุกฺขานํ. อากาสา\-นญฺจา\-ยตนสมาปตฺติยา รูปสญฺญาย ปฏิฆสญฺญาย นานตฺตสญฺญาย. วิญฺญาณญฺจายตน\-สมาปตฺติยา อากาสา\-นญฺจา\-ยตนสญฺญาย. อากิญฺจญฺญายตน\-สมาปตฺติยา วิญฺญณญฺจายตนสญฺญาย. เนว\-สญฺญา\-นา\-สญฺญา\-ยตนสมาปตฺติยา อากิญฺจญฺญายตน\-สญฺญาย.}

\prose{\csromanfolio{46}อนิจฺจา\-นุปสฺสนาย นิจฺจสญฺญาย. ทุกฺขา\-นุปสฺสนาย สุขสญฺญาย, อนตฺตา\-นุปสฺสนาย อตฺตสญฺญาย. นิพฺพิทา\-นุปสฺสนาย นนฺทิยา. วิราคา\-นุปสฺสนาย ราคสฺส. นิโรธา\-นุปสฺสนาย สมุทยสฺส. ปฏินิสฺสคฺคา\-นุปสฺสนาย อาทานสฺส. ขยานุปสฺสนาย ฆนสญฺญาย. วยา\-นุปสฺสนาย อายูหนสฺส. วิปริณามา\-นุปสฺสนาย ธุวสญฺญาย. อนิมิตฺตา\-นุปสฺสนาย นิมิตฺตสฺส. อปฺปณิหิตา\-นุปสฺสนาย ปณิธิยา. สุญฺญตา\-นุปสฺสนาย อภินิเวสสฺส. อธิปญฺญา\-ธมฺม\-วิปสฺสนาย สาราคา\-ภินิเวสสฺส. ยถา\-ภูต\-ญาณ\-ทสฺสเนน สมฺโมหา\-ภินิเวสสฺส. อาทีนวา\-นุปสฺสนาย อาลยา\-ภินิเวสสฺส. ปฏิสงฺขา\-นุปสฺสนาย อปฺปฏิสงฺขาย. วิวฏฺฏนา\-นุปสฺสนาย สญฺโญคา\-ภินิเวสสฺส.}

\prose{โสตาปตฺติ\-มคฺเคน ทิฏฺเฐ\-กฏฺฐานํ กิเลสานํ. สกทาคามิ\-มคฺเคน โอฬาริกานํ กิเลสานํ. อนาคามิมคฺเคน อณุสหคตานํ กิเลสานํ. อรหตฺต\-มคฺเคน สพฺพกิเลสานํ ปหานํ สีลํ, เวรมณี สีลํ, เจตนา สีลํ, สํวโร สีลํ, อวีติกฺกโม สีลํ. เอวรูปานิ สีลานิ จิตฺตสฺส อวิปฺปฏิสาราย สํวตฺตนฺติ, ปาโมชฺชาย สํวตฺตนฺติ, ปีติยา สํวตฺตนฺติ, ปสฺสทฺธิยา สํวตฺตนฺติ, โสมนสฺสาย สํวตฺตนฺติ, อาเสวนาย สํวตฺตนฺติ, ภาวนาย สํวตฺตนฺติ, พหุลีกมฺมาย สํวตฺตนฺติ, อลงฺการาย สํวตฺตนฺติ, ปริกฺขาราย สํวตฺตนฺติ, ปริวาราย สํวตฺตนฺติ, ปาริปูริยา สํวตฺตนฺติ, เอกนฺต\-นิพฺพิทาย วิราคาย นิโรธาย อุปสมาย อภิญฺญาย สมฺโพธาย นิพฺพานาย สํวตฺตนฺติ.}

\proseitem{42}{เอวรูปานํ สีลานํ สํวร\-ปาริสุทฺธิ อธิสีลํ. สํวร\-ปาริสุทฺธิยา ฐิตํ จิตฺตํ อวิกฺเขปํ คจฺฉติ, อวิกฺเขป\-ปาริสุทฺธิ อธิจิตฺตํ. สํวร\-ปาริสุทฺธิํ สมฺมา ปสฺสติ, อวิกฺเขป\-ปาริสุทฺธิํ สมฺมา ปสฺสติ. ทสฺสน\-ปาริสุทฺธิ อธิปญฺญา. โย ตตฺถ สํวรฏฺโฐ, อยํ อธิสีลสิกฺขา. โย ตตฺถ อวิกฺเขปฏฺโฐ, อยํ อธิจิตฺต\-สิกฺขา. โย ตตฺถ ทสฺสนฏฺโฐ, อยํ อธิปญฺญา\-สิกฺขา.}

\prosewithcloserruled{อิมา ติสฺโส สิกฺขาโย อาวชฺชนฺโต สิกฺขติ, ชานนฺโต สิกฺขติ, ปสฺสนฺโต สิกฺขติ, ปจฺจเวกฺขนฺโต สิกฺขติ, จิตฺตํ อธิฏฺฐหนฺโต สิกฺขติ, สทฺธาย อธิมุจฺจนฺโต สิกฺขติ, วีริยํ ปคฺคณฺหนฺโต สิกฺขติ, สติํ อุปฏฺฐเปนฺโต สิกฺขติ, จิตฺตํ สมาทหนฺโต สิกฺขติ, ปญฺญาย ปชานนฺโต สิกฺขติ, อภิญฺเญยฺยํ อภิชานนฺโต สิกฺขติ, ปริญฺเญยฺยํ ปริชานนฺโต สิกฺขติ, ปหาตพฺพํ ปชหนฺโต สิกฺขติ, สจฺฉิกาตพฺพํ สจฺฉิกโรนฺโต สิกฺขติ, \csromanfolio{47}ภาเวตพฺพํ ภาเวนฺโต สิกฺขติ, ตํญาตฏฺเฐน ญาณํ, ปชานนฏฺเฐน ปญฺญา, เตน วุจฺจติ สุตฺวาน สํวเร ปญฺญา สีลมเย ญาณํ.}{สีลมย\-ญาณ\-นิทฺเทโส ทุติโย.}
\csromanmatikaanchor{mtk.5}
\csromantocmarkref{subsection}{3. สมาธิภาวนามยญาณนิทฺเทส}{mtk.5}
\bookmark[dest={mtk.5},level=2]{3. สมาธิภาวนามยญาณนิทฺเทส}
\csromanheader{2}{3. สมาธิ\-ภาวนามย\-ญาณ\-นิทฺเทส}
\proseitem{43}{กถํ สํวริตฺวา สมาทหเน ปญฺญา สมาธิ\-ภาวนามเย ญาณํ\csromandash{}เอโก สมาธิ จิตฺตสฺส เอกคฺคตา.  เทฺว สมาธี, โลกิโย สมาธิ โลกุตฺตโร สมาธิ.  ตโย สมาธี, สวิตกฺก\-สวิจาโร สมาธิ อวิตกฺก\-วิจารมตฺโต สมาธิ อวิตกฺก\-อวิจาโร สมาธิ.  จตฺตาโร สมาธี, หานภาคิโย สมาธิ ฐิติภาคิโย สมาธิ วิเสสภาคิโย สมาธิ นิพฺเพธ\-ภาคิโย สมาธิ.  ปญฺจ สมาธี, ปีติผรณตา สุขผรณตา เจโตผรณตา อาโลกผรณตา ปจฺจเวกฺขณา\-นิมิตฺตํ\csromansharedfootnote{413d2eee79f0eeaa}{ปจฺจเวกฺขนานิมิตฺตํ (ก)}.  ฉ สมาธี, พุทฺธานุสฺสติ\-วเสน จิตฺตสฺส เอกคฺคตา อวิกฺเขโป สมาธิ, ธมฺมานุสฺสติ\-วเสน จิตฺตสฺส เอกคฺคตา อวิกฺเขโป สมาธิ, สํฆานุสฺสติวเสน จิตฺตสฺส เอกคฺคตา อวิกฺเขโป สมาธิ, สีลานุสฺสติ\-วเสน จิตฺตสฺส เอกคฺคตา อวิกฺเขโป สมาธิ, จาคานุสฺสติ\-วเสน จิตฺตสฺส เอกคฺคตา อวิกฺเขโป สมาธิ, เทวตานุสฺสติ\-วเสน จิตฺตสฺส เอกคฺคตา อวิกฺเขโป สมาธิ.  สตฺต สมาธี, สมาธิ\-กุสลตา สมาธิสฺส สมาปตฺติ\-กุสลตา สมาธิสฺส ฐิติกุสลตา สมาธิสฺส วุฏฺฐาน\-กุสลตา สมาธิสฺส กลฺลตา\-กุสลตา\csromansharedfootnote{2dd59b44144c12c7}{กลฺลิตกุสลตา (สฺยา, ก)} สมาธิสฺส โคจรกุสลตา สมาธิสฺส อภินีหาร\-กุสลตา.  อฏฺฐ สมาธี, ปถวีกสิณ\-วเสน จิตฺตสฺส เอกคฺคตา อวิกฺเขโป สมาธิ, อาโปกสิณ\-วเสน ฯเปฯ. เตโชกสิณ\-วเสน. วาโยกสิณ\-วเสน. นีลกสิณ\-วเสน. ปีตกสิณ\-วเสน. โลหิตกสิณ\-วเสน. โอทาตกสิณ\-วเสน จิตฺตสฺส เอกคฺคตา อวิกฺเขโป สมาธิ.  นว สมาธี, รูปาวจโร สมาธิ อตฺถิ หีโน อตฺถิ มชฺโฌ อตฺถิ ปณีโต, อรูปาวจโร สมาธิ อตฺถิ หีโน อตฺถิ มชฺโฌ อตฺถิ ปณีโต, สุญฺญโต สมาธิ อนิมิตฺโต สมาธิ อปฺปณิหิโต สมาธิ.  ทส สมาธี, อุทฺธุมาตก\-สญฺญาวเสน \csromanfolio{48}จิตฺตสฺส เอกคฺคตา อวิกฺเขโป สมาธิ, วินีลกสญฺญา\-วเสน. วิปุพฺพกสญฺญา\-วเสน. วิจฺฉิทฺทก\-สญฺญาวเสน. วิกฺขายิตก\-สญฺญา\-วเสน. วิกฺขิตฺตก\-สญฺญา\-วเสน. หตวิกฺขิตฺตก\-สญฺญา\-วเสน. โลหิตก\-สญฺญา\-วเสน. ปุฬวกสญฺญา\-วเสน\csromansharedfootnote{377a42e4964b8935}{ปุฬุวก สญฺญาวเสน (ก)}. อฏฺฐิกสญฺญา\-วเสน จิตฺตสฺส เอกคฺคตา อวิกฺเขโป สมาธิ. อิเม ปญฺจปญฺญาส สมาธี.}

\proseitemwithcloserruled{44}{อปิ จ ปญฺจวีสติ สมาธิสฺส สมาธิฏฺฐา\csromandash{}ปริคฺคห\-ฏฺเฐน สมาธิ, ปริวารฏฺเฐน สมาธิ, ปริปูรฏฺเฐน สมาธิ, เอกคฺคฏฺเฐน สมาธิ, อวิกฺเขป\-ฏฺเฐน สมาธิ, อนาวิลฏฺเฐน สมาธิ, อนิญฺชนฏฺเฐน สมาธิ, วิมุตฺตฏฺเฐน สมาธิ, เอกตฺตุ\-ปฏฺฐาน\-วเสน จิตฺตสฺส ฐิตตฺตา สมาธิ. สมํ เอสตีติ สมาธิ, วิสมํ เนสตีติ สมาธิ, สมํ เอสิตตฺตา สมาธิ, วิสมํ เนสิตตฺตา สมาธิ, สมํ อาทิยตีติ สมาธิ, วิสมํ นาทิยตีติ สมาธิ, สมํ อาทินฺนตฺตา สมาธิ, วิสมํ อนาทินฺนตฺตา สมาธิ, สมํ ปฏิปชฺชตีติ สมาธิ, วิสมํ นปฺปฏิปชฺชตีติ สมาธิ, สมํ ปฏิปนฺนตฺตา สมาธิ, วิสมํ นปฺปฏิปนฺนตฺตา สมาธิ, สมํ ฌายตีติ สมาธิ, วิสมํ ฌาเปตีติ สมาธิ, สมํ ฌาตตฺตา สมาธิ, วิสมํ ฌาปิตตฺตา สมาธิ. สโม จ หิโต จ สุโข จาติ สมาธิ. อิเม ปญฺจวีสติ สมาธิสฺส สมาธิฏฺฐา. ตํญาตฏฺเฐน ญาณํ, ปชานนฏฺเฐน ปญฺญา, เตน วุจฺจติ สํวริตฺวา สมาทหเน ปญฺญา สมาธิ\-ภาวนามเย ญาณํ.}{สมาธิ\-ภาวนามย\-ญาณ\-นิทฺเทโส ตติโย.}
\csromanmatikaanchor{mtk.6}
\csromantocmarkref{subsection}{4. ธมฺมฏฺฐิติญาณนิทฺเทส}{mtk.6}
\bookmark[dest={mtk.6},level=2]{4. ธมฺมฏฺฐิติญาณนิทฺเทส}
\csromanheader{2}{4. ธมฺมฏฺฐิติญาณ\-นิทฺเทส}
\proseitem{45}{กถํ ปจฺจยปริคฺคเห ปญฺญา ธมฺมฏฺฐิติ\-ญาณํ\csromandash{}อวิชฺชา สงฺขารานํ อุปฺปาทฏฺฐิติ จ ปวตฺตฏฺฐิติ จ นิมิตฺตฏฺฐิติ จ อายูหนฏฺฐิติ จ สญฺโญคฏฺฐิติ จ ปลิโพธ\-ฏฺฐิติ จ สมุทยฏฺฐิติ จ เหตุฏฺฐิติ จ ปจฺจยฏฺฐิติ จ, อิเมหิ นวหากาเรหิ อวิชฺชา ปจฺจโย, สงฺขารา ปจฺจยสมุปฺปนฺนา, อุโภเปเต ธมฺมา ปจฺจยสมุปฺปนฺนาติ ปจฺจยปริคฺคเห ปญฺญา ธมฺมฏฺฐิติ\-ญาณํ. อตีตมฺปิ อทฺธานํ, อนาคตมฺปิ อทฺธานํ อวิชฺชา สงฺขารานํ อุปฺปาทฏฺฐิติ จ ปวตฺตฏฺฐิติ จ นิมิตฺตฏฺฐิติ จ อายูหนฏฺฐิติ จ สญฺโญคฏฺฐิติ จ ปลิโพธ\-ฏฺฐิติ จ สมุทยฏฺฐิติ จ \csromanfolio{49}เหตุฏฺฐิติ จ ปจฺจยฏฺฐิติ จ, อิเมหิ นวหากาเรหิ อวิชฺชา ปจฺจโย, สงฺขารา ปจฺจยสมุปฺปนฺนา, อุโภเปเต ธมฺมา ปจฺจยสมุปฺปนฺนาติ ปจฺจยปริคฺคเห ปญฺญา ธมฺมฏฺฐิติ\-ญาณํ.}

\prose{สงฺขารา วิญฺญาณสฺส ฯเปฯ. วิญฺญาณํ นามรูปสฺส. นามรูปํ สฬายตนสฺส, สฬายตนํ ผสฺสสฺส. ผสฺโส เวทนาย. เวทนา ตณฺหาย. ตณฺหา อุปาทานสฺส. อุปาทานํ ภวสฺส. ภโว ชาติยา. ชาติ ชรามรณสฺส อุปฺปาทฏฺฐิติ จ ปวตฺตฏฺฐิติ จ นิมิตฺตฏฺฐิติ จ อายูหนฏฺฐิติ จ สญฺโญคฏฺฐิติ จ ปลิโพธ\-ฏฺฐิติ จ สมุทยฏฺฐิติ จ เหตุฏฺฐิติ จ ปจฺจยฏฺฐิติ จ, อิเมหิ นวหากาเรหิ ชาติ ปจฺจโย, ชรามรณํ ปจฺจยสมุปฺปนฺนํ, อุโภเปเต ธมฺมา ปจฺจยสมุปฺปนฺนาติ ปจฺจยปริคฺคเห ปญฺญา ธมฺมฏฺฐิติ\-ญาณํ. อตีตมฺปิ อทฺธานํ. อนาคตมฺปิ อทฺธานํ ชาติ ชรามรณสฺส อุปฺปาทฏฺฐิติ จ ปวตฺตฏฺฐิติ จ นิมิตฺตฏฺฐิติ จ อายูหนฏฺฐิติ จ สญฺโญคฏฺฐิติ จ ปลิโพธ\-ฏฺฐิติ จ สมุทยฏฺฐิติ จ เหตุฏฺฐิติ จ ปจฺจยฏฺฐิติ จ, อิเมหิ นวหากาเรหิ ชาติ ปจฺจโย, ชรามรณ ปจฺจยสมุปฺปนฺนํ, อุโภเปเต ธมฺมา ปจฺจยสมุปฺปนฺนาติ ปจฺจยปริคฺคเห ปญฺญา ธมฺมฏฺฐิติ\-ญาณํ.}

\proseitem{46}{อวิชฺชา เหตุ, สงฺขารา เหตุสมุปฺปนฺนา, อุโภเปเต ธมฺมา เหตุสมุปฺปนฺนาติ ปจฺจยปริคฺคเห ปญฺญา ธมฺม\-ฐิติ\-ญาณํ. อตีตมฺปิ อทฺธานํ. อนาคตมฺปิ อทฺธานํ อวิชฺชา เหตุ, สงฺขารา เหตุสมุปฺปนฺนา, อุโภเปเต ธมฺมา เหตุสมุปฺปนฺนาติ ปจฺจยปริคฺคเห ปญฺญา ธมฺมฏฺฐิติ\-ญาณํ.}

\prose{สงฺขารา เหตุ, วิญฺญาณํ เหตุ\-สมุปฺปนฺนํ ฯเปฯ. วิญฺญาณํ เหตุ, นามรูปํ เหตุ\-สมุปฺปนฺนํ. นามรูปํ เหตุ, สฬายตนํ เหตุ\-สมุปฺปนฺนํ. สฬายตนํ เหตุ, ผสฺโส เหตุสมุปฺปนฺโน. ผสฺโส เหตุ, เวทนา เหตุสมุปฺปนฺนา. เวทนา เหตุ, ตณฺหา เหตุสมุปฺปนฺนา. ตณฺหา เหตุ, อุปาทานํ เหตุ\-สมุปฺปนฺนํ. อุปาทานํ เหตุ, ภโว เหตุสมุปฺปนฺโน. ภโว เหตุ, ชาติ เหตุสมุปฺปนฺนา. ชาติ เหตุ, ชรามรณํ เหตุ\-สมุปฺปนฺนํ, อุโภเปเต ธมฺมา เหตุสมุปฺปนฺนาติ ปจฺจยปริคฺคเห ปญฺญา ธมฺมฏฺฐิติ\-ญาณํ. อตีตมฺปิ อทฺธานํ. อนาคตมฺปิ อทฺธานํ ชาติ เหตุ, ชรามรณํ เหตุ\-สมุปฺปนฺนํ, อุโภเปเต ธมฺมา เหตุสมุปฺปนฺนาติ ปจฺจยปริคฺคเห ปญฺญา ธมฺมฏฺฐิติ\-ญาณํ.}

\prose{อวิชฺชา ปฏิจฺจา, สงฺขารา ปฏิจฺจ\-สมุปฺปนฺนา, อุโภเปเต ธมฺมา ปฏิจฺจ\-สมุปฺปนฺนาติ ปจฺจยปริคฺคเห ปญฺญา ธมฺมฏฺฐิติ\-ญาณํ. อตีตมฺปิ อทฺธานํ. อนาคตมฺปิ อทฺธานํ \csromanfolio{50}อวิชฺชา ปฏิจฺจา, สงฺขารา ปฏิจฺจ\-สมุปฺปนฺนา, อุโภเปเต ธมฺมา ปฏิจฺจ\-สมุปฺปนฺนาติ ปจฺจยปริคฺคเห ปญฺญา ธมฺมฏฺฐิติ\-ญาณํ.}

\prose{สงฺขารา ปฏิจฺจา, วิญฺญาณํ ปฏิจฺจ\-สมุปฺปนฺนํ ฯเปฯ. วิญฺญาณํ ปฏิจฺจํ, นามรูปํ ปฏิจฺจ\-สมุปฺปนฺนํ. นามรูปํ ปฏิจฺจํ, สฬายตนํ ปฏิจฺจ\-สมุปฺปนฺนํ. สฬายตนํ ปฏิจฺจํ, ผสฺโส ปฏิจฺจ\-สมุปฺปนฺโน. ผสฺโส ปฏิจฺโจ, เวทนา ปฏิจฺจ\-สมุปฺปนฺนา. เวทนา ปฏิจฺจา, ตณฺหา ปฏิจฺจ\-สมุปฺปนฺนา. ตณฺหา ปฏิจฺจา, อุปาทานํ ปฏิจฺจ\-สมุปฺปนฺนํ. อุปาทานํ ปฏิจฺจํ, ภโว ปฏิจฺจ\-สมุปฺปนฺโน. ภโว ปฏิจฺโจ, ชาติ ปฏิจฺจ\-สมุปฺปนฺนา. ชาติ ปฏิจฺจา, ชรามรณํ ปฏิจฺจ\-สมุปฺปนฺนํ, อุโภเปเต ธมฺมา ปฏิจฺจ\-สมุปฺปนฺนาติ ปจฺจยปริคฺคเห ปญฺญา ธมฺมฏฺฐิติ\-ญาณํ. อตีตมฺปิ อทฺธานํ. อนาคตมฺปิ อทฺธานํ ชาติ ปฏิจฺจา, ชรามรณํ ปฏิจฺจ\-สมุปฺปนฺนํ, อุโภเปเต ธมฺมา ปฏิจฺจ\-สมุปฺปนฺนาติ ปจฺจยปริคฺคเห ปญฺญา ธมฺมฏฺฐิติ\-ญาณํ.}

\prose{อวิชฺชา ปจฺจโย, สงฺขารา ปจฺจยสมุปฺปนฺนา, อุโภเปเต ธมฺมา ปจฺจยสมุปฺปนฺนาติ ปจฺจยปริคฺคเห ปญฺญา ธมฺมฏฺฐิติ\-ญาณํ. อตีตมฺปิ อทฺธานํ. อนาคตมฺปิ อทฺธานํ อวิชฺชา ปจฺจโย, สงฺขารา ปจฺจยสมุปฺปนฺนา, อุโภเปเต ธมฺมา ปจฺจยสมุปฺปนฺนาติ ปจฺจยปริคฺคเห ปญฺญา ธมฺมฏฺฐิติ\-ญาณํ.}

\prose{สงฺขารา ปจฺจยา, วิญฺญาณํ ปจฺจยสมุปฺปนฺนํ ฯเปฯ. วิญฺญาณํ ปจฺจโย, นามรูปํ ปจฺจยสมุปฺปนฺนํ. นามรูปํ ปจฺจโย, สฬายตนํ ปจฺจยสมุปฺปนฺนํ. สฬายตนํ ปจฺจโย, ผสฺโส ปจฺจยสมุปฺปนฺโน. ผสฺโส ปจฺจโย, เวทนา ปจฺจยสมุปฺปนฺนา. เวทนา ปจฺจโย, ตณฺหา ปจฺจยสมุปฺปนฺนา. ตณฺหา ปจฺจโย, อุปาทานํ ปจฺจยสมุปฺปนฺนํ. อุปาทานํ ปจฺจโย, ภโว ปจฺจยสมุปฺปนฺโน. ภโว ปจฺจโย, ชาติ ปจฺจยสมุปฺปนฺนา. ชาติ ปจฺจโย, ชรามรณํ ปจฺจยสมุปฺปนฺนํ, อุโภเปเต ธมฺมา ปจฺจยสมุปฺปนฺนาติ ปจฺจยปริคฺคเห ปญฺญา ธมฺมฏฺฐิติ\-ญาณํ. อตีตมฺปิ อทฺธานํ. อนาคตมฺปิ อทฺธานํ ชาติ ปจฺจโย, ชรามรณํ ปจฺจยสมุปฺปนฺนํ, อุโภเปเต ธมฺมา ปจฺจยสมุปฺปนฺนาติ ปจฺจยปริคฺคเห ปญฺญา ธมฺมฏฺฐิติ\-ญาณํ.}

\proseitemwithcloserruled{47}{ปุริม\-กมฺมภวสฺมิํ โมโห อวิชฺชา, อายูหนา สงฺขารา, นิกนฺติ ตณฺหา, อุปคมนํ อุปาทานํ, เจตนา ภโว, อิเม ปญฺจ ธมฺมา ปุริม\-กมฺมภวสฺมิํ อิธ ปฏิสนฺธิยา ปจฺจยา. อิธ ปฏิสนฺธิ วิญฺญาณํ, โอกฺกนฺติ นามรูปํ, ปสาโท อายตนํ, ผุฏฺโฐ ผสฺโส, เวทยิตํ เวทนา, อิเม ปญฺจ ธมฺมา อิธุปปตฺติภวสฺมิํ ปุเรกตสฺส กมฺมสฺส ปจฺจยา. อิธ ปริปกฺกตฺตา อายตนานํ โมโห อวิชฺชา, อายูหนา สงฺขารา, นิกนฺติ ตณฺหา, อุปคมนํ อุปาทานํ, \csromanfolio{51}เจตนา ภโว, อิเม ปญฺจ ธมฺมา อิธ กมฺมภวสฺมิํ อายติํ ปฏิสนฺธิยา ปจฺจยา. อายติํ ปฏิสนฺธิ วิญฺญาณํ, โอกฺกนฺติ นามรูปํ, ปสาโท อายตนํ, ผุฏฺโฐ ผสฺโส, เวทยิตํ เวทนา, อิเม ปญฺจ ธมฺมา อายติํ อุปปตฺติภวสฺมิํ อิธ กตสฺส กมฺมสฺส ปจฺจยา. อิติเม จตุสงฺเขเป ตโย อทฺเธ ติสนฺธิํ วีสติยา อากาเรหิ ปฏิจฺจ\-สมุปฺปาทํ ชานาติ ปสฺสติ อญฺญาติ ปฏิวิชฺฌติ, ตํญาตฏฺเฐน ญาณํ, ปชานนฏฺเฐน ปญฺญา, เตน วุจฺจติ ปจฺจยปริคฺคเห ปญฺญา ธมฺมฏฺฐิติ\-ญาณํ.}{ธมฺมฏฺฐิติญาณ\-นิทฺเทโส จตุตฺโถ.}
\csromanmatikaanchor{mtk.7}
\csromantocmarkref{subsection}{5. สมฺมสนญาณนิทฺเทส}{mtk.7}
\bookmark[dest={mtk.7},level=2]{5. สมฺมสนญาณนิทฺเทส}
\csromanheader{2}{5. สมฺมสน\-ญาณ\-นิทฺเทส}
\proseitem{48}{กถํ อตีตา\-นาคต\-ปจฺจุปฺปนฺนานํ ธมฺมานํ สงฺขิปิตฺวา ววตฺถาเน ปญฺญา สมฺมสเน ญาณํ\csromandash{}}

\prose{ยํ กิญฺจิ รูปํ อตีตา\-นาคต\-ปจฺจุปฺปนฺนํ อชฺฌตฺตํ วา พหิทฺธา วา โอฬาริกํ วา สุขุมํ วา หีนํ วา ปณีตํ วา ยํ ทูเร สนฺติเก วา, สพฺพํ รูปํ อนิจฺจโต ววตฺเถติ เอกํ สมฺมสนํ, ทุกฺขโต ววตฺเถติ เอกํ สมฺมสนํ, อนตฺตโต ววตฺเถติ เอกํ สมฺมสนํ.}

\prose{ยา กาจิ เวทนา ฯเปฯ. ยา กาจิ สญฺญา. เย เกจิ สงฺขารา. ยํ กิญฺจิ วิญฺญาณํ อตีตา\-นาคต\-ปจฺจุปฺปนฺนํ อชฺฌตฺตํ วา พหิทฺธา วา โอฬาริกํ วา สุขุมํ วา หีนํ วา ปณีตํ วา ยํ ทูเร สนฺติเก วา, สพฺพํ วิญฺญาณํ อนิจฺจโต ววตฺเถติ เอกํ สมฺมสนํ, ทุกฺขโต ววตฺเถติ เอกํ สมฺมสนํ, อนตฺตโต ววตฺเถติ เอกํ สมฺมสนํ.}

\prose{จกฺขุํ ฯเปฯ. ชรามรณํ อตีตา\-นาคต\-ปจฺจุปฺปนฺนํ อนิจฺจโต ววตฺเถติ เอกํ สมฺมสนํ, ทุกฺขโต ววตฺเถติ เอกํ สมฺมสนํ, อนตฺตโต ววตฺเถติ เอกํ สมฺมสนํ.}

\prose{“รูปํ อตีตา\-นาคต\-ปจฺจุปฺปนฺนํ อนิจฺจํ ขยฏฺเฐน, ทุกฺขํ ภยฏฺเฐน, อนตฺตา อสารกฏฺเฐนา”ติ สงฺขิปิตฺวา ววตฺถาเน ปญฺญา สมฺมสเน ญาณํ. “เวทนา ฯเปฯ. สญฺญา ฯเปฯ. สงฺขารา ฯเปฯ. วิญฺญาณํ ฯเปฯ. จกฺขุ ฯเปฯ. ชรามรณํ อตีตา\-นาคต\-ปจฺจุปฺปนฺนํ อนิจฺจํ ขยฏฺเฐน ทุกฺขํ ภยฏฺเฐน อนตฺตา อสารกฏฺเฐนา”ติ สงฺขิปิตฺวา ววตฺถาเน ปญฺญา สมฺมสเน ญาณํ.}

\prose{\csromanfolio{52}“รูปํ อตีตา\-นาคต\-ปจฺจุปฺปนฺนํ อนิจฺจํ สงฺขตํ ปฏิจฺจ\-สมุปฺปนฺนํ ขยธมฺมํ วยธมฺมํ วิราคธมฺมํ นิโรธธมฺมนฺ”ติ สงฺขิปิตฺวา ววตฺถาเน ปญฺญา สมฺมสเน ญาณํ. “เวทนา ฯเปฯ. สญฺญา. สงฺขารา. วิญฺญาณํ.  จกฺขุ ฯเปฯ. ชรามรณํ อตีตา\-นาคต\-ปจฺจุปฺปนฺนํ อนิจฺจํ สงฺขตํ ปฏิจฺจ\-สมุปฺปนฺนํ ขยธมฺมํ วยธมฺมํ วิราคธมฺมํ นิโรธธมฺมนฺ”ติ สงฺขิปิตฺวา ววตฺถาเน ปญฺญา สมฺมสเน ญาณํ.}

\prosewithcloserruled{“ชาติปจฺจยา ชรามรณํ, อสติ ชาติยา นตฺถิ ชรามรณนฺ”ติ สงฺขิปิตฺวา ววตฺถาเน ปญฺญา สมฺมสเน ญาณํ. อตีตมฺปิ อทฺธานํ. อนาคตมฺปิ อทฺธานํ “ชาติปจฺจยา ชรามรณํ, อสติ ชาติยา นตฺถิ ชรามรณนฺ”ติ สงฺขิปิตฺวา ววตฺถาเน ปญฺญา สมฺมสเน ญณํ.  ภวปฺปจฺจยา ชาติ, อสติ ฯเปฯ. อุปาทานปจฺจยา ภโว, อสติ ฯเปฯ. ตณฺหาปจฺจยา อุปาทานํ, อสติ ฯเปฯ. เวทนาปจฺจยา ตณฺหา, อสติ ฯเปฯ. ผสฺสปจฺจยา เวทนา, อสติ ฯเปฯ. สฬายตน\-ปจฺจยา ผสฺโส, อสติ ฯเปฯ. นามรูป\-ปจฺจยา สฬายตนํ, อสติ ฯเปฯ. วิญฺญาณปจฺจยา นามรูปํ, อสติ ฯเปฯ. สงฺขาร\-ปจฺจยา วิญฺญาณํ, อสติ ฯเปฯ. “อวิชฺชาปจฺจยา สงฺขารา, อสติ อวิชฺชาย นตฺถิ สงฺขารา”ติ สงฺขิปิตฺวา ววตฺถาเน ปญฺญา สมฺมสเน ญาณํ. “อตีตมฺปิ อทฺธานํ. อนาคตมฺปิ อทฺธานํ อวิชฺชาปจฺจยา สงฺขารา, อสติ อวิชฺชาย นตฺถิ สงฺขารา”ติ สงฺขิปิตฺวา ววตฺถาเน ปญฺญา สมฺมสเน ญาณํ, ตํญาตฏฺเฐน ญาณํ, ปชานนฏฺเฐน ปญฺญา, เตน วุจฺจติ อตีตา\-นาคต\-ปจฺจุปฺปนฺนานํ ธมฺมานํ สงฺขิปิตฺวา ววตฺถาเน ปญฺญา สมฺมสเน ญาณํ.}{สมฺมสน\-ญาณ\-นิทฺเทโส ปญฺจโม.}
\csromanmatikaanchor{mtk.8}
\csromantocmarkref{subsection}{6. อุทยพฺพยญาณนิทฺเทส}{mtk.8}
\bookmark[dest={mtk.8},level=2]{6. อุทยพฺพยญาณนิทฺเทส}
\csromanheader{2}{6. อุทยพฺพยญาณ\-นิทฺเทส}
\proseitem{49}{กถํ ปจฺจุปฺปนฺนานํ ธมฺมานํ วิปริณามา\-นุปสฺสเน ปญฺญา อุทยพฺพยา\-นุปสฺสเน ญาณํ\csromandash{}ชาตํ รูปํ ปจฺจุปฺปนฺนํ, ตสฺส นิพฺพตฺติ\-ลกฺขณํ อุทโย, วิปริณาม\-ลกฺขณํ วโย, อนุปสฺสนา ญาณํ. ชาตา เวทนา. ชาตา สญฺญา. ชาตา สงฺขารา. ชาตํ วิญฺญาณํ. ชาตํ จกฺขุ ฯเปฯ. ชาโต ภโว ปจฺจุปฺปนฺโน, ตสฺส นิพฺพตฺติ\-ลกฺขณํ อุทโย, วิปริณาม\-ลกฺขณํ วโย, อนุปสฺสนา ญาณํ.}

\proseitem{50}{ปญฺจนฺนํ ขนฺธานํ อุทยํ ปสฺสนฺโต กติ ลกฺขณานิ ปสฺสติ, วยํ ปสฺสนฺโต กติ ลกฺขณานิ ปสฺสติ, อุทยพฺพยํ ปสฺสนฺโต กติ \csromanfolio{53}ลกฺขณานิ ปสฺสติ? ปญฺจนฺนํ ขนฺธานํ อุทยํ ปสฺสนฺโต ปญฺจวีสติ ลกฺขณานิ ปสฺสติ, วยํ ปสฺสนฺโต ปญฺจวีสติ ลกฺขณานิ ปสฺสติ, อุทยพฺพยํ ปสฺสนฺโต ปญฺญาส ลกฺขณานิ ปสฺสติ.}

\prose{รูปกฺขนฺธสฺส อุทยํ ปสฺสนฺโต กติ ลกฺขณานิ ปสฺสติ, วยํ ปสฺสนฺโต กติ ลกฺขณานิ ปสฺสติ, อุทยพฺพยํ ปสฺสนฺโต กติ ลกฺขณานิ ปสฺสติ? เวทนา\-กฺขนฺธสฺส ฯเปฯ. สญฺญา\-กฺขนฺธสฺส ฯเปฯ. สงฺขาร\-กฺขนฺธสฺส ฯเปฯ. วิญฺญาณ\-กฺขนฺธสฺส อุทยํ ปสฺสนฺโต กติ ลกฺขณานิ ปสฺสติ, วยํ ปสฺสนฺโต กติ ลกฺขณานิ ปสฺสติ, อุทยพฺพยํ ปสฺสนฺโต กติ ลกฺขณานิ ปสฺสติ? รูปกฺขนฺธสฺส อุทยํ ปสฺสนฺโต ปญฺจ ลกฺขณานิ ปสฺสติ, วยํ ปสฺสนฺโต ปญฺจ ลกฺขณานิ ปสฺสติ, อุทยพฺพยํ ปสฺสนฺโต ทส ลกฺขณานิ ปสฺสติ.  เวทนา\-กฺขนฺธสฺส ฯเปฯ. สญฺญา\-กฺขนฺธสฺส. สงฺขาร\-กฺขนฺธสฺส. วิญฺญาณ\-กฺขนฺธสฺส อุทยํ ปสฺสนฺโต ปญฺจ ลกฺขณานิ ปสฺสติ, วยํ ปสฺสนฺโต ปญฺจ ลกฺขณานิ ปสฺสติ, อุทยพฺพยํ ปสฺสนฺโต ทส ลกฺขณานิ ปสฺสติ.}

\prose{รูปกฺขนฺธสฺส อุทยํ ปสฺสนฺโต กตมานิ ปญฺจ ลกฺขณานิ ปสฺสติ? “อวิชฺชาสมุทยา รูปสมุทโย”ติ ปจฺจยสมุทย\-ฏฺเฐน รูปกฺขนฺธสฺส อุทยํ ปสฺสติ, “ตณฺหาสมุทยา รูปสมุทโย”ติ ปจฺจยสมุทย\-ฏฺเฐน รูปกฺขนฺธสฺส อุทยํ ปสฺสติ, “กมฺมสมุทยา รูปสมุทโย”ติ ปจฺจยสมุทย\-ฏฺเฐน รูปกฺขนฺธสฺส อุทยํ ปสฺสติ, “อาหารสมุทยา รูปสมุทโย”ติ ปจฺจยสมุทย\-ฏฺเฐน รูปกฺขนฺธสฺส อุทยํ ปสฺสติ. นิพฺพตฺติ\-ลกฺขณํ ปสฺสนฺโตปิ รูปกฺขนฺธสฺส อุทยํ ปสฺสติ. รูปกฺขนฺธสฺส อุทยํ ปสฺสนฺโต อิมานิ ปญฺจ ลกฺขณานิ ปสฺสติ.}

\prose{วยํ ปสฺสนฺโต กตมานิ ปญฺจ ลกฺขณานิ ปสฺสติ? “อวิชฺชานิโรธา รูปนิโรโธ”ติ ปจฺจยนิโรธ\-ฏฺเฐน รูปกฺขนฺธสฺส วยํ ปสฺสติ, “ตณฺหา นิโรธา รูปนิโรโธ”ติ ปจฺจยนิโรธ\-ฏฺเฐน รูปกฺขนฺธสฺส วยํ ปสฺสติ, “กมฺมนิโรธา รูปนิโรโธ”ติ ปจฺจยนิโรธ\-ฏฺเฐน รูปกฺขนฺธสฺส วยํ ปสฺสติ, “อาหารนิโรธา รูปนิโรโธ”ติ ปจฺจยนิโรธ\-ฏฺเฐน รูปกฺขนฺธสฺส วยํ ปสฺสติ, วิปริณาม\-ลกฺขณํ ปสฺสนฺโตปิ รูปกฺขนฺธสฺส วยํ ปสฺสติ. รูปกฺขนฺธสฺส วยํ ปสฺสนฺโต อิมานิ ปญฺจ ลกฺขณานิ ปสฺสติ. อุทยพฺพยํ ปสฺสนฺโต อิมานิ ทส ลกฺขณานิ ปสฺสติ.}

\prose{\csromanfolio{54}เวทนา\-กฺขนฺธสฺส อุทยํ ปสฺสนฺโต กตมานิ ปญฺจ ลกฺขณานิ ปสฺสติ? “อวิชฺชาสมุทยา เวทนาสมุทโย”ติ ปจฺจยสมุทย\-ฏฺเฐน เวทนา\-กฺขนฺธสฺส อุทยํ ปสฺสติ, “ตณฺหาสมุทยา เวทนาสมุทโย”ติ ปจฺจยสมุทย\-ฏฺเฐน เวทนา\-กฺขนฺธสฺส อุทยํ ปสฺสติ, “กมฺมสมุทยา เวทนาสมุทโย”ติ ปจฺจยสมุทย\-ฏฺเฐน เวทนา\-กฺขนฺธสฺส อุทยํ ปสฺสติ, “ผสฺสสมุทยา เวทนาสมุทโย”ติ ปจฺจยสมุทย\-ฏฺเฐน เวทนา\-กฺขนฺธสฺส อุทยํ ปสฺสติ. นิพฺพตฺติ\-ลกฺขณํ ปสฺสนฺโตปิ เวทนา\-กฺขนฺธสฺส อุทยํ ปสฺสติ. เวทนา\-กฺขนฺธสฺส อุทยํ ปสฺสนฺโต อิมานิ ปญฺจ ลกฺขณานิ ปสฺสติ.}

\prose{วยํ ปสฺสนฺโต กตมานิ ปญฺจ ลกฺขณานิ ปสฺสติ? “อวิชฺชานิโรธา เวทนานิโรโธ”ติ ปจฺจยนิโรธ\-ฏฺเฐน เวทนา\-กฺขนฺธสฺส วยํ ปสฺสติ, “ตณฺหานิโรธา เวทนานิโรโธ”ติ ปจฺจยนิโรธ\-ฏฺเฐน เวทนา\-กฺขนฺธสฺส วยํ ปสฺสติ, “กมฺมนิโรธา เวทนานิโรโธ”ติ ปจฺจยนิโรธ\-ฏฺเฐน เวทนา\-กฺขนฺธสฺส วยํ ปสฺสติ, “ผสฺสนิโรธา เวทนานิโรโธ”ติ ปจฺจยนิโรธ\-ฏฺเฐน เวทนา\-กฺขนฺธสฺส วยํ ปสฺสติ. วิปริณาม\-ลกฺขณํ ปสฺสนฺโตปิ เวทนา\-กฺขนฺธสฺส วยํ ปสฺสติ. เวทนา\-กฺขนฺธสฺส วยํ ปสฺสนฺโต อิมานิ ปญฺจ ลกฺขณานิ ปสฺสติ.  อุทยพฺพยํ ปสฺสนฺโต อิมานิ ทส ลกฺขณานิ ปสฺสติ.}

\prose{สญฺญา\-กฺขนฺธสฺส ฯเปฯ. สงฺขาร\-กฺขนฺธสฺส ฯเปฯ. วิญฺญาณ\-กฺขนฺธสฺส อุทยํ ปสฺสนฺโต กตมานิ ปญฺจ ลกฺขณานิ ปสฺสติ? “อวิชฺชาสมุทยา วิญฺญาณ\-สมุทโย”ติ ปจฺจยสมุทย\-ฏฺเฐน วิญฺญาณ\-กฺขนฺธสฺส อุทยํ ปสฺสติ, “ตณฺหาสมุทยา วิญฺญาณ\-สมุทโย”ติ ปจฺจยสมุทย\-ฏฺเฐน วิญฺญาณ\-กฺขนฺธสฺส อุทยํ ปสฺสติ, “กมฺมสมุทยา วิญฺญาณ\-สมุทโย”ติ ปจฺจยสมุทย\-ฏฺเฐน วิญฺญาณ\-กฺขนฺธสฺส อุทยํ ปสฺสติ, “นามรูป\-สมุทยา วิญฺญาณ\-สมุทโย”ติ ปจฺจยสมุทย\-ฏฺเฐน วิญฺญาณ\-กฺขนฺธสฺส อุทยํ ปสฺสติ. นิพฺพตฺติ\-ลกฺขณํ ปสฺสนฺโตปิ วิญฺญาณ\-กฺขนฺธสฺส อุทยํ ปสฺสติ. วิญฺญาณ\-กฺขนฺธสฺส อุทยํ ปสฺสนฺโต อิมานิ ปญฺจ ลกฺขณานิ ปสฺสติ.}

\prose{วยํ ปสฺสนฺโต กตมานิ ปญฺจ ลกฺขณานิ ปสฺสติ? “อวิชฺชานิโรธา วิญฺญาณนิโรโธ”ติ ปจฺจยนิโรธ\-ฏฺเฐน วิญฺญาณ\-กฺขนฺธสฺส วยํ ปสฺสติ, “ตณฺหานิโรธา วิญฺญาณนิโรโธ”ติ ปจฺจยนิโรธ\-ฏฺเฐน วิญฺญาณ\-กฺขนฺธสฺส วยํ ปสฺสติ, “กมฺมนิโรธา วิญฺญาณนิโรโธ”ติ ปจฺจยนิโรธ\-ฏฺเฐน วิญฺญาณ\-กฺขนฺธสฺส วยํ ปสฺสติ, “นามรูป\-นิโรธา วิญฺญาณนิโรโธ”ติ ปจฺจยนิโรธ\-ฏฺเฐน วิญฺญาณ\-กฺขนฺธสฺส วยํ ปสฺสติ, วิปริณาม\-ลกฺขณํ ปสฺสนฺโตปิ \csromanfolio{55}วิญฺญาณ\-กฺขนฺธสฺส วยํ ปสฺสติ. วิญฺญาณ\-กฺขนฺธสฺส วยํ ปสฺสนฺโต อิมานิ ปญฺจ ลกฺขณานิ ปสฺสติ.  อุทยพฺพยํ ปสฺสนฺโต อิมานิ ทส ลกฺขณานิ ปสฺสติ.}

\prosewithcloserruled{ปญฺจนฺนํ ขนฺธานํ อุทยํ ปสฺสนฺโต อิมานิ ปญฺจวีสติ ลกฺขณานิ ปสฺสติ, วยํ ปสฺสนฺโต อิมานิ ปญฺจวีสติ ลกฺขณานิ ปสฺสติ, อุทยพฺพยํ ปสฺสนฺโต อิมานิ ปญฺญาส ลกฺขณานิ ปสฺสติ. ตํญาตฏฺเฐน ญาณํ, ปชานนฏฺเฐน ปญฺญา, เตน วุจฺจติ ปจฺจุปฺปนฺนานํ ธมฺมานํ วิปริณามา\-นุปสฺสเน ปญฺญา อุทยพฺพยา\-นุปสฺสเน ญาณํ.  รูปกฺขนฺโธ\csromansharedfootnote{c2f739198f04f6f2}{รูปกฺขนฺธา (สฺยา)} อาหารสมุทโย, เวทนาสญฺญา สงฺขารา ตโย\csromansharedfootnote{f0d0dab0feebfff0}{สงฺขาราติ เสสา (สฺยา)} ขนฺธา ผสฺสสมุทยา. วิญฺญาณ\-กฺขนฺโธ นามรูป\-สมุทโย.}{อุทยพฺพยญาณนิทฺเทสฺโส ฉฏฺโฐ.}
\csromanmatikaanchor{mtk.9}
\csromantocmarkref{subsection}{7. ภงฺคานุปสฺสนาญาณนิทฺเทส}{mtk.9}
\bookmark[dest={mtk.9},level=2]{7. ภงฺคานุปสฺสนาญาณนิทฺเทส}
\csromanheader{2}{7. ภงฺคานุปสฺสนา\-ญาณ\-นิทฺเทส}
\proseitem{51}{กถํ อารมฺมณํ ปฏิสงฺขา ภงฺคานุปสฺสเน ปญฺญา วิปสฺสเน ญาณํ\csromandash{}รูปารมฺมณตา จิตฺตํ อุปฺปชฺชิตฺวา ภิชฺชติ, ตํ อารมฺมณํ ปฏิสงฺขา ตสฺส จิตฺตสฺส ภงฺค อนุปสฺสติ.  \textbf{อนุปสฺสตี}ติ กถํ อนุปสฺสติ? อนิจฺจโต อนุปสฺสติ, โน นิจฺจโต. ทุกฺขโต อนุปสฺสติ, โน สุขโต. อนตฺตโต อนุปสฺสติ, โน อตฺตโต. นิพฺพินฺทติ, โน นนฺทติ. วิรชฺชติ, โน รชฺชติ. นิโรเธติ, โน สมุเทติ. ปฏินิสฺสชฺชติ, โน อาทิยติ.}

\proseitem{52}{อนิจฺจโต อนุปสฺสนฺโต นิจฺจสญฺญํ ปชหาติ, ทุกฺขโต อนุปสฺสนฺโต สุขสญฺญํ ปชหาติ, อนตฺตโต อนุปสฺสนฺโต อตฺตสญฺญํ ปชหาติ, นิพฺพินฺทนฺโต นนฺทิํ ปชหาติ, วิรชฺชนฺโต ราคํ ปชหาติ, นิโรเธนฺโต สมุทยํ ปชหาติ, ปฏินิสฺสชฺชนฺโต อาทานํ ปชหาติ.}

\prose{เวทนา\-รมฺมณตา ฯเปฯ. สญฺญารมฺมณตา. สงฺขารารมฺมณตา. วิญฺญาณา\-รมฺมณตา.  จกฺขุ ฯเปฯ. ชรามรณา\-รมฺมณตา จิตฺตํ อุปฺปชฺชิตฺวา ภิชฺชติ, ตํ อารมฺมณํ ปฏิสงฺขา ตสฺส จิตฺตสฺส ภงฺคํ อนุปสฺสติ.  อนุปสฺสตีติ กถํ อนุปสฺสติ? อนิจฺจโต อนุปสฺสติ, โน นิจฺจโต. ทุกฺขโต อนุปสฺสติ, \csromanfolio{56}โน สุขโต. อนตฺตโต อนุปสฺสติ, โน อตฺตโต. นิพฺพินฺทติ, โน นนฺทติ. วิรชฺชติ, โน รชฺชติ. นิโรเธติ, โน สมุเทติ. ปฏินิสฺสชฺชติ, โน อาทิยติ.}

\prose{อนิจฺจโต อนุปสฺสนฺโต นิจฺจสญฺญํ ปชหาติ, ทุกฺขโต อนุปสฺสนฺโต สุขสญฺญํ ปชหาติ, อนตฺตโต อนุปสฺสนฺโต อตฺตสญฺญํ ปชหาติ, นิพฺพินฺทนฺโต นนฺทิํ ปชหาติ, วิรชฺชนฺโต ราคํ ปชหาติ, นิโรเธนฺโต สมุทยํ ปชหาติ, ปฏินิสฺสชฺชนฺโต อาทานํ ปชหาติ.}

\setlength{\csromangathapagewidth}{0pt}
\setlength{\csromangathawakwidth}{0pt}
\csromangathameasuregroup{วตฺถุสงฺกมนา เจว, \\ อาวชฺชนาพลญฺเจว, \\ อารมฺมณอนฺวเยน, \\ นิโรเธ อธิมุตฺตตา, \\ อารมฺมณญฺจ ปฏิสงฺขา, \\ สุญฺญโต จ อุปฏฺฐานํ, \\ กุสโล ตีสุ อนุปสฺสนาสุ, \\ ตโย อุปฏฺฐาเน กุสลตา,}{\csromangathabat{วตฺถุสงฺกมนา เจว,}{ปญฺญาย จ วิวฏฺฏนา.} \\ \csromangathabat{อาวชฺชนาพลญฺเจว,}{ปฏิสงฺขา วิปสฺสนา.} \\ \csromangathabat{อารมฺมณอนฺวเยน,}{อุโภ เอกววตฺถนา\textsuperscript{1}.} \\ \csromangathabat{นิโรเธ อธิมุตฺตตา,}{วยลกฺขณวิปสฺสนา.} \\ \csromangathabat{อารมฺมณญฺจ ปฏิสงฺขา,}{ภงฺคญฺจ อนุปสฺสติ.} \\ \csromangathabat{สุญฺญโต จ อุปฏฺฐานํ,}{อธิปญฺญา วิปสฺสนา.} \\ \csromangathabat{กุสโล ตีสุ อนุปสฺสนาสุ,}{จตสฺโส จ\textsuperscript{1} วิปสฺสนาสุ.} \\ \csromangathabat{ตโย อุปฏฺฐาเน กุสลตา,}{นานาทิฏฺฐีสุ น กมฺปตีติ.}}
\csromangathasetleft{วตฺถุสงฺกมนา เจว, \\ อาวชฺชนาพลญฺเจว, \\ อารมฺมณอนฺวเยน, \\ นิโรเธ อธิมุตฺตตา, \\ อารมฺมณญฺจ ปฏิสงฺขา, \\ สุญฺญโต จ อุปฏฺฐานํ, \\ กุสโล ตีสุ อนุปสฺสนาสุ, \\ ตโย อุปฏฺฐาเน กุสลตา,}
\csromangathagroup{\csromangathastanza{\csromangathabat{วตฺถุสงฺกมนา เจว,}{ปญฺญาย จ วิวฏฺฏนา.} \\ \csromangathabat{อาวชฺชนาพลญฺเจว,}{ปฏิสงฺขา วิปสฺสนา.}}\csromangathastanza{\csromangathabat{อารมฺมณอนฺวเยน,}{อุโภ เอกววตฺถนา\csromansharedfootnote{edb98764657ee704}{เอกววตฺถานา (ก)}.} \\ \csromangathabat{นิโรเธ อธิมุตฺตตา,}{วยลกฺขณวิปสฺสนา.}}\csromangathastanza{\csromangathabat{อารมฺมณญฺจ ปฏิสงฺขา,}{ภงฺคญฺจ อนุปสฺสติ.} \\ \csromangathabat{สุญฺญโต จ อุปฏฺฐานํ,}{อธิปญฺญา วิปสฺสนา.}}\csromangathastanza{\csromangathabat{กุสโล ตีสุ อนุปสฺสนาสุ,}{จตสฺโส จ\csromansharedfootnote{1bfa87c3f5a8f268}{จตูสุ จ (สฺยา)} วิปสฺสนาสุ.} \\ \csromangathabat{ตโย อุปฏฺฐาเน กุสลตา,}{นานาทิฏฺฐีสุ น กมฺปตีติ.}}}

\prosewithcloserruled{ตํญาตฏฺเฐน ญาณํ, ปชานนฏฺเฐน ปญฺญา, เตน วุจฺจติ อารมฺมณํ ปฏิสงฺขา ภงฺคานุปสฺสเน ปญฺญา วิปสฺสเน ญาณํ.}{ภงฺคานุปสฺสนา\-ญาณ\-นิทฺเทโส สตฺตโม.}
\csromanmatikaanchor{mtk.10}
\csromantocmarkref{subsection}{8. อาทีนวญาณนิทฺเทส}{mtk.10}
\bookmark[dest={mtk.10},level=2]{8. อาทีนวญาณนิทฺเทส}
\csromanheader{2}{8. อาทีนว\-ญาณ\-นิทฺเทส}
\proseitem{53}{กถํ ภยตุปฏฺฐาเน ปญฺญา อาทีนเว ญาณํ\csromandash{}อุปฺปาโท ภยนฺติ ภยตุปฏฺฐาเน ปญฺญา อาทีนเว ญาณํ, ปวตฺตํ ภยนฺติ ภยตุปฏฺฐาเน ปญฺญา อาทีนเว ญาณํ, นิมิตฺตํ ภยนฺติ ฯเปฯ, อายูหนา ภยนฺติ ฯเปฯ, ปฏิสนฺธิ ภยนฺติ. คติ ภยนฺติ. นิพฺพตฺติ ภยนฺติ. อุปปตฺติ ภยนฺติ. ชาติ ภยนฺติ. ชรา ภยนฺติ. พฺยาธิ ภยนฺติ. มรณํ ภยนฺติ. โสโก ภยนฺติ. ปริเทโว ภยนฺติ. อุปายาโส ภยนฺติ ภยตุปฏฺฐาเน ปญฺญา อาทีนเว ญาณํ.}

\prose{\csromanfolio{57}อนุปฺปาโท เขมนฺติ สนฺติปเท ญาณํ, อปฺปวตฺตํ เขมนฺติ สนฺติปเท ญาณํ ฯเปฯ อนุปายาโส เขมนฺติ สนฺติปเท ญาณํ.}

\prose{อุปฺปาโท ภยํ, อนุปฺปาโท เขมนฺติ สนฺติปเท ญาณํ, ปวตฺตํ ภยํ, อปฺปวตฺตํ เขมนฺติ สนฺติปเท ญาณํ ฯเปฯ อุปายาโส ภยํ, อนุปายาโส เขมนฺติ สนฺติปเท ญาณํ.}

\prose{อุปฺปาโท ทุกฺขนฺติ ภยตุปฏฺฐาเน ปญฺญา อาทีนเว ญาณํ, ปวตฺตํ ทุกฺขนฺติ ภยตุปฏฺฐาเน ปญฺญา อาทีนเว ญาณํ ฯเปฯ อุปายาโส ทุกฺขนฺติ ภยตุปฏฺฐาเน ปญฺญา อาทีนเว ญาณํ.}

\prose{อนุปฺปาโท สุขนฺติ สนฺติปเท ญาณํ, อปฺปวตฺตํ สุขนฺติ สนฺติปเท ญาณํ ฯเปฯ อนุปายาโส สุขนฺติ สนฺติปเท ญาณํ.}

\prose{อุปฺปาโท ทุกฺขํ, อนุปฺปาโท สุขนฺติ สนฺติปเท ญาณํ, ปวตฺตํ ทุกฺขํ, อปฺปวตฺตํ สุขนฺติ สนฺติปเท ญาณํ ฯเปฯ อุปายาโส ทุกฺขํ, อนุปายาโส สุขนฺติ สนฺติปเท ญาณํ.}

\prose{อุปฺปาโท สามิสนฺติ ภยตุปฏฺฐาเน ปญฺญา อาทีนเว ญาณํ, ปวตฺตํ สามิสนฺติ ภยตุปฏฺฐาเน ปญฺญา อาทีนเว ญาณํ ฯเปฯ อุปายาโส สามิสนฺติ ภยตุปฏฺฐาเน ปญฺญา อาทีนเว ญาณํ.}

\prose{อนุปฺปาโท นิรามิสนฺติ สนฺติปเท ญาณํ, อปฺปวตฺตํ นิรามิสนฺติ สนฺติปเท ญาณํ ฯเปฯ อนุปายาโส นิรามิสนฺติ สนฺติปเท ญาณํ.}

\prose{อุปฺปาโท สามิสํ, อนุปฺปาโท นิรามิสนฺติ สนฺติปเท ญาณํ, ปวตฺตํ สามิสํ, อปฺปวตฺตํ นิรามิสนฺติ สนฺติปเท ญาณํ ฯเปฯ อุปายาโส สามิสํ, อนุปายาโส นิรามิสนฺติ สนฺติปเท ญาณํ.}

\prose{อุปฺปาโท สงฺขาราติ ภยตุปฏฺฐาเน ปญฺญา อาทีนเว ญาณํ, ปวตฺตํ สงฺขาราติ ภยตุปฏฺฐาเน ปญฺญา อาทีนเว ญาณํ ฯเปฯ อุปายาโส สงฺขาราติ ภยตุปฏฺฐาเน ปญฺญา อาทีนเว ญาณํ.}

\prose{อนุปฺปาโท นิพฺพานนฺติ สนฺติปเท ญาณํ, อปฺปวตฺตํ นิพฺพานนฺติ สนฺติปเท ญาณํ ฯเปฯ อนุปายาโส นิพฺพานนฺติ สนฺติปเท ญาณํ.}

\prose{\csromanfolio{58}อุปฺปาโท สงฺขารา, อนุปฺปาโท นิพฺพานนฺติ สนฺติปเท ญาณํ, ปวตฺตํ สงฺขารา, อปฺปวตฺตํ นิพฺพานนฺติ สนฺติปเท ญาณํ ฯเปฯ อุปายาโส สงฺขารา, อนุปายาโส นิพฺพานนฺติ สนฺติปเท ญาณํ.}

\setlength{\csromangathapagewidth}{0pt}
\setlength{\csromangathawakwidth}{0pt}
\csromangathameasuregroup{อุปฺปาทญฺจ ปวตฺตญฺจ, \\ อายูหนํ ปฏิสนฺธิํ, \\ อนุปฺปาทํ อปฺปวตฺตํ, \\ อนายูหนํ อปฺปฏิสนฺธิํ, \\ อิทํ อาทีนเว ญาณํ, \\ ปญฺจฐาเน สนฺติปเท, \\ ทฺวินฺนํ ญาณานํ กุสลตา,}{\csromangathabat{อุปฺปาทญฺจ ปวตฺตญฺจ,}{นิมิตฺตํ ทุกฺขนฺติ ปสฺสติ.} \\ \csromangathabat{อายูหนํ ปฏิสนฺธิํ,}{ญาณํ อาทีนเว อิทํ.} \\ \csromangathabat{อนุปฺปาทํ อปฺปวตฺตํ,}{อนิมิตฺตํ สุขนฺติ จ.} \\ \csromangathabat{อนายูหนํ อปฺปฏิสนฺธิํ,}{ญาณํ สนฺติปเท อิทํ.} \\ \csromangathabat{อิทํ อาทีนเว ญาณํ,}{ปญฺจฐาเนสุ ชายติ.} \\ \csromangathabat{ปญฺจฐาเน สนฺติปเท,}{ทส ญาเณ ปชานาติ.} \\ \csromangathabat{ทฺวินฺนํ ญาณานํ กุสลตา,}{นานาทิฏฺฐีสุ น กมฺปตีติ.}}
\csromangathasetleft{อุปฺปาทญฺจ ปวตฺตญฺจ, \\ อายูหนํ ปฏิสนฺธิํ, \\ อนุปฺปาทํ อปฺปวตฺตํ, \\ อนายูหนํ อปฺปฏิสนฺธิํ, \\ อิทํ อาทีนเว ญาณํ, \\ ปญฺจฐาเน สนฺติปเท, \\ ทฺวินฺนํ ญาณานํ กุสลตา,}
\csromangathagroup{\csromangathastanza{\csromangathabat{อุปฺปาทญฺจ ปวตฺตญฺจ,}{นิมิตฺตํ ทุกฺขนฺติ ปสฺสติ.} \\ \csromangathabat{อายูหนํ ปฏิสนฺธิํ,}{ญาณํ อาทีนเว อิทํ.}}\csromangathastanza{\csromangathabat{อนุปฺปาทํ อปฺปวตฺตํ,}{อนิมิตฺตํ สุขนฺติ จ.} \\ \csromangathabat{อนายูหนํ อปฺปฏิสนฺธิํ,}{ญาณํ สนฺติปเท อิทํ.}}\csromangathastanza{\csromangathabat{อิทํ อาทีนเว ญาณํ,}{ปญฺจฐาเนสุ ชายติ.} \\ \csromangathabat{ปญฺจฐาเน สนฺติปเท,}{ทส ญาเณ ปชานาติ.} \\ \csromangathabat{ทฺวินฺนํ ญาณานํ กุสลตา,}{นานาทิฏฺฐีสุ น กมฺปตีติ.}}}

\prosewithcloserruled{ตํญาตฏฺเฐน ญาณํ, ปชานนฏฺเฐน ปญฺญา, เตน วุจฺจติ ภยตุปฏฺฐาเน ปญฺญา อาทีนเว ญาณํ.}{อาทีนว\-ญาณ\-นิทฺเทโส อฏฺฐโม.}
\csromanmatikaanchor{mtk.11}
\csromantocmarkref{subsection}{9. สงฺขารุเปกฺขาญาณนิทฺเทส}{mtk.11}
\bookmark[dest={mtk.11},level=2]{9. สงฺขารุเปกฺขาญาณนิทฺเทส}
\csromanheader{2}{9. สงฺขารุเปกฺขาญาณ\-นิทฺเทส}
\proseitem{54}{กถํ มุญฺจิตุกมฺยตา\-ปฏิสงฺขา\-สนฺติฏฺฐนา\csromansharedfootnote{d82d92723d7fa679}{มุจฺจิตุกมฺยตาปฏิสงฺขาสนฺติฏฺฐนา (ก)} ปญฺญา สงฺขารุ\-เปกฺขาสุ ญาณํ\csromandash{}อุปฺปาทํ มุญฺจิตุกมฺยตา\-ปฏิสงฺขา\-สนฺติฏฺฐนา\csromansharedfootnote{5b0fae0f9ad42358}{อุปฺปาทมุญฺจิตุกมฺยตา... (สฺยา)} ปญฺญา สงฺขารุ\-เปกฺขาสุ ญาณํ, ปวตฺตํ มุญฺจิตุกมฺยตา\-ปฏิสงฺขา\-สนฺติฏฺฐนา ปญฺญา สงฺขารุ\-เปกฺขาสุ ญาณํ, นิมิตฺตํ มุญฺจตุกมฺยตา ฯเปฯ อายูหนํ มุญฺจิตุกมฺยตา. ปฏิสนฺธิํ มุญฺจิตุกมฺยตา. คติํ มุญฺจิตุกมฺยตา. นิพฺพตฺตํ มุญฺจิตุกมฺยตา. อุปปตฺติํ มุญฺจิตุกมฺยตา. ชาติํ มุญฺจิตุกมฺยตา. ชรํ มุญฺจิตุกมฺยตา. พฺยาธิํ มุญฺจิตุกมฺยตา. มรณํ มุญฺจิตุกมฺยตา. โสกํ มุญฺจิตุกมฺยตา. ปริเทวํ มุญฺจิตุกมฺยตา ฯเปฯ อุปายาสํ มุญฺจิตุกมฺยตา\-ปฏิสงฺขา\-สนฺติฏฺฐนา ปญฺญา สงฺขารุ\-เปกฺขาสุ ญาณํ.}

\prose{อุปฺปาโท ทุกฺขนฺติ มุญฺจิตุกมฺยตา\-ปฏิสงฺขา\-สนฺติฏฺฐนา ปญฺญา สงฺขารุ\-เปกฺขาสุ ญาณํ, ปวตฺตํ ทุกฺขนฺติ มุญฺจิตุกมฺยตา\-ปฏิสงฺขา\-สนฺติฏฺฐนา ปญฺญา สงฺขารุ\-เปกฺขาสุ ญาณํ ฯเปฯ อุปายาโส ทุกฺขนฺติ มุญฺจิตุกมฺยตา\-ปฏิสงฺขา\-สนฺติฏฺฐนา ปญฺญา สงฺขารุ\-เปกฺขาสุ ญาณํ.}

\prose{\csromanfolio{59}อุปฺปาโท ภยนฺติ มุญฺจิตุกมฺยตา\-ปฏิสงฺขา\-สนฺติฏฺฐนา ปญฺญา สงฺขารุ\-เปกฺขาสุ ญาณํ, ปวตฺตํ ภยนฺติ มุญฺจิตุกมฺยตา\-ปฏิสงฺขา\-สนฺติฏฺฐนา ปญฺญา สงฺขารุ\-เปกฺขาสุ ญาณํ ฯเปฯ อุปายาโส ภยนฺติ มุญฺจิตุกมฺยตา\-ปฏิสงฺขา\-สนฺติฏฺฐนา ปญฺญา สงฺขารุ\-เปกฺขาสุ ญาณํ.}

\prose{อุปฺปาโท สามิสนฺติ มุญฺจิตุกมฺยตา\-ปฏิสงฺขา\-สนฺติฏฺฐนา ปญฺญา สงฺขารุ\-เปกฺขาสุ ญาณํ, ปวตฺตํ สามิสนฺติ มุญฺจิตุกมฺยตา\-ปฏิสงฺขา\-สนฺติฏฺฐนา ปญฺญา สงฺขารุ\-เปกฺขาสุ ญาณํ ฯเปฯ อุปายาโส สามิสนฺติ มุญฺจิตุกมฺยตา\-ปฏิสงฺขา\-สนฺติฏฺฐนา ปญฺญา สงฺขารุ\-เปกฺขาสุ ญาณํ.}

\prose{อุปฺปาโท สงฺขาราติ มุญฺจิตุกมฺยตา\-ปฏิสงฺขา\-สนฺติฏฺฐนา ปญฺญา สงฺขารุ\-เปกฺขาสุ ญาณํ, ปวตฺตํ สงฺขาราติ มุญฺจิตุกมฺยตา\-ปฏิสงฺขา\-สนฺติฏฺฐนา ปญฺญา สงฺขารุ\-เปกฺขาสุ ญาณํ ฯเปฯ อุปายาโส สงฺขาราติ มุญฺจิตุกมฺยตา\-ปฏิสงฺขา\-สนฺติฏฺฐนา ปญฺญา สงฺขารุ\-เปกฺขาสุ ญาณํ.}

\prose{อุปฺปาโท สงฺขารา, เต สงฺขาเร อชฺฌุเปกฺขตีติ สงฺขารุเปกฺขา. เย จ สงฺขารา ยา จ อุเปกฺขา, อุโภเปเต สงฺขารา, เต สงฺขาเร อชฺฌุเปกฺขตีติ สงฺขารุเปกฺขา. ปวตฺตํ สงฺขารา ฯเปฯ. นิมิตฺตํ สงฺขารา. อายูหนา สงฺขารา. ปฏิสนฺธิ สงฺขารา. คติ สงฺขารา. นิพฺพตฺติ สงฺขารา. อุปปตฺติ สงฺขารา. ชาติ สงฺขารา. ชรา สงฺขารา. พฺยาธิ สงฺขารา. มรณํ สงฺขารา. โสโก สงฺขารา. ปริเทโว สงฺขารา ฯเปฯ. อุปายาโส สงฺขารา, เต สงฺขาเร อชฺฌุเปกฺขตีติ สงฺขารุเปกฺขา. เย จ สงฺขารา ยา จ อุเปกฺขา, อุโภเปเต สงฺขารา, เต สงฺขาเร อชฺฌุเปกฺขตีติ สงฺขารุเปกฺขา.}

\proseitem{55}{กติหากาเรหิ สงฺขารุ\-เปกฺขาย จิตฺตสฺส อภินีหาโร โหติ? อฏฺฐหากาเรหิ สงฺขารุ\-เปกฺขาย จิตฺตสฺส อภินีหาโร โหติ.  ปุถุชฺชนสฺส กติหากาเรหิ สงฺขารุ\-เปกฺขาย จิตฺตสฺส อภินีหาโร โหติ? เสกฺขสฺส กติหากาเรหิ สงฺขารุ\-เปกฺขาย จิตฺตสฺส อภินีหาโร โหติ? วีตราคสฺส กติหากาเรหิ สงฺขารุ\-เปกฺขาย จิตฺตสฺส อภินีหาโร โหติ? ปุถุชฺชนสฺส ทฺวีหากาเรหิ สงฺขารุ\-เปกฺขาย จิตฺตสฺส อภินีหาโร โหติ, เสกฺขสฺส ตีหากาเรหิ สงฺขารุ\-เปกฺขาย จิตฺตสฺส อภินีหาโร โหติ, วีตราคสฺส ตีหากาเรหิ สงฺขารุ\-เปกฺขาย จิตฺตสฺส อภินีหาโร โหติ.}

\prose{\csromanfolio{60}ปุถุชฺชนสฺส กตเมหิ ทฺวีหากาเรหิ สงฺขารุ\-เปกฺขาย จิตฺตสฺส อภินีหาโร โหติ? ปุถุชฺชโน สงฺขารุ\-เปกฺขํ อภินนฺทติ วา วิปสฺสติ วา. ปุถุชฺชนสฺส อิเมหิ ทฺวีหากาเรหิ สงฺขารุ\-เปกฺขาย จิตฺตสฺส อภินีหาโร โหติ.  เสกฺขสฺส กตเมหิ ตีหากาเรหิ สงฺขารุ\-เปกฺขาย จิตฺตสฺส อภินีหาโร โหติ? เสกฺโข สงฺขารุ\-เปกฺขํ อภินนฺทติ วา วิปสฺสติ วา ปฏิสงฺขาย วา ผลสมาปตฺติํ สมาปชฺชติ เสกฺขสฺส อิเมหิ ตีหากาเรหิ สงฺขารุ\-เปกฺขาย จิตฺตสฺส อภินีหาโร โหติ.  วีตราคสฺส กตเมหิ ตีหากาเรหิ สงฺขารุ\-เปกฺขาย จิตฺตสฺส อภินีหาโร โหติ? วีตราโค สงฺขารุ\-เปกฺขํ วิปสฺสติ วา ปฏิสงฺขาย วา ผลสมาปตฺติํ สมาปชฺชติ, ตทชฺฌุเปกฺขิตฺวา สุญฺญต\-วิหาเรน วา อนิมิตฺต\-วิหาเรน วา อปฺปณิหิตหาเรน วา วิหรติ. วีตราคสฺส อิเมหิ ตีหากาเรหิ สงฺขารุ\-เปกฺขาย จิตฺตสฺส อภินีหาโร โหติ.}

\proseitem{56}{กถํ ปุถุชฺชนสฺส จ เสกฺขสฺส จ สงฺขารุ\-เปกฺขาย จิตฺตสฺส อภินีหาโร เอกตฺตํ โหติ? ปุถุชฺชนสฺส สงฺขารุ\-เปกฺขํ อภินนฺทโต จิตฺตํ กิลิสฺสติ, ภาวนาย ปริปนฺโถ โหติ, ปฏิเวธสฺส อนฺตราโย โหติ, อายติํ ปฏิสนฺธิยา ปจฺจโย โหติ. เสกฺขสฺสปิ สงฺขารุ\-เปกฺขํ อภินนฺทโต จิตฺตํ กิลิสฺสติ, ภาวนาย ปริปนฺโถ โหติ, อุตฺตริ\-ปฏิเวธสฺส อนฺตราโย โหติ, อายติํ ปฏิสนฺธิยา ปจฺจโย โหติ. เอวํ ปุถุชฺชนสฺส จ เสกฺขสฺส จ สงฺขารุ\-เปกฺขาย จิตฺตสฺส อภินีหาโร เอกตฺตํ โหติ อภินนฺท\-ฏฺเฐน.}

\prose{กถํ ปุถุชฺชนสฺส จ เสกฺขสฺส จ วีตราคสฺส จ สงฺขารุ\-เปกฺขาย จิตฺตสฺส อภินีหาโร เอกตฺตํ โหติ? ปุถุชฺชโน สงฺขารุ\-เปกฺขํ อนิจฺจโตปิ ทุกฺขโตปิ อนตฺตโตปิ วิปสฺสติ, เสกฺโขปิ สงฺขารุ\-เปกฺขํ อนิจฺจโตปิ ทุกฺขโตปิ อนตฺตโตปิ วิปสฺสติ, วีตราโคปิ สงฺขารุ\-เปกฺขํ อนิจฺจโตปิ ทุกฺขโตปิ อนตฺตโตปิ วิปสฺสติ. เอวํ ปุถุชฺชนสฺส จ เสกฺขสฺส จ วีตราคสฺส จ สงฺขารุ\-เปกฺขาย จิตฺตสฺส อภินีหาโร เอกตฺตํ โหติ อนุปสฺสนฏฺเฐน.}

\prose{กถํ ปุถุชฺชนสฺส จ เสกฺขสฺส จ วีตราคสฺส จ สงฺขารุ\-เปกฺขาย จิตฺตสฺส อภินีหาโร นานตฺตํ โหติ? ปุถุชฺชนสฺส สงฺขารุเปกฺขา กุสลา โหติ, เสกฺขสฺสปิ สงฺขารุเปกฺขา กุสลา โหติ, วีตราคสฺส \csromanfolio{61}สงฺขารุเปกฺขา อพฺยากตา โหติ. เอวํ ปุถุชฺชนสฺส จ เสกฺขสฺส จ วีตราคสฺส จ สงฺขารุ\-เปกฺขาย จิตฺตสฺส อภินีหาโร นานตฺตํ โหติ กุสลา\-พฺยากต\-ฏฺเฐน.}

\prose{กถํ ปุถุชฺชนสฺส จ เสกฺขสฺส จ วีตราคสฺส จ สงฺขารุ\-เปกฺขาย จิตฺตสฺส อภินีหาโร นานตฺตํ โหติ? ปุถุชฺชนสฺส สงฺขารุเปกฺขา กิญฺจิ กาเล สุวิทิตา โหติ, กิญฺจิ กาเล น สุวิทิตา โหติ. เสกฺขสฺสปิ สงฺขารุเปกฺขา กิญฺจิ กาเล สุวิทิตา โหติ, กิญฺจิ กาเล น สุวิทิตา โหติ. วีตราคสฺส สงฺขารุเปกฺขา อจฺจนฺตํ สุวิทิตา โหติ. เอวํ ปุถุชฺชนสฺส จ เสกฺขสฺส จ วีตราคสฺส จ สงฺขารุ\-เปกฺขาย จิตฺตสฺส อภินีหาโร นานตฺตํ โหติ วิทิตฏฺเฐน จ อวิทิตฏฺเฐน จ.}

\prose{กถํ ปุถุชฺชนสฺส จ เสกฺขสฺส จ วีตราคสฺส จ สงฺขารุ\-เปกฺขาย จิตฺตสฺส อภินีหาโร นานตฺตํ โหติ? ปุถุชฺชโน สงฺขารุ\-เปกฺขํ อติตฺตตฺตา วิปสฺสติ, เสกฺโขปิ สงฺขารุ\-เปกฺขํ อติตฺตตฺตา วิปสฺสติ, วีตราโค สงฺขารุ\-เปกฺขํ ติตฺตตฺตา วิปสฺสติ. เอวํ ปุถุชฺชนสฺส จ เสกฺขสฺส จ วีตราคสฺส จ สงฺขารุ\-เปกฺขาย จิตฺตสฺส อภินีหาโร นานตฺตํ โหติ ติตฺตฏฺเฐน จ อติตฺตฏฺเฐน จ.}

\prose{กถํ ปุถุชฺชนสฺส จ เสกฺขสฺส จ วีตราคสฺส จ สงฺขารุ\-เปกฺขาย จิตฺตสฺส อภินีหาโร นานตฺตํ โหติ? ปุถุชฺชโน สงฺขารุ\-เปกฺขํ ติณฺณํ สํโยชนานํ ปหานาย โสตาปตฺติ\-มคฺคํ ปฏิลาภ\-ตฺถาย วิปสฺสติ, เสกฺโข สงฺขารุ\-เปกฺขํ ติณฺณํ สํโยชนานํ ปหีนตฺตา อุตฺตริ\-ปฏิลาภ\-ตฺถาย วิปสฺสติ, วีตราโค สงฺขารุ\-เปกฺขํ สพฺพกิเลสานํ ปหีนตฺตา ทิฏฺฐธมฺมสุขวิหาร\-ตฺถาย วิปสฺสติ. เอวํ ปุถุชฺชนสฺส จ เสกฺขสฺส จ วีตราคสฺส จ สงฺขารุ\-เปกฺขาย จิตฺตสฺส อภินีหาโร นานตฺตํ โหติ ปหีนฏฺเฐน จ อปฺปหีนฏฺเฐน จ.}

\prose{กถํ เสกฺขสฺส จ วีตราคสฺส จ สงฺขารุ\-เปกฺขาย จิตฺตสฺส อภินีหาโร นานตฺตํ โหติ? เสกฺโข สงฺขารุ\-เปกฺขํ อภินนฺทติ วา วิปสฺสติ วา ปฏิสงฺขาย วา ผลสมาปตฺติํ สมาปชฺชติ, วีตราโค สงฺขารุ\-เปกฺขํ วิปสฺสติ วา ปฏิสงฺขาย วา ผลสมาปตฺติํ สมาปชฺชติ, ตทชฺฌุเปกฺขิตฺวา สุญฺญต\-วิหาเรน วา อนิมิตฺต\-วิหาเรน วา อปฺปณิหิต\-วิหาเรน วา วิหรติ. เอวํ เสกฺขสฺส จ วีตราคสฺส จ สงฺขารุ\-เปกฺขาย จิตฺตสฺส อภินีหาโร นานตฺตํ โหติ วิหารสมาปตฺต\-ฏฺเฐน.}

\proseitem{57}{\csromanfolio{62}กติ สงฺขารุเปกฺขา สมถวเสน อุปฺปชฺชนฺติ, กติ สงฺขารุเปกฺขา วิปสฺสนา\-วเสน อุปฺปชฺชนฺติ? อฏฺฐ สงฺขารุเปกฺขา สมถวเสน อุปฺปชฺชนฺติ, ทส สงฺขารุเปกฺขา วิปสฺสนา\-วเสน อุปฺปชฺชนฺติ.}

\prose{กตมา อฏฺฐ สงฺขารุเปกฺขา สมถวเสน อุปฺปชฺชนฺติ? ปฐมํ ฌานํ ปฏิลาภ\-ตฺถาย นีวรเณ ปฏิสงฺขา สนฺติฏฺฐนา ปญฺญา สงฺขารุ\-เปกฺขาสุ ญาณํ. ทุติยํ ฌานํ ปฏิลาภ\-ตฺถาย วิตกฺกวิจาเร ปฏิสงฺขา สนฺติฏฺฐนา ปญฺญา สงฺขารุ\-เปกฺขาสุ ญาณํ. ตติยํ ฌานํ ปฏิลาภ\-ตฺถาย ปีติํ ปฏิสงฺขา สนฺติฏฺฐนา ปญฺญา สงฺขารุ\-เปกฺขาสุ ญาณํ. จตุตฺถํ ฌานํ ปฏิลาภ\-ตฺถาย สุขทุกฺเข ปฏิสงฺขา สนฺติฏฺฐนา ปญฺญา สงฺขารุ\-เปกฺขาสุ ญาณํ.  อากาสา\-นญฺจา\-ยตนสมาปตฺติํ ปฏิลาภ\-ตฺถาย รูปสญฺญํ ปฏิฆสญฺญํ นานตฺตสญฺญํ ปฏิสงฺขา สนฺติฏฺฐนา ปญฺญา สงฺขารุ\-เปกฺขาสุ ญาณํ. วิญฺญาณญฺจา\-ยตนสมาปตฺติํ ปฏิลาภ\-ตฺถาย อากาสา\-นญฺจา\-ยตนสญฺญํ ปฏิสงฺขา สนฺติฏฺฐนา ปญฺญา สงฺขารุ\-เปกฺขาสุ ญาณํ. อากิญฺจญฺญายตน\-สมาปตฺติํ ปฏิลาภ\-ตฺถาย วิญฺญาณญฺจา\-ยตนสญฺญํ ปฏิสงฺขา สนฺติฏฺฐนา ปญฺญา สงฺขารุ\-เปกฺขาสุ ญาณํ. เนว\-สญฺญา\-นา\-สญฺญา\-ยตนสมาปตฺติํ ปฏิลาภ\-ตฺถาย อากิญฺจญฺญายตน\-สญฺญํ ปฏิสงฺขา สนฺติฏฺฐนา ปญฺญา สงฺขารุ\-เปกฺขาสุ ญาณํ. อิมา อฏฺฐ สงฺขารุเปกฺขา สมถวเสน อุปฺปชฺชนฺติ.}

\prose{กตมา ทส สงฺขารุเปกฺขา วิปสฺสนา\-วเสน อุปฺปชฺชนฺติ? โสตาปตฺติ\-มคฺคํ ปฏิลาภ\-ตฺถาย อุปฺปาทํ ปวตฺตํ นิมิตฺตํ อายูหนํ ปฏิสนฺธิํ คติํ นิพฺพตฺติํ อุปปตฺติํ ชาติํ ชรํ พฺยาธิํ มรณํ โสกํ ปริเทวํ อุปายาสํ ปฏิสงฺขา สนฺติฏฺฐนา ปญฺญา สงฺขารุ\-เปกฺขาสุ ญาณํ. โสตาปตฺติผล\-สมา\-ปตฺตตฺถาย อุปฺปาทํ ปวตฺตํ นิมิตฺตํ อายูหนํ ปฏิสนฺธิํ ฯเปฯ ปฏิสงฺขา สนฺติฏฺฐนา ปญฺญา สงฺขารุ\-เปกฺขาสุ ญาณํ. สกทาคามิ\-มคฺคํ ปฏิลาภ\-ตฺถาย ฯเปฯ. สกทาคามิผล\-สมา\-ปตฺตตฺถาย. อนาคามิมคฺคํ ปฏิลาภ\-ตฺถาย ฯเปฯ. อนาคามิผล\-สมา\-ปตฺตตฺถาย. อรหตฺตมคฺคํ ปฏิลาภ\-ตฺถาย อุปฺปาทํ ปวตฺตํ นิมิตฺตํ อายูหนํ ปฏิสนฺธิํ คติํ นิพฺพตฺติํ อุปปตฺติํ ชาติํ ชรํ พฺยาธิํ มรณํ โสกํ ปริเทวํ อุปายาสํ ปฏิสงฺขา สนฺติฏฺฐนา ปญฺญา สงฺขารุ\-เปกฺขาสุ ญาณํ. อรหตฺตผลสมาปตฺต\-ตฺถาย ฯเปฯ. สุญฺญต\-วิหารสมาปตฺต\-ตฺถาย ฯเปฯ. อนิมิตฺต\-วิหารสมาปตฺต\-ตฺถาย อุปฺปาทํ ปวตฺตํ นิมิตฺตํ อายูหนํ ปฏิสนฺธิํ ฯเปฯ. ปฏิสงฺขา สนฺติฏฺฐนา ปญฺญา สงฺขารุ\-เปกฺขาสุ ญาณํ. อิมา ทส สงฺขารุเปกฺขา วิปสฺสนา\-วเสน อุปฺปชฺชนฺติ.}

\proseitem{58}{\csromanfolio{63}กติ สงฺขารุเปกฺขา กุสลา, กติ อกุสลา, กติ อพฺยากตา? ปนฺนรส สงฺขารุเปกฺขา กุสลา, ติสฺโส สงฺขารุเปกฺขา อพฺยากตา, นตฺถิ สงฺขารุเปกฺขา อกุสลา.}

\setlength{\csromangathapagewidth}{0pt}
\setlength{\csromangathawakwidth}{0pt}
\csromangathameasuregroup{ปฏิสงฺขา สนฺติฏฺฐนา ปญฺญา, \\ ปุถุชฺชนสฺส เทฺว โหนฺติ, \\ ตโย จ วีตราคสฺส, \\ อฏฺฐ สมาธิสฺส ปจฺจยา, \\ อฏฺฐารส สงฺขารุเปกฺขา, \\ อิเม อฏฺฐารสาการา, \\ กุสโล สงฺขารุเปกฺขาสุ,}{\csromangathabat{ปฏิสงฺขา สนฺติฏฺฐนา ปญฺญา,}{อฏฺฐ จิตฺตสฺส โคจรา.} \\ \csromangathabat{ปุถุชฺชนสฺส เทฺว โหนฺติ,}{ตโย เสกฺขสฺส โคจรา.} \\ \csromangathabat{ตโย จ วีตราคสฺส,}{เยหิ จิตฺตํ วิวฏฺฏติ.} \\ \csromangathabat{อฏฺฐ สมาธิสฺส ปจฺจยา,}{ทส ญาณสฺส โคจรา.} \\ \csromangathabat{อฏฺฐารส สงฺขารุเปกฺขา,}{ติณฺณํ วิโมกฺขาน ปจฺจยา.} \\ \csromangathabat{อิเม อฏฺฐารสาการา,}{ปญฺญา ยสฺส ปริจฺจิตา.} \\ \csromangathabat{กุสโล สงฺขารุเปกฺขาสุ,}{นานาทิฏฺฐีสุ น กมฺปตีติ.}}
\csromangathasetleft{ปฏิสงฺขา สนฺติฏฺฐนา ปญฺญา, \\ ปุถุชฺชนสฺส เทฺว โหนฺติ, \\ ตโย จ วีตราคสฺส, \\ อฏฺฐ สมาธิสฺส ปจฺจยา, \\ อฏฺฐารส สงฺขารุเปกฺขา, \\ อิเม อฏฺฐารสาการา, \\ กุสโล สงฺขารุเปกฺขาสุ,}
\csromangathagroup{\csromangathastanza{\csromangathabat{ปฏิสงฺขา สนฺติฏฺฐนา ปญฺญา,}{อฏฺฐ จิตฺตสฺส โคจรา.} \\ \csromangathabat{ปุถุชฺชนสฺส เทฺว โหนฺติ,}{ตโย เสกฺขสฺส โคจรา.}}\csromangathastanza{\csromangathabat{ตโย จ วีตราคสฺส,}{เยหิ จิตฺตํ วิวฏฺฏติ.} \\ \csromangathabat{อฏฺฐ สมาธิสฺส ปจฺจยา,}{ทส ญาณสฺส โคจรา.}}\csromangathastanza{\csromangathabat{อฏฺฐารส สงฺขารุเปกฺขา,}{ติณฺณํ วิโมกฺขาน ปจฺจยา.} \\ \csromangathabat{อิเม อฏฺฐารสาการา,}{ปญฺญา ยสฺส ปริจฺจิตา.} \\ \csromangathabat{กุสโล สงฺขารุ\-เปกฺขาสุ,}{นานาทิฏฺฐีสุ น กมฺปตีติ.}}}

\prosewithcloserruled{ตํญาตฏฺเฐน ญาณํ, ปชานนฏฺเฐน ปญฺญา, เตน วุจฺจติ มุญฺจิตุกมฺยตา\-ปฏิสงฺขา\-สนฺติฏฺฐนา ปญฺญา สงฺขารุ\-เปกฺขาสุ ญาณํ.}{สงฺขารุเปกฺขาญาณ\-นิทฺเทโส นวโม.}
\csromanmatikaanchor{mtk.12}
\csromantocmarkref{subsection}{10. โคตฺรภุญาณนิทฺเทส}{mtk.12}
\bookmark[dest={mtk.12},level=2]{10. โคตฺรภุญาณนิทฺเทส}
\csromanheader{2}{10. โคตฺรภุญาณ\-นิทฺเทส}
\proseitem{59}{กถํ พหิทฺธา วุฏฺฐาน\-วิวฏฺฏเน ปญฺญา โคตฺรภุญาณํ\csromandash{} อุปฺปาทํ อภิภุยฺยตีติ โคตฺรภุ, ปวตฺตํ อภิภุยฺยตีติ โคตฺรภุ, นิมิตฺตํ อภิภุยฺยตีติ โคตฺรภุ, อายูหนํ อภิภุยฺยตีติ โคตฺรภุ, ปฏิสนฺธิํ อภิภุยฺยตีติ โคตฺรภุ, คติํ อภิภุยฺยตีติ โคตฺรภุ, นิพฺพตฺติํ อภิภุยฺยตีติ โคตฺรภุ, อุปปตฺติํ อภิภุยฺยตีติ โคตฺรภุ, ชาติํ อภิภุยฺยตีติ โคตฺรภุ, ชรํ อภิภุยฺยตีติ โคตฺรภุ, พฺยาธิํ อภิภุยฺยตีติ โคตฺรภุ, มรณํ อภิภุยฺยตีติ โคตฺรภุ, โสกํ อภิภุยฺยตีติ โคตฺรภุ, ปริเทวํ อภิภุยฺยตีติ โคตฺรภุ, อุปายาสํ อภิภุยฺยตีติ โคตฺรภุ, พหิทฺธา สงฺขาร\-นิมิตฺตํ อภิภุยฺยตีติ โคตฺรภุ.  อนุปฺปาทํ ปกฺขนฺทตีติ โคตฺรภุ, อปฺปวตฺตํ ปกฺขนฺทตีติ โคตฺรภุ ฯเปฯ นิโรธํ นิพฺพานํ ปกฺขนฺทตีติ โคตฺรภุ.}

\prose{อุปฺปาทํ อภิภุยฺยิตฺวา อนุปฺปาทํ ปกฺขนฺทตีติ โคตฺรภุ, ปวตฺตํ อภิภุยฺยิตฺวา อปฺปวตฺตํ ปกฺขนฺทตีติ โคตฺรภุ, นิมิตฺตํ อภิภุยฺยิตฺวา อนิมิตฺตํ ปกฺขนฺทตีติ \csromanfolio{64}โคตฺรภุ ฯเปฯ พหิทฺธา สงฺขาร\-นิมิตฺตํ อภิภุยฺยิตฺวา นิโรธํ นิพฺพานํ ปกฺขนฺทตีติ โคตฺรภุ.}

\prose{อุปฺปาทา วุฏฺฐาตีติ โคตฺรภุ, ปวตฺตา วุฏฺฐาตีติ โคตฺรภุ, นิมิตฺตา วุฏฺฐาตีติ โคตฺรภุ, อายูหนา วุฏฺฐาตีติ โคตฺรภุ, ปฏิสนฺธิยา วุฏฺฐาตีติ โคตฺรภุ, คติยา วุฏฺฐาตีติ โคตฺรภุ, นิพฺพตฺติยา วุฏฺฐาตีติ โคตฺรภุ, อุปปตฺติยา วุฏฺฐาตีติ โคตฺรภุ, ชาติยา วุฏฺฐาตีติ โคตฺรภุ, ชราย วุฏฺฐตีติ โคตฺรภุ, พฺยาธิมฺหา วุฏฺฐาตีติ โคตฺรภุ, มรณา วุฏฺฐาตีติ โคตฺรภุ, โสกา วุฏฺฐาตีติ โคตฺรภุ, ปริเทวา วุฏฺฐาตีติ โคตฺรภุ, อุปายาสา วุฏฺฐาตีติ โคตฺรภุ, พหิทฺธา สงฺขาร\-นิมิตฺตา วุฏฺฐาตีติ โคตฺรภุ.  อนุปฺปาทํ ปกฺขนฺทตีติ โคตฺรภุ, อปฺปวตฺตํ ปกฺขนฺทตีติ โคตฺรภุ ฯเปฯ นิโรธํ นิพฺพานํ ปกฺขนฺทตีติ โคตฺรภุ.}

\prose{อุปฺปาทา วุฏฺฐหิตฺวา\csromansharedfootnote{4d3e1bf9262be44f}{วุฏฺฐิตฺวา (สฺยา, ก)} อนุปฺปาทํ ปกฺขนฺทตีติ โคตฺรภุ, ปวตฺตา วุฏฺฐหิตฺวา อปฺปวตฺตํ ปกฺขนฺทตีติ โคตฺรภุ, นิมิตฺตา วุฏฺฐหิตฺวา อนิมิตฺตํ ปกฺขนฺทตีติ โคตฺรภุ, อายูหนา วุฏฺฐหิตฺวา อนายูหนํ ปกฺขนฺทตีติ โคตฺรภุ, ปฏิสนฺธิยา วุฏฺฐหิตฺวา อปฺปฏิสนฺธิํ ปกฺขนฺทตีติ โคตฺรภุ, คติยา วุฏฺฐหิตฺวา อคติํ ปกฺขนฺทตีติ โคตฺรภุ, นิพฺพตฺติยา วุฏฺฐหิตฺวา อนิพฺพตฺติํ ปกฺขนฺทตีติ โคตฺรภุ, อุปปตฺติยา วุฏฺฐหิตฺวา อนุปปตฺติํ ปกฺขนฺทตีติ โคตฺรภุ, ชาติยา วุฏฺฐหิตฺวา อชาติํ ปกฺขนฺทตีติ โคตฺรภุ, ชราย วุฏฺฐหิตฺวา อชรํ ปกฺขนฺทตีติ โคตฺรภุ, พฺยาธิมฺหา วุฏฺฐหิตฺวา อพฺยาธิํ ปกฺขนฺทตีติ โคตฺรภุ, มรณา วุฏฺฐหิตฺวา อมตํ ปกฺขนฺทตีติ โคตฺรภุ, โสกา วุฏฺฐหิตฺวา อโสกํ ปกฺขนฺทตีติ โคตฺรภุ, ปริเทวา วุฏฺฐหิตฺวา อปริเทวํ ปกฺขนฺทตีติ โคตฺรภุ, อุปายาสา วุฏฺฐหิตฺวา อนุปายาสํ ปกฺขนฺทตีติ โคตฺรภุ, พหิทฺธา สงฺขาร\-นิมิตฺตา วุฏฺฐหิตฺวา นิโรธํ นิพฺพานํ ปกฺขนฺทตีติ โคตฺรภุ.}

\prose{อุปฺปาทา วิวฏฺฏตีติ โคตฺรภุ, ปวตฺตา วิวฏฺฏตีติ โคตฺรภุ ฯเปฯ พหิทฺธา สงฺขาร\-นิมิตฺตา วิวฏฺฏตีติ โคตฺรภุ.  อนุปฺปาทํ ปกฺขนฺทตีติ โคตฺรภุ, อปฺปวตฺตํ ปกฺขนฺทตีติ โคตฺรภุ ฯเปฯ นิโรธํ นิพฺพานํ ปกฺขนฺทตีติ โคตฺรภุ.}

\prose{อุปฺปาทา วิวฏฺฏิตฺวา อนุปฺปาทํ ปกฺขนฺทตีติ โคตฺรภุ, ปวตฺตา วิวฏฺฏิตฺวา อปฺปวตฺตํ ปกฺขนฺทตีติ โคตฺรภุ ฯเปฯ พหิทฺธา สงฺขาร\-นิมิตฺตา วิวฏฺฏิตฺวา นิโรธํ นิพฺพานํ ปกฺขนฺทตีติ โคตฺรภุ.}

\proseitem{60}{\csromanfolio{65}กติ โคตฺรภู ธมฺมา สมถวเสน อุปฺปชฺชนฺติ, กติ โคตฺรภู ธมฺมา วิปสฺสนา\-วเสน อุปฺปชฺชนฺติ? อฏฺฐ โคตฺรภู ธมฺมา สมถวเสน อุปฺปชฺชนฺติ, ทส โคตฺรภู ธมฺมา วิปสฺสนา\-วเสน อุปฺปชฺชนฺติ.}

\prose{กตเม อฏฺฐ โคตฺรภู ธมฺมา สมถวเสน อุปฺปชฺชนฺติ? ปฐมํ ฌานํ ปฏิลาภ\-ตฺถาย นีวรเณ อภิภุยฺยตีติ โคตฺรภุ, ทุติยํ ฌานํ ปฏิลาภ\-ตฺถาย วิตกฺกวิจาเร อภิภุยฺยตีติ โคตฺรภุ, ตติยํ ฌานํ ปฏิลาภ\-ตฺถาย ปีติํ อภิภุยฺยตีติ โคตฺรภุ, จตุตฺถํ ฌานํ ปฏิลาภ\-ตฺถาย สุขทุกฺเข อภิภุยฺยตีติ โคตฺรภุ, อากาสา\-นญฺจา\-ยตนสมาปตฺติํ ปฏิลาภ\-ตฺถาย รูปสญฺญํ ปฏิฆสญฺญํ นานตฺตสญฺญํ อภิภุยฺยตีติ โคตฺรภุ, วิญฺญาณญฺจา\-ยตนสมาปตฺติํ ปฏิลาภ\-ตฺถาย อากาสา\-นญฺจา\-ยตนสญฺญํ อภิภุยฺยตีติ โคตฺรภุ, อากิญฺจญฺญายตน\-สมาปตฺติํ ปฏิลาภ\-ตฺถาย วิญฺญาณญฺจา\-ยตนสญฺญํ อภิภุยฺยตีติ โคตฺรภุ, เนว\-สญฺญา\-นา\-สญฺญา\-ยตนสมาปตฺติํ ปฏิลาภ\-ตฺถาย อากิญฺจญฺญายตน\-สญฺญํ อภิภุยฺยตีติ โคตฺรภุ. อิเม อฏฺฐ โคตฺรภู ธมฺมา สมถวเสน อุปฺปชฺชนฺติ.}

\prose{กตเม ทส โคตฺรภู ธมฺมา วิปสฺสนา\-วเสน อุปฺปชฺชนฺติ? โสตาปตฺติ\-มคฺคํ ปฏิลาภ\-ตฺถาย อุปฺปาทํ ปวตฺตํ นิมิตฺตํ อายูหนํ ปฏิสนฺธิํ คติํ นิพฺพตฺติํ อุปปตฺติํ ชาติํ ชรํ พฺยาธิํ มรณํ โสกํ ปริเทวํ อุปายาสํ พหิทฺธา สงฺขาร\-นิมิตฺตํ อภิภุยฺยตีติ โคตฺรภุ, โสตาปตฺติผล\-สมา\-ปตฺตตฺถาย อุปฺปาทํ ปวตฺตํ นิมิตฺตํ อายูหนํ ปฏิสนฺธิํ ฯเปฯ อภิภุยฺยตีติ โคตฺรภุ. สกทาคามิ\-มคฺคํ ปฏิลาภ\-ตฺถาย ฯเปฯ สกทาคามิผล\-สมา\-ปตฺตตฺถาย. อนาคามิมคฺคํ ปฏิลาภ\-ตฺถาย. อนาคามิผล\-สมา\-ปตฺตตฺถาย. อรหตฺตมคฺคํ ปฏิลาภ\-ตฺถาย อุปฺปาทํ ปวตฺตํ นิมิตฺตํ อายูหนํ ปฏิสนฺธิํ คติํ นิพฺพตฺติํ อุปปตฺติํ ชาติํ ชรํ พฺยาธิํ มรณํ โสกํ ปริเทวํ อุปายาสํ พหิทฺธา สงฺขาร\-นิมิตฺตํ อภิภุยฺยตีติ โคตฺรภุ, อรหตฺตผลสมาปตฺต\-ตฺถาย. สุญฺญต\-วิหารสมาปตฺต\-ตฺถาย. อนิมิตฺต\-วิหารสมาปตฺต\-ตฺถาย อุปฺปาทํ ปวตฺตํ นิมิตฺตํ อายูหนํ ปฏิสนฺธิํ ฯเปฯ อภิภุยฺยตีติ โคตฺรภุ. อิเม ทส โคตฺรภู ธมฺมา วิปสฺสนา\-วเสน อุปฺปชฺชนฺติ.}

\prose{กติ โคตฺรภู ธมฺมา กุสลา, กติ อกุสลา, กติ อพฺยากตา? ปนฺนรส โคตฺรภู ธมฺมา กุสลา, ตโย โคตฺรภู ธมฺมา อพฺยากตา, นตฺถิ โคตฺรภู ธมฺมา อกุสลาติ.}

\setlength{\csromangathapagewidth}{0pt}
\setlength{\csromangathawakwidth}{0pt}
\csromangathameasuregroup{สามิสญฺจ นิรามิสํ, \\ สญฺญุตฺตญฺจ วิสญฺญุตฺตํ, \\ อฏฺฐ สมาธิสฺส ปจฺจยา, \\ อฏฺฐารส โคตฺรภู ธมฺมา, \\ อิเม อฏฺฐารสาการา, \\ กุสโล วิวฏฺเฏ วุฏฺฐาเน,}{\csromangathabat{สามิสญฺจ นิรามิสํ,}{ปณิหิตญฺจ อปฺปณิหิตํ.} \\ \csromangathabat{สญฺญุตฺตญฺจ วิสญฺญุตฺตํ,}{วุฏฺฐิตญฺจ อวุฏฺฐิตํ.} \\ \csromangathabat{อฏฺฐ สมาธิสฺส ปจฺจยา,}{ทส ญาณสฺส โคจรา.} \\ \csromangathabat{อฏฺฐารส โคตฺรภู ธมฺมา,}{ติณฺณํ วิโมกฺขาน ปจฺจยา.} \\ \csromangathabat{อิเม อฏฺฐารสาการา,}{ปญฺญา ยสฺส ปริจฺจิตา.} \\ \csromangathabat{กุสโล วิวฏฺเฏ วุฏฺฐาเน,}{นานาทิฏฺฐีสุ น กมฺปตีติ.}}
\csromangathasetleft{สามิสญฺจ นิรามิสํ, \\ สญฺญุตฺตญฺจ วิสญฺญุตฺตํ, \\ อฏฺฐ สมาธิสฺส ปจฺจยา, \\ อฏฺฐารส โคตฺรภู ธมฺมา, \\ อิเม อฏฺฐารสาการา, \\ กุสโล วิวฏฺเฏ วุฏฺฐาเน,}
\csromangathagroup{\csromangathastanza{\csromanfolio{66}\csromangathabat{สามิสญฺจ นิรามิสํ,}{ปณิหิตญฺจ อปฺปณิหิตํ.} \\ \csromangathabat{สญฺญุตฺตญฺจ วิสญฺญุตฺตํ,}{วุฏฺฐิตญฺจ อวุฏฺฐิตํ.}}\csromangathastanza{\csromangathabat{อฏฺฐ สมาธิสฺส ปจฺจยา,}{ทส ญาณสฺส โคจรา.} \\ \csromangathabat{อฏฺฐารส โคตฺรภู ธมฺมา,}{ติณฺณํ วิโมกฺขาน ปจฺจยา.}}\csromangathastanza{\csromangathabat{อิเม อฏฺฐารสาการา,}{ปญฺญา ยสฺส ปริจฺจิตา.} \\ \csromangathabat{กุสโล วิวฏฺเฏ วุฏฺฐาเน,}{นานาทิฏฺฐีสุ น กมฺปตีติ.}}}

\prosewithcloserruled{ตํญาตฏฺเฐน ญาณํ, ปชานนฏฺเฐน ปญฺญา, เตน วุจฺจติ พหิทฺธา วุฏฺฐาน\-วิวฏฺฏเน ปญฺญา โคตฺรภุญาณํ.}{โคตฺรภุญาณ\-นิทฺเทโส ทสโม.}
\csromanmatikaanchor{mtk.13}
\csromantocmarkref{subsection}{11. มคฺคญาณนิทฺเทส}{mtk.13}
\bookmark[dest={mtk.13},level=2]{11. มคฺคญาณนิทฺเทส}
\csromanheader{2}{11. มคฺคญาณ\-นิทฺเทส}
\proseitem{61}{กถํ ทุภโต วุฏฺฐาน\-วิวฏฺฏเน ปญฺญา มคฺเค ญาณํ\csromandash{} โสตาปตฺติมคฺค\-กฺขเณ ทสฺสนฏฺเฐน สมฺมาทิฏฺฐิ มิจฺฉาทิฏฺฐิยา วุฏฺฐาติ, ตทนุวตฺตก\-กิเลเสหิ จ ขนฺเธหิ จ วุฏฺฐาติ, พหิทฺธา จ สพฺพนิมิตฺเตหิ วุฏฺฐาติ. เตน วุจฺจติ ทุภโต วุฏฺฐาน\-วิวฏฺฏเน ปญฺญา มคฺเค ญาณํ.}

\prose{อภินิโรปน\-ฏฺเฐน สมฺมาสงฺกปฺโป มิจฺฉาสงฺกปฺปา วุฏฺฐาติ, ตทนุวตฺตา\-กกิเลเสหิ จ ขนฺเธหิ จ วุฏฺฐาติ, พหิทฺธา จ สพฺพนิมิตฺเตหิ วุฏฺฐาติ. เตน วุจฺจติ ทุภโต วุฏฺฐาน\-วิวฏฺฏเน ปญฺญา มคฺเค ญาณํ.}

\prose{ปริคฺคห\-ฏฺเฐน สมฺมาวาจา มิจฺฉาวาจาย วุฏฺฐาติ, ตทนุวตฺตก\-กิเลเสหิ จ ขนฺเธหิ จ วุฏฺฐาติ, พหิทฺธา จ สพฺพนิมิตฺเตหิ วุฏฺฐาติ. เตน วุจฺจติ ทุภโต วุฏฺฐาน\-วิวฏฺฏเน ปญฺญา มคฺเค ญาณํ.}

\prose{สมุฏฺฐาน\-ฏฺเฐน สมฺมากมฺมนฺโต มิจฺฉากมฺมนฺตา วุฏฺฐาติ, ตทนุวตฺตก\-กิเลเสหิ จ ขนฺเธหิ จ วุฏฺฐาติ, พหิทฺธา จ สพฺพนิมิตฺเตหิ วุฏฺฐาติ. เตน วุจฺจติ ทุภโต วุฏฺฐาน\-วิวฏฺฏเน ปญฺญา มคฺเค ญาณํ.}

\prose{โวทานฏฺเฐน สมฺมาอาชีโว มิจฺฉาอาชีวา วุฏฺฐาติ, ตทนุวตฺตก\-กิเลเสหิ จ ขนฺเธหิ จ วุฏฺฐาติ, พหิทฺธา จ สพฺพนิมิตฺเตหิ วุฏฺฐาติ. เตน วุจฺจติ ทุภโต วุฏฺฐาน\-วิวฏฺฏเน ปญฺญา มคฺเค ญาณํ.}

\prose{\csromanfolio{67}ปคฺคหฏฺเฐน สมฺมาวายาโม มิจฺฉาวายามา วุฏฺฐาติ, ตทนุวตฺตก\-กิเลเสหิ จ ขนฺเธหิ จ วุฏฺฐาติ, พหิทฺธา จ สพฺพนิมิตฺเตหิ วุฏฺฐาติ. เตน วุจฺจติ ทุภโต วุฏฺฐาน\-วิวฏฺฏเน ปญฺญา มคฺเค ญาณํ.}

\prose{อุปฏฺฐาน\-ฏฺเฐน สมฺมาสติ มิจฺฉาสติยา วุฏฺฐาติ, ตทนุวตฺตก\-กิเลเสหิ จ ขนฺเธหิ จ วุฏฺฐาติ, พหิทฺธา จ สพฺพนิมิตฺเตหิ วุฏฺฐาติ. เตน วุจฺจติ ทุภโต วุฏฺฐาน\-วิวฏฺฏเน ปญฺญา มคฺเค ญาณํ.}

\prose{อวิกฺเขป\-ฏฺเฐน สมฺมาสมาธิ มิจฺฉา\-สมาธิโต วุฏฺฐาติ, ตทนุวตฺตก\-กิเลเสหิ จ ขนฺเธหิ จ วุฏฺฐาติ, พหิทฺธา จ สพฺพนิมิตฺเตหิ วุฏฺฐาติ. เตน วุจฺจติ ทุภโต วุฏฺฐาน\-วิวฏฺฏเน ปญฺญา มคฺเค ญาณํ.}

\prose{สกทาคามิ\-มคฺคกฺขเณ ทสฺสนฏฺเฐน สมฺมาทิฏฺฐิ ฯเปฯ. อวิกฺเขป\-ฏฺเฐน สมฺมาสมาธิ โอฬาริกา กามราค\-สญฺโญชนา ปฏิฆ\-สญฺโญชนา โอฬาริกา กามราคานุสยา ปฏิฆานุสยา วุฏฺฐาติ, ตทนุวตฺตก\-กิเลเสหิ จ ขนฺเธหิ จ วุฏฺฐาติ, พหิทฺธา จ สพฺพนิมิตฺเตหิ วุฏฺฐาติ. เตน วุจฺจติ ทุภโต วุฏฺฐาน\-วิวฏฺฏเน ปญฺญา มคฺเค ญาณํ.}

\prose{อนาคามิ\-มคฺคกฺขเณ ทสฺสนฏฺเฐน สมฺมาทิฏฺฐิ ฯเปฯ. อวิกฺเขป\-ฏฺเฐน สมฺมาสมาธิ อณุสหคตา กามราค\-สญฺโญชนา ปฏิฆ\-สญฺโญชนา อณุสหคตา กามราคานุสยา ปฏิฆานุสยา วุฏฺฐาติ, ตทนุวตฺตก\-กิเลเสหิ จ ขนฺเธหิ จ วุฏฺฐาติ, พหิทฺธา จ สพฺพนิมิตฺเตหิ วุฏฺฐาติ. เตน วุจฺจติ ทุภโต วุฏฺฐาน\-วิวฏฺฏเน ปญฺญา มคฺเค ญาณํ.}

\prose{อรหตฺต\-มคฺคกฺขเณ ทสฺสนฏฺเฐน สมฺมาทิฏฺฐิ ฯเปฯ. อวิกฺเขป\-ฏฺเฐน สมฺมาสมาธิ รูปราคา อรูปราคา มานา อุทฺธจฺจา อวิชฺชาย มานานุสยา ภวราคา\-นุสยา อวิชฺชานุสยา วุฏฺฐาติ, ตทนุวตฺตก\-กิเลเสหิ จ ขนฺเธหิ จ วุฏฺฐาติ, พหิทฺธา จ สพฺพนิมิตฺเตหิ วุฏฺฐาติ. เตน วุจฺจติ ทุภโต วุฏฺฐาน\-วิวฏฺฏเน ปญฺญา มคฺเค ญาณํ.}

\setlength{\csromangathapagewidth}{0pt}
\setlength{\csromangathawakwidth}{0pt}
\csromangathameasuregroup{สมาทหิตฺวา ยถา เจ วิปสฺสติ, \\ วิปสฺสนา จ สมโถ ตทา อหุ, \\ ทุกฺขา สงฺขารา สุโข, \\ ทุภโต วุฏฺฐิตา ปญฺญา, \\ วิโมกฺขจริยํ ชานาติ, \\ ทฺวินฺนํ ญาณานํ กุสลตา,}{อชาตํ ฌาเปติ ชาเตน, ฌานํ เตน ปวุจฺจติ. \\ ฌานวิโมกฺเข กุสลตา, นานาทิฏฺฐีสุ น กมฺปติ. \\ \csromangathabat{สมาทหิตฺวา ยถา เจ วิปสฺสติ,}{วิปสฺสมาโน ตถา เจ สมาทเห.} \\ \csromangathabat{วิปสฺสนา จ สมโถ ตทา อหุ,}{สมานภาคา ยุคนทฺธา วตฺตเร.} \\ \csromangathabat{ทุกฺขา สงฺขารา สุโข,}{นิโรโธ อิติ ทสฺสนํ\textsuperscript{1}.} \\ \csromangathabat{ทุภโต วุฏฺฐิตา ปญฺญา,}{ผสฺเสติ อมตํ ปทํ.} \\ \csromangathabat{วิโมกฺขจริยํ ชานาติ,}{นานตฺเตกตฺตโกวิโท.} \\ \csromangathabat{ทฺวินฺนํ ญาณานํ กุสลตา,}{นานาทิฏฺฐีสุ น กมฺปตีติ.}}
\csromangathasetleft{สมาทหิตฺวา ยถา เจ วิปสฺสติ, \\ วิปสฺสนา จ สมโถ ตทา อหุ, \\ ทุกฺขา สงฺขารา สุโข, \\ ทุภโต วุฏฺฐิตา ปญฺญา, \\ วิโมกฺขจริยํ ชานาติ, \\ ทฺวินฺนํ ญาณานํ กุสลตา,}
\csromangathagroup{\csromangathastanza{\csromangathaitemline{62}{อชาตํ ฌาเปติ ชาเตน, ฌานํ เตน ปวุจฺจติ.} \\ \csromangathacontline{ฌานวิโมกฺเข กุสลตา, นานาทิฏฺฐีสุ น กมฺปติ.} \\ \csromangathacontbat{สมาทหิตฺวา ยถา เจ วิปสฺสติ,}{วิปสฺสมาโน ตถา เจ สมาทเห.} \\ \csromangathacontbat{วิปสฺสนา จ สมโถ ตทา อหุ,}{สมานภาคา ยุคนทฺธา วตฺตเร.}}\csromangathastanza{\csromanfolio{68}\csromangathabat{ทุกฺขา สงฺขารา สุโข,}{นิโรโธ อิติ ทสฺสนํ\csromansharedfootnote{05701f1cec977f79}{นิโรโธติ ทสฺสนํ(สฺยา, ก)}.} \\ \csromangathabat{ทุภโต วุฏฺฐิตา ปญฺญา,}{ผสฺเสติ อมตํ ปทํ.}}\csromangathastanza{\csromangathabat{วิโมกฺขจริยํ ชานาติ,}{นานตฺเตกตฺตโกวิโท.} \\ \csromangathabat{ทฺวินฺนํ ญาณานํ กุสลตา,}{นานาทิฏฺฐีสุ น กมฺปตีติ.}}}

\prosewithcloserruled{ตํญาตฏฺเฐน ญาณํ, ปชานนฏฺเฐน ปญฺญา. เตน วุจฺจติ ทุภโต วุฏฺฐาน\-วิวฏฺฏเน ปญฺญา มคฺเค ญาณํ.}{มคฺคญาณ\-นิทฺเทโส เอกาทสโม.}
\csromanmatikaanchor{mtk.14}
\csromantocmarkref{subsection}{12. ผลญาณนิทฺเทส}{mtk.14}
\bookmark[dest={mtk.14},level=2]{12. ผลญาณนิทฺเทส}
\csromanheader{2}{12. ผลญาณ\-นิทฺเทส}
\proseitem{63}{กถํ ปโยค\-ปฺปฏิปฺปสฺสทฺธิ\-ปญฺญา ผเล ญาณํ\csromandash{} โสตาปตฺติมคฺค\-กฺขเณ ทสฺสนฏฺเฐน สมฺมาทิฏฺฐิ มิจฺฉาทิฏฺฐิยา วุฏฺฐาติ, ตทนุวตฺตก\-กิเลเสหิ จ ขนฺเธหิ จ วุฏฺฐาติ, พหิทฺธา จ สพฺพนิมิตฺเตหิ วุฏฺฐาติ. ตํปโยค\-ปฺปฏิปฺปสฺสทฺธ\-ตฺตา อุปฺปชฺชติ สมฺมาทิฏฺฐิ, มคฺคสฺเสตํ ผลํ.}

\prose{อภินิโรปน\-ฏฺเฐน สมฺมาสงฺกปฺโป มิจฺฉาสงฺกปฺปา วุฏฺฐาติ, ตทนุวตฺตก\-กิเลเสหิ จ ขนฺเธหิ จ วุฏฺฐาติ, พหิทฺธา จ สพฺพนิมิตฺเตหิ วุฏฺฐาติ. ตํปโยค\-ปฺปฏิปฺปสฺสทฺธ\-ตฺตา อุปฺปชฺชติ สมฺมาสงฺกปฺโป, มคฺคสฺเสตํ ผลํ.}

\prose{ปริคฺคห\-ฏฺเฐน สมฺมาวาจา มิจฺฉาวาจาย วุฏฺฐาติ, ตทนุวตฺตก\-กิเลเสหิ จ ขนฺเธหิ จ วุฏฺฐาติ, พหิทฺธา จ สพฺพนิมิตฺเตหิ วุฏฺฐาติ. ตํปโยค\-ปฺปฏิปฺปสฺสทฺธ\-ตฺตา อุปฺปชฺชติ สมฺมาวาจา, มคฺคสฺเสตํ ผลํ.}

\prose{สมุฏฺฐาน\-ฏฺเฐน สมฺมากมฺมนฺโต มิจฺฉากมฺมนฺตา วุฏฺฐาติ, ตทนุวตฺตก\-กิเลเสหิ จ ขนฺเธหิ จ วุฏฺฐาติ, พหิทฺธา จ สพฺพนิมิตฺเตหิ วุฏฺฐาติ. ตํปโยค\-ปฺปฏิปฺปสฺสทฺธ\-ตฺตา อุปฺปชฺชติ สมฺมากมฺมนฺโต, มคฺคสฺเสตํ ผลํ.}

\prose{โวทานฏฺเฐน สมฺมาอาชีโว มิจฺฉาอาชีวา วุฏฺฐาติ, ตทนุวตฺตก\-กิเลเสหิ จ ขนฺเธหิ จ วุฏฺฐาติ, พหิทฺธา จ สพฺพนิมิตฺเตหิ วุฏฺฐาติ. ตํปโยค\-ปฺปฏิปฺปสฺสทฺธ\-ตฺตา อุปฺปชฺชติ สมฺมาอาชีโว, มคฺคสฺเสตํ ผลํ.}

\prose{\csromanfolio{69}ปคฺคหฏฺเฐน สมฺมาวายาโม มิจฺฉาวายามา วุฏฺฐาติ, ตทนุวตฺตก\-กิเลเสหิ จ ขนฺเธหิ จ วุฏฺฐาติ, พหิทฺธา จ สพฺพนิมิตฺเตหิ วุฏฺฐาติ. ตํปโยค\-ปฺปฏิปฺปสฺสทฺธ\-ตฺตา อุปฺปชฺชติ สมฺมาวายาโม, มคฺคสฺเสตํ ผลํ.}

\prose{อุปฏฺฐาน\-ฏฺเฐน สมฺมาสติ มิจฺฉาสติยา วุฏฺฐาติ, ตทนุวตฺตก\-กิเลเสหิ จ ขนฺเธหิ จ วุฏฺฐาติ, พหิทฺธา จ สพฺพนิมิตฺเตหิ วุฏฺฐาติ. ตํปโยค\-ปฺปฏิปฺปสฺสทฺธ\-ตฺตา อุปฺปชฺชติ สมฺมาสติ, มคฺคสฺเสตํ ผลํ.}

\prose{อวิกฺเขป\-ฏฺเฐน สมฺมาสมาธิ มิจฺฉา\-สมาธิโต วุฏฺฐาติ, ตทนุวตฺตก\-กิเลเสหิ จ ขนฺเธหิ จ วุฏฺฐาติ, พหิทฺธา จ สพฺพนิมิตฺเตหิ วุฏฺฐาติ. ตํปโยค\-ปฺปฏิปฺปสฺสทฺธ\-ตฺตา อุปฺปชฺชติ สมฺมาสมาธิ, มคฺคสฺเสตํ ผลํ.}

\prose{สกทาคามิ\-มคฺคกฺขเณ ทสฺสนฏฺเฐน สมฺมาทิฏฺฐิ ฯเปฯ. อวิกฺเขป\-ฏฺเฐน สมฺมาสมาธิ โอฬาริกา กามราค\-สญฺโญชนา ปฏิฆ\-สญฺโญชนา, โอฬาริกา กามราคานุสยา ปฏิฆานุสยา วุฏฺฐาติ, ตทนุวตฺตก\-กิเลเสหิ จ ขนฺเธหิ จ วุฏฺฐาติ, พหิทฺธา จ สพฺพนิมิตฺเตหิ วุฏฺฐาติ. ตํปโยค\-ปฺปฏิปฺปสฺสทฺธ\-ตฺตา อุปฺปชฺชติ สมฺมาสมาธิ, มคฺคสฺเสตํ ผลํ.}

\prose{อนาคามิ\-มคฺคกฺขเณ ทสฺสนฏฺเฐน สมฺมาทิฏฺฐิ ฯเปฯ. อวิกฺเขป\-ฏฺเฐน สมฺมาสมาธิ อณุสหคตา กามราค\-สญฺโญชนา ปฏิฆ\-สญฺโญชนา, อณุสหคตา กามราคานุสยา ปฏิฆานุสยา วุฏฺฐาติ, ตทนุวตฺตก\-กิเลเสหิ จ ขนฺเธหิ จ วุฏฺฐาติ, พหิทฺธา จ สพฺพนิมิตฺเตหิ วุฏฺฐาติ. ตํปโยค\-ปฺปฏิปฺปสฺสทฺธ\-ตฺตา อุปฺปชฺชติ สมฺมาสมาธิ, มคฺคสฺเสตํ ผลํ.}

\prose{อรหตฺต\-มคฺคกฺขเณ ทสฺสนฏฺเฐน สมฺมาทิฏฺฐิ ฯเปฯ. อวิกฺเขป\-ฏฺเฐน สมฺมาสมาธิ รูปราคา อรูปราคา มานา อุทฺธจฺจา อวิชฺชาย, มานานุสยา ภวราคา\-นุสยา อวิชฺชานุสยา วุฏฺฐาติ, ตทนุวตฺตก\-กิเลเสหิ จ ขนฺเธหิ จ วุฏฺฐาติ, พหิทฺธา จ สพฺพนิมิตฺเตหิ วุฏฺฐาติ. ตํปโยค\-ปฺปฏิปฺปสฺสทฺธ\-ตฺตา อุปฺปชฺชติ สมฺมาสมาธิ, มคฺคสฺเสตํ ผลํ. ตํญาตฏฺเฐน ญาณํ, ปชานนฏฺเฐน ปญฺญา. เตน วุจฺจติ ปโยค\-ปฺปฏิปฺปสฺสทฺธิ\-ปญฺญา ผเล ญาณํ.}

\vspace{\tipitakaparskip}
\csromancenter{ผลญาณ\-นิทฺเทโส ทฺวาทสโม.}
\csromanmatikaanchor{mtk.15}
\csromantocmarkref{subsection}{13. วิมุตฺติญาณนิทฺเทส}{mtk.15}
\bookmark[dest={mtk.15},level=2]{13. วิมุตฺติญาณนิทฺเทส}
\csromanheader{2}{13. วิมุตฺติญาณ\-นิทฺเทส}
\proseitem{64}{\csromanfolio{70}กถํ ฉินฺนวฏุมา\-นุปสฺสเน ปญฺญา วิมุตฺติญาณํ\csromandash{} โสตาปตฺติ\-มคฺเคน สกฺกายทิฏฺฐิ วิจิกิจฺฉา สีลพฺพต\-ปรามาโส, ทิฏฺฐานุสโย วิจิกิจฺฉา\-นุสโย อตฺตโน จิตฺตสฺส อุปกฺกิเลสา สมฺมา สมุจฺฉินฺนา โหนฺติ, อิเมหิ ปญฺจหิ อุปกฺกิเลเสหิ สปริยุฏฺฐาเนหิ จิตฺตํ วิมุตฺตํ โหติ สุวิมุตฺตํ. ตํวิมุตฺติ\-ญาตฏฺเฐน ญาณํ, ปชานนฏฺเฐน ปญฺญา. เตน วุจฺจติ ฉินฺนวฏุมา\-นุปสฺสเน ปญฺญา วิมุตฺติญาณํ.}

\prose{สกทาคามิ\-มคฺเคน โอฬาริกํ กามราค\-สญฺโญชนํ ปฏิฆ\-สญฺโญชนํ, โอฬาริโก กามราคานุสโย ปฏิฆานุสโย อตฺตโน จิตฺตสฺส อุปกฺกิเลสา สมฺมา สมุจฺฉินฺนา โหนฺติ, อิเมหิ จตูหิ อุปกฺกิเลเสหิ สปริยุฏฺฐาเนหิ จิตฺตํ วิมุตฺตํ โหติ สุวิมุตฺตํ. ตํวิมุตฺติ\-ญาตฏฺเฐน ญาณํ, ปชานนฏฺเฐน ปญฺญา. เตน วุจฺจติ ฉินฺนวฏุมา\-นุปสฺสเน ปญฺญา วิมุตฺติญาณํ.}

\prose{อนาคามิมคฺเคน อณุสหคตํ กามราค\-สญฺโญชนํ ปฏิฆ\-สญฺโญชนํ, อณุสหคโต กามราคานุสโย ปฏิฆานุสโย อตฺตโน จิตฺตสฺส อุปกฺกิเลสา สมฺมา สมุจฺฉินฺนา โหนฺติ, อิเมหิ จตูหิ อุปกฺกิเลเสหิ สปริยุฏฺฐาเนหิ จิตฺตํ วิมุตฺตํ โหติ สุวิมุตฺตํ. ตํวิมุตฺติ\-ญาตฏฺเฐน ญาณํ, ปชานนฏฺเฐน ปญฺญา. เตน วุจฺจติ ฉินฺนวฏุมา\-นุปสฺสเน ปญฺญา วิมุตฺติญาณํ.}

\prosewithcloserruled{อรหตฺต\-มคฺเคน รูปราโค อรูปราโค มาโน อุทฺธจฺจํ อวิชฺชา, มานานุสโย ภวราคา\-นุสโย อวิชฺชานุสโย อตฺตโน จิตฺตสฺส อุปกฺกิเลสา สมฺมา สมุจฺฉินฺนา โหนฺติ, อิเมหิ อฏฺฐหิ อุปกฺกิเลเสหิ สปริยุฏฺฐาเนหิ จิตฺตํ วิมุตฺตํ โหติ สุวิมุตฺตํ. ตํวิมุตฺติ\-ญาตฏฺเฐน ญาณํ, ปชานนฏฺเฐน ปญฺญา. เตน วุจฺจติ ฉินฺนวฏุมา\-นุปสฺสเน ปญฺญา วิมุตฺติญาณํ. ตํญาตฏฺเฐน ญาณํ, ปชานนฏฺเฐน ปญฺญา. เตน วุจฺจติ ฉินฺนวฏุมา\-นุปสฺสเน ปญฺญา วิมุตฺติญาณํ.}{วิมุตฺติญาณ\-นิทฺเทโส เตรสโม.}
\csromanmatikaanchor{mtk.16}
\csromantocmarkref{subsection}{14. ปจฺจเวกฺขณญาณนิทฺเทส}{mtk.16}
\bookmark[dest={mtk.16},level=2]{14. ปจฺจเวกฺขณญาณนิทฺเทส}
\csromanheader{2}{14. ปจฺจเวกฺขณญาณ\-นิทฺเทส}
\proseitem{65}{กถํ ตทา สมุทาคเต ธมฺเม ปสฺสเน ปญฺญา ปจฺจเวกฺขเณ ญาณํ\csromandash{}โสตาปตฺติมคฺค\-กฺขเณ ทสฺสนฏฺเฐน สมฺมาทิฏฺฐิ ตทา สมุทาคตา. \csromanfolio{71}อภินิโรปน\-ฏฺเฐน สมฺมาสงฺกปฺโป ตทา สมุทาคโต. ปริคฺคห\-ฏฺเฐน สมฺมาวาจา ตทา สมุทาคตา. สมุฏฺฐาน\-ฏฺเฐน สมฺมากมฺมนฺโต ตทา สมุทาคโต. โวทานฏฺเฐน สมฺมาอาชีโว ตทา สมุทาคโต. ปคฺคหฏฺเฐน สมฺมาวายาโม ตทา สมุทาคโต. อุปฏฺฐาน\-ฏฺเฐน สมฺมาสติ ตทา สมุทาคตา. อวิกฺเขป\-ฏฺเฐน สมฺมาสมาธิ ตทา สมุทาคโต.}

\prose{อุปฏฺฐาน\-ฏฺเฐน สติสมฺโพชฺฌงฺโค ตทา สมุทาคโต. ปวิจยฏฺเฐน ธมฺมวิจย\-สมฺโพชฺฌงฺโค ตทา สมุทาคโต. ปคฺคหฏฺเฐน วีริย\-สมฺโพชฺฌงฺโค ตทา สมุทาคโต. ผรณฏฺเฐน ปีติสมฺโพชฺฌงฺโค ตทา สมุทาคโต. อุปสมฏฺเฐน ปสฺสทฺธิ\-สมฺโพชฺฌงฺโค ตทา สมุทาคโต. อวิกฺเขป\-ฏฺเฐน สมาธิ\-สมฺโพชฺฌงฺโค ตทา สมุทาคโต. ปฏิสงฺขา\-นฏฺเฐน อุเปกฺขา\-สมฺโพชฺฌงฺโค ตทา สมุทาคโต.}

\prose{อสฺสทฺธิเย อกมฺปิยฏฺเฐน สทฺธาพลํ ตทา สมุทาคตํ, โกสชฺเช อกมฺปิยฏฺเฐน วีริยพลํ ตทา สมุทาคตํ. ปมาเท อกมฺปิยฏฺเฐน สติพลํ ตทา สมุทาคตํ. อุทฺธจฺเจ อกมฺปิยฏฺเฐน สมาธิพลํ ตทา สมุทาคตํ. อวิชฺชาย อกมฺปิยฏฺเฐน ปญฺญาพลํ ตทา สมุทาคตํ.}

\prose{อธิโมกฺข\-ฏฺเฐน สทฺธินฺทฺริยํ ตทา สมุทาคตํ. ปคฺคหฏฺเฐน วีริยินฺทฺริยํ ตทา สมุทาคตํ. อุปฏฺฐาน\-ฏฺเฐน สตินฺทฺริยํ ตทา สมุทาคตํ. อวิกฺเขป\-ฏฺเฐน สมาธินฺทฺริยํ ตทา สมุทาคตํ. ทสฺสนฏฺเฐน ปญฺญินฺทฺริยํ ตทา สมุทาคตํ.}

\prose{อาธิปเตยฺย\-ฏฺเฐน อินฺทฺริยา ตทา สมุทาคตา. อกมฺปิยฏฺเฐน พลา ตทา สมุทาคตา. นิยฺยานฏฺเฐน สมฺโพชฺฌงฺคา ตทา สมุทาคตา. เหตุฏฺเฐน มคฺโค ตทา สมุทาคโต. อุปฏฺฐาน\-ฏฺเฐน สติปฏฺฐานา ตทา สมุทาคตา. ปทหนฏฺเฐน สมฺมปฺปธานา ตทา สมุทาคตา. อิชฺฌนฏฺเฐน อิทฺธิปาทา ตทา สมุทาคตา. ตถฏฺเฐน สจฺจา ตทา สมุทาคตา. อวิกฺเขป\-ฏฺเฐน สมโถ ตทา สมุทาคโต. อนุปสฺสนฏฺเฐน วิปสฺสนา ตทา สมุทาคตา. เอกรสฏฺเฐน สมถ\-วิปสฺสนา ตทา สมุทาคตา. อนติวตฺตน\-ฏฺเฐน ยุคนทฺธํ ตทา สมุทาคตํ. สํวรฏฺเฐน สีลวิสุทฺธิ ตทา สมุทาคตา. อวิกฺเขป\-ฏฺเฐน จิตฺตวิสุทฺธิ ตทา สมุทาคตา. ทสฺสนฏฺเฐน ทิฏฺฐิวิสุทฺธิ ตทา สมุทาคตา. วิมุตฺตฏฺเฐน วิโมกฺโข ตทา สมุทาคโต. ปฏิเวธ\-ฏฺเฐน วิชฺชา ตทา สมุทาคตา. ปริจฺจาค\-ฏฺเฐน วิมุตฺติ ตทา สมุทาคตา. สมุจฺเฉท\-ฏฺเฐน ขเย ญาณํ ตทา สมุทาคตํ.}

\prose{\csromanfolio{72}ฉนฺโท มูลฏฺเฐน ตทา สมุทาคโต. มนสิกาโร สมุฏฺฐาน\-ฏฺเฐน ตทา สมุทาคโต. ผสฺโส สโมธาน\-ฏฺเฐน ตทา สมุทาคโต. เวทนา สโมสรณ\-ฏฺเฐน ตทา สมุทาคตา. สมาธิ ปมุขฏฺเฐน ตทา สมุทาคโต. สติ อาธิปเตยฺย\-ฏฺเฐน ตทา สมุทาคตา. ปญฺญา ตทุตฺต\-รฏฺเฐน ตทา สมุทาคตา. วิมุตฺติ สารฏฺเฐน ตทา สมุทาคตา. อมโตคธํ นิพฺพานํ ปริโยสาน\-ฏฺเฐน ตทา สมุทาคตํ. วุฏฺฐหิตฺวา ปจฺจเวกฺขติ, อิเม ธมฺมา ตทา สมุทาคตา.}

\prose{โสตาปตฺติ\-ผลกฺขเณ ทสฺสนฏฺเฐน สมฺมาทิฏฺฐิ ตทา สมุทาคตา. อภินิโรปน\-ฏฺเฐน สมฺมาสงฺกปฺโป ตทา สมุทาคโต ฯเปฯ. ปฏิปฺปสฺสทฺธ\-ฏฺเฐน อนุปฺปาเท ญาณํ ตทา สมุทาคตํ.  ฉนฺโท มูลฏฺเฐน ตทา สมุทาคโต. มนสิกาโร สมุฏฺฐาน\-ฏฺเฐน ตทา สมุทาคโต. ผสฺโส สโมธาน\-ฏฺเฐน ตทา สมุทาคโต. เวทนา สโมสรณ\-ฏฺเฐน ตทา สมุทาคตา. สมาธิ ปมุขฏฺเฐน ตทา สมุทาคโต. สติ อาธิปเตยฺย\-ฏฺเฐน ตทา สมุทาคตา. ปญฺญา ตทุตฺต\-รฏฺเฐน ตทา สมุทาคตา. วิมุตฺติ สารฏฺเฐน ตทา สมุทาคตา. อมโตคธํ นิพฺพานํ ปริโยสาน\-ฏฺเฐน ตทา สมุทาคตํ. วุฏฺฐหิตฺวา ปจฺจเวกฺขติ, อิเม ธมฺมา ตทา สมุทาคตา.}

\prose{สกทาคามิ\-มคฺคกฺขเณ ฯเปฯ. สกทาคามิ\-ผลกฺขเณ ฯเปฯ. อนาคามิ\-มคฺคกฺขเณ ฯเปฯ. อนาคามิ\-ผลกฺขเณ ฯเปฯ. อรหตฺต\-มคฺคกฺขเณ ทสฺสนฏฺเฐน สมฺมาทิฏฺฐิ ตทา สมุทาคตา ฯเปฯ. สมุจฺเฉท\-ฏฺเฐน ขเย ญาณํ ตทา สมุทาคตํ.  ฉนฺโท มูลฏฺเฐน ตทา สมุทาคโต ฯเปฯ อมโตคธํ นิพฺพานํ ปริโยสาน\-ฏฺเฐน ตทา สมุทาคตํ. วุฏฺฐหิตฺวา ปจฺจเวกฺขติ, อิเม ธมฺมา ตทา สมุทาคตา.}

\prose{อรหตฺต\-ผลกฺขเณ ทสฺสนฏฺเฐน สมฺมาทิฏฺฐิ ตทา สมุทาคตา ฯเปฯ. ปฏิปฺปสฺสทฺธ\-ฏฺเฐน อนุปฺปาเท ญาณํ ตทา สมุทาคตํ. ฉนฺโท มูลฏฺเฐน ตทา สมุทาคโต ฯเปฯ อมโตคธํ นิพฺพานํ ปริโยสาน\-ฏฺเฐน ตทา สมุทาคตํ. วุฏฺฐหิตฺวา ปจฺจเวกฺขติ, อิเม ธมฺมา ตทา สมุทาคตา, ตํญาตฏฺเฐน ญาณํ, ปชานนฏฺเฐน ปญฺญา. เตน วุจฺจติ ตทา สมุทาคเต ธมฺเม ปสฺสเน ปญฺญา ปจฺจเวกฺขเณ ญาณํ.}

\vspace{\tipitakaparskip}
\csromancenter{ปจฺจเวกฺขณญาณ\-นิทฺเทโส จุทฺทสโม.}
\csromanmatikaanchor{mtk.17}
\csromantocmarkref{subsection}{15. วตฺถุนานตฺตญาณนิทฺเทส}{mtk.17}
\bookmark[dest={mtk.17},level=2]{15. วตฺถุนานตฺตญาณนิทฺเทส}
\csromanheader{2}{15. วตฺถุนานตฺต\-ญาณ\-นิทฺเทส}
\proseitem{66}{\csromanfolio{73}กถํ อชฺฌตฺต\-ววตฺถาเน ปญฺญา วตฺถุนานตฺเต ญาณํ, กถํ อชฺฌตฺตธมฺเม ววตฺเถติ\csromandash{}จกฺขุํ อชฺฌตฺตํ ววตฺเถติ, โสตํ อชฺฌตฺตํ ววตฺเถติ, ฆานํ อชฺฌตฺตํ ววตฺเถติ, ชิวฺหํ อชฺฌตฺตํ ววตฺเถติ, กายํ อชฺฌตฺตํ ววตฺเถติ, มนํ อชฺฌตฺตํ ววตฺเถติ.}

\prose{กถํ จกฺขุํ อชฺฌตฺตํ ววตฺเถติ? จกฺขุ อวิชฺชาสมฺภูตนฺติ ววตฺเถติ, จกฺขุ ตณฺหา\-สมฺภูตนฺติ ววตฺเถติ, จกฺขุ กมฺมสมฺภูตนฺติ ววตฺเถติ, จกฺขุ อาหาร\-สมฺภูตนฺติ ววตฺเถติ, จกฺขุ จตุนฺนํ มหาภูตานํ อุปาทายาติ ววตฺเถติ, จกฺขุ อุปฺปนฺนนฺติ ววตฺเถติ, จกฺขุ สมุทาคตนฺติ ววตฺเถติ. จกฺขุ อหุตฺวา สมฺภูตํ, หุตฺวา น ภวิสฺสตีติ ววตฺเถติ, จกฺขุํ อนฺตวนฺตโต ววตฺเถติ, จกฺขุ อทฺธุคํ อสสฺสตํ วิปริณาม\-ธมฺมนฺติ ววตฺเถติ, จกฺขุ อนิจฺจํ สงฺขตํ ปฏิจฺจ\-สมุปฺปนฺนํ ขยธมฺมํ วยธมฺมํ วิราคธมฺมํ นิโรธ\-ธมฺมนฺติ ววตฺเถติ, จกฺขุํ อนิจฺจโต ววตฺเถติ, โน นิจฺจโต, ทุกฺขโต ววตฺเถติ, โน สุขโต, อนตฺตโต ววตฺเถติ, โน อตฺตโต, นิพฺพินฺทติ, โน นนฺทติ, วิรชฺชติ, โน รชฺชติ, นิโรเธติ, โน สมุเทติ, ปฏินิสฺสชฺชติ, โน อาทิยติ.  อนิจฺจโต ววตฺเถนฺโต นิจฺจสญฺญํ ปชหาติ, ทุกฺขโต ววตฺเถนฺโต สุขสญฺญํ ปชหาติ, อนตฺตโต ววตฺเถนฺโต อตฺตสญฺญํ ปชหาติ, นิพฺพินฺทนฺโต นนฺทิํ ปชหาติ, วิรชฺชนฺโต ราคํ ปชหาติ, นิโรเธนฺโต สมุทยํ ปชหาติ, ปฏินิสฺสชฺชนฺโต อาทานํ ปชหาติ. เอวํ จกฺขุํ อชฺฌตฺตํ ววตฺเถติ.}

\prose{กถํ โสตํ อชฺฌตฺตํ ววตฺเถติ? โสตํ อวิชฺชาสมฺภูตนฺติ ววตฺเถติ ฯเปฯ เอวํ โสตํ อชฺฌตฺตํ ววตฺเถติ.  กถํ ฆานํ อชฺฌตฺตํ ววตฺเถติ? ฆานํ อวิชฺชาสมฺภูตนฺติ ววตฺเถติ ฯเปฯ เอวํ ฆานํ อชฺฌตฺตํ ววตฺเถติ.  กถํ ชิวฺหํ อชฺฌตฺตํ ววตฺเถติ? ชิวฺหา อวิชฺชา\-สมฺภูตาติ ววตฺเถติ, ชิวฺหา ตณฺหาสมฺภูตาติ ววตฺเถติ, ชิวฺหา กมฺม\-สมฺภูตาติ ววตฺเถติ, ชิวฺหา อาหารสมฺภูตาติ ววตฺเถติ, ชิวฺหา จตุนฺนํ มหาภูตานํ อุปาทายาติ ววตฺเถติ, ชิวฺหา อุปฺปนฺนาติ ววตฺเถติ, ชิวฺหา สมุทาคตาติ ววตฺเถติ, ชิวฺหา อหุตฺวา สมฺภูตา, หุตฺวา น ภวิสฺสตีติ ววตฺเถติ, ชิวฺหํ อนฺตวนฺตโต ววตฺเถติ, ชิวฺหา อทฺธุวา อสสฺสตา วิปริณาม\-ธมฺมาติ ววตฺเถติ, ชิวฺหา อนิจฺจา สงฺขตา ปฏิจฺจ\-สมุปฺปนฺนา ขยธมฺมา วยธมฺมา วิราคธมฺมา นิโรธธมฺมาติ ววตฺเถติ, ชิวฺหํ อนิจฺจโต ววตฺเถติ, \csromanfolio{74}โน นิจฺจโต ฯเปฯ ปฏินิสฺสชฺชติ, โน อาทิยติ.  อนิจฺจโต ววตฺเถนฺโต นิจฺจสญฺญํ ปชหาติ ฯเปฯ ปฏินิสฺสชฺชนฺโต อาทานํ ปชหาติ. เอวํ ชิวฺหํ อชฺฌตฺตํ ววตฺเถติ.}

\prose{กถํ กายํ อชฺฌตฺตํ ววตฺเถติ? กาโย อวิชฺชาสมฺภูโตติ ววตฺเถติ, กาโย ตณฺหา\-สมฺภูโตติ ววตฺเถติ, กาโย กมฺมสมฺภูโตติ ววตฺเถติ, กาโย อาหาร\-สมฺภูโตติ ววตฺเถติ, กาโย จตุนฺนํ มหาภูตานํ อุปาทายาติ ววตฺเถติ, กาโย อุปฺปนฺโนติ ววตฺเถติ, กาโย สมุทาคโตติ ววตฺเถติ, กาโย อหุตฺวา สมฺภูโต, หุตฺวา น ภวิสฺสตีติ ววตฺเถติ, กายํ อนฺตวนฺตโต ววตฺเถติ, กาโย อทฺธุโว อสสฺสโต วิปริณาม\-ธมฺโมติ ววตฺเถติ, กาโย อนิจฺโจ สงฺขโต ปฏิจฺจ\-สมุปฺปนฺโน ขยธมฺโม วยธมฺโม วิราคธมฺโม นิโรธ\-ธมฺโมติ ววตฺเถติ, กายํ อนิจฺจโต ววตฺเถติ, โน นิจฺจโต, ทุกฺขโต ววตฺเถติ, โน สุขโต ฯเปฯ ปฏินิสฺสชฺชติ, โน อาทิยติ.  อนิจฺจโต ววตฺเถนฺโต นิจฺจสญฺญํ ปชหาติ, ทุกฺขโต ววตฺเถนฺโต สุขสญฺญํ ปชหาติ ฯเปฯ ปฏินิสฺสชฺชนฺโต อาทานํ ปชหาติ. เอวํ กายํ อชฺฌตฺตํ ววตฺเถติ.}

\prosewithcloserruled{กถํ มนํ อชฺฌตฺตํ ววตฺเถติ? มโน อวิชฺชาสมฺภูโตติ ววตฺเถติ, มโน ตณฺหา\-สมฺภูโตติ ววตฺเถติ, มโน กมฺมสมฺภูโตติ ววตฺเถติ, มโน อาหาร\-สมฺภูโตติ ววตฺเถติ, มโน อุปฺปนฺโนติ ววตฺเถติ, มโน สมุทาคโตติ ววตฺเถติ, มโน อหุตฺวา สมฺภูโต, หุตฺวา น ภวิสฺสตีติ ววตฺเถติ, มนํ อนฺตวนฺตโต ววตฺเถติ, มโน อทฺธุโว อสสฺสโต วิปริณาม\-ธมฺโมติ ววตฺเถติ, มโน อนิจฺโจ สงฺขโต ปฏิจฺจ\-สมุปฺปนฺโน ขยธมฺโม วยธมฺโม วิราคธมฺโม นิโรธ\-ธมฺโมติ ววตฺเถติ. มนํ อนิจฺจโต ววตฺเถติ, โน นิจฺจโต, ทุกฺขโต ววตฺเถติ, โน สุขโต, อนตฺตโต ววตฺเถติ, โน อตฺตโต นิพฺพินฺทติ, โน นนฺทติ, วิรชฺชติ, โน รชฺชติ, นิโรเธติ, โน สมุเทติ, ปฏินิสฺสชฺชติ, โน อาทิยติ.  อนิจฺจโต ววตฺเถนฺโต นิจฺจสญฺญํ ปชหาติ. ทุกฺขโต ววตฺเถนฺโต สุขสญฺญํ ปชหาติ, อนตฺตโต ววตฺเถนฺโต อตฺตสญฺญํ ปชหาติ, นิพฺพินฺทนฺโต นนฺทิํ ปชหาติ, วิรชฺชนฺโต ราคํ ปชหาติ, นิโรเธนฺโต สมุทยํ ปชหาติ, ปฏินิสฺสชฺชนฺโต อาทานํ ปชหาติ. เอวํ มนํ อชฺฌตฺตํ ววตฺเถติ. เอวํ อชฺฌตฺตธมฺเม ววตฺเถติ. ตํญาตฏฺเฐน \csromanfolio{75}ญาณํ ปชานนฏฺเฐน ปญฺญา. เตน วุจฺจติ อชฺฌตฺต\-ววตฺถาเน ปญฺญา วตฺถุนานตฺเต ญาณํ.}{วตฺถุนานตฺต\-ญาณ\-นิทฺเทโส ปนฺนรสโม.}
\csromanmatikaanchor{mtk.18}
\csromantocmarkref{subsection}{16. โคจรนานตฺตญาณนิทฺเทส}{mtk.18}
\bookmark[dest={mtk.18},level=2]{16. โคจรนานตฺตญาณนิทฺเทส}
\csromanheader{2}{16. โคจร\-นานตฺต\-ญาณ\-นิทฺเทส}
\proseitem{67}{กถํ พหิทฺธา ววตฺถาเน ปญฺญา โคจรนานตฺเต ญาณํ\csromandash{}กถํ พหิทฺธา ธมฺเม ววตฺเถติ, รูเป พหิทฺธา ววตฺเถติ, สทฺเท พหิทฺธา ววตฺเถติ, คนฺเธ พหิทฺธา ววตฺเถติ, รเส พหิทฺธา ววตฺเถติ, โผฏฺฐพฺเพ พหิทฺธา ววตฺเถติ, ธมฺเม พหิทฺธา ววตฺเถติ?}

\prose{กถํ รูเป พหิทฺธา ววตฺเถติ? รูปา อวิชฺชา\-สมฺภูตาติ ววตฺเถติ, รูปา ตณฺหาสมฺภูตาติ ววตฺเถติ, รูปา กมฺม\-สมฺภูตาติ ววตฺเถติ, รูปา อาหารสมฺภูตาติ ววตฺเถติ. รูปา จตุนฺนํ มหาภูตานํ อุปาทายาติ ววตฺเถติ, รูปา อุปฺปนฺนาติ ววตฺเถติ, รูปา สมุทาคตาติ ววตฺเถติ, รูปา อหุตฺวา สมฺภูตา, หุตฺวา น ภวิสฺสนฺตีติ ววตฺเถติ, รูเป อนฺตวนฺตโต ววตฺเถติ, รูปา อทฺธุวา อสสฺสตา วิปริณาม\-ธมฺมาติ ววตฺเถติ. รูปา อนิจฺจา สงฺขตา ปฏิจฺจ\-สมุปฺปนฺนา ขยธมฺมา วยธมฺมา วิราคธมฺมา นิโรธธมฺมาติ ววตฺเถติ, รูเป อนิจฺจโต ววตฺเถติ, โน นิจฺจโต, ทุกฺขโต ววตฺเถติ โน สุขโต, อนตฺตโต ววตฺเถติ, โน อตฺตโต, นิพฺพินฺทติ, โน นนฺทติ, วิรชฺชติ โน รชฺชติ, นิโรเธติ, โน สมุเทติ, ปฏินิสฺสชฺชติ, โน อาทิยติ.  อนิจฺจโต ววตฺเถนฺโต นิจฺจสญฺญํ ปชหาติ, ทุกฺขโต ววตฺเถนฺโต สุขสญฺญํ ปชหาติ, อนตฺตโต ววตฺเถนฺโต อตฺตสญฺญํ ปชหาติ, นิพฺพินฺทนฺโต นนฺทิํ ปชหาติ, วิรชฺชนฺโต ราคํ ปชหาติ, นิโรเธนฺโต สมุทยํ ปชหาติ, ปฏินิสฺสชฺชนฺโต อาทานํ ปชหาติ. เอวํ รูเป พหิทฺธา ววตฺเถติ.}

\prose{กถํ สทฺเท พหิทฺธา ววตฺเถติ? สทฺทา จตุนฺนํ มหาภูตานํ อุปาทายาติ ววตฺเถติ, สทฺทา อุปฺปนฺนาติ ววตฺเถติ, สทฺทา สมุทาคตาติ ววตฺเถติ. สทฺธา อหุตฺวา สมฺภูตา, หุตฺวา น ภวิสฺสนฺตีติ ววตฺเถติ, สทฺเท อนฺตวนฺตโต ววตฺเถติ, สทฺทา \csromanfolio{76}อทฺธุวา อสสฺสตา วิปริณาม\-ธมฺมาติ ววตฺเถติ, สทฺทา อนิจฺจา สงฺขตา ปฏิจฺจ\-สมุปฺปนฺนา ขยธมฺมา วยธมฺมา วิราคธมฺมา นิโรธธมฺมาติ ววตฺเถติ, สทฺเท อนิจฺจโต ววตฺเถติ, โน นิจฺจโต ฯเปฯ. เอวํ สทฺเท พหิทฺธา ววตฺเถติ.}

\prose{กถํ คนฺเธ พหิทฺธา ววตฺเถติ? คนฺธา อวิชฺชา\-สมฺภูตาติ ววตฺเถติ, คนฺธา ตณฺหาสมฺภูตาติ ววตฺเถติ ฯเปฯ. เอวํ คนฺเธ พหิทฺธา ววตฺเถติ.  กถํ รเส พหิทฺธา ววตฺเถติ? รสา อวิชฺชา\-สมฺภูตาติ ววตฺเถติ, รสา ตณฺหาสมฺภูตาติ ววตฺเถติ ฯเปฯ. เอวํ รเส พหิทฺธา ววตฺเถติ.  กถํ โผฏฺฐพฺเพ พหิทฺธา ววตฺเถติ? โผฏฺฐพฺพา อวิชฺชา\-สมฺภูตาติ ววตฺเถติ, โผฏฺฐพฺพา ตณฺหาสมฺภูตาติ ววตฺเถติ, โผฏฺฐพฺพา กมฺม\-สมฺภูตาติ ววตฺเถติ, โผฏฺฐพฺพา อาหารสมฺภูตาติ ววตฺเถติ, โผฏฺฐพฺพา อุปฺปนฺนาติ ววตฺเถติ. โผฏฺฐพฺพา สมุทาคตาติ ววตฺเถติ ฯเปฯ. เอวํ โผฏฺฐพฺเพ พหิทฺธา ววตฺเถติ.}

\prose{กถํ ธมฺเม พหิทฺธา ววตฺเถติ? ธมฺมา อวิชฺชา\-สมฺภูตาติ ววตฺเถติ, ธมฺมา ตณฺหาสมฺภูตาติ ววตฺเถติ, ธมฺมา กมฺม\-สมฺภูตาติ ววตฺเถติ, ธมฺมา อาหารสมฺภูตาติ ววตฺเถติ, ธมฺมา อุปฺปนฺนาติ ววตฺเถติ, ธมฺมา สมุทาคตาติ ววตฺเถติ, ธมฺมา อหุตฺวา สมฺภูตา, หุตฺวา น ภวิสฺสนฺตีติ ววตฺเถติ, ธมฺเม อนฺตวนฺตโต ววตฺเถติ, ธมฺมา อทฺธุวา อสสฺสตา วิปริณาม\-ธมฺมาติ ววตฺเถติ, ธมฺมา อนิจฺจา สงฺขตา ปฏิจฺจ\-สมุปฺปนฺนา ขยธมฺมา วยธมฺมา วิราคธมฺมา นิโรธธมฺมาติ ววตฺเถติ, ธมฺเม อนิจฺจโต ววตฺเถติ, โน นิจฺจโต, ทุกฺขโต ววตฺเถติ, โน สุขโต, อนตฺตโต ววตฺเถติ, โน อตฺตโต, นิพฺพินฺทติ, โน นนฺทติ, วิรชฺชติ, โน รชฺชติ, นิโรเธติ, โน สมุเทติ, ปฏินิสฺสชฺชติ, โน อาทิยติ.  อนิจฺจโต ววตฺเถนฺโต นิจฺจสญฺญํ ปชหาติ, ทุกฺขโต ววตฺเถนฺโต สุขสญฺญํ ปชหาติ, อนตฺตโต ววตฺเถนฺโต อตฺตสญฺญํ ปชหาติ, นิพฺพินฺทนฺโต นนฺทิํ ปชหาติ, วิรชฺชนฺโต ราคํ ปชหาติ, นิโรเธนฺโต สมุทยํ ปชหาติ, ปฏินิสฺสชฺชนฺโต อาทานํ ปชหาติ. เอวํ ธมฺเม พหิทฺธา ววตฺเถติ. ตํญาตฏฺเฐน ญาณํ, ปชานนฏฺเฐน ปญฺญา. เตน วุจฺจติ พหิทฺธา ววตฺถาเน ปญฺญา โคจรนานตฺเต ญาณํ.}

\vspace{\tipitakaparskip}
\csromancenter{โคจร\-นานตฺต\-ญาณ\-นิทฺเทโส โสฬสโม.}
\csromanmatikaanchor{mtk.19}
\csromantocmarkref{subsection}{17. จริยานานตฺตญาณนิทฺเทส}{mtk.19}
\bookmark[dest={mtk.19},level=2]{17. จริยานานตฺตญาณนิทฺเทส}
\csromanheader{2}{17. จริยา\-นานตฺต\-ญาณ\-นิทฺเทส}
\proseitem{68}{\csromanfolio{77}กถํ จริยา\-ววตฺถาเน ปญฺญา จริยานานตฺเต ญาณํ\csromandash{}จริยาติ ติสฺโส จริยาโย วิญฺญาณจริยา, อญฺญาณจริยา, ญาณจริยา.}

\prose{กตมา วิญฺญาณจริยา? ทสฺสนตฺถาย อาวชฺชน\-กิริยา\-พฺยากตา วิญฺญาณจริยา รูเปสุ, ทสฺสนฏฺโฐ จกฺขุวิญฺญาณํ วิญฺญาณจริยา รูเปสุ, ทิฏฺฐตฺตา อภินิโรปนา วิปาก\-มโนธาตุ\csromansharedfootnote{50cf3766aef18f09}{อภินิโรปนวิปากมโนธาตุ (สฺยา)} วิญฺญาณจริยา รูเปสุ, อภินิโรปิต\-ตฺตา วิปาก\-มโนวิญฺญาณ\-ธาตุ วิญฺญาณจริยา รูเปสุ.  สวนตฺถาย อาวชฺชน\-กิริยา\-พฺยากตา วิญฺญาณจริยา สทฺเทสุ, สวนฏฺโฐ โสตวิญฺญาณํ วิญฺญาณจริยา สทฺเทสุ. สุตตฺตา อภินิโรปนา วิปาก\-มโนธาตุ วิญฺญาณจริยา สทฺเทสุ, อภินิโรปิต\-ตฺตา วิปาก\-มโนวิญฺญาณ\-ธาตุ วิญฺญาณจริยา สทฺเทสุ.  ฆายนตฺถาย อาวชฺชน\-กิริยา\-พฺยากตา วิญฺญาณจริยา คนฺเธสุ, ฆายนฏฺโฐ\csromansharedfootnote{1fcd9aa9da49f4f0}{ฆายนตฺโถ (สฺยา, ก)} ฆานวิญฺญาณํ วิญฺญาณจริยา คนฺเธสุ, ฆายิตตฺตา อภินิโรปนา วิปาก\-มโนธาตุ วิญฺญาณจริยา คนฺเธสุ, อภินิโรปิต\-ตฺตา วิปาก\-มโนวิญฺญาณ\-ธาตุ วิญฺญาณจริยา คนฺเธสุ.  สายนตฺถาย อาวชฺชน\-กิริยา\-พฺยากตา วิญฺญาณจริยา รเสสุ, สายนฏฺโฐ ชิวฺหาวิญฺญาณํ วิญฺญาณจริยา รเสสุ, สายิตตฺตา อภินิโรปนา วิปาก\-มโนธาตุ วิญฺญาณจริยา รเสสุ, อภินิโรปิต\-ตฺตา วิปาก\-มโนวิญฺญาณ\-ธาตุ วิญฺญาณจริยา รเสสุ.  ผุสนตฺถาย อาวชฺชน\-กิริยา\-พฺยากตา วิญฺญาณจริยา โผฏฺฐพฺเพสุ, ผุสนฏฺโฐ กายวิญฺญาณํ วิญฺญาณจริยา โผฏฺฐพฺเพสุ, ผุฏฺฐตฺตา อภินิโรปนา วิปาก\-มโนธาตุ วิญฺญาณจริยา โผฏฺฐพฺเพสุ, อภินิโรปิต\-ตฺตา วิปาก\-มโนวิญฺญาณ\-ธาตุ วิญฺญาณจริยา โผฏฺฐพฺเพสุ.  วิชานนตฺถาย อาวชฺชน\-กิริยา\-พฺยากตา วิญฺญาณจริยา ธมฺเมสุ, วิชานนฏฺโฐ มโนวิญฺญาณํ วิญฺญาณจริยา ธมฺเมสุ. วิญฺญาตตฺตา อภินิโรปนา วิปาก\-มโนธาตุ วิญฺญาณจริยา ธมฺเมสุ, อภินิโรปิต\-ตฺตา วิปาก\-มโนวิญฺญาณ\-ธาตุ วิญฺญาณจริยา ธมฺเมสุ.}

\proseitem{69}{\textbf{วิญฺญาณจริยา}ติ เกนฏฺเฐน วิญฺญาณจริยา? นีราคา จรตีติ วิญฺญาณจริยา, นิทฺโทสา จรตีติ วิญฺญาณจริยา, นิมฺโมหา จรตีติ วิญฺญาณจริยา, นิมฺมานา จรตีติ วิญฺญาณจริยา, นิทฺทิฏฺฐิ จรตีติ \csromanfolio{78}วิญฺญาณจริยา, นิอุทฺธจฺจา จรตีติ วิญฺญาณจริยา. นิพฺพิจิกิจฺฉา จรตีติ วิญฺญาณจริยา, นานุสยา จรตีติ วิญฺญาณจริยา. ราควิปฺปยุตฺตา จรตีติ วิญฺญาณจริยา, โทสวิปฺปยุตฺตา จรตีติ วิญฺญาณจริยา, โมหวิปฺปยุตฺตา จรตีติ วิญฺญาณจริยา, มานวิปฺปยุตฺตา จรตีติ วิญฺญาณจริยา, ทิฏฺฐิ\-วิปฺปยุตฺตา จรตีติ วิญฺญาณจริยา, อุทฺธจฺจ\-วิปฺปยุตฺตา จรตีติ วิญฺญาณจริยา, วิจิกิจฺฉา\-วิปฺปยุตฺตา จรตีติ วิญฺญาณจริยา, อนุสย\-วิปฺปยุตฺตา จรตีติ วิญฺญาณจริยา, กุสเลหิ กมฺเมหิ สมฺปยุตฺตา จรตีติ วิญฺญาณจริยา, อกุสเลหิ กมฺเมหิ วิปฺปยุตฺตา จรตีติ วิญฺญาณจริยา, สาวชฺเชหิ กมฺเมหิ วิปฺปยุตฺตา จรตีติ วิญฺญาณจริยา, อนวชฺเชหิ กมฺเมหิ สมฺปยุตฺตา จรตีติ วิญฺญาณจริยา, กเณฺหหิ กมฺเมหิ วิปฺปยุตฺตา จรตีติ วิญฺญาณจริยา, สุกฺเกหิ กมฺเมหิ สมฺปยุตฺตา จรตีติ วิญฺญาณจริยา, สุขุทฺรเยหิ กมฺเมหิ สมฺปยุตฺตา จรตีติ วิญฺญาณจริยา, ทุกฺขุทฺรเยหิ กมฺเมหิ วิปฺปยุตฺตา จรตีติ วิญฺญาณจริยา, สุขวิปาเกหิ กมฺเมหิ สมฺปยุตฺตา จรตีติ วิญฺญาณจริยา, ทุกฺขวิปาเกหิ กมฺเมหิ วิปฺปยุตฺตา จรตีติ วิญฺญาณจริยา, วิญฺญาเต จรตีติ วิญฺญาณจริยา, วิญฺญาณสฺส เอวรูปา จริยา โหตีติ วิญฺญาณจริยา, ปกติ\-ปริสุทฺธมิทํ จิตฺตํ นิกฺกิเลส\-ฏฺเฐนาติ วิญฺญาณจริยา. อยํ วิญฺญาณจริยา.}

\prose{กตมา อญฺญาณจริยา? มนาปิเยสุ รูเปสุ ราคสฺส ชวนตฺถาย อาวชฺชน\-กิริยา\-พฺยากตา วิญฺญาณจริยา, ราคสฺส ชวนา อญฺญาณจริยา. อมนาปิเยสุ รูเปสุ โทสสฺส ชวนตฺถาย อาวชฺชน\-กิริยา\-พฺยากตา วิญฺญาณจริยา, โทสสฺส ชวนา อญฺญาณจริยา. ตทุภเยน อสมเปกฺขนสฺมิํ วตฺถุสฺมิํ โมหสฺส ชวนตฺถาย อาวชฺชน\-กิริยา\-พฺยากตา วิญฺญาณจริยา, โมหสฺส ชวนา อญฺญาณจริยา, วินิพนฺธสฺส มานสฺส ชวนตฺถาย อาวชฺชน\-กิริยา\-พฺยากตา วิญฺญาณจริยา, มานสฺส ชวนา อญฺญาณจริยา. ปรามฏฺฐาย ทิฏฺฐิยา ชวนตฺถาย อาวชฺชน\-กิริยา\-พฺยากตา วิญฺญาณจริยา, ทิฏฺฐิยา ชวนา อญฺญาณจริยา. วิกฺเขป\-คตสฺส อุทฺธจฺจสฺส ชวนตฺถาย อาวชฺชน\-กิริยา\-พฺยากตา วิญฺญาณจริยา, อุทฺธจฺจสฺส ชวนา อญฺญาณจริยา. อนิฏฺฐงฺคตาย วิจิกิจฺฉาย ชวนตฺถาย อาวชฺชน\-กิริยา\-พฺยากตา วิญฺญาณจริยา, วิจิกิจฺฉาย ชวนา อญฺญาณจริยา. ถามคตสฺส อนุสยสฺส ชวนตฺถาย อาวชฺชน\-กิริยา\-พฺยากตา วิญฺญาณจริยา, อนุสยสฺส ชวนา อญฺญาณจริยา.}

\prose{\csromanfolio{79}มนาปิเยสุ สทฺเทสุ ฯเปฯ. มนาปิเยสุ คนฺเธสุ ฯเปฯ. มนาปิเยสุ รเสสุ ฯเปฯ. มนาปิเยสุ โผฏฺฐพฺเพสุ ฯเปฯ. มนาปิเยสุ ธมฺเมสุ ราคสฺส ชวนตฺถาย อาวชฺชน\-กิริยา\-พฺยากตา วิญฺญาณจริยา, ราคสฺส ชวนา อญฺญาณจริยา. อมนาปิเยสุ ธมฺเมสุ โทสสฺส ชวนตฺถาย อาวชฺชน\-กิริยา\-พฺยากตา วิญฺญาณจริยา, โทสสฺส ชวนา อญฺญาณจริยา. ตทุภเยน อสมเปกฺขนสฺมิํ วตฺถุสฺมิํ โมหสฺส ชวนตฺถาย อาวชฺชน\-กิริยา\-พฺยากตา วิญฺญาณจริยา, โมหสฺส ชวนา อญฺญาณจริยา. วินิพนฺธสฺส มานสฺส ชวนตฺถาย อาวชฺชน\-กิริยา\-พฺยากตา วิญฺญาณจริยา, มานสฺส ชวนา อญฺญาณจริยา. ปรามฏฺฐาย ทิฏฺฐิยา ชวนตฺถาย อาวชฺชน\-กิริยา\-พฺยากตา วิญฺญาณจริยา, ทิฏฺฐิยา ชวนา อญฺญาณจริยา. วิกฺเขป\-คตสฺส อุทฺธจฺจสฺส ชวนตฺถาย อาวชฺชน\-กิริยา\-พฺยากตา วิญฺญาณจริยา, อุทฺธจฺจสฺส ชวนา อญฺญาณจริยา. อนิฏฺฐงฺคตาย วิจิกิจฺฉาย ชวนตฺถาย อาวชฺชน\-กิริยา\-พฺยากตา วิญฺญาณจริยา, วิจิกิจฺฉาย ชวนา อญฺญาณจริยา. ถามคตสฺส อนุสยสฺส ชวนตฺถาย อาวชฺชน\-กิริยา\-พฺยากตา วิญฺญาณจริยา, อนุสยสฺส ชวนา อญฺญาณจริยา.}

\proseitem{70}{อญฺญาณจริยาติ เกนฏฺเฐน อญฺญาณจริยา? สราคา จรตีติ อญฺญาณจริยา, สโทสา จรตีติ อญฺญาณจริยา, สโมหา จรตีติ อญฺญาณจริยา, สมานา จรตีติ อญฺญาณจริยา, สทิฏฺฐิ จรตีติ อญฺญาณจริยา, สฺเอาทฺธจฺจา จรตีติ อญฺญาณจริยา, สวิจิกิจฺฉา จรตีติ อญฺญาณจริยา, สานุสยา จรตีติ อญฺญาณจริยา. ราคสมฺปยุตฺตา จรตีติ อญฺญาณจริยา, โทสสมฺปยุตฺตา จรตีติ อญฺญาณจริยา, โมหสมฺปยุตฺตา จรตีติ อญฺญาณจริยา, มานสมฺปยุตฺตา จรตีติ อญฺญาณจริยา, ทิฏฺฐิ\-สมฺปยุตฺตา จรตีติ อญฺญาณจริยา, อุทฺธจฺจ\-สมฺปยุตฺตา จรตีติ อญฺญาณจริยา. วิจิกิจฺฉา\-สมฺปยุตฺตา จรตีติ อญฺญาณจริยา, อนุสย\-สมฺปยุตฺตา จรตีติ อญฺญาณจริยา. กุสเลหิ กมฺเมหิ วิปฺปยุตฺตา จรตีติ อญฺญาณจริยา, อกุสเลหิ กมฺเมหิ สมฺปยุตฺตา จรตีติ อญฺญาณจริยา, สาวชฺเชหิ กมฺเมหิ สมฺปยุตฺตา จรตีติ อญฺญาณจริยา, อนวชฺเชหิ กมฺเมหิ วิปฺปยุตฺตา จรตีติ อญฺญาณจริยา, กเณฺหหิ กมฺเมหิ สมฺปยุตฺตา จรตีติ อญฺญาณจริยา, สุกฺเกหิ กมฺเมหิ วิปฺปยุตฺตา จรตีติ อญฺญาณจริยา, สุขุทฺรเยหิ กมฺเมหิ วิปฺปยุตฺตา จรตีติ อญฺญาณจริยา, ทุกฺขุทฺรเยหิ \csromanfolio{80}กมฺเมหิ สมฺปยุตฺตา จรตีติ อญฺญาณจริยา, สุขวิปาเกหิ กมฺเมหิ วิปฺปยุตฺตา จรตีติ อญฺญาณจริยา, ทุกฺขวิปาเกหิ กมฺเมหิ สมฺปยุตฺตา จรตีติ อญฺญาณจริยา, อญฺญาเต จรตีติ อญฺญาณจริยา, อญฺญาณสฺส เอวรูปา จริยา โหตีติ อญฺญาณจริยา. อยํ อญฺญาณจริยา.}

\proseitem{71}{กตมา ญาณจริยา? อนิจฺจานุปสฺสน\-ตฺถาย อาวชฺชน\-กิริยา\-พฺยากตา วิญฺญาณจริยา, อนิจฺจา\-นุปสฺสนา ญาณจริยา. ทุกฺขานุปสฺสน\-ตฺถาย อาวชฺชน\-กิริยา\-พฺยากตา วิญฺญาณจริยา, ทุกฺขา\-นุปสฺสนา ญาณจริยา. อนตฺตานุปสฺสน\-ตฺถาย อาวชฺชน\-กิริยา\-พฺยากตา วิญฺญาณจริยา, อนตฺตา\-นุปสฺสนา ญาณจริยา. นิพฺพิทานุปสฺสน\-ตฺถาย ฯเปฯ. วิราคานุปสฺสน\-ตฺถาย. นิโรธานุปสฺสน\-ตฺถาย. ปฏินิสฺสคฺคานุปสฺสน\-ตฺถาย. ขยานุปสฺสน\-ตฺถาย. วยา\-นุปสฺสน\-ตฺถาย. วิปริณามา\-นุปสฺสน\-ตฺถาย. อนิมิตฺตา\-นุปสฺสน\-ตฺถาย. อปฺปณิหิตา\-นุปสฺสน\-ตฺถาย. สุญฺญตานุปสฺสน\-ตฺถาย. อธิปญฺญา\-ธมฺมานุปสฺสน\-ตฺถาย. ยถา\-ภูต\-ญาณ\-ทสฺสน\-ตฺถาย. อาทีนวานุปสฺสน\-ตฺถาย. ปฏิสงฺขานุปสฺสน\-ตฺถาย. อาวชฺชน\-กิริยา\-พฺยากตา วิญฺญาณจริยา, ปฏิสงฺขา\-นุปสฺสนา ญาณจริยา. วิวฏฺฏนา\-นุปสฺสนา ญาณจริยา. โสตาปตฺติมคฺโค ญาณจริยา. โสตาปตฺติ\-ผลสมาปตฺติ ญาณจริยา. สกทาคามิ\-มคฺโค ญาณจริยา. สกทาคามิ\-ผลสมาปตฺติ ญาณจริยา. อนาคามิมคฺโค ญาณจริยา. อนาคามิ\-ผลสมาปตฺติ ญาณจริยา. อรหตฺตมคฺโค ญาณจริยา. อรหตฺตผล\-สมาปตฺติ ญาณจริยา.}

\prosewithcloserruled{ญาณจริยาติ เกนฏฺเฐน ญาณจริยา? นีราคา จรตีติ ญาณจริยา นิทฺโทสา จรตีติ ญาณจริยา ฯเปฯ นานุสยา จรตีติ ญาณจริยา, ราควิปฺปยุตฺตา จรตีติ ญาณจริยา, โทสวิปฺปยุตฺตา จรตีติ ญาณจริยา, โมหวิปฺปยุตฺตา จรตีติ ญาณจริยา, มานวิปฺปยุตฺตา ฯเปฯ ทิฏฺฐิ\-วิปฺปยุตฺตา. อุทฺธจฺจ\-วิปฺปยุตฺตา. วิจิกิจฺฉา\-วิปฺปยุตฺตา. อนุสย\-วิปฺปยุตฺตา. กุสเลหิ กมฺเมหิ สมฺปยุตฺตา. อกุสเลหิ กมฺเมหิ วิปฺปยุตฺตา. สาวชฺเชหิ กมฺเมหิ วิปฺปยุตฺตา. อนวชฺเชหิ กมฺเมหิ สมฺปยุตฺตา. กเณฺหหิ กมฺเมหิ วิปฺปยุตฺตา. สุกฺเกหิ กมฺเมหิ สมฺปยุตฺตา. สุขุทฺรเยหิ กมฺเมหิ สมฺปยุตฺตา. ทุกฺขุทฺรเยหิ กมฺเมหิ วิปฺปยุตฺตา. สุขวิปาเกหิ กมฺเมหิ สมฺปยุตฺตา จรตีติ ญาณจริยา, ทุกฺขวิปาเกหิ กมฺเมหิ วิปฺปยุตฺตา \csromanfolio{81}จรตีติ ญาณจริยา, ญาเต จรตีติ ญาณจริยา, ญาณสฺส เอวรูปา จริยา โหตีติ ญาณจริยา, อยํ ญาณจริยา, อญฺญา วิญฺญาณจริยา, อญฺญา อญฺญาณจริยา, อญฺญา ญาณจริยาติ. ตํญาตฏฺเฐน ญาณํ, ปชานนฏฺเฐน ปญฺญา. เตน วุจฺจติ จริยา\-ววตฺถาเน ปญฺญา จริยานานตฺเต ญาณํ.}{จริยา\-นานตฺต\-ญาณ\-นิทฺเทโส สตฺตรสโม.}
\csromanmatikaanchor{mtk.20}
\csromantocmarkref{subsection}{18. ภูมินานตฺตญาณนิทฺเทส}{mtk.20}
\bookmark[dest={mtk.20},level=2]{18. ภูมินานตฺตญาณนิทฺเทส}
\csromanheader{2}{18. ภูมินานตฺต\-ญาณ\-นิทฺเทส}
\proseitem{72}{กถํ จตุธมฺม\-ววตฺถาเน ปญฺญา ภูมินานตฺเต ญาณํ\csromandash{} จตสฺโส ภูมิโย กามาวจรา ภูมิ, รูปาวจรา ภูมิ, อรูปาวจรา ภูมิ, อปริยาปนฺนา ภูมิ. กตมา กามาวจรา ภูมิ? เหฏฺฐโต อวีจินิรยํ ปริยนฺตํ กริตฺวา อุปริโต ปรนิมฺมิต\-วส\-วตฺตี เทเว\csromansharedfootnote{eff10c429f37797f}{ปรนิมฺมิตวสวตฺติเทเว (ก)} อนฺโต กริตฺวา ยํ เอตสฺมิํ อนฺตเร เอตฺถาวจรา เอตฺถ ปริยาปนฺนา ขนฺธ\-ธาตุ\-อายตนา รูปํ เวทนา สญฺญา สงฺขารา วิญฺญาณํ. อยํ กามาวจรา ภูมิ.}

\prose{กตมา รูปาวจรา ภูมิ? เหฏฺฐโต พฺรหฺมโลกํ ปริยนฺตํ กริตฺวา อุปริโต อกนิฏฺเฐ เทเว อนฺโต กริตฺวา ยํ เอตสฺมิํ อนฺตเร เอตฺถาวจรา เอตฺถ ปริยาปนฺนา สมาปนฺนสฺส วา อุปปนฺนสฺส วา ทิฏฺฐธมฺม\-สุขวิหาริสฺส วา จิตฺตเจตสิกา ธมฺมา. อยํ รูปาวจรา ภูมิ.}

\prose{กตมา อรูปาวจรา ภูมิ? เหฏฺฐโต อากาสานญฺจายตนู\-ปเค เทเว ปริยนฺตํ กริตฺวา อุปริโต เนว\-สญฺญา\-นา\-สญฺญา\-ยตนู\-ปเค เทเว อนฺโต กริตฺวา ยํ เอตสฺมิํ อนฺตเร เอตฺถาวจรา เอตฺถ ปริยาปนฺนา สมาปนฺนสฺส วา อุปปนฺนสฺส วา ทิฏฺฐธมฺม\-สุขวิหาริสฺส วา จิตฺตเจตสิกา ธมฺมา. อยํ อรูปาวจรา ภูมิ.}

\prose{กตมา อปริยาปนฺนา ภูมิ? อปริยาปนฺนา มคฺคา จ มคฺคผลานิ จ อสงฺขตา จ ธาตุ. อยํ อปริยาปนฺนา ภูมิ. อิมา จตสฺโส ภูมิโย.}

\prosewithcloserruled{อปราปิ จตสฺโส ภูมิโย จตฺตาโร สติปฏฺฐานา, จตฺตาโร สมฺมปฺปธานา, จตฺตาโร อิทฺธิปาทา, จตฺตาริ ฌานานิ, จตสฺโส อปฺปมญฺญาโย, จตสฺโส อรูป\-สมาปตฺติโย, จตสฺโส ปฏิสมฺภิทา, \csromanfolio{82}จตสฺโส ปฏิปทา, จตฺตาริ อารมฺมณานิ, จตฺตาโร อริยวํสา, จตฺตาริ สงฺคห\-วตฺถูนิ, จตฺตาริ จกฺกานิ, จตฺตาริ ธมฺมปทานิ. อิมา จตสฺโส ภูมิโย. ตํญาตฏฺเฐน ญาณํ, ปชานนฏฺเฐน ปญฺญา. เตน วุจฺจติ จตุธมฺม\-ววตฺถาเน ปญฺญา ภูมินานตฺเต ญาณํ.}{ภูมินานตฺต\-ญาณ\-นิทฺเทโส อฏฺฐารสโม.}
\csromanmatikaanchor{mtk.21}
\csromantocmarkref{subsection}{19. ธมฺมนานตฺตญาณนิทฺเทส}{mtk.21}
\bookmark[dest={mtk.21},level=2]{19. ธมฺมนานตฺตญาณนิทฺเทส}
\csromanheader{2}{19. ธมฺม\-นานตฺต\-ญาณ\-นิทฺเทส}
\proseitem{73}{กถํ นวธมฺม\-ววตฺถาเน ปญฺญา ธมฺมนานตฺเต ญาณํ\csromandash{} กถํ ธมฺเม ววตฺเถติ? กามาวจเร ธมฺเม กุสลโต ววตฺเถติ, อกุสลโต ววตฺเถติ, อพฺยากตโต ววตฺเถติ. รูปาวจเร ธมฺเม กุสลโต ววตฺเถติ, อพฺยากตโต ววตฺเถติ. อรูปาวจเร ธมฺเม กุสลโต ววตฺเถติ, อพฺยากตโต ววตฺเถติ. อปริยาปนฺเน ธมฺเม กุสลโต ววตฺเถติ, อพฺยากตโต ววตฺเถติ.}

\prose{กถํ กามาวจเร ธมฺเม กุสลโต ววตฺเถติ, อกุสลโต ววตฺเถติ, อพฺยากตโต ววตฺเถติ? ทส กุสล\-กมฺม\-ปเถ กุสลโต ววตฺเถติ. ทส อกุสล\-กมฺมปเถ อกุสลโต ววตฺเถติ, รูปญฺจ วิปากญฺจ กิริยญฺจ อพฺยากตโต ววตฺเถติ. เอวํ กามาวจเร ธมฺเม กุสลโต ววตฺเถติ อกุสลโต ววตฺเถติ อพฺยากตโต ววตฺเถติ.}

\prose{กถํ รูปาวจเร ธมฺเม กุสลโต ววตฺเถติ, อพฺยากตโต ววตฺเถติ? อิธฏฺฐสฺส จตฺตาริ ฌานานิ กุสลโต ววตฺเถติ, ตตฺรู\-ปปนฺนสฺส จตฺตาริ ฌานานิ อพฺยากตโต ววตฺเถติ. เอวํ รูปาวจเร ธมฺเม กุสลโต ววตฺเถติ, อพฺยากตโต ววเถติ.}

\prose{กถํ อรูปาวจเร ธมฺเม กุสลโต ววตฺเถติ, อพฺยากตโต ววตฺเถติ? อิธฏฺฐสฺส จตสฺโส อรูปาวจร\-สมาปตฺติโย กุสลโต ววตฺเถติ, ตตฺรู\-ปปนฺนสฺส จตสฺโส อรูปาวจร\-สมาปตฺติโย อพฺยากตโต ววตฺเถติ. เอวํ อรูปาวจเร ธมฺเม กุสลโต ววตฺเถติ อพฺยากตโต ววตฺเถติ.}

\prose{\csromanfolio{83}กถํ อปริยาปนฺเน ธมฺเม กุสลโต ววตฺเถติ, อพฺยากตโต ววตฺเถติ? จตฺตาโร อริยมคฺเค กุสลโต ววตฺเถติ, จตฺตาริ จ สามญฺญผลานิ นิพฺพานญฺจ อพฺยากตโต ววตฺเถติ. เอวํ อปริยาปนฺเน ธมฺเม กุสลโต ววตฺเตติ อพฺยากตโต ววตฺเถติ. เอวํ ธมฺเม ววตฺเถติ.}

\prose{นว ปาโมชฺชมูลกา ธมฺมา, อนิจฺจโต มนสิกโรโต ปาโมชฺชํ ชายติ, ปมุทิตสฺส ปีติ ชายติ, ปีติมนสฺส กาโย ปสฺสมฺภติ, ปสฺสทฺธกาโย ทุขํ เวเทติ, สุขิโน จิตฺตํ สมาธิยติ, สมาหิเต จิตฺเต ยถาภูตํ ปชานาติ ปสฺสติ, ยถาภูตํ ชานํ ปสฺสํ นิพฺพินฺทติ. นิพฺพินฺทํ วิรชฺชติ, วิราคา วิมุจฺจติ. ทุกฺขโต มนสิกโรโต ปาโมชฺชํ ชายติ ฯเปฯ. อนตฺตโต มนสิกโรโต ปาโมชฺชํ ชายติ ฯเปฯ.}

\prose{รูปํ อนิจฺจโต มนสิกโรโต ปาโมชฺชํ ชายติ ฯเปฯ. รูปํ ทุกฺขโต มนสิกโรโต ฯเปฯ. เวทนํ. สญฺญํ. สงฺขาเร. วิญฺญาณํ. จกฺขุํ ฯเปฯ. ชรามรณํ อนิจฺจโต มนสิกโรโต ปาโมชฺชํ ชายติ ฯเปฯ. ชรามรณํ ทุกฺขโต มนสิกโรโต ปาโมชฺชํ ชายติ ฯเปฯ. ชรามรณํ อนตฺตโต มนสิกโรโต ปาโมชฺชํ ชายติ, ปมุทิตสฺส ปีติ ชายติ, ปีติมนสฺส กาโย ปสฺสมฺภติ, ปสฺสทฺธกาโย สุขํ เวเทติ, สุขิโน จิตฺตํ สมาธิยติ, สมาหิเต จิตฺเต ยถาภูตํ ปชานาติ ปสฺสติ, ยถาภูตํ ชานํ ปสฺสํ นิพฺพินฺทติ, นิพฺพินฺทํ วิรชฺชติ, วิราคา วิมุจฺจติ, อิเม นว ปาโมชฺชมูลกา ธมฺมา.}

\proseitem{74}{นว โยนิโสมนสิการ\-มูลกา ธมฺมา, อนิจฺจโต โยนิโส มนสิกโรโต ปาโมชฺชํ ชายติ, ปมุทิตสฺส ปีติ ชายติ, ปีติมนสฺส กาโย ปสฺสมฺภติ, ปสฺสทฺธกาโย สุขํ เวเทติ, สุขิโน จิตฺตํ สมาธิยติ, สมาหิเตน จิตฺเตน “อิทํ ทุกฺขนฺ”ติ ยถาภูตํ ปชานาติ, “อยํ ทุกฺขสมุทโย”ติ ยถาภูตํ ปชานาติ, “อยํ ทุกฺขนิโรโธ”ติ ยยาภูตํ ปชานาติ, “อยํ ทุกฺขนิโรธ\-คามินี ปฏิปทา”ติ ยถาภูตํ ปชานาติ.  ทุกฺขโต โยนิโส มนสิกโรโต ปาโมชฺชํ ชายติ, ปมุทิตสฺส ปีติ ชายติ, ปีติมนสฺส กาโย ปสฺสมฺภติ, ปสฺสทฺธกาโย สุขํ เวเทติ, สุขิโน จิตฺตํ สมาธิยติ, สมาหิเตน จิตฺเตน “อิทํ ทุกฺขนฺ”ติ ยถาภูตํ ปชานาติ, “อยํ ทุกฺขสมุทโย”ติ ยถาภูตํ ปชานาติ, “อยํ ทุกฺขนิโรโธ”ติ ยถาภูตํ ปชานาติ, \csromanfolio{84}“อยํ ทุกฺขนิโรธ\-คามินี ปฏิปทา”ติ ยถาภูตํ ปชานาติ. อนตฺตโต โยนิโส มนสิกโรโต ปาโมชฺชํ ชายติ ฯเปฯ.}

\prose{รูปํ อนิจฺจโต โยนิโส มนสิกโรโต ปาโมชฺชํ ชายติ ฯเปฯ. รูปํ ทุกฺขโต โยนิโส มนสิกโรโต ปาโมชฺชํ ชายติ ฯเปฯ. รูปํ อนตฺตโต โยนิโส มนสิกโรโต ปาโมชฺชํ ชายติ ฯเปฯ. เวทนํ. สญฺญํ. สงฺขาเร. วิญฺญาณํ. จกฺขุํ ฯเปฯ. ชรามรณํ อนิจฺจโต โยนิโส มนสิกโรโต ปาโมชฺชํ ชายติ ฯเปฯ. ชรามรณํ ทุกฺขโต โยนิโส มนสิกโรโต ปาโมชฺชํ ชายติ ฯเปฯ. ชรามรณํ อนตฺตโต โยนิโส มนสิกโรโต ปาโมชฺชํ ชายติ, ปมุทิตสฺส ปีติ ชายติ, ปีติมนสฺส กาโย ปสฺสมฺภติ, ปสฺสทฺธกาโย สุขํ เวเทติ, สุขิโน จิตฺตํ สมาธิยติ, สมาหิเตน จิตฺเตน “อิทํ ทุกฺขนฺ”ติ ยถาภูตํ ปชานาติ, “อยํ ทุกฺขสมุทโย”ติ ยถาภูตํ ปชานาติ, “อยํ ทุกฺขนิโรโธ”ติ ยถาภูตํ ปชานาติ, “อยํ ทุกฺขนิโรธ\-คามินี ปฏิปทา”ติ ยถาภูตํ ปชานาติ. อิเม นว โยนิโสมนสิการ\-มูลกา ธมฺมา.}

\prosewithcloserruled{นว นานตฺตา, ธาตุนานตฺตํ ปฏิจฺจ อุปฺปชฺชติ ผสฺสนานตฺตํ, ผสฺสนานตฺตํ ปฏิจฺจ อุปฺปชฺชติ เวทนานานตฺตํ, เวทนานานตฺตํ ปฏิจฺจ อุปฺปชฺชติ สญฺญานานตฺตํ, สญฺญานานตฺตํ ปฏิจฺจ อุปฺปชฺชติ สงฺกปฺป\-นานตฺตํ, สงฺกปฺป\-นานตฺตํ ปฏิจฺจ อุปฺปชฺชติ ฉนฺทนานตฺตํ, ฉนฺทนานตฺตํ ปฏิจฺจ อุปฺปชฺชติ ปริฬาห\-นานตฺตํ, ปริฬาห\-นานตฺตํ ปฏิจฺจ อุปฺปชฺชติ ปริเยสนา\-นานตฺตํ, ปริเยสนา\-นานตฺตํ ปฏิจฺจ อุปฺปชฺชติ ลาภนานตฺตํ. อิเม นว นานตฺตา. ตํญาตฏฺเฐน ญาณํ, ปชานนฏฺเฐน ปญฺญา. เตน วุจฺจติ นวธมฺม\-ววตฺถาเน ปญฺญา ธมฺมนานตฺเต ญาณํ.}{ธมฺม\-นานตฺต\-ญาณ\-นิทฺเทโส เอกูนวีสติโม.}
\csromanmatikaanchor{mtk.22}
\csromantocmarkref{subsection}{20-24. ญาณปญฺจกนิทฺเทส}{mtk.22}
\bookmark[dest={mtk.22},level=2]{20-24. ญาณปญฺจกนิทฺเทส}
\csromanheader{2}{20-24. ญาณ\-ปญฺจก\-นิทฺเทส}
\proseitemwithcloserruled{75}{กถํ อภิญฺญาปญฺญา ญาตฏฺเฐ ญาณํ, ปริญฺญาปญฺญา ตีรณฏฺเฐ ญาณํ, ปหาเน ปญฺญา ปริจฺจาคฏฺเฐ ญาณํ, ภาวนาปญฺญา เอกรสฏฺเฐ ญาณํ, สจฺฉิกิริยา\-ปญฺญา ผสฺสนฏฺเฐ ญาณํ\csromandash{}เย เย ธมฺมา อภิญฺญาตา โหนฺติ, เต เต ธมฺมา ญาตา โหนฺติ. เย เย ธมฺมา ปริญฺญาตา โหนฺติ, \csromanfolio{85}เต เต ธมฺมา ตีริตา โหนฺติ. เย เย ธมฺมา ปหีนา โหนฺติ, เต เต ธมฺมา ปริจฺจตฺตา โหนฺติ. เย เย ธมฺมา ภาวิตา โหนฺติ, เต เต ธมฺมา เอกรสา โหนฺติ. เย เย ธมฺมา สจฺฉิกตา โหนฺติ, เต เต ธมฺมา ผสฺสิตา โหนฺติ. ตํญาตฏฺเฐน ญาณํ, ปชานนฏฺเฐน ปญฺญา. เตน วุจฺจติ อภิญฺญาปญฺญา ญาตฏฺเฐ ญาณํ, ปริญฺญาปญฺญา ตีรณฏฺเฐ ญาณํ. ปหาเน ปญฺญา ปริจฺจาคฏฺเฐ ญาณํ. ภาวนาปญฺญา เอกรสฏฺเฐ ญาณํ. สจฺฉิกิริยา\-ปญฺญา ผสฺสนฏฺเฐ ญาณํ.}{ญาณ\-ปญฺจก\-นิทฺเทโส จตุวีสติโม.}
\csromanmatikaanchor{mtk.23}
\csromantocmarkref{subsection}{25-28. ปฏิสมฺภิทาญาณนิทฺเทส}{mtk.23}
\bookmark[dest={mtk.23},level=2]{25-28. ปฏิสมฺภิทาญาณนิทฺเทส}
\csromanheader{2}{25-28. ปฏิสมฺภิทา\-ญาณ\-นิทฺเทส}
\proseitem{76}{กถํ อตฺถนานตฺเต ปญฺญา อตฺถ\-ปฏิสมฺภิเท ญาณํ, ธมฺมนานตฺเต ปญฺญา ธมฺม\-ปฏิสมฺภิเท ญาณํ, นิรุตฺตินานตฺเต ปญฺญา นิรุตฺติ\-ปฏิสมฺภิเท ญาณํ, ปฏิภาน\-นานตฺเต ปญฺญา ปฏิภาน\-ปฏิสมฺภิเท ญาณํ\csromandash{}สทฺธินฺทฺริยํ ธมฺโม, วีริยินฺทฺริยํ ธมฺโม, สตินฺทฺริยํ ธมฺโม, สมาธินฺทฺริยํ ธมฺโม, ปญฺญินฺทฺริยํ ธมฺโม. อญฺโญ สทฺธินฺทฺริยํ ธมฺโม, อญฺโญ วีริยินฺทฺริยํ ธมฺโม, อญฺโญ สตินฺทฺริยํ ธมฺโม, อญฺโญ สมาธินฺทฺริยํ ธมฺโม, อญฺโญ ปญฺญินฺทฺริยํ ธมฺโม. เยน ญาเณน อิเม นานา ธมฺมา ญาตา, เตเนว ญาเณน อิเม นานา ธมฺมา ปฏิวิทิตาติ. เตน วุจฺจติ ธมฺมนานตฺเต ปญฺญา ธมฺม\-ปฏิสมฺภิเท ญาณํ. (1)}

\prose{อธิโมกฺขฏฺโฐ อตฺโถ, ปคฺคหฏฺโฐ อตฺโถ, อุปฏฺฐานฏฺโฐ อตฺโถ, อวิกฺเขปฏฺโฐ อตฺโถ, ทสฺสนฏฺโฐ อตฺโถ. อญฺโญ อธิโมกฺขฏฺโฐ อตฺโถ, อญฺโญ ปคฺคหฏฺโฐ อตฺโถ, อญฺโญ อุปฏฺฐานฏฺโฐ อตฺโถ, อญฺโญ อวิกฺเขปฏฺโฐ อตฺโถ, อญฺโญ ทสฺสนฏฺโฐ อตฺโถ. เยน ญาเณน อิเม นานา อตฺถา ญาตา, เตเนว ญาเณน อิเม นานา อตฺถา ปฏิวิทิตาติ. เตน วุจฺจติ อตฺถนานตฺเถ ปญฺญา อตฺถ\-ปฏิสมฺภิเท ญาณํ. (2)}

\prose{ปญฺจ ธมฺเม สนฺทสฺเสตุํ พฺยญฺชน\-นิรุตฺตา\-ภิลาปา, ปญฺจ อตฺเถ สนฺทสฺเสตุํ พฺยญฺชน\-นิรุตฺตา\-ภิลาปา. อญฺญา ธมฺม\-นิรุตฺติโย, อญฺญา อตฺถนิรุตฺติโย. เยน ญาเณน อิมา นานา นิรุตฺติโย ญาตา, เตเนว ญาเณน อิมา นานา นิรุตฺติโย ปฏิวิทิตาติ. เตน วุจฺจติ นิรุตฺตินานตฺเต ปญฺญา นิรุตฺติ\-ปฏิสมฺภิเท ญาณํ. (3)}

\prose{\csromanfolio{86}ปญฺจสุ ธมฺเมสุ ญาณานิ, ปญฺจสุ อตฺเถสุ ญาณานิ, ทสสุ นิรุตฺตีสุ ญาณานิ. อญฺญานิ ธมฺเมสุ ญาณานิ, อญฺญานิ อตฺเถสุ ญาณานิ, อญฺญานิ นิรุตฺตีสุ ญาณานิ. เยน ญาเณน อิเม นานา ญาณา ญาตา, เตเนว ญาเณน อิเม นานา ญาณา ปฏิวิทิตาติ. เตน วุจฺจติ ปฏิภาน\-นานตฺเต ปญฺญา ปฏิภาน\-ปฏิสมฺภิเท ญาณํ. (4)}

\proseitem{77}{สทฺธาพลํ ธมฺโม, วีริยพลํ ธมฺโม, สติพลํ ธมฺโม, สมาธิพลํ ธมฺโม, ปญฺญาพลํ ธมฺโม. อญฺโญ สทฺธาพลํ ธมฺโม, อญฺโญ วีริยพลํ ธมฺโม, อญฺโญ สติพลํ ธมฺโม, อญฺโญ สมาธิพลํ ธมฺโม, อญฺโญ ปญฺญาพลํ ธมฺโม. เยน ญาเณน อิเม นานา ธมฺมา ญาตา, เตเนว ญาเณน อิเม นานา ธมฺมา ปฏิวิทิตาติ. เตน วุจฺจติ ธมฺมนานตฺเต ปญฺญา ธมฺม\-ปฏิสมฺภิเท ญาณํ. (1)}

\prose{อสฺสทฺธิเย อกมฺปิยฏฺโฐ อตฺโถ, โกสชฺเช อกมฺปิยฏฺโฐ อตฺโถ, ปมาเท อกมฺปิยฏฺโฐ อตฺโถ, อุทฺธจฺเจ อกมฺปิยฏฺโฐ อตฺโถ, อวิชฺชาย อกมฺปิยฏฺโฐ อตฺโถ. อญฺโญ อสฺสทฺธิเย อกมฺปิยฏฺโฐ อตฺโถ, อญฺโญ โกสชฺเช อกมฺปิยฏฺโฐ อตฺโถ, อญฺโญ ปมาเท อกมฺปิยฏฺโฐ อตฺโถ, อญฺโญ อุทฺธจฺเจ อกมฺปิยฏฺโฐ อตฺโถ, อญฺโญ อวิชฺชาย อกมฺปิยฏฺโฐ อตฺโถ. เยน ญาเณน อิเม นานา อตฺถา ญาตา, เตเนว ญาเณน อิเม นานา อตฺถา ปฏิวิทิตาติ, เตน วุจฺจติ อตฺถนานตฺเต ปญฺญา อตฺถ\-ปฏิสมฺภิเท ญาณํ. (2)}

\prose{ปญฺจ ธมฺเม สนฺทสฺเสตุํ พฺยญฺชน\-นิรุตฺตา\-ภิลาปา, ปญฺจ อตฺเถ สนฺทสฺเสตุํ พฺยญฺชน\-นิรุตฺตา\-ภิลาปา. อญฺญา ธมฺม\-นิรุตฺติโย, อญฺญา อตฺถนิรุตฺติโย. เยน ญาเณน อิมา นานา นิรุตฺติโย ญาตา, เตเนว ญาเณน อิมา นานา นิรุตฺติโย ปฏิวิทิตาติ, เตน วุจฺจติ นิรุตฺตินานตฺเต ปญฺญา นิรุตฺติ\-ปฏิสมฺภิเท ญาณํ. (3)}

\prose{ปญฺจสุ\csromansharedfootnote{c78ddf09a78056d2}{ปญฺจ (?)} ธมฺเมสุ ญาณานิ, ปญฺจสุ\csromansharedfootnote{819819c26f1d0fbf}{ทส (?)} อตฺเถสุ ญาณานิ, ทสสุ\csromansharedfootnote{c78ddf09a78056d2}{ปญฺจ (?)} นิรุตฺตีสุ ญาณานิ. อญฺญานิ ธมฺเมสุ ญาณานิ, อญฺญานิ อตฺเถสุ ญาณานิ, อญฺญานิ นิรุตฺตีสุ ญาณานิ. เยน ญาเณน อิเม นานา ญาณา ญาตา, เตเนว ญาเณน อิเม นานา ญาณา ปฏิวิทิตาติ, เตน วุจฺจติ ปฏิภาน\-นานตฺเต ปญฺญา ปฏิภาน\-ปฏิสมฺภิเท ญาณํ. (4)}

\prose{\csromanfolio{87}สติสมฺโพชฺฌงฺโค ธมฺโม, ธมฺมวิจย\-สมฺโพชฺฌงฺโค ธมฺโม, วีริย\-สมฺโพชฺฌงฺโค ธมฺโม, ปีติสมฺโพชฺฌงฺโค ธมฺโม, ปสฺสทฺธิ\-สมฺโพชฺฌงฺโค ธมฺโม, สมาธิ\-สมฺโพชฺฌงฺโค ธมฺโม, อุเปกฺขา\-สมฺโพชฺฌงฺโค ธมฺโม. อญฺโญ สติสมฺโพชฺฌงฺโค ธมฺโม, อญฺโญ ธมฺมวิจย\-สมฺโพชฺฌงฺโค ธมฺโม, อญฺโญ วีริย\-สมฺโพชฺฌงฺโค ธมฺโม, อญฺโญ ปีติสมฺโพชฺฌงฺโค ธมฺโม, อญฺโญ ปสฺสทฺธิ\-สมฺโพชฺฌงฺโค ธมฺโม, อญฺโญ สมาธิ\-สมฺโพชฺฌงฺโค ธมฺโม, อญฺโญ อุเปกฺขา\-สมฺโพชฺฌงฺโค ธมฺโม, เยน ญาเณน อิเม นานา ธมฺมา ญาตา, เตเนว ญาเณน อิเม นานา ธมฺมา ปฏิวิทิตาติ, เตน วุจฺจติ ธมฺมนานตฺเต ปญฺญา ธมฺม\-ปฏิสมฺภิเท ญาณํ. (1)}

\prose{อุปฏฺฐานฏฺโฐ อตฺโถ, ปวิจยฏฺโฐ อตฺโถ, ปคฺคหฏฺโฐ อตฺโถ, ผรณฏฺโฐ อตฺโถ, อุปสมฏฺโฐ อตฺโถ, อวิกฺเขปฏฺโฐ อตฺโถ, ปฏิสงฺขา\-นฏฺโฐ อตฺโถ. อญฺโญ อุปฏฺฐานฏฺโฐ อตฺโถ, อญฺโญ ปวิจยฏฺโฐ อตฺโถ, อญฺโญ ปคฺคหฏฺโฐ อตฺโต, อญฺโญ ผรณฏฺโฐ อตฺโถ, อญฺโญ อุปสมฏฺโฐ อตฺโถ, อญฺโญ อวิกฺเขปฏฺโฐ อตฺโถ, อญฺโญ ปฏิสงฺขา\-นฏฺโฐ อตฺโถ. เยน ญาเณน อิเม นานา อตฺถา ญาตา. เตเนว ญาเณน อิเม นานา อตฺถา ปฏิวิทิตาติ. เตน วุจฺจติ อตฺถนานตฺเต ปญฺญา อตฺถ\-ปฏิสมฺภิเท ญาณํ. (2)}

\prose{สตฺต ธมฺเม สนฺทสฺเสตุํ พฺยญฺชน\-นิรุตฺตา\-ภิลาปา, สตฺต อตฺเถ สนฺทสฺเสตุํ พฺยญฺชน\-นิรุตฺตา\-ภิลาปา. อญฺญา ธมฺม\-นิรุตฺติโย, อญฺญา อตฺถนิรุตฺติโย. เยน ญาเณน อิมา นานา นิรุตฺติโย ญาตา, เตเนว ญาเณน อิมา นานา นิรุตฺติโย ปฏิวิทิตาติ, เตน วุจฺจติ นิรุตฺตินานตฺเต ปญฺญา นิรุตฺติ\-ปฏิสมฺภิเท ญาณํ. (3)}

\prose{สตฺตสุ ธมฺเมสุ ญาณานิ, สตฺตสุ อตฺเถสุ ญาณานิ, จุทฺทสสุ นิรุตฺตีสุ ญาณานิ. อญฺญานิ ธมฺเมสุ ญาณานิ, อญฺญานิ อตฺเถสุ ญาณานิ, อญฺญานิ นิรุตฺตีสุ ญาณานิ, เยน ญาเณน อิเม นานา ญาณา ญาตา, เตเนว ญาเณน อิเม นานา ญาณา ปฏิวิทิตาติ, เตน วุจฺจติ ปฏิภาน\-นานตฺเต ปญฺญา ปฏิภาน\-ปฏิสมฺภิเท ญาณํ. (4)}

\prose{สมฺมาทิฏฺฐิ ธมฺโม, สมฺมาสงฺกปฺโป ธมฺโม, สมฺมาวาจา ธมฺโม, สมฺมากมฺมนฺโต ธมฺโม, สมฺมาอาชีโว ธมฺโม, สมฺมาวายาโม ธมฺโม, สมฺมาสติ ธมฺโม, สมฺมาสมาธิ ธมฺโม. อญฺโญ สมฺมาทิฏฺฐิ ธมฺโม, อญฺโญ \csromanfolio{88}สมฺมาสงฺกปฺโป ธมฺโม, อญฺโญ สมฺมาวาจา ธมฺโม, อญฺโญ สมฺมากมฺมนฺโต ธมฺโม, อญฺโญ สมฺมาอาชีโว ธมฺโม, อญฺโญ สมฺมาวายาโม ธมฺโม, อญฺโญ สมฺมาสติ ธมฺโม, อญฺโญ สมฺมาสมาธิ ธมฺโม. เยน ญาเณน อิเม นานา ธมฺมา ญาตา, เตเนว ญาเณน อิเม นานา ธมฺมา ปฏิวิทิตาติ, เตน วุจฺจติ ธมฺมนานตฺเต ปญฺญา ธมฺม\-ปฏิสมฺภิเท ญาณํ.}

\prose{ทสฺสนฏฺโฐ อตฺโถ, อภินิโรปน\-ฏฺโฐ อตฺโถ, ปริคฺคหฏฺโฐ อตฺโถ, สมุฏฺฐานฏฺโฐ อตฺโถ, โวทานฏฺโฐ อตฺโถ, ปคฺคหฏฺโฐ อตฺโถ, อุปฏฺฐานฏฺโฐ อตฺโถ, อวิกฺเขปฏฺโฐ อตฺโถ. อญฺโญ ทสฺสนฏฺโฐ อตฺโถ, อญฺโญ อภินิโรปน\-ฏฺโฐ อตฺโถ, อญฺโญ ปริคฺคหฏฺโฐ อตฺโถ, อญฺโญ สมุฏฺฐานฏฺโฐ อตฺโถ, อญฺโญ โวทานฏฺโฐ อตฺโถ, อญฺโญ ปคฺคหฏฺโฐ อตฺโถ, อญฺโญ อุปฏฺฐานฏฺโฐ อตฺโถ, อญฺโญ อวิกฺเขปฏฺโฐ อตฺโถ. เยน ญาเณน อิเม นานา อตฺถา ญาตา, เตเนว ญาเณน อิเม นานา อตฺถา ปฏิวิทิตาติ, เตน วุจฺจติ อตฺถนานตฺเต ปญฺญา อตฺถ\-ปฏิสมฺภิเท ญาณํ.}

\prose{อฏฺฐ ธมฺเม สนฺทสฺเสตุํ พฺยญฺชน\-นิรุตฺตา\-ภิลาปา, อฏฺฐ อตฺเถ สนฺทสฺเสตุํ พฺยญฺชน\-นิรุตฺตา\-ภิลาปา. อญฺญา ธมฺม\-นิรุตฺติโย, อญฺญา อตฺถนิรุตฺติโย เยน ญาเณน อิมา นานา นิรุตฺติโย ญาตา, เตเนว ญาเณน อิมา นานา นิรุตฺติโย ปฏิวิทิตาติ, เตน วุจฺจติ นิรุตฺตินานตฺเต ปญฺญา นิรุตฺติ\-ปฏิสมฺภิเท ญาณํ.}

\prose{อฏฺฐสุ ธมฺเมสุ ญาณานิ, อฏฺฐสุ อตฺเถสุ ญาณานิ, โสฬสสุ นิรุตฺตีสุ ญาณานิ. อญฺญานิ ธมฺเมสุ ญาณานิ, อญฺญานิ อตฺเถสุ ญาณานิ อญฺญานิ นิรุตฺตีสุ ญาณานิ. เยน ญาเณน อิเม นานา ญาณา ญาตา, เตเนว ญาเณน อิเม นานา ญาณา ปฏิวิทิตาติ, เตน วุจฺจติ ปฏิภาน\-นานตฺเต ปญฺญา ปฏิภาน\-ปฏิสมฺภิเท ญาณํ. ตํญาตฏฺเฐน ญาณํ, ปชานนฏฺเฐน ปญฺญา, เตน วุจฺจติ อตฺถนานตฺเต ปญฺญา อตฺถ\-ปฏิสมฺภิเท ญาณํ, ธมฺมนานตฺเต ปญฺญา ธมฺม\-ปฏิสมฺภิเท ญาณํ, นิรุตฺตินานตฺเต ปญฺญา นิรุตฺติ\-ปฏิสมฺภิเท ญาณํ, ปฏิภาน\-นานตฺเต ปญฺญา ปฏิภาน\-ปฏิสมฺภิเท ญาณํ.}

\vspace{\tipitakaparskip}
\csromancenter{ปฏิสมฺภิทา\-ญาณ\-นิทฺเทโส อฏฺฐวีสติโม.}
\csromanmatikaanchor{mtk.24}
\csromantocmarkref{subsection}{29-31. ญาณตฺตยนิทฺเทส}{mtk.24}
\bookmark[dest={mtk.24},level=2]{29-31. ญาณตฺตยนิทฺเทส}
\csromanheader{2}{29-31. ญาณ\-ตฺตย\-นิทฺเทส}
\proseitem{78}{\csromanfolio{89}กถํ วิหารนานตฺเต ปญฺญา วิหารฏฺเฐ ญาณํ, สมาปตฺติ\-นานตฺเต ปญฺญา สมาปตฺตฏฺเฐ ญาณํ, วิหารสมาปตฺติ\-นานตฺเต ปญฺญา วิหารสมาปตฺต\-ฏฺเฐ ญาณํ\csromandash{}นิมิตฺตํ ภยโต สมฺปสฺสมาโน อนิมิตฺเต อธิมุตฺตตฺตา ผุสฺส ผุสฺส วยํ ปสฺสติ, อนิมิตฺโต วิหาโร. ปณิธิํ ภยโต สมฺปสฺสมาโน อปฺปณิหิเต อธิมุตฺตตฺตา ผุสฺส ผุสฺส วยํ ปสฺสติ, อปฺปณิหิโต วิหาโร. อภินิเวสํ ภยโต สมฺปสฺสมาโน สุญฺญเต อธิมุตฺตตฺตา ผุสฺส ผุสฺส วยํ ปสฺสติ, สุญฺญโต วิหาโร.}

\prose{นิมิตฺตํ ภยโต สมฺปสฺสมาโน อนิมิตฺเต อธิมุตฺตตฺตา ปวตฺตํ อชฺฌุเปกฺขิตฺวา นิโรธํ นิพฺพานํ อนิมิตฺตํ อาวชฺชิตฺวา สมาปชฺชติ, อนิมิตฺตา สมาปตฺติ. ปณิธิํ ภยโต สมฺปสฺสมาโน อปฺปณิหิเต อธิมุตฺตตฺตา ปวตฺตํ อชฺฌุเปกฺขิตฺวา นิโรธํ นิพฺพานํ อปฺปณิหิตํ อาวชฺชิตฺวา สมาปชฺชติ, อปฺปณิหิตา สมาปตฺติ. อภินิเวสํ ภยโต สมฺปสฺสมาโน สุญฺญเต อธิมุตฺตตฺตา ปวตฺตํ อชฺฌุเปกฺขิตฺวา นิโรธํ นิพฺพานํ สุญฺญตํ อาวชฺชิตฺวา สมาปชฺชติ, สุญฺญตา สมาปตฺติ.}

\prose{นิมิตฺตํ ภยโต สมฺปสฺสมาโน อนิมิตฺเต อธิมุตฺตตฺตา ผุสฺส ผุสฺส วยํ ปสฺสติ, ปวตฺตํ อชฺฌุเปกฺขิตฺวา นิโรธํ นิพฺพานํ อนิมิตฺตํ อาวชฺชิตฺวา สมาปชฺชติ, อนิมิตฺต\-วิหารสมาปตฺติ. ปณิธิํ ภยโต สมฺปสฺสมาโน อปฺปณิหิเต อธิมุตฺตตฺตา ผุสฺส ผุสฺส วยํ ปสฺสติ, ปวตฺตํ อชฺฌุเปกฺขิตฺวา นิโรธํ นิพฺพานํ อปฺปณิหิตํ อาวชฺชิตฺวา สมาปชฺชติ, อปฺปณิหิต\-วิหารสมาปตฺติ. อภินิเวสํ ภยโต สมฺปสฺสมาโน สุญฺญเต อธิมุตฺตตฺตา ผุสฺส ผุสฺส วยํ ปสฺสติ, ปวตฺตํ อชฺฌุเปกฺขิตฺวา นิโรธํ นิพฺพานํ สุญฺญตํ อาวชฺชิตฺวา สมาปชฺชติ, สุญฺญต\-วิหารสมาปตฺติ.}

\proseitem{79}{รูปนิมิตฺตํ ภยโต สมฺปสฺสมาโน อนิมิตฺเต อธิมุตฺตตฺตา ผุสฺส ผุสฺส วยํ ปสฺสติ, อนิมิตฺโต วิหาโร. รูปปณิธํ ภยโต สมฺปสฺสมาโน อปฺปณิหิเต อธิมุตฺตตฺตา ผุสฺส ผุสฺส วยํ ปสฺสติ, อปฺปณิหิโต วิหาโร. รูปาภินิเวสํ ภยโต สมฺปสฺสมาโน สุญฺญเต อธิมุตฺตตฺตา ผุสฺส ผุสฺส วยํ ปสฺสติ, สุญฺญโต วิหาโร.}

\prose{รูปนิมิตฺตํ ภยโต สมฺปสฺสมาโน อนิมิตฺเต อธิมุตฺตตฺตา ปวตฺตํ อชฺฌุเปกฺขิตฺวา นิโรธํ นิพฺพานํ อนิมิตฺตํ อาวชฺชิตฺวา สมาปชฺชติ, อนิมิตฺตา \csromanfolio{90}สมาปตฺติ. รูปปณิธิํ ภยโต สมฺปสฺสมาโน อปฺปณิหิเต อธิมุตฺตตฺตา ปวตฺตํ อชฺฌุเปกฺขิตฺวา นิโรธํ นิพฺพานํ อปฺปณิหิตํ อาวชฺชิตฺวา สมาปชฺชติ, อปฺปณิหิตา สมาปตฺติ. รูปาภินิเวสํ ภยโต สมฺปสฺสมาโน สุญฺญเต อธิมุตฺตตฺตา ปวตฺตํ อชฺฌุเปกฺขิตฺวา นิโรธํ นิพฺพานํ สุญฺญตํ อาวชฺชิตฺวา สมาปชฺชติ, สุญฺญตา สมาปตฺติ.}

\prose{รูปนิมิตฺตํ ภยโต สมฺปสฺสมาโน อนิมิตฺเต อธิมุตฺตตฺตา ผุสฺส ผุสฺส วยํ ปสฺสติ, ปวตฺตํ อชฺฌุเปกฺขิตฺวา นิโรธํ นิพฺพานํ อนิมิตฺตํ อาวชฺชิตฺวา สมาปชฺชติ, อนิมิตฺต\-วิหารสมาปตฺติ. รูปปณิธิํ ภยโต สมฺปสฺสมาโน อปฺปณิหิเต อธิมุตฺตตฺตา ผุสฺส ผุสฺส วยํ ปสฺสติ, ปวตฺตํ อชฺฌุเปกฺขิตฺวา นิโรธํ นิพฺพานํ อปฺปณิหิตํ อาวชฺชิตฺวา สมาปชฺชติ, อปฺปณิหิต\-วิหารสมาปตฺติ. รูปาภินิเวสํ ภยโต สมฺปสฺสมาโน สุญฺญเต อธิมุตฺตตฺตา ผุสฺส ผุสฺส วยํ ปสฺสติ, ปวตฺตํ อชฺฌุเปกฺขิตฺวา นิโรธํ นิพฺพานํ สุญฺญตํ อาวชฺชิตฺวา สมาปชฺชติ, สุญฺญต\-วิหารสมาปตฺติ.}

\prose{เวทนานิมิตฺตํ ฯเปฯ. สญฺญานิมิตฺตํ. สงฺขาร\-นิมิตฺตํ. วิญฺญาณ\-นิมิตฺตํ. จกฺขุนิมิตฺตํ ฯเปฯ. ชรามรณ\-นิมิตฺตํ ภยโต สมฺปสฺสมาโน อนิมิตฺเต อธิมุตฺตตฺตา ผุสฺส ผุสฺส วยํ ปสฺสติ, อนิมิตฺโต วิหาโร. ชรามรณ\-ปณิธิํ ภยโต สมฺปสฺสมาโน อปฺปณิหิเต อธิมุตฺตตฺตา ผุสฺส ผุสฺส วยํ ปสฺสติ, อปฺปณิหิโต วิหาโร. ชรามรณา\-ภินิเวสํ ภยโต สมฺปสฺสมาโน สุญฺญเต อธิมุตฺตตฺตา ผุสฺส ผุสฺส วยํ ปสฺสติ, สุญฺญโต วิหาโร.}

\prose{ชรามรณ\-นิมิตฺตํ ภยโต สมฺปสฺสมาโน อนิมิตฺเต อธิมุตฺตตฺตา ปวตฺตํ อชฺฌุเปกฺขิตฺวา นิโรธํ นิพฺพานํ อนิมิตฺตํ อาวชฺชิตฺวา สมาปชฺชติ, อนิมิตฺตา สมาปตฺติ. ชรามรณ\-ปณิธิํ ภยโต สมฺปสฺสมาโน อปฺปณิหิเต อธิมุตฺตตฺตา ปวตฺตํ อชฺฌุเปกฺขิตฺวา นิโรธํ นิพฺพานํ อปฺปณิหิตํ อาวชฺชิตฺวา สมาปชฺชติ, อปฺปณิหิตา สมาปตฺติ. ชรามรณา\-ภินิเวสํ ภยโต สมฺปสฺสมาโน สุญฺญเต อธิมุตฺตตฺตา ปวตฺตํ อชฺฌุเปกฺขิตฺวา นิโรธํ นิพฺพานํ สุญฺญตํ อาวชฺชิตฺวา สมาปชฺชติ, สุญฺญตา สมาปตฺติ.}

\prosewithcloserruled{ชรามรณ\-นิมิตฺตํ ภยโต สมฺปสฺสมาโน อนิมิตฺเต อธิมุตฺตตฺตา ผุสฺส ผุสฺส วยํ ปสฺสติ, ปวตฺตํ อชฺฌุเปกฺขิตฺวา นิโรธํ นิพฺพานํ อนิมิตฺตํ อาวชฺชิตฺวา สมาปชฺชติ, อนิมิตฺต\-วิหารสมาปตฺติ. ชรามรณ\-ปณิธิํ \csromanfolio{91}ภยโต สมฺปสฺสมาโน อปฺปณิหิเต อธิมุตฺตตฺตา ผุสฺส ผุสฺส วยํ ปสฺสติ, ปวตฺตํ อชฺฌุเปกฺขิตฺวา นิโรธํ นิพฺพานํ อปฺปณิหิตํ อาวชฺชิตฺวา สมาปชฺชติ, อปฺปณิหิต\-วิหารสมาปตฺติ. ชรามรณา\-ภินิเวสํ ภยโต สมฺปสฺสมาโน สุญฺญเต อธิมุตฺตตฺตา ผุสฺส ผุสฺส วยํ ปสฺสติ, ปวตฺตํ อชฺฌุเปกฺขิตฺวา นิโรธํ นิพฺพานํ สุญฺญตํ อาวชฺชิตฺวา สมาปชฺชติ, สุญฺญต\-วิหารสมาปตฺติ. อญฺโญ อนิมิตฺโต วิหาโร, อญฺโญ อปฺปณิหิโต วิหาโร, อญฺโญ สุญฺญโต วิหาโร. อญฺญา อนิมิตฺตา สมาปตฺติ, อญฺญา อปฺปณิหิตา สมาปตฺติ, อญฺญา สุญฺญตา สมาปตฺติ. อญฺญา อนิมิตฺต\-วิหารสมาปตฺติ, อญฺญา อปฺปณิหิต\-วิหารสมาปตฺติ, อญฺญา สุญฺญต\-วิหารสมาปตฺติ. ตํญาตฏฺเฐน ญาณํ, ปชานนฏฺเฐน ปญฺญา. เตน วุจฺจติ วิหารนานตฺเต ปญฺญา วิหารฏฺเฐ ญาณํ. สมาปตฺติ\-นานตฺเต ปญฺญา สมาปตฺตฏฺเฐ ญาณํ. วิหารสมาปตฺติ\-นานตฺเต ปญฺญา วิหารสมาปตฺต\-ฏฺเฐ ญาณํ.}{ญาณ\-ตฺตย\-นิทฺเทโส เอกติํสติโม.}
\csromanmatikaanchor{mtk.25}
\csromantocmarkref{subsection}{32. อานนฺตริก สมาธิญาณนิทฺเทส}{mtk.25}
\bookmark[dest={mtk.25},level=2]{32. อานนฺตริก สมาธิญาณนิทฺเทส}
\csromanheader{2}{32. อานนฺตริก\-สมาธิ\-ญาณ\-นิทฺเทส}
\proseitem{80}{กถํ อวิกฺเขป\-ปริสุทฺธตฺตา อาสว\-สมุจฺเฉเท ปญฺญา อานนฺตริก\-สมาธิมฺหิ ญาณํ\csromandash{}เนกฺขมฺม\-วเสน จิตฺตสฺส เอกคฺคตา อวิกฺเขโป สมาธิ, ตสฺส สมาธิสฺส วเสน อุปฺปชฺชติ ญาณํ, เตน ญาเณน อาสวา ขียนฺติ, อิติ ปฐมํ สมโถ, ปจฺฉา ญาณํ, เตน ญาเณน อาสวานํ ขโย โหติ. เตน วุจฺจติ อวิกฺเขป\-ปริสุทฺธตฺตา อาสว\-สมุจฺเฉเท ปญฺญา อานนฺตริก\-สมาธิมฺหิ ญาณํ.}

\prose{อาสวาติ กตเม เต อาสวา? กามาสโว ภวาสโว ทิฏฺฐาสโว อวิชฺชาสโว. กตฺเถเต อาสวา ขียนฺติ? โสตาปตฺติ\-มคฺเคน อนวเสโส ทิฏฺฐาสโว ขียติ, อปายคมนีโย กามาสโว ขียติ, อปายคมนีโย ภวาสโว ขียติ, อปายคมนีโย อวิชฺชาสโว ขียติ, เอตฺเถเต อาสวา ขียนฺติ. สกทาคามิ\-มคฺเคน โอฬาริโก กามาสโว ขียติ, ตเทกฏฺโฐ ภวาสโว ขียติ, ตเทกฏฺโฐ อวิชฺชาสโว ขียติ, เอตฺเถเต อาสวา ขียนฺติ. อนาคามิมคฺเคน อนวเสโส กามาสโว ขียติ, ตเทกฏฺโฐ ภวาสโว ขียติ, \csromanfolio{92}ตเทกฏฺโฐ อวิชฺชาสโว ขียติ, เอตฺเถเต อาสวา ขียนฺติ. อรหตฺต\-มคฺเคน อนวเสโส ภวาสโว ขียติ, อนวเสโส อวิชฺชาสโว ขียติ, เอตฺเถเต อาสวา ขียนฺติ.}

\prose{อพฺยาปาทวเสน. อาโลกสญฺญา\-วเสน. อวิกฺเขป\-วเสน. ธมฺมววตฺถาน\-วเสน. ญาณวเสน. ปาโมชฺชวเสน.  ปฐมชฺฌาน\-วเสน. ทุติยชฺฌาน\-วเสน. ตติยชฺฌาน\-วเสน. จตุตฺถชฺฌาน\-วเสน.  อากาสา\-นญฺจา\-ยตนสมาปตฺติวเสน. วิญฺญาณญฺจายตน\-สมาปตฺติวเสน. อากิญฺจญฺญายตน\-สมาปตฺติวเสน. เนว\-สญฺญา\-นา\-สญฺญา\-ยตนสมาปตฺติวเสน.  ปถวีกสิณ\-วเสน. อาโปกสิณ\-วเสน. เตโชกสิณ\-วเสน. วาโยกสิณ\-วเสน. นีลกสิณ\-วเสน. ปีตกสิณ\-วเสน. โลหิตกสิณ\-วเสน. โอทาตกสิณ\-วเสน. อากาสกสิณ\-วเสน. วิญฺญาณกสิณ\-วเสน.  พุทฺธานุสฺสติ\-วเสน. ธมฺมานุสฺสติ\-วเสน. สํฆานุสฺสติวเสน. สีลานุสฺสติ\-วเสน. จาคานุสฺสติ\-วเสน. เทวตานุสฺสติ\-วเสน. อานาปานสฺสติ\-วเสน. มรณสฺสติ\-วเสน. กายคตาสติ\-วเสน. อุปสมานุสฺสติ\-วเสน.  อุทฺธุมาตก\-สญฺญาวเสน. วินีลกสญฺญา\-วเสน. วิปุพฺพกสญฺญา\-วเสน. วิจฺฉิทฺทก\-สญฺญาวเสน. วิกฺขายิตก\-สญฺญา\-วเสน. วิกฺขิตฺตก\-สญฺญา\-วเสน. หตวิกฺขิตฺตก\-สญฺญา\-วเสน. โลหิตก\-สญฺญา\-วเสน. ปุฬวกสญฺญา\-วเสน. อฏฺฐิกสญฺญา\-วเสน.}

\proseitem{81}{ทีฆํ อสฺสาสวเสน, ทีฆํ ปสฺสาสวเสน. รสฺสํ อสฺสาสวเสน, รสฺสํ ปสฺสาสวเสน. สพฺพ\-กาย\-ปฏิสํเวที อสฺสาสวเสน, สพฺพ\-กาย\-ปฏิสํเวที ปสฺสาสวเสน. ปสฺสมฺภยํ กายสงฺขารํ อสฺสาสวเสน, ปสฺสมฺภยํ กายสงฺขารํ ปสฺสาสวเสน.  ปีติปฏิสํเวที อสฺสาสวเสน, ปีติปฏิสํเวที ปสฺสาสวเสน. สุขปฏิสํเวที อสฺสาสวเสน, สุขปฏิสํเวที ปสฺสาสวเสน. จิตฺต\-สงฺขาร\-ปฏิสํเวที อสฺสาสวเสน, จิตฺต\-สงฺขาร\-ปฏิสํเวที ปสฺสาสวเสน. ปสฺสมฺภยํ จิตฺตสงฺขารํ อสฺสาสวเสน, ปสฺสมฺภยํ จิตฺตาสงฺขารํ ปสฺสาสวเสน.  จิตฺต\-ปฏิสํเวที อสฺสาสวเสน, จิตฺต\-ปฏิสํเวที ปสฺสาสวเสน. อภิปฺปโมทยํ จิตฺตํ อสฺสาสวเสน, อภิปฺปโมทยํ จิตฺตํ ปสฺสาสวเสน. สมาทหํ จิตฺตํ ฯเปฯ. วิโมจยํ จิตฺตํ.  อนิจฺจานุปสฺสี. วิราคานุปสฺสี. นิโรธานุปสฺสี. ปฏินิสฺสคฺคา\-นุปสฺสี อสฺสาสวเสน, ปฏินิสฺสคฺคา\-นุปสฺสี \csromanfolio{93}ปสฺสาสวเสน จิตฺตสฺส เอกคฺคตา อวิกฺเขโป สมาธิ. ตสฺส สมาธิสฺส วเสน อุปฺปชฺชติ ญาณํ. เตน ญาเณน อาสวา ขียนฺติ, อิติ ปฐมํ สมโถ, ปจฺฉา ญาณํ. เตน ญาเณน อาสวานํ ขโย โหติ, เตน วุจฺจติ อวิกฺเขป\-ปริสุทฺธตฺตา อาสว\-สมุจฺเฉเท ปญฺญา อานนฺตริก\-สมาธิมฺหิ ญาณํ.}

\prosewithcloserruled{อาสวาติ กตเม เต อาสวา? กามาสโว ภวาสโว ทิฏฺฐาสโว อวิชฺชาสโว. กตฺเถเต อาสวา ขียนฺติ? โสตาปตฺติ\-มคฺเคน อนวเสโส ทิฏฺฐาสโว ขียติ, อปายคมนีโย กามาสโว ขียติ, อปายคมนีโย ภวาสโว ขียติ, อปายคมนีโย อวิชฺชาสโว ขียติ, เอตฺเถเต อาสวา ขียนฺติ. สกทาคามิ\-มคฺเคน โอฬาริโก กามาสโว ขียติ, ตเทกฏฺโฐ ภวาสโว ขียติ, ตเทกฏฺโฐ อวิชฺชาสโว ขียติ, เอตฺเถเต อาสวา ขียนฺติ. อนาคามิมคฺเคน อนวเสโส กามาสโว ขียติ, ตเทกฏฺโฐ ภวาสโว ขียติ, ตเทกฏฺโฐ อวิชฺชาสโว ขียติ, เอตฺเถเต อาสวา ขียนฺติ. อรหตฺต\-มคฺเคน อนวเสโส ภวาสโว ขียติ, อนวเสโส อวิชฺชาสโว ขียติ, เอตฺเถเต อาสวา ขียนฺติ. ตํญาตฏฺเฐน ญาณํ, ปชานนฏฺเฐน ปญฺญา, เตน วุจฺจติ อวิกฺเขป\-ปริสุทฺธตฺตา อาสว\-สมุจฺเฉเท ปญฺญา อานนฺตริก\-สมาธิมฺหิ ญาณํ.}{อานนฺตริก\-สมาธิ\-ญาณ\-นิทฺเทโส ทฺวตฺติํสติโม.}
\csromanmatikaanchor{mtk.26}
\csromantocmarkref{subsection}{33. อรณวิหารญาณนิทฺเทส}{mtk.26}
\bookmark[dest={mtk.26},level=2]{33. อรณวิหารญาณนิทฺเทส}
\csromanheader{2}{33. อรณวิหาร\-ญาณ\-นิทฺเทส}
\proseitem{82}{กถํ ทสฺสนา\-ธิปเตยฺยํ สนฺโต จ วิหาราธิคโม ปณีตา\-ธิมุตฺตตา ปญฺญา อรณวิหาเร ญาณํ\csromandash{}ทสฺสนาธิปเตยฺยนฺติ อนิจฺจา\-นุปสฺสนา ทสฺสนา\-ธิปเตยฺยํ, ทุกฺขา\-นุปสฺสนา ทสฺสนา\-ธิปเตยฺยํ, อนตฺตา\-นุปสฺสนา ทสฺสนา\-ธิปเตยฺยํ, รูเป อนิจฺจา\-นุปสฺสนา ทสฺสนา\-ธิปเตยฺยํ, รูเป ทุกฺขา\-นุปสฺสนา ทสฺสนา\-ธิปเตยฺยํ, รูเป อนตฺตา\-นุปสฺสนา ทสฺสนา\-ธิปเตยฺยํ, เวทนาย ฯเปฯ สญฺญาย. สงฺขาเรสุ. วิญฺญาเณ. จกฺขุสฺมิํ ฯเปฯ. ชรามรเณ อนิจฺจา\-นุปสฺสนา ทสฺสนา\-ธิปเตยฺยํ, ชรามรเณ ทุกฺขา\-นุปสฺสนา ทสฺสนา\-ธิปเตยฺยํ, ชรามรเณ อนตฺตา\-นุปสฺสนา ทสฺสนา\-ธิปเตยฺยํ.}

\prose{\csromanfolio{94}\textbf{สนฺโต จ วิหาราธิคโม}ติ สุญฺญโต วิหาโร สนฺโต วิหาราธิคโม, อนิมิตฺโต วิหาโร สนฺโต วิหาราธิคโม, อปฺปณิหิโต วิหาโร สนฺโต วิหาราธิคโม.}

\prose{ปณีตา\-ธิมุตฺต\-ตาติ สุญฺญเต อธิมุตฺตตา ปณีตา\-ธิมุตฺตตา, อนิมิตฺเต อธิมุตฺตตา ปณีตา\-ธิมุตฺตตา, อปฺปณิหิเต อธิมุตฺตตา ปณีตา\-ธิมุตฺตตา.}

\prose{อรณวิหาโรติ ปฐมํ ฌานํ อรณวิหาโร, ทุติยํ ฌานํ อรณวิหาโร, ตติยํ ฌานํ อรณวิหาโร, จตุตฺถํ ฌานํ อรณวิหาโร อากาสา\-นญฺจา\-ยตนสมาปตฺติ อรณวิหาโร ฯเปฯ เนว\-สญฺญา\-นา\-สญฺญา\-ยตน- สมาปตฺติ อรณวิหาโร.}

\prosewithcloserruled{อรณวิหาโรติ เกนฏฺเฐน อรณวิหาโร, ปฐเมน ฌาเนน นีวรเณ หรตีติ อรณวิหาโร, ทุติเยน ฌาเนน วิตกฺกวิจาเร หรตีติ อรณวิหาโร, ตติเยน ฌาเนน ปีติํ หรตีติ อรณวิหาโร, จตุตฺเถน ฌาเนน สุขทุกฺเข หรตีติ อรณวิหาโร.  อากาสา\-นญฺจา\-ยตนสมาปตฺติยา รูปสญฺญํ ปฏิฆสญฺญํ นานตฺตสญฺญํ หรตีติ อรณวิหาโร, วิญฺญาณญฺจายตน\-สมาปตฺติยา อากาสา\-นญฺจา\-ยตนสญฺญํ หรตีติ อรณวิหาโร, อากิญฺจญฺญายตน\-สมาปตฺติยา วิญฺญาณญฺจา\-ยตนสญฺญํ หรตีติ อรณวิหาโร, เนว\-สญฺญา\-นา\-สญฺญา\-ยตนสมาปตฺติยา อากิญฺจญฺญายตน\-สญฺญํ หรตีติ อรณวิหาโร, อยํ อรณวิหาโร. ตํญาตฏฺเฐน ญาณํ, ปชานนฏฺเฐน ปญฺญา เตน วุจฺจติ ทสฺสนา\-ธิปเตยฺยํ สนฺโต จ วิหาราธิคโม ปณีตา\-ธิมุตฺตตา ปญฺญา อรณวิหาเร ญาณํ.}{อรณวิหาร\-ญาณ\-นิทฺเทโส เตตฺติํสติโม.}
\csromanmatikaanchor{mtk.27}
\csromantocmarkref{subsection}{34. นิโรธสมาปตฺติญาณนิทฺเทส}{mtk.27}
\bookmark[dest={mtk.27},level=2]{34. นิโรธสมาปตฺติญาณนิทฺเทส}
\csromanheader{2}{34. นิโรธสมาปตฺติ\-ญาณ\-นิทฺเทส}
\proseitem{83}{กถํ ทฺวีหิ พเลหิ สมนฺนาคตตฺตา ตโย จ สงฺขารานํ ปฏิปฺปสฺสทฺธิยา โสฬสหิ ญาณจริยาหิ นวหิ สมาธิจริยาหิ วสิภาวตา ปญฺญา นิโรธ\-สมาปตฺติยา ญาณํ\csromandash{}}

\prose{\csromanfolio{95}ทฺวีหิ พเลหีติ เทฺว พลานิ สมถพลํ วิปสฺสนาพลํ. กตมํ สมถพลํ? เนกฺขมฺม\-วเสน จิตฺตสฺเส\-กคฺคตา อวิกฺเขโป สมถพลํ, อพฺยาปาทวเสน จิตฺตสฺเส\-กคฺคตา อวิกฺเขโป สมถพลํ, อาโลกสญฺญา\-วเสน จิตฺตสฺเส\-กคฺคตา อวิกฺเขโป สมถพลํ, อวิกฺเขป\-วเสน จิตฺตสฺเส\-กคฺคตา อวิกฺเขโป สมถพลํ ฯเปฯ ปฏินิสฺสคฺคา\-นุปสฺสี อสฺสาสวเสน จิตฺตสฺเส\-กคฺคตา อวิกฺเขโป สมถพลํ, ปฏินิสฺสคฺคา\-นุปสฺสี ปสฺสาสวเสน จิตฺตสฺเส\-กคฺคตา อวิกฺเขโป สมถพลํ.}

\prose{สมถพลนฺติ เกนฏฺเฐน สมถพลํ? ปฐเมน ฌาเนน นีวรเณ น กมฺปตีติ สมถพลํ, ทุติเยน ฌาเนน วิตกฺกวิจาเร น กมฺปตีติ สมถพลํ, ตติเยน ฌาเนน ปีติยา น กมฺปตีติ สมถพลํ, จตุตฺเถน ฌาเนน สุขทุกฺเข น กมฺปตีติ สมถพลํ.  อากาสา\-นญฺจา\-ยตนสมาปตฺติยา รูปสญฺญาย ปฏิฆสญฺญาย นานตฺตสญฺญาย น กมฺปตีติ สมถพลํ, วิญฺญาณญฺจายตน\-สมาปตฺติยา อากาสา\-นญฺจา\-ยตนสญฺญาย น กมฺปตีติ สมถพลํ, อากิญฺจญฺญายตน\-สมาปตฺติยา วิญฺญาณญฺจา\-ยตนสญฺญาย น กมฺปตีติ สมถพลํ, เนว\-สญฺญา\-นา\-สญฺญา\-ยตนสมาปตฺติยา อากิญฺจญฺญายตน\-สญฺญาย น กมฺปตีติ สมถพลํ, อุทฺธจฺเจ จ อุทฺธจฺจสหคต\-กิเลเส จ ขนฺเธ จ น กมฺปติ น จลติ น เวธตีติ สมถพลํ. อิทํ สมถพลํ.}

\prose{กตมํ วิปสฺสนาพลํ? อนิจฺจา\-นุปสฺสนา วิปสฺสนาพลํ, ทุกฺขา\-นุปสฺสนา วิปสฺสนาพลํ, อนตฺตา\-นุปสฺสนา วิปสฺสนาพลํ, นิพฺพิทา\-นุปสฺสนา วิปสฺสนาพลํ, วิราคา\-นุปสฺสนา วิปสฺสนาพลํ, นิโรธา\-นุปสฺสนา วิปสฺสนาพลํ, ปฏินิสฺสคฺคา\-นุปสฺสนา วิปสฺสนาพลํ. รูเป อนิจฺจา\-นุปสฺสนา วิปสฺสนาพลํ ฯเปฯ รูเป ปฏินิสฺสคฺคา\-นุปสฺสนา วิปสฺสนาพลํ.  เวทนาย ฯเปฯ. สญฺญาย. สงฺขาเรสุ. วิญฺญาเณ. จกฺขุสฺมิํ ฯเปฯ. ชรามรเณ อนิจฺจา\-นุปสฺสนา วิปสฺสนาพลํ ฯเปฯ ชรามรเณ ปฏินิสฺสคฺคา\-นุปสฺสนา วิปสฺสนาพลํ.}

\prose{วิปสฺสนา\-พลนฺติ เกนฏฺเฐน วิปสฺสนาพลํ? อนิจฺจา\-นุปสฺสนาย นิจฺจสญฺญาย น กมฺปตีติ วิปสฺสนาพลํ, ทุกฺขา\-นุปสฺสนาย สุขสญฺญาย น กมฺปตีติ วิปสฺสนาพลํ, อนตฺตา\-นุปสฺสนาย อตฺตสญฺญาย น กมฺปตีติ วิปสฺสนาพลํ, นิพฺพิทา\-นุปสฺสนาย นนฺทิยา น กมฺปตีติ วิปสฺสนาพลํ, วิราคา\-นุปสฺสนาย ราเค น กมฺปตีติ วิปสฺสนาพลํ, นิโรธา\-นุปสฺสนาย สมุทเย น กมฺปตีติ วิปสฺสนาพลํ, ปฏินิสฺสคฺคา\-นุปสฺสนาย อาทาเน น \csromanfolio{96}กมฺปตีติ วิปสฺสนาพลํ, อวิชฺชาย จ อวิชฺชา\-สหคต\-กิเลเส จ ขนฺเธ จ น กมฺปติ น จลติ น เวธตีติ วิปสฺสนาพลํ. อิทํ วิปสฺสนาพลํ.}

\prose{ตโย จ สงฺขารานํ ปฏิปฺปสฺสทฺธิยาติ กตเมสํ ติณฺณนฺนํ สงฺขารานํ ปฏิปฺปสฺสทฺธิยา? ทุติยํ ฌานํ สมาปนฺนสฺส วิตกฺกวิจารา วจีสงฺขารา ปฏิปฺปสฺสทฺธา โหนฺติ, จตุตฺถํ ฌานํ สมาปนฺนสฺส อสฺสาสปสฺสาสา กายสงฺขารา ปฏิปฺปสฺสทฺธา โหนฺติ, สญฺญาเวทยิต\-นิโรธํ สมาปนฺนสฺส สญฺญา จ เวทนา จ จิตฺตสงฺขารา ปฏิปฺปสฺสทฺธา โหนฺติ. อิเมสํ ติณฺณนฺนํ สงฺขารานํ ปฏิปฺปสฺสทฺธิยา.}

\proseitem{84}{โสฬสหิ ญาณจริยาหีติ กตมาหิ โสฬสหิ ญาณจริยาหิ? อนิจฺจา\-นุปสฺสนา ญาณจริยา, ทุกฺขา\-นุปสฺสนา ญาณจริยา, อนตฺตา\-นุปสฺสนา ญาณจริยา, นิพฺพิทา\-นุปสฺสนา ญาณจริยา, วิราคา\-นุปสฺสนา ญาณจริยา, นิโรธา\-นุปสฺสนา ญาณจริยา, ปฏินิสฺสคฺคา\-นุปสฺสนา ญาณจริยา, วิวฏฺฏนา\-นุปสฺสนา ญาณจริยา. โสตาปตฺติมคฺโค ญาณจริยา, โสตาปตฺติ\-ผลสมาปตฺติ ญาณจริยา, สกทาคามิ\-มคฺโค ญาณจริยา, สกทาคามิ\-ผลสมาปตฺติ ญาณจริยา, อนาคามิมคฺโค ญาณจริยา, อนาคามิ\-ผลสมาปตฺติ ญาณจริยา, อรหตฺตมคฺโค ญาณจริยา, อรหตฺตผล\-สมาปตฺติ ญาณจริยา. อิมาหิ โสฬสหิ ญาณจริยาหิ.}

\proseitem{85}{นวหิ สมาธิจริยาหีติ กตมาหิ นวหิ สมาธิจริยาหิ? ปฐมํ ฌานํ สมาธิจริยา, ทุติยํ ฌานํ สมาธิจริยา, ตติยํ ฌานํ สมาธิจริยา, จตุตฺถํ ฌานํ สมาธิจริยา.  อากาสา\-นญฺจา\-ยตนสมาปตฺติ ฯเปฯ. วิญฺญาณญฺจา\-ยตนสมาปตฺติ. อากิญฺจญฺญายตน\-สมาปตฺติ. เนว\-สญฺญา\-นา\-สญฺญา\-ยตนสมาปตฺติ สมาธิจริยา.  ปฐมํ ฌานํ ปฏิลาภ\-ตฺถาย วิตกฺโก จ วิจาโร จ ปีติ จ สุขญฺจ จิตฺเตกคฺคตา จ ฯเปฯ. เนว\-สญฺญา\-นา\-สญฺญา\-ยตนสมาปตฺติํ ปฏิลาภ\-ตฺถาย วิตกฺโก จ วิจาโร จ ปีติ จ สุขญฺจ จิตฺเตกคฺคตา จ. อิมาหิ นวหิ สมาธิจริยาหิ.}

\prose{\textbf{วสี}ติ ปญฺจ วสิโย อาวชฺชนวสี สมาปชฺชนวสี อธิฏฺฐานวสี วุฏฺฐานวสี ปจฺจเวกฺขณาวสี. ปฐมํ ฌานํ ยตฺถิจฺฉกํ ยทิจฺฉกํ ยาวติจฺฉกํ อาวชฺชติ, อาวชฺชนาย ทนฺธายิตตฺตํ นตฺถีติ อาวชฺชนวสี. ปฐมํ ฌานํ ยตฺถิจฺฉกํ ยทิจฺฉกํ ยาวติจฺฉกํ สมาปชฺชติ, สมาปชฺชนาย ทนฺธายิตตฺตํ นตฺถีติ สมาปชฺชนวสี. ปฐมํ ฌานํ ยตฺถิจฺฉกํ ยทิจฺฉกํ ยาวติจฺฉกํ \csromanfolio{97}อธิฏฺฐาติ, อธิฏฺฐาเน ทนฺธายิตตฺตํ นตฺถีติ อธิฏฺฐานวสี. ปฐมํ ฌานํ ยตฺถิจฺฉกํ ยทิจฺฉกํ ยาวติจฺฉกํ วุฏฺฐาติ, วุฏฺฐาเน ทนฺธายิตตฺตํ นตฺถีติ วุฏฺฐานวสี. ปฐมํ ฌานํ ยตฺถิจฺฉกํ ยทิจฺฉกํ ยาวติจฺฉกํ ปจฺจเวกฺขติ, ปจฺจเวกฺขณาย ทนฺธายิตตฺตํ นตฺถีติ ปจฺจเวกฺขณาวสี.}

\prosewithcloserruled{ทุติยํ ฌานํ ฯเปฯ. เนว\-สญฺญา\-นา\-สญฺญา\-ยตนสมาปตฺติํ ยตฺถิจฺฉกํ ยทิจฺฉกํ ยาวติจฺฉกํ อาวชฺชติ, อาวชฺชนาย ทนฺธายิตตฺตํ นตฺถีติ อาวชฺชนวสี. เนว\-สญฺญา\-นา\-สญฺญา\-ยตนสมาปตฺติํ ยตฺถิจฺฉกํ ยทิจฺฉกํ ยาวติจฺฉกํ สมาปชฺชติ ฯเปฯ อธิฏฺฐาติ, วุฏฺฐาติ, ปจฺจเวกฺขติ, ปจฺจเวกฺขณาย ทนฺธายิตตฺตํ นตฺถีติ ปจฺจเวกฺขณาวสี. อิมา ปญฺจ วสิโย. ตํญาตฏฺเฐน ญาณํ ปชานนฏฺเฐน ปญฺญา. เตน วุจฺจติ ทฺวีหิ พเลหิ สมนฺนาคตตฺตา ตโย จ สงฺขารานํ ปฏิปสฺสทฺธิยา โสฬสหิ ญาณจริยาหิ นวหิ สมาธิจริยาหิ วสีภาวตา ปญฺญา นิโรธ\-สมาปตฺติยา ญาณํ.}{นิโรธสมาปตฺติ\-ญาณ\-นิทฺเทโส จตุตฺติํสติโม.}
\csromanmatikaanchor{mtk.28}
\csromantocmarkref{subsection}{35. ปรินิพฺพานญาณนิทฺเทส}{mtk.28}
\bookmark[dest={mtk.28},level=2]{35. ปรินิพฺพานญาณนิทฺเทส}
\csromanheader{2}{35. ปรินิพฺพาน\-ญาณ\-นิทฺเทส}
\proseitem{86}{กถํ สมฺปชานสฺส ปวตฺต\-ปริยาทาเน ปญฺญา ปรินิพฺพาเน ญาณํ\csromandash{}อิธ สมฺปชาโน เนกฺขมฺเมน กามจฺฉนฺทสฺส ปวตฺตํ ปริยาทิยติ. อพฺยาปาเทน พฺยาปาทสฺส ปวตฺตํ ปริยาทิยติ. อาโลกสญฺญาย ถินมิทฺธสฺส ปวตฺตํ ปริยาทิยติ. อวิกฺเขเปน อุทฺธจฺจสฺส ปวตฺตํ ปริยาทิยติ. ธมฺม\-ววตฺถาเนน วิจิกิจฺฉาย ฯเปฯ. ญาเณน อวิชฺชาย. ปาโมชฺเชน อรติยา. ปฐเมน ฌาเนน นีวรณานํ ปวตฺตํ ปริยาทิยติ ฯเปฯ. อรหตฺต\-มคฺเคน สพฺพกิเลสานํ ปวตฺตํ ปริยาทิยติ.}

\prosewithcloserruled{อถวา ปน สมฺปชานสฺส อนุปาทิเสสาย นิพฺพานธาตุยา ปรินิพฺพายนฺตสฺส อิทํ เจว จกฺขุปวตฺตํ ปริยาทิยติ, อญฺญญฺจ จกฺขุปวตฺตํ น อุปฺปชฺชติ. อิทํ เจว โสตปวตฺตํ ฯเปฯ. ฆานปวตฺตํ. ชิวฺหาปวตฺตํ. กายปวตฺตํ. มโนปวตฺตํ ปริยาทิยติ, อญฺญญฺจ มโนปวตฺตํ น อุปฺปชฺชติ. อิทํ สมฺปชานสฺส \csromanfolio{98}ปวตฺต\-ปริยาทาเน ปญฺญา ปรินิพฺพาเน ญาณํ. ตํญาตฏฺเฐน ญาณํ, ปชานนฏฺเฐน ปญฺญา. เตน วุจฺจติ สมฺปชานสฺส ปวตฺต\-ปริยาทาเน ปญฺญา ปรินิพฺพาเน ญาณํ.}{ปรินิพฺพาน\-ญาณ\-นิทฺเทโส ปญฺจติํสติโม.}
\csromanmatikaanchor{mtk.29}
\csromantocmarkref{subsection}{36. สมสีสฏฺฐญาณนิทฺเทส}{mtk.29}
\bookmark[dest={mtk.29},level=2]{36. สมสีสฏฺฐญาณนิทฺเทส}
\csromanheader{2}{36. สมสีส\-ฏฺฐญาณ\-นิทฺเทส}
\proseitem{87}{กถํ สพฺพธมฺมานํ สมฺมา สมุจฺเฉเท นิโรเธ จ อนุปฏฺฐานตา ปญฺญา สมสีสฏฺเฐ ญาณํ\csromandash{}สพฺพ\-ธมฺมานนฺติ ปญฺจกฺขนฺธา ทฺวาทสา\-ยตนานิ อฏฺฐารส ธาตุโย กุสลา ธมฺมา อกุสลา ธมฺมา อพฺยากตา ธมฺมา กามาวจรา ธมฺมา รูปาวจรา ธมฺมา อรูปาวจรา ธมฺมา อปริยาปนฺนา ธมฺมา. สมฺมา สมุจฺเฉเทติ เนกฺขมฺเมน กามจฺฉนฺทํ สมฺมา สมุจฺฉินฺทติ, อพฺยาปาเทน พฺยาปาทํ สมฺมา สมุจฺฉินฺทติ. อาโลกสญฺญาย ถินมิทฺธํ สมฺมา สมุจฺฉินฺทติ. อวิกฺเขเปน อุทฺธจฺจํ สมฺมา สมุจฺฉินฺทติ. ธมฺม\-ววตฺถาเนน วิจิกิจฺฉํ สมฺมา สมุจฺฉินฺทติ. ญาเณน อวิชฺชํ สมฺมา สมุจฺฉินฺทติ. ปาโมชฺเชน อรติํ สมฺมา สมุจฺฉินฺทติ. ปฐเมน ฌาเนน นีวรเณ สมฺมา สมุจฺฉินฺทติ ฯเปฯ. อรหตฺต\-มคฺเคน สพฺพกิเลเส สมฺมา สมุจฺฉินฺทติ.}

\prose{นิโรเธติ เนกฺขมฺเมน กามจฺฉนฺทํ นิโรเธติ, อพฺยาปาเทน พฺยาปาทํ นิโรเธติ. อาโลกสญฺญาย ถินมิทฺธํ นิโรเธติ. อวิกฺเขเปน อุทฺธจฺจํ นิโรเธติ, ธมฺม\-ววตฺถาเนน วิจิกิจฺฉํ นิโรเธติ. ญาเณน อวิชฺชํ นิโรเธติ. ปาโมชฺเชน อรติํ นิโรเธติ. ปฐเมน ฌาเนน นีวรเณ นิโรเธติ ฯเปฯ อรหตฺต\-มคฺเคน สพฺพกิเลเส นิโรเธติ.}

\prose{อนุป\-ฏฺฐานตาติ เนกฺขมฺมํ ปฏิลทฺธสฺส กามจฺฉนฺโท น อุปฏฺฐาติ. อพฺยาปาทํ ปฏิลทฺธสฺส พฺยาปาโท น อุปฏฺฐาติ. อาโลกสญฺญํ ปฏิลทฺธสฺส ถินมิทฺธํ น อุปฏฺฐาติ. อวิกฺเขปํ ปฏิลทฺธสฺส อุทฺธจฺจํ น อุปฏฺฐาติ. ธมฺม\-ววตฺถานํ ปฏิลทฺธสฺส วิจิกิจฺฉา น อุปฏฺฐาติ. ญาณํ ปฏิลทฺธสฺส อวิชฺชา น อุปฏฺฐาติ. ปาโมชฺชํ ปฏิลทฺธสฺส อรติ น อุปฏฺฐาติ. ปฐมํ ฌานํ ปฏิลทฺธสฺส นีวรณา น อุปฏฺฐหนฺติ ฯเปฯ. อรหตฺตมคฺคํ ปฏิลทฺธสฺส สพฺพกิเลสา น อุปฏฺฐหนฺติ.}

\prose{\textbf{สมนฺ}ติ กามจฺฉนฺทสฺส ปหีนตฺตา เนกฺขมฺมํ สมํ. พฺยาปาทสฺส ปหีนตฺตา อพฺยาปาโท สมํ. ถินมิทฺธสฺส ปหีนตฺตา อาโลกสญฺญา สมํ. \csromanfolio{99}อุทฺธจฺจสฺส ปหีนตฺตา อวิกฺเขโป สมํ. วิจิกิจฺฉาย ปหีนตฺตา ธมฺม\-ววตฺถานํ สมํ. อวิชฺชาย ปหีนตฺตา ญาณํ สมํ. อรติยา ปหีนตฺตา ปาโมชฺชํ สมํ. นีวรณานํ ปหีนตฺตา ปฐมํ ฌานํ สมํ ฯเปฯ. สพฺพกิเลสานํ ปหีนตฺตา อรหตฺตมคฺโค สมํ.}

\prosewithcloserruled{สีสนฺติ เตรส สีสานิ. ปลิโพธ\-สีสญฺจ ตณฺหา, วินิพนฺธน\-สีสญฺจ มาโน, ปรามาส\-สีสญฺจ ทิฏฺฐิ, วิกฺเขป\-สีสญฺจ อุทฺธจฺจํ, สํกิเลส\-สีสญฺจ อวิชฺชา, อธิโมกฺข\-สีสญฺจ สทฺธา, ปคฺคหสีสญฺจ วีริยํ, อุปฏฺฐาน\-สีสญฺจ สติ, อวิกฺเขป\-สีสญฺจ สมาธิ, ทสฺสนสีสญฺจ ปญฺญา, ปวตฺตสีสญฺจ ชีวิตินฺทฺริยํ, โคจรสีสญฺจ วิโมกฺโข, สงฺขาร\-สีสญฺจ นิโรโธ. ตํญาตฏฺเฐน ญาณํ ปชานนฏฺเฐน ปญฺญา. เตน วุจฺจติ สพฺพธมฺมานํ สมฺมา สมุจฺเฉเท นิโรเธ จ อนุปฏฺฐานตา ปญฺญา สมสีสฏฺเฐ ญาณํ.}{สมสีส\-ฏฺฐญาณ\-นิทฺเทโส ฉตฺติํสติโม.}
\csromanmatikaanchor{mtk.30}
\csromantocmarkref{subsection}{37. สลฺเลขฏฺฐญาณนิทฺเทส}{mtk.30}
\bookmark[dest={mtk.30},level=2]{37. สลฺเลขฏฺฐญาณนิทฺเทส}
\csromanheader{2}{37. สลฺเลข\-ฏฺฐญาณ\-นิทฺเทส}
\proseitem{88}{กถํ ปุถุนานตฺต\-เตช\-ปริยาทาเน ปญฺญา สลฺเลขฏฺเฐ\csromansharedfootnote{fcc8d8b27bb4f062}{ปุถุนานตฺเตกตฺตเตชปริยาทาเน (ก)} ญาณํ\csromandash{} ปุถูติ ราโค ปุถุ, โทโส ปุถุ, โมโห ปุถุ, โกโธ ฯเปฯ. อุปนาโห. มกฺโข. ปฬาโส. อิสฺสา. มจฺฉริยํ. มายา. สาเฐยฺยํ. ถมฺโภ. สารมฺโภ. มาโน. อติมาโน. มโท. ปมาโท. สพฺเพ กิเลสา. สพฺเพ ทุจฺจริตา. สพฺเพ อภิสงฺขารา. สพฺเพ ภวคามิกมฺมา.}

\prose{นานตฺเตกตฺตนฺติ กามจฺฉนฺโท นานตฺตํ, เนกฺขมฺมํ เอกตฺตํ. พฺยาปาโท นานตฺตํ อพฺยาปาโท เอกตฺตํ. ถินมิทฺธํ นานตฺตํ, อาโลกสญฺญา เอกตฺตํ. อุทฺธจฺจํ นานตฺตํ, อวิกฺเขโป เอกตฺตํ. วิจิกิจฺฉา นานตฺตํ, ธมฺม\-ววตฺถานํ เอกตฺตํ. อวิชฺชา นานตฺตํ, ญาณํ เอกตฺตํ. อรติ นานตฺตํ, ปาโมชฺชํ เอกตฺตํ. นีวรณา นานตฺตํ, ปฐมํ ฌานํ เอกตฺตํ ฯเปฯ สพฺเพ กิเลสา นานตฺตํ, อรหตฺตมคฺโค เอกตฺตํ.}

\prose{\textbf{เตโช}ติ ปญฺจ เตชา จรณเตโช คุณเตโช ปญฺญาเตโช ปุญฺญเตโช ธมฺมเตโช. จรณเตเชน เตชิตตฺตา ทุสฺสีลฺยเตชํ ปริยาทิยติ, คุณเตเชน เตชิตตฺตา อคุณเตชํ ปริยาทิยติ, ปญฺญาเตเชน เตชิตตฺตา ทุปฺปญฺญเตชํ ปริยาทิยติ, ปุญฺญเตเชน \csromanfolio{100}เตชิตตฺตา อปุญฺญเตชํ ปริยาทิยติ, ธมฺมเตเชน เตชิตตฺตา อธมฺมเตชํ ปริยาทิยติ.}

\prosewithcloserruled{สลฺเลโขติ กามจฺฉนฺโท อสลฺเลโข, เนกฺขมฺมํ สลฺเลโข. พฺยาปาโท อสลฺเลโข, อพฺยาปาโท สลฺเลโข. ถินมิทฺธํ อสลฺเลโข, อาโลกสญฺญา สลฺเลโข. อุทฺธจฺจํ อสลฺเลโข, อวิกฺเขโป สลฺเลโข. วิจิกิจฺฉา อสลฺเลโข, ธมฺม\-ววตฺถานํ สลฺเลโข. อวิชฺชา อสลฺเลโข, ญาณํ สลฺเลโข. อรติ อสลฺเลโข, ปาโมชฺชํ สลฺเลโข. นีวรณา อสลฺเลโข, ปฐมํ ฌานํ สลฺเลโข ฯเปฯ. สพฺพกิเลสา อสลฺเลโข, อรหตฺตมคฺโค สลฺเลโข. ตํญาตฏฺเฐน ญาณํ, ปชานนฏฺเฐน ปญฺญา. เตน วุจฺจติ ปุถุนานตฺต\-เตช\-ปริยาทาเน ปญฺญา สลฺเลขฏฺเฐ ญาณํ.}{สลฺเลข\-ฏฺฐญาณ\-นิทฺเทโส สตฺตติํสติโม.}
\csromanmatikaanchor{mtk.31}
\csromantocmarkref{subsection}{38. วีริยารมฺภญาณนิทฺเทส}{mtk.31}
\bookmark[dest={mtk.31},level=2]{38. วีริยารมฺภญาณนิทฺเทส}
\csromanheader{2}{38. วีริยารมฺภ\-ญาณ\-นิทฺเทส}
\proseitem{89}{กถํ อสลฺลีน\-ตฺต\-ปหิตตฺต\-ปคฺคหฏฺเฐ ปญฺญา วีริยารมฺเภ ญาณํ\csromandash{}อนุปฺปนฺนานํ ปาปกานํ อกุสลานํ ธมฺมานํ อนุปฺปาทาย อสลฺลีน\-ตฺต\-ปหิตตฺต\-ปคฺคหฏฺเฐ ปญฺญา วีริยารมฺเภ ญาณํ. อุปฺปนฺนานํ ปาปกานํ อกุสลานํ ธมฺมานํ ปหานาย อสลฺลีน\-ตฺต\-ปหิตตฺต\-ปคฺคหฏฺเฐ ปญฺญา วีริยารมฺเภ ญาณํ. อนุปฺปนฺนานํ กุสลานํ ธมฺมานํ อุปฺปาทาย อสลฺลีน\-ตฺต\-ปหิตตฺต\-ปคฺคหฏฺเฐ ปญฺญา วีริยารมฺเภ ญาณํ. อุปฺปนฺนานํ กุสลานํ ธมฺมานํ ฐิติยา อสมฺโมสาย ภิโยภาวาย เวปุลฺลาย ภาวนาย ปาริปูริยา อสลฺลีน\-ตฺต\-ปหิตตฺต\-ปคฺคหฏฺเฐ ปญฺญา วีริยารมฺเภ ญาณํ.}

\prose{อนุปฺปนฺนสฺส กามจฺฉนฺทสฺส อนุปฺปาทาย อสลฺลีน\-ตฺต\-ปหิตตฺต\-ปคฺคหฏฺเฐ ปญฺญา วีริยารมฺเภ ญาณํ. อุปฺปนฺนสฺส กามจฺฉนฺทสฺส ปหานาย อสลฺลีน\-ตฺต\-ปหิตตฺต\-ปคฺคหฏฺเฐ ปญฺญา วีริยารมฺเภ ญาณํ. อนุปฺปนฺนสฺส เนกฺขมฺมสฺส อุปฺปาทาย อสลฺลีน\-ตฺต\-ปหิตตฺต\-ปคฺคหฏฺเฐ ปญฺญา วีริยารมฺเภ ญาณํ. อุปฺปนฺนสฺส เนกฺขมฺมสฺส ฐิติยา อสมฺโมสาย ภิยฺโยภาวาย เวปุลฺลาย ภาวนาย ปาริปูริยา อสลฺลีน\-ตฺต\-ปหิตตฺต\-ปคฺคหฏฺเฐ ปญฺญา วีริยารมฺเภ ญาณํ ฯเปฯ.}

\prosewithcloserruled{อนุปฺปนฺนานํ สพฺพกิเลสานํ อนุปฺปาทาย อสลฺลีน\-ตฺต\-ปหิตตฺต\-ปคฺคหฏฺเฐ ปญฺญา วีริยารมฺเภ ญาณํ. อุปฺปนฺนานํ สพฺพกิเลสานํ ปหานาย \csromanfolio{101}อสลฺลีน\-ตฺต\-ปหิตตฺต\-ปคฺคหฏฺเฐ ปญฺญา วีริยารมฺเภ ญาณํ ฯเปฯ. อนุปฺปนฺนสฺส อรหตฺต\-มคฺคสฺส อุปฺปาทาย อสลฺลีน\-ตฺต\-ปหิตตฺต\-ปคฺคหฏฺเฐ ปญฺญา วีริยารมฺเภ ญาณํ. อุปฺปนฺนสฺส อรหตฺต\-มคฺคสฺส ฐิติยา อสมฺโมสาย ภิยฺโยภาวาย เวปุลฺลาย ภาวนาย ปาริปูริยา อสลฺลีน\-ตฺต\-ปหิตตฺต\-ปคฺคหฏฺเฐ ปญฺญา วีริยารมฺเภ ญาณํ. ตํญาตฏฺเฐน ญาณํ, ปชานนฏฺเฐน ปญฺญา. เตน วุจฺจติ อสลฺลีน\-ตฺต\-ปหิตตฺต\-ปคฺคหฏฺเฐ ปญฺญา วีริยารมฺเภ ญาณํ.}{วีริยารมฺภ\-ญาณ\-นิทฺเทโส อฏฺฐติํสติโม.}
\csromanmatikaanchor{mtk.32}
\csromantocmarkref{subsection}{39. อตฺถสนฺทสฺสนญาณนิทฺเทส}{mtk.32}
\bookmark[dest={mtk.32},level=2]{39. อตฺถสนฺทสฺสนญาณนิทฺเทส}
\csromanheader{2}{39. อตฺถ\-สนฺทสฺสน\-ญาณ\-นิทฺเทส}
\proseitem{90}{กถํ นานา\-ธมฺมปฺปกาสนตา ปญฺญา อตฺถ\-สนฺทสฺสเน ญาณํ\csromandash{} นานาธมฺมาติ ปญฺจกฺขนฺธา ทฺวาทสา\-ยตนานิ อฏฺฐารส ธาตุโย กุสลา ธมฺมา อกุสลา ธมฺมา อพฺยากตา ธมฺมา กามาวจรา ธมฺมา รูปาวจรา ธมฺมา อรูปาวจรา ธมฺมา อปริยาปนฺนา ธมฺมา.}

\prose{\textbf{ปกาสนตา}ติ รูปํ อนิจฺจโต ปกาเสติ, รูปํ ทุกฺขโต ปกาเสติ, รูปํ อนตฺตโต ปกาเสติ. เวทนํ ฯเปฯ. สญฺญํ. สงฺขาร. วิญฺญาณํ. จกฺขุํ ฯเปฯ. ชรามรณํ อนิจฺจโต ปกาเสติ, ชรามรณํ ทุกฺขโต ปกาเสติ, ชรามรณํ อนตฺตโต ปกาเสติ.}

\prose{อตฺถ\-สนฺทสฺส\-เนติ กามจฺฉนฺทํ ปชหนฺโต เนกฺขมฺมตฺถํ สนฺทสฺเสติ, พฺยาปาทํ ปชหนฺโต อพฺยาปาทตฺถํ สนฺทสฺเสติ, ถินมิทฺธํ ปชหนฺโต อาโลกสญฺญตฺถํ สนฺทสฺเสติ, อุทฺธจฺจํ ปชหนฺโต อวิกฺเขปตฺถํ สนฺทสฺเสติ, วิจิกิจฺฉํ ปชหนฺโต ธมฺมววตฺถานตฺถํ สนฺทสฺเสติ, อวิชฺชํ ปชหนฺโต ญาณตฺถํ สนฺทสฺเสติ, อรติํ ปชหนฺโต ปาโมชฺชตฺถํ สนฺทสฺเสติ, นีวรเณ ปชหนฺโต ปฐม\-ฌา\-นตฺถํ สนฺทสฺเสติ ฯเปฯ สพฺพกิเลเส ปชหนฺโต อรหตฺตมคฺคตฺถํ สนฺทสฺเสติ. ตํญาตฏฺเฐน ญาณํ, ปชานนฏฺเฐน ปญฺญา. เตน วุจฺจติ นานา\-ธมฺม\-ปกาส\-นตา ปญฺญา อตฺถ\-สนฺทสฺสเน ญาณํ.}

\vspace{\tipitakaparskip}
\csromancenter{อตฺถ\-สนฺทสฺสน\-ญาณ\-นิทฺเทโส นวติํสติโม.}
\csromanmatikaanchor{mtk.33}
\csromantocmarkref{subsection}{40. ทสฺสนวิสุทฺธิญาณนิทฺเทส}{mtk.33}
\bookmark[dest={mtk.33},level=2]{40. ทสฺสนวิสุทฺธิญาณนิทฺเทส}
\csromanheader{2}{40. ทสฺสน\-วิสุทฺธิญาณ\-นิทฺเทส}
\proseitem{91}{\csromanfolio{102}กถํ สพฺพธมฺมานํ เอก\-สงฺค\-หตา\-นานตฺเต\-กตฺต\-ปฏิเวเธ ปญฺญา ทสฺสน\-วิสุทฺธิญาณํ\csromandash{}สพฺพ\-ธมฺมานนฺติ ปญฺจกฺขนฺธา ฯเปฯ อปริยาปนฺนา ธมฺมา.}

\prose{เอกสงฺคหตาติ ทฺวาทสหิ อากาเรหิ สพฺเพ ธมฺมา เอกสงฺคหิตา, ตถฏฺเฐน อนตฺตฏฺเฐน สจฺจฏฺเฐน ปฏิเวธ\-ฏฺเฐน อภิชาน\-นฏฺเฐน ปริชาน\-นฏฺเฐน ธมฺมฏฺเฐน ธาตุฏฺเฐน ญาตฏฺเฐน สจฺฉิกิริยฏฺเฐน ผุสนฏฺเฐน อภิสมย\-ฏฺเฐน. อิเมหิ ทฺวาทสหิ อากาเรหิ สพฺเพ ธมฺมา เอกสงฺคหิตา.}

\prose{\textbf{นานตฺเตกตฺตนฺ}ติ กามจฺฉนฺโท นานตฺตํ, เนกฺขมฺมํ เอกตฺตํ ฯเปฯ สพฺพกิเลสา นานตฺตํ, อรหตฺตมคฺโค เอกตฺตํ.}

\prose{ปฏิเวเธติ ทุกฺขสจฺจํ ปริญฺญา\-ปฏิเวธํ ปฏิวิชฺฌติ, สมุทยสจฺจํ ปหาน\-ปฏิเวธํ ปฏิวิชฺฌติ, นิโรธสจฺจํ สจฺฉิกิริยา\-ปฏิเวธํ ปฏิวิชฺฌติ, มคฺคสจฺจํ ภาวนา\-ปฏิเวธํ ปฏิวิชฺฌติ.}

\prosewithcloserruled{ทสฺสน\-วิสุทฺธีติ โสตาปตฺติมคฺค\-กฺขเณ ทสฺสนํ วิสุชฺฌติ, โสตาปตฺติ\-ผลกฺขเณ ทสฺสนํ วิสุทฺธํ. สกทาคามิ\-มคฺคกฺขเณ ทสฺสนํ วิสุชฺฌติ, สกทาคามิ\-ผลกฺขเณ ทสฺสนํ วิสุทฺธํ, อนาคามิ\-มคฺคกฺขเณ ทสฺสนํ วิสุชฺฌติ, อนาคามิ\-ผลกฺขเณ ทสฺสนํ วิสุทฺธํ. อรหตฺต\-มคฺคกฺขเณ ทสฺสนํ วิสุชฺฌติ, อรหตฺต\-ผลกฺขเณ ทสฺสนํ วิสุทฺธํ. ตํญาตฏฺเฐน ญาณํ ปชานนฏฺเฐน ปญฺญา. เตน วุจฺจติ สพฺพธมฺมานํ เอก\-สงฺค\-หตา\-นานตฺเต\-กตฺต\-ปฏิเวเธ ปญฺญา ทสฺสน\-วิสุทฺธิญาณํ.}{ทสฺสนวิสุทฺธิญาณนนิทฺเทโส จตฺตาลีสโม.}
\csromanmatikaanchor{mtk.34}
\csromantocmarkref{subsection}{41. ขนฺติญาณนิทฺเทส}{mtk.34}
\bookmark[dest={mtk.34},level=2]{41. ขนฺติญาณนิทฺเทส}
\csromanheader{2}{41. ขนฺติ\-ญาณ\-นิทฺเทส}
\proseitem{92}{กถํ วิทิตตฺตา ปญฺญา ขนฺติญาณํ\csromandash{}รูปํ อนิจฺจโต วิทิตํ, รูปํ ทุกฺขโต วิทิตํ, รูปํ อนตฺตโต วิทิตํ. ยํ ยํ วิทิตํ, ตํ ตํ ขมตีติ วิทิตตฺตา ปญฺญา ขนฺติญาณํ. เวทนา ฯเปฯ. สญฺญา. สงฺขารา. วิญฺญาณํ. จกฺขุ ฯเปฯ. ชรามรณํ อนิจฺจโต วิทิตํ, ชรามรณํ ทุกฺขโต วิทิตํ, อรามรณํ อนตฺตโต วิทิตํ. ยํ ยํ วิทิตํ, ตํ ตํ ขมตีติ วิทิตตฺตา ปญฺญา ขนฺติญาณํ. \csromanfolio{103}ตํญาตฏฺเฐน ญาณํ, ปชานนฏฺเฐน ปญฺญา. เตน วุจฺจติ วิทิตตฺตา ปญฺญา ขนฺติญาณํ. ขนฺติ\-ญาณ\-นิทฺเทโส เอกจตฺตาลีสโม.}

\csromanmatikaanchor{mtk.35}
\csromantocmarkref{subsection}{42. ปริโยคาหณญาณนิทฺเทส}{mtk.35}
\bookmark[dest={mtk.35},level=2]{42. ปริโยคาหณญาณนิทฺเทส}
\csromanheader{2}{42. \textbf{ปริโยคาหณญาณนิทฺเทส}}
\proseitemwithcloserruled{93}{กถํ ผุฏฺฐตฺตา ปญฺญา ปริโยคาหเณ ญาณํ\csromandash{}รูปํ อนิจฺจโต ผุสติ, รูปํ ทุกฺขโต ผุสติ, รูปํ อนตฺตโต ผุสติ. ยํ ยํ ผุสติ, ตํ ตํ ปริโยคหตีติ ผุฏฺฐตฺตา ปญฺญา ปริโยคาหเณ ญาณํ. เวทนํ ฯเปฯ. สญฺญํ. สงฺขาเร. วิญฺญาณํ. จกฺขุํ ฯเปฯ. ชรามรณํ อนิจฺจโต ผุสติ, ทุกฺขโต ผุสติ, อนตฺตโต ผุสติ. ยํ ยํ ผุสติ, ตํ ตํ ปริโยคหตีติ ผุฏฺฐตฺตา ปญฺญา ปริโยคาหเณ ญาณํ. ตํญาตฏฺเฐน ญาณํ. ปชานนฏฺเฐน ปญฺญา. เตน วุจฺจติ ผุฏฺฐตฺตา ปญฺญา ปริโยคาหเณ ญาณํ.}{ปริโยคาหณญาณนิทฺเทโส เทฺวจตฺตาลีสโม.}
\csromanmatikaanchor{mtk.36}
\csromantocmarkref{subsection}{43. ปเทสวิหารญาณนิทฺเทส}{mtk.36}
\bookmark[dest={mtk.36},level=2]{43. ปเทสวิหารญาณนิทฺเทส}
\csromanheader{2}{43. ปเทส\-วิหาร\-ญาณ\-นิทฺเทส}
\proseitem{94}{กถํ สโมทหเน ปญฺญา ปเทสวิหาเร ญาณํ\csromandash{}มิจฺฉาทิฏฺฐิ ปจฺจยาปิ เวทยิตํ, มิจฺฉาทิฏฺฐิวูปสมปจฺจยาปิ เวทยิตํ. สมฺมาทิฏฺฐิ\-ปจฺจยาปิ เวทยิตํ, สมฺมาทิฏฺฐิ\-วูปสม\-ปจฺจยาปิ เวทยิตํ. มิจฺฉาสงฺกปฺป\-ปจฺจยาปิ เวทยิตํ, มิจฺฉาสงฺกปฺป\-วูปสม\-ปจฺจยาปิ เวทยิตํ. สมฺมาสงฺกปฺป\-ปจฺจยาปิ เวทยิตํ, สมฺมาสงฺกปฺป\-วูปสม\-ปจฺจยาปิ เวทยิตํ ฯเปฯ. มิจฺฉาวิมุตฺติ\-ปจฺจยาปิ เวทยิตํ, มิจฺฉาวิมุตฺติ\-วูปสม\-ปจฺจยาปิ เวทยิตํ. สมฺมาวิมุตฺติ\-ปจฺจยาปิ เวทยิตํ, สมฺมาวิมุตฺติ\-วูปสม\-ปจฺจยาปิ เวทยิตํ. ฉนฺทปจฺจยาปิ เวทยิตํ, ฉนฺท\-วูปสม\-ปจฺจยาปิ เวทยิตํ. วิตกฺกปจฺจยาปิ เวทยิตํ. วิตกฺก\-วูปสม\-ปจฺจยาปิ เวทยิตํ. สญฺญาปจฺจยาปิ เวทยิตํ, สญฺญา\-วูปสม\-ปจฺจยาปิ เวทยิตํ.}

\prosewithcloserruled{ฉนฺโท จ อวูปสนฺโต โหติ, วิตกฺโก จ อวูปสนฺโต โหติ, สญฺญา จ อวูปสนฺตา โหติ, ตปฺปจฺจยาปิ เวทยิตํ. ฉนฺโท จ \csromanfolio{104}วูปสนฺโต โหติ, วิตกฺโก จ อวูปสนฺโต โหติ, สญฺญา จ อวูปสนฺตา โหติ, ตปฺปจฺจยาปิ เวทยิตํ. ฉนฺโท จ วูปสนฺโต โหติ, วิตกฺโก จ วูปสนฺโต โหติ, สญฺญา จ อวูปสนฺตา โหติ, ตปฺปจฺจยาปิ เวทยิตํ. ฉนฺโท จ วูปสนฺโต โหติ, วิตกฺโก จ วูปสนฺโต โหติ, สญฺญา จ วูปสนฺตา โหติ. ตปฺปจฺจยาปิ เวทยิตํ. อปฺปตฺตสฺส ปตฺติยา อตฺถิ อายวํ, ตสฺมิมฺปิ ฐาเน อนุปฺปตฺเต ตปฺปจฺจยาปิ เวทยิตํ. ตํญาตฏฺเฐน ญาณํ, ปชานนฏฺเฐน ปญฺญา. เตน วุจฺจติ สโมทหเน ปญฺญา ปเทสวิหาเร ญาณํ.}{ปเทส\-วิหาร\-ญาณ\-นิทฺเทโส เตจตฺตาลีสโม.}
\csromanmatikaanchor{mtk.37}
\csromantocmarkref{subsection}{44-49. ฉวิวฏฺฏญาณนิทฺเทส}{mtk.37}
\bookmark[dest={mtk.37},level=2]{44-49. ฉวิวฏฺฏญาณนิทฺเทส}
\csromanheader{2}{44-49. ฉวิวฏฺฏ\-ญาณ\-นิทฺเทส}
\proseitem{95}{กถํ อธิปตตฺตา ปญฺญา สญฺญาวิวฏฺเฏ ญาณํ\csromandash{} เนกฺขมฺมา\-ธิป\-ตตฺตา ปญฺญา กามจฺฉนฺทโต สญฺญาย วิวฏฺฏตีติ อธิปตตฺตา ปญฺญา สญฺญาวิวฏฺเฏ ญาณํ, อพฺยาปาทา\-ธิป\-ตตฺตา ปญฺญา พฺยาปาทโต สญฺญาย วิวฏฺฏตีติ อธิปตตฺตา ปญฺญา สญฺญาวิวฏฺเฏ ญาณํ, อาโลกสญฺญา\-ธิป\-ตตฺตา ปญฺญา ถินมิทฺธโต สญฺญาย วิวฏฺฏตีติ อธิปตตฺตา ปญฺญา สญฺญาวิวฏฺเฏ ญาณํ, อวิกฺเขปา\-ธิป\-ตตฺตา ปญฺญา อุทฺธจฺจโต สญฺญาย วิวฏฺฏตีติ อธิปตตฺตา ปญฺญา สญฺญาวิวฏฺเฏ ญาณํ, ธมฺมววตฺถานา\-ธิป\-ตตฺตา ปญฺญา วิจิกิจฺฉาย สญฺญาย วิวฏฺฏตีติ อธิปตตฺตา ปญฺญา สญฺญาวิวฏฺเฏ ญาณํ, ญาณาธิปตตฺตา ปญฺญา อวิชฺชาย สญฺญาย วิวฏฺฏตีติ อธิปตตฺตา ปญฺญา สญฺญาวิวฏฺเฏ ญาณํ. ปาโมชฺชา\-ธิป\-ตตฺตา ปญฺญา อรติยา สญฺญาย วิวฏฺฏตีติ อธิปตตฺตา ปญฺญา สญฺญาวิวฏฺเฏ ญาณํ, ปฐมชฺฌานา\-ธิป\-ตตฺตา ปญฺญา นีวรเณหิ สญฺญาย วิวฏฺฏตีติ อธิปตตฺตา ปญฺญา สญฺญาวิวฏฺเฏ ญาณํ ฯเปฯ อรหตฺตมคฺคา\-ธิป\-ตตฺตา ปญฺญา สพฺพกิเลเสหิ สญฺญาย วิวฏฺฏตีติ อธิปตตฺตา ปญฺญา สญฺญาวิวฏฺเฏ ญาณํ. ตํญาตฏฺเฐน ญาณํ, ปชานนฏฺเฐน ปญฺญา. เตน วุจฺจติ อธิปตตฺตา ปญฺญา สญฺญาวิวฏฺเฏ ญาณํ.}

\proseitem{96}{กถํ นานตฺเต ปญฺญา เจโตวิวฏฺเฏ ญาณํ\csromandash{}กามจฺฉนฺโท นานตฺตํ, เนกฺขมฺมํ เอกตฺตํ. เนกฺขมฺเม\-กตฺตํ เจตยโต กามจฺฉนฺทโต จิตฺตํ วิวฏฺฏตีติ \csromanfolio{105}นานตฺเต ปญฺญา เจโตวิวฏฺเฏ ญาณํ. พฺยาปาโท นานตฺตํ, อพฺยาปาโท เอกตฺตํ. อพฺยาปาเทกตฺตํ เจตยโต พฺยาปาทโต จิตฺตํ วิวฏฺฏตีติ นานตฺเต ปญฺญา เจโตวิวฏฺเฏ ญาณํ. ถินมิทฺธํ นานตฺตํ, อาโลกสญฺญา เอกตฺตํ. อาโลกสญฺเญ\-กตฺตํ เจตยโต ถินมิทฺธโต จิตฺตํ วิวฏฺฏตีติ นานตฺเต ปญฺญา เจโตวิวฏฺเฏ ญาณํ ฯเปฯ. สพฺพกิเลสา นานตฺตํ, อรหตฺตมคฺโค เอกตฺตํ. อรหตฺตมคฺเค\-กตฺตํ เจตยโต สพฺพกิเลเสหิ จิตฺตํ วิวฏฺฏตีติ นานตฺเต ปญฺญา เจโตวิวฏฺเฏ ญาณํ, ตํญาตฏฺเฐน ญาณํ, ปชานนฏฺเฐน ปญฺญา. เตน วุจฺจติ นานตฺเต ปญฺญา เจโตวิวฏฺเฏ ญาณํ.}

\proseitem{97}{กถํ อธิฏฺฐาเน ปญฺญา จิตฺตวิวฏฺเฏ ญาณํ\csromandash{}กามจฺฉนฺทํ ปชหนฺโต เนกฺขมฺม\-วเสน จิตฺตํ อธิฏฺฐาตีติ อธิฏฺฐาเน ปญฺญา จิตฺตวิวฏฺเฏ ญาณํ. พฺยาปาทํ ปชหนฺโต อพฺยาปาทวเสน จิตฺตํ อธิฏฺฐาตีติ อธิฏฺฐาเน ปญฺญา จิตฺตวิวฏฺเฏ ญาณํ. ถินมิทฺธํ ปชหนฺโต อาโลกสญฺญา\-วเสน จิตฺตํ อธิฏฺฐาตีติ อธิฏฺฐาเน ปญฺญา จิตฺตวิวฏฺเฏ ญาณํ ฯเปฯ. สพฺพกิเลเส ปชหนฺโต อรหตฺตมคฺค\-วเสน จิตฺตํ อธิฏฺฐาตีติ อธิฏฺฐาเน ปญฺญา จิตฺตวิวฏฺเฏ ญาณํ, ตํ ญาตฏฺเฐน ญาณํ, ปชานนฏฺเฐน ปญฺญา. เตน วุจฺจติ อธิฏฺฐาเน ปญฺญา จิตฺตวิวฏฺเฏ ญาณํ.}

\proseitem{98}{กถํ สุญฺญเต ปญฺญา ญาณวิวฏฺเฏ ญาณํ\csromandash{}จกฺขุ สุญฺญํ อตฺเตน วา อตฺตนิเยน วา นิจฺเจน วา ธุเวน วา สสฺสเตน วา อวิปริณาม\-ธมฺเมน วาติ ยถาภูตํ ชานโต\csromansharedfootnote{313ceeae43e8dd11}{ปชานโต (สฺยา)} ปสฺสโต จกฺขา\-ภินิเวสโต\csromansharedfootnote{2b40c0ed467be222}{กามาภินิเวสโต (สฺยา) อฏฺฐกถา โอโลเกตพฺพา.} ญาณํ วิวฏฺฏตีติ สุญฺญเต ปญฺญา ญาณวิวฏฺเฏ ญาณํ. โสตํ สุญฺญํ ฯเปฯ. ฆานํ สุญฺญํ. ชิวฺหา สุญฺญา. กาโย สุญฺโญ. มโน สุญฺโญ อตฺเตน วา อตฺตนิเยน วา นิจฺเจน วา ธุเวน วา สสฺสเตน วา อวิปริณาม\-ธมฺเมน วาติ ยถาภูตํ ชานโต ปสฺสโต มนาภิ\-นิเวสโต ญาณํ วิวฏฺฏตีติ สุญฺญเต ปญฺญา ญาณวิวฏฺเฏ ญาณํ, ตํญาตฏฺเฐน ญาณํ, ปชานนฏฺเฐน ปญฺญา. เตน วุจฺจติ สุญฺญเต ปญฺญา ญาณวิวฏฺเฏ ญาณํ.}

\proseitem{99}{กถํ โวสคฺเค ปญฺญา วิโมกฺข\-วิวฏฺเฏ ญาณํ\csromandash{}เนกฺขมฺเมน กามจฺฉนฺทํ โวสชฺชตีติ โวสคฺเค ปญฺญา วิโมกฺข\-วิวฏฺเฏ ญาณํ, อพฺยาปาเทน พฺยาปาทํ โวสชฺชตีติ โวสคฺเค ปญฺญา วิโมกฺข\-วิวฏฺเฏ ญาณํ, อาโลกสญฺญาย ถินมิทฺธํ โวสชฺชตีติ โวสคฺเค ปญฺญา วิโมกฺข\-วิวฏฺเฏ \csromanfolio{106}ญาณํ, อวิกฺเขเปน อุทฺธจฺจํ โวสชฺชตีติ โวสคฺเค ปญฺญา วิโมกฺข\-วิวฏฺเฏ ญาณํ, ธมฺม\-ววตฺถาเนน วิจิกิจฺฉํ โวสชฺชตีติ โวสคฺเค ปญฺญา วิโมกฺข\-วิวฏฺเฏ ญาณํ ฯเปฯ อรหตฺต\-มคฺเคน สพฺพกิเลเส โวสชฺชตีติ โวสคฺเค ปญฺญา วิโมกฺข\-วิวฏฺเฏ ญาณํ, ตํญาตฏฺเฐน ญาณํ, ปชานนฏฺเฐน ปญฺญา. เตน วุจฺจติ โวสคฺเค ปญฺญา วิโมกฺข\-วิวฏฺเฏ ญาณํ.}

\proseitem{100}{กถํ ตถฏฺเฐ ปญฺญา สจฺจวิวฏฺเฏ ญาณํ\csromandash{}ทุกฺขสฺส ปีฬนฏฺฐํ สงฺขตฏฺฐํ สนฺตาปฏฺฐํ วิปริณามฏฺฐํ ปริชานนฺโต วิวฏฺฏตีติ ตถฏฺเฐ ปญฺญา สจฺจวิวฏฺเฏ ญาณํ, สมุทยสฺส อายูหนฏฺฐํ นิทานฏฺฐํ สญฺโญคฏฺฐํ ปลิโพธฏฺฐํ ปชหนฺโต วิวฏฺฏตีติ ตถฏฺเฐ ปญฺญา สจฺจวิวฏฺเฏ ญาณํ, นิโรธสฺส นิสฺสรณฏฺฐํ วิเวกฏฺฐํ อสงฺขตฏฺฐํ อมตฏฺฐํ สจฺฉิกโรนฺโต วิวฏฺฏตีติ ตถฏฺเฐ ปญฺญา สจฺจวิวฏฺเฏ ญาณํ, มคฺคสฺส นิยฺยานฏฺฐํ เหตุฏฺฐํ ทสฺสนฏฺฐํ อาธิปเตยฺยฏฺฐํ ภาเวนฺโต วิวฏฺฏตีติ ตถฏฺเฐ ปญฺญา สจฺจวิวฏฺเฏ ญาณํ.}

\prose{สญฺญาวิวฏฺโฏ เจโตวิวฏฺโฏ จิตฺตวิวฏฺโฏ ญาณวิวฏฺโฏ วิโมกฺข\-วิวฏฺโฏ สจฺจวิวฏฺโฏ. สญฺชานนฺโต วิวฏฺฏตีติ สญฺญาวิวฏฺโฏ, เจตยนฺโต วิวฏฺฏตีติ เจโตวิวฏฺโฏ, วิชานนฺโต วิวฏฺฏตีติ จิตฺตวิวฏฺโฏ, ญาณํ กโรนฺโต วิวฏฺฏตีติ ญาณวิวฏฺโฏ, โวสชฺชนฺโต วิวฏฺฏตีติ วิโมกฺข\-วิวฏฺโฏ, ตถฏฺเฐ วิวฏฺฏตีติ สจฺจวิวฏฺโฏ.}

\prose{ยตฺถ สญฺญาวิวฏฺโฏ, ตตฺถ เจโตวิวฏฺโฏ. ยตฺถ เจโตวิวฏฺโฏ, ตตฺถ สญฺญาวิวฏฺโฏ. ยตฺถ สญฺญาวิวฏฺโฏ เจโตวิวฏฺโฏ, ตตฺถ จิตฺตวิวฏฺโฏ. ยตฺถ จิตฺตวิวฏฺโฏ, ตตฺถ สญฺญาวิวฏฺโฏ เจโตวิวฏฺโฏ. ยตฺถ สญฺญาวิวฏฺโฏ เจโตวิวฏฺโฏ จิตฺตวิวฏฺโฏ, ตตฺถ ญาณวิวฏฺโฏ. ยตฺถ ญาณวิวฏฺโฏ, ตตฺถ สญฺญาวิวฏฺโฏ เจโตวิวฏฺโฏ จิตฺตวิวฏฺโฏ. ยตฺถ สญฺญาวิวฏฺโฏ เจโตวิวฏฺโฏ จิตฺตวิวฏฺโฏ ญาณวิวฏฺโฏ, ตตฺถ วิโมกฺข\-วิวฏฺโฏ. ยตฺถ วิโมกฺข\-วิวฏฺโฏ, ตตฺถ สญฺญาวิวฏฺโฏ เจโตวิวฏฺโฏ จิตฺตวิวฏฺโฏ ญาณวิวฏฺโฏ. ยตฺถ สญฺญาวิวฏฺโฏ เจโตวิวฏฺโฏ จิตฺตวิวฏฺโฏ ญาณวิวฏฺโฏ วิโมกฺข\-วิวฏฺโฏ, ตตฺถ สจฺจวิวฏฺโฏ. ยตฺถ สจฺจวิวฏฺโฏ, ตตฺถ สญฺญาวิวฏฺโฏ เจโตวิวฏฺโฏ จิตฺตวิวฏฺโฏ ญาณวิวฏฺโฏ วิโมกฺข\-วิวฏฺโฏ. ตํญาตฏฺเฐน ญาณํ, ปชานนฏฺเฐน ปญฺญา. เตน วุจฺจติ ตถฏฺเฐ ปญฺญา สจฺจวิวฏฺเฏ ญาณํ.}

\vspace{\tipitakaparskip}
\csromancenter{ฉวิวฏฺฏ\-ญาณ\-นิทฺเทโส นวจตฺตาลีสโม.}
\csromanmatikaanchor{mtk.38}
\csromantocmarkref{subsection}{50. อิทฺธิวิธญาณนิทฺเทส}{mtk.38}
\bookmark[dest={mtk.38},level=2]{50. อิทฺธิวิธญาณนิทฺเทส}
\csromanheader{2}{50. อิทฺธิวิธญาณ\-นิทฺเทส}
\proseitem{101}{\csromanfolio{107}กถํ กายมฺปิ จิตฺตมฺปิ เอก\-ววตฺถานตา สุขสญฺญญฺจ ลหุสญฺญญฺจ อธิฏฺฐาน\-วเสน อิชฺฌนฏฺเฐ ปญฺญา อิทฺธิวิเธ ญาณํ\csromandash{}อิธ ภิกฺขุ ฉนฺท\-สมาธิ\-ปธาน\-สงฺขาร\-สมนฺนาคตํ อิทฺธิปาทํ ภาเวติ, วีริย\-สมาธิ\-ปธาน\-สงฺขาร\-สมนฺนาคตํ อิทฺธิปาทํ ภาเวติ, จิตฺต\-สมาธิ\-ปธาน\-สงฺขาร\-สมนฺนาคตํ อิทฺธิปาทํ ภาเวติ, วีมํสา\-สมาธิ\-ปธาน\-สงฺขาร\-สมนฺนาคตํ อิทฺธิปาทํ ภาเวติ. โส อิเมสุ จตูสุ อิทฺธิปาเทสุ จิตฺตํ ปริภาเวติ ปริทเมติ, มุทุํ กโรติ กมฺมนิยํ. โส อิเมสุ จตูสุ อิทฺธิปาเทสุ จิตฺตํ ปริภาเวตฺวา ปริทเมตฺวา มุทุํ กริตฺวา กมฺมนิยํ กายมฺปิ จิตฺเต สโมทหติ, จิตฺตมฺปิ กาเย สโมทหติ, กายวเสน จิตฺตํ ปริณาเมติ, จิตฺตวเสน กายํ ปริณาเมติ, กายวเสน จิตฺตํ อธิฏฺฐาติ, จิตฺตวเสน กายํ อธิฏฺฐาติ, กายวเสน จิตฺตํ ปริณาเมตฺวา จิตฺตวเสน กายํ ปริณาเมตฺวา กายวเสน จิตฺตํ อธิฏฺฐหิตฺวา จิตฺตวเสน กายํ อธิฏฺฐหิตฺวา สุขสญฺญญฺจ ลหุสญฺญญฺจ กาเย โอกฺกมิตฺวา วิหรติ, โส ตถาภาวิเตน จิตฺเตน ปริสุทฺเธน ปริโยทาเตน อิทฺธิวิธ\-ญาณาย จิตฺตํ อภินีหรติ อภินินฺนาเมติ, โส อเนกวิหิตํ อิทฺธิวิธํ ปจฺจนุโภติ.}

\proseitem{102}{เอโกปิ หุตฺวา พหุธา โหติ, พหุธาปิ หุตฺวา เอโก โหติ, อาวิภาวํ ติโรภาวํ ติโรกุฏฺฏํ\csromansharedfootnote{d504a26cf43fb05d}{ติโรกุฑฺฑํ (สฺยา)} ติโรปาการํ ติโรปพฺพตํ อสชฺชมาโน คจฺฉติ เสยฺยถาปิ อากาเส. ปถวิยาปิ อุมฺมุชฺช\-นิมุชฺชํ กโรติ เสยฺยถาปิ อุทเก. อุทเกปิ อภิชฺชมาเน คจฺฉติ เสยฺยถาปิ ปถวิยํ. อากาเสปิ ปลฺลงฺเกน กมติ\csromansharedfootnote{27dd9f299094b980}{จงฺกมติ (สฺยา) ที 1. 74 ปิฏฺเฐ ปสฺสิตพฺพา.} เสยฺยถาปิ ปกฺขี สกุโณ. อิเมปิ จนฺทิมสูริเย เอวํมหิทฺธิเก เอวํ\-มหา\-นุภาเว ปาณินา ปรามสติ\csromansharedfootnote{d7cf39cf20e401b0}{ปริมสติ (สฺยา)} ปริมชฺชติ, ยาว พฺรหฺมโลกาปิ กาเยน วสํ วตฺเตติ. ตํญาตฏฺเฐน ญาณํ, ปชานนฏฺเฐน ปญฺญา. เตน วุจฺจติ กามมฺปิ จิตฺตมฺปิ เอก\-ววตฺถานตา สุขสญฺญญฺจ ลหุสญฺญญฺจ อธิฏฺฐาน\-วเสน อิชฺฌนฏฺเฐ ปญฺญา อิทฺธิวิเธ ญาณํ.}

\vspace{\tipitakaparskip}
\csromancenter{อิทฺธิวิธญาณ\-นิทฺเทโส ปญฺญาสโม.}
\csromanmatikaanchor{mtk.39}
\csromantocmarkref{subsection}{51. โสตธาตุวิสุทฺธิญาณนิทฺเทส}{mtk.39}
\bookmark[dest={mtk.39},level=2]{51. โสตธาตุวิสุทฺธิญาณนิทฺเทส}
\csromanheader{2}{51. โสตธาตุ\-วิสุทฺธิญาณ\-นิทฺเทส}
\proseitemwithcloserruled{103}{\csromanfolio{108}กถํ วิตกฺก\-วิปฺผาร\-วเสน นานตฺเต\-กตฺต\-สทฺทนิมิตฺตานํ ปริโยคาหเณ ปญฺญา โสตธาตุ\-วิสุทฺธิญาณํ\csromandash{}อิธ ภิกฺขุ ฉนฺทสมาธิ ฯเปฯ วีริยสมาธิ. จิตฺตสมาธิ. วีมํสาสมาธิ\-ปธานสงฺขาร- สมนฺนาคตํ อิทฺธิปาทํ ภาเวติ. โส อิเมสุ จตูสุ อิทฺธิปาเทสุ จิตฺตํ ปริภาเวติ ปริทเมติ, มุทุํ กโรติ กมฺมนิยํ. โส อิเมสุ จตูสุ อิทฺธิปาเทสุ จิตฺตํ ปริภาเวตฺวา ปริทเมตฺวา มุทุํ กริตฺวา กมฺมนิยํ ทูเรปิ สทฺทานํ สทฺทนิมิตฺตํ มนสิ กโรติ, สนฺติเกปิ สทฺทานํ สทฺทนิมิตฺตํ มนสิ กโรติ, โอฬาริกานมฺปิ สทฺทานํ สทฺทนิมิตฺตํ มนสิ กโรติ, สุขุมานมฺปิ สทฺทานํ สทฺทนิมิตฺตํ มนสิ กโรติ, สณฺหสณฺหานมฺปิ สทฺทานํ สทฺทนิมิตฺตํ มนสิ กโรติ, ปุรตฺถิมายปิ ทิสาย สทฺทานํ สทฺทนิมิตฺตํ มนสิ กโรติ, ปจฺฉิมายปิ ทิสาย สทฺทานํ สทฺทนิมิตฺตํ มนสิ กโรติ, อุตฺตรายปิ ทิสาย สทฺทานํ สทฺทนิมิตฺตํ มนสิ กโรติ, ทกฺขิณายปิ ทิสาย สทฺทานํ สทฺทนิมิตฺตํ มนสิ กโรติ, ปุรตฺถิมายปิ อนุทิสาย สทฺทานํ สทฺทนิมิตฺตํ มนสิ กโรติ, ปจฺฉิมายปิ อนุทิสาย สทฺทานํ สทฺทนิมิตฺตํ มนสิ กโรติ, อุตฺตรายปิ อนุทิสาย สทฺทานํ สทฺทนิมิตฺตํ มนสิ กโรติ, ทกฺขิณายปิ อนุทิสาย สทฺทานํ สทฺทนิมิตฺตํ มนสิ กโรติ, เหฏฺฐิมยปิ ทิสาย สทฺทานํ สทฺทนิมิตฺตํ มนสิ กโรติ, อุปริมายปิ ทิสาย สทฺทานํ สทฺทนิมิตฺตํ มนสิ กโรติ. โส ตถาภาวิเตน จิตฺเตน ปริสุทฺเธน ปริโยทาเตน โสตธาตุ\-วิสุทฺธิญาณาย จิตฺตํ อภินีหรติ อภินินฺนาเมติ, โส ทิพฺพาย โสตธาตุยา วิสุทฺธาย อติกฺกนฺตมานุสิกาย อุโภ สทฺเท สุณาติ ทิพฺเพ จ มานุเส จ เย ทูเร สนฺติเก จ. ตํญาตฏฺเฐน ญาณํ, ปชานนฏฺเฐน ปญฺญา. เตน วุจฺจติ วิตกฺก\-วิปฺผาร\-วเสน นานตฺเต\-กตฺต\-สทฺทนิมิตฺตานํ ปริโยคาหเณ ปญฺญา โสตธาตุ\-วิสุทฺธิญาณํ.}{โสตธาตุ\-วิสุทฺธิญาณ\-นิทฺเทโส เอกปญฺญาสโม.}
\csromanmatikaanchor{mtk.40}
\csromantocmarkref{subsection}{52. เจโตปริยญาณนิทฺเทส}{mtk.40}
\bookmark[dest={mtk.40},level=2]{52. เจโตปริยญาณนิทฺเทส}
\csromanheader{2}{52. เจโตปริยญาณ\-นิทฺเทส}
\proseitemwithcloserruled{104}{กถํ ติณฺณํ จิตฺตานํ วิปฺผารตฺตา อินฺทฺริยานํ ปสาทวเสน นานตฺเต\-กตฺต\-วิญฺญาณจริยา\-ปริโยคาหเณ ปญฺญา เจโตปริย\-ญาณํ\csromandash{}อิธ \csromanfolio{109}ภิกฺขุ ฉนฺท\-สมาธิ\-ปธาน\-สงฺขาร\-สมนฺนาคตํ อิทฺธิปาทํ ภาเวติ, วีริยสมาธฺ ฯเปฯ จิตฺตสมาธิ ฯเปฯ วีมํสาสมาธิ\-ปธานสงฺขาร- สมนฺนาคตํ อิทฺธิปาทํ ภาเวติ. โส อิเมสุ จตูสุ อิทฺธิปาเทสุ จิตฺตํ ปริภาเวติ ปริทเมติ, มุทุํ กโรติ กมฺมนิยํ. โส อิเมสุ จตูสุ อิทฺธิปาเทสุ จิตฺตํ ปริภาเวตฺวา ปริทเมตฺวา มุทุํ กริตฺวา กมฺมนิยํ เอวํ ปชานาติ “อิทํ รูปํ โสมนสฺสินฺทฺริย\-สมุฏฺฐิตํ, อิทํ รูปํ โทมนสฺสินฺทฺริย\-สมุฏฺฐิตํ, อิทํ รูปํ อุเปกฺขินฺทฺริยสมุฏฺฐิตนฺ”ติ. โส ตถาภาวิเตน จิตฺเตน ปริสุทฺเธน ปริโยทาเตน เจโตปริย\-ญาณาย จิตฺตํ อภินีหรติ อภินินฺนาเมติ. โส ปรสตฺตานํ ปรปุคฺคลานํ เจตสา เจโต ปริจฺจ ปชานาติ สราคํ วา จิตฺตํ “สราคํ จิตฺตนฺ”ติ ปชานาติ, วีตราคํ วา จิตฺตํ “วีตราคํ จิตฺตนฺ”ติ ปชานาติ, สโทสํ วา จิตฺตํ ฯเปฯ วีตโทสํ วา จิตฺตํ. สโมหํ วา จิตฺตํ. วีตโมหํ วา จิตฺตํ. สํขิตฺตํ วา จิตฺตํ. วิกฺขิตฺตํ วา จิตฺตํ. มหคฺคตํ วา จิตฺตํ. อมหคฺคตํ วา จิตฺตํ. สอุตฺตรํ วา จิตฺตํ. อนุตฺตรํ วา จิตฺตํ. สมาหิตํ วา จิตฺตํ. อสมาหิตํ วา จิตฺตํ. วิมุตฺตํ วา จิตฺตํ. อวิมุตฺตํ วา จิตฺตํ “อวิมุตฺตํ จิตฺตนฺ”ติ ปชานาติ. ตํญาตฏฺเฐน ญาณํ, ปชานนฏฺเฐน ปญฺญา. เตน วุจฺจติ ติณฺณํ จิตฺตานํ วิปฺผารตฺตา อินฺทฺริยานํ ปสาทวเสน นานตฺเต\-กตฺต\-วิญฺญาณจริยา\-ปริโยคาหเณ ปญฺญา เจโตปริย\-ญาณํ.}{เจโตปริยญาณ\-นิทฺเทโส เทฺวปญฺญาสโม.}
\csromanmatikaanchor{mtk.41}
\csromantocmarkref{subsection}{53. ปุพฺเพนิวาสานุสฺสติญาณนิทฺเทส}{mtk.41}
\bookmark[dest={mtk.41},level=2]{53. ปุพฺเพนิวาสานุสฺสติญาณนิทฺเทส}
\csromanheader{2}{53. ปุพฺเพ\-นิวาสา\-นุสฺสติ\-ญาณนิทฺเทส}
\proseitem{105}{กถํ ปจฺจย\-ปวตฺตานํ ธมฺมานํ นานตฺเต\-กตฺต\-กมฺม\-วิปฺผาร\-วเสน ปริโยคาหเณ ปญฺญา ปุพฺเพ\-นิวาสา\-นุสฺสติ\-ญาณํ\csromandash{}อิธ ภิกฺขุ ฉนฺทสมาธิ ฯเปฯ มุทุํ กริตฺวา กมฺมนิยํ เอวํ ปชานาติ “อิมสฺมิํ สติ อิทํ โหติ, อิมสฺสุปฺปาทา อิทํ อุปฺปชฺชติ ยทิทํ อวิชฺชาปจฺจยา สงฺขารา, สงฺขาร\-ปจฺจยา วิญฺญาณํ, วิญฺญาณปจฺจยา นามรูปํ, นามรูป\-ปจฺจยา สฬายตนํ, สฬายตน\-ปจฺจยา ผสฺโส, ผสฺสปจฺจยา เวทนา, เวทนาปจฺจยา ตณฺหา, ตณฺหาปจฺจยา อุปาทานํ, อุปาทานปจฺจยา ภโว, ภวปจฺจยา ชาติ, ชาติปจฺจยา ชรามรณํ โสก\-ปริเทว\-ทุกฺข\-โทมนสฺสุ\-ปายาสา สมฺภวนฺติ, เอวเมตสฺส เกวลสฺส ทุกฺข\-กฺขนฺธสฺส สมุทโย โหติ.}

\prosewithcloserruled{\csromanfolio{110}โส ตถาภาวิเตน จิตฺเตน ปริสุทฺเธน ปริโยทาเตน ปุพฺเพ\-นิวาสา\-นุสฺสติ\-ญาณาย จิตฺตํ อภินีหรติ อภินินฺนาเมติ. โส อเนกวิหิตํ ปุพฺเพนิวาสํ อนุสฺสรติ. เสยฺยถิทํ, เอกมฺปิ ชาติํ เทฺวปิ ชาติโย ติสฺโสปิ ชาติโย จตสฺโสปิ ชาติโย ปญฺจปิ ชาติโย ทสปิ ชาติโย วีสมฺปิ ชาติโย ติํสมฺปิ ชาติโย จตฺตาลีสมฺปิ ชาติโย ปญฺญาสมฺปิ ชาติโย, ชาติสตมฺปิ ชาติสหสฺสมฺปิ ชาติ\-สต\-สหสฺสมฺปิ, อเนเกปิ สํวฏฺฏกปฺเป อเนเกปิ วิวฏฺฏกปฺเป อเนเกปิ สํวฏฺฏ\-วิวฏฺฏ\-กปฺเป, อมุตฺราสิํ เอวํนาโม เอวํโคตฺโต เอวํวณฺโณ เอวมาหาโร เอวํ\-สุขทุกฺข\-ปฺปฏิสํเวที เอวมา\-ยุปริยนฺโต. โส ตโต จุโต อมุตฺร อุทปาทิํ. ตตฺราปาสิํ เอวํนาโม เอวํโคตฺโต เอวํวณฺโณ เอวมาหาโร เอวํ\-สุขทุกฺข\-ปฺปฏิสํเวที เอวมา\-ยุปริยนฺโต. โส ตโต จุโต อิธูปปนฺโน”ติ, อิติ สาการํ สอุทฺเทสํ อเนกวิหิตํ ปุพฺเพนิวาสํ อนุสฺสรติ. ตํญาตฏฺเฐน ญาณํ, ปชานนฏฺเฐน ปญฺญา. เตน วุจฺจติ ปจฺจย\-ปวตฺตานํ ธมฺมานํ นานตฺเต\-กตฺต\-กมฺม\-วิปฺผาร\-วเสน ปริโยคาหเณ ปญฺญา ปุพฺเพ\-นิวาสา\-นุสฺสติ\-ญาณํ.}{ปุพฺเพ\-นิวาสา\-นุสฺสติ\-ญาณนิทฺเทโส เตปญฺญาสโม.}
\csromanmatikaanchor{mtk.42}
\csromantocmarkref{subsection}{54. ทิพฺพจกฺขุญาณนิทฺเทส}{mtk.42}
\bookmark[dest={mtk.42},level=2]{54. ทิพฺพจกฺขุญาณนิทฺเทส}
\csromanheader{2}{54. ทิพฺพจกฺขุ\-ญาณนิทฺเทส}
\proseitemwithcloserruled{106}{กถํ โอภาสวเสน นานตฺเต\-กตฺต\-รูปนิมิตฺตานํ ทสฺสนฏฺเฐ ปญฺญา ทิพฺพจกฺขุ\-ญาณํ\csromandash{}อิธ ภิกฺขุ ฉนฺท\-สมาธิ\-ปธาน\-สงฺขาร\-สมนฺนาคตํ อิทฺธิปาทํ ภาเวติ. วีริยสมาธิ ฯเปฯ. จิตฺตสมาธิ ฯเปฯ. วีมํสา\-สมาธิ\-ปธาน\-สงฺขาร\-สมนฺนาคตํ อิทฺธิปาทํ ภาเวติ. โส อิเมสุ จตูสุ อิทฺธิปาเทสุ จิตฺตํ ปริภาเวติ ปริทเมติ, มุทุํ กโรติ กมฺมนิยํ. โส อิเมสุ จตูสุ อิทฺธิปาเทสุ จิตฺตํ ปริภาเวตฺวา ปริทเมตฺวา มุทุํ กริตฺวา กมฺมนิยํ อาโลกสญฺญํ มนสิ กโรติ, ทิวาสญฺญํ อธิฏฺฐาติ, ยถา ทิวา ตถา รตฺติํ, ยถา รตฺติํ ตถา ทิวา. อิติ วิวเฏน เจตสา อปริโยนทฺเธน สปฺปภาสํ จิตฺตํ ภาเวติ. โส ตถาภาวิเตน จิตฺเตน ปริสุทฺเธน ปริโยทาเตน สตฺตานํ จุตูปปาต\-ญาณาย จิตฺตํ อภินีหรติ อภินินฺนาเมติ. โส ทิพฺเพน จกฺขุนา วิสุทฺเธน อติกฺกนฺต\-มานุสเกน สตฺเต ปสฺสติ จวมาเน อุปปชฺชมาเน หีเน ปณีเต สุวณฺเณ \csromanfolio{111}ทุพฺพณฺเณ สุคเต ทุคฺคเต, ยถากมฺมูปเค สตฺเต ปชานาติ “อิเม วต โภนฺโต สตฺตา กาย\-ทุจฺจริเตน สมนฺนาคตา, วจีทุจฺจริเตน สมนฺนาคตา, มโน\-ทุจฺจริเตน สมนฺนาคตา, อริยานํ อุปวาทกา, มิจฺฉาทิฏฺฐิกา มิจฺฉาทิฏฺฐิ\-กมฺม\-สมาทานา, เต กายสฺส เภทา ปรํ มรณา อปายํ ทุคฺคติํ วินิปาตํ นิรยํ อุปปนฺนา. อิเม วา ปน โภนฺโต สตฺตา กายสุจริเตน สมนฺนาคตา, วจีสุจริเตน สมนฺนาคตา, มโนสุจริเตน สมนฺนาคตา, อริยานํ อนุปวาทกา, สมฺมาทิฏฺฐิกา สมฺมา\-ทิฏฺฐิ\-กมฺม\-สมาทานา, เต กายสฺส เภทา ปรํ มรณา สุคติํ สคฺคํ โลกํ อุปปนฺนา”ติ, อิติ ทิพฺเพน จกฺขุนา วิสุทฺเธน อติกฺกนฺต\-มานุสเกน สตฺเต ปสฺสติ จวมาเน อุปปชฺชมาเน หีเน ปณีเต สุวณฺเณ ทุพฺพณฺเณ สุคเต ทุคฺคเต, ยถากมฺมูปเค สตฺเต ปชานาติ. ตํญาตฏฺเฐน ญาณํ, ปชานนฏฺเฐน ปญฺญา. เตน วุจฺจติ โอภาสวเสน นานตฺเต\-กตฺต\-รูปนิมิตฺตานํ ทสฺสงฏฺเฐ ปญฺญา ทิพฺพจกฺขุ\-ญาณํ.}{ทิพฺพจกฺขุญาณ\-นิทฺเทโส จตุปญฺญาสโม.}
\csromanmatikaanchor{mtk.43}
\csromantocmarkref{subsection}{55. อาสวกฺขยญาณนิทฺเทส}{mtk.43}
\bookmark[dest={mtk.43},level=2]{55. อาสวกฺขยญาณนิทฺเทส}
\csromanheader{2}{55. อาสวกฺขย\-ญาณนิทฺเทส}
\proseitem{107}{กถํ จตุสฏฺฐิยา อากาเรหิ ติณฺณนฺนํ อินฺทฺริยานํ วสิภาวตา ปญฺญา อาสวานํ ขเย ญาณํ\csromandash{}กตเมสํ ติณฺณนฺนํ อินฺทฺริยานํ? อนญฺญาต\-ญฺญสฺสามี\-ตินฺทฺริยสฺส อญฺญินฺทฺริยสฺส อญฺญาตาวิ\-นฺทฺริยสฺส.}

\prose{อนญฺญาต\-ญฺญสฺสามี\-ตินฺทฺริยํ กติ ฐานานิ คจฺฉติ, อญฺญินฺทฺริยํ กติ ฐานานิ คจฺฉติ, อญฺญาตาวิ\-นฺทฺริยํ กติ ฐานานิ คจฺฉติ? อนญฺญาต\-ญฺญสฺสามี\-ตินฺทฺริยํ เอกํ ฐานํ คจฺฉติ, โสตาปตฺติ\-มคฺคํ. อญฺญินฺทฺริยํ ฉ ฐานานิ คจฺฉติ, โสตาปตฺติผลํ สกทาคามิ\-มคฺคํ สกทาคามิผลํ อนาคามิมคฺคํ อนาคามิผลํ อรหตฺตมคฺคํ. อญฺญาตาวิ\-นฺทฺริยํ เอกํ ฐานํ คจฺฉติ, อรหตฺตผลํ.}

\prose{โสตาปตฺติมคฺค\-กฺขเณ อนญฺญาต\-ญฺญสฺสามี\-ตินฺทฺริยสฺส สทฺธินฺทฺริยํ อธิโมกฺข\-ปริวารํ โหติ, วีริยินฺทฺริยํ ปคฺคห\-ปริวารํ โหติ, สตินฺทฺริยํ อุปฏฺฐาน\-ปริวารํ โหติ, สมาธินฺทฺริยํ อวิกฺเขป\-ปริวารํ โหติ, ปญฺญินฺทฺริยํ ทสฺสน\-ปริวารํ โหติ, มนินฺทฺริยํ วิชานน\-ปริวารํ โหติ, โสมนสฺสิ\-นฺทฺริยํ \csromanfolio{112}อภิสนฺทน\-ปริวารํ โหติ, ชีวิตินฺทฺริยํ ปวตฺต\-สนฺตตา\-ธิปเตยฺย\-ปริวารํ โหติ. โสตาปตฺติมคฺค\-กฺขเณ ชาตา ธมฺมา ฐเปตฺวา จิตฺต\-สมุฏฺฐานํ รูปํ สพฺเพว กุสลา โหนฺติ, สพฺเพว อนาสวา โหนฺติ, สพฺเพว นิยฺยานิกา โหนฺติ, สพฺเพว อปจยคามิโน โหนฺติ, สพฺเพว โลกุตฺตรา โหนฺติ, สพฺเพว นิพฺพานารมฺมณา โหนฺติ. โสตาปตฺติมคฺค\-กฺขเณ อนญฺญาต\-ญฺญสฺสามี\-ตินฺทฺริยสฺส อิมานิ อฏฺฐินฺทฺริยานิ สหชาต\-ปริวารา โหนฺติ, อญฺญมญฺญ\-ปริวารา โหนฺติ, นิสฺสย\-ปริวารา โหนฺติ, สมฺปยุตฺต\-ปริวารา โหนฺติ, สหคตา โหนฺติ, สหชาตา โหนฺติ, สํสฏฺฐา โหนฺติ, สมฺปยุตฺตา โหนฺติ. เตว ตสฺส อาการา เจว โหนฺติ ปริวารา จ.}

\prose{โสตาปตฺติ\-ผลกฺขเณ อญฺญินฺทฺริยสฺส สทฺธินฺทฺริยํ อธิโมกฺข\-ปริวารํ โหติ, วีริยินฺทฺริยํ ปคฺคห\-ปริวารํ โหติ, สตินฺทฺริยํ อุปฏฺฐาน\-ปริวารํ โหติ, สมาธินฺทฺริยํ อวิกฺเขป\-ปริวารํ โหติ, ปญฺญินฺทฺริยํ ทสฺสน\-ปริวารํ โหติ, มนินฺทฺริยํ วิชานน\-ปริวารํ โหติ, โสมนสฺสิ\-นฺทฺริยํ อภิสนฺทน\-ปริวารํ โหติ, ชีวิตินฺทฺริยํ ปวตฺต\-สนฺตตา\-ธิปเตยฺย\-ปริวารํ โหติ. โสตาปตฺติ\-ผลกฺขเณ ชาตา ธมฺมา สพฺเพว อพฺยากตา โหนฺติ, ฐเปตฺวา จิตฺต\-สมุฏฺฐานํ รูปํ สพฺเพว อนาสวา โหนฺติ, สพฺเพว โลกุตฺตรา โหนฺติ, สพฺเพว นิพฺพานารมฺมณา โหนฺติ. โสตาปตฺติ\-ผลกฺขเณ อญฺญินฺทฺริยสฺส อิมานิ อฏฺฐินฺทฺริยานิ สหชาต\-ปริวารา โหนฺติ, อญฺญมญฺญ\-ปริวารา โหนฺติ, นิสฺสย\-ปริวารา โหนฺติ, สมฺปยุตฺต\-ปริวารา โหนฺติ, สหคตา โหนฺติ, สหชาตา โหนฺติ, สํสฏฺฐา โหนฺติ, สมฺปยุตฺตา โหนฺติ, เตว ตสฺส อาการา เจว โหนฺติ ปริวารา จ.}

\prose{สกทาคามิ\-มคฺคกฺขเณ ฯเปฯ. สกทาคามิ\-ผลกฺขเณ ฯเปฯ. อนาคามิ\-มคฺคกฺขเณ ฯเปฯ. อนาคามิ\-ผลกฺขเณ ฯเปฯ. อรหตฺต\-มคฺคกฺขเณ อญฺญินฺทฺริยสฺส สทฺธินฺทฺริยํ อธิโมกฺข\-ปริวารํ โหติ ฯเปฯ.ชีวิตินฺทฺริยํ ปวตฺต\-สนฺตตา\-ธิปเตยฺย\-ปริวารํ โหติ. อรหตฺต\-มคฺคกฺขเณ ชาตา ธมฺมา ฐเปตฺวา จิตฺต\-สมุฏฺฐานํ รูปํ สพฺเพว กุสลา โหนฺติ, สพฺเพว อนาสวา โหนฺติ, สพฺเพว นิยฺยานิกา โหนฺติ, สพฺเพว อปจยคามิโน โหนฺติ, สพฺเพว โลกุตฺตรา โหนฺติ, สพฺเพว นิพฺพานารมฺมณา โหนฺติ. อรหตฺต\-มคฺคกฺขเณ อญฺญินฺทฺริยสฺส อิมานิ อฏฺฐินฺทฺริยานิ สหชาต\-ปริวารา โหนฺติ, อญฺญมญฺญ\-ปริวารา โหนฺติ, นิสฺสย\-ปริวารา โหนฺติ, สมฺปยุตฺต\-ปริวารา โหนฺติ, \csromanfolio{113}\textbf{สหคตา โหนฺติ, สหชาตา โหนฺติ, สํสฏฺฐา โหนฺติ, สมฺปยุตฺตา โหนฺติ, เตว ตสฺส} อาการา เจว โหนฺติ ปริวารา จ.}

\prose{อรหตฺต\-ผลกฺขเณ อญฺญาตาวิ\-นฺทฺริยสฺส สทฺธินฺทฺริยํ อธิโมกฺข\-ปริวารํ โหติ, วีริยินฺทฺริยํ ปคฺคห\-ปริวารํ โหติ, สตินฺทฺริยํ อุปฏฺฐาน\-ปริวารํ โหติ, สมาธินฺทฺริยํ อวิกฺเขป\-ปริวารํ โหติ, ปญฺญินฺทฺริยํ ทสฺสน\-ปริวารํ โหติ, มนินฺทฺริยํ วิชานน\-ปริวารํ โหติ, โสมนสฺสิ\-นฺทฺริยํ อภิสนฺทน\-ปริวารํ โหติ, ชีวิตินฺทฺริยํ ปวตฺต\-สนฺตตา\-ธิปเตยฺย\-ปริวารํ โหติ. อรหตฺต\-ผลกฺขเณ ชาตา ธมฺมา สพฺเพว อพฺยากตา โหนฺติ, ฐเปตฺวา จิตฺต\-สมุฏฺฐานํ รูปํ สพฺเพว อนาสวา โหนฺติ, สพฺเพว โลกุตฺตรา โหนฺติ, สพฺเพว นิพฺพานารมฺมณา โหนฺติ. อรหตฺต\-ผลกฺขเณ อญฺญาตาวิ\-นฺทฺริยสฺส อิมานิ อฏฺฐินฺทฺริยานิ สหชาต\-ปริวารา โหนฺติ, อญฺญมญฺญ\-ปริวารา โหนฺติ, นิสฺสย\-ปริวารา โหนฺติ, สมฺปยุตฺต\-ปริวารา โหนฺติ, สหคตา โหนฺติ, สหชาตา โหนฺติ, สํสฏฺฐา โหนฺติ, สมฺปยุตฺตา โหนฺติ, เตว ตสฺส อาการา เจว โหนฺติ ปริวารา จ. อิติ อิมานิ อฏฺฐฏฺฐกานิ จตุสฏฺฐิ โหนฺติ.}

\prose{อาสวาติ กตเม เต อาสวา? กามาสโว ภวาสโว ทิฏฺฐาสโว อวิชฺชาสโว. กตฺเถเต อาสวา ขียนฺติ? โสตาปตฺติ\-มคฺเคน อนวเสโส ทิฏฺฐาสโว ขียติ, อปายคมนีโย กามาสโว ขียติ, อปายคมนีโย ภวาสโว ขียติ, อปายคมนีโย อวิชฺชาสโว ขียติ, เอตฺเถเต อาสวา ขียนฺติ. สกทาคามิ\-มคฺเคน โอฬาริโก กามาสโว ขียติ, ตเทกฏฺโฐ ภวาสโว ขียติ, ตเทกฏฺโฐ อวิชฺชาสโว ขียติ, เอตฺเถเต อาสวา ขียนฺติ. อนาคามิมคฺเคน อนวเสโส กามาสโว ขียติ, ตเทกฏฺโฐ ภวาสโว ขียติ, ตเทกฏฺโฐ อวิชฺชาสโค ขียติ, เอตฺเถเต อาสวา ขียนฺติ. อรหตฺต\-มคฺเคน อนวเสโส ภวาสโว ขียติ, อนวเสโส อวิชฺชาสโว ขียติ, เอตฺเถเต อาสวา ขียนฺติ. ตํญาตฏฺเฐน ญาณํ, ปชานนฏฺเฐน ปญฺญา, เตน วุจฺจติ จตุสฏฺฐิยา อากาเรหิ ติณฺณนฺนํ อินฺทฺริยานํ วสิภาวตา ปญฺญา อาสวานํ ขเย ญาณํ.}

\vspace{\tipitakaparskip}
\csromancenter{อาสวกฺขย\-ญาณนิทฺเทโส ปญฺจปญฺญาสโม.}
\csromanmatikaanchor{mtk.44}
\csromantocmarkref{subsection}{56-63. สจฺจญาณจตุกฺกทฺวยนิทฺเทส}{mtk.44}
\bookmark[dest={mtk.44},level=2]{56-63. สจฺจญาณจตุกฺกทฺวยนิทฺเทส}
\csromanheader{2}{56-63. สจฺจญาณ\-จตุกฺก\-ทฺวย\-นิทฺเทส}
\proseitem{108}{\csromanfolio{114}กถํ ปริญฺญฏฺเฐ ปญฺญา ทุกฺเข ญาณํ, ปหานฏฺเฐ ปญฺญา สมุทเย ญาณํ, สจฺฉิกิริยฏฺเฐ ปญฺญา นิโรเธ ญาณํ, ภาวนฏฺเฐ ปญฺญา มคฺเค ญาณํ\csromandash{}ทุกฺขสฺส ปีฬนฏฺโฐ สงฺขตฏฺโฐ สนฺตาปฏฺโฐ วิปริณามฏฺโฐ ปริญฺญาตฏฺโฐ. สมุทยสฺส อายูหนฏฺโฐ นิทานฏฺโฐ สญฺโญคฏฺโฐ ปลิโพธฏฺโฐ ปหานฏฺโฐ. นิโรธสฺส นิสฺสรณฏฺโฐ วิเวกฏฺโฐ อสงฺขตฏฺโฐ อมตฏฺโฐ สจฺฉิกิริยฏฺโฐ. มคฺคสฺส นิยฺยานฏฺโฐ เหตุฏฺโฐ ทสฺสนฏฺโฐ อาธิปเตยฺย\-ฏฺโฐ ภาวนฏฺโฐ. ตํญาตฏฺเฐน ญาณํ, ปชานนฏฺเฐน ปญฺญา, เตน วุจฺจติ ปริญฺญฏฺเฐ ปญฺญา ทุกฺเข ญาณํ. ปหานฏฺเฐ ปญฺญา สมุทเย ญาณํ. สจฺฉิกิริยฏฺเฐ ปญฺญา นิโรเธ ญาณํ. ภาวนฏฺเฐ ปญฺญา มคฺเค ญาณํ.}

\proseitem{109}{กถํ ทุกฺเข ญาณํ. ทุกฺขสมุทเย ญาณํ. ทุกฺขนิโรเธ ญาณํ. ทุกฺขนิโรธ\-คามินิยา ปฏิปทาย ญาณํ\csromandash{}มคฺค\-สมงฺคิสฺส ญาณํ ทุกฺเขเปตํ ญาณํ, ทุกฺขสมุทเยเปตํ ญาณํ, ทุกฺขนิโรเธเปตํ ญาณํ, ทุกฺขนิโรธ\-คามินิยา ปฏิปทา\-ยเปตํ ญาณํ.}

\prose{ตตฺถ กตมํ ทุกฺเข ญาณํ? ทุกฺขํ อารพฺภ ยา อุปฺปชฺชติ ปญฺญา ปชานนา วิจโย ปวิจโย ธมฺมวิจโย สลฺลกฺขณา อุปลกฺขณา ปจฺจุ\-ปลกฺขณา ปณฺฑิจฺจํ โกสลฺลํ เนปุญฺญํ เวภพฺยา จินฺตา อุปปริกฺขา ภูริ เมธา ปริณายิกา วิปสฺสนา สมฺปชญฺญํ ปโตโท ปญฺญา ปญฺญินฺทฺริยํ ปญฺญาพลํ ปญฺญาสตฺถํ ปญฺญาปาสาโท ปญฺญาอาโลโก ปญฺญาโอภาโส ปญฺญาปชฺโชโต ปญฺญารตนํ อโมโห ธมฺมวิจโย สมฺมาทิฏฺฐิ, อิทํ วุจฺจติ ทุกฺเข ญาณํ. ทุกฺข\-สมุทยํ อารพฺภ ฯเปฯ ทุกฺขนิโรธํ อารพฺภ ฯเปฯ ทุกฺขนิโรธ\-คามินิํ ปฏิปทํ อารพฺภ ยา อุปฺปชฺชติ ปญฺญา ปชานนา ฯเปฯ อโมโห ธมฺมวิจโย สมฺมาทิฏฺฐิ, อิทํ วุจฺจติ ทุกฺขนิโรธ\-คามินิยา ปฏิปทาย ญาณํ, ตํญาตฏฺเฐน ญาณํ, ปชานนฏฺเฐน ปญฺญา, เตน วุจฺจติ ทุกฺเข ญาณํ. ทุกฺขสมุทเย ญาณํ. ทุกฺขนิโรเธ ญาณํ. ทุกฺขนิโรธ\-คามินิยา ปฏิปทาย ญาณํ.}

\vspace{\tipitakaparskip}
\csromancenter{สจฺจญาณ\-จตุกฺก\-ทฺวย\-นิทฺเทโส เตสฏฺฐิโม.}
\csromanmatikaanchor{mtk.45}
\csromantocmarkref{subsection}{64-67. สุทฺธิกปฏิสมฺภิทาญาณนิทฺเทส}{mtk.45}
\bookmark[dest={mtk.45},level=2]{64-67. สุทฺธิกปฏิสมฺภิทาญาณนิทฺเทส}
\csromanheader{2}{64-67. สุทฺธิก\-ปฏิสมฺภิทา\-ญาณ\-นิทฺเทส}
\proseitem{110}{\csromanfolio{115}กถํ อตฺถ\-ปฏิสมฺภิเท ญาณํ, ธมฺม\-ปฏิสมฺภิเท ญาณํ, นิรุตฺติ\-ปฏิสมฺภิเท ญาณํ, ปฏิภาน\-ปฏิสมฺภิเท ญาณํ\csromandash{}อตฺเถสุ ญาณํ อตฺถ\-ปฏิสมฺภิทา, ธมฺเมสุ ญาณํ ธมฺม\-ปฏิสมฺภิทา, นิรุตฺตีสุ ญาณํ นิรุตฺติ\-ปฏิสมฺภิทา, ปฏิภาเนสุ ญาณํ ปฏิภาน\-ปฏิสมฺภิทา. อตฺถนานตฺเต ปญฺญา อตฺถ\-ปฏิสมฺภิเท ญาณํ, ธมฺมนานตฺเต ปญฺญา ธมฺม\-ปฏิสมฺภิเท ญาณํ, นิรุตฺตินานตฺเต ปญฺญา นิรุตฺติ\-ปฏิสมฺภิเท ญาณํ, ปฏิภาน\-นานตฺเต ปญฺญา ปฏิภาน\-ปฏิสมฺภิเท ญาณํ.  อตฺถ\-ววตฺถาเน ปญฺญา อตฺถ\-ปฏิสมฺภิเท ญาณํ, ธมฺม\-ววตฺถาเน ปญฺญา ธมฺม\-ปฏิสมฺภิเท ญาณํ, นิรุตฺติ\-ววตฺถาเน ปญฺญา นิรุตฺติ\-ปฏิสมฺภิเท ญาณํ, ปฏิภาน\-ววตฺถาเน ปญฺญา ปฏิภาน\-ปฏิสมฺภิเท ญาณํ.}

\prose{อตฺถ\-สลฺลกฺขเณ ปญฺญา อตฺถ\-ปฏิสมฺภิเท ญาณํ, ธมฺม\-สลฺลกฺขเณ ปญฺญา ธมฺม\-ปฏิสมฺภิเท ญาณํ, นิรุตฺติ\-สลฺลกฺขเณ ปญฺญา นิรุตฺติ\-ปฏิสมฺภิเท ญาณํ, ปฏิภาน\-สลฺลกฺขเณ ปญฺญา ปฏิภาน\-ปฏิสมฺภิเท ญาณํ.  อตฺถู\-ปลกฺขเณ ปญฺญา อตฺถ\-ปฏิสมฺภิเท ญาณํ, ธมฺมู\-ปลกฺขเณ ปญฺญา ธมฺม\-ปฏิสมฺภิเท ญาณํ, นิรุตฺตู\-ปลกฺขเณ ปญฺญา นิรุตฺติ\-ปฏิสมฺภิเท ญาณํ. ปฏิภานู\-ปลกฺขเณ ปญฺญา ปฏิภาน\-ปฏิสมฺภิเท ญาณํ.}

\prose{อตฺถปฺปเภเท ปญฺญา อตฺถ\-ปฏิสมฺภิเท ญาณํ, ธมฺมปฺปเภเท ปญฺญา ธมฺม\-ปฏิสมฺภิเท ญาณํ, นิรุตฺติ\-ปฺปเภเท ปญฺญา นิรุตฺติ\-ปฏิสมฺภิเท ญาณํ, ปฏิภาน\-ปฺปเภเท ปญฺญา ปฏิภาน\-ปฏิสมฺภิเท ญาณํ.  อตฺถ\-ปฺปภาวเน ปญฺญา อตฺถ\-ปฏิสมฺภิเท ญาณํ, ธมฺม\-ปฺปภาวเน ปญฺญา ธมฺม\-ปฏิสมฺภิเท ญาณํ, นิรุตฺติ\-ปฺปภาวเน ปญฺญา นิรุตฺติ\-ปฏิสมฺภิเท ญาณํ, ปฏิภาน\-ปฺปภาวเน ปญฺญา ปฏิภาน\-ปฏิสมฺภิเท ญาณํ.}

\prosewithcloserruled{อตฺถโชตเน ปญฺญา อตฺถ\-ปฏิสมฺภิเท ญาณํ, ธมฺมโชตเน ปญฺญา ธมฺม\-ปฏิสมฺภิเท ญาณํ, นิรุตฺติโชตเน ปญฺญา นิรุตฺติ\-ปฏิสมฺภิเท ญาณํ, ปฏิภาน\-โชตเน ปญฺญา ปฏิภาน\-ปฏิสมฺภิเท ญาณํ.  อตฺถวิโรจเน ปญฺญา อตฺถ\-ปฏิสมฺภิเท ญาณํ, ธมฺมวิโรจเน ปญฺญา ธมฺม\-ปฏิสมฺภิเท ญาณํ, นิรุตฺติ\-วิโรจเน ปญฺญา นิรุตฺติ\-ปฏิสมฺภิเท ญาณํ, ปฏิภาน\-วิโรจเน ปญฺญา ปฏิภาน\-ปฏิสมฺภิเท ญาณํ.  อตฺถปฺปกาสเน ปญฺญา อตฺถ\-ปฏิสมฺภิเท ญาณํ, ธมฺม\-ปฺปกาสเน ปญฺญา ธมฺม\-ปฏิสมฺภิเท ญาณํ, นิรุตฺถิปฺปกาสเน ปญฺญา นิรุตฺติ\-ปฏิสมฺภิเท ญาณํ, ปฏิภาน\-ปฺปกาสเน ปญฺญา ปฏิภาน\-ปฏิสมฺภิเท ญาณํ. \csromanfolio{116}ตํญาตฏฺเฐน ญาณํ, ปชานนฏฺเฐน ปญฺญา, เตน วุจฺจติ อตฺถ\-ปฏิสมฺภิเท ญาณํ, ธมฺม\-ปฏิสมฺภิเท ญาณํ, นิรุตฺติ\-ปฏิสมฺภิเท ญาณํ, ปฏิภาน\-ปฏิสมฺภิเท ญาณํ.}{สุทฺธิก\-ปฏิสมฺภิทา\-ญาณ\-นิทฺเทโส สตฺตสฏฺฐิโม.}
\csromanmatikaanchor{mtk.46}
\csromantocmarkref{subsection}{68. อินฺทฺริยปโรปริยตฺตญาณนิทฺเทส}{mtk.46}
\bookmark[dest={mtk.46},level=2]{68. อินฺทฺริยปโรปริยตฺตญาณนิทฺเทส}
\csromanheader{2}{68. อินฺทฺริยปโรปริยตฺต\-ญาณ\-นิทฺเทส}
\proseitem{111}{กตมํ ตถาคตสฺส อินฺทฺริยปโรปริยตฺต\-ญาณํ? อิธ ตถาคโต สตฺเต ปสฺสติ อปฺปรชกฺเข มหารชกฺเข ติกฺขินฺทฺริเย มุทินฺทฺริเย สฺวากาเร ทฺวากาเร สุวิญฺญาปเย ทุวิญฺญาปเย อปฺเปกจฺเจ ปรโลก\-วชฺช\-ภยทสฺสาวิโน อปฺเปกจฺเจ นปรโลก\-วชฺช\-ภยทสฺสาวิโน.}

\prose{\textbf{อปฺปรชกฺเข มหารชกฺเข}ติ สทฺโธ ปุคฺคโล อปฺปรชกฺโข, อสฺสทฺโธ ปุคฺคโล มหารชกฺโข. อารทฺธวีริโย ปุคฺคโล อปฺปรชกฺโข, กุสีโต ปุคฺคโล มหารชกฺโข. อุปฏฺฐิตสฺสติ ปุคฺคโล อปฺปรชกฺโข, มุฏฺฐสฺสติ ปุคฺคโล มหารชกฺโข. สมาหิโต ปุคฺคโล อปฺปรชกฺโข, อสมาหิโต ปุคฺคโล มหารชกฺโข. ปญฺญวา ปุคฺคโล อปฺปรชกฺโข, ทุปฺปญฺโญ ปุคฺคโล มหารชกฺโข.}

\prose{\textbf{ติกฺขินฺทฺริเย มุทินฺทฺริเย}ติ สทฺโธ ปุคฺคโล ติกฺขินฺทฺริโย, อสฺสทฺโธ ปุคฺคโล มุทินฺทฺริโย. อารทฺธวีริโย ปุคฺคโล ติกฺขินฺทฺริโย, กุสีโต ปุคฺคโล มุทินฺทฺริโย. อุปฏฺฐิตสฺสติ ปุคฺคโล ติกฺขินฺทฺริโย, มุฏฺฐสฺสติ ปุคฺคโล มุทินฺทฺริโย. สมาหิโต ปุคฺคโล ติกฺขินฺทฺริโย, อสมาหิโต ปุคฺคโล มุทินฺทฺริโย. ปญฺญวา ปุคฺคโล ติกฺขินฺทฺริโย, ทุปฺปญฺโญ ปุคฺคโล มุทินฺทฺริโย.}

\prose{\textbf{สฺวากาเร ทฺวากาเร}ติ สทฺโธ ปุคฺคโล สฺวากาโร, อสฺสทฺโธ ปุคฺคโล ทฺวากาโร. อารทฺธวีริโย ปุคฺคโล สฺวากาโร, กุสีโต ปุคฺคโล ทฺวากาโร. อุปฏฺฐิตสฺสติ ปุคฺคโล สฺวากาโร, มุฏฺฐสฺสติ ปุคฺคโล ทฺวากาโร. สมาหิโต ปุคฺคโล สฺวากาโร, อสมาหิโต ปุคฺคโล ทฺวากาโร. ปญฺญวา ปุคฺคโล สฺวากาโร, ทุปฺปญฺโญ ปุคฺคโล ทฺวากาโร.}

\prose{\textbf{สุวิญฺญาปเย ทุวิญฺญาปเย}ติ สทฺโธ ปุคฺคโล สุวิญฺญาปโย, อสฺสทฺโธ ปุคฺคโล ทุวิญฺญาปโย. อารทฺธวีริโย ปุคฺคโล สุวิญฺญาปโย, กุสีโต ปุคฺคโล ทุวิญฺญาปโย. อุปฏฺฐิตสฺสติ ปุคฺคโล สุวิญฺญาปโย, มุฏฺฐสฺสติ ปุคฺคโล ทุวิญฺญาปโย. สมาหิโต ปุคฺคโล \csromanfolio{117}สุวิญฺญาปโย, อสมาหิโต ปุคฺคโล ทุวิญฺญาปโย. ปญฺญวา ปุคฺคโล สุวิญฺญาปโย, ทุปฺปญฺโญ ปุคฺคโล ทุวิญฺญาปโย.}

\prose{อปฺเปกจฺเจ ปรโลก\-วชฺช\-ภยทสฺสาวิโน อปฺเปกจฺเจ นปรโลก\-วชฺช\-ภยทสฺสาวิโนติ สทฺโธ ปุคฺคโล ปรโลก\-วชฺช\-ภยทสฺสาวี, อสฺสทฺโธ ปุคฺคโล นปรโลก\-วชฺช\-ภยทสฺสาวี. อารทฺธวีริโย ปุคฺคโล ปรโลก\-วชฺช\-ภยทสฺสาวี, กุสีโต ปุคฺคโล นปรโลก\-วชฺช\-ภยทสฺสาวี. อุปฏฺฐิตสฺสติ ปุคฺคโล ปรโลก\-วชฺช\-ภยทสฺสาวี, มุฏฺฐสฺสติ ปุคฺคโล นปรโลก\-วชฺช\-ภยทสฺสาวี. สมาหิโต ปุคฺคโล ปรโลก\-วชฺช\-ภยทสฺสาวี, อสมาหิโต ปุคฺคโล นปรโลก\-วชฺช\-ภยทสฺสาวี. ปญฺญวา ปุคฺคโล ปรโลก\-วชฺช\-ภยทสฺสาวี, ทุปฺปญฺโญ ปุคฺคโล นปรโลก\-วชฺช\-ภยทสฺสาวี.}

\proseitem{112}{โลโกติ ขนฺธโลโก ธาตุโลโก อายตนโลโก วิปตฺติ\-ภว\-โลโก วิปตฺติ\-สมฺภว\-โลโก สมฺปตฺติ\-ภว\-โลโก สมฺปตฺติ\-สมฺภว\-โลโก.}

\prose{เอโก โลโก, สพฺเพ สตฺตา อาหารฏฺฐิติกา. เทฺว โลกา, นามญฺจ รูปญฺจ. ตโย โลกา, ติสฺโส เวทนา. จตฺตาโร โลกา, จตฺตาโร อาหารา. ปญฺจ โลกา, ปญฺจุ\-ปาทาน\-กฺขนฺธา. ฉ โลกา, ฉ อชฺฌตฺติกานิ อายตนานิ. สตฺต โลกา, สตฺต วิญฺญาณ\-ฏฺฐิติโย. อฏฺฐ โลกา, อฏฺฐ โลกธมฺมา. นว โลกา, นว สตฺตาวาสา. ทส โลกา, ทสายตนานิ. ทฺวาทส โลกา, ทฺวาทสา\-ยตนานิ. อฏฺฐารส โลกา, อฏฺฐารส ธาตุโย.}

\prosewithcloserruled{วชฺชนฺติ สพฺเพ กิเลสา วชฺชา, สพฺเพ ทุจฺจริตา วชฺชา, สพฺเพ อภิสงฺขารา วชฺชา, สพฺเพ ภวคามิกมฺมา วชฺชา, อิติ อิมสฺมิญฺจ โลเก อิมสฺมิญฺจ วชฺเช ติพฺพา ภยสญฺญา ปจฺจุปฏฺฐิตา โหติ เสยฺยถาปิ อุกฺขิตฺตาสิเก วธเก. อิเมหิ ปญฺญาสาย อากาเรหิ อิมานิ ปญฺจินฺทฺริยานิ ชานาติ ปสฺสติ อญฺญาติ ปฏิวิชฺฌติ, อิทํ ตถาคตสฺส อินฺทฺริยปโรปริยตฺต\-ญาณํ.}{อินฺทฺริยปโรปริยตฺต\-ญาณ\-นิทฺเทโส อฏฺฐสฏฺฐิโม.}
\csromanmatikaanchor{mtk.47}
\csromantocmarkref{subsection}{69. อาสยานุสยญาณนิทฺเทส}{mtk.47}
\bookmark[dest={mtk.47},level=2]{69. อาสยานุสยญาณนิทฺเทส}
\csromanheader{2}{69. อาสยา\-นุสย\-ญาณ\-นิทฺเทส}
\proseitem{113}{กตมํ \textbf{ตถาคตสฺส สตฺตานํ อาสยานุสเย ญาณํ}? อิธ ตถาคโต สตฺตานํ อาสยํ ชานาติ, อนุสยํ ชานาติ, จริตํ \csromanfolio{118}ชานาติ, อธิมุตฺติํ ชานาติ, ภพฺพาภพฺเพ สตฺเต ปชานาติ. กตโม\csromansharedfootnote{9f56e9234ddbfb2e}{กตโม จ (สฺยา, ก)} สตฺตานํ อาสโย? “สสฺสโต โลโก”ติ วา “อสสฺสโต โลโก”ติ วา “อนฺตวา โลโก”ติ วา “อนนฺตวา โลโก”ติ วา “ตํ ชีวํ ตํ สรีรนฺ”ติ วา “อญฺญํ ชีวํ อญฺญํ สรีรนฺ”ติ วา “โหติ ตถาคโต ปรํ มรณา”ติ วา “น โหติ ตถาคโต ปรํ มรณา”ติ วา “โหติ จ น จ โหติ ตถาคโต ปรํ มรณา”ติ วา “เนว โหติ น น โหติ ตถาคโต ปรํ มรณา”ติ วา. อิติ ภวทิฏฺฐิ\-สนฺนิสฺสิตา วา สตฺตา โหนฺติ วิภวทิฏฺฐิ\-สนฺนิสฺสิตา วา.}

\prose{เอเต วา ปน อุโภ อนฺเต อนุปคมฺม อิทปฺปจฺจยตา\-ปฏิจฺจ\-สมุปฺปนฺเนสุ ธมฺเมสุ อนุโลมิกา ขนฺติ ปฏิลทฺธา โหติ, ยถาภูตํ วา ญาณํ. กามํ เสวนฺตญฺเญว ชานาติ “อยํ ปุคฺคโล กามครุโก กามาสโย กามาธิมุตฺโต”ติ. กามํ เสวนฺตญฺเญว ชานาติ “อยํ ปุคฺคโล เนกฺขมฺม\-ครุโก เนกฺขมฺมาสโย เนกฺขมฺมา\-ธิมุตฺโต”ติ. เนกฺขมฺมํ เสวนฺตญฺเญว ชานาติ “อยํ ปุคฺคโล (เนกฺขมฺม\-ครุโก เนกฺขมฺมาสโย เนกฺขมฺมา\-ธิมุตฺโต”ติ. เนกฺขมฺมํ เสวนฺตญฺเญว ชานาติ “อยํ ปุคฺคโล กามครุโก กามาสโย กามาธิมุตฺโต”ติ. พฺยาปาทํ เสวนฺตญฺเญว ชานาติ)\csromansharedfootnote{a775221e2e8c2b67}{( ) เอตฺถนฺตเร ปาฐา นตฺถิ สฺยามโปตฺถเก, เอวมุปริปิ.} “อยํ ปุคฺคโล พฺยาปาทครุโก พฺยาปาทาสโย พฺยาปาทาธิมุตฺโต”ติ. พฺยาปาทํ เสวนฺตญฺเญว ชานาติ “อยํ ปุคฺคโล อพฺยาปาทครุโก อพฺยาปาทาสโย อพฺยาปาทาธิมุตฺโต”ติ. อพฺยาปาทํ เสวนฺตญฺเญว ชานาติ “อยํ ปุคฺคโล (อพฺยาปาทครุโก อพฺยาปาทาสโย อพฺยาปาทาธิมุตฺโต”ติ. อพฺยาปาทํ เสวนฺตญฺเญว ชานาติ “อยํ ปุคฺคโล พฺยาปาทครุโก พฺยาปาทาสโย พฺยาปาทาธิมุตฺโต”ติ. ถินมิทฺธํ เสวนฺตญฺเญว ชานาติ) “อยํ ปุคฺคโล ถินมิทฺธ\-ครุโก ถินมิทฺธา\-สโย ถินมิทฺธาธิมุตฺโต”ติ. ถินมิทฺธํ เสวนฺตญฺเญว ชานาติ “อยํ ปุคฺคโล อาโลกสญฺญา\-ครุโก อาโลกสญฺญาสโย อาโลกสญฺญาธิมุตฺโต”ติ. อาโลกสญฺญํ เสวนฺตญฺเญว ชานาติ “อยํ ปุคฺคโล (อาโลกสญฺญา\-ครุโก อาโลกสญฺญาสโย อาโลกสญฺญาธิมุตฺโต”ติ. อาโลกสญฺญํ เสวนฺตญฺเญว ชานาติ “อยํ ปุคฺคโล ถินมิทฺธ\-ครุโก ถินมิทฺธา\-สโย ถินมิทฺธาธิมุตฺโต”ติ,) อยํ สตฺตานํ อาสโย.}

\proseitem{114}{\csromanfolio{119}กตโม จ สตฺตานํ อนุสโย? สตฺตานุสยา กามราคานุสโย ปฏิฆานุสโย มานานุสโย ทิฏฺฐานุสโย วิจิกิจฺฉา\-นุสโย ภวราคา\-นุสโย อวิชฺชานุสโย. ยํ โลเก ปิยรูปํ สาตรูปํ, เอตฺถ สตฺตานํ กามราคานุสโย อนุเสติ. ยํ โลเก อปฺปิยรูปํ อสาตรูปํ, เอตฺถ สตฺตานํ ปฏิฆานุสโย อนุเสติ. อิติ อิเมสุ ทฺวีสุ ธมฺเมสุ อวิชฺชา อนุปติตา, ตเทกฏฺโฐ มาโน จ ทิฏฺฐิ จ วิจิกิจฺฉา จ ทฏฺฐพฺพา. อยํ สตฺตานํ อนุสโย.}

\prose{กตมญฺจ สตฺตานํ จริตํ? ปุญฺญา\-ภิสงฺขาโร อปุญฺญา\-ภิสงฺขาโร อาเนญฺชา\-ภิสงฺขาโร ปริตฺตภูมโก วา มหาภูมโก วา. อิทํ สตฺตานํ จริตํ.}

\proseitem{115}{กตมา จ สตฺตานํ อธิมุตฺติ? สนฺติ สตฺตา หีนาธิมุตฺติกา, สนฺติ สตฺตา ปณีตา\-ธิมุตฺติกา.  หีนาธิมุตฺติกา สตฺตา หีนาธิมุตฺติเก สตฺเต เสวนฺติ ภชนฺติ ปยิรุปาสนฺติ, ปณีตา\-ธิมุตฺติกา สตฺตา ปณีตา\-ธิมุตฺติเก สตฺเต เสวนฺติ ภชนฺติ ปยิรุปาสนฺติ. อตีตมฺปิ อทฺธานํ หีนาธิมุตฺติกา สตฺตา หีนาธิมุตฺติเก สตฺเต เสวิํสุ ภชิํสุ ปยิรุปาสิํสุ ปณีตา\-ธิมุตฺติกา สตฺตา ปณีตา\-ธิมุตฺติเก สตฺเต เสวิํสุ ภชิํสุ ปยิรุปาสิํสุ. อนาคตมฺปิ อทฺธานํ หีนาธิมุตฺติกา สตฺตา หีนาธิมุตฺติเก สตฺเต เสวิสฺสนฺติ ภชิสฺสนฺติ ปยิรุปาสสฺสนฺติ, ปณีตา\-ธิมุตฺติกา สตฺตา ปณีภาธิมุตฺติเก สตฺเถ เสวิสฺสนฺติ ภชิสฺสนฺติ ปยิรุปาสิสฺสนฺติ. อยํ สตฺตานํ อธิมุตฺติ.}

\prose{กตเม สตฺตา อภพฺพา? เย เต สตฺตา กมฺมาวรเณน สมนฺนาคตา กิเลสาวรเณน สมนฺนาคตา วิปากาวรเณน สมนฺนาคตา อสฺสทฺธา อจฺฉนฺทิกา ทุปฺปญฺญา อภพฺพา นิยามํ โอกฺกมิตุํ กุสเลสุ ธมฺเมสุ สมฺมตฺตํ, อิเม เต สตฺตา อภพฺพา.}

\prose{กตเม สตฺตา ภพฺพา? เย เต สตฺตา น กมฺมาวรเณน สมนฺนาคตา น กิเลสาวรเณน สมนฺนาคตา น วิปากาวรเณน สมนฺนาคตา สทฺธา ฉนฺทิกา ปญฺญวนฺโต ภพฺพา นิยามํ โอกฺกมิตุํ กุสเลสุ ธมฺเมสุ สมฺมตฺตํ. อิเม เต สตฺตา ภพฺพา, อิทํ ตถาคตสฺส สตฺตานํ อาสยานุสเย ญาณํ.}

\vspace{\tipitakaparskip}
\csromancenter{อาสยา\-นุสย\-ญาณ\-นิทฺเทโส นวสฏฺฐิโม.}
\csromanmatikaanchor{mtk.48}
\csromantocmarkref{subsection}{70. ยมกปาฏิหีรญาณนิทฺเทส}{mtk.48}
\bookmark[dest={mtk.48},level=2]{70. ยมกปาฏิหีรญาณนิทฺเทส}
\csromanheader{2}{70. ยมกปาฏิหีร\-ญาณ\-นิทฺเทส}
\proseitem{116}{\csromanfolio{120}กตมํ ตถาคตสฺส ยมกปาฏิหีเร ญาณํ? อิธ ตถาคโต ยมกปาฏิหีรํ กโรติ อสาธารณํ สาวเกหิ, อุปริมกายโต อคฺคิกฺขนฺโธ ปวตฺตติ, เหฏฺฐิมกายโต อุทกธารา ปวตฺตติ. เหฏฺฐิมกายโต อคฺคิกฺขนฺโธ ปวตฺตติ, อุปริมกายโต อุทกธารา ปวตฺตติ.  ปุรตฺถิม\-กายโต อคฺคิกฺขนฺโธ ปวตฺตติ, ปจฺฉิมกายโต อุทกธารา ปวตฺตติ. ปจฺฉิมกายโต อคฺคิกฺขนฺโธ ปวตฺตติ, ปุรตฺติมกายโต อุทกธารา ปวตฺตติ.  ทกฺขิณ\-อกฺขิโต อคฺคิกฺขนฺโธ ปวตฺตติ, วามอกฺขิโต อุทกธารา ปวตฺตติ. วามอกฺขิโต อคฺคิกฺขนฺโธ ปวตฺตติ, ทกฺขิณ\-อกฺขิโต อุทกธารา ปวตฺตติ.  ทกฺขิณ\-กณฺณโสตโต อคฺคิกฺขนฺโธ ปวตฺตติ, วาม\-กณฺณโสตโต อุทกธารา ปวตฺตติ, วาม\-กณฺณโสตโต อคฺคิกฺขนฺโธ ปวตฺตติ, ทกฺขิณ\-กณฺณโสตโต อุทกธารา ปวตฺตติ.  ทกฺขิณ\-นาสิกาโสตโต อคฺคิกฺขนฺโธ ปวตฺตติ, วาม\-นาสิกาโสตโต อุทกธารา ปวตฺตติ. วาม\-นาสิกาโสตโต อคฺคิกฺขนฺโธ ปวตฺตติ, ทกฺขิณ\-นาสิกาโสตโต อุทกธารา ปวตฺตติ.  ทกฺขิณ\-อํสกูฏโต อคฺคิกฺขนฺโธ ปวตฺตติ, วามอํสกูฏโต อุทกธารา ปวตฺตติ. วามอํสกูฏโต อคฺคิกฺขนฺโธ ปวตฺตติ, ทกฺขิณ\-อํสกูฏโต อุทกธารา ปวตฺตติ.  ทกฺขิณ\-หตฺถโต อคฺคิกฺขนฺโธ ปวตฺตติ, วามหตฺถโต อุทกธารา ปวตฺตติ. วามหตฺถโต อคฺคิกฺขนฺโธ ปวตฺตติ, ทกฺขิณ\-หตฺถโต อุทกธารา ปวตฺตติ.  ทกฺขิณ\-ปสฺสโต อคฺคิกฺขนฺโธ ปวตฺตติ, วามปสฺสโต อุทกธารา ปวตฺตติ. วามปสฺสโต อคฺคิกฺขนฺโธ ปวตฺตติ, ทกฺขิณ\-ปสฺสโต อุทกธารา ปวตฺตติ.  ทกฺขิณปาทโต อคฺคิกฺขนฺโธ ปวตฺตติ, วามปาทโต อุทกธารา ปวตฺตติ. วามปาทโต อคฺคิกฺขนฺโธ ปวตฺตติ, ทกฺขิณปาทโต อุทกธารา ปวตฺตติ.  องฺคุลงฺคุเลหิ อคฺคิกฺขนฺโธ ปวตฺตติ, องฺคุล\-นฺตริกาหิ อุทกธารา ปวตฺตติ, องฺคุล\-นฺตริกาหิ อคฺคิกฺขนฺโธ ปวตฺตติ, องฺคุลงฺคุเลหิ อุทกธารา ปวตฺตติ.  เอเกกโลมโต อคฺคิกฺขนฺโธ ปวตฺตติ, เอเกกโลมโต อุทกธารา ปวตฺตติ. โลมกูปโต โลมกูปโต อคฺคิกฺขนฺโธ ปวตฺตติ, โลมกูปโต โลมกูปโต อุทกธารา ปวตฺตติ. \csromanfolio{121}ฉนฺนํ วณฺณานํ นีลานํ ปีตกานํ โลหิตกานํ โอทาตานํ มญฺชิฏฺฐานํ\csromansharedfootnote{e4d0c69913a23938}{มญฺเชฏฺฐานํ (สฺยา, ก)} ปภสฺสรานํ.}

\prosewithcloserruled{ภควา จงฺกมติ, นิมฺมิโต ติฏฺฐติ วา นิสีทติ วา เสยฺยํ วา กปฺเปติ. ภควา ติฏฺฐติ, นิมฺมิโต จงฺกมติ วา นิสีทติ วา เสยฺยํ วา กปฺเปติ. ภควา นิสีทติ, นิมฺมิโต จงฺกมติ วา ติฏฺฐติ วา เสยฺยํ วา กปฺเปติ. ภควา เสยฺยํ กปฺเปติ, นิมฺมิโต จงฺกมติ วา ติฏฺฐติ วา นิสีทติ วา. นิมฺมิโต จงฺกมติ, ภควา ติฏฺฐติ วา นิสีทติ วา เสยฺยํ วา กปฺเปติ. นิมฺมิโต ติฏฺฐติ, ภควา จงฺกมติ วา นิสีทติ วา เสยฺยํ วา กปฺเปติ. นิมฺมิโต นิสีทติ, ภควา จงฺกมติ วา ติฏฺฐติ วา เสยฺยํ วา กปฺเปติ. นิมฺมิโต เสยฺยํ กปฺเปติ, ภควา จงฺกมติ วา ติฏฺฐติ วา นิสีทติ วา. อิทํ ตถาคตสฺส ยมกปาฏิหีเร ญาณํ.}{ยมกปาฏิหีร\-ญาณ\-นิทฺเทโส สตฺตติโม.}
\csromanmatikaanchor{mtk.49}
\csromantocmarkref{subsection}{71. มหากรุณาญาณนิทฺเทส}{mtk.49}
\bookmark[dest={mtk.49},level=2]{71. มหากรุณาญาณนิทฺเทส}
\csromanheader{2}{71. มหากรุณา\-ญาณ\-นิทฺเทส}
\proseitem{117}{กตมํ ตถาคตสฺส มหาครุณาสมาปตฺติยา ญาณํ? พหุเกหิ อากาเรหิ ปสฺสนฺตานํ พุทฺธานํ ภควนฺตานํ สตฺเตตุ มหากรุณา โอกฺกมติ, “อาทิตฺโต โลกสนฺนิวาโส”ติ ปสฺสนฺตานํ พุทฺธานํ ภควนฺตานํ สตฺเตสุ มหากรุณา โอกฺกมติ. “อุยฺยุตฺโต โลกสนฺนิวาโส”ติ ปสฺสนฺตานํ พุทฺธานํ ภควนฺตานํ สตฺเตสุ มหากรุณา โอกฺกมติ, “ปยาโต โลกสนฺนิวาโส”ติ ปสฺสนฺตานํ พุทฺธานํ ภควนฺตานํ สตฺเตสุ มหากรุณา โอกฺกมติ. “กุมฺมคฺค\-ปฺปฏิปนฺโน\csromansharedfootnote{e488cba17824729a}{กุมฺมคฺคํ ปฏิปนฺโน (สฺยา)} โลกสนฺนิวาโส”ติ ปสฺสนฺตานํ พุทฺธานํ ภควนฺตานํ สตฺเตสุ มหากรุณา โอกฺกมติ, “อุปนียติ โลโก อทฺธุโว”ติ ปสฺสนฺตานํ พุทฺธานํ ภควนฺตานํ สตฺเตสุ มหากรุณา โอกฺกมติ. “อตาโณ\csromansharedfootnote{b0f9ad759fac1822}{อตฺตาโณ (สฺยา)} โลโก อนภิสฺสโร”ติ ปสฺสนฺตานํ พุทฺธานํ ภควนฺตานํ สตฺเตสุ มหากรุณา โอกฺกมติ, “อสฺสโก โลโก สพฺพํ ปหาย คมนียนฺ”ติ ปสฺสนฺตานํ พุทฺธานํ ภควนฺตานํ สตฺเตสุ มหากรุณา โอกฺกมติ, “อูโน โลโก อติตฺโต ตณฺหาทาโส”ติ ปสฺสนฺตานํ พุทฺธานํ ภควนฺตานํ สตฺเตสุ มหากรุณา โอกฺกมติ, “อตายโน โลกสนฺนิวาโส”ติ \csromanfolio{122}ปสฺสนฺตานํ ฯเปฯ “อเลโณ โลกสนฺนิวาโส”ติ ปสฺสนฺตานํ ฯเปฯ “อสรโณ โลกสนฺนิวาโส”ติ ปสฺสนฺตานํ ฯเปฯ “อสรณีภูโต โลกสนฺนิวาโส”ติ ปสฺสนฺตานํ ฯเปฯ.}

\prose{“อุทฺธโต โลโก อวูปสนฺโต”ติ ปสฺสนฺตานํ ฯเปฯ “สสลฺโลโลกสนฺนิวาโส วิทฺโธ ปุถุสลฺเลหิ, ตสฺส นตฺถญฺโญ โกจิ สลฺลานํ อุทฺธตา อญฺญตฺร มยา”ติ ปสฺสนฺตานํ ฯเปฯ “อวิชฺชนฺธการา\-วรโณ โลกสนฺนิวาโส อณฺฑภูโต กิเลส\-ปญฺชร\-ปกฺขิตฺโต, ตสฺส นตฺถญฺโญ โกจิ อาโลกํ ทสฺเสตา อญฺญตฺร มยา”ติ ปสฺสนฺตานํ ฯเปฯ “อวิชฺชาคโต โลกสนฺนิวาโส อณฺฑภูโต ปริโยนทฺโธ ตนฺตา\-กุลก\-ชาโต\csromansharedfootnote{b135c7953cd12fe9}{ตนฺตากุลชาโต (สฺยา)} กุลาคณฺฐิกชาโต\csromansharedfootnote{713672b48883909e}{คุฬาคุณฺฐิกชาโต (สฺยา), คุลาคุณฺฑิกชาโต (ก, สีฏฺฐ) ที 2. 47 ปิฏฺเฐ ปสฺสิตพฺพา.} มุญฺช\-ปพฺพช\-ภูโต อปายํ ทุคฺคติํ วินิปาตํ สํสารํ นาติวตฺตตี”ติ ปสฺสนฺตานํ ฯเปฯ “อวิชฺชาวิสโทส\-สํลิตฺโต โลกสนฺนิวาโส กิเลส\-กลลีภูโต”ติ ปสฺสนฺตานํ ฯเปฯ “ราคโทสโมห- ชฏาชฏิโต โลกสนฺนิวาโส, ตสฺส นตฺถญฺโญ โกจิ ชฏํ วิชเฏตา อญฺญตฺร มยา”ติ ปสฺสนฺตานํ ฯเปฯ.}

\prose{“ตณฺหาสํฆาฏปฏิมุกฺโก โลกสนฺนิวาโส”ติ ปสฺสนฺตานํ ฯเปฯ “ตณฺหาชาเลน โอตฺถโฏ\csromansharedfootnote{99f19a78d4b67832}{โอตฺถโต (ก)} โลกสนฺนิวาโส”ติ ปสฺสนฺตานํ ฯเปฯ “ตณฺหาโสเตน วุยฺหติ โลกสนฺนิวาโส”ติ ปสฺสนฺตานํ ฯเปฯ “ตณฺหา\-สญฺโญชเนน สญฺญุตฺโต โลกสนฺนิวาโส”ติ ปสฺสนฺตานํ ฯเปฯ “ตณฺหานุสเยน อนุสโฏ โลกสนฺนิวาโส”ติ ปสฺสนฺตานํ ฯเปฯ “ตณฺหาสนฺตาเปน สนฺตปฺปติ โลกสนฺนิวาโส”ติ ปสฺสนฺตานํ ฯเปฯ “ตณฺหา\-ปริฬาเหน ปริฑยฺหติ โลกสนฺนิวาโส”ติ ปสฺสนฺตานํ ฯเปฯ.}

\prose{“ทิฏฺฐิสํฆาฏปฏิมุกฺโก โลกสนฺนิวาโส”ติ ปสฺสนฺตานํ ฯเปฯ “ทิฏฺฐิชาเลน โอตฺถโฏ โลกสนฺนิวาโส”ติ ปสฺสนฺตานํ ฯเปฯ “ทิฏฺฐิโสเตน วุยฺหติ โลกสนฺนิวาโส”ติ ปสฺสนฺตานํ ฯเปฯ “ทิฏฺฐิ\-สญฺโญชเนน สญฺญุตฺโต โลกสนฺนิวาโส”ติ ปสฺสนฺตานํ ฯเปฯ “ทิฏฺฐานุสเยน อนุสโฏ โลกสนฺนิวาโส”ติ ปสฺสนฺตานํ ฯเปฯ “ทิฏฺฐิ\-สนฺตาเปน สนฺตปฺปติ โลกสนฺนิวาโส”ติ ปสฺสนฺตานํ ฯเปฯ “ทิฏฺฐิ\-ปริฬาเหน ปริฑยฺหติ โลกสนฺนิวาโส”ติ ปสฺสนฺตานํวฺ ฯเปฯ. \csromanfolio{123}“ชาติยา อนุคโต โลกสนฺนิวาโส”ติ ปสฺสนฺตานํ ฯเปฯ “ชราย อนุสโฏ โลกสนฺนิวาโส”ติ ปสฺสนฺตานํ ฯเปฯ “พฺยาธินา อภิภูโต โลกสนฺนิวาโส”ติ ปสฺสนฺตานํ ฯเปฯ “มรเณน อพฺภาหโต โลกสนฺนิวาโส”ติ ปสฺสนฺตานํ ฯเปฯ “ทุกฺเข ปติฏฺฐิโต โลกสนฺนิวาโส”ติ ปสฺสนฺตานํ ฯเปฯ.}

\prose{“ตณฺหาย อุฑฺฑิโต โลกสนฺนิวาโส”ติ ปสฺสนฺตานํ ฯเปฯ “ชรา\-ปาการปริกฺขิตฺโต โลกสนฺนิวาโส”ติ ปสฺสนฺตานํ ฯเปฯ “มจฺจุปาเสน ปริกฺขิตฺโต โลกสนฺนิวาโส”ติ ปสฺสนฺตานํ ฯเปฯ “มหาพนฺธน\-พทฺโธ โลกสนฺนิวาโส ราคพนฺธเนน โทสพนฺธเนน โมหพนฺธเนน มานพนฺธเนน ทิฏฺฐิ\-พนฺธเนน กิเลส\-พนฺธเนน ทุจฺจริต\-พนฺธเนน, ตสฺส นตฺถญฺโญ โกจิ พนฺธนํ โมเจตา อญฺญตฺร มยา”ติ ปสฺสนฺตานํ ฯเปฯ “มหา\-สมฺพาธ\-ปฺปฏิปนฺโน โลกสนฺนิวาโส, ตสฺส นตฺถญฺโญ โกจิ โอกาสํ ทสฺเสตา อญฺญตฺร มยา”ติ ปสฺสนฺตานํ. “มหา\-ปลิโพเธน ปลิพุทฺโธ โลกสนฺนิวาโส, ตสฺส นตฺถญฺโญ โกจิ ปลิโพธํ เฉตา อญฺญตฺร มยา”ติ ปสฺสนฺตานํ ฯเปฯ “มหาปปาเต ปติโต โลกสนฺนิวาโส, ตสฺส นตฺถญฺโญ โกจิ ปปาตา อุทฺธตา อญฺญตฺร มยา”ติ ปสฺสนฺตานํ ฯเปฯ “มหา\-กนฺตาร\-ปฺปฏิปนฺโน โลกสนฺนิวาโส, ตสฺส นตฺถญฺโญ โกจิ กนฺตารํ ตาเรตา อญฺญตฺร มยา”ติ ปสฺสนฺตานํ ฯเปฯ “มหา\-สํสาร\-ปฺปฏิปนฺโน โลกสนฺนิวาโส, ตสฺส นตฺถญฺโญ โกจิ สํสารา โมเจตา อญฺญตฺร มยา”ติ ปสฺสนฺตานํ ฯเปฯ “มหาวิทุคฺเค สมฺปริวตฺตติ โลกสนฺนิวาโส, ตสฺส นตฺถญฺโญ โกจิ วิทุคฺคา อุทฺธตา อญฺญตฺร มยา”ติ ปสฺสนฺตานํ ฯเปฯ “มหาปลิเป\csromansharedfootnote{3c201226260d106b}{มหาปลฺเลเป (สฺยา)} ปลิปนฺโน โลกสนฺนิวาโส, ตสฺส นตฺถญฺโญ โกจิ ปลิปา อุทฺธตา อญฺญตฺร มยา”ติ ปสฺสนฺตานํ ฯเปฯ.}

\prose{“อพฺภาหโต โลกสนฺนิวาโส”ติ ปสฺสนฺตานํ ฯเปฯ “อาทิตฺโต โลกสนฺนิวาโส ราคคฺคินา โทสคฺคินา โมหคฺคินา ชาติยา ชราย มรเณน โสเกหิ ปริเทเวหิ ทุกฺเขหิ โทมนสฺเสหิ อุปายาเสหิ, ตสฺส นตฺถญฺโญ โกจิ นิพฺพาเปตา อญฺญตฺร มยา”ติ ปสฺสนฺตานํ ฯเปฯ “อุนฺนีตโก โลกสนฺนิวาโส หญฺญติ นิจฺจมตาโณ ปตฺตทณฺโฑ ตกฺกโร”ติ ปสฺสนฺตานํ ฯเปฯ “วชฺช\-พนฺธนพทฺโธ โลกสนฺนิวาโส อาฆาตน\-ปจฺจุปฏฺฐิโต, ตสฺส นตฺถญฺโญ โกจิ พนฺธนํ โมเจตา \csromanfolio{124}อญฺญตฺร มยา”ติ ปสฺสนฺตานํ ฯเปฯ “อนาโถ โลกสนฺนิวาโส ปรม\-การุญฺญ\-ปฺปตฺโต, ตสฺส นตฺถญฺโญ โกจิ ตาเยตา อญฺญตฺร มยา”ติ ปสฺสนฺตานํ ฯเปฯ “ทุกฺขาภิตุนฺโน\csromansharedfootnote{97a19570b89cffac}{ทุกฺขาภิตุณฺโณ (ก)} โลกสนฺนิวาโส จิรรตฺตํ ปีฬิโต”ติ ปสฺสนฺตานํ ฯเปฯ “คธิโต โลกสนฺนิวาโส นิจฺจํ ปิปาสิโต”ติ ปสฺสนฺตานํ ฯเปฯ.}

\prose{“อนฺโธ โลกสนฺนิวาโส อจกฺขุโก”ติ ปสฺสนฺตานํ ฯเปฯ “หตเนตฺโต โลกสนฺนิวาโส อปริณายโก”ติ ปสฺสนฺตานํ ฯเปฯ “วิปถ\-ปกฺขนฺโท\csromansharedfootnote{90b34c55a351c8f3}{วิปถํ ปกฺขนฺโต (สฺยา)} โลกสนฺนิวาโส อญฺชสาปรทฺโธ, ตสฺส นตฺถญฺโญ โกจิ อริยปถํ อาเนตา อญฺญตฺร มยา”ติ ปสฺสนฺตานํ ฯเปฯ “มโหฆ\-ปกฺขนฺโท โลกสนฺนิวาโส, ตสฺส นตฺถญฺโญ โกจิ โอฆา อุทฺธตา อญฺญตฺร มยา”ติ ปสฺสนฺตานํ ฯเปฯ.}

\proseitem{118}{“ทฺวีหิ ทิฏฺฐิคเตหิ ปริยุฏฺฐิโต โลกสนฺนิวาโส”ติ ปสฺสนฺตานํ ฯเปฯ “ตีหิ ทุจฺจริเตหิ วิปฺปฏิปนฺโน โลกสนฺนิวาโส”ติ ปสฺสนฺตานํ ฯเปฯ “จตูหิ โยเคหิ ยุตฺโต โลกสนฺนิวาโส จตุโยคโยชิโต”ติ ปสฺสนฺตานํ ฯเปฯ “จตูหิ คนฺเถหิ\csromansharedfootnote{a7e2c17842af9cb3}{คณฺเฐหิ (สฺยา)} คนฺถิโต โลกสนฺนิวาโส”ติ ปสฺสนฺตานํ ฯเปฯ “จตูหิ อุปาทาเนหิ อุปาทียติ โลกสนฺนิวาโส”ติ ปสฺสนฺตานํ ฯเปฯ “ปญฺจ\-คติ\-สมารุโฬฺห โลกสนฺนิวาโส”ติ ปสฺสนฺตานํ ฯเปฯ “ปญฺจหิ กามคุเณหิ รชฺชติ โลกสนฺนิวาโส”ติ ปสฺสนฺตานํ ฯเปฯ “ปญฺจหิ นีวรเณหิ โอตฺถโฏ โลกสนฺนิวาโส”ติ ปสฺสนฺตานํ ฯเปฯ “ฉหิ วิวาทมูเลหิ วิวทติ โลกสนฺนิวาโส”ติ ปสฺสนฺตานํ ฯเปฯ “ฉหิ ตณฺหากาเยหิ รชฺชติ โลกสนฺนิวาโส”ติ ปสฺสนฺตานํ ฯเปฯ “ฉหิ ทิฏฺฐิคเตหิ ปริยุฏฺฐิโต โลกสนฺนิวาโส”ติ ปสฺสนฺตานํ ฯเปฯ “สตฺตหิ อนุสเยหิ อนุสโฏ โลกสนฺนิวาโส”ติ ปสฺสนฺตานํ ฯเปฯ “สตฺตหิ สญฺโญชเนหิ สญฺญุตฺโต โลกสนฺนิวาโส”ติ ปสฺสนฺตานํ ฯเปฯ “สตฺตหิ มาเนหิ อุนฺนโต\csromansharedfootnote{c7572fb77d32cd77}{อุณฺณโต (สฺยา, ก)} โลกสนฺนิวาโส”ติ ปสฺสนฺตานํ ฯเปฯ “อฏฺฐหิ โลกธมฺเมหิ สมฺปริวตฺตติ โลกสนฺนิวาโส”ติ ปสฺสนฺตานํ ฯเปฯ “อฏฺฐหิ มิจฺฉตฺเตหิ นิยฺยาโต โลกสนฺนิวาโส”ติ ปสฺสนฺตานํ ฯเปฯ “อฏฺฐหิ ปุริสโทเสหิ ทุสฺสติ โลกสนฺนิวาโส”ติ \csromanfolio{125}ปสฺสนฺตานํ ฯเปฯ “นวหิ อาฆาตวตฺถูหิ อาฆาติโต โลกสนฺนิวาโส”ติ ปสฺสนฺตานํ ฯเปฯ “นววิธ\-มาเนหิ อุนฺนโต โลกสนฺนิวาโส”ติ ปสฺสนฺตานํ ฯเปฯ “นวหิ ตณฺหามูลเกหิ ธมฺเมหิ รชฺชติ โลกสนฺนีวาโส”ติ ปสฺสนฺตานํ ฯเปฯ “ทสหิ กิเลสวตฺถูหิ กิลิสฺสติ โลกสนฺนิวาโส”ติ ปสฺสนฺตานํ ฯเปฯ “ทสหิ อาฆาตวตฺถูหิ อาฆาติโต โลกสนฺนิวาโส”ติ ปสฺสนฺตานํ ฯเปฯ “ทสหิ อกุสล\-กมฺมปเถหิ สมนฺนาคโต โลกสนฺนิวาโส”ติ ปสฺสนฺตานํ ฯเปฯ “ทสหิ สญฺโญชเนหิ สญฺญุตฺโต โลกสนฺนิวาโส”ติ ปสฺสนฺตานํ ฯเปฯ “ทสหิ มิจฺฉตฺเตหิ นิยฺยาโต โลกสนฺนิวาโส”ติ ปสฺสนฺตานํ ฯเปฯ “ทสวตฺถุกาย มิจฺฉาทิฏฺฐิยา สมนฺนาคโต โลกสนฺนิวาโส”ติ ปสฺสนฺตานํ ฯเปฯ “ทสวตฺถุกาย อนฺตคฺคาหิกาย ทิฏฺฐิยา สมนฺนาคโต โลกสนฺนิวาโส”ติ ปสฺสนฺตานํ ฯเปฯ “อฏฺฐสต\-ตณฺหา\-ปปญฺจ\-สเตหิ ปปญฺจิโต โลกสนฺนิวาโส”ติ ปสฺสนฺตานํ พุทฺธานํ ภควนฺตานํ สตฺเตสุ มหากรุณา โอกฺกมติ. ทฺวาสฏฺฐิยา ทิฏฺฐิคเตหิ ปริยุฏฺฐิโต โลก\-สนฺนิวาโสติ ปสฺสนฺตานํ พุทฺธานํ ภควนฺตานํ สตฺเตสุ มหากรุณา โอกฺกมติ.}

\prosewithcloserruled{“อหญฺจมฺหิ ติณฺโณ, โลโก จ อติณฺโณ. อหญฺจมฺหิ มุตฺโต, โลโก จ อมุตฺโต. อหญฺจมฺหิ ทนฺโต, โลโก จ อทนฺโต. อหญฺจมฺหิ สนฺโต, โลโก จ อสนฺโต. อหญฺจมฺหิ อสฺสตฺโถ, โลโก จ อนสฺสตฺโถ. อหญฺจมฺหิ ปรินิพฺพุโต, โลโก จ อปรินิพฺพุโต. ปโหมิ ขฺวาหํ ติณฺโณ ตาเรตุํ, มุตฺโต โมเจตุํ, ทนฺโต ทเมตุํ, สนฺโต สเมตุํ, อสฺสตฺโถ อสฺสาเสตุํ, ปรินิพฺพุโต ปเร จ ปรินิพฺพาเปตุนฺ”ติ ปสฺสนฺตานํ พุทฺธานํ ภควนฺตานํ สตฺเตสุ มหากรุณา โอกฺกมติ. อิทํ ตถาคตสฺส มหากรุณา\-สมาปตฺติยา ญาณํ.}{มหากรุณา\-ญาณ\-นิทฺเทโส เอกสตฺตติโม.}
\csromanmatikaanchor{mtk.50}
\csromantocmarkref{subsection}{72-73. สพฺพญฺญุตญฺญาณนิทฺเทส}{mtk.50}
\bookmark[dest={mtk.50},level=2]{72-73. สพฺพญฺญุตญฺญาณนิทฺเทส}
\csromanheader{2}{72-73. สพฺพญฺญุตญฺญาณนิทฺเทส}
\proseitem{119}{กตมํ ตถาคตสฺส สพฺพญฺญุต\-ญฺญาณํ? สพฺพํ สงฺขตม\-สงฺขตํ อนวเสสํ ชานาตีติ สพฺพญฺญุต\-ญฺญาณํ, ตตฺถ อาวรณํ นตฺถีติ อนาวรณญาณํ\csromansharedfootnote{9852563c5624dd68}{อนาวรณํ ญาณํ (สฺยา) เอวมุปริปิ.}.}

\proseitem{120}{\csromanfolio{126}อตีตํ สพฺพํ ชานาตีติ สพฺพญฺญุต\-ญฺญาณํ, ตตฺถ อาวรณํ นตฺถีติ อนาวรณญาณํ. อนาคตํ สพฺพํ ชานาตีติ สพฺพญฺญุต\-ญฺญาณํ, ตตฺถ อาวรณํ นตฺถีติ อนาวรณญาณํ. ปจฺจุปฺปนฺนํ สพฺพํ ชานาตีติ สพฺพญฺญุต\-ญฺญาณํ, ตตฺถ อาวรณํ นตฺถีติ อนาวรณญาณํ.}

\prose{จกฺขุ เจว รูปา จ, เอวํ ตํ สพฺพํ ชานาตีติ สพฺพญฺญุต\-ญฺญาณํ, ตตฺถ อาวรณํ นตฺถีติ อนาวรณญาณํ. โสตญฺเจว สทฺทา จ ฯเปฯ. ฆานญฺเจว คนฺธา จ, ชิวฺหา เจว รสา จ. กาโย เจว โผฏฺฐพฺพา จ. มโน เจว ธมฺมา จ, เอวํ ตํ สพฺพํ ชานาตีติ สพฺพญฺญุต\-ญฺญาณํ, ตตฺถ อาวรณํ นตฺถีติ อนาวรณญาณํ.}

\prose{ยาวตา อนิจฺจฏฺฐํ ทุกฺขฏฺฐํ อนตฺตฏฺฐํ, ตํ สพฺพํ ชานาตีติ สพฺพญฺญุต\-ญฺญาณํ, ตตฺถ อาวรณํ นตฺถีติ อนาวรณญาณํ. ยาวตา รูปสฺส อนิจฺจฏฺฐํ ทุกฺขฏฺฐํ อนตฺตฏฺฐํ, ตํ สพฺพํ ชานาตีติ สพฺพญฺญุต\-ญฺญาณํ, ตตฺถ อาวรณํ นตฺถีติ อนาวรณญาณํ. ยาวตา เวทนาย ฯเปฯ. ยาวตา สญฺญาย ฯเปฯ. ยาวตา สงฺขารานํ ฯเปฯ. ยาวตา วิญฺญาณสฺส ฯเปฯ. ยาวตา จกฺขุสฺส ฯเปฯ. ชรามรณสฺส อนิจฺจฏฺฐํ ทุกฺขฏฺฐํ อนตฺตฏฺฐํ, ตํ สพฺพํ ชานาตีติ สพฺพญฺญุต\-ญฺญาณํ, ตตฺถ อาวรณํ นตฺถีติ อนาวรณญาณํ.}

\prose{ยาวตา อภิญฺญาย อภิญฺญฏฺฐํ, ตํ สพฺพํ ชานาตีติ สพฺพญฺญุต\-ญฺญาณํ, ตตฺถ อาวรณํ นตฺถีติ อนาวรณญาณํ. ยาวตา ปริญฺญาย ปริญฺญฏฺฐํ ฯเปฯ. ยาวตา ปหานสฺส\csromansharedfootnote{d457060fb2c7dcda}{ปหานาย (สฺยา)} ปหานฏฺฐํ ฯเปฯ. ยาวตา ภาวนาย ภาวนฏฺฐํ ฯเปฯ. ยาวตา สจฺฉิกิริยาย สจฺฉิกิริยฏฺฐํ, ตํ สพฺพํ ชานาตีติ สพฺพญฺญุต\-ญฺญาณํ, ตตฺถ อาวรณํ นตฺถีติ อนาวรณญาณํ.}

\prose{ยาวตา ขนฺธานํ ขนฺธฏฺฐํ, ตํ สพฺพํ ชานาตีติ สพฺพญฺญุต\-ญฺญาณํ, ตตฺถ อาวรณํ นตฺถีติ อนาวรณญาณํ. ยาวตา ธาตูนํ ธาตุฏฺฐํ ฯเปฯ. ยาวตา อายตนานํ อายตนฏฺฐํ ฯเปฯ. ยาวตา สงฺขตานํ สงฺขตฏฺฐํ ฯเปฯ. ยาวตา อสงฺขตสฺส อสงฺขตฏฺฐํ, ตํ สพฺพํ ชานาตีติ สพฺพญฺญุต\-ญฺญาณํ, ตตฺถ อาวรณํ นตฺถีติ อนาวรณญาณํ.}

\prose{ยาวตา กุสเล ธมฺเม สพฺพํ\csromansharedfootnote{c7c77b14f3fd10d9}{สพฺเพ (อฏฺฐกถา)} ชานาตีติ สพฺพญฺญุต\-ญฺญาณํ, ตตฺถ อาวรณํ นตฺถีติ อนาวรณญาณํ. ยาวตา อกุสเล ธมฺเม ฯเปฯ. ยาวตา อพฺยากเต ธมฺเม. ยาวตา กามาวจเร ธมฺเม. ยาวตา \csromanfolio{127}รูปาวจเร ธมฺเม. ยาวตา อรูปาวจเร ธมฺเม. ยาวตา อปริยาปนฺเน ธมฺเม สพฺพํ\csromansharedfootnote{c7c77b14f3fd10d9}{สพฺเพ (อฏฺฐกถา)} ชานาตีติ สพฺพญฺญุต\-ญฺญาณํ, ตตฺถ อาวรณํ นตฺถีติ อนาวรณญาณํ.}

\prose{ยาวตา ทุกฺขสฺส ทุกฺขฏฺฐํ, ตํ สพฺพํ ชานาตีติ สพฺพญฺญุต\-ญฺญาณํ, ตตฺถ อาวรณํ นตฺถีติ อนาวรณญาณํ. ยาวตา สมุทยสฺส สมุทยฏฺฐํ ฯเปฯ. ยาวตา นิโรธสฺส นิโรธฏฺฐํ ฯเปฯ. ยาวตา มคฺคสฺส มคฺคฏฺฐํ, ตํ สพฺพํ ชานาตีติ สพฺพญฺญุต\-ญฺญาณํ. ตตฺถ อาวรณํ นตฺถีติ อนาวรณญาณํ.}

\prose{ยาวตา อตฺถ\-ปฏิสมฺภิทาย อตฺถปฏิสมฺภิทฏฺฐํ, ตํ สพฺพํ ชานาตีติ สพฺพญฺญุต\-ญฺญาณํ, ตตฺถ อาวรณํ นตฺถีติ อนาวรณญาณํ. ยาวตา ธมฺม\-ปฏิสมฺภิทาย ธมฺมปฏิสมฺภิทฏฺฐํ ฯเปฯ. ยาวตา นิรุตฺติ\-ปฏิสมฺภิทาย นิรุตฺติปฏิสมฺภิทฏฺฐํ ฯเปฯ. ยาวตา ปฏิภาน\-ปฏิสมฺภิทาย ปฏิภานปฏิสมฺภิทฏฺฐํ, ตํ สพฺพํ ชานาตีติ สพฺพญฺญุต\-ญฺญาณํ, ตตฺถ อาวรณํ นตฺถีติ อนาวรณญาณํ.}

\prose{ยาวตา อินฺทฺริยปโรปริยตฺต\-ญาณํ, ตํ สพฺพํ ชานาตีติ สพฺพญฺญุต\-ญฺญาณํ, ตตฺถ อาวรณํ นตฺถีติ อนาวรณญาณํ. ยาวตา สตฺตานํ อาสยานุสเย ญาณํ ฯเปฯ. ยาวตา ยมกปาฏิหีเร ญาณํ ฯเปฯ. ยาวตา มหากรุณา\-สมาปตฺติยา ญาณํ, ตํ สพฺพํ ชานาตีติ สพฺพญฺญุต\-ญฺญาณํ, ตตฺถ อาวรณํ นตฺถีติ อนาวรณญาณํ.}

\prose{ยาวตา สเทวกสฺส โลกสฺส สมารกสฺส สพฺรหฺมกสฺส สสฺสมณพฺรหฺมณิยา ปชาย สเทว\-มนุสฺสาย ทิฏฺฐํ สุตํ มุตํ วิญฺญาตํ ปตฺตํ ปริเยสิตํ อนุวิจริตํ มนสา, ตํ สพฺพํ ชานาตีติ สพฺพญฺญุต\-ญฺญาณํ, ตตฺถ อาวรณํ นตฺถีติ อนาวรณญาณํ.}

\setlength{\csromangathapagewidth}{0pt}
\setlength{\csromangathawakwidth}{0pt}
\csromangathameasuregroup{}{น ตสฺส อทฺทิฏฺฐมิธตฺถิ กิญฺจิ, \\ อโถ อวิญฺญาตมชานิตพฺพํ. \\ สพฺพํ อภิญฺญาสิ ยทตฺถิ เนตฺยํ, \\ ตถาคโต เตน สมนฺตจกฺขูติ\textsuperscript{1}.}
\csromangathameasurewak{น ตสฺส อทฺทิฏฺฐมิธตฺถิ กิญฺจิ, \\ อโถ อวิญฺญาตมชานิตพฺพํ. \\ สพฺพํ อภิญฺญาสิ ยทตฺถิ เนตฺยํ, \\ ตถาคโต เตน สมนฺตจกฺขูติ\textsuperscript{1}.}
\setlength{\csromangathaleftcol}{0pt}
\csromangathagroup{\csromangathastanza{\csromangathawak{\csromangathaitemline{121}{น ตสฺส อทฺทิฏฺฐมิธตฺถิ กิญฺจิ,} \\ \csromangathacontline{อโถ อวิญฺญาตมชานิตพฺพํ.} \\ \csromangathacontline{สพฺพํ อภิญฺญาสิ ยทตฺถิ เนตฺยํ,} \\ \csromangathacontline{ตถาคโต เตน สมนฺต\-จกฺขูติ\csromansharedfootnote{ddd7e3274fd824bd}{ขุ 7. 280; ขุ 8. 178 ปิฏฺเฐสุปิ.}.}}}}

\prose{สมนฺต\-จกฺขูติ เกนฏฺเฐน สมนฺตจกฺขุ? จุทฺทส พุทฺธญาณานิ ทุกฺเข ญาณํ พุทฺธญาณํ, ทุกฺขสมุทเย ญาณํ พุทฺธญาณํ, ทุกฺขนิโรเธ ญาณํ พุทฺธญาณํ, \csromanfolio{128}ทุกฺขนิโรธ\-คามินิยา ปฏิปทาย ญาณํ พุทฺธญาณํ, อตฺถ\-ปฏิสมฺภิเท ญาณํ พุทฺธญาณํ, ธมฺม\-ปฏิสมฺภิเท ญาณํ พุทฺธญาณํ, นิรุตฺติ\-ปฏิสมฺภิเท ญาณํ พุทฺธญาณํ, ปฏิภาน\-ปฏิสมฺภิเท ญาณํ พุทฺธญาณํ, อินฺทฺริยปโรปริยตฺต\-ญาณํ พุทฺธญาณํ, สตฺตานํ อาสยานุสเย ญาณํ พุทฺธญาณํ, ยมกปาฏิหีเร ญาณํ พุทฺธญาณํ, มหากรุณา\-สมาปตฺติยา ญาณํ พุทฺธญาณํ, สพฺพญฺญุต\-ญฺญาณํ พุทฺธญาณํ, อนาวรณญาณํ พุทฺธญาณํ. อิมานิ จุทฺทส พุทฺธญาณานิ. อิเมสํ จุทฺทสนฺนํ พุทฺธญาณานํ อฏฺฐ ญาณานิ สาวก\-สาธารณานิ, ฉ ญาณานิ อสาธารณานิ สาวเกหิ.}

\prose{ยาวตา ทุกฺขสฺส ทุกฺขฏฺโฐ, สพฺโพ ญาโต, อญฺญาโต ทุกฺขฏฺโฐ นตฺถีติ สพฺพญฺญุต\-ญฺญาณํ, ตตฺถ อาวรณํ นตฺถีติ อนาวรณญาณํ. ยาวตา ทุกฺขสฺส ทุกฺขฏฺโฐ, สพฺโพ ญาโต สพฺโพ ทิฏฺโฐ สพฺโพ วิทิโต สพฺโพ สจฺฉิกโต สพฺโพ ผสฺสิโต ปญฺญาย, อผสฺสิโต ปญฺญาย ทุกฺขฏฺโฐ นตฺถีติ สพฺพญฺญุต\-ญฺญาณํ, ตตฺถ อาวรณํ นตฺถีติ อนาวรณญาณํ. ยาวตา สมุทยสฺส สมุทยฏฺโฐ. ยาวตา นิโรธสฺส นิโรธฏฺโฐ. ยาวตา มคฺคสฺส มคฺคฏฺโฐ ฯเปฯ. ยาวตา อตฺถ\-ปฏิสมฺภิทาย อตฺถปฏิสมฺภิท\-ฏฺโฐ. ยาวตา ธมฺม\-ปฏิสมฺภิทาย ธมฺมปฏิสมฺภิท\-ฏฺโฐ. ยาวตา นิรุตฺติ\-ปฏิสมฺภิทาย นิรุตฺติปฏิสมฺภิท\-ฏฺโฐ. ยาวตา ปฏิภาน\-ปฏิสมฺภิทาย ปฏิภานปฏิสมฺภิท\-ฏฺโฐ, สพฺโพ ญาโต สพฺโพ ทิฏฺโฐ สพฺโพ วิทิโต สพฺโพ สจฺฉิกโต สพฺโพ ผสฺสิโต ปญฺญาย, อผสฺสิโต ปญฺญาย ปฏิภานปฏิสมฺภิท\-ฏฺโฐ นตฺถีติ สพฺพญฺญุต\-ญฺญาณํ, ตตฺถ อาวรณํ นตฺถีติ อนาวรณญาณํ.}

\prose{ยาวตา อินฺทฺริยปโรปริยตฺต\-ญาณํ. ยาวตา สตฺตานํ อาสยานุสเย ญาณํ. ยาวตา ยมกปาฏิหีเร ญาณํ. ยาวตา มหากรุณา\-สมาปตฺติยา ญาณํ, สพฺพํ ญาตํ สพฺพํ ทิฏฺฐํ สพฺพํ วิทิตํ สพฺพํ สจฺฉิกตํ สพฺพํ ผสฺสิตํ ปญฺญาย, อผสฺสิตํ ปญฺญาย มหากรุณา\-สมาปตฺติยา ญาณํ นตฺถีติ สพฺพญฺญุต\-ญฺญาณํ, ตตฺถ อาวรณํ นตฺถีติ อนาวรณญาณํ.}

\prose{ยาวตา สเทวกสฺส โลกสฺส สมารกสฺส สพฺรหฺมกสฺส สสฺสมณ\-พฺราหฺมณิยา ปชาย สเทว\-มนุสฺสาย ทิฏฺฐํ สุตํ มุตํ วิญฺญาตํ ปตฺตํ ปริเยสิตํ อนุวิจริตํ มนสา, สพฺพํ ญาตํ สพฺพํ ทิฏฺฐํ สพฺพํ วิทิตํ สพฺพํ สจฺฉิกตํ สพฺพํ ผสฺสิตํ ปญฺญาย, อผสฺสิตํ ปญฺญาย นตฺถีติ สพฺพญฺญุต\-ญฺญาณํ, ตตฺถ อาวรณํ นตฺถีติ อนาวรณญาณํ.}

\setlength{\csromangathapagewidth}{0pt}
\setlength{\csromangathawakwidth}{0pt}
\csromangathameasuregroup{}{น ตสฺส อทฺทิฏฺฐมิธตฺถิ กิญฺจิ, \\ อโถ อวิญฺญาตมชานิตพฺพํ. \\ สพฺพํ อภิญฺญาสิ ยทตฺถิ เนยฺยํ, \\ ตถาคโต เตน สมนฺตจกฺขูติ.}
\csromangathameasurewak{น ตสฺส อทฺทิฏฺฐมิธตฺถิ กิญฺจิ, \\ อโถ อวิญฺญาตมชานิตพฺพํ. \\ สพฺพํ อภิญฺญาสิ ยทตฺถิ เนยฺยํ, \\ ตถาคโต เตน สมนฺตจกฺขูติ.}
\setlength{\csromangathaleftcol}{0pt}
\csromangathagroupwithcloserruled{\csromangathastanza{\csromangathawak{\csromanfolio{129}น ตสฺส อทฺทิฏฺฐมิธตฺถิ กิญฺจิ, \\ อโถ อวิญฺญาตมชานิตพฺพํ. \\ สพฺพํ อภิญฺญาสิ ยทตฺถิ เนยฺยํ, \\ ตถาคโต เตน สมนฺต\-จกฺขูติ.}}}{ญาณกถา นิฏฺฐิตา.}
\proseitem{122}{กา ทิฏฺฐิ, กติ ทิฏฺฐิฏฺฐานานิ, กติ ทิฏฺฐิ\-ปริยุฏฺฐานานิ, กติ ทิฏฺฐิโย กติ ทิฏฺฐา\-ภินิเวสา, กตโม ทิฏฺฐิฏฺฐานสมุคฺฆาโตติ.}

\prose{กา ทิฏฺฐีติ อภินิเวส\-ปรามาโส ทิฏฺฐิ.  กติ ทิฏฺฐิฏฺฐานานีติ อฏฺฐ ทิฏฺฐิฏฺฐานานิ.  กติ ทิฏฺฐิปริยุฏฺฐานานีติ อฏฺฐารส ทิฏฺฐิ\-ปริยุฏฺฐานานิ.  กติ ทิฏฺฐิโยติ โสฬส ทิฏฺฐิโย.  กติ ทิฏฺฐา\-ภินิเวสาติ ตีณิ สตํ ทิฏฺฐา\-ภินิเวสา.  กตโม ทิฏฺฐิฏฺฐาน\-สมุคฺฆาโตติ โสตาปตฺติมคฺโค ทิฏฺฐิฏฺฐาน\-สมุคฺฆาโต.}

\proseitem{123}{กถํ อภินิเวส\-ปรามาโส ทิฏฺฐิ\csromansharedfootnote{72772e7cc62f33cf}{กติ อภินิเวสปรามาโส ทิฏฺฐีติ (สฺยา)}? รูปํ “เอตํ มม, เอโสหมสฺมิ, เอโส เม อตฺตา”ติ อภินิเวส\-ปรามาโส ทิฏฺฐิ. เวทนํ เอตํ มม ฯเปฯ. สญฺญํ เอตํ มม ฯเปฯ. สงฺขาเร เอตํ มม ฯเปฯ. วิญฺญาณํ “เอตํ มม, เอโสหมสฺมิ, เอโส เม อตฺตา”ติ อภินิเวส\-ปรามาโส ทิฏฺฐิ. จกฺขุํ เอตํ มม. โสตํ เอตํ มม. ฆานํ เอตํ มม. ชิวฺหํ เอตํ มม. กายํ เอตํ มม. มนํ “เอตํ มม, เอโสหมสฺมิ, เอโส เม อตฺตา”ติ อภินิเวส\-ปรามาโส ทิฏฺฐิ.  รูเป\csromansharedfootnote{3bc42fa63555e067}{รูปํ (สฺยา) ตถา ปญฺจสุ อารมฺมเณสุ เอกวจเนน.} เอตํ มม. สทฺเท เอตํ มม. คนฺเธ เอตํ มม. รเส เอตํ มม. โผฏฺฐพฺเพ เอตํ มม. ธมฺเม “เอตํ มม, เอโสหมสฺมิ, เอโส เม อตฺตา”ติ อภินิเวส\-ปรามาโส ทิฏฺฐิ.  จกฺขุวิญฺญาณํ เอตํ มม. โสตวิญฺญาณํ เอตํ มม. ฆานวิญฺญาณํ เอตํ มม. ชิวฺหาวิญฺญาณํ เอตํ มม. กายวิญฺญาณํ เอตํ มม. มโนวิญฺญาณํ “เอตํ มม, เอโสหมสฺมิ, เอโส เม อตฺตา”ติ อภินิเวส\-ปรามาโส ทิฏฺฐิ.  จกฺขุ\-สมฺผสฺสํ เอตํ มม. โสตสมฺผสฺสํ เอตํ มม. ฆาน\-สมฺผสฺสํ เอตํ มม. ชิวฺหา\-สมฺผสฺสํ เอตํ มม. \csromanfolio{130}กายสมฺผสฺสํ เอตํ มม. มโนสมฺผสฺสํ “เอตํ มม, เอโสหมสฺมิ, เอโส เม อตฺตา”ติ อภินิเวส\-ปรามาโส ทิฏฺฐิ.  จกฺขุ\-สมฺผสฺสชํ เวทนํ. โสต\-สมฺผสฺสชํ เวทนํ. ฆาน\-สมฺผสฺสชํ เวทนํ. ชิวฺหา\-สมฺผสฺสชํ เวทนํ. กาย\-สมฺผสฺสชํ เวทนํ. มโน\-สมฺผสฺสชํ เวทนํ “เอตํ มม, เอโสหมสฺมิ, เอโส เม อตฺตา”ติ อภินิเวส\-ปรามาโส ทิฏฺฐิ.}

\prose{รูปสญฺญํ เอตํ มม. สทฺทสญฺญํ เอตํ มม. คนฺธสญฺญํ เอตํ มม. รสสญฺญํ เอตํ มม. โผฏฺฐพฺพ\-สญฺญํ เอตํ มม. ธมฺมสญฺญํ “เอตํ มม, เอโสหมสฺมิ, เอโส เม อตฺตา”ติ อภินิเวส\-ปรามาโส ทิฏฺฐิ.  รูปสญฺเจตนํ เอตํ มม. สทฺท\-สญฺเจตนํ เอตํ มม. คนฺธ\-สญฺเจตนํ เอตํ มม. รสสญฺเจตนํ เอตํ มม. โผฏฺฐพฺพ\-สญฺเจตนํ เอตํ มม. ธมฺม\-สญฺเจตนํ “เอตํ มม, เอโสหมสฺมิ, เอโส เม อตฺตา”ติ อภินิเวส\-ปรามาโส ทิฏฺฐิ.  รูปตณฺหํ เอตํ มม. สทฺทตณฺหํ เอตํ มม. คนฺธตณฺหํ เอตํ มม. รสตณฺหํ เอตํ มม. โผฏฺฐพฺพ\-ตณฺหํ เอตํ มม. ธมฺมตณฺหํ “เอตํ มม, เอโสหมสฺมิ, เอโส เม อตฺตา”ติ อภินิเวส\-ปรามาโส ทิฏฺฐิ.  รูปวิตกฺกํ เอตํ มม. สทฺทวิตกฺกํ เอตํ มม. คนฺธวิตกฺกํ เอตํ มม. รสวิตกฺกํ เอตํ มม. โผฏฺฐพฺพ\-วิตกฺกํ เอตํ มม. ธมฺมวิตกฺกํ “เอตํ มม, เอโสหมสฺมิ, เอโส เม อตฺตา”ติ อภินิเวส\-ปรามาโส ทิฏฺฐิ.  รูปวิจารํ เอตํ มม. สทฺทวิจารํ เอตํ มม. คนฺธวิจารํ เอตํ มม. รสวิจารํ เอตํ มม. โผฏฺฐพฺพ\-วิจารํ เอตํ มม. ธมฺมวิจารํ “เอตํ มม, เอโสหมสฺมิ, เอโส เม อตฺตา”ติ อภินิเวส\-ปรามาโส ทิฏฺฐิ.}

\prose{ปถวีธาตุํ เอตํ มม. อาโปธาตุํ เอตํ มม. เตโชธาตุํ เอตํ มม. วาโยธาตุํ เอตํ มม. อากาสธาตุํ เอตํ มม. วิญฺญาณธาตุํ “เอตํ มม, เอโสหมสฺมิ, เอโส เม อตฺตา”ติ อภินิเวส\-ปรามาโส ทิฏฺฐิ.  ปถวีกสิณํ เอตํ มม. อาโปกสิณํ. เตโชกสิณํ. วาโยกสิณํ. นีลกสิณํ. ปีตกสิณํ. โลหิตกสิณํ. โอทาตกสิณํ. อากาสกสิณํ. วิญฺญาณกสิณํ “เอตํ มม, เอโสหมสฺมิ, เอโส เม อตฺตา”ติ อภินิเวส\-ปรามาโส ทิฏฺฐิ.}

\prose{เกสํ เอตํ มม. โลมํ เอตํ มม. นขํ เอตํ มม. ทนฺตํ เอตํ มม. ตจํ เอตํ มม. มํสํ เอตํ มม. นฺหารุํ เอตํ มม. อฏฺฐิํ เอตํ มม. อฏฺฐิมิญฺชํ เอตํ มม. วกฺกํ เอตํ มม. หทยํ เอตํ มม. ยกนํ เอตํ มม. กิโลมกํ เอตํ มม. ปิหกํ เอตํ มม. ปปฺผาสํ เอตํ มม. อนฺตํ เอตํ มม. อนฺตคุณํ เอตํ มม. อุทริยํ เอตํ มม. กรีสํ เอตํ มม. ปิตฺตํ เอตํ มม. เสมฺหํ เอตํ มม. ปุพฺพํ เอตํ \csromanfolio{131}มม. โลหิตํ เอตํ มม. เสทํ เอตํ มม. เมทํ เอตํ มม. อสฺสุํ เอตํ มม. วสํ เอตํ มม. เขฬํ เอตํ มม. สิงฺฆาณิกํ เอตํ มม. ลสิกํ เอตํ มม. มุตฺตํ เอตํ มม. มตฺถลุงฺคํ “เอตํ มม, เอโสหมสฺมิ, เอโส เม อตฺตา”ติ อภินิเวส\-ปรามาโส ทิฏฺฐิ.}

\prose{จกฺขายตนํ เอตํ มม. รูปายตนํ เอตํ มม. โสตายตนํ เอตํ มม. สทฺทายตนํ เอตํ มม. ฆานายตนํ เอตํ มม. คนฺธายตนํ เอตํ มม. ชิวฺหายตนํ เอตํ มม. รสายตนํ เอตํ มม. กายายตนํ เอตํ มม. โผฏฺฐพฺพา\-ยตนํ เอตํ มม. มนายตนํ เอตํ มม. ธมฺมายตนํ เอตํ มม.}

\prose{จกฺขุธาตุํ เอตํ มม. รูปธาตุํ เอตํ มม. จกฺขุวิญฺญาณ\-ธาตุํ เอตํ มม. โสตธาตุํ เอตํ มม. สทฺทธาตุํ เอตํ มม. โสตวิญฺญาณ\-ธาตุํ เอตํ มม. ฆานธาตุํ เอตํ มม. คนฺธธาตุํ เอตํ มม. ฆานวิญฺญาณ\-ธาตุํ เอตํ มม. ชิวฺหาธาตุํ เอตํ มม. รสธาตุํ เอตํ มม. ชิวฺหาวิญฺญาณ\-ธาตุํ เอตํ มม. กายธาตุํ เอตํ มม. โผฏฺฐพฺพ\-ธาตุํ เอตํ มม. กายวิญฺญาณ\-ธาตุํ เอตํ มม. มโนธาตุํ เอตํ มม. ธมฺมธาตุํ เอตํ มม. มโนวิญฺญาณ\-ธาตุํ “เอตํ มม, เอโสหมสฺมิ, เอโส เม อตฺตา”ติ อภินิเวส\-ปรามาโส ทิฏฺฐิ.}

\prose{จกฺขุนฺทฺริยํ เอตํ มม. โสตินฺทฺริยํ เอตํ มม. ฆานินฺทฺริยํ เอตํ มม. ชิวฺหินฺทฺริยํ เอตํ มม. กายินฺทฺริยํ เอตํ มม, มนินฺทฺริยํ เอตํ มม. ชีวิตินฺทฺริยํ เอตํ มม. อิตฺถินฺทฺริยํ เอตํ มม. ปุริสินฺทฺริยํ เอตํ มม. สุขินฺทฺริยํ เอตํ มม. ทุกฺขินฺทฺริยํ เอตํ มม. โสมนสฺสิ\-นฺทฺริยํ เอตํ มม. โทมนสฺสิ\-นฺทฺริยํ เอตํ มม. อุเปกฺขินฺทฺริยํ เอตํ มม. สทฺธินฺทฺริยํ เอตํ มม. วีริยินฺทฺริยํ เอตํ มม. สตินฺทฺริยํ เอตํ มม. สมาธินฺทฺริยํ เอตํ มม. ปญฺญินฺทฺริยํ “เอตํ มม, เอโสหมสฺมิ, เอโส เม อตฺตา”ติ อภินิเวส\-ปรามาโส ทิฏฺฐิ.}

\prose{กามธาตุํ เอตํ มม. รูปธาตุํ เอตํ มม. อรูปธาตุํ เอตํ มม. กามภวํ เอตํ มม. รูปภวํ เอตํ มม. อรูปภวํ เอตํ มม. สญฺญาภวํ เอตํ มม. อสญฺญาภวํ เอตํ มม. เนว\-สญฺญา\-นา\-สญฺญาภวํ เอตํ มม. เอกโวการภวํ เอตํ มม. จตุโวการภวํ เอตํ มม. ปญฺจโวการภวํ เอตํ มม.  ปฐมชฺฌานํ เอตํ มม. ทุติยชฺฌานํ เอตํ มม. ตติยชฺฌานํ เอตํ มม. จตุตฺถ\-ชฺฌานํ เอตํ มม. เมตฺตํ เจโตวิมุตฺติํ เอตํ มม. กรุณํ เจโตวิมุตฺติํ เอตํ มม. มุทิตํ เจโตวิมุตฺติํ เอตํ มม. อุเปกฺขํ เจโตวิมุตฺติํ เอตํ มม.  อากาสา\-นญฺจา\-ยตนสมาปตฺติํ เอตํ มม. \csromanfolio{132}วิญฺญาณญฺจา\-ยตนสมาปตฺติํ เอตํ มม. อากิญฺจญฺญายตน\-สมาปตฺติํ เอตํ มม. เนว\-สญฺญา\-นา\-สญฺญา\-ยตนสมาปตฺติํ “เอตํ มม, เอโสหมสฺมิ, เอโส เม อตฺตา”ติ อภินิเวส\-ปรามาโส ทิฏฺฐิ.}

\proseitem{2}{อวิชฺชํ เอตํ มม. สงฺขาเร เอตํ มม. วิญฺญาณํ เอตํ มม. นามรูปํ เอตํ มม. สฬายตนํ เอตํ มม. ผสฺสํ เอตํ มม. เวทนํ เอตํ มม. ตณฺหํ เอตํ มม. อุปาทานํ เอตํ มม. ภวํ เอตํ มม. ชาติํ เอตํ มม. ชรามรณํ “เอตํ มม, เอโสหมสฺมิ, เอโส เม อตฺตา”ติ อภินิเวส\-ปรามาโส ทิฏฺฐิ. เอวํ อภินิเวส\-ปรามาโส ทิฏฺฐิ.}

\proseitem{124}{กตมานิ \textbf{อฏฺฐ ทิฏฺฐิฏฺฐานานิ?} ขนฺธาปิ ทิฏฺฐิฏฺฐานํ, อวิชฺชาปิ ทิฏฺฐิฏฺฐานํ, ผสฺโสปิ ทิฏฺฐิฏฺฐานํ, สญฺญาปิ ทิฏฺฐิฏฺฐานํ, วิตกฺโกปิ\csromansharedfootnote{283e8f85fb54c13b}{วิตกฺกาปิ (สฺยา)} ทิฏฺฐิฏฺฐานํ, อโยนิโส มนสิกาโรปิ ทิฏฺฐิฏฺฐานํ, ปาปมิตฺโตปิ ทิฏฺฐิฏฺฐานํ, ปรโตโฆโสปิ ทิฏฺฐิฏฺฐานํ.}

\prose{ขนฺธา เหตุ ขนฺธา ปจฺจโย ทิฏฺฐิฏฺฐานํ อุปาทาย สมุฏฺฐาน\-ฏฺเฐน, เอวํ ขนฺธาปิ ทิฏฺฐิฏฺฐานํ. อวิชฺชา เหตุ อวิชฺชา ปจฺจโย ทิฏฺฐิฏฺฐานํ อุปาทาย สมุฏฺฐาน\-ฏฺเฐน, เอวํ อวิชฺชาปิ ทิฏฺฐิฏฺฐานํ. ผสฺโส เหตุ ผสฺโส ปจฺจโย ทิฏฺฐิฏฺฐานํ อุปาทาย สมุฏฺฐาน\-ฏฺเฐน, เอวํ ผสฺโสปิ ทิฏฺฐิฏฺฐานํ. สญฺญา เหตุ สญฺญา ปจฺจโย ทิฏฺฐิฏฺฐานํ อุปาทาย สมุฏฺฐาน\-ฏฺเฐน, เอวํ สญฺญาปิ ทิฏฺฐิฏฺฐานํ. วิตกฺโก เหตุ วิตกฺโก ปจฺจโย ทิฏฺฐิฏฺฐานํ อุปาทาย สมุฏฺฐาน\-ฏฺเฐน, เอวํ วิตกฺโกปิ ทิฏฺฐิฏฺฐานํ. อโยนิโส มนสิกาโร เหตุ อโยนิโส มนสิกาโร ปจฺจโย ทิฏฺฐิฏฺฐานํ อุปาทาย สมุฏฺฐาน\-ฏฺเฐน, เอวํ อโยนิโส มนสิกาโรปิ ทิฏฺฐิฏฺฐานํ. ปาปมิตฺโต เหตุ ปาปมิตฺโต ปจฺจโย ทิฏฺฐิฏฺฐานํ อุปาทาย สมุฏฺฐาน\-ฏฺเฐน, เอวํ ปาปมิตฺโตปิ ทิฏฺฐิฏฺฐานํ. ปรโตโฆโส เหตุ ปรโตโฆโส ปจฺจโย ทิฏฺฐิฏฺฐานํ อุปาทาย สมุฏฺฐาน\-ฏฺเฐน, เอวํ ปรโตโฆโสปิ ทิฏฺฐิฏฺฐานํ. อิมานิ อฏฺฐ ทิฏฺฐิฏฺฐานานิ.}

\proseitem{125}{กตมานิ อฏฺฐารส ทิฏฺฐิ\-ปริยุฏฺฐานานิ? ยา ทิฏฺฐิ ทิฏฺฐิคตํ ทิฏฺฐิคหนํ ทิฏฺฐิกนฺตารํ ทิฏฺฐิวิสูกํ ทิฏฺฐิ\-วิปฺผนฺทิตํ ทิฏฺฐิ\-สญฺโญชนํ ทิฏฺฐิสลฺลํ ทิฏฺฐิสมฺพาโธ ทิฏฺฐิ\-ปลิโพโธ ทิฏฺฐิ\-พนฺธนํ ทิฏฺฐิปปาโต ทิฏฺฐานุสโย ทิฏฺฐิสนฺตาโป \csromanfolio{133}ทิฏฺฐิปริฬาโห ทิฏฺฐิคนฺโถ ทิฏฺฐุปาทานํ ทิฏฺฐา\-ภินิเวโส ทิฏฺฐิปรามาโส อิมานิ อฏฺฐารส ทิฏฺฐิ\-ปริยุฏฺฐานานิ.}

\proseitem{126}{กตมา โสฬส ทิฏฺฐิโย? อสฺสาททิฏฺฐิ อตฺตานุทิฏฺฐิ มิจฺฉาทิฏฺฐิ สกฺกายทิฏฺฐิ สกฺกายวตฺถุกา สสฺสตทิฏฺฐิ สกฺกายวตฺถุกา อุจฺเฉททิฏฺฐิ อนฺตคฺคาหิกา ทิฏฺฐิ ปุพฺพนฺตา\-นุทิฏฺฐิ อปรนฺตา\-นุทิฏฺฐิ สญฺโญชนิกา ทิฏฺฐิ “อหนฺ”ติ มานวินิพนฺธาทฺùฐิ “มมนฺ”ติ มานวินิพนฺธา ทิฏฺฐิ อตฺตวาท\-ปฏิสํยุตฺตา ทิฏฺฐิ โลกวาท\-ปฏิสํยุตฺตา ทิฏฺฐิ ภวทิฏฺฐิ วิภวทิฏฺฐิ, อิมา โสฬส ทิฏฺฐิโย.}

\proseitem{127}{กตเม ตีณิ สตํ ทิฏฺฐา\-ภินิเวสา? อสฺสาททิฏฺฐิยา กติหากาเรหิ อภินิเวโส โหติ, อตฺตานุทิฏฺฐิยา กติหากาเรหิ อภินิเวโส โหติ, มิจฺฉาทิฏฺฐิยา กติหากาเรหิ อภินิเวโส โหติ, สกฺกาย\-ทิฏฺฐิยา กติหากาเรหิ อภินิเวโส โหติ, สกฺกายวตฺถุกาย สสฺสต\-ทิฏฺฐิยา กติหากาเรหิ อภินิเวโส โหติ, สกฺกายวตฺถุกาย อุจฺเฉท\-ทิฏฺฐิยา กติหากาเรหิ อภินิเวโส โหติ, อนฺตคฺคาหิกาย ทิฏฺฐิยา กติหากาเรหิ อภินิเวโส โหติ, ปุพฺพนฺตา\-นุทิฏฺฐิยา กติหากาเรหิ อภินิเวโส โหติ, อปรนฺตา\-นุทิฏฺฐิยา กติหากาเรหิ อภินิเวโส โหติ, สญฺโญชนิกาย ทิฏฺฐิยา กติหากาเรหิ อภินิเวโส โหติ, “อหนฺ”ติ มาน\-วินิพนฺธาย ทิฏฺฐิยา กติหากาเรหิ อภินิเวโส โหติ, “มมนฺ”ติ มาน\-วินิพนฺธาย ทิฏฺฐิยา กติหากาเรหิ อภินิเวโส โหติ, อตฺตวาท\-ปฏิสํยุตฺตาย ทิฏฺฐิยา กติหากาเรหิ อภินิเวโส โหติ, โลกวาท\-ปฏิสํยุตฺตาย ทิฏฺฐิยา กติหากาเรหิ อภินิเวโส โหติ, ภวทิฏฺฐิยา กติหากาเรหิ อภินิเวโส โหติ, วิภว\-ทิฏฺฐิยา กติหากาเรหิ อภินิเวโส โหติ?}

\prose{อสฺสาททิฏฺฐิยา ปญฺจติํสาย อากาเรหิ อภินิเวโส โหติ, อตฺตานุทิฏฺฐิยา วีสติยา อากาเรหิ อภินิเวโส โหติ, มิจฺฉาทิฏฺฐิยา ทสหากาเรหิ อภินิเวโส โหติ, สกฺกาย\-ทิฏฺฐิยา วีสติยา อากาเรหิ อภินิเวโส โหติ, สกฺกายวตฺถุกาย สสฺสต\-ทิฏฺฐิยา ปนฺนรสหิ อากาเรหิ อภินิเวโส โหติ, สกฺกายวตฺถุกาย อุจฺเฉท\-ทิฏฺฐิยา ปญฺจหากาเรหิ อภินิเวโส โหติ, อนฺตคฺคาหิกาย ทิฏฺฐิยา ปญฺญาสาย อากาเรหิ อภินิเวโส \csromanfolio{134}โหติ, ปุพฺพนฺตา\-นุทิฏฺฐิยา อฏฺฐารสหิ อากาเรหิ อภินิเวโส โหติ, อปรนฺตา\-นุทิฏฺฐิยา จตุจตฺตาลีสาย อากาเรหิ อภินิเวโส โหติ, สญฺโญชนิกาย ทิฏฺฐิยา อฏฺฐารสหิ อากาเรหิ อภินิเวโส โหติ, “อหนฺ”ติ มาน\-วินิพนฺธาย ทิฏฺฐิยา อฏฺฐารสหิ อากาเรหิ อภินิเวโส โหติ, “มมนฺ”ติ มาน\-วินิพนฺธาย ทิฏฺฐิยา อฏฺฐารสหิ อากาเรหิ อภินิเวโส โหติ, อตฺตวาท\-ปฏิสํยุตฺตาย ทิฏฺฐิยา วีสติยา อากาเรหิ อภินิเวโส โหติ, โลกวาท\-ปฏิสํยุตฺตาย ทิฏฺฐิยา อฏฺฐหิ อากาเรหิ อภินิเวโส โหติ, ภวทิฏฺฐิยา เอเกน อากาเรน\csromansharedfootnote{7ff9fb6891cbfcbf}{เอกูนวีสติยา อากาเรหิ (สฺยา)} อภินิเวโส โหติ, วิภว\-ทิฏฺฐิยา เอเกน อากาเรน\csromansharedfootnote{7ff9fb6891cbfcbf}{เอกูนวีสติยา อากาเรหิ (สฺยา)} อภินิเวโส โหติ.}

\csromanmatikaanchor{mtk.51}
\csromanmatikaanchor{mtk.52}
\csromantocmarknopage{section}{2. ทิฏฺฐิกถา}{mtk.51}
\bookmark[dest={mtk.51},level=1]{2. ทิฏฺฐิกถา}
\csromantocmarkref{subsection}{1. อสฺสาททิฏฺฐินิทฺเทส}{mtk.52}
\bookmark[dest={mtk.52},level=2]{1. อสฺสาททิฏฺฐินิทฺเทส}
\csromanheader{2}{1. อสฺสาททิฏฺฐิ\-นิทฺเทส}
\proseitem{128}{อสฺสาททิฏฺฐิยา กตเมหิ ปญฺจติํสาย อากาเรหิ อภินิเวโส โหติ? ยํ รูปํ ปฏิจฺจ อุปฺปชฺชติ สุขํ โสมนสฺสํ, “อยํ รูปสฺส อสฺสาโท”ติ อภินิเวส\-ปรามาโส ทิฏฺฐิ. ทิฏฺฐิ น อสฺสาโท, อสฺสาโท น ทิฏฺฐิ. อญฺญา ทิฏฺฐิ, อญฺโญ อสฺสาโท. ยา จ ทิฏฺฐิ โย จ อสฺสาโท, อยํ วุจฺจติ อสฺสาททิฏฺฐิ.}

\prose{อสฺสาททิฏฺฐิ มิจฺฉาทิฏฺฐิ ทิฏฺฐิวิปตฺติ, ตาย ทิฏฺฐิ\-วิปตฺติยา สมนฺนาคโต ปุคฺคโล ทิฏฺฐิวิปนฺโน. ทิฏฺฐิวิปนฺโน ปุคฺคโล น เสวิตพฺโพ น ภชิตพฺโพ น ปยิรุปาสิตพฺโพ, ตํ กิสฺส เหตุ, ทิฏฺฐิ หิสฺส ปาปิกา. โย ทิฏฺฐิยา ราโค\csromansharedfootnote{1c5a23557fc00ddc}{ยา ทิฏฺฐิ โย ราโค (สฺยา)}, โส น ทิฏฺฐิ, ทิฏฺฐิ น ราโค. อญฺญา ทิฏฺฐิ, อญฺโญ ราโค. ยา จ ทิฏฺฐิ โย จ ราโค, อยํ วุจฺจติ ทิฏฺฐิราโค. ตาย จ ทิฏฺฐิยา เตน จ ราเคน สมนฺนาคโต ปุคฺคโล ทิฏฺฐิ\-ราครตฺโต. ทิฏฺฐิ\-ราครตฺเต ปุคฺคเล ทินฺนํ ทานํ น มหปฺผลํ โหติ น มหานิสํสํ. ตํ กิสฺส เหตุ, ทิฏฺฐิ หิสฺส ปาปิกา, อสฺสาททิฏฺฐิ มิจฺฉาทิฏฺฐิ.}

\prose{มิจฺฉา\-ทิฏฺฐิกสฺส ปุริส\-ปุคฺคลสฺส เทฺวว คติโย นิรโย วา ติรจฺฉานโยนิ วา. มิจฺฉา\-ทิฏฺฐิกสฺส ปุริส\-ปุคฺคลสฺส ยญฺเจว กายกมฺมํ ยถาทิฏฺฐิ สมตฺตํ สมาทินฺนํ, ยญฺจ วจีกมฺมํ ฯเปฯ ยญฺจ มโนกมฺมํ ยถาทิฏฺฐิ สมตฺตํ สมาทินฺนํ, ยา จ เจตนา ยา จ ปตฺถนา โย จ ปณิธิ เย จ สงฺขารา, \csromanfolio{135}สพฺเพ เต ธมฺมา อนิฏฺฐาย อกนฺตาย อมนาปาย อหิตาย ทุกฺขาย สํวตฺตนฺติ. ตํ กิสฺส เหตุ, ทิฏฺฐิ หิสฺส ปาปิกา. เสยฺยถาปิ นิมฺพพีชํ วา โกสาตกีพีชํ วา ติตฺตกาลาพุพีชํ วา อลฺลาย ปถวิยา นิกฺขิตฺตํ ยญฺเจว ปถวิรสํ อุปาทิยติ, ยญฺจ อาโปรสํ อุปาทิยติ, สพฺพํ ตํ ติตฺตกตฺตาย กฏุกตฺตาย อสาตตฺตาย\csromansharedfootnote{9c79c28e2f48118e}{อสารตาย (สฺยา)} สํวตฺตติ. ตํ กิสฺส เหตุ, พีชํ หิสฺส ปาปิกํ. เอวเมวํ มิจฺฉา\-ทิฏฺฐิกสฺส ปุริส\-ปุคฺคลสฺส ยญฺเจว กายกมฺมํ ยถาทิฏฺฐิ สมตฺตํ สมาทินฺนํ, ยญฺจ วจีกมฺมํ ฯเปฯ ยญฺจ มโนกมฺมํ ยถาทิฏฺฐิ สมตฺตํ สมาทินฺนํ, ยา จ เจตนา ยา จ ปตฺถนา โย จ ปณิธิ เย จ สงฺขารา, สพฺเพ เต ธมฺมา อนิฏฺฐาย อกนฺตาย อมนาปาย อหิตาย ทุกฺขาย สํวตฺตนฺติ. ตํ กิสฺส เหตุ, ทิฏฺฐิ หิสฺส ปาปิกา. อสฺสาททิฏฺฐิ มิจฺฉาทิฏฺฐิ.}

\prose{มิจฺฉาทิฏฺฐิ ทิฏฺฐิคตํ ทิฏฺฐิคหนํ ทิฏฺฐิกนฺตารํ ทิฏฺฐิวิสูกํ ทิฏฺฐิ\-วิปฺผนฺทิตํ ทิฏฺฐิ\-สญฺโญชนํ ทิฏฺฐิสลฺลํ ทิฏฺฐิสมฺพาโธ ทิฏฺฐิ\-ปลิโพโธ ทิฏฺฐิ\-พนฺธนํ ทิฏฺฐิปปาโต ทิฏฺฐานุสโย ทิฏฺฐิสนฺตาโป ทิฏฺฐิปริฬาโห ทิฏฺฐิคนฺโถ ทิฏฺฐุปาทานํ ทิฏฺฐา\-ภินิเวโส ทิฏฺฐิปรามาโส, อิเมหิ อฏฺฐารสหิ อากาเรหิ ปริยุฏฺฐิต\-จิตฺตสฺส สญฺโญโค.}

\proseitem{129}{อตฺถิ สญฺโญชนานิ เจว ทิฏฺฐิโย จ, อตฺถิ สญฺโญชนานิ น จ ทิฏฺฐิโย. กตมานิ สญฺโญชนานิ เจว ทิฏฺฐิโย จ? สกฺกายทิฏฺฐิ สีลพฺพต\-ปรามาโส, อิมานิ สญฺโญชนานิ เจว ทิฏฺฐิโย จ.  กตมานิ สญฺโญชนานิ น จ ทิฏฺฐิโย? กามราค\-สญฺโญชนํ ปฏิฆ\-สญฺโญชนํ มานสญฺโญชนํ วิจิกิจฺฉา\-สญฺโญชนํ ภวราค\-สญฺโญชนํ อิสฺสาสญฺโญชนํ มจฺฉริย\-สญฺโญชนํ อนุนย\-สญฺโญชนํ อวิชฺชา\-สญฺโญชนํ, อิมานิ สญฺโญชนานิ น จ ทิฏฺฐิโย.}

\prose{ยํ เวทนํ ปฏิจฺจ ฯเปฯ. ยํ สญฺญํ ปฏิจฺจ ฯเปฯ. ยํ สงฺขาเร\csromansharedfootnote{df9d96d09c29f165}{เย สงฺขาเร (ก) สํ 2. 24 ปิฏฺเฐ ปสฺสิตพฺพา.} ปฏิจฺจ ฯเปฯ. ยํ วิญฺญาณํ ปฏิจฺจ.  ยํ จกฺขุํ ปฏิจฺจ. ยํ โสตํ ปฏิจฺจ. ยํ ฆานํ ปฏิจฺจ. ยํ ชิวฺหํ ปฏิจฺจ. ยํ กายํ ปฏิจฺจ. ยํ มนํ ปฏิจฺจ. ยํ รูเป ปฏิจฺจ. ยํ สทฺเท ปฏิจฺจ. ยํ คนฺเธ ปฏิจฺจ. ยํ รเส ปฏิจฺจ. ยํ โผฏฺฐพฺเพ ปฏิจฺจ. ยํ ธมฺเม ปฏิจฺจ.  ยํ จกฺขุวิญฺญาณํ ปฏิจฺจ. ยํ โสตวิญฺญาณํ ปฏิจฺจ. ยํ ฆานวิญฺญาณํ ปฏิจฺจ. ยํ ชิวฺหาวิญฺญาณํ ปฏิจฺจ. ยํ กายวิญฺญาณํ ปฏิจฺจ. ยํ มโนวิญฺญาณํ \csromanfolio{136}ปฏิจฺจ.  ยํ จกฺขุ\-สมฺผสฺสํ ปฏิจฺจ. ยํ โสตสมฺผสฺสํ ปฏิจฺจ. ยํ ฆาน\-สมฺผสฺสํ ปฏิจฺจ. ยํ ชิวฺหา\-สมฺผสฺสํ ปฏิจฺจ. ยํ กายสมฺผสฺสํ ปฏิจฺจ. ยํ มโนสมฺผสฺสํ ปฏิจฺจ.  ยํ จกฺขุ\-สมฺผสฺสชํ เวทนํ ปฏิจฺจ. ยํ โสต\-สมฺผสฺสชํ เวทนํ ปฏิจฺจ. ยํ ฆาน\-สมฺผสฺสชํ เวทนํ ปฏิจฺจ. ยํ ชิวฺหา\-สมฺผสฺสชํ เวทนํ ปฏิจฺจ. ยํ กาย\-สมฺผสฺสชํ เวทนํ ปฏิจฺจ. ยํ มโน\-สมฺผสฺสชํ เวทนํ ปฏิจฺจ อุปฺปชฺชติ สุขํ โสมนสฺสํ “อยํ มโน\-สมฺผสฺสชาย เวทนาย อสฺสาโท”ติ อภินิเวส\-ปรามาโส ทิฏฺฐิ. ทิฏฺฐิ น อสฺสาโท, อสฺสาโท น ทิฏฺฐิ. อญฺญา ทิฏฺฐิ, อญฺโญ อสฺสาโท. ยา จ ทิฏฺฐิ โย จ อสฺสาโท, อยํ วุจฺจติ อสฺสาททิฏฺฐิ.}

\prose{อสฺสาททิฏฺฐิ มิจฺฉาทิฏฺฐิ ทิฏฺฐิวิปตฺติ, ตาย ทิฏฺฐิ\-วิปตฺติยา สมานฺนาคโต ปุคฺคโล ทิฏฺฐิวิปนฺโน. ทิฏฺฐิวิปนฺโน ปุคฺคโล น เสวิตพฺโพ น ภชิตพฺโพ น ปยิรุปาสิตพฺโพ. ตํ กิสฺส เหตุ, ทิฏฺฐิ หิสฺส ปาปิกา. โย ทิฏฺฐิยา ราโค, โส น ทิฏฺฐิ, ทิฏฺฐิ น ราโค. อญฺญา ทิฏฺฐิ, อญฺโญ ราโค. ยา จ ทิฏฺฐิ โย จ ราโค, อยํ วุจฺจติ ทิฏฺฐิราโค. ตาย จ ทิฏฺฐิยา เตน จ ราเคน สมนฺนาคโต ปุคฺคโล ทิฏฺฐิ\-ราครตฺโต. ทิฏฺฐิ\-ราครตฺเต ปุคฺคเล ทินฺนํ ทานํ น มหปฺผลํ โหติ น มหานิสํสํ. ตํ กิสฺส เหสุ, ทิฏฺฐิ หิสฺส ปาปิกา. อสฺสาททิฏฺฐิ มิจฺฉาทิฏฺฐิ.}

\prose{มิจฺฉา\-ทิฏฺฐิกสฺส ปุริส\-ปุคฺคลสฺส เทฺวว คติโย นิรโย วา ติรจฺฉานโยนิ วา. มิจฺฉา\-ทิฏฺฐิกสฺส ปุริส\-ปุคฺคลสฺส ยญฺเจว กายกมฺมํ ยถาทิฏฺฐิ สมตฺตํ สมาทินฺนํ, ยญฺจ วจีกมฺมํ ฯเปฯ ยญฺจ มโนกมฺมํ ยถาทิฏฺฐิ สมตฺตํ สมาทินฺนํ, ยา จ เจตนา ยา จ ปตฺถนา โย จ ปณิธิ เย จ สงฺขารา, สพฺเพ เต ธมฺมา อนิฏฺฐาย อกนฺตาย อมนาปาย อหิตาย ทุกฺขาย สํวตฺตนฺติ. ตํ กิสฺส เหตุ, ทิฏฺฐิ หิสฺส ปาปิกา. เสยฺยถาปิ นิมฺพพีชํ วา โกสาตกีพีชํ วา ติตฺตกาลาพุพีชํ วา อลฺลาย ปถวิยา นิกฺขิตฺตํ ยญฺเจว ปถวิรสํ อุปาทิยติ, ยญฺจ อาโปรสํ อุปาทิยติ, สพฺพํ ตํ ติตฺตกตฺตาย กฏุกตฺตาย อสาตตฺตาย สํวตฺตติ. ตํ กิสฺส เหตุ, พีชํ หิสฺส ปาปิกํ. เอวเมวํ มิจฺฉา\-ทิฏฺฐิกสฺส ปุริส\-ปุคฺคลสฺส ยญฺเจว กายกมฺมํ ยถาทิฏฺฐิ สมตฺตํ สมาทินฺนํ, ยญฺจ วจีกมฺมํ ฯเปฯ ยญฺจ มโนกมฺมํ ยถาทิฏฺฐิ สมตฺตํ สมาทินฺนํ, ยา จ เจตนา ยา จ ปตฺถนา โย จ ปณิธิ เย จ สงฺขารา, สพฺเพ เต ธมฺมา อนิฏฺฐาย อกนฺตาย \csromanfolio{137}อมนาปาย อหิตาย ทุกฺขาย สํวตฺตนฺติ. ตํ กิสฺส เหตุ, ทิฏฺฐิ หิสฺส ปาปิกา. อสฺสาททิฏฺฐิ มิจฺฉาทิฏฺฐิ.}

\prose{มิจฺฉาทิฏฺฐิ ทิฏฺฐิคตํ ทิฏฺฐิคหนํ ฯเปฯ ทิฏฺฐา\-ภินิเวโส ทิฏฺฐิปรามาโส, อิเมหิ อฏฺฐารสหิ อากาเรหิ ปริยุฏฺฐิต\-จิตฺตสฺส สญฺโญโค.}

\prosewithcloserruled{อตฺถิ สญฺโญชนานิ เจว ทิฏฺฐิโย จ, อตฺถิ สญฺโญชนานิ น จ ทิฏฺฐิโย. กตมานิ สญฺโญชนานิ เจว ทิฏฺฐิโย จ? สกฺกายทิฏฺฐิ สีลพฺพต\-ปรามาโส, อิมานิ สญฺโญชนานิ เจว ทิฏฺฐิโย จ.  กตมานิ สญฺโญชนานิ น จ ทิฏฺฐิโย? กามราค\-สญฺโญชนํ ปฏิฆ\-สญฺโญชนํ มานสญฺโญชนํ วิจิกิจฺฉา\-สญฺโญชนํ ภวราค\-สญฺโญชนํ อิสฺสาสญฺโญชนํ มจฺฉริย\-สญฺโญชนํ อนุนย\-สญฺโญชนํ อวิชฺชา\-สญฺโญชนํ, อิมานิ สญฺโญชนานิ น จ ทิฏฺฐิโย. อสฺสาททิฏฺฐิยา อิเมหิ ปญฺจติํสาย อากาเรหิ อภินิเวโส โหติ.}{อสฺสาททิฏฺฐิ\-นิทฺเทโส ปฐโม.}
\csromanmatikaanchor{mtk.53}
\csromantocmarkref{subsection}{2. อตฺตานุทิฏฺฐินิทฺเทส}{mtk.53}
\bookmark[dest={mtk.53},level=2]{2. อตฺตานุทิฏฺฐินิทฺเทส}
\csromanheader{2}{2. อตฺตานุทิฏฺฐิ\-นิทฺเทส}
\proseitem{130}{อตฺตานุทิฏฺฐิยา กตเมหิ วีสติยา อากาเรหิ อภินิเวโส โหติ? อิธ อสฺสุตวา ปุถุชฺชโน อริยานํ อทสฺสาวี อริยธมฺมสฺส อโกวิโท อริยธมฺเม อวินีโต สปฺปุริสานํ อทสฺสาวี สปฺปุริส\-ธมฺมสฺส อโกวิโท สปฺปุริส\-ธมฺเม อวินีโต รูปํ อตฺตโต สมนุปสฺสติ, รูปวนฺตํ วา อตฺตานํ, อตฺตนิ วา รูปํ, รูปสฺมิํ วา อตฺตานํ.  เวทนํ ฯเปฯ. สญฺญํ ฯเปฯ. สงฺขาเร ฯเปฯ. วิญฺญาณํ อตฺตโต สมนุปสฺสติ, วิญฺญาณวนฺตํ วา อตฺตานํ, อตฺตนิ วา วิญฺญาณํ, วิญฺญาณสฺมิํ วา อตฺตานํ.}

\proseitem{131}{กถํ \textbf{รูปํ อตฺตโต สมนุปสฺสติ}? อิเธกจฺโจ ปถวีกสิณํ อตฺตโต สมนุปสฺสติ “ยํ ปถวีกสิณํ, โส อหํ. โย อหํ, ตํ ปถวีกสิณนฺ”ติ ปถวีกสิณญฺจ อตฺตญฺจ อทฺวยํ สมนุปสฺสติ. เสยฺยถาปิ เตลปฺปทีปสฺส ฌายโต ยา อจฺจิ, โส วณฺโณ. โย วณฺโณ, สา อจฺจีติ อจฺจิญฺจ วณฺณญฺจ อทฺวยํ สมนุปสฺสติ. เอวเมวํ อิเธกจฺโจ ปถวีกสิณํ อตฺตโต สมนุปสฺสติ “ยํ ปถวีกสิณํ, โส อหํ. โย \csromanfolio{138}อหํ, ตํ ปถวีกสิณนฺ”ติ ปถวีกสิณญฺจ อตฺตญฺจ อทฺวยํ สมนุปสฺสติ. อภินิเวส\-ปรามาโส ทิฏฺฐิ, ทิฏฺฐิ น วตฺถุ, วตฺถุ น ทิฏฺฐิ. อญฺญา ทิฏฺฐิ, อญฺญํ วตฺถุ. ยา จ ทิฏฺฐิ ยญฺจ วตฺถุ, อยํ ปฐมา รูปวตฺถุกา อตฺตานุทิฏฺฐิ. อตฺตานุทิฏฺฐิ มิจฺฉาทิฏฺฐิ ทิฏฺฐิวิปตฺติ ฯเปฯ อตฺตานุทิฏฺฐิ มิจฺฉาทิฏฺฐิ. มิจฺฉา\-ทิฏฺฐิกสฺส ปุริส\-ปุคฺคลสฺส เทฺวว คติโย ฯเปฯ อิมานิ สญฺโญชนานิ น จ ทิฏฺฐิโย.}

\prose{อิเธกจฺโจ อาโปกสิณํ. เตโชกสิณํ. วาโยกสิณํ. นีลกสิณํ. ปีตกสิณํ. โลหิตกสิณํ. โอทาตกสิณํ อตฺตโต สมนุปสฺสติ “ยํ โอทาตกสิณํ, โส อหํ. โย อหํ, ตํ โอทาตกสิณนฺ”ติ โอทาตกสิณญฺจ, อตฺตญฺจ อทฺวยํ สมนุปสฺสติ. เสยฺยถาปิ เตลปฺปทีปสฺส ฌายโต ยา อจฺจิ, โส วณฺโณ. โย วณฺโณ, สา อจฺจีติ อจฺจิญฺจ วณฺณญฺจ อทฺวยํ สมนุปสฺสติ. เอวเมวํ อิเธกจฺโจ ฯเปฯ. โอทาตกสิณญฺจ อตฺตญฺจ อทฺวยํ สมนุปสฺสติ. อภินิเวส\-ปรามาโส ทิฏฺฐิ, ทิฏฺฐิ น วตฺถุ, วตฺถุ น ทิฏฺฐิ. อญฺญา ทิฏฺฐิ, อญฺญํ วตฺถุ. ยา จ ทิฏฺฐิ ยญฺจ วตฺถุ, อยํ ปฐมา รูปวตฺถุกา อตฺตานุทิฏฺฐิ. อตฺตานุทิฏฺฐิ มิจฺฉาทิฏฺฐิ ทิฏฺฐิวิปตฺติ ฯเปฯ อิมานิ สญฺโญชนานิ น จ ทิฏฺฐิโย. เอวํ รูปํ อตฺตโต สมนุปสฺสติ. (1)}

\prose{กถํ รูปวนฺตํ อตฺตานํ สมนุปสฺสติ? อิเธกจฺโจ เวทนํ สญฺญํ สงฺขาเร วิญฺญาณํ อตฺตโต สมนุปสฺสติ, ตสฺส เอวํ โหติ “อยํ โข เม อตฺตา, โส โข ปน เม อยํ อตฺตา อิมินา รูเปน รูปวา”ติ รูปวนฺตํ อตฺตานํ สมนุปสฺสติ. เสยฺยถาปิ รุกฺโข ฉายาสมฺปนฺโน อสฺส. ตเมนํ ปุริโส เอวํ วเทยฺย “อยํ รุกฺโข, อยํ ฉายา. อญฺโญ รุกฺโข, อญฺญา ฉายา. โส โข ปนายํ รุกฺโข อิมาย ฉายาย ฉายาวา”ติ ฉายาวนฺตํ รุกฺขํ สมนุปสฺสติ. เอวเมวํ อิเธกจฺโจ เวทนํ. สญฺญํ. สงฺขาเร. วิญฺญาณํ อตฺตโต สมนุปสฺสติ, ตสฺส เอวํ โหติ “อยํ โข เม อตฺตา, โส โข ปน อยํ อตฺตา อิมินา รูเปน รูปวา”ติ รูปวนฺตํ อตฺตานํ สมนุปสฺสติ. อภินิเวส\-ปรามาโส ทิฏฺฐิ, ทิฏฺฐิ น วตฺถุ, วตฺถุ น ทิฏฺฐิ. อญฺญา ทิฏฺฐิ, อญฺญํ วตฺถุ. ยา จ ทิฏฺฐิ ยญฺจ วตฺถุ, อยํ ทุติยา รูปวตฺถุกา อตฺตานุทิฏฺฐิ. อตฺตานุทิฏฺฐิ มิจฺฉาทิฏฺฐิ ทิฏฺฐิวิปตฺติ ฯเปฯ อิมานิ สญฺโญชนานิ น จ ทิฏฺฐิโย. เอวํ รูปวนฺตํ อตฺตานํ สมนุปสฺสติ. (2)}

\prose{\csromanfolio{139}กถํ อตฺตนิ รูปํ สมนุปสฺสติ? อิเธกจฺโจ เวทนํ สญฺญํ สงฺขาเร วิญฺญาณํ อตฺตโต สมนุปสฺสติ, ตสฺส เอวํ โหติ “อยํ โข เม อตฺตา, อิมสฺมิํ จ ปน อตฺตนิ อิทํ รูปนฺ”ติ อตฺตนิ รูปํ สมนุปสฺสติ. เสยฺยถาปิ ปุปฺผํ คนฺธ\-สมฺปนฺนํ อสฺส. ตเมนํ ปุริโส เอวํ วเทยฺย “อิทํ ปุปฺผํ, อยํ คนฺโธ. อญฺญํ ปุปฺผํ, อญฺโญ คนฺโธ. โส โข ปนายํ คนฺโธ อิมสฺมิํ ปุปฺเผ”ติ ปุปฺผสฺมิํ คนฺธํ สมนุปสฺสติ. เอวเมวํ อิเธกจฺโจ เวทนํ สญฺญํ สงฺขาเร วิญฺญาณํ อตฺตโต สมนุปสฺสติ. ตสฺส เอวํ โหติ “อยํ โข เม อตฺตา, อิมสฺมิํ จ ปน อตฺตนิ อิทํ รูปนฺ”ติ อตฺตนิ รูปํ สมนุปสฺสติ. อภินิเวส\-ปรามาโส ทิฏฺฐิ. ทิฏฺฐิ น วตฺถุ, วตฺถุ น ทิฏฺฐิ, อญฺญา ทิฏฺฐิ, อญฺญํ วตฺถุ. ยา จ ทิฏฺฐิ ยญฺจ วตฺถุ, อยํ ตติยา รูปวตฺถุกา อตฺตานุทิฏฺฐิ. อตฺตานุทิฏฺฐิ มิจฺฉาทิฏฺฐิ ทิฏฺฐิวิปตฺติ ฯเปฯ อิมานิ สญฺโญชนานิ น จ ทิฏฺฐิโย. เอวํ อตฺตนิ รูปํ สมนุปสฺสติ. (3)}

\prose{กถํ รูปสฺมิํ อตฺตานํ สมนุปสฺสติ? อิเธกจฺโจ เวทนํ สญฺญํ สงฺขาเร วิญฺญาณํ อตฺตโต สมนุปสฺสติ, ตสฺส เอวํ โหติ “อยํ โข เม อตฺตา, โส โข ปน เม อยํ อตฺตา อิมสฺมิํ รูเป”ติ รูปสฺมิํ อตฺตานํ สมนุปสฺสติ. เสยฺยถาปิ มณิ กรณฺฑเก ปกฺขิตฺโต อสฺส. ตเมนํ ปุริโส เอวํ วเทยฺย “อยํ มณิ, อยํ กรณฺฑโก. อญฺโญ มณิ, อญฺโญ กรณฺฑโก. โส โข ปนายํ มณิ อิมสฺมิํ กรณฺฑเก”ติ กรณฺฑกสฺมิํ มณิํ สมนุปสฺสติ. เอวเมวํ อิเธกจฺโจ เวทนํ สญฺญํ สงฺขาเร วิญฺญาณํ อตฺตโต สมนุปสฺสติ. ตสฺส เอวํ โหติ “อยํ โข เม อตฺตา, โส โข ปน เม อยํ อตฺตา อิมสฺมิํ รูเป”ติ รูปสฺมิํ อตฺตานํ สมนุปสฺสติ, อภินิเวส\-ปรามาโส ทิฏฺฐิ, ทิฏฺฐิ น วตฺถุ, วตฺถุ น ทิฏฺฐิ. อญฺญา ทิฏฺฐิ, อญฺญํ วตฺถุ. ยา จ ทิฏฺฐิ ยญฺจ วตฺถุ. อยํ จตุตฺถา รูปวตฺถุกา อตฺตานุทิฏฺฐิ. อตฺตานุทิฏฺฐิ มิจฺฉาทิฏฺฐิ ทิฏฺฐิวิปตฺติ ฯเปฯ อิมานิ สญฺโญชนานิ น จ ทิฏฺฐิโย. เอวํ รูปสฺมิํ อตฺตานํ สมนุปสฺสติ. (4)}

\proseitem{132}{กถํ เวทนํ อตฺตโต สมนุปสฺสติ? อิเธกจฺโจ จกฺขุ\-สมฺผสฺสชํ เวทนํ. โสต\-สมฺผสฺสชํ เวทนํ. ฆาน\-สมฺผสฺสชํ เวทนํ. ชิวฺหา\-สมฺผสฺสชํ เวทนํ. กาย\-สมฺผสฺสชํ เวทนํ. มโน\-สมฺผสฺสชํ เวทนํ อตฺตโต สมนุปสฺสติ “ยา มโน\-สมฺผสฺสชา เวทนา, โส อหํ. โย อหํ, สา มโน\-สมฺผสฺสชา เวทนา”ติ มโน\-สมฺผสฺสชํ เวทนญฺจ \csromanfolio{140}อตฺตญฺจ อทฺวยํ สมนุปสฺสติ. เสยฺยถาปิ เตลปฺปทีปสฺส ฌายโต ยา อจฺจิ, โส วณฺโณ. โย วณฺโณ, สา อจฺจีติ อจฺจิญฺจ วณฺณญฺจ อทฺวยํ สมนุปสฺสติ. เอวเมวํ อิเธกจฺโจ มโน\-สมฺผสฺสชํ เวทนํ อตฺตโต สมนุปสฺสติ “ยา มโน\-สมฺผสฺสชา เวทนา, โส อหํ. โย อหํ, สา มโน\-สมฺผสฺสชา เวทนา”ติ มโน\-สมฺผสฺสชํ เวทนญฺจ อตฺตญฺจ อทฺวยํ สมนุปสฺสติ. อภินิเวส\-ปรามาโส ทิฏฺฐิ, ทิฏฺฐิ น วตฺถุ, วตฺถุ น ทิฏฺฐิ. อญฺญา ทิฏฺฐิ, อญฺญํ วตฺถุ. ยา จ ทิฏฺฐิ ยญฺจ วตฺถุ. อยํ ปฐมา เวทนาวตฺถุกา อตฺตานุทิฏฺฐิ. อตฺตานุทิฏฺฐิ มิจฺฉาทิฏฺฐิ ทิฏฺฐิวิปตฺติ ฯเปฯ อิมานิ สญฺโญชนานิ น จ ทิฏฺฐิโย. เอวํ เวทนํ อตฺตโต สมนุปสฺสติ. \mbox{(1, 5)}}

\prose{กถํ เวทนาวนฺตํ อตฺตานํ สมนุปสฺสติ? อิเธกจฺโจ สญฺญํ สงฺขาเร วิญฺญาณํ รูปํ อตฺตโต สมนุปสฺสติ. ตสฺส เอวํ โหติ “อยํ โข เม อตฺตา, โส โข ปน เม อยํ อตฺตา, อิมาย เวทนาย เวทนาวา”ติ เวทนาวนฺตํ อตฺตานํ สมนุปสฺสติ. เสยฺยถาปิ รุกฺโข ฉายาสมฺปนฺโน อสฺส, ตเมนํ ปุริโส เอวํ วเทยฺย “อยํ รุกฺโข, อยํ ฉายา. อญฺโญ รุกฺโข, อญฺญา ฉายา. โส โข ปนายํ รุกฺโข อิมาย ฉายาย ฉายาวา”ติ ฉายาวนฺตํ รุกฺขํ สมนุปสฺสติ. เอวเมวํ อิเธกจฺโจ สญฺญํ สงฺขาเร วิญฺญาณํ รูปํ อตฺตโต สมนุปสฺสติ, ตสฺส เอวํ โหติ “อยํ โข เม อตฺตา, โส โข ปน เม อยํ อตฺตา อิมาย เวทนาย เวทนาวา”ติ เวทนาวนฺตํ อตฺตานํ สมนุปสฺสติ. อภินิเวส\-ปรามาโส ทิฏฺฐิ, ทิฏฺฐิ น วตฺถุ, วตฺถุ น ทิฏฺฐิ. อญฺญา ทิฏฺฐิ, อญฺญํ วตฺถุ, ยา จ ทิฏฺฐิ ยญฺจ วตฺถุ. อยํ ทุติยา เวทนาวตฺถุกา อตฺตานุทิฏฺฐิ. อตฺตานุทิฏฺฐิ มิจฺฉาทิฏฺฐิ ทิฏฺฐิวิปตฺติ ฯเปฯ อิมานิ สญฺโญชนานิ น จ ทิฏฺฐิโย. เอวํ เวทนาวนฺตํ อตฺตานํ สมนุปสฺสติ. \mbox{(2, 6)}}

\prose{กถํ อตฺตนิ เวทนํ สมนุปสฺสติ? อิเธกจฺโจ สญฺญํ สงฺขาเร วิญฺญาณํ รูปํ อตฺตโต สมนุปสฺสติ, ตสฺส เอวํ โหติ “อยํ โข เม อตฺตา, อิมสฺมิญฺจ ปน อตฺตนิ อยํ เวทนา”ติ อตฺตนิ เวทนํ สมนุปสฺสติ. เสยฺยถาปิ ปุปฺผํ คนฺธ\-สมฺปนฺนํ อสฺส, ตเมนํ ปุริโส เอวํ วเทยฺย “อิทํ ปุปฺผํ, อยํ คนฺโธ. อญฺญํ ปุปฺผํ, อญฺโญ คนฺโธ. โส โข ปนายํ คนฺโธ อิมสฺมิํ ปุปฺเผ”ติ ปุปฺผสฺมิํ คนฺธํ สมนุปสฺสติ. เอวเมวํ อิเธกจฺโจ สญฺญํ สงฺขาเร วิญฺญาณํ รูปํ อตฺตโต สมนุปสฺสติ, ตสฺส เอวํ โหติ “อยํ โข เม อตฺตา, อิมสฺมิํ จ ปน อตฺตนิ อยํ เวทนา”ติ อตฺตนิ เวทนํ \csromanfolio{141}สมนุปสฺสติ, อภินิเวส\-ปรามาโส ทิฏฺฐิ, ทิฏฺฐิ น วตฺถุ, วตฺถุ น ทิฏฺฐิ. อญฺญา ทิฏฺฐิ, อญฺญํ วตฺถุ. ยา จ ทิฏฺฐิ ยญฺจ วตฺถุ. อยํ ตติยา เวทนาวตฺถุกา อตฺตานุทิฏฺฐิ. อตฺตานุทิฏฺฐิ มิจฺฉาทิฏฺฐิ ทิฏฺฐิวิปตฺติ ฯเปฯ อิมานิ สญฺโญชนานิ น จ ทิฏฺฐิโย. เอวํ อตฺตนิ เวทนํ สมนุปสฺสติ. \mbox{(3, 7)}}

\prose{กถํ เวทนาย อตฺตานํ สมนุปสฺสติ? อิเธกจฺโจ สญฺญํ สงฺขาเร วิญฺญาณํ รูปํ อตฺตโต สมนุปสฺสติ, ตสฺส เอวํ โหติ “อยํ โข เม อตฺตา, โส โข ปน เม อยํ อตฺตา อิมาย เวทนายา”ติ เวทนาย อตฺตานํ สมนุปสฺสติ. เสยฺยถาปิ มณิ กรณฺฑเก ปกฺขิตฺโต อสฺส. ตเมนํ ปุริโส เอวํ วเทยฺย “อยํ มณิ, อยํ กรณฺฑโก. อญฺโญ มณิ, อญฺโญ กรณฺฑโก. โส โข ปนายํ มณิ อิมสฺมิํ กรณฺฑเก”ติ กรณฺฑกสฺมิํ มณิํ สมนุปสฺสติ. เอวเมวํ อิเธกจฺโจ สญฺญํ สงฺขาเร วิญฺญาณํ รูปํ อตฺตโต สมนุปสฺสติ, ตสฺส เอวํ โหติ “อยํ โข เม อตฺตา, โส โข ปน เม อยํ อตฺตา อิมาย เวทนายา”ติ เวทนาย อตฺตานํ สมนุปสฺสติ. อภินิเวส\-ปรามาโส ทิฏฺฐิ, ทิฏฺฐิ น วตฺถุ, วตฺถุ น ทิฏฺฐิ. อญฺญา ทิฏฺฐิ, อญฺญํ วตฺถุ. ยา จ ทิฏฺฐิ ยญฺจ วตฺถุ. อยํ จตุตฺถา เวทนาวตฺถุกา อตฺตานุทิฏฺฐิ. อตฺตานุทิฏฺฐิ มิจฺฉาทิฏฺฐิ ทิฏฺฐิวิปตฺติ ฯเปฯ อิมานิ สญฺโญชนานิ น จ ทิฏฺฐิโย. เอวํ เวทนาย อตฺตานํ สมนุปสฺสติ. \mbox{(4, 8)}}

\proseitem{133}{กถํ สญฺญํ อตฺตโต สมนุปสฺสติ? อิเธกจฺโจ จกฺขุ\-สมฺผสฺสชํ สญฺญํ. โสต\-สมฺผสฺสชํ สญฺญํ. ฆาน\-สมฺผสฺสชํ สญฺญํ. ชิวฺหา\-สมฺผสฺสชํ สญฺญํ. กาย\-สมฺผสฺสชํ สญฺญํ. มโน\-สมฺผสฺสชํ สญฺญํ อตฺตโต สมนุปสฺสติ “ยา มโน\-สมฺผสฺสชา สญฺญา, โส อหํ. โย อหํ, สา มโน\-สมฺผสฺสชา สญฺญา”ติ มโน\-สมฺผสฺสชํ สญฺญญฺจ อตฺตญฺจ อทฺวยํ สมนุปสฺสติ. เสยฺยถาปิ เตลปฺปทีปสฺส ฌายโต ยา อจฺจิ, โส วณฺโณ. โย วณฺโณ, สา อจฺจีติ อจฺจิญฺจ วณฺณญฺจ อทฺวยํ สมนุปสฺสติ. เอวเมวํ อิเธกจฺโจ มโน\-สมฺผสฺสชํ สญฺญํ อตฺตโต สมนุปสฺสติ “ยา มโน\-สมฺผสฺสชา สญฺญา, โส อหํ. โย อหํ, สา มโน\-สมฺผสฺสชา สญฺญา”ติ มโน\-สมฺผสฺสชํ สญฺญญฺจ อตฺตญฺจ อทฺวยํ สมนุปสฺสติ, อภินิเวส\-ปรามาโส ทิฏฺฐิ, ทิฏฺฐิ น วตฺถุ, วตฺถุ น ทิฏฺฐิ. อญฺญา ทิฏฺฐิ, อญฺญํ วตฺถุ. ยา จ ทิฏฺฐิ ยญฺจ วตฺถุ. อยํ ปฐมา สญฺญาวตฺถุกา อตฺตานุทิฏฺฐิ. \csromanfolio{142}อตฺตานุทิฏฺฐิ มิจฺฉาทิฏฺฐิ. ทิฏฺฐิวิปตฺติ ฯเปฯ อิมานิ สญฺโญชนานิ น จ ทิฏฺฐิโย. เอวํ สญฺญํ อตฺตโต สมนุปสฺสติ. \mbox{(1, 9)}}

\prose{กถํ สญฺญาวนฺตํ อตฺตานํ สมนุปสฺสติ? อิเธกจฺโจ สงฺขาเร วิญฺญาณํ รูปํ เวทนํ อตฺตโต สมนุปสฺสติ, ตสฺส เอวํ โหติ “อยํ โข เม อตฺตา, โส โข ปน เม อยํ อตฺตา อิมาย สญฺญาย สญฺญาวา”ติ สญฺญาวนฺตํ อตฺตานํ สมนุปสฺสติ. เสยฺยถาปิ รุกฺโข ฉายาสมฺปนฺโน อสฺส, ตเมนํ ปุริโส เอวํ วเทยฺย “อยํ รุกฺโข, อยํ ฉายา. อญฺโญ รุกฺโข, อญฺญา ฉายา. โส โข ปนายํ รุกฺโข อิมาย ฉายาย ฉายาวา”ติ ฉายาวนฺตํ รุกฺขํ สมนุปสฺสติ. เอวเมวํ อิเธกจฺโจ สงฺขาเร วิญฺญาณํ รูปํ เวทนํ อตฺตโต สมนุปสฺสติ, ตสฺส เอวํ โหติ “อยํ โข เม อตฺตา, โส โข ปน เม อยํ อตฺตา อิมาย สญฺญาย สญฺญาวา”ติ สญฺญาวนฺตํ อตฺตานํ สมนุปสฺสติ, อภินิเวส\-ปรามาโส ทิฏฺฐิ, ทิฏฺฐิ น วตฺถุ, วตฺถุ น ทิฏฺฐิ. อญฺญา ทิฏฺฐิ, อญฺญํ วตฺถุ. ยา จ ทิฏฺฐิ ยญฺจ วตฺถุ. อยํ ทุติยา สญฺญาวตฺถุกา อตฺตานุทิฏฺฐิ. อตฺตานุทิฏฺฐิ มิจฺฉาทิฏฺฐิ ฯเปฯ อิมานิ สญฺโญชนานิ น จ ทิฏฺฐิโย. เอวํ สญฺญาวนฺตํ อตฺตานํ สมนุปสฺสติ. \mbox{(2, 10)}}

\prose{กถํ อตฺตนิ สญฺญํ สมนุปสฺสติ? อิเธกจฺโจ สงฺขาเร วิญฺญาณํ รูปํ เวทนํ อตฺตโต สมนุปสฺสติ, ตสฺส เอวํ โหติ “อยํ โข เม อตฺตา, อิมสฺมิํ จ ปน อตฺตนิ อยํ สญฺญา”ติ อตฺตนิ สญฺญํ สมนุปสฺสติ. เสยฺยถาปิ ปุปฺผํ คนฺธ\-สมฺปนฺนํ อสฺส, ตเมนํ ปุริโส เอวํ วเทยฺย “อิทํ ปุปฺผํ อยํ คนฺโธ. อญฺญํ ปุปฺผํ, อญฺโญ คนฺโธ. โส โข ปนายํ คนฺโธ อิมสฺมิํ ปุปฺเผ”ติ ปุปฺผสฺมิํ คนฺธํ สมนุปสฺสติ. เอวเมวํ อิเธกจฺโจ สงฺขาเร วิญฺญาณํ รูปํ เวทนํ อตฺตโต สมนุปสฺสติ, ตสฺส เอวํ โหติ “อยํ โข เม อตฺตา, อิมสฺมิํ จ ปน อตฺตนิ อยํ สญฺญา”ติ อตฺตนิ สญฺญํ สมนุปสฺสติ, อภินิเวส\-ปรามาโส ทิฏฺฐิ, ทิฏฺฐิ น วตฺถุ, วตฺถุ น ทิฏฺฐิ. อญฺญา ทิฏฺฐิ, อญฺญํ วตฺถุ. ยา จ ทิฏฺฐิ ยญฺจ วตฺถุ. อยํ ตติยา สญฺญาวตฺถุกา อตฺตานุทิฏฺฐิ. อตฺตานุทิฏฺฐิ มิจฺฉาทิฏฺฐิ ฯเปฯ อิมานิ สญฺโญชนานิ น จ ทิฏฺฐิโย. เอวํ อตฺตนิ สญฺญํ สมนุปสฺสติ. \mbox{(3, 11)}}

\prose{กถํ \textbf{สญฺญาย อตฺตานํ สมนุปสฺสติ?} อิเธกจฺโจ สงฺขาเร วิญฺญาณํ รูปํ เวทนํ อตฺตโต สมนุปสฺสติ, ตสฺส เอวํ โหติ “อยํ โข เม อตฺตา, โส โข ปน เม อยํ อตฺตา อิมาย สญฺญายา”ติ \csromanfolio{143}สญฺญาย อตฺตานํ สมนุปสฺสติ. เสยฺยถาปิ มณิ กรณฺฑเก ปกฺขิตฺโต อสฺส, ตเมนํ ปุริโส เอวํ วเทยฺย “อยํ มณิ, อยํ กรณฺฑโก. อญฺโญ มณิ, อญฺโญ กรณฺฑโก. โส โข ปนายํ มณิ อิมสฺมิํ กรณฺฑเก”ติ กรณฺฑกสฺมิํ, มณิํ สมนุปสฺสติ. เอวเมวํ อิเธกจฺโจ สงฺขาเร วิญฺญาณํ รูปํ เวทนํ อตฺตโต สมนุปสฺสติ, ตสฺส เอวํ โหติ “อยํ โข เม อตฺตา, โส โข ปน เม อยํ อตฺตา อิมาย สญฺญายา”ติ สญฺญาย อตฺตานํ สมนุปสฺสติ, อภินิเวส\-ปรามาโส ทิฏฺฐิ, ทิฏฺฐิ น วตฺถุ, วตฺถุ น ทิฏฺฐิ. อญฺญา ทิฏฺฐิ, อญฺญํ วตฺถุ. ยา จ ทิฏฺฐิ ยญฺจ วตฺถุ. อยํ จตุตฺถา สญฺญาวตฺถุกา อตฺตานุทิฏฺฐิ. อตฺตานุทิฏฺฐิ มิจฺฉาทิฏฺฐิ ฯเปฯ อิมานิ สญฺโญชนานิ น จ ทิฏฺฐิโย. เอวํ สญฺญาย อตฺตานํ สมนุปสฺสติ. \mbox{(4, 12)}}

\proseitem{134}{กถํ สงฺขาเร อตฺตโต สมนุปสฺสติ? อิเธกจฺโจ จกฺขุ\-สมฺผสฺสชํ เจตนํ. โสต\-สมฺผสฺสชํ เจตนํ. ฆาน\-สมฺผสฺสชํ เจตนํ. ชิวฺหา\-สมฺผสฺสชํ เจตนํ. กาย\-สมฺผสฺสชํ เจตนํ. มโน\-สมฺผสฺสชํ เจตนํ อตฺตโต สมนุปสฺสติ “ยา มโน\-สมฺผสฺสชา เจตนา, โส อหํ. โย อหํ, สา มโน\-สมฺผสฺสชา เจตนา”ติ มโน\-สมฺผสฺสชํ เจตนญฺจ อตฺตญฺจ อทฺวยํ สมนุปสฺสติ. เสยฺยถาปิ เตลปฺปทีปสฺส ฌายโต ยา อจฺจิ, โส วณฺโณ. โย วณฺโณ, สา อจฺจีติ อจฺจิญฺจ วณฺณญฺจ อทฺวยํ สมนุปสฺสติ. เอวเมวํ อิเธกจฺโจ มโน\-สมฺผสฺสชํ เจตนํ อตฺตโต สมนุปสฺสติ “ยา มโน\-สมฺผสฺสชา เจตนา, โส อหํ. โย อหํ, สา มโน\-สมฺผสฺสชา เจตนา”ติ มโน\-สมฺผสฺสชํ เจตนญฺจ อตฺตญฺจ อทฺวยํ สมนุปสฺสติ, อภินิเวส\-ปรามาโส ทิฏฺฐิ, ทิฏฺฐิ น วตฺถุ, วตฺถุ น ทิฏฺฐิ. อญฺญา ทิฏฺฐิ, อญฺญํ วตฺถุ. ยา จ ทิฏฺฐิ ยญฺจ วตฺถุ. อยํ ปฐมา สงฺขาร\-วตฺถุกา อตฺตานุทิฏฺฐิ. อตฺตานุทิฏฺฐิ มิจฺฉาทิฏฺฐิ ฯเปฯ อิมานิ สญฺโญชนานิ น จ ทิฏฺฐิโย. เอวํ สงฺขาเร อตฺตโต สมนุปสฺสติ. \mbox{(1, 13)}}

\prose{กถํ \textbf{สงฺขารวนฺตํ อตฺตานํ สมนุปสฺสติ}? อิเธกจฺโจ วิญฺญาณํ รูปํ เวทนํ สญฺญํ อตฺตโต สมนุปสฺสติ, ตสฺส เอวํ โหติ “อยํ โข เม อตฺตา, โส โข ปน เม อยํ อตฺตา อิเมหิ สงฺขาเรหิ สงฺขารวา”ติ สงฺขารวนฺตํ อตฺตานํ สมนุปสฺสติ. เสยฺยถาปิ รุกฺโข ฉายาสมฺปนฺโน อสฺส, ตเมนํ ปุริโส เอวํ วเทยฺย “อยํ รุกฺโข, อยํ ฉายา. อญฺโญ รุกฺโข, อญฺญา ฉายา. โส โข ปนายํ \csromanfolio{144}รุกฺโข อิมาย ฉายาย ฉายาวา”ติ ฉายาวนฺตํ รุกฺขํ สมนุปสฺสติ. เอวเมวํ อิเธกจฺโจ วิญฺญาณํ รูปํ เวทนํ สญฺญํ อตฺตโต สมนุปสฺสติ, ตสฺส เอวํ โหติ “อยํ โข เม อตฺตา, โส โข ปน เม อยํ อตฺตา อิเมหิ สงฺขาเรหิ สงฺขารวา”ติ สงฺขารวนฺตํ อตฺตานํ สมนุปสฺสติ, อภินิเวส\-ปรามาโส ทิฏฺฐิ, ทิฏฺฐิ น วตฺถุ, วตฺถุ น ทิฏฺฐิ. อญฺญา ทิฏฺฐิ, อญฺญํ วตฺถุ. ยา จ ทิฏฺฐิ ยญฺจ วตฺถุ. อยํ ทุติยา สงฺขาร\-วตฺถุกา อตฺตานุทิฏฺฐิ. อตฺตานุทิฏฺฐิ มิจฺฉาทิฏฺฐิ ฯเปฯ อิมานิ สญฺโญชนานิ น จ ทิฏฺฐิโย. เอวํ สงฺขารวนฺตํ อตฺตานํ สมนุปสฺสติ. \mbox{(1, 14)}}

\prose{กถํ อตฺตนิ สงฺขาเร สมนุปสฺสติ? อิเธกจฺโจ วิญฺญาณํ รูปํ เวทนํ สญฺญํ อตฺตโต สมนุปสฺสติ, ตสฺส เอวํ โหติ “อยํ โข เม อตฺตา, อิมสฺมิํ จ ปน อตฺตนิ อิเม สงฺขารา”ติ อตฺตนิ สงฺขาเร สมนุปสฺสติ. เสยฺยถาปิ, ปุปฺผํ คนฺธ\-สมฺปนฺนํ อสฺส, ตเมนํ ปุริโส เอวํ วเทยฺย “อิทํ ปุปฺผํ, อยํ คนฺโธ. อญฺญํ ปุปฺผํ, อญฺโญ คนฺโธ. โส โข ปนายํ คนฺโธ อิมสฺมิํ ปุปฺเผ”ติ ปุปฺผสฺมิํ คนฺธํ สมนุปสฺสติ. เอวเมวํ อิเธกจฺโจ วิญฺญาณํ รูปํ เวทนํ สญฺญํ อตฺตโต สมนุปสฺสติ, ตสฺส เอวํ โหติ “อยํ โข เม อตฺตา, อิมสฺมิํ จ ปน อตฺตนิ อิเม สงฺขารา”ติ อตฺตนิ สงฺขาเร สมนุปสฺสติ, อภินิเวส\-ปรามาโส ทิฏฺฐิ, ทิฏฺฐิ น วตฺถุ. วตฺถุ น ทิฏฺฐิ. อญฺญา ทิฏฺฐิ, อญฺญํ วตฺถุ. ยา จ ทิฏฺฐิ ยญฺจ วตฺถุ. อยํ ตติยา สงฺขาร\-วตฺถุกา อตฺตานุทิฏฺฐิ. อตฺตานุทิฏฺฐิ มิจฺฉาทิฏฺฐิ ฯเปฯ อิมานิ สญฺโญชนานิ น จ ทิฏฺฐิโย. เอวํ อตฺตนิ สงฺขาเร สมนุปสฺสติ. \mbox{(3, 15)}}

\prose{กถํ สขาเรสุ อตฺตานํ สมนุปสฺสติ? อิเธกจฺโจ วิญฺญาณํ รูปํ เวทนํ สญฺญํ อตฺตโต สมนุปสฺสติ, ตสฺส เอวํ โหติ “อยํ โข เม อตฺตา, โส โข ปน เม อยํ อตฺตา อิเมสุ สงฺขาเรสู”ติ สงฺขาเรสุ อตฺตานํ สมนุปสฺสติ. เสยฺยถาปิ มณิ กรณฺฑเก ปกฺขิตฺโต อสฺส, ตเมนํ ปุริโส เอวํ วเทยฺย “อยํ มณิ, อยํ กรณฺฑโก. อญฺโญ มณิ, อญฺโญ กรณฺฑโก. โส โข ปนายํ มณิ อิมสฺมิํ กรณฺฑเก”ติ กรณฺฑกสฺมิํ มณิํ สมนุปสฺสติ. เอวเมวํ อิเธกจฺโจ วิญฺญาณํ รูปํ เวทนํ สญฺญํ อตฺตโต สมนุปสฺสติ, ตสฺส เอวํ โหติ “อยํ โข เม อตฺตา, โส โข ปน เม อยํ อตฺตา อิเมสุ สงฺขาเรสูติ สงฺขาเรสุ อตฺตานํ สมนุปสฺสติ, อภินิเวส\-ปรามาโส ทิฏฺฐิ, ทิฏฺฐิ น วตฺถุ, วตฺถุ น ทิฏฺฐิ. อญฺญา ทิฏฺฐิ, อญฺญํ วตฺถุ. ยา จ ทิฏฺฐิ ยญฺจ \csromanfolio{145}วตฺถุ อยํ จตุตฺถา สงฺขาร\-วตฺถุกา อตฺตานุทิฏฺฐิ. อตฺตานุทิฏฺฐิ มิจฺฉาทิฏฺฐิ ฯเปฯ อิมานิ สญฺโญชนานิ น จ ทิฏฺฐิโย. เอวํ สงฺขาเรสุ อตฺตานํ สมนุปสฺสติ. \mbox{(4, 16)}}

\proseitem{135}{กถํ วิญฺญาณํ อตฺตโต สมนุปสฺสติ? อิเธกจฺโจ จกฺขุวิญฺญาณํ. โสตวิญฺญาณํ. ฆานวิญฺญาณํ. ชิวฺหาวิญฺญาณํ. กายวิญฺญาณํ. มโนวิญฺญาณํ อตฺตโต สมนุปสฺสติ “ยํ มโนวิญฺญาณํ, โส อหํ. โย อหํ, ตํ มโนวิญฺญาณนฺ”ติ มโนวิญฺญาณญฺจ, อตฺตญฺจ อทฺวยํ สมนุปสฺสติ. เสยฺยถาปิ เตลปฺปทีปสฺส ฌายโต ยา อจฺจิ, โส วณฺโณ. โย วณฺโณ, สา อจฺจีติ อจฺจิญฺจ วณฺณญฺจ อทฺวยํ สมนุปสฺสติ. เอวเมวํ อิเธกจฺโจ มโนวิญฺญาณํ อตฺตโต สมนุปสฺสติ “ยํ มโนวิญฺญาณํ, โส อหํ. โย อหํ, ตํ มโนวิญฺญาณนฺ”ติ มโนวิญฺญาณญฺจ อตฺตญฺจ อทฺวยํ สมนุปสฺสติ, อภินิเวส\-ปรามาโส ทิฏฺฐิ, ทิฏฺฐิ น วตฺถุ, วตฺถุ น ทิฏฺฐิ, อญฺญา ทิฏฺฐิ, อญฺญํ วตฺถุ, ยา จ ทิฏฺฐิ, ยญฺจ วตฺถุ. อยํ ปฐมา วิญฺญาณ\-วตฺถุกา อตฺตานุทิฏฺฐิ. อตฺตานุทิฏฺฐิ มิจฺฉาทิฏฺฐิ ฯเปฯ อิมานิ สญฺโญชนานิ น จ ทิฏฺฐิโย. เอวํ วิญฺญาณํ อตฺตโต สมนุปสฺสติ. \mbox{(1, 17)}}

\prose{กถํ วิญฺญาณวนฺตํ อตฺตานํ สมนุปสฺสติ? อิเธกจฺโจ รูปํ เวทนํ สญฺญํ สงฺขาเร อตฺตโต สมนุปสฺสติ, ตสฺส เอวํ โหติ “อยํ โข เม อตฺตา, โส โข ปน เม อยํ อตฺตา อิมินา วิญฺญาเณน วิญฺญาณวา”ติ วิญฺญาณวนฺตํ อตฺตานํ สมนุปสฺสติ. เสยฺยถาปิ รุกฺโข ฉายาสมฺปนฺโน อสฺส, ตเมนํ ปุริโส เอวํ วเทยฺย “อยํ รุกฺโข, อยํ ฉายา. อญฺโญ รุกฺโข, อญฺญา ฉายา. โส โข ปนายํ รุกฺโข อิมาย ฉายาย ฉายาวา”ติ ฉายาวนฺตํ รุกฺขํ สมนุปสฺสติ. เอวเมวํ อิเธกจฺโจ รูปํ เวทนํ สญฺญํ สงฺขาเร อตฺตโต สมนุปสฺสติ, ตสฺส เอวํ โหติ “อยํ โข เม อตฺตา, โส โข ปน เม อยํ อตฺตา อิมินา วิญฺญาเณน วิญฺญาณวา”ติ วิญฺญาณวนฺตํ อตฺตานํ สมนุปสฺสติ, อภินิเวส\-ปรามาโส ทิฏฺฐิ, ทิฏฺฐิ น วตฺถุ, วตฺถุ น ทิฏฺฐิ, อญฺญา ทิฏฺฐิ, อญฺญํ วตฺถุ, ยา จ ทิฏฺฐิ ยญฺจ วตฺถุ. อยํ ทุติยา วิญฺญาณ\-วตฺถุกา อตฺตานุทิฏฺฐิ, อตฺตานุทิฏฺฐิ มิจฺฉาทิฏฺฐิ ฯเปฯ อิมานิ สญฺโญชนานิ น จ ทิฏฺฐิโย. เอวํ วิญฺญาณวนฺตํ อตฺตานํ สมนุปสฺสติ. \mbox{(2, 18)}}

\prose{กถํ \textbf{อตฺตนิ วิญฺญาณํ สมนุปสฺสติ}? อิเธกจฺโจ รูปํ เวทนํ สญฺญํ สงฺขาเร อตฺตโต สมนุปสฺสติ, ตสฺส เอว โหติ “อยํ โข เม \csromanfolio{146}อตฺตา, อิมสฺมิญฺจ ปน อตฺตนิ อิทํ วิญฺญาณนฺ”ติ อตฺตนิ วิญฺญาณํ สมนุปสฺสติ. เสยฺยถาปิ ปุปฺผํ คนฺธ\-สมฺปนฺนํ อสฺส, ตเมนํ ปุริโส เอวํ วเทยฺย “อิทํ ปุปฺผํ, อยํ คนฺโธ. อญฺญํ ปุปฺผํ, อญฺโญ คนฺโธ. โส โข ปนายํ คนฺโธ อิมสฺมิํ ปุปฺเผ”ติ ปุปฺผสฺมิํ คนฺธํ สมนุปสฺสติ. เอวเมวํ อิเธกจฺโจ รูปํ เวทนํ สญฺญํ สงฺขาเร อตฺตโต สมนุปสฺสติ, ตสฺส เอวํ โหติ “อยํ โข เม อตฺตา, อิมสฺมิญฺจ ปน อตฺตนิ อิทํ วิญฺญาณนฺ”ติ อตฺตนิ วิญฺญาณํ สมนุปสฺสติ. อภินิเวส\-ปรามาโส ทิฏฺฐิ, ทิฏฺฐิ น วตฺถุ, วตฺถุ น ทิฏฺฐิ, อญฺญา ทิฏฺฐิ, อญฺญํ วตฺถุ, ยา จ ทิฏฺฐิ ยญฺจ วตฺถุ. อยํ ตติยา วิญฺญาณ\-วตฺถุกา อตฺตานุทิฏฺฐิ. อตฺตานุทิฏฺฐิ มิจฺฉาทิฏฺฐิ ฯเปฯ อิมานิ สญฺโญชนานิ น จ ทิฏฺฐิโย. เอวํ อตฺตนิ วิญฺญาณํ สมนุปสฺสติ. \mbox{(3, 19)}}

\prosewithcloserruled{กถํ วิญฺญาณสฺมิํ อตฺตานํ สมนุปสฺสติ? อิเธกจฺโจ รูปํ เวทนํ สญฺญํ สงฺขาเร อตฺตโต สมนุปสฺสติ, ตสฺส เอวํ โหติ “อยํ โข เม อตฺตา, โส โข ปน เม อยํ อตฺตา อิมสฺมิํ วิญฺญาเณ”ติ วิญฺญาณสฺมิํ อตฺตานํ สมนุปสฺสติ. เสยฺยถาปิ มณิ กรณฺฑเก ปกฺขิตฺโต อสฺส, ตเมนํ ปุริโส เอวํ วเทยฺย “อยํ มณิ, อยํ กรณฺฑโก. อญฺโญ มณิ, อญฺโญ กรณฺฑโก. โส โข ปนายํ มณิ อิมสฺมิํ กรณฺฑเก”ติ กรณฺฑกสฺมิํ มณิํ สมนุปสฺสติ. เอวเมวํ อิเธกจฺโจ รูปํ เวทนํ สญฺญํ สงฺขาเร อตฺตโต สมนุปสฺสติ, ตสฺส เอวํ โหติ “อยํ โข เม อตฺตา, โส โข ปน เม อยํ อตฺตา อิมสฺมิํ วิญฺญาเณ”ติ วิญฺญาณสฺมิํ อตฺตานํ สมนุปสฺสติ, อภินิเวส\-ปรามาโส ทิฏฺฐิ, ทิฏฺฐิ น วตฺถุ, วตฺถุ น ทิฏฺฐิ, อญฺญา ทิฏฺฐิ, อญฺญํ วตฺถุ, ยา จ ทิฏฺฐิ ยญฺจ วตฺถุ. อยํ จตุตฺถา วิญฺญาณ\-วตฺถุกา อตฺตานุทิฏฺฐิ. อตฺตานุทิฏฺฐิ มิจฺฉาทิฏฺฐิ ฯเปฯ อิมานิ สญฺโญชนานิ น จ ทิฏฺฐิโย. เอวํ วิญฺญาณสฺมิํ อตฺตานํ สมนุปสฺสติ. อตฺตานุทิฏฺฐิยา อิเมหิ วีสติยา อากาเรหิ อภินิเวโส โหติ. \mbox{(4, 20)}}{อตฺตานุทิฏฺฐิ\-นิทฺเทโส ทุติโย.}
\csromanmatikaanchor{mtk.54}
\csromantocmarkref{subsection}{3. มิจฺฉาทิฏฺฐินิทฺเทส}{mtk.54}
\bookmark[dest={mtk.54},level=2]{3. มิจฺฉาทิฏฺฐินิทฺเทส}
\csromanheader{2}{3. มิจฺฉาทิฏฺฐิ\-นิทฺเทส}
\proseitem{136}{มิจฺฉาทิฏฺฐิยา กตเมหิ ทสหากาเรหิ อภินิเวโส โหติ? “นตฺถิ ทินฺนนฺ”ติ วตฺถุ\csromansharedfootnote{af398228508721bf}{วตฺถุํ (สฺยา) เอวมุปริปิ.} เอวํวาโท มิจฺฉา\-ภินิเวสปรามาโส\csromansharedfootnote{6be1081a3ece28bf}{มิจฺฉาทิฏฺฐาภินิเวสปรามาโส (สฺยา)} ทิฏฺฐิ, \csromanfolio{147}ทิฏฺฐิ น วตฺถุ, วตฺถุ น ทิฏฺฐิ, อญฺญา ทิฏฺฐิ, อญฺญํ วตฺถุ, ยา จ ทิฏฺฐิ ยญฺจ วตฺถุ. อยํ ปฐมา มิจฺฉาวตฺถุกา มิจฺฉาทิฏฺฐิ. มิจฺฉาทิฏฺฐิ ทิฏฺฐิวิปตฺติ ฯเปฯ อิมานิ สญฺโญชนานิ น จ ทิฏฺฐิโย. “นตฺถิ ยิฏฺฐนฺ”ติ วตฺถุ ฯเปฯ “นตฺถิ หุตนฺ”ติ วตฺถุ. “นตฺถิ สุกต\-ทุกฺกฏานํ กมฺมานํ ผลํ วิปาโก”ติ วตฺถุ. “นตฺถิ อยํ โลโก”ติ วตฺถุ. “นตฺถิ ปโร โลโก”ติ วตฺถุ. “นตฺถิ มาตา”ติ วตฺถุ. “นตฺถิ ปิตา”ติ วตฺถุ. “นตฺถิ สตฺตา โอปปาติกา”ติ วตฺถุ. “นตฺถิ โลเก สมณพฺราหฺมณา สมฺมคฺคตา\csromansharedfootnote{44029fb9de358ac7}{สมคฺคตา (ก)} สมฺมาปฏิปนฺนา เย อิมญฺจ โลกํ ปรญฺจ โลกํ สยํ อภิญฺญา สจฺฉิกตฺวา ปเวเทนฺตี”ติ วตฺถุ เอวํวาโท มิจฺฉา\-ภินิเวสปรามาโส ทิฏฺฐิ, ทิฏฺฐิ น วตฺถุ, วตฺถุ น ทิฏฺฐิ, อญฺญา ทิฏฺฐิ, อญฺญํ วตฺถุ, ยา จ ทิฏฺฐิ, ยญฺจ วตฺถุ, อยํ ทสมา มิจฺฉาวตฺถุกา มิจฺฉาทิฏฺฐิ มิจฺฉาทิฏฺฐิ ทิฏฺฐิวิปตฺติ ฯเปฯ. มิจฺฉา\-ทิฏฺฐิกสฺส ปุริส\-ปุคฺคลสฺส เทฺวว คติโย ฯเปฯ อิมานิ สญฺโญชนานิ น จ ทิฏฺฐิโย. มิจฺฉาทิฏฺฐิยา อิเมหิ ทสหากาเรหิ อภินิเวโส โหติ.}

\csromanmatikaanchor{mtk.55}
\csromantocmarkref{subsection}{4. สกฺกายทิฏฺฐินิทฺเทส}{mtk.55}
\bookmark[dest={mtk.55},level=2]{4. สกฺกายทิฏฺฐินิทฺเทส}
\csromanheader{2}{4. สกฺกายทิฏฺฐิ\-นิทฺเทส}
\proseitem{137}{สกฺกาย\-ทิฏฺฐิยา กตเมหิ วีสติยา อากาเรหิ อภินิเวโส โหติ? อิธ อสฺสุตวา ปุถุชฺชโน อริยานํ อทสฺสาวี อริยธมฺมสฺส อโกวิโท อริยธมฺเม อวินีโต สปฺปุริสานํ อทสฺสาวี สปฺปุริส\-ธมฺมสฺส อโกวิโท สปฺปุริส\-ธมฺเม อวินีโต รูปํ อตฺตโต สมนุปสฺสติ, รูปวนฺตํ วา อตฺตานํ, อตฺตนิ วา รูปํ, รูปสฺมิํ วา อตฺตานํ. เวทนํ. สญฺญํ. สงฺขาเร. วิญฺญาณํ อตฺตโต สมนุปสฺสติ, วิญฺญาณวนฺตํ วา อตฺตานํ, อตฺตนิ วา วิญฺญาณํ วิญฺญาณสฺมิํ วา อตฺตานํ.}

\prose{กถํ รูปํ อตฺตโต สมนุปสฺสติ? อิเธกจฺโจ ปถวีกสิณํ ฯเปฯ. โอทาตกสิณํ อตฺตโต สมนุปสฺสติ “ยํ โอทาตกสิณํ, โส อหํ. โย อหํ, ตํ โอทาตกสิณนฺ”ติ โอทาตกสิณญฺจ อตฺตญฺจ อทฺวยํ สมนุปสฺสติ. เสยฺยถาปิ เตลปฺปทีปสฺส ฌายโต ฯเปฯ. เอวเมวํ อิเธกจฺโจ โอทาตกสิณํ อตฺตโต สมนุปสฺสติ, อภินิเวส\-ปรามาโส ทิฏฺฐิ ฯเปฯ. อยํ ปฐมา รูปวตฺถุกา สกฺกายทิฏฺฐิ, สกฺกายทิฏฺฐิ มิจฺฉาทิฏฺฐิ ฯเปฯ อิมานิ สญฺโญชนานิ น จ ทิฏฺฐิโย. เอวํ รูปํ อตฺตโต สมนุปสฺสติ ฯเปฯ. สกฺกาย\-ทิฏฺฐิยา อิเมหิ วีสติยา อากาเรหิ อภินิเวโส โหติ.}

\csromanmatikaanchor{mtk.56}
\csromantocmarkref{subsection}{5. สสฺสตทิฏฺฐินิทฺเทส}{mtk.56}
\bookmark[dest={mtk.56},level=2]{5. สสฺสตทิฏฺฐินิทฺเทส}
\csromanheader{2}{5. สสฺสตทิฏฺฐิ\-นิทฺเทส}
\proseitem{138}{\csromanfolio{148}สกฺกายวตฺถุกาย สสฺสต\-ทิฏฺฐิยา กตเมหิ ปนฺนรสหิ อากาเรหิ อภินิเวโส โหติ? อิธ อสฺสุตวา ปุถุชฺชโน อริยานํ อทสฺสาวี อริยธมฺมสฺส อโกวิโท อริยธมฺเม อวินีโต สปฺปุริสานํ อทสฺสาวี สปฺปุริส\-ธมฺมสฺส อโกวิโท สปฺปุริส\-ธมฺเม อวินีโต รูปวนฺตํ วา อตฺตานํ, อตฺตนิ วา รูปํ, รูปสฺมิํ วา อตฺตานํ. เวทนาวนฺตํ วา อตฺตานํ. สญฺญาวนฺตํ วา อตฺตานํ. สงฺขารวนฺตํ วา อตฺตานํ. วิญฺญาณวนฺตํ วา อตฺตานํ. อตฺตนิ วา วิญฺญาณํ, วิญฺญาณสฺมิํ วา อตฺตานํ.}

\prose{กถํ รูปวนฺตํ อตฺตานํ สมนุปสฺสติ? อิเธกจฺโจ เวทนํ สญฺญํ สงฺขาเร วิญฺญาณํ อตฺตโต สมนุปสฺสติ, ตสฺส เอวํ โหติ “อยํ โข เม อตฺตา, โส โข ปน เม อยํ อตฺตา อิมินา รูเปน รูปวา”ติ รูปวนฺตํ อตฺตานํ สมนุปสฺสติ. เสยฺยถาปิ รุกฺโข ฉายาสมฺปนฺโน อสฺส, ตเมนํ ปุริโส เอวํ วเทยฺย “อยํ รุกฺโข, อยํ ฉายา. อญฺโญ รุกฺโข, อญฺญา ฉายา. โส โข ปนายํ รุกฺโข อิมาย ฉายาย ฉายาวา”ติ ฉายาวนฺตํ รุกฺขํ สมนุปสฺสติ. เอวเมวํ อิเธกจฺโจ เวทนํ ฯเปฯ. อยํ ปฐมา สกฺกายวตฺถุกา สสฺสตทิฏฺฐิ, สสฺสตทิฏฺฐิ มิจฺฉาทิฏฺฐิ ฯเปฯ อิมานิ สญฺโญชนานิ น จ ทิฏฺฐิโย. เอวํ รูปวนฺตํ อตฺตานํ สมนุปสฺสติ ฯเปฯ. สกฺกายวตฺถุกาย สสฺสต\-ทิฏฺฐิยา อิเมหิ ปนฺนรสหิ อากาเรหิ อภินิเวโส โหติ.}

\csromanmatikaanchor{mtk.57}
\csromantocmarkref{subsection}{6. อุจฺเฉททิฏฺฐินิทฺเทส}{mtk.57}
\bookmark[dest={mtk.57},level=2]{6. อุจฺเฉททิฏฺฐินิทฺเทส}
\csromanheader{2}{6. อุจฺเฉททิฏฺฐิ\-นิทฺเทส}
\proseitem{139}{สกฺกายวตฺถุกาย อุจฺเฉท\-ทิฏฺฐิยา กตเมหิ ปญฺจหิ อากาเรหิ อภินิเวโส โหติ? อิธ อสฺสุตวา ปุถุชฺชโน อริยานํ อทสฺสวี อริยธมฺมสฺส อโกวิโท อริยธมฺเม อวินีโต สปฺปุริสานํ อทสฺสาวี สปฺปุริส\-ธมฺมสฺส อโกวิโท สปฺปุริส\-ธมฺเม อวินีโต รูปํ อตฺตโต สมนุปสฺสติ. เวทนํ อตฺตโต สมนุปสฺสติ. สญฺญํ อตฺตโต สมนุปสฺสติ. สงฺขาเร อตฺตโต สมนุปสฺสติ. วิญฺญาณํ อตฺตโต สมนุปสฺสติ.}

\prose{กถํ \textbf{รูปํ อตฺตโต สมนุปสฺสติ?} อิเธกจฺโจ ปถวีกสิณํ ฯเปฯ. โอทาตกสิณํ อตฺตโต สมนุปสฺสติ “ยํ โอทาตกสิณํ, โส \csromanfolio{149}อหํ. โย อหํ, ตํ โอทาตกสิณนฺ”ติ โอทาตกสิณญฺจ อตฺตญฺจ อทฺวยํ สมนุปสฺสติ. เสยฺยถาปิ เตลปฺปทีปสฺส ฌายโต ฯเปฯ. อยํ ปฐมา สกฺกายวตฺถุกา อุจฺเฉททิฏฺฐิ, อุจฺเฉททิฏฺฐิ มิจฺฉาทิฏฺฐิ ฯเปฯ อิมานิ สญฺโญชนานิ น จ ทิฏฺฐิโย. เอวํ รูปํ อตฺตโต สมนุปสฺสติ ฯเปฯ. สกฺกายวตฺถุกาย อุจฺเฉท\-ทิฏฺฐิยา อิเมหิ ปญฺจหิ อากาเรหิ อภินิเวโส โหติ.}

\csromanmatikaanchor{mtk.58}
\csromantocmarkref{subsection}{7. อนฺตคฺคาหิกาทิฏฺฐินิทฺเทส}{mtk.58}
\bookmark[dest={mtk.58},level=2]{7. อนฺตคฺคาหิกาทิฏฺฐินิทฺเทส}
\csromanheader{2}{7. อนฺตคฺคาหิกา\-ทิฏฺฐินิทฺเทส}
\proseitem{140}{อนฺตคฺคาหิกาย ทิฏฺฐิยา กตเมหิ ปญฺญาสาย อากาเรหิ อภินิเวโส โหติ? “สสฺสโต โลโก”ติ อนฺตคฺคาหิกาย ทิฏฺฐิยา กติหากาเรหิ อภินิเวโส โหติ, “อสสฺสโต โลโก”ติ อนฺตคฺคาหิกาย ทิฏฺฐิยา กติหากาเรหิ อภินิเวโส โหติ, “อนฺตวา โลโก”ติ อนฺตคฺคาหิกาย ทิฏฺฐิยา. “อนนฺตวา โลโก”ติ อนฺตคฺคาหิกาย ทิฏฺฐิยา. “ตํ ชีวํ ตํ สรีรนฺ”ติ อนฺตคฺคาหิกาย ทิฏฺฐิยา. “อญฺญํ ชีวํ อญฺญํ สรีรนฺ”ติ อนฺตคฺคาหิกาย ทิฏฺฐิยา. “โหติ ตถาคโต ปรํ มรณา”ติ ฯเปฯ “น โหติ ตถาคโต ปรํ มรณา”ติ. “โหติ จ น จ โหติ ตถาคโต ปรํ มรณา”ติ. “เนว โหติ น น โหติ ตถาคโต ปรํ มรณา”ติ อนฺตคฺคาหิกาย ทิฏฺฐิยา กติหากาเรหิ อภินิเวโส โหติ?}

\prose{“สสฺสโต โลโก”ติ อนฺตคฺคาหิกาย ทิฏฺฐิยา ปญฺจหากาเรหิ อภินิเวโส โหติ ฯเปฯ “เนว โหติ น น โหติ ตถาคโต ปรํ มรณา”ติ อนฺตคฺคาหิกาย ทิฏฺฐิยา ปญฺจหากาเรหิ อภินิเวโส โหติ.}

\prose{(ก) “สสฺสโต โลโก”ติ อนฺตคฺคาหิกาย ทิฏฺฐิยา กตเมหิ ปญฺจหากาเรหิ อภินิเวโส โหติ? “รูปํ โลโก เจว สสฺสตญฺจา”ติ อภินิเวส\-ปรามาโส ทิฏฺฐิ, ตาย ทิฏฺฐิยา โส อนฺโต คหิโตติ อนฺตคฺคาหิกา ทิฏฺฐิ, ทิฏฺฐิ น วตฺถุ, วตฺถุ น ทิฏฺฐิ, อญฺญา ทิฏฺฐิ, อญฺญํ วตฺถุ, ยา จ ทิฏฺฐิ ยญฺจ วตฺถุ. อยํ ปฐมา “สสฺสโต โลโก”ติ อนฺตคฺคาหิกา ทิฏฺฐิ, อนฺตคฺคาหิกา ทิฏฺฐิ มิจฺฉาทิฏฺฐิ ฯเปฯ อิมานิ สญฺโญชนานิ น จ ทิฏฺฐิโย.}

\prose{“เวทนา โลโก เจว สสฺสตา จา”ติ ฯเปฯ “สญฺญา โลโก เจว อสฺสตา จา”ติ ฯเปฯ “สงฺขารา โลโก เจว สสฺสตา จา”ติ ฯเปฯ \csromanfolio{150}“วิญฺญาณํ โลโก เจว สสฺสตญฺจา”ติ อภินิเวส\-ปรามาโส ทิฏฺฐิ, ตาย ทิฏฺฐิยา โส อนฺโต คหิโตติ อนฺตคฺคาหิกา ทิฏฺฐิ, ทิฏฺฐิ น วตฺถุ, วตฺถุ น ทิฏฺฐิ, อญฺญา ทิฏฺฐิ, อญฺญํ วตฺถุ, ยา จ ทิฏฺฐิ, ยญฺจ วตฺถุ. อยํ ปญฺจมี “สสฺสโต โลโก”ติ อนฺตคฺคาหิกา ทิฏฺฐิ. อนฺตคฺคาหิกา ทิฏฺฐิ มิจฺฉาทิฏฺฐิ ฯเปฯ อิมานิ สญฺโญชนานิ น จ ทิฏฺฐิโย. “สสฺสโต โลโก”ติ อนฺตคฺคาหิกาย ทิฏฺฐิยา อิเมหิ ปญฺจหากาเรหิ อภินิเวโส โหติ. (5)}

\prose{(ข) “อสสฺสโต โลโก”ติ อนฺตคฺคาหิกาย ทิฏฺฐิยา กตเมหิ ปญฺจหากาเรหิ อภินิเวโส โหติ? “รูปํ โลโก เจว อสสฺสตญฺจา”ติ อภินิเวส\-ปรามาโส ทิฏฺฐิ, ตาย ทิฏฺฐิยา โส อนฺโต คหิโตติ อนฺตคฺคาหิกา ทิฏฺฐิ ฯเปฯ. อยํ ปฐมา “อสสฺสโต โลโก”ติ อนฺตคฺคาหิกา ทิฏฺฐิ, อนฺตคฺคาหิกา ทิฏฺฐิ มิจฺฉาทิฏฺฐิ ทิฏฺฐิวิปตฺติ ฯเปฯ อิมานิ สญฺโญชนานิ น จ ทิฏฺฐิโย.}

\prose{“เวทนา โลโก เจว อสสฺสตา จา”ติ ฯเปฯ “สญฺญา โลโก เจว อสสฺสตา จา”ติ ฯเปฯ “สงฺขารา โลโก เจว อสสฺสตา จา”ติ ฯเปฯ “วิญฺญาณํ โลโก เจว อสสฺสตญฺจา”ติ อภินิเวส\-ปรามาโส ทิฏฺฐิ. ตาย ทิฏฺฐิยา โส อนฺโต คหิโตติ อนฺตคฺคาหิกา ทิฏฺฐิ, อนฺตคฺคาหิกา ทิฏฺฐิ มิจฺฉาทิฏฺฐิ ฯเปฯ อิมานิ สญฺโญชนานิ น จ ทิฏฺฐิโย. “อสสฺสโต โลโก”ติ อนฺตคฺคาหิกาย ทิฏฺฐิยา อิเมหิ ปญฺจหากาเรหิ อภินิเวโส โหติ. \mbox{(5, 10)}}

\prose{(ค) “อนฺตวา โลโก”ติ อนฺตคฺคาหิกาย ทิฏฺฐิยา กตเมหิ ปญฺจหากาเรหิ อภินิเวโส โหติ? อิเธกจฺโจ ปริตฺตํ โอกาสํ นีลกโต ผรติ, ตสฺส เอวํ โหติ “อนฺตวา อยํ โลโก ปริวฏุโม”ติ อนฺตสญฺญี โหติ. “ยํ ผรติ, ตํ วตฺถุ เจว โลโก จ. เยน ผรติ, โส อตฺตา เจว โลโก จา”ติ อภินิเวส\-ปรามาโส ทิฏฺฐิ, ตาย ทิฏฺฐิยา โส อนฺโต คหิโตติ อนฺตคฺคาหิกา ทิฏฺฐิ, ทิฏฺฐิ น วตฺถุ, วตฺถุ น ทิฏฺฐิ, อญฺญา ทิฏฺฐิ, อญฺญํ วตฺถุ, ยา จ ทิฏฺฐิ ยญฺจ วตฺถุ. อยํ ปฐมา “อนฺตวา โลโก”ติ อนฺตคฺคาหิกา ทิฏฺฐิ, อนฺตคฺคาหิกา ทิฏฺฐิ มิจฺฉาทิฏฺฐิ ฯเปฯ อิมานิ สญฺโญชนานิ น จ ทิฏฺฐิโย.}

\prose{\csromanfolio{151}อิเธกจฺโจ ปริตฺตํ โอกาสํ ปีตกโต ผรติ. โลหิตกโต ผรติ. โอทาตกโต ผรติ. โอภาสกโต ผรติ, ตสฺส เอวํ โหติ “อนฺตวา อยํ โลโก ปริวฏุโม”ติ อนฺตสญฺญี โหติ. “ยํ ผรติ, ตํ วตฺถุ เจว โลโก จ. เยน ผรติ, โส อตฺตา เจว โลโก จา”ติ อภินิเวส\-ปรามาโส ทิฏฺฐิ, ตาย ทิฏฺฐิยา โส อนฺโต คหิโตติ อนฺตคฺคาหิกา ทิฏฺฐิ ฯเปฯ “อนฺตวา โลโก”ติ อนฺตคฺคาหิกาย ทิฏฺฐิยา อิเมหิ ปญฺจหากาเรหิ อภินิเวโส โหติ. \mbox{(5, 15)}}

\prose{(ฆ) “อนนฺตวา โลโก”ติ อนฺตคฺคาหิกาย ทิฏฺฐิยา กตเมหิ ปญฺจหากาเรหิ อภินิเวโส โหติ\csromandash{}อิเธกจฺโจ วิปุลํ โอกาสํ นีลกโต ผรติ, ตสฺส เอวํ โหติ “อนนฺตวา อยํ โลโก อปริยนฺโต”ติ อนนฺตสญฺญี โหติ. “ยํ ผรติ, ตํ วตฺถุ เจว โลโก จ. เยน ผรติ, โส อตฺตา เจว โลโก จา”ติ อภินิเวส\-ปรามาโส ทิฏฺฐิ, ตาย ทิฏฺฐิยา โส อนฺโต คหิโตติ อนฺตคฺคาหิกา ทิฏฺฐิ, ทิฏฺฐิ น วตฺถุ, วตฺถุ น ทิฏฺฐิ, อญฺญา ทิฏฺฐิ, อญฺญํ วตฺถุ, ยา จ ทิฏฺฐิ ยญฺจ วตฺถุ. อยํ ปฐมา “อนนฺตวา โลโก”ติ อนฺตคฺคาหิกา ทิฏฺฐิ, อนฺตคฺคาหิกา ทิฏฺฐิ มิจฺฉาทิฏฺฐิ ฯเปฯ อิมานิ สญฺโญชนานิ น จ ทิฏฺฐิโย.}

\prose{อิเธกจฺโจ วิปุลํ โอกาสํ ปีตกโต ผรติ. โลหิตกโต ผรติ. โอทาตกโต ผรติ. โอภาสกโต ผรติ, ตสฺส เอวํ โหติ “อนนฺตวา อยํ โลโก อปริยนฺโต”ติ อนนฺตสญฺญี โหติ. “ยํ ผรติ, ตํ วตฺถุ เจว โลโก จ. เยน ผรติ, โส อตฺตา เจว โลโก จา”ติ อภินิเวส\-ปรามาโส ทิฏฺฐิ, ตาย ทิฏฺฐิยา โส อนฺโต คหิโตติ อนฺตคฺคาหิกา ทิฏฺฐิ ฯเปฯ “อนนฺตวา โลโก”ติ อนฺตคฺคาหิกาย ทิฏฺฐิยา อิเมหิ ปญฺจหากาเรหิ อภินิเวโส โหติ. \mbox{(5, 20)}}

\prose{(\^{}อ) “ตํ ชีวํ ตํ สรีรนฺ”ติ อนฺตคฺคาหิกาย ทิฏฺฐิยา กตเมหิ ปญฺจหากาเรหิ อภินิเวโส โหติ? “รูปํ ชีวญฺเจว สรีรญฺจ, ยํ ชีวํ, ตํ สรีรํ. ยํ สรีรํ, ตํ ชีวนฺ”ติ อภินิเวส\-ปรามาโส ทิฏฺฐิ, ตาย ทิฏฺฐิยา โส อนฺโต คหิโตติ อนฺตคฺคาหิกา ทิฏฺฐิ, ทิฏฺฐิ น วตฺถุ, วตฺถุ น ทิฏฺฐิ, อญฺญา ทิฏฺฐิ, อญฺญํ วตฺถุ, ยา จ ทิฏฺฐิ ยญฺจ วตฺถุ. อยํ ปฐมา “ตํ ชีวํ ตํ สรีรนฺ”ติ อนฺตคฺคาหิกา ทิฏฺฐิ, อนฺตคฺคาหิกา ทิฏฺฐิ มิจฺฉาทิฏฺฐิ ฯเปฯ อิมานิ สญฺโญชนานิ น จ ทิฏฺฐิโย.}

\prose{\csromanfolio{152}“เวทนา ชีวา เจว สรีรญฺจ. สญฺญา ชีวา เจว สรีรญฺจ. สงฺขารา ชีวา เจว สรีรญฺจ. วิญฺญาณํ ชีวํ เจว สรีรญฺจ. ยํ ชีวํ ตํ สรีรํ, ยํ สรีรํ ตํ ชีวนฺ”ติ อภินิเวส\-ปรามาโส ทิฏฺฐิ, ตาย ทิฏฺฐิยา โส อนฺโต คหิโตติ อนฺตคฺคาหิกา ทิฏฺฐิ ฯเปฯ “ตํ ชีวํ ตํ สรีรนฺ”ติ อนฺตคฺคาหิกาย ทิฏฺฐิยา อิเมหิ ปญฺจหากาเรหิ อภินิเวโส โหติ. \mbox{(5, 25)}}

\prose{(จ) “อญฺญํ ชีวํ อญฺญํ สรีรนฺ”ติ อนฺตคฺคาหิกาย ทิฏฺฐิยา กตเมหิ ปญฺจหากาเรหิ อภินิเวโส โหติ? “รูปํ สรีรํ น ชีวํ, ชีวํ น สรีรํ. อญฺญํ ชีวํ อญฺญํ สรีรนฺ”ติ อภินิเวส\-ปรามาโส ทิฏฺฐิ, ตาย ทิฏฺฐิยา โส อนฺโต คหิโตติ อนฺตคฺคาหิกา ทิฏฺฐิ, ทิฏฺฐิ น วตฺถุ, วตฺถุ น ทิฏฺฐิ, อญฺญา ทิฏฺฐิ, อญฺญํ วตฺถุ, ยา จ ทิฏฺฐิ ยญฺจ วตฺถุ. อยํ ปฐมา “อญฺญํ ชีวํ อญฺญํ สรีรนฺ”ติ อนฺตคฺคาหิกา ทิฏฺฐิ, อนฺตคฺคาหิกา ทิฏฺฐิ มิจฺฉาทิฏฺฐิ ฯเปฯ อิมานิ สญฺโญชนานิ น จ ทิฏฺฐิโย.}

\prose{“เวทนา สรีรํ น ชีวํ. สญฺญา สรีรํ น ชีวํ. สงฺขารา สรีรํ น ชีวํ. วิญฺญาณํ สรีรํ น ชีวํ. ชีวํ น สรีรํ. อญฺญํ ชีวํ, อญฺญํ สรีรนฺ”ติ อภินิเวส\-ปรามาโส ทิฏฺฐิ, ตาย ทิฏฺฐิยา โส อนฺโต คหิโตติ อนฺตคฺคาหิกา ทิฏฺฐิ ฯเปฯ “อญฺญํ ชีวํ อญฺญํ สรีรนฺ”ติ อนฺตคฺคาหิกาย ทิฏฺฐิยา อิเมหิ ปญฺจหากาเรหิ อภินิเวโส โหติ. \mbox{(5, 30)}}

\prose{(ฉ) “โหติ ตถาคโต ปรํ มรณา”ติ อนฺตคฺคาหิกาย ทิฏฺฐิยา กตเมหิ ปญฺจหากาเรหิ อภินิเวโส โหติ? “รูปํ อิเธว มรณธมฺมํ, ตถาคโต กายสฺส เภทา โหติปิ ติฏฺฐติปิ อุปฺปชฺชติปิ นิพฺพตฺตติปี”ติ อภินิเวส\-ปรามาโส ทิฏฺฐิ, ตาย ทิฏฺฐิยา โส อนฺโต คหิโตติ อนฺตคฺคาหิกา ทิฏฺฐิ, ทิฏฺฐิ น วตฺถุ, วตฺถุ น ทิฏฺฐิ, อญฺญา ทิฏฺฐิ, อญฺญํ วตฺถุ, ยา จ ทิฏฺฐิ ยญฺจ วตฺถุ. อยํ ปฐมา “โหติ ตถาคโต ปรํ มรณา”ติ อนฺตคฺคาหิกา ทิฏฺฐิ, อนฺตคฺคาหิกา ทิฏฺฐิ มิจฺฉาทิฏฺฐิ ฯเปฯ อิมานิ สญฺโญชนานิ น จ ทิฏฺฐิโย.}

\prose{“เวทนา อิเธว มรณธมฺมา. สญฺญา อิเธว มรณธมฺมา. สงฺขารา อิเธว มรณธมฺมา. วิญฺญาณํ อิเธว มรณธมฺมํ, ตถาคโต กายสฺส เภทา โหติปิ ติฏฺฐติปิ อุปฺปชฺชติปิ นิพฺพตฺตติปี”ติ อภินิเวส\-ปรามาโส ทิฏฺฐิ, ตาย ทิฏฺฐิยา โส อนฺโต คหิโตติ อนฺตคฺคาหิกา ทิฏฺฐิ ฯเปฯ “โหติ \csromanfolio{153}ตถาคโต ปรํ มรณา”ติ อนฺตคฺคาหิกาย ทิฏฺฐิยา อิเมหิ ปญฺจหากาเรหิ อภินิเวโส โหติ. \mbox{(5, 35)}}

\proseitem{140}{(ช) “น โหติ ตถาคโต ปรํ มรณา”ติ อนฺตคฺคาหิกาย ทิฏฺฐิยา กตเมหิ ปญฺจหากาเรหิ อภินิเวโส โหติ? “รูปํ อิเธว มรณธมฺมํ, ตถาคโตปิ กายสฺส เภทา อุจฺฉิชฺชติ วินสฺสติ น โหติ ตถาคโต ปรํ มรณา”ติ อภินิเวส\-ปรามาโส ทิฏฺฐิ, ตาย ทิฏฺฐิยา โส อนฺโต คหิโตติ อนฺตคฺคาหิกา ทิฏฺฐิ, ทิฏฺฐิ น วตฺถุ, วตฺถุ น ทิฏฺฐิ, อญฺญา ทิฏฺฐิ, อญฺญํ วตฺถุ. ยา จ ทิฏฺฐิ ยญฺจ วตฺถุ. อยํ ปฐมา “น โหติ ตถาคโต ปรํ มรณา”ติ อนฺตคฺคาหิกา ทิฏฺฐิ, อนฺตคฺคาหิกา ทิฏฺฐิ มิจฺฉาทิฏฺฐิ ฯเปฯ อิมานิ สญฺโญชนานิ น จ ทิฏฺฐิโย.}

\prose{“เวทนา อิเธว มรณธมฺมา. สญฺญา อิเธว มรณธมฺมา. สงฺขารา อิเธว มรณธมฺมา. วิญฺญาณํ อิเธว มรณธมฺมํ, ตถาคโตปิ กายสฺส เภทา อุจฺฉิชฺชติ วินสฺสติ น โหติ ตถาคโต ปรํ มรณา”ติ อภินิเวส\-ปรามาโส ทิฏฺฐิ. ตาย ทิฏฺฐิยา โส อนฺโต คหิโตติ อนฺตคฺคาหิกา ทิฏฺฐิ ฯเปฯ “น โหติ ตถาคโต ปรํ มรณา”ติ อนฺตคฺคาหิกาย ทิฏฺฐิยา อิเมหิ ปญฺจหากาเรหิ อภินิเวโส โหติ. \mbox{(5, 40)}}

\prose{(ฌ) “โหติ จ น จ โหติ ตถาคโต ปรํ มรณา”ติ อนฺตคฺคาหิกาย ทิฏฺฐิยา กตเมหิ ปญฺจหากาเรหิ อภินิเวโส โหติ? “รูปํ อิเธว มรณธมฺมํ, ตถาคโต กายสฺส เภทา โหติ จ น จ โหตี”ติ อภินิเวส\-ปรามาโส ทิฏฺฐิ, ตาย ทิฏฺฐิยา โส อนฺโต คหิโตติ อนฺตคฺคาหิกา ทิฏฺฐิ, ทิฏฺฐิ น วตฺถุ, วตฺถุ น ทิฏฺฐิ, อญฺญา ทิฏฺฐิ, อญฺญํ วตฺถุ, ยา จ ทิฏฺฐิ ยญฺจ วตฺถุ. อยํ ปฐมา “โหติ จ น จ โหติ ตถาคโต ปรํ มรณา”ติ อนฺตคฺคาหิกา ทิฏฺฐิ. อนฺตคฺคาหิกา ทิฏฺฐิ มิจฺฉาทิฏฺฐิ ฯเปฯ อิมานิ สญฺโญชนานิ น จ ทิฏฺฐิโย.}

\prose{“เวทนา อิเธว มรณธมฺมา. สญฺญา อิเธว มรณธมฺมา. สงฺขารา อิเธว มรณธมฺมา. วิญฺญาณํ อิเธว มรณธมฺมํ, ตถาคโต กายสฺส เภทา โหติ จ น จ โหตี”ติ อภินิเวส\-ปรามาโส ทิฏฺฐิ, ตาย ทิฏฺฐิยา โส อนฺโต คหิโตติ อนฺตคฺคาหิกา ทิฏฺฐิ ฯเปฯ “โหติ จ น จ โหติ ตถาคโต ปรํ มรณา”ติ อนฺตคฺคาหิกาย ทิฏฺฐิยา อิเมหิ ปญฺจหากาเรหิ อภินิเวโส โหติ. \mbox{(5, 45)}}

\proseitem{2}{\csromanfolio{154}(ญ) “เนว โหติ น น โหติ ตถาคโต ปรํ มรณา”ติ อนฺตคฺคาหิกาย ทิฏฺฐิยา กตเมหิ ปญฺจหากาเรหิ อภินิเวโส โหติ? “รูปํ อิเธว มรณธมฺมํ, ตถาคโต กายสฺส เภทา ปรํ มรณา เนว โหติ น น โหตี”ติ อภินิเวส\-ปรามาโส ทิฏฺฐิ, ตาย ทิฏฺฐิยา โส อนฺโต คหิโตติ อนฺตคฺคาหิกา ทิฏฺฐิ, ทิฏฺฐิ น วตฺถุ, วตฺถุ น ทิฏฺฐิ, อญฺญา ทิฏฺฐิ, อญฺญํ วตฺถุ, ยา จ ทิฏฺฐิ, ยญฺจ วตฺถุ. อยํ ปฐมา “เนว โหติ น น โหติ ตถาคโต ปรํ มรณา”ติ อนฺตคฺคาหิกา ทิฏฺฐิ, อนฺตคฺคาหิกา ทิฏฺฐิ มิจฺฉาทิฏฺฐิ ฯเปฯ อิมานิ สญฺโญชนานิ น จ ทิฏฺฐิโย.}

\prosewithcloserruled{“เวทนา อิเธว มรณธมฺมา. สญฺญา อิเธว มรณธมฺมา. สงฺขารา อิเธว มรณธมฺมา. วิญฺญาณํ อิเธว มรณธมฺมํ, ตถาคโต กายสฺส เภทา ปรํ มรณา เนว โหติ น น โหตี”ติ อภินิเวส\-ปรามาโส ทิฏฺฐิ, ตาย ทิฏฺฐิยา โส อนฺโต คหิโตติ อนฺตคฺคาหิกา ทิฏฺฐิ, ทิฏฺฐิ น วตฺถุ, วตฺถุ น ทิฏฺฐิ, อญฺญา ทิฏฺฐิ, อญฺญํ วตฺถุ, ยา จ ทิฏฺฐิ ยญฺจ วตฺถุ. อยํ ปญฺจมี “เนว โหติ น น โหติ ตถาคโต ปรํ มรณา”ติ อนฺตคฺคาหิกา ทิฏฺฐิ, อนฺตคฺคาหิกา ทิฏฺฐิ มิจฺฉาทิฏฺฐิ ฯเปฯ อิมานิ สญฺโญชนานิ น จ ทิฏฺฐิโย. “เนว โหติ น น โหติ ตถาคโต ปรํ มรณา”ติ อนฺตคฺคาหิกาย ทิฏฺฐิยา อิเมหิ ปญฺจหากาเรหิ อภินิเวโส โหติ. อนฺตคฺคาหิกาย ทิฏฺฐิยา อิเมหิ ปญฺญาสาย อากาเรหิ อภินิเวโส โหติ. (50)}{อนฺตคฺคาหิกา\-ทิฏฺฐินิทฺเทโส สตฺตโม.}
\csromanmatikaanchor{mtk.59}
\csromantocmarkref{subsection}{8. ปุพฺพนฺตานุทิฏฺฐินิทฺเทส}{mtk.59}
\bookmark[dest={mtk.59},level=2]{8. ปุพฺพนฺตานุทิฏฺฐินิทฺเทส}
\csromanheader{2}{8. ปุพฺพนฺตา\-นุทิฏฺฐินิทฺเทส}
\proseitem{141}{ปุพฺพนฺตา\-นุทิฏฺฐิยา กตเมหิ อฏฺฐารสหิ อากาเรหิ อภินิเวโส โหติ? จตฺตาโร สสฺสตวาทา, จตฺตาโร เอกจฺจ\-สสฺสติกา, จตฺตาโร อนฺตานนฺติกา, จตฺตาโร อมรา\-วิกฺเขปิกา, เทฺว อธิจฺจ\-สมุปฺปนฺนิกา. ปุพฺพนฺตา\-นุทิฏฺฐิยา อิเมหิ อฏฺฐารสหิ อากาเรหิ อภินิเวโส โหติ.}

\csromanmatikaanchor{mtk.60}
\csromantocmarkref{subsection}{9. อปรนฺตานุทิฏฺฐินิทฺเทส}{mtk.60}
\bookmark[dest={mtk.60},level=2]{9. อปรนฺตานุทิฏฺฐินิทฺเทส}
\csromanheader{2}{9. อปรนฺตา\-นุทิฏฺฐินิทฺเทส}
\proseitem{142}{อปรนฺตา\-นุทิฏฺฐิยา กตเมหิ จตุจตฺตาลีสาย อากาเรหิ อภินิเวโส โหติ? โสฬส สญฺญีวาทา, อฏฺฐ อสญฺญีวาทา, อฏฺฐ \csromanfolio{155}เนวสญฺญี\-นาสญฺญี\-วาทา, สตฺต อุจฺเฉทวาทา, ปญฺจ ทิฏฺฐธมฺมนิพฺพาน\-วาทา. อปรนฺตา\-นุทิฏฺฐิยา อิเมหิ จตุจตฺตาลีสาย อากาเรหิ อภินิเวโส โหติ.}

\csromanmatikaanchor{mtk.61}
\csromantocmarkref{subsection}{10-12. สญฺโญชนิกาทิทิฏฺฐินิทฺเทส}{mtk.61}
\bookmark[dest={mtk.61},level=2]{10-12. สญฺโญชนิกาทิทิฏฺฐินิทฺเทส}
\csromanheader{2}{10-12. สญฺโญชนิกาทิทิฏฺฐินิทฺเทส}
\proseitem{143}{สญฺโญชนิกาย ทิฏฺฐิยา กตเมหิ อฏฺฐารสหิ อากาเรหิ อภินิเวโส โหติ? ยา ทิฏฺฐิ ทิฏฺฐิคตํ ทิฏฺฐิคหนํ ฯเปฯ ทิฏฺฐา\-ภินิเวโส ทิฏฺฐิปรามาโส. สญฺโญชนิกาย ทิฏฺฐิยา อิเมหิ อฏฺฐารสหิ อากาเรหิ อภินิเวโส โหติ.}

\proseitem{144}{“อหนฺ”ติ มาน\-วินิพนฺธาย ทิฏฺฐิยา กตเมหิ อฏฺฐารสหิ อากาเรหิ อภินิเวโส โหติ? “จกฺขุ อหนฺ”ติ อภินิเวส\-ปรามาโส “อหนฺ”ติ มานวินิพนฺธา ทิฏฺฐิ, ทิฏฺฐิ น วตฺถุ, วตฺถุ น ทิฏฺฐิ, อญฺญา ทิฏฺฐิ, อญฺญํ วตฺถุ, ยา จ ทิฏฺฐิ, ยญฺจ วตฺถุ, อยํ ปฐมา “อหนฺ”ติ มานวินิพนฺธา ทิฏฺฐิ. มานวินิพนฺธา ทิฏฺฐิ มิจฺฉาทิฏฺฐิ ฯเปฯ อิมานิ สญฺโญชนานิ น จ ทิฏฺฐิโย.}

\prose{“โสตํ อหนฺ”ติ ฯเปฯ “ฆานํ อหนฺ”ติ ฯเปฯ “ชิวฺหา อหนฺ”ติ ฯเปฯ “กาโย อหนฺ”ติ ฯเปฯ “มโน อหนฺ”ติ ฯเปฯ “รูปา อหนฺ”ติ ฯเปฯ “ธมฺมา อหนฺ”ติ. “จกฺขุวิญฺญาณํ อหนฺ”ติ ฯเปฯ “มโนวิญฺญาณํ อหนฺ”ติ อภินิเวส\-ปรามาโส “อหนฺ”ติ มานวินิพนฺธา ทิฏฺฐิ, ทิฏฺฐิ น วตฺถุ, วตฺถุ น ทิฏฺฐิ, อญฺญา ทิฏฺฐิ, อญฺญํ วตฺถุ, ยา จ ทิฏฺฐิ ยญฺจ วตฺถุ, อยํ อฏฺฐารสมี “อหนฺ”ติ มานวินิพนฺธา ทิฏฺฐิ, มานวินิพนฺธา ทิฏฺฐิ มิจฺฉาทิฏฺฐิ ฯเปฯ อิมานิ สญฺโญชนานิ น จ ทิฏฺฐิโย. “อหนฺ”ติ มาน\-วินิพนฺธาย ทิฏฺฐิยา อิเมหิ อฏฺฐารสหิ อากาเรหิ อภินิเวโส โหติ.}

\proseitem{145}{“มมนฺ”ติ มาน\-วินิพนฺธาย ทิฏฺฐิยา กตเมหิ อฏฺฐารสหิ อากาเรหิ อภินิเวโส โหติ? “จกฺขุ มมนฺ”ติ อภินิเวส\-ปรามาโส “มมนฺ”ติ มานวินิพนฺธา ทิฏฺฐิ, ทิฏฺฐิ น วตฺถุ, วตฺถุ น ทิฏฺฐิ, อญฺญา ทิฏฺฐิ, อญฺญํ วตฺถุ, ยา จ ทิฏฺฐิ ยญฺจ วตฺถุ, อยํ ปฐมา “มมนฺ”ติ มานวินิพนฺธา ทิฏฺฐิ. มานวินิพนฺธา ทิฏฺฐิ มิจฺฉาทิฏฺฐิ ฯเปฯ อิมานิ สญฺโญชนานิ น จ ทิฏฺฐิโย.}

\prose{“โสตํ มมนฺ”ติ ฯเปฯ “ฆานํ มมนฺ”ติ ฯเปฯ “ชิวฺหา มมนฺ”ติ ฯเปฯ “กาโย มมนฺ”ติ ฯเปฯ “มโน มมนฺ”ติ ฯเปฯ “รูปา มมนฺ”ติ ฯเปฯ “ธมฺมา มมนฺ”ติ ฯเปฯ “จกฺขุวิญฺญาณํ มมนฺ”ติ ฯเปฯ “มโนวิญฺญาณํ มมนฺ”ติ อภินิเวส\-ปรามาโส \csromanfolio{156}“มมนฺ”ติ มานวินิพนฺธา ทิฏฺฐิ, ทิฏฺฐิ น วตฺถุ, วตฺถุ น ทิฏฺฐิ, อญฺญา ทิฏฺฐิ, อญฺญํ วตฺถุ, ยา จ ทิฏฺฐิ, ยญฺจ วตฺถุ, อยํ อฏฺฐารสมี “มมนฺ”ติ มานวินิพนฺธา ทิฏฺฐิ, มานวินิพนฺธา ทิฏฺฐิ มิจฺฉาทิฏฺฐิ ฯเปฯ อิมานิ สญฺโญชนานิ น จ ทิฏฺฐิโย. “มมนฺ”ติ มาน\-วินิพนฺธาย ทิฏฺฐิยา อิเมหิ อฏฺฐารสหิ อากาเรหิ อภินิเวโส โหติ.}

\csromanmatikaanchor{mtk.62}
\csromantocmarkref{subsection}{13. อตฺตวาทปฏิสํยุตฺตทิฏฺฐินิทฺเทส}{mtk.62}
\bookmark[dest={mtk.62},level=2]{13. อตฺตวาทปฏิสํยุตฺตทิฏฺฐินิทฺเทส}
\csromanheader{2}{13. อตฺตวาทปฏิสํยุตฺต\-ทิฏฺฐิ\-นิทฺเทส}
\proseitem{146}{อตฺตวาท\-ปฏิสํยุตฺตาย ทิฏฺฐิยา กตเมหิ วีสติยา อากาเรหิ อภินิเวโส โหติ? อิธ อสฺสุตวา ปุถุชฺชโน อริยานํ อทสฺสาวี อริยธมฺมสฺส อโกวิโท อริยธมฺเม อวินีโต สปฺปุริสานํ อทสฺสาวี สปฺปุริส\-ธมฺมสฺส อโกวิโท สปฺปุริส\-ธมฺเม อวินีโต รูปํ อตฺตโต สมนุปสฺสติ, รูปวนฺตํ วา อตฺตานํ, อตฺตนิ วา รูปํ, รูปสฺมิํ วา อตฺตานํ ฯเปฯ เวทนํ. สญฺญํ. สงฺขาเร. วิญฺญาณํ อตฺตโต สมนุปสฺสติ, วิญฺญาณวนฺตํ วา อตฺตานํ, อตฺตนิ วา วิญฺญาณํ, วิญฺญาณสฺมิํ วา อตฺตานํ ฯเปฯ.}

\prose{กถํ รูปํ อตฺตโต สมนุปสฺสติ? อิเธกจฺโจ ปถวีกสิณํ ฯเปฯ. โอทาตกสิณํ อตฺตโต สมนุปสฺสติ. “ยํ โอทาตกสิณํ, โส อหํ. โย อหํ, ตํ โอทาตกสิณนฺ”ติ โอทาตกสิณญฺจ อตฺตญฺจ อทฺวยํ สมนุปสฺสติ. เสยฺยถาปิ เตลปฺปทีปสฺส ฌายโต ยา อจฺจิ, โส วณฺโณ. โย วณฺโณ, สา อจฺจีติ อจฺจิญฺจ วณฺณญฺจ อทฺวยํ สมนุปสฺสติ. เอวเมวํ อิเธกจฺโจ โอทาตกสิณํ อตฺตโต สมนุปสฺสติ ฯเปฯ อยํ ปฐมา รูปวตฺถุกา อตฺตวาท\-ปฏิสํยุตฺตา ทิฏฺฐิ. อตฺตวาท\-ปฏิสํยุตฺตา ทิฏฺฐิ มิจฺฉาทิฏฺฐิ ฯเปฯ อิมานิ สญฺโญชนานิ น จ ทิฏฺฐิโย, เอวํ รูปํ อตฺตโต สมนุปสฺสติ ฯเปฯ. อตฺตวาท\-ปฏิสํยุตฺตาย ทิฏฺฐิยา อิเมหิ วีสติยา อากาเรหิ อภินิเวโส โหติ.}

\csromanmatikaanchor{mtk.63}
\csromantocmarkref{subsection}{14. โลกวาทปฏิสํยุตฺตทิฏฺฐินิทฺเทส}{mtk.63}
\bookmark[dest={mtk.63},level=2]{14. โลกวาทปฏิสํยุตฺตทิฏฺฐินิทฺเทส}
\csromanheader{2}{14. โลกวาทปฏิสํยุตฺตทิฏฺฐิ\-นิทฺเทส}
\proseitem{147}{โลกวาท\-ปฏิสํยุตฺตาย ทิฏฺฐิยา กตเมหิ อฏฺฐหิ อากาเรหิ อภินิเวโส โหติ? “สสฺสโต อตฺตา จ โลโก จา”ติ อภินิเวส\-ปรามาโส โลกวาท\-ปฏิสํยุตฺตา ทิฏฺฐิ, ทิฏฺฐิ น วตฺถุ, วตฺถุ น ทิฏฺฐิ, อญฺญา ทิฏฺฐิ, อญฺญํ วตฺถุ, ยา จ ทิฏฺฐิ ยญฺจ วตฺถุ. อยํ ปฐมา โลกวาท\-ปฏิสํยุตฺตา ทิฏฺฐิ. โลกวาท\-ปฏิสํยุตฺตา ทิฏฺฐิ มิจฺฉาทิฏฺฐิ ฯเปฯ อิมานิ สญฺโญชนานิ น จ ทิฏฺฐิโย.}

\csromanmatikaanchor{mtk.64}
\csromantocmarkref{subsection}{15-16. ภววิภวทิฏฺฐินิทฺเทส}{mtk.64}
\bookmark[dest={mtk.64},level=2]{15-16. ภววิภวทิฏฺฐินิทฺเทส}
\prose{\csromanfolio{157}“อสสฺสโต อตฺตา จ โลโก จา”ติ ฯเปฯ “สสฺสโต จ อสสฺสโต จ อตฺตา จ โลโก จา”ติ ฯเปฯ “เนว สสฺสโต นาสสฺสโต อตฺตา จ โลโก จา”ติ. “อนฺตวา อตฺตา จ โลโก จา”ติ. “อนนฺตวา อตฺตา จ โลโก จา”ติ. “อนฺตวา จ อนนฺตวา จ อตฺตา จ โลโก จา”ติ. “เนว อนฺตวา น อนนฺตวา อตฺตา จ โลโก จา”ติ อภินิเวส\-ปรามาโส โลกวาท\-ปฏิสํยุตฺตา ทิฏฺฐิ. ทิฏฺฐิ น วตฺถุ, วตฺถุ น ทิฏฺฐิ, อญฺญา ทิฏฺฐิ, อญฺญํ วตฺถุ, ยา จ ทิฏฺฐิ ยญฺจ วตฺถุ. อยํ อฏฺฐมี โลกวาท\-ปฏิสํยุตฺตา ทิฏฺฐิ. โลกวาท\-ปฏิสํยุตฺตา ทิฏฺฐิ มิจฺฉาทิฏฺฐิ ฯเปฯ อิมานิ สญฺโญชนานิ น จ ทิฏฺฐิโย. โลกวาท\-ปฏิสํยุตฺตาย ทิฏฺฐิยา อิเมหิ อฏฺฐหิ อากาเรหิ อภินิเวโส โหติ.}

\prose{15-16. ภววิภวทิฏฺฐิ\-นิทฺเทส}

\proseitem{148}{โอลียนาภิ\-นิเวโส ภวทิฏฺฐิ, อติธาวนา\-ภินิเวโส วิภวทิฏฺฐิ, อสฺสาททิฏฺฐิยา ปญฺจติํสาย อากาเรหิ อภินิเวโส กติ ภวทิฏฺฐิโย, กติ วิภว\-ทิฏฺฐิโย. อตฺตานุทิฏฺฐิยา วีสติยา อากาเรหิ อภินิเวโส กติ ภวทิฏฺฐิโย, กติ วิภว\-ทิฏฺฐิโย ฯเปฯ. โลกวาท\-ปฏิสํยุตฺตาย ทิฏฺฐิยา อฏฺฐหิ อากาเรหิ อภินิเวโส กติ ภวทิฏฺฐิโย, กติ วิภว\-ทิฏฺฐิโย?}

\prose{อสฺสาททิฏฺฐิยา ปญฺจาติํสาย อากาเรหิ อภินิเวโส สิยา ภวทิฏฺฐิโย, สิยา วิภว\-ทิฏฺฐิโย. อตฺตานุทิฏฺฐิยา วีสติยา อากาเรหิ อภินิเวโส ปนฺนรสภวทิฏฺฐิโย, ปญฺจ วิภว\-ทิฏฺฐิโย. มิจฺฉาทิฏฺฐิยา ทสหิ อากาเรหิ อภินิเวโส, สพฺพาว ตา วิภว\-ทิฏฺฐิโย. สกฺกาย\-ทิฏฺฐิยา วีสติยา อากาเรหิ อภินิเวโส ปนฺนรส ภวทิฏฺฐิโย, ปญฺจ วิภว\-ทิฏฺฐิโย. สกฺกายวตฺถุกาย สสฺสต\-ทิฏฺฐิยา ปนฺนรสหิ อากาเรหิ อภินิเวโส, สพฺพาว ตา ภวทิฏฺฐิโย. สกฺกายวตฺถุกาย อุจฺเฉท\-ทิฏฺฐิยา ปญฺจหากาเรหิ อภินิเวโส, สพฺพาว ตา วิภว\-ทิฏฺฐิโย.}

\prose{“สสฺสโต โลโก”ติ อนฺตคฺคาหิกาย ทิฏฺฐิยา ปญฺจหากาเรหิ อภินิเวโส, สพฺพาว ตา ภวทิฏฺฐิโย, “อสสฺสโต โลโก”ติ อนฺตคฺคาหิกาย ทิฏฺฐิยา ปญฺจหากาเรหิ อภินิเวโส, สพฺพาว ตา \csromanfolio{158}วิภว\-ทิฏฺฐิโย. “อนฺตวา โลโก”ติ อนฺตคฺคาหิกาย ทิฏฺฐิยา ปญฺจหากาเรหิ อภินิเวโส สิยา ภวทิฏฺฐิโย, สิยา วิภว\-ทิฏฺฐิโย. “อนนฺตวา โลโก”ติ อนฺตคฺคาหิกาย ทิฏฺฐิยา ปญฺจหากาเรหิ อภินิเวโส สิยา ภวทิฏฺฐิโย, สิยา วิภว\-ทิฏฺฐิโย. “ตํ ชีวํ ตํ สรีรนฺ”ติ อนฺตคฺคาหิกาย ทิฏฺฐิยา ปญฺจหากาเรหิ อภินิเวโส, สพฺพาว ตา วิภว\-ทิฏฺฐิโย. “อญฺญํ ชีวํ อญฺญํ สรีรนฺ”ติ อนฺตคฺคาหิกาย ทิฏฺฐิยา ปญฺจหากาเรหิ อภินิเวโส, สพฺพาว ตา ภวทิฏฺฐิโย. “โหติ ตถาคโต ปรํ มรณา”ติ อนฺตคฺคาหิกาย ทิฏฺฐิยา ปญฺจหากาเรหิ อภินิเวโส, สพฺพาว ตา ภวทิฏฺฐิโย. “น โหติ ตถาคโต ปรํ มรณา”ติ อนฺตคฺคาหิกาย ทิฏฺฐิยา ปญฺจหากาเรหิ อภินิเวโส, สพฺพาว ตา วิภว\-ทิฏฺฐิโย. “โหติ จ น จ โหติ ตถาคโต ปรํ มรณา”ติ อนฺตคฺคาหิกาย ทิฏฺฐิยา ปญฺจหากาเรหิ อภินิเวโส สิยา ภวทิฏฺฐิโย, สิยา วิภว\-ทิฏฺฐิโย. “เนว โหติ น น โหติ ตถาคโต ปรํ มรณา”ติ อนฺตคฺคาหิกาย ทิฏฺฐิยา ปญฺจหากาเรหิ อภินิเวโส สิยา ภวทิฏฺฐิโย, สิยา วิภว\-ทิฏฺฐิโย.}

\prose{ปุพฺพนฺตา\-นุทิฏฺฐิยา อฏฺฐารสหิ อากาเรหิ อภินิเวโส สิยา ภวทิฏฺฐิโย, สิยา วิภว\-ทิฏฺฐิโย. อปรนฺตา\-นุทิฏฺฐิยา จตุจตฺตารีสาย อากาเรหิ อภินิเวโส สิยา ภวทิฏฺฐิโย, สิยา วิภว\-ทิฏฺฐิโย. สญฺโญชนิกาย ทิฏฺฐิยา อฏฺฐารสหิ อากาเรหิ อภินิเวโส สิยา ภวทิฏฺฐิโย, สิยา วิภว\-ทิฏฺฐิโย. “อหนฺ”ติ มาน\-วินิพนฺธาย ทิฏฺฐิยา อฏฺฐารสหิ อากาเรหิ อภินิเวโส สพฺพาว ตา วิภว\-ทิฏฺฐิโย. “มมนฺ”ติ มาน\-วินิพนฺธาย ทิฏฺฐิยา อฏฺฐารสหิ อากาเรหิ อภินิเวโส สพฺพาว ตา ภวทิฏฺฐิโย. อตฺตวาท\-ปฏิสํยุตฺตาย ทิฏฺฐิยา วีสติยา อากาเรหิ อภินิเวโส ปนฺนรส ภวทิฏฺฐิโย, ปญฺจ วิภว\-ทิฏฺฐิโย. โลกวาท\-ปฏิสํยุตฺตาย ทิฏฺฐิยา อฏฺฐหิ อากาเรหิ อภินิเวโส สิยา ภวทิฏฺฐิโย, สิยา วิภว\-ทิฏฺฐิโย.}

\prose{สพฺพาว ตา ทิฏฺฐิโย อสฺสาททิฏฺฐิโย. สพฺพาว ตา ทิฏฺฐิโย อตฺตานุทิฏฺฐิโย. สพฺพาว ตา ทิฏฺฐิโย มิจฺฉาทิฏฺฐิโย. สพฺพาว ตา ทิฏฺฐิโย สกฺกาย\-ทิฏฺฐิโย สพฺพาว ตา ทิฏฺฐิโย อนฺตคฺคาหิกา ทิฏฺฐิโย. สพฺพาว ตา ทิฏฺฐิโย สญฺโญชนิกา ทิฏฺฐิโย. สพฺพาว ตา ทิฏฺฐิโย อตฺตวาท\-ปฏิสํยุตฺตา ทิฏฺฐิโย.}

\setlength{\csromangathapagewidth}{0pt}
\setlength{\csromangathawakwidth}{0pt}
\csromangathameasuregroup{}{ภวญฺจ ทิฏฺฐิํ วิภวญฺจ ทิฏฺฐิํ, \\ เอตํ ทฺวยํ ตกฺกิกา นิสฺสิตาเส. \\ เตสํ นิโรธมฺหิ น ห’ตฺถิ ญาณํ, \\ ยตฺถายํ โลโก วิปรีตสญฺญีติ.}
\csromangathameasurewak{ภวญฺจ ทิฏฺฐิํ วิภวญฺจ ทิฏฺฐิํ, \\ เอตํ ทฺวยํ ตกฺกิกา นิสฺสิตาเส. \\ เตสํ นิโรธมฺหิ น ห’ตฺถิ ญาณํ, \\ ยตฺถายํ โลโก วิปรีตสญฺญีติ.}
\setlength{\csromangathaleftcol}{0pt}
\csromangathagroup{\csromangathastanza{\csromangathawak{\csromanfolio{159}ภวญฺจ ทิฏฺฐิํ วิภวญฺจ ทิฏฺฐิํ, \\ เอตํ ทฺวยํ ตกฺกิกา นิสฺสิตาเส. \\ เตสํ นิโรธมฺหิ น ห’ตฺถิ ญาณํ, \\ ยตฺถายํ โลโก วิปรีตสญฺญีติ.}}}

\proseitem{149}{ทฺวีหิ ภิกฺขเว ทิฏฺฐิคเตหิ ปริยุฏฺฐิตา เทวมนุสฺสา โอลียนฺติ\csromansharedfootnote{17ed8e51169e47d6}{โอลิยนฺติ (สฺยา, ก) ขุ 1. 224 ปิฏฺเฐ ปสฺสิตพฺพา.} เอเก, อติธาวนฺติ เอเก, จกฺขุมนฺโต จ ปสฺสนฺติ. กถญฺจ ภิกฺขเว โอลียนฺติ เอเก? ภวารามา ภิกฺขเว เทวมนุสฺสา ภวรตา ภวสมฺมุทิตา, เตสํ ภวนิโรธาย ธมฺเม เทสิยมาเน จิตฺตํ น ปกฺขนฺทติ น ปสีทติ น สนฺติฏฺฐติ นาธิมุจฺจติ. เอวํ โข ภิกฺขเว โอลียนฺติ เอเก.}

\prose{กถญฺจ ภิกฺขเว อติธาวนฺติ เอเก? ภเวเนว โข ปเนเก อฏฺฏียมานา\csromansharedfootnote{365e9e1ec10ebfb5}{อฏฺฏิยมานา (สฺยา, ก) ขุ 1. 224 ปิฏฺเฐ ปสฺสิตพฺพา.} หรายมานา ชิคุจฺฉมานา วิภวํ อภินนฺทนฺติ, ยโต กิร โภ อยํ อตฺตา กายสฺส เภทา ปรํ มรณา อุจฺฉิชฺชติ วินสฺสติ น โหติ ปรํ มรณา, เอตํ สนฺตํ เอตํ ปณีตํ เอตํ ยาถาวนฺติ. เอวํ โข ภิกฺขเว อติธาวนฺติ เอเก.}

\prose{กถญฺจ ภิกฺขเว จกฺขุมนฺโต จ ปสฺสนฺติ? อิธ ภิกฺขเว ภิกฺขุ ภูตํ ภูตโต ปสฺสติ, ภูตํ ภูตโต ทิสฺวา ภูตสฺส นิพฺพิทาย วิราคาย นิโรธาย ปฏิปนฺโน โหติ. เอวํ โข ภิกฺขเว จกฺขุมนฺโต จ ปสฺสนฺติ.}

\setlength{\csromangathapagewidth}{0pt}
\setlength{\csromangathawakwidth}{0pt}
\csromangathameasuregroup{“โย\textsuperscript{1} ภูตํ ภูตโต ทิสฺวา, \\ ยถาภูเต’ธิมุจฺจติ, \\ ส เว ภูตปริญฺญาโต, \\ ภูตสฺส วิภวา ภิกฺขุ,}{\csromangathabat{“โย\textsuperscript{1} ภูตํ ภูตโต ทิสฺวา,}{ภูตสฺส จ อติกฺกมํ.} \\ \csromangathabat{ยถาภูเต’ธิมุจฺจติ,}{ภวตณฺหา ปริกฺขยา.} \\ \csromangathabat{ส เว ภูตปริญฺญาโต,}{วีตตโณฺห ภวาภเว.} \\ \csromangathabat{ภูตสฺส วิภวา ภิกฺขุ,}{นาคจฺฉติ ปุนพฺภวนฺ”ติ.}}
\csromangathasetleft{“โย\textsuperscript{1} ภูตํ ภูตโต ทิสฺวา, \\ ยถาภูเต’ธิมุจฺจติ, \\ ส เว ภูตปริญฺญาโต, \\ ภูตสฺส วิภวา ภิกฺขุ,}
\csromangathagroup{\csromangathastanza{\csromangathabat{“โย\csromansharedfootnote{b2b7268fe1253db2}{อยํ คาถา ขุ 1. 225 ปิฏฺเฐ ทิสฺสติ.} ภูตํ ภูตโต ทิสฺวา,}{ภูตสฺส จ อติกฺกมํ.} \\ \csromangathabat{ยถาภูเต’ธิมุจฺจติ,}{ภวตณฺหา ปริกฺขยา.}}\csromangathastanza{\csromangathabat{ส เว ภูตปริญฺญาโต,}{วีตตโณฺห ภวาภเว.} \\ \csromangathabat{ภูตสฺส วิภวา ภิกฺขุ,}{นาคจฺฉติ ปุนพฺภวนฺ”ติ.}}}

\proseitem{150}{ตโย ปุคฺคลา วิปนฺนทิฏฺฐี, ตโย ปุคฺคลา สมฺปนฺนทิฏฺฐี. กตเม ตโย ปุคฺคลา วิปนฺนทิฏฺฐี? ติตฺถิโย จ ติตฺถิยสาวโก จ โย จ มิจฺฉาทิฏฺฐิโก. อิเม ตโย ปุคฺคลา วิปนฺนทิฏฺฐี.}

\prose{กตเม ตโย ปุคฺคลา สมฺปนฺนทิฏฺฐี? ตถาคโต จ ตถาคต\-สาวโก จ โย จ สมฺมาทิฏฺฐิโก. อิเม ตโย ปุคฺคลา สมฺปนฺนทิฏฺฐี.}

\setlength{\csromangathapagewidth}{0pt}
\setlength{\csromangathawakwidth}{0pt}
\csromangathameasuregroup{โกธโน อุปนาหี จ, \\ วิปนฺนทิฏฺฐิ มายาวี, \\ อกฺโกธโน อนุปนาหี, \\ สมฺปนฺนทิฏฺฐิ เมธาวี,}{\csromangathabat{โกธโน อุปนาหี จ,}{ปาปมกฺขี จ โย นโร.} \\ \csromangathabat{วิปนฺนทิฏฺฐิ มายาวี,}{ตํ ชญฺญา วสโล อิติ\textsuperscript{1}.} \\ \csromangathabat{อกฺโกธโน อนุปนาหี,}{วิสุทฺโธ\textsuperscript{1} สุทฺธตํ คโต.} \\ \csromangathabat{สมฺปนฺนทิฏฺฐิ เมธาวี,}{ตํ ชญฺญา อริโย อิตีติ.}}
\csromangathasetleft{โกธโน อุปนาหี จ, \\ วิปนฺนทิฏฺฐิ มายาวี, \\ อกฺโกธโน อนุปนาหี, \\ สมฺปนฺนทิฏฺฐิ เมธาวี,}
\csromangathagroup{\csromangathastanza{\csromanfolio{160}\csromangathabat{โกธโน อุปนาหี จ,}{ปาปมกฺขี จ โย นโร.} \\ \csromangathabat{วิปนฺนทิฏฺฐิ มายาวี,}{ตํ ชญฺญา วสโล อิติ\csromansharedfootnote{ebe911b8658f1f05}{อิเม (สฺยา) ขุ 1. 298 ปิฏฺเฐ ปสฺสิตพฺพา.}.}}\csromangathastanza{\csromangathabat{อกฺโกธโน อนุปนาหี,}{วิสุทฺโธ\csromansharedfootnote{d61cdb7dae96b5d0}{อมกฺขี (สฺยา) อฏฺฐกถา โอโลเกตพฺพา. 3. สกฺกายทิฏฺฐิปฺปมุเขน (สฺยา)} สุทฺธตํ คโต.} \\ \csromangathabat{สมฺปนฺนทิฏฺฐิ เมธาวี,}{ตํ ชญฺญา อริโย อิตีติ.}}}

\prose{ติสฺโส วิปนฺน\-ทิฏฺฐิโย, ติสฺโส สมฺปนฺน\-ทิฏฺฐิโย. กตมา ติสฺโส วิปนฺน\-ทิฏฺฐิโย? “เอตํ มมา”ติ วิปนฺนทิฏฺฐิ, “เอโสหมสฺมี”ติ วิปนฺนทิฏฺฐิ, “เอโส เม อตฺตา”ติ วิปนฺนทิฏฺฐิ. อิมา ติสฺโส วิปนฺน\-ทิฏฺฐิโย.}

\prose{กตมา ติสฺโส สมฺปนฺน\-ทิฏฺฐิโย? “เนตํ มมา”ติ สมฺปนฺนทิฏฺฐิ, “เนโสหมสฺมี”ติ สมฺปนฺนทิฏฺฐิ, “น เมโส อตฺตา”ติ สมฺปนฺนทิฏฺฐิ. อิมา ติสฺโส สมฺปนฺน\-ทิฏฺฐิโย.}

\prose{“เอตํ มมา”ติ กา ทิฏฺฐิ, กติ ทิฏฺฐิโย, กตมนฺตา\-นุคฺคหิตา ตา ทิฏฺฐิโย. “เอโสหมสฺมี”ติ กา ทิฏฺฐิ, กติ ทิฏฺฐิโย, กตมนฺตา\-นุคฺคหิตา ตา ทิฏฺฐิโย. “เอโส เม อตฺตา”ติ กา ทิฏฺฐิ, กติ ทิฏฺฐิโย, กตมนฺตา\-นุคฺคหิตา ตา ทิฏฺฐิโย?}

\prose{“เอตํ มมา”ติ ปุพฺพนฺตา\-นุทิฏฺฐิ, อฏฺฐารส ทิฏฺฐิโย, ปุพฺพนฺตา\-นุคฺคหิตา ตา ทิฏฺฐิโย. “เอโสหมสฺมี”ติ อปรนฺตา\-นุทิฏฺฐิ, จตุจตฺตารีสํ ทิฏฺฐิโย, อปรนฺตา\-นุคฺคหิตา ตา ทิฏฺฐิโย. “เอโส เม อตฺตา”ติ วีสติวตฺถุกา อตฺตานุทิฏฺฐิ วีสติวตฺถุกา สกฺกายทิฏฺฐิ, สกฺกายทิฏฺฐิ\-ปฺปมุขานิ\csromansharedfootnote{9a4550b4ef050c68}{สกฺกายทิฏฺฐิปฺปมุเขน (สฺยา)} ทฺวาสฏฺฐิ ทิฏฺฐิคตานิ, ปุพฺพนฺตาปรนฺตา\-นุคฺคหิตา ตา ทิฏฺฐิโย.}

\proseitem{151}{เย เกจิ ภิกฺขเว มยิ นิฏฺฐํ คตา, สพฺเพ เต ทิฏฺฐิสมฺปนฺนา. เตสํ ทิฏฺฐิ\-สมฺปนฺนานํ ปญฺจนฺนํ อิธ นิฏฺฐา, ปญฺจนฺนํ อิธ วิหาย นิฏฺฐา. กตเมสํ ปญฺจนฺนํ อิธ นิฏฺฐา? สตฺตกฺขตฺตุ\-ปรมสฺส, โกลํโกลสฺส, เอกพีชิสฺส สกทาคามิสฺส โย จ ทิฏฺเฐว ธมฺเม อรหา. อิเมสํ ปญฺจนฺนํ อิธ นิฏฺฐา.}

\prose{กตเมสํ ปญฺจนฺนํ อิธ วิหาย นิฏฺฐา? อนฺตรา\-ปรินิพฺพายิสฺส, อุปหจฺจ\-ปรินิพฺพายิสฺส, อสงฺขารปรินิพฺพายิสฺส, สสงฺขาร\-ปรินิพฺพายิสฺส, อุทฺธํโสตสฺส อกนิฏฺฐ\-คามิโน. อิเมสํ ปญฺจนฺนํ อิธ วิหาย นิฏฺฐา.}

\prose{\csromanfolio{161}เย เกจิ ภิกฺขเว มยิ นิฏฺฐํ คตา, สพฺเพ เต ทิฏฺฐิสมฺปนฺนา, เตสํ ทิฏฺฐิ\-สมฺปนฺนานํ อิเมสํ ปญฺจนฺนํ อิธ นิฏฺฐา, อิเมสํ ปญฺจนฺนํ อิธ วิหาย นิฏฺฐา.}

\prose{เย เกจิ ภิกฺขเว มยิ อเวจฺจปฺปสนฺนา, สพฺเพ เต โสตาปนฺนา, เตสํ โสตาปนฺนานํ ปญฺจนฺนํ อิธ นิฏฺฐา, ปญฺจนฺนํ อิธ วิหาย นิฏฺฐา, กตเมสํ ปญฺจนฺนํ อิธ นิฏฺฐา? สตฺตกฺขตฺตุ\-ปรมสฺส โกลํโกลสฺส เอกพีชิสฺส สกทาคามิสฺส โย จ ทิฏฺเฐว ธมฺเม อรหา, อิเมสํ ปญฺจนฺนํ อิธ นิฏฺฐา.}

\prose{กตเมสํ ปญฺจนฺนํ อิธ วิหาย นิฏฺฐา? อนฺตรา\-ปรินิพฺพายิสฺส อุปหจฺจ\-ปรินิพฺพายิสฺส อสงฺขารปรินิพฺพายิสฺส สสงฺขาร\-ปรินิพฺพายิสฺส อุทฺธํโสตสฺส อกนิฏฺฐ\-คามิโน, อิเมสํ ปญฺจนฺนํ อิธ วิหาย นิฏฺฐา.}

\prose{เย เกจิ ภิกฺขเว มยิ อเวจฺจปฺปสนฺนา, สพฺเพ เต โสตาปนฺนา, เตสํ โสตาปนฺนานํ อิเมสํ ปญฺจนฺนํ อิธ นิฏฺฐา, อิเมสํ ปญฺจนฺนํ อิธ วิหาย นิฏฺฐาติ.}

\vspace{\tipitakaparskip}
\csromancenter{ภววิภวทิฏฺฐิ\-นิทฺเทโส โสฬสโม.}
\nitthitamruled{ทิฏฺฐิกถา นิฏฺฐิตา.}
\csromanmatikaanchor{mtk.65}
\csromantocmarknopage{section}{3. อานาปานสฺสติกถา}{mtk.65}
\bookmark[dest={mtk.65},level=1]{3. อานาปานสฺสติกถา}
\csromanheader{1}{3. อานาปานสฺสติ\-กถา}
\csromanmatikaanchor{mtk.66}
\csromantocmarkref{subsection}{1. คณนวาร}{mtk.66}
\bookmark[dest={mtk.66},level=2]{1. คณนวาร}
\csromanheader{2}{1. \textbf{คณนวาร}}
\proseitem{152}{โสฬส\-วตฺถุกํ อานาปาน\-สฺสติ\-สมาธิํ\csromansharedfootnote{992e1ecc59128f2b}{อานาปานสติสมาธิํ (สีฏฺฐ)} ภาวยโต สมธิกานิ เทฺว ญาณสตานิ อุปฺปชฺชนฺติ, อฏฺฐ ปริปนฺเถ\csromansharedfootnote{04edc6c95358c2d5}{ปริพนฺเธ (ก)} ญาณานิ, อฏฺฐ จ อุปกาเร ญาณานิ, อฏฺฐารส อุปกฺกิเลเส ญาณานิ, เตรส โวทาเน ญาณานิ, พาตฺติํส สโตการิสฺส\csromansharedfootnote{fad8b7852414fa8c}{สโตการีสุ (สฺยา)} ญาณานิ, จตุวีสติ สมาธิวเสน ญาณานิ, เทฺวสตฺตติ วิปสฺสนา\-วเสน ญาณานิ, อฏฺฐ นิพฺพิทาญาณานิ, อฏฺฐ นิพฺพิทา\-นุโลมญาณานิ, อฏฺฐ นิพฺพิทา\-ปฏิปฺปสฺสทฺธิ\-ญาณานิ, เอกวีสติ วิมุตฺติสุเข ญาณานิ.}

\prosewithcloserruled{กตมานิ อฏฺฐ ปริปนฺเถ ญาณานิ, อฏฺฐ จ อุปกาเร ญาณานิ? กามจฺฉนฺโท สมาธิสฺส ปริปนฺโถ, เนกฺขมฺมํ สมาธิสฺส อุปการํ. พฺยาปาโท \csromanfolio{162}สมาธิสฺส ปริปนฺโถ, อพฺยาปาโท สมาธิสฺส อุปการํ. ถินมิทฺธํ สมาธิสฺส ปริปนฺโถ, อาโลกสญฺญา สมาธิสฺส อุปการํ. อุทฺธจฺจํ สมาธิสฺส ปริปนฺโถ, อวิกฺเขโป สมาธิสฺส อุปการํ. วิจิกิจฺฉา สมาธิสฺส ปริปนฺโถ, ธมฺม\-ววตฺถานํ สมาธิสฺส อุปการํ. อวิชฺชา สมาธิสฺส ปริปนฺโถ, ญาณํ สมาธิสฺส อุปการํ. อรติ สมาธิสฺส ปริปนฺโถ, ปาโมชฺชํ สมาธิสฺส อุปการํ. สพฺเพปิ อกุสลา ธมฺมา สมาธิสฺส ปริปนฺถา, สพฺเพปิ กุสลา ธมฺมา สมาธิสฺส อุปการา. อิมานิ อฏฺฐ ปริปนฺเถ ญาณานิ, อฏฺฐ จ อุปกาเร ญาณานิ.}{คณนวาโร ปฐโม.}
\csromanmatikaanchor{mtk.67}
\csromantocmarkref{subsection}{2. โสฬสญาณนิทฺเทส}{mtk.67}
\bookmark[dest={mtk.67},level=2]{2. โสฬสญาณนิทฺเทส}
\csromanheader{2}{2. โสฬส\-ญาณ\-นิทฺเทส}
\proseitem{153}{อิเมหิ โสฬสหิ อากาเรหิ อุทุจิตํ จิตฺตํ สมุทุจิตํ จิตฺตํ เอกตฺเต สนฺติฏฺฐติ, นีวรเณหิ วิสุชฺฌติ. กตเม เต เอกตฺตา? เนกฺขมฺมํ เอกตฺตํ, อพฺยาปาโท เอกตฺตํ, อาโลกสญฺญา เอกตฺตํ, อวิกฺเขโป เอกตฺตํ, ธมฺม\-ววตฺถานํ เอกตฺตํ, ญาณํ เอกตฺตํ, ปาโมชฺชํ เอกตฺตํ, สพฺเพปิ กุสลา ธมฺมา เอกตฺตา.}

\prose{นีวรณาติ กตเม เต นีวรณา? กามจฺฉนฺโท นีวรณํ, พฺยาปาโท นีวรณํ, ถินมิทฺธํ นีวรณํ, อุทฺธจฺจ\-กุกฺกุจฺจํ นีวรณํ, วิจิกิจฺฉา นีวรณํ, อวิชฺชา นีวรณํ, อรติ นีวรณํ, สพฺเพปิ อกุสลา ธมฺมา นีวรณา.}

\prosewithcloserruled{นีวรณาติ เกนฏฺเฐน นีวรณา? นิยฺยานา\-วรณ\-ฏฺเฐน นีวรณา. กตเม เต นิยฺยานา? เนกฺขมฺมํ อริยานํ นิยฺยานํ, เตน จ เนกฺขมฺเมน อริยา นิยฺยนฺติ, กามจฺฉนฺโท นิยฺยานาวรณํ, เตน จ กามจฺฉนฺเทน นิวุตตฺตา เนกฺขมฺมํ อริยานํ นิยฺยานํ นปฺปชานาตีติ กามจฺฉนฺโท นิยฺยานาวรณํ.  อพฺยาปาโท อริยานํ นิยฺยานํ, เตน จ อพฺยาปาเทน อริยา นิยฺยนฺติ, พฺยาปาโท นิยฺยานาวรณํ, เตน จ พฺยาปาเทน นิวุตตฺตา อพฺยาปาทํ อริยานํ นิยฺยานํ นปฺปชานาตีติ พฺยาปาโท นิยฺยานาวรณํ.  อาโลกสญฺญา อริยานํ นิยฺยานํ, ตาย จ อาโลกสญฺญาย อริยา นิยฺยนฺติ, ถินมิทฺธํ นิยฺยานาวรณํ, เตน จ ถินมิทฺเธน นิวุตตฺตา อาโลกสญฺญํ อริยานํ นิยฺยานํ นปฺปชานาตีติ ถินมิทฺธํ นิยฺยานาวรณํ.  อวิกฺเขโป อริยานํ นิยฺยานํ, เตน จ อวิกฺเขเปน อริยา นิยฺยนฺติ, อุทฺธจฺจํ \csromanfolio{163}นิยฺยานาวรณํ, เตน จ อุทฺธจฺเจน นิวุตตฺตา อวิกฺเขปํ อริยานํ นิยฺยานํ นปฺปชานาตีติ อุทฺธจฺจํ นิยฺยานาวรณํ.  ธมฺม\-ววตฺถานํ อริยานํ นิยฺยานํ, เตน จ ธมฺม\-ววตฺถาเนน อริยา นิยฺยนฺติ, วิจิกิจฺฉา นิยฺยานาวรณํ, ตาย จ วิจิกิจฺฉาย นิวุตตฺตา ธมฺม\-ววตฺถานํ อริยานํ นิยฺยานํ นปฺปชานาตีติ วิจิกิจฺฉา นิยฺยานาวรณํ.  ญาณํ อริยานํ นิยฺยานํ, เตน จ ญาเณน อริยา นิยฺยนฺติ, อวิชฺชา นิยฺยานาวรณํ, ตาย จ อวิชฺชาย นิวุตตฺตา ญาณํ อริยานํ นิยฺยานํ นปฺปชานาตีติ อวิชฺชา นิยฺยานาวรณํ.  ปาโมชฺชํ อริยานํ นิยฺยานํ, เตน จ ปาโมชฺเชน อริยา นิยฺยนฺติ, อรติ นิยฺยานาวรณํ, ตาย จ อรติยา นิวุตตฺตา ปาโมชฺชํ อริยานํ นิยฺยานํ นปฺปชานาตีติ อรติ นิยฺยานาวรณํ.  สพฺเพปิ กุสลา ธมฺมา อริยานํ นิยฺยานํ, เตหิ จ กุสเลหิ ธมฺเมหิ อริยา นิยฺยนฺติ. สพฺเพปิ อกุสลา ธมฺมา นิยฺยานาวรณา, เตหิ จ อกุสเลหิ ธมฺเมหิ นิวุตตฺตา กุสเล ธมฺเม อริยานํ นิยฺยานํ นปฺปชานาตีติ สพฺเพปิ อกุสลา ธมฺมา นิยฺยานาวรณา.}{โสฬส\-ญาณ\-นิทฺเทโส ทุติโย.}
\csromanmatikaanchor{mtk.68}
\csromantocmarkref{subsection}{3. อุปกฺกิเลสญาณนิทฺเทส}{mtk.68}
\bookmark[dest={mtk.68},level=2]{3. อุปกฺกิเลสญาณนิทฺเทส}
\csromanheader{2}{3. อุปกฺกิเลส\-ญาณ\-นิทฺเทส}
\csromanheader{2}{\textbf{ปฐมจฺฉกฺก}}
\proseitem{154}{อิเมหิ จ ปน นีวรเณหิ วิสุทฺธ\-จิตฺตสฺส โสฬส\-วตฺถุกํ อานาปาน\-สฺสติ\-สมาธิํ ภาวยโต ขณิก\-สโมธานา กตเม อฏฺฐารส อุปกฺกิเลสา อุปฺปชฺชนฺติ? อสฺสาสา\-ทิมชฺฌปริโยสานํ สติยา อนุคจฺฉโต อชฺฌตฺต\-วิกฺเขปคตํ จิตฺตํ สมาธิสฺส ปริปนฺโถ,}

\prose{ปสฺสาสา\-ทิมชฺฌปริโยสานํ สติยา อนุคจฺฉโต พหิทฺธา\-วิกฺเขปคตํ จิตฺตํ สมาธิสฺส ปริปนฺโถ, อสฺสาส\-ปฏิก\-งฺขนา นิกนฺติ\-ตณฺหา\-จริยา สมาธิสฺส ปริปนฺโถ, ปสฺสาส\-ปฏิก\-งฺขนา นิกนฺติ\-ตณฺหา\-จริยา สมาธิสฺส ปริปนฺโถ, อสฺสาเสนา\-ภิตุนฺนสฺส ปสฺสาส\-ปฏิลาเภ มุจฺฉนา สมาธิสฺส ปริปนฺโถ, ปสฺสาเสนา\-ภิตุนฺนสฺส อสฺสาส\-ปฏิลาเภ มุจฺฉนา สมาธิสฺส ปริปนฺโถ.}

\setlength{\csromangathapagewidth}{0pt}
\setlength{\csromangathawakwidth}{0pt}
\csromangathameasuregroup{อนุคจฺฉนา จ อสฺสาสํ,}{\csromangathabat{อนุคจฺฉนา จ อสฺสาสํ,}{ปสฺสาสํ, อนุคจฺฉนา.}}
\csromangathasetleft{อนุคจฺฉนา จ อสฺสาสํ,}
\csromangathagroup{\csromangathastanza{\csromangathaitembat{154}{อนุคจฺฉนา จ อสฺสาสํ,}{ปสฺสาสํ, อนุคจฺฉนา.}}}

\prose{\csromanfolio{164}สติ อชฺฌตฺต\-วิกฺเขปา\-กงฺขนา พหิทฺธา\-วิกฺเขป\-ปตฺถนา\csromansharedfootnote{92624d34a73a4026}{วิกฺเขปปนฺถนา (สฺยา)}.}

\setlength{\csromangathapagewidth}{0pt}
\setlength{\csromangathawakwidth}{0pt}
\csromangathameasuregroup{อสฺสาเสนาภิตุนฺนสฺส, \\ ปสฺสาเสนาภิตุนฺนสฺส, \\ ฉ เอเต อุปกฺกิเลสา, \\ เยหิ วิกฺขิปฺปมานสฺส\textsuperscript{1}, \\ วิโมกฺขํ อปฺปชานนฺตา,}{\csromangathabat{อสฺสาเสนาภิตุนฺนสฺส,}{ปสฺสาสปฏิลาเภ มุจฺฉนา.} \\ \csromangathabat{ปสฺสาเสนาภิตุนฺนสฺส,}{อสฺสาสปฏิลาเภ มุจฺฉนา.} \\ \csromangathabat{ฉ เอเต อุปกฺกิเลสา,}{อานาปานสฺสติสมาธิสฺส.} \\ \csromangathabat{เยหิ วิกฺขิปฺปมานสฺส\textsuperscript{1},}{โน จ จิตฺตํ วิมุจฺจติ.} \\ \csromangathabat{วิโมกฺขํ อปฺปชานนฺตา,}{เต โหนฺติ ปรปตฺติยาติ.}}
\csromangathasetleft{อสฺสาเสนาภิตุนฺนสฺส, \\ ปสฺสาเสนาภิตุนฺนสฺส, \\ ฉ เอเต อุปกฺกิเลสา, \\ เยหิ วิกฺขิปฺปมานสฺส\textsuperscript{1}, \\ วิโมกฺขํ อปฺปชานนฺตา,}
\csromangathagroup{\csromangathastanza{\csromangathabat{อสฺสาเสนา\-ภิตุนฺนสฺส,}{ปสฺสาส\-ปฏิลาเภ มุจฺฉนา.} \\ \csromangathabat{ปสฺสาเสนา\-ภิตุนฺนสฺส,}{อสฺสาส\-ปฏิลาเภ มุจฺฉนา.}}\csromangathastanza{\csromangathabat{ฉ เอเต อุปกฺกิเลสา,}{อานาปาน\-สฺสติ\-สมาธิสฺส.} \\ \csromangathabat{เยหิ วิกฺขิปฺปมานสฺส\csromansharedfootnote{546e2f21eafcc599}{วิกมฺปมานสฺส (สฺยา)},}{โน จ จิตฺตํ วิมุจฺจติ.} \\ \csromangathabat{วิโมกฺขํ อปฺปชานนฺตา,}{เต โหนฺติ ปรปตฺติยาติ.}}}

\csromanheader{2}{\textbf{ทุติยจฺฉกฺก}}
\proseitem{155}{นิมิตฺตํ อาวชฺชโต อสฺสาเส จิตฺตํ วิกมฺปติ, สมาธิสฺส ปริปนฺโถ. อสฺสาสํ อาวชฺชโต นิมิตฺเต จิตฺตํ วิกมฺปติ, สมาธิสฺส ปริปนฺโถ. นิมิตฺตํ อาวชฺชโต ปสฺสาเส จิตฺตํ วิกมฺปติ, สมาธิสฺส ปริปนฺโถ. ปสฺสาสํ อาวชฺชโต นิมิตฺเต จิตฺตํ วิกมฺปติ, สมาธิสฺส ปริปนฺโถ. อสฺสาสํ อาวชฺชโต ปสฺสาเส จิตฺตํ วิกมฺปติ, สมาธิสฺส ปริปนฺโถ. ปสฺสาสํ อาวชฺชโต อสฺสาเส จิตฺตํ วิกมฺปติ, สมาธิสฺส ปริปนฺโถ.}

\setlength{\csromangathapagewidth}{0pt}
\setlength{\csromangathawakwidth}{0pt}
\csromangathameasuregroup{นิมิตฺตํ อาวชฺชมานสฺส, \\ อสฺสาสํ อาวชฺชมานสฺส, \\ นิมิตฺตํ อาวชฺชมานสฺส, \\ ปสฺสาสํ อาวชฺชมานสฺส, \\ อสฺสาสํ อาวชฺชมานสฺส, \\ ปสฺสาสํ อาวชฺชมานสฺส, \\ ฉ เอเต อุปกฺกิเลสา, \\ เยหิ วิกฺขิปฺปมานสฺส, \\ วิโมกฺขํ อปฺปชานนฺตา,}{\csromangathabat{นิมิตฺตํ อาวชฺชมานสฺส,}{อสฺสาเส วิกฺขิปฺปเต มโน.} \\ \csromangathabat{อสฺสาสํ อาวชฺชมานสฺส,}{นิมิตฺเต จิตฺตํ วิกมฺปติ.} \\ \csromangathabat{นิมิตฺตํ อาวชฺชมานสฺส,}{ปสฺสาเส วิกฺขิปฺปเต มโน.} \\ \csromangathabat{ปสฺสาสํ อาวชฺชมานสฺส,}{นิมิตฺเต จิตฺตํ วิกมฺปติ.} \\ \csromangathabat{อสฺสาสํ อาวชฺชมานสฺส,}{ปสฺสาเส วิกฺขิปฺปเต มโน.} \\ \csromangathabat{ปสฺสาสํ อาวชฺชมานสฺส,}{อสฺสาเส จิตฺตํ วิกมฺปติ.} \\ \csromangathabat{ฉ เอเต อุปกฺกิเลสา,}{อานาปานสฺสติสมาธิสฺส.} \\ \csromangathabat{เยหิ วิกฺขิปฺปมานสฺส,}{โน จ จิตฺตํ วิมุจฺจติ.} \\ \csromangathabat{วิโมกฺขํ อปฺปชานนฺตา,}{เต โหนฺติ ปรปตฺติยาติ.}}
\csromangathasetleft{นิมิตฺตํ อาวชฺชมานสฺส, \\ อสฺสาสํ อาวชฺชมานสฺส, \\ นิมิตฺตํ อาวชฺชมานสฺส, \\ ปสฺสาสํ อาวชฺชมานสฺส, \\ อสฺสาสํ อาวชฺชมานสฺส, \\ ปสฺสาสํ อาวชฺชมานสฺส, \\ ฉ เอเต อุปกฺกิเลสา, \\ เยหิ วิกฺขิปฺปมานสฺส, \\ วิโมกฺขํ อปฺปชานนฺตา,}
\csromangathagroup{\csromangathastanza{\csromangathabat{นิมิตฺตํ อาวชฺชมานสฺส,}{อสฺสาเส วิกฺขิปฺปเต มโน.} \\ \csromangathabat{อสฺสาสํ อาวชฺชมานสฺส,}{นิมิตฺเต จิตฺตํ วิกมฺปติ.}}\csromangathastanza{\csromangathabat{นิมิตฺตํ อาวชฺชมานสฺส,}{ปสฺสาเส วิกฺขิปฺปเต มโน.} \\ \csromangathabat{ปสฺสาสํ อาวชฺชมานสฺส,}{นิมิตฺเต จิตฺตํ วิกมฺปติ.}}\csromangathastanza{\csromangathabat{อสฺสาสํ อาวชฺชมานสฺส,}{ปสฺสาเส วิกฺขิปฺปเต มโน.} \\ \csromangathabat{ปสฺสาสํ อาวชฺชมานสฺส,}{อสฺสาเส จิตฺตํ วิกมฺปติ.}}\csromangathastanza{\csromangathabat{ฉ เอเต อุปกฺกิเลสา,}{อานาปาน\-สฺสติ\-สมาธิสฺส.} \\ \csromangathabat{เยหิ วิกฺขิปฺปมานสฺส,}{โน จ จิตฺตํ วิมุจฺจติ.} \\ \csromangathabat{วิโมกฺขํ อปฺปชานนฺตา,}{เต โหนฺติ ปรปตฺติยาติ.}}}

\csromanheader{2}{\textbf{ตติยจฺฉกฺก}}
\proseitem{156}{อตีตา\-นุ\-ธาวนํ จิตฺตํ วิกฺเขปา\-นุปติตํ สมาธิสฺส ปริปนฺโถ. อนาคต\-ปฏิก\-งฺขนํ จิตฺตํ วิกมฺปิตํ สมาธิสฺส ปริปนฺโถ. ลีนํ จิตฺตํ \csromanfolio{165}โกสชฺชา\-นุปติตํ สมาธิสฺส ปริปนฺโถ. อติปคฺคหิตํ จิตฺตํ อุทฺธจฺจา\-นุปติตํ สมาธิสฺส ปริปนฺโถ. อภินตํ จิตฺตํ ราคานุปติตํ สมาธิสฺส ปริปนฺโถ. อปนตํ จิตฺตํ พฺยาปาทา\-นุปติตํ สมาธิสฺส ปริปนฺโถ.}

\setlength{\csromangathapagewidth}{0pt}
\setlength{\csromangathawakwidth}{0pt}
\csromangathameasuregroup{อตีตานุธาวนํ จิตฺตํ, \\ อติปคฺคหิตํ อภินตํ, \\ ฉ เอเต อุปกฺกิเลสา, \\ เยหิ อุปกฺกิลิฏฺฐสงฺกปฺโป,}{\csromangathabat{อตีตานุธาวนํ จิตฺตํ,}{อนาคตปฏิกงฺขนํ ลีนํ.} \\ \csromangathabat{อติปคฺคหิตํ อภินตํ,}{อปนตํ จิตฺตํ น สมาธิยติ.} \\ \csromangathabat{ฉ เอเต อุปกฺกิเลสา,}{อานาปานสฺสติสมาธิสฺส.} \\ \csromangathabat{เยหิ อุปกฺกิลิฏฺฐสงฺกปฺโป,}{อธิจิตฺตํ นปฺปชานาตีติ.}}
\csromangathasetleft{อตีตานุธาวนํ จิตฺตํ, \\ อติปคฺคหิตํ อภินตํ, \\ ฉ เอเต อุปกฺกิเลสา, \\ เยหิ อุปกฺกิลิฏฺฐสงฺกปฺโป,}
\csromangathagroup{\csromangathastanza{\csromangathabat{อตีตา\-นุ\-ธาวนํ จิตฺตํ,}{อนาคต\-ปฏิก\-งฺขนํ ลีนํ.} \\ \csromangathabat{อติปคฺคหิตํ อภินตํ,}{อปนตํ จิตฺตํ น สมาธิยติ.}}\csromangathastanza{\csromangathabat{ฉ เอเต อุปกฺกิเลสา,}{อานาปาน\-สฺสติ\-สมาธิสฺส.} \\ \csromangathabat{เยหิ อุปกฺกิลิฏฺฐสงฺกปฺโป,}{อธิจิตฺตํ นปฺปชานาตีติ.}}}

\proseitem{157}{อสฺสาสา\-ทิมชฺฌปริโยสานํ สติยา อนุคจฺฉโต อชฺฌตฺตํ วิกฺเขปคเตน จิตฺเตน กาโยปิ จิตฺตมฺปิ สารทฺธา จ โหนฺติ อิญฺชิตา จ ผนฺทิตา จ, ปสฺสาสา\-ทิมชฺฌปริโยสานํ สติยา อนุคจฺฉโต พหิทฺธา\-วิกฺเขปคเตน จิตฺเตน กาโยปิ จิตฺตมฺปิ สารทฺธา จ โหนฺติ อิญฺชิตา จ ผนฺทิตา จ, อสฺสาส\-ปฏิกงฺข\-นาย นิกนฺติยา ตณฺหาจริยาย กาโยปิ จิตฺตมฺปิ สารทฺธา จ โหนฺติ อิญฺชิตา จ ผนฺทิตา จ, ปสฺสาส\-ปฏิกงฺข\-นาย นิกนฺติยา ตณฺหาจริยาย กาโยปิ จิตฺตมฺปิ สารทฺธา จ โหนฺติ อิญฺชิตา จ ผนฺทิตา จ, อสฺสาเสนา\-ภิตุนฺนสฺส ปสฺสาส\-ปฏิลาเภ มุจฺฉิตตฺตา กาโยปิ จิตฺตมฺปิ สารทฺธา จ โหนฺติ อิญฺชิตา จ ผนฺทิตา จ, ปสฺสาเสนา\-ภิตุนฺนสฺส อสฺสาส\-ปฏิลาเภ มุจฺฉิตตฺตา กาโยปิ จิตฺตมฺปิ สารทฺธา จ โหนฺติ อิญฺชิตา จ ผนฺทิตา จ.}

\prose{นิมิตฺตํ อาวชฺชโต อสฺสาเส จิตฺตํ วิกมฺปิตตฺตา กาโยปิ จิตฺตมฺปิ สารทฺธา จ โหนฺติ อิญฺชิตา จ ผนฺทิตา จ, อสฺสาสํ อาวชฺชโต นิมิตฺเต จิตฺตํ วิกมฺปิตตฺตา กาโยปิ จิตฺตมฺปิ สารทฺธา จ โหนฺติ อิญฺชิตา จ ผนฺทิตา จ. นิมิตฺตํ อาวชฺชโต ปสฺสาเส จิตฺตํ วิกมฺปิตตฺตา กาโยปิ จิตฺตมฺปิ สารทฺธา จ โหนฺติ อิญฺชิตา จ ผนฺทิตา จ, ปสฺสาสํ อาวชฺชโต นิมิตฺเต จิตฺตํ วิกมฺปิตตฺตา กาโยปิ จิตฺตมฺปิ สารทฺธา จ โหนฺติ อิญฺชิตา จ ผนฺทิตา จ. อสฺสาสํ อาวชฺชโต ปสฺสาเส จิตฺตํ วิกมฺปิตตฺตา กาโยปิ จิตฺตมฺปิ สารทฺธา จ โหนฺติ อิญฺชิตา จ ผนฺทิตา จ, ปสฺสาสํ อาวชฺชโต อสฺสาเส จิตฺตํ วิกมฺปิตตฺตา กาโยปิ จิตฺตมฺปิ สารทฺธา จ โหนฺติ อิญฺชิตา จ ผนฺทิตา จ.}

\prose{อตีตา\-นุ\-ธาวเนน จิตฺเตน วิกฺเขปา\-นุปติเตน กาโยปิ จิตฺตมฺปิ สารทฺธา จ โหนฺติ อิญฺชิตา จ ผนฺทิตา จ, อนาคต\-ปฏิกงฺข\-เนน จิตฺเตน วิกมฺปิเตน กาโยปิ จิตฺตมฺปิ สารทฺธา จ โหนฺติ อิญฺชิตา จ ผนฺทิตา จ. ลีเนน \csromanfolio{166}จิตฺเตน โกสชฺชา\-นุปติเตน กาโยปิ จิตฺตมฺปิ สารทฺธา จ โหนฺติ อิญฺชิตา จ ผนฺทิตา จ, อติปคฺคหิเตน จิตฺเตน อุทฺธจฺจา\-นุปติเตน กาโยปิ จิตฺตมฺปิ สารทฺธา จ โหนฺติ อิญฺชิตา จ ผนฺทิตา จ. อภินเตน จิตฺเตน ราคานุปติเตน กาโยปิ จิตฺตมฺปิ สารทฺธา จ โหนฺติ อิญฺชิตา จ ผนฺทิตา จ, อปนเตน จิตฺเตน พฺยาปาทา\-นุปติเตน กาโยปิ จิตฺตมฺปิ สารทฺธา จ โหนฺติ อิญฺชิตา จ ผนฺทิตา จ.}

\setlength{\csromangathapagewidth}{0pt}
\setlength{\csromangathawakwidth}{0pt}
\csromangathameasuregroup{อานาปานสฺสติ ยสฺส, \\ กาโยปิ อิญฺชิโต โหติ, \\ กาโยปิ ผนฺทิโต โหติ, \\ อานาปานสฺสติ ยสฺส, \\ กาโยปิ อนิญฺชิโต โหติ, \\ กาโยปิ อผนฺทิโต โหติ,}{\csromangathabat{อานาปานสฺสติ ยสฺส,}{ปริปุณฺณา อภาวิตา.} \\ \csromangathabat{กาโยปิ อิญฺชิโต โหติ,}{จิตฺตมฺปิ โหติ อิญฺชิตํ.} \\ \csromangathabat{กาโยปิ ผนฺทิโต โหติ,}{จิตฺตมฺปิ โหติ ผนฺทิตํ.} \\ \csromangathabat{อานาปานสฺสติ ยสฺส,}{ปริปุณฺณา สุภาวิตา.} \\ \csromangathabat{กาโยปิ อนิญฺชิโต โหติ,}{จิตฺตมฺปิ โหติ อนิญฺชิตํ.} \\ \csromangathabat{กาโยปิ อผนฺทิโต โหติ,}{จิตฺตมฺปิ โหติ อผนฺทิตนฺติ.}}
\csromangathasetleft{อานาปานสฺสติ ยสฺส, \\ กาโยปิ อิญฺชิโต โหติ, \\ กาโยปิ ผนฺทิโต โหติ, \\ อานาปานสฺสติ ยสฺส, \\ กาโยปิ อนิญฺชิโต โหติ, \\ กาโยปิ อผนฺทิโต โหติ,}
\csromangathagroup{\csromangathastanza{\csromangathabat{อานาปานสฺสติ ยสฺส,}{ปริปุณฺณา อภาวิตา.} \\ \csromangathabat{กาโยปิ อิญฺชิโต โหติ,}{จิตฺตมฺปิ โหติ อิญฺชิตํ.}}\csromangathastanza{\csromangathabat{กาโยปิ ผนฺทิโต โหติ,}{จิตฺตมฺปิ โหติ ผนฺทิตํ.} \\ \csromangathabat{อานาปานสฺสติ ยสฺส,}{ปริปุณฺณา สุภาวิตา.}}\csromangathastanza{\csromangathabat{กาโยปิ อนิญฺชิโต โหติ,}{จิตฺตมฺปิ โหติ อนิญฺชิตํ.} \\ \csromangathabat{กาโยปิ อผนฺทิโต โหติ,}{จิตฺตมฺปิ โหติ อผนฺทิตนฺติ.}}}

\prosewithcloserruled{อิเมหิ จ ปน นีวรเณหิ วิสุทฺธ\-จิตฺตสฺส โสฬส\-วตฺถุกํ อานาปาน\-สฺสติ\-สมาธิํ ภาวยโต ขณิก\-สโมธานา อิเม อฏฺฐารส อุปกฺกิเลสา อุปฺปชฺชนฺติ.}{อุปกฺกิเลส\-ญาณ\-นิทฺเทโส ตติโย.}
\csromanmatikaanchor{mtk.69}
\csromantocmarkref{subsection}{4. โวทานญาณนิทฺเทส}{mtk.69}
\bookmark[dest={mtk.69},level=2]{4. โวทานญาณนิทฺเทส}
\csromanheader{2}{4. โวทาน\-ญาณ\-นิทฺเทส}
\proseitem{158}{กตมานิ เตรส โวทาเน ญาณานิ? อตีตา\-นุ\-ธาวนํ จิตฺตํ วิกฺเขปา\-นุปติตํ, ตํ วิวชฺชยิตฺวา เอกฏฺฐาเน สมาทหติ, เอวมฺปิ จิตฺตํ น วิกฺเขปํ คจฺฉติ. อนาคต\-ปฏิก\-งฺขนํ จิตฺตํ วิกมฺปิตํ, ตํ วิวชฺชยิตฺวา ตตฺเถว อธิโมเจติ, เอวมฺปิ จิตฺตํ น วิกฺเขปํ คจฺฉติ. ลีนํ จิตฺตํ โกสชฺชา\-นุปติตํ, ตํ ปคฺคณฺหิตฺวา โกสชฺชํ ปชหติ, เอวมฺปิ จิตฺตํ น วิกฺเขปํ คจฺฉติ. อติปคฺคหิตํ จิตฺตํ อุทฺธจฺจา\-นุปติตํ, ตํ วินิคฺคณฺหิตฺวา อุทฺธจฺจํ ปชหติ, เอวมฺปิ จิตฺตํ น วิกฺเขปํ คจฺฉติ. อภินตํ จิตฺตํ ราคานุปติตํ, ตํ สมฺปชาโน หุตฺวา ราคํ ปชหติ, เอวมฺปิ จิตฺตํ น วิกฺเขปํ คจฺฉติ. อปนตํ จิตฺตํ พฺยาปาทา\-นุปติตํ, ตํ สมฺปชาโน หุตฺวา พฺยาปาทํ ปชหติ, เอวมฺปิ จิตฺตํ น วิกฺเขปํ คจฺฉติ. อิเมหิ ฉหิ ฐาเนหิ ปริสุทฺธํ จิตฺตํ ปริโยทาตํ เอกตฺตคตํ โหติ.}

\prose{\csromanfolio{167}กตเม เต เอกตฺตา? ทาน\-โวสคฺคุ\-ปฏฺฐาเน\-กตฺตํ, สมถนิมิตฺตุ\-ปฏฺฐาเน\-กตฺตํ, วยลกฺขณุ\-ปฏฺฐาเน\-กตฺตํ, นิโรธุ\-ปฏฺฐาเน\-กตฺตํ. ทาน\-โวสคฺคุ\-ปฏฺฐาเน\-กตฺตํ จาคา\-ธิมุตฺตานํ. สมถนิมิตฺตุ\-ปฏฺฐาเน\-กตฺตญฺจ อธิจิตฺตม\-นุยุตฺตานํ. วยลกฺขณุ\-ปฏฺฐาเน\-กตฺตญฺจ วิปสฺสกานํ. นิโรธุ\-ปฏฺฐาเน\-กตฺตญฺจ อริย\-ปุคฺคลานํ. อิเมหิ จตูหิ ฐาเนหิ เอกตฺตคตํ จิตฺตํ ปฏิปทาวิสุทฺธิ\-ปกฺขนฺท\-ญฺเจว โหติ, อุเปกฺขานุพฺรูหิตญฺจ ญาเณน จ สมฺปหํสิตํ.}

\prose{ปฐมสฺส ฌานสฺส โก อาทิ, กิํ มชฺเฌ, กิํ ปริโยสานํ? ปฐมสฺส ฌานสฺส ปฏิปทา\-วิสุทฺธิ อาทิ, อุเปกฺขา\-นุพฺรูหนา มชฺเฌ, สมฺปหํสนา ปริโยสานํ.  ปฐมสฺส ฌานสฺส ปฏิปทา\-วิสุทฺธิ อาทิ. อาทิสฺส กติ ลกฺขณานิ? อาทิสฺส ตีณิ ลกฺขณานิ, โย ตสฺส ปริปนฺโถ, ตโต จิตฺตํ วิสุชฺฌติ, วิสุทฺธตฺตา จิตฺตํ มชฺฌิมํ สมถ\-นิมิตฺตํ ปฏิปชฺชติ, ปฏิปนฺนตฺตา ตตฺถ จิตฺตํ ปกฺขนฺทติ, ยญฺจ ปริปนฺถโต จิตฺตํ วิสุชฺฌติ, ยญฺจ วิสุทฺธตฺตา จิตฺตํ มชฺฌิมํ สมถ\-นิมิตฺตํ ปฏิปชฺชติ, ยญฺจ ปฏิปนฺนตฺตา ตตฺถ จิตฺตํ ปกฺขนฺทติ, ปฐมสฺส ฌานสฺส ปฏิปทา\-วิสุทฺธิ อาทิ, อาทิสฺส อิมานิ ตีณิ ลกฺขณานิ, เตน วุจฺจติ ปฐมํ ฌานํ อาทิกลฺยาณ\-ญฺเจว โหติ ลกฺขณสมฺปนฺนญฺจ.}

\prose{ปฐมสฺส ฌานสฺส อุเปกฺขา\-นุพฺรูหนา มชฺเฌ. มชฺฌสฺส กติ ลกฺขณานิ? มชฺฌสฺส ตีณิ ลกฺขณานิ, วิสุทฺธํ จิตฺตํ อชฺฌุเปกฺขติ, สมถ\-ปฏิปนฺนํ อชฺฌุเปกฺขติ, เอกตฺตุ\-ปฏฺฐานํ อชฺฌุเปกฺขติ, ยญฺจ วิสุทฺธํ จิตฺตํ อชฺฌุเปกฺขติ, ยญฺจ สมถ\-ปฏิปนฺนํ อชฺฌุเปกฺขติ, ยญฺจ เอกตฺตุ\-ปฏฺฐานํ อชฺฌุเปกฺขติ, ปฐมสฺส ฌานสฺส อุเปกฺขา\-นุพฺรูหนา มชฺเฌ, มชฺฌสฺส อิมานิ ตีณิ ลกฺขณานิ, เตน วุจฺจติ ปฐมํ ฌานํ มชฺเฌกลฺยาณ\-ญฺเจว โหติ ลกฺขณสมฺปนฺนญฺจ.}

\prose{ปฐมสฺส ฌานสฺส สมฺปหํสนา ปริโยสานํ. ปริโยสานสฺส กติ ลกฺขณานิ? ปริโยสานสฺส จตฺตาริ ลกฺขณานิ, ตตฺถ ชาตานํ ธมฺมานํ อนติวตฺตน\-ฏฺเฐน สมฺปหํสนา, อินฺทฺริยานํ เอกรสฏฺเฐน สมฺปหํสนา, ตทุปค\-วีริย\-วาห\-นฏฺเฐน สมฺปหํสนา, อาเสวนฏฺเฐน สมฺปหํสนา, ปฐมสฺส ฌานสฺส สมฺปหํสนา ปริโยสานํ, ปริโยสานสฺส อิมานิ จตฺตาริ ลกฺขณานิ, เตน วุจฺจติ ปฐมํ ฌานํ ปริโยสานกลฺยาณ\-ญฺเจว โหติ ลกฺขณสมฺปนฺนญฺจ, เอวํ ติวตฺตคตํ จิตฺตํ ติวิธ\-กลฺยาณกํ ทสลกฺขณสมฺปนฺนํ วิตกฺก\-สมฺปนฺน\-ญฺเจว โหติ วิจาร\-สมฺปนฺนญฺจ ปีติสมฺปนฺนญฺจ สุขสมฺปนฺนญฺจ จิตฺตสฺส อธิฏฺฐาน\-สมฺปนฺนญฺจ สทฺธาสมฺปนฺนญฺจ วีริย\-สมฺปนฺนญฺจ สติสมฺปนฺนญฺจ สมาธิสมฺปนฺนญฺจ ปญฺญาสมฺปนฺนญฺจ.}

\prose{\csromanfolio{168}ทุติยสฺส ฌานสฺส โก อาทิ, กิํ มชฺเฌ, กิํ ปริโยสานํ? ทุติยสฺส ฌานสฺส ปฏิปทา\-วิสุทฺธิ อาทิ, อุเปกฺขา\-นุพฺรูหนา มชฺเฌ, สมฺปหํสนา ปริโยสานํ ฯเปฯ เอวํ ติวตฺตคตํ จิตฺตํ ติวิธ\-กลฺยาณกํ ทสลกฺขณสมฺปนฺนํ ปีติสมฺปนฺน\-ญฺเจว โหติ สุขสมฺปนฺนญฺจ จิตฺตสฺส อธิฏฺฐาน\-สมฺปนฺนญฺจ สทฺธาสมฺปนฺนญฺจ วีริย\-สมฺปนฺนญฺจ สติสมฺปนฺนญฺจ สมาธิสมฺปนฺนญฺจ ปญฺญาสมฺปนฺนญฺจ.}

\prose{ตติยสฺส ฌานสฺส โก อาทิ, กิํ มชฺเฌ, กิํ ปริโยสานํ ฯเปฯ เอวํ ติวตฺตคตํ จิตฺตํ ติวิธ\-กลฺยาณกํ ทสลกฺขณสมฺปนฺนํ สุขสมฺปนฺน\-ญฺเจว โหติ จิตฺตสฺส อธิฏฺฐาน\-สมฺปนฺนญฺจ สทฺธาสมฺปนฺนญฺจ วีริย\-สมฺปนฺนญฺจ สติสมฺปนฺนญฺจ สมาธิสมฺปนฺนญฺจ ปญฺญาสมฺปนฺนญฺจ.}

\prose{จตุตฺถสฺส ฌานสฺส โก อาทิ, กิํ มชฺเฌ, กิํ ปริโยสานํ ฯเปฯ เอวํ ติวตฺตคตํ จิตฺตํ ติวิธ\-กลฺยาณกํ ทสลกฺขณสมฺปนฺนญฺจ อุเปกฺขาสมฺปนญฺเจว โหติ จิตฺตสฺส อธิฏฺฐาน\-สมฺปนฺนญฺจ สทฺธาสมฺปนฺนญฺจ วีริย\-สมฺปนฺนญฺจ สติสมฺปนฺนญฺจ สมาธิสมฺปนฺนญฺจ ปญฺญาสมฺปนฺนญฺจ.}

\prose{อากาสา\-นญฺจา\-ยตนสมาปตฺติยา ฯเปฯ. วิญฺญาณญฺจายตน\-สมาปตฺติยา. อากิญฺจญฺญายตน\-สมาปตฺติยา. เนว\-สญฺญา\-นา\-สญฺญา\-ยตนสมาปตฺติยา โก อาทิ, กิํ มชฺเฌ, กิํปริโยสานํ ฯเปฯ เอวํ ติวตฺตคตํ จิตฺตํ ติวิธ\-กลฺยาณกํ ทสลกฺขณสมฺปนฺนํ อุเปกฺขา\-สมฺปนฺน\-ญฺเจว โหติ จิตฺตสฺส อธิฏฺฐาน\-สมฺปนฺนญฺจ ฯเปฯ ปญฺญาสมฺปนฺนญฺจ.}

\prose{อนิจฺจา\-นุปสฺสนาย โก อาทิ, กิํ มชฺเฌ, กิํ ปริโยสานํ ฯเปฯ เอวํ ติวตฺตคตํ จิตฺตํ ติวิธ\-กลฺยาณกํ ทสลกฺขณสมฺปนฺนํ วิตกฺก\-สมฺปนฺน\-ญฺเจว โหติ วิจาร\-สมฺปนฺนญฺจ ปีติสมฺปนฺนญฺจ สุขสมฺปนฺนญฺจ จิตฺตสฺส อธิฏฺฐาน\-สมฺปนฺนญฺจ สทฺธาสมฺปนฺนญฺจ วีริย\-สมฺปนฺนญฺจ สติสมฺปนฺนญฺจ สมาธิสมฺปนฺนญฺจ ปญฺญาสมฺปนฺนญฺจ.  ทุกฺขา\-นุปสฺสนาย ฯเปฯ. อนตฺตา\-นุปสฺสนาย. นิพฺพิทา\-นุปสฺสนาย. วิราคา\-นุปสฺสนาย. นิโรธา\-นุปสฺสนาย. ปฏินิสฺสคฺคา\-นุปสฺสนาย. ขยานุปสฺสนาย. วยา\-นุปสฺสนาย. วิปริณามา\-นุปสฺสนาย. อนิมิตฺตา\-นุปสฺสนาย. อปฺปณิหิตา\-นุปสฺสนาย. สุญฺญตา\-นุปสฺสนาย. อธิปญฺญา\-ธมฺม\-วิปสฺสนาย. ยถา\-ภูต\-ญาณ\-ทสฺสนาย. อาทีนวา\-นุปสฺสนาย. ปฏิสงฺขา\-นุปสฺสนาย. วิวฏฺฏนา\-นุปสฺสนาย.}

\prose{โสตาปตฺติ\-มคฺคสฺส. สกทาคามิ\-มคฺคสฺส. อนาคามิ\-มคฺคสฺส. อรหตฺต\-มคฺคสฺส โก อาทิ, กิํ มชฺเฌ, กิํ ปริโยสานํ? อรหตฺต\-มคฺคสฺส \csromanfolio{169}ปฏิปทา\-วิสุทฺธิ อาทิ, อุเปกฺขา\-นุพฺรูหนา มชฺเฌ, สมฺปหํสนา ปริโยสานํ. อรหตฺต\-มคฺคสฺส ปฏิปทา\-วิสุทฺธิ อาทิ. อาทิสฺส กติ ลกฺขณานิ? อาทิสฺส ตีณิ ลกฺขณานิ, โย ตสฺส ปริปนฺโถ, ตโต จิตฺตํ วิสุชฺฌติ, วิสุทฺธตฺตา จิตฺตํ มชฺฌิมํ สมถ\-นิมิตฺตํ ปฏิปชฺชติ. ปฏิปนฺนตฺตา ตตฺถ จิตฺตํ ปกฺขนฺทติ, ยญฺจ ปริปนฺถโต จิตฺตํ วิสุชฺฌติ, ยญฺจ วิสุทฺธตฺตา จิตฺตํ มชฺฌิมํ สมถ\-นิมิตฺตํ ปฏิปชฺชติ, ยญฺจ ปฏิปนฺนตฺตา ตตฺถ จิตฺตํ ปกฺขนฺทติ, อรหตฺต\-มคฺคสฺส ปฏิปทา\-วิสุทฺธิ อาทิ, อาทิสฺส อิมานิ ตีณิ ลกฺขณานิ. เตน วุจฺจติ อรหตฺตมคฺโค อาทิกลฺยาโณ เจว โหติ ลกฺขณสมฺปนฺโน จ.}

\proseitem{158}{อรหตฺต\-มคฺคสฺส อุเปกฺขา\-นุพฺรูหนา มชฺเฌ. มชฺฌสฺส กติ ลกฺขณานิ? มชฺฌสฺส ตีณิ ลกฺขณานิ, วิสุทฺธํ จิตฺตํ อชฺฌูเปกฺขติ. สมถ\-ปฏิปนฺนํ อชฺฌุเปกฺขติ, เอกตฺตุ\-ปฏฺฐานํ อชฺฌุเปกฺขติ, ยญฺจ วิสุทฺธํ จิตฺตํ อชฺฌุเปกฺขติ, ยญฺจ สมถ\-ปฏิปนฺนํ อชฺฌุเปกฺขติ, ยญฺจ เอกตฺตุ\-ปฏฺฐานํ อชฺฌุเปกฺขติ. เตน วุจฺจติ อรหตฺตมคฺโค มชฺเฌกลฺยาโณ เจว โหติ ลกฺขณสมฺปนฺโน จ.}

\prose{อรหตฺต\-มคฺคสฺส สมฺปหํสนา ปริโยสานํ. ปริโยสานสฺส กติ ลกฺขณานิ? ปริโยสานสฺส จตฺตาริ ลกฺขณานิ, ตตฺถ ชาตานํ ธมฺมานํ อนติวตฺตน\-ฏฺเฐน สมฺปหํสนา, อินฺทฺริยานํ เอกรสฏฺเฐน สมฺปหํสนา, ตทุปค\-วีริย\-วาห\-นฏฺเฐน สมฺปหํสนา, อาเสวนฏฺเฐน สมฺปหํสนา. อรหตฺต\-มคฺคสฺส สมฺปหํสนา ปริโยสานํ, ปริโยสานสฺส อิมานิ จตฺตาริ ลกฺขณานิ, เตน วุจฺจติ อรหตฺตมคฺโค ปริโยสาน\-กลฺยาโณ เจว โหติ ลกฺขณสมฺปนฺโน จ. เอวํ ติวตฺตคตํ จิตฺตํ ติวิธ\-กลฺยาณกํ ทสลกฺขณสมฺปนฺนํ วิตกฺก\-สมฺปนฺน\-ญฺเจว โหติ วิจาร\-สมฺปนฺนญฺจ ปีติสมฺปนฺนญฺจ สุขสมฺปนฺนญฺจ จิตฺตสฺส อธิฏฺฐาน\-สมฺปนฺนญฺจ สทฺธาสมฺปนฺนญฺจ วีริย\-สมฺปนฺนญฺจ สติสมฺปนฺนญฺจ สมาธิสมฺปนฺนญฺจ ปญฺญาสมฺปนฺนญฺจ.}

\setlength{\csromangathapagewidth}{0pt}
\setlength{\csromangathawakwidth}{0pt}
\csromangathameasuregroup{นิมิตฺตํ อสฺสาสปสฺสาสา, \\ อชานโต จ ตโย ธมฺเม, \\ นิมิตฺตํ อสฺสาสปสฺสาสา, \\ ชานโต จ ตโย ธมฺเม,}{\csromangathabat{นิมิตฺตํ อสฺสาสปสฺสาสา,}{อนารมฺมณเมกจิตฺตสฺส.} \\ \csromangathabat{อชานโต จ ตโย ธมฺเม,}{ภาวนา นุปลพฺภติ.} \\ \csromangathabat{นิมิตฺตํ อสฺสาสปสฺสาสา,}{อนารมฺมณเมกจิตฺตสฺส.} \\ \csromangathabat{ชานโต จ ตโย ธมฺเม,}{ภาวนา อุปลพฺภตีติ.}}
\csromangathasetleft{นิมิตฺตํ อสฺสาสปสฺสาสา, \\ อชานโต จ ตโย ธมฺเม, \\ นิมิตฺตํ อสฺสาสปสฺสาสา, \\ ชานโต จ ตโย ธมฺเม,}
\csromangathagroup{\csromangathastanza{\csromangathaitembat{159}{นิมิตฺตํ อสฺสาสปสฺสาสา,}{อนารมฺมณเม\-กจิตฺตสฺส.} \\ \csromangathacontbat{อชานโต จ ตโย ธมฺเม,}{ภาวนา นุปลพฺภติ.} \\ \csromangathacontbat{นิมิตฺตํ อสฺสาสปสฺสาสา,}{อนารมฺมณเม\-กจิตฺตสฺส.} \\ \csromangathacontbat{ชานโต จ ตโย ธมฺเม,}{ภาวนา อุปลพฺภตีติ.}}}

\prose{กถํ อิเม ตโย ธมฺมา เอกจิตฺตสฺส อารมฺมณา น โหนฺติ, น จิเม ตโย ธมฺมา อวิทิตา โหนฺติ, น จ จิตฺตํ วิกฺเขปํ คจฺฉติ, ปธานญฺจ ปญฺญายติ, ปโยคญฺจ สาเธติ, วิเสสม\-ธิคจฺฉติ? เสยฺยถาปิ \csromanfolio{170}รุกฺโข สเม ภูมิภาเค นิกฺขิตฺโต, ตเมนํ ปุริโส กกเจน ฉินฺเทยฺย, รุกฺเข ผุฏฺฐ\-กกจทนฺตานํ วเสน ปุริสสฺส สติ อุปฏฺฐิตา โหติ, น อาคเต วา คเต วา กกจทนฺเต มนสิ กโรติ, น อาคตา วา คตา วา กกจทนฺตา อวิทิตา โหนฺติ, ปธานญฺจ ปญฺญายติ, ปโยคญฺจ สาเธติ. ยถา รุกฺโข สเม ภูมิภาเค นิกฺขิตฺโต, เอวํ อุปนิพนฺธนา นิมิตฺตํ. ยถา กกจทนฺตา, เอวํ อสฺสาสปสฺสาสา. ยถา รุกฺเข ผุฏฺฐ\-กกจทนฺตานํ วเสน ปุริสสฺส สติ อุปฏฺฐิตา โหติ, น อาคเต วา คเต วา กกจทนฺเต มนสิ กโรติ, น อาคตา วา คตา วา กกจทนฺตา อวิทิตา โหนฺติ, ปธานญฺจ ปญฺญายติ, ปโยคญฺจ สาเธติ. เอวเมวํ ภิกฺขุ นาสิกคฺเค วา มุขนิมิตฺเต วา สติํ อุปฏฺฐเปตฺวา นิสินฺโน โหติ. น อาคเต วา คเต วา อสฺสาสปสฺสาเส มนสิ กโรติ, น อาคตา วา คตา วา อสฺสาสปสฺสาสา อวิทิตา โหนฺติ, ปธานญฺจ ปญฺญายติ, ปโยคญฺจ สาเธติ, วิเสสม\-ธิคจฺฉติ ปธานญฺจ.}

\prose{กตมํ ปธานํ? อารทฺธ\-วีริยสฺส กาโยปิ จิตฺตมฺปิ กมฺมนิยํ โหติ. อิทํ ปธานํ.  กตโม ปโยโค? อารทฺธ\-วีริยสฺส อุปกฺกิเลสา ปหียนฺติ, วิตกฺกา วูปสมฺมนฺติ. อยํ ปโยโค.  กตโม วิเสโส? อารทฺธ\-วีริยสฺส สญฺโญชนา ปหียนฺติ, อนุสยา พฺยนฺตีโหนฺติ\csromansharedfootnote{3df7ff4f6376e8e9}{อนุสยา พฺยาสนฺติ (สฺยา)}. อยํ วิเสโส. เอวํ อิเม ตโย ธมฺมา เอกจิตฺตสฺส อารมฺมณา น โหนฺติ, น จิเม ตโย ธมฺมา อวิทิตา โหนฺติ, น จ จิตฺตํ วิกฺเขปํ คจฺฉติ, ปธานญฺจ ปญฺญายติ, ปโยคญฺจ สาเธติ. วิเสสม\-ธิคจฺฉติ.}

\setlength{\csromangathapagewidth}{0pt}
\setlength{\csromangathawakwidth}{0pt}
\csromangathameasuregroup{อานาปานสฺสติ ยสฺส, \\ อนุปุพฺพํ ปริจิตา, \\ โส อิมํ โลกํ ปภาเสติ,}{\csromangathabat{อานาปานสฺสติ ยสฺส,}{ปริปุณฺณา สุภาวิตา.} \\ \csromangathabat{อนุปุพฺพํ ปริจิตา,}{ยถา พุทฺเธน เทสิตา.} \\ \csromangathabat{โส อิมํ โลกํ ปภาเสติ,}{อพฺภา มุตฺโตว จนฺทิมาติ.}}
\csromangathasetleft{อานาปานสฺสติ ยสฺส, \\ อนุปุพฺพํ ปริจิตา, \\ โส อิมํ โลกํ ปภาเสติ,}
\csromangathagroup{\csromangathastanza{\csromangathaitembat{160}{อานาปานสฺสติ ยสฺส,}{ปริปุณฺณา สุภาวิตา.} \\ \csromangathacontbat{อนุปุพฺพํ ปริจิตา,}{ยถา พุทฺเธน เทสิตา.} \\ \csromangathacontbat{โส อิมํ โลกํ ปภาเสติ,}{อพฺภา มุตฺโตว จนฺทิมาติ.}}}

\prose{\textbf{อานนฺ}ติ อสฺสาโส, โน ปสฺสาโส. \textbf{อาปานนฺ}ติ\csromansharedfootnote{1088bd4872155f19}{อปานนฺติ (ก) 69 ปิฏฺเฐ อฏฺฐกถา โอโลเกตพฺพา.} ปสฺสาโส, โน อสฺสาโส. อสฺสาสวเสน อุปฏฺฐานํ สติ, ปสฺสาสวเสน อุปฏฺฐานํ สติ.}

\prose{โย อสฺสสติ, ตสฺสุปฏฺฐาติ. โย ปสฺสสติ ตสฺสุปฏฺฐาติ. ปริปุณฺณาติ ปริคฺคห\-ฏฺเฐน ปริปุณฺณา, ปริวารฏฺเฐน ปริปุณฺณา, ปริปูรฏฺเฐน ปริปุณฺณา. สุภาวิตาติ จตสฺโส ภาวนา. ตตฺถ ชาตานํ ธมฺมานํ \csromanfolio{171}อนติวตฺตน\-ฏฺเฐน ภาวนา, อินฺทฺริยานํ เอกรสฏฺเฐน ภาวนา, ตทุปค\-วีริย\-วาห\-นฏฺเฐน ภาวนา. อาเสวนฏฺเฐน ภาวนา. ตสฺสิเม จตฺตาโร ภาวนฏฺฐา ยานีกตา โหนฺติ วตฺถุกตา อนุฏฺฐิตา ปริจิตา สุสมารทฺธา.}

\prose{ยานีกตาติ ยตฺถ ยตฺถ อากงฺขติ, ตตฺถ ตตฺถ วสิปฺปตฺโต โหติ พลปฺปตฺโต เวสารชฺช\-ปฺปตฺโต. ตสฺส เม เต ธมฺมา อาวชฺชนปฏิพทฺธา โหนฺติ อากงฺข\-ปฏิพทฺธา มนสิการ\-ปฏิพทฺธา จิตฺตุปฺปาท\-ปฏิพทฺธา. เตน วุจฺจติ “ยานีกตา”ติ.  วตฺถุกตาติ ยสฺมิํ ยสฺมิํ วตฺถุสฺมิํ จิตฺตํ สฺวาธิฏฺฐิตํ โหติ, ตสฺมิํ ตสฺมิํ วตฺถุสฺมิํ สติ สูปฏฺฐิตา\csromansharedfootnote{92ebfb123559fb6c}{สุปฏฺฐิตา (ก)} โหติ. ยสฺมิํ ยสฺมิํ วา ปน วตฺถุสฺมิํ สติ สูปฏฺฐิตา โหติ, ตสฺมิํ ตสฺมิํ วตฺถุสฺมิํ จิตฺตํ สฺวาธิฏฺฐิตํ โหติ. เตน วุจฺจติ “วตฺถุกตา”ติ.  อนุฏฺฐิตาติ วตฺถุสฺมิํ เยน เยน จิตฺตํ อภินีหรติ, เตน เตน สติ อนุปริวตฺตติ. เยน เยน วา ปน สติ อนุปริวตฺตติ, เตน เตน จิตฺตํ อภินีหรติ. เตน วุจฺจติ “อนุฏฺฐิตา”ติ.  ปริจิตาติ ปริคฺคห\-ฏฺเฐน ปริจิตา, ปริวารฏฺเฐน ปริจิตา, ปริปูรฏฺเฐน ปริจิตา, สติยา ปริคฺคณฺหนฺโต ชินาติ ปาปเก อกุสเล ธมฺเม. เตน วุจฺจติ “ปริจิตา”ติ.  สุสมารทฺธาติ จตฺตาโร สุสมารทฺธา. ตตฺถ ชาตานํ ธมฺมานํ อนติวตฺตน\-ฏฺเฐน สุสมารทฺธา, อินฺทฺริยานํ เอกรสฏฺเฐน สุสมารทฺธา, ตทุปค\-วีริย\-วาห\-นฏฺเฐน สุสมารทฺธา, ตปฺปจฺจนีกานํ กิเลสานํ สุสมูหตตฺตา\csromansharedfootnote{5faccb15cd15103d}{สุสมุหตตฺตา (ก)} สุสมารทฺธา.}

\proseitem{161}{\textbf{สุสมนฺ}ติ อตฺถิ สมํ อตฺถิ สุสมํ. กตมํ สมํ? เย ตตฺถ ชาตา อนวชฺชา กุสลา โพธิปกฺขิยา, อิทํ สมํ. กตมํ สุสมํ? ยํ เตสํ เตสํ ธมฺมานํ อารมฺมณํ นิโรโธ นิพฺพานํ, อิทํ สุสมํ. อิติ อิทญฺจ สมํ, อิทญฺจ สุสมํ ญาตํ โหติ ทิฏฺฐิ วิทิตํ สจฺฉิกตํ ผสฺสิตํ ปญฺญาย, อารทฺธํ โหติ วีริยํ อสลฺลีนํ, อุปฏฺฐิตา สติ อสมฺมุฏฺฐา\csromansharedfootnote{930e052c730284af}{อปมุฏฺฐา (สฺยา)}, ปสฺสทฺโธ กาโย อสารทฺโธ, สมาหิตํ จิตฺตํ เอกคฺคํ. เตน วุจฺจติ “สุสมารทฺธา”ติ.}

\prose{\textbf{อนุปุพฺพํ ปริจิตา}ติ ทีฆํ อสฺสาสวเสน ปุริมา ปุริมา ปริจิตา, ปจฺฉิมา ปจฺฉิมา อนุปริจิตา. ทีฆํ ปสฺสาสวเสน ปุริมา ปุริมา ปริจิตา, ปจฺฉิมา ปจฺฉิมา อนุปริจิตา. รสฺสํ อสฺสาสวเสน ปุริมา ปุริมา \csromanfolio{172}ปริจิตา, ปจฺฉิมา ปจฺฉิมา อนุปริจิตา. รสฺสํ ปสฺสาสวเสน ปุริมา ปุริมา ปริจิตา, ปจฺฉิมา ปจฺฉิมา อนุปริจิตา ฯเปฯ. ปฏินิสฺสคฺคา\-นุปสฺสี อสฺสาสวเสน ปุริมา ปุริมา ปริจิตา, ปจฺฉิมา ปจฺฉิมา อนุปริจิตา. ปฏินิสฺสคฺคา\-นุปสฺสี ปสฺสาสวเสน ปุริมา ปุริมา ปริจิตา, ปจฺฉิมา ปจฺฉิมา อนุปริจิตา. สพฺพาปิ โสฬสวตฺถุกา อานาปานสฺสติโย อญฺญมญฺญํ ปริจิตา เจว โหนฺติ อนุปริจิตา จ. เตน วุจฺจติ “อนุปุพฺพปริจิตา”ติ.}

\prose{ยถาติ ทส ยถตฺถา. อตฺต\-ทมถ\-ตฺโถ ยถตฺโถ, อตฺต\-สมถ\-ตฺโถ ยถตฺโถ, อตฺต\-ปรินิพฺพาปน\-ตฺโถ ยถตฺโถ, อภิญฺญตฺโถ ยถตฺโถ, ปริญฺญตฺโถ ยถตฺโถ, ปหานตฺโถ ยถตฺโถ, ภาวนตฺโถ ยถตฺโถ, สจฺฉิกิริย\-ตฺโถ ยถตฺโถ, สจฺจาภิสมย\-ตฺโถ ยถตฺโถ, นิโรเธ ปติฏฺฐาปก\-ตฺโถ ยถตฺโถ.}

\prose{\textbf{พุทฺโธ}ติ โย โส ภควา สยมฺภู อนาจริยโก ปุพฺเพ อนนุสฺสุเตสุ ธมฺเมสุ สามํ สจฺจานิ อภิสมฺพุชฺฌิ, ตตฺถ จ สพฺพญฺญุตํ ปาปุณิ, พเลสุ จ วสีภาวํ.}

\proseitem{162}{พุทฺโธติ เกนฏฺเฐน พุทฺโธ? พุชฺฌิตา สจฺจานีติ พุทฺโธ. โพเธตา ปชายาติ พุทฺโธ. สพฺพญฺญุตาย พุทฺโธ. สพฺพทสฺสาวิ\-ตาย พุทฺโธ. อนญฺญเนยฺย\-ตาย พุทฺโธ. วิสวิตาย\csromansharedfootnote{0161e45bfae910f8}{วิกติตาย (สฺยา)} พุทฺโธ. ขีณาสว\-สงฺขาเตน พุทฺโธ. นิรุปเลป\-สงฺขาเตน\csromansharedfootnote{cb7b5566b968f32b}{นิรุปกฺกิเลสสงฺขาเตน (สฺยา)} พุทฺโธ. เอกนฺตวีตราโคติ พุทฺโธ. เอกนฺตวีตโทโสติ พุทฺโธ. เอกนฺตวีตโมโหติ พุทฺโธ. เอกนฺตนิกฺกิเลโสติ พุทฺโธ. เอกายนมคฺคํ คโตติ พุทฺโธ. เอโก อนุตฺตรํ สมฺมาสมฺโพธิํ อภิสมฺพุทฺโธติ พุทฺโธ. อพุทฺธิ\-วิห\-ตตฺตา พุทฺธิ\-ปฏิลาภา พุทฺโธ. “พุทฺโธ”ติ เนตํ นามํ มาตรา กตํ, น ปิตรา กตํ, น ภาตรา กตํ, น ภคินิยา กตํ, น มิตฺตามจฺเจหิ กตํ, น ญาติสาโลหิเตหิ กตํ, น สมณ\-พฺราหฺมเณหิ กตํ, น เทวตาหิ กตํ, วิโมกฺข\-นฺติกเมตํ พุทฺธานํ ภควนฺตานํ โพธิยา มุเล สห สพฺพญฺญุต\-ญาณสฺส ปฏิลาภา สจฺฉิกา ปญฺญตฺติ, ยทิทํ พุทฺโธติ.  เทสิตาติ อตฺต\-ทมถ\-ตฺโถ ยถตฺโถ ยถา พุทฺเธน เทสิโต. อตฺต\-สมถ\-ตฺโถ ยถตฺโถ ยถา พุทฺเทน เทสิโต. อตฺต\-ปรินิพฺพาปน\-ตฺโถ ยถตฺโถ ยถา พุทฺเธน เทสิโต ฯเปฯ นิโรเธ ปติฏฺฐาปก\-ตฺโถ ยถตฺโถ ยถา พุทฺเธน เทสิโต.}

\prose{\csromanfolio{173}โสติ คหฏฺโฐ วา โหติ ปพฺพชิโต วา.  โลโกติ ขนฺธโลโก ธาตุโลโก อายตนโลโก วิปตฺติ\-ภว\-โลโก วิปตฺติ\-สมฺภว\-โลโก สมฺปตฺติ\-ภว\-โลโก สมฺปตฺติ\-สมฺภว\-โลโก. เอโก โลโก, สพฺเพ สตฺตา อาหารฏฺฐิติกา ฯเปฯ. อฏฺฐารส โลกา อฏฺฐารส ธาตุโย.  ปภาเสตีติ อตฺตทมตฺถตฺถํ ยถตฺถํ อภิสมฺพุทฺธตฺตา โส อิมํ โลกํ โอภาเสติ ภาเสติ ปภาเสติ. อตฺต\-สมถตฺถํ ยถตฺถํ อภิสมฺพุทฺธตฺตา โส อิมํ โลกํ โอภาเสติ ภาเสติ ปภาเสติ. อตฺต\-ปรินิพฺพาปนตฺถํ ยถตฺถํ อภิสมฺพุทฺธตฺตา โส อิมํ โลกํ โอภาเสติ ภาเสติ ปภาเสติ ฯเปฯ นิโรเธ ปติฏฺฐาปกตฺถํ ยถตฺถํ อภิสมฺพุทฺธตฺตา โส อิมํ โลกํ โอภาเสติ ภาเสติ ปภาเสติ.}

\prose{\textbf{อพฺภา มุตฺโตว จนฺทิมา}ติ ยถา อพฺภา, เอวํ กิเลสา. ยถา จนฺโท เอวํ อริยญาณํ. ยถา จนฺทิมา เทวปุตฺโต, เอวํ ภิกฺขุ. ยถา จนฺโท อพฺภา มุตฺโต มหิกา มุตฺโต ธูมรชา มุตฺโต ราหุคหณา\csromansharedfootnote{54fdb146f46f1075}{ราหุปาณา (สฺยา)} วิปฺปมุตฺโต ภาสเต จ ตปเต จ วิโรจเต\csromansharedfootnote{2390845d169c2e6a}{วิโรจติ (สฺยา)} จ. เอวเมวํ ภิกฺขุ สพฺพกิเลเสหิ วิปฺปมุตฺโต ภาสเต จ ตปเต จ วิโรจเต จ, เตน วุจฺจติ “อพฺภา มุตฺโตว จนฺทิมา”ติ. อิมานิ เตรส โวทาเน ญาณานิ.}

\vspace{\tipitakaparskip}
\csromancenter{โวทาน\-ญาณ\-นิทฺเทโส จตุตฺโถ.}
\nitthitamruled{ภาณวาโร.}
\csromanmatikaanchor{mtk.70}
\csromantocmarkref{subsection}{5. สโตการิญาณนิทฺเทส}{mtk.70}
\bookmark[dest={mtk.70},level=2]{5. สโตการิญาณนิทฺเทส}
\csromanheader{2}{5. \textbf{สโตการิญาณนิทฺเทส}}
\proseitem{163}{กตมานิ พาตฺติํส สโตการิสฺส ญาณานิ? อิธ ภิกฺขุ อรญฺญคโต วา รุกฺขมูลคโต วา สุญฺญาคารคโต วา นิสีทติ ปลฺลงฺกํ อาภุชิตฺวา อุชุํ กายํ ปณิธาย ปริมุขํ สติํ อุปฏฺฐเปตฺวา, โส สโตว อสฺสสติ, สโต ปสฺสสติ\csromandash{}ทีฆํ วา อสฺสสนฺโต “ทีฆํ อสฺสสามี”ติ ปชานาติ, ทีฆํ วา ปสฺสสนฺโต “ทีฆํ ปสฺสสามี”ติ ปชานาติ. รสฺสํ วา อสฺสสนฺโต “รสฺสํ อสฺสสามี”ติ ปชานาติ, รสฺสํ วา ปสฺสสนฺโต “รสฺสํ ปสฺสสามี”ติ ปชานาติ. “สพฺพ\-กาย\-ปฏิสํเวที อสฺสสิสฺสามี”ติ สิกฺขติ, “สพฺพ\-กาย\-ปฏิสํเวที ปสฺสสิสฺสามี”ติ \csromanfolio{174}สิกฺขติ. “ปสฺสมฺภยํ กายสงฺขารํ อสฺสสิสฺสามี”ติ สิกฺขติ, “ปสฺสมฺภยํ กายสงฺขารํ ปสฺสสิสฺสามี”ติ สิกฺขติ. “ปีติปฏิสํเวที ฯเปฯ. สุขปฏิสํเวที. จิตฺต\-สงฺขาร\-ปฏิสํเวที. ปสฺสมฺภยํ จิตฺตสงฺขารํ. จิตฺต\-ปฏิสํเวที. อภิปฺปโมทยํ จิตฺตํ. สมาทหํ จิตฺตํ. วิโมจยํ จิตฺตํ. อนิจฺจานุปสฺสี. วิราคานุปสฺสี. นิโรธานุปสฺสี. ปฏินิสฺสคฺคา\-นุปสฺสี อสฺสสิสฺสามี”ติ สิกฺขติ, “ปฏินิสฺสคฺคา\-นุปสฺสี ปสฺสสิสฺสามี”ติ สิกฺขติ.}

\proseitem{164}{อิธาติ อิมิสฺสา ทิฏฺฐิยา อิมิสฺสา ขนฺติยา อิมิสฺสา รุจิยา อิมสฺมิํ อาทาเย อิมสฺมิํ ธมฺเม อิมสฺมิํ วินเย อิมสฺมิํ ธมฺมวินเย อิมสฺมิํ ปาวจเน อิมสฺมิํ พฺรหฺมจริเย อิมสฺมิํ สตฺถุสาสเน, เตน วุจฺจติ “อิธา”ติ. พฺหิกฺขูติ ปุถุชฺชน\-กลฺยาณโก วา โหติ ภิกฺขุ เสกฺโข วา อรหา วา อกุปฺปธมฺโม. อรญฺญนฺติ นิกฺขมิตฺวา พหิ อินฺทขีลา, สพฺพเมตํ อรญฺญํ. รุกฺขมูลนฺติ ยตฺถ ภิกฺขุโน อาสนํ ปญฺญตฺตํ โหติ มญฺโจ วา ปีฐํ วา ภิสิ วา ตฏฺฏิกา วา จมฺมขณฺโฑ วา ติณสนฺถโร\csromansharedfootnote{ec5bf7c6ae4769a4}{ติณสณฺฐโร (สฺยา)} วา ปณฺณสนฺถโร วา ปลาลสนฺถโร\csromansharedfootnote{675fc3124b1a76b3}{ปลาสสณฺฐโร (สฺยา)} วา, ตตฺถ ภิกฺขุ จงฺกมติ วา ติฏฺฐติ วา นิสีทติ วา เสยฺยํ วา กปฺเปติ. สุญฺญนฺติ เกนจิ อนากิณฺณํ โหติ คหฏฺเฐหิ วา ปพฺพชิเตหิ วา. อคารนฺติ\csromansharedfootnote{c4e815c108689203}{อาคารนฺติ (สฺยา)} วิหาโร อฑฺฒโยโค ปาสาโท หมฺมิยํ คุหา. นิสีทติ ปลฺลงฺกํ อาภุชิตฺวาติ นิสินฺโน โหติ ปลฺลงฺกํ อาภุชิตฺวา. อุชุํ กายํ ปณิธายาติ อุชุโก โหติ กาโย ฐิโต\csromansharedfootnote{08ccda66835307fb}{ปนิธิโต (สฺยา)} สุปณิหิโต. ปริมุขํ สติํ อุปฏฺฐเปตฺวาติ ปรีติ ปริคฺคหฏฺโฐ. มฺอุขนฺติ นิยฺยานฏฺโฐ. สตีติ อุปฏฺฐานฏฺโฐ, เตน วุจฺจติ “ปริมุขํ สติํ อุปฏฺฐเปตฺวา”ติ.}

\proseitem{165}{\textbf{สโตว อสฺสสติ, สโต ปสฺสสตี}ติ พาตฺติํสาย อากาเรหิ สโต การี โหติ. ทีฆํ อสฺสาสวเสน จิตฺตสฺส เอกคฺคตํ อวิกฺเขปํ ปชานโต สติ อุปฏฺฐิตา โหติ, ตาย สติยา เตน ญาเณน สโต การี โหติ. ทีฆํ ปสฺสาสวเสน จิตฺตสฺส เอกคฺคตํ อวิกฺเขปํ ปชานโต สติ อุปฏฺฐิตา โหติ, ตาย สติยา เตน ญาเณน สโต การี โหติ. รสฺสํ อสฺสาสวเสน จิตฺตสฺส เอกคฺคตํ อวิกฺเขปํ ปชานโต สติ อุปฏฺฐิตา โหติ, ตาย สติยา เตน ญาเณน สโต การี โหติ. รสฺสํ ปสฺสาสวเสน จิตฺตสฺส \csromanfolio{175}เอกคฺคตํ อวิกฺเขปํ ปชานโต สติ อุปฏฺฐิตา โหติ, ตาย สติยา เตน ญาเณน สโต การี โหติ ฯเปฯ. ปฏินิสฺสคฺคา\-นุปสฺสี อสฺสาสวเสน. ปฏินิสฺสคฺคา\-นุปสฺสี ปสฺสาสวเสน จิตฺตสฺส เอกคฺคตํ อวิกฺเขปํ ปชานโต สติ อุปฏฺฐิตา โหติ, ตาย สติยา เตน ญาเณน สโต การี โหติ.}

\csromanheader{2}{ปฐม\-จตุกฺก\-นิทฺเทส}
\proseitem{166}{กถํ ทีฆํ อสฺสสนฺโต “ทีฆํ อสฺสสามี”ติ ปชานาติ, ทีฆํ ปสฺสสนฺโต “ทีฆํ ปสฺสสามี”ติ ปชานาติ? ทีฆํ อสฺสาสํ อทฺธาน\-สงฺขาเต อสฺสสติ, ทีฆํ ปสฺสาสํ อทฺธาน\-สงฺขาเต ปสฺสสติ, ทีฆํ อสฺสาสปสฺสาสํ อทฺธาน\-สงฺขาเต อสฺสสติปิ อสฺสสติปิ, ทีฆํ อสฺสาสปสฺสาสํ อทฺธาน\-สงฺขาเต อสฺสสโตปิ ปสฺสสโตปิ ฉนฺโท อุปฺปชฺชติ. ฉนฺทวเสน ตโต สุขุมตรํ ทีฆํ อสฺสาสํ อทฺธาน\-สงฺขาเต อสฺสสติ. ฉนฺทวเสน ตโต สุขุมตรํ ทีฆํ ปสฺสาสํ อทฺธาน\-สงฺขาเต ปสฺสสติ. ฉนฺทวเสน ตโต สุขุมตรํ ทีฆํ อสฺสาสปสฺสาสํ อทฺธาน\-สงฺขาเต อสฺสสติปิ ปสฺสสติปิ. ฉนฺทวเสน ตโต สุขุมตรํ ทีฆํ อสฺสาสปสฺสาสํ อทฺธาน\-สงฺขาเต อสฺสสโตปิ ปสฺสสโตปิ ปาโมชฺชํ อุปฺปชฺชติ. ปาโมชฺชวเสน ตโต สุขุมตรํ ทีฆํ อสฺสาสํ อทฺธาน\-สงฺขาเต อสฺสสติ. ปาโมชฺชวเสน ตโต สุขุมตรํ ทีฆํ ปสฺสาสํ อทฺธาน\-สงฺขาเต ปสฺสสติ. ปาโมชฺชวเสน ตโต สุขุมตรํ ทีฆํ อสฺสาสปสฺสาสํ อทฺธาน\-สงฺขาเต อสฺสสติปิ ปสฺสสติปิ. ปาโมชฺชวเสน ตโต สุขุมตรํ ทีฆํ อสฺสาสปสฺสาสํ อทฺธาน\-สงฺขาเต อสฺสสโตปิ ปสฺสสโตปิ ทีฆํ อสฺสาสปสฺสาสาปิ จิตฺตํ วิวตฺตติ, อุเปกฺขา สณฺฐาติ. อิเมหิ นวหากาเรหิ ทีฆํ อสฺสาสปสฺสาสา กาโย, อุปฏฺฐานํ สติ, อนุปสฺสนา ญาณํ. กาโย อุปฏฺฐานํ, โน สติ. สติ อุปฏฺฐานญฺเจว สติ จ. ตาย สติยา เตน ญาเณน ตํ กายํ อนุปสฺสติ, เตน วุจฺจติ “กาเย กายานุปสฺสนา\-สติปฏฺฐาน\-ภาวนา”ติ.}

\proseitem{167}{อ\textbf{นุปสฺสตี}ติ กถํ ตํ กายํ อนุปสฺสติ? อนิจฺจโต อนุปสฺสติ, โน นิจฺจโต. ทุกฺขโต อนุปสฺสติ, โน สุขโต. อนตฺตโต อนุปสฺสติ, โน อตฺตโต. นิพฺพินฺทติ, โน นนฺทติ. วิรชฺชติ, โน รชฺชติ. นิโรเธติ, โน สมุเทติ. ปฏินิสฺสชฺชติ, โน อาทิยติ. \csromanfolio{176}อนิจฺจโต อนุปสฺสนฺโต นิจฺจสญฺญํ ปชหติ. ทุกฺขโต อนุปสฺสนฺโต สุขสญฺญํ ปชหติ. อนตฺตโต อนุปสฺสนฺโต อตฺตสญฺญํ ปชหติ. นิพฺพินฺทนฺโต นนฺทิํ ปชหติ. วิรชฺชนฺโต ราคํ ปชหติ. นิโรเธนฺโต สมุทยํ ปชหติ. ปฏินิสฺสชฺชนฺโต อาทานํ ปชหติ. เอวํ ตํ กายํ อนุปสฺสติ.}

\prose{ภาวนาติ จตสฺโส ภาวนา. ตตฺถ ชาตานํ ธมฺมานํ อนติวตฺตน\-ฏฺเฐน ภาวนา, อินฺทฺริยานํ เอกรสฏฺเฐน ภาวนา, ตทุปค\-วีริย\-วาห\-นฏฺเฐน ภาวนา, อาเสวนฏฺเฐน ภาวนา, ทีฆํ อสฺสาสปสฺสาส\-วเสน จิตฺตสฺส เอกคฺคตํ อวิกฺเขปํ ปชานโต วิทิตา เวทนา อุปฺปชฺชนฺติ, วิทิตา อุปฏฺฐหนฺติ, วิทิตา อพฺภตฺถํ คจฺฉนฺติ, วิทิตา สญฺญา อุปฺปชฺชนฺติ, วิทิตา อุปฏฺฐหนฺติ, วิทิตา อพฺภตฺถํ คจฺฉนฺติ, วิทิตา วิตกฺกา อุปฺปชฺชนฺติ, วิทิตา อุปฏฺฐนฺติ, วิทิตา อพฺภตฺถํ คจฺฉนฺติ.}

\prose{กถํ วิทิตา เวทนา อุปฺปชฺชนฺติ, วิทิตา อุปฏฺฐหนฺติ, วิทิตา อพฺภตฺถํ คจฺฉนฺติ? กถํ เวทนาย อุปฺปาโท วิทิโต โหติ? “อวิชฺชาสมุทยา เวทนาสมุทโย”ติ ปจฺจยสมุทย\-ฏฺเฐน เวทนาย อุปฺปาโท วิทิโต โหติ. “ตณฺหาสมุทยา เวทนาสมุทโย”ติ. “กมฺมสมุทยา เวทนาสมุทโย”ติ. “ผสฺสสมุทยา เวทนาสมุทโย”ติ ปจฺจยสมุทย\-ฏฺเฐน เวทนาย อุปฺปาโท วิทิโต โหติ. นิพฺพตฺติ\-ลกฺขณํ ปสฺสโตปิ เวทนาย อุปฺปาโท วิทิโต โหติ. เอวํ เวทนาย อุปฺปาโท วิทิโต โหติ.}

\prose{กถํ เวทนาย อุปฏฺฐานํ วิทิตํ โหติ? อนิจฺจโต มนสิกโรโต ขยตุ\-ปฏฺฐานํ วิทิตํ โหติ, ทุกฺขโต มนสิกโรโต ภยตุปฏฺฐานํ วิทิตํ โหติ, อนตฺตโต มนสิกโรโต สุญฺญตุ\-ปฏฺฐานํ วิทิตํ โหติ. เอวํ เวทนาย อุปฏฺฐานํ วิทิตํ โหติ.}

\prose{กถํ เวทนาย อตฺถงฺคโม วิทิโต โหติ? “อวิชฺชานิโรธา เวทนานิโรโธ”ติ ปจฺจยนิโรธ\-ฏฺเฐน เวทนาย อตฺถงฺคโม วิทิโต โหติ, “ตณฺหานิโรธา เวทนานิโรโธ”ติ ฯเปฯ “กมฺมนิโรธา เวทนานิโรโธ”ติ ฯเปฯ “ผสฺสนิโรธา เวทนานิโรโธ”ติ ปจฺจยนิโรธ\-ฏฺเฐน เวทนาย อตฺถงฺคโม วิทิโต โหติ, วิปริณาม\-ลกฺขณํ ปสฺสโตปิ เวทนาย อตฺถงฺคโม วิทิโต โหติ. เอวํ เวทนาย อตฺถงฺคโม วิทิโต โหติ. เอวํ วิทิตา เวทนา อุปฺปชฺชนฺติ, วิทิตา อุปฏฺฐหนฺติ, วิทิตา อพฺภตฺถํ คจฺฉนฺติ.}

\prose{\csromanfolio{177}กถํ วิทิตา สญฺญา อุปฺปชฺชนฺติ, วิทิตา อุปฏฺฐหนฺติ, วิทิตา อพฺภตฺถํ คจฺฉนฺติ, กถํ สญฺญาย อุปฺปาโท วิทิโต โหติ? “อวิชฺชาสมุทยา สญฺญาสมุทโย”ติ ปจฺจยสมุทย\-ฏฺเฐน สญฺญาย อุปฺปาโท วิทิโต โหติ, “ตณฺหาสมุทยา สญฺญาสมุทโย”ติ ฯเปฯ “กมฺมสมุทยา สญฺญาสมุทโย”ติ ฯเปฯ “ผสฺสสมุทยา สญฺญาสมุทโย”ติ ปจฺจยสมุทย\-ฏฺเฐน สญฺญาย อุปฺปาโท วิทิโต โหติ, นิพฺพตฺติ\-ลกฺขณํ ปสฺสโตปิ สญฺญาย อุปฺปโท วิทิโต โหติ. เอวํ สญฺญาย อุปฺปาโท วิทิโต โหติ.}

\prose{กถํ สญฺญาย อุปฏฺฐานํ วิทิตํ โหติ? อนิจฺจโต มนสิกโรโต ขยตุ\-ปฏฺฐานํ วิทิตํ โหติ, ทุกฺขโต มนสิกโรโต ภยตุปฏฺฐานํ วิทิตํ โหติ, อนตฺตโต มนสิกโรโต สุญฺญาตุปฏฺฐานํ วิทิตํ โหติ. เอวํ สญฺญาย อุปฏฺฐานํ วิทิตํ โหติ.}

\prose{กถํ สญฺญาย อตฺถงฺคโม วิทิโต โหติ? “อวิชฺชานิโรธา สญฺญานิโรโธ”ติ ปจฺจยนิโรธ\-ฏฺเฐน สญฺญาย อตฺถงฺคโม วิทิโต โหติ, “ตณฺหานิโรธา สญฺญานิโรโธ”ติ ฯเปฯ “กมฺมนิโรธา สญฺญานิโรโธ”ติ ฯเปฯ “ผสฺสนิโรธา สญฺญานิโรโธ”ติ ปจฺจยนิโรธ\-ฏฺเฐน สญฺญาย อตฺถงฺคโม วิทิโต โหติ, วิปริณาม\-ลกฺขณํ ปสฺสโตปิ สญฺญาย อตฺถงฺคโม วิทิโต โหติ. เอวํ สญฺญาย อตฺถงฺคโม วิทิโต โหติ. เอวํ วิทิตา สญฺญา อุปฺปชฺชนฺติ, วิทิตา อุปฏฺฐหนฺติ, วิทิตา อพฺภตฺถํ คจฺฉนฺติ.}

\prose{กถํ วิทิตา วิตกฺกา อุปฺปชฺชนฺติ, วิทิตา อุปฏฺฐหนฺติ, วิทิตา อพฺภตฺถํ คจฺฉนฺติ, กถํ วิตกฺกานํ อุปฺปาโท วิทิโต โหติ? “อวิชฺชาสมุทยา วิตกฺก\-สมุทโย”ติ ปจฺจยสมุทย\-ฏฺเฐน วิตกฺกานํ อุปฺปาโท วิทิโต โหติ, “ตณฺหาสมุทยา วิตกฺก\-สมุทโย”ติ ฯเปฯ “กมฺมสมุทยา วิตกฺก\-สมุทโย”ติ ฯเปฯ “สญฺญาสมุทยา วิตกฺก\-สมุทโย”ติ ปจฺจยสมุทย\-ฏฺเฐน วิตกฺกานํ อุปฺปาโท วิทิโต โหติ, นิพฺพตฺติ\-ลกฺขณํ ปสฺสโตปิ วิตกฺกานํ อุปฺปาโท วิทิโต โหติ. เอวํ วิตกฺกานํ อุปฺปาโท วิทิโต โหติ.}

\prose{กถํ วิตกฺกานํ อุปฏฺฐานํ วิทิตํ โหติ? อนิจฺจโต มนสิกโรโต ขยตุ\-ปฏฺฐานํ วิทิตํ โหติ, ทุกฺขโต มนสิกโรโต ภยตุปฏฺฐานํ วิทิตํ โหติ, อนตฺตโต มนสิกโรโต สุญฺญตุ\-ปฏฺฐานํ วิทิตํ โหติ. เอวํ วิตกฺกานํ อุปฏฺฐานํ วิทิตํ โหติ.}

\prose{\csromanfolio{178}กถํ วิตกฺกานํ อตฺถงฺคโม วิทิโต โหติ? “อวิชฺชานิโรธา วิตกฺกนิโรโธ”ติ ปจฺจยนิโรธ\-ฏฺเฐน วิตกฺกานํ อตฺถงฺคโม วิทิโต โหติ, “ตณฺหานิโรธา วิตกฺกนิโรโธ”ติ ฯเปฯ “กมฺมนิโรธา วิตกฺกนิโรโธ”ติ ฯเปฯ “สญฺญานิโรธา วิตกฺกนิโรโธ”ติ ปจฺจยนิโรธ\-ฏฺเฐน วิตกฺกานํ อตฺถงฺคโม วิทิโต โหติ, วิปริณาม\-ลกฺขณํ ปสฺสโตปิ วิตกฺกานํ อตฺถงฺคโม วิทิโต โหติ. เอวํ วิตกฺกานํ อตฺถงฺคโม วิทิโต โหติ. เอวํ วิทิตา วิตกฺกา อุปฺปชฺชนฺติ, วิทิตา อุปฏฺฐหนฺติ, วิทิตา อพฺภตฺถํ คจฺฉนฺติ.}

\proseitem{168}{ทีฆํ อสฺสาสปสฺสาส\-วเสน จิตฺตสฺส เอกคฺคตํ อวิกฺเขปํ ปชานนฺโต อินฺทฺริยานิ สโมธาเนติ, โคจรญฺจ ปชานาติ, สมตฺถญฺจ ปฏิวิชฺฌติ ฯเปฯ มคฺคํ สโมธาเนติ, ธมฺเม สโมธาเนติ, โคจรญฺจ ปชานาติ, สมตฺตญฺจ ปฏิวิชฺฌติ.}

\prose{อินฺทฺริยานิ สโมธาเนตีติ กถํ อินฺทฺริยานิ สโมธาเนติ? อธิโมกฺข\-ฏฺเฐน สทฺธินฺทฺริยํ สโมธาเนติ, ปคฺคหฏฺเฐน วีริยินฺทฺริยํ สโมธาเนติ, อุปฏฺฐาน\-ฏฺเฐน สตินฺทฺริยํ สโมธาเนติ, อวิกฺเขป\-ฏฺเฐน สมาธินฺทฺริยํ สโมธาเนติ, ทสฺสนฏฺเฐน ปญฺญินฺทฺริยํ สโมธาเนติ. อยํ ปุคฺคโล อิมานิ อินฺทฺริยานิ อิมสฺมิํ อารมฺมเณ สโมธาเนติ. เตน วุจฺจติ “อินฺทฺริยานิ สโมธาเนตี”ติ.}

\prose{\textbf{โคจรญฺจ ปชานาตี}ติ ยํ ตสฺส อารมฺมณํ, ตํ ตสฺส โคจรํ. ยํ ตสฺส โคจรํ, ตํ ตสฺส อารมฺมณํ. \textbf{ปชานาตี}ติ ปุคฺคโล ปชานนา ปญฺญา.}

\prose{\textbf{สมนฺ}ติ อารมฺมณสฺส อุปฏฺฐานํ สมํ, จิตฺตสฺส อวิกฺเขโป สมํ, จิตฺตสฺส อธิฏฺฐานํ สมํ, จิตฺตสฺส โวทานํ สมํ. \textbf{อตฺโถ}ติ อนวชฺชฏฺโฐ นิเกฺลสฏฺโฐ โวทานฏฺโฐ ปรมฏฺโฐ. \textbf{ปฏิวิชฺฌตี}ติ อารมฺมณสฺส อุปฏฺฐานฏฺฐํ ปฏิวิชฺฌติ, จิตฺตสฺส อวิกฺเขปฏฺฐํ ปฏิวิชฺฌติ, จิตฺตสฺส อธิฏฺฐานฏฺฐํ ปฏิวิชฺฌติ, จิตฺตสฺส โวทานฏฺฐํ ปฏิวิชฺฌติ, เตน วุจฺจติ “สมตฺถญฺจ ปฏิวิชฺฌตี”ติ.}

\prose{\textbf{พลานิ สโมธาเนตี}ติ กถํ พลานิ สโมธาเนติ? อสฺสทฺธิเย อกมฺปิยฏฺเฐน สทฺธาพลํ สโมธาเนติ, โกสชฺเช อกมฺปิยฏฺเฐน วิริยพลํ สโมธาเนติ, ปมาเท อกมฺปิยฏฺเฐน สติพลํ สโมธาเนติ, อุทฺธจฺเจ อกมฺปิยฏฺเฐน สมาธิพลํ สโมธาเนติ, อวิชฺชาย อกมฺปิยฏฺเฐน ปญฺญาพลํ สโมธาเนติ. อยํ ปุคฺคโล อิมานิ พลานิ อิมสฺมิํ อารมฺมเณ \csromanfolio{179}สโมธาเนติ, เตน วุจฺจติ “พลานิ สโมธาเนตี”ติ.  โคจรญฺจ ปชานาตีติ ฯเปฯ. เตน วุจฺจติ “สมตฺถญฺจ ปฏิวิชฺฌตี”ติ.}

\prose{โพชฺฌงฺเค สโมธาเนตีติ กถํ โพชฺฌงฺค สโมธาเนติ? อุปฏฺฐาน\-ฏฺเฐน สติสมฺโพชฺฌงฺค สโมธาเนติ, ปวิจยฏฺเฐน ธมฺมวิจย\-สมฺโพชฺฌงฺคํ สโมธาเนติ, ปคฺคหฏฺเฐน วีริย\-สมฺโพชฺฌงฺคํ สโมธาเนติ, ผรณฏฺเฐน ปีติสมฺโพชฺฌงฺคํ สโมธาเนติ, อุปสมฏฺเฐน ปสฺสทฺธิ\-สมฺโพชฺฌงฺคํ สโมธาเนติ, อวิกฺเขป\-ฏฺเฐน สมาธิ\-สมฺโพชฺฌงฺค สโมธาเนติ, ปฏิสงฺขา\-นฏฺเฐน อุเปกฺขา\-สมฺโพชฺฌงฺคํ สโมธาเนติ. อยํ ปุคฺคโล อิเม โพชฺฌงฺเค อิมสฺมิํ อารมฺมเณ สโมธาเนติ. เตน วุจฺจติ “โพชฺฌงฺเค สโมธาเนตี”ติ.  โคจรญฺจ ปชานาตีติ ฯเปฯ. เตน วุจฺจติ “สมตฺถญฺจ ปฏิวิชฺฌตี”ติ.}

\prose{มคฺคํ สโมธาเนตีติ กถํ มคฺคํ สโมธาเนติ? ทสฺสนฏฺเฐน สมฺมาทิฏฺฐิํ สโมธาเนติ, อภินิโรปน\-ฏฺเฐน สมฺมาสงฺกปฺปํ สโมธาเนติ, ปริคฺคห\-ฏฺเฐน สมฺมาวาจํ สโมธาเนติ, สมุฏฺฐาน\-ฏฺเฐน สมฺมากมฺมนฺตํ สโมธาเนติ, โวทานฏฺเฐน สมฺมาอาชีวํ สโมธาเนติ, ปคฺคหฏฺเฐน สมฺมาวายามํ สโมธาเนติ, อุปฏฺฐาน\-ฏฺเฐน สมฺมาสติํ สโมธาเนติ, อวิกฺเขป\-ฏฺเฐน สมฺมาสมาธิํ สโมธาเนติ. อยํ ปุคฺคโล อิมํ มคฺคํ อิมสฺมิํ อารมฺมเณ สโมธาเนติ, เตน วุจฺจติ “มคฺคํ สโมธาเนตี”ติ.  โคจรญฺจ ปชานาตีติ ฯเปฯ. เตน วุจฺจติ “สมตฺถญฺจ ปฏิวิชฺฌตี”ติ.}

\prose{ธมฺเม สโมธาเนตีติ กถํ ธมฺเม สโมธาเนติ? อาธิปเตยฺย\-ฏฺเฐน อินฺทฺริยานิ สโมธาเนติ, อกมฺปิยฏฺเฐน พลานิ สโมธาเนติ, นิยฺยานฏฺเฐน โพชฺฌงฺเค สโมธาเนติ, เหตุฏฺเฐน มคฺคํ สโมธาเนติ, อุปฏฺฐาน\-ฏฺเฐน สติปฏฺฐานํ สโมธาเนติ, ปทหนฏฺเฐน สมฺมปฺปธานํ สโมธาเนติ, อิชฺฌนฏฺเฐน อิทฺธิปาทํ สโมธาเนติ, ตถฏฺเฐน สจฺจํ สโมธาเนติ, อวิกฺเขป\-ฏฺเฐน สมถํ สโมธาเนติ, อนุปสฺสนฏฺเฐน วิปสฺสนํ สโมธาเนติ, เอกรสฏฺเฐน สมถ\-วิปสฺสนํ สโมธาเนติ, อนติวตฺตน\-ฏฺเฐน ยุคนทฺธํ สโมธาเนติ, สํวรฏฺเฐน สีลวิสุทฺธิํ สโมธาเนติ, อวิกฺเขป\-ฏฺเฐน จิตฺตวิสุทฺธิํ สโมธาเนติ, ทสฺสนฏฺเฐน ทิฏฺฐิ\-วิสุทฺธิํ สโมธาเนติ, วิมุตฺตฏฺเฐน วิโมกฺขํ สโมธาเนติ, ปฏิเวธ\-ฏฺเฐน วิชฺชํ สโมธาเนติ, ปริจฺจาค\-ฏฺเฐน วิมุตฺติํ สโมธาเนติ, สมุจฺเฉท\-ฏฺเฐน ขเย ญาณํ สโมธาเนติ, ปฏิปฺปสฺสทฺธ\-ฏฺเฐน อนุปฺปาเท ญาณํ สโมธาเนติ, \csromanfolio{180}ฉนฺทํ มูลฏฺเฐน สโมธาเนติ, มนสิการํ สมุฏฺฐาน\-ฏฺเฐน สโมธาเนติ, ผสฺสํ สโมธาน\-ฏฺเฐน สโมธาเนติ, เวทนํ สโมสรณ\-ฏฺเฐน สโมธาเนติ, สมาธิํ ปมุขฏฺเฐน สโมธาเนติ, สติํ อาธิปเตยฺย\-ฏฺเฐน สโมธาเนติ, ปญฺญํ ตทุตฺต\-รฏฺเฐน สโมธาเนติ, วิมุตฺติํ สารฏฺเฐน สโมธาเนติ, อมโตคธํ นิพฺพานํ ปริโยสาน\-ฏฺเฐน สโมธาเนติ. อยํ ปุคฺคโล อิเม ธมฺเม อิมสฺมิํ อารมฺมเณ สโมธาเนติ, เตน วุจฺจติ “ธมฺเม สโมธาเนตี”ติ.}

\proseitem{3}{\textbf{โคจรญฺจ ปชานาตี}ติ ยํ ตสฺส อารมฺมณํ, ตํ ตสฺส โคจรํ. ยํ ตสฺส โคจรํ, ตํ ตสฺส อารมฺมณํ. \textbf{ปชานาตี}ติ ปุคฺคโล ปชานนา ปญฺญา. \textbf{สมนฺ}ติ อารมฺมณสฺส อุปฏฺฐานํ สมํ, จิตฺตสฺส อวิกฺเขโป สมํ, จิตฺตสฺส อธิฏฺฐานํ สมํ, จิตฺตสฺส โวทานํ สมํ. \textbf{อตฺโถ}ติ อนวชฺชฏฺโฐ นิเกฺลสฏฺโฐ โวทานฏฺโฐ ปรมฏฺโฐ. \textbf{ปฏิวิชฺฌตี}ติ อารมฺมณสฺส อุปฏฺฐานฏฺฐํ ปฏิวิชฺฌติ, จิตฺตสฺส อวิกฺเขปฏฺฐํ ปฏิวิชฺฌติ, จิตฺตสฺส อธิฏฺฐานฏฺฐํ ปฏิวิชฺฌติ, จิตฺตสฺส โวทานฏฺฐํ ปฏิวิชฺฌติ, เตน วุจฺจติ “สมตฺถญฺจ ปฏิวิชฺฌตี”ติ. (1)}

\proseitem{169}{กถํ รสฺสํ อสฺสสนฺโต “รสฺสํ อสฺสสามี”ติ ปชานาติ, รสฺสํ ปสฺสสนฺโต “รสฺสํ ปสฺสสามี”ติ ปชานาติ? รสฺสํ อสฺสาสํ อิตฺตรสงฺขาเต อสฺสสติ, รสฺสํ ปสฺสาสํ อิตฺตรสงฺขาเต ปสฺสสติ, รสฺสํ อสฺสาสปสฺสาสํ อิตฺตรสงฺขาเต อสฺสสติปิ ปสฺสสติปิ. รสฺสํ อสฺสาสปสฺสาสํ อิตฺตรสงฺขาเต อสฺสสโตปิ ปสฺสสโตปิ ฉนฺโท อุปฺปชฺชติ, ฉนฺทวเสน ตโต สุขุมตรํ รสฺสํ อสฺสาสํ อิตฺตรสงฺขาเต อสฺสสติ, ฉนฺทวเสน ตโต สุขุมตรํ รสฺสํ ปสฺสาสํ อิตฺตรสงฺขาเต ปสฺสสติ, ฉนฺทวเสน ตโต สุขุมตรํ รสฺสํ อสฺสาสปสฺสาสํ อิตฺตรสงฺขาเต อสฺสสติปิ ปสฺสสติปิ. ฉนฺทวเสน ตโต สุขุมตรํ รสฺสํ อสฺสาสปสฺสาสํ อิตฺตรสงฺขาเต อสฺสสโตปิ ปสฺสสโตปิ ปาโมชฺชํ อุปฺปชฺชติ, ปาโมชฺชวเสน ตโต สุขุมตรํ รสฺสํ อสฺสาสํ อิตฺตรสงฺขาเต อสฺสสติ, ปาโมชฺชวเสน ตโต สุขุมตรํ รสฺสํ ปสฺสาสํ อิตฺตรสงฺขาเต ปสฺสสติ, ปาโมชฺชวเสน ตโต สุขุมตรํ รสฺสํ อสฺสาสปสฺสาสํ อิตฺตรสงฺขาเต อสฺสสติปิ ปสฺสสติปิ. ปาโมชฺชวเสน ตโต สุขุมตรํ รสฺสํ อสฺสาสปสฺสาสํ อิตฺตรสงฺขาเต อสฺสสโตปิ ปสฺสสโตปิ, รสฺสา อสฺสาสปสฺสาสา จิตฺตํ วิวตฺตติ, อุเปกฺขา สณฺฐาติ. อิเมหิ นวหากาเรหิ รสฺสา \csromanfolio{181}อสฺสาสปสฺสาสา กาโย, อุปฏฺฐานํ สติ, อนุปสฺสนา ญาณํ. กาโย อุปฏฺฐานํ โน สติ, สติ อุปฏฺฐานญฺเจว สติ จ. ตาย สติยา เตน ญาเณน ตํ กายํ อนุปสฺสติ, เตน วุจฺจติ “กาเย กายานุปสฺสนา\-สติปฏฺฐาน\-ภาวนา”ติ.}

\prose{อนุปสฺสตีติ กถํ ตํ กายํ อนุปสฺสติ ฯเปฯ. เอวํ ตํ กายํ อนุปสฺสติ. พฺหาวนาติ จตสฺโส ภาวนา ฯเปฯ อาเสวนฏฺเฐน ภาวนา. รสฺสํ อสฺสาสปสฺสาส\-วเสน จิตฺตสฺส เอกคฺคตํ อวิกฺเขปํ ปชานโต วิทิตา เวทนา อุปฺปชฺชนฺติ ฯเปฯ. รสฺสํ อสฺสาสปสฺสาส\-วเสน จิตฺตสฺส เอกคฺคตํ อวิกฺเขปํ ปชานนฺโต อินฺทฺริยานิ สโมธาเนติ ฯเปฯ เตน วุจฺจติ “สมตฺถญฺจ ปฏิวิชฺฌตี”ติ. (2)}

\proseitem{170}{กถํ “สพฺพ\-กาย\-ปฏิสํเวที อสฺสสิสฺสามี”ติ สิกฺขติ, “สพฺพ\-กาย\-ปฏิสํเวที ปสฺสสิสฺสามี”ติ สิกฺขติ? กาโยติ เทฺว กายา นามกาโย จ รูปกาโย จ. กตโม นามกาโย? เวทนา สญฺญา เจตนา ผสฺโส มนสิกาโร นามญฺจ นามกาโย จ, เย จ วุจฺจนฺติ จิตฺตสงฺขารา, อยํ นามกาโย.  กตโม รูปกาโย? จตฺตาโร จ มหาภูตา จตุนฺนญฺจ มหาภูตานํ อุปาทายรูปํ อสฺสาโส จ ปสฺสาโส จ นิมิตฺตญฺจ อุปนิพนฺธนา, เย จ วุจฺจนฺติ กายสงฺขารา, อยํ รูปกาโย.}

\prose{กถํ เต กายา ปฏิวิทิตา โหนฺติ? ทีฆํ อสฺสาสวเสน จิตฺตสฺส เอกคฺคตํ อวิกฺเขปํ ปชานโต สติ อุปฏฺฐิตา โหติ, ตาย สติยา เตน ญาเณน เต กายา ปฏิวิทิตา โหนฺติ. ทีฆํ ปสฺสาสวเสน จิตฺตสฺส เอกคฺคตํ อวิกฺเขปํ ปชานโต สติ อุปฏฺฐิตา โหติ, ตาย สติยา เตน ญาเณน เต กายา ปฏิวิทิตา โหนฺติ. รสฺสํ อสฺสาสวเสน จิตฺตสฺส เอกคฺคตํ อวิกฺเขปํ ปชานโต สติ อุปฏฺฐิตา โหติ, ตาย สติยา เตน ญาเณน เต กายา ปฏิวิทิตา โหนฺติ. รสฺสํ ปสฺสาสวเสน จิตฺตสฺส เอกคฺคตํ อวิกฺเขปํ ปชานโต สติ อุปฏฺฐิตา โหติ, ตาย สติยา เตน ญาเณน เต กายา ปฏิวิทิตา โหนฺติ.}

\prose{อาวชฺชโต เต กายา ปฏิวิทิตา โหนฺติ, ปชานโต เต กายา ปฏิวิทิตา โหนฺติ, ปสฺสโต เต กายา ปฏิวิทิตา โหนฺติ, \csromanfolio{182}ปจฺจเวกฺขโต เต กายา ปฏิวิทิตา โหนฺติ, จิตฺตํ อธิฏฺฐหโต เต กายา ปฏิวิทิตา โหนฺติ, สทฺธาย อธิมุจฺจโต เต กายา ปฏิวิทิตา โหนฺติ, วีริยํ ปคฺคณฺหโต เต กายา ปฏิวิทิตา โหนฺติ, สติํ อุปฏฺฐาปยโต เต กายา ปฏิวิทิตา โหนฺติ, จิตฺตํ สมาทหโต เต กายา ปฏิวิทิตา โหนฺติ, ปญฺญาย ปชานโต เต กายา ปฏิวิทิตา โหนฺติ, อภิญฺเญยฺยํ อภิชานโต เต กายา ปฏิวิทิตา โหนฺติ, ปริญฺเญยฺยํ ปริชานโต เต กายา ปฏิวิทิตา โหนฺติ, ปหาตพฺพํ ปชหโต เต กายา ปฏิวิทิตา โหนฺติ, ภาเวตพฺพํ ภาวยโต เต กายา ปฏิวิทิตา โหนฺติ, สจฺฉิกาตพฺพํ สจฺฉิกโรโต เต กายา ปฏิวิทิตา โหนฺติ. เอวํ เต กายา ปฏิวิทิตา โหนฺติ, สพฺพ\-กาย\-ปฏิสํเวที อสฺสาสปสฺสาสา กาโย, อุปฏฺฐานํ สติ, อนุปสฺสนา ญาณํ. กาโย อุปฏฺฐานํ, โน สติ, สติ อุปฏฺฐานญฺเจว สติ จ. ตาย สติยา เตน ญาเณน ตํ กายํ อนุปสฺสติ, เตน วุจฺจติ “กาเย กายานุปสฺสนา\-สติปฏฺฐาน\-ภาวนา”ติ.}

\prose{\textbf{อนุปสฺสตี}ติ กถํ ตํ กายํ อนุปสฺสติ ฯเปฯ. เอวํ ตํ กายํ อนุปสฺสติ. ภาวนาติ จตสฺโส พฺ\textbf{หาวนา} ฯเปฯ. อาเสวนฏฺเฐน ภาวนา.}

\prose{สพฺพ\-กาย\-ปฏิสํเวที อสฺสาส\-ปสฺสาสานํ\csromansharedfootnote{b9c998ee299ba1e8}{อสฺสาสปสฺสาสา (สฺยา)} สํวรฏฺเฐน สีลวิสุทฺธิ, อวิกฺเขป\-ฏฺเฐน จิตฺตวิสุทฺธิ, ทสฺสนฏฺเฐน ทิฏฺฐิวิสุทฺธิ. โย ตตฺถ สํวรฏฺโฐ, อยํ อธิสีลสิกฺขา. โย ตตฺถ อวิกฺเขปฏฺโฐ, อยํ อธิจิตฺต\-สิกฺขา. โย ตตฺถ ทสฺสนฏฺโฐ, อยํ อธิปญฺญา\-สิกฺขา. อิมา ติสฺโส สิกฺขาโย อาวชฺชนฺโต สิกฺขติ, ชานนฺโต สิกฺขติ, ปสฺสนฺโต สิกฺขติ, ปจฺจเวกฺขนฺโต สิกฺขติ, จิตฺตํ อธิฏฺฐหนฺโต สิกฺขติ, สทฺธาย อธิมุจฺจนฺโต สิกฺขติ, วีริยํ ปคฺคณฺหนฺโต สิกฺขติ, สติํ อุปฏฺฐเปนฺโต สิกฺขติ, จิตฺตํ สมาทหนฺโต สิกฺขติ, ปญฺญาย ปชานนฺโต สิกฺขติ, อภิญฺเญยฺยํ อภิชานนฺโต สิกฺขติ, ปริญฺเญยฺยํ ปริชานนฺโต สิกฺขติ, ปหาตพฺพํ ปชหนฺโต สิกฺขติ, ภาเวตพฺพํ ภาเวนฺโต สิกฺขติ, สิจฺฉิกาตพฺพํ สจฺฉิกโรนฺโต สิกฺขติ.}

\prose{\csromanfolio{183}สพฺพ\-กาย\-ปฏิสํเวที อสฺสาสปสฺสาส\-วเสน จิตฺตสฺส เอกคฺคตํ อวิกฺเขปํ ปชานโต วิทิตา เวทนา อุปฺปชฺชนฺติ ฯเปฯ. สพฺพ\-กาย\-ปฏิสํเวที อสฺสาสปสฺสาส\-วเสน จิตฺตสฺส เอกคฺคตํ อวิกฺเขปํ ปชานนฺโต อินฺทฺริยานิ สโมธาเนติ ฯเปฯ เตน วุจฺจติ “สมตฺถญฺจ ปฏิวิชฺฌตี”ติ. (3)}

\proseitem{171}{กถํ “ปสฺสมฺภยํ กายสงฺขารํ อสฺสสิสฺสามี”ติ สิกฺขติ, “ปสฺสมฺภยํ กายสงฺขารํ ปสฺสสิสฺสามี”ติ สิกฺขติ? กตโม กายสงฺขาโร? ทีฆํ อสฺสาสา กายิกา เอเต ธมฺมา กายปฏิพทฺธา กายสงฺขารา, เต กายสงฺขาเร ปสฺสมฺเภนฺโต นิโรเธนฺโต วูปสเมนฺโต สิกฺขติ. ทีฆํ ปสฺสาสา กายิกา เอเต ธมฺมา กายปฏิพทฺธา กายสงฺขารา, เต กายสงฺขาเร ปสฺสมฺเภนฺโต นิโรเธนฺโต วูปสเมนฺโต สิกฺขติ. รสฺสํ อสฺสาสา. รสฺสํ ปสฺสาสา. สพฺพ\-กาย\-ปฏิสํเวที อสฺสาสา. สพฺพ\-กาย\-ปฏิสํเวที ปสฺสาสา กายิกา เอเต ธมฺมา กายปฏิพทฺธา กายสงฺขารา, เต กายสงฺขาเร ปสฺสมฺเภนฺโต นิโรเธนฺโต วูปสเมนฺโต สิกฺขติ.}

\prose{ยถารูเปหิ กายสงฺขาเรหิ ยา กายสฺส อานมนา วินมนา สนฺนมนา ปณมนา อิญฺชนา ผนฺทนา จลนา ปกมฺปนา, “ปสฺสมฺภยํ กายสงฺขารํ อสฺสสิสฺสามี”ติ สิกฺขติ, “ปสฺสมฺภยํ กายสงฺขารํ ปสฺสสิสฺสามี”ติ สิกฺขติ. ยถารูเปหิ กายสงฺขาเรหิ ยา กายสฺส น อานมนา น วินมนา น สนฺนมนา น ปณมนา อนิญฺชนา อผนฺทนา อจลนา อกมฺปนา, “สนฺตํ สุขุมํ ปสฺสมฺภยํ กายสงฺขารํ อสฺสสิสฺสามี”ติ สิกฺขติ, “ปสฺสมฺภยํ กายสงฺขารํ ปสฺสสิสฺสามี”ติ สิกฺขติ.}

\prose{อิติ กิร “ปสฺสมฺภยํ กายสงฺขารํ อสฺสสิสฺสามี”ติ สิกฺขติ, “ปสฺสมฺภยํ กายสงฺขารํ ปสฺสสิสฺสามี”ติ สิกฺขติ. เอวํ สนฺเต วาตูปลทฺธิยา จ ปภาวนา น โหติ, อสฺสาสปสฺสาสานญฺจ ปภาวนา น โหติ, อานาปานสฺสติยา จ ปภาวนา น โหติ, อานาปาน\-สฺสติ\-สมาธิสฺส จ ปภาวนา น โหติ, น จ นํ ตํ สมาปตฺติํ ปณฺฑิตา สมาปชฺชนฺติปิ วุฏฺฐหนฺติปิ.}

\prose{อิติ กิร “ปสฺสมฺภยํ กายสงฺขารํ อสฺสสิสฺสามี”ติ สิกฺขติ, “ปสฺสมฺภยํ กายสงฺขารํ ปสฺสสิสฺสามี”ติ สิกฺขติ. เอวํ สนฺเต วาตูปลทฺธิยา จ ปภาวนา โหติ, อสฺสาสปสฺสาสานญฺจ ปภาวนา โหติ, อานาปานสฺสติยา จ ปภาวนา โหติ, อานาปาน\-สฺสติ\-สมาธิสฺส \csromanfolio{184}จ ปพฺหาวนา โหติ, ตญฺจ นํ สมาปตฺติํ ปณฺฑิตา สมาปชฺชนฺติปิ วุฏฺฐหนฺติปิ. ยถา กถํ วิย เสยฺยถาปิ กํเส อาโกฏิเต ปฐมํ โอฬาริกา สทฺทา ปวตฺตนฺติ, โอฬาริกานํ สทฺทานํ นิมิตฺตํ สุคฺคหิตตฺตา สุมนสิกต\-ตฺตา สูปธาริตตฺตา นิรุทฺเธปิ โอฬาริเก สทฺเท อถ ปจฺฉา สุขุมกา สทฺทา ปวตฺตนฺติ. สุขุมกานํ สทฺทานํ นิมิตฺตํ สุคฺคหิตตฺตา สุมนสิกต\-ตฺตา สูปธาริตตฺตา นิรุทฺเธปิ สุขุมเก สทฺเท อถ ปจฺฉา สุขุม\-สทฺทนิมิตฺตา\-รมฺมณตาปิ จิตฺตํ ปวตฺตติ. เอวเมวํ ปฐมํ โอฬาริกา อสฺสาสปสฺสาสา ปวตฺตนฺติ, โอฬาริกานํ อสฺสาส\-ปสฺสาสานํ นิมิตฺตํ สุคฺคหิตตฺตา สุมนสิกต\-ตฺตา สูปธาริตตฺตา นิรุทฺเธปิ โอฬาริเก อสฺสาสปสฺสาเส อถ ปจฺฉา สุขุมกา อสฺสาสปสฺสาสา ปวตฺตนฺติ, สุขุมกานํ อสฺสาส\-ปสฺสาสานํ นิมิตฺตํ สุคฺคหิตตฺตา สุมนสิกต\-ตฺตา สูปธาริตตฺตา นิรุทฺเธปิ สุขุมเก อสฺสาสปสฺสาเส อถ ปจฺฉา สุขุมก\-อสฺสาส\-ปสฺสาสานํ นิมิตฺตา\-รมฺมณตาปิ จิตฺตํ น วิกฺเขปํ คจฺฉติ.}

\prose{เอวํ สนฺเต วาตูปลทฺธิยา จ ปภาวนา โหติ, อสฺสาสปสฺสาสานญฺจ ปภาวนา โหติ, อานาปานสฺสติยา จ ปภาวนา โหติ, อานาปาน\-สฺสติ\-สมาธิสฺส จ ปภาวนา โหติ, ตญฺจ นํ สมาปตฺติํ ปณฺฑิตา สมาปชฺชนฺติปิ วุฏฺฐหนฺติปิ. ปสฺสมฺภยํ กายสงฺขารํ อสฺสาสปสฺสาสา กาโย, อุปฏฺฐานํ สติ, อนุปสฺสนา ญาณํ. กาโย อุปฏฺฐานํ, โน สติ, สติ อุปฏฺฐานญฺเจว สติ จ. ตาย สติยา เตน ญาเณน ตํ กายํ อนุปสฺสติ, เตน วุจฺจติ “กาเย กายานุปสฺสนา\-สติปฏฺฐาน\-ภาวนา”ติ.}

\prose{อนุปสฺสตีติ กถํ ตํ กายํ อนุปสฺสติ ฯเปฯ. เอวํ ตํ กายํ อนุปสฺสติ. ภาวนาติ จตสฺโส พฺหาวนา ฯเปฯ อาเสวนฏฺเฐน ภาวนา. ปสฺสมฺภยํ กายสงฺขารํ อสฺสาส\-ปสฺสาสานํ สํวรฏฺเฐน สีลวิสุทฺธิ, อวิกฺเขป\-ฏฺเฐน จิตฺตวิสุทฺธิ, ทสฺสนฏฺเฐน ทิฏฺฐิวิสุทฺธิ. โย ตตฺถ สํวรฏฺโฐ, อยํ อธิสีลสิกฺขา. โย ตตฺถ อวิกฺเขปฏฺโฐ, อยํ อธิจิตฺต\-สิกฺขา. โย ตตฺถ ทสฺสนฏฺโฐ, อยํ อธิปญฺญา\-สิกฺขา. อิมา ติสฺโส สิกฺขาโย อาวชฺชนฺโต สิกฺขติ ฯเปฯ สจฺฉิกาตพฺพํ สจฺฉิกโรนฺโต สิกฺขติ, ปสฺสมฺภยํ กายสงฺขารํ อสฺสาสปสฺสาส\-วเสน จิตฺตสฺส เอกคฺคตํ อวิกฺเขปํ ปชานโต วิทิตา เวทนา อุปฺปชฺชนฺติ ฯเปฯ. ปสฺสมฺภยํ กายสงฺขารํ อสฺสาสปสฺสาส\-วเสน จิตฺตสฺส เอกคฺคตํ อวิกฺเขปํ \csromanfolio{185}ปชานนฺโต อินฺทฺริยานิ สโมธาเนติ ฯเปฯ เตน วุจฺจติ “สมตฺถญฺจ ปฏิวิชฺฌตี”ติ. (4)}

\prosewithcloserruled{อฏฺฐ อนุปสฺสนา\-ญาณานิ, อฏฺฐ จ อุปฏฺฐานา\-นุสฺสติโย, จตฺตาริ สุตฺตนฺติก\-วตฺถูนิ กาเย กายา\-นุปสฺสนาย.}{ภาณวาโร.}
\csromanheader{2}{ทุติย\-จตุกฺก\-นิทฺเทส}
\proseitem{172}{กถํ “ปีติปฏิสํเวที อสฺสสิสฺสามี”ติ สิกฺขติ, “ปีติปฏิสํเวที ปสฺสสิสฺสามี”ติ สิกฺขติ? กตมา ปีติ? ทีฆํ อสฺสาสวเสน จิตฺตสฺส เอกคฺคตํ อวิกฺเขปํ ปชานโต อุปฺปชฺชติ ปีติ ปาโมชฺชํ, ยา ปีติ ปาโมชฺชํ อาโมทนา ปโมทนา หาโส ปหาโส วิตฺติ โอทคฺยํ อตฺตมนตา จิตฺตสฺส. ทีฆํ ปสฺสาสวเสน จิตฺตสฺส เอกคฺคตํ อวิกฺเขปํ ปชานโต อุปฺปชฺชติ ปีติ ปาโมชฺชํ ฯเปฯ. รสฺสํ อสฺสาสวเสน. รสฺสํ ปสฺสาสวเสน. สพฺพ\-กาย\-ปฏิสํเวที อสฺสาสวเสน. สพฺพ\-กาย\-ปฏิสํเวที ปสฺสาสวเสน. ปสฺสมฺภยํ กายสงฺขารํ อสฺสาสวเสน. ปสฺสมฺภยํ กายสงฺขารํ ปสฺสาสวเสน จิตฺตสฺส เอกคฺคตํ อวิกฺเขปํ ปชานโต อุปฺปชฺชติ ปีติ ปาโมชฺชํ. ยา ปีติ ปาโมชฺชํ อาโมทนา ปโมทนา หาโส ปหาโส วิตฺติ โอทคฺยํ อตฺตมนตา จิตฺตสฺส, อยํ ปีติ.}

\prose{กถํ สา ปีติ ปฏิวิทิตา โหติ? ทีฆํ อสฺสาสปสฺสาส\-วเสน จิตฺตสฺส เอกคฺคตํ อวิกฺเขปํ ปชานโต สติ อุปฏฺฐิตา โหติ, ตาย สติยา เตน ญาเณน สา ปีติ ปฏิวิทิตา โหติ. ทีฆํ ปสฺสาสวเสน จิตฺตสฺส เอกคฺคตํ อวิกฺเขปํ ปชานโต สติ อุปฏฺฐิตา โหติ, ตาย สติยา เตน ญาเณน สา ปีติ ปฏิวิทิตา โหติ. รสฺสํ อสฺสาสวเสน ฯเปฯ. รสฺสํ ปสฺสาสวเสน. สพฺพ\-กาย\-ปฏิสํเวที อสฺสาสวเสน. สพฺพ\-กาย\-ปฏิสํเวที ปสฺสาสวเสน. ปสฺสมฺภยํ กายสงฺขารํ อสฺสาสวเสน. ปสฺสมฺภยํ กายสงฺขารํ ปสฺสาสวเสน จิตฺตสฺส เอกคฺคตํ อวิกฺเขปํ ปชานโต สติ อุปฏฺฐิตา โหติ, ตาย สติยา เตน ญาเณน สา ปีติ ปฏิวิทิตา โหติ. อาวชฺชโต สา ปีติ ปฏิวิทิตา โหติ, ชานโต. ปสฺสโต. ปจฺจเวกฺขโต. จิตฺตํ อธิฏฺฐหโต. \csromanfolio{186}สทฺธาย อธิมุจฺจโต. วีริยํ ปคฺคณฺหโต. สติํ อุปฏฺฐาปยโต. จิตฺตํ สมาทหโต. ปญฺญาย ปชานโต. อภิญฺเญยฺยํ อภิชานโต. ปริญฺเญยฺยํ ปริชานโต. ปหาตพฺพํ ปชหโต. ภาเวตพฺพํ ภาวยโต. สจฺฉิกาตพฺพํ สจฺฉิกโรโต สา ปีติ ปฏิวิทิตา โหติ. เอวํ สา ปีติ ปฏิวิทิตา โหติ.}

\prose{ปีติปฏิสํเวที อสฺสาสปสฺสาส\-วเสน เวทนา, อุปฏฺฐานํ สติ, อนุปสฺสนา ญาณํ. เวทนา อุปฏฺฐานํ, โน สติ, สติ อุปฏฺฐานญฺเจว สติ จ.ตาย สติยา เตน ญาเณน ตํ เวทนํ อนุปสฺสติ, เตน วุจฺจติ “เวทนาสุ เวทนานุปสฺสนาสติปฏฺฐานภาวนา”ติ.}

\prose{อนุปสฺสตีติ กถํ ตํ เวทนํ อนุปสฺสติ, อนิจฺจโต อนุปสฺสติ ฯเปฯ. เอวํ ตํ เวทนํ อนุปสฺสติ. ภาวนาติ จตสฺโส ภาวนา ฯเปฯ อาเสวนฏฺเฐน ภาวนา. ปีติปฏิสํเวที อสฺสาส\-ปสฺสาสานํ สํวรฏฺเฐน สีลวิสุทฺธิ ฯเปฯ ปีติปฏิสํเวที อสฺสาสปสฺสาส\-วเสน จิตฺตสฺส เอกคฺคตํ อวิกฺเขปํ ปชานโต ฯเปฯ ปชานนฺโต อินฺทฺริยานิ สโมธาเนติ, เตน วุจฺจติ “สมตฺถญฺจ ปฏิวิชฺฌตี”ติ. (1)}

\proseitem{173}{กถํ “สุขปฏิสํเวที อสฺสสิสฺสามี”ติ สิกฺขติ, “สุขปฏิสํเวที ปสฺสสิสฺสามี”ติ สิกฺขติ? สุขนฺติ เทฺว สุขานิ กายิกญฺจ สุขํ, เจตสิกญฺจ สุขํ. กตมํ กายิกํ สุขํ? ยํ กายิกํ สาตํ กายิกํ สุขํ กาย\-สมฺผสฺสชํ สาตํ สุขํ เวทยิตํ กาย\-สมฺผสฺสชา สาตา สุขา เวทนา, อิทํ กายิกํ สุขํ. กตมํ เจตสิกํ สุขํ? ยํ เจตสิกํ สาตํ เจตสิกํ สุขํ เจโต\-สมฺผสฺสชํ สาตํ สุขํ เวทยิตํ เจโต\-สมฺผสฺสชา สาตา สุขา เวทนา, อิทํ เจตสิกํ สุขํ.}

\prose{กถํ เต สุขา ปฏิวิทิตา โหนฺติ? ทีฆํ อสฺสาสวเสน จิตฺตสฺส เอกคฺคตํ อวิกฺเขปํ ปชานโต สติ อุปฏฺฐิตา โหติ, ตาย สติยา เตน ญาเณน เต สุขา ปฏิวิทิตา โหนฺติ. ทีฆํ ปสฺสาสวเสน จิตฺตสฺส เอกคฺคตํ อวิกฺเขปํ ปชานโต สติ อุปฏฺฐิตา โหติ, ตาย สติยา เตน ญาเณน เต สุขา ปฏิวิทิตา โหนฺติ ฯเปฯ สจฺฉิกาตพฺพํ สจฺฉิกโรโต เต สุขา ปฏิวิทิตา โหนฺติ. เอวํ เต สุขา ปฏิวิทิตา โหนฺติ. สุขปฏิสํเวที อสฺสาสปสฺสาส\-วเสน เวทนา, อุปฏฺฐานํ สติ, อนุปสฺสนา ญาณํ. เวทนา อุปฏฺฐานํ, โน สติ, สติ อุปฏฺฐานญฺเจว \csromanfolio{187}สติ จ, ตาย สติยา เตน ญาเณน ตํ เวทนํ อนุปสฺสติ, เตน วุจฺจติ “เวทนาสุ เวทนานุปสฺสนาสติปฏฺฐานภาวนา”ติ.}

\prose{อนุปสฺสตีติ กถํ ตํ เวทนํ อนุปสฺสติ, อนิจฺจโต อนุปสฺสติ ฯเปฯ. เอวํ ตํ เวทนํ อนุปสฺสติ. ภาวนาติ จตสฺโส ภาวนา ฯเปฯ อาเสวนฏฺเฐน ภาวนา. สุขปฏิสํเวที อสฺสาส\-ปสฺสาสานํ สํวรฏฺเฐน สีลวิสุทฺธิ ฯเปฯ. สุขปฏิสํเวที อสฺสาสปสฺสาส\-วเสน จิตฺตสฺส เอกคฺคตํ อวิกฺเขปํ ปชานโต ฯเปฯ. ปชานนฺโต อินฺทฺริยานิ สโมธาเนติ, เตน วุจฺจติ “สมตฺถญฺจ ปฏิวิชฺฌตี”ติ. (2)}

\proseitem{174}{กถํ “จิตฺต\-สงฺขาร\-ปฏิสํเวที อสฺสสิสฺสามี”ติ สิกฺขติ, “จิตฺต\-สงฺขาร\-ปฏิสํเวที ปสฺสสิสฺสามี”ติ สิกฺขติ? กตโม จิตฺตสงฺขาโร? ทีฆํ อสฺสาสวเสน สญฺญา จ เวทนา จ เจตสิกา, เอเต ธมฺมา จิตฺต\-ปฏิพทฺธา จิตฺตสงฺขารา. ทีฆํ ปสฺสาสวเสน สญฺญา จ เวทนา จ เจตสิกา, เอเต ธมฺมา จิตฺต\-ปฏิพทฺธา จิตฺตสงฺขารา ฯเปฯ. สุขปฏิสํเวที อสฺสาสวเสน. สุขปฏิสํเวที ปสฺสาสวเสน สญฺญา จ เวทนา จ เจตสิกา, เอเต ธมฺมา จิตฺต\-ปฏิพทฺธา จิตฺตสงฺขารา, อยํ จิตฺตสงฺขาโร.}

\prose{กถํ เต จิตฺตสงฺขารา ปฏิวิทิตา โหนฺติ? ทีฆํ อสฺสาสวเสน จิตฺตสฺส เอกคฺคตํ อวิกฺเขปํ ปชานโต สติ อุปฏฺฐิตา โหติ, ตาย สติยา เตน ญาเณน เต จิตฺตสงฺขารา ปฏิวิทิตา โหนฺติ. ทีฆํ ปสฺสาสวเสน จิตฺตสฺส เอกคฺคตํ อวิกฺเขปํ ปชานโต สติ อุปฏฺฐิตา โหติ, ตาย สติยา เตน ญาเณน เต จิตฺตสงฺขารา ปฏิวิทิตา โหนฺติ ฯเปฯ. สจฺฉิกาตพฺพํ สจฺฉิกโรโต เต จิตฺตสงฺขารา ปฏิวิทิตา โหนฺติ. เอวํ เต จิตฺตสงฺขารา ปฏิวิทิตา โหนฺติ. จิตฺต\-สงฺขาร\-ปฏิสํเวที อสฺสาสปสฺสาส\-วเสน เวทนา, อุปฏฺฐานํ สติ, อนุปสฺสนา ญาณํ. เวทนา อุปฏฺฐานํ, โน สติ, สติ อุปฏฺฐานญฺเจว สติ จ, ตาย สติยา เตน ญาเณน ตํ เวทนํ อนุปสฺสติ, เตน วุจฺจติ “เวทนาสุ เวทนานุปสฺสนาสติปฏฺฐานภาวนา”ติ.}

\prose{อนุปสฺสตีติ กถํ ตํ เวทนํ อนุปสฺสติ, อนิจฺจโต อนุปสฺสติ ฯเปฯ. เอวํ ตํ เวทนํ อนุปสฺสติ. ภาวนาติ จตสฺโส ภาวนา ฯเปฯ อาเสวนฏฺเฐน ภาวนา. จิตฺต\-สงฺขาร\-ปฏิสํเวที อสฺสาส\-ปสฺสาสานํ สํวรฏฺเฐน สีลวิสุทฺธิ ฯเปฯ. จิตฺต\-สงฺขาร\-ปฏิสํเวที อสฺสาสปสฺสาส\-วเสน \csromanfolio{188}จิตฺตสฺส เอกคฺคตํ อวิกฺเขปํ ปชานโต ฯเปฯ ปชานนฺโต อินฺทฺริยานิ สโมธาเนติ, เตน วุจฺจติ “สมตฺถญฺจ ปฏิวิชฺฌตี”ติ. (3)}

\proseitem{175}{กถํ “ปสฺสมฺภยํ จิตฺตสงฺขารํ อสฺสสิสฺสามี”ติ สิกฺขติ, “ปสฺสมฺภยํ จิตฺตสงฺขารํ ปสฺสสิสฺสามี”ติ สิกฺขติ? กตโม จิตฺตสงฺขาโร? ทีฆํ อสฺสาสวเสน สญฺญา จ เวทนา จ เจตสิกา, เอเต ธมฺมา จิตฺต\-ปฏิพทฺธา จิตฺตสงฺขารา, เต จิตฺตสงฺขาเร ปสฺสมฺเภนฺโต นิโรเธนฺโต วูปสเมนฺโต สิกฺขติ. ทีฆํ ปสฺสาสวเสน สญฺญา จ เวทนา จ เจตสิกา, เอเต ธมฺมา จิตฺต\-ปฏิพทฺธา จิตฺตสงฺขารา, เต จิตฺตสงฺขาเร ปสฺสมฺเภนฺโต นิโรเธนฺโต วูปสเมนฺโต สิกฺขติ. จิตฺต\-สงฺขาร\-ปฏิสํเวที อสฺสาสวเสน. จิตฺต\-สงฺขาร\-ปฏิสํเวที ปสฺสาสวเสน สญฺญา จ เวทนา จ เจตสิกา, เอเต ธมฺมา จิตฺต\-ปฏิพทฺธา จิตฺตสงฺขารา, เต จิตฺตสงฺขาเร ปสฺสมฺเภนฺโต นิโรเธนฺโต วูปสเมนฺโต สิกฺขติ. ปสฺสมฺภยํ จิตฺตสงฺขารํ อสฺสาสปสฺสาส\-วเสน เวทนา, อุปฏฺฐานํ สติ, อนุปสฺสนา ญาณํ. เวทนา อุปฏฺฐานํ, โน สติ, สติ อุปฏฺฐานญฺเจว สติ จ. ตาย สติยา เตน ญาเณน ตํ เวทนํ อนุปสฺสติ, เตน วุจฺจติ “เวทนาสุ เวทนานุปสฺสนาสติปฏฺฐานภาวนา”ติ.}

\prose{อนุปสฺสตีติ กถํ ตํ เวทนํ อนุปสฺสติ ฯเปฯ. เอวํ ตํ เวทนํ อนุปสฺสติ, ภาวนาติ จตสฺโส ภาวนา ฯเปฯ อาเสวนฏฺเฐน ภาวนา. ปสฺสมฺภยํ จิตฺตสงฺขารํ อสฺสาส\-ปสฺสาสานํ สํวรฏฺเฐน สีลวิสุทฺธิ ฯเปฯ. ปสฺสมฺภยํ จิตฺตสงฺขารํ อสฺสาสปสฺสาส\-วเสน จิตฺตสฺส เอกคฺคตํ อวิกฺเขปํ ปชานโต ฯเปฯ. ปชานนฺโต อินฺทฺริยานิ สโมธาเนติ, เตน วุจฺจติ “สมตฺถญฺจ ปฏิวิชฺฌตี”ติ. (4)}

\prose{อฏฺฐ อนุปสฺสนา\-ญาณานิ, อฏฺฐ จ อุปฏฺฐานา\-นุสฺสติโย, จตฺตาริ สุตฺตนฺติก\-วตฺถูนิ เวทนาสุ เวทนา\-นุปสฺสนาย.}

\csromanheader{2}{ตติย\-จตุกฺก\-นิทฺเทส}
\proseitem{176}{กถํ “จิตฺต\-ปฏิสํเวที อสฺสสิสฺสามี”ติ สิกฺขติ, “จิตฺต\-ปฏิสํเวที ปสฺสสิสฺสามี”ติ สิกฺขติ? กตมํ ตํ จิตฺตํ? ทีฆํ อสฺสาสวเสน วิญฺญาณํ จิตฺตํ. ยํ จิตฺตํ มโน มานสํ หทยํ ปณฺฑรํ มโน มนายตนํ มนินฺทฺริยํ วิญฺญาณํ วิญฺญาณ\-กฺขนฺโธ ตชฺชา มโนวิญฺญาณ\-ธาตุ. ทีฆํ ปสฺสาสวเสน ฯเปฯ. ปสฺสมฺภยํ จิตฺตสงฺขารํ อสฺสาสวเสน. ปสฺสมฺภยํ จิตฺตสงฺขารํ \csromanfolio{189}ปสฺสาสวเสน วิญฺญาณํ จิตฺตํ. ยํ จิตฺตํ มโน มานสํ หทยํ ปณฺฑรํ มโน มนายตนํ มนินฺทฺริยํ วิญฺญาณํ วิญฺญาณ\-กฺขนฺโธ ตชฺชา มโนวิญฺญาณ\-ธาตุ, อิทํ จิตฺตํ.}

\prose{กถํ ตํ จิตฺตํ ปฏิวิทิตํ โหติ? ทีฆํ อสฺสาสวเสน จิตฺตสฺส เอกคฺคตํ อวิกฺเขปํ ปชานโต สติ อุปฏฺฐิตา โหติ, ตาย สติยา เตน ญาเณน ตํ จิตฺตํ ปฏิวิทิตํ โหติ. ทีฆํ ปสฺสาสวเสน จิตฺตสฺส เอกคฺคตํ อวิกฺเขปํ ปชานโต สติ อุปฏฺฐิตา โหติ, ตาย สติยา เตน ญาเณน ตํ จิตฺตํ ปฏิวิทิตํ โหติ ฯเปฯ สจฺฉิกาตพฺพํ สจฺฉิกโรโต ตํ จิตฺตํ ปฏิวิทิตํ โหติ. เอวํ ตํ จิตฺตํ ปฏิวิทิตํ โหติ. จิตฺต\-ปฏิสํเวที อสฺสาสปสฺสาส\-วเสน วิญฺญาณํ จิตฺตํ, อุปฏฺฐานํ สติ, อนุปสฺสนา ญาณํ. จิตฺตํ อุปฏฺฐานํ, โน สติ, สติ อุปฏฺฐานญฺเจว สติ จ. ตาย สติยา เตน ญาเณน ตํ จิตฺตํ อนุปสฺสติ, เตน วุจฺจติ “จิตฺเต จิตฺตานุปสฺสนาสติปฏฺฐานภาวนา”ติ.}

\prose{อนุปสฺสตีติ กถํ ตํ จิตฺตํ อนุปสฺสติ ฯเปฯ. เอวํ ตํ จิตฺตํ อนุปสฺสติ. ภาวนาติ จตสฺโส พฺหาวนา ฯเปฯ อาเสวนฏฺเฐน ภาวนา. จิตฺต\-ปฏิสํเวที อสฺสาส\-ปสฺสาสานํ สํวรฏฺเฐน สีลวิสุทฺธิ ฯเปฯ. จิตฺต\-ปฏิสํเวที อสฺสาสปสฺสาส\-วเสน จิตฺตสฺส เอกคฺคตํ อวิกฺเขปํ ปชานโต ฯเปฯ ปชานนฺโต อินฺทฺริยานิ สโมธาเนติ, เตน วุจฺจติ “สมตฺถญฺจ ปฏิวิชฺฌตี”ติ. (1)}

\proseitem{177}{กถํ “อภิปฺปโมทยํ จิตฺตํ อสฺสสิสฺสามี”ติ สิกฺขติ, “อภิปฺปโมทยํ จิตฺตํ ปสฺสสิสฺสามี”ติ สิกฺขติ? กตโม จิตฺตสฺส อภิปฺปโมโท? ทีฆํ อสฺสาสวเสน จิตฺตสฺส เอกคฺคตํ อวิกฺเขปํ ปชานโต อุปฺปชฺชติ จิตฺตสฺส อภิปฺปโมโท, ยา จิตฺตสฺส อาโมทนา ปโมทนา หาโส ปหาโส วิตฺติ โอทคฺยํ อตฺตมนตา จิตฺตสฺส. ทีฆํ อสฺสาสปสฺสาส\-วเสน จิตฺตสฺส เอกคฺคตํ อวิกฺเขปํ ปชานโต อุปฺปชฺชติ จิตฺตสฺส อภิปฺปโมโท, ยา จิตฺตสฺส อาโมทนา ปโมทนา หาโส ปหาโส วิตฺติ โอทคฺยํ อตฺตมนตา จิตฺตสฺส ฯเปฯ. จิตฺต\-ปฏิสํเวที อสฺสาสวเสน. จิตฺต\-ปฏิสํเวที ปสฺสาสวเสน จิตฺตสฺส เอกคฺคตํ อวิกฺเขปํ ปชานโต อุปฺปชฺชติ จิตฺตสฺส อภิปฺปโมโท, ยา จิตฺตสฺส อาโมทนา ปโมทนา หาโส ปหาโส วิตฺติ โอทคฺยํ อตฺตมนตา จิตฺตสฺส, อยํ จิตฺตสฺส อภิปฺปโมโท. อภิปฺปโมทยํ จิตฺตํ อสฺสาสปสฺสาส\-วเสน วิญฺญาณํ จิตฺตํ, อุปฏฺฐานํ \csromanfolio{190}สติ, อนุปสฺสนา ญาณํ. จิตฺตํ อุปฏฺฐานํ, โน สติ, สติ อุปฏฺฐานญฺเจว สติ จ. ตาย สติยา เตน ญาเณน ตํ จิตฺตํ อนุปสฺสติ, เตน วุจฺจติ “จิตฺเต จิตฺตานุปสฺสนาสติปฏฺฐานภาวนา”ติ.}

\prose{อนุปสฺสตีติ กถํ ตํ จิตฺตํ อนุปสฺสติ ฯเปฯ. เอวํ ตํ จิตฺตํ อนุปสฺสตีติ. ภาวนาติ จตสฺโส พฺหาวนา ฯเปฯ อาเสวนฏฺเฐน ภาวนา. อภิปฺปโมทยํ จิตฺตํ อสฺสาส\-ปสฺสาสานํ สํวรฏฺเฐน สีลวิสุทฺธิ ฯเปฯ. อภิปฺปโมทยํ จิตฺตํ อสฺสาสปสฺสาส\-วเสน จิตฺตสฺส เอกคฺคตํ อวิกฺเขปํ ปชานโต ฯเปฯ ปชานนฺโต อินฺทฺริยานิ สโมธาเนติ, เตน วุจฺจติ “สมตฺถญฺจ ปฏิวิชฺฌตี”ติ. (2)}

\proseitem{178}{กถํ “สมาทหํ จิตฺตํ อสฺสสิสฺสามี”ติ สิกฺขติ, “สมาทหํ จิตฺตํ ปสฺสสิสฺสามี”ติ สิกฺขติ? กตโม สมาธิ\csromansharedfootnote{a9a3e3b77a51ca68}{กตมญฺจ สมาธินฺทฺริยํ (สฺยา)}? ทีฆํ อสฺสาสวเสน จิตฺตสฺส เอกคฺคตา อวิกฺเขโป สมาธิ, ยา จิตฺตสฺส ฐิติ สณฺฐิติ อวฏฺฐิติ\csromansharedfootnote{00dc4d8c2e02237a}{อธิฏฺฐิติ (สฺยา)} อวิสาหาโร อวิกฺเขโป อวิสาหฏ\-มานสตา\csromansharedfootnote{1dc199891782ac32}{อวิสาหตมานสตา (สฺยา, ก)} สมโถ สมาธินฺทฺริยํ สมาธิพลํ สมฺมาสมาธิ. ทีฆํ ปสฺสาสวเสน จิตฺตสฺส เอกคฺคตา อวิกฺเขโป สมาธิ ฯเปฯ สมาทหํ จิตฺตํ อสฺสาสวเสน ฯเปฯ สมาทหํ จิตฺตํ ปสฺสาสวเสน จิตฺตสฺส เอกคฺคตา อวิกฺเขโป สมาธิ, ยา จิตฺตสฺส ฐิติ สณฺฐิติ อวฏฺฐิติ อวิสาหาโร อวิกฺเขโป อวิสาหฏ\-มานสตา สมโถ สมาธินฺทฺริยํ สมาธิพลํ สมฺมาสมาธิ, อยํ สมาธิ. สมาทหํ จิตฺตํ อสฺสาสปสฺสาส\-วเสน วิญฺญาณํ จิตฺตํ, อุปฏฺฐานํ สติ, อนุปสฺสนา ญาณํ. จิตฺตํ อุปฏฺฐานํ, โน สติ. สติ อุปฏฺฐานญฺเจว สติ จ. ตาย สติยา เตน ญาเณน ตํ จิตฺตํ อนุปสฺสติ, เตน วุจฺจติ “จิตฺเต จิตฺตานุปสฺสนาสติปฏฺฐานภาวนา”ติ.}

\prose{อนุปสฺสตีติ กถํ ตํ จิตฺตํ อนุปสฺสติ ฯเปฯ. เอวํ ตํ จิตฺตํ อนุปสฺสติ. ภาวนาติ จตสฺโส ภาวนา ฯเปฯ อาเสวนฏฺเฐน ภาวนา. สมาทหํ จิตฺตํ อสฺสาส\-ปสฺสาสานํ สํวรฏฺเฐน สีลวิสุทฺธิ ฯเปฯ สมาทหํ จิตฺตํ อสฺสาสปสฺสาส\-วเสน จิตฺตสฺส เอกคฺคตํ อวิกฺเขปํ ปชานโต ฯเปฯ ปชานนฺโต อินฺทฺริยานิ สโมธาเนติ, เตน วุจฺจติ “สมตฺถญฺจ ปฏิวิชฺฌตี”ติ. (3)}

\proseitem{179}{\csromanfolio{191}กถํ “วิโมจยํ จิตฺตํ อสฺสสิสฺสามี”ติ สิกฺขติ, “วิโมจยํ จิตฺตํ ปสฺสสิสฺสามี”ติ สิกฺขติ? “ราคโต วิโมจยํ จิตฺตํ อสฺสสิสฺสามี”ติ สิกฺขติ, “ราคโต วิโมจยํ จิตฺตํ ปสฺสสิสฺสามี”ติ สิกฺขติ. “โทสโต วิโมจยํ จิตฺตํ อสฺสสิสฺสามี”ติ สิกฺขติ, “โทสโต วิโมจยํ จิตฺตํ ปสฺสสิสฺสามี”ติ สิกฺขติ. “โมหโต วิโมจยํ จิตฺตํ อสฺสสิสฺสามี”ติ สิกฺขติ ฯเปฯ “มานโต วิโมจยํ จิตฺตํ. ทิฏฺฐิยา วิโมจยํ จิตฺตํ. วิจิกิจฺฉาย วิโมจยํ จิตฺตํ. ถินโต วิโมจยํ จิตฺตํ. อุทฺธจฺจโต วิโมจยํ จิตฺตํ. อหิริกโต วิโมจยํ จิตฺตํ. อโนตฺตปฺปโต วิโมจยํ จิตฺตํ อสฺสสิสฺสามี”ติ สิกฺขติ, “อโนตฺตปฺปโต วิโมจยํ จิตฺตํ ปสฺสสิสฺสามี”ติ สิกฺขติ. วิโมจยํ จิตฺตํ อสฺสาสปสฺสาส\-วเสน วิญฺญาณํ จิตฺตํ, อุปฏฺฐานํ สติ ฯเปฯ.}

\prose{อนุปสฺสตีติ กถํ ตํ จิตฺตํ อนุปสฺสติ ฯเปฯ. เอวํ ตํ จิตฺตํ อนุปสฺสติ. ภาวนาติ จตสฺโส พฺหาวนา ฯเปฯ อาเสวนฏฺเฐน ภาวนา. วิโมจยํ จิตฺตํ อสฺสาส\-ปสฺสาสานํ สํวรฏฺเฐน สีลวิสุทฺธิ ฯเปฯ. วิโมจยํ จิตฺตํ อสฺสาสปสฺสาส\-วเสน จิตฺตสฺส เอกคฺคตํ อวิกฺเขปํ ปชานโต ฯเปฯ ปชานนฺโต อินฺทฺริยานิ สโมธาเนติ, เตน วุจฺจติ “สมตฺถญฺจ ปฏิวิชฺฌตี”ติ. (4)}

\prose{อฏฺฐ อนุปสฺสนา\-ญาณานิ, อฏฺฐ จ อุปฏฺฐานา\-นุสฺสติโย, จตฺตาริ สุตฺตนฺติก\-วตฺถูนิ จิตฺเต จิตฺตา\-นุปสฺสนาย.}

\csromanheader{2}{จตุตฺถ\-จตุกฺก\-นิทฺเทส}
\proseitem{180}{กถํ “อนิจฺจานุปสฺสี อสฺสสิสฺสามี”ติ สิกฺขติ, “อนิจฺจานุปสฺสี ปสฺสสิสฺสามี”ติ สิกฺขติ? อนิจฺจนฺติ กิํ อนิจฺจํ? ปญฺจกฺขนฺธา อนิจฺจา. เกนฏฺเฐน อนิจฺจา? อุปฺปาท\-วยฏฺเฐน อนิจฺจา. ปญฺจนฺนํ ขนฺธานํ อุทยํ ปสฺสนฺโต กติ ลกฺขณานิ ปสฺสติ, วยํ ปสฺสนฺโต กติ ลกฺขณานิ ปสฺสติ, อุทยพฺพยํ ปสฺสนฺโต กติ ลกฺขณานิ ปสฺสติ? ปญฺจนฺนํ ขนฺธานํ อุทยํ ปสฺสนฺโต ปญฺจวีสติ ลกฺขณานิ ปสฺสติ, วยํ ปสฺสนฺโต ปญฺจวีสติ ลกฺขณานิ ปสฺสติ, ปญฺจนฺนํ ขนฺธานํ อุทยพฺพยํ ปสฺสนฺโต อิมานิ ปญฺญาส ลกฺขณานิ ปสฺสติ.}

\prose{“รูเป อนิจฺจานุปสฺสี อสฺสสิสฺสามี”ติ สิกฺขติ, “รูเป อนิจฺจานุปสฺสี ปสฺสสิสฺสามี”ติ สิกฺขติ, “เวทนาย ฯเปฯ. สญฺญาย. สงฺขาเรสุ. วิญฺญาเณ. จกฺขุสฺมิํ ฯเปฯ. ชรามรเณ อนิจฺจานุปสฺสี อสฺสสิสฺสามี”ติ สิกฺขติ, \csromanfolio{192}“ชรามรเณ อนิจฺจานุปสฺสี ปสฺสสิสฺสามี”ติ สิกฺขติ. อนิจฺจานุปสฺสี อสฺสาสปสฺสาส\-วเสน ธมฺมา, อุปฏฺฐานํ สติ, อนุปสฺสนา ญาณํ. ธมฺมา อุปฏฺฐานํ, โน สติ, สติ อุปฏฺฐานญฺเจว สติ จ. ตาย สติยา เตน ญาเณน เต ธมฺเม อนุปสฺสติ, เตน วุจฺจติ “ธมฺเมสุ ธมฺมานุปสฺสนา\-สติปฏฺฐาน\-ภาวนา”ติ.}

\prose{อนุปสฺสตีติ กถํ เต ธมฺเม อนุปสฺสติ ฯเปฯ. เอวํ เต ธมฺเม อนุปสฺสติ. ภาวนาติ จตสฺโส พฺหาวนา ฯเปฯ อาเสวนฏฺเฐน ภาวนา อนิจฺจานุปสฺสี อสฺสาส\-ปสฺสาสานํ สํวรฏฺเฐน สีลวิสุทฺธิ ฯเปฯ. อนิจฺจานุปสฺสี อสฺสาสปสฺสาส\-วเสน จิตฺตสฺส เอกคฺคตํ อวิกฺเขปํ ปชานโต ฯเปฯ. ปชานนฺโต อินฺทฺริยานิ สโมธาเนติ, เตน วุจฺจติ “สมตฺถญฺจ ปฏิวิชฺฌตี”ติ. (1)}

\prose{กถํ “วิราคานุปสฺสี อสฺสสิสฺสามี”ติ สิกฺขติ, “วิราคานุปสฺสี ปสฺสสิสฺสามี”ติ สิกฺขติ? รูเป อาทีนวํ ทิสฺวา รูปวิราเค ฉนฺทชาโต โหติ สทฺธาธิมุตฺโต, จิตฺตญฺจสฺส สฺวาธิฏฺฐิตํ. “รูเป วิราคานุปสฺสี อสฺสสิสฺสามี”ติ สิกฺขติ, “รูเป วิราคานุปสฺสี ปสฺสสิสฺสามี”ติ สิกฺขติ, เวทนาย ฯเปฯ สญฺญาย. สงฺขาเรสุ. วิญฺญาเณ. จกฺขุสฺมิํ ฯเปฯ. ชรามรเณ อาทีนวํ ทิสฺวา ชรามรณ\-วิราเค ฉนฺทชาโต โหติ สทฺธาธิมุตฺโต, จิตฺตญฺจสฺส สฺวาธิฏฺฐิตํ. “ชรามรเณ วิราคานุปสฺสี อสฺสสิสฺสามี”ติ สิกฺขติ, “ชรามรเณ วิราคานุปสฺสี ปสฺสสิสฺสามี”ติ สิกฺขติ. วิราคานุปสฺสี อสฺสาสปสฺสาส\-วเสน ธมฺมา, อุปฏฺฐานํ สติ, อนุปสฺสนา ญาณํ. ธมฺมา อุปฏฺฐานํ, โน สติ, สติ อุปฏฺฐานญฺเจว สติ จ. ตาย สติยา เตน ญาเณน เต ธมฺเม อนุปสฺสติ, เตน วุจฺจติ “ธมฺเมสุ ธมฺมานุปสฺสนา\-สติปฏฺฐาน\-ภาวนา”ติ.}

\prose{อนุปสฺสตีติ กถํ เต ธมฺเม อนุปสฺสติ ฯเปฯ. เอวํ เต ธมฺเม อนุปสฺสติ. ภาวนาติ จตสฺโส พฺหาวนา ฯเปฯ อาเสวนฏฺเฐน ภาวนา, วิราคานุปสฺสี อสฺสาส\-ปสฺสาสานํ สํวรฏฺเฐน สีลวิสุทฺธิ ฯเปฯ. วิราคานุปสฺสี อสฺสาสปสฺสาส\-วเสน จิตฺตสฺส เอกคฺคตํ อวิกฺเขปํ ปชานโต ฯเปฯ ปชานนฺโต อินฺทฺริยานิ สโมธาเนติ, เตน วุจฺจติ “สมตฺถญฺจ ปฏิวิชฺฌตี”ติ. (2)}

\prose{\csromanfolio{193}กถํ “นิโรธานุปสฺสี อสฺสสิสฺสามี”ติ สิกฺขติ, “นิโรธานุปสฺสี ปสฺสสิสฺสามี”ติ สิกฺขติ? รูเป อาทีนวํ ทิสฺวา รูปนิโรเธ ฉนฺทชาโต โหติ สทฺธาธิมุตฺโต, จิตฺตญฺจสฺส สฺวาธิฏฺฐิตํ. “รูเป นิโรธานุปสฺสี อสฺสสิสฺสามี”ติ สิกฺขติ, “รูเป นิโรธานุปสฺสี ปสฺสสิสฺสามี”ติ สิกฺขติ, เวทนาย ฯเปฯ สญฺญาย. สงฺขาเรสุ. วิญฺญาเณ. จกฺขุสฺมิํ ฯเปฯ. ชรามรเณ อาทีนวํ ทิสฺวา ชรามรณ\-นิโรเธ ฉนฺทชาโต โหติ สทฺธาธิมุตฺโต, จิตฺตญฺจสฺส สฺวาธิฏฺฐิตํ. “ชรามรเณ นิโรธานุปสฺสี อสฺสสิสฺสามี”ติ สิกฺขติ, “ชรามรเณ นิโรธานุปสฺสี ปสฺสสิสฺสามี”ติ สิกฺขติ.}

\proseitem{181}{กติหากาเรหิ อวิชฺชาย อาทีนโว โหติ, กติหากาเรหิ อวิชฺชา นิรุชฺฌติ? ปญฺจหากาเรหิ อวิชฺชาย อาทีนโว โหติ, อฏฺฐหากาเรหิ อวิชฺชา นิรุชฺฌติ.}

\prose{กตเมหิ ปญฺจหากาเรหิ อวิชฺชาย อาทีนโว โหติ? อนิจฺจฏฺเฐน อวิชฺชาย อาทีนโว โหติ, ทุกฺขฏฺเฐน อวิชฺชาย อาทีนโว โหติ, อนตฺตฏฺเฐน อวิชฺชาย อาทีนโว โหติ, สนฺตาปฏฺเฐน อวิชฺชาย อาทีนโว โหติ, วิปริณาม\-ฏฺเฐน อวิชฺชาย อาทีนโว โหติ, อิเมหิ ปญฺจหากาเรหิ อวิชฺชาย อาทีนโว โหติ.}

\prose{กตเมหิ อฏฺฐหากาเรหิ อวิชฺชา นิรุชฺฌติ? นิทาน\-นิโรเธน อวิชฺชา นิรุชฺฌติ, สมุทย\-นิโรเธน อวิชฺชา นิรุชฺฌติ, ชาตินิโรเธน อวิชฺชา นิรุชฺฌติ, ปภว\-นิโรเธน\csromansharedfootnote{78528e18bf382bfb}{อาหารนิโรเธน (สฺยา)} อวิชฺชา นิรุชฺฌติ, เหตุนิโรเธน อวิชฺชา นิรุชฺฌติ, ปจฺจย\-นิโรเธน อวิชฺชา นิรุชฺฌติ, ญาณุปฺปาเทน อวิชฺชา นิรุชฺฌติ, นิโรธุ\-ปฏฺฐาเนน อวิชฺชา นิรุชฺฌติ, อิเมหิ อฏฺฐหากาเรหิ อวิชฺชา นิรุชฺฌติ. อิเมหิ ปญฺจหากาเรหิ อวิชฺชาย อาทีนวํ ทิสฺวา อิเมหิ อฏฺฐหากาเรหิ อวิชฺชานิโรเธ ฉนฺทชาโต โหติ สทฺธาธิมุตฺโต, จิตฺตญฺจสฺส สฺวาธิฏฺฐิตํ. “อวิชฺชาย นิโรธานุปสฺสี อสฺสสิสฺสามี”ติ สิกฺขติ, “อวิชฺชาย นิโรธานุปสฺสี ปสฺสสิสฺสามี”ติ สิกฺขติ.}

\prose{กติหากาเรหิ สงฺขาเรสุ อาทีนโว โหติ, กติหากาเรหิ สงฺขารา นิรุชฺฌนฺติ ฯเปฯ. กติหากาเรหิ วิญฺญาเณ อาทีนโว โหติ, กติหากาเรหิ วิญฺญาณํ นิรุชฺฌติ. กติหากาเรหิ นามรูเป \csromanfolio{194}อาทีนโว โหติ, กติหากาเรหิ นามรูปํ นิรุชฺฌติ. กติหากาเรหิ สฬายตเน อาทีนโว โหติ, กติหากาเรหิ สฬายตนํ นิรุชฺฌติ. กติหากาเรหิ ผสฺเส อาทีนโว โหติ, กติหากาเรหิ ผสฺโส นิรุชฺฌติ. กติหากาเรหิ เวทนาย อาทีนโว โหติ, กติหากาเรหิ เวทนา นิรุชฺฌติ. กติหากาเรหิ ตณฺหาย อาทีนโว โหติ, กติหากาเรหิ ตณฺหา นิรุชฺฌติ. กติหากาเรหิ อุปาทาเน อาทีนโว โหติ, กติหากาเรหิ อุปาทานํ นิรุชฺฌติ. กติหากาเรหิ ภเว อาทีนโว โหติ, กติหากาเรหิ ภโว นิรุชฺฌติ. กติหากาเรหิ ชาติยา อาทีนโว โหติ, กติหากาเรหิ ชาติ นิรุชฺฌติ. กติหากาเรหิ ชรามรเณ อาทีนโว โหติ, กติหากาเรหิ ชรามรณํ นิรุชฺฌติ? ปญฺจหากาเรหิ ชรามรเณ อาทีนโว โหติ, อฏฺฐหากาเรหิ ชรามรณํ นิรุชฺฌติ.}

\prose{กตเมหิ ปญฺจหากาเรหิ ชรามรเณ อาทินโว โหติ? อนิจฺจฏฺเฐน ชรามรเณ อาทีนโว โหติ, ทุกฺขฏฺเฐน ฯเปฯ อนตฺตฏฺเฐน ฯเปฯ สนฺตาปฏฺเฐน ฯเปฯ วิปริณาม\-ฏฺเฐน ชรามรเณ อาทีนโว โหติ. อิเมหิ ปญฺจหากาเรหิ ชรามรเณ อาทีนโว โหติ.}

\prose{กตเมหิ อฏฺฐหากาเรหิ ชรามรณํ นิรุชฺฌติ? นิทาน\-นิโรเธน ชรามรณํ นิรุชฺฌติ, สมุทย\-นิโรเธน ฯเปฯ ชาตินิโรเธน ฯเปฯ ปภว\-นิโรเธน\csromansharedfootnote{d456a48095298172}{ภวนิโรเธน (สฺยา)}. เหตุนิโรเธน. ปจฺจย\-นิโรเธน. ญาณุปฺปาเทน ฯเปฯ. นิโรธุ\-ปฏฺฐาเนน ชรามรณํ นิรุชฺฌติ. อิเมหิ อฏฺฐหากาเรหิ ชรามรณํ นิรุชฺฌติ. อิเมหิ ปญฺจหากาเรหิ ชรามรเณ อาทีนวํ ทิสฺวา อิเมหิ อฏฺฐหากาเรหิ ชรามรณ\-นิโรเธ ฉนฺทชาโต โหติ สทฺธาธิมุตฺโต, จิตฺตญฺจสฺส สฺวาธิฏฺฐิตํ. “ชรามรเณ นิโรธานุปสฺสี อสฺสสิสฺสามี”ติ สิกฺขติ, “ชรามรเณ นิโรธานุปสฺสี ปสฺสสิสฺสามี”ติ สิกฺขติ. นิโรธานุปสฺสี อสฺสาสปสฺสาส\-วเสน ธมฺมา, อุปฏฺฐานํ สติ, อนุปสฺสนา ญาณํ. ธมฺมา อุปฏฺฐานํ, โน สติ, สติ อุปฏฺฐานญฺเจว สติ จ. ตาย สติยา เตน ญาเณน เต ธมฺเม อนุปสฺสติ, เตน วุจฺจติ “ธมฺเมสุ ธมฺมานุปสฺสนา\-สติปฏฺฐาน\-ภาวนา”ติ.}

\prose{\csromanfolio{195}อนุปสฺสตีติ กถํ เต ธมฺเม อนุปสฺสติ ฯเปฯ. เอวํ เต ธมฺเม อนุปสฺสติ. ภาวนาติ จตสฺโส พฺหาวนา ฯเปฯ อาเสวนฏฺเฐน ภาวนา. นิโรธานุปสฺสี อสฺสาส\-ปสฺสาสานํ สํวรฏฺเฐน สีลวิสุทฺธิ ฯเปฯ นิโรธานุปสฺสี อสฺสาสปสฺสาส\-วเสน จิตฺตสฺส เอกคฺคตํ อวิกฺเขปํ ปชานโต ฯเปฯ ปชานนฺโต อินฺทฺริยานิ สโมธาเนติ, เตน วุจฺจติ “สมตฺถญฺจ ปฏิวิชฺฌตี”ติ. (3)}

\proseitem{182}{กถํ “ปฏินิสฺสคฺคา\-นุปสฺสี อสฺสสิสฺสามี”ติ สิกฺขติ, “ปฏินิสฺสคฺคา\-นุปสฺสี ปสฺสสิสฺสามี”ติ สิกฺขติ? ปฏินิสฺสคฺคาติ เทฺว ปฏินิสฺสคฺคา ปริจฺจาค\-ปฏินิสฺสคฺโค จ ปกฺขนฺทน\-ปฏินิสฺสคฺโค จ. รูปํ ปริจฺจชตีติ ปริจฺจาค\-ปฏินิสฺสคฺโค, รูปนิโรเธ นิพฺพาเน จิตฺตํ ปกฺขนฺทตีติ ปกฺขนฺทน\-ปฏินิสฺสคฺโค. “รูเป ปฏินิสฺสคฺคา\-นุปสฺสี อสฺสสิสฺสามี”ติ สิกฺขติ, “รูเป ปฏินิสฺสคฺคา\-นุปสฺสี ปสฺสสิสฺสามี”ติ สิกฺขติ. เวทนํ ฯเปฯ. สญฺญํ. สงฺขาเร. วิญฺญาณํ. จกฺขุํ ฯเปฯ. ชรามรณํ ปริจฺจชตีติ ปริจฺจาค\-ปฏินิสฺสคฺโค, ชรามรณ\-นิโรเธ นิพฺพาเน จิตฺตํ ปกฺขนฺทตีติ ปกฺขนฺทน\-ปฏินิสฺสคฺโค, “ชรามรเณ ปฏินิสฺสคฺคา\-นุปสฺสี อสฺสสิสฺสามี”ติ สิกฺขติ, “ชรามรเณ ปฏินิสฺสคฺคา\-นุปสฺสี ปสฺสสิสฺสามี”ติ สิกฺขติ, ปฏินิสฺสคฺคา\-นุปสฺสี อสฺสาสปสฺสาส\-วเสน ธมฺมา, อุปฏฺฐานํ สติ, อนุปสฺสนา ญาณํ. ธมฺมา อุปฏฺฐานํ, โน สติ, สติ อุปฏฺฐานญฺเจว สติ จ. ตาย สติยา เตน ญาเณน เต ธมฺเม อนุปสฺสติ, เตน วุจฺจติ “ธมฺเมสุ ธมฺมานุปสฺสนา\-สติปฏฺฐาน\-ภาวนา”ติ.}

\prose{อนุปสฺสตีติ กถํ เต ธมฺเม อนุปสฺสติ อนิจฺจโต อนุปสฺสติ, โน นิจฺจโต ฯเปฯ ปฏินิสฺสชฺชติ, โน อาทิยติ. อนิจฺจโต อนุปสฺสนฺโต นิจฺจสญฺญํ ปชหติ ฯเปฯ ปฏินิสฺสชฺชนฺโต อาทานํ ปชหติ, เอวํ เต ธมฺเม อนุปสฺสติ. พฺหาวนาติ จตสฺโส ภาวนา. ตตฺถ ชาตานํ ธมฺมานํ อนติวตฺตน\-ฏฺเฐน ภาวนา ฯเปฯ อาเสวนฏฺเฐน ภาวนา, ปฏินิสฺสคฺคา\-นุปสฺสี อสฺสาส\-ปสฺสาสานํ สํวรฏฺเฐน สีลวิสุทฺธิ, อวิกฺเขป\-ฏฺเฐน จิตฺตวิสุทฺธิ, ทสฺสนฏฺเฐน ทิฏฺฐิวิสุทฺธิ, โย ตตฺถ สํวรฏฺโฐ, อยํ อธิสีลสิกฺขา. โย ตตฺถ อวิกฺเขปฏฺโฐ, อยํ อธิจิตฺต\-สิกฺขา. โย ตตฺถ ทสฺสนฏฺโฐ, อยํ อธิปญฺญา\-สิกฺขา. อิมา ติสฺโส สิกฺขาโย อาวชฺชนฺโต สิกฺขติ, ชานนฺโต สิกฺขติ ฯเปฯ สจฺฉิกาตพฺพํ สจฺฉิกโรนฺโต สิกฺขติ.}

\prose{ปฏินิสฺสคฺคา\-นุปสฺสี อสฺสาสปสฺสาส\-วเสน จิตฺตสฺส เอกคฺคตํ อวิกฺเขปํ ปชานโต วิทิตา เวทนา อุปชฺชนฺติ, วิทิตา อุปฏฺฐหนฺติ, วิทิตา อพฺภตฺถํ \csromanfolio{196}คจฺฉนฺติ ฯเปฯ. ปฏินิสฺสคฺคา\-นุปสฺสี อสฺสาสปสฺสาส\-วเสน จิตฺตสฺส เอกคฺคตํ อวิกฺเขปํ ปชานนฺโต อินฺทฺริยานิ สโมธาเนติ, โคจรญฺจ ปชานาติ, สมตฺถญฺจ ปฏิวิชฺฌติ, พลานิ สโมธาเนติ, โพชฺฌงฺเค สโมธาเนติ, มคฺคํ สโมธาเนติ, ธมฺเม สโมธาเนติ, โคจรญฺจ ปชานาติ, สมตฺถญฺจ ปฏิวิชฺฌติ.}

\prose{อินฺทฺริยานิ สโมธาเนตีติ กถํ อินฺทฺริยานิ สโมธาเนติ, อธิโมกฺข\-ฏฺเฐน สทฺธินฺทฺริยํ สโมธาเนติ ฯเปฯ เตน วุจฺจติ “สมตฺถญฺจ ปฏิวิชฺฌตี”ติ. (4)}

\prosewithcloserruled{อฏฺฐ อนุปสฺสเน ญาณานิ, อฏฺฐ จ อุปฏฺฐานา\-นุสฺสติโย, จตฺตาริ สุตฺตนฺติก\-วตฺถูนิ ธมฺเมสุ ธมฺมา\-นุปสฺสนาย. อิมานิ พาตฺติํส สโตการิสฺส ญาณานิ.}{สโตการิญาณนิทฺเทโส ปญฺจโม.}
\csromanmatikaanchor{mtk.71}
\csromantocmarkref{subsection}{6. ญาณราสิฉกฺกนิทฺเทส}{mtk.71}
\bookmark[dest={mtk.71},level=2]{6. ญาณราสิฉกฺกนิทฺเทส}
\csromanheader{2}{6. ญาณ\-ราสิ\-ฉกฺก\-นิทฺเทส}
\proseitem{183}{กตมานิ จตุวีสติ สมาธิวเสน ญาณานิ? ทีฆํ อสฺสาสวเสน จิตฺตสฺส เอกคฺคตา อวิกฺเขโป สมาธิ, ทีฆํ ปสฺสาสวเสน จิตฺตสฺส เอกคฺคตา อวิกฺเขโป สมาธิ ฯเปฯ. วิโมจยํ จิตฺตํ อสฺสาสวเสน จิตฺตสฺส เอกคฺคตา อวิกฺเขโป สมาธิ. วิโมจยํ จิตฺตํ ปสฺสาสวเสน จิตฺตสฺส เอกคฺคตา อวิกฺเขโป สมาธิ. อิมานิ จตุวีสติ สมาธิวเสน ญาณานิ. (1)}

\prose{กตมานิ เทฺวสตฺตติ วิปสฺสนา\-วเสน ญาณานิ? ทีฆํ อสฺสาสํ อนิจฺจโต อนุปสฺสนฏฺเฐน วิปสฺสนา, ทุกฺขโต อนุปสฺสนฏฺเฐน วิปสฺสนา, อนตฺตโต อนุปสฺสนฏฺเฐน วิปสฺสนา. ทีฆํ ปสฺสาสํ อนิจฺจโต อนุปสฺสนฏฺเฐน วิปสฺสนา, ทุกฺขโต อนุปสฺสนฏฺเฐน วิปสฺสนา, อนตฺตโต อนุปสฺสนฏฺเฐน วิปสฺสนา ฯเปฯ วิโมจยํ จิตฺตํ อสฺสาสํ. วิโมจยํ จิตฺตํ ปสฺสาสํ อนิจฺจโต อนุปสฺสนฏฺเฐน วิปสฺสนา, ทุกฺขโต อนุปสฺสนฏฺเฐน วิปสฺสนา, อนตฺตโต อนุปสฺสนฏฺเฐน วิปสฺสนา. อิมานิ เทฺวสตฺตติ วิปสฺสนา\-วเสน ญาณานิ. (2)}

\prose{\csromanfolio{197}กตมานิ อฏฺฐ นิพฺพิทาญาณานิ? อนิจฺจานุปสฺสี อสฺสสํ ยถาภูตํ ชานาติ ปสฺสตีติ นิพฺพิทาญาณํ, อนิจฺจานุปสฺสี ปสฺสสํ ยถาภูตํ ชานาติ ปสฺสตีติ นิพฺพิทาญาณํ ฯเปฯ ปฏินิสฺสคฺคา\-นุปสฺสี อสฺสสํ ยถาภูตํ ชานาติ ปสฺสตีติ นิพฺพิทาญาณํ, ปฏินิสฺสคฺคา\-นุปสฺสี ปสฺสสํ ยถาภูตํ ชานาติ ปสฺสตีติ นิพฺพิทาญาณํ. อิมานิ อฏฺฐ นิพฺพิทาญาณานิ. (3)}

\prose{กตมานิ อฏฺฐ นิพฺพิทานุโลเม ญาณานิ? อนิจฺจานุปสฺสี อสฺสสํ ภยตุปฏฺฐาเน ปญฺญา นิพฺพิทานุโลเม ญาณํ. อนิจฺจานุปสฺสี ปสฺสสํ ภยตุปฏฺฐาเน ปญฺญา นิพฺพิทานุโลเม ญาณํ ฯเปฯ. ปฏินิสฺสคฺคา\-นุปสฺสี อสฺสสํ ภยตุปฏฺฐาเน ปญฺญา นิพฺพิทานุโลเม ญาณํ. ปฏินิสฺสคฺคา\-นุปสฺสี ปสฺสสํ ภยตุปฏฺฐาเน ปญฺญา นิพฺพิทานุโลเม ญาณํ. อิมานิ อฏฺฐ นิพฺพิทานุโลเม ญาณานิ. (4)}

\prose{กตมานิ อฏฺฐ นิพฺพิทา\-ปฏิปฺปสฺสทฺธิ\-ญาณานิ? อนิจฺจานุปสฺสี อสฺสสํ ปฏิสงฺขา สนฺติฏฺฐนา ปญฺญา นิพฺพิทา\-ปฏิปฺปสฺสทฺธิ\-ญาณํ. อนิจฺจานุปสฺสี ปสฺสสํ ปฏิสงฺขา สนฺติฏฺฐนา ปญฺญา นิพฺพิทา\-ปฏิปฺปสฺสทฺธิ\-ญาณํ ฯเปฯ. ปฏินิสฺสคฺคา\-นุปสฺสี อสฺสสํ ปฏิสงฺขา สนฺติฏฺฐนา ปญฺญา นิพฺพิทา\-ปฏิปฺปสฺสทฺธิ\-ญาณํ. ปฏินิสฺสคฺคา\-นุปสฺสี ปสฺสสํ ปฏิสงฺขา สนฺติฏฺฐนา ปญฺญา นิพฺพิทา\-ปฏิปฺปสฺสทฺธิ\-ญาณํ. อิมานิ อฏฺฐ นิพฺพิทา\-ปฏิปฺปสฺสทฺธิ\-ญาณานิ. (5)}

\prose{กตมานิ เอกวีสติ วิมุตฺติสุเข ญาณานิ? โสตาปตฺติ\-มคฺเคน สกฺกาย\-ทิฏฺฐิยา ปหีนตฺตา สมุจฺฉินฺนตฺตา อุปฺปชฺชติ วิมุตฺติสุเข ญาณํ. วิจิกิจฺฉาย ปหีนตฺตา สมุจฺฉินฺนตฺตา อุปฺปชฺชติ วิมุตฺติสุเข ญาณํ. สีลพฺพต\-ปรามาสสฺส ฯเปฯ. ทิฏฺฐา\-นุสยสฺส วิจิกิจฺฉา\-นุสยสฺส ปหีนตฺตา สมุจฺฉินฺนตฺตา อุปฺปชฺชติ วิมุตฺติสุเข ญาณํ. สกทาคามิ\-มคฺเคน โอฬาริกสฺส กามราค\-สญฺโญชนสฺส ฯเปฯ. ปฏิฆ\-สญฺโญชนสฺส โอฬาริกสฺส กาม\-ราคา\-นุสยสฺส ปฏิฆา\-นุสยสฺส ปหีนตฺตา สมุจฺฉินฺนตฺตา อุปฺปชฺชติ วิมุตฺติสุเข ญาณํ. อนาคามิมคฺเคน อณุสหคตสฺส กามราค\-สญฺโญชนสฺส ฯเปฯ. ปฏิฆ\-สญฺโญชนสฺส อณุสหคตสฺส กาม\-ราคา\-นุสยสฺส ปฏิฆา\-นุสยสฺส ปหีนตฺตา สมุจฺฉินฺนตฺตา อุปฺปชฺชติ วิมุตฺติสุเข ญาณํ. อรหตฺต\-มคฺเคน รูปราคสฺส ฯเปฯ. อรูปราคสฺส มานสฺส อุทฺธจฺจสฺส อวิชฺชาย มานานุสยสฺส ภวราคา\-นุสยสฺส อวิชฺชา\-นุสยสฺส ปหีนตฺตา สมุจฺฉินฺนตฺตา อุปฺปชฺชติ}

\prosewithcloserruled{\csromanfolio{198}วิมุตฺติสุเข ญาณํ. อิมานิ เอกวีสติ วิมุตฺติสุเข ญาณานิ. โสฬส\-วตฺถุกํ อานาปาน\-สฺสติ\-สมาธิํ ภาวยโต สมธิกานิ อิมานิ เทฺว ญาณสตานิ อุปฺปชฺชนฺติ. (6)}{อานาปานสฺสติ\-กถา\csromansharedfootnote{575e9a2e1d1e29b6}{อานาปานกถา (สฺยา, ก)} นิฏฺฐิตา.}
\csromanmatikaanchor{mtk.72}
\csromantocmarknopage{section}{4. อินฺทฺริยกถา}{mtk.72}
\bookmark[dest={mtk.72},level=1]{4. อินฺทฺริยกถา}
\csromanheader{1}{4. \textbf{อินฺทฺริยกถา}}
\csromanmatikaanchor{mtk.73}
\csromantocmarkref{subsection}{1. ปฐมสุตฺตนฺตนิทฺเทส}{mtk.73}
\bookmark[dest={mtk.73},level=2]{1. ปฐมสุตฺตนฺตนิทฺเทส}
\csromanheader{2}{1. ปฐม\-สุตฺตนฺต\-นิทฺเทส}
\proseitem{184}{เอวํ เม สุตํ\csromandash{}เอกํ สมยํ ภควา สาวตฺถิยํ วิหรติ เชตวเน อนาถ\-ปิณฺฑิกสฺส อาราเม. ตตฺร โข ภควา ภิกฺขู อามนฺเตสิ “ภิกฺขโว”ติ. “ภทนฺเต”ติ\csromansharedfootnote{c532c828f0cc1381}{“ภทฺทนฺเต”ติ (ก)} เต ภิกฺขู ภควโต ปจฺจสฺโสสุํ. ภควา เอตทโวจ\csromandash{}}

\prose{ปญฺจิมานิ ภิกฺขเว อินฺทฺริยานิ. กตมานิ ปญฺจ? สทฺธินฺทฺริยํ วีริยินฺทฺริยํ สตินฺทฺริยํ สมาธินฺทฺริยํ ปญฺญินฺทฺริยํ. อิมานิ โข ภิกฺขเว ปญฺจินฺทฺริยานิ\csromansharedfootnote{e5cc754ca7b750ba}{สํ 3. 170 ปิฏฺเฐปิ.}.}

\proseitem{185}{อิมานิ ปญฺจินฺทฺริยานิ กติหากาเรหิ วิสุชฺฌนฺติ? อิมานิ ปญฺจินฺทฺริยานิ ปนฺนรสหิ อากาเรหิ วิสุชฺฌนฺติ. อสฺสทฺเธ ปุคฺคเล ปริวชฺชยโต สทฺเธ ปุคฺคเล เสวโต ภชโต ปยิรุปาสโต ปสาทนีเย สุตฺตนฺเต ปจฺจเวกฺขโต อิเมหิ ตีหากาเรหิ สทฺธินฺทฺริยํ วิสุชฺฌติ.  กุสีเต ปุคฺคเล ปริวชฺชยโต อารทฺธวีริเย ปุคฺคเล เสวโต ภชโต ปยิรุปาสโต สมฺมปฺปธาเน ปจฺจเวกฺขโต อิเมหิ ตีหากาเรหิ วีริยินฺทฺริยํ วิสุชฺฌติ.  มุฏฺฐสฺสตี ปุคฺคเล ปริวชฺชยโต อุปฏฺฐิตสฺสตี ปุคฺคเล เสวโต ภชโต ปยิรุปาสโต สติปฏฺฐาเน ปจฺจเวกฺขโต อิเมหิ ตีหากาเรหิ สตินฺทฺริยํ วิสุชฺฌติ.  อสมาหิเต ปุคฺคเล ปริวชฺชยโต สมาหิเต ปุคฺคเล เสวโต ภชโต ปยิรุปาสโต ฌานวิโมกฺเข ปจฺจเวกฺขโต อิเมหิ ตีหากาเรหิ สมาธินฺทฺริยํ วิสุชฺฌติ.  ทุปฺปญฺเญ ปุคฺคเล ปริวชฺชยโต ปญฺญวนฺเต ปุคฺคเล เสวโต ภชโต ปยิรุปาสโต คมฺภีร\-ญาณจริยํ ปจฺจเวกฺขโต อิเมหิ ตีหากาเรหิ ปญฺญินฺทฺริยํ วิสุชฺฌติ. อิติ อิเม ปญฺจ ปุคฺคเล ปริวชฺชยโต ปญฺจ}

\prose{\csromanfolio{199}ปุคฺคเล เสวโต ภชโต ปยิรุปาสโต ปญฺจ สุตฺตนฺต\-กฺขนฺเธ ปจฺจเวกฺขโต อิเมหิ ปนฺนรสหิ อากาเรหิ อิมานิ ปญฺจินฺทฺริยานิ วิสุชฺฌนฺติ.}

\prose{กติหากาเรหิ ปญฺจินฺทฺริยานิ ภาวิยนฺติ, กติหากาเรหิ ปญฺจนฺนํ อินฺทฺริยานํ ภาวนา โหติ? ทสหากาเรหิ ปญฺจินฺทฺริยานิ ภาวิยนฺติ, ทสหากาเรหิ ปญฺจนฺนํ อินฺทฺริยานํ ภาวนา โหติ. อสฺสทฺธิยํ ปชหนฺโต สทฺธินฺทฺริยํ ภาเวติ, สทฺธินฺทฺริยํ ภาเวนฺโต อสฺสทฺธิยํ ปชหติ, โกสชฺชํ ปชหนฺโต วีริยินฺทฺริยํ ภาเวติ, วีริยินฺทฺริยํ ภาเวนฺโต โกสชฺชํ ปชหติ, ปมาทํ ปชหนฺโต สตินฺทฺริยํ ภาเวติ, สตินฺทฺริยํ ภาเวนฺโต ปมาทํ ปชหติ, อุทฺธจฺจํ ปชหนฺโต สมาธินฺทฺริยํ ภาเวติ, สมาธินฺทฺริยํ ภาเวนฺโต อุทฺธจฺจํ ปชหติ, อวิชฺชํ ปชหนฺโต ปญฺญินฺทฺริยํ ภาเวติ, ปญฺญินฺทฺริยํ ภาเวนฺโต อวิชฺชํ ปชหติ. อิเมหิ ทสหากาเรหิ ปญฺจินฺทฺริยานิ ภาวิยนฺติ, อิเมหิ ทสหากาเรหิ ปญฺจนฺนํ อินฺทฺริยานํ ภาวนา โหติ.}

\prose{กติหากาเรหิ ปญฺจินฺทฺริยานิ ภาวิตานิ โหนฺติ สุภาวิตานิ. ทสหากาเรหิ ปญฺจินฺทฺริยานิ ภาวิตานิ โหนฺติ สุภาวิตานิ. อสฺสทฺธิยสฺส ปหีนตฺตา สุปฺปหีนตฺตา สทฺธินฺทฺริยํ ภาวิตํ โหติ สุภาวิตํ, สทฺธิ\-นฺทฺริยสฺส ภาวิตตฺตา สุภาวิตตฺตา อสฺสทฺธิยํ ปหีนํ โหติ สุปฺปหีนํ, โกสชฺชสฺส ปหีนตฺตา สุปฺปหีนตฺตา วีริยินฺทฺริยํ ภาวิตํ โหติ สุภาวิตํ, วีริยิ\-นฺทฺริยสฺส ภาวิตตฺตา สุภาวิตตฺตา โกสชฺชํ ปหีนํ โหติ สุปฺปหีนํ, ปมาทสฺส ปหีนตฺตา สุปฺปหีนตฺตา สตินฺทฺริยํ ภาวิตํ โหติ สุภาวิตํ, สตินฺทฺริยสฺส ภาวิตตฺตา สุภาวิตตฺตา ปมาโท ปหีโน โหติ สุปฺปหีโน, อุทฺธจฺจสฺส ปหีนตฺตา สุปฺปหีนตฺตา สมาธินฺทฺริยํ ภาวิตํ โหติ สุภาวิตํ, สมาธิ\-นฺทฺริยสฺส ภาวิตตฺตา สุภาวิตตฺตา อุทฺธจฺจํ ปหีนํ โหติ สุปฺปหีนํ, อวิชฺชาย ปหีนตฺตา ปญฺญินฺทฺริยํ ภาวิตํ โหติ สุภาวิตํ, ปญฺญินฺทฺริยสฺส ภาวิตตฺตา สุภาวิตตฺตา อวิชฺชา ปหีนา โหติ สุปฺปหีนา. อิเมหิ ทสหากาเรหิ ปญฺจินฺทฺริยานิ ภาวิตานิ โหนฺติ สุภาวิตานิ.}

\proseitem{186}{กติหากาเรหิ ปญฺจินฺทฺริยานิ ภาวิยนฺติ, กติหากาเรหิ ปญฺจินฺทฺริยานิ ภาวิตานิ เจว โหนฺติ สุภาวิตานิ จ ปฏิปฺปสฺสทฺธานิ จ สุปฺปฏิปฺปสฺสทฺธานิ จ? จตูหากาเรหิ ปญฺจินฺทฺริยานิ ภาวิยนฺติ, จตูหากาเรหิ ปญฺจินฺทฺริยานิ ภาวิตานิ เจว โหนฺติ สุภาวิตานิ จ ปฏิปฺปสฺสทฺธานิ จ สุปฺปฏิปฺปสฺสทฺธานิ จ. โสตาปตฺติมคฺค\-กฺขเณ ปญฺจินฺทฺริยานิ ภาวิยนฺติ, \csromanfolio{200}โสตาปตฺติ\-ผลกฺขเณ ปญฺจินฺทฺริยานิ ภาวิตานิ เจว โหนฺติ สุภาวิตานิ จ ปฏิปฺปสฺสทฺธานิ จ สุปฺปฏิปฺปสฺสทฺธานิ จ. สกทาคามิ\-มคฺคกฺขเณ ปญฺจินฺทฺริยานิ ภาวิยนฺติ, สกทาคามิ\-ผลกฺขเณ ปญฺจินฺทฺริยานิ ภาวิตานิ เจว โหนฺติ สุภาวิตานิ จ ปฏิปฺปสฺสทฺธานิ จ สุปฺปฏิปฺปสฺสทฺธานิ จ. อนาคามิ\-มคฺคกฺขเณ ปญฺจินฺทฺริยานิ ภาวิยนฺติ, อนาคามิ\-ผลกฺขเณ ปญฺจินฺทฺริยานิ ภาวิตานิ เจว โหนฺติ สุภาวิตานิ จ ปฏิปฺปสฺสทฺธานิ จ สุปฺปฏิปฺปสฺสทฺธานิ จ. อรหตฺต\-มคฺคกฺขเณ ปญฺจินฺทฺริยานิ ภาวิยนฺติ, อรหตฺต\-ผลกฺขเณ ปญฺจินฺทฺริยานิ ภาวิตานิ เจว โหนฺติ สุภาวิตานิ จ ปฏิปฺปสฺสทฺธานิ จ สุปฺปฏิปฺปสฺสทฺธานิ จ. อิติ จตสฺโส มคฺค\-วิสุทฺธิโย จตสฺโส ผลวิสุทฺธิโย จตสฺโส สมุจฺเฉท\-วิสุทฺธิโย จตสฺโส ปฏิปฺปสฺสทฺธิ\-วิสุทฺธิโย. อิเมหิ จตูหากาเรหิ ปญฺจินฺทฺริยานิ ภาวิยนฺติ, อิเมหิ จตูหากาเรหิ ปญฺจินฺทฺริยานิ ภาวิตานิ เจว โหนฺติ สุภาวิตานิ จ ปฏิปฺปสฺสทฺธานิ จ สุปฺปฏิปฺปสฺสทฺธานิ จ.}

\prosewithcloserruled{กตินํ ปุคฺคลานํ อินฺทฺริยภาวนา, กติ ปุคฺคลา ภาวิตินฺทฺริยา? อฏฺฐนฺนํ ปุคฺคลานํ อินฺทฺริยภาวนา, ตโย ปุคฺคลา ภาวิตินฺทฺริยา.  กตเมสํ อฏฺฐนฺนํ ปุคฺคลานํ อินฺทฺริยภาวนา? สตฺตนฺนญฺจ เสกฺขานํ ปุถุชฺชน\-กลฺยาณกสฺส จ, อิเมสํ อฏฺฐนฺนํ ปุคฺคลานํ อินฺทฺริยภาวนา.  กตเม ตโย ปุคฺคลา ภาวิตินฺทฺริยา, สวเนน พุทฺโธ ตถาคตสฺส สาวโก ขีณาสโว ภาวิตินฺทฺริโย, สยํ ภูตฏฺเฐน ปจฺเจก\-สมฺพุทฺโธ ภาวิตินฺทฺริโย, อปฺปเมยฺย\-ฏฺเฐน ตถาคโต อรหํ สมฺมาสมฺพุทฺโธ ภาวิตินฺทฺริโย. อิเม ตโย ปุคฺคลา ภาวิตินฺทฺริยา. อิติ อิเมสํ อฏฺฐนฺนํ ปุคฺคลานํ อินฺทฺริยภาวนา. อิเม ตโย ปุคฺคลา ภาวิตินฺทฺริยา.}{สุตฺตนฺต\-นิทฺเทโส ปฐโม.}
\csromanmatikaanchor{mtk.74}
\csromantocmarkref{subsection}{2. ทุติยสุตฺตนฺตนิทฺเทส}{mtk.74}
\bookmark[dest={mtk.74},level=2]{2. ทุติยสุตฺตนฺตนิทฺเทส}
\csromanheader{2}{2. ทุติย\-สุตฺตนฺต\-นิทฺเทส}
\proseitem{187}{สาวตฺถิ\-นิทานํ.  ปญฺจิมานิ ภิกฺขเว อินฺทฺริยานิ. กตมานิ ปญฺจ? สทฺธินฺทฺริยํ วีริยินฺทฺริยํ สตินฺทฺริยํ สมาธินฺทฺริยํ ปญฺญินฺทฺริยํ. เย หิ เกจิ ภิกฺขเว สมณา วา พฺราหฺมณา วา อิเมสํ ปญฺจนฺนํ อินฺทฺริยานํ สมุทยญฺจ อตฺถงฺคมญฺจ อสฺสาทญฺจ อาทีนวญฺจ นิสฺสรณญฺจ ยถาภูตํ นปฺปชานนฺติ, น เม เต ภิกฺขเว สมณา วา พฺราหฺมณา วา สมเณสุ วา\csromansharedfootnote{92e6ab4cec54c8e2}{เจว (ก) สํ 3. 171 ปิฏฺเฐ ปสฺสิตพฺพา.} สมณสมฺมตา,}

\prose{\csromanfolio{201}พฺราหฺมเณสุ วา พฺราหฺมณ\-สมฺมตา, น จ ปเนเต อายสฺมนฺตา\csromansharedfootnote{13cb9116dbd4a4e4}{อายสฺมนฺโต สํ 1. 433 ปิฏฺเฐ; สํ 3. 171 ปิฏฺเฐปิ.} สามญฺญตฺถํ วา พฺรหฺมญฺญตฺถํ วา ทิฏฺเฐว ธมฺเม สยํ อภิญฺญา สจฺฉิกตฺวา อุปสมฺปชฺช วิหรนฺติ. เย จ โข เกจิ\csromansharedfootnote{15240d6da46de396}{เย เกจิ โข (สฺยา), เย หิ เกจิ โข (ก) สํ 3. 171 ปิฏฺเฐ ปสฺสิตพฺพา.} ภิกฺขเว สมณา วา พฺราหฺมณา วา อิเมสํ ปญฺจนฺนํ อินฺทฺริยานํ สมุทยญฺจ อตฺถงฺคมญฺจ อสฺสาทญฺจ อาทีนวญฺจ นิสฺสรณญฺจ ยถาภูตํ ปชานนฺติ, เต โข เม ภิกฺขเว สมณา วา พฺราหฺมณา วา สมเณสุ วา สมณสมฺมตา, พฺราหฺมเณสุ วา พฺราหฺมณ\-สมฺมตา, เต จ ปนายสฺมนฺตา สามญฺญตฺถญฺจ พฺราหฺมญฺญตฺถญฺจ ทิฏฺเฐว ธมฺเม สยํ อภิญฺญา สจฺฉิกตฺวา อุปสมฺปชฺช วิหรนฺตีติ.}

\proseitem{188}{กติหากาเรหิ ปญฺจนฺนํ อินฺทฺริยานํ สมุทโย โหติ, กติหากาเรหิ ปญฺจนฺนํ อินฺทฺริยานํ สมุทยํ ปชานาติ. กติหากาเรหิ ปญฺจนฺนํ อินฺทฺริยานํ อตฺถงฺคโม โหติ, กติหากาเรหิ ปญฺจนฺนํ อินฺทฺริยานํ อตฺถงฺคมํ ปชานาติ. กติหากาเรหิ ปญฺจนฺนํ อินฺทฺริยานํ อสฺสาโท โหติ, กติหากาเรหิ ปญฺจนฺนํ อินฺทฺริยานํ อสฺสาทํ ปชานาติ. กติหากาเรหิ ปญฺจนฺนํ อินฺทฺริยานํ อาทีนโว โหติ, กติหากาเรหิ ปญฺจนฺนํ อินฺทฺริยานํ อาทีนวํ ปชานาติ. กติหากาเรหิ ปญฺจนฺนํ อินฺทฺริยานํ นิสฺสรณํ โหติ, กติหากาเรหิ ปญฺจนฺนํ อินฺทฺริยานํ นิสฺสรณํ ปชานาติ?}

\prose{จตฺตารีสาย อากาเรหิ ปญฺจนฺนํ อินฺทฺริยานํ สมุทโย โหติ, จตฺตารีสาย อากาเรหิ ปญฺจนฺนํ อินฺทฺริยานํ สมุทยํ ปชานาติ. จตฺตารีสาย อากาเรหิ ปญฺจนฺนํ อินฺทฺริยานํ อตฺถงฺคโม โหติ, จตฺตารีสาย อากาเรหิ ปญฺจนฺนํ อินฺทฺริยานํ อตฺถงฺคมํ ปชานาติ. ปญฺจวีสติยา อากาเรหิ ปญฺจนฺนํ อินฺทฺริยานํ อสฺสาโท โหติ, ปญฺจวีสติยา อากาเรหิ ปญฺจนฺนํ อินฺทฺริยานํ อสฺสาทํ ปชานาติ. ปญฺจวีสติยา อากาเรหิ ปญฺจนฺนํ อินฺทฺริยานํ อาทีนโว โหติ, ปญฺจวีสติยา อากาเรหิ ปญฺจนฺนํ อินฺทฺริยานํ อาทีนวํ ปชานาติ. อสีติสตํ อากาเรหิ ปญฺจนฺนํ อินฺทฺริยานํ นิสฺสรณํ โหติ, อสีติสตํ อากาเรหิ\csromansharedfootnote{a774aaa4dbe48876}{อสีติสตากาเรหิ (สฺยา)} ปญฺจนฺนํ อินฺทฺริยานํ นิสฺสรณํ ปชานาติ.}

\prose{\csromanfolio{202}กตเมหิ จตฺตารีสาย อากาเรหิ ปญฺจนฺนํ อินฺทฺริยานํ สมุทโย โหติ, กตเมหิ จตฺตารีสาย อากาเรหิ ปญฺจนฺนํ อินฺทฺริยานํ สมุทยํ ปชานาติ? อธิโมกฺข\-ตฺถาย อาวชฺชนายา สมุทโย สทฺธิ\-นฺทฺริยสฺส สมุทโย โหติ, อธิโมกฺข\-วเสน ฉนฺทสฺส สมุทโย สทฺธิ\-นฺทฺริยสฺส สมุทโย โหติ, อธิโมกฺข\-วเสน มนสิการสฺส สมุทโย สทฺธิ\-นฺทฺริยสฺส สมุทโย โหติ. สทฺธิ\-นฺทฺริยสฺส วเสน เอกตฺตุ\-ปฏฺฐานํ สทฺธิ\-นฺทฺริยสฺส สมุทโย โหติ.  ปคฺคหตฺถาย อาวชฺชนาย สมุทโย วีริยิ\-นฺทฺริยสฺส สมุทโย โหติ, ปคฺคหวเสน ฉนฺทสฺส สมุทโย วีริยิ\-นฺทฺริยสฺส สมุทโย โหติ, ปคฺคหวเสน มนสิการสฺส สมุทโย วีริยิ\-นฺทฺริยสฺส สมุทโย โหติ. วีริยิ\-นฺทฺริยสฺส วเสน เอกตฺตุ\-ปฏฺฐานํ วีริยิ\-นฺทฺริยสฺส สมุทโย โหติ.  อุปฏฺฐาน\-ตฺถาย อาวชฺชนาย สมุทโย สตินฺทฺริยสฺส สมุทโย โหติ, อุปฏฺฐาน\-วเสน ฉนฺทสฺส สมุทโย สตินฺทฺริยสฺส สมุทโย โหติ, อุปฏฺฐาน\-วเสน มนสิการสฺส สมุทโย สตินฺทฺริยสฺส สมุทโย โหติ. สตินฺทฺริยสฺส วเสน เอกตฺตุ\-ปฏฺฐานํ สตินฺทฺริยสฺส สมุทโย โหติ.  อวิกฺเขป\-ตฺถาย อาวชฺชนาย สมุทโย สมาธิ\-นฺทฺริยสฺส สมุทโย โหติ, อวิกฺเขป\-วเสน ฉนฺทสฺส สมุทโย สมาธิ\-นฺทฺริยสฺส สมุทโย โหติ, อวิกฺเขป\-วเสน มนสิการสฺส สมุทโย สมาธิ\-นฺทฺริยสฺส สมุทโย โหติ. สมาธิ\-นฺทฺริยสฺส วเสน เอกตฺตุ\-ปฏฺฐานํ สมาธิ\-นฺทฺริยสฺส สมุทโย โหติ.  ทสฺสนตฺถาย อาวชฺชนาย สมุทโย ปญฺญินฺทฺริยสฺส สมุทโย โหติ, ทสฺสนวเสน ฉนฺทสฺส สมุทโย ปญฺญินฺทฺริยสฺส สมุทโย โหติ, ทสฺสนวเสน มนสิการสฺส สมุทโย ปญฺญินฺทฺริยสฺส สมุทโย โหติ. ปญฺญินฺทฺริยสฺส วเสน เอกตฺตุ\-ปฏฺฐานํ ปญฺญินฺทฺริยสฺส สมุทโย โหติ.}

\prose{อธิโมกฺข\-ตฺถาย อาวชฺชนาย สมุทโย สทฺธิ\-นฺทฺริยสฺส สมุทโย โหติ, ปคฺคหตฺถาย อาวชฺชนาย สมุทโย วีริยิ\-นฺทฺริยสฺส สมุทโย โหติ, อุปฏฺฐาน\-ตฺถาย อาวชฺชนาย สมุทโย สตินฺทฺริยสฺส สมุทโย โหติ. อวิกฺเขป\-ตฺถาย อาวชฺชนาย สมุทโย สมาธิ\-นฺทฺริยสฺส สมุทโย โหติ, ทสฺสนตฺถาย อาวชฺชนาย สมุทโย ปญฺญินฺทฺริยสฺส สมุทโย โหติ.  อธิโมกฺข\-วเสน ฉนฺทสฺส สมุทโย สทฺธิ\-นฺทฺริยสฺส สมุทโย โหติ, ปคฺคหวเสน ฉนฺทสฺส สมุทโย วีริยิ\-นฺทฺริยสฺส สมุทโย โหติ, อุปฏฺฐาน\-วเสน ฉนฺทสฺส สมุทโย สตินฺทฺริยสฺส}

\prose{\csromanfolio{203}สมุทโย โหติ, อวิกฺเขป\-วเสน ฉนฺทสฺส สมุทโย สมาธิ\-นฺทฺริยสฺส สมุทโย โหติ, ทสฺสนวเสน ฉนฺทสฺส สมุทโย ปญฺญินฺทฺริยสฺส สมุทโย โหติ.  อธิโมกฺข\-วเสน มนสิการสฺส สมุทโย สทฺธิ\-นฺทฺริยสฺส สมุทโย โหติ, ปคฺคหวเสน มนสิการสฺส สมุทโย วีริยิ\-นฺทฺริยสฺส สมุทโย โหติ, อุปฏฺฐาน\-วเสน มนสิการสฺส สมุทโย สตินฺทฺริยสฺส สมุทโย โหติ, อวิกฺเขป\-วเสน มนสิการสฺส สมุทโย สมาธิ\-นฺทฺริยสฺส สมุทโย โหติ, ทสฺสนวเสน มนสิการสฺส สมุทโย ปญฺญินฺทฺริยสฺส สมุทโย โหติ, สทฺธิ\-นฺทฺริยสฺส วเสน เอกตฺตุ\-ปฏฺฐานํ สทฺธิ\-นฺทฺริยสฺส สมุทโย โหติ, วีริยิ\-นฺทฺริยสฺส วเสน เอกตฺตุ\-ปฏฺฐานํ วีริยิ\-นฺทฺริยสฺส สมุทโย โหติ, สตินฺทฺริยสฺส วเสน เอกตฺตุ\-ปฏฺฐานํ สตินฺทฺริยสฺส สมุทโย โหติ, สมาธิ\-นฺทฺริยสฺส วเสน เอกตฺตุ\-ปฏฺฐานํ สมาธิ\-นฺทฺริยสฺส สมุทโย โหติ, ปญฺญินฺทฺริยสฺส วเสน เอกตฺตุ\-ปฏฺฐานํ ปญฺญินฺทฺริยสฺส สมุทโย โหติ.}

\prose{อิเมหิ จตฺตารีสาย อากาเรหิ ปญฺจนฺนํ อินฺทฺริยานํ สมุทโย โหติ, อิเมหิ จตฺตารีสาย อากาเรหิ ปญฺจนฺนํ อินฺทฺริยานํ สมุทยํ ปชานาติ.}

\prose{กตเมหิ จตฺตารีสาย อากาเรหิ ปญฺจนฺนํ อินฺทฺริยานํ อตฺถงฺคโม โหติ, กตเมหิ จตฺตารีสาย อากาเรหิ ปญฺจนฺนํ อินฺทฺริยานํ อตฺถงฺคมํ ปชานาติ? อธิโมกฺข\-ตฺถาย อาวชฺชนาย อตฺถงฺคโม สทฺธิ\-นฺทฺริยสฺส อตฺถงฺคโม โหติ, อธิโมกฺข\-วเสน ฉนฺทสฺส อตฺถงฺคโม สทฺธิ\-นฺทฺริยสฺส อตฺถงฺคโม โหติ, อธิโมกฺข\-วเสน มนสิการสฺส อตฺถงฺคโม สทฺธิ\-นฺทฺริยสฺส อตฺถงฺคโม โหติ. สทฺธิ\-นฺทฺริยสฺส วเสน เอกตฺต\-อนุ\-ปฏฺฐานํ\csromansharedfootnote{06ee186022390ef3}{เอกตฺตํ อนุปฏฺฐานํ (สฺยา)} สทฺธิ\-นฺทฺริยสฺส วเสน อตฺถงฺคโม โหติ.  ปคฺคหตฺถาย อาวชฺชนาย อตฺถงฺคโม วีริยิ\-นฺทฺริยสฺส อตฺถงฺคโม โหติ, ปคฺคหวเสน ฉนฺทสฺส อตฺถงฺคโม วีริยิ\-นฺทฺริยสฺส อตฺถงฺคโม โหติ. ปคฺคหวเสน มนสิการสฺส อตฺถงฺคโม วีริยิ\-นฺทฺริยสฺส อตฺถงฺคโม โหติ, วีริยิ\-นฺทฺริยสฺส วเสน เอกตฺต\-อนุ\-ปฏฺฐานํ วีริยิ\-นฺทฺริยสฺส อตฺถงฺคโม โหติ.  อุปฏฺฐาน\-ตฺถาย อาวชฺชนาย อตฺถงฺคโม สตินฺทฺริยสฺส อตฺถงฺคโม โหติ, อุปฏฺฐาน\-วเสน ฉนฺทสฺส อตฺถงฺคโม สตินฺทฺริยสฺส อตฺถงฺคโม โหติ, อุปฏฺฐาน\-วเสน มนสิการสฺส อตฺถงฺคโม สตินฺทฺริยสฺส อตฺถงฺคโม โหติ. สตินฺทฺริยสฺส}

\prose{\csromanfolio{204}วเสน เอกตฺต\-อนุ\-ปฏฺฐานํ สตินฺทฺริยสฺส อตฺถงฺคโม โหติ.  อวิกฺเขป\-ตฺถาย อาวชฺชนาย อตฺถงฺคโม สมาธิ\-นฺทฺริยสฺส อตฺถงฺคโม โหติ, อวิกฺเขป\-วเสน ฉนฺทสฺส อตฺถงฺคโม สมาธิ\-นฺทฺริยสฺส อตฺถงฺคโม โหติ, อวิกฺเขป\-วเสน มนสิการสฺส อตฺถงฺคโม สมาธิ\-นฺทฺริยสฺส อตฺถงฺคโม โหติ. สมาธิ\-นฺทฺริยสฺส วเสน เอกตฺต\-อนุ\-ปฏฺฐานํ สมาธิ\-นฺทฺริยสฺส อตฺถงฺคโม โหติ.  ทสฺสนตฺถาย อาวชฺชนาย อตฺถงฺคโม ปญฺญินฺทฺริยสฺส อตฺถงฺคโม โหติ, ทสฺสนวเสน ฉนฺทสฺส อตฺถงฺคโม ปญฺญินฺทฺริยสฺส อตฺถงฺคโม โหติ, ทสฺสนวเสน มนสิการสฺส อตฺถงฺคโม ปญฺญินฺทฺริยสฺส อตฺถงฺคโม โหติ. ปญฺญินฺทฺริยสฺส วเสน เอกตฺต\-อนุ\-ปฏฺฐานํ ปญฺญินฺทฺริยสฺส อตฺถงฺคโม โหติ.}

\prose{อธิโมกฺข\-ตฺถาย อาวชฺชนาย อตฺถงฺคโม สทฺธิ\-นฺทฺริยสฺส อตฺถงฺคโม โหติ, ปคฺคหตฺถาย อาวชฺชนาย อตฺถงฺคโม วีริยิ\-นฺทฺริยสฺส อตฺถงฺคโม โหติ, อุปฏฺฐาน\-ตฺถาย อาวชฺชนาย อตฺถงฺคโม สตินฺทฺริยสฺส อตฺถงฺคโม โหติ, อวิกฺเขป\-ตฺถาย อาวชฺชนาย อตฺถงฺคโม สมาธิ\-นฺทฺริยสฺส อตฺถงฺคโม โหติ, ทสฺสนตฺถาย อาวชฺชนาย อตฺถงฺคโม ปญฺญินฺทฺริยสฺส อตฺถงฺคโม โหติ.  อธิโมกฺข\-วเสน ฉนฺทสฺส อตฺถงฺคโม สทฺธิ\-นฺทฺริยสฺส อตฺถงฺคโม โหติ, ปคฺคหวเสน ฉนฺทสฺส อตฺถงฺคโม วีริยิ\-นฺทฺริยสฺส อตฺถงฺคโม โหติ, อุปฏฺฐาน\-วเสน ฉนฺทสฺส อตฺถงฺคโม สตินฺทฺริยสฺส อตฺถงฺคโม โหติ, อวิกฺเขป\-วเสน ฉนฺทสฺส อตฺถงฺคโม สมาธิ\-นฺทฺริยสฺส อตฺถงฺคโม โหติ, ทสฺสนวเสน ฉนฺทสฺส อตฺถงฺคโม ปญฺญินฺทฺริยสฺส อตฺถงฺคโม โหติ.  อธิโมกฺข\-วเสน มนสิการสฺส อตฺถงฺคโม สทฺธิ\-นฺทฺริยสฺส อตฺถงฺคโม โหติ, ปคฺคหวเสน มนสิการสฺส อตฺถงฺคโม วีริยิ\-นฺทฺริยสฺส อตฺถงฺคโม โหติ, อุปฏฺฐาน\-วเสน มนสิการสฺส อตฺถงฺคโม สตินฺทฺริยสฺส อตฺถงฺคโม โหติ, อวิกฺเขป\-วเสน มนสิการสฺส อตฺถงฺคโม สมาธิ\-นฺทฺริยสฺส อตฺถงฺคโม โหติ, ทสฺสนวเสน มนสิการสฺส อตฺถงฺคโม ปญฺญินฺทฺริยสฺส อตฺถงฺคโม โหติ.  สทฺธิ\-นฺทฺริยสฺส วเสน เอกตฺต\-อนุ\-ปฏฺฐานํ สทฺธิ\-นฺทฺริยสฺส อตฺถงฺคโม โหติ. วีริยิ\-นฺทฺริยสฺส วเสน เอกตฺต\-อนุ\-ปฏฺฐานํ วีริยิ\-นฺทฺริยสฺส อตฺถงฺคโม โหติ. สตินฺทฺริยสฺส วเสน เอกตฺต\-อนุ\-ปฏฺฐานํ สตินฺทฺริยสฺส อตฺถงฺคโม โหติ. สมาธิ\-นฺทฺริยสฺส วเสน เอกตฺต\-อนุ\-ปฏฺฐานํ สมาธิ\-นฺทฺริยสฺส อตฺถงฺคโม โหติ. ปญฺญินฺทฺริยสฺส วเสน เอกตฺต\-อนุ\-ปฏฺฐานํ ปญฺญินฺทฺริยสฺส อตฺถงฺคโม โหติ.}

\prose{\csromanfolio{205}อิเมหิ จตฺตารีสาย อากาเรหิ ปญฺจนฺนํ อินฺทฺริยานํ อตฺถงฺคโม โหติ. อิเมหิ จตฺตารีสาย อากาเรหิ ปญฺจนฺนํ อินฺทฺริยานํ อตฺถงฺคมํ ปชานาติ.}

\csromanheader{2}{\textbf{ก. อสฺสาทนิทฺเทส}}
\proseitem{189}{กตเมหิ ปญฺจวีสติยา อากาเรหิ ปญฺจนฺนํ อินฺทฺริยานํ อสฺสาโท โหติ, กตเมหิ ปญฺจวีสติยา อากาเรหิ ปญฺจนฺนํ อินฺทฺริยานํ อสฺสาทํ ปชานาติ? อสฺสทฺธิยสฺส อนุปฏฺฐานํ สทฺธิ\-นฺทฺริยสฺส อสฺสาโท โหติ. อสฺสทฺธิย\-ปริฬาหสฺส อนุปฏฺฐานํ สทฺธิ\-นฺทฺริยสฺส อสฺสาโท โหติ. อธิโมกฺข\-จริยาย เวสารชฺชํ สทฺธิ\-นฺทฺริยสฺส อสฺสาโท โหติ. สนฺโต จ วิหาราธิคโม สทฺธิ\-นฺทฺริยสฺส อสฺสาโท โหติ. ยํ สทฺธินฺทฺริยํ ปฏิจฺจ อุปฺปชฺชติ สุขํ โสมนสฺสํ, อยํ สทฺธิ\-นฺทฺริยสฺส อสฺสาโท.}

\prose{โกสชฺชสฺส อนุปฏฺฐานํ วีริยิ\-นฺทฺริยสฺส อสฺสาโท โหติ. โกสชฺช\-ปริฬาหสฺส อนุปฏฺฐานํ วีริยิ\-นฺทฺริยสฺส อสฺสาโท โหติ. ปคฺคห\-จริยาย เวสารชฺชํ วีริยิ\-นฺทฺริยสฺส อสฺสาโท โหติ. สนฺโต จ วิหาราธิคโม วีริยิ\-นฺทฺริยสฺส อสฺสาโท โหติ. ยํ วีริยินฺทฺริยํ ปฏิจฺจ อุปฺปชฺชติ สุขํ โสมนสฺสํ, อยํ วีริยิ\-นฺทฺริยสฺส อสฺสาโท.}

\prose{ปมาทสฺส อนุปฏฺฐานํ สตินฺทฺริยสฺส อสฺสาโท โหติ. ปมาท\-ปริฬาหสฺส อนุปฏฺฐานํ สตินฺทฺริยสฺส อสฺสาโท โหติ. อุปฏฺฐาน\-จริยาย เวสารชฺชํ สตินฺทฺริยสฺส อสฺสาโท โหติ. สนฺโต จ วิหาราธิคโม สตินฺทฺริยสฺส อสฺสาโท โหติ. ยํ สตินฺทฺริยํ ปฏิจฺจ อุปฺปชฺชติ สุขํ โสมนสฺสํ, อยํ สตินฺทฺริยสฺส อสฺสาโท.}

\prose{อุทฺธจฺจสฺส อนุปฏฺฐานํ สมาธิ\-นฺทฺริยสฺส อสฺสาโท โหติ. อุทฺธจฺจ\-ปริฬาหสฺส อนุปฏฺฐานํ สมาธิ\-นฺทฺริยสฺส อสฺสาโท โหติ. อวิกฺเขป\-จริยาย เวสารชฺชํ สมาธิ\-นฺทฺริยสฺส อสฺสาโท โหติ. สนฺโต จ วิหาราธิคโม สมาธิ\-นฺทฺริยสฺส อสฺสาโท โหติ. ยํ สมาธินฺทฺริยํ ปฏิจฺจ อุปฺปชฺชติ สุขํ โสมนสฺส, อยํ สมาธิ\-นฺทฺริยสฺส อสฺสาโท.}

\prose{อวิชฺชาย อนุปฏฺฐานํ ปญฺญินฺทฺริยสฺส อสฺสาโท โหติ. อวิชฺชา\-ปริฬาหสฺส อนุปฏฺฐานํ ปญฺญินฺทฺริยสฺส อสฺสาโท โหติ. ทสฺสน\-จริยาย เวสารชฺชํ ปญฺญินฺทฺริยสฺส อสฺสาโท โหติ. สนฺโต จ วิหาราธิคโม \csromanfolio{206}ปญฺญินฺทฺริยสฺส อสฺสาโท โหติ. ยํ ปญฺญินฺทฺริยํ ปฏิจฺจ อุปฺปชฺชติ สุขํ โสมนสฺสํ, อยํ ปญฺญินฺทฺริยสฺส อสฺสาโท.}

\prose{อิเมหิ ปญฺจวีสติยา อากาเรหิ ปญฺจนฺนํ อินฺทฺริยานํ อสฺสาโท โหติ, อิเมหิ ปญฺจวีสติยา อากาเรหิ ปญฺจนฺนํ อินฺทฺริยานํ อสฺสาทํ ปชานาติ.}

\csromanheader{2}{\textbf{ข. อาทีนวนิทฺเทส}}
\proseitem{190}{กตเมหิ ปญฺจวีสติยา อากาเรหิ ปญฺจนฺนํ อินฺทฺริยานํ อาทีนโว โหติ, กตเมหิ ปญฺจวีสติยา อากาเรหิ ปญฺจนฺนํ อินฺทฺริยานํ อาทีนวํ ปชานาติ? อสฺสทฺธิยสฺส อุปฏฺฐานํ สทฺธิ\-นฺทฺริยสฺส อาทีนโว โหติ. อสฺสทฺธิย\-ปริฬาหสฺส อุปฏฺฐานํ สทฺธิ\-นฺทฺริยสฺส อาทีนโว โหติ. อนิจฺจฏฺเฐน สทฺธิ\-นฺทฺริยสฺส อาทีนโว โหติ. ทุกฺขฏฺเฐน สทฺธิ\-นฺทฺริยสฺส อาทีนโว โหติ. อนตฺตฏฺเฐน สทฺธิ\-นฺทฺริยสฺส อาทีนโว โหติ.}

\prose{โกสชฺชสฺส อุปฏฺฐานํ วีริยิ\-นฺทฺริยสฺส อาทีนโว โหติ. โกสชฺช\-ปริฬาหสฺส อุปฏฺฐานํ วีริยิ\-นฺทฺริยสฺส อาทีนโว โหติ. อนิจฺจฏฺเฐน วีริยิ\-นฺทฺริยสฺส อาทีนโว โหติ. ทุกฺขฏฺเฐน ฯเปฯ. อนตฺตฏฺเฐน วีริยินฺทฺริสฺส อาทีนโว โหติ.}

\prose{ปมาทสฺส อุปฏฺฐานํ สตินฺทฺริยสฺส อาทีนโว โหติ. ปมาท\-ปริฬาหสฺส อุปฏฺฐานํ สตินฺทฺริยสฺส อาทีนโว โหติ. อนิจฺจฏฺเฐน สตินฺทฺริยสฺส อาทีนโว โหติ. ทุกฺขฏฺเฐน ฯเปฯ. อนตฺตฏฺเฐน สตินฺทฺริยสฺส อาทีนโว โหติ.}

\prose{อุทฺธจฺจสฺส อุปฏฺฐานํ สมาธิ\-นฺทฺริยสฺส อาทีนโว โหติ. อุทฺธจฺจ\-ปริฬาหสฺส อุปฏฺฐานํ สมาธิ\-นฺทฺริยสฺส อาทีนโว โหติ. อนิจฺจฏฺเฐน สมาธิ\-นฺทฺริยสฺส อาทีนโว โหติ. ทุกฺขฏฺเฐน ฯเปฯ. อนตฺตฏฺเฐน สมาธิ\-นฺทฺริยสฺส อาทีนโว โหติ.}

\prose{อวิชฺชาย อุปฏฺฐานํ ปญฺญินฺทฺริยสฺส อาทีนโว โหติ. อวิชฺชา\-ปริฬาหสฺส อุปฏฺฐานํ ปญฺญินฺทฺริยสฺส อาทีนโว โหติ. อนิจฺจฏฺเฐน ปญฺญินฺทฺริยสฺส อาทีนโว โหติ. ทุกฺขฏฺเฐน ปญฺญินฺทฺริยสฺส อาทีนโว โหติ. อนตฺตฏฺเฐน ปญฺญินฺทฺริยสฺส อาทีนโว โหติ.}

\prose{\csromanfolio{207}อิเมหิ ปญฺจวีสติยา อากาเรหิ ปญฺจนฺนํ อินฺทฺริยานํ อาทีนโว โหติ, อิเมหิ ปญฺจวีสติยา อากาเรหิ ปญฺจนฺนํ อินฺทฺริยานํ อาทีนวํ ปชานาติ.}

\csromanheader{2}{8. นิสฺสรณ\-นิทฺเทส}
\proseitem{191}{กตเมหิ อสีติสตํ อากาเรหิ ปญฺจนฺนํ อินฺทฺริยานํ นิสฺสรณํ โหติ, กตเมหิ อสีติสตํ อากาเรหิ ปญฺจนฺนํ อินฺทฺริยานํ นิสฺสรณํ ปชานาติ? อธิโมกฺข\-ฏฺเฐน สทฺธินฺทฺริยํ อสฺสทฺธิยา นิสฺสฏํ โหติ. อสฺสทฺธิย\-ปริฬาหา นิสฺสฏํ โหติ, ตทนุวตฺตก\-กิเลเสหิ จ ขนฺเธหิ จ นิสฺสฏํ โหติ, พหิทฺธา จ สพฺพนิมิตฺเตหิ นิสฺสฏํ โหติ, ตโต ปณีตตร\-สทฺธิ\-นฺทฺริยสฺส ปฏิลาภา ปุริมตร\-สทฺธินฺทฺริยา นิสฺสฏํ โหติ. ปคฺคหฏฺเฐน วีริยินฺทฺริยํ โกสชฺชา นิสฺสฏํ โหติ, โกสชฺช\-ปริฬาหา นิสฺสฏํ โหติ, ตทนุวตฺตก\-กิเลเสหิ จ ขนฺเธหิ จ นิสฺสฏํ โหติ, พหิทฺธา จ สพฺพนิมิตฺเตหิ นิสฺสฏํ โหติ, ตโต ปณีตตร\-วีริยิ\-นฺทฺริยสฺส ปฏิลาภา ปุริมตร\-วีริยินฺทฺริยา นิสฺสฏํ โหติ. อุปฏฺฐาน\-ฏฺเฐน สตินฺทฺริยํ ปมาทา นิสฺสฏํ โหติ, ปมาทปริฬาหา นิสฺสฏํ โหติ, ตทนุวตฺตก\-กิเลเสหิ จ ขนฺเธหิ จ นิสฺสฏํ โหติ, พหิทฺธา จ สพฺพนิมิตฺเตหิ นิสฺสฏํ โหติ, ตโต ปณีตตร\-สตินฺทฺริยสฺส ปฏิลาภา ปุริมตร\-สตินฺทฺริยา นิสฺสฏํ โหติ. อวิกฺเขป\-ฏฺเฐน สมาธินฺทฺริยํ อุทฺธจฺจา นิสฺสฏํ โหติ, อุทฺธจฺจ\-ปริฬาหา นิสฺสฏํ โหติ, ตทนุวตฺตก\-กิเลเสหิ จ ขนฺเธหิ จ นิสฺสฏํ โหติ, พหิทฺธา จ สพฺพนิมิตฺเตหิ นิสฺสฏํ โหติ, ตโต ปณีตตร\-สมาธิ\-นฺทฺริยสฺส ปฏิลาภา ปุริมตร\-สมาธินฺทฺริยา นิสฺสฏํ โหติ. ทสฺสนฏฺเฐน ปญฺญินฺทฺริยํ อวิชฺชาย นิสฺสฏํ โหติ, อวิชฺชาปริฬาหา นิสฺสฏํ โหติ, ตทนุวตฺตก\-กิเลเสหิ จ ขนฺเธหิ จ นิสฺสฏํ โหติ, พหิทฺธา จ สพฺพนิมิตฺเตหิ นิสฺสฏํ โหติ, ตโต ปณีตตร\-ปญฺญินฺทฺริยสฺส ปฏิลาภา ปุริมตร\-ปญฺญินฺทฺริยา นิสฺสฏํ โหติ.}

\proseitem{192}{ปุพฺพภาเค ปญฺจหิ อินฺทฺริเยหิ ปฐมชฺฌาน\-วเสน ปญฺจินฺทฺริยานิ นิสฺสฏานิ โหนฺติ, ปฐเม ฌาเน ปญฺจหิ\-นฺทฺริเยหิ ทุติยชฺฌาน\-วเสน ปญฺจินฺทฺริยานิ นิสฺสฏานิ โหนฺติ. ทุติเย ฌาเน ปญฺจหิ\-นฺทฺริเยหิ ตติยชฺฌาน\-วเสน ปญฺจินฺทฺริยานิ นิสฺสฏานิ โหนฺติ, ตติเย ฌาเน ปญฺจหิ\-นฺทฺริเยหิ จตุตฺถชฺฌาน\-วเสน ปญฺจินฺทฺริยานิ นิสฺสาฏานิ โหนฺติ, จตุตฺเถ ฌาเน ปญฺจหิ\-นฺทฺริเยหิ อากาสา\-นญฺจา\-ยตนสมาปตฺติวเสน ปญฺจินฺทฺริยานิ นิสฺสฏานิ \csromanfolio{208}โหนฺติ, อากาสา\-นญฺจา\-ยตนสมาปตฺติยา ปญฺจหิ\-นฺทฺริเยหิ วิญฺญาณญฺจายตน\-สมาปตฺติวเสน ปญฺจินฺทฺริยานิ นิสฺสฏานิ โหนฺติ, วิญฺญาณญฺจายตน\-สมาปตฺติยา ปญฺจหิ\-นฺทฺริเยหิ อากิญฺจญฺญายตน\-สมาปตฺติวเสน ปญฺจินฺทฺริยานิ นิสฺสฏานิ โหนฺติ, อากิญฺจญฺญายตน\-สมาปตฺติยา ปญฺจหิ\-นฺทฺริเยหิ เนว\-สญฺญา\-นา\-สญฺญา\-ยตนสมาปตฺติวเสน ปญฺจินฺทฺริยานิ นิสฺสฏานิ โหนฺติ, เนว\-สญฺญา\-นา\-สญฺญา\-ยตนสมาปตฺติยา ปญฺจหิ\-นฺทฺริเยหิ อนิจฺจา\-นุปสฺสนาวเสน ปญฺจินฺทฺริยานิ นิสฺสฏานิ โหนฺติ. (8)}

\prose{อนิจฺจา\-นุปสฺสนาย ปญฺจหิ\-นฺทฺริเยหิ ทุกฺขานุปสฺสนา\-วเสน ปญฺจินฺทฺริยานิ นิสฺสฏานิ โหนฺติ, ทุกฺขา\-นุปสฺสนาย ปญฺจหิ\-นฺทฺริเยหิ อนตฺตา\-นุปสฺสนาวเสน ปญฺจินฺทฺริยานิ นิสฺสฏานิ โหนฺติ, อนตฺตา\-นุปสฺสนาย ปญฺจหิ\-นฺทฺริเยหิ นิพฺพิทา\-นุปสฺสนาวเสน ปญฺจินฺทฺริยานิ นิสฺสฏานิ โหนฺติ, นิพฺพิทา\-นุปสฺสนาย ปญฺจหิ\-นฺทฺริเยหิ วิราคา\-นุปสฺสนาวเสน ปญฺจินฺทฺริยานิ นิสฺสฏานิ โหนฺติ, วิราคา\-นุปสฺสนาย ปญฺจหิ\-นฺทฺริเยหิ นิโรธา\-นุปสฺสนาวเสน ปญฺจินฺทฺริยานิ นิสฺสฏานิ โหนฺติ, นิโรธา\-นุปสฺสนาย ปญฺจหิ\-นฺทฺริเยหิ ปฏินิสฺสคฺคา\-นุปสฺสนาวเสน ปญฺจินฺทฺริยานิ นิสฺสฏานิ โหนฺติ, ปฏินิสฺสคฺคา\-นุปสฺสนาย ปญฺจหิ\-นฺทฺริเยหิ ขยานุปสฺสนา\-วเสน ปญฺจินฺทฺริยานิ นิสฺสฏานิ โหนฺติ, ขยานุปสฺสนาย ปญฺจหิ\-นฺทฺริเยหิ วยา\-นุปสฺสนา\-วเสน ปญฺจินฺทฺริยานิ นิสฺสฏานิ โหนฺติ, วยา\-นุปสฺสนาย ปญฺจหิ\-นฺทฺริเยหิ วิปริณามา\-นุปสฺสนา\-วเสน ปญฺจินฺทฺริยานิ นิสฺสฏานิ โหนฺติ, วิปริณามา\-นุปสฺสนาย ปญฺจหิ\-นฺทฺริเยหิ อนิมิตฺตา\-นุปสฺสนา\-วเสน ปญฺจินฺทฺริยานิ นิสฺสฏานิ โหนฺติ, อนิมิตฺตา\-นุปสฺสนาย ปญฺจหิ\-นฺทฺริเยหิ อปฺปณิหิตา\-นุปสฺสนา\-วเสน ปญฺจินฺทฺริยานิ นิสฺสฏานิ โหนฺติ, อปฺปณิหิตา\-นุปสฺสนาย ปญฺจหิ\-นฺทฺริเยหิ สุญฺญตา\-นุปสฺสนาวเสน ปญฺจินฺทฺริยานิ นิสฺสฏานิ โหนฺติ, สุญฺญตา\-นุปสฺสนาย ปญฺจหิ\-นฺทฺริเยหิ อธิปญฺญา\-ธมฺมวิปสฺสนาวเสน ปญฺจินฺทฺริยานิ นิสฺสฏานิ โหนฺติ. อธิปญฺญา\-ธมฺม\-วิปสฺสนาย ปญฺจหิ\-นฺทฺริเยหิ ยถาภูตญาณ\-ทสฺสนวเสน ปญฺจินฺทฺริยานิ นิสฺสฏานิ โหนฺติ, ยถา\-ภูต\-ญาณ\-ทสฺสเน ปญฺจหิ\-นฺทฺริเยหิ อาทีนวา\-นุปสฺสนาวเสน ปญฺจินฺทฺริยานิ นิสฺสฏานิ โหนฺติ, อาทีนวา\-นุปสฺสนาย ปญฺจหิ\-นฺทฺริเยหิ ปฏิสงฺขา\-นุปสฺสนาวเสน ปญฺจินฺทฺริยานิ นิสฺสฏานิ โหนฺติ, ปฏิสงฺขา\-นุปสฺสนาย ปญฺจหิ\-นฺทฺริเยหิ วิวฏฺฏนา\-นุปสฺสนา\-วเสน ปญฺจินฺทฺริยานิ นิสฺสฏานิ โหนฺติ. วิวฏฺฏนา\-นุปสฺสนาย ปญฺจหิ\-นฺทฺริเยหิ โสตาปตฺติมคฺค\-วเสน ปญฺจินฺทฺริยานิ นิสฺสฏานิ โหนฺติ. \mbox{(18, 26)}}

\prose{\csromanfolio{209}โสตาปตฺติมคฺเค ปญฺจหิ\-นฺทฺริเยหิ โสตาปตฺติผล\-สมาปตฺติวเสน ปญฺจินฺทฺริยานิ นิสฺสฏานิ โหนฺติ, โสตาปตฺติ\-ผล\-สมาปตฺติยา ปญฺจหิ\-นฺทฺริเยหิ สกทาคามิมคฺค\-วเสน ปญฺจินฺทฺริยานิ นิสฺสฏานิ โหนฺติ, สกทาคามิ\-มคฺเค ปญฺจหิ\-นฺทฺริเยหิ สกทาคามิผล\-สมาปตฺติวเสน ปญฺจินฺทฺริยานิ นิสฺสฏานิ โหนฺติ, สกทามิผลสมาปตฺติยา ปญฺจหิ\-นฺทฺริเยหิ อนาคามิมคฺค\-วเสน ปญฺจินฺทฺริยานิ นิสฺสฏานิ โหนฺติ, อนาคามิมคฺเค ปญฺจหิ\-นฺทฺริเยหิ อนาคามิผล\-สมาปตฺติวเสน ปญฺจินฺทฺริยานิ นิสฺสฏานิ โหนฺติ, อนาคามิ\-ผล\-สมาปตฺติยา ปญฺจหิ\-นฺทฺริเยหิ อรหตฺตมคฺค\-วเสน ปญฺจินฺทฺริยานิ นิสฺสฏานิ โหนฺติ, อรหตฺตมคฺเค ปญฺจหิ\-นฺทฺริเยหิ อรหตฺตผล\-สมาปตฺติวเสน ปญฺจินฺทฺริยานิ นิสฺสฏานิ โหนฺติ. \mbox{(8, 34)}}

\prose{เนกฺขมฺเม ปญฺจินฺทฺริยานิ กามจฺฉนฺทโต นิสฺสฏานิ โหนฺติ, อพฺยาปาเท ปญฺจินฺทฺริยานิ พฺยาปาทโต นิสฺสฏานิ โหนฺติ, อาโลกสญฺญาย ปญฺจินฺทฺริยานิ ถินมิทฺธโต นิสฺสฏานิ โหนฺติ, อวิกฺเขเป ปญฺจินฺทฺริยานิ อุทฺธจฺจโต นิสฺสฏานิ โหนฺติ, ธมฺม\-ววตฺถาเน ปญฺจินฺทฺริยานิ วิจิกิจฺฉาย นิสฺสฏานิ โหนฺติ, ญาเณ ปญฺจินฺทฺริยานิ อวิชฺชาย นิสฺสฏานิ โหนฺติ, ปาโมชฺเช ปญฺจินฺทฺริยานิ อรติยา นิสฺสฏานิ โหนฺติ.}

\proseitem{193}{ปฐเม ฌาเน ปญฺจินฺทฺริยานิ นีวรเณหิ นิสฺสฏานิ โหนฺติ, ทุติเย ฌาเน ปญฺจินฺทฺริยานิ วิตกฺกวิจาเรหิ นิสฺสฏานิ โหนฺติ, ตติเย ฌาเน ปญฺจินฺทฺริยานิ ปีติยา นิสฺสฏานิ โหนฺติ, จตุตฺเถ ฌาเน ปญฺจินฺทฺริยานิ สุขทุกฺเขหิ นิสฺสฏานิ โหนฺติ, อากาสา\-นญฺจา\-ยตนสมาปตฺติยา ปญฺจินฺทฺริยานิ รูปสญฺญาย ปฏิฆสญฺญาย นานตฺตสญฺญาย นิสฺสฏานิ โหนฺติ, วิญฺญาณญฺจายตน\-สมาปตฺติยา ปญฺจินฺทฺริยานิ อากาสา\-นญฺจา\-ยตนสญฺญาย นิสฺสฏานิ โหนฺติ, อากิญฺจญฺญายตน\-สมาปตฺติยา ปญฺจินฺทฺริยานิ วิญฺญาณญฺจา\-ยตนสญฺญาย นิสฺสฏานิ โหนฺติ, เนว\-สญฺญา\-นา\-สญฺญา\-ยตนสมาปตฺติยา ปญฺจินฺทฺริยานิ อากิญฺจญฺญายตน\-สญฺญาย นิสฺสฏานิ โหนฺติ.}

\prose{อนิจฺจา\-นุปสฺสนาย ปญฺจินฺทฺริยานิ นิจฺจสญฺญาย นิสฺสฏานิ โหนฺติ, ทุกฺขา\-นุปสฺสนาย ปญฺจินฺทฺริยานิ สุขสญฺญาย นิสฺสฏานิ โหนฺติ, อนตฺตา\-นุปสฺสนาย ปญฺจินฺทฺริยานิ อตฺตสญฺญาย นิสฺสฏานิ โหนฺติ, นิพฺพิทา\-นุปสฺสนาย ปญฺจินฺทฺริยานิ นนฺทิยา นิสฺสฏานิ โหนฺติ, วิราคา\-นุปสฺสนาย ปญฺจินฺทฺริยานิ ราคโต นิสฺสฏานิ โหนฺติ, นิโรธา\-นุปสฺสนาย ปญฺจินฺทฺริยานิ สมุทยโต \csromanfolio{210}นิสฺสฏานิ โหนฺติ, ปฏินิสฺสคฺคา\-นุปสฺสนาย ปญฺจินฺทฺริยานิ อาทานโต นิสฺสฏานิ โหนฺติ, ขยานุปสฺสนาย ปญฺจินฺทฺริยานิ ฆนสญฺญาย นิสฺสฏานิ โหนฺติ, วยา\-นุปสฺสนาย ปญฺจินฺทฺริยานิ อายูหนโต นิสฺสฏานิ โหนฺติ, วิปริณามา\-นุปสฺสนาย ปญฺจินฺทฺริยานิ ธุวสญฺญาย นิสฺสฏานิ โหนฺติ, อนิมิตฺตา\-นุปสฺสนาย ปญฺจินฺทฺริยานิ นิมิตฺตโต นิสฺสฏานิ โหนฺติ, อปฺปณิหิตา\-นุปสฺสนาย ปญฺจินฺทฺริยานิ ปณิธิยา นิสฺสฏานิ โหนฺติ, สุญฺญตา\-นุปสฺสนาย ปญฺจินฺทฺริยานิ อภินิเวสโต นิสฺสฏานิ โหนฺติ. อธิปญฺญา\-ธมฺม\-วิปสฺสนาย ปญฺจินฺทฺริยานิ สาราทานา\-ภินิเวสโต นิสฺสฏานิ โหนฺติ, ยถา\-ภูต\-ญาณ\-ทสฺสเน ปญฺจินฺทฺริยานิ สมฺโมหา\-ภินิเวสโต นิสฺสฏานิ โหนฺติ, อาทีนวา\-นุปสฺสนาย ปญฺจินฺทฺริยานิ อาลยา\-ภินิเวสโต นิสฺสฏานิ โหนฺติ, ปฏิสงฺขา\-นุปสฺสนาย ปญฺจินฺทฺริยานิ อปฺปฏิสงฺขาย นิสฺสฏานิ โหนฺติ, วิวฏฺฏนา\-นุปสฺสนาย ปญฺจินฺทฺริยานิ สญฺโญคา\-ภินิเวสโต นิสฺสฏานิ โหนฺติ.}

\prose{โสตาปตฺติมคฺเค ปญฺจินฺทฺริยานิ ทิฏฺเฐกฏฺเฐหิ กิเลเสหิ นิสฺสฏานิ โหนฺติ, สกทาคามิ\-มคฺเค ปญฺจินฺทฺริยานิ โอฬาริเกหิ กิเลเสหิ นิสฺสฏานิ โหนฺติ, อนาคามิมคฺเค ปญฺจินฺทฺริยานิ อณุสหคเตหิ กิเลเสหิ นิสฺสฏานิ โหนฺติ, อรหตฺตมคฺเค ปญฺจินฺทฺริยานิ สพฺพกิเลเสหิ นิสฺสฏานิ โหนฺติ, สพฺเพสญฺเญว ขีณาสวานํ ตตฺถ ตตฺถ ปญฺจินฺทฺริยานิ นิสฺสฏานิ เจว โหนฺติ สุนิสฺสฏานิ จ ปฏิปฺปสฺสทฺธานิ จ สุปฺปฏิปฺปสฺสทฺธานิ จ.}

\prose{อิเมหิ อสีติสตํ อากาเรหิ ปญฺจนฺนํ อินฺทฺริยานํ นิสฺสรณํ โหติ, อิเมหิ อสีติสตํ อากาเรหิ ปญฺจนฺนํ อินฺทฺริยานํ นิสฺสรณํ ปชานาติ.}

\vspace{\tipitakaparskip}
\csromancenter{สุตฺตนฺต\-นิทฺเทโส ทุติโย.}
\nitthitamruled{ปฐม\-ภาณวาโร.}
\csromanmatikaanchor{mtk.75}
\csromantocmarkref{subsection}{3. ตติยสุตฺตนฺตนิทฺเทส}{mtk.75}
\bookmark[dest={mtk.75},level=2]{3. ตติยสุตฺตนฺตนิทฺเทส}
\csromanheader{2}{3. ตติย\-สุตฺตนฺต\-นิทฺเทส}
\proseitem{194}{สาวตฺถิ\-นิทานํ\csromansharedfootnote{95a9fff84238f432}{สํ 3. 172 ปิฏฺเฐ ปสฺสิตพฺพา.}.  ปญฺจิมานิ ภิกฺขเว อินฺทฺริยานิ. กตมานิ ปญฺจ? สทฺธินฺทฺริยํ วีริยินฺทฺริยํ สตินฺทฺริยํ สมาธินฺทฺริยํ ปญฺญินฺทฺริยํ. กตฺถ จ ภิกฺขเว สทฺธินฺทฺริยํ ทฏฺฐพฺพํ? จตูสุ โสตาปตฺติยงฺคสุ เอตฺถ สทฺธินฺทฺริยํ ทฏฺฐพฺพํ. กตฺถ จ ภิกฺขเว วีริยินฺทฺริยํ ทฏฺฐพฺพํ? จตูสุ สมฺม\-ปฺปธาเนสุ เอตฺถ วีริยินฺทฺริยํ ทฏฺฐพฺพํ. กตฺถ}

\prose{\csromanfolio{211}จ ภิกฺขเว สตินฺทฺริยํ ทฏฺฐพฺพํ? จตูสุ สติปฏฺฐาเนสุ เอตฺถ สตินฺทฺริยํ ทฏฺฐพฺพํ, กตฺถ จ ภิกฺขเว สมาธินฺทฺริยํ ทฏฺฐพฺพํ? จตูสุ ฌาเนสุ เอตฺถ สมาธินฺทฺริยํ ทฏฺฐพฺพํ. กตฺถ จ ภิกฺขเว ปญฺญินฺทฺริยํ ทฏฺฐพฺพํ? จตูสุ อริยสจฺเจสุ เอตฺถ ปญฺญินฺทฺริยํ ทฏฺฐพฺพํ.}

\prose{จตูสุ โสตาปตฺติยงฺเคสุ สทฺธิ\-นฺทฺริยสฺส วเสน กติหากาเรหิ ปญฺจินฺทฺริยานิ ทฏฺฐพฺพานิ, จตูสุ สมฺม\-ปฺปธาเนสุ วีริยิ\-นฺทฺริยสฺส วเสน กติหากาเรหิ ปญฺจินฺทฺริยานิ ทฏฺฐพฺพานิ, จตูสุ สติปฏฺฐาเนสุ สตินฺทฺริยสฺส วเสน กติหากาเรหิ ปญฺจินฺทฺริยานิ ทฏฺฐพฺพานิ, จตูสุ ฌาเนสุ สมาธิ\-นฺทฺริยสฺส วเสน กติหากาเรหิ ปญฺจินฺทฺริยานิ ทฏฺฐพฺพานิ, จตูสุ อริยสจฺเจสุ ปญฺญินฺทฺริยสฺส วเสน กติหากาเรหิ ปญฺจินฺทฺริยานิ ทฏฺฐพฺพานิ?}

\prose{จตูสุ โสตาปตฺติยงฺเคสุ สทฺธิ\-นฺทฺริยสฺส วเสน วีสติยา อากาเรหิ ปญฺจินฺทฺริยานิ ทฏฺฐพฺพานิ, จตูสุ สมฺม\-ปฺปธาเนสุ วีริยิ\-นฺทฺริยสฺส วเสน วีสติยา อากาเรหิ ปญฺจินฺทฺริยานิ ทฏฺฐพฺพานิ, จตูสุ สติปฏฺฐาเนสุ สตินฺทฺริยสฺส วเสน วีสติยา อากาเรหิ ปญฺจินฺทฺริยานิ ทฏฺฐพฺพานิ, จตูสุ ฌาเนสุ สมาธิ\-นฺทฺริยสฺส วเสน วีสติยา อากาเรหิ ปญฺจินฺทฺริยานิ ทฏฺฐพฺพานิ, จตูสุ อริยสจฺเจสุ ปญฺญินฺทฺริยสฺส วเสน วีสติยา อากาเรหิ ปญฺจินฺทฺริยานิ ทฏฺฐพฺพานิ.}

\csromanheader{2}{ก. ปเภท\-คณน\-นิทฺเทส}
\proseitem{195}{จตูสุ โสตาปตฺติยงฺเคสุ สทฺธิ\-นฺทฺริยสฺส วเสน กตเมหิ วีสติยา อากาเรหิ ปญฺจินฺทฺริยานิ ทฏฺฐพฺพานิ? สปฺปุริส\-สํเสเว โสตาปตฺติยงฺเค อธิโมกฺขา\-ธิปเตยฺยฏฺเฐน สทฺธินฺทฺริยํ ทฏฺฐพฺพํ, สทฺธิ\-นฺทฺริยสฺส วเสน ปคฺคหฏฺเฐน วีริยินฺทฺริยํ ทฏฺฐพฺพํ, อุปฏฺฐาน\-ฏฺเฐน สตินฺทฺริยํ ทฏฺฐพฺพํ, อวิกฺเขป\-ฏฺเฐน สมาธินฺทฺริยํ ทฏฺฐพฺพํ, ทสฺสนฏฺเฐน ปญฺญินฺทฺริยํ ทฏฺฐพฺพํ, สทฺธมฺม\-สวเน โสตาปตฺติยงฺเค โยนิโส\-มนสิกาเร โสตาปตฺติยงฺเค ธมฺมานุธมฺม\-ปฏิปตฺติยา โสตาปตฺติยงฺเค อธิโมกฺขา\-ธิปเตยฺยฏฺเฐน สทฺธินฺทฺริยํ ทฏฺฐพฺพํ, สทฺธิ\-นฺทฺริยสฺส วเสน ปคฺคหฏฺเฐน วีริยินฺทฺริยํ ทฏฺฐพฺพํ, อุปฏฺฐาน\-ฏฺเฐน สตินฺทฺริยํ ทฏฺฐพฺพํ, อวิกฺเขป\-ฏฺเฐน สมาธินฺทฺริยํ ทฏฺฐพฺพํ, ทสฺสนฏฺเฐน ปญฺญินฺทฺริยํ ทฏฺฐพฺพํ. จตูสุ โสตาปตฺติยงฺเคสุ สทฺธิ\-นฺทฺริยสฺส วเสน อิเมหิ วีสติยา อากาเรหิ ปญฺจินฺทฺริยานิ ทฏฺฐพฺพานิ.}

\prose{\csromanfolio{212}จตูสุ สมฺม\-ปฺปธาเนสุ วีริยิ\-นฺทฺริยสฺส วเสน กตเมหิ วีสติยา อากาเรหิ ปญฺจินฺทฺริยานิ ทฏฺฐพฺพานิ? อนุปฺปนฺนานํ ปาปกานํ อกุสลานํ ธมฺมานํ อนุปฺปาทาย สมฺมปฺปธาเน ปคฺคหา\-ธิปเตยฺยฏฺเฐน วีริยินฺทฺริยํ ทฏฺฐพฺพํ, วีริยิ\-นฺทฺริยสฺส วเสน อุปฏฺฐาน\-ฏฺเฐน สตินฺทฺริยํ ทฏฺฐพฺพํ, อวิกฺเขป\-ฏฺเฐน สมาธินฺทฺริยํ ทฏฺฐพฺพํ, ทสฺสนฏฺเฐน ปญฺญินฺทฺริยํ ทฏฺฐพฺพํ, อธิโมกฺข\-ฏฺเฐน สทฺธินฺทฺริยํ ทฏฺฐพฺพํ. อุปฺปนฺนานํ ปาปกานํ อกุสลานํ ธมฺมานํ ปหานาย สมฺมปฺปธาเน ฯเปฯ. อนุปฺปนฺนานํ กุสลานํ ธมฺมานํ อุปฺปาทาย สมฺมปฺปธาเน ฯเปฯ. อุปฺปนฺนานํ กุสลานํ ธมฺมานํ ฐิติยา อสมฺโมสาย ภิยฺโยภาวาย เวปุลฺลาย ภาวนาย ปาริปูริยา สมฺมปฺปธาเน ปคฺคหา\-ธิปเตยฺยฏฺเฐน วีริยินฺทฺริยํ ทฏฺฐพฺพํ, วีริยิ\-นฺทฺริยสฺส วเสน อุปฏฺฐาน\-ฏฺเฐน สตินฺทฺริยํ ทฏฺฐพฺพํ, อวิกฺเขป\-ฏฺเฐน สมาธินฺทฺริยํ ทฏฺฐพฺพํ, ทสฺสนฏฺเฐน ปญฺญินฺทฺริยํ ทฏฺฐพฺพํ, อธิโมกฺข\-ฏฺเฐน สทฺธินฺทฺริยํ ทฏฺฐพฺพํ. จตูสุ สมฺม\-ปฺปธาเนสุ วีริยิ\-นฺทฺริยสฺส วเสน อิเมหิ วีสติยา อากาเรหิ ปญฺจินฺทฺริยานิ ทฏฺฐพฺพานิ.}

\prose{จตูสุ สติปฏฺฐาเนสุ สตินฺทฺริยสฺส วเสน กตเมหิ วีสติยา อากาเรหิ ปญฺจินฺทฺริยานิ ทฏฺฐพฺพานิ? กาเย กายานุปสฺสนา\-สติปฏฺฐาเน อุปฏฺฐานา\-ธิปเตยฺยฏฺเฐน สตินฺทฺริยํ ทฏฺฐพฺพํ, สตินฺทฺริยสฺส วเสน อวิกฺเขป\-ฏฺเฐน สมาธินฺทฺริยํ ทฏฺฐพฺพํ, ทสฺสนฏฺเฐน ปญฺญินฺทฺริยํ ทฏฺฐพฺพํ, อธิโมกฺข\-ฏฺเฐน สทฺธินฺทฺริยํ ทฏฺฐพฺพํ, ปคฺคหฏฺเฐน วีริยินฺทฺริยํ ทฏฺฐพฺพํ. เวทนาสุ เวทนานุปสฺสนา\-สติปฏฺฐาเน ฯเปฯ. จิตฺเต จิตฺตานุปสฺสนา\-สติปฏฺฐาเน. ธมฺเมสุ ธมฺมานุปสฺสนา\-สติปฏฺฐาเน อุปฏฺฐานา\-ธิปเตยฺยฏฺเฐน สตินฺทฺริยํ ทฏฺฐพฺพํ, สตินฺทฺริยสฺส วเสน อวิกฺเขป\-ฏฺเฐน สมาธินฺทฺริยํ ทฏฺฐพฺพํ, ทสฺสนฏฺเฐน ปญฺญินฺทฺริยํ ทฏฺฐพฺพํ, อธิโมกฺข\-ฏฺเฐน สทฺธินฺทฺริยํ ทฏฺฐพฺพํ, ปคฺคหฏฺเฐน วีริยินฺทฺริยํ ทฏฺฐพฺพํ. จตูสุ สติปฏฺฐาเนสุ สตินฺทฺริยสฺส วเสน อิเมหิ วีสติยา อากาเรหิ ปญฺจินฺทฺริยานิ ทฏฺฐพฺพานิ.}

\prose{จตูสุ ฌาเนสุ สมาธิ\-นฺทฺริยสฺส วเสน กตเมหิ วีสติยา อากาเรหิ ปญฺจินฺทฺริยานิ ทฏฺฐพฺพานิ? ปฐเม ฌาเน อวิกฺเขปา\-ธิปเตยฺยฏฺเฐน สมาธินฺทฺริยํ ทฏฺฐพฺพํ, สมาธิ\-นฺทฺริยสฺส วเสน ทสฺสนฏฺเฐน ปญฺญินฺทฺริยํ ทฏฺฐพฺพํ, อธิโมกฺข\-ฏฺเฐน สทฺธินฺทฺริยํ ทฏฺฐพฺพํ, ปคฺคหฏฺเฐน วีริยินฺทฺริยํ ทฏฺฐพฺพํ, อุปฏฺฐาน\-ฏฺเฐน สตินฺทฺริยํ ทฏฺฐพฺพํ. ทุติเย ฌาเน ฯเปฯ. ตติเย ฌาเน. จตุตฺเถ ฌาเน อวิกฺเขปา\-ธิปเตยฺยฏฺเฐน สมาธินฺทฺริยํ ทฏฺฐพฺพํ, สมาธิ\-นฺทฺริยสฺส วเสน ทสฺสนฏฺเฐน ปญฺญินฺทฺริยํ ทฏฺฐพฺพํ, อธิโมกฺข\-ฏฺเฐน สทฺธินฺทฺริยํ ทฏฺฐพฺพํ, ปคฺคหฏฺเฐน}

\prose{\csromanfolio{213}วีริยินฺทฺริยํ ทฏฺฐพฺพํ, อุปฏฺฐาน\-ฏฺเฐน สตินฺทฺริยํ ทฏฺฐพฺพํ. จตูสุ ฌาเนสุ สมาธิ\-นฺทฺริยสฺส วเสน อิเมหิ วีสติยา อากาเรหิ ปญฺจินฺทฺริยานิ ทฏฺฐพฺพานิ.}

\prose{จตูสุ อริยสจฺเจสุ ปญฺญินฺทฺริยสฺส วเสน กตเมหิ วีสติยา อากาเรหิ ปญฺจินฺทฺริยานิ ทฏฺฐพฺพานิ? ทุกฺเข อริยสจฺเจ ทสฺสนา\-ธิปเตยฺยฏฺเฐน ปญฺญินฺทฺริยํ ทฏฺฐพฺพํ, ปญฺญินฺทฺริยสฺส วเสน อธิโมกฺข\-ฏฺเฐน สทฺธินฺทฺริยํ ทฏฺฐพฺพํ, ปคฺคหฏฺเฐน วีริยินฺทฺริยํ ทฏฺฐพฺพํ. อุปฏฺฐาน\-ฏฺเฐน สตินฺทฺริยํ ทฏฺฐพฺพํ. อวิกฺเขป\-ฏฺเฐน สมาธินฺทฺริยํ ทฏฺฐพฺพํ. ทุกฺขสมุทเย อริยสจฺเจ ฯเปฯ. ทุกฺขนิโรเธ อริยสจฺเจ. ทุกฺขนิโรธ\-คามินิยา ปฏิปทาย อริยสจฺเจ ทสฺสนา\-ธิปเตยฺยฏฺเฐน ปญฺญินฺทฺริยํ ทฏฺฐพฺพํ, ปญฺญินฺทฺริยสฺส วเสน อธิโมกฺข\-ฏฺเฐน สทฺธินฺทฺริยํ ทฏฺฐพฺพํ, ปคฺคหฏฺเฐน วีริยินฺทฺริยํ ทฏฺฐพฺพํ, อุปฏฺฐาน\-ฏฺเฐน สตินฺทฺริยํ ทฏฺฐพฺพํ, อวิกฺเขป\-ฏฺเฐน สมาธินฺทฺริยํ ทฏฺฐพฺพํ. จตูสุ อริยสจฺเจสุ ปญฺญินฺทฺริยสฺส วเสน อิเมหิ วีสติยา อากาเรหิ ปญฺจินฺทฺริยานิ ทฏฺฐพฺพานิ.}

\csromanheader{2}{\textbf{ข. จริยวาร}}
\proseitem{196}{จตูสุ โสตาปตฺติยงฺเคสุ สทฺธิ\-นฺทฺริยสฺส วเสน กติหากาเรหิ ปญฺจนฺนํ อินฺทฺริยานํ จริยา ทฏฺฐพฺพา? จตูสุ สมฺม\-ปฺปธาเนสุ ฯเปฯ. จตูสุ สติปฏฺฐาเนสุ. จตูสุ ฌาเนสุ. จตูสุ อริยสจฺเจสุ ปญฺญินฺทฺริยสฺส วเสน กติหากาเรหิ ปญฺจนฺนํ อินฺทฺริยานํ จริยา ทฏฺฐพฺพา.  จตูสุ โสตาปตฺติยงฺเคสุ สทฺธิ\-นฺทฺริยสฺส วเสน วีสติยา อากาเรหิ ปญฺจนฺนํ อินฺทฺริยานํ จริยา ทฏฺฐพฺพา. จตูสุ สมฺม\-ปฺปธาเนสุ ฯเปฯ. จตูสุ สติปฏฺฐาเนสุ จตูสุ ฌาเนสุ จตูสุ อริยสจฺเจสุ ปญฺญินฺทฺริยสฺส วเสน วีสติยา อากาเรหิ ปญฺจนฺนํ อินฺทฺริยานํ จริยา ทฏฺฐพฺพา.}

\prose{จตูสุ โสตาปตฺติยงฺเคสุ สทฺธิ\-นฺทฺริยสฺส วเสน กตเมหิ วีสติยา อากาเรหิ ปญฺจนฺนํ อินฺทฺริยานํ จริยา ทฏฺฐพฺพา? สปฺปุริส\-สํเสเว โสตาปตฺติยงฺเค อธิโมกฺขา\-ธิปเตยฺยฏฺเฐน สทฺธิ\-นฺทฺริยสฺส จริยา ทฏฺฐพฺพา, สทฺธิ\-นฺทฺริยสฺส วเสน ปคฺคหฏฺเฐน วีริยิ\-นฺทฺริยสฺส จริยา ทฏฺฐพฺพา, อุปฏฺฐาน\-ฏฺเฐน สตินฺทฺริยสฺส จริยา ทฏฺฐพฺพา, อวิกฺเขป\-ฏฺเฐน สมาธิ\-นฺทฺริยสฺส จริยา ทฏฺฐพฺพา. ทสฺสนฏฺเฐน ปญฺญินฺทฺริยสฺส จริยา ทฏฺฐพฺพา. สทฺธมฺม\-สวเน โสตาปตฺติยงฺเค ฯเปฯ. โยนิโส มนสิกาเร โสตาปตฺติยงฺเค ธมฺมานุธมฺม\-ปฏิปตฺติยา โสตาปตฺติยงฺเค อธิโมกฺขา\-ธิปเตยฺยฏฺเฐน สทฺธิ\-นฺทฺริยสฺส จริยา ทฏฺฐพฺพา, สทฺธิ\-นฺทฺริยสฺส วเสน ปคฺคหฏฺเฐน วีริยิ\-นฺทฺริยสฺส \csromanfolio{214}จริยา ทฏฺฐพฺพา, อุปฏฺฐาน\-ฏฺเฐน สตินฺทฺริยสฺส จริยา ทฏฺฐพฺพา, อวิกฺเขป\-ฏฺเฐน สมาธิ\-นฺทฺริยสฺส จริยา ทฏฺฐพฺพา, ทสฺสนฏฺเฐน ปญฺญินฺทฺริยสฺส จริยา ทฏฺฐพฺพา. จตูสุ โสตาปตฺติยงฺเคสุ สทฺธิ\-นฺทฺริยสฺส วเสน อิเมหิ วีสติยา อากาเรหิ ปญฺจนฺนํ อินฺทฺริยานํ จริยา ทฏฺฐพฺพา.}

\prose{จตูสุ สมฺม\-ปฺปธาเนสุ วีริยิ\-นฺทฺริยสฺส วเสน กตเมหิ วีสติยา อากาเรหิ ปญฺจนฺนํ อินฺทฺริยานํ จริยา ทฏฺฐพฺพา? อนุปฺปนฺนานํ ปาปกานํ อกุสลานํ ธมฺมานํ อนุปฺปาทาย สมฺมปฺปธาเน ปคฺคหา\-ธิปเตยฺยฏฺเฐน วีริยิ\-นฺทฺริยสฺส จริยา ทฏฺฐพฺพา, วีริยิ\-นฺทฺริยสฺส วเสน อุปฏฺฐาน\-ฏฺเฐน สตินฺทฺริยสฺส จริยา ทฏฺฐพฺพา, อวิกฺเขป\-ฏฺเฐน สมาธิ\-นฺทฺริยสฺส จริยา ทฏฺฐพฺพา, ทสฺสนฏฺเฐน ปญฺญินฺทฺริยสฺส จริยา ทฏฺฐพฺพา, อธิโมกฺข\-ฏฺเฐน สทฺธิ\-นฺทฺริยสฺส จริยา ทฏฺฐพฺพา. อุปฺปนฺนานํ ปาปกานํ อกุสลานํ ธมฺมานํ ปหานาย สมฺมปฺปธาเน ฯเปฯ. อนุปฺปนฺนานํ กุสลานํ ธมฺมานํ อุปฺปาทาย สมฺมปฺปธาเน ฯเปฯ. อุปฺปนฺนานํ กุสลานํ ธมฺมานํ ฐิติยา อสมฺโมสาย ภิยฺโยภาวาย เวปุลฺลาย ภาวนาย ปาริปูริยา สมฺมปฺปธาเน ปคฺคหา\-ธิปเตยฺยฏฺเฐน วีริยิ\-นฺทฺริยสฺส จริยา ทฏฺฐพฺพา ฯเปฯ. จตูสุ สมฺม\-ปฺปธาเนสุ วีริยิ\-นฺทฺริยสฺสวเสน อิเมหิ วีสติยา อากาเรหิ ปญฺจนฺนํ อินฺทฺริยานํ จริยา ทฏฺฐพฺพา.}

\prose{จตูสุ สติปฏฺฐาเนสุ สตินฺทฺริยสฺส วเสน กตเมหิ วีสติยา อากาเรหิ ปญฺจนฺนํ อินฺทฺริยานํ จริยา ทฏฺฐพฺพา? กาเย กายานุปสฺสนา\-สติปฏฺฐาเน อุปฏฺฐานา\-ธิปเตยฺยฏฺเฐน สตินฺทฺริยสฺส จริยา ทฏฺฐพฺพา, สตินฺทฺริยสฺส วเสน อวิกฺเขป\-ฏฺเฐน สมาธิ\-นฺทฺริยสฺส จริยา ทฏฺฐพฺพา, ทสฺสนฏฺเฐน ปญฺญินฺทฺริยสฺส จริยา ทฏฺฐพฺพา, อธิโมกฺข\-ฏฺเฐน สทฺธิ\-นฺทฺริยสฺส จริยา ทฏฺฐพฺพา, ปคฺคหฏฺเฐน วีริยิ\-นฺทฺริยสฺส จริยา ทฏฺฐพฺพา. เวทนาสุ เวทนานุปสฺสนา\-สติปฏฺฐาเน ฯเปฯ. จิตฺเต จิตฺตานุปสฺสนา\-สติปฏฺฐาเน ฯเปฯ. ธมฺเมสุ ธมฺมา\-นุปสฺสนา สติปฏฺฐาเน อุปฏฺฐานา\-ธิปเตยฺยฏฺเฐน สตินฺทฺริยสฺส จริยา ทฏฺฐพฺพา ฯเปฯ. จตูสุ สติปฏฺฐาเนสุ สตินฺทฺริยสฺส วเสน อิเมหิ วีสติยา อากาเรหิ ปญฺจนฺนํ อินฺทฺริยานํ จริยา ทฏฺฐพฺพา.}

\prose{จตูสุ ฌาเนสุ สมาธิ\-นฺทฺริยสฺส วเสน กตเมหิ วีสติยา อากาเรหิ ปญฺจนฺนํ อินฺทฺริยานํ จริยา ทฏฺฐพฺพา? ปฐเม ฌาเน อวิกฺเขปา\-ธิปเตยฺยฏฺเฐน สมาธิ\-นฺทฺริยสฺส จริยา ทฏฺฐพฺพา, สมาธิ\-นฺทฺริยสฺส วเสน ทสฺสนฏฺเฐน ปญฺญินฺทฺริยสฺส จริยา ทฏฺฐพฺพา, อธิโมกฺข\-ฏฺเฐน สทฺธิ\-นฺทฺริยสฺส จริยา ทฏฺฐพฺพา, ปคฺคหฏฺเฐน วีริยิ\-นฺทฺริยสฺส จริยา ทฏฺฐพฺพา, อุปฏฺฐาน\-ฏฺเฐน}

\prose{\csromanfolio{215}สตินฺทฺริยสฺส จริยา ทฏฺฐพฺพา. ทุติเย ฌาเน ฯเปฯ. ตติเย ฌาเน ฯเปฯ. จตุตฺเถ ฌาเน อวิกฺเขปา\-ธิปเตยฺยฏฺเฐน สมาธิ\-นฺทฺริยสฺส จริยา ทฏฺฐพฺพา ฯเปฯ. จตูสุ ฌาเนสุ สมาธิ\-นฺทฺริยสฺส วเสน อิเมหิ วีสติยา อากาเรหิ ปญฺจนฺนํ อินฺทฺริยานํ จริยา ทฏฺฐพฺพา.}

\prose{จตูสุ อริยสจฺเจสุ ปญฺญินฺทฺริยสฺส วเสน กตเมหิ วีสติยา อากาเรหิ ปญฺจนฺนํ อินฺทฺริยานํ จริยา ทฏฺฐพฺพา? ทุกฺเข อริยสจฺเจ ทสฺสนา\-ธิปเตยฺยฏฺเฐน ปญฺญินฺทฺริยสฺส จริยา ทฏฺฐพฺพา, ปญฺญินฺทฺริยสฺส วเสน อธิโมกฺข\-ฏฺเฐน สทฺธิ\-นฺทฺริยสฺส จริยา ทฏฺฐพฺพา, ปคฺคหฏฺเฐน วีริยิ\-นฺทฺริยสฺส จริยา ทฏฺฐพฺพา, อุปฏฺฐาน\-ฏฺเฐน สตินฺทฺริยสฺส จริยา ทฏฺฐพฺพา, อวิกฺเขป\-ฏฺเฐน สมาธิ\-นฺทฺริยสฺส จริยา ทฏฺฐพฺพา. ทุกฺขสมุทเย อริยสจฺเจ ฯเปฯ. ทุกฺขนิโรเธ อริยสจฺเจ. ทุกฺขนิโรธ\-คามินิยา ปฏิปทาย อริยสจฺเจ ทสฺสนา\-ธิปเตยฺยฏฺเฐน ปญฺญินฺทฺริยสฺส จริยา ทฏฺฐพฺพา, ปญฺญินฺทฺริยสฺส วเสน อธิโมกฺข\-ฏฺเฐน สทฺธิ\-นฺทฺริยสฺส จริยา ทฏฺฐพฺพา, ปคฺคหฏฺเฐน วีริยิ\-นฺทฺริยสฺส จริยา ทฏฺฐพฺพา, อุปฏฺฐาน\-ฏฺเฐน สตินฺทฺริยสฺส จริยา ทฏฺฐพฺพา, อวิกฺเขป\-ฏฺเฐน สมาธิ\-นฺทฺริยสฺส จริยา ทฏฺฐพฺพา. จตูสุ อริยสจฺเจสุ ปญฺญินฺทฺริยสฺส วเสน อิเมหิ วีสติยา อากาเรหิ ปญฺจนฺนํ อินฺทฺริยานํ จริยา ทฏฺฐพฺพา.}

\csromanheader{2}{จาร\-วิหาร\-นิทฺเทส}
\proseitem{197}{จาโร จ วิหาโร จ อนุพุทฺโธ โหติ ปฏิวิทฺโธ. ยถาจรนฺตํ ยถา\-วิหรนฺตํ วิญฺญู สพฺรหฺมจารี คมฺภีเรสุ ฐาเนสุ โอกปฺเปยฺยุํ อทฺธา อยมายสฺมา ปตฺโต วา ปาปุณิสฺสติ วา.}

\prose{จริยาติ อฏฺฐ จริยาโย อิริยาปถ\-จริยา, อายตนจริยา, สติจริยา, สมาธิจริยา, ญาณจริยา, มคฺคจริยา, ปตฺติจริยา, โลกตฺถจริยาติ. อิริยาปถ\-จริยาติ จตูสุ อิริยาปเถสุ. อายตน\-จริยาติ ฉสุ อชฺฌตฺติก\-พาหิเรสุ อายตเนสุ. สติจริยาติ จตูสุ สติปฏฺฐาเนสุ. สมาธิจริยาติ จตูสุ ฌาเนสุ. ญาณจริยาติ จตูสุ อริยสจฺเจสุ. มคฺคจริยาติ จตูสุ อริยมคฺเคสุ. ปตฺติจริยาติ จตูสุ สามญฺญผเลสุ. โลกตฺถจริยาติ ตถาคเตสุ อรหนฺเตสุ สมฺมา\-สมฺพุทฺเธสุ, ปเทเส ปจฺเจก\-พุทฺเธสุ, ปเทเส สาวเกสุ. อิริยาปถ\-จริยา จ ปณิธิ\-สมฺปนฺนานํ. อายตนจริยา จ \csromanfolio{216}อินฺทฺริเยสุ คุตฺตทฺวารานํ. สติจริยา จ อปฺปมาท\-วิหารีนํ. สมาธิจริยา จ อธิจิตฺตม\-นุยุตฺตานํ. ญาณจริยา จ พุทฺธิ\-สมฺปนฺนานํ. มคฺคจริยา จ สมฺมา\-ปฏิปนฺนานํ. ปตฺติจริยา จ อธิคต\-ผลานํ. โลกตฺถจริยา จ ตถาคตานํ อรหนฺตานํ สมฺมา\-สมฺพุทฺธานํ, ปเทเส ปจฺเจก\-พุทฺธานํ, ปเทเส สาวกานํ. อิมา อฏฺฐ จริยาโย.}

\prose{อปราปิ อฏฺฐ จริยาโย อธิมุจฺจนฺโต สทฺธาย จรติ, ปคฺคณฺหนฺโต วีริเยน จรติ, อุปฏฺฐาเปนฺโต สติยา จรติ, อวิกฺเขปํ กโรนฺโต สมาธินา จรติ, ปชานนฺโต ปญฺญาย จรติ, วิชานนฺโต วิญฺญาณ\-จริยาย จรติ, “เอวํ ปฏิปนฺนสฺส กุสลา ธมฺมา อายาเปนฺตี”ติ อายตน\-จริยาย จรติ, “เอวํ ปฏิปนฺโน วิเสสมธิคจฺฉตี”ติ วิเสสจริยาย จรติ. อิมา อฏฺฐ จริยาโย.}

\prose{อปราปิ อฏฺฐ จริยาโย ทสฺสนจริยา จ สมฺมาทิฏฺฐิยา, อภินิโรปน\-จริยา จ สมฺมา\-สงฺกปฺปสฺส, ปริคฺคห\-จริยา จ สมฺมาวาจาย, สมุฏฺฐาน\-จริยา จ สมฺมา\-กมฺมนฺตสฺส, โวทานจริยา จ สมฺมาอาชีวสฺส, ปคฺคหจริยา จ สมฺมาวายามสฺส, อุปฏฺฐาน\-จริยา จ สมฺมาสติยา, อวิกฺเขป\-จริยา จ สมฺมา\-สมาธิสฺส. อิมา อฏฺฐ จริยาโย.}

\prose{\textbf{วิหาโร}ติ อธิมุจฺจนฺโต สทฺธาย วิหรติ, ปคฺคณฺหนฺโต วีริเยน วิหรติ, อุปฏฺฐาเปนฺโต สติยา วิหรติ, อวิกฺเขปํ กโรนฺโต สมาธินา วิหรติ, ปชานนฺโต ปญฺญาย วิหรติ.}

\prose{อนุพุทฺโธติ สทฺธิ\-นฺทฺริยสฺส อธิโมกฺขฏฺโฐ อนุพุทฺโธ โหติ, วีริยิ\-นฺทฺริยสฺส ปคฺคหฏฺโฐ อนุพุทฺโธ โหติ, สตินฺทฺริยสฺส อุปฏฺฐานฏฺโฐ อนุพุทฺโธ โหติ, สมาธิ\-นฺทฺริยสฺส อวิกฺเขปฏฺโฐ อนุพุทฺโธ โหติ, ปญฺญินฺทฺริยสฺส ทสฺสนฏฺโฐ อนุพุทฺโธ โหติ.}

\prose{ปฏิวิทฺโธติ สทฺธิ\-นฺทฺริยสฺส อธิโมกฺขฏฺโฐ ปฏิวิทฺโธ โหติ, วีริยิ\-นฺทฺริยสฺส ปคฺคหฏฺโฐ ปฏิวิทฺโธ โหติ, สตินฺทฺริยสฺส อุปฏฺฐานฏฺโฐ ปฏิวิทฺโธ โหติ, สมาธิ\-นฺทฺริยสฺส อวิกฺเขปฏฺโฐ ปฏิวิทฺโธ โหติ, ปญฺญินฺทฺริยสฺส ทสฺสนฏฺโฐ ปฏิวิทฺโธ โหติ.  ยถาจรนฺตนฺติ เอวํ สทฺธาย จรนฺตํ เอวํ วีริเยน จรนฺตํ เอวํ สติยา จรนฺตํ เอวํ สมาธินา จรนฺตํ เอวํ ปญฺญาย จรนฺตํ.  ยถาวิหรนฺตนฺติ เอวํ สทฺธาย วิหรนฺตํ เอวํ วีริเยน วิหรนฺตํ เอวํ สติยา วิหรนฺตํ เอวํ สมาธินา วิหรนฺตํ เอวํ ปญฺญาย วิหรนฺตํ. วิญฺญูติ}

\prosewithcloserruled{\csromanfolio{217}วิญฺญู วิภาวี เมธาวี ปณฺฑิตา พุทฺธิสมฺปนฺนา. สพฺรหฺมจารีติ เอกกมฺมํ เอกุทฺเทโส สมสิกฺขตา. คมฺภีเรสุ ฐาเนสูติ คมฺภีรานิ ฐานานิ วุจฺจนฺติ ฌานา จ วิโมกฺขา จ สมาธี จ สมาปตฺติโย จ มคฺคา จ ผลานิ จ อภิญฺญาโย จ ปฏิสมฺภิทา จ. โอกปฺเปยฺยุนฺติ สทฺทเหยฺยุํ อธิมุจฺเจยฺยุํ. อทฺธาติ เอกํสวจนเมตํ นิสฺสํสยวจนเมตํ นิกฺกงฺข\-วจนเมตํ อเทฺวชฺฌวจนเมตํ อเทฺวฬฺหก\-วจนเมตํ นิโยค\-วจนเมตํ อปณฺณก\-วจนเมตํ อวตฺถาปน\-วจนเมตํ “อทฺธา”ติ. อายสฺมาติ ปิยวจนเมตํ ครุวจนเมตํ สคารวสปฺปติสฺสาธิวจนเมตํ “อายสฺมา”ติ. ปตฺโต วาติ อธิคโต วา. ปาปุณิสฺสติ วาติ อธิคมิสฺสติ วา.}{สุตฺตนฺต\-นิทฺเทโส ตติโย.}
\csromanmatikaanchor{mtk.76}
\csromantocmarkref{subsection}{4. จตุตฺถสุตฺตนฺตนิทฺเทส}{mtk.76}
\bookmark[dest={mtk.76},level=2]{4. จตุตฺถสุตฺตนฺตนิทฺเทส}
\csromanheader{2}{4. จตุตฺถ\-สุตฺตนฺต\-นิทฺเทส}
\proseitem{198}{ปุริมนิทานํ\csromansharedfootnote{3dd02d192165dcc8}{ปริปุณฺณนิทานํ (สฺยา, ก)}.  ปญฺจิมานิ ภิกฺขเว อินฺทฺริยานิ. กตมานิ ปญฺจ? สทฺธินฺทฺริยํ, วีริยินฺทฺริยํ, สตินฺทฺริยํ, สมาธินฺทฺริยํ, ปญฺญินฺทฺริยํ. อิมานิ โข ภิกฺขเว ปญฺจินฺทฺริยานิ. อิมานิ ปญฺจินฺทฺริยานิ กติหากาเรหิ เกนฏฺเฐน ทฏฺฐพฺพานิ? อิมานิ ปญฺจินฺทฺริยานิ ฉหากาเรหิ เตนฏฺเฐน ทฏฺฐพฺพานิ อาธิปเตยฺย\-ฏฺเฐน, อาทิ\-วิโสธ\-นฏฺเฐน, อธิมตฺต\-ฏฺเฐน, อธิฏฺฐาน\-ฏฺเฐน, ปริยาทาน\-ฏฺเฐน, ปติฏฺฐา\-ปกฏฺเฐน.}

\csromanheader{2}{ก. อาธิปเตยฺยฏฺฐ\-นิทฺเทส}
\proseitem{199}{กถํ อาธิปเตยฺย\-ฏฺเฐน อินฺทฺริยานิ ทฏฺฐพฺพานิ? อสฺสทฺธิยํ ปชหโต อธิโมกฺขา\-ธิปเตยฺยฏฺเฐน สทฺธินฺทฺริยํ ทฏฺฐพฺพํ, สทฺธิ\-นฺทฺริยสฺส วเสน ปคฺคหฏฺเฐน วีริยินฺทฺริยํ ทฏฺฐพฺพํ, อุปฏฺฐาน\-ฏฺเฐน สตินฺทฺริยํ ทฏฺฐพฺพํ, อวิกฺเขป\-ฏฺเฐน สมาธินฺทฺริยํ ทฏฺฐพฺพํ, ทสฺสนฏฺเฐน ปญฺญินฺทฺริยํ ทฏฺฐพฺพํ. โกสชฺชํ ปชหโต ปคฺคหา\-ธิปเตยฺยฏฺเฐน วีริยินฺทฺริยํ ทฏฺฐพฺพํ, วีริยิ\-นฺทฺริยสฺส วเสน อุปฏฺฐาน\-ฏฺเฐน สตินฺทฺริยํ ทฏฺฐพฺพํ, อวิกฺเขป\-ฏฺเฐน สมาธินฺทฺริยํ ทฏฺฐพฺพํ, ทสฺสนฏฺเฐน ปญฺญินฺทฺริยํ ทฏฺฐพฺพํ, อธิโมกฺข\-ฏฺเฐน สทฺธินฺทฺริยํ ทฏฺฐพฺพํ. ปมาทํ ปชหโต อุปฏฺฐานา\-ธิปเตยฺยฏฺเฐน สตินฺทฺริยํ ทฏฺฐพฺพํ, สตินฺทฺริยสฺส วเสน อวิกฺเขป\-ฏฺเฐน สมาธินฺทฺริยํ ทฏฺฐพฺพํ, ทสฺสนฏฺเฐน ปญฺญินฺทฺริยํ ทฏฺฐพฺพํ, อธิโมกฺข\-ฏฺเฐน สทฺธินฺทฺริยํ ทฏฺฐพฺพํ. ปคฺคหฏฺเฐน \csromanfolio{218}วีริยินฺทฺริยํ ทฏฺฐพฺพํ. อุทฺธจฺจํ ปชหโต อวิกฺเขปา\-ธิปเตยฺยฏฺเฐน สมาธินฺทฺริยํ ทฏฺฐพฺพํ, สมาธิ\-นฺทฺริยสฺส วเสน ทสฺสนฏฺเฐน ปญฺญินฺทฺริยํ ทฏฺฐพฺพํ อธิโมกฺข\-ฏฺเฐน สทฺธินฺทฺริยํ ทฏฺฐพฺพํ, ปคฺคหฏฺเฐน วีริยินฺทฺริยํ ทฏฺฐพฺพํ. อุปฏฺฐาน\-ฏฺเฐน สตินฺทฺริยํ ทฏฺฐพฺพํ. อวิชฺชํ ปชหโต ทสฺสนา\-ธิปเตยฺยฏฺเฐน ปญฺญินฺทฺริยํ ทฏฺฐพฺพํ. ปญฺญินฺทฺริยสฺส วเสน อธิโมกฺข\-ฏฺเฐน สทฺธินฺทฺริยํ ทฏฺฐพฺพํ. ปคฺคหฏฺเฐน วีริยินฺทฺริยํ ทฏฺฐพฺพํ. อุปฏฺฐาน\-ฏฺเฐน สตินฺทฺริยํ ทฏฺฐพฺพํ. อวิกฺเขป\-ฏฺเฐน สมาธินฺทฺริยํ ทฏฺฐพฺพํ. ทสฺสนฏฺเฐน ปญฺญินฺทฺริยํ ทฏฺฐพฺพํ. กามจฺฉนฺทํ ปชหโต เนกฺขมฺม\-วเสน อธิโมกฺขา\-ธิปเตยฺยฏฺเฐน สทฺธินฺทฺริยํ ทฏฺฐพฺพํ, สทฺธิ\-นฺทฺริยสฺส วเสน ปคฺคหฏฺเฐน วีริยินฺทฺริยํ ทฏฺฐพฺพํ. อุปฏฺฐาน\-ฏฺเฐน สตินฺทฺริยํ ทฏฺฐพฺพํ, อวิกฺเขป\-ฏฺเฐน สมาธินฺทฺริยํ ทฏฺฐพฺพํ, ทสฺสนฏฺเฐน ปญฺญินฺทฺริยํ ทฏฺฐพฺพํ. กามจฺฉนฺทํ ปชหโต เนกฺขมฺม\-วเสน ปคฺคหา\-ธิปเตยฺยฏฺเฐน วีริยินฺทฺริยํ ทฏฺฐพฺพํ, วีริยิ\-นฺทฺริยสฺส วเสน อุปฏฺฐาน\-ฏฺเฐน สตินฺทฺริยํ ทฏฺฐพฺพํ, อวิกฺเขป\-ฏฺเฐน สมาธินฺทฺริยํ ทฏฺฐพฺพํ, ทสฺสนฏฺเฐน ปญฺญินฺทฺริยํ ทฏฺฐพฺพํ, อธิโมกฺข\-ฏฺเฐน สทฺธินฺทฺริยํ ทฏฺฐพฺพํ. กามจฺฉนฺทํ ปชหโต เนกฺขมฺม\-วเสน อุปฏฺฐานา\-ธิปเตยฺยฏฺเฐน สตินฺทฺริยํ ทฏฺฐพฺพํ, สตินฺทฺริยสฺส วเสน อวิกฺเขป\-ฏฺเฐน สมาธินฺทฺริยํ ทฏฺฐพฺพํ, ทสฺสนฏฺเฐน ปญฺญินฺทฺริยํ ทฏฺฐพฺพํ, อธิโมกฺข\-ฏฺเฐน สทฺธินฺทฺริยํ ทฏฺฐพฺพํ, ปคฺคหฏฺเฐน วีริยินฺทฺริยํ ทฏฺฐพฺพํ. กามจฺฉนฺทํ ปชหโต เนกฺขมฺม\-วเสน อวิกฺเขปา\-ธิปเตยฺยฏฺเฐน สมาธินฺทฺริยํ ทฏฺฐพฺพํ, สมาธิ\-นฺทฺริยสฺส วเสน ทสฺสนฏฺเฐน ปญฺญินฺทฺริยํ ทฏฺฐพฺพํ, อธิโมกฺข\-ฏฺเฐน สทฺธินฺทฺริยํ ทฏฺฐพฺพํ, ปคฺคหฏฺเฐน วีริยินฺทฺริยํ ทฏฺฐพฺพํ, อุปฏฺฐาน\-ฏฺเฐน สตินฺทฺริยํ ทฏฺฐพฺพํ. กามจฺฉนฺทํ ปชหโต เนกฺขมฺม\-วเสน ทสฺสนา\-ธิปเตยฺยฏฺเฐน ปญฺญินฺทฺริยํ ทฏฺฐพฺพํ, ปญฺญินฺทฺริยสฺส วเสน อธิโมกฺข\-ฏฺเฐน สทฺธินฺทฺริยํ ทฏฺฐพฺพํ, ปคฺคหฏฺเฐน วีริยินฺทฺริยํ ทฏฺฐพฺพํ, อุปฏฺฐาน\-ฏฺเฐน สตินฺทฺริยํ ทฏฺฐพฺพํ, อวิกฺเขป\-ฏฺเฐน สมาธินฺทฺริยํ ทฏฺฐพฺพํ.}

\prose{พฺยาปาทํ ปชหโต อพฺยาปาทวเสน ฯเปฯ. ถนมิทฺธํ ปชหโต อาโลกสญฺญา\-วเสน ฯเปฯ. สพฺพกิเลเส ปชหโต อรหตฺตมคฺค\-วเสน อธิโมกฺขา\-ธิปเตยฺยฏฺเฐน สทฺธินฺทฺริยํ ทฏฺฐพฺพํ, สทฺธิ\-นฺทฺริยสฺส วเสน ปคฺคหฏฺเฐน วีริยินฺทฺริยํ ทฏฺฐพฺพํ, อุปฏฺฐาน\-ฏฺเฐน สตินฺทฺริยํ ทฏฺฐพฺพํ, อวิกฺเขป\-ฏฺเฐน สมาธินฺทฺริยํ ทฏฺฐพฺพํ, ทสฺสนฏฺเฐน ปญฺญินฺทฺริยํ ทฏฺฐพฺพํ ฯเปฯ สพฺพกิเลเส ปชหโต อรหตฺตมคฺค\-วเสน ทสฺสนา\-ธิปเตยฺยฏฺเฐน ปญฺญินฺทฺริยํ ทฏฺฐพฺพํ, ปญฺญินฺทฺริยสฺส วเสน อธิโมกฺข\-ฏฺเฐน สทฺธินฺทฺริยํ ทฏฺฐพฺพํ, ปคฺคหฏฺเฐน วีริยินฺทฺริยํ ทฏฺฐพฺพํ,}

\prose{\csromanfolio{219}อุปฏฺฐาน\-ฏฺเฐน สตินฺทฺริยํ ทฏฺฐพฺพํ, อวิกฺเขป\-ฏฺเฐน สมาธินฺทฺริยํ ทฏฺฐพฺพํ. เอวํ อาธิปเตยฺย\-ฏฺเฐน อินฺทฺริยานิ ทฏฺฐพฺพานิ. (1)}

\csromanheader{2}{ข. อาทิ\-วิโสธ\-นฏฺฐ\-นิทฺเทส}
\proseitem{200}{กถํ อาทิ\-วิโสธ\-นฏฺเฐน อินฺทฺริยานิ ทฏฺฐพฺพานิ? อธิโมกฺข\-ฏฺเฐน สทฺธินฺทฺริยํ อสฺสทฺธิย\-สํวรฏฺเฐน\csromansharedfootnote{ae2cb31c07462c88}{อสฺสทฺธิยํ สํวรฏฺเฐน (สฺยา, ก) เอวมีทิเสสุ ปเทสุปิ.} สีลวิสุทฺธิ สทฺธิ\-นฺทฺริยสฺส อาทิวิโสธนา. ปคฺคหฏฺเฐน วีริยินฺทฺริยํ โกสชฺช\-สํวรฏฺเฐน สีลวิสุทฺธิ วีริยิ\-นฺทฺริยสฺส อาทิวิโสธนา. อุปฏฺฐาน\-ฏฺเฐน สตินฺทฺริยํ ปมาท\-สํวรฏฺเฐน สีลวิสุทฺธิ สตินฺทฺริยสฺส อาทิวิโสธนา. อวิกฺเขป\-ฏฺเฐน สมาธินฺทฺริยํ อุทฺธจฺจ\-สํวรฏฺเฐน สีลวิสุทฺธิ สมาธิ\-นฺทฺริยสฺส อาทิวิโสธนา. ทสฺสนฏฺเฐน ปญฺญินฺทฺริยํ อวิชฺชา\-สํวรฏฺเฐน สีลวิสุทฺธิ ปญฺญินฺทฺริยสฺส อาทิวิโสธนา. เนกฺขมฺเม ปญฺจินฺทฺริยานิ กามจฺฉนฺท\-สํวรฏฺเฐน สีลวิสุทฺธิ ปญฺจนฺนํ อินฺทฺริยานํ อาทิวิโสธนา. อพฺยาปาเท ปญฺจินฺทฺริยานิ พฺยาปาท\-สํวรฏฺเฐน สีลวิสุทฺธิ ปญฺจนฺนํ อินฺทฺริยานํ อาทิวิโสธนา ฯเปฯ. อรหตฺตมคฺเค ปญฺจินฺทฺริยานิ สพฺพกิเลส\-สํวรฏฺเฐน สีลวิสุทฺธิ ปญฺจนฺนํ อินฺทฺริยานํ อาทิวิโสธนา. เอวํ อาทิ\-วิโสธ\-นฏฺเฐน อินฺทฺริยานิ ทฏฺฐพฺพานิ. (2)}

\csromanheader{2}{ค. อธิมตฺต\-ฏฺฐ\-นิทฺเทส}
\proseitem{201}{กถํ อธิมตฺต\-ฏฺเฐน อินฺทฺริยานิ ทฏฺฐพฺพานิ? สทฺธิ\-นฺทฺริยสฺส ภาวนาย ฉนฺโท อุปฺปชฺชติ, ฉนฺทวเสน สทฺธาวเสน สทฺธินฺทฺริยํ อธิมตฺตํ โหติ. ฉนฺทวเสน ปาโมชฺชํ อุปฺปชฺชติ, ปาโมชฺชวเสน สทฺธาวเสน สทฺธินฺทฺริยํ อธิมตฺตํ โหติ. ปาโมชฺชวเสน ปีติ อุปฺปชฺชติ, ปีติวเสน สทฺธาวเสน สทฺธินฺทฺริยํ อธิมตฺตํ โหติ. ปีติวเสน ปสฺสทฺธิ อุปฺปชฺชติ, ปสฺสทฺธิ\-วเสน สทฺธาวเสน สทฺธินฺทฺริยํ อธิมตฺตํ โหติ. ปสฺสทฺธิ\-วเสน สุขํ อุปฺปชฺชติ, สุขวเสน สทฺธาวเสน สทฺธินฺทฺริยํ อธิมตฺตํ โหติ. สุขวเสน โอภาโส อุปฺปชฺชติ, โอภาสวเสน สทฺธาวเสน สทฺธินฺทฺริยํ อธิมตฺตํ โหติ. โอภาสวเสน สํเวโค อุปฺปชฺชติ, สํเวควเสน สทฺธาวเสน สทฺธินฺทฺริยํ อธิมตฺตํ โหติ. สํเวเชตฺวา จิตฺตํ สมาทหติ, สมาธิวเสน สทฺธาวเสน สทฺธินฺทฺริยํ อธิมตฺตํ โหติ. ตถาสมาหิตํ จิตฺตํ สาธุกํ ปคฺคณฺหาติ, ปคฺคหวเสน สทฺธาวเสน สทฺธินฺทฺริยํ อธิมตฺตํ โหติ. ตถา\-ปคฺคหิตํ จิตฺตํ สาธุกํ อชฺฌุเปกฺขติ, อุเปกฺขาวเสน สทฺธาวเสน สทฺธินฺทฺริยํ อธิมตฺตํ โหติ. อุเปกฺขาวเสน \csromanfolio{220}นานตฺต\-กิเลเสหิ จิตฺตํ วิมุจฺจติ, วิโมกฺขวเสน สทฺธาวเสน สทฺธินฺทฺริยํ อธิมตฺตํ โหติ. วิมุตฺตตฺตา เต ธมฺมา เอกรสา โหนฺติ, เอกรสฏฺเฐน ภาวนาวเสน สทฺธาวเสน สทฺธินฺทฺริยํ อธิมตฺตํ โหติ. ภาวิตตฺตา ตโต ปณีตตเร วิวฏฺฏนฺติ, วิวฏฺฏนา\-วเสน สทฺธาวเสน สทฺธินฺทฺริยํ อธิมตฺตํ โหติ. วิวฏฺฏิตตฺตา ตโต โวสชฺชติ\csromansharedfootnote{eaaa066ff9d635ea}{โวสฺสชฺชติ (สฺยา, ก)}, โวสคฺควเสน สทฺธาวเสน สทฺธินฺทฺริยํ อธิมตฺตํ โหติ. โวสชฺชิตตฺตา ตโต นิรุชฺฌนฺติ, นิโรธวเสน สทฺธาวเสน สทฺธินฺทฺริยํ อธิมตฺตํ โหติ. นิโรธวเสน เทฺว โวสคฺคา ปริจฺจาค\-โวสคฺโค จ ปกฺขนฺทน\-โวสคฺโค จ. กิเลเส จ ขนฺเธ จ ปริจฺจชตีติ ปริจฺจาค\-โวสคฺโค. นิโรธ\-นิพฺพานธาตุยา จิตฺตํ ปกฺขนฺทตีติ ปกฺขนฺทน\-โวสคฺโค. นิโรธวเสน อิเม เทฺว โวสคฺคา.}

\prose{อสฺสทฺธิยสฺส ปหานาย ฉนฺโท อุปฺปชฺชติ ฯเปฯ. อสฺสทฺธิย\-ปริฬาหสฺส ปหานาย ฉนฺโท อุปฺปชฺชติ. ทิฏฺเฐ\-กฏฺฐานํ กิเลสานํ ปหานาย ฉนฺโท อุปฺปชฺชติ. โอฬาริกานํ กิเลสานํ ปหานาย ฉนฺโท อุปฺปชฺชติ. อณุสหคตานํ กิเลสานํ ปหานาย ฉนฺโท อุปฺปชฺชติ. สพฺพกิเลสานํ ปหานาย ฉนฺโท อุปฺปชฺชติ, ฉนฺทวเสน สทฺธาวเสน สทฺธินฺทฺริยํ อธิมตฺตํ โหติ ฯเปฯ. วีริยิ\-นฺทฺริยสฺส ภาวนาย ฉนฺโท อุปฺปชฺชติ ฯเปฯ. โกสชฺชสฺส ปหานาย ฉนฺโท อุปฺปชฺชติ. โกสชฺช\-ปริฬาหสฺส ปหานาย ฉนฺโท อุปฺปชฺชติ. ทิฏฺเฐ\-กฏฺฐานํ กิเลสานํ ปหานาย ฉนฺโท อุปฺปชฺชติ ฯเปฯ. สพฺพกิเลสานํ ปหานาย ฉนฺโท อุปฺปชฺชติ. สตินฺทฺริยสฺส ภาวนาย ฉนฺโท อุปฺปชฺชติ ฯเปฯ. ปมาทสฺส ปหานาย ฉนฺโท อุปฺปชฺชติ. ปมาท\-ปริฬาหสฺส ปหานาย ฉนฺโท อุปฺปชฺชติ ฯเปฯ. สพฺพกิเลสานํ ปหานาย ฉนฺโท อุปฺปชฺชติ. สมาธิ\-นฺทฺริยสฺส ภาวนาย ฉนฺโท อุปฺปชฺชติ ฯเปฯ. อุทฺธจฺจสฺส ปหานาย ฉนฺโท อุปฺปชฺชติ. อุทฺธจฺจ\-ปริฬาหสฺส ปหานาย ฉนฺโท อุปฺปชฺชติ ฯเปฯ. สพฺพกิเลสานํ ปหานาย ฉนฺโท อุปฺปชฺชติ. ปญฺญินฺทฺริยสฺส ภาวนาย ฉนฺโท อุปฺปชฺชติ ฯเปฯ. อวิชฺชาย ปหานาย ฉนฺโท อุปฺปชฺชติ. อวิชฺชา\-ปริฬาหสฺส ปหานาย ฉนฺโท อุปฺปชฺชติ. ทิฏฺเฐ\-กฏฺฐานํ กิเลสานํ ปหานาย ฉนฺโท อุปฺปชฺชติ. โอฬาริกานํ กิเลสานํ ปหานาย ฉนฺโท อุปฺปชฺชติ. อณุสหคตานํ กิเลสานํ ปหานาย ฉนฺโท อุปฺปชฺชติ. สพฺพกิเลสานํ ปหานาย ฉนฺโท อุปฺปชฺชติ. ฉนฺทวเสน ปญฺญาวเสน ปญฺญินฺทฺริยํ อธิมตฺตํ โหติ. ฉนฺทวเสน ปาโมชฺชํ อุปฺปชฺชติ, ปาโมชฺชวเสน}

\prosewithcloserruled{\csromanfolio{221}ปญฺญาวเสน ปญฺญินฺทฺริยํ อธิมตฺตํ โหติ. ปาโมชฺชวเสน ปีติ อุปฺปชฺชติ, ปีติวเสน ปญฺญาวเสน ปญฺญินฺทฺริยํ อธิมตฺตํ โหติ. ปีติวเสน ปสฺสทฺธิ อุปฺปชฺชติ, ปสฺสทฺธิ\-วเสน ปญฺญาวเสน ปญฺญินฺทฺริยํ อธิมตฺตํ โหติ. ปสฺสทฺธิ\-วเสน สุขํ อุปฺปชฺชติ, สุขวเสน ปญฺญาวเสน ปญฺญินฺทฺริยํ อธิมตฺตํ โหติ. สุขวเสน โอภาโส อุปฺปชฺชติ, โอภาสวเสน ปญฺญาวเสน ปญฺญินฺทฺริยํ อธิมตฺตํ โหติ. โอภาสวเสน สํเวโค อุปฺปชฺชติ, สํเวควเสน ปญฺญาวเสน ปญฺญินฺทฺริยํ อธิมตฺตํ โหติ. สํเวเชตฺวา จิตฺตํ สมาทหติ, สมาธิวเสน ปญฺญาวเสน ปญฺญินฺทฺริยํ อธิมตฺตํ โหติ. ตถาสมาหิตํ จิตฺตํ สาธุกํ ปคฺคณฺหาติ, ปคฺคหวเสน ปญฺญาวเสน ปญฺญินฺทฺริยํ อธิมตฺตํ โหติ. ตถา\-ปคฺคหิตํ จิตฺตํ สาธุกํ อชฺฌุเปกฺขติ, อุเปกฺขาวเสน ปญฺญาวเสน ปญฺญินฺทฺริยํ อธิมตฺตํ โหติ. อุเปกฺขาวเสน นานตฺต\-กิเลเสหิ จิตฺตํ วิมุจฺจติ, วิโมกฺขวเสน ปญฺญาวเสน ปญฺญินฺทฺริยํ อธิมตฺตํ โหติ. วิมุตฺตตฺตา เต ธมฺมา เอกรสา โหนฺติ, ภาวนาวเสน\csromansharedfootnote{43012462123e79ae}{เอกรสฏฺเฐน ภาวนาวเสน (สฺยา, ก) อฏฺฐกถา โอโลเกตพฺโพ.} ปญฺญาวเสน ปญฺญินฺทฺริยํ อธิมตฺตํ โหติ. ภาวิตตฺตา ตโต ปณีตตเร วิวฏฺฏนฺติ, วิวฏฺฏนา\-วเสน ปญฺญาวเสน ปญฺญินฺทฺริยํ อธิมตฺตํ โหติ. วิวฏฺฏิตตฺตา ตโต โวสชฺชติ, โวสคฺควเสน ปญฺญาวเสน ปญฺญินฺทฺริยํ อธิมตฺตํ โหติ. โวสชฺชิตตฺตา ตโต นิรุชฺฌนฺติ, นิโรธวเสน ปญฺญาวเสน ปญฺญินฺทฺริยํ อธิมตฺตํ โหติ. นิโรธวเสน เทฺว โวสคฺคา ปริจฺจาค\-โวสคฺโค จ ปกฺขนฺทน\-โวสคฺโค จ. กิเลเส จ ขนฺเธ จ ปริจฺจชตีติ ปริจฺจาค\-โวสคฺโค. นิโรธ\-นิพฺพานธาตุยา จิตฺตํ ปกฺขนฺทตีติ ปกฺขนฺทน\-โวสคฺโค. นิโรธวเสน อิเม เทฺว โวสคฺคา. เอวํ อธิมตฺต\-ฏฺเฐน อินฺทฺริยานิ ทฏฺฐพฺพานิ.}{ทุติย\-ภาณวาโร.}
\csromanheader{2}{ฆ. อธิฏฺฐานฏฺฐ\-นิทฺเทส}
\proseitem{202}{กถํ อธิฏฺฐาน\-ฏฺเฐน อินฺทฺริยานิ ทฏฺฐพฺพานิ? สทฺธิ\-นฺทฺริยสฺส ภาวนาย ฉนฺโท อุปฺปชฺชติ, ฉนฺทวเสน สทฺธาวเสน สทฺธินฺทฺริยํ อธิฏฺฐาติ. ฉนฺทวเสน ปาโมชฺชํ อุปฺปชฺชติ, ปาโมชฺชวเสน สทฺธาวเสน สทฺธินฺทฺริยํ อธิฏฺฐาติ ฯเปฯ. เอวํ อธิฏฺฐาน\-ฏฺเฐน อินฺทฺริยานิ ทฏฺฐพฺพานิ.}

\vspace{\tipitakaparskip}
\csromancenter{\^{}อ. ปริยาทาน\-ฏฺฐ\-นิทฺเทส}
\prose{\csromanfolio{222}กถํ ปริยาทาน\-ฏฺเฐน อินฺทฺริยานิ ทฏฺฐพฺพานิ? อธิโมกฺข\-ฏฺเฐน สทฺธินฺทฺริยํ อสฺสทฺธิยํ ปริยาทิยติ, อสฺสทฺธิย\-ปริฬาหํ ปริยาทิยติ. ปคฺคหฏฺเฐน วีริยินฺทฺริยํ โกสชฺชํ ปริยาทิยติ, โกสชฺช\-ปริฬาหํ ปริยาทิยติ. อุปฏฺฐาน\-ฏฺเฐน สตินฺทฺริยํ ปมาทํ ปริยาทิยติ, ปมาท\-ปริฬาหํ ปริยาทิยติ. อวิกฺเขป\-ฏฺเฐน สมาธินฺทฺริยํ อุทฺธจฺจํ ปริยาทิยติ, อุทฺธจฺจ\-ปริฬาหํ ปริยาทิยติ. ทสฺสนฏฺเฐน ปญฺญินฺทฺริยํ อวิชฺชํ ปริยาทิยติ, อวิชฺชา\-ปริฬาหํ ปริยาทิยติ. เนกฺขมฺเม ปญฺจินฺทฺริยานิ กามจฺฉนฺทํ ปริยาทิยนฺติ. อพฺยาปาเท ปญฺจินฺทฺริยานิ พฺยาปาทํ ปริยาทิยนฺติ. อาโลกสญฺญาย ปญฺจินฺทฺริยานิ ถินมิทฺธํ ปริยาทิยนฺติ. อวิกฺเขเป ปญฺจินฺทฺริยานิ อุทฺธจฺจํ ปริยาทิยนฺติ ฯเปฯ. อรหตฺตมคฺเค ปญฺจินฺทฺริยานิ สพฺพกิเลเส ปริยาทิยนฺติ. เอวํ ปริยาทาน\-ฏฺเฐน อินฺทฺริยานิ ทฏฺฐพฺพานิ.}

\csromanheader{2}{จ. ปติฏฺฐา\-ปกฏฺฐ\-นิทฺเทส}
\proseitem{203}{กถํ ปติฏฺฐา\-ปกฏฺเฐน อินฺทฺริยานิ ทฏฺฐพฺพานิ? สทฺโธ สทฺธินฺทฺริยํ อธิโมกฺเข ปติฏฺฐาเปติ, สทฺธสฺส สทฺธินฺทฺริยํ อธิโมกฺเข ปติฏฺฐาเปติ. วีริยวา วีริยินฺทฺริยํ ปคฺคเห ปติฏฺฐาเปติ, วีริยวโต วีริยินฺทฺริยํ ปคฺคเห ปติฏฺฐาเปติ. สติมา สตินฺทฺริยํ อุปฏฺฐาเน ปติฏฺฐาเปติ, สติมโต สตินฺทฺริยํ อุปฏฺฐาเน ปติฏฺฐาเปติ. สมาหิโต สมาธินฺทฺริยํ อวิกฺเขเป ปติฏฺฐาเปติ, สมาหิตสฺส สมาธินฺทฺริยํ อวิกฺเขเป ปติฏฺฐาเปติ. ปญฺญวา ปญฺญินฺทฺริยํ ทสฺสเน ปติฏฺฐาเปติ, ปญฺญวโต ปญฺญินฺทฺริยํ ทสฺสเน ปติฏฺฐาเปติ. โยคาวจโร ปญฺจินฺทฺริยานิ เนกฺขมฺเม ปติฏฺฐาเปติ, โยคาวจรสฺส ปญฺจินฺทฺริยานิ เนกฺขมฺเม ปติฏฺฐาเปนฺติ. โยคาวจโร ปญฺจินฺทฺริยานิ อพฺยาปาเท ปติฏฺฐาเปติ, โยคาวจรสฺส ปญฺจินฺทฺริยานิ อพฺยาปาเท ปติฏฺฐาเปนฺติ. โยคาวจโร ปญฺจินฺทฺริยานิ อาโลกสญฺญาย ปติฏฺฐาเปติ, โยคาวจรสฺส ปญฺจินฺทฺริยานิ อาโลกสญฺญาย ปติฏฺฐาเปนฺติ. โยคาวจโร ปญฺจินฺทฺริยานิ อวิกฺเขเป ปติฏฺฐาเปติ, โยคาวจรสฺส ปญฺจินฺทฺริยานิ อวิกฺเขเป ปติฏฺฐาเปนฺติ ฯเปฯ. โยคาวจโร ปญฺจินฺทฺริยานิ อรหตฺตมคฺเค ปติฏฺฐาเปติ, โยคาวจรสฺส ปญฺจินฺทฺริยานิ อรหตฺตมคฺเค ปติฏฺฐาเปนฺติ. เอวํ ปติฏฺฐา\-ปกฏฺเฐน อินฺทฺริยานิ ทฏฺฐพฺพานิ. สุตฺตนฺต\-นิทฺเทโส จตุตฺโถ.}

\csromanmatikaanchor{mtk.77}
\csromantocmarkref{subsection}{5. อินฺทฺริยสโมธาน}{mtk.77}
\bookmark[dest={mtk.77},level=2]{5. อินฺทฺริยสโมธาน}
\csromanheader{2}{5. อินฺทฺริย\-สโมธาน}
\proseitem{204}{\csromanfolio{223}ปุถุชฺชโน สมาธิํ ภาเวนฺโต กติหากาเรหิ อุปฏฺฐาน\-กุสโล โหติ, เสกฺโข สมาธิํ ภาเวนฺโต กติหากาเรหิ อุปฏฺฐาน\-กุสโล โหติ, วีตราโค สมาธิํ ภาเวนฺโต กติหากาเรหิ อุปฏฺฐาน\-กุสโล โหติ? ปุถุชฺชโน สมาธิํ ภาเวนฺโต สตฺตหิ อากาเรหิ อุปฏฺฐาน\-กุสโล โหติ, เสกฺโข สมาธิํ ภาเวนฺโต อฏฺฐหิ อากาเรหิ อุปฏฺฐาน\-กุสโล โหติ, วีตราโค สมาธิํ ภาเวนฺโต ทสหิ อากาเรหิ อุปฏฺฐาน\-กุสโล โหติ.}

\prose{ปุถุชฺชโน สมาธิํ ภาเวนฺโต กตเมหิ สตฺตหิ อากาเรหิ อุปฏฺฐาน\-กุสโล โหติ\csromandash{}อาวชฺชิตตฺตา อารมฺมณู\-ปฏฺฐานกุสโล โหติ, สมถนิมิตฺตู\-ปฏฺฐานกุสโล โหติ, ปคฺคหนิมิตฺตู\-ปฏฺฐานกุสโล โหติ, อวิกฺเขปู\-ปฏฺฐานกุสโล โหติ, โอภาสู\-ปฏฺฐานกุสโล โหติ, สมฺปหํสนู\-ปฏฺฐานกุสโล โหติ, อุเปกฺขู\-ปฏฺฐานกุสโล โหติ. ปุถุชฺชโน สมาธิํ ภาเวนฺโต อิเมหิ สตฺตหิ อากาเรหิ อุปฏฺฐาน\-กุสโล โหติ.}

\prose{เสกฺโข สมาธิํ ภาเวนฺโต กตเมหิ อฏฺฐหิ อากาเรหิ อุปฏฺฐาน\-กุสโล โหติ\csromandash{}อาวชฺชิตตฺตา อารมฺมณู\-ปฏฺฐานกุสโล โหติ, สมถนิมิตฺตู\-ปฏฺฐานกุสโล โหติ, ปคฺคหนิมิตฺตู\-ปฏฺฐานกุสโล โหติ, อวิกฺเขปู\-ปฏฺฐานกุสโล โหติ, โอภาสู\-ปฏฺฐานกุสโล โหติ, สมฺปหํสนู\-ปฏฺฐานกุสโล โหติ, อุเปกฺขู\-ปฏฺฐานกุสโล โหติ, เอกตฺตู\-ปฏฺฐานกุสโล โหติ. เสกฺโข สมาธิํ ภาเวนฺโต อิเมหิ อฏฺฐหิ อากาเรหิ อุปฏฺฐาน\-กุสโล โหติ.}

\prose{วีตราโค สมาธิํ ภาเวนฺโต กตเมหิ ทสหิ อากาเรหิ อุปฏฺฐาน\-กุสโล โหติ\csromandash{}อาวชฺชิตตฺตา อารมฺมณู\-ปฏฺฐานกุสโล โหติ ฯเปฯ เอกตฺตู\-ปฏฺฐานกุสโล โหติ, ญาณู\-ปฏฺฐานกุสโล โหติ, วิมุตฺตู\-ปฏฺฐานกุสโล โหติ. วีตราโค สมาธิํ ภาเวนฺโต อิเมหิ ทสหิ อากาเรหิ อุปฏฺฐาน\-กุสโล โหติ.}

\proseitem{205}{ปุถุชฺชโน วิปสฺสนํ ภาเวนฺโต กติหากาเรหิ อุปฏฺฐาน\-กุสโล โหติ, กติหากาเรหิ อนุป\-ฏฺฐานกุสโล โหติ. เสกฺโข วิปสฺสนํ ภาเวนฺโต กติหากาเรหิ อุปฏฺฐาน\-กุสโล \csromanfolio{224}โหติ, กติหาเรหิ อนุป\-ฏฺฐานกุสโล โหติ. วีตราโค วิปสฺสนํ ภาเวนฺโต กติหากาเรหิ อุปฏฺฐาน\-กุสโล โหติ, กติหากาเรหิ อนุป\-ฏฺฐานกุสโล โหติ?}

\prose{ปุถุชฺชโน วิปสฺสนํ ภาเวนฺโต นวหิ อากาเรหิ อุปฏฺฐาน\-กุสโล โหติ, นวหิ อากาเรหิ อนุป\-ฏฺฐานกุสโล โหติ. เสกฺโข วิปสฺสนํ ภาเวนฺโต ทสหิ อากาเรหิ อุปฏฺฐาน\-กุสโล โหติ, ทสหิ อากาเรหิ อนุป\-ฏฺฐานกุสโล โหติ. วีตราโค วิปสฺสนํ ภาเวนฺโต ทฺวาทสหิ อากาเรหิ อุปฏฺฐาน\-กุสโล โหติ, ทฺวาทสหิ อากาเรหิ อนุป\-ฏฺฐานกุสโล โหติ.}

\prose{ปุถุชฺชโน วิปสฺสนํ ภาเวนฺโต กตเมหิ นวหิ อากาเรหิ อุปฏฺฐาน\-กุสโล โหติ, กตเมหิ นวหิ อากาเรหิ อนุป\-ฏฺฐานกุสโล โหติ? อนิจฺจโต อุปฏฺฐาน\-กุสโล โหติ, นิจฺจโต อนุป\-ฏฺฐานกุสโล โหติ. ทุกฺขโต อุปฏฺฐาน\-กุสโล โหติ, สุขโต อนุป\-ฏฺฐานกุสโล โหติ. อนตฺตโต อุปฏฺฐาน\-กุสโล โหติ, อตฺตโต อนุป\-ฏฺฐานกุสโล โหติ. ขยโต อุปฏฺฐาน\-กุสโล โหติ, ฆนโต อนุป\-ฏฺฐานกุสโล โหติ. วยโต อุปฏฺฐาน\-กุสโล โหติ, อายูหนานุ\-ปฏฺฐานกุสโล โหติ. วิปริณามู\-ปฏฺฐานกุสโล โหติ, ธุวโต อนุป\-ฏฺฐานกุสโล โหติ. อนิมิตฺตู\-ปฏฺฐานกุสโล โหติ, นิมิตฺตานุ\-ปฏฺฐานกุสโล โหติ. อปฺปณิหิตู\-ปฏฺฐานกุสโล โหติ, ปณิธิ\-อนุป\-ฏฺฐานกุสโล\csromansharedfootnote{4c6316b2f17f70ac}{ปณิธานุปฏฺฐานกุสโล (สฺยา)} โหติ, สุญฺญตู\-ปฏฺฐานกุสโล โหติ, อภินิเวสานุ\-ปฏฺฐานกุสโล โหติ. ปุถุชฺชโน วิปสฺสนํ ภาเวนฺโต อิเมหิ นวหิ อากาเรหิ อุปฏฺฐาน\-กุสโล โหติ, อิเมหิ นวหิ อากาเรหิ อนุป\-ฏฺฐานกุสโล โหติ.}

\proseitem{206}{เสกฺโข วิปสฺสนํ ภาเวนฺโต กตเมหิ ทสหิ อากาเรหิ อุปฏฺฐาน\-กุสโล โหติ, กตเมหิ ทสหิ อากาเรหิ อนุป\-ฏฺฐานกุสโล โหติ? อนิจฺจโต อุปฏฺฐาน\-กุสโล โหติ, นิจฺจโต อนุป\-ฏฺฐานกุสโล โหติ ฯเปฯ. สุญฺญตู\-ปฏฺฐานกุสโล โหติ, อภินิเวสานุ\-ปฏฺฐานกุสโล โหติ. ญาณู\-ปฏฺฐานกุสโล โหติ อญาณ\-อนุป\-ฏฺฐานกุสโล โหติ, เสกฺโข วิปสฺสนํ ภาเวนฺโต}

\prose{\csromanfolio{225}อิเมหิ ทสหิ อากาเรหิ อุปฏฺฐาน\-กุสโล โหติ, อิเมหิ ทสหิ อากาเรหิ อนุป\-ฏฺฐานกุสโล โหติ.}

\prose{วีตราโค วิปสฺสนํ ภาเวนฺโต กตเมหิ ทฺวาทสหิ อากาเรหิ อุปฏฺฐาน\-กุสโล โหติ, กตเมหิ ทฺวาทสหิ อากาเรหิ อนุป\-ฏฺฐานกุสโล โหติ? อนิจฺจโต อุปฏฺฐาน\-กุสโล โหติ, นิจฺจโต อนุป\-ฏฺฐานกุสโล โหติ ฯเปฯ. ญาณู\-ปฏฺฐานกุสโล โหติ, อญาณานุ\-ปฏฺฐานกุสโล โหติ. วิสญฺโญคู\-ปฏฺฐานกุสโล โหติ, สญฺโญคานุ\-ปฏฺฐานกุสโล โหติ. นิโรธู\-ปฏฺฐานกุสโล โหติ, สงฺขารานุ\-ปฏฺฐานกุสโล โหติ. วีตราโค วิปสฺสนํ ภาเวนฺโต อิเมหิ ทฺวาทสหิ อากาเรหิ อุปฏฺฐาน\-กุสโล โหติ, อิเมหิ ทฺวาทสหิ อากาเรหิ อนุป\-ฏฺฐานกุสโล โหติ.}

\prose{อาวชฺชิตตฺตา อารมฺมณู\-ปฏฺฐานกุสล\-วเสน อินฺทฺริยานิ สโมธาเนติ โคจรญฺจ ปชานาติ. สมตฺถญฺจ ปฏิวิชฺฌติ ฯเปฯ ธมฺเม สโมธาเนติ. โคจรญฺจ ปชานาติ สมตฺถญฺจ ปฏิวิชฺฌติ. อินฺทฺริยานิ สโมธาเนตีติ กถํ อินฺทฺริยานิ สโมธาเนติ. อธิโมกฺข\-ฏฺเฐน สทฺธินฺทฺริยํ สโมธาเนติ ฯเปฯ. สมถนิมิตฺตู\-ปฏฺฐานกุสล\-วเสน. ปคฺคหนิมิตฺตู\-ปฏฺฐานกุสล\-วเสน. อวิกฺเขปู\-ปฏฺฐานกุสล\-วเสน. โอภาสู\-ปฏฺฐานกุสล\-วเสน. สมฺปหํสนู\-ปฏฺฐานกุสล\-วเสน. อุเปกฺขู\-ปฏฺฐานกุสล\-วเสน. เอกตฺตู\-ปฏฺฐานกุสล\-วเสน. ญาณู\-ปฏฺฐานกุสล\-วเสน. วิมุตฺตู\-ปฏฺฐานกุสล\-วเสน.}

\prose{อนิจฺจโต อุปฏฺฐานกุสล\-วเสน. นิจฺจโต อนุป\-ฏฺฐานกุสล\-วเสน. ทุกฺขโต อุปฏฺฐานกุสล\-วเสน. สุขโต อนุป\-ฏฺฐานกุสล\-วเสน. อนตฺตโต อุปฏฺฐานกุสล\-วเสน. อตฺตโต อนุป\-ฏฺฐานกุสล\-วเสน, ขยโต อุปฏฺฐานกุสล\-วเสน. ฆนโต อนุป\-ฏฺฐานกุสล\-วเสน. วยโต อุปฏฺฐานกุสล\-วเสน. อายูหนานุ\-ปฏฺฐานกุสล\-วเสน. วิปริณามู\-ปฏฺฐานกุสล\-วเสน. ธุวโต อนุป\-ฏฺฐานกุสล\-วเสน. อนิมิตฺตู\-ปฏฺฐานกุสล\-วเสน. นิมิตฺตานุ\-ปฏฺฐานกุสล\-วเสน. อปฺปณิหิตู\-ปฏฺฐานกุสล\-วเสน. ปณิธิ\-อนุ\-ปฏฺฐาน\-กุสลวเสน. สุญฺญตู\-ปฏฺฐานกุสล\-วเสน. อภินิเวสานุ\-ปฏฺฐานกุสล\-วเสน. ญาณู\-ปฏฺฐานกุสล\-วเสน. อญาณานุ\-ปฏฺฐานกุสล\-วเสน. วิสญฺโญ รูปฏฺฐานกุสลวเสน.}

\prose{\csromanfolio{226}สญฺโญคานุ\-ปฏฺฐานกุสล\-วเสน. นิโรธู\-ปฏฺฐานกุสล\-วเสน. สงฺขารานุ\-ปฏฺฐานกุสล\-วเสน. อินฺทฺริยานิ สโมธาเนติ. โคจรญฺจ ปชานาติ. สมตฺถญฺจ ปฏิวิชฺฌติ.}

\proseitem{207}{จตุสฏฺฐิยา อากาเรหิ ติณฺณนฺนํ อินฺทฺริยานํ วสิภาวตา ปญฺญา อาสวานํ ขเย ญาณํ\csromandash{}กตเมสํ ติณฺณนฺนํ อินฺทฺริยานํ? อนญฺญาต\-ญฺญสฺสามี\-ตินฺทฺริยสฺส อญฺญินฺทฺริยสฺส อญฺญาตาวิ\-นฺทฺริยสฺส.}

\prose{อนญฺญาต\-ญฺญสฺสามี\-ตินฺทฺริยํ กติ ฐานานิ คจฺฉติ, อญฺญินฺทฺริยํ กติ ฐานานิ คจฺฉติ, อญฺญาตาวิ\-นฺทฺริยํ กติ ฐานานิ คจฺฉติ? อนญฺญาต\-ญฺญสฺสามี\-ตินฺทฺริยํ เอกํ ฐานํ คจฺฉติ, โสตาปตฺติ\-มคฺคํ. อญฺญินฺทฺริยํ ฉ ฐานานิ คจฺฉติ, โสตาปตฺติผลํ สกทาคามิ\-มคฺคํ สกทามิผลํ อนาคามิมคฺคํ อนาคามิผลํ อรหตฺตมคฺคํ. อญฺญาตาวิ\-นฺทฺริยํ เอกํ ฐานํ คจฺฉติ, อรหตฺตผลํ.}

\prose{โสตาปตฺติมคฺค\-กฺขเณ อนญฺญาต\-ญฺญสฺสามี\-ตินฺทฺริยสฺส สทฺธินฺทฺริยํ อธิโมกฺข\-ปริวารํ โหติ. วีริยินฺทฺริยํ ปคฺคห\-ปริวารํ โหติ. สตินฺทฺริยํ อุปฏฺฐาน\-ปริวารํ โหติ. สมาธินฺทฺริยํ อวิกฺเขป\-ปริวารํ โหติ. ปญฺญินฺทฺริยํ ทสฺสน\-ปริวารํ โหติ. มนินฺทฺริยํ วิชานน\-ปริวารํ โหติ. โสมนสฺสิ\-นฺทฺริยํ อภิสนฺทน\-ปริวารํ โหติ. ชีวิตินฺทฺริยํ ปวตฺต\-สนฺตตา\-ธิปเตยฺย\-ปริวารํ โหติ. โสตาปตฺติมคฺค\-กฺขเณ ชาตา ธมฺมา ฐเปตฺวา จิตฺต\-สมุฏฺฐานํ รูปํ สพฺเพว กุสลา โหนฺติ. สพฺเพว อนาสวา โหนฺติ. สพฺเพว นิยฺยานิกา โหนฺติ สพฺเพว อปจยคามิโน โหนฺติ. สพฺเพว โลกุตฺตรา โหนฺติ. สพฺเพว นิพฺพานารมฺมณา โหนฺติ. โสตาปตฺติมคฺค\-กฺขเณ อนญฺญาต\-ญฺญสฺสามี\-ตินฺทฺริยสฺส อิมานิ อฏฺฐินฺทฺริยานิ สหชาต\-ปริวารา โหนฺติ, อญฺญมญฺญ\-ปริวารา โหนฺติ, นิสฺสย\-ปริวารา โหนฺติ, สมฺปยุตฺต\-ปริวารา โหนฺติ, สหคตา โหนฺติ, สหชาตา โหนฺติ, สํสฏฺฐา โหนฺติ, สมฺปยุตฺตา โหนฺติ, เตว ตสฺส อาการา เจว โหนฺติ ปริวารา จ.}

\prose{โสตาปตฺติ\-ผลกฺขเณ ฯเปฯ. อรหตฺต\-ผลกฺขเณ อญฺญาตาวิ\-นฺทฺริยสฺส สทฺธินฺทฺริยํ อธิโมกฺข\-ปริวารํ โหติ ฯเปฯ ชีวิตินฺทฺริยํ ปวตฺต\-สนฺตตา\-ธิปเตยฺย\-ปริวารํ โหติ, อรหตฺต\-ผลกฺขเณ ชาตา ธมฺมา สพฺเพว อพฺยากตา โหนฺติ, ฐเปตฺวา จิตฺต\-สมุฏฺฐานํ รูปํ สพฺเพว อนาสวา โหนฺติ สพฺเพว โลกุตฺตรา โหนฺติ สพฺเพว นิพฺพานารมฺมณา โหนฺติ. อรหตฺต\-ผลกฺขเณ อญฺญาตาวิ\-นฺทฺริยสฺส อิมานิ อฏฺฐินฺทฺริยานิ สหชาต\-ปริวารา}

\prose{\csromanfolio{227}โหนฺติ, เตว ตสฺส อาการา เจว โหนฺติ ปริวารา จ. อิติ อิมานิ อฏฺฐ อฏฺฐกานิ จตุสฏฺฐิ โหนฺติ.}

\prose{อาสวาติ กตเม เต อาสวา? กามาสโว ภวาสโว ทิฏฺฐาสโว อวิชฺชาสโว. กตฺเถเต อาสวา ขียนฺติ\csromansharedfootnote{e868022c001f068d}{ขิยฺยนฺติ (ก)}? โสตาปตฺติ\-มคฺเคน อนวเสโส ทิฏฺฐาสโว ขียติ, อปายคมนีโย กามาสโว ขียติ, อปายคมนีโย ภวาสโว ขียติ, อปายคมนีโย อวิชฺชาสโว ขียติ, เอตฺเถเต อาสวา ขียนฺติ. สกทาคามิ\-มคฺเคน โอฬาริโก กามาสโว ขียติ, ตเทกฏฺโฐ ภวาสโว ขียติ, ตเทกฏฺโฐ อวิชฺชาสโว ขียติ, เอตฺเถเต อาสวา ขียนฺติ. อนาคามิมคฺเคน อนวเสโส กามาสโว ขียติ, ตเทกฏฺโฐ ภวาสโว ขียติ, ตเทกฏฺโฐ อวิชฺชาสโว ขียติ, เอตฺเถเต อาสวา ขียนฺติ. อรหตฺต\-มคฺเคน อนวเสโส ภวาสโว ขียติ, อนวเสโส อวิชฺชาสโว ขียติ, เอตฺเถเต อาสวา ขียนฺติ.}

\setlength{\csromangathapagewidth}{0pt}
\setlength{\csromangathawakwidth}{0pt}
\csromangathameasuregroup{}{น ตสฺส อทฺทิฏฺฐมิธตฺถิกิญฺจิ, \\ อโถ อวิญฺญาตมชานิตพฺพํ. \\ สพฺพํ อภิญฺญาสิ ยทตฺถิ เนยฺยํ, \\ ตถาคโต เตน สมนฺตจกฺขูติ.}
\csromangathameasurewak{น ตสฺส อทฺทิฏฺฐมิธตฺถิกิญฺจิ, \\ อโถ อวิญฺญาตมชานิตพฺพํ. \\ สพฺพํ อภิญฺญาสิ ยทตฺถิ เนยฺยํ, \\ ตถาคโต เตน สมนฺตจกฺขูติ.}
\setlength{\csromangathaleftcol}{0pt}
\csromangathagroup{\csromangathastanza{\csromangathawak{\csromangathaitemline{208}{น ตสฺส อทฺทิฏฺฐมิธตฺถิกิญฺจิ,} \\ \csromangathacontline{อโถ อวิญฺญาตมชานิตพฺพํ.} \\ \csromangathacontline{สพฺพํ อภิญฺญาสิ ยทตฺถิ เนยฺยํ,} \\ \csromangathacontline{ตถาคโต เตน สมนฺต\-จกฺขูติ.}}}}

\prose{สมนฺต\-จกฺขูติ เกนฏฺเฐน สมนฺตจกฺขุ? จุทฺทส พุทฺธญาณานิ ทุกฺเข ญาณํ พุทฺธญาณํ, ทุกฺขสมุทเย ญาณํ พุทฺธญาณํ ฯเปฯ สพฺพญฺญุต\-ญฺญาณํ พุทฺธญาณํ, อนาวรณญาณํ พุทฺธญาณํ, อิมานิ จุทฺทส พุทฺธญาณานิ. อิเมสํ จุทฺทสนฺนํ พุทฺธญาณานํ อฏฺฐ ญาณานิ สาวก\-สาธารณานิ ฉ ญาณานิ อสาธารณานิ สาวเกหิ.}

\prose{ยาวตา ทุกฺขสฺส ทุกฺขฏฺโฐ ญาโต, อญฺญาโต ทุกฺขฏฺโฐ นตฺถีติ สมนฺตจกฺขุ. ยํ สมนฺตจกฺขุ ตํ ปญฺญินฺทฺริยํ, ปญฺญินฺทฺริยสฺส วเสน อธิโมกฺข\-ฏฺเฐน สทฺธินฺทฺริยํ, วคฺคหฏฺเฐน วีริยินฺทฺริยํ, อุปฏฺฐาน\-ฏฺเฐน สตินฺทฺริยํ, อวิกฺเขป\-ฏฺเฐน สมาธินฺทฺริยํ. ยาวตา ทุกฺขสฺส ทุกฺขฏฺโฐ ทิฏฺโฐ วิทิโต สจฺฉิกโต ผสฺสิโต ปญฺญาย, อผสฺสิโต ปญฺญาย ทุกฺขฏฺโฐ นตฺถีติ สมนฺตจกฺขุ. ยํ สมนฺตจกฺขุ, ตํ ปญฺญินฺทฺริยํ, ปญฺญินฺทฺริยสฺส วเสน อธิโมกฺข\-ฏฺเฐน สทฺธินฺทฺริยํ, ปคฺคหฏฺเฐน วีริยินฺทฺริยํ, อุปฏฺฐาน\-ฏฺเฐน สตินฺทฺริยํ, อวิกฺเขป\-ฏฺเฐน สมาธินฺทฺริยํ. \csromanfolio{228}ยาวตา สมุทยสฺส สมุทยฏฺโฐ ฯเปฯ. ยาวตา นิโรธสฺส นิโรธฏฺโฐ. ยาวตา มคฺคสฺส มคฺคฏฺโฐ. ยาวตา อตฺถ\-ปฏิสมฺภิทาย อตฺถปฏิสมฺภิท\-ฏฺโฐ. ยาวตา ธมฺม\-ปฏิสมฺภิทาย ธมฺมปฏิสมฺภิท\-ฏฺโฐ. ยาวตา นิรุตฺติ\-ปฏิสมฺภิทาย นิรุตฺติปฏิสมฺภิท\-ฏฺโฐ. ยาวตา ปฏิภาน\-ปฏิสมฺภิทาย ปฏิภานปฏิสมฺภิท\-ฏฺโฐ. ยาวตา อินฺทฺริยปโรปริยตฺต\-ญาณํ. ยาวตา สตฺตานํ อาสยานุสเย ญาณํ. ยาวตา ยมกปาฏิหีเร ญาณํ. ยาวตา มหากรุณา\-สมาปตฺติยา ญาณํ. ยาวตา สเทวกสฺส โลกสฺส สมารกสฺส สพฺรหฺมกสฺส สสฺสมณ\-พฺราหฺมณิยา ปชาย สเทว\-มนุสฺสาย ทิฏฺฐํ สุตํ มุตํ วิญฺญาตํ ปตฺตํ ปริเยสิตํ อนุวิจริตํ มนสา, ตํ ญาตํ ทิฏฺฐํ วิทิตํ สจฺฉิกตํ ผสฺสิตํ ปญฺญาย, อผสฺสิตํ ปญฺญาย นตฺถีติ สมนฺตจกฺขุ. ยํ สมนฺตจกฺขุ, ตํ ปญฺญินฺทฺริยํ, ปญฺญินฺทฺริยสฺส วเสน อธิโมกฺข\-ฏฺเฐน สทฺธินฺทฺริยํ, ปคฺคหฏฺเฐน วีริยินฺทฺริยํ, อุปฏฺฐาน\-ฏฺเฐน สตินฺทฺริยํ, อวิกฺเขป\-ฏฺเฐน สมาธินฺทฺริยํ.}

\prose{สทฺทหนฺโต ปคฺคณฺหาติ, ปคฺคณฺหนฺโต สทฺทหติ, สทฺทหนฺโต อุปฏฺฐาเปติ, อุปฏฺฐาเปนฺโต สทฺทหติ, สทฺทหนฺโต สมาทหติ, สมาทหนฺโต สทฺทหติ, สทฺทหนฺโต ปชานาติ, ปชานนฺโต สทฺทหติ. ปคฺคณฺหนฺโต อุปฏฺฐาเปติ, อุปฏฺฐาเปนฺโต ปคฺคณฺหาติ, ปคฺคณฺหนฺโต สมาทหติ, สมาทหนฺโต ปคฺคณฺหาติ, ปคฺคณฺหนฺโต ปชานาติ, ปชานนฺโต ปคฺคณฺหาติ, ปคฺคณฺหนฺโต สทฺทหติ, สทฺทหนฺโต ปคฺคณฺหาติ. อุปฏฺฐาเปนฺโต สมาทหติ, สมาทหนฺโต อุปฏฺฐาเปติ, อุปฏฺฐาเปนฺโต ปชานาติ, ปชานนฺโต อุปฏฺฐาเปติ, อุปฏฺฐาเปนฺโต สทฺทหติ, สทฺทหนฺโต อุปฏฺฐาเปติ, อุปฏฺฐาเปนฺโต ปคฺคณฺหาติ, ปคฺคณฺหนฺโต อุปฏฺฐาเปติ. สมาทหนฺโต ปชานาติ, ปชานนฺโต สมาทหติ, สมาทหนฺโต สทฺทหติ, สทฺทหนฺโต สมาทหติ, สมาทหนฺโต ปคฺคณฺหาติ, ปคฺคณฺหนฺโต สมาทหติ, สมาทหนฺโต อุปฏฺฐาเปติ, อุปฏฺฐาเปนฺโต สมาทหติ. ปชานนฺโต สทฺทหติ, สทฺทหนฺโต ปชานาติ, ปชานนฺโต ปคฺคณฺหาติ, ปคฺคณฺหนฺโต ปชานาติ, ปชานนฺโต อุปฏฺฐาเปติ, อุปฏฺฐาเปนฺโต ปชานาติ, ปชานนฺโต สมาทหติ, สมาทหนฺโต ปชานาติ.}

\prose{สทฺทหิตตฺตา ปคฺคหิตํ, ปคฺคหิตตฺตา สทฺทหิตํ, สทฺทหิตตฺตา อุปฏฺฐาปิตํ, อุปฏฺฐาปิตตฺตา สทฺทหิตํ, สทฺทหิตตฺตา สมาทหิตํ, สมาทหิตตฺตา สทฺทหิตํ สทฺทหิตตฺตา ปชานิตํ, ปชานิตตฺตา สทฺทหิตํ. ปคฺคหิตตฺตา}

\prose{\csromanfolio{229}อุปฏฺฐาปิตํ, อุปฏฺฐาปิตตฺตา ปคฺคหิตํ, ปคฺคติตตฺตา สมาทหิตํ, สมาทหิตตฺตา ปคฺคหิตํ, ปคฺคหิตตฺตา ปชานิตํ, ปชานิตตฺตา ปคฺคหิตํ, ปคฺคหิตตฺตา สทฺทหิตํ, สทฺทหิตตฺตา ปคฺคหิตํ. อุปฏฺฐาปิตตฺตา สมาทหิตํ, สมาทหิตตฺตา อุปฏฺฐาปิตํ, อุปฏฺฐาปิตตฺตา ปชานิตํ, ปชานิตตฺตา อุปฏฺฐาปิตํ, อุปฏฺฐาปิตตฺตา สทฺทหิตํ, สทฺทหิตตฺตา อุปฏฺฐาปิตํ, อุปฏฺฐาปิตตฺตา ปคฺคหิตํ, ปคฺคหิตตฺตา อุปฏฺฐาปิตํ. สมาทหิตตฺตา ปชานิตํ, ปชานิตตฺตา สมาทหิตํ, สมาทหิตตฺตา สทฺทหิตํ, สทฺทหิตตฺตา สมาทหิตํ, สมาทหิตตฺตา ปคฺคหิตํ, ปคฺคหิตตฺตา สมาทหิตํ, สมาทหิตตฺตา อุปฏฺฐาปิตํ, อุปฏฺฐาปิตตฺตา สมาทหิตํ. ปชานิตตฺตา สทฺทหิตํ, สทฺทหิตตฺตา ปชานิตํ, ปชานิตตฺตา ปคฺคหิตํ, ปคฺคหิตตฺตา ปชานิตํ, ปชานิตตฺตา อุปฏฺฐาปิตํ, อุปฏฺฐาปิตตฺตา ปชานิตํ, ปชานิตตฺตา สมาทหิตํ, สมาทหิตตฺตา ปชานิตํ.}

\prose{ยํ พุทฺธจกฺขุ, ตํ พุทฺธญาณํ. ยํ พุทฺธญาณํ, ตํ พุทฺธจกฺขุ. เยน จกฺขุนา ตถาคโต สตฺเต ปสฺสติ อปฺปรชกฺเข มหารชกฺเข ติกฺขินฺทฺริเย มุทินฺทฺริเย สฺวากาเร ทฺวากาเร สุวิญฺญาปเย ทุวิญฺญาปเย อปฺเปกจฺเจ ปรโลก\-วชฺช\-ภยทสฺสาวิโน\csromansharedfootnote{03c59bb8c660ed57}{ปรโลกวชฺชภยทสฺสาวิเน (สฺยา, ก)} อปฺเปกจฺเจ นปรโลก\-วชฺช\-ภยทสฺสาวิโน.}

\prose{\textbf{อปฺปรชกฺเข มหารชกฺเข}ติ สทฺโธ ปุคฺคโล อปฺปรชกฺโข, อสฺสทฺโธ ปุคฺคโล มหารชกฺโข. อารทฺธวีริโย ปุคฺคโล อปฺปรชกฺโข, กุสีโต ปุคฺคโล มหารชกฺโข. อุปฏฺฐิตสฺสติ ปุคฺคโล อปฺปรชกฺโข, มุฏฺฐสฺสติ ปุคฺคโล มหารชกฺโข. สมาหิโต ปุคฺคโล อปฺปรชกฺโข, อสมาหิโต ปุคฺคโล มหารชกฺโข. ปญฺญวา ปุคฺคโล อปฺปรชกฺโข, ทุปฺปญฺโญ ปุคฺคโล มหารชกฺโข.}

\prose{\textbf{ติกฺขินฺทฺริเย มุทินฺทฺริเย}ติ สทฺโธ ปุคฺคโล ติกฺขินฺทฺริโย, อสฺสทฺโธ ปุคฺคโล มุทินฺทฺริโย ฯเปฯ. ปญฺญวา ปุคฺคโล ติกฺขินฺทฺริโย, ทุปฺปญฺโญ ปุคฺคโล มุทินฺทฺริโย. \textbf{สฺวากาเร ทฺวากาเร}ติ สทฺโธ ปุคฺคโล สฺวากาโร, อสฺสทฺโธ ปุคฺคโล ทฺวากาโร ฯเปฯ. ปญฺญวา ปุคฺคโล สฺวากาโร, ทุปฺปญฺโญ ปุคฺคโล ทฺวากาโร. \textbf{สุวิญฺญาปเย ทุวิญฺญาปเย}ติ สทฺโธ ปุคฺคโล สุวิญฺญาปโย, อสฺสทฺโธ ปุคฺคโล ทุวิญฺญาปโย ฯเปฯ. ปญฺญวา ปุคฺคโล สุวิญฺญาปโย, ทุปฺปญฺโญ ปุคฺคโล ทุวิญฺญาปโย.}

\prose{\csromanfolio{230}อปฺเปกจฺเจ ปรโลก\-วชฺช\-ภยทสฺสาวิโน อปฺเปกจฺเจ นปรโลก\-วชฺช\-ภยทสฺสาวิโนติ สทฺโธ ปุคฺคโล ปรโลก\-วชฺช\-ภยทสฺสาวี, อสฺสทฺโธ ปุคฺคโล นปรโลก\-วชฺช\-ภยทสฺสาวี. อารทฺธวีริโย ปุคฺคโล ปรโลก\-วชฺช\-ภยทสฺสาวี, กุสีโต ปุคฺคโล นปรโลก\-วชฺช\-ภยทสฺสาวี ฯเปฯ. ปญฺญวา ปุคฺคโล ปรโลก\-วชฺช\-ภยทสฺสาวี, ทุปฺปญฺโญ ปุคฺคโล นปรโลก\-วชฺช\-ภยทสฺสาวี.}

\prose{โลโกติ ขนฺธโลโก ธาตุโลโก อายตนโลโก วิปตฺติ\-ภว\-โลโก วิปตฺติ\-สมฺภว\-โลโก สมฺปตฺติ\-ภว\-โลโก สมฺปตฺติ\-สมฺภว\-โลโก.}

\prose{เอโก โลโก, สพฺเพ สตฺตา อาหารฏฺฐิติกา. เทฺว โลกา, นามญฺจ รูปญฺจ. ตโย โลกา, ติสฺโส เวทนา. จตฺตาโร โลกา, จตฺตาโร อาหารา. ปญฺจ โลกา, ปญฺจุ\-ปาทาน\-กฺขนฺธา. ฉ โลกา, ฉ อชฺฌตฺติกานิ อายตนานิ. สตฺต โลกา, สตฺต วิญฺญาณ\-ฏฺฐิติโย. อฏฺฐ โลกา, อฏฺฐ โลกธมฺมา. นว โลกา, นว สตฺตาวาสา. ทส โลกา, ทสายตนานิ. ทฺวาทส โลกา, ทฺวาทสา\-ยตนานิ. อฏฺฐารส โลกา, อฏฺฐารส ธาตุโย.}

\prose{\textbf{วชฺชนฺ}ติ สพฺเพ กิเลสา วชฺชา, สพฺเพ ทุจฺจริตา วชฺชา, สพฺเพ อภิสงฺขารา วชฺชา, สพฺเพ ภวคามิกมฺมา วชฺชา. อิติ อิมสฺมิญฺจ โลเก อิมสฺมิญฺจ วชฺเช ติพฺพา ภยสญฺญา ปจฺจุปฏฺฐิตา โหนฺติ เสยฺยถาปิ อุกฺขิตฺตาสิเก วธเก. อิเมหิ ปญฺญาสาย อากาเรหิ อิมานิ ปญฺจินฺทฺริยานิ ชานาติ ปสฺสติ อญฺญาติ ปฏิวิชฺฌตีติ.}

\vspace{\tipitakaparskip}
\csromancenter{ตติย\-ภาณวาโร.}
\nitthitamruled{อินฺทฺริยกถา นิฏฺฐิตา.}
\csromanmatikaanchor{mtk.78}
\csromantocmarkref{section}{5. วิโมกฺขกถา}{mtk.78}
\bookmark[dest={mtk.78},level=1]{5. วิโมกฺขกถา}
\csromanheader{1}{5. \textbf{วิโมกฺขกถา}}
\csromanheader{2}{\textbf{อุทฺเทส}}
\proseitem{209}{ปุริมนิทานํ, ตโยเม ภิกฺขเว วิโมกฺขา. กตเม ตโย? สุญฺญโต วิโมกฺโข อนิมิตฺโต วิโมกฺโข อปฺปณิหิโต วิโมกฺโข, อิเม โข ภิกฺขเว ตโย วิโมกฺขา.}

\prose{\csromanfolio{231}อปิ จ อฏฺฐสฏฺฐิ วิโมกฺขา\csromandash{}สุญฺญโต วิโมกฺโข, อนิมิตฺโต วิโมกฺโข, อปฺปณิหิโต วิโมกฺโข. อชฺฌตฺต\-วุฏฺฐาโน วิโมกฺโข, พหิทฺธา\-วุฏฺฐาโน วิโมกฺโข, ทุภโต วุฏฺฐาโน วิโมกฺโข. อชฺฌตฺต\-วุฏฺฐานา จตฺตาโร วิโมกฺขา, พหิทฺธา\-วุฏฺฐานา จตฺตาโร วิโมกฺขา, ทุภโต วุฏฺฐานา จตฺตาโร วิโมกฺขา. อชฺฌตฺต\-วุฏฺฐานานํ อนุโลมา จตฺตาโร วิโมกฺขา, พหิทฺธา\-วุฏฺฐานานํ อนุโลมา จตฺตาโร วิโมกฺขา, ทุภโต วุฏฺฐานานํ อนุโลมา จตฺตาโร วิโมกฺขา. อชฺฌตฺตวุฏฺฐาน\-ปฏิปฺปสฺสทฺธี\csromansharedfootnote{19407d2c105fbc9e}{อชฺฌตฺตวุฏฺฐานา ปฏิปฺปสฺสทฺธี (สฺยา)} จตฺตาโร วิโมกฺขา, พหิทฺธาวุฏฺฐาน\-ปฏิปฺปสฺสทฺธี จตฺตาโร วิโมกฺขา, ทุภโต วุฏฺฐาน\-ปฏิปฺปสฺสทฺธี จตฺตาโร วิโมกฺขา. รูปี รูปานิ ปสฺสตีติ วิโมกฺโข, อชฺฌตฺตํ อรูปสญฺญี พหิทฺธา รูปานิ ปสฺสตีติ วิโมกฺโข, “สุภนฺ”เตว อธิมุตฺโต โหตีติ วิโมกฺโข. อากาสา\-นญฺจา\-ยตนสมาปตฺติ วิโมกฺโข, วิญฺญาณญฺจา\-ยตนสมาปตฺติ วิโมกฺโข, อากิญฺจญฺญายตน\-สมาปตฺติ วิโมกฺโข. เนว\-สญฺญา\-นา\-สญฺญา\-ยตนสมาปตฺติ วิโมกฺโข, สญฺญา\-เวทยิต\-นิโรธ\-สมาปตฺติ วิโมกฺโข. สมยวิโมกฺโข อสมย\-วิโมกฺโข. สามยิโก วิโมกฺโข, อสามยิโก วิโมกฺโข. กุปฺโป วิโมกฺโข, อกุปฺโป วิโมกฺโข. โลกิโก วิโมกฺโข, โลกุตฺตโร วิโมกฺโข. สาสโว วิโมกฺโข, อนาสโว วิโมกฺโข. สามิโส วิโมกฺโข, นิรามิโส วิโมกฺโข, นิรามิสา นิรามิสตโร วิโมกฺโข, ปณิหิโต วิโมกฺโข, อปฺปณิหิโต วิโมกฺโข, ปณิหิต\-ปฺปฏิปฺปสฺสทฺธิ วิโมกฺโข. สญฺญุตฺโต วิโมกฺโข, วิสญฺญุตฺโต วิโมกฺโข. เอกตฺตวิโมกฺโข, นานตฺต\-วิโมกฺโข, สญฺญาวิโมกฺโข. ญาณวิโมกฺโข, สีติสิยา\-วิโมกฺโข\csromansharedfootnote{f831f6cadda98b6a}{สีตีสิยาวิโมกฺโข (สีฏฺฐ)}, ฌานวิโมกฺโข, อนุปาทา จิตฺตสฺส วิโมกฺโข. (75)}

\csromanheader{2}{\textbf{นิทฺเทส}}
\proseitem{210}{กตโม \textbf{สุญฺญโต วิโมกฺโข?} อิธ ภิกฺขุ อรญฺญคโต วา รุกฺขมูลคโต วา สุญฺญาคารคโต วา อิติ ปฏิสญฺจิกฺขติ “สุญฺญมิทํ อตฺเตน วา อตฺตนิเยน วา”ติ. โส ตตฺถ อภินิเวสํ น กโรตีติ สุญฺญโต วิโมกฺโข. อยํ สุญฺญโต วิโมกฺโข.}

\prose{\csromanfolio{232}กตโม \textbf{อนิมิตฺโต วิโมกฺโข?} อิธ ภิกฺขุ อรญฺญคโต วา รุกฺขมูลคโต วา สุญฺญาคารคโต วา อิติ ปฏิสญฺจิกฺขติ “สุญฺญมิทํ อตฺเตน วา อตฺตนิเยน วา”ติ, โส ตตฺถ นิมิตฺตํ น กโรตีติ อนิมิตฺโต วิโมกฺโข. อยํ อนิมิตฺโต วิโมกฺโข.}

\prose{กตโม \textbf{อปฺปณิหิโต วิโมกฺโข?} อิธ ภิกฺขุ อรญฺญคโต วา รุกฺขมูลคโต วา สุญฺญาคารคโต วา อิติ ปฏิสญฺจิกฺขติ “สุญฺญมิทํ อตฺเตน วา อตฺตนิเยน วา”ติ, โส ตตฺถ ปณิธิํ น กโรตีติ อปฺปณิหิโต วิโมกฺโข. อยํ อปฺปณิหิโต วิโมกฺโข.}

\prose{กตโม อชฺฌตฺต\-วุฏฺฐาโน วิโมกฺโข? จตฺตาริ ฌานานิ. อยํ อชฺฌตฺต\-วุฏฺฐาโน วิโมกฺโข. กตโม พหิทฺธา\-วุฏฺฐาโน วิโมกฺโข? จตสฺโส อรูป\-สมาปตฺติโย. อยํ พหิทฺธา\-วุฏฺฐาโน วิโมกฺโข. กตโม ทุภโต วุฏฺฐาโน วิโมกฺโข? จตฺตาโร อริยมคฺคา. อยํ ทุภโต วุฏฺฐาโน วิโมกฺโข.}

\prose{กตเม อชฺฌตฺต\-วุฏฺฐานา จตฺตาโร วิโมกฺขา? ปฐมํ ฌานํ นีวรเณหิ วุฏฺฐาติ. ทุติยํ ฌานํ วิตกฺกวิจาเรหิ วุฏฺฐาติ. ตติยํ ฌานํ ปีติยา วุฏฺฐาติ. จตุตฺถํ ฌานํ สุขทุกฺเขหิ วุฏฺฐาติ. อิเม อชฺฌตฺต\-วุฏฺฐานา จตฺตาโร วิโมกฺขา. (4)}

\prose{กตเม พหิทฺธา\-วุฏฺฐานา จตฺตาโร วิโมกฺขา? อากาสา\-นญฺจา\-ยตนสมาปตฺติ รูปสญฺญาย ปฏิฆสญฺญาย นานตฺตสญฺญาย วุฏฺฐาติ. วิญฺญาณญฺจา\-ยตนสมาปตฺติ อากาสา\-นญฺจา\-ยตนสญฺญาย วุฏฺฐาติ. อากิญฺจญฺญายตน\-สมาปตฺติ วิญฺญาณญฺจา\-ยตนสญฺญาย วุฏฺฐาติ. เนว\-สญฺญา\-นา\-สญฺญา\-ยตนสมาปตฺติ อากิญฺจญฺญายตน\-สญฺญาย วุฏฺฐาติ. อิเม พหิทฺธา\-วุฏฺฐานา จตฺตาโร วิโมกฺขา. \mbox{(4, 8)}}

\prose{กตเม ทุภโต วุฏฺฐานา จตฺตาโร วิโมกฺขา? โสตาปตฺติมคฺโค สกฺกายทิฏฺฐิ\-วิจิกิจฺฉา\-สีลพฺพตปรามาสา ทิฏฺฐานุสยา วิจิกิจฺฉา\-นุสยา วุฏฺฐาภิ, ตทนุวตฺตก\-กิเลเสหิ จ ขนฺเธหิ จ วุฏฺฐาติ, พหิทฺธา จ สพฺพนิมิตฺเตหิ วุฏฺฐาติ. สกทาคามิ\-มคฺโค โอฬาริกา กามราค\-สญฺโญชนา ปฏิฆ\-สญฺโญชนา โอฬาริกา กามราคานุสยา ปฏิฆานุสยา วุฏฺฐาติ, ตทนุวตฺตก\-กิเลเสหิ จ ขนฺเธหิ จ วุฏฺฐาติ, พหิทฺธา จ สพฺพนิมิตฺเตหิ}

\prose{\csromanfolio{233}วุฏฺฐาติ. อนาคามิมคฺโค อณุสหคตา กามราค\-สญฺโญชนา ปฏิฆ\-สญฺโญชนา อณุสหคตา กามราคานุสยา ปฏิฆานุสยา วุฏฺฐาติ, ตทนุวตฺตก\-กิเลเสหิ จ ขนฺเธหิ จ วุฏฺฐาติ, พหิทฺธา จ สพฺพนิมิตฺเตหิ วุฏฺฐาติ. อรหตฺตมคฺโค รูปราคา อรูปราคา มานา อุทฺธจฺจา อวิชฺชาย มานานุสยา ภวราคา\-นุสยา อวิชฺชานุสยา วุฏฺฐาติ, ตทนุวตฺตก\-กิเลเสหิ จ ขนฺเธหิ จ วุฏฺฐาติ, พหิทฺธา จ สพฺพนิมิตฺเตหิ วุฏฺฐาติ, อิเม ทุภโต วุฏฺฐานา จตฺตาโร วิโมกฺขา. \mbox{(4, 12)}}

\proseitem{211}{กตเม อชฺฌตฺต\-วุฏฺฐานานํ อนุโลมา จตฺตาโร วิโมกฺขา? ปฐมํ ฌานํ ปฏิลาภ\-ตฺถาย วิตกฺโก จ วิจาโร จ ปีติ จ สุขญฺจ จิตฺเตกคฺคตา จ. ทุติยํ ฌานํ ปฏิลาภ\-ตฺถาย ฯเปฯ. ตติยํ ฌานํ ปฏิลาภ\-ตฺถาย ฯเปฯ. จตุตฺถํ ฌานํ ปฏิลาภ\-ตฺถาย วิตกฺโก จ วิจาโร จ ปีติ จ สุขญฺจ จิตฺเตกคฺคตา จ. อิเม อชฺฌตฺต\-วุฏฺฐานานํ อนุโลมา จตฺตาโร วิโมกฺขา. \mbox{(4, 16)}}

\prose{กตเม พหิทฺธา\-วุฏฺฐานานํ อนุโลมา จตฺตาโร วิโมกฺขา? อากาสา\-นญฺจา\-ยตนสมาปตฺติํ ปฏิลาภ\-ตฺถาย วิตกฺโก จ วิจาโร จ ปีติ จ สุขญฺจ จิตฺเตกคฺคตา จ. วิญฺญาณญฺจา\-ยตนสมาปตฺติํ ฯเปฯ. อากิญฺจญฺญายตน\-สมาปตฺติํ ฯเปฯ. เนว\-สญฺญา\-นา\-สญฺญา\-ยตนสมาปตฺติํ ปฏิลาภ\-ตฺถาย วิตกฺโก จ วิจาโร จ ปีติ จ สุขญฺจ จิตฺเตกคฺคตา จ. อิเม พหิทฺธา\-วุฏฺฐานานํ อนุโลมา จตฺตาโร วิโมกฺขา. \mbox{(4, 20)}}

\prose{กตเม ทุภโต วุฏฺฐานานํ อนุโลมา จตฺตาโร วิโมกฺขา? โสตาปตฺติ\-มคฺคํ ปฏิลาภ\-ตฺถาย อนิจฺจา\-นุปสฺสนา, ทุกฺขา\-นุปสฺสนา, อนตฺตา\-นุปสฺสนา. สกทาคามิ\-มคฺคํ ปฏิลาภ\-ตฺถาย ฯเปฯ. อนาคามิมคฺคํ ปฏิลาภ\-ตฺถาย ฯเปฯ. อรหตฺตมคฺคํ ปฏิลาภ\-ตฺถาย อนิจฺจา\-นุปสฺสนา, ทุกฺขา\-นุปสฺสนา, อนตฺตา\-นุปสฺสนา, อิเม ทุภโต วุฏฺฐานานํ อนุโลมา จตฺตาโร วิโมกฺขา. \mbox{(4, 24)}}

\prose{กตเม อชฺฌตฺตวุฏฺฐาน\-ปฏิปฺปสฺสทฺธิ จตฺตาโร วิโมกฺขา? ปฐมสฺส ฌานสฺส ปฏิลาโภ วา วิปาโก วา. ทุติยสฺส ฌานสฺส ฯเปฯ. ตติยสฺส ฌานสฺส ฯเปฯ. จตุตฺถสฺส ฌานสฺส ปฏิลาโภ วา วิปาโก วา. อิเม อชฺฌตฺตวุฏฺฐาน\-ปฏิปฺปสฺสทฺธี จตฺตาโร วิโมกฺขา. \mbox{(4, 28)}}

\prose{\csromanfolio{234}กตเม พหิทฺธาวุฏฺฐาน\-ปฏิปฺปสฺสทฺธี จตฺตาโร วิโมกฺขา? อากาสา\-นญฺจา\-ยตนสมาปตฺติยา ปฏิลาโภ วา วิปาโก วา. วิญฺญาณญฺจายตน\-สมาปตฺติยา ฯเปฯ. อากิญฺจญฺญายตน\-สมาปตฺติยา ฯเปฯ. เนว\-สญฺญา\-นา\-สญฺญา\-ยตนสมาปตฺติยา ปฏิลาโภ วา วิปาโก วา. อิเม พหิทฺธาวุฏฺฐาน\-ปฏิปฺปสฺสทฺธี จตฺตาโร วิโมกฺขา. \mbox{(4, 32)}}

\prose{กตเม ทุภโต วุฏฺฐาน\-ปฏิปฺปสฺสทฺธิ จตฺตาโร วิโมกฺขา? โสตาปตฺติ\-มคฺคสฺส โสตาปตฺติผลํ, สกทาคามิ\-มคฺคสฺส สกทาคามิผลํ, อนาคามิ\-มคฺคสฺส อนาคามิผลํ, อรหตฺต\-มคฺคสฺส อรหตฺตผลํ. อิเม ทุภโต วุฏฺฐาน\-ปฏิปฺปสฺสทฺธี จตฺตาโร วิโมกฺขา. \mbox{(4, 36)}}

\proseitem{212}{กถํ \textbf{รูปี รูปานิ ปสฺสตีติ วิโมกฺโข?} อิเธกจฺโจ อชฺฌตฺตํ ปจฺจตฺตํ นีลนิมิตฺตํ มนสิ กโรติ, นีลสญฺญํ ปฏิลภติ, โส ตํ นิมิตฺตํ สุคฺคหิตํ กโรติ, สูปธาริตํ\csromansharedfootnote{0211ed2151f21587}{สุปธาริตํ (ก)} อุปธาเรติ, สฺวาวตฺถิตํ อวตฺถาเปติ. โส ตํ นิมิตฺตํ สุคฺคหิตํ กตฺวา สูปธาริตํ อุปธาเรตฺวา สฺวาวตฺถิตํ อวตฺถาเปตฺวา พหิทฺธา นีลนิมิตฺเต จิตฺตํ อุปสํหรติ, นีลสญฺญํ ปฏิลภติ. โส ตํ นิมิตฺตํ สุคฺคหิตํ กโรติ, สูปธาริตํ อุปธาเรติ, สฺวาวตฺถิตํ อวตฺถาเปติ. โส ตํ นิมิตฺตํ สุคฺคหิตํ กตฺวา อุปธาเรตฺวา สฺวาวตฺถิตํ อวตฺถาเปตฺวา อาเสวติ ภาเวติ พหุลีกโรติ. ตสฺส เอวํ โหติ “อชฺฌตฺตญฺจ พหิทฺธา จ อุภยมิทํ รูปนฺ”ติ รูปสญฺญี โหติ. อิเธกจฺโจ อชฺฌตฺตํ ปจฺจตฺตํ ปีตนิมิตฺตํ ฯเปฯ โลหิตนิมิตฺตํ ฯเปฯ โอทาตนิมิตฺตํ มนสิ กโรติ, โอทาตสญฺญํ ปฏิลภติ. โส ตํ นิมิตฺตํ สุคฺคหิตํ กโรติ, สูปธาริตํ อุปธาเรติ, สฺวาวตฺถิตํ อวตฺถาเปติ. โส ตํ นิมิตฺตํ สุคฺคหิตํ กตฺวา สูปธาริตํ อุปธาเรตฺวา สฺวาวตฺถิตํ อวตฺถาเปตฺวา พหิทฺธา โอทาตนิมิตฺเต จิตฺตํ อุปสํหรติ, โอทาตสญฺญํ ปฏิลภติ. โส ตํ นิมิตฺตํ สุคฺคหิตํ กโรติ, สูปธาริตํ อุปธาเรติ, สฺวาวตฺถิตํ อวตฺถาเปติ. โส ตํ นิมิตฺตํ สุคฺคหิตํ กตฺวา สูปธาริตํ อุปธาเรตฺวา สฺวาวตฺถิตํ อวตฺถาเปตฺวา อาเสวติ ภาเวติ พหุลีกโรติ. ตสฺส เอวํ โหติ “อชฺฌตฺตญฺจ พหิทฺธา จ อุภยมิทํ รูปนฺ”ติ รูปสญฺญี โหติ. เอวํ รูปี รูปานิ ปสฺสตีติ วิโมกฺโข. (37)}

\prose{\csromanfolio{235}กถํ \textbf{อชฺฌตฺตํ อรูปสญฺญี พหิทฺธา รูปานิ ปสฺสตีติ วิโมกฺโข?} อิเธกจฺโจ อชฺฌตฺตํ ปจฺจตฺตํ นีลนิมิตฺตํ น มนสิ กโรติ, นีลสญฺญํ น ปฏิลภติ, พหิทฺธา นีลนิมิตฺเต จิตฺตํ อุปสํหรติ, นีลสญฺญํ ปฏิลภติ. โส ตํ นิมิตฺตํ สุคฺคหิตํ กโรติ, สูปธาริตํ อุปธาเรติ, สฺวาวตฺถิตํ อวตฺถาเปติ. โส ตํ นิมิตฺตํ สุคฺคหิตํ กตฺวา สูปธาริตํ อุปธาเรตฺวา สฺวาวตฺถิตํ อวตฺถาเปตฺวา อาเสวติ ภาเวติ พหุลีกโรติ. ตสฺส เอวํ โหติ “อชฺฌตฺตํ อรูปํ พหิทฺธา รูปมิทนฺ”ติ รูปสญฺญี โหติ. อิเธกจฺโจ อชฺฌตฺตํ ปจฺจตฺตํ ปีตนิมิตฺตํ ฯเปฯ โลหิตนิมิตฺตํ ฯเปฯ โอทาตนิมิตฺตํ น มนสิ กโรติ, โอทาตสญฺญํ น ปฏิลภติ. พหิทฺธา โอทาตนิมิตฺเต จิตฺตํ อุปสํหรติ, โอทาตสญฺญํ ปฏิลภติ. โส ตํ นิมิตฺตํ สุคฺคหิตํ กโรติ, สูปธาริตํ อุปธาเรติ, สฺวาวตฺถิตํ อวตฺถาเปติ. โส ตํ นิมิตฺตํ สุคฺคหิตํ กตฺวา สูปธาริตํ อุปธาเรตฺวา สฺวาวตฺถิตํ อวตฺถาเปตฺวา อาเสวติ ภาเวติ พหุลีกโรติ. ตสฺส เอวํ โหติ “อชฺฌตฺตํ อรูปํ พหิทฺธา รูปมิทนฺ”ติ รูปสญฺญี โหติ. เอวํ อชฺฌตฺตํ อรูปสญฺญี พหิทฺธา รูปานิ ปสฺสตีติ วิโมกฺโข. (38)}

\prose{กถํ “สุภนฺ”เตว อธิมุตฺโต โหตีติ วิโมกฺโข? อิธ ภิกฺขุ เมตฺตา\-สหคเตน เจตสา เอกํ ทิสํ ผริตฺวา วิหรติ, ตถา ทุติยํ. ตถา ตติยํ. ตถา จตุตฺถํ. อิติ อุทฺธมโธ ติริยํ สพฺพธิ สพฺพตฺตตาย สพฺพาวนฺตํ โลกํ เมตฺตา\-สหคเตน เจตสา วิปุเลน มหคฺคเตน อปฺปมาเณน อเวเรน อพฺยาปชฺเชน\csromansharedfootnote{078f1c26e84c46a3}{อพฺยาปชฺเฌน (สฺยา)} ผริตฺวา วิหรติ, เมตฺตาย ภาวิตตฺตา สตฺตา อปฺปฏิกูลา โหนฺติ. กรุณา\-สหคเตน เจเตสา ฯเปฯ กรุณาย ภาวิตตฺตา สตฺตา อปฺปฏิกูลา โหนฺติ. มุทิตา\-สหคเตน เจตสา ฯเปฯ มุทิตาย ภาวิตตฺตา สตฺตา อปฺปฏิกูลา โหนฺติ. อุเปกฺขา\-สหคเตน เจตสา เอกํ ทิสํ ผริตฺวา วิหรติ ฯเปฯ อุเปกฺขาย ภาวิตตฺตา สตฺตา อปฺปฏิกูลา โหนฺติ. เอวํ “สุภนฺ”เตว อธิมุตฺโต โหตีติ วิโมกฺโข. (39)}

\proseitem{213}{กตโม อากาสา\-นญฺจา\-ยตนสมาปตฺติ วิโมกฺโข? อิธ ภิกฺขุ สพฺพโส รูปสญฺญานํ สมติกฺกมา ปฏิฆ\-สญฺญานํ อตฺถงฺคมา \csromanfolio{236}นานตฺต\-สญฺญานํ อมนสิการา “อนนฺโต อากาโส”ติ อากาสา\-นญฺจา\-ยตนํ อุปสมฺปชฺช วิหรติ. อยํ อากาสา\-นญฺจา\-ยตนสมาปตฺติ วิโมกฺโข. (40)}

\prose{กตโม วิญฺญาณญฺจา\-ยตนสมาปตฺติ วิโมกฺโข? อิธ ภิกฺขุ สพฺพโส อากาสา\-นญฺจา\-ยตนํ สมติกฺกมฺม “อนนฺตํ วิญฺญาณนฺ”ติ วิญฺญาณญฺจา\-ยตนํ อุปสมฺปชฺช วิหรติ. อยํ วิญฺญาณญฺจา\-ยตนสมาปตฺติ วิโมกฺโข. (41)}

\prose{กตโม อากิญฺจญฺญายตน\-สมาปตฺติ วิโมกฺโข? อิธ ภิกฺขุ สพฺพโส วิญฺญาณญฺจา\-ยตนํ สมติกฺกมฺม “นตฺถิ กิญฺจี”ติ อากิญฺจญฺญา\-ยตนํ อุปสมฺปชฺช วิหรติ. อยํ อากิญฺจญฺญายตน\-สมาปตฺติ วิโมกฺโข. (42)}

\prose{กตโม เนว\-สญฺญา\-นา\-สญฺญา\-ยตนสมาปตฺติ วิโมกฺโข? อิธ ภิกฺขุ สพฺพโส อากิญฺจญฺญา\-ยตนํ สมติกฺกมฺม เนว\-สญฺญา\-นา\-สญฺญา\-ยตนํ อุปสมฺปชฺช วิหรติ. อยํ เนว\-สญฺญา\-นา\-สญฺญา\-ยตนสมาปตฺติ วิโมกฺโข. (43)}

\prose{กตโม สญฺญา\-เวทยิต\-นิโรธ\-สมาปตฺติ วิโมกฺโข? อิธ ภิกฺขุ สพฺพโส เนว\-สญฺญา\-นา\-สญฺญา\-ยตนํ สมติกฺกมฺม สญฺญาเวทยิต\-นิโรธํ อุปสมฺปชฺช วิหรติ. อยํ สญฺญา\-เวทยิต\-นิโรธ\-สมาปตฺติ วิโมกฺโข. (44)}

\prose{กตโม สมยวิโมกฺโข? จตฺตาริ จ ฌานานิ, จตสฺโส จ อรูป\-สมาปตฺติโย. อยํ สมยวิโมกฺโข. (45)}

\prose{กตโม อสมย\-วิโมกฺโข? จตฺตาโร จ อริยมคฺคา, จตฺตาริ จ สามญฺญผลานิ, นิพฺพานญฺจ. อยํ อสมย\-วิโมกฺโข. (46)}

\prose{กตโม สามยิโก วิโมกฺโข? จตฺตาริ จ ฌานานิ, จตสฺโส จ อรูป\-สมาปตฺติโย. อยํ สามยิโก วิโมกฺโข. (47)}

\prose{กตโม \textbf{อสามยิโก วิโมกฺโข?} จตฺตาโร จ อริยมคฺคา, จตฺตาริ จ สามญฺญผลานิ, นิพฺพานญฺจ. อยํ อสามยิโก วิโมกฺโข. (48)}

\prose{กตโม กุปฺโป วิโมกฺโข? จตฺตาริ จ ฌานานิ, จตสฺโส จ อรูป\-สมาปตฺติโย. อยํ กุปฺโป วิโมกฺโข. (49)}

\prose{\csromanfolio{237}กตโม \textbf{อกุปฺโป วิโมกฺโข?} จตฺตาโร จ อริยมคฺคา, จตฺตาริ จ สามญฺญผลานิ, นิพฺพานญฺจ. อยํ อกุปฺโป วิโมกฺโข. (50)}

\prose{กตโม โลกิโก วิโมกฺโข? จตฺตาริ จ ฌานานิ, จตสฺโส จ อรูป\-สมาปตฺติโย. อยํ โลกิโก วิโมกฺโข. (51)}

\prose{กตโม \textbf{โลกุตฺตโร วิโมกฺโข?} จตฺตาโร จ อริยมคฺคา, จตฺตาริ จ สามญฺญผลานิ, นิพฺพานญฺจ. อยํ โลกุตฺตโร วิโมกฺโข. (52)}

\prose{กตโม สาสโว วิโมกฺโข? จตฺตาริ จ ฌานานิ, จตสฺโส จ อรูป\-สมาปตฺติโย. อยํ สาสโว วิโมกฺโข. (53)}

\prose{กตโม \textbf{อนาสโว วิโมกฺโข?} จตฺตาโร จ อริยมคฺคา, จตฺตาริ จ สามญฺญผลานิ, นิพฺพานญฺจ. อยํ อนาสโว วิโมกฺโข. (54)}

\prose{กตโม \textbf{สามิโส วิโมกฺโข}? รูปปฺปฏิสญฺญุตฺโต วิโมกฺโข. อยํ สามิโส วิโมกฺโข. (55)}

\prose{กตโม นิรามิโส วิโมกฺโข? อรูป\-ปฺปฏิสญฺญุตฺโต วิโมกฺโข. อยํ นิรามิโส วิโมกฺโข. (56)}

\prose{กตโม \textbf{นิรามิสา นิรามิสตโร วิโมกฺโข?} จตฺตาโร จ อริยมคฺคา, จตฺตาริ จ สามญฺญผลานิ, นิพฺพานญฺจ. อยํ นิรามิสา นิรามิสตโร วิโมกฺโข. (57)}

\prose{กตโม ปณิหิโต วิโมกฺโข? จตฺตาริ จ ฌานานิ, จตสฺโส จ อรูป\-สมาปตฺติโย. อยํ ปณิหิโต วิโมกฺโข. กตโม อปฺปณิหิโต วิโมกฺโข? จตฺตาโร จ อริยมคฺคา, จตฺตาริ จ สามญฺญผลานิ, นิพฺพานญฺจ. อยํ อปฺปณิหิโต วิโมกฺโข. (58)}

\prose{กตโม ปณิหิต\-ปฺปฏิปฺปสฺสทฺธิ วิโมกฺโข? ปฐมสฺส ฌานสฺส ปฏิลาโภ วา วิปาโก วา ฯเปฯ. เนว\-สญฺญา\-นา\-สญฺญา\-ยตนสมาปตฺติยา ปฏิลาโภ วา วิปาโก วา. อยํ ปณิหิต\-ปฺปฏิปฺปสฺสทฺธิ วิโมกฺโข. (59)}

\prose{กตโม สญฺญุตฺโต วิโมกฺโข? จตฺตาริ จ ฌานานิ, จตสฺโส จ อรูป\-สมาปตฺติโย. อยํ สญฺญุตฺโต วิโมกฺโข. (60)}

\prose{กตโม \textbf{วิสญฺญุตฺโต วิโมกฺโข?} จตฺตาโร จ อริยมคฺคา, จตฺตาริ จ สามญฺญผลานิ, นิพฺพานญฺจ. อยํ วิสญฺญุตฺโต วิโมกฺโข. (61)}

\prose{\csromanfolio{238}กตโม \textbf{เอกตฺตวิโมกฺโข?} จตฺตาโร จ อริยมคฺคา, จตฺตาริ จ สามญฺญผลานิ, นิพฺพานญฺจ. อยํ เอกตฺตวิโมกฺโข. (62)}

\prose{กตโม นานตฺต\-วิโมกฺโข? จตฺตาริ จ ฌานานิ, จตสฺโส จ อรูป\-สมาปตฺติโย. อยํ นานตฺต\-วิโมกฺโข. (63)}

\proseitem{214}{กตโม \textbf{สญฺญาวิโมกฺโข?} สิยา เอโก สญฺญาวิโมกฺโข ทส สญฺญาวิโมกฺขา โหนฺติ. ทส สญฺญาวิโมกฺขา เอโก สญฺญาวิโมกฺโข โหติ วตฺถุวเสน ปริยาเยน.}

\prose{สิยาติ กถญฺจ สิยา? อนิจฺจานุปสฺสน\-ญาณํ นิจฺจโต สญฺญาย มุจฺจตีติ สญฺญาวิโมกฺโข, ทุกฺขานุปสฺสน\-ญาณํ สุขโต สญฺญาย มุจฺจตีติ สญฺญาวิโมกฺโข, อนตฺตานุปสฺสน\-ญาณํ อตฺตโต สญฺญาย มุจฺจตีติ สญฺญาวิโมกฺโข, นิพฺพิทานุปสฺสน\-ญาณํ นนฺทิยา สญฺญาย มุจฺจตีติ สญฺญาวิโมกฺโข, วิราคานุปสฺสน\-ญาณํ ราคโต สญฺญาย มุจฺจตีติ สญฺญาวิโมกฺโข, นิโรธานุปสฺสน\-ญาณํ สมุทยโต สญฺญาย มุจฺจตีติ สญฺญาวิโมกฺโข, ปฏินิสฺสคฺคานุปสฺสน\-ญาณํ อาทานโต สญฺญาย มุจฺจตีติ สญฺญาวิโมกฺโข, อนิมิตฺตา\-นุปสฺสน\-ญาณํ นิมิตฺตโต สญฺญาย มุจฺจตีติ สญฺญาวิโมกฺโข, อปฺปณิหิตา\-นุปสฺสน\-ญาณํ ปณิธิยา สญฺญาย มุจฺจตีติ สญฺญาวิโมกฺโข, สุญฺญตานุปสฺสน\-ญาณํ อภินิเวสโต สญฺญาย มุจฺจตีติ สญฺญาวิโมกฺโข. เอวํ สิยา เอโก สญฺญาวิโมกฺโข ทส สญฺญาวิโมกฺขา โหนฺติ, ทส สญฺญาวิโมกฺขา เอโก สญฺญาวิโมกฺโข โหติ วตฺถุวเสน ปริยาเยน.}

\prose{รูเป อนิจฺจานุปสฺสน\-ญาณํ นิจฺจโต สญฺญาย มุจฺจตีติ สญฺญาวิโมกฺโข ฯเปฯ รูเป สุญฺญตานุปสฺสน\-ญาณํ อภินิเวสโต สญฺญาย มุจฺจตีติ สญฺญาวิโมกฺโข. เอวํ สิยา เอโก สญฺญาวิโมกฺโข ทส สญฺญาวิโมกฺขา โหนฺติ. ทส สญฺญาวิโมกฺขา เอโก สญฺญาวิโมกฺโข โหติ วตฺถุวเสน ปริยาเยน.}

\prose{เวทนาย ฯเปฯ. สญฺญาย. สงฺขาเรสุ. วิญฺญาเณ. จกฺขุสฺมิํ ฯเปฯ. ชรามรเณ อนิจฺจานุปสฺสน\-ญาณํ นิจฺจโต สญฺญาย มุจฺจตีติ สญฺญาวิโมกฺโข ฯเปฯ ชรามรเณ สุญฺญตานุปสฺสน\-ญาณํ อภินิเวสโต สญฺญาย มุจฺจตีติ}

\prose{\csromanfolio{239}สญฺญาวิโมกฺโข. เอวํ สิยา เอโก สญฺญาวิโมกฺโข ทส สญฺญาวิโมกฺขา โหนฺติ. ทส สญฺญาวิโมกฺขา เอโก สญฺญาวิโมกฺโข โหติ วตฺถุวเสน ปริยาเยน. อยํ สญฺญาวิโมกฺโข. (64)}

\proseitem{215}{กตโม \textbf{ญาณวิโมกฺโข?} สิยา เอโก ญาณวิโมกฺโข ทส ญาณวิโมกฺขา โหนฺติ. ทส ญาณวิโมกฺขา เอโก ญาณวิโมกฺโข โหติ วตฺถุวเสน ปริยาเยน.}

\prose{สิยาติ กถญฺจ สิยา? อนิจฺจา\-นุปสฺสนา ยถาภูตํ ญาณํ นิจฺจโต สมฺโมหา อญฺญาณา มุจฺจตีติ ญาณวิโมกฺโข, ทุกฺขา\-นุปสฺสนา ยถาภูตํ ญาณํ สุขโต สมฺโมหา อญฺญาณา มุจฺจตีติ ญาณวิโมกฺโข, อนตฺตา\-นุปสฺสนา ยถาภูตํ ญาณํ อตฺตโต สมฺโมหา อญฺญาณา มุจฺจตีติ ญาณวิโมกฺโข, นิพฺพิทา\-นุปสฺสนา ยถาภูตํ ญาณํ นนฺทิยา สมฺโมหา อญฺญาณา มุจฺจตีติ ญาณวิโมกฺโข, วิราคา\-นุปสฺสนา ยถาภูตํ ญาณํ ราคโต สมฺโมหา อญฺญาณา มุจฺจตีติ ญาณวิโมกฺโข, นิโรธา\-นุปสฺสนา ยถาภูตํ ญาณํ สมุทยโต สมฺโมหา อญฺญาณา มุจฺจตีติ ญาณวิโมกฺโข, ปฏินิสฺสคฺคา\-นุปสฺสนา ยถาภูตํ ญาณํ อาทานโต สมฺโมหา อญฺญาณา มุจฺจตีติ ญาณวิโมกฺโข, อนิมิตฺตา\-นุปสฺสนา ยถาภูตํ ญาณํ นิมิตฺตโต สมฺโมหา อญฺญาณา มุจฺจตีติ ญาณวิโมกฺโข, อปฺปณิหิตา\-นุปสฺสนา ยถาภูตํ ญาณํ ปณิธิยา สมฺโมหา อญฺญาณา มุจฺจตีติ ญาณวิโมกฺโข, สุญฺญตา\-นุปสฺสนา ยถาภูตํ ญาณํ อภินิเวสโต สมฺโมหา อญฺญาณา มุจฺจตีติ ญาณวิโมกฺโข. เอวํ สิยา เอโก ญาณวิโมกฺโข ทส ญาณวิโมกฺขา โหนฺติ, ทส ญาณวิโมกฺขา เอโก ญาณวิโมกฺโข โหติ วตฺถุวเสน ปริยาเยน.}

\prose{รูเป อนิจฺจา\-นุปสฺสนา ยถาภูตํ ญาณํ นิจฺจโต สมฺโมหา อญฺญาณา มุจฺจตีติ ญาณวิโมกฺโข ฯเปฯ รูเป สุญฺญตา\-นุปสฺสนา ยถาภูตํ ญาณํ อภินิเวสโต สมฺโมหา อญฺญาณา มุจฺจตีติ ญาณวิโมกฺโข. เอวํ สิยา เอโก ญาณวิโมกฺโข ทส ญาณวิโมกฺขา โหนฺติ, ทส ญาณวิโมกฺขา เอโก ญาณวิโมกฺโข โหติ วตฺถุวเสน ปริยาเยน.}

\prose{\csromanfolio{240}เวทนาย ฯเปฯ. สญฺญาย. สงฺขาเรสุ. วิญฺญาเณ. จกฺขุสฺมิํ ฯเปฯ. ชรามรเณ อนิจฺจา\-นุปสฺสนา ยถาภูตํ ญาณํ นิจฺจโต สมฺโมหา อญฺญาณา มุจฺจตีติ ญาณวิโมกฺโข ฯเปฯ ชรามรเณ สุญฺญตา\-นุปสฺสนา ยถาภูตํ ญาณํ อภินิเวสโต สมฺโมหา อญฺญาณา มุจฺจตีติ ญาณวิโมกฺโข. เอวํ สิยา เอโก ญาณวิโมกฺโข ทส ญาณวิโมกฺขา โหนฺติ, ทส ญาณวิโมกฺขา เอโก ญาณวิโมกฺโข โหติ วตฺถุวเสน ปริยาเยน. อยํ ญาณวิโมกฺโข. (65)}

\proseitem{216}{กตโม สีติสิยา\-วิโมกฺโข? สิยา เอโก สีติสิยา\-วิโมกฺโข ทส สีติสิยา\-วิโมกฺขา โหนฺติ, ทส สีติสิยา\-วิโมกฺขา เอโก สีติสิยา\-วิโมกฺโข โหติ วตฺถุวเสน ปริยาเยน.}

\prose{สิยาติ กถญฺจ สิยา? อนิจฺจา\-นุปสฺสนา อนุตฺตรํ สีติภาวญาณํ นิจฺจโต สนฺตาป\-ปริฬาห\-ทรถา มุจฺจตีติ สีติสิยา\-วิโมกฺโข. ทุกฺขา\-นุปสฺสนา อนุตฺตรํ สีติภาวญาณํ สุขโต สนฺตาป\-ปริฬาห\-ทรถา มุจฺจตีติ สีติสิยา\-วิโมกฺโข. อนตฺตา\-นุปสฺสนา อนุตฺตรํ สีติภาวญาณํ อตฺตโต สนฺตาป\-ปริฬาห\-ทรถา มุจฺจตีติ สีติสิยา\-วิโมกฺโข. นิพฺพิทา\-นุปสฺสนา อนุตฺตรํ สีติภาวญาณํ นนฺทิยา สนฺตาป\-ปริฬาห\-ทรถา มุจฺจตีติ สีติสิยา\-วิโมกฺโข. วิราคา\-นุปสฺสนา อนุตฺตรํ สีติภาวญาณํ ราคโต สนฺตาป\-ปริฬาห\-ทรถา มุจฺจตีติ สีติสิยา\-วิโมกฺโข. นิโรธา\-นุปสฺสนา อนุตฺตรํ สีติภาวญาณํ สมุทยโต สนฺตาป\-ปริฬาห\-ทรถา มุจฺจตีติ สีติสิยา\-วิโมกฺโข. ปฏินิสฺสคฺคา\-นุปสฺสนา อนุตฺตรํ สีติภาวญาณํ อาทานโต สนฺตาป\-ปริฬาห\-ทรถา มุจฺจตีติ สีกิสิยาวิโมกฺโข. อนิมิตฺตา\-นุปสฺสนา อนุตฺตรํ สีติภาวญาณํ นิมิตฺตโต สนฺตาป\-ปริฬาห\-ทรถา มุจฺจตีติ สีติสิยา\-วิโมกฺโข. อปฺปณิหิตา\-นุปสฺสนา อนุตฺตรํ สีติภาวญาณํ ปณิธิยา สนฺตาป\-ปริฬาห\-ทรถา มุจฺจตีติ สีติสิยา\-วิโมกฺโข. สุญฺญตา\-นุปสฺสนา อนุตฺตรํ สีติภาวญาณํ อภินิเวสโต สนฺตาป\-ปริฬาห\-ทรถา มุจฺจตีติ สีติสิยา\-วิโมกฺโข. เอวํ สิยา เอโก สีติสิยา\-วิโมกฺโข ทส สีติสิยา\-วิโมกฺขา โหนฺติ, ทส สีติสิยา\-วิโมกฺขา เอโก สีติสิยา\-วิโมกฺโข โหติ วตฺถุวเสน ปริยาเยน.}

\prose{\csromanfolio{241}รูเป อนิจฺจา\-นุปสฺสนา อนุตฺตรํ สีติภาวญาณํ นิจฺจโต สนฺตาป\-ปริฬาห\-ทรถา มุจฺจตีติ สีติสิยา\-วิโมกฺโข ฯเปฯ รูเป สุญฺญาตานุปสฺสนา อนุตฺตรํ สีติภาวญาณํ อภินิเวสโต สนฺตาป\-ปริฬาห\-ทรถา มุจฺจตีติ สีติสิยา\-วิโมกฺโข. เอวํ สิยา เอโก สีติสิยา\-วิโมกฺโข ทส สีติสิยา\-วิโมกฺขา โหนฺติ, ทส สีติสิยา\-วิโมกฺขา เอโก สีติสิยา\-วิโมกฺโข โหติ วตฺถุวเสน ปริยาเยน.}

\prose{เวทนาย ฯเปฯ. สญฺญาย. สงฺขาเรสุ. วิญฺญาเณ. จกฺขุสฺมิํ ฯเปฯ. ชรามรเณ อนิจฺจา\-นุปสฺสนา อนุตฺตรํ สีติภาวญาณํ นิจฺจโต สนฺตาป\-ปริฬาห\-ทรถา มุจฺจตีติ สีติสิยา\-วิโมกฺโข ฯเปฯ ชรามรเณ สุญฺญตา\-นุปสฺสนา อนุตฺตรํ สีติภาวญาณํ อภินิเวสโต สนฺตาป\-ปริฬาห\-ทรถา มุจฺจตีติ สีติสิยา\-วิโมกฺโข. เอวํ สิยา เอโก สีติสิยา\-วิโมกฺโข ทส สีติสิยา\-วิโมกฺขา โหนฺติ, ทส สีติสิยา\-วิโมกฺขา เอโก สีติสิยา\-วิโมกฺโข โหติ วตฺถุวเสน ปริยาเยน, อยํ สีติสียาวิโมกฺโข. (66)}

\proseitem{217}{กตโม ฌานวิโมกฺโข? เนกฺขมฺมํฌายตีติ\csromansharedfootnote{62795bfbccc92995}{ชายตีติ (สฺยา)} ฌานํ, กามจฺฉนฺทํ ฌาเปตีติ ฌานํ, ฌายนฺโต\csromansharedfootnote{d3f4d81bb916b1ea}{ชายนฺโต (สฺยา)} มุจฺจตีติ ฌานวิโมกฺโข, ฌาเปนฺโต มุจฺจตีติ ฌานวิโมกฺโข. ฌายนฺตีติ ธมฺมา, ฌาเปตีติ กิเลเส, ฌาเต จ ฌาเป จ ชานาตีติ ฌานวิโมกฺโข. อพฺยาปาโท ฌายตีติ ฌานํ พฺยาปาทํ ฌาเปตีติ ฌานํ. ฌายนฺโต มุจฺจตีติ ฌานวิโมกฺโข, ฌาเปนฺโต มุจฺจตีติ ฌานวิโมกฺโข. ฌายนฺตีติ ธมฺมา, ฌาเปตีติ กิเลเส, ฌาเต จ ฌาเป จ ชานาตีติ ฌานวิโมกฺโข. อาโลกสญฺญา ฌายตีติ ฌานํ. ถินมิทฺธํ ฌาเปตีติ ฌานํ ฯเปฯ. อวิกฺเขโป ฌายตีติ ฌานํ, อุทฺธจฺจํ ฌาเปตีติ ฌานํ. ธมฺม\-ววตฺถานํ ฌายตีติ ฌานํ, วิจิกิจฺฉํ ฌาเปตีติ ฌานํ. ญาณํ ฌายตีติ ฌานํ, อวิชฺชํ ฌาเปตีติ ฌานํ. ปาโมชฺชํ ฌายตีติ ฌานํ, อรติํ ฌาเปตีติ ฌานํ. ปฐมํ ฌานํ ฌายตีติ ฌานํ, นีวรเณ ฌาเปตีติ ฌานํ ฯเปฯ. อรหตฺตมคฺโค ฌายตีติ ฌานํ สพฺพกิเลเส ฌาเปตีติ ฌานํ. ฌายนฺโต มุจฺจตีติ ฌานวิโมกฺโข, ฌาเปนฺโต มุจฺจตีติ ฌานวิโมกฺโข. ฌายนฺตีติ ธมฺมา, ฌาเปตีติ กิเลเส, ฌาเต จ ฌาเป จ ชานาตีติ ฌานวิโมกฺโข. อยํ ฌานวิโมกฺโข. (67)}

\proseitem{218}{\csromanfolio{242}กตโม \textbf{อนุปาทา จิตฺตสฺส วิโมกฺโข?} สิยา เอโก อนุปาทา จิตฺตสฺส วิโมกฺโข ทส อนุปาทา จิตฺตสฺส วิโมกฺขา โหนฺติ, ทส อนุปาทา จิตฺตสฺส วิโมกฺขา เอโก อนุปาทา จิตฺตสฺส วิโมกฺโข โหติ วตฺถุวเสน ปริยาเยน.}

\prose{สิยาติ กถญฺจ สิยา? อนิจฺจานุปสฺสน\-ญาณํ นิจฺจโต อุปาทานา มุจฺจตีติ อนุปาทา จิตฺตสฺส วิโมกฺโข. ทุกฺขานุปสฺสน\-ญาณํ สุขโต อุปาทานา มุจฺจตีติ อนุปาทา จิตฺตสฺส วิโมกฺโข. อนตฺตานุปสฺสน\-ญาณํ อตฺตโต อุปาทานา มุจฺจตีติ อนุปาทา จิตฺตสฺส วิโมกฺโข. นิพฺพิทานุปสฺสน\-ญาณํ นนฺทิยา อุปาทานา มุจฺจตีติ อนุปาทา จิตฺตสฺส วิโมกฺโข. วิราคานุปสฺสน\-ญาณํ ราคโต อุปาทานา มุจฺจตีติ อนุปาทา จิตฺตสฺส วิโมกฺโข. นิโรธานุปสฺสน\-ญาณํ สมุทยโต อุปาทานา มุจฺจตีติ อนุปาทา จิตฺตสฺส วิโมกฺโข. ปฏินิสฺสคฺคานุปสฺสน\-ญาณํ อาทานโต อุปาทานา มุจฺจตีติ อนุปาทา จิตฺตสฺส วิโมกฺโข. อนิมิตฺตา\-นุปสฺสน\-ญาณํ นิมิตฺตโต อุปาทานา มุจฺจตีติ อนุปาทา จิตฺตสฺส วิโมกฺโข. อปฺปณิหิตา\-นุปสฺสน\-ญาณํ ปณิธิยา อุปาทานา มุจฺจตีติ อนุปาทา จิตฺตสฺส วิโมกฺโข. สุญฺญตานุปสฺสน\-ญาณํ อภินิเวสโต อุปาทานา มุจฺจตีติ อนุปาทา จิตฺตสฺส วิโมกฺโข. เอวํ สิยา เอโก อนุปาทา จิตฺตสฺส วิโมกฺโข ทส อนุปาทา จิตฺตสฺส วิโมกฺขา โหนฺติ, ทส อนุปาทา จิตฺตสฺส วิโมกฺขา เอโก อนุปาทา จิตฺตสฺส วิโมกฺโข โหติ วตฺถุวเสน ปริยาเยน.}

\prose{รูเป อนิจฺจานุปสฺสน\-ญาณํ นิจฺจโต อุปาทานา มุจฺจตีติ อนุปาทา จิตฺตสฺส วิโมกฺโข ฯเปฯ รูเป สุญฺญตานุปสฺสน\-ญาณํ อภินิเวสโต อุปาทานา มุจฺจตีติ อนุปาทา จิตฺตสฺส วิโมกฺโข. เอวํ สิยา เอโก อนุปาทา จิตฺตสฺส วิโมกฺโข ทส อนุปาทา จิตฺตสฺส วิโมกฺขา โหนฺติ, ทส อนุปาทา จิตฺตสฺส วิโมกฺขา เอโก อนุปาทา จิตฺตสฺส วิโมกฺโข โหติ วตฺถุวเสน ปริยาเยน.}

\prose{เวทนาย ฯเปฯ. สญฺญาย. สงฺขาเรสุ. วิญฺญาเณ. จกฺขุสฺมิํ ฯเปฯ. ชรามรเณ อนิจฺจานุปสฺสน\-ญาณํ นิจฺจโต อุปาทานา มุจฺจตีติ อนุปาทา จิตฺตสฺส วิโมกฺโข ฯเปฯ ชรามรเณ สุญฺญตานุปสฺสน\-ญาณํ อภินิเวสโต อุปาทานา มุจฺจตีติ อนุปาทา จิตฺตสฺส วิโมกฺโข. เอวํ สิยา เอโก อนุปาทา จิตฺตสฺส วิโมกฺโข ทส อนุปาทา จิตฺตสฺส วิโมกฺขา}

\prose{\csromanfolio{243}โหนฺติ, ทส อนุปาทา จิตฺตสฺส วิโมกฺขา เอโก อนุปาทา จิตฺตสฺส วิโมกฺโข โหติ วตฺถุวเสน ปริยาเยน.}

\prose{อนิจฺจานุปสฺสน\-ญาณํ กติหุปาทาเนหิ มุจฺจติ, ทุกฺขานุปสฺสน\-ญาณํ กติหุปาทาเนหิ มุจฺจติ, อนตฺตานุปสฺสน\-ญาณํ กติหุปาทาเนหิ มุจฺจติ, นิพฺพิทานุปสฺสน\-ญาณํ ฯเปฯ วิราคานุปสฺสน\-ญาณํ. นิโรธานุปสฺสน\-ญาณํ. ปฏินิสฺสคฺคานุปสฺสน\-ญาณํ. อนิมิตฺตา\-นุปสฺสน\-ญาณํ. อปฺปณิหิตา\-นุปสฺสน\-ญาณํ. สุญฺญตานุปสฺสน\-ญาณํ กติหุปาทาเนหิ มุจฺจตีติ?}

\prose{อนิจฺจานุปสฺสน\-ญาณํ ตีหุปาทาเนหิ มุจฺจติ, ทุกฺขานุปสฺสน\-ญาณํ เอกุปาทานา มุจฺจติ, อนตฺตานุปสฺสน\-ญาณํ ตีหุปาทาเนหิ มุจฺจติ, นิพฺพิทานุปสฺสน\-ญาณํ เอกุปาทานา มุจฺจติ, วิราคานุปสฺสน\-ญาณํ เอกุปาทานา มุจฺจติ, นิโรธานุปสฺสน\-ญาณํ จตูหุปาทาเนหิ มุจฺจติ, ปฏินิสฺสคฺคานุปสฺสน\-ญาณํ จตูหุปาทาเนหิ มุจฺจติ, อนิมิตฺตา\-นุปสฺสน\-ญาณํ ตีหุปาทาเนหิ มุจฺจติ, อปฺปณิหิตา\-นุปสฺสน\-ญาณํ เอกุปาทานา มุจฺจติ, สุญฺญตานุปสฺสน\-ญาณํ ตีหุปาทาเนหิ มุจฺจติ.}

\prose{อนิจฺจานุปสฺสน\-ญาณํ กตเมหิ ตีหุปาทาเนหิ มุจฺจติ, ทิฏฺฐุปาทานา สีลพฺพตุ\-ปาทานา อตฺตวาทุปาทานา, อนิจฺจานุปสฺสน\-ญาณํ อิเมหิ ตีหุปาทาเนหิ มุจฺจติ.  ทุกฺขานุปสฺสน\-ญาณํ กตมา เอกุปาทานา\csromansharedfootnote{98e479c438a89946}{เอกูปาทานา (สีฏฺฐ)} มุจฺจติ, กามุปาทานา. ทุกฺขานุปสฺสน\-ญาณํ อิทํ เอกุปาทานา\csromansharedfootnote{914ceda28cd1c1c8}{อิมา เอกุปาทานา (สฺยา) 159 ปิฏฺเฐ อฏฺฐกถา โอโลเกตพฺพา.} มุจฺจติ.  อนตฺตานุปสฺสน\-ญาณํ กตเมหิ ตีหุปาทาเนหิ มุจฺจติ, ทิฏฺฐุปาทานา สีลพฺพตุ\-ปาทานา อตฺตวาทุปาทานา. อนตฺตานุปสฺสน\-ญาณํ อิเมหิ ตีหุปาทาเนหิ มุจฺจติ.  นิพฺพิทานุปสฺสน\-ญาณํ กตมา เอกุปาทานา มุจฺจติ, กามุปาทานา. นิพฺพิทานุปสฺสน\-ญาณํ อิทํ เอกุปาทานา มุจฺจติ.  วิราคานุปสฺสน\-ญาณํ กตมา เอกุปาทานา มุจฺจติ, กามุปาทานา. วิราคานุปสฺสน\-ญาณํ อิทํ เอกุปาทานา มุจฺจติ.  นิโรธานุปสฺสน\-ญาณํ กตเมหิ จตูหุปาทาเนหิ มุจฺจติ, กามุปาทานา ทิฏฺฐุปาทานา สีลพฺพตุ\-ปาทานา อตฺตวาทุปาทานา. นิโรธานุปสฺสน\-ญาณํ อิเมหิ จตูหุปาทาเนหิ มุจฺจติ.  ปฏินิสฺสคฺคานุปสฺสน\-ญาณํ กตเมหิ จตูหุปาทาเนหิ มุจฺจติ, กามุปาทานา ทิฏฺฐุปาทานา สีลพฺพตุ\-ปาทานา อตฺตวาทุปาทานา. ปฏินิสฺสคฺคานุปสฺสน\-ญาณํ อิเมหิ จตูหุปาทาเนหิ มุจฺจติ.  อนิมิตฺตา\-นุปสฺสน\-ญาณํ กตเมหิ ตีหุปาทาเนหิ มุจฺจติ, ทิฏฺฐุปาทานา สีลพฺพตุ\-ปาทานา}

\prose{\csromanfolio{244}อตฺตวาทุปาทานา. อนิมิตฺตา\-นุปสฺสน\-ญาณํ อิเมหิ ตีหุปาทาเนหิ มุจฺจติ.  อปฺปณิหิตา\-นุปสฺสน\-ญาณํ กตมา เอกุปาทานา มุจฺจติ, กามุปาทานา. อปฺปณิหิตา\-นุปสฺสน\-ญาณํ อิทํ เอกุปาทานา มุจฺจติ.  สุญฺญตานุปสฺสน\-ญาณํ กตเมหิ ตีหุปาทาเนหิ มุจฺจติ, ทิฏฺฐุปาทานา สีลพฺพตุ\-ปาทานา อตฺตวาทุปาทานา. สุญฺญตานุปสฺสน\-ญาณํ อิเมหิ ตีหุปาทาเนหิ มุจฺจติ.}

\prose{ยญฺจ อนิจฺจานุปสฺสน\-ญาณํ ยญฺจ อนตฺตานุปสฺสน\-ญาณํ ยญฺจ อนิมิตฺตา\-นุปสฺสน\-ญาณํ ยญฺจ สุญฺญตานุปสฺสน\-ญาณํ อิมานิ จตฺตาริ ญาณานิ ตีหุปาทาเนหิ มุจฺจนฺติ, ทิฏฺฐุปาทานา สีลพฺพตุ\-ปาทานา อตฺตวาทุปาทานา.  ยญฺจ ทุกฺขานุปสฺสน\-ญาณํ ยญฺจ นิพฺพิทานุปสฺสน\-ญาณํ ยญฺจ วิราคานุปสฺสน\-ญาณํ ยญฺจ อปฺปณิหิตา\-นุปสฺสน\-ญาณํ อิมานิ จตฺตาริ ญาณานิ เอกุปาทานา มุจฺจนฺติ กามุปาทานา.  ยญฺจ นิโรธานุปสฺสน\-ญาณํ ยญฺจ ปฏินิสฺสคฺคานุปสฺสน\-ญาณํ อิมานิ เทฺว ญาณานิ จตูหุปาทาเนหิ มุจฺจนฺติ กามุปาทานา ทิฏฺฐุปาทานา สีลพฺพตุ\-ปาทานา อตฺตวาทุปาทานา. อยํ อนุปาทา จิตฺตสฺส วิโมกฺโข. (68) วิโมกฺข\-กถาย ปฐม\-ภาณวาโร.}

\proseitem{219}{ตีณิ โข ปนิมานิ วิโมกฺขมุขานิ โลกนิยฺยานาย สํวตฺตนฺติ, สพฺพสงฺขาเร ปริจฺเฉท\-ปริวฏุมโต\csromansharedfootnote{538a2a471839409d}{ปริจฺเฉทปริวฏฺฏุมโต (สฺยา)} สมนุปสฺส\-นตาย อนิมิตฺตาย จ ธาตุยา จิตฺตสมฺปกฺขนฺทนตาย, สพฺพ\-สงฺขาเรสุ มโน\-สมุตฺเตชน\-ตาย อปฺปณิหิตาย จ ธาตุยา จิตฺตสมฺปกฺขนฺทนตาย, สพฺพธมฺเม ปรโต สมนุปสฺส\-นตาย สุญฺญตาย จ ธาตุยา จิตฺตสมฺปกฺขนฺทนตาย, อิมานิ ตีณิ วิโมกฺขมุขานิ โลกนิยฺยานาย สํวตฺตนฺติ.}

\prose{อนิจฺจโต มนสิกโรโต กถํ สงฺขารา อุปฏฺฐนฺติ, ทุกฺขโต มนสิกโรโต กถํ สงฺขารา อุปฏฺฐนฺติ, อนตฺตโต มนสิกโรโต กถํ สงฺขารา อุปฏฺฐนฺติ? อนิจฺจโต มนสิกโรโต ขยโต สงฺขารา อุปฏฺฐนฺติ, ทุกฺขโต มนสิกโรโต ภยโต สงฺขารา อุปฏฺฐนฺติ, อนตฺตโต มนสิกโรโต สุญฺญโต สงฺขารา อุปฏฺฐนฺติ.}

\prose{\csromanfolio{245}อนิจฺจโต มนสิกโรโต กิํ พหุลํ จิตฺตํ โหติ, ทุกฺขโต มนสิกโรโต กิํ พหุลํ จิตฺตํ โหติ, อนตฺตโต มนสิกโรโต กิํ พหุลํ จิตฺตํ โหติ? อนิจฺจโต มนสิกโรโต อธิโมกฺข\-พหุลํ จิตฺตํ โหติ, ทุกฺขโต มนสิกโรโต ปสฺสทฺธิ\-พหุลํ จิตฺตํ โหติ, อนตฺตโต มนสิกโรโต เวทพหุลํ จิตฺตํ โหติ.}

\prose{อนิจฺจโต มนสิกโรนฺโต อธิโมกฺข\-พหุโล กตมินฺทฺริยํ ปฏิลภติ, ทุกฺขโต มนสิกโรนฺโต ปสฺสทฺธิ\-พหุโล กตมินฺทฺริยํ ปฏิลภติ, อนตฺตโต มนสิกโรนฺโต เวทพหุโล กตมินฺทฺริยํ ปฏิลภติ? อนิจฺจโต มนสิกโรนฺโต อธิโมกฺข\-พหุโล สทฺธินฺทฺริยํ ปฏิลภติ, ทุกฺขโต มนสิกโรนฺโต ปสฺสทฺธิ\-พหุโล สมาธินฺทฺริยํ ปฏิลภติ, อนตฺตโต มนสิกโรนฺโต เวทพหุโล ปญฺญินฺทฺริยํ ปฏิลภติ.}

\prose{อนิจฺจโต มนสิกโรโต อธิโมกฺข\-พหุลสฺส กตมินฺทฺริยํ อาธิปเตยฺยํ โหติ, ภาวนาย กตินฺทฺริยานิ ตทนฺวยา\csromansharedfootnote{4dba3316017503c1}{ตทนฺวยานิ (สฺยา)} โหนฺติ, สหชาต\-ปจฺจยา โหนฺติ, อญฺญมญฺญ\-ปจฺจยา โหนฺติ, นิสฺสยปจฺจยา โหนฺติ, สมฺปยุตฺต\-ปจฺจยา โหนฺติ, เอกรสา โหนฺติ, เกนฏฺเฐน ภาวนา, โก ภาเวติ. ทุกฺขโต มนสิกโรโต ปสฺสทฺธิ\-พหุลสฺส กตมินฺทฺริยํ อาธิปเตยฺยํ โหติ, ภาวนาย กตินฺทฺริยานิ ตทนฺวยา โหนฺติ, สหชาต\-ปจฺจยา โหนฺติ, อญฺญมญฺญ\-ปจฺจยา โหนฺติ, นิสฺสยปจฺจยา โหนฺติ, สมฺปยุตฺต\-ปจฺจยา โหนฺติ, เอกรสา โหนฺติ, เกนฏฺเฐน ภาวนา, โก ภาเวติ. อนตฺตโต มนสิกโรโต เวทพหุลสฺส กตมินฺทฺริยํ อาธิปเตยฺยํ โหติ, ภาวนาย กตินฺทฺริยานิ ตทนฺวยา โหนฺติ, สหชาต\-ปจฺจยา โหนฺติ, อญฺญมญฺญ\-ปจฺจยา โหนฺติ, นิสฺสยปจฺจยา โหนฺติ, สมฺปยุตฺต\-ปจฺจยา โหนฺติ, เอกรสา โหนฺติ, เกนฏฺเฐน ภาวนา, โก ภาเวติ? อนิจฺจโต มนสิกโรโห อธิโมกฺข\-พหุลสฺส สทฺธินฺทฺริยํ อาธิปเตยฺยํ โหติ, ภาวนาย จตฺตาริ\-นฺทฺริยานิ ตทนฺวยา โหนฺติ, สหชาต\-ปจฺจยา โหนฺติ, อญฺญมญฺญ\-ปจฺจยา โหนฺติ, นิสฺสยปจฺจยา โหนฺติ, สมฺปยุตฺต\-ปจฺจยา โหนฺติ, เอกรสา โหนฺติ, เอกรสฏฺเฐน ภาวนา, โย สมฺมาปฏิปนฺโน, โส ภาเวติ, นตฺถิ มิจฺฉา\-ปฏิปนฺนสฺส อินฺทฺริยภาวนา.  ทุกฺขโต มนสิกโรโต ปสฺสทฺธิ\-พหุลสฺส สมาธินฺทฺริยํ อาธิปเตยฺยํ โหติ, ภาวนาย จตฺตาริ\-นฺทฺริยานิ ตทนฺวยา}

\prose{\csromanfolio{246}โหนฺติ, สหชาต\-ปจฺจยา โหนฺติ, อญฺญมญฺญ\-ปจฺจยา โหนฺติ, นิสฺสยปจฺจยา โหนฺติ, สมฺปยุตฺต\-ปจฺจยา โหนฺติ, เอกรสา โหนฺติ, เอกรสฏฺเฐน ภาวนา, โย สมฺมาปฏิปนฺโน, โส ภาเวติ, นตฺถิ มิจฺฉา\-ปฏิปนฺนสฺส อินฺทฺริยภาวนา.  อนตฺตโต มนสิกโรโต เวทพหุลสฺส ปญฺญินฺทฺริยํ อาธิปเตยฺยํ โหติ, ภาวนาย จตฺตาริ\-นฺทฺริยานิ ตทนฺวยา โหนฺติ, สหชาต\-ปจฺจยา โหนฺติ, อญฺญมญฺญ\-ปจฺจยา โหนฺติ, นิสฺสยปจฺจยา โหนฺติ, สมฺปยุตฺต\-ปจฺจยา โหนฺติ, เอกรสา โหนฺติ, เอกรสฏฺเฐน ภาวนา, โย สมฺมาปฏิปนฺโน, โส ภาเวติ, นตฺถิ มิจฺฉา\-ปฏิปนฺนสฺส อินฺทฺริยภาวนา.}

\proseitem{220}{อนิจฺจโต มนสิกโรโต อธิโมกฺข\-พหุลสฺส กตมินฺทฺริยํ อาธิปเตยฺยํ โหติ, ภาวนาย กตินฺทฺริยานิ ตทนฺวยา โหนฺติ, สหชาต\-ปจฺจยา โหนฺติ, อญฺญมญฺญ\-ปจฺจยา โหนฺติ, นิสฺสยปจฺจยา โหนฺติ, สมฺปยุตฺต\-ปจฺจยา โหนฺติ. ปฏิเวธกาเล กตมินฺทฺริยํ อาธิปเตยฺยํ โหติ, ปฏิเวธาย กตินฺทฺริยานิ ตทนฺวยา โหนฺติ. สหชาต\-ปจฺจยา โหนฺติ, อญฺญมญฺญ\-ปจฺจยา โหนฺติ, นิสฺสยปจฺจยา โหนฺติ, สมฺปยุตฺต\-ปจฺจยา โหนฺติ, เอกรสา โหนฺติ, เกนฏฺเฐน ภาวนา, เกนฏฺเฐน ปฏิเวโธ.  ทุกฺขโต มนสิกโรโต ปสฺสทฺธิ\-พหุลสฺส กตมินฺทฺริยํ อาธิปเตยฺยํ โหติ, ภาวนาย กตินฺทฺริยานิ ตทนฺวยา โหนฺติ, สหชาต\-ปจฺจยา โหนฺติ, อญฺญมญฺญ\-ปจฺจยา โหนฺติ, นิสฺสยปจฺจยา โหนฺติ, สมฺปยุตฺต\-ปจฺจยา โหนฺติ. ปฏิเวธกาเล กตมินฺทฺริยํ อาธิปเตยฺยํ โหติ, ปฏิเวธาย กตินฺทฺริยานิ ตทนฺวยา โหนฺติ, สหชาต\-ปจฺจยา โหนฺติ, อญฺญมญฺญ\-ปจฺจยา โหนฺติ, นิสฺสยปจฺจยา โหนฺติ, สมฺปยุตฺต\-ปจฺจยา โหนฺติ, เอกรสา โหนฺติ, เกนฏฺเฐน ภาวนา, เกนฏฺเฐน ปฏิเวโธ.  อนตฺตโต มนสิกโรโต เวทพหุลสฺส กตมินฺทฺริยํ อาธิปเตยฺยํ โหติ, ภาวนาย กตินฺทฺริยานิ ตทนฺวยา โหนฺติ, สหชาต\-ปจฺจยา โหนฺติ, อญฺญมญฺญ\-ปจฺจยา โหนฺติ, นิสฺสยปจฺจยา โหนฺติ, สมฺปยุตฺต\-ปจฺจยา โหนฺติ, ปฏิเวธกาเล กตมินฺทฺริยํ อาธิปเตยฺยํ โหติ, ปฏิเวธาย กตินฺทฺริยานิ ตทนฺวยา โหนฺติ. สหชาต\-ปจฺจยา โหนฺติ, อญฺญมญฺญ\-ปจฺจยา โหนฺติ, นิสฺสยปจฺจยา โหนฺติ, สมฺปยุตฺต\-ปจฺจยา โหนฺติ, เอกรสา โหนฺติ, เกนฏฺเฐน ภาวนา, เกนฏฺเฐน ปฏิเวโธ?}

\prose{\csromanfolio{247}อนิจฺจโต มนสิกโรโต อธิโมกฺข\-พหุลสฺส สทฺธินฺทฺริยํ อาธิปเตยฺยํ โหติ, ภาวนาย จตฺตาริ\-นฺทฺริยานิ ตทนฺวยา โหนฺติ, สหชาต\-ปจฺจยา โหนฺติ, อญฺญมญฺญ\-ปจฺจยา โหนฺติ, นิสฺสยปจฺจยา โหนฺติ, สมฺปยุตฺต\-ปจฺจยา โหนฺติ, ปฏิเวธกาเล ปญฺญินฺทฺริยํ อาธิปเตยฺยํ โหติ, ปฏิเวธาย จตฺตาริ\-นฺทฺริยานิ ตทนฺวยา โหนฺติ, สหชาต\-ปจฺจยา โหนฺติ, อญฺญมญฺญ\-ปจฺจยา โหนฺติ, นิสฺสยปจฺจยา โหนฺติ, สมฺปยุตฺต\-ปจฺจยา โหนฺติ, เอกรสา โหนฺติ, เอกรสฏฺเฐน ภาวนา, ทสฺสนฏฺเฐน ปฏิเวโธ. เอวํ ปฏิวิชฺฌนฺโตปิ ภาเวติ, ภาเวนฺโตปิ ปฏิวิชฺฌติ.  ทุกฺขโต มนสิกโรโต ปสฺสทฺธิ\-พหุลสฺส สมาธินฺทฺริยํ อาธิปเตยฺยํ โหติ, ภาวนาย จตฺตาริ\-นฺทฺริยานิ ตทนฺวยา โหนฺติ, สหชาต\-ปจฺจยา โหนฺติ, อญฺญมญฺญ\-ปจฺจยา โหนฺติ นิสฺสยปจฺจยา โหนฺติ, สมฺปยุตฺต\-ปจฺจยา โหนฺติ, ปฏิเวธกาเล ปญฺญินฺทฺริยํ อาธิปเตยฺยํ โหติ, ปฏิเวธาย จตฺตาริ\-นฺทฺริยานิ ตทนฺวยา โหนฺติ, สหชาต\-ปจฺจยา โหนฺติ, อญฺญมญฺญ\-ปจฺจยา โหนฺติ, นิสฺสยปจฺจยา โหนฺติ, สมฺปยุตฺต\-ปจฺจยา โหนฺติ, เอกรสา โหนฺติ, เอกรสฏฺเฐน ภาวนา, ทสฺสนฏฺเฐน ปฏิเวโธ. เอวํ ปฏิวิชฺฌนฺโตปิ ภาเวติ, ภาเวนฺโตปิ ปฏิวิชฺฌติ.  อนตฺตโต มนสิกโรโต เวทพหุลสฺส ปญฺญินฺทฺริยํ อาธิปเตยฺยํ โหติ, ภาวนาย จตฺตาริ\-นฺทฺริยานิ ตทนฺวยา โหนฺติ, สหชาต\-ปจฺจยา โหนฺติ, อญฺญมญฺญ\-ปจฺจยา โหนฺติ, นิสฺสยปจฺจยา โหนฺติ, สมฺปยุตฺต\-ปจฺจยา โหนฺติ, ปฏิเวธกาเลปิ ปญฺญินฺทฺริยํ อาธิปเตยฺยํ โหติ, ปฏิเวธาย จตฺตาริ\-นฺทฺริยานิ ตทนฺวยา โหนฺติ, สหชาต\-ปจฺจยา โหนฺติ, อญฺญมญฺญ\-ปจฺจยา โหนฺติ, นิสฺสยปจฺจยา โหนฺติ, สมฺปยุตฺต\-ปจฺจยา โหนฺติ, เอกรสา โหนฺติ, เอกรสฏฺเฐน ภาวนา, ทสฺสนฏฺเฐน ปฏิเวโธ. เอวํ ปฏิวิชฺฌนฺโตปิ ภาเวติ, ภาเวนฺโตปิ ปฏิวิชฺฌติ.}

\proseitem{221}{อนิจฺจโต มนสิกโรโต กตมินฺทฺริยํ อธิมตฺตํ โหติ, กตมิ\-นฺทฺริยสฺส อธิมตฺตตฺตา สทฺธาวิมุตฺโต\csromansharedfootnote{0c7a8232b65ba0bf}{สทฺธาธิมุตฺโต (สฺยา)} โหติ.  ทุกฺขโต มนสิกโรโต กตมินฺทฺริยํ อธิมตฺตํ โหติ, กตมิ\-นฺทฺริยสฺส อธิมตฺตตฺตา กายสกฺขี\csromansharedfootnote{6d0fd622bc8bee18}{กายสกฺขิ (สฺยา, ก) อภิ 3. 132 ปิฏฺเฐ ปสฺสิตพฺพา.} โหติ.  อนตฺตโต มนสิกโรโต กตมินฺทฺริยํ อธิมตฺตํ โหติ. กตมิ\-นฺทฺริยสฺส อธิมตฺตตฺตา ทิฏฺฐิปฺปตฺโต โหติ?}

\prose{\csromanfolio{248}อนิจฺจโต มนสิกโรโต สทฺธินฺทฺริยํ อธิมตฺตํ โหติ, สทฺธิ\-นฺทฺริยสฺส อธิมตฺตตฺตา สทฺธาวิมุตฺโต โหติ.  ทุกฺขโต มนสิกโรโต สมาธินฺทฺริยํ อธิมตฺตํ โหติ, สมาธิ\-นฺทฺริยสฺส อธิมตฺตตฺตา กายสกฺขี โหติ.  อนตฺตโต มนสิกโรโต ปญฺญินฺทฺริยํ อธิมตฺตํ โหติ, ปญฺญินฺทฺริยสฺส อธิมตฺตตฺตา ทิฏฺฐิปฺปตฺโต โหติ.}

\prose{สทฺทหนฺโต วิมุตฺโตติ สทฺธาวิมุตฺโต, ผุฏฺฐตฺตา สจฺฉิกโตติ\csromansharedfootnote{38e75862411f5e5c}{สจฺฉิกโรตีติ (สฺยา, อิ)} กายสกฺขี, ทิฏฺฐตฺตา ปตฺโตติ ทิฏฺฐิปฺปตฺโต. สทฺทหนฺโต วิมุจฺจตีติ สทฺธาวิมุตฺโต. ฌานผสฺสํ ปฐมํ ผุสติ, ปจฺฉา นิโรธํ นิพฺพานํ สจฺฉิกโรตีติ กายสกฺขี. “ทุกฺขา สงฺขารา, สุโข นิโรโธ”ติ ญาตํ โหติ ทิฏฺฐํ วิทิตํ สจฺฉิกตํ ผสฺสิตํ ปญฺญายาติ ทิฏฺฐิปฺปตฺโต. โย จายํ ปุคฺคโล สทฺธาวิมุตฺโต, โย จ กายสกฺขี, โย จ ทิฏฺฐิปฺปตฺโต, สิยา อิเม ตโย ปุคฺคลา สทฺธาวิมุตฺตาปิ กายสกฺขีปิ ทิฏฺฐิปฺปตฺตาปิ วตฺถุวเสน ปริยาเยน.}

\prose{สิยาติ กถญฺจ สิยา? อนิจฺจโต มนสิกโรโต สทฺธินฺทฺริยํ อธิมตฺตํ โหติ, สทฺธิ\-นฺทฺริยสฺส อธิมตฺตตฺตา สทฺธาวิมุตฺโต โหติ. ทุกฺขโต มนสิกโรโต สทฺธินฺทฺริยํ อธิมตฺตํ โหติ, สทฺธิ\-นฺทฺริยสฺส อธิมตฺตตฺตา สทฺธาวิมุตฺโต โหติ. อนตฺตโต มนสิกโรโต สทฺธินฺทฺริยํ อธิมตฺตํ โหติ, สทฺธิ\-นฺทฺริยสฺส อธิมตฺตตฺตา สทฺธาวิมุตฺโต โหติ. เอวํ อิเม ตโย ปุคฺคลา สทฺธิ\-นฺทฺริยสฺส วเสน สทฺธาวิมุตฺตา.}

\prose{ทุกฺขโต มนสิกโรโต สมาธินฺทฺริยํ อธิมตฺตํ โหติ, สมาธิ\-นฺทฺริยสฺส อธิมตฺตตฺตา กายสกฺขี โหติ. อนตฺตโต มนสิกโรโต สมาธินฺทฺริยํ อธิมตฺตํ โหติ, สมาธิ\-นฺทฺริยสฺส อธิมตฺตตฺตา กายสกฺขี โหติ. อนิจฺจโต มนสิกโรโต สมาธินฺทฺริยํ อธิมตฺตํ โหติ, สมาธิ\-นฺทฺริยสฺส อธิมตฺตตฺตา กายสกฺขี โหติ. เอวํ อิเม ตโย ปุคฺคลา สมาธิ\-นฺทฺริยสฺส วเสน กายสกฺขี.}

\prose{อนตฺตโต มนสิกโรโต ปญฺญินฺทฺริยํ อธิมตฺตํ โหติ, ปญฺญินฺทฺริยสฺส อธิมตฺตตฺตา ทิฏฺฐิปฺปตฺโต โหติ. อนิจฺจโต มนสิกโรโต ปญฺญินฺทฺริยํ อธิมตฺตํ โหติ. ปญฺญินฺทฺริยสฺส อธิมตฺตตฺตา ทิฏฺฐิปฺปตฺโต โหติ. ทุกฺขโต มนสิกโรโต ปญฺญินฺทฺริยํ อธิมตฺตํ โหติ, ปญฺญินฺทฺริยสฺส อธิมตฺตตฺตา}

\prose{\csromanfolio{249}ทิฏฺฐิปฺปตฺโต โหติ. เอวํ อิเม ตโย ปุคฺคลา ปญฺญินฺทฺริยสฺส วเสน ทิฏฺฐิปฺปตฺตา.}

\prose{โย จายํ ปุคฺคโล สทฺธาวิมุตฺโต, โย จ กายสกฺขี โย จ ทิฏฺฐิปฺปตฺโต, เอวํ สิยา อิเม ตโย ปุคฺคลา สทฺธาวิมุตฺตาปิ กายสกฺขีปิ ทิฏฺฐิปฺปตฺตาปิ วตฺถุวเสน ปริยาเยน.  โย จายํ ปุคฺคโล สทฺธาวิมุตฺโต, โย จ กายสกฺขี, โย จ ทิฏฺฐิปฺปตฺโต, สิยา อิเม ตโย ปุคฺคลา ฯเปฯ. อญฺโญเยว สทฺธาวิมุตฺโต, อญฺโญ กายสกฺขี, อญฺโญ ทิฏฺฐิปฺปตฺโต.}

\prose{สิยาติ กถญฺจ สิยา? อนิจฺจโต มนสิกโรโต สทฺธินฺทฺริยํ อธิมตฺตํ โหติ, สทฺธิ\-นฺทฺริยสฺส อธิมตฺตตฺตา สทฺธาวิมุตฺโต โหติ. ทุกฺขโต มนสิ กโรโต สมาธินฺทฺริยํ อธิมตฺตํ โหติ. สมาธิ\-นฺทฺริยสฺส อธิมตฺตตฺตา กายสกฺขี โหติ. อนตฺตโต มนสิกโรโต ปญฺญินฺทฺริยํ อธิมตฺตํ โหติ, ปญฺญินฺทฺริยสฺส อธิมตฺตตฺตา ทิฏฺฐิปฺปตฺโต โหติ. โย จายํ ปุคฺคโล สทฺธาวิมุตฺโต, โย จ กายสกฺขี, โย จ ทิฏฺฐิปฺปตฺโต, เอวํ สิยา อิเม ตโย ปุคฺคลา สทฺธาวิมุตฺตาปิ กายสกฺขีปิ ทิฏฺฐิปฺปตฺตาปิ วตฺถุวเสน ปริยาเยน. อญฺโญเยว สทฺธาวิมุตฺโต, อญฺโญ กายสกฺขี อญฺโญ ทิฏฺฐิปฺปตฺโต.}

\prose{อนิจฺจโต มนสิกโรโต สทฺธินฺทฺริยํ อธิมตฺตํ โหติ, สทฺธิ\-นฺทฺริยสฺส อธิมตฺตตฺตา โสตาปตฺติ\-มคฺคํ ปฏิลภติ, เตน วุจฺจติ สทฺธานุสารี. จตฺตาริ\-นฺทฺริยานิ ตทนฺวยา โหนฺติ, สหชาต\-ปจฺจยา โหนฺติ, อญฺญมญฺญ\-ปจฺจยา โหนฺติ, นิสฺสยปจฺจยา โหนฺติ, สมฺปยุตฺต\-ปจฺจยา โหนฺติ. สทฺธิ\-นฺทฺริยสฺส วเสน จตุนฺนํ อินฺทฺริยานํ ภาวนา โหติ. เย หิ เกจิ สทฺธิ\-นฺทฺริยสฺส วเสน โสตาปตฺติ\-มคฺคํ ปฏิลภนฺติ, สพฺเพ เต สทฺธานุสาริโน.}

\prose{อนิจฺจโต มนสิกโรโต สทฺธินฺทฺริยํ อธิมตฺตํ โหติ, สทฺธิ\-นฺทฺริยสฺส อธิมตฺตตฺตา โสตาปตฺติผลํ สจฺฉิกตํ โหติ, เตน วุจฺจติ สทฺธาวิมุตฺโต. จตฺตาริ\-นฺทฺริยานิ ตทนฺวยา โหนฺติ, สหชาต\-ปจฺจยา โหนฺติ, อญฺญมญฺญ\-ปจฺจยา โหนฺติ, นิสฺสยปจฺจยา โหนฺติ, สมฺปยุตฺต\-ปจฺจยา โหนฺติ, สทฺธิ\-นฺทฺริยสฺส วเสน จตฺตาริ\-นฺทฺริยานิ ภาวิตานิ โหนฺติ}

\prose{\csromanfolio{250}สุภาวิตานิ. เย หิ เกจิ สทฺธิ\-นฺทฺริยสฺส วเสน โสตาปตฺติผลํ สจฺฉิกตา, สพฺเพ เต สทฺธาวิมุตฺตา.}

\prose{อนิจฺจโต มนสิกโรโต สทฺธินฺทฺริยํ อธิมตฺตํ โหติ. สทฺธิ\-นฺทฺริยสฺส อธิมตฺตตฺตา สกทาคามิ\-มคฺคํ ปฏิลภติ ฯเปฯ สกทาคามิผลํ สจฺฉิกตํ โหติ. อนาคามิมคฺคํ ปฏิลภติ, อนาคามิผลํ สจฺฉิกตํ โหติ. อรหตฺตมคฺคํ ปฏิลภติ, อรหตฺตํ\csromansharedfootnote{911b2d14715a50a6}{อรหตฺตผลํ (สฺยา, ก) อฏฺฐกถา โอโลเกตพฺพา.} สจฺฉิกตํ โหติ, เตน วุจฺจติ สทฺธาวิมุตฺโต. จตฺตาริ\-นฺทฺริยานิ ตทนฺวยา โหนฺติ ฯเปฯ สมฺปยุตฺต\-ปจฺจยา โหนฺติ, สทฺธิ\-นฺทฺริยสฺส วเสน จตฺตาริ\-นฺทฺริยานิ ภาวิตานิ โหนฺติ สุภาวิตานิ. เย หิ เกจิ สทฺธิ\-นฺทฺริยสฺส วเสน อรหตฺตํ สจฺฉิกตา, สพฺเพ เต สทฺธาวิมุตฺตา.}

\prose{ทุกฺขโต มนสิกโรโต สมาธินฺทฺริยํ อธิมตฺตํ โหติ, สมาธิ\-นฺทฺริยสฺส อธิมตฺตตฺตา โสตาปตฺติ\-มคฺคํ ปฏิลภติ, เตน วุจฺจติ กายสกฺขี. จตฺตาริ\-นฺทฺริยานิ ตทนฺวยา โหนฺติ, สหชาต\-ปจฺจยา โหนฺติ, อญฺญมญฺญ\-ปจฺจยา โหนฺติ, นิสฺสยปจฺจยา โหนฺติ. สมฺปยุตฺต\-ปจฺจยา โหนฺติ, สมาธิ\-นฺทฺริยสฺส วเสน จตุนฺนํ อินฺทฺริยานํ ภาวนา โหติ, เย หิ เกจิ สมาธิ\-นฺทฺริยสฺส วเสน โสตาปตฺติ\-มคฺคํ ปฏิลภนฺติ, สพฺเพ เต กายสกฺขี.}

\prose{ทุกฺขโต มนสิกโรโต สมาธินฺทฺริยํ อธิมตฺตํ โหติ, สมาธิ\-นฺทฺริยสฺส อธิมตฺตตฺตา โสตาปตฺติผลํ สจฺฉิกตํ โหติ. สกทาคามิ\-มคฺคํ ปฏิลภติ, สกทาคามิผลํ สจฺฉิกตํ โหติ. อนาคามิมคฺคํ ปฏิลภติ, อนาคามิผลํ สจฺฉิกตํ โหติ. อรหตฺตมคฺคํ ปฏิลภติ, อรหตฺตํ สจฺฉิกตํ โหติ, เตน วุจฺจติ กายสกฺขี. จตฺตาริ\-นฺทฺริยานิ ตทนฺวยา โหนฺติ, สหชาต\-ปจฺจยา โหนฺติ, อญฺญมญฺญ\-ปจฺจยา โหนฺติ, นิสฺสยปจฺจยา โหนฺติ, สมฺปยุตฺต\-ปจฺจยา โหนฺติ. สมาธิ\-นฺทฺริยสฺส วเสน จตฺตาริ\-นฺทฺริยานิ ภาวิตานิ โหนฺติ สุภาวิตานิ. เย หิ เกจิ สมาธิ\-นฺทฺริยสฺส วเสน อรหตฺตํ สจฺฉิกตํ, สพฺเพ เต กายสกฺขี.}

\prose{อนตฺตโต มนสิกโรโต ปญฺญินฺทฺริยํ อธิมตฺตํ โหติ, ปญฺญินฺทฺริยสฺส อธิมตฺตตฺตา โสตาปตฺติ\-มคฺคํ ปฏิลภติ, เตน วุจฺจติ ธมฺมานุสารี. จตฺตาริ\-นฺทฺริยานิ ตทนฺวยา โหนฺติ ฯเปฯ สมฺปยุตฺต\-ปจฺจยา โหนฺติ. ปญฺญินฺทฺริยสฺส}

\prose{\csromanfolio{251}วเสน จตุนฺนํ อินฺทฺริยานํ ภาวนา โหติ. เย หิ เกจิ ปญฺญินฺทฺริยสฺส วเสน โสตาปตฺติ\-มคฺคํ ปฏิลภนฺติ, สพฺเพ เต ธมฺมานุสาริโน.}

\prose{อนตฺตโต มนสิกโรโต ปญฺญินฺทฺริยํ อธิมตฺตํ โหติ, ปญฺญินฺทฺริยสฺส อธิมตฺตตฺตา โสตาปตฺติผลํ สจฺฉิกตํ โหติ, เตน วุจฺจติ ทิฏฺฐิปฺปตฺโต. จตฺตาริ\-นฺทฺริยานิ ตทนฺวยา โหนฺติ ฯเปฯ สมฺปยุตฺต\-ปจฺจยา โหนฺติ, ปญฺญินฺทฺริยสฺส วเสน จตฺตาริ\-นฺทฺริยานิ ภาวิตานิ โหนฺติ สุภาวิตานิ. เย หิ เกจิ ปญฺญินฺทฺริยสฺส วเสน โสตาปตฺติผลํ สจฺฉิกตา, สพฺเพ เต ทิฏฺฐิปฺปตฺตา.}

\prose{อนตฺตโต มนสิกโรโต ปญฺญินฺทฺริยํ อธิมตฺตํ โหติ, ปญฺญินฺทฺริยสฺส อธิมตฺตตฺตา สกทาคามิ\-มคฺคํ ปฏิลภติ, สกทาคามิผลํ สจฺฉิกตํ โหติ. อนาคามิมคฺคํ ปฏิลภติ, อนาคามิผลํ สจฺฉิกตํ โหติ. อรหตฺตมคฺคํ ปฏิลภติ, อรหตฺตํ สจฺฉิกตํ โหติ, เตน วุจฺจติ ทิฏฺฐิปฺปตฺโต. จตฺตาริ\-นฺทฺริยานิ ตทนฺวยา โหนฺติ, สหชาต\-ปจฺจยา โหนฺติ, อญฺญมญฺญ\-ปจฺจยา โหนฺติ, นิสฺสยปจฺจยา โหนฺติ, สมฺปยุตฺต\-ปจฺจยา โหนฺติ, ปญฺญินฺทฺริยสฺส วเสน จตฺตาริ\-นฺทฺริยานิ ภาวิตานิ โหนฺติ สุภาวิตานิ. เย หิ เกจิ ปญฺญินฺทฺริยสฺส วเสน อรหตฺตํ สจฺฉิกตา, สพฺเพ เต ทิฏฺฐิปฺปตฺตา.}

\proseitem{222}{เย หิ เกจิ เนกฺขมฺมํ ภาวิตา วา ภาเวนฺติ วา ภาวิสฺสนฺติ วา, อธิคตา วา อธิคจฺฉนฺติ วา อธิคมิสฺสนฺติ วา, ปตฺตา วา ปาปุณนฺติ วา ปาปุณิสฺสนฺติ วา, ปฏิลทฺธา วา ปฏิลภนฺติ วา ปฏิลภิสฺสนฺติ วา, ปฏิวิทฺธา วา ปฏิวิชฺฌนฺติ วา ปฏิวิชฺฌิสฺสนฺติ วา, สจฺฉิกตา วา สจฺฉิกโรนฺติ วา สจฺฉิกริสฺสนฺติ วา, ผสฺสิตา วา ผสฺสนฺติ วา ผสฺสิสฺสนฺติ วา, วสิปฺปตฺตา วา ปาปุณนฺติ วา ปาปุณิสฺสนฺติ วา, ปารมิปฺปตฺตา วา ปาปุณนฺติ วา ปาปุณิสฺสนฺติ วา, เวสารชฺช\-ปฺปตฺตา วา ปาปุณนฺติ วา ปาปุณิสฺสนฺติ วา, สพฺเพ เต สทฺธิ\-นฺทฺริยสฺส วเสน สทฺธาวิมุตฺตา, สมาธิ\-นฺทฺริยสฺส วเสน กายสกฺขี, ปญฺญินฺทฺริยสฺส วเสน ทิฏฺฐิปฺปตฺตา.}

\prose{เย หิ เกจิ อพฺยาปาทํ ฯเปฯ อาโลกสญฺญํ. อวิกฺเขปํ. ธมฺม\-ววตฺถานํ. ญาณํ. ปาโมชฺชํ. ปฐมํ ฌานํ. ทุติยํ ฌานํ. ตติยํ ฌานํ. จตุตฺถํ ฌานํ. อากาสา\-นญฺจา\-ยตนสมาปตฺติํ. วิญฺญาณญฺจา\-ยตนสมาปตฺติํ. อากิญฺจญฺญายตนาสมาปตฺติํ. เนว\-สญฺญา\-นา\-สญฺญา\-ยตนสมาปตฺติํ. อนิจฺจา\-นุปสฺสนํ. ทุกฺขา\-นุปสฺสนํ. อนตฺตา\-นุปสฺสนํ. นิพฺพิทา\-นุปสฺสนํ. \csromanfolio{252}วิราคา\-นุปสฺสนํ. นิโรธา\-นุปสฺสนํ. ปฏินิสฺสคฺคา\-นุปสฺสนํ. ขยานุปสฺสนํ. วยานุปสฺสนํ. วิปริณามา\-นุปสฺสนํ. อนิมิตฺตา\-นุปสฺสนํ. อปฺปณิหิตา\-นุปสฺสนํ. สุญฺญตา\-นุปสฺสนํ. อธิปญฺญา\-ธมฺมวิปสฺสนํ. ยถา\-ภูต\-ญาณ\-ทสฺสนํ. อาทีนวา\-นุปสฺสนํ. ปฏิสงฺขา\-นุปสฺสนํ. วิวฏฺฏนา\-นุปสฺสนํ. โสตาปตฺติ\-มคฺคํ. สกทาคามิ\-มคฺคํ. อนาคามิมคฺคํ. อรหตฺตมคฺคํ.}

\prose{เย หิ เกจิ จตฺตาโร สติปฏฺฐาเน. จตฺตาโร สมฺมปฺปธาเน. จตฺตาโร อิทฺธิปาเท. ปญฺจินฺทฺริยานิ. ปญฺจ พลานิ. สตฺต โพชฺฌงฺเค. อริยํ อฏฺฐงฺคิกํ มคฺคํ. เย หิ เกจิ อฏฺฐ วิโมกฺเข ภาวิตา วา ภาเวนฺติ วา ภาวิสฺสนฺติ วา, อธิคตา วา อธิคจฺฉนฺติ วา อธิคมิสฺสนฺติ วา, ปตฺตา วา ปาปุณนฺติ วา ปาปุณิสฺสนฺติ วา, ปฏิลทฺธา วา ปฏิลภนฺติ วา ปฏิลภิสฺสนฺติ วา, ปฏิวิทฺธา วา ปฏิวิชฺฌนฺติ วา ปฏิวิชฺฌิสฺสนฺติ วา, สจฺฉิกตา วา สจฺฉิกโรนฺติ วา สจฺฉิกริสฺสนฺติ วา, ผสฺสิตา วา ผสฺสนฺติ วา ผสฺสิสฺสนฺติ วา, วสิปฺปตฺตา วา ปาปุณนฺติ วา ปาปุณิสฺสนฺติ วา, ปารมิปฺปตฺตา วา ปาปุณนฺติ วา ปาปุณิสฺสนฺติ วา, เวสารชฺช\-ปฺปตฺตา วา ปาปุณนฺติ วา ปาปุณิสฺสนฺติ วา, สพฺเพ เต สทฺธิ\-นฺทฺริยสฺส วเสน สทฺธาวิมุตฺตา, สมาธิ\-นฺทฺริยสฺส วเสน กายสกฺขี, ปญฺญินฺทฺริยสฺส วเสน ทิฏฺฐิปฺปตฺตา.}

\prose{เย หิ เกจิ จตสฺโส ปฏิสมฺภิทา ปตฺตา วา ปาปุณนฺติ วา ปาปุณิสฺสนฺติ วา ฯเปฯ สพฺเพ เต สทฺธิ\-นฺทฺริยสฺส วเสน สทฺธาวิมุตฺตา, สมาธิ\-นฺทฺริยสฺส วเสน กายสกฺขี, ปญฺญินฺทฺริยสฺส วเสน ทิฏฺฐิปฺปตฺตา.}

\prose{เย หิ เกจิ ติสฺโส วิชฺชา ปฏิวิทฺธา วา ปฏิวิชฺฌนฺติ วา ปฏิวิชฺฌิสฺสนฺติ วา ฯเปฯ สพฺเพ เต สทฺธิ\-นฺทฺริยสฺส วเสน สทฺธาวิมุตฺตา, สมาธิ\-นฺทฺริยสฺส วเสน กายสกฺขี, ปญฺญินฺทฺริยสฺส วเสน ทิฏฺฐิปฺปตฺตา.}

\prose{เย หิ เกจิ ติสฺโส สิกฺขา สิกฺขิตา วา สิกฺขนฺติ วา สิกฺขิสฺสนฺติ วา, สจฺฉิกตา วา สจฺฉิกโรนฺติ วา สจฺฉิกริสฺสนฺติ วา, ผสฺสิตา วา ผสฺสนฺติ วา ผสฺสิสฺสนฺติ วา, วสิปฺปตฺตา วา ปาปุณนฺติ วา ปาปุณิสฺสนฺติ วา, ปารมิปฺปตฺตา วา ปาปุณนฺติ วา ปาปุณิสฺสนฺติ วา, เวสารชฺช\-ปฺปตฺตา วา ปาปุณนฺติ วา ปาปุณิสฺสนฺติ วา, สพฺเพ เต สทฺธิ\-นฺทฺริยสฺส วเสน สทฺธาวิมุตฺตา, สมาธิ\-นฺทฺริยสฺส วเสน กายสกฺขี, ปญฺญินฺทฺริยสฺส วเสน ทิฏฺฐิปฺปตฺตา.}

\prose{\csromanfolio{253}เย หิ เกจิ ทุกฺขํ ปริชานนฺติ. สมุทยํ ปชหนฺติ. นิโรธํ สจฺฉิกโรนฺติ. มคฺคํ ภาเวนฺติ. สพฺเพ เต สทฺธิ\-นฺทฺริยสฺส วเสน สทฺธาวิมุตฺตา, สมาธิ\-นฺทฺริยสฺส วเสน กายสกฺขี, ปญฺญินฺทฺริยสฺส วเสน ทิฏฺฐิปฺปตฺตา.}

\prose{กติหากาเรหิ สจฺจ\-ปฺปฏิเวโธ โหติ, กติหากาเรหิ สจฺจานิ ปฏิวิชฺฌติ? จตูหากาเรหิ สจฺจ\-ปฺปฏิเวโธ โหติ, จตูหากาเรหิ สจฺจานิ ปฏิวิชฺฌติ. ทุกฺขสจฺจํ ปริญฺญา\-ปฏิเวธํ ปฏิวิชฺฌติ, สมุทยสจฺจํ ปหาน\-ปฺปฏิเวธํ ปฏิวิชฺฌติ, นิโรธสจฺจํ สจฺฉิกิริยา\-ปฏิเวธํ ปฏิวิชฺฌติ, มคฺคสจฺจํ ภาวนา\-ปฏิเวธํ ปฏิวิชฺฌติ. อิเมหิ จตูหากาเรหิ สจฺจ\-ปฺปฏิเวโธ โหติ.  อิเมหิ จตูหากาเรหิ สจฺจานิ ปฏิวิชฺฌนฺโต สทฺธิ\-นฺทฺริยสฺส วเสน สทฺธาวิมุตฺโต, สมาธิ\-นฺทฺริยสฺส วเสน กายสกฺขี, ปญฺญินฺทฺริยสฺส วเสน ทิฏฺฐิปฺปตฺโต.}

\prose{กติหากาเรหิ สจฺจ\-ปฺปฏิเวโธ โหติ, กติหากาเรหิ สจฺจานิ ปฏิวิชฺฌติ? นวหากาเรหิ สจฺจ\-ปฺปฏิเวโธ โหติ, นวหากาเรหิ สจฺจานิ ปฏิวิชฺฌติ, ทุกฺขสจฺจํ ปริญฺญา\-ปฏิเวธํ ปฏิวิชฺฌติ, สมุทยสจฺจํ ปหาน\-ปฺปฏิเวธํ ปฏิวิชฺฌติ. นิโรธสจฺจํ สจฺฉิกิริยา\-ปฏิเวธํ ปฏิวิชฺฌติ, มคฺคสจฺจํ ภาวนา\-ปฏิเวธํ ปฏิวิชฺฌติ. อภิญฺญา\-ปฏิเวโธ จ สพฺพธมฺมานํ, ปริญฺญา\-ปฏิเวโธ จ สพฺพ\-สงฺขารานํ, ปหาน\-ปฺปฏิเวโธ จ สพฺพากุสลานํ, ภาวนา\-ปฏิเวโธ จ จตุนฺนํ มคฺคานํ, สจฺฉิกิริยา\-ปฏิเวโธ จ นิโรธสฺส. อิเมหิ นวหากาเรหิ สจฺจ\-ปฺปฏิเวโธ โหติ.  อิเมหิ นวหากาเรหิ สจฺจานิ ปฏิวิชฺฌนฺโต สทฺธิ\-นฺทฺริยสฺส วเสน สทฺธาวิมุตฺโต, สมาธิ\-นฺทฺริยสฺส วเสน กายสกฺขี, ปญฺญินฺทฺริยสฺส วเสน ทิฏฺฐิปฺปตฺโต. ทุติย\-ภาณวาโร.}

\proseitem{223}{อนิจฺจโต มนสิกโรโต กถํ สงฺขารา อุปฏฺฐนฺติ, ทุกฺขโต มนสิกโรโต กถํ สงฺขารา อุปฏฺฐนฺติ, อนตฺตโต มนสิกโรโต กถํ สงฺขารา อุปฏฺฐนฺติ? อนิจฺจโต มนสิกโรโต ขยโต สงฺขารา อุปฏฺฐนฺติ, ทุกฺขโต มนสิกโรโต ภยโต สงฺขารา อุปฏฺฐนฺติ, อนตฺตโต มนสิกโรโต สุญฺญโต สงฺขารา อุปฏฺฐนฺติ.}

\prose{อนิจฺจโต มนสิกโรโต กิํ พหุลํ จิตฺตํ โหติ, ทุกฺขโต มนสิกโรโต กิํ พหุลํ จิตฺตํ โหติ, อนตฺตโต มนสิกโรโต \csromanfolio{254}กิํ พหุลํ จิตฺตํ โหติ? อนิจฺจโต มนสิกโรโต อธิโมกฺข\-พหุลํ จิตฺตํ โหติ, ทุกฺขโต มนสิกโรโต ปสฺสทฺธิ\-พหุลํ จิตฺตํ โหติ, อนตฺตโต มนสิกโรโต เวทพหุลํ จิตฺตํ โหติ.}

\prose{อนิจฺจโต มนสิกโรนฺโต อธิโมกฺข\-พหุโล กตมํ วิโมกฺขํ ปฏิลภติ. ทุกฺขโต มนสิกโรนฺโต ปสฺสทฺธิ\-พหุโล กตมํ วิโมกฺขํ ปฏิลภติ. อนตฺตโต มนสิกโรนฺโต เวทพหุโล กตมํ วิโมกฺขํ ปฏิลภติ.  อนิจฺจโต มนสิกโรนฺโต อธิโมกฺข\-พหุโล อนิมิตฺตํ วิโมกฺขํ ปฏิลภติ. ทุกฺขโต มนสิกโรนฺโต ปสฺสทฺธิ\-พหุโล อปฺปณิหิตํ วิโมกฺขํ ปฏิลภติ. อนตฺตโต มนสิกโรนฺโต เวทพหุโล สุญฺญตํ วิโมกฺขํ ปฏิลภติ.}

\proseitem{224}{อนิจฺจโต มนสิกโรโต อธิโมกฺข\-พหุลสฺส กตโม วิโมกฺโข อาธิปเตยฺโย โหติ, ภาวนาย กติ วิโมกฺขา ตทนฺวยา โหนฺติ, สหชาต\-ปจฺจยา โหนฺติ, อญฺญมญฺญ\-ปจฺจยา โหนฺติ, นิสฺสยปจฺจยา โหนฺติ, สมฺปยุตฺต\-ปจฺจยา โหนฺติ, เอกรสา โหนฺติ. เกนฏฺเฐน ภาวนา, โก ภาเวติ. ทุกฺขโต มนสิกโรโต ปสฺสทฺธิ\-พหุลสฺส กตโม วิโมกฺโข อาธิปเตยฺโย โหติ, ภาวนาย กติ วิโมกฺขา ตทนฺวยา โหนฺติ, สหชาต\-ปจฺจยา โหนฺติ, อญฺญมญฺญ\-ปจฺจยา โหนฺติ, นิสฺสยปจฺจยา โหนฺติ, สมฺปยุตฺต\-ปจฺจยา โหนฺติ, เอกรสา โหนฺติ. เกนฏฺเฐน ภาวนา, โก ภาเวติ. อนตฺตโต มนสิกโรโต เวทพหุลสฺส กตโม วิโมกฺโข อาธิปเตยฺโย โหติ, ภาวนาย กติ วิโมกฺขา ตทนฺวยา โหนฺติ, สหชาต\-ปจฺจยา โหนฺติ, อญฺญมญฺญ\-ปจฺจยา โหนฺติ, นิสฺสยปจฺจยา โหนฺติ, สมฺปยุตฺต\-ปจฺจยา โหนฺติ, เอกรสา โหนฺติ. เกนฏฺเฐน ภาวนา, โก ภาเวติ?}

\prose{อนิจฺจโต มนสิกโรโต อธิโมกฺข\-พหุลสฺส อนิมิตฺโต วิโมกฺโข อาธิปเตยฺโย โหติ, ภาวนาย เทฺว วิโมกฺขา ตทนฺวยา โหนฺติ, สหชาต\-ปจฺจยา โหนฺติ, อญฺญมญฺญ\-ปจฺจยา โหนฺติ, นิสฺสยปจฺจยา โหนฺติ, สมฺปยุตฺต\-ปจฺจยา โหนฺติ, เอกรสา โหนฺติ. เอกรสฏฺเฐน ภาวนา. โย สมฺมาปฏิปนฺโน, โส ภาเวติ. นตฺถิ มิจฺฉา\-ปฏิปนฺนสฺส วิโมกฺข\-ภาวนา. ทุกฺขโต มนสิกโรโต ปสฺสทฺธิ\-พหุลสฺส อปฺปณิหิโต วิโมกฺโข อาธิปเตยฺโย โหติ, ภาวนาย เทฺว วิโมกฺขา}

\prose{\csromanfolio{255}ตทนฺวยา โหนฺติ, สหชาต\-ปจฺจยา โหนฺติ, อญฺญมญฺญ\-ปจฺจยา โหนฺติ, นิสฺสยปจฺจยา โหนฺติ, สมฺปยุตฺต\-ปจฺจยา โหนฺติ, เอกรสา โหนฺติ. เอกรสฏฺเฐน ภาวนา. โย สมฺมาปฏิปนฺโน, โส ภาเวติ. นตฺถิ มิจฺฉา\-ปฏิปนฺนสฺส วิโมกฺข\-ภาวนา. อนตฺตโต มนสิกโรโต เวทพหุลสฺส สุญฺญโต วิโมกฺโข อาธิปเตยฺโย โหติ, ภาวนาย เทฺว วิโมกฺขา ตทนฺวยา โหนฺติ, สหชาต\-ปจฺจยา โหนฺติ, อญฺญมญฺญ\-ปจฺจยา โหนฺติ, นิสฺสยปจฺจยา โหนฺติ, สมฺปยุตฺต\-ปจฺจยา โหนฺติ, เอกรสา โหนฺติ, เอกรสฏฺเฐน ภาวนา. โย สมฺมาปฏิปนฺโน, โส ภาเวติ. นตฺถิ มิจฺฉา\-ปฏิปนฺนสฺส วิโมกฺข\-ภาวนา.}

\prose{อนิจฺจโต มนสิกโรโต อธิโมกฺข\-พหุลสฺส กตโม วิโมกฺโข อาธิปเตยฺโย โหติ, ภาวนาย กติ วิโมกฺขา ตทนฺวยา โหนฺติ, สหชาต\-ปจฺจยา โหนฺติ, อญฺญมญฺญ\-ปจฺจยา โหนฺติ, นิสฺสยปจฺจยา โหนฺติ, สมฺปยุตฺต\-ปจฺจยา โหนฺติ, เอกรสา โหนฺติ. ปฏิเวธกาเล กตโม วิโมกฺโข อาธิปเตยฺโย โหติ, ปฏิเวธาย กติ วิโมกฺขา ตทนฺวยา โหนฺติ, สหชาต\-ปจฺจยา โหนฺติ, อญฺญมญฺญ\-ปจฺจยา โหนฺติ, นิสฺสยปจฺจยา โหนฺติ, สมฺปยุตฺต\-ปจฺจยา โหนฺติ, เอกรสา โหนฺติ. เกนฏฺเฐน ภาวนา, เกนฏฺเฐน ปฏิเวโธ.  ทุกฺขโต มนสิกโรโต ปสฺสทฺธิ\-พหุลสฺส กตโม วิโมกฺโข อาธิปเตยฺโย โหติ, ภาวนาย กติ วิโมกฺขา ตทนฺวยา โหนฺติ, สหชาต\-ปจฺจยา โหนฺติ, อญฺญมญฺญ\-ปจฺจยา โหนฺติ, นิสฺสยปจฺจยา โหนฺติ, สมฺปยุตฺต\-ปจฺจยา โหนฺติ, เอกรสา โหนฺติ. ปฏิเวธกาเล กตโม วิโมกฺโข อาธิปเตยฺโย โหติ, ปฏิเวธาย กติ วิโมกฺขา ตทนฺวยา โหนฺติ, สหชาต\-ปจฺจยา โหนฺติ, อญฺญมญฺญ\-ปจฺจยา โหนฺติ, นิสฺสยปจฺจยา โหนฺติ, สมฺปยุตฺต\-ปจฺจยา โหนฺติ, เอกรสา โหนฺติ. เกนฏฺเฐน ภาวนา, เกนฏฺเฐน ปฏิเวโธ.  อนตฺตโต มนสิกโรโต เวทพหุลสฺส กตโม วิโมกฺโข อาธิปเตยฺโย โหติ, ภาวนาย กติ วิโมกฺขา ตทนฺวยา โหนฺติ, สหชาต\-ปจฺจยา โหนฺติ, อญฺญมญฺญ\-ปจฺจยา โหนฺติ, นิสฺสยปจฺจยา โหนฺติ, สมฺปยุตฺต\-ปจฺจยา โหนฺติ, เอกรสา โหนฺติ. ปฏิเวธกาเล กตโม วิโมกฺโข อาธิปเตยฺโย โหติ, ปฏิเวธาย กติวิโมกฺขา ตทนฺวยา โหนฺติ, สหชาต\-ปจฺจยา โหนฺติ, อญฺญมญฺญ\-ปจฺจยา}

\prose{\csromanfolio{256}โหนฺติ, นิสฺสยปจฺจยา โหนฺติ, สมฺปยุตฺต\-ปจฺจยา โหนฺติ, เอกรสา โหนฺติ. เกนฏฺเฐน ภาวนา, เกนฏฺเฐน ปฏิเวโธ?}

\proseitem{225}{อนิจฺจโต มนสิกโรโต อธิโมกฺข\-พหุลสฺส อนิมิตฺโต วิโมกฺโข อาธิปเตยฺโย โหติ, ภาวนาย เทฺว วิโมกฺขา ตทนฺวยา โหนฺติ, สหชาต\-ปจฺจยา โหนฺติ, อญฺญมญฺญ\-ปจฺจยา โหนฺติ, นิสฺสยปจฺจยา โหนฺติ, สมฺปยุตฺต\-ปจฺจยา โหนฺติ, เอกรสา โหนฺติ. ปฏิเวธกาเลปิ อนิมิตฺโต วิโมกฺโข อาธิปเตยฺโย โหติ, ปฏิเวธาย เทฺว วิโมกฺขา ตทนฺวยา โหนฺติ, สหชาต\-ปจฺจยา โหนฺติ, อญฺญมญฺญ\-ปจฺจยา โหนฺติ, นิสฺสยปจฺจยา โหนฺติ, สมฺปยุตฺต\-ปจฺจยา โหนฺติ, เอกรสา โหนฺติ. เอกรสฏฺเฐน ภาวนา, ทสฺสนฏฺเฐน ปฏิเวโธ. เอวํ ปฏิวิชฺฌนฺโตปิ ภาเวติ, ภาเวนฺโตปิ ปฏิวิชฺฌติ.  ทุกฺขโต มนสิกโรโต ปสฺสทฺธิ\-พหุลสฺส อปฺปณิหิโต วิโมกฺโข อาธิปเตยฺโย โหติ, ภาวนาย เทฺว วิโมกฺขา ตทนฺวยา โหนฺติ, สหชาต\-ปจฺจยา โหนฺติ, อญฺญมญฺญ\-ปจฺจยา โหนฺติ, นิสฺสยปจฺจยา โหนฺติ, สมฺปยุตฺต\-ปจฺจยา โหนฺติ, เอกรสา โหนฺติ. ปฏิเวธกาเลปิ อปฺปณิหิโต วิโมกฺโข อาธิปเตยฺโย โหติ, ปฏิเวธาย เทฺว วิโมกฺขา ตทนฺวยา โหนฺติ, สหชาต\-ปจฺจยา โหนฺติ, อญฺญมญฺญ\-ปจฺจยา โหนฺติ, นิสฺสยปจฺจยา โหนฺติ, สมฺปยุตฺต\-ปจฺจยา โหนฺติ, เอกรสา โหนฺติ. เอกรสฏฺเฐน ภาวนา, ทสฺสนฏฺเฐน ปฏิเวโธ. เอวํ ปฏิวิชฺฌนฺโตปิ ภาเวติ, ภาเวนฺโตปิ ปฏิวิชฺฌติ.  อนตฺตโต มนสิกโรโต เวทพหุลสฺส สุญฺญโต วิโมกฺโข อาธิปเตยฺโย โหติ, ภาวนาย เทฺว วิโมกฺขา ตทนฺวยา โหนฺติ, สหชาต\-ปจฺจยา โหนฺติ, อญฺญมญฺญ\-ปจฺจยา โหนฺติ, นิสฺสยปจฺจยา โหนฺติ, สมฺปยุตฺต\-ปจฺจยา โหนฺติ. เอกรสา โหนฺติ. ปฏิเวธกาเลปิ สุญฺญโต วิโมกฺโข อาธิปเตยฺโย โหติ, ปฏิเวธาย เทฺว วิโมกฺขา ตทนฺวยา โหนฺติ, สหชาต\-ปจฺจยา โหนฺติ, อญฺญมญฺญ\-ปจฺจยา โหนฺติ, นิสฺสยปจฺจยา โหนฺติ, สมฺปยุตฺต\-ปจฺจยา โหนฺติ, เอกรสา โหนฺติ, เอกรสฏฺเฐน ภาวนา. ทสฺสนฏฺเฐน ปฏิเวธา. เอวํ ปฏิวิชฺฌนฺโตปิ ภาเวติ, ภาเวนฺโตปิ ปฏิวิชฺฌติ.}

\proseitem{226}{อนิจฺจโต มนสิกโรโต กตโม วิโมกฺโข อธิมตฺโต โหติ, กตม\-วิโมกฺขสฺส อธิมตฺตตฺตา สทฺธาวิมุตฺโต โหติ.}

\prose{\csromanfolio{257}ทุกฺขโต มนสิกโรโต กตโม วิโมกฺโข อธิมตฺโต โหติ, กตม\-วิโมกฺขสฺส อธิมตฺตตฺตา กายสกฺขี โหติ. อนตฺตโต มนสิกโรโต กตโม วิโมกฺโข อธิมตฺโต โหติ, กตม\-วิโมกฺขสฺส อธิมตฺตตฺตา ทิฏฺฐิปฺปตฺโต โหติ?}

\prose{อนิจฺจโต มนสิกโรโต อนิมิตฺโต วิโมกฺโข อธิมตฺโต โหติ, อนิมิตฺต\-วิโมกฺขสฺส อธิมตฺตตฺตา สทฺธาวิมุตฺโต โหติ. ทุกฺขโต มนสิกโรโต อปฺปณิหิโต วิโมกฺโข อธิมตฺโต โหติ, อปฺปณิหิต\-วิโมกฺขสฺส อธิมตฺตตฺตา กายสกฺขี โหติ. อนตฺตโต มนสิกโรโต สุญฺญโต วิโมกฺโข อธิมตฺโต โหติ, สุญฺญต\-วิโมกฺขสฺส อธิมตฺตตฺตา ทิฏฺฐิปฺปตฺโต โหติ.}

\prose{สทฺทหนฺโต วิมุตฺโตติ สทฺธาวิมุตฺโต, ผุฏฺฐตฺตา สจฺฉิกโตติ กายสกฺขี, ทิฏฺฐตฺตา ปตฺโตติ ทิฏฺฐิปฺปตฺโต. สทฺทหนฺโต วิมุจฺจตีติ สทฺธาวิมุตฺโต. ฌานผสฺสํ ปฐมํ ผุสติ, ปจฺฉา นิโรธํ นิพฺพานํ สจฺฉิกโรตีติ กายสกฺขี. “ทุกฺขา สงฺขารา, สุโข นิโรโธ”ติ วิญฺญาตํ โหติ ทิฏฺฐํ วิทิตํ สจฺฉิกตํ ผสฺสิตํ ปญฺญายาติ ทิฏฺฐิปฺปตฺโต ฯเปฯ.}

\prose{เย หิ เกจิ เนกฺขมฺมํ ภาวิตา วา ภาเวนฺติ วา ภาวิสฺสนฺติ วา ฯเปฯ สฺพฺเพ เต อนิมิตฺต\-วิโมกฺขสฺส วเสน สทฺธาวิมุตฺตา, อปฺปณิหิต\-วิโมกฺขสฺส วเสน กายสกฺขี, สุญฺญต\-วิโมกฺขสฺส วเสน ทิฏฺฐิปฺปตฺตา.}

\prose{เย หิ เกจิ อพฺยาปาทํ ฯเปฯ. อาโลกสญฺญํ. อวิกฺเขปํ ฯเปฯ. เย หิ เกจิ ทุกฺขํ ปริชานนฺติ, สมุทยํ ปชหนฺติ, นิโรธํ สจฺฉิกโรนฺติ, มคฺคํ ภาเวนฺติ. สพฺเพ เต อนิมิตฺต\-วิโมกฺขสฺส วเสน สทฺธาวิมุตฺตา, อปฺปณิหิต\-วิโมกฺขสฺส วเสน กายสกฺขี, สุญฺญต\-วิโมกฺขสฺส วเสน ทิฏฺฐิปฺปตฺตา.}

\prose{กติหากาเรหิ สจฺจ\-ปฺปฏิเวโธ โหติ, กติหากาเรหิ สจฺจานิ ปฏิวิชฺฌติ? จตูหากาเรหิ สจฺจ\-ปฺปฏิเวโธ โหติ, จตูหากาเรหิ สจฺจานิ ปฏิวิชฺฌติ. ทุกฺขสจฺจํ ปริญฺญา\-ปฏิเวธํ ปฏิวิชฺฌติ, สมุทยสจฺจํ ปหาน\-ปฺปฏิเวธํ ปฏิวิชฺฌติ. นิโรธสจฺจํ สจฺฉิกิริยา\-ปฏิเวธํ ปฏิวิชฺฌติ. มคฺคสจฺจํ ภาวนา\-ปฏิเวธํ ปฏิวิชฺฌติ. อิเมหิ จตูหากาเรหิ สจฺจ\-ปฺปฏิเวโธ โหติ.  อิเมหิ จตูหากาเรหิ สจฺจานิ ปฏิวิชฺฌนฺโต อนิมิตฺต\-วิโมกฺขสฺส วเสน สทฺธาวิมุตฺโต, อปฺปณิหิต\-วิโมกฺขสฺส วเสน กายสกฺขี, สุญฺญต\-วิโมกฺขสฺส วเสน ทิฏฺฐิปฺปตฺโต.}

\prose{\csromanfolio{258}กติหากาเรหิ สจฺจ`ปฏิเวโธ โหติ, กติหากาเรหิ สจฺจานิ ปฏิวิชฺฌติ? นวหากาเรหิ สจฺจ\-ปฺปฏิเวโธ โหติ, นวหากาเรหิ สจฺจานิ ปฏิวิชฺฌติ. ทุกฺขสจฺจํ ปริญฺญา\-ปฏิเวธํ ปฏิวิชฺฌติ ฯเปฯ. สจฺฉิกิริยา\-ปฏิเวโธ จ นิโรธสฺส. อิเมหิ นวหากาเรหิ สจฺจ\-ปฺปฏิเวโธ โหติ.  อิเมหิ นวหากาเรหิ สจฺจานิ ปฏิวิชฺฌนฺโต อนิมิตฺต\-วิโมกฺขสฺส วเสน สทฺธาวิมุตฺโต, อปฺปณิหิต\-วิโมกฺขสฺส วเสน กายสกฺขี, สุญฺญต\-วิโมกฺขสฺส วเสน ทิฏฺฐิปฺปตฺโต.}

\proseitem{227}{อนิจฺจโต มนสิกโรนฺโต กตเม ธมฺเม ยถาภูตํ ชานาติ ปสฺสติ, กถํ สมฺมาทสฺสนํ โหติ, กถํ ตทนฺวเยน สพฺเพ สงฺขารา อนิจฺจโต สุทิฏฺฐา โหนฺติ, กตฺถ กงฺขา ปหียติ. ทุกฺขโต มนสิกโรนฺโต กตเม ธมฺเม ยถาภูตํ ชานาติ ปสฺสติ, กถํ สมฺมาทสฺสนํ โหติ, กถํ ตทนฺวเยน สพฺเพ สงฺขารา ทุกฺขโต สุทิฏฺฐา โหนฺติ, กตฺถ กงฺขา ปหียติ. อนตฺตโต มนสิกโรนฺโต กตเม ธมฺเม ยถาภูตํ ชานาติ ปสฺสติ, กถํ สมฺมาทสฺสนํ โหติ, กถํ ตทนฺวเยน สพฺเพ ธมฺมา อนตฺตโต สุทิฏฺฐา โหนฺติ, กตฺถ กงฺขา ปหียติ?}

\prose{อนิจฺจโต มนสิกโรนฺโต นิมิตฺตํ ยถาภูตํ ชานาติ ปสฺสติ, เตน วุจฺจติ สมฺมาทสฺสนํ. เอวํ ตทนฺวเยน สพฺเพ สงฺขารา อนิจฺจโต สุทิฏฺฐา โหนฺติ, เอตฺถ กงฺขา ปหียติ. ทุกฺขโต มนสิกโรนฺโต ปวตฺตํ ยถาภูตํ ชานาติ ปสฺสติ, เตน วุจฺจติ สมฺมาทสฺสนํ. เอวํ ตทนฺวเยน สพฺเพ สงฺขารา ทุกฺขโต สุทิฏฺฐา โหนฺติ, เอตฺถ กงฺขา ปหียติ. อนตฺตโต มนสิกโรนฺโต นิมิตฺตญฺจ ปวตฺตญฺจ ยถาภูตํ ชานาติ ปสฺสติ, เตน วุจฺจติ สมฺมาทสฺสนํ. เอวํ ตทนฺวเยน สพฺเพ ธมฺมา อนตฺตโต สุทิฏฺฐา โหนฺติ, เอตฺถ กงฺขา ปหียติ.}

\prose{ยญฺจ ยถาภูตํ ญาณํ ยญฺจ สมฺมาทสฺสนํ ยา จ กงฺขาวิตรณา, อิเม ธมฺมา นานตฺถา เจว นานาพฺยญฺชนา จ, อุทาหุ เอกตฺถา, พฺยญฺชนเมว นานนฺติ? ยญฺจ ยถาภูตํ ญาณํ ยญฺจ สมฺมาทสฺสนํ ยา จ กงฺขาวิตรณา, อิเม ธมฺมา เอกตฺถา, พฺยญฺชนเมว นานํ.}

\prose{อนิจฺจโต มนสิกโรโต กิํ ภยโต อุปฏฺฐาติ, ทุกฺขโต มนสิกโรโต กิํ ภยโต อุปฏฺฐาติ, อนตฺตโต มนสิกโรโต กิํ ภยโต อุปฏฺฐาติ? อนิจฺจโต มนสิกโรโต นิมิตฺตํ ภยโต}

\prose{\csromanfolio{259}อุปฏฺฐาติ, ทุกฺขโต มนสิกโรโต ปวตฺตํ ภยโต อุปฏฺฐาติ, อนตฺตโต มนสิกโรโต นิมิตฺตญฺจ ปวตฺตญฺจ ภยโต อุปฏฺฐาติ.}

\prose{ยา จ ภยตุปฏฺฐาเน ปญฺญา ยญฺจ อาทีนเว ญาณํ ยา จ นิพฺพิทา, อิเม ธมฺมา นานตฺถา เจว นานาพฺยญฺชนา จ. อุทาหุ เอกตฺถา, พฺยญฺชนเมว นานนฺติ. ยา จ ภยตุปฏฺฐาเน ปญฺญา ยญฺจ อาทีนเว ญาณํ ยา จ นิพฺพิทา, อิเม ธมฺมา เอกตฺถา, พฺยญฺชนเมว นานํ.}

\prose{ยา จ อนตฺตา\-นุปสฺสนา ยา จ สุญฺญตา\-นุปสฺสนา, อิเม ธมฺมา นานตฺถา เจว นานาพฺยญฺชนา จ. อุทาหุ เอกตฺถา, พฺยญฺชนเมว นานนฺติ? ยา จ อนตฺตา\-นุปสฺสนา ยา จ สุญฺญตา\-นุปสฺสนา, อิเม ธมฺมา เอกตฺถา, พฺยญฺชนเมว นานํ.}

\prose{อนิจฺจโต มนสิกโรโต กิํ ปฏิสงฺขา ญาณํ อุปฺปชฺชติ, ทุกฺขโต มนสิกโรโต กิํ ปฏิสงฺขา ญาณํ อุปฺปชฺชติ, อนตฺตโต มนสิกโรโต กิํ ปฏิสงฺขา ญาณํ อุปฺปชฺชติ? อนิจฺจโต มนสิกโรโต นิมิตฺตํ ปฏิสงฺขา ญาณํ อุปฺปชฺชติ, ทุกฺขโต มนสิกโรโต ปวตฺตํ ปฏิสงฺขา ญาณํ อุปฺปชฺชติ, อนตฺตโต มนสิกโรโต นิมิตฺตญฺจ ปวตฺตญฺจ ปฏิสงฺขา ญาณํ อุปฺปชฺชติ.}

\prose{ยา จ มุจฺจิตุกมฺยตา ยา จ ปฏิสงฺขา\-นุปสฺสนา ยา จ สงฺขารุเปกฺขา, อิเม ธมฺมา นานตฺถา เจว นานาพฺยญฺชนา จ. อุทาหุ เอกตฺถา, พฺยญฺชนเมว นานนฺติ? ยา จ มุจฺจิตุกมฺยตา ยา จ ปฏิสงฺขา\-นุปสฺสนา ยา จ สงฺขารุเปกฺขา อิเม ธมฺมา เอกตฺถา, พฺยญฺชนเมว นานํ.}

\prose{อนิจฺจโต มนสิกโรโต กุโต จิตฺตํ วุฏฺฐาติ, กตฺถ จิตฺตํ ปกฺขนฺทติ. ทุกฺขโต มนสิกโรโต กุโต จิตฺตํ วุฏฺฐาติ, กตฺถ จิตฺตํ ปกฺขนฺทติ. อนตฺตโต มนสิกโรโต กุโต จิตฺตํ วุฏฺฐาติ, กตฺถ จิตฺตํ ปกฺขนฺทติ? อนิจฺจโต มนสิกโรโต นิมิตฺตา จิตฺตํ วุฏฺฐาติ, อนิมิตฺเต จิตฺตํ ปกฺขนฺทติ. ทุกฺขโต มนสิกโรโต ปวตฺตา จิตฺตํ วุฏฺฐาติ, อปฺปวตฺเต จิตฺตํ ปกฺขนฺทติ. อนตฺตโต มนสิกโรโต นิมิตฺตา จ ปวตฺตา จ จิตฺตํ วุฏฺฐาติ, อนิมิตฺเต อปฺปวตฺเต นิโรเธ นิพฺพานธาตุยา\csromansharedfootnote{5d3895473c58a39d}{นิโรธนิพฺพานธาตุยา (สี)} จิตฺตํ ปกฺขนฺทติ.}

\prose{ยา จ พหิทฺธาวุฏฺฐาน\-วิวฏฺฏเน ปญฺญา เย จ โคตฺรภู ธมฺมา, อิเม ธมฺมา นานตฺถา เจว นานาพฺยญฺชนา จ. อุทาหุ เอกตฺถา, พฺยญฺชนเมว นานนฺติ? ยา จ}

\prose{\csromanfolio{260}พหิทฺธาวุฏฺฐาน\-วิวฏฺฏเน ปญฺญา เย จ โคตฺรภู ธมฺมา, อิเม ธมฺมา เอกตฺถา, พฺยญฺชนเมว นานํ.}

\prose{อนิจฺจโต มนสิกโรนฺโต กตเมน วิโมกฺเขน วิมุจฺจติ, ทุกฺขโต มนสิกโรนฺโต กตเมน วิโมกฺเขน วิมุจฺจติ, อนตฺตโต มนสิกโรนฺโต กตเมน วิโมกฺเขน วิมุจฺจติ? อนิจฺจโต มนสิกโรนฺโต อนิมิตฺต\-วิโมกฺเขน วิมุจฺจติ, ทุกฺขโต มนสิกโรนฺโต อปฺปณิหิต\-วิโมกฺเขน วิมุจฺจติ, อนตฺตโต มนสิกโรนฺโต สุญฺญต\-วิโมกฺเขน วิมุจฺจติ.}

\prose{ยา จ ทุภโต วุฏฺฐาน\-วิวฏฺฏเน ปญฺญา ยญฺจ มคฺเค ญาณํ, อิเม ธมฺมา นานตฺถา เจว นานาพฺยญฺชนา จ. อุทาหุ เอกตฺถา, พฺยญฺชนเมว นานนฺติ? ยา จ ทุภโต วุฏฺฐาน\-วิวฏฺฏเน ปญฺญา ยญฺจ มคฺเค ญาณํ, อิเม ธมฺมา เอกตฺถา, พฺยญฺชนเมว นานํ.}

\proseitem{228}{กติหากาเรหิ ตโย วิโมกฺขา นานากฺขเณ\csromansharedfootnote{334619900cde0bb8}{นานาขเณ (สฺยา, ก)} โหนฺติ กติหากาเรหิ ตโย วิโมกฺขา เอกกฺขเณ โหนฺติ? จตูหากาเรหิ ตโย วิโมกฺขา นานากฺขเณ โหนฺติ, สตฺตหากาเรหิ ตโย วิโมกฺขา เอกกฺขเณ โหนฺติ.}

\prose{กตเมหิ จตูหากาเรหิ ตโย วิโมกฺขา นานากฺขเณ โหนฺติ? อาธิปเตยฺย\-ฏฺเฐน อธิฏฺฐาน\-ฏฺเฐน อภินีหาร\-ฏฺเฐน นิยฺยานฏฺเฐน. กถํ อาธิปเตยฺย\-ฏฺเฐน ตโย วิโมกฺขา นานากฺขเณ โหนฺติ? อนิจฺจโต มนสิกโรโต อนิมิตฺโต วิโมกฺโข อาธิปเตยฺโย โหติ, ทุกฺขโต มนสิกโรโต อปฺปณิหิโต วิโมกฺโข อาธิปเตยฺโย โหติ, อนตฺตโต มนสิกโรโต สุญฺญโต วิโมกฺโข อาธิปเตยฺโย โหติ. เอวํ อาธิปเตยฺย\-ฏฺเฐน ตโย วิโมกฺขา นานากฺขเณ โหนฺติ.}

\prose{กถํ อธิฏฺฐาน\-ฏฺเฐน ตโย วิโมกฺขา นานากฺขเณ โหนฺติ? อนิจฺจโต มนสิกโรนฺโต อนิมิตฺต\-วิโมกฺขสฺส วเสน จิตฺตํ อธิฏฺฐาติ, ทุกฺขโต มนสิกโรนฺโต อปฺปณิหิต\-วิโมกฺขสฺส วเสน จิตฺตํ อธิฏฺฐาติ, อนตฺตโต มนสิกโรนฺโต สุญฺญต\-วิโมกฺขสฺส วเสน จิตฺตํ อธิฏฺฐาติ. เอวํ อธิฏฺฐาน\-ฏฺเฐน ตโย วิโมกฺขา นานากฺขเณ โหนฺติ.}

\prose{\csromanfolio{261}กถํ อภินีหาร\-ฏฺเฐน ตโย วิโมกฺขา นานากฺขเณ โหนฺติ? อนิจฺจโต มนสิกโรนฺโต อนิมิตฺต\-วิโมกฺขสฺส วเสน จิตฺตํ อภินีหรติ, ทุกฺขโต มนสิกโรนฺโต อปฺปณิหิต\-วิโมกฺขสฺส วเสน จิตฺตํ อภินีหรติ, อนตฺตโต มนสิกโรนฺโต สุญฺญต\-วิโมกฺขสฺส วเสน จิตฺตํ อภินีหรติ. เอวํ อภินีหาร\-ฏฺเฐน ตโย วิโมกฺขา นานากฺขเณ โหนฺติ.}

\prose{กถํ นิยฺยานฏฺเฐน ตโย วิโมกฺขา นานากฺขเณ โหนฺติ? อนิจฺจโต มนสิกโรนฺโต อนิมิตฺต\-วิโมกฺขสฺส วเสน นิโรธํ นิพฺพานํ นิยฺยาติ, ทุกฺขโต มนสิกโรนฺโต อปฺปณิหิต\-วิโมกฺขสฺส วเสน นิโรธํ นิพฺพานํ นิยฺยาติ, อนตฺตโต มนสิกโรนฺโต สุญฺญต\-วิโมกฺขสฺส วเสน นิโรธํ นิพฺพานํ นิยฺยาติ. เอวํ นิยฺยานฏฺเฐน ตโย วิโมกฺขา นานากฺขเณ โหนฺติ.  อิเมหิ จตูหากาเรหิ ตโย วิโมกฺขา นานากฺขเณ โหนฺติ.}

\prose{กตเมหิ สตฺตหากาเรหิ ตโย วิโมกฺขา เอกกฺขเณ โหนฺติ? สโมธาน\-ฏฺเฐน อธิคม\-นฏฺเฐน ปฏิลาภ\-ฏฺเฐน ปฏิเวธ\-ฏฺเฐน สจฺฉิกิริยฏฺเฐน ผสฺสนฏฺเฐน อภิสมย\-ฏฺเฐน. กถํ สโมธาน\-ฏฺเฐน อธิคม\-นฏฺเฐน ปฏิลาภ\-ฏฺเฐน ปฏิเวธ\-ฏฺเฐน สจฺฉิกิริยฏฺเฐน ผสฺสนฏฺเฐน อภิสมย\-ฏฺเฐน ตโย วิโมกฺขา เอกกฺขเณ โหนฺติ? อนิจฺจโต มนสิกโรนฺโต นิมิตฺตา มุจฺจตีติ อนิมิตฺโต วิโมกฺโข. ยโต มุจฺจติ, ตตฺถ น ปณิทหตีติ อปฺปณิหิโต วิโมกฺโข. ยตฺถ น ปณิทหติ, เตน สุญฺโญติ สุญฺญโต วิโมกฺโข. เยน สุญฺโญ, เตน นิมิตฺเตน อนิมิตฺโตติ อนิมิตฺโต วิโมกฺโข. เอวํ สโมธาน\-ฏฺเฐน อธิคม\-นฏฺเฐน ปฏิลาภ\-ฏฺเฐน ปฏิเวธ\-ฏฺเฐน สจฺฉิกิริยฏฺเฐน ผสฺสนฏฺเฐน อภิสมย\-ฏฺเฐน ตโย วิโมกฺขา เอกกฺขเณ โหนฺติ.}

\prose{ทุกฺขโต มนสิกโรนฺโต ปณิธิยา มุจฺจตีติ อปฺปณิหิโต วิโมกฺโข. ยตฺถ น ปณิทหติ, เตน สุญฺโญติ สุญฺญโต วิโมกฺโข. เยน สุญฺโญ, เตน นิมิตฺเตน อนิมิตฺโตติ อนิมิตฺโต วิโมกฺโข. เยน นิมิตฺเตน อนิมิตฺโต, ตตฺถ น ปณิทหตีติ อปฺปณิหิโต วิโมกฺโข. เอวํ สโมธาน\-ฏฺเฐน อธิคม\-นฏฺเฐน ปฏิลาภ\-ฏฺเฐน ปฏิเวธ\-ฏฺเฐน สจฺฉิกิริยฏฺเฐน ผสฺสนฏฺเฐน อภิสมย\-ฏฺเฐน ตโย วิโมกฺขา เอกกฺขเณ โหนฺติ.}

\prose{\csromanfolio{262}อนตฺตโต มนสิกโรนฺโต อภินิเวสา มุจฺจตีติ สุญฺญโต วิโมกฺโข. เยน สุญฺโญ, เตน นิมิตฺเตน อนิมิตฺโตติ อนิมิตฺโต วิโมกฺโข. เยน นิมิตฺเตน อนิมิตฺโต, ตตฺถ น ปณิทหตีติ อปฺปณิหิโต วิโมกฺโข. ยตฺถ น ปณิทหติ, เตน สุญฺโญติ สุญฺญโต วิโมกฺโข. เอวํ สโมธาน\-ฏฺเฐน อธิคม\-นฏฺเฐน ปฏิลาภ\-ฏฺเฐน ปฏิเวธ\-ฏฺเฐน สจฺฉิกิริยฏฺเฐน ผสฺสนฏฺเฐน อภิสมย\-ฏฺเฐน ตโย วิโมกฺขา เอกกฺขเณ โหนฺติ. อิเมหิ สตฺตหากาเรหิ ตโย วิโมกฺขา เอกกฺขเณ โหนฺติ.}

\proseitem{229}{อตฺถิ วิโมกฺโข, อตฺถิ มุขํ, อตฺถิ วิโมกฺขมุขํ, อตฺถิ วิโมกฺข\-ปจฺจนีกํ, อตฺถิ วิโมกฺขา\-นุโลมํ, อตฺถิ วิโมกฺข\-วิวฏฺโฏ, อตฺถิ วิโมกฺข\-ภาวนา, อตฺถิ วิโมกฺข\-ปฺปฏิปฺปสฺสทฺธิ.}

\prose{กตโม วิโมกฺโข? สุญฺญโต วิโมกฺโข อนิมิตฺโต วิโมกฺโข อปฺปณิหิโต วิโมกฺโข.  กตโม สุญฺญโต วิโมกฺโข? อนิจฺจานุปสฺสน\-ญาณํ นิจฺจโต อภินิเวสา มุจฺจตีติ สุญฺญโต วิโมกฺโข. ทุกฺขานุปสฺสน\-ญาณํ สุขโต อภินิเวสา มุจฺจตีติ สุญฺญโต วิโมกฺโข. อนตฺตานุปสฺสน\-ญาณํ อตฺตโต อภินิเวสา มุจฺจตีติ สุญฺญโต วิโมกฺโข. นิพฺพิทานุปสฺสน\-ญาณํ นนฺทิยา อภินิเวสา มุจฺจตีติ สุญฺญโต วิโมกฺโข. วิราคานุปสฺสน\-ญาณํ ราคโต อภินิเวสา มุจฺจตีติ สุญฺญโต วิโมกฺโข. นิโรธานุปสฺสน\-ญาณํ สมุทยโต อภินิเวสา มุจฺจตีติ สุญฺญโต วิโมกฺโข. ปฏินิสฺสคฺคานุปสฺสน\-ญาณํ อาทานโต อภินิเวสา มุจฺจตีติ สุญฺญโต วิโมกฺโข. อนิมิตฺตา\-นุปสฺสน\-ญาณํ นิมิตฺตโต อภินิเวสา มุจฺจตีติ สุญฺญโต วิโมกฺโข. อปฺปณิหิตา\-นุปสฺสน\-ญาณํ ปณิธิยา อภินิเวสา มุจฺจตีติ สุญฺญโต วิโมกฺโข. สุญฺญตานุปสฺสน\-ญาณํ สพฺพาภิ\-นิเวเสหิ มุจฺจตีติ สุญฺญโต วิโมกฺโข.}

\prose{รูเป อนิจฺจานุปสฺสน\-ญาณํ นิจฺจโต อภินิเวสา มุจฺจตีติ สุญฺญโต วิโมกฺโข ฯเปฯ รุเป สุญฺญตานุปสฺสน\-ญาณํ สพฺพาภิ\-นิเวเสหิ มุจฺจตีติ สุญฺญโต วิโมกฺโข.  เวทนาย ฯเปฯ. สญฺญาย. สงฺขาเรสุ. วิญฺญาเณ. จกฺขุสฺมิํ ฯเปฯ. ชรามรเณ อนิจฺจานุปสฺสน\-ญาณํ นิจฺจโต อภินิเวสา มุจฺจตีติ สุญฺญโต วิโมกฺโข ฯเปฯ ชรามรเณ สุญฺญตานุปสฺสน\-ญาณํ}

\prose{\csromanfolio{263}สพฺพาภิ\-นิเวเสหิ มุจฺจตีติ สุญฺญโต วิโมกฺโข. อยํ สุญฺญโต วิโมกฺโข.}

\prose{กตโม อนิมิตฺโต วิโมกฺโข? อนิจฺจานุปสฺสน\-ญาณํ นิจฺจโต นิมิตฺตา มุจฺจตีติ อนิมิตฺโต วิโมกฺโข. ทุกฺขานุปสฺสน\-ญาณํ สุขโต นิมิตฺตา มุจฺจตีติ อนิมิตฺโต วิโมกฺโข. อนตฺตานุปสฺสน\-ญาณํ อตฺตโต นิมิตฺตา มุจฺจตีติ อนิมิตฺโต วิโมกฺโข. นิพฺพิทานุปสฺสน\-ญาณํ นนฺทิยา นิมิตฺตา มุจฺจตีติ อนิมิตฺโต วิโมกฺโข. วิราคานุปสฺสน\-ญาณํ ราคโต นิมิตฺตา มุจฺจตีติ อนิมิตฺโต วิโมกฺโข. นิโรธานุปสฺสน\-ญาณํ สมุทยโต นิมิตฺตา มุจฺจตีติ อนิมิตฺโต วิโมกฺโข. ปฏินิสฺสคฺคานุปสฺสน\-ญาณํ อาทานโต นิมิตฺตา มุจฺจตีติ อนิมิตฺโต วิโมกฺโข. อนิมิตฺตา\-นุปสฺสน\-ญาณํ สพฺพนิมิตฺเตหิ มุจฺจตีติ อนิมิตฺโต วิโมกฺโข. อปฺปณิหิตา\-นุปสฺสน\-ญาณํ ปณิธิยา นิมิตฺตา มุจฺจตีติ อนิมิตฺโต วิโมกฺโข. สุญฺญตานุปสฺสน\-ญาณํ อภินิเวสโต นิมิตฺตา มุจฺจตีติ อนิมิตฺโต วิโมกฺโข.}

\prose{รูเป อนิจฺจานุปสฺสน\-ญาณํ นิจฺจโต นิมิตฺตา มุจฺจตีติ อนิมิตฺโต วิโมกฺโข ฯเปฯ รูเป อนิมิตฺตา\-นุปสฺสน\-ญาณํ สพฺพนิมิตฺเตหิ มุจฺจตีติ อนิมิตฺโต วิโมกฺโข. รูเป อปฺปณิหิตา\-นุปสฺสน\-ญาณํ ปณิธิยา นิมิตฺตา มุจฺจตีติ อนิมิตฺโต วิโมกฺโข. รูเป สุญฺญตานุปสฺสน\-ญาณํ อภินิเวสโต นิมิตฺตา มุจฺจตีติ อนิมิตฺโต วิโมกฺโข.  เวทนาย ฯเปฯ. สญฺญาย. สงฺขาเรสุ. วิญฺญาเณ. จกฺขุสฺมิํ ฯเปฯ. ชรามรเณ อนิจฺจานุปสฺสน\-ญาณํ นิจฺจโต นิมิตฺตา มุจฺจตีติ อนิมิตฺโต วิโมกฺโข ฯเปฯ ชรามรเณ อนิจฺจานุปสฺสน\-ญาณํ สพฺพนิมิตฺเตหิ มุจฺจตีติ อนิมิตฺโต วิโมกฺโข. ชรามรเณ อปฺปณิหิตา\-นุปสฺสน\-ญาณํ ปณิธิยา นิมิตฺตา มุจฺจตีติ อนิมิตฺโต วิโมกฺโข. ชรามรเณ สุญฺญตานุปสฺสน\-ญาณํ อภินิเวสโต นิมิตฺตา มุจฺจตีติ อนิมิตฺโต วิโมกฺโข. อยํ อนิมิตฺโต วิโมกฺโข.}

\prose{กตโม อปฺปณิหิโต วิโมกฺโข? อนิจฺจานุปสฺสน\-ญาณํ นิจฺจโต ปณิธิยา มุจฺจตีติ อปฺปณิหิโต วิโมกฺโข. ทุกฺขานุปสฺสน\-ญาณํ สุขโต ปณิธิ มุจฺจตีติ อปฺปณิหิโต วิโมกฺโข. อนตฺตานุปสฺสน\-ญาณํ อตฺตโต ปณิธิ มุจฺจตีติ อปฺปณิหิโต วิโมกฺโข. นิพฺพิทานุปสฺสน\-ญาณํ นนฺทิยา ปณิธิ มุจฺจตีติ อปฺปณิหิโต วิโมกฺโข.}

\prose{\csromanfolio{264}วิราคานุปสฺสน\-ญาณํ ราคโต ปณิธิ มุจฺจตีติ อปฺปณิหิโต วิโมกฺโข. นิโรธานุปสฺสน\-ญาณํ สมุทยโต ปณิธิ มุจฺจตีติ อปฺปณิหิโต วิโมกฺโข. ปฏินิสฺสคฺคานุปสฺสน\-ญาณํ อาทานโต ปณิธิ มุจฺจตีติ อปฺปณิหิโต วิโมกฺโข. อนิมิตฺตา\-นุปสฺสน\-ญาณํ นิมิตฺตโต ปณิธิ มุจฺจตีติ อปฺปณิหิโต วิโมกฺโข. อปฺปณิหิตา\-นุปสฺสน\-ญาณํ สพฺพปณิธีหิ มุจฺจตีติ อปฺปณิหิโต วิโมกฺโข. สุญฺญตานุปสฺสน\-ญาณํ อภินิเวสโต ปณิธิ มุจฺจตีติ อปฺปณิหิโต วิโมกฺโข.}

\prose{รูเป อนิจฺจานุปสฺสน\-ญาณํ นิจฺจโต ปณิธิ มุจฺจตีติ อปฺปณิหิโต วิโมกฺโข ฯเปฯ รูเป อปฺปณิหิตา\-นุปสฺสน\-ญาณํ สพฺพปณิธีหิ มุจฺจตีติ อปฺปณิหิโต วิโมกฺโข. รูเป สุญฺญตานุปสฺสน\-ญาณํ อภินิเวสโต ปณิธิ มุจฺจตีติ อปฺปณิหิโต วิโมกฺโข.  เวทนาย ฯเปฯ. สญฺญาย. สงฺขาเรสุ. วิญฺญาเณ. จกฺขุสฺมิํ ฯเปฯ. ชรามรเณ อนิจฺจานุปสฺสน\-ญาณํ นิจฺจโต ปณิธิ มุจฺจตีติ อปฺปณิหิโต วิโมกฺโข ฯเปฯ ชรามรเณ อปฺปณิหิตา\-นุปสฺสน\-ญาณํ สพฺพปณิธีหิ มุจฺจตีติ อปฺปณิหิโต วิโมกฺโข. ชรามรเณ สุญฺญตานุปสฺสน\-ญาณํ อภินิเวสโต ปณิธิ มุจฺจตีติ อปฺปณิหิโต วิโมกฺโข. อยํ อปฺปณิหิโต วิโมกฺโข. อยํ วิโมกฺโข. (1)}

\proseitem{230}{กตมํ \textbf{มุขํ?} เย ตตฺถ ชาตา อนวชฺชกุสลา โพธิปกฺขิยา ธมฺมา, อิทํ มุขํ. (2)}

\prose{กตมํ \textbf{วิโมกฺขมุขํ?} ยํ เตสํ ธมฺมานํ อารมฺมณํ นิโรโธ นิพฺพานํ, อิทํ วิโมกฺขมุขํ.  วิโมกฺขญฺจ มุขญฺจ วิโมกฺขมุขํ, อิทํ วิโมกฺขมุขํ. (3)}

\prose{กตมํ วิโมกฺข\-ปจฺจนีกํ? ตีณิ อกุสลมูลานิ วิโมกฺข\-ปจฺจนีกานิ, ตีณิ ทุจฺจริตานิ วิโมกฺข\-ปจฺจนีกานิ, สพฺเพปิ อกุสลา ธมฺมา วิโมกฺข\-ปจฺจนีกา, อิทํ วิโมกฺข\-ปจฺจนีกํ. (4)}

\prose{กตมํ วิโมกฺขา\-นุโลมํ? ตีณิ กุสลมูลานิ วิโมกฺขา\-นุโลมานิ, ตีณิ สุจริตานิ วิโมกฺขา\-นุโลมานิ, สพฺเพปิ กุสลา ธมฺมา วิโมกฺขานุโลมา, อิทํ วิโมกฺขา\-นุโลมํ. (5)}

\prose{\csromanfolio{265}กตโม วิโมกฺข\-วิวฏฺโฏ? สญฺญาวิวฏฺโฏ เจโตวิวฏฺโฏ จิตฺตวิวฏฺโฏ ญาณวิวฏฺโฏ วิโมกฺข\-วิวฏฺโฏ สจฺจวิวฏฺโฏ. สญฺชานนฺโต วิวฏฺฏตีติ สญฺญาวิวฏฺโฏ, เจตยนฺโต วิวฏฺฏตีติ เจโตวิวฏฺโฏ, วิชานนฺโต วิวฏฺฏตีติ จิตฺตวิวฏฺโฏ, ญาณํ กโรนฺโต วิวฏฺฏตีติ ญาณวิวฏฺโฏ, โวสชฺชนฺโต วิวฏฺฏตีติ วิโมกฺข\-วิวฏฺโฏ, ตถฏฺเฐ วิวฏฺฏตีติ สจฺจวิวฏฺโฏ.}

\prose{ยตฺถ สญฺญาวิวฏฺโฏ, ตตฺถ เจโตวิวฏฺโฏ. ยตฺถ เจโตวิวฏฺโฏ, ตตฺถ สญฺญาวิวฏฺโฏ. ยตฺถ สญฺญาวิวฏฺโฏ เจโตวิวฏฺโฏ, ตตฺถ จิตฺตวิวฏฺโฏ. ยตฺถ จิตฺตวิวฏฺโฏ, ตตฺถ สญฺญาวิวฏฺโฏ เจโตวิวฏฺโฏ. ยตฺถ สญฺญาวิวฏฺโฏ เจโตวิวฏฺโฏ จิตฺตวิวฏฺโฏ, ตตฺถ ญาณวิวฏฺโฏ. ยตฺถ ญาณวิวฏฺโฏ, ตตฺถ สญฺญาวิวฏฺโฏ เจโตวิวฏฺโฏ จิตฺตวิวฏฺโฏ. ยตฺถ สญฺญาวิวฏฺโฏ เจโตวิวฏฺโฏ จิตฺตวิวฏฺโฏ ญาณวิวฏฺโฏ, ตตฺถ วิโมกฺข\-วิวฏฺโฏ. ยตฺถ วิโมกฺข\-วิวฏฺโฏ, ตตฺถ สญฺญาวิวฏฺโฏ เจโตวิวฏฺโฏ จิตฺตวิวฏฺโฏ ญาณวิวฏฺโฏ. ยตฺถ สญฺญาวิวฏฺโฏ เจโตวิวฏฺโฏ จิตฺตวิวฏฺโฏ ญาณวิวฏฺโฏ วิโมกฺข\-วิวฏฺโฏ, ตตฺถ สจฺจวิวฏฺโฏ. ยตฺถ สจฺจวิวฏฺโฏ, ตตฺถ สญฺญาวิวฏฺโฏ เจโตวิวฏฺโฏ จิตฺตวิวฏฺโฏ ญาณวิวฏฺโฏ วิโมกฺข\-วิวฏฺโฏ. อยํ วิโมกฺข\-วิวฏฺโฏ. (6)}

\prose{กตมา วิโมกฺข\-ภาวนา? ปฐมสฺส ฌานสฺส อาเสวนา ภาวนา พหุลีกมฺมํ, ทุติยสฺส ฌานสฺส อาเสวนา ภาวนา พหุลีกมฺมํ, ตติยสฺส ฌานสฺส อาเสวนา ภาวนา พหุลีกมฺมํ, จตุตฺถสฺส ฌานสฺส อาเสวนา ภาวนา พหุลีกมฺมํ, อากาสา\-นญฺจา\-ยตนสมาปตฺติยา อาเสวนา ภาวนา พหุลีกมฺมํ, วิญฺญาณญฺจายตน\-สมาปตฺติยา ฯเปฯ อากิญฺจญฺญายตน\-สมาปตฺติยา. เนว\-สญฺญา\-นา\-สญฺญา\-ยตนสมาปตฺติยา อาเสวนา ภาวนา พหุลีกมฺมํ, โสตาปตฺติ\-มคฺคสฺส อาเสวนา ภาวนา พหุลีกมฺมํ, สกทาคามิ\-มคฺคสฺส อาเสวนา ภาวนา พหุลีกมฺมํ, อนาคามิ\-มคฺคสฺส อาเสวนา ภาวนา พหุลีกมฺมํ, อรหตฺต\-มคฺคสฺส อาเสวนา ภาวนา พหุลีกมฺมํ. อยํ วิโมกฺข\-ภาวนา. (7)}

\prose{กตมา วิโมกฺข\-ปฺปฏิปฺปสฺสทฺธิ? ปฐมสฺส ฌานสฺส ปฏิลาโภ วา วิปาโก วา, ทุติยสฺส ฌานสฺส ปฏิลาโภ วา วิปาโก วา, ตติยสฺส ฌานสฺส ฯเปฯ จตุตฺถสฺส ฌานสฺส. อากาสา\-นญฺจา\-ยตนสมาปตฺติยา. วิญฺญาณญฺจายตน\-สมาปตฺติยา. อากิญฺจญฺญายตน\-สมาปตฺติยา. เนว\-สญฺญา\-นา\-สญฺญา\-ยตนสมาปตฺติยา ปฏิลาโภ วา}

\prosewithcloser{\csromanfolio{266}วิปาโก วา, โสตาปตฺติ\-มคฺคสฺส โสตาปตฺติผลํ, สกทาคามิ\-มคฺคสฺส สกทาคามิผลํ, อนาคามิ\-มคฺคสฺส อนาคามิผลํ, อรหตฺต\-มคฺคสฺส อรหตฺตผลํ. อยํ วิโมกฺข\-ปฺปฏิปฺปสฺสทฺธีติ. (8)}{วิโมกฺขกถา นิฏฺฐิตา.}
\nitthitamruled{ตติย\-ภาณวาโร.}
\csromanmatikaanchor{mtk.79}
\csromantocmarkref{section}{6. คติกถา}{mtk.79}
\bookmark[dest={mtk.79},level=1]{6. คติกถา}
\csromanheader{1}{6. \textbf{คติกถา}}
\proseitem{231}{คติสมฺปตฺติยา ญาณสมฺปยุตฺเต กตินํ เหตูนํ ปจฺจยา อุปปตฺติ โหติ. ขตฺติย\-มหาสาลานํ พฺราหฺมณ\-มหาสาลานํ คหปติ\-มหาสาลานํ กามาวจรานํ เทวานํ ญาณสมฺปยุตฺเต กตินํ เหตูนํ ปจฺจยา อุปปตฺติ โหติ. รูปาวจรานํ เทวานํ กตินํ เหตูนํ ปจฺจยา อุปปตฺติ โหติ. อรูปาวจรานํ เทวานํ กตินํ เหตูนํ ปจฺจยา อุปปตฺติ โหติ?}

\prose{คติสมฺปตฺติยา ญาณสมฺปยุตฺเต อฏฺฐนฺนํ เหตูนํ ปจฺจยา อุปปตฺติ โหติ. ขตฺติย\-มหาสาลานํ พฺราหฺมณ\-มหาสาลานํ คหปติ\-มหาสาลานํ กามาวจรานํ เทวานํ ญาณสมฺปยุตฺเต อฏฺฐนฺนํ เหตูนํ ปจฺจยา อุปปตฺติ โหติ. รูปาวจรานํ เทวานํ อฏฺฐนฺนํ เหตูนํ ปจฺจยา อุปปตฺติ โหติ, อรูปาวจรานํ เทวานํ อฏฺฐนฺนํ เหตูนํ ปจฺจยา อุปปตฺติ โหติ.}

\proseitem{232}{คติสมฺปตฺติยา ญาณสมฺปยุตฺเต กตเมสํ อฏฺฐนฺนํ เหตูนํ ปจฺจยา อุปปตฺติ โหติ? กุสลกมฺมสฺส ชวนกฺขเณ ตโย เหตู กุสลา ตสฺมิํ ขเณ ชาตเจตนาย สหชาต\-ปจฺจยา โหนฺติ, เตน วุจฺจติ กุสลมูล\-ปจฺจยาปิ สงฺขารา. นิกนฺติกฺขเณ เทฺว เหตู อกุสลา ตสฺมิํ ขเณ ชาตเจตนาย สหชาต\-ปจฺจยา โหนฺติ, เตน วุจฺจติ อกุสลมูล\-ปจฺจยาปิ สงฺขารา. ปฏิสนฺธิ\-กฺขเณ ตโย เหตู อพฺยากตา ตสฺมิํ ขเณ ชาตเจตนาย สหชาต\-ปจฺจยา โหนฺติ, เตน วุจฺจติ นามรูป\-ปจฺจยาปิ วิญฺญาณํ, วิญฺญาณ\-ปจฺจยาปิ นามรูปํ.}

\prose{ปฏิสนฺธิ\-กฺขเณ ปญฺจกฺขนฺธา สหชาต\-ปจฺจยา โหนฺติ, อญฺญมญฺญ\-ปจฺจยา โหนฺติ, นิสฺสยปจฺจยา โหนฺติ, วิปฺปยุตฺต\-ปจฺจยา โหนฺติ. ปฏิสนฺธิ\-กฺขเณ}

\csromanheader{2}{6. คติกถา}
\prose{\csromanfolio{267}จตฺตาโร มหาภูตา สหชาต\-ปจฺจยา โหนฺติ, อญฺญมญฺญ\-ปจฺจยา โหนฺติ, นิสฺสยปจฺจยา โหนฺติ. ปฏิสนฺธิ\-กฺขเณ ตโย ชีวิตสงฺขารา สหชาต\-ปจฺจยา โหนฺติ, อญฺญมญฺญ\-ปจฺจยา โหนฺติ, นิสฺสยปจฺจยา โหนฺติ, วิปฺปยุตฺต\-ปจฺจยา โหนฺติ. ปฏิสนฺธิ\-กฺขเณ นามญฺจ รูปญฺจ สหชาต\-ปจฺจยา โหนฺติ, อญฺญมญฺญ\-ปจฺจยา โหนฺติ, นิสฺสยปจฺจยา โหนฺติ, วิปฺปยุตฺต\-ปจฺจยา โหนฺติ. ปฏิสนฺธิ\-กฺขเณ อิเม จุทฺทส ธมฺมา สหชาต\-ปจฺจยา โหนฺติ, อญฺญมญฺญ\-ปจฺจยา โหนฺติ, นิสฺสยปจฺจยา โหนฺติ, วิปฺปยุตฺต\-ปจฺจยา โหนฺติ. ปฏิสนฺธิ\-กฺขเณ จตฺตาโร ขนฺธา อรูปิโน สหชาต\-ปจฺจยา โหนฺติ, อญฺญมญฺญ\-ปจฺจยา โหนฺติ, นิสฺสยปจฺจยา โหนฺติ, สมฺปยุตฺต\-ปจฺจยา โหนฺติ. ปฏิสนฺธิ\-กฺขเณ ปญฺจินฺทฺริยานิ สหชาต\-ปจฺจยา โหนฺติ, อญฺญมญฺญ\-ปจฺจยา โหนฺติ, นิสฺสยปฺปจฺจยา โหนฺติ, สมฺปยุตฺต\-ปจฺจยา โหนฺติ. ปฏิสนฺธิ\-กฺขเณ ตโย เหตู สหชาต\-ปจฺจยา โหนฺติ, อญฺญมญฺญ\-ปจฺจยา โหนฺติ, นิสฺสยปจฺจยา โหนฺติ, สมฺปยุตฺต\-ปจฺจยา โหนฺติ. ปฏิสนฺธิ\-กฺขเณ นามญฺจ วิญฺญาณญฺจ สหชาต\-ปจฺจยา โหนฺติ, อญฺญมญฺญ\-ปจฺจยา โหนฺติ, นิสฺสยปจฺจยา โหนฺติ, สมฺปยุตฺต\-ปจฺจยา โหนฺติ. ปฏิสนฺธิ\-กฺขเณ อิเม จุทฺทส ธมฺมา สหชาต\-ปจฺจยา โหนฺติ, อญฺญมญฺญ\-ปจฺจยา โหนฺติ, นิสฺสยปจฺจยา โหนฺติ, สมฺปยุตฺต\-ปจฺจยา โหนฺติ. ปฏิสนฺธิ\-กฺขเณ อิเม อฏฺฐวีสติ ธมฺมา สหชาต\-ปจฺจยา โหนฺติ, อญฺญมญฺญ\-ปจฺจยา โหนฺติ, นิสฺสยปจฺจยา โหนฺติ, วิปฺปยุตฺต\-ปจฺจยา โหนฺติ. คติสมฺปตฺติยา ญาณสมฺปยุตฺเต อิเมสํ อฏฺฐนฺนํ เหตูนฺนํ ปจฺจยา อุปปตฺติ โหติ. (1)}

\prose{ขตฺติย\-มหาสาลานํ พฺราหฺมณ\-มหาสาลานํ คหปติ\-มหาสาลานํ กามาวจรานํ เทวานํ ญาณสมฺปยุตฺเต กตเมสํ อฏฺฐนฺนํ เหตูนํ ปจฺจยา อุปปตฺติ โหติ? กุสลกมฺมสฺส ชวนกฺขเณ ตโย เหตู กุสลา ตสฺมิํ ขเณ ชาตเจตนาย สหชาต\-ปจฺจยา โหนฺติ, เตน วุจฺจติ กุสลมูล\-ปจฺจยาปิ สงฺขารา. นิกนฺติกฺขเณ เทฺว เหตู อกุสลา ตสฺมิํ ขเณ ชาตเจตนาย สหชาต\-ปจฺจยา โหนฺติ, เตน วุจฺจติ อกุสลมูล\-ปจฺจยาปิ สงฺขารา. ปฏิสนฺธิ\-กฺขเณ ตโย เหตู อพฺยากตา ตสฺมิํ ขเณ ชาตเจตนาย สหชาต\-ปจฺจยา โหนฺติ, เตน วุจฺจติ นามรูป\-ปจฺจยาปิ วิญฺญาณํ, วิญฺญาณ\-ปจฺจยาปิ นามรูปํ.}

\prose{\csromanfolio{268}ปฏิสนฺธิ\-กฺขเณ ปญฺจกฺขนฺธา สหชาต\-ปจฺจยา โหนฺติ, อญฺญมญฺญ\-ปจฺจยา โหนฺติ, นิสฺสยปจฺจยา โหนฺติ, วิปฺปยุตฺต\-ปจฺจยา โหนฺติ, ปฏิสนฺธิ\-กฺขเณ จตฺตาโร มหาภูตา สหชาต\-ปจฺจยา โหนฺติ. อญฺญมญฺญ\-ปจฺจยา โหนฺติ, นิสฺสยปจฺจยา โหนฺติ. ปฏิสนฺธิ\-กฺขเณ ตโย ชีวิตสงฺขารา สหชาต\-ปจฺจยา โหนฺติ, อญฺญมญฺญ\-ปจฺจยา โหนฺติ, นิสฺสยปจฺจยา โหนฺติ, วิปฺปยุตฺต\-ปจฺจยา โหนฺติ. ปฏิสนฺธิ\-กฺขเณ นามญฺจ รูปญฺจ สหชาต\-ปจฺจยา โหนฺติ, อญฺญมญฺญ\-ปจฺจยา โหนฺติ, นิสฺสยปจฺจยา โหนฺติ, วิปฺปยุตฺต\-ปจฺจยา โหนฺติ. ปฏิสนฺธิ\-กฺขเณ อิเม จุทฺทส ธมฺมา สหชาต\-ปจฺจยา โหนฺติ, อญฺญมญฺญ\-ปจฺจยา โหนฺติ, นิสฺสยปจฺจยา โหนฺติ, วิปฺปยุตฺต\-ปจฺจยา โหนฺติ. ปฏิสนฺธิ\-กฺขเณ จตฺตาโร ขนฺธา อรูปิโน สหชาต\-ปจฺจยา โหนฺติ, อญฺญมญฺญ\-ปจฺจยา โหนฺติ, นิสฺสยปจฺจยา โหนฺติ, สมฺปยุตฺต\-ปจฺจยา โหนฺติ. ปฏิสนฺธิ\-กฺขเณ ปญฺจินฺทฺริยานิ สหชาต\-ปจฺจยา โหนฺติ, อญฺญมญฺญ\-ปจฺจยา โหนฺติ, นิสฺสยปจฺจยา โหนฺติ, สมฺปยุตฺต\-ปจฺจยา โหนฺติ. ปฏิสนฺธิ\-กฺขเณ ตโย เหตู สหชาต\-ปจฺจยา โหนฺติ, อญฺญมญฺญ\-ปจฺจยา โหนฺติ, นิสฺสยปจฺจยา โหนฺติ, สมฺปยุตฺต\-ปจฺจยา โหนฺติ. ปฏิสนฺธิ\-กฺขเณ นามญฺจ วิญฺญาณญฺจ สหชาต\-ปจฺจยา โหนฺติ, อญฺญมญฺญ\-ปจฺจยา โหนฺติ, นิสฺสยปจฺจยา โหนฺติ, สมฺปยุตฺต\-ปจฺจยา โหนฺติ. ปฏิสนฺธิ\-กฺขเณ อิเม จุทฺทส ธมฺมา สหชาต\-ปจฺจยา โหนฺติ, อญฺญมญฺญ\-ปจฺจยา โหนฺติ, นิสฺสยปจฺจยา โหนฺติ, สมฺปยุตฺต\-ปจฺจยา โหนฺติ. ปฏิสนฺธิ\-กฺขเณ อิเม อฏฺฐวีสติ ธมฺมา สหชาต\-ปจฺจยา โหนฺติ, อญฺญมญฺญ\-ปจฺจยา โหนฺติ, นิสฺสยปจฺจยา โหนฺติ, วิปฺปยุตฺต\-ปจฺจยา โหนฺติ, ขตฺติย\-มหาสาลานํ พฺรหฺมณมหาสาลานํ คหปติ\-มหาสาลานํ กามาวจรานํ เทวานํ ญาณสมฺปยุตฺเต อิเมสํ อฏฺฐนฺนํ เหตูนํ ปจฺจยา อุปปตฺติ โหติ. (2)}

\prose{รูปาวจรานํ เทวานํ กตเมสํ อฏฺฐนฺนํ เหตูนํ ปจฺจยา อุปปตฺติ โหติ? กุสลกมฺมสฺส ชวนกฺขเณ ตโย เหตู กุสลา ฯเปฯ รูปาวจรานํ เทวานํ อิเมสํ อฏฺฐนฺนํ เหตูนํ ปจฺจยา อุปปตฺติ โหติ. (3)}

\prose{อรูปาวจรานํ เทวานํ กตเมสํ อฏฺฐนฺนํ เหตูนํ ปจฺจยา อุปปตฺติ โหติ? กุสลกมฺมสฺส ชวนกฺขเณ ตโย เหตู กุสลา ตสฺมิํ ขเณ ชาตเจตนาย สหชาต\-ปจฺจยา โหนฺติ, เตน วุจฺจติ กุสลมูล\-ปจฺจยาปิ สงฺขารา. นิกนฺติกฺขเณ เทฺว เหตู อกุสลา ตสฺมิํ ขเณ}

\csromanheader{2}{6. คติกถา}
\prose{\csromanfolio{269}ชาตเจตนาย สหชาต\-ปจฺจยา โหนฺติ, เตน วุจฺจติ อกุสลมูล\-ปจฺจยาปิ สงฺขารา. ปฏิสนฺธิ\-กฺขเณ ตโย เหตู อพฺยากตา ตสฺมิํ ขเณ ชาตเจตนาย สหชาต\-ปจฺจยา โหนฺติ, เตน วุจฺจติ นามปจฺจยาปิ วิญฺญาณํ วิญฺญาณ\-ปจฺจยาปิ นามํ.}

\prose{ปฏิสนฺธิ\-กฺขเณ จตฺตาโร ขนฺธา อรูปิโน สหชาต\-ปจฺจยา โหนฺติ, อญฺญมญฺญ\-ปจฺจยา โหนฺติ, นิสฺสยปจฺจยา โหนฺติ, สมฺปยุตฺต\-ปจฺจยา โหนฺติ. ปฏิสนฺธิ\-กฺขเณ ปญฺจินฺทฺริยานิ สหชาต\-ปจฺจยา โหนฺติ, อญฺญมญฺญ\-ปจฺจยา โหนฺติ, นิสฺสยปจฺจยา โหนฺติ, สมฺปยุตฺต\-ปจฺจยา โหนฺติ. ปฏิสนฺธิ\-กฺขเณ ตโย เหตู สหชาต\-ปจฺจยา โหนฺติ, อญฺญมญฺญ\-ปจฺจยา โหนฺติ, นิสฺสยปจฺจยา โหนฺติ, สมฺปยุตฺต\-ปจฺจยา โหนฺติ. ปฏิสนฺธิ\-กฺขเณ นามญฺจ วิญฺญาณญฺจ สหชาต\-ปจฺจยา โหนฺติ, อญฺญมญฺญ\-ปจฺจยา โหนฺติ, นิสฺสยปจฺจยา โหนฺติ, สมฺปยุตฺต\-ปจฺจยา โหนฺติ. ปฏิสนฺธิ\-กฺขเณ อิเม จุทฺทส ธมฺมา สหชาต\-ปจฺจยา โหนฺติ, อญฺญมญฺญ\-ปจฺจยา โหนฺติ, นิสฺสยปจฺจยา โหนฺติ, สมฺปยุตฺต\-ปจฺจยา โหนฺติ, อรูปาวจรานํ เทวานํ อิเมสํ อฏฺฐนฺนํ เหตูนํ ปจฺจยา อุปปตฺติ โหติ. (4)}

\proseitem{233}{คติสมฺปตฺติยา ญาณวิปฺปยุตฺเต กตินํ เหตูนํ ปจฺจยา อุปปตฺติ โหติ. ขตฺติย\-มหาสาลานํ พฺราหฺมณ\-มหาสาลานํ คหปติ\-มหาสาลานํ กามาวจรานํ เทวานํ ญาณวิปฺปยุตฺเต กตินํ เหตูนํ ปจฺจยา อุปปตฺติ โหติ?}

\prose{คติสมฺปตฺติยา ญาณวิปฺปยุตฺเต ฉนฺนํ เหตูนํ ปจฺจยา อุปปตฺติ โหติ. ขตฺติย\-มหาสาลานํ พฺราหฺมณ\-มหาสาลานํ คหปติ\-มหาสาลานํ กามาวจรานํ เทวานํ ญาณวิปฺปยุตฺเต ฉนฺนํ เหตูนํ ปจฺจยา อุปปตฺติ โหติ.}

\prose{คติสมฺปตฺติยา ญาณวิปฺปยุตฺเต กตเมสํ ฉนฺนํ เหตูนํ ปจฺจยา อุปปตฺติ โหติ? กุสลกมฺมสฺส ชวนกฺขเณ เทฺว เหตู กุสลา ตสฺมิํ ขเณ ชาตเจตนาย สหชาต\-ปจฺจยา โหนฺติ, เตน วุจฺจติ กุสลมูล\-ปจฺจยาปิ สงฺขารา. นิกนฺติกฺขเณ เทฺว เหตู อกุสลา ตสฺมิํ ขเณ ชาตเจตนาย สหชาต\-ปจฺจยา โหนฺติ, เตน วุจฺจติ อกุสลมูล\-ปจฺจยาปิ สงฺขารา. ปฏิสนฺธิ\-กฺขเณ เทฺว เหตู อพฺยากตา ตสฺมิํ ขเณ \csromanfolio{270}ชาตเจตนาย สหชาต\-ปจฺจยา โหนฺติ, เตน วุจฺจติ นามรูป\-ปจฺจยาปิ วิญฺญาณํ, วิญฺญาณ\-ปจฺจยาปิ นามรูปํ.}

\prose{ปฏิสนฺธิ\-กฺขเณ ปญฺจกฺขนฺธา สหชาต\-ปจฺจยา โหนฺติ, อญฺญมญฺญ\-ปจฺจยา โหนฺติ, นิสฺสยปจฺจยา โหนฺติ, วิปฺปยุตฺต\-ปจฺจยา โหนฺติ. ปฏิสนฺธิ\-กฺขเณ จตฺตาโร มหาภูตา สหชาต\-ปจฺจยา โหนฺติ, อญฺญมญฺญ\-ปจฺจยา โหนฺติ, นิสฺสยปจฺจยา โหนฺติ. ปฏิสนฺธิ\-กฺขเณ ตโย ชีวิตสงฺขารา สหชาต\-ปจฺจยา โหนฺติ, อญฺญมญฺญ\-ปจฺจยา โหนฺติ, นิสฺสยปจฺจยา โหนฺติ, วิปฺปยุตฺต\-ปจฺจยา โหนฺติ. ปฏิสนฺธิ\-กฺขเณ นามญฺจ รูปญฺจ สหชาต\-ปจฺจยา โหนฺติ, อญฺญมญฺญ\-ปจฺจยา โหนฺติ, นิสฺสยปจฺจยา โหนฺติ, วิปฺปยุตฺต\-ปจฺจยา โหนฺติ. ปฏิสนฺธิ\-กฺขเณ อิเม จุทฺทส ธมฺมา สหชาต\-ปจฺจยา โหนฺติ, อญฺญมญฺญ\-ปจฺจยา โหนฺติ, นิสฺสยปจฺจยา โหนฺติ, วิปฺปยุตฺต\-ปจฺจยา โหนฺติ. ปฏิสนฺธิ\-กฺขเณ จตฺตาโร ขนฺธา อรูปิโน สหชาต\-ปจฺจยา โหนฺติ, อญฺญมญฺญ\-ปจฺจยา โหนฺติ, นิสฺสยปจฺจยา โหนฺติ, สมฺปยุตฺต\-ปจฺจยา โหนฺติ. ปฏิสนฺธิ\-กฺขเณ จตฺตาริ อินฺทฺริยานิ สหชาต\-ปจฺจยา โหนฺติ, อญฺญมญฺญ\-ปจฺจยา โหนฺติ, นิสฺสยปจฺจยา โหนฺติ, สมฺปยุตฺต\-ปจฺจยา โหนฺติ. ปฏิสนฺธิ\-กฺขเณ เทฺว เหตู สหชาต\-ปจฺจยา โหนฺติ, อญฺญมญฺญ\-ปจฺจยา โหนฺติ, นิสฺสยปจฺจยา โหนฺติ, สมฺปยุตฺต\-ปจฺจยา โหนฺติ. ปฏิสนฺธิ\-กฺขเณ นามญฺจ วิญฺญาณญฺจ สหชาต\-ปจฺจยา โหนฺติ, อญฺญมญฺญ\-ปจฺจยา โหนฺติ, นิสฺสยปจฺจยา โหนฺติ, สมฺปยุตฺต\-ปจฺจยา โหนฺติ. ปฏิสนฺธิ\-กฺขเณ อิเม ทฺวาทส ธมฺมา สหชาต\-ปจฺจยา โหนฺติ, อญฺญมญฺญ\-ปจฺจยา โหนฺติ, นิสฺสยปจฺจยา โหนฺติ, สมฺปยุตฺต\-ปจฺจยา โหนฺติ. ปฏิสนฺธิ\-กฺขเณ อิเม ฉพฺพีสติ ธมฺมา สหชาต\-ปจฺจยา โหนฺติ, อญฺญมญฺญ\-ปจฺจยา โหนฺติ, นิสฺสยปจฺจยา โหนฺติ, วิปฺปยุตฺต\-ปจฺจยา โหนฺติ, คติสมฺปตฺติยา ญาณวิปฺปยุตฺเต อิเมสํ ฉนฺนํ เหตูนํ ปจฺจยา อุปปตฺติ โหติ. (1)}

\prose{ขตฺติย\-มหาสาลานํ พฺราหฺมณ\-มหาสาลานํ คหปติ\-มหาสาลานํ กามาวจรานํ เทวานํ ญาณวิปฺปยุตฺเต กตเมสํ ฉนฺนํ เหตูนํ ปจฺจยา อุปปตฺติ โหติ? กุสลกมฺมสฺส ชวนกฺขเณ เทฺว เหตู กุสลา ตสฺมิํ ขเณ ชาตเจตนาย สหชาต\-ปจฺจยา โหนฺติ, เตน วุจฺจติ กุสลมูล\-ปจฺจยาปิ สงฺขารา ฯเปฯ ขตฺติย\-มหาสาลานํ พฺราหฺมณ\-มหาสาลานํ}

\csromanmatikaanchor{mtk.80}
\csromantocmarkref{section}{7. กมฺมกถา}{mtk.80}
\bookmark[dest={mtk.80},level=1]{7. กมฺมกถา}
\csromanheader{1}{7. กมฺมกถา}
\prosewithcloserruled{\csromanfolio{271}คหปติ\-มหาสาลานํ กามาวจรานํ เทวานํ ญาณวิปฺปยุตฺเต อิเมสํ ฉนฺนํ เหตูนํ ปจฺจยา อุปปตฺติ โหตีติ. (2)}{คติกถา นิฏฺฐิตา.}
\csromanheader{2}{7. กมฺมกถา}
\proseitem{234}{อโหสิ กมฺมํ อโหสิ กมฺมวิปาโก. อโหสิ กมฺมํ นาโหสิ กมฺมวิปาโก. อโหสิ กมฺมํ อตฺถิ กมฺมวิปาโก. อโหสิ กมฺมํ นตฺถิ กมฺมวิปาโก. อโหสิ กมฺมํ ภวิสฺสติ กมฺมวิปาโก. อโหสิ กมฺมํ น ภวิสฺสติ กมฺมวิปาโก. (อตีตกมฺมํ) (6)}

\prose{อตฺถิ กมฺมํ อตฺถิ กมฺมวิปาโก. อตฺถิ กมฺมํ นตฺถิ กมฺมวิปาโก. อตฺถิ กมฺมํ ภวิสฺสติ กมฺมวิปาโก. อตฺถิ กมฺมํ น ภวิสฺสติ กมฺมวิปาโก. (ปจฺจุปฺปนฺน\-กมฺมํ) \mbox{(4, 10)}}

\prose{ภวิสฺสติ กมฺมํ ภวิสฺสติ กมฺมวิปาโก. ภวิสฺสติ กมฺมํ น ภวิสฺสติ กมฺมวิปาโก. (อนาคตกมฺมํ) \mbox{(2, 12)}}

\proseitem{235}{อโหติ กุสลํ กมฺมํ อโหสิ กุสลสฺส กมฺมสฺส วิปาโก, อโหสิ กุสลํ กมฺมํ นาโหสิ กุสลสฺส กมฺมสฺส วิปาโก, อโหสิ กุสลํ กมฺมํ อตฺถิ กุสลสฺส กมฺมสฺส วิปาโก, อโหสิ กุสลํ กมฺมํ นตฺถิ กุสลสฺส กมฺมสฺส วิปาโก, อโหสิ กุสลํ กมฺมํ ภวิสฺสติ กุสลสฺส กมฺมสฺส วิปาโก, อโหสิ กุสลํ กมฺมํ น ภวิสฺสติ กุสลสฺส กมฺมสฺส วิปาโก.}

\prose{อตฺถิ กุสลํ กมฺมํ อตฺถิ กุสลสฺส กมฺมสฺส วิปาโก, อตฺถิ กุสลํ กมฺมํ นตฺถิ กุสลสฺส กมฺมสฺส วิปาโก, อตฺถิ กุสลํ กมฺมํ ภวิสฺสติ กุสลสฺส กมฺมสฺส วิปาโก, อตฺถิ กุสลํ กมฺมํ น ภวิสฺสติ กุสลสฺส กมฺมสฺส วิปาโก.}

\prose{ภวิสฺสติ กุสลํ กมฺมํ ภวิสฺสติ กุสลสฺส กมฺมสฺส วิปาโก, ภวิสฺสติ กุสลํ กมฺมํ น ภวิสฺสติ กุสลสฺส กมฺมสฺส วิปาโก.}

\prose{อโหสิ อกุสลํ กมฺมํ อโหสิ อกุสลสฺส กมฺมสฺส วิปาโก, อโหสิ อกุสลํ กมฺมํ นาโหสิ อกุสลสฺส กมฺมสฺส วิปาโก, \csromanfolio{272}อโหสิ อกุสลํ กมฺมํ อตฺถิ อกุสลสฺส กมฺมสฺส วิปาโก, อโหสิ อกุสลํ กมฺมํ นตฺถิ อกุสลสฺส กมฺมสฺส วิปาโก. อโหสิ อกุสลํ กมฺมํ ภวิสฺสติ อกุสลสฺส กมฺมสฺส วิปาโก, อโหสิ อกุสลํ กมฺมํ น ภวิสฺสติ อกุสลสฺส กมฺมสฺส วิปาโก.}

\prose{อตฺถิ อกุสลํ กมฺมํ อตฺถิ อกุสลสฺส กมฺมสฺส วิปาโก, อตฺถิ อกุสลํ กมฺมํ นตฺถิ อกุสลสฺส กมฺมสฺส วิปาโก, อตฺถิ อกุสลํ กมฺมํ ภวิสฺสติ อกุสลสฺส กมฺมสฺส วิปาโก, อตฺถิ อกุสลํ กมฺมํ น ภวิสฺสติ อกุสลสฺส กมฺมสฺส วิปาโก.}

\prose{ภวิสฺสติ อกุสลํ กมฺมํ ภวิสฺสติ อกุสลสฺส กมฺมสฺส วิปาโก, ภวิสฺสติ อกุสลํ กมฺมํ น ภวิสฺสติ อกุสลสฺส กมฺมสฺส วิปาโก.}

\prose{อโหสิ สาวชฺชํ กมฺมํ ฯเปฯ. อโหสิ อนวชฺชํ กมฺมํ ฯเปฯ. อโหสิ กณฺหํ กมฺมํ ฯเปฯ. อโหสิ สุกฺกํ กมฺมํ ฯเปฯ. อโหสิ สุขุทฺรยํ กมฺมํ ฯเปฯ. อโหสิ ทุกฺขุทฺรยํ กมฺมํ ฯเปฯ. อโหสิ สุขวิปากํ กมฺมํ ฯเปฯ. อโหสิ ทุกฺขวิปากํ กมฺมํ อโหสิ ทุกฺข\-วิปากสฺส กมฺมสฺส วิปาโก, อโหสิ ทุกฺขวิปากํ กมฺมํ นาโหสิ ทุกฺข\-วิปากสฺส กมฺมสฺส วิปาโก, อโหสิ ทุกฺขวิปากํ กมฺมํ อตฺถิ ทุกฺข\-วิปากสฺส กมฺมสฺส วิปาโก, อโหสิ ทุกฺขวิปากํ กมฺมํ นตฺถิ ทุกฺข\-วิปากสฺส กมฺมสฺส วิปาโก, อโหสิ ทุกฺขวิปากํ กมฺมํ ภวิสฺสติ ทุกฺข\-วิปากสฺส กมฺมสฺส วิปาโก, อโหสิ ทุกฺขวิปากํ กมฺมํ น ภวิสฺสติ ทุกฺข\-วิปากสฺส กมฺมสฺส วิปาโก.}

\prose{อตฺถิ ทุกฺขวิปากํ กมฺมํ อตฺถิ ทุกฺข\-วิปากสฺส กมฺมสฺส วิปาโก, อตฺถิ ทุกฺขวิปากํ กมฺมํ นตฺถิ ทุกฺข\-วิปากสฺส กมฺมสฺส วิปาโก, อตฺถิ ทุกฺข วิปากํ กมฺมํ ภวิสฺสติ ทุกฺข\-วิปากสฺส กมฺมสฺส วิปาโก, อตฺถิ ทุกฺขวิปากํ กมฺมํ น ภวิสฺสติ ทุกฺข\-วิปากสฺส กมฺมสฺส วิปาโก.}

\prosewithcloser{ภวิสฺสติ ทุกฺขวิปากํ กมฺมํ ภวิสฺสติ ทุกฺข\-วิปากสฺส กมฺมสฺส วิปาโก, ภวิสฺสติ ทุกฺขวิปากํ กมฺมํ น ภวิสฺสติ ทุกฺข\-วิปากสฺส กมฺมสฺส วิปาโกติ.}{กมฺมกถา นิฏฺฐิตา.}
\csromanmatikaanchor{mtk.81}
\csromantocmarkref{section}{8. วิปลฺลาสกถา}{mtk.81}
\bookmark[dest={mtk.81},level=1]{8. วิปลฺลาสกถา}
\csromanheader{1}{8. วิปลฺลาสกถา}
\csromanheader{2}{8. วิปลฺลาสกถา}
\proseitem{236}{\csromanfolio{273}ปุริมนิทานํ.  จตฺตาโรเม ภิกฺขเว\csromansharedfootnote{954371fb65eaf87b}{อํ 1. 361 ปิฏฺเฐ ปสฺสิตพฺพา.} สญฺญาวิปลฺลาสา จิตฺตวิปลฺลาสา ทิฏฺฐิ\-วิปลฺลาสา. กตเม จตฺตาโร? อนิจฺเจ ภิกฺขเว “นิจฺจนฺ”ติ สญฺญาวิปลฺลาโส จิตฺตวิปลฺลาโส ทิฏฺฐิ\-วิปลฺลาโส, ทุกฺเข ภิกฺขเว “สุขนฺ”ติ สญฺญาวิปลฺลาโส จิตฺตวิปลฺลาโส ทิฏฺฐิ\-วิปลฺลาโส, อนตฺตนิ ภิกฺขเว “อตฺตา”ติ สญฺญาวิปลฺลาโส จิตฺตวิปลฺลาโส ทิฏฺฐิ\-วิปลฺลาโส, อสุเภ ภิกฺขเว “สุภนฺ”ติ สญฺญาวิปลฺลาโส จิตฺตวิปฺปลฺลาโส ทิฏฺฐิ\-วิปลฺลาโส, อิเม โข ภิกฺขเว จตฺตาโร สญฺญาวิปลฺลาสา จิตฺตวิปลฺลาสา ทิฏฺฐิ\-วิปลฺลาสา.}

\prose{จตฺตาโรเม ภิกฺขเว นสญฺญา\-วิปลฺลาสา นจิตฺต\-วิปลฺลาสา นทิฏฺฐิปลฺลาสา. กตเม จตฺตาโร? อนิจฺเจ ภิกฺขเว “อนิจฺจนฺ”ติ นสญฺญา\-วิปลฺลาโส นจิตฺต\-วิปลฺลาโส นทิฏฺฐิ\-วิปลฺลาโส, ทุกฺเข ภิกฺขเว “ทุกฺขนฺ”ติ นสญฺญา\-วิปลฺลาโส นจิตฺต\-วิปลฺลาโส นทิฏฺฐิ\-วิปลฺลาโส, อนตฺตนิ ภิกฺขเว “อนตฺตา”ติ นสญฺญา\-วิปลฺลาโส นจิตฺต\-วิปลฺลาโส นทิฏฺฐิ\-วิปลฺลาโส, อสุเภ ภิกฺขเว “อสุภนฺ”ติ นสญฺญา\-วิปลฺลาโส นจิตฺต\-วิปลฺลาโส นทิฏฺฐิ\-วิปลฺลาโส, อิเม โข ภิกฺขเว จตฺตโร นสญฺญา\-วิปลฺลาสา นจิตฺต\-วิปลฺลาสา นทิฏฺฐิ\-วิปลฺลาสาติ.}

\setlength{\csromangathapagewidth}{0pt}
\setlength{\csromangathawakwidth}{0pt}
\csromangathameasuregroup{อนิจฺเจ นิจฺจสญฺญิโน, \\ อนตฺตนิ จ “อตฺตา”ติ\textsuperscript{1}, \\ มิจฺฉาทิฏฺฐิหตา สตฺตา, \\ เต โยคยุตฺตา มารสฺส, \\ สตฺตา คจฺฉนฺติ สํสารํ, \\ ยทา จ พุทฺธา โลกสฺมิํ, \\ เต อิมํ ธมฺมํ ปกาเสนฺติ, \\ เตสํ สุตฺวาน สปฺปญฺญา, \\ อนิจฺจํ อนิจฺจโต ทกฺขุํ, \\ อนตฺตนิ “อนตฺตา”ติ, \\ สมฺมาทิฏฺฐิสมาทานา,}{\csromangathabat{อนิจฺเจ นิจฺจสญฺญิโน,}{ทุกฺเข จ สุขสญฺญิโน.} \\ \csromangathabat{อนตฺตนิ จ “อตฺตา”ติ\textsuperscript{1},}{อสุเภ สุภสญฺญิโน.} \\ \csromangathabat{มิจฺฉาทิฏฺฐิหตา สตฺตา,}{ขิตฺตจิตฺตา วิสญฺญิโน.} \\ \csromangathabat{เต โยคยุตฺตา มารสฺส,}{อโยคกฺเขมิโน ชนา.} \\ \csromangathabat{สตฺตา คจฺฉนฺติ สํสารํ,}{ชาติมรณคามิโน.} \\ \csromangathabat{ยทา จ พุทฺธา โลกสฺมิํ,}{อุปฺปชฺชนฺติ ปภงฺกรา.} \\ \csromangathabat{เต อิมํ ธมฺมํ ปกาเสนฺติ,}{ทุกฺขูปสมคามินํ.} \\ \csromangathabat{เตสํ สุตฺวาน สปฺปญฺญา,}{สจิตฺตํ ปจฺจลทฺธุ\textsuperscript{1} เต.} \\ \csromangathabat{อนิจฺจํ อนิจฺจโต ทกฺขุํ,}{ทุกฺขมทฺทกฺขุ ทุกฺขโต.} \\ \csromangathabat{อนตฺตนิ “อนตฺตา”ติ,}{อสุภํ อสุภตทฺทสุํ.} \\ \csromangathabat{สมฺมาทิฏฺฐิสมาทานา,}{สพฺพํ ทุกฺขํ อุปจฺจคุนฺติ.}}
\csromangathasetleft{อนิจฺเจ นิจฺจสญฺญิโน, \\ อนตฺตนิ จ “อตฺตา”ติ\textsuperscript{1}, \\ มิจฺฉาทิฏฺฐิหตา สตฺตา, \\ เต โยคยุตฺตา มารสฺส, \\ สตฺตา คจฺฉนฺติ สํสารํ, \\ ยทา จ พุทฺธา โลกสฺมิํ, \\ เต อิมํ ธมฺมํ ปกาเสนฺติ, \\ เตสํ สุตฺวาน สปฺปญฺญา, \\ อนิจฺจํ อนิจฺจโต ทกฺขุํ, \\ อนตฺตนิ “อนตฺตา”ติ, \\ สมฺมาทิฏฺฐิสมาทานา,}
\csromangathagroup{\csromangathastanza{\csromangathabat{อนิจฺเจ นิจฺจสญฺญิโน,}{ทุกฺเข จ สุขสญฺญิโน.} \\ \csromangathabat{อนตฺตนิ จ “อตฺตา”ติ\csromansharedfootnote{89ead441242ca62a}{อตฺตสญฺญิโน (สฺยา) อํ 1. 362 ปิฏฺเฐ ปสฺสิตพฺพา. 3. ปจฺจลทฺธุํ (สฺยา), ปจฺจลทฺธา (อํ 1. 362 ปิฏฺเฐ.)},}{อสุเภ สุภสญฺญิโน.}}\csromangathastanza{\csromangathabat{มิจฺฉาทิฏฺฐิหตา สตฺตา,}{ขิตฺตจิตฺตา วิสญฺญิโน.} \\ \csromangathabat{เต โยคยุตฺตา มารสฺส,}{อโยคกฺเขมิโน ชนา.}}\csromangathastanza{\csromangathabat{สตฺตา คจฺฉนฺติ สํสารํ,}{ชาติมรณคามิโน.} \\ \csromangathabat{ยทา จ พุทฺธา โลกสฺมิํ,}{อุปฺปชฺชนฺติ ปภงฺกรา.}}\csromangathastanza{\csromangathabat{เต อิมํ ธมฺมํ ปกาเสนฺติ,}{ทุกฺขู\-ปสม\-คามินํ.} \\ \csromangathabat{เตสํ สุตฺวาน สปฺปญฺญา,}{สจิตฺตํ ปจฺจลทฺธุ\csromansharedfootnote{2a174a2c956b9b08}{ปจฺจลทฺธุํ (สฺยา), ปจฺจลทฺธา (อํ 1. 362 ปิฏฺเฐ.)} เต.}}\csromangathastanza{\csromangathabat{อนิจฺจํ อนิจฺจโต ทกฺขุํ,}{ทุกฺขมทฺทกฺขุ ทุกฺขโต.} \\ \csromangathabat{อนตฺตนิ “อนตฺตา”ติ,}{อสุภํ อสุภตทฺทสุํ.} \\ \csromangathabat{สมฺมาทิฏฺฐิสมาทานา,}{สพฺพํ ทุกฺขํ อุปจฺจคุนฺติ.}}}

\prosewithcloserruled{\csromanfolio{274}อิเม จตฺตาโร วิปลฺลาสา ทิฏฺฐิ\-สมฺปนฺนสฺส ปุคฺคลสฺส ปหีนา อปฺปหีนาติ? เกจิ ปหีนา เกจิ อปฺปหีนา. อนิจฺเจ “นิจฺจนฺ”ติ สญฺญาวิปลฺลาโส จิตฺตวิปลฺลาโส ทิฏฺฐิ\-วิปลฺลาโส ปหีโน. ทุกฺเข “สุขนฺ”ติ สญฺญา อุปฺปชฺชติ, จิตฺตํ อุปฺปชฺชติ, ทิฏฺฐิ\-วิปลฺลาโส ปหีโน. อนตฺตนิ “อตฺตา”ติ สญฺญาวิปลฺลาโส จิตฺตวิปลฺลาโส ทิฏฺฐิ\-วิปลฺลาโส ปหีโน, อสุเภ “สุภนฺ”ติ สญฺญา อุปฺปชฺชติ, จิตฺตํ อุปฺปชฺชติ, ทิฏฺฐิ\-วิปลฺลาโส ปหีโน. ทฺวีสุ วตฺถูสุ ฉ วิปลฺลาสา ปหีนา, ทฺวีสุ วตฺถูสุ เทฺว วิปลฺลาสา ปหีนา, จตฺตาโร วิปลฺลาสา อปฺปหีนา, จตูสุ วตฺถูสุ อฏฺฐ วิปลฺลาสา ปหีนา, จตฺตาโร วิปลฺลาสา อปฺปหีนาติ.}{วิปลฺลาสกถา นิฏฺฐิตา.}
\csromanmatikaanchor{mtk.82}
\csromantocmarkref{section}{9. มคฺคกถา}{mtk.82}
\bookmark[dest={mtk.82},level=1]{9. มคฺคกถา}
\csromanheader{1}{9. \textbf{มคฺคกถา}}
\proseitem{237}{มคฺโคติ เกนฏฺเฐน มฺอคฺโค? โสตาปตฺติมคฺค\-กฺขเณ ทสฺสนฏฺเฐน สมฺมาทิฏฺฐิ มิจฺฉาทิฏฺฐิยา ปหานาย มคฺโค เจว เหตุ จ, สหชาตานํ ธมฺมานํ อุปตฺถมฺภนาย มคฺโค เจว เหตุ จ, กิเลสานํ ปริยาทานาย มคฺโค เจว เหตุ จ, ปฏิเวธา\-ทิวิ\-โสธนาย มคฺโค เจว เหตุ จ, จิตฺตสฺส อธิฏฺฐานาย มคฺโค เจว เหตุ จ, จิตฺตสฺส โวทานาย มคฺโค เจว เหตุ จ, วิเสสา\-ธิคมาย มคฺโค เจว เหตุ จ, อุตฺตริ\-ปฏิเวธาย\csromansharedfootnote{4824d6730c5138d8}{อุตฺตริํ ปฏิเวธาย (สฺยา)} มคฺโค เจว เหตุ จ, สจฺจา\-ภิสมยาย มคฺโค เจว เหตุ จ, นิโรเธ ปติฏฺฐาปนาย มคฺโค เจว เหตุ จ.}

\prose{อภินิโรปน\-ฏฺเฐน สมฺมาสงฺกปฺโป มิจฺฉา\-สงฺกปฺปสฺส ปหานาย มคฺโค เจว เหตุ จ, สหชาตานํ ธมฺมานํ อุปตฺถมฺภนาย มคฺโค เจว เหตุ จ, กิเลสานํ ปริยาทานาย มคฺโค เจว เหตุ จ, ปฏิเวธา\-ทิวิ\-โสธนาย มคฺโค เจว เหตุ จ, จิตฺตสฺส อธิฏฺฐานาย มคฺโค เจว เหตุ จ, จิตฺตสฺส โวทานาย มคฺโค เจว เหตุ จ, วิเสสา\-ธิคมาย มคฺโค เจว เหตุ จ,}

\csromanheader{2}{9. มคฺคกถา}
\prose{\csromanfolio{275}อุตฺตริ\-ปฏิเวธาย มคฺโค เจว เหตุ จ, สจฺจา\-ภิสมยาย มคฺโค เจว เหตุ จ, นิโรเธ ปติฏฺฐาปนาย มคฺโค เจว เหตุ จ.}

\prose{ปริคฺคห\-ฏฺเฐน สมฺมาวาจา มิจฺฉาวาจาย ปหานาย มคฺโค เจว เหตุ จ, สหชาตานํ ธมฺมานํ อุปตฺถมฺภนาย มคฺโค เจว เหตุ จ, กิเลสานํ ปริยาทานาย มคฺโค เจว เหตุ จ, ปฏิเวธา\-ทิวิ\-โสธนาย มคฺโค เจว เหตุ จ, จิตฺตสฺส อธิฏฺฐานาย มคฺโค เจว เหตุ จ, จิตฺตสฺส โวทานาย มคฺโค เจว เหตุ จ, วิเสสา\-ธิคมาย มคฺโค เจว เหตุ จ, อุตฺตริ\-ปฏิเวธาย มคฺโค เจว เหตุ จ, สจฺจา\-ภิสมยาย มคฺโค เจว เหตุ จ, นิโรเธ ปติฏฺฐาปนาย มคฺโค เจว เหตุ จ.}

\prose{สมุฏฺฐาน\-ฏฺเฐน สมฺมากมฺมนฺโต มิจฺฉา\-กมฺมนฺตสฺส ปหานาย มคฺโค เจว เหตุ จ, สหชาตานํ ธมฺมานํ อุปตฺถมฺภนาย มคฺโค เจว เหตุ จ, กิเลสานํ ปริยาทานาย มคฺโค เจว เหตุ จ, ปฏิเวธา\-ทิวิ\-โสธนาย มคฺโค เจว เหตุ จ, จิตฺตสฺส อธิฏฺฐานาย มคฺโค เจว เหตุ จ, จิตฺตสฺส โวทานาย มคฺโค เจว เหตุ จ. วิเสสา\-ธิคมาย มคฺโค เจว เหตุ จ, อุตฺตริ\-ปฏิเวธาย มคฺโค เจว เหตุ จ, สจฺจา\-ภิสมยาย มคฺโค เจว เหตุ จ, นิโรเธ ปติฏฺฐาปนาย มคฺโค เจว เหตุ จ.}

\prose{โวทานฏฺเฐน สมฺมาอาชีโว มิจฺฉาอาชีวสฺส ปหานาย มคฺโค เจว เหตุ จ ฯเปฯ. ปคฺคหฏฺเฐน สมฺมาวายาโม มิจฺฉา\-วายามสฺส ปหานาย มคฺโค เจว เหตุ จ ฯเปฯ. อุปฏฺฐาน\-ฏฺเฐน สมฺมาสติ มิจฺฉาสติยา ปหานาย มคฺโค เจว เหตุ จ ฯเปฯ. อวิกฺเขป\-ฏฺเฐน สมฺมาสมาธิ มิจฺฉา\-สมาธิสฺส ปหานาย มคฺโค เจว เหตุ จ, สหชาตานํ ธมฺมานํ อุปตฺถมฺภนาย มคฺโค เจว เหตุ จ, กิเลสานํ ปริยาทานาย มคฺโค เจว เหตุ จ, ปฏิเวธา\-ทิวิ\-โสธนาย มคฺโค เจว เหตุ จ, จิตฺตสฺส อธิฏฺฐานาย มคฺโค เจว เหตุ จ, จิตฺตสฺส โวทานาย มคฺโค เจว เหตุ จ, วิเสสา\-ธิคมาย มคฺโค เจว เหตุ จ, อุตฺตริ\-ปฏิเวธาย มคฺโค เจว เหตุ จ สจฺจา\-ภิสมยาย มคฺโค เจว เหตุ จ, นิโรเธ ปติฏฺฐาปนาย มคฺโค เจว เหตุ จ.}

\prose{สกทาคามิ\-มคฺคกฺขเณ ทสฺสนฏฺเฐน สมฺมาทิฏฺฐิ ฯเปฯ. อวิกฺเขป\-ฏฺเฐน สมฺมาสมาธิ โอฬาริกสฺส กามราค\-สญฺโญชนสฺส ปฏิฆ\-สญฺโญชนสฺส โอฬาริกสฺส กาม\-ราคา\-นุสยสฺส ปฏิฆา\-นุสยสฺส ปหานาย มคฺโค}

\prose{\csromanfolio{276}เจว เหตุ จ, สหชาตานํ ธมฺมานํ อุปตฺถมฺภนาย มคฺโค เจว เหตุ จ, กิเลสานํ ปริยาทานาย มคฺโค เจว เหตุ จ, ปฏิเวธา\-ทิวิ\-โสธนาย มคฺโค เจว เหตุ จ, จิตฺตสฺส อธิฏฺฐานาย มคฺโค เจว เหตุ จ, จิตฺตสฺส โวทานาย มคฺโค เจว เหตุ จ, วิเสสา\-ธิคมาย มคฺโค เจว เหตุ จ, อุตฺตริ\-ปฏิเวธาย มคฺโค เจว เหตุ จ, สจฺจา\-ภิสมยาย มคฺโค เจว เหตุ จ, นิโรเธ ปติฏฺฐาปนาย มคฺโค เจว เหตุ จ.}

\prose{อนาคามิ\-มคฺคกฺขเณ ทสฺสนฏฺเฐน สมฺมาทิฏฺฐิ ฯเปฯ. อวิกฺเขป\-ฏฺเฐน สมฺมาสมาธิ อณุสหคตสฺส กามราค\-สญฺโญชนสฺส ปฏิฆ\-สญฺโญชนสฺส, อณุสหคตสฺส กาม\-ราคา\-นุสยสฺส ปฏิฆา\-นุสยสฺส ปหานาย มคฺโค เจว เหตุ จ, สหชาตานํ ธมฺมานํ อุปตฺถมฺภนาย มคฺโค เจว เหตุ จ, กิเลสานํ ปริยาทานาย มคฺโค เจว เหตุ จ, ปฏิเวธา\-ทิวิ\-โสธนาย มคฺโค เจว เหตุ จ, จิตฺตสฺส อธิฏฺฐานาย มคฺโค เจว เหตุ จ, จิตฺตสฺส โวทานาย มคฺโค เจว เหตุ จ, วิเสสา\-ธิคมาย มคฺโค เจว เหตุ จ, อุตฺตริ\-ปฏิเวธาย มคฺโค เจว เหตุ จ, สจฺจา\-ภิสมยาย มคฺโค เจว เหตุ จ, นิโรเธ ปติฏฺฐาปนาย มคฺโค เจว เหตุ จ.}

\prose{อรหตฺต\-มคฺคกฺขเณ ทสฺสนฏฺเฐน สมฺมาทิฏฺฐิ ฯเปฯ. อวิกฺเขป\-ฏฺเฐน สมฺมาสมาธิ รูปราคสฺส อรูปราคสฺส มานสฺส อุทฺธจฺจสฺส อวิชฺชาย, มานานุสยสฺส ภวราคา\-นุสยสฺส อวิชฺชา\-นุสยสฺส ปหานาย มคฺโค เจว เหตุ จ, สหชาตานํ ธมฺมานํ อุปตฺถมฺภนาย มคฺโค เจว เหตุ จ, กิเลสานํ ปริยาทานาย มคฺโค เจว เหตุ จ, ปฏิเวธา\-ทิวิ\-โสธนาย มคฺโค เจว เหตุ จ, จิตฺตสฺส อธิฏฺฐานาย มคฺโค เจว เหตุ จ, จิตฺตสฺส โวทานาย มคฺโค เจว เหตุ จ, วิเสสา\-ธิคมาย มคฺโค เจว เหตุ จ, อุตฺตริ\-ปฏิเวธาย มคฺโค เจว เหตุ จ, สจฺจา\-ภิสมยาย มคฺโค เจว เหตุ จ, นิโรเธ ปติฏฺฐาปนาย มคฺโค เจว เหตุ จ.}

\prose{ทสฺสนมคฺโค สมฺมาทิฏฺฐิ, อภินิโรปน\-มคฺโค สมฺมาสงฺกปฺโป, ปริคฺคหมคฺโค สมฺมาวาจา, สมุฏฺฐาน\-มคฺโค สมฺมากมฺมนฺโต, โวทานมคฺโค สมฺมาอาชีโว, ปคฺคหมคฺโค สมฺมาวายาโม, อุปฏฺฐานมคฺโค สมฺมาสติ, อวิกฺเขปมคฺโค สมฺมาสมาธิ. อุปฏฺฐานมคฺโค สติสมฺโพชฺฌงฺโค, ปวิจยมคฺโค ธมฺมวิจย\-สมฺโพชฺฌงฺโค, ปคฺคหมคฺโค วีริย\-สมฺโพชฺฌงฺโค, ผรณมคฺโค ปีติสมฺโพชฺฌงฺโค,}

\csromanmatikaanchor{mtk.83}
\csromantocmarkref{section}{10. มณฺฑเปยฺยกถา}{mtk.83}
\bookmark[dest={mtk.83},level=1]{10. มณฺฑเปยฺยกถา}
\prose{\csromanfolio{277}อุปสมมคฺโค ปสฺสทฺธิ\-สมฺโพชฺฌงฺโค, อวิกฺเขปมคฺโค สมาธิ\-สมฺโพชฺฌงฺโค, ปฏิสงฺขาน\-มคฺโค อุเปกฺขา\-สมฺโพชฺฌงฺโค.}

\prose{อสฺสทฺธิเย อกมฺปิยมคฺโค สทฺธาพลํ, โกสชฺเช อกมฺปิยมคฺโค วีริยพลํ, ปมาเท อกมฺปิยมคฺโค สติพลํ, อุทฺธจฺเจ อกมฺปิยมคฺโค สมาธิพลํ, อวิชฺชาย อกมฺปิยมคฺโค ปญฺญาพลํ, อธิโมกฺขมคฺโค สทฺธินฺทฺริยํ, ปคฺคหมคฺโค วีริยินฺทฺริยํ, อุปฏฺฐานมคฺโค สตินฺทฺริยํ, อวิกฺเขปมคฺโค สมาธินฺทฺริยํ, ทสฺสนมคฺโค ปญฺญินฺทฺริยํ.}

\prosewithcloserruled{อาธิปเตยฺย\-ฏฺเฐน อินฺทฺริยา มคฺโค, อกมฺปิยฏฺเฐน พลา มคฺโค, นิยฺยานฏฺเฐน โพชฺฌงฺคา มคฺโค, เหตุฏฺเฐน มคฺคงฺคา มคฺโค, อุปฏฺฐาน\-ฏฺเฐน สติปฏฺฐานา มคฺโค, ปทหนฏฺเฐน สมฺมปฺปธานา มคฺโค, อิชฺฌนฏฺเฐน อิทฺธิปาทา มคฺโค, ตถฏฺเฐน สจฺจานิ มคฺโค, อวิกฺเขป\-ฏฺเฐน สมโถ มคฺโค, อนุปสฺสนฏฺเฐน วิปสฺสนา มคฺโค, เอกรสฏฺเฐน สมถ\-วิปสฺสนา มคฺโค. อนติวตฺตน\-ฏฺเฐน ยุคนทฺธา มคฺโค. สํวรฏฺเฐน สีลวิสุทฺธิ มคฺโค, อวิกฺเขป\-ฏฺเฐน จิตฺตวิสุทฺธิ มคฺโค, ทสฺสนฏฺเฐน ทิฏฺฐิวิสุทฺธิ มคฺโค, มุตฺตฏฺเฐน วิโมกฺโข มคฺโค, ปฏิเวธ\-ฏฺเฐน วิชฺชา มคฺโค, ปริจฺจาค\-ฏฺเฐน วิมุตฺติ มคฺโค, สมุจฺเฉท\-ฏฺเฐน ขเย ญาณํ มคฺโค, ฉนฺโท มูลฏฺเฐน มคฺโค, มนสิกาโร สมุฏฺฐาน\-ฏฺเฐน มคฺโค, ผสฺโส สโมธาน\-ฏฺเฐน มคฺโค, เวทนา สโมสรณ\-ฏฺเฐน มคฺโค, สมาธิ ปมุขฏฺเฐน มคฺโค, สติ อาธิปเตยฺย\-ฏฺเฐน มคฺโค, ปญฺญา ตทุตฺต\-รฏฺเฐน มคฺโค, วิมุตฺติ สารฏฺเฐน มคฺโค, อมโตคธํ นิพฺพานํ ปริโยสาน\-ฏฺเฐน มคฺโคติ.}{มคฺคกถา นิฏฺฐิตา.}
\proseitem{238}{มณฺฑเปยฺยมิทํ ภิกฺขเว พฺรหฺมจริยํ, สตฺถา สมฺมุขีภูโต. ติธตฺตมณฺโฑ\csromansharedfootnote{bf934cc989b0ce2f}{ติวิโธ มณฺโฑ (พหูสุ) อฏฺฐกถา โอโลเกตพฺพา.} สตฺถริ สมฺมุขีภูเต เทสนามณฺโฑ ปฏิคฺคหมณฺโฑ พฺรหฺมจริย\-มณฺโฑ.}

\prose{กตโม เทสนามณฺโฑ? จตุนฺนํ อริยสจฺจานํ อาจิกฺขนา เทสนา ปญฺญาปนา ปฏฺฐปนา วิวรณา วิภชนา อุตฺตานีกมฺมํ\csromansharedfootnote{93855b27a579489b}{อุตฺตานิกมฺมํ (ก)}. จตุนฺนํ สติปฏฺฐานานํ ฯเปฯ. จตุนฺนํ สมฺม\-ปฺปธานานํ. จตุนฺนํ อิทฺธิปาทานํ. ปญฺจนฺนํ อินฺทฺริยานํ.}

\prose{\csromanfolio{278}ปญฺจนฺนํ พลานํ. สตฺตนฺนํ โพชฺฌงฺคานํ. อริยสฺส อฏฺฐงฺคิกสฺส มคฺคสฺส อาจิกฺขนา เทสนา ปญฺญาปนา ปฏฺฐปนา วิวรณา วิภชนา อุตฺตานีกมฺมํ. อยํ เทสนามณฺโฑ. (1)}

\prose{กตโม \textbf{ปฏิคฺคหมณฺโฑ?} ภิกฺขู ภิกฺขุนิโย อุปาสกา อุปาสิกาโย เทวา มนุสฺสา, เย วา ปนญฺเญปิ เกจิ วิญฺญาตาโร. อยํ ปฏิคฺคหมณฺโฑ. (2)}

\prose{กตโม พฺรหฺมจริย\-มณฺโฑ? อยเมว อริโย อฏฺฐงฺคิโก มคฺโค. เสยฺยถิทํ\csromandash{}สมฺมาทิฏฺฐิ สมฺมาสงฺกปฺโป สมฺมาวาจา สมฺมากมฺมนฺโต สมฺมาอาชีโว สมฺมาวายาโม สมฺมาสติ สมฺมาสมาธิ. อยํ พฺรหฺมจริย\-มณฺโฑ. (3)}

\proseitem{239}{อธิโมกฺข\-มณฺโฑ สทฺธินฺทฺริยํ, อสฺสทฺธิยํ กสโฏ, อสฺสทฺธิยํ กสฏํ ฉฑฺเฑตฺวา สทฺธิ\-นฺทฺริยสฺส อธิโมกฺข\-มณฺฑํ ปิวตีติ มณฺฑเปยฺยํ. ปคฺคหมณฺโฑ วีริยินฺทฺริยํ, โกสชฺชํ กสโฏ, โกสชฺชํ กสฏํ ฉฑฺเฑตฺวา วีริยิ\-นฺทฺริยสฺส ปคฺคหมณฺฑํ ปิวตีติ มณฺฑเปยฺยํ. อุปฏฺฐานมณฺโฑ สตินฺทฺริยํ, ปมาโท กสโฏ, ปมาทํ กสฏํ ฉฑฺเฑตฺวา สตินฺทฺริยสฺส อุปฏฺฐาน\-มณฺฑํ ปิวตีติ มณฺฑเปยฺยํ. อวิกฺเขปมณฺโฑ สมาธินฺทฺริยํ, อุทฺธจฺจํ กสโฏ, อุทฺธจฺจํ กสฏํ ฉฑฺเฑตฺวา สมาธิ\-นฺทฺริยสฺส อวิกฺเขป\-มณฺฑํ ปิวตีติ มณฺฑเปยฺยํ. ทสฺสนมณฺโฑ ปญฺญินฺทฺริยํ, อวิชฺชา กสโฏ, อวิชฺชํ กสฏํ ฉฑฺเฑตฺวา ปญฺญินฺทฺริยสฺส ทสฺสนมณฺฑํ ปิวตีติ มณฺฑเปยฺยํ.}

\prose{อสฺสทฺธิเย อกมฺปิยมณฺโฑ สทฺธาพลํ, อสฺสทฺธิยํ กสโฏ, อสฺสทฺธิยํ กสฏํ ฉฑฺเฑตฺวา สทฺธาพลสฺส อสฺสทฺธิเย อกมฺปิยมณฺฑํ ปิวตีติ มณฺฑเปยฺยํ. โกสชฺเช อกมฺปิยมณฺโฑ วิริยพลํ, โกสชฺชํ กสโฏ, โกสชฺชํ กสฏํ ฉฑฺเฑตฺวา วีริยพลสฺส โกสชฺเช อกมฺปิยมณฺฑํ ปิวตีติ มณฺฑเปยฺยํ. ปมาเท อกมฺปิยมณฺโฑ สติพลํ, ปมาโท กสโฏ, ปมาทํ กสฏํ ฉฑฺเฑตฺวา สติพลสฺส ปมาเท อกมฺปิยมณฺฑํ ปิวตีติ มณฺฑเปยฺยํ. อุทฺธจฺเจ อกมฺปิยมณฺโฑ สมาธิพลํ, อุทฺธจฺจํ กสโฏ, อุทฺธจฺจํ กสฏํ ฉฑฺเฑตฺวา สมาธิพลสฺส อุทฺธจฺเจ อกมฺปิยมณฺฑํ ปิวตีติ มณฺฑเปยฺยํ. อวิชฺชาย อกมฺปิยมณฺโฑ ปญฺญาพลํ, อวิชฺชา กสโฏ, อวิชฺชํ กสฏํ ฉฑฺเฑตฺวา ปญฺญาพลสฺส อวิชฺชาย อกมฺปิยมณฺฑํ ปิวตีติ มณฺฑเปยฺยํ.}

\prose{\csromanfolio{279}อุปฏฺฐานมณฺโฑ สติสมฺโพชฺฌงฺคา, ปมาโท กสโฏ, ปมาทํ กสฏํ ฉฑฺเฑตฺวา สติสมฺโพชฺฌงฺคสฺส อุปฏฺฐาน\-มณฺฑํ ปิวตีติ มณฺฑเปยฺยํ. ปวิจยมณฺโฑ ธมฺมวิจย\-สมฺโพชฺฌงฺโค, อวิชฺชา กสโฏ, อวิชฺชํ กสฏํ ฉฑฺเฑตฺวา ธมฺมวิจย\-สมฺโพชฺฌงฺคสฺส ปวิจยมณฺฑํ ปิวตีติ มณฺฑเปยฺยํ. ปคฺคหมณฺโฑ วีริย\-สมฺโพชฺฌงฺโค, โกสชฺชํ กสโฏ, โกสชฺชํ กสฏํ ฉฑฺเฑตฺวา วีริย\-สมฺโพชฺฌงฺคสฺส ปคฺคหมณฺฑํ ปิวตีติ มณฺฑเปยฺยํ. ผรณมณฺโฑ ปีติสมฺโพชฺฌงฺโค, ปริฬาโห กสโฏ, ปริฬาหํ กสฏํ ฉฑฺเฑตฺวา ปีติสมฺโพชฺฌงฺคสฺส ผรณมณฺฑํ ปิวตีติ มณฺฑเปยฺยํ. อุปสมมณฺโฑ ปสฺสทฺธิ\-สมฺโพชฺฌงฺโค, ทุฏฺฐุลฺลํ กสโฏ, ทุฏฺฐุลฺลํ กสฏํ ฉฑฺเฑตฺวา ปสฺสทฺธิ\-สมฺโพชฺฌงฺคสฺส อุปสมมณฺฑํ ปิวตีติ มณฺฑเปยฺยํ. อวิกฺเขปมณฺโฑ สมาธิ\-สมฺโพชฺฌงฺโค, อุทฺธจฺจํ กสโฏ, อุทฺธจฺจํ กสฏํ ฉฑฺเฑตฺวา สมาธิ\-สมฺโพชฺฌงฺคสฺส อวิกฺเขป\-มณฺฑํ ปิวตีติ มณฺฑเปยฺยํ. ปฏิสงฺขาน\-มณฺโฑ อุเปกฺขา\-สมฺโพชฺฌงฺโค, อปฺปฏิสงฺขา กสโฏ, อปฺปฏิสงฺขํ กสฏํ ฉฑฺเฑตฺวา อุเปกฺขา\-สมฺโพชฺฌงฺคสฺส ปฏิสงฺขาน\-มณฺฑํ ปิวตีติ มณฺฑเปยฺยํ.}

\prose{ทสฺสนมณฺโฑ สมฺมาทิฏฺฐิ, มิจฺฉาทิฏฺฐิ กสโฏ, มิจฺฉาทิฏฺฐิํ กสฏํ ฉฑฺเฑตฺวา สมฺมาทิฏฺฐิยา ทสฺสนมณฺฑํ ปิวตีติ มณฺฑเปยฺยํ. อภินิโรปน\-มณฺโฑ สมฺมาสงฺกปฺโป, มิจฺฉาสงฺกปฺโป กสโฏ, มิจฺฉา\-สงฺกปฺปํ กสฏํ ฉฑฺเฑตฺวา สมฺมา\-สงฺกปฺปสฺส อภินิโรปน\-มณฺฑํ ปิวตีติ มณฺฑเปยฺยํ. ปริคฺคหมณฺโฑ สมฺมาวาจา, มิจฺฉาวาจา กสโฏ, มิจฺฉาวาจํ กสฏํ ฉฑฺเฑตฺวา สมฺมาวาจาย ปริคฺคห\-มณฺฑํ ปิวตีติ มณฺฑเปยฺยํ. สมุฏฺฐาน\-มณฺโฑ สมฺมากมฺมนฺโต, มิจฺฉากมฺมนฺโต กสโฏ, มิจฺฉา\-กมฺมนฺตํ กสฏํ ฉฏฺเฏตฺวา สมฺมา\-กมฺมนฺตสฺส สมุฏฺฐาน\-มณฺฑํ ปิวตีติ มณฺฑเปยฺยํ. โวทานมณฺโฑ สมฺมาอาชีโว, มิจฺฉาอาชีโว กสโฏ, มิจฺฉาอาชีวํ กสฏํ ฉฑฺเฑตฺว ฉฑฺเฑตฺวา สมฺมาอาชีวสฺส โวทานมณฺฑํ ปิวตีติ มณฺฑเปยฺยํ. ปคฺคหมณฺโฑ สมฺมาวายาโม, มิจฺฉาวายาโม กสโฏ. มิจฺฉาวายามํ กสฏํ ฉฑฺเฑตฺวา สมฺมาวายามสฺส ปคฺคหมณฺฑํ ปิวตีติ มณฺฑเปยฺยํ. อุปฏฺฐานมณฺโฑ สมฺมาสติ, มิจฺฉาสติ กสโฏ, มิจฺฉาสติํ กสฏํ ฉฑฺเฑตฺวา สมฺมาสติยา อุปฏฺฐาน\-มณฺฑํ ปิวตีติ มณฺฑเปยฺยํ. อวิกฺเขปมณฺโฑ สมฺมาสมาธิ, มิจฺฉาสมาธิ กสโฏ, มิจฺฉาสมาธิํ กสฏํ ฉฑฺเฑตฺวา สมฺมา\-สมาธิสฺส อวิกฺเขป\-มณฺฑํ ปิวตีติ มณฺฑเปยฺยํ.}

\proseitem{240}{อตฺถิ มณฺโฑ, อตฺถิ เปยฺยํ, อตฺถิ กสโฏ. อธิโมกฺข\-มณฺโฑ สทฺธินฺทฺริยํ, อสฺสทฺธิยํ กสโฏ, โย ตตฺถ อตฺถรโส ธมฺมรโส \csromanfolio{280}วิมุตฺติรโส, อิทํ เปยฺยํ. ปคฺคหมณฺโฑ วีริยินฺทฺริยํ, โกสชฺชํ กสโฏ, โย ตตฺถ อตฺถรโส ธมฺมรโส วิมุตฺติรโส, อิทํ เปยฺยํ. อุปฏฺฐานมณฺโฑ สตินฺทฺริยํ, ปมาโท กสโฏ, โย ตตฺถ อตฺถรโส ธมฺมรโส วิมุตฺติรโส, อิทํ เปยฺยํ. อวิกฺเขปมณฺโฑ สมาธินฺทฺริยํ, อุทฺธจฺจํ กสโฏ, โย ตตฺถ อตฺถรโส ธมฺมรโส วิมุตฺติรโส, อิทํ เปยฺยํ. ทสฺสนมณฺโฑ ปญฺญินฺทฺริยํ, อวิชฺชา กสโฏ, โย ตตฺถ อตฺถรโส ธมฺมรโส วิมุตฺติรโส, อิทํ เปยฺยํ.}

\prose{อสฺสทฺธิเย อกมฺปิยมณฺโฑ สทฺธาพลํ, อสฺสทฺธิยํ กสโฏ, โย ตตฺถ อตฺถรโส ธมฺมรโส วิมุตฺติรโส, อิทํ เปยฺยํ. โกสชฺเช อกมฺปิยมณฺโฑ วีริยพลํ, โกสชฺชํ กสโฏ, โย ตตฺถ อตฺถรโส ธมฺมรโส วิมุตฺติรโส, อิทํ เปยฺยํ. ปมาเท อกมฺปิยมณฺโฑ สติพลํ, ปมาโท กสโฏ, โย ตตฺถ อตฺถรโส ธมฺมรโส วิมุตฺติรโส, อิทํ เปยฺยํ. อุทฺธจฺเจ อกมฺปิยมณฺโฑ สมาธิพลํ, อุทฺธจฺจํ กสโฏ, โย ตตฺถ อตฺถรโส ธมฺมรโส วิมุตฺติรโส, อิทํ เปยฺยํ. อวิชฺชาย อกมฺปิยมณฺโฑ ปญฺญาพลํ, อวิชฺชา กสโฏ, โย ตตฺถ อตฺถรโส ธมฺมรโส วิมุตฺติรโส, อิทํ เปยฺยํ.}

\prose{อุปฏฺฐานมณฺโฑ สติสมฺโพชฺฌงฺโค, ปมาโท กสโฏ, โย ตตฺถ อตฺถรโส ธมฺมรโส วิมุตฺติรโส, อิทํ เปยฺยํ. ปวิจยมณฺโฑ ธมฺมวิจย\-สมฺโพชฺฌงฺโค, อวิชฺชา กสโฏ, โย ตตฺถ อตฺถรโส ธมฺมรโส วิมุตฺติรโส, อิทํ เปยฺยํ. ปคฺคหมณฺโฑ วีริย\-สมฺโพชฺฌงฺโค, โกสชฺชํ กสโฏ, โย ตตฺถ อตฺถรโส ธมฺมรโส วิมุตฺติรโส, อิทํ เปยฺยํ. ผรณมณฺโฑ ปีติสมฺโพชฺฌงฺโค, ปริฬาโห กสโฏ, โย ตตฺถ อตฺถรโส ธมฺมรโส วิมุตฺติรโส, อิทํ เปยฺยํ. อุปสมมณฺโฑ ปสฺสทฺธิ\-สมฺโพชฺฌงฺโค, ทุฏฺฐุลฺลํ กสโฏ, โย ตตฺถ อตฺถรโส ธมฺมรโส วิมุตฺติรโส, อิทํ เปยฺยํ. อวิกฺเขปมณฺโฑ สมาธิ\-สมฺโพชฺฌงฺโค, อุทฺธจฺจํ กสโฏ, โย ตตฺถ อตฺถรโส ธมฺมรโส วิมุตฺติรโส, อิทํ เปยฺยํ. ปฏิสงฺขาน\-มณฺโฑ อุเปกฺขา\-สมฺโพชฺฌงฺโค, อปฏิสงฺขา กสโฏ, โย ตตฺถ อตฺถรโส ธมฺมรโส วิมุตฺติรโส, อิทํ เปยฺยํ.}

\prose{ทสฺสนมณฺโฑ สมฺมาทิฏฺฐิ, มิจฺฉาทิฏฺฐิ กสโฏ, โย ตตฺถ อตฺถรโส ธมฺมรโส วิมุตฺติรโส, อิทํ เปยฺยํ. อภินิโรปน\-มณฺโฑ สมฺมาสงฺกปฺโป, มิจฺฉาสงฺกปฺโป กสโฏ, โย ตตฺถ อตฺถรโส ธมฺมรโส วิมุตฺติรโส,}

\prose{\csromanfolio{281}อิทํ เปยฺยํ. ปริคฺคหมณฺโฑ สมฺมาวาจา, มิจฺฉาวาจา กสโฏ, โย ตตฺถ อตฺถรโส ธมฺมรโส วิมุตฺติรโส, อิทํ เปยฺยํ. สมุฏฺฐาน\-มณฺโฑ สมฺมากมฺมนฺโต, มิจฺฉากมฺมนฺโต กสโฏ, โย ตตฺถ อตฺถรโส ธมฺมรโส วิมุตฺติรโส, อิทํ เปยฺยํ. โวทานมณฺโฑ สมฺมาอาชีโว, มิจฺฉาอาชีโว กสโฏ, โย ตตฺถ อตฺถรโส ธมฺมรโส วิมุตฺติรโส, อิทํ เปยฺยํ. ปคฺคหมณฺโฑ สมฺมาวายาโม, มิจฺฉาวายาโม กสโฏ, โย ตตฺถ อตฺถรโส ธมฺมรโส วิมุตฺติรโส, อิทํ เปยฺยํ. อุปฏฺฐานมณฺโฑ สมฺมาสติ, มิจฺฉาสติ กสโฏ, โย ตตฺถ อตฺถรโส ธมฺมรโส วิมุตฺติรโส, อิทํ เปยฺยํ. อวิกฺเขปมณฺโฑ สมฺมาสมาธิ, มิจฺฉาสมาธิ กสโฏ, โย ตตฺถ อตฺถรโส ธมฺมรโส วิมุตฺติรโส, อิทํ เปยฺยํ.}

\prose{ทสฺสนมณฺโฑ สมฺมาทิฏฺฐิ. อภินิโรปน\-มณฺโฑ สมฺมาสงฺกปฺโป. ปริคฺคหมณฺโฑ สมฺมาวาจา. สมุฏฺฐาน\-มณฺโฑ สมฺมากมฺมนฺโต. โวทานมณฺโฑ สมฺมาอาชีโว. ปคฺคหมณฺโฑ สมฺมาวายาโม. อุปฏฺฐานมณฺโฑ สมฺมาสติ. อวิกฺเขปมณฺโฑ สมฺมาสมาธิ.}

\prose{อุปฏฺฐานมณฺโฑ สติสมฺโพชฺฌงฺโค, ปวิจยมณฺโฑ ธมฺมวิจย\-สมฺโพชฺฌงฺโค, ปคฺคหมณฺโฑ วีริย\-สมฺโพชฺฌงฺโค, ผรณมณฺโฑ ปีติสมฺโพชฺฌงฺโค, อุปสมมณฺโฑ ปสฺสทฺธิ\-สมฺโพชฺฌงฺโค, อวิกฺเขปมณฺโฑ สมาธิ\-สมฺโพชฺฌงฺโค, ปฏิสงฺขาน\-มณฺโฑ อุเปกฺขา\-สมฺโพชฺฌงฺโค.}

\prose{อสฺสทฺธิเย อกมฺปิยมณฺโฑ สทฺธาพลํ, โกสชฺเช อกมฺปิยมณฺโฑ วีริยพลํ, ปมาเท อกมฺปิยมณฺโฑ สติพลํ, อุทฺธจฺเจ อกมฺปิยมณฺโฑ สมาธิพลํ, อวิชฺชาย อกมฺปิยมณฺโฑ ปญฺญาพลํ.}

\prose{อธิโมกฺข\-มณฺโฑ สทฺธินฺทฺริยํ, ปคฺคหมณฺโฑ วีริยินฺทฺริยํ, อุปฏฺฐานมณฺโฑ สตินฺทฺริยํ, อวิกฺเขปมณฺโฑ สมาธินฺทฺริยํ, ทสฺสนมณฺโฑ ปญฺญินฺทฺริยํ.}

\prose{อาธิปเตยฺย\-ฏฺเฐน อินฺทฺริยา มณฺโฑ, อกมฺปิยฏฺเฐน พลา มณฺโฑ, นิยฺยานฏฺเฐน โพชฺฌงฺคา มณฺโฑ, เหตุฏฺเฐน มคฺโค มณฺโฑ, อุปฏฺฐาน\-ฏฺเฐน สติปฏฺฐานา มณฺโฑ, ปทหนฏฺเฐน สมฺมปฺปธานา มณฺโฑ, อิชฺฌนฏฺเฐน อิทฺธิปาทา มณฺโฑ, ตถฏฺเฐน สจฺจา มณฺโฑ, อวิกฺเขป\-ฏฺเฐน สมโถ มณฺโฑ, อนุปสฺสนฏฺเฐน วิปสฺสนา มณฺโฑ, เอกรสฏฺเฐน สมถ\-วิปสฺสนา มณฺโฑ, อนติวตฺตน\-ฏฺเฐน ยุคนทฺธา มณฺโฑ, สํวรฏฺเฐน สีลวิสุทฺธิ มณฺโฑ, อวิกฺเขป\-ฏฺเฐน จิตฺตวิสุทฺธิ มณฺโฑ, ทสฺสนฏฺเฐน ทิฏฺฐิวิสุทฺธิ มณฺโฑ, มุตฺตฏฺเฐน}

\prosewithcloser{\csromanfolio{282}วิโมกฺโข มณฺโฑ, ปฏิเวธ\-ฏฺเฐน วิชฺชา มณฺโฑ, ปริจฺจาค\-ฏฺเฐน วิมุตฺติ มณฺโฑ, สมุจฺเฉท\-ฏฺเฐน ขเย ญาณํ มณฺโฑ, ปฏิปฺปสฺสทฺธ\-ฏฺเฐน อนุปฺปาเท ญาณํ มณฺโฑ. ฉนฺโท มูลฏฺเฐน มณฺโฑ, มนสิกาโร สมุฏฺฐาน\-ฏฺเฐน มณฺโฑ, ผสฺโส สโมธาน\-ฏฺเฐน มณฺโฑ, เวทนา สโมสรณ\-ฏฺเฐน มณฺโฑ, สมาธิ ปมุขฏฺเฐน มณฺโฑ, สติ อาธิปเตยฺย\-ฏฺเฐน มณฺโฑ, ปญฺญา ตทุตฺต\-รฏฺเฐน มณฺโฑ, วิมุตฺติ สารฏฺเฐน มณฺโฑ, อมโตคธํ นิพฺพานํ ปริโยสาน\-ฏฺเฐน มณฺโฑติ.}{มณฺฑเปยฺย\-กถา นิฏฺฐิตา.}
\csromancenter{จตุตฺถ\-ภาณวาโร.}
\nitthitammidruled{\textbf{มหาวคฺโค ปฐโม.}}
\csromancenter{\textbf{ตสฺสุทฺทานํ}}
\setlength{\csromangathapagewidth}{0pt}
\setlength{\csromangathawakwidth}{0pt}
\csromangathameasuregroup{ญาณทิฏฺฐี จ อสฺสาสา, \\ คติกมฺมวิปลฺลาสา,}{\csromangathabat{ญาณทิฏฺฐี จ อสฺสาสา,}{อินฺทฺริยํ วิโมกฺขปญฺจมา.} \\ \csromangathabat{คติกมฺมวิปลฺลาสา,}{มคฺโค มณฺเฑน เต ทสาติ.}}
\csromangathasetleft{ญาณทิฏฺฐี จ อสฺสาสา, \\ คติกมฺมวิปลฺลาสา,}
\csromangathagroup{\csromangathastanza{\csromangathabat{ญาณทิฏฺฐี จ อสฺสาสา,}{อินฺทฺริยํ วิโมกฺขปญฺจมา.} \\ \csromangathabat{คติกมฺมวิปลฺลาสา,}{มคฺโค มณฺเฑน เต ทสาติ.}}}

\prose{เอส นิกายธเรหิ ฐปิโต อสโม\csromansharedfootnote{338dd9d64fd7ee1e}{[มิสฺสินฺคฺ โนเต 0]} ปฐโม ปวโร วรวคฺโคติ\csromansharedfootnote{61785d03c494c2ae}{[มิสฺสินฺคฺ โนเต 1]}.}

\csromanmatikaanchor{mtk.84}
\csromantocmarknopage{chapter}{2. ยุคนทฺธวคฺค}{mtk.84}
\bookmark[dest={mtk.84},level=0]{2. ยุคนทฺธวคฺค}
\chapterheadpageread{2. ยุคนทฺธ\-วคฺค}
\proseitem{1}{\csromanfolio{283}เอวํ เม สุตํ, เอกํ สมยํ อายสฺมา อานนฺโท โกสมฺพิยํ วิหรติ โฆสิตาราเม. ตตฺร โข อายสฺมา อานนฺโท ภิกฺขู อามนฺเตสิ “อาวุโส ภิกฺขโว”ติ. “อาวุโส”ติ โข เต ภิกฺขู อายสฺมโต อานนฺทสฺส ปจฺจสฺโสสุํ. อายสฺมา อานนฺโท เอตทโวจ\csromandash{}}

\prose{โย หิ โกจิ อาวุโส ภิกฺขุ วา ภิกฺขุนี วา มม สนฺติเก อรหตฺต\-ปฺปตฺตํ\csromansharedfootnote{38377dcff48d3ef0}{อรหตฺตํ (สฺยา), อรหตฺตปตฺติํ (อํ 1. 475 ปิฏฺเฐ.)} พฺยากโรติ สพฺพโส จตูหิ มคฺเคหิ, เอเตสํ วา อญฺญตเรน. กตเมหิ จตูหิ?}

\prose{อิธาวุโส ภิกฺขุ สมถ\-ปุพฺพงฺคมํ วิปสฺสนํ ภาเวติ, ตสฺส สมถ\-ปุพฺพงฺคมํ วิปสฺสนํ ภาวยโต มคฺโค สญฺชายติ, โส ตํ มคฺคํ อาเสวติ ภาเวติ พหุลีกโรติ\csromansharedfootnote{a6f9453652f58a9d}{พหุลิํ กโรติ (ก) อํ 1. 475 ปิฏฺเฐ ปสฺสิตพฺพา.}, ตสฺส ตํ มคฺคํ อาเสวโต ภาวยโต พหุลี กโรโต สญฺโญชนานิ ปหียนฺติ, อนุสยา พฺยนฺตีโหนฺติ.}

\prose{ปุน จปรํ อาวุโส ภิกฺขุ วิปสฺสนา\-ปุพฺพงฺคมํ สมถํ ภาเวติ, ตสฺส วิปสฺสนา\-ปุพฺพงฺคมํ สมถํ ภาวยโต มคฺโค สญฺชายติ, โส ตํ มคฺคํ อาเสวติ ภาเวติ พหุลีกโรติ, ตสฺส ตํ มคฺคํ อาเสวโต ภาวยโต พหุลีกโรโต สญฺโญชนานิ ปหียนฺติ, อนุสยา พฺยนฺตีโหนฺติ.}

\prose{ปุน จปรํ อาวุโส ภิกฺขุ สมถ\-วิปสฺสนํ ยุคนทฺธํ\csromansharedfootnote{5095021e6a593337}{ยุคนนฺธํ (ก, สีฏฺฐ)} ภาเวติ, ตสฺส สมถ\-วิปสฺสนํ ยุคนทฺธํ ภาวยโต มคฺโค สญฺชายติ, โส ตํ มคฺคํ อาเสวติ ภาเวติ พหุลีกโรติ, ตสฺส ตํ มคฺคํ อาเสวโต ภาวยโต พหุลีกโรโต สญฺโญชนานิ ปหียนฺติ, อนุสยา พฺยนฺตีโหนฺติ.}

\prose{ปุน จปรํ อาวุโส ภิกฺขุโน ธมฺมุ\-ทฺธจฺจ\-วิคฺคหิตํ มานสํ โหติ, โหติ โส อาวุโส สมโย ยํ ตํ จิตฺตํ อชฺฌตฺตเมว\csromansharedfootnote{e19e28cdd64458cc}{อชฺฌตฺตญฺเญว (สฺยา, ก)} สนฺติฏฺฐติ สนฺนิสีทติ \csromanfolio{284}เอโกทิ โหติ สมาธิยติ, ตสฺส มคฺโค สญฺชายติ, โส ตํ มคฺคํ อาเสวติ ภาเวติ พหุลีกโรติ, ตสฺส ตํ มคฺคํ อาเสวโต ภาวยโต พหุลีกโรโต สญฺโญชนานิ ปหียนฺติ, อนุสยา พฺยนฺตีโหนฺติ.}

\prose{โย หิ โกจิ อาวุโส ภิกฺขุ วา ภิกฺขุนี วา มม สนฺติเก อรหตฺต\-ปฺปตฺตํ พฺยากโรติ สพฺพโส อิเมหิ จตูหิ มคฺเคหิ, เอเตสํ วา อญฺญตเรนาติ.}

\csromanmatikaanchor{mtk.85}
\csromanmatikaanchor{mtk.86}
\csromantocmarknopage{section}{1. ยุคนทฺธกถา}{mtk.85}
\bookmark[dest={mtk.85},level=1]{1. ยุคนทฺธกถา}
\csromantocmarkref{subsection}{1. สุตฺตนฺตนิทฺเทส}{mtk.86}
\bookmark[dest={mtk.86},level=2]{1. สุตฺตนฺตนิทฺเทส}
\csromanheader{2}{1. สุตฺตนฺต\-นิทฺเทส}
\proseitem{1}{กถํ สมถ\-ปุพฺพงฺคมํ วิปสฺสนํ ภาเวติ? เนกฺขมฺม\-วเสน จิตฺตสฺส เอกคฺคตา อวิกฺเขโป สมาธิ, ตตฺถ ชาเต ธมฺเม อนิจฺจโต อนุปสฺสนฏฺเฐน วิปสฺสนา, ทุกฺขโต อนุปสฺสนฏฺเฐน วิปสฺสนา, อนตฺตโต อนุปสฺสฏฺเฐน วิปสฺสนา. อิติ ปฐมํ สมโถ, ปจฺฉา วิปสฺสนา. เตน วุจฺจติ “สมถ\-ปุพฺพงฺคมํ วิปสฺสนํ ภาเวตี”ติ.  ภาเวตีติ จตสฺโส ภาวนา, ตตฺถ ชาตานํ ธมฺมานํ อนติวตฺตน\-ฏฺเฐน ภาวนา, อินฺทฺริยานํ เอกรสฏฺเฐน ภาวนา, ตทุปค\-วีริย\-วาห\-นฏฺเฐน ภาวนา, อาเสวนฏฺเฐน ภาวนา.}

\prose{มคฺโค สญฺชายตีติ กถํ มคฺโค สญฺชายติ? ทสฺสนฏฺเฐน สมฺมาทิฏฺฐิ มคฺโค สญฺชายติ, อภินิโรปน\-ฏฺเฐน สมฺมาสงฺกปฺโป มคฺโค สญฺชายติ, ปริคฺคห\-ฏฺเฐน สมฺมาวาจา มคฺโค สญฺชายติ, สมุฏฺฐาน\-ฏฺเฐน สมฺมากมฺมนฺโต มคฺโค สญฺชายติ, โวทานฏฺเฐน สมฺมาอาชีโว มคฺโค สญฺชายติ, ปคฺคหฏฺเฐน สมฺมาวายาโม มคฺโค สญฺชายติ, อุปฏฺฐาน\-ฏฺเฐน สมฺมาสติ มคฺโค สญฺชายติ, อวิกฺเขป\-ฏฺเฐน สมฺมาสมาธิ มคฺโค สญฺชายติ. เอวํ มคฺโค สญฺชายติ.}

\prose{\textbf{โส ตํ มคฺคํ อาเสวติ ภาเวติ พหุลีกโรติ อาเสวตี}ติ กถํ อาเสวติ? อาวชฺชนฺโต อาเสวติ, ชานนฺโต อาเสวติ, ปสฺสนฺโต อาเสวติ, ปจฺจเวกฺขนฺโต อาเสวติ, จิตฺตํ อธิฏฺฐหนฺโต อาเสวติ, สทฺธาย อธิมุจฺจนฺโต อาเสวติ, วีริยํ ปคฺคณฺหนฺโต อาเสวติ, สติํ อุปฏฺฐาเปนฺโต อาเสวติ, จิตฺตํ สมาทหนฺโต อาเสวติ, ปญฺญาย ปชานนฺโต อาเสวติ, อภิญฺเญยฺยํ อภิชานนฺโต อาเสวติ, ปริญฺเญยฺยํ ปริชานนฺโต อาเสวติ, ปหาตพฺพํ ปชหนฺโต \csromanfolio{285}อาเสวติ, ภาเวตพฺพํ ภาเวนฺโต อาเสวติ, สจฺฉิกาตพฺพํ สจฺฉิกโรนฺโต อาเสวติ, เอวํ อาเสวติ.}

\proseitem{2}{\textbf{ภาเวตี}ติ กถํ ภาเวติ? อาวชฺชนฺโต ภาเวติ, ชนนฺโต ภาเวติ, ปสฺสนฺโต ภาเวติ, ปจฺจเวกฺขนฺโต ภาเวติ, จิตฺตํ อธิฏฺฐหนฺโต ภาเวติ, สทฺธาย อธิมุจฺจนฺโต ภาเวติ วีริยํ ปคฺคณฺหนฺโต ภาเวติ, สติํ อุปฏฺฐาเปนฺโต ภาเวติ, จิตฺตํ สมาทหนฺโต ภาเวติ, ปญฺญาย ปชานนฺโต ภาเวติ. อภิญฺเญยฺยํ อภิชานนฺโต ภาเวติ, ปริญฺเญยฺยํ ปริชานนฺโต ภาเวติ, ปหาตพฺพํ ปชหนฺโต ภาเวติ, ภาเวตพฺพํ ภาเวนฺโต ภาเวติ, สจฺฉิกาตพฺพํ สจฺฉิกโรนฺโต ภาเวติ, เอวํ ภาเวติ.}

\prose{\textbf{พหุลีกโรตี}ติ กถํ พหุลีกโรติ? อาวชฺชนฺโต พหุลีกโรติ, ชานนฺโต พหุลีกโรติ, ปสฺสนฺโต พหุลีกโรติ, ปจฺจเวกฺขนฺโต พหุลีกโรติ, จิตฺตํ อธิฏฺฐหนฺโต พหุลีกโรติ, สทฺธาย อธิมุจฺจนฺโต พหุลีกโรติ, วีริยํ ปคฺคณฺหนฺโต พหุลีกโรติ, สติํ อุปฏฺฐาเปนฺโต พหุลีกโรติ, จิตฺตํ สมาทหนฺโต พหุลีกโรติ, ปญฺญาย ปชานนฺโต พหุลีกโรติ, อภิญฺเญยฺยํ อภิชานนฺโต พหุลีกโรติ, ปริญฺเญยฺยํ ปริชานนฺโต พหุลีกโรติ, ปหาตพฺพํ ปชหนฺโต พหุลีกโรติ, ภาเวตพฺพํ ภาเวนฺโต พหุลีกโรติ, สจฺฉิกาตพฺพํ สจฺฉิกโรนฺโต พหุลีกโรติ, เอวํ พหุลีกโรติ.}

\prose{ตสฺส ตํ มคฺคํ อาเสวโต ภาวยโต พหุลีกโรโต สญฺโญชนานิ ปหียนฺติ, อนุสยา พฺยนฺตีโหนฺตีติ กถํ สญฺโญชนานิ ปหียนฺติ, อนุสยา พฺยนฺตีโหนฺติ? โสตาปตฺติ\-มคฺเคน สกฺกายทิฏฺฐิ วิจิกิจฺฉา สีลพฺพต\-ปรามาโส, อิมานิ ตีณิ สญฺโญชนานิ ปหียนฺติ, ทิฏฺฐานุสโย วิจิกิจฺฉา\-นุสโย อิเม เทฺว อนุสยา พฺยนฺตีโหนฺติ. สกทาคามิ\-มคฺเคน โอฬาริกํ กามราค\-สญฺโญชนํ ปฏิฆ\-สญฺโญชนํ อิมานิ เทฺว สญฺโญชนานิ ปหียนฺติ, โอฬาริโก กามราคานุสโย ปฏิฆานุสโย อิเม เทฺว อนุสยา พฺยนฺตีโหนฺติ. อนาคามิมคฺเคน อณุสหคตํ กามราค\-สญฺโญชนํ ปฏิฆ\-สญฺโญชนํ อิมานิ เทฺว สญฺโญชนานิ ปหียนฺติ, อณุสหคโต กามราคานุสโย ปฏิฆานุสโย อิเม เทฺว อนุสยา พฺยนฺตีโหนฺติ. อรหตฺต\-มคฺเคน รูปราโค อรูปราโค มาโน อุทฺธจฺจํ อวิชฺชา อิมานิ ปญฺจ สญฺโญชนานิ ปหียนฺติ, มานานุสโย \csromanfolio{286}ภวราคา\-นุสโย อวิชฺชานุสโย อิเม ตโย อนุสยา พฺยนฺตีโหนฺติ, เอวํ สญฺโญชนานิ ปหียนฺติ, อนุสยา พฺยนฺตีโหนฺติ.}

\proseitem{1}{อพฺยาปาทวเสน จิตฺตสฺส เอกคฺคตา อวิกฺเขโป สมาธิ ฯเปฯ. อาโลกสญฺญา\-วเสน จิตฺตสฺส เอกคฺคตา อวิกฺเขโป สมาธิ ฯเปฯ. ปฏินิสฺสคฺคา\-นุปสฺสี อสฺสาสวเสน ปฏินิสฺสคฺคา\-นุปสฺสี ปสฺสาสวเสน จิตฺตสฺส เอกคฺคตา อวิกฺเขโป สมาธิ. ตตฺถ ชาเต ธมฺเม อนิจฺจโต อนุปสฺสนฏฺเฐน วิปสฺสนา, ทุกฺขโต อนุปสฺสนฏฺเฐน วิปสฺสนา, อนตฺตโต อนุปสฺสนฏฺเฐน วิปสฺสนา. อิติ ปฐมํ สมโถ, ปจฺฉา วิปสฺสนา. เตน วุจฺจติ “สมถ\-ปุพฺพงฺคมํ วิปสฺสนํ ภาเวตี”ติ.  ภาเวตีติ จตสฺโส ภาวนา. ตตฺถ ชาตานํ ธมฺมานํ อนติวตฺตน\-ฏฺเฐน ภาวนา, อินฺทฺริยานํ เอกรสฏฺเฐน ภาวนา, ตทุปค\-วีริย\-วาห\-นฏฺเฐน ภาวนา, อาเสวนฏฺเฐน ภาวนา.}

\prose{มคฺโค สญฺชายตีติ กถํ มคฺโค สญฺชายติ? ทสฺสนฏฺเฐน สมฺมาทิฏฺฐิ มคฺโค สญฺชายติ, อภินิโรปน\-ฏฺเฐน สมฺมาสงฺกปฺโป มคฺโค สญฺชายติ ฯเปฯ อวิกฺเขป\-ฏฺเฐน สมฺมาสมาธิ มคฺโค สญฺชายติ, เอวํ มคฺโค สญฺชายติ.}

\prose{\textbf{โส ตํ มคฺคํ อาเสวติ ภาเวติ พหุลีกโรติ อาเสวตี}ติ กถํ อาเสวติ? อาวชฺชนฺโต อาเสวติ ฯเปฯ สจฺฉิกาตพฺพํ สจฺฉิกโรนฺโต อาเสวติ, เอวํ อาเสวติ.  ภาเวตีติ กถํ ภาเวติ? อาวชฺชนฺโต ภาเวติ, ชานนฺโต ภาเวติ ฯเปฯ สจฺฉิกาตพฺพํ สจฺฉิกโรนฺโต ภาเวติ, เอวํ ภาเวติ.  พหุลีกโรตีติ กถํ พหุลีกโรติ? อาวชฺชนฺโต พหุลีกโรติ, ชานนฺโต พหุลีกโรติ ฯเปฯ สจฺฉิกาตพฺพํ สจฺฉิกโรนฺโต พหุลีกโรติ, เอวํ พหุลีกโรติ.}

\prose{ตสฺส ตํ มคฺคํ อาเสวโต ภาวยโต พหุลีกโรโต สญฺโญชนานิ ปหียนฺติ, อนุสยา พฺยนฺตีโหนฺตีติ กถํ สญฺโญชนา ปหียนฺติ, อนุสยา พฺยนฺตีโหนฺติ? โสตาปตฺติ\-มคฺเคน สกฺกายทิฏฺฐิ วิจิกิจฺฉา สีลพฺพต\-ปรามาโส, อิมานิ ตีณิ สญฺโญชนานิ ปหียนฺติ, ทิฏฺฐานุสโย วิจิกิจฺฉา\-นุสโย อิเม เทฺว อนุสยา พฺยนฺตีโหนฺติ. สกทาคามิ\-มคฺเคน โอฬาริกํ กามราค\-สญฺโญชนํ ปฏิฆ\-สญฺโญชนํ อิมานิ เทฺว สญฺโญชนานิ ปหียนฺติ, โอฬาริโก กามราคานุสโย \csromanfolio{287}ปฏิฆานุสโย อิเม เทฺว อนุสยา พฺยนฺตีโหนฺติ. อนาคามิมคฺเคน อณุสหคตํ กามราค\-สญฺโญชนํ ปฏิฆ\-สญฺโญชนํ อิมานิ เทฺว สญฺโญชนานิ ปหียนฺติ, อณุสหคโต กามราคานุสโย ปฏิฆานุสโย, อิเม เทฺว อนุสยา พฺยนฺตีโหนฺติ. อรหตฺต\-มคฺเคน รูปราโค อรูปราโค มาโน อุทฺธจฺจํ อวิชฺชา อิมานิ ปญฺจ สญฺโญชนานิ ปหียนฺติ, มานานุสโย ภวราคา\-นุสโย อวิชฺชานุสโย อิเม ตโย อนุสยา พฺยนฺตีโหนฺติ. เอวํ สญฺโญชนานิ ปหียนฺติ, อนุสยา พฺยนฺตีโหนฺติ. เอวํ สมถ\-ปุพฺพงฺคมํ วิปสฺสนํ ภาเวติ.}

\proseitem{3}{กถํ วิปสฺสนา\-ปุพฺพงฺคมํ สมถํ ภาเวติ? อนิจฺจโต อนุปสฺสนฏฺเฐน วิปสฺสนา, ทุกฺขโต อนุปสฺสนฏฺเฐน วิปสฺสนา, อนตฺตโต อนุปสฺสนฏฺเฐน วิปสฺสนา. ตตฺถ ชาตานํ ธมฺมานญฺจ โวสคฺคา\-รมฺมณตา\csromansharedfootnote{a6744c3ab5485f88}{โวสฺสคฺคารมฺมณตา (สฺยา, ก)} จิตฺตสฺส เอกคฺคตา อวิกฺเขโป สมาธิ. อิติ ปฐมํ วิปสฺสนา, ปจฺฉา สมโถ. เตน วุจฺจติ “วิปสฺสนา\-ปุพฺพงฺคมํ สมถํ ภาเวตี”ติ.  ภาเวตีติ จตสฺโส ภาวนา, อาเสวนฏฺเฐน ภาวนา ฯเปฯ. มคฺโค สญฺชายตีติ กถํ มคฺโค สญฺชายติ ฯเปฯ เอวํ มคฺโค สญฺชายติ. เอวํ สญฺโญชนานิ ปหียนฺติ, อนุสยา พฺยนฺตีโหนฺติ.}

\prose{รูปํ อนิจฺจโต อนุปสฺสนฏฺเฐน วิปสฺสนา, รูปํ ทุกฺขโต อนุปสฺสนฏฺเฐน วิปสฺสนา, รูปํ อนตฺตโต อนุปสฺสนฏฺเฐน วิปสฺสนา. ตตฺถ ชาตานํ ธมฺมานญฺจ โวสคฺคา\-รมฺมณตา จิตฺตสฺส เอกคฺคตา อวิกฺเขโป สมาธิ. อิติ ปฐมํ วิปสฺสนา, ปจฺฉา สมโถ. เตน วุจฺจติ “วิปสฺสนา\-ปุพฺพงฺคมํ สมถํ ภาเวตี”ติ. ภาเวตีติ จตสฺโส ภาวนา, อาเสวนฏฺเฐน ภาวนา ฯเปฯ. มคฺโค สญฺชายตีติ กถํ มคฺโค สญฺชายติ ฯเปฯ เอวํ มคฺโค สญฺชายติ. เอวํ สญฺโญชนานิ ปหียนฺติ, อนุสยา พฺยนฺตีโหนฺติ.}

\prose{เวทนํ ฯเปฯ. สญฺญํ. สงฺขาเร. วิญฺญาณํ. จกฺขุํ ฯเปฯ. ชรามรณํ อนิจฺจโต อนุปสฺสนฏฺเฐน วิปสฺสนา, ชรามรณํ ทุกฺขโต ฯเปฯ. อนตฺตโต อนุปสฺสนฏฺเฐน วิปสฺสนา. ตตฺถ ชาตานํ ธมฺมานญฺจ โวสคฺคา\-รมฺมณตา จิตฺตสฺส เอกคฺคตา อวิกฺเขโป สมาธิ. อิติ ปฐมํ วิปสฺสนา, ปจฺฉา สมโถ. เตน วุจฺจติ “วิปสฺสนา\-ปุพฺพงฺคมํ สมถํ ภาเวตี”ติ.  ภาเวตีติ จตสฺโส ภาวนา, อาเสวนฏฺเฐน ภาวนา ฯเปฯ. มคฺโค สญฺชายตีติ \csromanfolio{288}กถํ มคฺโค สญฺชายติ ฯเปฯ เอวํ มคฺโค สญฺชายติ. เอวํ สญฺโญชนานิ ปหียนฺติ, อนุสยา พฺยนฺตีโหนฺติ. เอวํ วิปสฺสนา\-ปุพฺพงฺคมํ สมถํ ภาเวติ.}

\proseitem{4}{กถํ สมถ\-วิปสฺสนํ ยุคนทฺธํ ภาเวติ? โสฬสหิ อากาเรหิ สมถ\-วิปสฺสนํ ยุคนทฺธํ ภาเวติ, อารมฺมณฏฺเฐน โคจรฏฺเฐน ปหานฏฺเฐน ปริจฺจาค\-ฏฺเฐน วุฏฺฐานฏฺเฐน วิวฏฺฏ\-นฏฺเฐน สนฺตฏฺเฐน ปณีตฏฺเฐน วิมุตฺตฏฺเฐน อนาสวฏฺเฐน ตรณฏฺเฐน อนิมิตฺตฏฺเฐน อปฺปณิหิต\-ฏฺเฐน สุญฺญตฏฺเฐน เอกรสฏฺเฐน อนติวตฺตน\-ฏฺเฐน ยุคนทฺธ\-ฏฺเฐน.}

\prose{กถํ อารมฺมณฏฺเฐน สมถ\-วิปสฺสนํ ยุคนทฺธํ ภาเวติ? อุทฺธจฺจํ ปชหโต จิตฺตสฺส เอกคฺคตา อวิกฺเขโป สมาธิ นิโรธารมฺมโณ, อวิชฺชํ ปชหโต อนุปสฺสนฏฺเฐน วิปสฺสนา นิโรธารมฺมณา. อิติ อารมฺมณฏฺเฐน สมถ\-วิปสฺสนา เอกรสา โหนฺติ, ยุคนทฺธา โหนฺติ, อญฺญมญฺญํ นาติวตฺตนฺตีติ. เตน วุจฺจติ “อารมฺมณฏฺเฐน สมถ\-วิปสฺสนํ ยุคนทฺธํ ภาเวตี”ติ. ภาเวตีติ จตสฺโส ภาวนา. อาเสวนฏฺเฐน ภาวนา ฯเปฯ. มคฺโค สญฺชายตีติ กถํ มคฺโค สญฺชายติ ฯเปฯ เอวํ มคฺโค สญฺชายติ. เอวํ สญฺโญชนานิ ปหียนฺติ, อนุสยา พฺยนฺตีโหนฺติ. เอวํ อารมฺมณฏฺเฐน สมถ\-วิปสฺสนํ ยุคนทฺธํ ภาเวติ. (1)}

\prose{กถํ โคจรฏฺเฐน สมถ\-วิปสฺสนํ ยุคนทฺธํ ภาเวติ? อุทฺธจฺจํ ปชหโต จิตฺตสฺส เอกคฺคตา อวิกฺเขโป สมาธิ นิโรธโคจโร, อวิชฺชํ ปชหโต อนุปสฺสนฏฺเฐน วิปสฺสนา นิโรธโคจรา. อิติ โคจรฏฺเฐน สมถ\-วิปสฺสนา เอกรสา โหนฺติ, ยุคนทฺธา โหนฺติ, อญฺญมญฺญํ นาติวตฺตนฺตีติ. เตน วุจฺจติ “โคจรฏฺเฐน สมถ\-วิปสฺสนํ ยุคนทฺธํ ภาเวตี”ติ. (2)}

\prose{กถํ ปหานฏฺเฐน สมถ\-วิปสฺสนํ ยุคนทฺธํ ภาเวติ? อุทฺธจฺจสหคต\-กิเลเส จ ขนฺเธ จ ปชหโต จิตฺตสฺส เอกคฺคตา อวิกฺเขโป สมาธิ นิโรธโคจโร, อวิชฺชา\-สหคต\-กิเลเส จ ขนฺเธ จ ปชหโต อนุปสฺสนฏฺเฐน วิปสฺสนา นิโรธโคจรา. อิติ ปหานฏฺเฐน สมถ\-วิปสฺสนา เอกรสา โหนฺติ, ยุคนทฺธา โหนฺติ, อญฺญมญฺญํ นาติวตฺตนฺตีติ. เตน วุจฺจติ “ปหานฏฺเฐน สมถ\-วิปสฺสนํ ยุคนทฺธํ ภาเวตี”ติ. (3)}

\prose{กถํ ปริจฺจาค\-ฏฺเฐน สมถ\-วิปสฺสนํ ยุคนทฺธํ ภาเวติ? อุทฺธจฺจสหคต\-กิเลเส จ ขนฺเธ จ ปริจฺจชโต จิตฺตสฺส เอกคฺคตา อวิกฺเขโป สมาธิ นิโรธโคจโร, อวิชฺชา\-สหคต\-กิเลเส จ ขนฺเธ จ ปริจฺจชโต}

\proseitem{5}{\csromanfolio{289}อนุปสฺสนฏฺเฐน วิปสฺสนา นิโรธโคจรา. อิติ ปริจฺจาค\-ฏฺเฐน สมถ\-วิปสฺสนา เอกรสา โหนฺติ, ยุคนทฺธา โหนฺติ, อญฺญมญฺญํ นาติวตฺตนฺตีติ. เตน วุจฺจติ “ปริจฺจาค\-ฏฺเฐน สมถ\-วิปสฺสนํ ยุคนทฺธํ ภาเวตี”ติ. (4)}

\prose{กถํ วุฏฺฐานฏฺเฐน สมถ\-วิปสฺสนํ ยุคนทฺธํ ภาเวติ? อุทฺธจฺจสหคต\-กิเลเสหิ จ ขนฺเธหิ จ วุฏฺฐหโต จิตฺตสฺส เอกคฺคตา อวิกฺเขโป สมาธิ นิโรธโคจโร, อวิชฺชา\-สหคต\-กิเลเสหิ จ ขนฺเธหิ จ วุฏฺฐหโต อนุปสฺสนฏฺเฐน วิปสฺสนา นิโรธโคจรา. อิติ วุฏฺฐานฏฺเฐน สมถ\-วิปสฺสนา เอกรสา โหนฺติ, ยุคนทฺธา โหนฺติ, อญฺญมญฺญํ นาติวตฺตนฺตีติ. เตน วุจฺจติ “วุฏฺฐานฏฺเฐน สมถ\-วิปสฺสนํ ยุคนทฺธํ ภาเวตี”ติ. (5)}

\prose{กถํ วิวฏฺฏ\-นฏฺเฐน สมถ\-วิปสฺสนํ ยุคนทฺธํ ภาเวติ? อุทฺธจฺจสหคต\-กิเลเสหิ จ ขนฺเธหิ จ วิวฏฺฏโต จิตฺตสฺส เอกคฺคตา อวิกฺเขโป สมาธิ นิโรธโคจโร, อวิชฺชา\-สหคต\-กิเลเสหิ จ ขนฺเธหิ จ วิวฏฺฏโต อนุปสฺสนฏฺเฐน วิปสฺสนา นิโรธโคจรา. อิติ วิวฏฺฏ\-นฏฺเฐน สมถ\-วิปสฺสนา เอกรสา โหนฺติ, ยุคนทฺธา โหนฺติ, อญฺญมญฺญํ นาติวตฺตนฺตีติ. เตน วุจฺจติ “วิวฏฺฏ\-นฏฺเฐน สมถ\-วิปสฺสนํ ยุคนทฺธํ ภาเวตี”ติ. (6)}

\prose{กถํ สนฺตฏฺเฐน สมถ\-วิปสฺสนํ ยุคนทฺธํ ภาเวติ? อุทฺธจฺจํ ปชหโต จิตฺตสฺส เอกคฺคตา อวิกฺเขโป สมาธิ สนฺโต โหติ นิโรธโคจโร, อวิชฺชํ ปชหโต อนุปสฺสนฏฺเฐน วิปสฺสนา สนฺตา โหติ นิโรธโคจรา. อิติ สนฺตฏฺเฐน สมถ\-วิปสฺสนา เอกรสา โหนฺติ, ยุคนทฺธา โหนฺติ, อญฺญมญฺญํ นาติวตฺตนฺตีติ. เตน วุจฺจติ “สนฺตฏฺเฐน สมถ\-วิปสฺสนํ ยุคนทฺธํ ภาเวตี”ติ. (7)}

\prose{กถํ ปณีตฏฺเฐน สมถ\-วิปสฺสนํ ยุคนทฺธํ ภาเวติ? อุทฺธจฺจํ ปชหโต จิตฺตสฺส เอกคฺคตา อวิกฺเขโป สมาธิ ปณีโต โหติ นิโรธโคจโร, อวิชฺชํ ปชหโต อนุปสฺสนฏฺเฐน วิปสฺสนา ปณีตา โหติ นิโรธโคจรา. อิติ ปณีตฏฺเฐน สมถ\-วิปสฺสนา เอกรสา โหนฺติ, ยุคนทฺธา โหนฺติ, อญฺญมญฺญํ นาติวตฺตนฺตีติ. เตน วุจฺจติ “ปณีตฏฺเฐน สมถ\-วิปสฺสนํ ยุคนทฺธํ ภาเวตี”ติ. (8)}

\prose{กถํ วิมุตฺตฏฺเฐน สมถ\-วิปสฺสนํ ยุคนทฺธํ ภาเวติ? อุทฺธจฺจํ ปชหโต จิตฺตสฺส เอกคฺคตา อวิกฺเขโป สมาธิ กามาสวา วิมุตฺโต โหติ นิโรธโคจโร, อวิชฺชํ ปชหโต อนุปสฺสนฏฺเฐน วิปสฺสนา อวิชฺชาสวา \csromanfolio{290}วิมุตฺตา โหติ นิโรธโคจรา. อิติ ราควิราคา เจโตวิมุตฺติ, อวิชฺชาวิราคา ปญฺญาวิมุตฺติ วิมุตฺตฏฺเฐน สมถ\-วิปสฺสนา เอกรสา โหนฺติ, ยุคนทฺธา โหนฺติ, อญฺญมญฺญํ นาติวตฺตนฺตีติ. เตน วุจฺจติ “วิมุตฺตฏฺเฐน สมถ\-วิปสฺสนํ ยุคนทฺธํ ภาเวตี”ติ. (9)}

\proseitem{1}{กถํ อนาสวฏฺเฐน สมถ\-วิปสฺสนํ ยุคนทฺธํ ภาเวติ? อุทฺธจฺจํ ปชหโต จิตฺตสฺส เอกคฺคตา อวิกฺเขโป สมาธิ กามาสเวน อนาสโว โหติ นิโรธโคจโร, อวิชฺชํ ปชหโต อนุปสฺสนฏฺเฐน วิปสฺสนา อวิชฺชาสเวน อนาสวา โหติ นิโรธโคจรา. อิติ อนาสวฏฺเฐน สมถ\-วิปสฺสนา เอกรสา โหนฺติ, ยุคนทฺธา โหนฺติ, อญฺญมญฺญํ นาติวตฺตนฺตีติ. เตน วุจฺจติ “อนาสวฏฺเฐน สมถ\-วิปสฺสนํ ยุคนทฺธํ ภาเวตี”ติ. (10)}

\prose{กถํ ตรณฏฺเฐน สมถ\-วิปสฺสนํ ยุคนทฺธํ ภาเวติ? อุทฺธจฺจสหคต\-กิเลเส จ ขนฺเธ จ ตรโต จิตฺตสฺส เอกคฺคตา อวิกฺเขโป สมาธิ นิโรธโคจโร, อวิชฺชา\-สหคต\-กิเลเส จ ขนฺเธ จ ตรโต อนุปสฺสนฏฺเฐน วิปสฺสนา นิโรธโคจรา. อิติ ตรณฏฺเฐน สมถ\-วิปสฺสนา เอกรสา โหนฺติ, ยุคนทฺธา โหนฺติ, อญฺญมญฺญํ นาติวตฺตนฺตีติ. เตน วุจฺจติ “ตรณฏฺเฐน สมถ\-วิปสฺสนํ ยุคนทฺธํ ภาเวตี”ติ. (11)}

\prose{กถํ อนิมิตฺตฏฺเฐน สมถ\-วิปสฺสนํ ยุคนทฺธํ ภาเวติ? อุทฺธจฺจํ ปชหโต จิตฺตสฺส เอกคฺคตา อวิกฺเขโป สมาธิ สพฺพนิมิตฺเตหิ อนิมิตฺโต โหติ นิโรธโคจโร, อวิชฺชํ ปชหโต อนุปสฺสนฏฺเฐน วิปสฺสนา สพฺพนิมิตฺเตหิ อนิมิตฺตา โหติ นิโรธโคจรา. อิติ อนิมิตฺตฏฺเฐน สมถ\-วิปสฺสนา เอกรสา โหนฺติ, ยุคนทฺธา โหนฺติ, อญฺญมญฺญํ นาติวตฺตนฺตีติ. เตน วุจฺจติ “อนิมิตฺตฏฺเฐน สมถ\-วิปสฺสนํ ยุคนทฺธํ ภาเวตี”ติ. (12)}

\prose{กถํ อปฺปณิหิต\-ฏฺเฐน สมถ\-วิปสฺสนํ ยุคนทฺธํ ภาเวติ? อุทฺธจฺจํ ปชหโต จิตฺตสฺส เอกคฺคตา อวิกฺเขโป สมาธิ สพฺพปณิธีหิ อปฺปณิหิโต โหติ นิโรธโคจโร, อวิชฺชํ ปชหโต อนุปสฺสนฏฺเฐน วิปสฺสนา สพฺพปณิธีหิ อปฺปณิหิตา โหติ นิโรธโคจรา. อิติ อปฺปณิหิต\-ฏฺเฐน สมถ\-วิปสฺสนา เอกรสา โหนฺติ, ยุคนทฺธา โหนฺติ, อญฺญมญฺญํ นาติวตฺตนฺตีติ. เตน วุจฺจติ “อปฺปณิหิต\-ฏฺเฐน สมถ\-วิปสฺสนํ ยุคนทฺธํ ภาเวตี”ติ. (13)}

\prose{\csromanfolio{291}กถํ สุญฺญตฏฺเฐน สมถ\-วิปสฺสนํ ยุคนทฺธํ ภาเวติ? อุทฺธจฺจํ ปชหโต จิตฺตสฺส เอกคฺคตา อวิกฺเขโป สมาธิ สพฺพาภิ\-นิเวเสหิ สุญฺโญ โหติ นิโรธโคจโร, อวิชฺชํ ปชหโต อนุปสฺสนฏฺเฐน วิปสฺสนา สพฺพาภิ\-นิเวเสหิ สุญฺญา โหติ นิโรธโคจรา. อิติ สุญฺญตฏฺเฐน สมถ\-วิปสฺสนา เอกรสา โหนฺติ, ยุคนทฺธา โหนฺติ, อญฺญมญฺญํ นาติวตฺตนฺตีติ. เตน วุจฺจติ “สุญฺญตฏฺเฐน สมถ\-วิปสฺสนํ ยุคนทฺธํ พฺหาเวตี”ติ. ภาเวตีติ จตสฺโส ภาวนา, อาเสวนฏฺเฐน ภาวนา ฯเปฯ. มฺอคฺโค สญฺชายตีติ กถํ มคฺโค สญฺชายติ ฯเปฯ เอวํ มคฺโค สญฺชายติ. เอวํ สญฺโญชนานิ ปหียนฺติ, อนุสยา พฺยนฺตีโหนฺติ. เอวํ สุญฺญตฏฺเฐน สมถ\-วิปสฺสนํ ยุคนทฺธํ ภาเวติ. อิเมหิ โสฬสหิ อากาเรหิ สมถ\-วิปสฺสนํ ยุคนทฺธํ ภาเวติ, เอวํ สมถ\-วิปสฺสนํ ยุคนทฺธํ ภาเวติ. (14) สุตฺตนฺต\-นิทฺเทโส.}
\csromansectionrule

\csromanmatikaanchor{mtk.87}
\csromantocmarkref{subsection}{2. ธมฺมุทฺธจฺจวารนิทฺเทส}{mtk.87}
\bookmark[dest={mtk.87},level=2]{2. ธมฺมุทฺธจฺจวารนิทฺเทส}
\csromanheader{2}{2. ธมฺมุทฺธจฺจ\-วาร\-นิทฺเทส}
\proseitem{6}{กถํ ธมฺมฺอุทฺธจฺจวิคฺคหิตํ มานสํ โหติ? อนิจฺจโต มนสิกโรโต โอภาโส อุปฺปชฺชติ, “โอภาโส ธมฺโม”ติ โอภาสํ อาวชฺชติ. ตโต วิกฺเขโป อุทฺธจฺจํ. เตน อุทฺธจฺเจน วิคฺคหิต\-มานโส อนิจฺจโต อุปฏฺฐานํ ยถาภูตํ นปฺปชานาติ, ทุกฺขโต อุปฏฺฐานํ ยถาภูตํ นปฺปชานาติ, อนตฺตโต อุปฏฺฐานํ ยถาภูตํ นปฺปชานาติ, เตน วุจฺจติ ธมฺมุทฺธจฺจวิคฺคหิต\-มานโส. โหติ โส สมโย, ยํ ตํ จิตฺตํ อชฺฌตฺตเมว สนฺติฏฺฐติ สนฺนิสีทติ เอโกทิ โหติ สมาธิยติ. ตสฺส มคฺโค สญฺชายตีติ กถํ มคฺโค สญฺชายติ ฯเปฯ เอวํ มคฺโค สญฺชายติ, เอวํ สญฺโญชนานิ ปหียนฺติ, อนุสยา พฺยนฺตีโหนฺติ.}

\prose{อนิจฺจโต มนสิกโรโต ญาณํ อุปฺปชฺชติ, ปีติ อุปฺปชฺชติ, ปสฺสทฺธิ อุปฺปชฺชติ, สุขํ อุปฺปชฺชติ, อธิโมกฺโข อุปฺปชฺชติ, ปคฺคโห อุปฺปชฺชติ, อุปฏฺฐานํ อุปฺปชฺชติ, อุเปกฺขา อุปฺปชฺชติ, นิกนฺติ อุปฺปชฺชติ, “นิกนฺติ ธมฺโม”ติ นิกนฺติํ อาวชฺชติ. ตโต วิกฺเขโป อุทฺธจฺจํ. เตน อุทฺธจฺเจน วิคฺคหิต\-มานโส อนิจฺจโต อุปฏฺฐานํ ยถาภูตํ นปฺปชานาติ, ทุกฺขโต อุปฏฺฐานํ ยถาภูตํ นปฺปชานาติ, อนตฺตโต อุปฏฺฐานํ ยถาภูตํ นปฺปชานาติ, เตน วุจฺจติ \csromanfolio{292}ธมฺมุทฺธจฺจวิคฺคหิต\-มานโส. โหติ โส สมโย, ยํ ตํ จิตฺตํ อชฺฌตฺตเมว สนฺติฏฺฐติ สนฺนิสีทติ เอโกทิ โหติ สมาธิยติ. ตสฺส มคฺโค สญฺชายตีติ กถํ มคฺโค สญฺชายติ ฯเปฯ เอวํ มคฺโค สญฺชายติ, เอวํ สญฺโญชนานิ ปหียนฺติ, อนุสยา พฺยนฺตีโหนฺติ.}

\prose{ทุกฺขโต มนสิกโรโต ฯเปฯ. อนตฺตโต มนสิกโรโต โอภาโส อุปฺปชฺชติ ฯเปฯ ญาณํ อุปฺปชฺชติ, ปีติ อุปฺปชฺชติ, ปสฺสทฺธิ อุปฺปชฺชติ, สุขํ อุปฺปชฺชติ, อธิโมกฺโข อุปฺปชฺชติ, ปคฺคโห อุปฺปชฺชติ, อุปฏฺฐานํ อุปฺปชฺชติ, อุเปกฺขา อุปฺปชฺชติ, นิกนฺติ อุปฺปชฺชติ, นิกนฺติ ธมฺโมติ นิกนฺติํ อาวชฺชติ. ตโต วิกฺเขโป อุทฺธจฺจํ. เตน อุทฺธจฺเจน วิคฺคหิต\-มานโส อนตฺตโต อุปฏฺฐานํ, อนิจฺจโต อุปฏฺฐานํ, ทุกฺขโต อุปฏฺฐานํ ยถาภูตํ นปฺปชานาติ, เตน วุจฺจติ ธมฺมุทฺธจฺจวิคฺคหิต\-มานโส ฯเปฯ เอวํ สญฺโญชนานิ ปหียนฺติ, อนุสยา พฺยนฺตีโหนฺติ.}

\prose{รูปํ อนิจฺจโต มนสิกโรโต ฯเปฯ รูปํ ทุกฺขโต มนสิกโรโต. รูปํ อนตฺตโต มนสิกโรโต. เวทนํ ฯเปฯ. สญฺญํ. สงฺขาเร, วิญฺญาณํ. จกฺขุํ ฯเปฯ. ชรามรณํ อนิจฺจโต มนสิกโรโต ฯเปฯ ชรามรณํ ทุกฺขโต มนสิกโรโต. ชรามรณํ อนตฺตโต มนสิกโรโต โอภาโส อุปฺปชฺชติ ฯเปฯ ญาณํ อุปฺปชฺชติ, ปีติ อุปฺปชฺชติ, ปสฺสทฺธิ อุปฺปชฺชติ, สุขํ อุปฺปชฺชติ, อธิโมกฺโข อุปฺปชฺชติ, ปคฺคโห อุปฺปชฺชติ, อุปฏฺฐานํ อุปฺปชฺชติ, อุเปกฺโข อุปฺปชฺชติ, นิกนฺติ อุปฺปชฺชติ, “นิกนฺติ ธมฺโม”ติ นิกนฺติํ อาวชฺชติ. ตโต วิกฺเขโป อุทฺธจฺจํ. เตน อุทฺธจฺเจน วิคฺคหิต\-มานโส. ชรามรณํ อนตฺตโต อุปฏฺฐานํ ยถาภูตํ นปฺปชานาติ. ชรามรณํ อนิจฺจโต อุปฏฺฐานํ ยถาภูตํ นปฺปชานาติ, ชรามรณํ ทุกฺขโต อุปฏฺฐานํ ยถาภูตํ นปฺปชานาติ. เตน วุจฺจติ ธมฺมุทฺธจฺจวิคฺคหิต\-มานโส. โหติ โส สมโย, ยํ ตํ จิตฺตํ อชฺฌตฺตเมว สนฺติฏฺฐติ สนฺนิสีทติ เอโกทิ โหติ สมาธิยติ. ตสฺส มคฺโค สญฺชายตีติ กถํ มคฺโค สญฺชายติ ฯเปฯ เอวํ มคฺโค สญฺชายติ. เอวํ สญฺโญชนานิ ปหียนฺติ, อนุสยา พฺยนฺตีโหนฺติ. เอวํ ธมฺมุ\-ทฺธจฺจ\-วิคฺคหิตํ มานสํ โหติ.}

\setlength{\csromangathapagewidth}{0pt}
\setlength{\csromangathawakwidth}{0pt}
\csromangathameasuregroup{โอภาเส เจว ญาเณ จ, \\ ปสฺสทฺธิยา สุเข เจว, \\ อธิโมกฺเข จ ปคฺคาเห, \\ อุเปกฺขาวชฺชนาย เจว, \\ อิมานิ ทส ฐานานิ, \\ ธมฺมุทฺธจฺจกุสโล โหติ, \\ วิกฺขิปฺปติ เจว กิลิสฺสติ จ, \\ วิกฺขิปติ น กิลิสฺสติ, \\ วิกฺขิปติ น กิลิสฺสติ, \\ น จ วิกฺขิปเต จิตฺตํ น กิลิสฺสติ,}{\csromangathabat{โอภาเส เจว ญาเณ จ,}{ปีติยา จ วิกมฺปติ.} \\ \csromangathabat{ปสฺสทฺธิยา สุเข เจว,}{เยหิ จิตฺตํ ปเวธติ.} \\ \csromangathabat{อธิโมกฺเข จ ปคฺคาเห,}{อุปฏฺฐาเน จ กมฺปติ.} \\ \csromangathabat{อุเปกฺขาวชฺชนาย เจว,}{อุเปกฺขาย จ นิกนฺติยา.} \\ \csromangathabat{อิมานิ ทส ฐานานิ,}{ปญฺญา ยสฺส ปริจฺจิตา.} \\ \csromangathabat{ธมฺมุทฺธจฺจกุสโล โหติ,}{น จ สมฺโมห คจฺฉติ.} \\ \csromangathabat{วิกฺขิปฺปติ เจว กิลิสฺสติ จ,}{จวติ จิตฺตภาวนา.} \\ \csromangathabat{วิกฺขิปติ น กิลิสฺสติ,}{ภาวนา ปริหายติ.} \\ \csromangathabat{วิกฺขิปติ น กิลิสฺสติ,}{ภาวนา น ปริหายติ.} \\ \csromangathabat{น จ วิกฺขิปเต จิตฺตํ น กิลิสฺสติ,}{น จวติ จิตฺตภาวนา.}}
\csromangathasetleft{โอภาเส เจว ญาเณ จ, \\ ปสฺสทฺธิยา สุเข เจว, \\ อธิโมกฺเข จ ปคฺคาเห, \\ อุเปกฺขาวชฺชนาย เจว, \\ อิมานิ ทส ฐานานิ, \\ ธมฺมุทฺธจฺจกุสโล โหติ, \\ วิกฺขิปฺปติ เจว กิลิสฺสติ จ, \\ วิกฺขิปติ น กิลิสฺสติ, \\ วิกฺขิปติ น กิลิสฺสติ, \\ น จ วิกฺขิปเต จิตฺตํ น กิลิสฺสติ,}
\csromangathagroup{\csromangathastanza{\csromangathaitembat{6}{โอภาเส เจว ญาเณ จ,}{ปีติยา จ วิกมฺปติ.} \\ \csromangathacontbat{ปสฺสทฺธิยา สุเข เจว,}{เยหิ จิตฺตํ ปเวธติ.}}\csromangathastanza{\csromanfolio{293}\csromangathabat{อธิโมกฺเข จ ปคฺคาเห,}{อุปฏฺฐาเน จ กมฺปติ.} \\ \csromangathabat{อุเปกฺขาวชฺชนาย เจว,}{อุเปกฺขาย จ นิกนฺติยา.}}\csromangathastanza{\csromangathabat{อิมานิ ทส ฐานานิ,}{ปญฺญา ยสฺส ปริจฺจิตา.} \\ \csromangathabat{ธมฺมุทฺธจฺจกุสโล โหติ,}{น จ สมฺโมห คจฺฉติ.}}\csromangathastanza{\csromangathabat{วิกฺขิปฺปติ เจว กิลิสฺสติ จ,}{จวติ จิตฺตภาวนา.} \\ \csromangathabat{วิกฺขิปติ น กิลิสฺสติ,}{ภาวนา ปริหายติ.}}\csromangathastanza{\csromangathabat{วิกฺขิปติ น กิลิสฺสติ,}{ภาวนา น ปริหายติ.} \\ \csromangathabat{น จ วิกฺขิปเต จิตฺตํ น กิลิสฺสติ,}{น จวติ จิตฺตภาวนา.}}}

\prosewithcloserruled{อิเมหิ จตูหิ ฐาเนหิ จิตฺตสฺส สงฺเขป\-วิกฺเขป\-วิคฺคหิตํ\csromansharedfootnote{a4fd89ca45460c7a}{สงฺเขปํ วิกฺเขปํ วิคฺคหิตํ (สฺยา)} ทส ฐาเน สมฺปชานาตีติ.}{ยุคนทฺธ\-กถา นิฏฺฐิตา.}
\proseitem{8}{ปุริมนิทานํ.  จตฺตาริมานิ ภิกฺขเว\csromansharedfootnote{4ae3b922c9571cb4}{สํ 3. 377 ปิฏฺเฐปิ.} ตถานิ อวิตถานิ อนญฺญถานิ. กตมานิ จตฺตาริ\csromandash{}“อิทํ ทุกฺขนฺ”ติ ภิกฺขเว ตถเมตํ อวิตถเมตํ อนญฺญถเมตํ, “อยํ ทุกฺขสมุทโย”ติ ตถเมตํ อวิตถเมตํ อนญฺญถเมตํ, “อยํ ทุกฺขนิโรโธ”ติ ตถเมตํ อวิตถเมตํ อนญฺญถเมตํ, “อยํ ทุกฺขนิโรธ\-คามินี ปฏิปทา”ติ ตถเมตํ อวิตถเมตํ อนญฺญถเมตํ. อิมานิ โข ภิกฺขเว จตฺตาริ ตถานิ อวิตถานิ อนญฺญถานิ.}

\csromanmatikaanchor{mtk.88}
\csromanmatikaanchor{mtk.89}
\csromantocmarknopage{section}{2. สจฺจกถา}{mtk.88}
\bookmark[dest={mtk.88},level=1]{2. สจฺจกถา}
\csromantocmarkref{subsection}{1. ปฐมสุตฺตนฺตนิทฺเทส}{mtk.89}
\bookmark[dest={mtk.89},level=2]{1. ปฐมสุตฺตนฺตนิทฺเทส}
\csromanheader{2}{1. ปฐม\-สุตฺตนฺต\-นิทฺเทส}
\prose{กถํ ทุกฺขํ ตถฏฺเฐน สจฺจํ? จตฺตาโร ทุกฺขสฺส ทุกฺขฏฺฐา ตถา อวิตถา อนญฺญถา, ทุกฺขสฺส ปีฬนฏฺโฐ สงฺขตฏฺโฐ สนฺตาปฏฺโฐ วิปริณามฏฺโฐ, อิเม จตฺตาโร ทุกฺขสฺส ทุกฺขฏฺฐา ตถา อวิตถา อนญฺญถา, เอวํ ทุกฺขํ ตถฏฺเฐน สจฺจํ.}

\prose{กถํ สมุทโย ตถฏฺเฐน สจฺจํ? จตฺตาโร สมุทยสฺส สมุทยฏฺฐา ตถา อวิตถา อนญฺญถา, สมุทยสฺส อายูหนฏฺโฐ นิทานฏฺโฐ}

\prose{\csromanfolio{294}สํโยคฏฺโฐ ปลิโพธฏฺโฐ, อิเม จตฺตาโร สมุทยสฺส สมุทยฏฺฐา ตถา อวิตถา อนญฺญถา, เอวํ สมุทโย ตถฏฺเฐน สจฺจํ.}

\prose{กถํ นิโรโธ ตถฏฺเฐน สจฺจํ? จตฺตาโร นิโรธสฺส นิโรธฏฺฐา ตถา อวิตถา อนญฺญถา, นิโรธสฺส นิสฺสรณฏฺโฐ วิเวกฏฺโฐ อสงฺขตฏฺโฐ อมตฏฺโฐ, อิเม จตฺตาโร นิโรธสฺส นิโรธฏฺฐา ตถา อวิตถา อนญฺญถา, เอวํ นิโรโธ ตถฏฺเฐน สจฺจํ.}

\prose{กถํ มคฺโค ตถฏฺเฐน สจฺจํ? จตฺตาโร มคฺคสฺส มคฺคฏฺฐา ตถา อวิตถา อนญฺญถา, มคฺคสฺส นิยฺยานฏฺโฐ เหตุฏฺโฐ ทสฺสนฏฺโฐ อาธิปเตยฺย\-ฏฺโฐ, อิเม จตฺตาโร มคฺคสฺส มคฺคฏฺฐา ตถา อวิตถา อนญฺญถา, เอวํ มคฺโค ตถฏฺเฐน สจฺจํ.}

\proseitem{9}{กติหากาเรหิ จตฺตาริ สจฺจานิ เอก\-ปฺปฏิเวธานิ? จตูหากาเรหิ จตฺตาริ สจฺจานิ เอก\-ปฺปฏิเวธานิ ตถฏฺเฐน อนตฺตฏฺเฐน สจฺจฏฺเฐน ปฏิเวธ\-ฏฺเฐน, อิเมหิ จตูหากาเรหิ จตฺตาริ สจฺจานิ เอกสงฺคหิตานิ, ยํ เอกสงฺคหิตํ ตํ เอกตฺตํ, เอกตฺตํ เอเกน ญาเณน ปฏิวิชฺฌตีติ จตฺตาริ สจฺจานิ เอก\-ปฺปฏิเวธานิ.}

\prose{กถํ ตถฏฺเฐน จตฺตาริ สจฺจานิ เอก\-ปฺปฏิเวธานิ? จตูหากาเรหิ ตถฏฺเฐน จตฺตาริ สจฺจานิ เอก\-ปฺปฏิเวธานิ ทุกฺขสฺส ทุกฺขฏฺโฐ ตถฏฺโฐ, สมุทยสฺส สมุทยฏฺโฐ ตถฏฺโฐ, นิโรธสฺส นิโรธฏฺโฐ ตถฏฺโฐ, มคฺคสฺส มคฺคฏฺโฐ ตถฏฺโฐ, อิเมหิ จตูหากาเรหิ ตถฏฺเฐน จตฺตาริ สจฺจานิ เอกสงฺคหิตานิ, ยํ เอกสงฺคหิตํ ตํ เอกตฺตํ, เอกตฺตํ เอเกน ญาเณน ปฏิวิชฺฌตีติ จตฺตาริ สจฺจานิ เอก\-ปฺปฏิเวธานิ.}

\prose{กถํ อนตฺตฏฺเฐน จตฺตาริ สจฺจานิ เอก\-ปฺปฏิเวธานิ? จตูหากาเรหิ อนตฺตฏฺเฐน จตฺตาริ สจฺจานิ เอก\-ปฺปฏิเวธานิ, ทุกฺขสฺส ทุกฺขฏฺโฐ อนตฺตฏฺโฐ, สมุทยสฺส สมุทยฏฺโฐ อนตฺตฏฺโฐ, นิโรธสฺส นิโรธฏฺโฐ อนตฺตฏฺโฐ, มคฺคสฺส มคฺคฏฺโฐ อนตฺตฏฺโฐ, อิเมหิ จตูหากาเรหิ อนตฺตฏฺเฐน จตฺตาริ สจฺจานิ เอกสงฺคหิตานิ, ยํ เอกสงฺคหิตํ ตํ เอกตฺตํ, เอกตฺตํ เอเกน ญาเณน ปฏิวิชฺฌตีติ จตฺตาริ สจฺจานิ เอก\-ปฺปฏิเวธานิ.}

\prose{กถํ สจฺจฏฺเฐน จตฺตาริ สจฺจานิ เอก\-ปฺปฏิเวธานิ? จตูหากาเรหิ สจฺจฏฺเฐน จตฺตาริ สจฺจานิ เอก\-ปฺปฏิเวธานิ. ทุกฺขสฺส ทุกฺขฏฺโฐ สจฺจฏฺโฐ. \csromanfolio{295}สมุทยสฺส สมุทยฏฺโฐ สจฺจฏฺโฐ. นิโรธสฺส นิโรธฏฺโฐ สจฺจฏฺโฐ. มคฺคสฺส มคฺคฏฺโฐ สจฺจฏฺโฐ. อิเมหิ จตูหากาเรหิ สจฺจฏฺเฐน จตฺตาริ สจฺจานิ เอกสงฺคหิตานิ. ยํ เอกสงฺคหิตํ ตํ เอกตฺตํ, เอกตฺตํ เอเกน ญาเณน ปฏิวิชฺฌตีติ จตฺตาริ สจฺจานิ เอก\-ปฺปฏิเวธานิ.}

\prose{กถํ ปฏิเวธ\-ฏฺเฐน จตฺตาริ สจฺจานิ เอก\-ปฺปฏิเวธานิ? จตูหากาเรหิ ปฏิเวธ\-ฏฺเฐน จตฺตาริ สจฺจานิ เอก\-ปฺปฏิเวธานิ. ทุกฺขสฺส ทุกฺขฏฺโฐ ปฏิเวธฏฺโฐ สมุทยสฺส สมุทยฏฺโฐ ปฏิเวธฏฺโฐ. นิโรธสฺส นิโรธฏฺโฐ ปฏิเวธฏฺโฐ. มคฺคสฺส มคฺคฏฺโฐ ปฏิเวธฏฺโฐ. อิเมหิ จตูหากาเรหิ ปฏิเวธ\-ฏฺเฐน จตฺตาริ สจฺจานิ เอกสงฺคหิตานิ. ยํ เอกสงฺคหิตํ ตํ เอกตฺตํ, เอกตฺตํ เอเกน ญาเณน ปฏิวิชฺฌตีติ จตฺตาริ สจฺจานิ เอก\-ปฺปฏิเวธานิ.}

\proseitem{10}{กติหากาเรหิ จตฺตาริ สจฺจานิ เอก\-ปฺปฏิเวธานิ? ยํ อนิจฺจํ, ตํ ทุกฺขํ. ยํ ทุกฺขํ, ตํ อนิจฺจํ. ยํ อนิจฺจญฺจ ทุกฺขญฺจ, ตํ อนตฺตา. ยํ อนิจฺจญฺจ ทุกฺขญฺจ อนตฺตา จ, ตํ ตถํ. ยํ อนิจฺจญฺจ ทุกฺขญฺจ อนตฺตา จ ตถญฺจ, ตํ สจฺจํ. ยํ อนิจฺจญฺจ ทุกฺขญฺจ อนตฺตา จ ตถญฺจ สจฺจญฺจ, ตํ เอกสงฺคหิตํ. ยํ เอกสงฺคหิตํ ตํ เอกตฺตํ, เอกตฺตํ เอเกน ญาเณน ปฏิวิชฺฌตีติ จตฺตาริ สจฺจานิ เอก\-ปฺปฏิเวธานิ.}

\prose{กติหากาเรหิ จตฺตาริ สจฺจานิ เอก\-ปฺปฏิเวธานิ? นวหากาเรหิ จตฺตาริ สจฺจานิ เอก\-ปฺปฏิเวธานิ, ตถฏฺเฐน อนตฺตฏฺเฐน สจฺจฏฺเฐน ปฏิเวธ\-ฏฺเฐน อภิญฺญฏฺเฐน ปริญฺญฏฺเฐน ปหานฏฺเฐน ภาวนฏฺเฐน สจฺฉิกิริยฏฺเฐน. อิเมหิ นวหากาเรหิ จตฺตาริ สจฺจานิ เอกสงฺคหิตานิ. ยํ เอกสงฺคหิตํ ตํ เอกตฺตํ, เอกตฺตํ เอเกน ญาเณน ปฏิวิชฺฌตีติ จตฺตาริ สจฺจานิ เอก\-ปฺปฏิเวธานิ.}

\prose{กถํ ตถฏฺเฐน จตฺตาริ สจฺจานิ เอกปฺปฏิเวธา? นวหากาเรหิ ตถฏฺเฐน จตฺตาริ สจฺจานิ เอก\-ปฺปฏิเวธานิ. ทุกฺขสฺส ทุกฺขฏฺโฐ ตถฏฺโฐ. สมุทยสฺส สมุทยฏฺโฐ ตถฏฺโฐ. นิโรธสฺส นิโรธฏฺโฐ ตถฏฺโฐ. มคฺคสฺส มคฺคฏฺโฐ ตถฏฺโฐ. อภิญฺญาย อภิญฺญฏฺโฐ ตถฏฺโฐ. ปริญฺญาย ปริญฺญฏฺโฐ ตถฏฺโฐ. ปหานสฺส ปหานฏฺโฐ ตถฏฺโฐ. ภาวนาย ภาวนฏฺโฐ ตถฏฺโฐ. สจฺฉิกิริยาย สจฺฉิกิริยฏฺโฐ ตถฏฺโฐ. อิเมหิ นวหากาเรหิ ตถฏฺเฐน จตฺตาริ สจฺจานิ เอกสงฺคหิตานิ. ยํ เอกสงฺคหิตํ ตํ เอกตฺตํ, เอกตฺตํ เอเกน ญาเณน ปฏิวิชฺฌตีติ จตฺตาริ สจฺจานิ เอก\-ปฺปฏิเวธานิ.}

\prose{\csromanfolio{296}กถํ อนตฺตฏฺเฐน สจฺจฏฺเฐน ปฏิเวธ\-ฏฺเฐน จตฺตาริ สจฺจานิ เอก\-ปฺปฏิเวธานิ? นวหากาเรหิ ปฏิเวธ\-ฏฺเฐน จตฺตาริ สจฺจานิ เอก\-ปฺปฏิเวธานิ. ทุกฺขสฺส ทุกฺขฏฺโฐ ปฏิเวธฏฺโฐ. สมุทยสฺส สมุทยฏฺโฐ ปฏิเวธฏฺโฐ. นิโรธสฺส นิโรธฏฺโฐ ปฏิเวธฏฺโฐ. มคฺคสฺส มคฺคฏฺโฐ ปฏิเวธฏฺโฐ. อภิญฺญาย อภิญฺญฏฺโฐ ปฏิเวธฏฺโฐ. ปริญฺญาย ปริญฺญฏฺโฐ ปฏิเวธฏฺโฐ. ปหานสฺส ปหานฏฺโฐ ปฏิเวธฏฺโฐ. ภาวนาย ภาวนฏฺโฐ ปฏิเวธฏฺโฐ. สจฺฉิกิริยาย สจฺฉิกิริยฏฺโฐ ปฏิเวธฏฺโฐ. อิเมหิ นวหากาเรหิ ปฏิเวธ\-ฏฺเฐน จตฺตาริ สจฺจานิ เอกสงฺคหิตานิ. ยํ เอกสงฺคหิตํ ตํ เอกตฺตํ, เอกตฺตํ เอเกน ญาเณน ปฏิวิชฺฌตีติ จตฺตาริ สจฺจานิ เอก\-ปฺปฏิเวธานิ.}

\proseitem{11}{กติหากาเรหิ จตฺตาริ สจฺจานิ เอก\-ปฺปฏิเวธานิ? ทฺวาทสหา\-กาเรหิ จตฺตาริ สจฺจานิ เอก\-ปฺปฏิเวธานิ, ตถฏฺเฐน อนตฺตฏฺเฐน สจฺจฏฺเฐน ปฏิเวธ\-ฏฺเฐน อภิชาน\-นฏฺเฐน ปริชาน\-นฏฺเฐน ธมฺมฏฺเฐน ตถฏฺเฐน ญาตฏฺเฐน สจฺฉิกิริยฏฺเฐน ผสฺสนฏฺเฐน อภิสมย\-ฏฺเฐน. อิเมหิ ทฺวาทสหิ อากาเรหิ จตฺตาริ สจฺจานิ เอกสงฺคหิตานิ. ยํ เอกสงฺคหิตํ ตํ เอกตฺตํ, เอกตฺตํ เอเกน ญาเณน ปฏิวิชฺฌตีติ จตฺตาริ สจฺจานิ เอก\-ปฺปฏิเวธานิ.}

\prose{กถํ ตถฏฺเฐน จตฺตาริ สจฺจานิ เอก\-ปฺปฏิเวธานิ? โสฬสหากาเรหิ ตถฏฺเฐน จตฺตาริ สจฺจานิ เอก\-ปฺปฏิเวธานิ. ทุกฺขสฺส ปีฬนฏฺโฐ สงฺขตฏฺโฐ สนฺตาปฏฺโฐ วิปริณามฏฺโฐ ตถฏฺโฐ. สมุทยสฺส อายูหนฏฺโฐ นิทานฏฺโฐ สํโยคฏฺโฐ ปลิโพธฏฺโฐ ตถฏฺโฐ. นิโรธสฺส นิสฺสรณฏฺโฐ วิเวกฏฺโฐ อสงฺขตฏฺโฐ อมตฏฺโฐ ตถฏฺโฐ. มคฺคสฺส นิยฺยานฏฺโฐ เหตุฏฺโฐ ทสฺสนฏฺโฐ อาธิปเตยฺย\-ฏฺโฐ ตถฏฺโฐ. อิเมหิ โสฬสหิ อากาเรหิ ตถฏฺเฐน จตฺตาริ สจฺจานิ เอกสงฺคหิตานิ. ยํ เอกสงฺคหิตํ ตํ เอกตฺตํ, เอกตฺตํ เอเกน ญาเณน ปฏิวิชฺฌตีติ จตฺตาริ สจฺจานิ เอก\-ปฺปฏิเวธานิ.}

\prose{กถํ อนตฺตฏฺเฐน ฯเปฯ. สจฺจฏฺเฐน ปฏิเวธ\-ฏฺเฐน อภิชาน\-นฏฺเฐน ปริชาน\-นฏฺเฐน ธมฺมฏฺเฐน ตถฏฺเฐน ญาตฏฺเฐน สจฺฉิกิริยฏฺเฐน ผสฺสนฏฺเฐน อภิสมย\-ฏฺเฐน จตฺตาริ สจฺจานิ เอก\-ปฺปฏิเวธานิ. โสฬสหิ อากาเรหิ อภิสมย\-ฏฺเฐน จตฺตาริ สจฺจานิ เอก\-ปฺปฏิเวธานิ. ทุกฺขสฺส ปีฬนฏฺโฐ สงฺขตฏฺโฐ สนฺตาปฏฺโฐ วิปริณามฏฺโฐ อภิสมยฏฺโฐ. สมุทยสฺส อายูหนฏฺโฐ นิทานฏฺโฐ สํโยคฏฺโฐ ปลิโพธฏฺโฐ อภิสมยฏฺโฐ. นิโรธสฺส \csromanfolio{297}นิสฺสรณฏฺโฐ วิเวกฏฺโฐ อสงฺขตฏฺโฐ อมตฏฺโฐ อภิสมยฏฺโฐ. มคฺคสฺส นิยฺยานฏฺโฐ เหตุฏฺโฐ ทสฺสนฏฺโฐ อาธิปเตยฺย\-ฏฺโฐ อภิสมยฏฺโฐ. อิเมหิ โสฬสหิ อากาเรหิ อภิสมย\-ฏฺเฐน จตฺตาริ สจฺจานิ เอกสงฺคหิตานิ. ยํ เอกสงฺคหิตํ ตํ เอกตฺตํ, เอกตฺตํ เอเกน ญาเณน ปฏิวิชฺฌตีติ จตฺตาริ สจฺจานิ เอก\-ปฺปฏิเวธานิ.}

\proseitem{12}{สจฺจานํ กติ ลกฺขณานิ? สจฺจานํ เทฺว ลกฺขณานิ สงฺขตลกฺขณญฺจ อสงฺขตลกฺขณญฺจ, สจฺจานํ อิมานิ เทฺว ลกฺขณานิ.}

\prose{สจฺจานํ กติ ลกฺขณานิ? สจฺจานํ ฉ ลกฺขณานิ. สงฺขตานํ สจฺจานํ อุปฺปาโท ปญฺญายติ, วโย ปญฺญายติ, ฐิตานํ อญฺญถตฺตํ ปญฺญายติ. อสงฺขตสฺส สจฺจสฺส น อุปฺปาโท ปญฺญายติ, น วโย ปญฺญายติ, น ฐิตสฺส อญฺญถตฺตํ ปญฺญายติ, สจฺจานํ อิมานิ ฉ ลกฺขณานิ.}

\prose{สจฺจานํ กติ ลกฺขณานิ? สจฺจานํ ทฺวาทส ลกฺขณานิ, ทุกฺขสจฺจสฺส อุปฺปาโท ปญฺญายติ, วโย ปญฺญายติ, ฐิตสฺส อญฺญถตฺตํ ปญฺญายติ, สมุทย\-สจฺจสฺส อุปฺปาโท ปญฺญายติ, วโย ปญฺญายติ, ฐิตสฺส อญฺญถตฺตํ ปญฺญายติ. มคฺคสจฺจสฺส อุปฺปาโท ปญฺญายติ, วโย ปญฺญายติ, ฐิตสฺส อญฺญถตฺตํ ปญฺญายติ. นิโรธ\-สจฺจสฺส น อุปฺปาโท ปญฺญายติ, น วโย ปญฺญายติ, น ฐิตสฺส อญฺญถตฺตํ ปญฺญายติ, สจฺจานํ อิมานิ ทฺวาทส ลกฺขณานิ.}

\prose{จตุนฺนํ สจฺจานํ กติ กุสลา, กติ อกุสลา, กติ อพฺยากตา? สมุทยสจฺจํ อกุสลํ, มคฺคสจฺจํ กุสลํ, นิโรธสจฺจํ อพฺยากตํ. ทุกฺขสจฺจํ สิยา กุสลํ, สิยา อกุสลํ, สิยา อพฺยากตํ.}

\prose{สิยา ตีณิ สจฺจานิ เอกสจฺเจน สงฺคหิตานิ. เอกสจฺจํ ตีหิ สจฺเจหิ สงฺคหิตํ วตฺถุวเสน ปริยาเยน. \textbf{สิยา}ติ กถญฺจ สิยา? ยํ ทุกฺขสจฺจํ อกุสลํ, สมุทยสจฺจํ อกุสลํ, เอวํ อกุสลฏฺเฐน เทฺว สจฺจานิ เอกสจฺเจน สงฺคหิตานิ. เอกสจฺจํ ทฺวีหิ สจฺเจหิ สงฺคหิตํ, ยํ ทุกฺขสจฺจํ กุสลํ, มคฺคสจฺจํ กุสลํ เอวํ กุสลฏฺเฐน เทฺว สจฺจานิ เอกสจฺเจน สงฺคหิตานิ. เอกสจฺจํ ทฺวีหิ สจฺเจหิ สงฺคหิตํ. ยํ ทุกฺขสจฺจํ อพฺยากตํ, นิโรธสจฺจํ อพฺยากตํ, เอวํ อพฺยากตฏฺเฐน เทฺว สจฺจานิ เอกสจฺเจน สงฺคหิตานิ. เอกสจฺจํ ทฺวีหิ สจฺเจหิ สงฺคหิตํ. เอวํ สิยา ตีณิ สจฺจานิ เอกสจฺเจน สงฺคหิตานิ, เอกสจฺจํ ตีหิ สจฺเจหิ สงฺคหิตํ วตฺถุวเสน ปริยาเยนาติ.}

\csromanpalirecto
\csromanmatikaanchor{mtk.90}
\csromantocmarkref{subsection}{2. ทุติยสุตฺตนฺตปาฬิ}{mtk.90}
\bookmark[dest={mtk.90},level=2]{2. ทุติยสุตฺตนฺตปาฬิ}
\csromanheader{2}{2. ทุติย\-สุตฺตนฺต\-ปาฬิ}
\proseitem{13}{\csromanfolio{298}ปุพฺเพ เม ภิกฺขเว\csromansharedfootnote{d752468903101992}{สํ 2. 23 ปิฏฺเฐปิ.} สมฺโพธา อนภิสมฺพุทฺธสฺส โพธิสตฺตสฺเสว สโต เอตทโหสิ “โก นุ โข รูปสฺส อสฺสาโท, โก อาทีนโว, กิํ นิสฺสรณํ. โก เวทนาย อสฺสาโท, โก อาทีนโว, กิํ นิสฺสรณํ. โก สญฺญาย อสฺสาโท, โก อาทีนโว, กิํ นิสฺสรณํ. โก สงฺขารานํ อสฺสาโท, โก อาทีนโว, กิํ นิสฺสรณํ. โก วิญฺญาณสฺส อสฺสาโท, โก อาทีนโว, กิํ นิสฺสรณนฺ”ติ. ตสฺส มยฺหํ ภิกฺขเว เอตทโหสิ “ยํ โข รูปํ ปฏิจฺจ อุปฺปชฺชติ สุขํ โสมนสฺสํ, อยํ รูปสฺส อสฺสาโท. ยํ รูปํ อนิจฺจํ, ตํ ทุกฺขํ วิปริณาม\-ธมฺมํ, อยํ รูปสฺส อาทีนโว. โย รูปสฺมิํ ฉนฺทราค\-วินโย ฉนฺท\-ราค\-ปฺปหานํ, อิทํ รูปสฺส นิสฺสรณํ.  ยํ เวทนํ ปฏิจฺจ ฯเปฯ. ยํ สญฺญํ ปฏิจฺจ. ยํ สงฺขาเร ปฏิจฺจ. ยํ วิญฺญาณํ ปฏิจฺจ อุปฺปชฺชติ สุขํ โสมนสฺสํ อยํ วิญฺญาณสฺส อสฺสาโท. ยํ วิญฺญาณํ อนิจฺจํ, ตํ ทุกฺขํ วิปริณาม\-ธมฺมํ, อยํ วิญฺญาณสฺส อาทีนโว. โย วิญฺญาณสฺมิํ ฉนฺทราค\-วินโย ฉนฺท\-ราค\-ปฺปหานํ, อิทํ วิญฺญาณสฺส นิสฺสรณํ.}

\prose{ยาว กีวญฺจาหํ ภิกฺขเว อิเมสํ ปญฺจนฺนํ อุปาทาน\-กฺขนฺธานํ เอวํ อสฺสาทญฺจ อสฺสาทโต, อาทีนวญฺจ อาทีนวโต, นิสฺสรณญฺจ นิสฺสรณโต ยถาภูตํ นาพฺภญฺญาสิํ. เนว ตาวาหํ ภิกฺขเว สเทวเก โลเก สมารเก สพฺรหฺมเก สสฺสมณ\-พฺราหฺมณิยา ปชาย สเทว\-มนุสฺสาย “อนุตฺตรํ สมฺมาสมฺโพธิํ อภิสมฺพุทฺโธ”ติ ปจฺจญฺญาสิํ\csromansharedfootnote{e1edf5fe7df39663}{อภิสมฺพุทฺโธ ปจฺจญฺญาสิํ (สฺยา) สํ 2. 24 ปิฏฺเฐ ปสฺสิตพฺพา.}. ยโต จ ขฺวาหํ ภิกฺขเว อิเมสํ ปญฺจนฺนํ อุปาทาน\-กฺขนฺธานํ เอวํ อสฺสาทญฺจ อสฺสาทโต, อาทีนวญฺจ อาทีนวโต, นิสฺสรณญฺจ นิสฺสรณโต ยถาภูตํ อพฺภญฺญาสิํ. อถาหํ ภิกฺขเว สเทวเก โลเก สมารเก สพฺรหฺมเก สสฺสมณ\-พฺราหฺมณิยา ปชาย สเทว\-มนุสฺสาย “อนุตฺตรํ สมฺมาสมฺโพธํ อภิสมฺพุทฺโธ”ติ ปจฺจญฺญาสิํ. ญาณญฺจ ปน เม ทสฺสนํ อุทปาทิ, อกุปฺปา เม วิมุตฺติ\csromansharedfootnote{122af127b1219dfa}{เจโตวิมุตฺติ (สฺยา, ก) สํ 2. 24 ปิฏฺเฐ ปสฺสิตพฺพา.}. อยมนฺติมา ชาติ, นตฺถิ ทานิ ปุนพฺภโวติ.}

\csromanmatikaanchor{mtk.91}
\csromantocmarkref{subsection}{3. ทุติยสุตฺตนฺตนิทฺเทส}{mtk.91}
\bookmark[dest={mtk.91},level=2]{3. ทุติยสุตฺตนฺตนิทฺเทส}
\csromanmarkh{2. สจฺจกถา}\csromanmarkhii{3. ทุติย\-สุตฺตนฺต\-นิทฺเทส}\csromanheadernomark{2}{2. สจฺจกถา 3. ทุติย\-สุตฺตนฺต\-นิทฺเทส}
\proseitem{14}{\csromanfolio{299}ยํ รูปํ ปฏิจฺจ อุปฺปชฺชติ สุขํ โสมนสฺสํ, “อยํ รูปสฺส อสฺสาโท”ติ ปหาน\-ปฺปฏิเวโธ สมุทยสจฺจํ. ยํ รูปํ อนิจฺจํ, ตํ ทุกฺขํ วิปริณาม\-ธมฺมํ, “อยํ รูปสฺส อาทีนโว”ติ ปริญฺญา\-ปฏิเวโธ ทุกฺขสจฺจํ. โย รูปสฺมิํ ฉนฺทราค\-วินโย ฉนฺท\-ราค\-ปฺปหานํ, “อิทํ รูปสฺส นิสฺสรณนฺ”ติ สจฺฉิกิริยา\-ปฏิเวโธ นิโรธสจฺจํ. ยา อิเมสุ ตีสุ ฐาเนสุ ทิฏฺฐิ สงฺกปฺโป วาจา กมฺมนฺโต อาชีโว วายาโม สติ สมาธิ ภาวนา\-ปฏิเวโธ มคฺคสจฺจํ.}

\prose{ยํ เวทนํ ปฏิจฺจ ฯเปฯ. ยํ สญฺญํ ปฏิจฺจ. ยํ สงฺขาเร ปฏิจฺจ. ยํ วิญฺญาณํ ปฏิจฺจ อุปฺปชฺชติ สุขํ โสมนสฺสํ, “อยํ วิญฺญาณสฺส อสฺสาโท”ติ ปหาน\-ปฺปฏิเวโธ สมุทยสจฺจํ. ยํ วิญฺญาณํ อนิจฺจํ, ตํ ทุกฺขํ วิปริณาม\-ธมฺมํ, “อยํ วิญฺญาณสฺส อาทีนโว”ติ ปริญฺญา\-ปฏิเวโธ ทุกฺขสจฺจํ. โย วิญฺญาณสฺมิํ ฉนฺทราค\-วินโย ฉนฺท\-ราค\-ปฺปหานํ, “อิทํ วิญฺญาณสฺส นิสฺสรณนฺ”ติ สจฺฉิกิริยา\-ปฏิเวโธ นิโรธสจฺจํ. ยา อิเมสุ ตีสุ ฐาเนสุ ทิฏฺฐิ สงฺกปฺโป วาจา กมฺมนฺโต อาชีโว วายาโม สติ สมาธิ ภาวนา\-ปฏิเวโธ มคฺคสจฺจํ.}

\proseitem{15}{สจฺจนฺติ กติหากาเรหิ สจฺจํ? เอสนฏฺเฐน ปริคฺคห\-ฏฺเฐน ปฏิเวธ\-ฏฺเฐน. กถํ เอสนฏฺเฐน สจฺจํ? ชรามรณํ กิํนิทานํ, กิํสมุทยํ, กิํชาติกํ, กิํปภวนฺติ เอวํ เอสนฏฺเฐน สจฺจํ. ชรามรณํ ชาตินิทานํ, ชาติสมุทยํ, ชาติชาติกํ, ชาติ\-ปฺปภวนฺติ เอวํ ปริคฺคห\-ฏฺเฐน สจฺจํ. ชรามรณญฺจ ปชานาติ, ชรามรณสมุทยญฺจ ปชานาติ, ชรามรณนิโรธญฺจ ปชานาติ, ชรามรณนิโรธ\-คามินิํ ปฏิปทญฺจ ปชานาติ, เอวํ ปฏิเวธ\-ฏฺเฐน สจฺจํ.}

\prose{ชาติ กิํนิทานา, กิํสมุทยา, กิํชาติกา, กิํปภวาติ เอวํ เอสนฏฺเฐน สจฺจํ. ชาติ ภวนิทานา, ภวสมุทยา, ภวชาติกา, ภวปฺปภวาติ เอวํ ปริคฺคห\-ฏฺเฐน สจฺจํ. ชาติญฺจ ปชานาติ, ชาติสมุทยญฺจ ปชานาติ, ชาตินิโรธญฺจ ปชานาติ, ชาตินิโรธ\-คามินิํ ปฏิปทญฺจ ปชานาติ, เอวํ ปฏิเวธ\-ฏฺเฐน สจฺจํ.}

\prose{ภโว กิํนิทาโน, กิํสมุทโย, กิํชาติโก, กิํปภโวติ เอวํ เอสนฏฺเฐน สจฺจํ, ภโว อุปาทานนิทาโน, อุปาทาน\-สมุทโย, อุปาทานชาติโก, อุปาทาน\-ปฺปภโวติ เอวํ ปริคฺคห\-ฏฺเฐน สจฺจํ. ภวญฺจ ปชานาติ, \csromanfolio{300}ภวสมุทยญฺจ ปชานาติ, ภวนิโรธญฺจ ปชานาติ, ภวนิโรธ\-คามินิํ ปฏิปทญฺจ ปชานาติ, เอวํ ปฏิเวธ\-ฏฺเฐน สจฺจํ.}

\prose{อุปาทานํ กิํนิทานํ, กิํสมุทยํ, กิํชาติกํ, กิํปภวนฺติ เอวํ เอสนฏฺเฐน สจฺจํ. อุปาทานํ ตณฺหานิทานํ, ตณฺหาสมุทยํ, ตณฺหาชาติกํ, ตณฺหา\-ปภวนฺติ เอวํ ปริคฺคห\-ฏฺเฐน สจฺจํ. อุปาทานญฺจ ปชานาติ, อุปาทานสมุทยญฺจ ปชานาติ, อุปาทานนิโรธญฺจ ปชานาติ, อุปาทานนิโรธ\-คามินิํ ปฏิปทญฺจ ปชานาติ, เอวํ ปฏิเวธ\-ฏฺเฐน สจฺจํ.}

\prose{ตณฺหา กิํ นิทานา, กิํสมุทยา, กิํชาติกา, กิํปภวาติ เอวํ เอสนฏฺเฐน สจฺจํ. ตณฺหา เวทนานิทานา, เวทนาสมุทยา, เวทนาชาติกา, เวทนาปภวาติ เอวํ ปริคฺคห\-ฏฺเฐน สจฺจํ. ตณฺหญฺจ ปชานาติ, ตณฺหาสมุทยญฺจ ปชานาติ, ตณฺหานิโรธญฺจ ปชานาติ, ตณฺหานิโรธ\-คามินิํ ปฏิปทญฺจ ปชานาติ, เอวํ ปฏิเวธ\-ฏฺเฐน สจฺจํ.}

\prose{เวทนา กิํนิทานา, กิํสมุทยา, กิํชาติกา, กิํปภวาติ เอวํ เอสนฏฺเฐน สจฺจํ. เวทนา ผสฺสนิทานา, ผสฺสสมุทยา, ผสฺสชาติกา, ผสฺส\-ปฺปภวาติ เอวํ ปริคฺคห\-ฏฺเฐน สจฺจํ. เวทนญฺจ ปชานาติ, เวทนาสมุทยญฺจ ปชานาติ, เวทนานิโรธญฺจ ปชานาติ, เวทนา\-นิโรธ\-คามินิํ ปฏิปทญฺจ ปชานาติ, เอวํ ปฏิเวธ\-ฏฺเฐน สจฺจํ.}

\prose{ผสฺโส กิํนิทาโน, กิํสมุทโย, กิํชาติโก, กิํปภโวติ เอวํ เอสนฏฺเฐน สจฺจํ. ผสฺโส สฬายต\-นนิทาโน, สฬายตน\-สมุทโย, สฬายตน\-ชาติโก, สฬายตน\-ปฺปภโวติ เอวํ ปริคฺคห\-ฏฺเฐน สจฺจํ. ผสฺสญฺจ ปชานาติ, ผสฺสสมุทยญฺจ ปชานาติ, ผสฺสนิโรธญฺจ ปชานาติ, ผสฺสนิโรธ\-คามินิํ ปฏิปทญฺจ ปชานาติ, เอวํ ปฏิเวธ\-ฏฺเฐน สจฺจํ.}

\prose{สฬายตนํ กิํนิทานํ, กิํสมุทยํ, กิํชาติกํ, กิํปภวนฺติ เอวํ เอสนฏฺเฐน สจฺจํ. สฬายตนํ นามรูป\-นิทานํ, นามรูป\-สมุทยํ, นามรูป\-ชาติกํ, นามรู\-ปปฺป\-ภวนฺติ เอวํ ปริคฺคห\-ฏฺเฐน สจฺจํ. สฬายตนญฺจ ปชานาติ, สฬายตนสมุทยญฺจ ปชานาติ, สฬายตนนิโรธญฺจ ปชานาติ, สฬายตนนิโรธ\-คามินิํ ปฏิปทญฺจ ปชานาติ, เอวํ ปฏิเวธ\-ฏฺเฐน สจฺจํ.}

\prose{นามรูปํ กิํนิทานํ, กิํสมุทยํ, กิํชาติกํ, กิํปภวนฺติ เอวํ เอสนฏฺเฐน สจฺจํ. นามรูปํ วิญฺญาณนิทานํ, วิญฺญาณ\-สมุทยํ, วิญฺญาณชาติกํ, \csromanfolio{301}วิญฺญาณ\-ปฺปภ\-วนฺติ เอวํ ปริคฺคห\-ฏฺเฐน สจฺจํ. นามรูปญฺจ ปชานาติ, นามรูปสมุทยญฺจ ปชานาติ, นามรูปนิโรธญฺจ ปชานาติ, นามรูปนิโรธ\-คามินิํ ปฏิปทญฺจ ปชานาติ, เอวํ ปฏิเวธ\-ฏฺเฐน สจฺจํ.}

\prose{วิญฺญาณํ กิํนิทานํ, กิํสมุทยํ, กิํชาติกํ, กิํปภวนฺติ เอวํ เอสนฏฺเฐน สจฺจํ. วิญฺญาณํ สงฺขาร\-นิทานํ, สงฺขาร\-สมุทยํ, สงฺขาร\-ชาติกํ, สงฺขาร\-ปฺปภ\-วนฺติ เอวํ ปริคฺคห\-ฏฺเฐน สจฺจํ. วิญฺญาณญฺจ ปชานาติ, วิญฺญาณสมุทยญฺจ ปชานาติ, วิญฺญาณนิโรธญฺจ ปชานาติ, วิญฺญาณ\-นิโรธ\-คามินิํ ปฏิปทญฺจ ปชานาติ, เอวํ ปฏิเวธ\-ฏฺเฐน สจฺจํ.}

\prose{สงฺขารา กิํนิทานา, กิํสมุทยา, กิํชาติกา, กิํปภวาติ เอวํ เอสนฏฺเฐน สจฺจํ. สงฺขารา อวิชฺชานิทานา, อวิชฺชาสมุทยา, อวิชฺชาชาติกา, อวิชฺชาปภาวาติ เอวํ ปริคฺคห\-ฏฺเฐน สจฺจํ. สงฺขาเร จ ปชานาติ, สงฺขารสมุทยญฺจ ปชานาติ, สงฺขาร\-นิโรธญฺจ ปชานาติ, สงฺขาร\-นิโรธ\-คามินิํ ปฏิปทญฺจ ปชานาติ, เอวํ ปฏิเวธ\-ฏฺเฐน สจฺจํ.}

\proseitem{16}{ชรามรณํ ทุกฺขสจฺจํ, ชาติ สมุทยสจฺจํ, อุภินฺนมฺปิ นิสฺสรณํ นิโรธสจฺจํ, นิโรธ\-ปฺปชานนา มคฺคสจฺจํ.  ชาติ ทุกฺขสจฺจํ, ภโว สมุทยสจฺจํ, อุภินฺนมฺปิ นิสฺสรณํ นิโรธสจฺจํ, นิโรธ\-ปฺปชานนา มคฺคสจฺจํ.  ภโว ทุกฺขสจฺจํ, อุปาทานํ สมุทยสจฺจํ, อุภินฺนมฺปิ นิสฺสรณํ นิโรธสจฺจํ, นิโรธ\-ปฺปชานนา มคฺคสจฺจํ.  อุปาทานํ ทุกฺขสจฺจํ, ตณฺหา สมุทยสจฺจํ, อุภินฺนมฺปิ นิสฺสรณํ นิโรธสจฺจํ, นิโรธ\-ปฺปชานนา มคฺคสจฺจํ.  ตณฺหา ทุกฺขสจฺจํ, เวทนา สมุทยสจฺจํ, อุภินฺนมฺปิ นิสฺสรณํ นิโรธสจฺจํ, นิโรธ\-ปฺปชานนา มคฺคสจฺจํ.  เวทนา ทุกฺขสจฺจํ, ผสฺโส สมุทยสจฺจํ, อุภินฺนมฺปิ นิสฺสรณํ นิโรธสจฺจํ, นิโรธ\-ปฺปชานนา มคฺคสจฺจํ.  ผสฺโส ทุกฺขสจฺจํ, สฬายตนํ สมุทยสจฺจํ, อุภินฺนมฺปิ นิสฺสรณํ นิโรธสจฺจํ, นิโรธ\-ปฺปชานนา มคฺคสจฺจํ.  สฬายตนํ ทุกฺขสจฺจํ, นามรูปํ สมุทยสจฺจํ, อุภินฺนมฺปิ นิสฺสรณํ นิโรธสจฺจํ, นิโรธ\-ปฺปชานนา มคฺคสจฺจํ.  นามรูปํ ทุกฺขสจฺจํ, วิญฺญาณํ สมุทยสจฺจํ, อุภินฺนมฺปิ นิสฺสรณํ นิโรธสจฺจํ, นิโรธ\-ปฺปชานนา มคฺคสจฺจํ.  วิญฺญาณํ ทุกฺขสจฺจํ, สงฺขารา สมุทยสจฺจํ, อุภินฺนมฺปิ นิสฺสรณํ นิโรธสจฺจํ, นิโรธ\-ปฺปชานนา มคฺคสจฺจํ.  สงฺขารา ทุกฺขสจฺจํ, อวิชฺชา สมุทยสจฺจํ, อุภินฺนมฺปิ นิสฺสรณํ นิโรธสจฺจํ, นิโรธ\-ปฺปชานนา มคฺคสจฺจํ.}

\prosewithcloser{\csromanfolio{302}ชรามรณํ สิยา ทุกฺขสจฺจํ, สิยา สมุทยสจฺจํ, อุภินฺนมฺปิ นิสฺสรณํ นิโรธสจฺจํ, นิโรธ\-ปฺปชานนา มคฺคสจฺจํ.  ชาติ สิยา ทุกฺขสจฺจํ, สิยา สมุทยสจฺจํ ฯเปฯ. ภโว สิยา ทุกฺขสจฺจํ, สิยา สมุทยสจฺจํ, อุภินฺนมฺปิ นิสฺสรณํ นิโรธสจฺจํ, นิโรธ\-ปฺปชานนา มคฺคสจฺจนฺติ.}{สจฺจกถา นิฏฺฐิตา.}
\nitthitamruled{\textbf{ภาณวาโร.}}
\proseitem{17}{สาวตฺถิ\-นิทานํ.  สตฺติเม ภิกฺขเว\csromansharedfootnote{22b64cd58d2756ff}{สํ 3. 75 ปิฏฺเฐ.} โพชฺฌงฺคา, กตเม สตฺต? สติสมฺโพชฺฌงฺโค ธมฺมวิจย\-สมฺโพชฺฌงฺโค วีริย\-สมฺโพชฺฌงฺโค ปีติสมฺโพชฺฌงฺโค ปสฺสทฺธิ\-สมฺโพชฺฌงฺโค สมาธิ\-สมฺโพชฺฌงฺโค อุเปกฺขา\-สมฺโพชฺฌงฺโค. อิเม โข ภิกฺขเว สตฺต โพชฺฌงฺคา.}

\prose{\textbf{โพชฺฌงฺคา}ติ เกนฏฺเฐน โพชฺฌงฺคา? โพธาย สํวตฺตนฺตีติ โพชฺฌงฺคา. พุชฺฌนฺตีติ โพชฺฌงฺคา, อนุพุชฺฌนฺตีติ โพชฺฌงฺคา, ปฏิพุชฺฌนฺตีติ โพชฺฌงฺคา, สมฺพุชฺฌนฺตีติ โพชฺฌงฺคา. (1)}

\prose{พุชฺฌนฏฺเฐน โพชฺฌงฺคา, อนุพุชฺฌ\-นฏฺเฐน โพชฺฌงฺคา, ปฏิพุชฺฌ\-นฏฺเฐน โพชฺฌงฺคา, สมฺพุชฺฌ\-นฏฺเฐน โพชฺฌงฺคา. (2)}

\prose{โพเธนฺตีติ โพชฺฌงฺคา, อนุโพเธนฺตีติ โพชฺฌงฺคา, ปฏิโพเธนฺตีติ โพชฺฌงฺคา, สมฺโพเธนฺตีติ โพชฺฌงฺคา. (3)}

\prose{โพธนฏฺเฐน โพชฺฌงฺคา, อนุโพธ\-นฏฺเฐน โพชฺฌงฺคา, ปฏิโพธ\-นฏฺเฐน โพชฺฌงฺคา, สมฺโพธ\-นฏฺเฐน โพชฺฌงฺคา. (4)}

\prose{โพธิปกฺขิย\-ฏฺเฐน โพชฺฌงฺคา, อนุโพธิปกฺขิย\-ฏฺเฐน โพชฺฌงฺคา, ปฏิโพธิปกฺขิย\-ฏฺเฐน โพชฺฌงฺคา, สมฺโพธิ\-ปกฺขิย\-ฏฺเฐน โพชฺฌงฺคา. (5) (ปญฺจม\-จตุกฺกํ)}

\prose{พุทฺธิ\-ลภ\-นฏฺเฐน โพชฺฌงฺคา, พุทฺธิ\-ปฏิลภ\-นฏฺเฐน โพชฺฌงฺคา, พุทฺธิ\-โรป\-นฏฺเฐน โพชฺฌงฺคา, พุทฺธิอภิโรปนฏฺเฐน โพชฺฌงฺคา, พุทฺธิ\-ปาป\-นฏฺเฐน โพชฺฌงฺคา, พุทฺธิ\-สมฺปาปน\-ฏฺเฐน โพชฺฌงฺคา. (ฉกฺกํ)}

\csromanmatikaanchor{mtk.92}
\csromanmatikaanchor{mtk.93}
\csromantocmarknopage{section}{3. โพชฺฌงฺคกถา}{mtk.92}
\bookmark[dest={mtk.92},level=1]{3. โพชฺฌงฺคกถา}
\csromantocmarkref{section}{มูลมูลกาทิทสก}{mtk.93}
\bookmark[dest={mtk.93},level=1]{มูลมูลกาทิทสก}
\csromanheader{1}{3. โพชฺฌงฺคกถา มูลมูลกา\-ทิ\-ทสก}
\proseitem{18}{\csromanfolio{303}มูลฏฺเฐน โพชฺฌงฺคา, มูลจริย\-ฏฺเฐน โพชฺฌงฺคา, มูลปริคฺคห\-ฏฺเฐน โพชฺฌงฺคา, มูลปริวาร\-ฏฺเฐน โพชฺฌงฺคา, มูลปริปูรณ\-ฏฺเฐน\csromansharedfootnote{1df16b690afe2ea6}{มูลปริปูรฏฺเฐน (สฺยา, ก) เอวมุปริปิ.} โพชฺฌงฺคา, มูลปริปาก\-ฏฺเฐน โพชฺฌงฺคา, มูลปฏิสมฺภิท\-ฏฺเฐน โพชฺฌงฺคา, มูลปฏิสมฺภิทา\-ปาป\-นฏฺเฐน โพชฺฌงฺคา, มูลปฏิสมฺภิทาย วสีภาว\-ฏฺเฐน โพชฺฌงฺคา, มูลปฏิสมฺภิทาย วสีภาวปฺปตฺตานมฺปิ โพชฺฌงฺคา. (1)}

\prose{เหตุฏฺเฐน โพชฺฌงฺคา, เหตุ\-จริย\-ฏฺเฐน โพชฺฌงฺคา, เหตุ\-ปริคฺคห\-ฏฺเฐน โพชฺฌงฺคา, เหตุ\-ปริวาร\-ฏฺเฐน โพชฺฌงฺคา, เหตุ\-ปริปูรณ\-ฏฺเฐน โพชฺฌงฺคา, เหตุ\-ปริปาก\-ฏฺเฐน โพชฺฌงฺคา, เหตุ\-ปฏิสมฺภิท\-ฏฺเฐน โพชฺฌงฺคา, เหตุ\-ปฏิสมฺภิทา\-ปาป\-นฏฺเฐน โพชฺฌงฺคา, เหตุ\-ปฏิสมฺภิทาย วสีภาว\-ฏฺเฐน โพชฺฌงฺคา, เหตุ\-ปฏิสมฺภิทาย วสีภาวปฺปตฺตานมฺปิ โพชฺฌงฺคา. (2)}

\prose{ปจฺจยฏฺเฐน โพชฺฌงฺคา, ปจฺจย\-จริย\-ฏฺเฐน โพชฺฌงฺคา, ปจฺจย\-ปริคฺคห\-ฏฺเฐน โพชฺฌงฺคา, ปจฺจย\-ปริวาร\-ฏฺเฐน โพชฺฌงฺคา, ปจฺจย\-ปริปูรณ\-ฏฺเฐน โพชฺฌงฺคา, ปจฺจย\-ปริปาก\-ฏฺเฐน โพชฺฌงฺคา, ปจฺจย\-ปฏิสมฺภิท\-ฏฺเฐน โพชฺฌงฺคา, ปจฺจย\-ปฏิสมฺภิทา\-ปาป\-นฏฺเฐน โพชฺฌงฺคา, ปจฺจย\-ปฏิสมฺภิทาย วสีภาว\-ฏฺเฐน โพชฺฌงฺคา, ปจฺจย\-ปฏิสมฺภิทาย วสีภาวปฺปตฺตานมฺปิ โพชฺฌงฺคา. (3)}

\prose{วิสุทฺธฏฺเฐน โพชฺฌงฺคา, วิสุทฺธิ\-จริย\-ฏฺเฐน โพชฺฌงฺคา, วิสุทฺธิ\-ปริคฺคห\-ฏฺเฐน โพชฺฌงฺคา, วิสุทฺธิ\-ปริวาร\-ฏฺเฐน โพชฺฌงฺคา, วิสุทฺธิ\-ปริปูรณ\-ฏฺเฐน โพชฺฌงฺคา, วิสุทฺธิ\-ปริปาก\-ฏฺเฐน โพชฺฌงฺคา, วิสุทฺธิ\-ปฏิสมฺภิท\-ฏฺเฐน โพชฺฌงฺคา, วิสุทฺธิ\-ปฏิสมฺภิทา\-ปาป\-นฏฺเฐน โพชฺฌงฺคา, วิสุทฺธิ\-ปฏิสมฺภิทาย วสีภาว\-ฏฺเฐน โพชฺฌงฺคา, วิสุทฺธิ\-ปฏิสมฺภิทาย วสีภาวปฺปตฺตานมฺปิ โพชฺฌงฺคา. (4)}

\prose{อนวชฺชฏฺเฐน โพชฺฌงฺคา, อนวชฺช\-จริย\-ฏฺเฐน โพชฺฌงฺคา, อนวชฺช\-ปริคฺคห\-ฏฺเฐน โพชฺฌงฺคา, อนวชฺช\-ปริวาร\-ฏฺเฐน โพชฺฌงฺคา, อนวชฺช\-ปริปูรณ\-ฏฺเฐน โพชฺฌงฺคา, อนวชฺช\-ปริปาก\-ฏฺเฐน โพชฺฌงฺคา, อนวชฺช\-ปฏิสมฺภิท\-ฏฺเฐน โพชฺฌงฺคา, อนวชฺช\-ปฏิสมฺภิทา\-ปาป\-นฏฺเฐน โพชฺฌงฺคา, อนวชฺช\-ปฏิสมฺภิทาย วสีภาว\-ฏฺเฐน โพชฺฌงฺคา, อนวชฺช\-ปฏิสมฺภิทาย วสีภาวปฺปตฺตานมฺปิ โพชฺฌงฺคา. (5)}

\prose{เนกฺขมฺม\-ฏฺเฐน โพชฺฌงฺคา, เนกฺขมฺม\-จริย\-ฏฺเฐน โพชฺฌงฺคา, เนกฺขมฺม\-ปริคฺคห\-ฏฺเฐน โพชฺฌงฺคา, เนกฺขมฺม\-ปริวาร\-ฏฺเฐน โพชฺฌงฺคา, เนกฺขมฺม\-ปริปูรณ\-ฏฺเฐน โพชฺฌงฺคา, \csromanfolio{304}เนกฺขมฺม\-ปริปาก\-ฏฺเฐน โพชฺฌงฺคา, เนกฺขมฺม\-ปฏิสมฺภิท\-ฏฺเฐน โพชฺฌงฺคา, เนกฺขมฺม\-ปฏิสมฺภิทา\-ปาป\-นฏฺเฐน โพชฺฌงฺคา, เนกฺขมฺม\-ปฏิสมฺภิทาย วสีภาว\-ฏฺเฐน โพชฺฌงฺคา, เนกฺขมฺม\-ปฏิสมฺภิทาย วสีภาวปฺปตฺตานมฺปิ โพชฺฌงฺคา. (6)}

\prose{วิมุตฺตฏฺเฐน โพชฺฌงฺคา, วิมุตฺติ\-จริย\-ฏฺเฐน โพชฺฌงฺคา, วิมุตฺติ\-ปริคฺคห\-ฏฺเฐน โพชฺฌงฺคา, วิมุตฺติ\-ปริวาร\-ฏฺเฐน โพชฺฌงฺคา, วิมุตฺติ\-ปริปูรณ\-ฏฺเฐน โพชฺฌงฺคา, วิมุตฺติ\-ปริปาก\-ฏฺเฐน โพชฺฌงฺคา, วิมุตฺติ\-ปฏิสมฺภิท\-ฏฺเฐน โพชฺฌงฺคา, วิมุตฺติ\-ปฏิสมฺภิทา\-ปาป\-นฏฺเฐน โพชฺฌงฺคา, วิมุตฺติ\-ปฏิสมฺภิทาย วสีภาว\-ฏฺเฐน โพชฺฌงฺคา, วิมุตฺติ\-ปฏิสมฺภิทาย วสีภาวปฺปตฺตานมฺปิ โพชฺฌงฺคา. (7)}

\prose{อนาสวฏฺเฐน โพชฺฌงฺคา, อนาสว\-จริย\-ฏฺเฐน โพชฺฌงฺคา, อนาสว\-ปริคฺคห\-ฏฺเฐน โพชฺฌงฺคา, อนาสว\-ปริวาร\-ฏฺเฐน โพชฺฌงฺคา, อนาสว\-ปริปูรณ\-ฏฺเฐน โพชฺฌงฺคา, อนาสว\-ปริปาก\-ฏฺเฐน โพชฺฌงฺคา, อนาสว\-ปฏิสมฺภิท\-ฏฺเฐน โพชฺฌงฺคา, อนาสว\-ปฏิสมฺภิทา\-ปาป\-นฏฺเฐน โพชฺฌงฺคา, อนาสว\-ปฏิสมฺภิทาย วสีภาว\-ฏฺเฐน โพชฺฌงฺคา, อนาสว\-ปฏิสมฺภิทาย วสีภาวปฺปตฺตานมฺปิ โพชฺฌงฺคา. (8)}

\prose{วิเวกฏฺเฐน โพชฺฌงฺคา, วิเวก\-จริย\-ฏฺเฐน โพชฺฌงฺคา, วิเวก\-ปริคฺคห\-ฏฺเฐน โพชฺฌงฺคา, วิเวก\-ปริวาร\-ฏฺเฐน โพชฺฌงฺคา, วิเวก\-ปริปูรณ\-ฏฺเฐน โพชฺฌงฺคา, วิเวก\-ปริปาก\-ฏฺเฐน โพชฺฌงฺคา, วิเวก\-ปฏิสมฺภิท\-ฏฺเฐน โพชฺฌงฺคา, วิเวก\-ปฏิสมฺภิทา\-ปาป\-นฏฺเฐน โพชฺฌงฺคา, วิเวก\-ปฏิสมฺภิทาย วสีภาว\-ฏฺเฐน โพชฺฌงฺคา, วิเวก\-ปฏิสมฺภิทาย วสีภาวปฺปตฺตานมฺปิ โพชฺฌงฺคา. (9)}

\prose{โวสคฺคฏฺเฐน โพชฺฌงฺคา, โวสคฺค\-จริย\-ฏฺเฐน โพชฺฌงฺคา, โวสคฺค\-ปริคฺคห\-ฏฺเฐน โพชฺฌงฺคา, โวสคฺค\-ปริวาร\-ฏฺเฐน โพชฺฌงฺคา, โวสคฺค\-ปริปูรณ\-ฏฺเฐน โพชฺฌงฺคา, โวสคฺค\-ปริปาก\-ฏฺเฐน โพชฺฌงฺคา, โวสคฺค\-ปฏิสมฺภิท\-ฏฺเฐน โพชฺฌงฺคา, โวสคฺค\-ปฏิสมฺภิทา\-ปาป\-นฏฺเฐน โพชฺฌงฺคา, โวสคฺค\-ปฏิสมฺภิทาย วสีภาว\-ฏฺเฐน โพชฺฌงฺคา, โวสคฺค\-ปฏิสมฺภิทาย วสีภาวปฺปตฺตานมฺปิ โพชฺฌงฺคา. (10)}

\proseitem{19}{มูลฏฺฐํ พุชฺฌนฺตีติ โพชฺฌงฺคา, เหตุฏฺฐํ พุชฺฌนฺตีติ โพชฺฌงฺคา, ปจฺจยฏฺฐํ พุชฺฌนฺตีติ โพชฺฌงฺคา, วิสุทฺธฏฺฐํ พุชฺฌนฺตีติ โพชฺฌงฺคา, อนวชฺชฏฺฐํ พุชฺฌนฺตีติ โพชฺฌงฺคา, เนกฺขมฺมฏฺฐํ พุชฺฌนฺตีติ โพชฺฌงฺคา, วิมุตฺตฏฺฐํ พุชฺฌนฺตีติ โพชฺฌงฺคา, อนาสวฏฺฐํ พุชฺฌนฺตีติ โพชฺฌงฺคา, วิเวกฏฺฐํ พุชฺฌนฺตีติ โพชฺฌงฺคา, โวสคฺคฏฺฐํ พุชฺฌนฺตีติ โพชฺฌงฺคา.}

\prose{\csromanfolio{305}มูลจริยฏฺฐํ พุชฺฌนฺตีติ โพชฺฌงฺคา, เหตุ\-จริยฏฺฐํ พุชฺฌนฺตีติ โพชฺฌงฺคา, ปจฺจย\-จริยฏฺฐํ พุชฺฌนฺตีติ โพชฺฌงฺคา, วิสุทฺธิ\-จริยฏฺฐํ พุชฺฌนฺตีติ โพชฺฌงฺคา, อนวชฺช\-จริยฏฺฐํ พุชฺฌนฺตีติ โพชฺฌงฺคา, เนกฺขมฺม\-จริยฏฺฐํ พุชฺฌนฺตีติ โพชฺฌงฺคา, วิมุตฺติ\-จริยฏฺฐํ พุชฺฌนฺตีติ โพชฺฌงฺคา, อนาสว\-จริยฏฺฐํ พุชฺฌนฺตีติ โพชฺฌงฺคา, วิเวก\-จริยฏฺฐํ พุชฺฌนฺตีติ โพชฺฌงฺคา, โวสคฺค\-จริยฏฺฐํ พุชฺฌนฺตีติ โพชฺฌงฺคา.}

\prose{มูลปริคฺคหฏฺฐํ พุชฺฌนฺตีติ โพชฺฌงฺคา ฯเปฯ โวสคฺค\-ปริคฺคหฏฺฐํ พุชฺฌนฺตีติ โพชฺฌงฺคา. มูลปริวารฏฺฐํ พุชฺฌนฺตีติ โพชฺฌงฺคา ฯเปฯ โวสคฺค\-ปริวารฏฺฐํ พุชฺฌนฺตีติ โพชฺฌงฺคา. มูลปริปูรณฏฺฐํ พุชฺฌนฺตีติ โพชฺฌงฺคา ฯเปฯ โวสคฺค\-ปริปูรฏฺฐํ พุชฺฌนฺตีติ โพชฺฌงฺคา. มูลปริปากฏฺฐํ พุชฺฌนฺตีติ โพชฺฌงฺคา ฯเปฯ โวสคฺค\-ปริปากฏฺฐํ พุชฺฌนฺตีติ โพชฺฌงฺคา. มูลปฏิสมฺภิทฏฺฐํ พุชฺฌนฺตีติ โพชฺฌงฺคา ฯเปฯ โวสคฺค\-ปฏิสมฺภิทฏฺฐํ พุชฺฌนฺตีติ โพชฺฌงฺคา.  มูลปฏิสมฺภิทา\-ปาป\-นฏฺฐํ พุชฺฌนฺตีติ โพชฺฌงฺคา ฯเปฯ โวสคฺค\-ปฏิสมฺภิทา\-ปาป\-นฏฺฐํ พุชฺฌนฺตีติ โพชฺฌงฺคา.  มูลปฏิสมฺภิทาย วสีภาวฏฺฐํ พุชฺฌนฺตีติ โพชฺฌงฺคา ฯเปฯ โวสคฺค\-ปฏิสมฺภิทาย วสีภาวฏฺฐํ พุชฺฌนฺตีติ โพชฺฌงฺคา ฯเปฯ.}

\prose{ปริคฺคหฏฺฐํ พุชฺฌนฺตีติ โพชฺฌงฺคา, ปริวารฏฺฐํ พุชฺฌนฺตีติ โพชฺฌงฺคา, ปริปูรณฏฺฐํ พุชฺฌนฺตีติ โพชฺฌงฺคา, เอกคฺคฏฺฐํ พุชฺฌนฺตีติ โพชฺฌงฺคา, อวิกฺเขปฏฺฐํ พุชฺฌนฺตีติ โพชฺฌงฺคา, ปคฺคหฏฺฐํ พุชฺฌนฺตีติ โพชฺฌงฺคา, อวิสารฏฺฐํ พุชฺฌนฺตีติ โพชฺฌงฺคา, อนาวิลฏฺฐํ พุชฺฌนฺตีติ โพชฺฌงฺคา, อนิญฺชนฏฺฐํ พุชฺฌนฺตีติ โพชฺฌงฺคา, เอกตฺตุ\-ปฏฺฐาน\-วเสน จิตฺตสฺส ฐิตฏฺฐํ พุชฺฌนฺตีติ โพชฺฌงฺคา, อารมฺมณฏฺฐํ พุชฺฌนฺตีติ โพชฺฌงฺคา, โคจรฏฺฐํ พุชฺฌนฺตีติ โพชฺฌงฺคา, ปหานฏฺฐํ พุชฺฌนฺตีติ โพชฺฌงฺคา, ปริจฺจาคฏฺฐํ พุชฺฌนฺตีติ โพชฺฌงฺคา, วุฏฺฐานฏฺฐํ พุชฺฌนฺตีติ โพชฺฌงฺคา, วิวฏฺฏนฏฺฐํ พุชฺฌนฺตีติ โพชฺฌงฺคา, สนฺตฏฺฐํ พุชฺฌนฺตีติ โพชฺฌงฺคา, ปณีตฏฺฐํ พุชฺฌนฺตีติ โพชฺฌงฺคา, วิมุตฺตฏฺฐํ พุชฺฌนฺตีติ โพชฺฌงฺคา, อนาสวฏฺฐํ พุชฺฌนฺตีติ โพชฺฌงฺคา, ตรณฏฺฐํ พุชฺฌนฺตีติ โพชฺฌงฺคา, อนิมิตฺตฏฺฐํ พุชฺฌนฺตีติ โพชฺฌงฺคา, อปฺปณิหิตฏฺฐํ พุชฺฌนฺตีติ โพชฺฌงฺคา, สุญฺญตฏฺฐํ พุชฺฌนฺตีติ โพชฺฌงฺคา, เอกรสฏฺฐํ พุชฺฌนฺตีติ โพชฺฌงฺคา, อนติวตฺตนฏฺฐํ พุชฺฌนฺตีติ โพชฺฌงฺคา, ยุคนทฺธฏฺฐํ พุชฺฌนฺตีติ โพชฺฌงฺคา, นิยฺยานฏฺฐํ พุชฺฌนฺตีติ โพชฺฌงฺคา, เหตุฏฺฐํ พุชฺฌนฺตีติ โพชฺฌงฺคา, ทสฺสนฏฺฐํ พุชฺฌนฺตีติ โพชฺฌงฺคา, อาธิปเตยฺยฏฺฐํ พุชฺฌนฺตีติ โพชฺฌงฺคา.}

\prose{สมถสฺส อวิกฺเขปฏฺฐํ พุชฺฌนฺตีติ โพชฺฌงฺคา, วิปสฺสนาย อนุปสฺสนฏฺฐํ พุชฺฌนฺตีติ โพชฺฌงฺคา, สมถ\-วิปสฺสนานํ เอกรสฏฺฐํ พุชฺฌนฺตีติ โพชฺฌงฺคา, ยุคนทฺธสฺส อนติวตฺตนฏฺฐํ พุชฺฌนฺตีติ โพชฺฌงฺคา.}

\prose{\csromanfolio{306}สิกฺขาย สมาทานฏฺฐํ พุชฺฌนฺตีติ โพชฺฌงฺคา, อารมฺมณสฺส โคจรฏฺฐํ พุชฺฌนฺตีติ โพชฺฌงฺคา, ลีนสฺส จิตฺตสฺส ปคฺคหฏฺฐํ พุชฺฌนฺตีติ โพชฺฌงฺคา, อุทฺธตสฺส จิตฺตสฺส นิคฺคหฏฺฐํ พุชฺฌนฺตีติ โพชฺฌงฺคา. อุโภ\-วิสุทฺธานํ อชฺฌุเปกฺข\-นฏฺฐํ พุชฺฌนฺตีติ โพชฺฌงฺคา. วิเสสาธิคมฏฺฐํ พุชฺฌนฺตีติ โพชฺฌงฺคา. อุตฺตริ\-ปฏิเวธฏฺฐํ พุชฺฌนฺตีติ โพชฺฌงฺคา. สจฺจาภิสมยฏฺฐํ พุชฺฌนฺตีติ โพชฺฌงฺคา. นิโรเธ ปติฏฺฐา\-ปกฏฺฐํ พุชฺฌนฺตีติ โพชฺฌงฺคา.}

\prose{สทฺธิ\-นฺทฺริยสฺส อธิโมกฺขฏฺฐํ พุชฺฌนฺตีติ โพชฺฌงฺคา ฯเปฯ ปญฺญินฺทฺริยสฺส ทสฺสนฏฺฐํ พุชฺฌนฺตีติ โพชฺฌงฺคา. สทฺธาพลสฺส อสฺสทฺธิเย อกมฺปิยฏฺฐํ พุชฺฌนฺตีติ โพชฺฌงฺคา ฯเปฯ ปญฺญาพลสฺส อวิชฺชาย อกมฺปิยฏฺฐํ พุชฺฌนฺตีติ โพชฺฌงฺคา. สติสมฺโพชฺฌงฺคสฺส อุปฏฺฐานฏฺฐํ พุชฺฌนฺตีติ โพชฺฌงฺคา ฯเปฯ อุเปกฺขา\-สมฺโพชฺฌงฺคสฺส ปฏิสงฺขา\-นฏฺฐํ พุชฺฌนฺตีติ โพชฺฌงฺคา. สมฺมาทิฏฺฐิยา ทสฺสนฏฺฐํ พุชฺฌนฺตีติ โพชฺฌงฺคา ฯเปฯ สมฺมา\-สมาธิสฺส อวิกฺเขปฏฺฐํ พุชฺฌนฺตีติ โพชฺฌงฺคา.}

\prose{อินฺทฺริยานํ อาธิปเตยฺยฏฺฐํ พุชฺฌนฺตีติ โพชฺฌงฺคา. พลานํ อกมฺปิยฏฺฐํ พุชฺฌนฺตีติ โพชฺฌงฺคา. นิยฺยานฏฺฐํ พุชฺฌนฺตีติ โพชฺฌงฺคา. มคฺคสฺส เหตุฏฺฐํ พุชฺฌนฺตีติ โพชฺฌงฺคา. สติปฏฺฐานานํ อุปฏฺฐานฏฺฐํ พุชฺฌนฺตีติ โพชฺฌงฺคา. สมฺม\-ปฺปธานานํ ปทหนฏฺฐํ พุชฺฌนฺตีติ โพชฺฌงฺคา. อิทฺธิปาทานํ อิชฺฌนฏฺฐํ พุชฺฌนฺตีติ โพชฺฌงฺคา. สจฺจานํ ตถฏฺฐํ พุชฺฌนฺตีติ โพชฺฌงฺคา. ปโยคานํ ปฏิปฺปสฺสทฺธฏฺฐํ พุชฺฌนฺตีติ โพชฺฌงฺคา. ผลานํ สจฺฉิกิริยฏฺฐํ พุชฺฌนฺตีติ โพชฺฌงฺคา.}

\prose{วิตกฺกสฺส อภินิโรปนฏฺฐํ พุชฺฌนฺตีติ โพชฺฌงฺคา. วิจารสฺส อุปวิจารฏฺฐํ พุชฺฌนฺตีติ โพชฺฌงฺคา. ปีติยา ผรณฏฺฐํ พุชฺฌนฺตีติ โพชฺฌงฺคา. สุขสฺส อภิสนฺท\-นฏฺฐํ พุชฺฌนฺตีติ โพชฺฌงฺคา. จิตฺตสฺส เอกคฺคฏฺฐํ พุชฺฌนฺตีติ โพชฺฌงฺคา.}

\prose{อาวชฺชนฏฺฐํ พุชฺฌนฺตีติ โพชฺฌงฺคา. วิชานนฏฺฐํ พุชฺฌนฺตีติ โพชฺฌงฺคา. ปชานนฏฺฐํ พุชฺฌนฺตีติ โพชฺฌงฺคา. สญฺชานนฏฺฐํ พุชฺฌนฺตีติ โพชฺฌงฺคา. เอโกทฏฺฐํ พุชฺฌนฺตีติ โพชฺฌงฺคา. อภิญฺญาย ญาตฏฺฐํ\csromansharedfootnote{1b6f2e69566b17c3}{อภิญฺเญยฺยตฺถํ (สฺยา) 17 ปิฏฺเฐ ปสฺสิตพฺพา.} พุชฺฌนฺตีติ โพชฺฌงฺคา. ปริญฺญาย ตีรณฏฺฐํ พุชฺฌนฺตีติ โพชฺฌงฺคา. ปหานสฺส ปริจฺจาคฏฺฐํ พุชฺฌนฺตีติ โพชฺฌงฺคา. ภาวนาย เอกรสฏฺฐํ พุชฺฌนฺตีติ โพชฺฌงฺคา. สจฺฉิกิริยาย ผสฺสนฏฺฐํ พุชฺฌนฺตีติ โพชฺฌงฺคา. ขนฺธานํ ขนฺธฏฺฐํ พุชฺฌนฺตีติ โพชฺฌงฺคา. ธาตูนํ ธาตุฏฺฐํ พุชฺฌนฺตีติ โพชฺฌงฺคา. อายตนานํ อายตนฏฺฐํ พุชฺฌนฺตีติ โพชฺฌงฺคา. สงฺขตานํ สงฺขตฏฺฐํ พุชฺฌตีติ โพชฺฌงฺคา. อสงฺขตสฺส อสงฺขตฏฺฐํ พุชฺฌนฺตีติ โพชฺฌงฺคา.}

\prose{\csromanfolio{307}จิตฺตฏฺฐํ พุชฺฌนฺตีติ โพชฺฌงฺคา. จิตฺตา\-นนฺตริยฏฺฐํ พุชฺฌนฺตีติ โพชฺฌงฺคา. จิตฺตสฺส วุฏฺฐานฏฺฐํ พุชฺฌนฺตีติ โพชฺฌงฺคา. จิตฺตสฺส วิวฏฺฏนฏฺฐํ\csromansharedfootnote{f7218b4853fa7c0d}{วิวชฺชนฏฺฐํ (สฺยา)} พุชฺฌนฺตีติ โพชฺฌงฺคา. จิตฺตสฺส เหตุฏฺฐํ พุชฺฌนฺตีติ โพชฺฌงฺคา. จิตฺตสฺส ปจฺจยฏฺฐํ พุชฺฌนฺตีติ โพชฺฌงฺคา. จิตฺตสฺส วตฺถุฏฺฐํ พุชฺฌนฺตีติ โพชฺฌงฺคา. จิตฺตสฺส ภูมฏฺฐํ\csromansharedfootnote{6f0032775c7f22e9}{ภุมฺมฏฺฐํ (สฺยา) 18 ปิฏฺเฐ ปสฺสิตพฺพา.} พุชฺฌนฺตีติ โพชฺฌงฺคา. จิตฺตสฺส อารมฺมณฏฺฐํ พุชฺฌนฺตีติ โพชฺฌงฺคา. จิตฺตสฺส โคจรฏฺฐํ พุชฺฌนฺตีติ โพชฺฌงฺคา. จิตฺตสฺส จริยฏฺฐํ พุชฺฌนฺตีติ โพชฺฌงฺคา. จิตฺตสฺส คตฏฺฐํ พุชฺฌนฺตีติ โพชฺฌงฺคา. จิตฺตสฺส อภินีหารฏฺฐํ พุชฺฌนฺตีติ โพชฺฌงฺคา. จิตฺตสฺส นิยฺยานฏฺฐํ พุชฺฌนฺตีติ โพชฺฌงฺคา. จิตฺตสฺส นิสฺสรณฏฺฐํ พุชฺฌนฺตีติ โพชฺฌงฺคา.}

\prose{เอกตฺเต อาวชฺชนฏฺฐํ พุชฺฌนฺตีติ โพชฺฌงฺคา. เอกตฺเต วิชานนฏฺฐํ พุชฺฌนฺตีติ โพชฺฌงฺคา. เอกตฺเต ปชานนฏฺฐํ พุชฺฌนฺตีติ โพชฺฌงฺคา. เอกตฺเต สญฺชานนฏฺฐํ พุชฺฌนฺตีติ โพชฺฌงฺคา. เอกตฺเต เอโกทฏฺฐํ พุชฺฌนฺตีติ โพชฺฌงฺคา. ( )\csromansharedfootnote{98e1091e2825076a}{(เอกตฺเต อุปนิพนฺธนฏฺฐํ พุชฺฌนฺตีติ โพชฺฌงฺคา) (สพฺพตฺถ)} เอกตฺเต ปกฺขนฺท\-นฏฺฐํ พุชฺฌนฺตีติ โพชฺฌงฺคา. เอกตฺเต ปสีทนฏฺฐํ พุชฺฌนฺตีติ โพชฺฌงฺคา. เอกตฺเต สนฺติฏฺฐ\-นฏฺฐํ พุชฺฌนฺตีติ โพชฺฌงฺคา. เอกตฺเต วิมุจฺจนฏฺฐํ พุชฺฌนฺตีติ โพชฺฌงฺคา. เอกตฺเต “เอตํ สนฺตนฺ”ติ ปสฺสนฏฺฐํ พุชฺฌนฺตีติ โพชฺฌงฺคา. เอกตฺเต ยานีกตฏฺฐํ พุชฺฌนฺตีติ โพชฺฌงฺคา. เอกตฺเต วตฺถุกตฏฺฐํ พุชฺฌนฺตีติ โพชฺฌงฺคา. เอกตฺเต อนุฏฺฐิตฏฺฐํ พุชฺฌนฺตีติ โพชฺฌงฺคา. เอกตฺเต ปริจิตฏฺฐํ พุชฺฌนฺตีติ โพชฺฌงฺคา. เอกตฺเต สุสมารทฺธฏฺฐํ พุชฺฌนฺตีติ โพชฺฌงฺคา. เอกตฺเต ปริคฺคหฏฺฐํ พุชฺฌนฺตีติ โพชฺฌงฺคา. เอกตฺเต ปริวารฏฺฐํ พุชฺฌนฺตีติ โพชฺฌงฺคา. เอกตฺเต ปริปูรณฏฺฐํ พุชฺฌนฺตีติ โพชฺฌงฺคา. เอกตฺเต สโมธานฏฺฐํ พุชฺฌนฺตีติ โพชฺฌงฺคา. เอกตฺเต อธิฏฺฐานฏฺฐํ พุชฺฌนฺตีติ โพชฺฌงฺคา. เอกตฺเต อาเสวนฏฺฐํ พุชฺฌนฺตีติ โพชฺฌงฺคา. เอกตฺเต ภาวนฏฺฐํ พุชฺฌนฺตีติ โพชฺฌงฺคา. เอกตฺเต พหุลีกมฺมฏฺฐํ พุชฺฌนฺตีติ โพชฺฌงฺคา. เอกตฺเต สุสมุคฺคตฏฺฐํ พุชฺฌนฺตีติ โพชฺฌงฺคา. เอกตฺเต สุวิมุตฺตฏฺฐํ พุชฺฌนฺตีติ โพชฺฌงฺคา. เอกตฺเต พุชฺฌนฏฺฐํ พุชฺฌนฺตีติ โพชฺฌงฺคา.เอกตฺเต อนุพุชฺฌ\-นฏฺฐํ พุชฺฌนฺตีติ โพชฺฌงฺคา. เอกตฺเต ปฏิพุชฺฌ\-นฏฺฐํ พุชฺฌนฺตีติ โพชฺฌงฺคา. เอกตฺเต สมฺพุชฺฌ\-นฏฺฐํ พุชฺฌนฺตีติ โพชฺฌงฺคา. เอกตฺเต โพธนฏฺฐํ พุชฺฌนฺตีติ โพชฺฌงฺคา. เอกตฺเต อนุโพธ\-นฏฺฐํ พุชฺฌนฺตีติ โพชฺฌงฺคา. เอกตฺเต ปฏิโพธ\-นฏฺฐํ พุชฺฌนฺตีติ โพชฺฌงฺคา. เอกตฺเต สมฺพุชฺฌ\-นฏฺฐํ พุชฺฌนฺตีติ โพชฺฌงฺคา. เอกตฺเต โพธิปกฺขิยฏฺฐํ พุชฺฌนฺตีติ โพชฺฌงฺคา. เอกตฺเต อนุโพธิปกฺขิยฏฺฐํ พุชฺฌนฺตีติ โพชฺฌงฺคา. เอกตฺเต ปฏิโพธิปกฺขิยฏฺฐํ พุชฺฌนฺตีติ โพชฺฌงฺคา. เอกตฺเต สมฺโพธิ\-ปกฺขิยฏฺฐํ พุชฺฌนฺตีติ โพชฺฌงฺคา. เอกตฺเต}

\prose{\csromanfolio{308}โชตนฏฺฐํ พุชฺฌนฺตีติ โพชฺฌงฺคา. เอกตฺเต อุชฺโชตนฏฺฐํ พุชฺฌนฺตีติ โพชฺฌงฺคา. เอกตฺเต อนุโชตนฏฺฐํ พุชฺฌนฺตีติ โพชฺฌงฺคา. เอกตฺเต ปฏิโชตนฏฺฐํ พุชฺฌนฺตีติ โพชฺฌงฺคา. เอกตฺเต สญฺโชตนฏฺฐํ พุชฺฌนฺตีติ โพชฺฌงฺคา.}

\prose{ปตาปนฏฺฐํ พุชฺฌนฺตีติ โพชฺฌงฺคา. วิโรจนฏฺฐํ พุชฺฌนฺตีติ โพชฺฌงฺคา. กิเลสานํ สนฺตาปนฏฺฐํ พุชฺฌนฺตีติ โพชฺฌงฺคา. อมลฏฺฐํ พุชฺฌนฺตีติ โพชฺฌงฺคา. วิมลฏฺฐํ พุชฺฌนฺตีติ โพชฺฌงฺคา. นิมฺมลฏฺฐํ พุชฺฌนฺตีติ โพชฺฌงฺคา. สมฏฺฐํ พุชฺฌนฺตีติ โพชฺฌงฺคา. สมยฏฺฐํ พุชฺฌนฺตีติ โพชฺฌงฺคา, วิเวกฏฺฐํ พุชฺฌนฺตีติ โพชฺฌงฺคา, วิเวก\-จริยฏฺฐํ พุชฺฌนฺตีติ โพชฺฌงฺคา, วิราคฏฺฐํ พุชฺฌนฺตีติ โพชฺฌงฺคา, วิราค\-จริยฏฺฐํ พุชฺฌนฺตีติ โพชฺฌงฺคา, นิโรธฏฺฐํ พุชฺฌนฺตีติ โพชฺฌงฺคา, นิโรธ\-จริยฏฺฐํ พุชฺฌนฺตีติ โพชฺฌงฺคา, โวสคฺคฏฺฐํ พุชฺฌนฺตีติ โพชฺฌงฺคา, โวสคฺค\-จริยฏฺฐํ พุชฺฌนฺตีติ โพชฺฌงฺคา, วิมุตฺตฏฺฐํ พุชฺฌนฺตีติ โพชฺฌงฺคา, วิมุตฺติ\-จริยฏฺฐํ พุชฺฌนฺตีติ โพชฺฌงฺคา.}

\prose{ฉนฺทฏฺฐํ พุชฺฌนฺตีติ โพชฺฌงฺคา, ฉนฺทสฺส มูลฏฺฐํ พุชฺฌนฺตีติ โพชฺฌงฺคา, ฉนฺทสฺส ปาทฏฺฐํ พุชฺฌนฺตีติ โพชฺฌงฺคา, ฉนฺทสฺส ปธานฏฺฐํ พุชฺฌนฺตีติ โพชฺฌงฺคา, ฉนฺทสฺส อิชฺฌนฏฺฐํ พุชฺฌนฺตีติ โพชฺฌงฺคา, ฉนฺทสฺส อธิโมกฺขฏฺฐํ พุชฺฌนฺตีติ โพชฺฌงฺคา, ฉนฺทสฺส ปคฺคหฏฺฐํ พุชฺฌนฺตีติ โพชฺฌงฺคา, ฉนฺทสฺส อุปฏฺฐานฏฺฐํ พุชฺฌนฺตีติ โพชฺฌงฺคา, ฉนฺทสฺส อวิกฺเขปฏฺฐํ พุชฺฌนฺตีติ โพชฺฌงฺคา, ฉนฺทสฺส ทสฺสนฏฺฐํ พุชฺฌนฺตีติ โพชฺฌงฺคา.}

\prose{วีริยฏฺฐํ พุชฺฌนฺตีติ โพชฺฌงฺคา ฯเปฯ. จิตฺตฏฺฐํ พุชฺฌนฺตีติ โพชฺฌงฺคา ฯเปฯ. วีมํสฏฺฐํ พุชฺฌนฺตีติ โพชฺฌงฺคา, วีมํสาย มูลฏฺฐํ พุชฺฌนฺตีติ โพชฺฌงฺคา, วีมํสาย ปาทฏฺฐํ พุชฺฌนฺตีติ โพชฺฌงฺคา, วีมํสาย ปธานฏฺฐํ พุชฺฌนฺตีติ โพชฺฌงฺคา. วีมํสาย อิชฺฌนฏฺฐํ พุชฺฌนฺตีติ โพชฺฌงฺคา, วีมํสาย อธิโมกฺขฏฺฐํ พุชฺฌนฺตีติ โพชฺฌงฺคา, วีมํสาย ปคฺคหฏฺฐํ พุชฺฌนฺตีติ โพชฺฌงฺคา, วีมํสาย อุปฏฺฐานฏฺฐํ พุชฺฌนฺตีติ โพชฺฌงฺคา, วีมํสาย อวิกฺเขปฏฺฐํ พุชฺฌนฺตีติ โพชฺฌงฺคา, วีมํสาย ทสฺสนฏฺฐํ พุชฺฌนฺตีติ โพชฺฌงฺคา.}

\prose{ทุกฺขฏฺฐํ พุชฺฌนฺตีติ โพชฺฌงฺคา. ทุกฺขสฺส ปีฬนฏฺฐํ พุชฺฌนฺตีติ โพชฺฌงฺคา, ทุกฺขสฺส สงฺขตฏฺฐํ พุชฺฌนฺตีติ โพชฺฌงฺคา, ทุกฺขสฺส สนฺตาปฏฺฐํ พุชฺฌนฺตีติ โพชฺฌงฺคา, ทุกฺขสฺส วิปริณามฏฺฐํ พุชฺฌนฺตีติ โพชฺฌงฺคา. สมุทยฏฺฐํ พุชฺฌนฺตีติ โพชฺฌงฺคา. สมุทยสฺส อายูหนฏฺฐํ. นิทานฏฺฐํ. สญฺโญคฏฺฐํ. ปลิโพธฏฺฐํ พุชฺฌนฺตีติ โพชฺฌงฺคา. นิโรธฏฺฐํ พุชฺฌนฺตีติ โพชฺฌงฺคา, นิโรธสฺส นิสฺสรณฏฺฐํ. วิเวกฏฺฐํ. อสงฺขตฏฺฐํ. อมตฏฺฐํ พุชฺฌนฺตีติ โพชฺฌงฺคา. มคฺคฏฺฐํ พุชฺฌนฺตีติ โพชฺฌงฺคา, มคฺคสฺส นิยฺยานฏฺฐํ. เหตุฏฺฐํ. ทสฺสนฏฺฐํ. อาธิปเตยฺยฏฺฐํ พุชฺฌนฺตีติ โพชฺฌงฺคา.}

\prose{\csromanfolio{309}ตถฏฺฐํ พุชฺฌนฺตีติ โพชฺฌงฺคา, อนตฺตฏฺฐํ พุชฺฌนฺตีติ โพชฺฌงฺคา, สจฺจฏฺฐํ พุชฺฌนฺตีติ โพชฺฌงฺคา, ปฏิเวธฏฺฐํ พุชฺฌนฺตีติ โพชฺฌงฺคา, อภิชาน\-นฏฺฐํ พุชฺฌนฺตีติ โพชฺฌงฺคา, ปริชาน\-นฏฺฐํ พุชฺฌนฺตีติ โพชฺฌงฺคา, ธมฺมฏฺฐํ พุชฺฌนฺตีติ โพชฺฌงฺคา, ธาตุฏฺฐํ พุชฺฌนฺตีติ โพชฺฌงฺคา, ญาตฏฺฐํ พุชฺฌนฺตีติ โพชฺฌงฺคา, สจฺฉิกิริยฏฺฐํ พุชฺฌนฺตีติ โพชฺฌงฺคา, ผสฺสนฏฺฐํ พุชฺฌนฺตีติ โพชฺฌงฺคา, อภิสมยฏฺฐํ พุชฺฌนฺตีติ โพชฺฌงฺคา.}

\prose{เนกฺขมฺมํ พุชฺฌนฺตีติ โพชฺฌงฺคา, อพฺยาปาทํ พุชฺฌนฺตีติ โพชฺฌงฺคา, อาโลกสญฺญํ พุชฺฌนฺตีติ โพชฺฌงฺคา, อวิกฺเขปํ พุชฺฌนฺตีติ โพชฺฌงฺคา, ธมฺม\-ววตฺถานํ พุชฺฌนฺตีติ โพชฺฌงฺคา, ญาณํ พุชฺฌนฺตีติ โพชฺฌงฺคา, ปาโมชฺชํ พุชฺฌนฺตีติ โพชฺฌงฺคา, ปฐมํ ฌานํ พุชฺฌนฺตีติ โพชฺฌงฺคา ฯเปฯ อรหตฺตมคฺคํ พุชฺฌนฺตีติ โพชฺฌงฺคา. อรหตฺตผล\-สมาปตฺติํ พุชฺฌนฺตีติ โพชฺฌงฺคา.}

\prose{อธิโมกฺข\-ฏฺเฐน สทฺธินฺทฺริยํ พุชฺฌนฺตีติ โพชฺฌงฺคา ฯเปฯ. ทสฺสนฏฺเฐน ปญฺญินฺทฺริยํ พุชฺฌนฺตีติ โพชฺฌงฺคา.  อสฺสทฺธิเย อกมฺปิยฏฺเฐน สทฺธาพลํ พุชฺฌนฺตีติ โพชฺฌงฺคา ฯเปฯ. อวิชฺชาย อกมฺปิยฏฺเฐน ปญฺญาพลํ พุชฺฌนฺตีติ โพชฺฌงฺคา.  อุปฏฺฐาน\-ฏฺเฐน สติสมฺโพชฺฌงฺคํ พุชฺฌนฺตีติ โพชฺฌงฺคา ฯเปฯ. ปฏิสงฺขา\-นฏฺเฐน อุเปกฺขา\-สมฺโพชฺฌงฺคํ พุชฺฌนฺตีติ โพชฺฌงฺคา.}

\prose{ทสฺสนฏฺเฐน สมฺมาทิฏฺฐิํ พุชฺฌนฺตีติ โพชฺฌงฺคา ฯเปฯ. อวิกฺเขป\-ฏฺเฐน สมฺมาสมาธิํ พุชฺฌนฺตีติ โพชฺฌงฺคา.  อาธิปเตยฺย\-ฏฺเฐน อินฺทฺริยํ พุชฺฌนฺตีติ โพชฺฌงฺคา, อกมฺปิยฏฺเฐน พลํ พุชฺฌนฺตีติ โพชฺฌงฺคา, นิยฺยานฏฺฐํ พุชฺฌนฺตีติ โพชฺฌงฺคา, เหตุฏฺเฐน มคฺคํ พุชฺฌนฺตีติ โพชฺฌงฺคา, อุปฏฺฐาน\-ฏฺเฐน สติปฏฺฐานํ พุชฺฌนฺตีติ โพชฺฌงฺคา, ปทหนฏฺเฐน สมฺมปฺปธานํ พุชฺฌนฺตีติ โพชฺฌงฺคา, อิชฺฌนฏฺเฐน อิทฺธิปาทํ พุชฺฌนฺตีติ โพชฺฌงฺคา, ตถฏฺเฐน สจฺจํ พุชฺฌนฺตีติ โพชฺฌงฺคา, อวิกฺเขป\-ฏฺเฐน สมถํ พุชฺฌนฺตีติ โพชฺฌงฺคา. อนุปสฺสนฏฺเฐน วิปสฺสนํ ฯเปฯ เอกรสฏฺเฐน สมถ\-วิปสฺสนํ พุชฺฌนฺตีติ โพชฺฌงฺคา, อนติวตฺตน\-ฏฺเฐน ยุคนทฺธํ พุชฺฌนฺตีติ โพชฺฌงฺคา, สํวรฏฺเฐน สีลวิสุทฺธิํ พุชฺฌนฺตีติ โพชฺฌงฺคา, อวิกฺเขป\-ฏฺเฐน จิตฺตวิสุทฺธิํ พุชฺฌนฺตีติ โพชฺฌงฺคา, ทสฺสนฏฺเฐน ทิฏฺฐิ\-วิสุทฺธิํ พุชฺฌนฺตีติ โพชฺฌงฺคา, มุตฺตฏฺเฐน วิโมกฺขํ พุชฺฌนฺตีติ โพชฺฌงฺคา, ปฏิเวธ\-ฏฺเฐน วิชฺชํ พุชฺฌนฺตีติ โพชฺฌงฺคา, ปริจฺจาค\-ฏฺเฐน วิมุตฺติํ พุชฺฌนฺตีติ โพชฺฌงฺคา, สมุจฺเฉท\-ฏฺเฐน ขเย ญาณํ พุชฺฌนฺตีติ โพชฺฌงฺคา, ปฏิปฺปสฺสทฺธ\-ฏฺเฐน อนุปฺปาเท ญาณํ พุชฺฌนฺตีติ โพชฺฌงฺคา.}

\prose{ฉนฺทํ มูลฏฺเฐน พุชฺฌนฺตีติ โพชฺฌงฺคา, มนสิการํ สมุฏฺฐาน\-ฏฺเฐน พุชฺฌนฺตีติ โพชฺฌงฺคา, ผสฺสํ สโมธาน\-ฏฺเฐน พุชฺฌนฺตีติ โพชฺฌงฺคา, เวทนํ สโมสรณ\-ฏฺเฐน พุชฺฌนฺตีติ โพชฺฌงฺคา, สมาธิํ ปมุขฏฺเฐน พุชฺฌนฺตีติ โพชฺฌงฺคา, สติํ}

\prose{\csromanfolio{310}อาธิปเตยฺย\-ฏฺเฐน พุชฺฌนฺตีติ โพชฺฌงฺคา, ปญฺญํ ตทุตฺต\-รฏฺเฐน พุชฺฌนฺตีติ โพชฺฌงฺคา, วิมุตฺติํ สารฏฺเฐน พุชฺฌนฺตีติ โพชฺฌงฺคา, อมโตคธํ นิพฺพานํ ปริโยสาน\-ฏฺเฐน พุชฺฌนฺตีติ โพชฺฌงฺคา.}

\proseitem{20}{สาวตฺถิ\-นิทานํ.  ตตฺร โข อายสฺมา สาริปุตฺโต ภิกฺขู อามนฺเตสิ “อาวุโส ภิกฺขโว”ติ\csromansharedfootnote{62456211b07a0e89}{อาวุโสติ (สฺยา) สํ 3. 64 ปิฏฺเฐ ปสฺสิตพฺพา.}. “อาวุโส”ติ โข เต ภิกฺขู อายสฺมโต สาริปุตฺตสฺส ปจฺจสฺโสสุํ. อายสฺมา สาริปุตฺโต เอตทโวจ\csromandash{}}

\prose{สตฺติเม อาวุโส โพชฺฌงฺคา. กตเม สตฺต? สติสมฺโพชฺฌงฺโค ธมฺมวิจย\-สมฺโพชฺฌงฺโค ฯเปฯ อุเปกฺขา\-สมฺโพชฺฌงฺโค, อิเม โข อาวุโส สตฺต โพชฺฌงฺคา. อิเมสํ ขฺวาหํ อาวุโส สตฺตนฺนํ โพชฺฌงฺคานํ เยน เยน โพชฺฌงฺเคน อากงฺขามิ ปุพฺพณฺหสมยํ วิหริตุํ, เตน เตน โพชฺฌงฺเคน ปุพฺพณฺหสมยํ วิหรามิ. เยน เยน โพชฺฌงฺเคน อากงฺขามิ มชฺฌนฺหิกสมยํ ฯเปฯ สายนฺหสมยํ วิหริตุํ, เตน เตน โพชฺฌงฺเคน สายนฺหสมยํ วิหรามิ. สติสมฺโพชฺฌงฺโค อิติ เจ เม อาวุโส โหติ, “อปฺปมาโณ”ติ เม โหติ, “สุสมารทฺโธ”ติ เม โหติ, “ติฏฺฐนฺตํ จ นํ\csromansharedfootnote{040138212259c15b}{ติฏฺฐนฺตํ จรํ (สฺยา) สํ 3. 64 ปิฏฺเฐ ปสฺสิตพฺพา.} ติฏฺฐตี”ติ ปชานามิ, สเจปิ เม จวติ, “อิทปฺปจฺจยา เม จวตี”ติ ปชานามิ. ธมฺมวิจย\-สมฺโพชฺฌงฺโค ฯเปฯ อุเปกฺขา\-สมฺโพชฺฌงฺโค อิติ เจ เม อาวุโส โหติ, “อปฺปมาโณ”ติ เม โหติ, “สุสมารทฺโธ”ติ เม โหติ, “ติฏฺฐนฺตํ จ นํ ติฏฺฐตี”ติ ปชานามิ. สเจปิ เม จวติ, “อิทปฺปจฺจยา เม จวตี”ติ ปชานามิ.}

\prose{เสยฺยถาปิ อาวุโส รญฺโญ วา ราช\-มหามตฺตสฺส วา นานารตฺตานํ ทุสฺสานํ ทุสฺสกรณฺฑโก ปูโร อสฺส, โส ยญฺญเทว ทุสฺสยุคํ อากงฺเขยฺย ปุพฺพณฺหสมยํ ปารุปิตุํ, ตํ ตเทว ทุสฺสยุคํ ปุพฺพณฺหสมยํ ปารุเปยฺย. ยญฺญเทว ทุสฺสยุคํ อากงฺเขยฺย มชฺฌนฺหิกสมยํ ฯเปฯ สายนฺหสมยํ ปารุปิตุํ, ตํ ตเทว ทุสฺสยุคํ สายนฺหสมยํ ปารุเปยฺย, เอวเมว ขฺวาหํ อาวุโส อิเมสํ สตฺตนฺนํ โพชฺฌงฺคานํ เยน เยน โพชฺฌงฺเคน อากงฺขามิ ปุพฺพณฺหสมยํ วิหริตุํ, เตน เตน โพชฺฌงฺเคน ปุพฺพณฺหสมยํ วิหรามิ. เยน เยน โพชฺฌงฺเคน อากงฺขามิ มชฺฌนฺหิกสมยํ ฯเปฯ \csromanfolio{311}สายนฺหสมยํ วิหริตุํ, เตน เตน โพชฺฌงฺเคน สายนฺหสมยํ วิหรามิ. สติสมฺโพชฺฌงฺโค อิติ เจ เม อาวุโส โหติ, “อปฺปมาโณ”ติ เม โหติ, “สุสมารทฺโธ”ติ เม โหติ, “ติฏฺฐนฺตํ จ นํ ติฏฺฐตี”ติ ปชานามิ, สเจปิ เม จวติ, “อิทปฺปจฺจยา เม จวตี”ติ ปชานามิ ฯเปฯ อุเปกฺขา\-สมฺโพชฺฌงฺโค อิติ เจ เม อาวุโส โหติ, “อปฺปมาโณ”ติ เม โหติ, “สุสมารทฺโธ”ติ เม โหติ, “ติฏฺฐนฺตํ จ นํ ติฏฺฐตี”ติ ปชานามิ, สเจปิ เม จวติ, “อิทปฺปจฺจยา เม จวตี”ติ ปชานามิ.}

\csromanmatikaanchor{mtk.94}
\csromantocmarkref{section}{สุตฺตนฺตนิทฺเทส}{mtk.94}
\bookmark[dest={mtk.94},level=1]{สุตฺตนฺตนิทฺเทส}
\csromanheader{1}{สุตฺตนฺต\-นิทฺเทส}
\proseitem{21}{กถํ \textbf{สติสมฺโพชฺฌงฺโค อิติ เจ เม โหตี}ติ โพชฺฌงฺโค? ยาวตา นิโรธู’ปฏฺฐาติ, ตาวตา สติสมฺโพชฺฌงฺโค, อิติ เจ เม โหตีติ โพชฺฌงฺโค. เสยฺยถาปิ เตลปฺปทีปสฺส ฌายโต ยาวตา อจฺจิ, ตาวตา วณฺโณ. ยาวตา วณฺโณ, ตาวตา อจฺจิ. เอวเมว ยาวตา นิโรธู’ปฏฺฐาติ, ตาวตา สติสมฺโพชฺฌงฺโค. อิติ เจ เม โหตีติ โพชฺฌงฺโค.}

\prose{กถํ อปฺปมาโณ อิติ เจ เม โหตีติ โพชฺฌงฺโค? ปมาณพทฺธา\csromansharedfootnote{6bc093930aabd778}{ปมาณวนฺตา (สฺยา), อฏฺฐกถา โอโลเกตพฺพา.} กิเลสา สพฺเพ จ ปริยุฏฺฐานา เย จ สงฺขารา โปโนภวิกา, อปฺปมาโณ นิโรโธ อจลฏฺเฐน อสงฺขต\-ฏฺเฐน ยาวตา นิโรธู’ปฏฺฐาติ, ตาวตา อปฺปมาโณ อิติ เจ เม โหตีติ โพชฺฌงฺโค.}

\prose{กถํ \textbf{สุสมารทฺโธ อิติ เจ เม โหตี}ติ โพชฺฌงฺโค? วิสมา กิเลสา สพฺเพ จ ปริยุฏฺฐานา เย จ สงฺขารา โปโนภวิกา, สมธมฺโม นิโรโธ สนฺตฏฺเฐน ปณีตฏฺเฐน. ยาวตา นิโรธู’ปฏฺฐาติ, ตาวตา สุสมารทฺโธ, อิติ เจ เม โหตีติ โพชฺฌงฺโค.}

\prose{กถํ \textbf{“ติฏฺฐนฺตํ จ นํ ติฏฺฐตี”ติ ปชานามิ, สเจปิ จวติ, “อิทปฺปจฺจยา จวตี”ติ ปชานามิ?} กติหากาเรหิ สติสมฺโพชฺฌงฺโค ติฏฺฐติ, กติหากาเรหิ สติสมฺโพชฺฌงฺโค จวติ? อฏฺฐหากาเรหิ สติสมฺโพชฺฌงฺโค ติฏฺฐติ, อฏฺฐหากาเรหิ สติสมฺโพชฺฌงฺโค จวติ.}

\prose{\csromanfolio{312}กตเมหิ \textbf{อฏฺฐหากาเรหิ สติสมฺโพชฺฌงฺโค ติฏฺฐติ?} อนุปฺปาทํ อาวชฺชิตตฺตา สติสมฺโพชฺฌงฺโค ติฏฺฐติ, อุปฺปาทํ อนาวชฺชิตตฺตา สติสมฺโพชฺฌงฺโค ติฏฺฐติ, อปฺปวตฺตํ อาวชฺชิตตฺตา สติสมฺโพชฺฌงฺโค ติฏฺฐติ, ปวตฺตํ อนาวชฺชิตตฺตา สติสมฺโพชฺฌงฺโค ติฏฺฐติ, อนิมิตฺตํ อาวชฺชิตตฺตา สติสมฺโพชฺฌงฺโค ติฏฺฐติ, นิมิตฺตํ อนาวชฺชิตตฺตา สติสมฺโพชฺฌงฺโค ติฏฺฐติ, นิโรธํ อาวชฺชิตตฺตา สติสมฺโพชฺฌงฺโค ติฏฺฐติ, สงฺขาเร อนาวชฺชิตตฺตา สติสมฺโพชฺฌงฺโค ติฏฺฐติ, อิเมหิ อฏฺฐหากาเรหิ สติสมฺโพชฺฌงฺโค ติฏฺฐติ.}

\prose{\textbf{กตเมหิ อฏฺฐหากาเรหิ} สติสมฺโพชฺฌงฺโค จวติ? อุปฺปาทํ อาวชฺชิตตฺตา สติสมฺโพชฺฌงฺโค จวติ, อนุปฺปาทํ อนาวชฺชิตตฺตา สติสมฺโพชฺฌงฺโค จวติ, ปวตฺตํ อาวชฺชิตตฺตา สติสมฺโพชฺฌงฺโค จวติ, อปฺปวตฺตํ อนาวชฺชิตตฺตา สติสมฺโพชฺฌงฺโค จวติ, นิมิตฺตํ อาวชฺชิตตฺตา สติสมฺโพชฺฌงฺโค จวติ, อนิมิตฺตํ อนาวชฺชิตตฺตา สติสมฺโพชฺฌงฺโค จวติ, สงฺขาเร อาวชฺชิตตฺตา สติสมฺโพชฺฌงฺโค จวติ, นิโรธํ อนาวชฺชิตตฺตา สติสมฺโพชฺฌงฺโค จวติ. อิเมหิ อฏฺฐหากาเรหิ สติสมฺโพชฺฌงฺโค จวติ.  เอวํ “ติฏฺฐนฺตํ จ นํ ติฏฺฐตี”ติ ปชานามิ. สเจปิ จวติ, “อิทปฺปจฺจยา เม จวตี”ติ ปชานามิ ฯเปฯ.}

\prose{กถํ อุเปกฺขา\-สมฺโพชฺฌงฺโค อิติ เจ เม โหตีติ โพชฺฌงฺโค? ยาวตา นิโรธู’ปฏฺฐาติ, ตาวตา อุเปกฺขา\-สมฺโพชฺฌงฺโค อิติ เจ เม โหตีติ โพชฺฌงฺโค. เสยฺยถาปิ เตลปฺปทีปสฺส ฌายโต ยาวตา อจฺจิ, ตาวตา วณฺโณ. ยาวตา วณฺโณ, ตาวตา อจฺจิ. เอวเมว ยาวตา นิโรธู’ปฏฺฐาติ, ตาวตา อุเปกฺขา\-สมฺโพชฺฌงฺโค อิติ เจ โหตีติ โพชฺฌงฺโค.}

\prose{กถํ อปฺปมาโณ อิติ เจ โหตีติ โพชฺฌงฺโค? ปมาณพทฺธา กิเลสา สพฺเพ จ ปริยุฏฺฐานา เย จ สงฺขารา โปโนภวิกา, อปฺปมาโณ นิโรโธ อจลฏฺเฐน อสงฺขต\-ฏฺเฐน. ยาวตา นิโรธู’ปฏฺฐาติ, ตาวตา อปฺปมาโณ อิติ เจ โหตีติ โพชฺฌงฺโค.}

\prose{กถํ \textbf{สุสมารทฺโธ อิติ เจ โหตี}ติ โพชฺฌงฺโค? วิสมา กิเลสา สพฺเพ จ ปริยุฏฺฐานา เย จ สงฺขารา โปโนภวิกา, สมธมฺโม \csromanfolio{313}นิโรโธ สนฺตฏฺเฐน ปณีตฏฺเฐน. ยาวตา นิโรธู\-ปฏฺฐาติ, ตาวตา สุสมารทฺโธ, อิติ เจ โหตีติ โพชฺฌงฺโค.}

\prose{กถํ “ติฏฺฐนฺตํ จ นํ ติฏฺฐตี”ติ ปชานามิ. สเจปิ จวติ, “อิทปฺปจฺจยา เม จวตี”ติ ปชานามิ? กติหากาเรหิ อุเปกฺขา\-สมฺโพชฺฌงฺโค ติฏฺฐติ, กติหากาเรหิ อุเปกฺขา\-สมฺโพชฺฌงฺโค จวติ? อฏฺฐหากาเรหิ อุเปกฺขา\-สมฺโพชฺฌงฺโค ติฏฺฐติ, อฏฺฐหากาเรหิ อุเปกฺขา\-สมฺโพชฺฌงฺโค จวติ.}

\prose{กตเมหิ อฏฺฐหากาเรหิ อุเปกฺขา\-สมฺโพชฺฌงฺโค ติฏฺฐติ? อนุปฺปาทํ อาวชฺชิตตฺตา อุเปกฺขา\-สมฺโพชฺฌงฺโค ติฏฺฐติ, อุปฺปาทํ อนาวชฺชิตตฺตา อุเปกฺขา\-สมฺโพชฺฌงฺโค ติฏฺฐติ, อปฺปวตฺตํ อาวชฺชิตตฺตา อุเปกฺขา\-สมฺโพชฺฌงฺโค ติฏฺฐติ, ปวตฺตํ อนาวชฺชิตตฺตา อุเปกฺขา\-สมฺโพชฺฌงฺโค ติฏฺฐติ, อนิมิตฺตํ อาวชฺชิตตฺตา อุเปกฺขา\-สมฺโพชฺฌงฺโค ติฏฺฐติ, นิมิตฺตํ อนาวชฺชิตตฺตา อุเปกฺขา\-สมฺโพชฺฌงฺโค ติฏฺฐติ, นิโรธํ อาวชฺชิตตฺตา อุเปกฺขา\-สมฺโพชฺฌงฺโค ติฏฺฐติ, สงฺขาเร อนาวชฺชิตตฺตา อุเปกฺขา\-สมฺโพชฺฌงฺโค ติฏฺฐติ. อิเมหิ อฏฺฐหากาเรหิ อุเปกฺขา\-สมฺโพชฺฌงฺโค ติฏฺฐติ.}

\prosewithcloserruled{กตเมหิ อฏฺฐหากาเรหิ อุเปกฺขา\-สมฺโพชฺฌงฺโค จวติ? อุปฺปาทํ อาวชฺชิตตฺตา อุเปกฺขา\-สมฺโพชฺฌงฺโค จวติ, อนุปฺปาทํ อนาวชฺชิตตฺตา อุเปกฺขา\-สมฺโพชฺฌงฺโค จวติ, ปวตฺตํ อาวชฺชิตตฺตา อุเปกฺขา\-สมฺโพชฺฌงฺโค จวติ, อปฺปวตฺตํ อนาวชฺชิตตฺตา อุเปกฺขา\-สมฺโพชฺฌงฺโค จวติ, นิมิตฺตํ อาวชฺชิตตฺตา อุเปกฺขา\-สมฺโพชฺฌงฺโค จวติ, อนิมิตฺตํ อนาวชฺชิตตฺตา อุเปกฺขา\-สมฺโพชฺฌงฺโค จวติ, สงฺขาเร อาวชฺชิตตฺตา อุเปกฺขา\-สมฺโพชฺฌงฺโค จวติ, นิโรธํ อนาวชฺชิตตฺตา อุเปกฺขา\-สมฺโพชฺฌงฺโค จวติ. อิเมหิ อฏฺฐหากาเรหิ อุเปกฺขา\-สมฺโพชฺฌงฺโค จวติ, เอวํ “ติฏฺฐนฺตํ จ นํ ติฏฺฐตี”ติ ปชานามิ. สเจปิ จวติ, “อิทปฺปจฺจยา เม จวตี”ติ ปชานามิ.}{โพชฺฌงฺคกถา นิฏฺฐิตา.}
\proseitem{22}{สาวตฺถิ\-นิทานํ.  เมตฺตาย ภิกฺขเว เจโตวิมุตฺติยา อาเสวิตาย ภาวิตาย พหุลีกตาย ยานีกตาย วตฺถุกตาย \csromanfolio{314}อนุฏฺฐิตาย ปริจิตาย สุสมารทฺธาย เอกาทสานิสํสา ปาฏิกงฺขา\csromandash{}กตเม เอกาทส? สุขํ สุปติ, สุขํ ปฏิพุชฺฌติ, น ปาปกํ สุปินํ ปสฺสติ, มนุสฺสานํ ปิโย โหติ, อมนุสฺสานํ ปิโย โหติ, เทวตา รกฺขนฺติ, นาสฺส อคฺคิ วา วิสํ วา สตฺถํ วา กมติ, ตุวฏํ จิตฺตํ สมาธิยติ, มุขวณฺโณ วิปฺปสีทติ, อสมฺมูโฬฺห กาลํ กโรติ, อุตฺตริ\csromansharedfootnote{ac751aaffc35d373}{อุตฺตริํ (สฺยา, อิ) อํ 3. 542 ปิฏฺเฐ ปสฺสิตพฺพา.} อปฺปฏิวิชฺฌนฺโต พฺรหฺมโลกูปโค โหติ. เมตฺตาย ภิกฺขเว เจโตวิมุตฺติยา อาเสวิตาย ภาวิตาย พหุลีกตาย ยานีกตาย วตฺถุกตาย อนุฏฺฐิตาย ปริจิตาย สุสมารทฺธาย อิเม เอกาทสานิสํสา ปาฏิกงฺขา.}

\prose{อตฺถิ อโนธิโส ผรณา เมตฺตา\-เจโตวิมุตฺติ, อตฺถิ โอธิโส ผรณา เมตฺตา\-เจโตวิมุตฺติ, อตฺถิ ทิสาผรณา เมตฺตา\-เจโตวิมุตฺติ. กติหากาเรหิ อโนธิโส ผรณา เมตฺตา\-เจโตวิมุตฺติ, กติหากาเรหิ โอธิโส ผรณา เมตฺตา\-เจโตวิมุตฺติ, กติหากาเรหิ ทิสาผรณา เมตฺตา\-เจโตวิมุตฺติ? ปญฺจหากาเรหิ อโนธิโส ผรณา เมตฺตา\-เจโตวิมุตฺติ, สตฺตหากาเรหิ โอธิโส ผรณา เมตฺตา\-เจโตวิมุตฺติ, ทสหากาเรหิ ทิสาผรณา เมตฺตา\-เจโตวิมุตฺติ.}

\prose{กตเมหิ ปญฺจหากาเรหิ อโนธิโส ผรณา เมตฺตา\-เจโตวิมุตฺติ? สพฺเพ สตฺตา อเวรา อพฺยาปชฺชา\csromansharedfootnote{33d64628ab896501}{อพฺยาปชฺฌา (สฺยา)} อนีฆา สุขี อตฺตานํ ปริหรนฺตุ. สพฺเพ ปาณา ฯเปฯ. สพฺเพ ภูตา ฯเปฯ. สพฺเพ ปุคฺคลา ฯเปฯ. สพฺเพ อตฺตภาว\-ปริยาปนฺนา อเวรา อพฺยาปชฺชา อนีฆา สุขี อตฺตานํ ปริหรนฺตูติ. อิเมหิ ปญฺจหากาเรหิ อโนธิโส ผรณา เมตฺตา\-เจโตวิมุตฺติ.}

\prose{กตเมหิ สตฺตหากาเรหิ โอธิโส ผรณา เมตฺตา\-เจโตวิมุตฺติ, สพฺพา อิตฺถิโย อเวรา อพฺยาปชฺชา อนีฆา สุขี อตฺตานํ ปริหรนฺตุ. สพฺเพ ปุริสา ฯเปฯ. สพฺเพ อริยา ฯเปฯ. สพฺเพ อนริยา ฯเปฯ. สพฺเพ เทวา ฯเปฯ. สพฺเพ มนุสฺสา ฯเปฯ. สพฺเพ วินิปาติกา อเวรา อพฺยาปชฺชา อนีฆา สุขี อตฺตานํ ปริหรนฺตูติ อิเมหิ สตฺตหากาเรหิ โอธิโส ผรณา เมตฺตา\-เจโตวิมุตฺติ.}

\prose{กตเมหิ ทสหากาเรหิ ทิสาผรณา เมตฺตา\-เจโตวิมุตฺติ, สพฺเพ ปุรตฺถิมาย ทิสาย สตฺตา อเวรา อพฺยาปชฺชา อนีฆา สุขี อตฺตานํ \csromanfolio{315}ปริหรนฺตุ. สพฺเพ ปจฺฉิมาย ทิสาย สตฺตา ฯเปฯ. สพฺเพ อุตฺตราย ทิสาย สตฺตา ฯเปฯ. สพฺเพ ทกฺขิณาย ทิสาย สตฺตา ฯเปฯ. สพฺเพ ปุรตฺถิมาย อนุทิสาย สตฺตา ฯเปฯ. สพฺเพ ปจฺฉิมาย อนุทิสาย สตฺตา ฯเปฯ. สพฺเพ อุตฺตราย อนุทิสาย สตฺตา ฯเปฯ. สพฺเพ ทกฺขิณาย อนุทิสาย สตฺตา ฯเปฯ. สพฺเพ เหฏฺฐิมาย ทิสาย สตฺตา ฯเปฯ. สพฺเพ อุปริมาย ทิสาย สตฺตา อเวรา อพฺยาปชฺชา อนีฆา สุขี อตฺตานํ ปริหรนฺตุ. สพฺเพ ปุรตฺถิมาย ทิสาย ปาณา ฯเปฯ. ภูตา. ปุคฺคลา. อตฺตภาว\-ปริยาปนฺนา. สพฺพา อิตฺถิโย. สพฺเพ ปุริสา. สพฺเพ อริยา. สพฺเพ อนริยา. สพฺเพ เทวา. สพฺเพ มนุสฺสา. สพฺเพ วินิปาติกา อเวรา อพฺยาปชฺชา อนีฆา สุขี อตฺตานํ ปริหรนฺตุ. สพฺเพ ปจฺฉิมาย ทิสาย วินิปาติกา ฯเปฯ. สพฺเพ อุตฺตราย ทิสาย วินิปาติกา. สพฺเพ ทกฺขิณาย ทิสาย วินิปาติกา. สพฺเพ ปุรตฺถิมาย อนุทิสาย วินิปาติกา. สพฺเพ ปจฺฉิมาย อนุทิสาย วินิปาติกา. สพฺเพ อุตฺตราย อนุทิสาย วินิปาติกา. สพฺเพ ทกฺขิณาย อนุทิสาย วินิปาติกา. สพฺเพ เหฏฺฐิมาย ทิสาย วินิปาติกา. สพฺเพ อุปริมาย ทิสาย วินิปาติกา อเวรา อพฺยาปชฺชา อนีฆา สุขี อตฺตานํ ปริหรนฺตูติ อิเมหิ ทสหากาเรหิ ทิสาผรณา เมตฺตา\-เจโตวิมุตฺติ.}

\csromanmatikaanchor{mtk.95}
\csromanmatikaanchor{mtk.96}
\csromantocmarknopage{section}{4. เมตฺตากถา}{mtk.95}
\bookmark[dest={mtk.95},level=1]{4. เมตฺตากถา}
\csromantocmarkref{subsection}{1. อินฺทฺริยวาร}{mtk.96}
\bookmark[dest={mtk.96},level=2]{1. อินฺทฺริยวาร}
\csromanheader{2}{1. \textbf{อินฺทฺริยวาร}}
\proseitem{23}{สพฺเพสํ สตฺตานํ ปีฬนํ วชฺเชตฺวา อปีฬนาย, อุปฆาตํ วชฺเชตฺวา อนุปฆาเตน, สนฺตาปํ วชฺเชตฺวา อสนฺตาเปน ปริยาทานํ วชฺเชตฺวา อปริยาทาเนน, วิเหสํ วชฺเชตฺวา อวิเหสาย สพฺเพ สตฺตา อเวริโน โหนฺตุ, มา เวริโน. สุขิโน โหนฺตุ, มา ทุกฺขิโน. สุขิตตฺตา โหนฺตุ, มา ทุกฺขิตตฺตาติ อิเมหิ อฏฺฐหากาเรหิ สพฺเพ สตฺเต เมตฺตายตีติ เมตฺตา. ตํ ธมฺมํ เจตยตีติ เจโต. สพฺพ\-พฺยาปาทปริยุฏฺฐาเนหิ วิมุจฺจตีติ วิมุตฺติ. เมตฺตา จ เจโต จ วิมุตฺติ จาติ เมตฺตา\-เจโตวิมุตฺติ.}

\prose{สพฺเพ สตฺตา อเวริโน โหนฺตุ, เขมิโน โหนฺตุ, สุขิโน โหนฺตูติ สทฺธาย อธิมุจฺจติ. สทฺธินฺทฺริย\-ปริภาวิตา โหติ เมตฺตา\-เจโตวิมุตฺติ.}

\prose{\csromanfolio{316}สพฺเพ สตฺตา อเวริโน โหนฺตุ, เขมิโน โหนฺตุ. สุขิโน โหนฺตูติ วีริยํ ปคฺคณฺหาติ. วีริยินฺทฺริย\-ปริภาวิตา โหติ เมตฺตา\-เจโตวิมุตฺติ.}

\prose{สพฺเพ สตฺตา อเวริโน โหนฺตุ, เขมิโน โหนฺตุ, สุขิโน โหนฺตูติ สติํ อุปฏฺฐาเปติ. สตินฺทฺริย\-ปริภาวิตา โหติ เมตฺตา\-เจโตวิมุตฺติ.}

\prose{สพฺเพ สตฺตา อเวริโน โหนฺตุ, เขมิโน โหนฺตุ, สุขิโน โหนฺตูติ จิตฺตํ สมาทหติ. สมาธินฺทฺริย\-ปริภาวิตา โหติ เมตฺตา\-เจโตวิมุตฺติ.}

\prose{สพฺเพ สตฺตา อเวริโน โหนฺตุ, เขมิโน โหนฺตุ, สุขิโน โหนฺตูติ ปญฺญาย ปชานาติ. ปญฺญินฺทฺริย\-ปริภาวิตา โหติ เมตฺตา\-เจโตวิมุตฺติ.}

\prose{อิมานิ ปญฺจินฺทฺริยานิ เมตฺตาย เจโตวิมุตฺติยา อาเสวนา โหนฺติ. อิเมหิ ปญฺจหิ อินฺทฺริเยหิ เมตฺตา\-เจโตวิมุตฺติ อาเสวียติ. อิมานิ ปญฺจินฺทฺริยานิ เมตฺตาย เจโตวิมุตฺติยา ภาวนา โหนฺติ. อิเมหิ ปญฺจหิ อินฺทฺริเยหิ เมตฺตา\-เจโตวิมุตฺติ ภาวียติ. อิมานิ ปญฺจินฺทฺริยานิ เมตฺตาย เจโตวิมุตฺติยา พหุลีกตา โหนฺติ. อิเมหิ ปญฺจหิ อินฺทฺริเยหิ เมตฺตา\-เจโตวิมุตฺติ พหุลีกรียติ. อิมานิ ปญฺจินฺทฺริยานิ เมตฺตาย เจโตวิมุตฺติยา อลงฺการา โหนฺติ. อิเมหิ ปญฺจหิ อินฺทฺริเยหิ เมตฺตา\-เจโตวิมุตฺติ สฺวาลงฺกตา โหติ. อิมานิ ปญฺจินฺทฺริยานิ เมตฺตาย เจโตวิมุตฺติยา ปริกฺขารา โหนฺติ. อิเมหิ ปญฺจหิ อินฺทฺริเยหิ เมตฺตา\-เจโตวิมุตฺติ สุปริกฺขตา โหติ. อิมานิ ปญฺจินฺทฺริยานิ เมตฺตาย เจโตวิมุตฺติยา ปริวารา โหนฺติ. อิเมหิ ปญฺจหิ อินฺทฺริเยหิ เมตฺตา\-เจโตวิมุตฺติ สุปริวุตา โหติ. อิมานิ ปญฺจินฺทฺริยานิ เมตฺตาย เจโตวิมุตฺติยา อาเสวนา โหนฺติ, ภาวนา โหนฺติ, พหุลีกตา โหนฺติ, อลงฺการา โหนฺติ, ปริกฺขารา โหนฺติ, ปริวารา โหนฺติ, ปาริปูรี โหนฺติ, สหคตา โหนฺติ, สหชาตา โหนฺติ, สํสฏฺฐา โหนฺติ, สมฺปยุตฺตา โหนฺติ, ปกฺขนฺทนา โหนฺติ, ปสีทนา\csromansharedfootnote{ba9d07747f083655}{สํสีทนา (ก, สีฏฺฐ)} โหนฺติ, สนฺติฏฺฐนา โหนฺติ, วิมุจฺจนา โหนฺติ, “เอตํ สนฺตนฺ”ติ ผสฺสนา โหนฺติ, ยานีกตา โหนฺติ, วตฺถุกตา โหนฺติ, \csromanfolio{317}อนุฏฺฐิตา โหนฺติ, ปริจิตา โหนฺติ, สุสมารทฺธา โหนฺติ, สุภาวิตา โหนฺติ, สฺวาธิฏฺฐิตา โหนฺติ, สุสมุคฺคตา โหนฺติ, สุวิมุตฺตา โหนฺติ นิพฺพตฺเตนฺติ โชเตนฺติ ปตาเปนฺติ.}

\csromanmatikaanchor{mtk.97}
\csromantocmarkref{subsection}{2. พลวาร}{mtk.97}
\bookmark[dest={mtk.97},level=2]{2. พลวาร}
\csromanheader{2}{2. \textbf{พลวาร}}
\proseitem{24}{สพฺเพ สตฺตา อเวริโน โหนฺตุ, เขมิโน โหนฺตุ, สุขิโน โหนฺตูติ อสฺสทฺธิเย น กมฺปติ. สทฺธาพล\-ปริภาวิตา โหติ เมตฺตา\-เจโตวิมุตฺติ.}

\prose{สพฺเพ สตฺตา อเวริโน โหนฺตุ, เขมิโน โหนฺตุ, สุขิโน โหนฺตูติ โกสชฺเชน กมฺปติ. วีริยพล\-ปริภาวิตา โหติ เมตฺตา\-เจโตวิมุตฺติ.}

\prose{สพฺเพ สตฺตา อเวริโน โหนฺตุ, เขมิโน โหนฺตุ, สุขิโน โหนฺตูติ ปมาเท น กมฺปติ. สติพล\-ปริภาวิตา โหติ เมตฺตาเจติวิมุตฺติ.}

\prose{สพฺเพ สตฺตา อเวริโน โหนฺตุ, เขมิโน โหนฺตุ, สุขิโน โหนฺตูติ อุทฺธจฺเจ น กมฺปติ. สมาธิพล\-ปริภาวิตา โหติ เมตฺตาเจติวิมุตฺติ.}

\prose{สพฺเพ สตฺตา อเวริโน โหนฺตุ, เขมิโน โหนฺตุ, สุขิโน โหนฺตูติ อวิชฺชาย น กมฺปติ. ปญฺญาพล\-ปริภาวิตา โหติ เมตฺตา\-เจโตวิมุตฺติ.}

\prose{อิมานิ ปญฺจ พลานิ เมตฺตาย เจโตวิมุตฺติยา อาเสวนา โหนฺติ. อิเมหิ ปญฺจหิ พเลหิ เมตฺตา\-เจโตวิมุตฺติ อาเสวียติ. อิมานิ ปญฺจ พลานิ เมตฺตาย เจโตวิมุตฺติยา ภาวนา โหนฺติ. อิเมหิ ปญฺจหิ พเลหิ เมตฺตา\-เจโตวิมุตฺติ ภาวียติ. อิมานิ ปญฺจ พลานิ เมตฺตาย เจโตวิมุตฺติยา พหุลีกตา โหนฺติ. อิเมหิ ปญฺจหิ พเลหิ เมตฺตา\-เจโตวิมุตฺติ พหุลีกรียติ. อิมานิ ปญฺจ พลานิ เมตฺตาย เจโตวิมุตฺติยา อลงฺการา โหนฺติ. อิเมหิ ปญฺจหิ พเลหิ เมตฺตา\-เจโตวิมุตฺติ สฺวาลงฺกตา โหติ. อิมานิ ปญฺจ พลานิ เมตฺตาย เจโตวิมุตฺติยา ปริกฺขารา โหนฺติ. อิเมหิ ปญฺจหิ พเลหิ เมตฺตา\-เจโตวิมุตฺติ สุปริกฺขตา โหติ. อิมานิ ปญฺจ พลานิ เมตฺตาย เจโตวิมุตฺติยา \csromanfolio{318}ปริวารา โหนฺติ. อิเมหิ ปญฺจหิ พเลหิ เมตฺตา\-เจโตวิมุตฺติ สุปริวุตา โหติ. อิมานิ ปญฺจ พลานิ เมตฺตาย เจโตวิมุตฺติยา อาเสวนา โหนฺติ, ภาวนา โหนฺติ, พหุลีกตา โหนฺติ, อลงฺการา โหนฺติ, ปริกฺขารา โหนฺติ, ปริวารา โหนฺติ, ปาริปูรี โหนฺติ, สหคตา โหนฺติ, สหชาตา โหนฺติ, สํสฏฺฐา โหนฺติ, สมฺปยุตฺตา โหนฺติ, ปกฺขนฺทนา โหนฺติ, ปสีทนา โหนฺติ, สนฺติฏฺฐนา โหนฺติ, วิมุจฺจนา โหนฺติ, “เอตํ สนฺตนฺ”ติ ผสฺสนา โหนฺติ, ยานีกตา โหนฺติ, วตฺถุกตา โหนฺติ, อนุฏฺฐิตา โหนฺติ, ปริจิตา โหนฺติ, สุสมารทฺธา โหนฺติ, สุภาวิตา โหนฺติ, สฺวาธิฏฺฐิตา โหนฺติ, สุสมุคฺคตา โหนฺติ, สุวิมุตฺตา โหนฺติ นิพฺพตฺเตนฺติ โชเตนฺติ ปตาเปนฺติ.}

\csromanmatikaanchor{mtk.98}
\csromantocmarkref{subsection}{3. โพชฺฌงฺควาร}{mtk.98}
\bookmark[dest={mtk.98},level=2]{3. โพชฺฌงฺควาร}
\csromanheader{2}{3. \textbf{โพชฺฌงฺควาร}}
\proseitem{25}{สพฺเพ สตฺตา อเวริโน โหนฺตุ, เขมิโน โหนฺตุ, สุขิโน โหนฺตูติ สติํ อุปฏฺฐาเปติ. สติสมฺโพชฺฌงฺค\-ปริภาวิตา โหติ เมตฺตา\-เจโตวิมุตฺติ.}

\prose{สพฺเพ สตฺตา อเวริโน โหนฺตุ, เขมิโน โหนฺตุ, สุขิโน โหนฺตูติ ปญฺญาย ปวิจินาติ. ธมฺมวิจยสมฺโพชฺฌงฺค\-ปริภาวิตา โหติ เมตฺตา\-เจโตวิมุตฺติ.}

\prose{สพฺเพ สตฺตา อเวริโน โหนฺตุ, เขมิโน โหนฺตุ, สุขิโน โหนฺตูติ วีริยํ ปคฺคณฺหาติ. วีริยสมฺโพชฺฌงฺค\-ปริภาวิตา โหติ เมตฺตา\-เจโตวิมุตฺติ.}

\prose{สพฺเพ สตฺตา อเวริโน โหนฺตุ, เขมิโน โหนฺตุ, สุขิโน โหนฺตูติ ปริฬาหํ ปฏิปฺปสฺสมฺเภติ. ปีติสมฺโพชฺฌงฺค\-ปริภาวิตา โหติ เมตฺตา\-เจโตวิมุตฺติ.}

\prose{สพฺเพ สตฺตา อเวริโน โหนฺตุ, เขมิโน โหนฺตุ, สุขิโน โหนฺตูติ ทุฏฺฐุลฺลํ ปฏิปฺปสฺสมฺเภติ. ปสฺสทฺธิสมฺโพชฺฌงฺค\-ปริภาวิตา โหติ เมตฺตา\-เจโตวิมุตฺติ.}

\prose{สพฺเพ สตฺตา อเวริโน โหนฺตุ, เขมิโน โหนฺตุ, สุขิโน โหนฺตูติ จิตฺตํ สมาทหติ. สมาธิสมฺโพชฺฌงฺค\-ปริภาวิตา โหติ เมตฺตา\-เจโตวิมุตฺติ.}

\prose{\csromanfolio{319}สพฺเพ สตฺตา อเวริโน โหนฺตุ, เขมิโน โหนฺตุ, สุขิโน โหนฺตูติ ญาเณน กิเลเส ปฏิสงฺขาติ. อุเปกฺขาสมฺโพชฺฌงฺค\-ปริภาวิตา โหติ เมตฺตา\-เจโตวิมุตฺติ.}

\prose{อิเม สตฺต โพชฺฌงฺคา เมตฺตาย เจโตวิมุตฺติยา อาเสวนา โหนฺติ. อิเมหิ สตฺตหิ โพชฺฌงฺเคหิ เมตฺตา\-เจโตวิมุตฺติ อาเสวียติ. อิเม สตฺต โพชฺฌงฺคา เมตฺตาย เจโตวิมุตฺติยา ภาวนา โหนฺติ. อิเมหิ สตฺตหิ โพชฺฌงฺเคหิ เมตฺตา\-เจโตวิมุตฺติ ภาวียติ. อิเม สตฺต โพชฺฌงฺคา เมตฺตาย เจโตวิมุตฺติยา พหุลีกตา โหนฺติ. อิเมหิ สตฺตหิ โพชฺฌงฺเคหิ เมตฺตา\-เจโตวิมุตฺติ พหุลีกรียติ. อิเม สตฺต โพชฺฌงฺคา เมตฺตาย เจโตวิมุตฺติยา อลงฺการา โหนฺติ. อิเมหิ สตฺตหิ โพชฺฌงฺเคหิ เมตฺตา\-เจโตวิมุตฺติ สฺวาลงฺกตา โหติ. อิเม สตฺต โพชฺฌงฺคา เมตฺตาย เจโตวิมุตฺติยา ปริกฺขารา โหนฺติ. อิเมหิ สตฺตหิ โพชฺฌงฺเคหิ เมตฺตา\-เจโตวิมุตฺติ สุปริกฺขตา โหติ. อิเม สตฺต โพชฺฌงฺคา เมตฺตาย เจโตวิมุตฺติยา ปริวารา โหนฺติ. อิเมหิ สตฺตหิ โพชฺฌงฺเคหิ เมตฺตา\-เจโตวิมุตฺติ สุปริวุตา โหติ. อิเม สตฺต โพชฺฌงฺคา เมตฺตาย เจโตวิมุตฺติยา อาเสวนา โหนฺติ. ภาวนา โหนฺติ, พหุลีกตา โหนฺติ, อลงฺการา โหนฺติ, ปริกฺขารา โหนฺติ, ปริวารา โหนฺติ, ปาริปูรี โหนฺตี, สหคตา โหนฺติ, สหชาตา โหนฺติ, สํสฏฺฐา โหนฺติ, สมฺปยุตฺตา โหนฺติ, ปกฺขนฺทนา โหนฺติ, ปสีทนา โหนฺติ, สนฺติฏฺฐนา โหนฺติ, วิมุจฺจนา โหนฺติ, “เอตํ สนฺตนฺ”ติ ผสฺสนา โหนฺติ, ยานีกตา โหนฺติ, วตฺถุกตา โหนฺติ, อนุฏฺฐิตา โหนฺติ, ปริจิตา โหนฺติ, สุสมารทฺธา โหนฺติ, สุภาวิตา โหนฺติ, สฺวาธิฏฺฐิตา โหนฺติ, สุสมุคฺคตา โหนฺติ, สุวิมุตฺตา โหนฺติ นิพฺพตฺเตนฺติ โชเตนฺติ ปตาเปนฺติ.}

\csromanmatikaanchor{mtk.99}
\csromantocmarkref{subsection}{4. มคฺคงฺควาร}{mtk.99}
\bookmark[dest={mtk.99},level=2]{4. มคฺคงฺควาร}
\csromanheader{2}{4. \textbf{มคฺคงฺควาร}}
\proseitem{26}{สพฺเพ สตฺตา อเวริโน โหนฺตุ, เขมิโน โหนฺตุ, สุขิโน โหนฺตูติ สมฺมา ปสฺสติ. สมฺมาทิฏฺฐิ\-ปริภาวิตา โหติ เมตฺตา\-เจโตวิมุตฺติ.}

\prose{สพฺเพ สตฺตา อเวริโน โหนฺตุ, เขมิโน โหนฺตุ, สุขิโน โหนฺตูติ สมฺมา อภินิโรเปติ. สมฺมาสงฺกปฺป\-ปริภาวิตา โหติ เมตฺตา\-เจโตวิมุตฺติ.}

\prose{\csromanfolio{320}สพฺเพ สตฺตา อเวริโน โหนฺตุ, เขมิโน โหนฺตุ, สุขิโน โหนฺตูติ สมฺมา ปริคฺคณฺหาติ. สมฺมาวาจา\-ปริภาวิตา โหติ เมตฺตา\-เจโตวิมุตฺติ.}

\prose{สพฺเพ สตฺตา อเวริโน โหนฺตุ, เขมิโน โหนฺตุ, สุขิโน โหนฺตูติ สมฺมา สมุฏฺฐาเปติ. สมฺมากมฺมนฺต\-ปริภาวิตา โหติ เมตฺตา\-เจโตวิมุตฺติ.}

\prose{สพฺเพ สตฺตา อเวริโน โหนฺตุ, เขมิโน โหนฺตุ, สุขิโน โหนฺตูติ สมฺมา โวทาเปติ. สมฺมาอาชีว\-ปริภาวิตา โหติ เมตฺตา\-เจโตวิมุตฺติ.}

\prose{สพฺเพ สตฺตา อเวริโน โหนฺตุ, เขมิโน โหนฺตุ, สุขิโน โหนฺตูติ สมฺมา ปคฺคณฺหาติ. สมฺมาวายาม\-ปริภาวิตา โหติ เมตฺตา\-เจโตวิมุตฺติ.}

\prose{สพฺเพ สตฺตา อเวริโน โหนฺตุ, เขมิโน โหนฺตุ, สุขิโน โหนฺตูติ สมฺมา อุปฏฺฐาเปติ. สมฺมาสติ\-ปริภาวิตา โหติ เมตฺตา\-เจโตวิมุตฺติ.}

\prose{สพฺเพ สตฺตา อเวริโน โหนฺตุ, เขมิโน โหนฺตุ, สุขิโน โหนฺตูติ สมฺมา สมาทหติ. สมฺมา\-สมาธิ\-ปริภาวิตา โหติ เมตฺตา\-เจโตวิมุตฺติ.}

\prose{อิเม อฏฺฐ มคฺคงฺคา เมตฺตาย เจโตวิมุตฺติยา อาเสวนา โหนฺติ. อิเมหิ อฏฺฐหิ มคฺคงฺเคหิ เมตฺตา\-เจโตวิมุตฺติ อาเสวียติ. อิเม อฏฺฐ มคฺคงฺคา เมตฺตาย เจโตวิมุตฺติยา ภาวนา โหนฺติ. อิเมหิ อฏฺฐหิ มคฺคงฺเคหิ เมตฺตา\-เจโตวิมุตฺติ ภาวียติ. อิเม อฏฺฐ มคฺคงฺคา เมตฺตาย เจโตวิมุตฺติยา พหุลีกตา โหนฺติ. อิเมหิ อฏฺฐหิ มคฺคงฺเคหิ เมตฺตา\-เจโตวิมุตฺติ พหุลีกรียติ. อิเม อฏฺฐ มคฺคงฺคา เมตฺตาย เจโตวิมุตฺติยา อลงฺการา โหนฺติ. อิเมหิ อฏฺฐหิ มคฺคงฺเคหิ เมตฺตา\-เจโตวิมุตฺติ สฺวาลงฺกตา โหนฺติ. อิเม อฏฺฐ มคฺคงฺคา เมตฺตาย เจโตวิมุตฺติยา ปริกฺขารา โหนฺติ. อิเมหิ อฏฺฐหิ มคฺคงฺเคหิ เมตฺตา\-เจโตวิมุตฺติ สุปริกฺขตา โหติ. อิเม อฏฺฐ มคฺคงฺคา เมตฺตาย เจโตวิมุตฺติยา ปริวารา โหนฺติ. อิเมหิ อฏฺฐหิ มคฺคงฺเคหิ เมตฺตา\-เจโตวิมุตฺติ สุปริวุตา โหติ. อิเม อฏฺฐ มคฺคงฺคา เมตฺตาย เจโตวิมุตฺติยา \csromanfolio{321}อาเสวนา โหนฺติ, ภาวนา โหนฺติ, พหุลีกตา โหนฺติ, อลงฺการา โหนฺติ, ปริกฺขารา โหนฺติ, ปริวารา โหนฺติ, ปาริปูรี โหนฺติ, สหคตา โหนฺติ, สหชาตา โหนฺติ, สํสฏฺฐา โหนฺติ, สมฺปยุตฺตา โหนฺติ, ปกฺขนฺทนา โหนฺติ, ปสีทนา โหนฺติ, สนฺติฏฺฐนา โหนฺติ, วิมุจฺจนา โหนฺติ, “เอตํ สนฺตนฺ”ติ ผสฺสนา โหนฺติ, ยานีกตา โหนฺติ, วตฺถุกตา โหนฺติ, อนุฏฺฐิตา โหนฺติ, ปริจิตา โหนฺติ, สุสมารทฺธา โหนฺติ, สุภาวิตา โหนฺติ, สฺวาธิฏฺฐิตา โหนฺติ, สุสมุคฺคตา โหนฺติ, สุวิมุตฺตา โหนฺติ นิพฺพตฺเตนฺติ โชเตนฺติ ปตาเปนฺติ.}

\proseitem{27}{สพฺเพสํ ปาณานํ ฯเปฯ. สพฺเพสํ ภูตานํ. สพฺเพสํ ปุคฺคลานํ. สพฺเพสํ อตฺตภาว\-ปริยาปนฺนานํ. สพฺพาสํ อิตฺถีนํ. สพฺเพสํ ปุริสานํ. สพฺเพสํ อริยานํ. สพฺเพสํ อนริยานํ. สพฺเพสํ เทวานํ. สพฺเพสํ มนุสฺสานํ. สพฺเพสํ วินิปาติกานํ ปีฬนํ วชฺเชตฺวา อปีฬนาย อุปฆาตํ วชฺเชตฺวา อนุปฆาเตน, สนฺตาปํ วชฺเชตฺวา อสนฺตาเปน, ปริยาทานํ วชฺเชตฺวา อปริยาทาเนน, วิเหสํ วชฺเชตฺวา อวิเหสาย สพฺเพ วินิปาติกา อเวริโน โหนฺตุ, มา เวริโน. สุขิโน โหนฺตุ, มา ทุกฺขิโน. สุขิตตฺตา โหนฺตุ, มา ทุกฺขิตตฺตาติ อิเมหิ อฏฺฐหากาเรหิ สพฺเพ วินิปาติเก เมตฺตายตีติ เมตฺตา, ตํ ธมฺมํ เจตยตีติ เจโต, สพฺพ\-พฺยาปาทปริยุฏฺฐาเนหิ วิมุจฺจตีติ วิมุตฺติ, เมตฺตา จ เจโต จ วิมุตฺติ จาติ เมตฺตา\-เจโตวิมุตฺติ.}

\prose{สพฺเพ วินิปาติกา อเวริโน โหนฺตุ, เขมิโน โหนฺตุ, สุขิโน โหนฺตูติ สทฺธาย อธิมุจฺจติ. สทฺธินฺทฺริย\-ปริภาวิตา โหติ เมตฺตา\-เจโตวิมุตฺติ ฯเปฯ นิพฺพตฺเตนฺติ โชเตนฺติ ปตาเปนฺติ.}

\prose{สพฺเพสํ ปุรตฺถิมาย ทิสาย สตฺตานํ ฯเปฯ. สพฺเพสํ ปจฺฉิมาย ทิสาย สตฺตานํ. สพฺเพสํ อุตฺตราย ทิสาย สตฺตานํ. สพฺเพสํ ทกฺขิณาย ทิสาย สตฺตานํ. สพฺเพสํ ปุรตฺถิมาย อนุทิสาย สตฺตานํ. สพฺเพสํ ปจฺฉิมาย อนุทิสาย สตฺตานํ. สพฺเพสํ อุตฺตราย อนุทิสาย สตฺตานํ. สพฺเพสํ ทกฺขิณาย อนุทิสาย สตฺตานํ. สพฺเพสํ เหฏฺฐิมาย ทิสาย สตฺตานํ สพฺเพสํ อุปริมาย ทิสาย สตฺตานํ ปีฬนํ วชฺเชตฺวา อปีฬนาย, อุปฆาตํ วชฺเชตฺวา อนุปฆาเตน, สนฺตาปํ วชฺเชตฺวา อสนฺตาเปน, ปริยาทานํ วชฺเชตฺวา อปริยาทาเนน, วิเหสํ วชฺเชตฺวา อวิเหสาย สพฺเพ อุปริมาย ทิสาย สตฺตา อเวริโน โหนฺตุ, มา เวริโน. สุขิโน \csromanfolio{322}โหนฺตุ, มา ทุกฺขิโน. สุขิตตฺตา โหนฺตุ, มา ทุกฺขิตตฺตาติ อิเมหิ อฏฺฐหากาเรหิ สพฺเพ อุปริมาย ทิสาย สตฺเต เมตฺตายตีติ เมตฺตา, ตํ ธมฺมํ เจตยตีติ เจโต, สพฺพ\-พฺยาปาทปริยุฏฺฐาเนหิ วิมุจฺจตีติ วิมุตฺติ, เมตฺตา จ เจโต จ วิมุตฺติ จาติ เมตฺตา\-เจโตวิมุตฺติ.}

\prose{สพฺเพ อุปริมาย ทิสาย สตฺตา อเวริโน โหนฺตุ, เขมิโน โหนฺตุ, สุขิโน โหนฺตูติ สทฺธาย อธิมุจฺจติ. สทฺธินฺทฺริย\-ปริภาวิตา โหติ, เมตฺตา\-เจโตวิมุตฺติ ฯเปฯ นิพฺพตฺเตนฺติ โชเตนฺติ ปตาเปนฺติ.}

\prose{สพฺเพสํ ปุรตฺถิมาย ทิสาย ปาณานํ. ภูตานํ. ปุคฺคลานํ. อตฺตภาว\-ปริยาปนฺนานํ. สพฺพาสํ อิตฺถีนํ. สพฺเพสํ ปุริสานํ. สพฺเพสํ อริยานํ. สพฺเพสํ อนริยานํ. สพฺเพสํ เทวานํ. สพฺเพสํ มนุสฺสานํ. สพฺเพสํ วินิปาติกานํ. สพฺเพสํ ปจฺฉิมาย ทิสาย วินิปาติกานํ. สพฺเพสํ อุตฺตราย ทิสาย วินิปาติกานํ. สพฺเพสํ ทกฺขิณาย ทิสาย วินิปาติกานํ. สพฺเพสํ ปุรตฺถิมาย อนุทิสาย วินิปาติกานํ. สพฺเพสํ ปจฺฉิมาย อนุทิสาย วินิปาติกานํ. สพฺเพสํ อุตฺตราย อนุทิสาย วินิปาติกานํ. สพฺเพสํ ทกฺขิณาย อนุทิสาย วินิปาติกานํ. สพฺเพสํ เหฏฺฐิมาย ทิสาย วินิปาติกานํ. สพฺเพสํ อุปริมาย ทิสาย วินิปาติกานํ ปีฬนํ วชฺเชตฺวา อปีฬนาย, อุปฆาตํ วชฺเชตฺวา อนุปฆาเตน, สนฺตาปํ วชฺเชตฺวา อสนฺตาเปน, ปริยาทานํ วชฺเชตฺวา อปริยาทาเนน, วิเหสํ วชฺเชตฺวา อวิเหสาย สพฺเพ อุปริมาย ทิสาย วินิปาติกา อเวริโน โหนฺตุ, มา เวริโน. สุขิโน โหนฺตุ, มา ทุกฺขิโน. สุขิตตฺตา โหนฺตุ, มา ทุกฺขิตตฺตาติ อิเมหิ อฏฺฐหากาเรหิ สพฺเพ อุปริมาย ทิสาย วินิปาติเก เมตฺตายตีติ เมตฺตา, ตํ ธมฺมํ เจตยตีติ เจโต, สพฺพ\-พฺยาปาทปริยุฏฺฐาเนหิ วิมุจฺจตีติ วิมุตฺติ, เมตฺตา จ เจโต จ วิมุตฺติ จาติ เมตฺตา\-เจโตวิมุตฺติ.}

\prose{สพฺเพ อุปริมาย ทิสาย วินิปาติกา อเวริโน โหนฺตุ, เขมิโน โหนฺตุ, สุขิโน โหนฺตูติ สทฺธาย อธิมุจฺจติ. สทฺธินฺทฺริย\-ปริภาวิตา โหติ เมตฺตา\-เจโตวิมุตฺติ.}

\prose{สพฺเพ อุปริมาย ทิสาย วินิปาติกา อเวริโน โหนฺตุ, เขมิโน โหนฺตุ, สุขิโน โหนฺตูติ วีริยํ ปคฺคณฺหาติ. วีริยินฺทฺริย\-ปริภาวิตา โหติ เมตฺตา\-เจโตวิมุตฺติ.}

\prose{\csromanfolio{323}สติํ อุปฏฺฐาเปติ, สตินฺทฺริย\-ปริภาวิตา โหติ เมตฺตา\-เจโตวิมุตฺติ. จิตฺตํ สมาทหติ, สมาธินฺทฺริย\-ปริภาวิตา โหติ เมตฺตา\-เจโตวิมุตฺติ. ปญฺญาย ปชานาติ, ปญฺญินฺทฺริย\-ปริภาวิตา โหติ เมตฺตา\-เจโตวิมุตฺติ.}

\prose{อิมานิ ปญฺจินฺทฺริยานิ เมตฺตาย เจโตวิมุตฺติยา อาเสวนา โหนฺติ. อิเมหิ ปญฺจหิ อินฺทฺริเยหิ เมตฺตา\-เจโตวิมุตฺติ อาเสวียติ ฯเปฯ นิพฺพตฺเตนฺติ โชเตนฺติ ปตาเปนฺติ.}

\prose{สพฺเพ อุปริมาย ทิสาย วินิปาติกา อเวริโน โหนฺตุ, เขมิโน โหนฺตุ, สุขิโน โหนฺตูติ อสฺสทฺธิเย น กมฺปติ, สทฺธาพล\-ปริภาวิตา โหติ เมตฺตา\-เจโตวิมุตฺติ. โกสชฺเช น กมฺปติ, วีริยพล\-ปริภาวิตา โหติ เมตฺตา\-เจโตวิมุตฺติ. ปมาเท น กมฺปติ, สติพล\-ปริภาวิตา โหติ เมตฺตา\-เจโตวิมุตฺติ. อุทฺธจฺเจ น กมฺปติ, สมาธิพล\-ปริภาวิตา โหติ เมตฺตา\-เจโตวิมุตฺติ. อวิชฺชาย น กมฺปติ, ปญฺญาพล\-ปริภาวิตา โหติ เมตฺตา\-เจโตวิมุตฺติ.}

\prose{อิมานิ ปญฺจ พลานิ เมตฺตาย เจโตวิมุตฺติยา อาเสวนา โหนฺติ. อิเมหิ ปญฺจหิ พเลหิ เมตฺตา\-เจโตวิมุตฺติ อาเสวียติ ฯเปฯ นิพฺพตฺเตนฺติ โชเตนฺติ ปตาเปนฺติ.}

\prose{สพฺเพ อุปริมาย ทิสาย วินิปาติกา อเวริโน โหนฺตุ, เขมิโน โหนฺตุ, สุขิโน โหนฺตูติ สติํ อุปฏฺฐาเปติ, สติสมฺโพชฺฌงฺค\-ปริภาวิตา โหติ เมตฺตา\-เจโตวิมุตฺติ.}

\prose{ปญฺญาย ปวิจินาติ, ธมฺมวิจยสมฺโพชฺฌงฺค\-ปริภาวิตา โหติ เมตฺตา\-เจโตวิมุตฺติ. วีริยํ ปคฺคณฺหาติ, วีริยสมฺโพชฺฌงฺค\-ปริภาวิตา โหติ เมตฺตา\-เจโตวิมุตฺติ. ปริฬาหํ ปฏิปฺปสฺสมฺเภติ, ปีติสมฺโพชฺฌงฺค\-ปริภาวิตา โหติ เมตฺตา\-เจโตวิมุตฺติ. ทุฏฺฐุลฺลํ ปฏิปฺปสฺสมฺเภติ, ปสฺสทฺธิสมฺโพชฺฌงฺค\-ปริภาวิตา โหติ เมตฺตา\-เจโตวิมุตฺติ. จิตฺตํ สมาทหติ, สมาธิสมฺโพชฺฌงฺค\-ปริภาวิตา โหติ เมตฺตา\-เจโตวิมุตฺติ. ญาเณน กิเลเส ปฏิสงฺขาติ, อุเปกฺขาสมฺโพชฺฌงฺค\-ปริภาวิตา โหติ เมตฺตา\-เจโตวิมุตฺติ.}

\prose{อิเม สตฺต โพชฺฌงฺคา เมตฺตาย เจโตวิมุตฺติยา อาเสวนา โหนฺติ. อิเมหิ สตฺตหิ โพชฺฌงฺเคหิ เมตฺตา\-เจโตวิมุตฺติ อาเสวียติ ฯเปฯ นิพฺพตฺเตนฺติ โชเตนฺติ ปตาเปนฺติ.}

\csromanmatikaanchor{mtk.100}
\csromantocmarkref{section}{5. วิราคกถา}{mtk.100}
\bookmark[dest={mtk.100},level=1]{5. วิราคกถา}
\prose{\csromanfolio{324}สพฺเพ อุปริมาย ทิสาย วินิปาติกา อเวริโน โหนฺตุ, เขมิโน โหนฺตุ, สุขิโน โหนฺตูติ สมฺมา ปสฺสติ, สมฺมาทิฏฺฐิ\-ปริภาวิตา โหติ เมตฺตา\-เจโตวิมุตฺติ. สมฺมา อภินิโรเปติ, สมฺมาสงฺกปฺป\-ปริภาวิตา โหติ เมตฺตา\-เจโตวิมุตฺติ. สมฺมา ปริคฺคณฺหาติ, สมฺมาวาจา\-ปริภาวิตา โหติ เมตฺตา\-เจโตวิมุตฺติ. สมฺมา สมุฏฺฐาเปติ, สมฺมากมฺมนฺต\-ปริภาวิตา โหติ เมตฺตา\-เจโตวิมุตฺติ. สมฺมา โวทาเปติ, สมฺมาอาชีว\-ปริภาวิตา โหติ เมตฺตา\-เจโตวิมุตฺติ. สมฺมา ปคฺคณฺหาติ, สมฺมาวายาม\-ปริภาวิตา โหติ เมตฺตา\-เจโตวิมุตฺติ. สมฺมา อุปฏฺฐาติ, สมฺมาสติ\-ปริภาวิตา โหติ เมตฺตา\-เจโตวิมุตฺติ. สมฺมา สมาทหติ, สมฺมา\-สมาธิ\-ปริภาวิตา โหติ เมตฺตา\-เจโตวิมุตฺติ.}

\prosewithcloserruled{อิเม อฏฺฐ มคฺคงฺคา เมตฺตาย เจโตวิมุตฺติยา อาเสวนา โหนฺติ. อิเมหิ อฏฺฐหิ มคฺคงฺเคหิ เมตฺตา\-เจโตวิมุตฺติ อาเสวียติ ฯเปฯ. อิเม อฏฺฐ มคฺคงฺคา เมตฺตาย เจโตวิมุตฺติยา ปริภาวิตา โหนฺติ. อิเมหิ อฏฺฐหิ มคฺคงฺเคหิ เมตฺตา\-เจโตวิมุตฺติ สุปริวุตา โหติ. อิเม อฏฺฐ มคฺคงฺคา เมตฺตาย เจโตวิมุตฺติยา อาเสวนา โหนฺติ, ภาวนา โหนฺติ, พหุลีกตา โหนฺติ, อลงฺการา โหนฺติ, ปริกฺขารา โหนฺติ, ปริวารา โหนฺติ, ปาริปูรี โหนฺติ, สหคตา โหนฺติ, สหชาตา โหนฺติ, สํสฏฺฐา โหนฺติ, สมฺปยุตฺตา โหนฺติ, ปกฺขนฺทนา โหนฺติ, ปสีทนา โหนฺติ, สนฺติฏฺฐนา โหนฺติ, วิมุจฺจนา โหนฺติ, “เอตํ สนฺตนฺ”ติ ผสฺสนา โหนฺติ, ยานีกตา โหนฺติ, วตฺถุกตา โหนฺติ, อนุฏฺฐิตา โหนฺติ, ปริจิตา โหนฺติ, สุสมารทฺธา โหนฺติ, สุภาวิตา โหนฺติ, สฺวาธิฏฺฐิตา โหนฺติ, สุสมุคฺคตา โหนฺติ, สุวิมุตฺตา โหนฺติ นิพฺพตฺเตนฺติ โชเตนฺติ ปตาเปนฺตีติ.}{เมตฺตากถา นิฏฺฐิตา.}
\csromanheader{2}{5. \textbf{วิราคตถา}}
\proseitem{28}{วิราโค มคฺโค, วิมุตฺติ ผลํ. กถํ วิราโค มคฺโค? โสตาปตฺติมคฺค\-กฺขเณ ทสฺสนฏฺเฐน สมฺมาทิฏฺฐิ มิจฺฉาทิฏฺฐิยา วิรชฺชติ, ตทนุวตฺตก\-กิเลเสหิ จ ขนฺเธหิ จ วิรชฺชติ. พหิทฺธา จ สพฺพนิมิตฺเตหิ วิรชฺชติ.}

\vspace{\tipitakaparskip}
\csromancenter{วิราคกถา วิราโค วิราคารมฺมโณ วิราคโคจโร วิราเค สมุทาคโต\csromansharedfootnote{1710181fb9e45489}{สมุปาคโต (สฺยา)} วิราเค ฐิโต วิราเค ปติฏฺฐิโต.}
\prose{\csromanfolio{325}วิราโคติ เทฺว วิราคา นิพฺพานญฺจ วิราโค, เย จ นิพฺพานารมฺมณตา\-ชาตา ธมฺมา สพฺเพ วิราคา โหนฺตีติ วิราคา. สหชาตานิ สตฺตงฺคานิ วิราคํ คจฺฉนฺตีติ วิราโค มคฺโค. เอเตน มคฺเคน พุทฺธา จ สาวกา จ อคตํ ทิสํ นิพฺพานํ คจฺฉนฺตีติ อฏฺฐงฺคิโก มคฺโค, ยาวตา ปุถุสมณ\-พฺราหฺมณานํ ปรปฺปวาทานํ มคฺคา, อยเมว อริโย อฏฺฐงฺคิโก มคฺโค อคฺโค จ เสฏฺโฐ จ ปาโมกฺโข\csromansharedfootnote{1e08688fdf89c250}{วิโมกฺโข (สฺยา), โมกฺโข (ฏฺฐ)} จ อุตฺตโม จ ปวโร จาติ มคฺคานํ อฏฺฐงฺคิโก เสฏฺโฐ.}

\prose{อภินิโรปน\-ฏฺเฐน สมฺมาสงฺกปฺโป มิจฺฉาสงฺกปฺปา วิรชฺชติ ปริคฺคห\-ฏฺเฐน สมฺมาวาจา มิจฺฉาวาจาย วิรชฺชติ. สมุฏฺฐาน\-ฏฺเฐน สมฺมากมฺมนฺโต มิจฺฉากมฺมนฺตา วิรชฺชติ. โวทานฏฺเฐน สมฺมาอาชีโว มิจฺฉาอาชีวา วิรชฺชติ. ปคฺคหฏฺเฐน สมฺมาวายาโม มิจฺฉาวายามา วิรชฺชติ. อุปฏฺฐาน\-ฏฺเฐน สมฺมาสติ มิจฺฉาสติยา วิรชฺชติ. อวิกฺเขป\-ฏฺเฐน สมฺมาสมาธิ มิจฺฉา\-สมาธิโต วิรชฺชติ. ตทนุวตฺตก\-กิเลเสหิ จ ขนฺเธหิ จ วิรชฺชติ. พหิทฺธา จ สพฺพนิมิตฺเตหิ วิรชฺชติ, วิราโค วิราคารมฺมโณ วิราคโคจโร วิราเค สมุทาคโต วิราเค ฐิโต วิราเค ปติฏฺฐิโต.}

\prose{วิราโคติ เทฺว วิราคา นิพฺพานญฺจ วิราโค, เย จ นิพฺพานารมฺมณตา\-ชาตา ธมฺมา สพฺเพ วิราคา โหนฺตีติ วิราคา, สหชาตานิ สตฺตงฺคานิ วิราคํ คจฺฉนฺตีติ วิราโค มคฺโค. เอเตน มคฺเคน พุทฺธา จ สาวกา จ อคตํ ทิสํ นิพฺพานํ คจฺฉนฺตีติ อฏฺฐงฺคิโก มคฺโค, ยาวตา ปุถุสมณ\-พฺราหฺมณานํ ปรปฺปวาทานํ มคฺคา, อยเมว อริโย อฏฺฐงฺคิโก มคฺโค อคฺโค จ เสฏฺโฐ จ ปาโมกฺโข จ อุตฺตโม จ ปวโร จาติ มคฺคานํ อฏฺฐงฺคิโก เสฏฺโฐ.}

\prose{สกทาคามิ\-มคฺคกฺขเณ ทสฺสนฏฺเฐน สมฺมาทิฏฺฐิ ฯเปฯ อวิกฺเขป\-ฏฺเฐน สมฺมาสมาธิ โอฬาริกา กามราค\-สญฺโญชนา ปฏิฆ\-สญฺโญชนา โอฬาริกา กามราคานุสยา ปฏิฆานุสยา วิรชฺชติ, ตทนุวตฺตก\-กิเลเสหิ จ \csromanfolio{326}ขนฺเธหิ จ วิรชฺชติ, พหิทฺธา จ สพฺพนิมิตฺเตหิ วิรชฺชติ, วิราโค วิราคารมฺมโณ วิราคโคจโร วิราเค สมุทาคโต วิราเค ฐิโต วิราเค ปติฏฺฐิโต.}

\prose{วิราโคติ เทฺว วิราคา นิพฺพานญฺจ วิราโค, เย จ นิพฺพานารมฺมณตา\-ชาตา ธมฺมา สพฺเพ วิราคา โหนฺตีติ วิราคา, สหชาตานิ สตฺตงฺคานิ วิราคํ คจฺฉนฺตีติ วิราโค มคฺโค, เอเตน มคฺเคน พุทฺธา จ สาวกา จ อคตํ ทิสํ นิพฺพานํ คจฺฉนฺตีติ อฏฺฐงฺคิโก มคฺโค, ยาวตา ปุถุสมณ\-พฺราหฺมณานํ ปรปฺปวาทานํ มคฺคา. อยเมว อริโย อฏฺฐงฺคิโก มคฺโค อคฺโค จ เสฏฺโฐ จ ปาโมกฺโข จ อุตฺตโม จ ปวโร จาติ มคฺคานํ อฏฺฐงฺคิโก เสฏฺโฐ.}

\prose{อนาคามิ\-มคฺคกฺขเณ ทสฺสนฏฺเฐน สมฺมาทิฏฺฐิ ฯเปฯ อวิกฺเขป\-ฏฺเฐน สมฺมาสมาธิ อณุสหคตา กามราค\-สญฺโญชนา ปฏิฆ\-สญฺโญชนา อณุสหคตา กามราคานุสยา ปฏิฆานุสยา วิรชฺชติ, ตทนุวตฺตก\-กิเลเสหิ จ ขนฺเธหิ จ วิรชฺชติ, พหิทฺธา จ สพฺพนิมิตฺเตหิ วิรชฺชติ, วิราโค วิราคารมฺมโณ ฯเปฯ มคฺคานํ อฏฺฐงฺคิโก เสฏฺโฐ.}

\prose{อรหตฺต\-มคฺคกฺขเณ ทสฺสนฏฺเฐน สมฺมาทิฏฺฐิ ฯเปฯ อวิกฺเขป\-ฏฺเฐน สมฺมาสมาธิ รูปราคา อรูปราคา มานา อุทฺธจฺจา อวิชฺชาย มานานุสยา ภวราคา\-นุสยา อวิชฺชานุสยา วิรชฺชติ, ตทนุวตฺตก\-กิเลเสหิ จ ขนฺเธหิ จ วิรชฺชติ, พหิทฺธา จ สพฺพนิมิตฺเตหิ วิรชฺชติ, วิราโค วิราคารมฺมโณ วิราคโคจโร วิราเค สมุทาคโต วิราเค ฐิโต วิราเค ปติฏฺฐิโต.}

\prose{วิราโคติ เทฺว วิราคา นิพฺพานญฺจ วิราโค, เย จ นิพฺพานารมฺมณตา\-ชาตา ธมฺมา สพฺเพ วิราคา โหนฺตีติ วิราคา, สหชาตานิ สตฺตงฺคานิ วิราคํ คจฺฉนฺตีติ วิราโค มคฺโค, เอเตน มคฺเคน พุทฺธา จ สาวกา จ อคตํ ทิสํ นิพฺพานํ คจฺฉนฺตีติ อฏฺฐงฺคิโก มคฺโค, ยาวตา ปุถุสมณ\-พฺราหฺมณานํ ปรปฺปวาทานํ มคฺคา, อยเมว อริโย อฏฺฐงฺคิโก มคฺโค อคฺโค จ เสฏฺโฐ จ ปาโมกฺโข จ อุตฺตโม จ ปวโร จาติ มคฺคานํ อฏฺฐงฺคิโก เสฏฺโฐ.}

\prose{ทสฺสนวิราโค สมฺมาทิฏฺฐิ, อภินิโรปน\-วิราโค สมฺมาสงฺกปฺโป, ปริคฺคห\-วิราโค สมฺมาวาจา, สมุฏฺฐาน\-วิราโค สมฺมากมฺมนฺโต, โวทานวิราโค}

\vspace{\tipitakaparskip}
\csromancenter{วิราคกถา สมฺมาอาชีโว, ปคฺคหวิราโค สมฺมาวายาโม, อุปฏฺฐาน\-วิราโค สมฺมาสติ, อวิกฺเขป\-วิราโค สมฺมาสมาธิ.  อุปฏฺฐาน\-วิราโค สติสมฺโพชฺฌงฺโค, ปวิจยวิราโค ธมฺมวิจย\-สมฺโพชฺฌงฺโค, ปคฺคหวิราโค วีริย\-สมฺโพชฺฌงฺโค, ผรณวิราโค ปีติสมฺโพชฺฌงฺโค, อุปสมวิราโค ปสฺสทฺธิ\-สมฺโพชฺฌงฺโค, อวิกฺเขป\-วิราโค สมาธิสโมชฺฌงฺโค, ปฏิสงฺขาน\-วิราโค อุเปกฺขา\-สมฺโพชฺฌงฺโค.  อสฺสทฺธิเย อกมฺปิยวิราโค สทฺธาพลํ, โกสชฺเช อกมฺปิยวิราโค วีริยพลํ, ปมาเท อกมฺปิยวิราโค สติพลํ, อุทฺธจฺเจ อกมฺปิยวิราโค สมาธิพลํ, อวิชฺชาย อกมฺปิยวิราโค ปญฺญาพลํ.  อธิโมกฺข\-วิราโค สทฺธินฺทฺริยํ, ปคฺคหวิราโค วีริยินฺทฺริยํ, อุปฏฺฐาน\-วิราโค สตินฺทฺริยํ, อวิกฺเขป\-วิราโค สมาธินฺทฺริยํ, ทสฺสนวิราโค ปญฺญินฺทฺริยํ.  อาธิปเตยฺย\-ฏฺเฐน อินฺทฺริยานิ วิราโค, อกมฺปิยฏฺเฐน พลํ วิราโค, นิยฺยานฏฺเฐน โพชฺฌงฺคา วิราโค, เหตุฏฺเฐน มคฺโค วิราโค, อุปฏฺฐาน\-ฏฺเฐน สติปฏฺฐานา วิราโค, ปทหนฏฺเฐน สมฺมปฺปธานา วิราโค, อิชฺฌนฏฺเฐน อิทฺธิปฺปาทา วิราโค, ตถฏฺเฐน สจฺจา วิราโค, อวิกฺเขป\-ฏฺเฐน สมโถ วิราโค, อนุปสฺสนฏฺเฐน วิปสฺสนา วิราโค, เอกรสฏฺเฐน สมถ\-วิปสฺสนา วิราโค, อนติวตฺตน\-ฏฺเฐน ยุคนทฺธํ วิราโค, สํวรฏฺเฐน สีลวิสุทฺธิ วิราโค, อวิกฺเขป\-ฏฺเฐน จิตฺตวิสุทฺธิ วิราโค, ทสฺสนฏฺเฐน ทิฏฺฐิวิสุทฺธิ วิราโค, วิมุตฺตฏฺเฐน วิโมกฺโข วิราโค, ปฏิเวธ\-ฏฺเฐน วิชฺชา วิราโค, ปริจฺจาค\-ฏฺเฐน วิมุตฺติ วิราโค, สมุจฺเฉท\-ฏฺเฐน ขเย ญาณํ วิราโค, ฉนฺโท มุลฏฺเฐน วิราโค, มนสิกาโร สมุฏฺฐาน\-ฏฺเฐน วิราโค, ผสฺโส สโมธาน\-ฏฺเฐน วิราโค, เวทนา สโมสรณ\-ฏฺเฐน วิราโค, สมาธิ ปมุขฏฺเฐน วิราโค, สติ อาธิปเตยฺย\-ฏฺเฐน วิราโค, ปญฺญา ตทุตฺต\-รฏฺเฐน วิราโค, วิมุตฺติ สารฏฺเฐน วิราโค, อมโตคธํ นิพฺพานํ ปริโยสาน\-ฏฺเฐน มคฺโค.}
\prose{\csromanfolio{327}ทสฺสนมคฺโค สมฺมาทิฏฺฐิ, อภินิโรปน\-มคฺโค สมฺมาสงฺกปฺโป ฯเปฯ อมโตคธํ นิพฺพานํ ปริโยสาน\-ฏฺเฐน มคฺโค, เอวํ วิราโค มคฺโค.}

\proseitem{29}{กถํ วิมุตฺติ ผลํ? โสตาปตฺติ\-ผลกฺขเณ ทสฺสนฏฺเฐน สมฺมาทิฏฺฐิ มิจฺฉาทิฏฺฐิยา วิมุตฺตา โหติ, ตทนุวตฺตก\-กิเลเสหิ จ ขนฺเธหิ จ วิมุตฺตา โหติ, พหิทฺธา จ สพฺพนิมิตฺเตหิ วิมุตฺตา โหติ, วิมุตฺติ วิมุตฺตารมฺมณา วิมุตฺติโคจรา วิมุตฺติยา สมุทาคตา วิมุตฺติยา ฐิตา วิมุตฺติยา ปติฏฺฐิตา. วิมุตฺตีติ เทฺว วิมุตฺติโย นิพฺพานญฺจ วิมุตฺติ, เย จ นิพฺพานารมฺมณตา\-ชาตา ธมฺมา สพฺเพ จ วิมุตฺตา โหนฺตีติ วิมุตฺติ ผลํ.}

\prose{\csromanfolio{328}อภินิโรปน\-ฏฺเฐน สมฺมาสงฺกปฺโป มิจฺฉาสงฺกปฺปา วิมุตฺโต โหติ, ตทนุวตฺตก\-กิเลเสหิ จ ขนฺเธหิ จ วิมุตฺโต โหติ, พหิทฺธา จ สพฺพนิมิตฺเตหิ วิมุตฺโต โหติ, วิมุตฺติ วิมุตฺตารมฺมณา วิมุตฺติโคจรา วิมุตฺติยา สมุทาคตา วิมุตฺติยา ฐิตา วิมุตฺติยา ปติฏฺฐิตา. วิมุตฺตีติ เทฺว วิมุตฺติโย นิพฺพานญฺจ วิมุตฺติ, เย จ นิพฺพานารมฺมณตา\-ชาตา ธมฺมา สพฺเพ วิมุตฺตา โหนฺตีติ วิมุตฺติ ผลํ.}

\prose{ปริคฺคห\-ฏฺเฐน สมฺมาวาจา มิจฺฉาวาจาย วิมุตฺตา โหติ, สมุฏฺฐาน\-ฏฺเฐน สมฺมากมฺมนฺโต มิจฺฉากมฺมนฺตา วิมุตฺโต โหติ. โวทานฏฺเฐน สมฺมาอาชีโว มิจฺฉาอาชีวา วิมุตฺโต โหติ, ปคฺคหฏฺเฐน สมฺมาวายาโม มิจฺฉาวายามา วิมุตฺโต โหติ, อุปฏฺฐาน\-ฏฺเฐน สมฺมาสติ มิจฺฉาสติยา วิมุตฺตา โหติ, อวิกฺเขป\-ฏฺเฐน สมฺมาสมาธิ มิจฺฉา\-สมาธิโต วิมุตฺโต โหติ, ตทนุวตฺตก\-กิเลเสหิ จ ขนฺเธหิ จ วิมุตฺโต โหติ, พหิทฺธา จ สพฺพนิมิตฺเตหิ วิมุตฺโต โหติ, วิมุตฺติ วิมุตฺตารมฺมณา วิมุตฺติโคจรา วิมุตฺติยา สมุทาคตา วิมุตฺติยา ฐิตา วิมุตฺติยา ปติฏฺฐิตา. วิมุตฺตีติ เทฺว วิมุตฺติโย, นิพฺพานญฺจ วิมุตฺติ, เย จ นิพฺพานารมฺมณตา\-ชาตา ธมฺมา สพฺเพ วิมุตฺตา โหนฺตีติ วิมุตฺติ ผลํ.}

\prose{สกทาคามิ\-ผลกฺขเณ ทสฺสนฏฺเฐน สมฺมาทิฏฺฐิ ฯเปฯ อวิกฺเขป\-ฏฺเฐน สมฺมาสมาธิ โอฬาริกา กามราค\-สญฺโญชนา ปฏิฆ\-สญฺโญชนา โอฬาริกา กามราคานุสยา ปฏิฆานุสยา วิมุตฺโต โหติ, ตทนุวตฺตก\-กิเลเสหิ จ ขนฺเธหิ จ วิมุตฺโต โหติ, พหิทฺธา จ สพฺพนิมิตฺเตหิ วิมุตฺโต โหติ, วิมุตฺติ วิมุตฺตารมฺมณา วิมุตฺติโคจรา วิมุตฺติยา สมุทาคตา วิมุตฺติยา ฐิตา วิมุตฺติยา ปติฏฺฐิตา. วิมุตฺตีติ เทฺว วิมุตฺติโย, นิพฺพานญฺจ วิมุตฺติ, เย จ นิพฺพานารมฺมณตา\-ชาตา ธมฺมา สพฺเพ วิมุตฺตา โหนฺตีติ วิมุตฺติ ผลํ.}

\prose{อนาคามิผลกฺเขเณ ทสฺสนฏฺเฐน สมฺมาทิฏฺฐิ ฯเปฯ อวิกฺเขป\-ฏฺเฐน สมฺมาสมาธิ อณุสหคตา กามราค\-สญฺโญชนา ปฏิฆ\-สญฺโญชนา อณุสหคตา กามราคานุสยา ปฏิฆานุสยา วิมุตฺโต โหติ, ตทนุวตฺตก\-กิเลเสหิ จ ขนฺเธหิ จ วิมุตฺโต โหติ, พหิทฺธา จ สพฺพนิมิตฺเตหิ วิมุตฺโต โหติ, วิมุตฺติ วิมุตฺตารมฺมณา วิมุตฺติโคจรา วิมุตฺติยา สมุทาคตา วิมุตฺติยา ฐิตา วิมุตฺติยา ปติฏฺฐิตา ฯเปฯ.}

\prose{\csromanfolio{329}อรหตฺต\-ผลกฺขเณ ทสฺสนฏฺเฐน สมฺมาทิฏฺฐิ ฯเปฯ อวิกฺเขป\-ฏฺเฐน สมฺมาสมาธิ รูปราคา อรูปราคา มานา อุทฺธจฺจา อวิชฺชาย มานานุสยา ภาวราคานุสยา อวิชฺชานุสยา วิมุตฺโต โหติ. ตทนุวตฺตก\-กิเลเสหิ จ ขนฺเธหิ จ วิมุตฺโต โหติ, พหิทฺธา จ สพฺพนิมิตฺเตหิ วิมุตฺโต โหติ, วิมุตฺติ วิมุตฺตารมฺมณา วิมุตฺติโคจรา วิมุตฺติยา สมุทาคตา วิมุตฺติยา ฐิตา วิมุตฺติยา ปติฏฺฐิตา. วิมุตฺตีติ เทฺว วิมุตฺติโย, นิพฺพานญฺจ วิมุตฺติ, เย จ นิพฺพานารมฺมณตา\-ชาตา ธมฺมา สพฺเพ วิมุตฺตา โหนฺตีติ วิมุตฺติ ผลํ.}

\prose{ทสฺสนวิมุตฺติ สมฺมาทิฏฺฐิ ฯเปฯ อวิกฺเขป\-วิมุตฺติ สมฺมาสมาธิ.  อุปฏฺฐาน\-วิมุตฺติ สติสมฺโพชฺฌงฺโค. ปฏิสงฺขาน\-วิมุตฺติ อุเปกฺขา\-สมฺโพชฺฌงฺโค.  อสฺสทฺธิเย อกมฺปิย\-วิมุตฺติ สทฺธาพลํ ฯเปฯ อวิชฺชาย อกมฺปิย\-วิมุตฺติ ปญฺญาพลํ.  อธิโมกฺข\-วิมุตฺติ สทฺธินฺทฺริยํ ฯเปฯ ทสฺสนวิมุตฺติ ปญฺญินฺทฺริยํ.}

\prosewithcloserruled{อาธิปเตยฺย\-ฏฺเฐน อินฺทฺริยา วิมุตฺติ. อกมฺปิยฏฺเฐน พลา วิมุตฺติ, นิยฺยานฏฺเฐน โพชฺฌงฺคา วิมุตฺติ, เหตุฏฺเฐน มคฺโค วิมุตฺติ, อุปฏฺฐาน\-ฏฺเฐน สติปฏฺฐานา วิมุตฺติ, ปทหนฏฺเฐน สมฺมปฺปธานา วิมุตฺติ, อิชฺฌนฏฺเฐน อิทฺธิปาทา วิมุตฺติ, ตถฏฺเฐน สจฺจา วิมุตฺติ, อวิกฺเขป\-ฏฺเฐน สมโถ วิมุตฺติ, อนุปสฺสนฏฺเฐน วิปสฺสนา วิมุตฺติ, เอกรสฏฺเฐน สมถ\-วิปสฺสนา วิมุตฺติ, อนติวตฺตน\-ฏฺเฐน ยุคนทฺธํ วิมุตฺติ, สํวรฏฺเฐน สีลวิสุทฺธิ วิมุตฺติ, อวิกฺเขป\-ฏฺเฐน จิตฺตวิสุทฺธิ วิมุตฺติ, ทสฺสนฏฺเฐน ทิฏฺฐิวิสุทฺธิ วิมุตฺติ, วิมุตฺตฏฺเฐน วิโมกฺโข วิมุตฺติ, ปฏิเวธ\-ฏฺเฐน วิชฺชา วิมุตฺติ, ปริจฺจาค\-ฏฺเฐน วิมุตฺติ วิมุตฺติ, ปฏิปฺปสฺสทฺธิย\-ฏฺเฐน อนุปฺปาเท ญาณํ วิมุตฺติ, ฉนฺโท มูลฏฺเฐน วิมุตฺติ, มนสิกาโร สมุฏฺฐาน\-ฏฺเฐน วิมุตฺติ, ผสฺโส สโมธาน\-ฏฺเฐน วิมุตฺติ, เวทนา สโมสรณ\-ฏฺเฐน วิมุตฺติ, สมาธิ ปมุขฏฺเฐน วิมุตฺติ, สติ อาธิปเตยฺย\-ฏฺเฐน วิมุตฺติ, ปญฺญา ตทุตฺต\-รฏฺเฐน วิมุตฺติ, วิมุตฺติ สารฏฺเฐน วิมุตฺติ, อมโตคธํ นิพฺพานํ ปริโยสาน\-ฏฺเฐน วิมุตฺติ, เอวํ วิมุตฺติ ผลํ. เอวํ วิราโค มคฺโค วิมุตฺติ ผลนฺติ.}{วิราคกถา นิฏฺฐิตา.}
\csromanmatikaanchor{mtk.101}
\csromantocmarknopage{section}{6. ปฏิสมฺภิทากถา}{mtk.101}
\bookmark[dest={mtk.101},level=1]{6. ปฏิสมฺภิทากถา}
\csromanheader{1}{6. ปฏิสมฺภิทา\-กถา}
\csromanmatikaanchor{mtk.102}
\csromantocmarkref{subsection}{1. ธมฺมจกฺกปวตฺตนวาร}{mtk.102}
\bookmark[dest={mtk.102},level=2]{1. ธมฺมจกฺกปวตฺตนวาร}
\csromanheader{2}{1. ธมฺมจกฺก\-ปวตฺตน\-วาร}
\proseitem{30}{เอวํ เม สุตํ\csromandash{}เอกํ สมยํ ภควา พาราณสิยํ วิหรติ อิสิปตเน มิคทาเย. ตตฺร โข ภควา ปญฺจวคฺคิเย ภิกฺขู อามนฺเตสิ}

\prose{\csromanfolio{330}เทฺวเม ภิกฺขเว อนฺตา ปพฺพชิเตน น เสวิตพฺพา. กตเม เทฺว? โย จายํ กาเมสุ กามสุขลฺลิกา\-นุโยโค หีโน คมฺโม โปถุชฺชนิโก อนริโย อนตฺถสํหิโต, โย จายํ อตฺต\-กิลมถา\-นุโยโค ทุกฺโข อนริโย อนตฺถสํหิโต. เอเต โข\csromansharedfootnote{9e7eb5c492968ae0}{เอเต เต (สฺยา, ก, สีฏฺฐ) สํ 3. 369; วิ 3. 14 ปิฏฺเฐสุ ปสฺสิตพฺพา.} ภิกฺขเว อุโภ อนฺเต อนุปคมฺม มชฺฌิมา ปฏิปทา ตถาคเตน อภิสมฺพุทฺธา จกฺขุกรณี ญาณกรณี อุปสมาย อภิญฺญาย สมฺโพธาย นิพฺพานาย สํวตฺตติ.}

\prose{กตมา จ สา ภิกฺขเว มชฺฌิมา ปฏิปทา ตถาคเตน อภิสมฺพุทฺธา จกฺขุกรณี ญาณกรณี อุปสมาย อภิญฺญาย สมฺโพธาย นิพฺพานาย สํวตฺตติ? อยเมว อริโย อฏฺฐงฺคิโก มคฺโค. เสยฺยถิทํ, สมฺมาทิฏฺฐิ ฯเปฯ สมฺมาสมาธิ. อยํ โข สา ภิกฺขเว มชฺฌิมา ปฏิปทา ตถาคเตน อภิสมฺพุทฺธา จกฺขุกรณี ญาณกรณี อุปสมาย อภิญฺญาย สมฺโพธาย นิพฺพานาย สํวตฺตติ.}

\prose{อิทํ โข ปน ภิกฺขเว ทุกฺขํ อริยสจฺจํ\csromandash{}ชาติปิ ทุกฺขา, ชราปิ ทุกฺขา, พฺยาธิปิ ทุกฺโข, มรณมฺปิ ทุกฺขํ, อปฺปิเยหิ สมฺปโยโค ทุกฺโข, ปิเยหิ วิปฺปโยโค ทุกฺโข, ยมฺปิจฺฉํ น ลภติ ตมฺปิ ทุกฺขํ, สํขิตฺเตน ปญฺจุ\-ปาทาน\-กฺขนฺธา ทุกฺขา. อิทํ โข ปน ภิกฺขเว ทุกฺข\-สมุทยํ อริยสจฺจํ\csromandash{}ยายํ ตณฺหา โปโนภวิกา นนฺทิ\-ราค\-สหคตา ตตฺร ตตฺรา\-ภินนฺทินี. เสยฺยถิทํ, กามตณฺหา ภวตณฺหา วิภวตณฺหา. อิทํ โข ปน ภิกฺขเว ทุกฺขนิโรธํ อริยสจฺจํ\csromandash{}โย ตสฺสาเยว ตณฺหาย อเสส\-วิราค\-นิโรโธ จาโค ปฏินิสฺสคฺโค มุตฺติ อนาลโย. อิทํ โข ปน ภิกฺขเว ทุกฺขนิโรธ\-คามินี ปฏิปทา อริยสจฺจํ, อยเมว อริโย อฏฺฐงฺคิโก มคฺโค. เสยฺยถิทํ, สมฺมาทิฏฺฐิ ฯเปฯ สมฺมาสมาธิ.}

\prose{“อิทํ ทุกฺขํ อริยสจฺจนฺ”ติ เม ภิกฺขเว ปุพฺเพ อนนุสฺสุเตสุ ธมฺเมสุ จกฺขุํ อุทปาทิ, ญาณํ อุทปาทิ, ปญฺญา อุทปาทิ, วิชฺชา อุทปาทิ, อาโลโก อุทปาทิ.  “ตํ โข ปนิทํ ทุกฺขํ อริยสจฺจํ ปริญฺเญยฺยนฺ”ติ เม ภิกฺขเว ฯเปฯ “ปริญฺญาตนฺ”ติ เม ภิกฺขเว ปุพฺเพ อนนุสฺสุเตสุ ธมฺเมสุ จกฺขุํ อุทปาทิ, ญาณํ อุทปาทิ, ปญฺญา อุทปาทิ, วิชฺชา อุทปาทิ, อาโลโก อุทปาทิ.}

\prose{\csromanfolio{331}“อิทํ ทุกฺข\-สมุทยํ อริยสจฺจนฺ”ติ เม ภิกฺขเว ปุพฺเพ อนนุสฺสุเตสุ ธมฺเมสุ จกฺขุํ อุทปาทิ, ญาณํ อุทปาทิ, ปญฺญา อุทปาทิ, วิชฺชา อุทปาทิ, อาโลโก อุทปาทิ.  “ตํ โข ปนิทํ ทุกฺข\-สมุทยํ อริยสจฺจํ ปหาตพฺพนฺ”ติ เม ภิกฺขเว ฯเปฯ “ปหีนนฺ”ติ เม ภิกฺขเว ปุพฺเพ อนนุสฺสุเตสุ ธมฺเมสุ จกฺขุํ อุทปาทิ, ญาณํ อุทปาทิ, ปญฺญา อุทปาทิ, วิชฺชา อุทปาทิ, อาโลโก อุทปาทิ.}

\prose{“อิทํ ทุกฺขนิโรธํ อริยสจฺจนฺ”ติ เม ภิกฺขเว ปุพฺเพ อนนุสฺสุเตสุ ธมฺเมสุ จกฺขุํ อุทปาทิ, ญาณํ อุทปาทิ, ปญฺญา อุทปาทิ, วิชฺชา อุทปาทิ, อาโลโก อุทปาทิ.  “ตํ โข ปนิทํ ทุกฺขนิโรธํ อริยสจฺจํ สจฺฉิกาตพฺพนฺ”ติ เม ภิกฺขเว ฯเปฯ “สจฺฉิกตนฺ”ติ เม ภิกฺขเว ปุพฺเพ อนนุสฺสุเตสุ ธมฺเมสุ จกฺขุํ อุทปาทิ, ญาณํ อุทปาทิ, ปญฺญา อุทปาทิ, วิชฺชา อุทปาทิ, อาโลโก อุทปาทิ.}

\prose{“อิทํ ทุกฺขนิโรธ\-คามินี ปฏิปทา อริยสจฺจนฺ”ติ เม ภิกฺขเว ปุพฺเพ อนนุสฺสุเตสุ ธมฺเมสุ จกฺขุํ อุทปาทิ, ญาณํ อุทปาทิ, ปญฺญา อุทปาทิ, วิชฺชา อุทปาทิ, อาโลโก อุทปาทิ.  “ตํ โข ปนิทํ ทุกฺขนิโรธ\-คามินี ปฏิปทา อริยสจฺจํ ภาเวตพฺพนฺ”ติ เม ภิกฺขเว ฯเปฯ “ภาวิตนฺ”ติ เม ภิกฺขเว ปุพฺเพ อนนุสฺสุเตสุ ธมฺเมสุ จกฺขุํ อุทปาทิ, ญาณํ อุทปาทิ, ปญฺญา อุทปาทิ, วิชฺชา อุทปาทิ, อาโลโก อุทปาทิ.}

\prose{ยาว กีวญฺจ เม ภิกฺขเว อิเมสุ จตูสุ อริยสจฺเจสุ เอวํ ติปริวฏฺฏํ ทฺวาทสาการํ ยถาภูตํ ญาณทสฺสนํ น สุวิสุทฺธํ อโหสิ, เนว ตาวาหํ ภิกฺขเว สเทวเก โลเก สมารเก สพฺรหฺมเก สสฺสมณ\-พฺราหฺมณิยา ปชาย สเทว\-มนุสฺสาย “อนุตฺตรํ สมฺมาสมฺโพธิํ อภิสมฺพุทฺโธ”ติ ปจฺจญฺญาสิํ. ยโต จ โข เม ภิกฺขเว อิเมสุ จตูสุ อริยสจฺเจสุ เอวํ ติปริวฏฺฏํ ทฺวาทสาการํ ยถาภูตํ ญาณทสฺสนํ สุวิสุทฺธํ อโหสิ, อถาหํ ภิกฺขเว สเทวเก โลเก สมารเก สพฺรหฺมเก สสฺสมณพฺรหฺมณิยา ปชาย สเทว\-มนุสฺสาย “อนุตฺตรํ สมฺมาสมฺโพธํ อภิสมฺพุทฺโธ”ติ ปจฺจญฺญาสิํ. ญาณญฺจ ปน เม ทสฺสนํ อุทปาทิ. อกุปฺปา เม วิมุตฺติ, อยมนฺติมา ชาติ, นตฺถิ ทานิ ปุนพฺภโวติ.}

\prose{อิทมโวจ ภควา, อตฺตมนา ปญฺจวคฺคิยา ภิกฺขู ภควโต ภาสิตํ อภินนฺทุนฺติ.}

\prose{\csromanfolio{332}อิมสฺมิํ จ ปน เวยฺยากรณสฺมิํ ภญฺญมาเน อายสฺมโต โกณฺฑญฺญสฺส วิรชํ วีตมลํ ธมฺมจกฺขุํ อุทปาทิ “ยํ กิญฺจิ สมุทย\-ธมฺมํ, สพฺพํ ตํ นิโรธธมฺมนฺ”ติ.}

\prose{ปวตฺติเต จ ปน ภควตา ธมฺมจกฺเก ภุมฺมา\csromansharedfootnote{37c882b93d956eaa}{ภูมา (ก)} เทวา สทฺทม\-นุสฺสาเวสุํ “เอตํ ภควตา พาราณสิยํ อิสิปตเน มิคทาเย อนุตฺตรํ ธมฺมจกฺกํ ปวตฺติตํ อปฺปฏิวตฺติยํ สมเณน วา พฺราหฺมเณน วา เทเวน วา มาเรน วา พฺรหฺมุนา วา เกนจิ วา โลกสฺมินฺ”ติ. ภุมฺมานํ เทวานํ สทฺทํ สุตฺวา จาตุมหาราชิกา\csromansharedfootnote{5402bf593da90e17}{จาตุมฺมหาราชิกา (สฺยา)} เทวา สทฺทมนุสฺสาเวสํ. จาตุมหาราชิกานํ เทวานํ สทฺทํ สุตฺวา ตาวติํสา เทวา ฯเปฯ. ยามา เทวา ฯเปฯ. ตุสิตา เทวา ฯเปฯ. นิมฺมานรตี เทวา ฯเปฯ, ปรนิมฺมิต\-วส\-วตฺตี เทวา ฯเปฯ. พฺรหฺมกายิกา เทวา สทฺทม\-นุสฺสาเวสุํ “เอตํ ภควตา พาราณสิยํ อิสิปตเน มิคทาเย อนุตฺตรํ ธมฺมจกฺกํ ปวตฺติกํ อปฺปฏิวตฺติยํ สมเณน วา พฺราหฺมเณน วา เทเวน วา มาเรน วา พฺรหฺมุนา วา เกนจิ วา โลกสฺมินฺ”ติ.}

\prose{อิติห เตน ขเณน เตน ลเยน เตน มุหุตฺเตน ยาว พฺรหฺมโลกา สทฺโท อพฺภุคฺคจฺฉิ. อยญฺจ ทสสหสฺสี โลกธาตุ สํกมฺปิ สมฺปกมฺปิ สมฺปเวธิ, อปฺปมาโณ จ อุฬาโร โอภาโส โลเก ปาตุรโหสิ อติกฺกมฺม\csromansharedfootnote{f9e8b2d2619eeb34}{อติกฺกมฺเมว (สฺยา, ก) สํ 3. 371 ปิฏฺเฐ ปสฺสิตพฺพา.} เทวานํ เทวานุภาวนฺติ.}

\prose{อถ โข ภควา อิมํ อุทานํ อุทาเนสิ “อญฺญาสิ วต โภ โกณฺฑญฺโญ, อญฺญาสิ วต โภ โกณฺฑญฺโญ”ติ. อิติ หิทํ อายสฺมโต โกณฺฑญฺญสฺส “อญฺญาสิโกณฺฑญฺโญ”เตฺวว\csromansharedfootnote{eaf769ce62be91fe}{อญฺญาโกณฺฑญฺโญเตฺวว (สฺยา)} นามํ อโหสิ.}

\prose{(ก) “อิทํ ทุกฺขํ อริยสจฺจนฺ”ติ ปุพฺเพ อนนุสฺสุเตสุ ธมฺเมสุ จกฺขุํ อุทปาทิ, ญาณํ อุทปาทิ, ปญฺญา อุทปาทิ, วิชฺชา อุทปาทิ, อาโลโก อุทปาทิ.}

\prose{จกฺขุํ อุทปาทีติ เกนฏฺเฐน, ญาณํ อุทปาทีติ เกนฏฺเฐน, ปญฺญา อุทปาทีติ เกนฏฺเฐน, วิชฺชา อุทปาทีติ เกนฏฺเฐน, อาโลโก อุทปาทีติ เกนฏฺเฐน? จกฺขุํ อุทปาทีติ ทสฺสนฏฺเฐน, ญาณํ อุทปาทีติ ญาตฏฺเฐน, ปญฺญา อุทปาทีติ ปชานนฏฺเฐน, วิชฺชา อุทปาทีติ ปฏิเวธ\-ฏฺเฐน, อาโลโก อุทปาทีติ โอภาสฏฺเฐน.}

\prose{\csromanfolio{333}จกฺขุํ ธมฺโม, ญาณํ ธมฺโม, ปญฺญา ธมฺโม, วิชฺชา ธมฺโม, อาโลโก ธมฺโม. อิเม ปญฺจ ธมฺมา ธมฺม\-ปฏิสมฺภิทาย อารมฺมณา เจว โหนฺติ โคจรา จ. เย ตสฺสา อารมฺมณา, เต ตสฺสา โคจรา. เย ตสฺสา โคจรา, เต ตสฺสา อารมฺมณา. เตน วุจฺจติ ธมฺเมสุ ญาณํ ธมฺม\-ปฏิสมฺภิทา.}

\prose{ทสฺสนฏฺโฐ อตฺโถ, ญาตฏฺโฐ อตฺโถ ปชานนฏฺโฐ อตฺโถ, ปฏิเวธฏฺโฐ อตฺโถ โอภาสฏฺโฐ อตฺโถ. อิเม ปญฺจ อตฺถา อตฺถ\-ปฏิสมฺภิทาย อารมฺมณา เจว โหนฺติ โคจรา จ, เย ตสฺสา อารมฺมณา, เต ตสฺสา โคจรา. เย ตสฺสา โคจรา, เต ตสฺสา อารมฺมณา. เตน วุจฺจติ อตฺเถสุ ญาณํ อตฺถ\-ปฏิสมฺภิทา.}

\prose{ปญฺจ ธมฺเม สนฺทสฺเสตุํ พฺยญฺชน\-นิรุตฺตา\-ภิลาปา, ปญฺจ อตฺเถ สนฺทสฺเสตุํ พฺยญฺชน\-นิรุตฺตา\-ภิลาปา. อิมา ทส นิรุตฺติโย นิรุตฺติ\-ปฏิสมฺภิทาย อารมฺมณา เจว โหนฺติ โคจรา จ, เย ตสฺสา อารมฺมณา, เต ตสฺสา โคจรา. เย ตสฺสา โคจรา, เต ตสฺสา อารมฺมณา. เตน วุจฺจติ นิรุตฺตีสุ ญาณํ นิรุตฺติ\-ปฏิสมฺภิทา.}

\prose{ปญฺจสุ ธมฺเมสุ ญาณานิ, ปญฺจสุ อตฺเถสุ ญาณานิ, ทสสุ นิรุตฺตีสุ ญาณานิ. อิมานิ วีสติ ญาณานิ ปฏิภาน\-ปฏิสมฺภิทาย อารมฺมณา เจว โหนฺติ โคจรา จ, เย ตสฺสา อารมฺมณา, เต ตสฺสา โคจรา. เย ตสฺสา โคจรา, เต ตสฺสา อารมฺมณา. เตน วุจฺจติ ปฏิภาเนสุ ญาณํ ปฏิภาน\-ปฏิสมฺภิทา.}

\prose{“ตํ โข ปนิทํ ทุกฺขํ อริยสจฺจํ ปริญฺเญยฺยนฺ”ติ ฯเปฯ “ปริญฺญาตนฺ”ติ ปุพฺเพ อนนุสฺสุเตสุ ธมฺเมสุ จกฺขุํ อุทปาทิ, ญาณํ อุทปาทิ, ปญฺญา อุทปาทิ, วิชฺชา อุทปาทิ, อาโลโก อุทปาทิ.}

\prose{จกฺขุํ อุทปาทีติ เกนฏฺเฐน, ญาณํ อุทปาทีติ เกนฏฺเฐน, ปญฺญา อุทปาทีติ เกนฏฺเฐน, วิชฺชา อุทปาทีติ เกนฏฺเฐน, อาโลโก อุทปาทีติ เกนฏฺเฐน? จกฺขุํ อุทปาทีติ ทสฺสนฏฺเฐน, ญาณํ อุทปาทีติ ญาตฏฺเฐน, ปญฺญา อุทปาทีติ ปชานนฏฺเฐน, วิชฺชา อุทปาทีติ ปฏิเวธ\-ฏฺเฐน, อาโลโก อุทปาทีติ โอภาสฏฺเฐน.}

\prose{จกฺขุํ ธมฺโม, ญาณํ ธมฺโม, ปญฺญา ธมฺโม, วิชฺชา ธมฺโม, อาโลโก ธมฺโม. อิเม ปญฺจ ธมฺมา ธมฺม\-ปฏิสมฺภิทาย อารมฺมณา เจว โหนฺติ โคจรา}

\prose{\csromanfolio{334}จ, เย ตสฺสา อารมฺมณา, เต ตสฺสา โคจรา. เย ตสฺสา โคจรา, เต ตสฺสา อารมฺมณา. เตน วุจฺจติ ธมฺเมสุ ญาณํ ธมฺม\-ปฏิสมฺภิทา.}

\prose{ทสฺสนฏฺโฐ อตฺโถ, ญาตฏฺโฐ อตฺโถ, ปชานนฏฺโฐ อตฺโถ, ปฏิเวธฏฺโฐ อตฺโถ, โอภาสฏฺโฐ อตฺโถ. อิเม ปญฺจ อตฺถา อตฺถ\-ปฏิสมฺภิทาย อารมฺมณา เจว โหนฺติ โคจรา จ, เย ตสฺสา อารมฺมณา, เต ตสฺสา โคจรา. เย ตสฺสา โคจรา, เต ตสฺสา อารมฺมณา. เตน วุจฺจติ อตฺเถสุ ญาณํ อตฺถ\-ปฏิสมฺภิทา.}

\prose{ปญฺจ ธมฺเม สนฺทสฺเสตุํ พฺยญฺชน\-นิรุตฺตา\-ภิลาปา, ปญฺจ อตฺเถ สนฺทสฺเสตุํ พฺยญฺชน\-นิรุตฺตา\-ภิลาปา, อิมา ทส นิรุตฺติโย นิรุตฺติ\-ปฏิสมฺภิทาย อารมฺมณา เจว โหนฺติ โคจรา จ, เย ตสฺสา อารมฺมณา, เต ตสฺสา โคจรา. เย ตสฺสา โคจรา, เต ตสฺสา อารมฺมณา. เตน วุจฺจติ นิรุตฺตีสุ ญาณํ นิรุตฺติ\-ปฏิสมฺภิทา.}

\prose{ปญฺจสุ ธมฺเมสุ ญาณานิ, ปญฺจสุ อตฺเถสุ ญาณานิ, ทสสุ นิรุตฺตีสุ ญาณานิ. อิมานิ วีสติ ญาณานิ ปฏิภาน\-ปฏิสมฺภิทาย อารมฺมณา เจว โหนฺติ โคจรา จ, เย ตสฺสา อารมฺมณา, เต ตสฺสา โคจรา. เย ตสฺสา โคจรา, เต ตสฺสา อารมฺมณา. เตน วุจฺจติ ปฏิภาเนสุ ญาณํ ปฏิภาน\-ปฏิสมฺภิทา.}

\prose{ทุกฺเข อริยสจฺเจ ปนฺนรส ธมฺมา, ปนฺนรส อตฺถา, ติํส นิรุตฺติโย, สฏฺฐิ ญาณานิ.}

\prose{(ข) “อิทํ ทุกฺข\-สมุทยํ อริยสจฺจนฺ”ติ ปุพฺเพ อนนุสฺสุเตสุ ธมฺเมสุ จกฺขุํ อุทปาทิ ฯเปฯ อาโลโก อุทปาทิ. “ตํ โข ปนิทํ ทุกฺข\-สมุทยํ อริยสจฺจํ ปหาตพฺพนฺ”ติ ฯเปฯ “ปหีนนฺ”ติ ปุพฺเพ อนนุสฺสุเตสุ ธมฺเมสุ จกฺขุํ อุทปาทิ ฯเปฯ อาโลโก อุทปาทิ ฯเปฯ.}

\prose{ทุกฺขสมุทเย อริยสจฺเจ ปนฺนรส ธมฺมา, ปนฺนรส อตฺถา, ติํส นิรุตฺติโย, สฏฺฐิ ญาณานิ.}

\prose{(ค) “อิทํ ทุกฺขนิโรธํ อริยสจฺจนฺ”ติ ปุพฺเพ อนนุสฺสุเตสุ ธมฺเมสุ จกฺขุํ อุทปาทิ ฯเปฯ อาโลโก อุทปาทิ. “ตํ โข ปนิทํ ทุกฺขนิโรธํ อริยสจฺจํ สจฺฉิกาตพฺพนฺ”ติ ฯเปฯ “สจฺฉิกตนฺ”ติ ปุพฺเพ อนนุสฺสุเตสุ ธมฺเมสุ จกฺขุํ อุทปาทิ ฯเปฯ อาโลโก อุทาปาทิ ฯเปฯ.}

\prose{\csromanfolio{335}ทุกฺขนิโรเธ อริยสจฺเจ ปนฺนรส ธมฺมา, ปนฺนรส อตฺถา, ติํส นิรุตฺติโย, สฏฺฐิ ญาณานิ.}

\prose{(ฆ) “อิทํ ทุกฺขนิโรธ\-คามินี ปฏิปทา อริยสจฺจนฺ”ติ ปุพฺเพ อนนุสฺสุเตสุ ธมฺเมสุ จกฺขุํ อุทปาทิ ฯเปฯ อาโลโก อุทปาทิ. “ตํ โข ปนิทํ ทุกฺขนิโรธ\-คามินี ปฏิปทา อริยสจฺจํ ภาเวตพฺพนฺ”ติ ฯเปฯ “ภาวิตนฺ”ติ ปุพฺเพ อนนุสฺสุเตสุ ธมฺเมสุ จกฺขุํ อุทปาทิ ฯเปฯ อาโลโก อุทปาทิ ฯเปฯ.}

\prose{ทุกฺขนิโรธ\-คามินิยา ปฏิปทาย อริยสจฺเจ ปนฺนรส ธมฺมา, ปนฺนรส อตฺถา, ติํส นิรุตฺติโย, สฏฺฐิ ญาณานิ.}

\prose{จตูสุ อริยสจฺเจสุ สฏฺฐิ ธมฺมา, สฏฺฐิ อตฺถา, วีสติสต\-นิรุตฺติโย, จตฺตาลีสญฺจ เทฺว จ ญาณสตานิ.}

\csromanmatikaanchor{mtk.103}
\csromantocmarkref{subsection}{2. สติปฏฺฐานวาร}{mtk.103}
\bookmark[dest={mtk.103},level=2]{2. สติปฏฺฐานวาร}
\csromanheader{2}{2. สติปฏฺฐาน\-วาร}
\proseitem{31}{“อยํ กาเย กายานุปสฺสนา”ติ เม ภิกฺขเว ปุพฺเพ อนนุสฺสุเตสุ ธมฺเมสุ จกฺขุํ อุทปาทิ ฯเปฯ อาโลโก อุทปาทิ, “สา โข ปนายํ กาเย กายานุปสฺสนา ภาเวตพฺพา”ติ เม ภิกฺขเว ฯเปฯ “ภาวิตา”ติ เม ภิกฺขเว ปุพฺเพ อนนุสฺสุเตสุ ธมฺเมสุ จกฺขุํ อุทปาทิ ฯเปฯ อาโลโก อุทปาทิ.}

\prose{“อยํ เวทนาสุ ฯเปฯ. อยํ จิตฺเต ฯเปฯ. อยํ ธมฺเมสุ ธมฺมา\-นุปสฺสนา”ติ เม ภิกฺขเว ปุพฺเพ อนนุสฺสุเตสุ ธมฺเมสุ จกฺขุํ อุทปาทิ ฯเปฯ อาโลโก อุทปาทิ, “สา โข ปนายํ ธมฺเมสุ ธมฺมา\-นุปสฺสนา ภาเวตพฺพา”ติ ฯเปฯ “ภาวิตา”ติ เม ภิกฺขเว ปุพฺเพ อนนุสฺสุเตสุ ธมฺเมสุ จกฺขุํ อุทปาทิ ฯเปฯ อาโลโก อุทปาทิ.}

\prose{(ก) “อยํ กาเย กายานุปสฺสนา”ติ ปุพฺเพ อนนุสฺสุเตสุ ธมฺเมสุ จกฺขุํ อุทปาทิ ฯเปฯ อาโลโก อุทปาทิ ฯเปฯ “สา โข ปนายํ กาเย กายานุปสฺสนา ภาเวตพฺพา”ติ ฯเปฯ “ภาวิตา”ติ ปุพฺเพ อนนุสฺสุเตสุ ธมฺเมสุ จกฺขุํ อุทปาทิ, ญาณํ อุทปาทิ, ปญฺญา อุทปาทิ, วิชฺชา อุทปาทิ, อาโลโก อุทปาทิ.}

\prose{จกฺขุํ อุทปาทีติ เกนฏฺเฐน, ญาณํ อุทปาทีติ เกนฏฺเฐน, ปญฺญา อุทปาทีติ เกนฏฺเฐน, วิชฺชา อุทปาทีติ เกนฏฺเฐน, อาโลโก อุทปาทีติ เกนฏฺเฐน? \textbf{จกฺขุํ อุทปาที}ติ ทสฺสนฏฺฐน, \textbf{ญาณํ อุทปาที}ติ ญาตฏฺเฐน, \csromanfolio{336}ปญฺญา อุทปาทีติ ปชานนฏฺเฐน, วิชฺชา อุทปาทีติ ปฏิเวธ\-ฏฺเฐน, อาโลโก อุทปาทีติ โอภาสฏฺเฐน.}

\prose{จกฺขุํ ธมฺโม, ญาณํ ธมฺโม, ปญฺญา ธมฺโม, วิชฺชา ธมฺโม, อาโลโก ธมฺโม. อิเม ปญฺจ ธมฺมา ธมฺม\-ปฏิสมฺภิทาย อารมฺมณา เจว โหนฺติ โคจรา จ. เย ตสฺสา อารมฺมณา, เต ตสฺสา โคจรา, เย ตสฺสา โคจรา, เต ตสฺสา อารมฺมณา, เตน วุจฺจติ ธมฺเมสุ ญาณํ ธมฺม\-ปฏิสมฺภิทา.}

\prose{ทสฺสนฏฺโฐ อตฺโถ, ญาตฏฺโฐ อตฺโถ, ปชานนฏฺโฐ อตฺโถ, ปฏิเวธฏฺโฐ อตฺโถ, โอภาสฏฺโฐ อตฺโถ. อิเม ปญฺจ อตฺถา อตฺถ\-ปฏิสมฺภิทาย อารมฺมณา เจว โหนฺติ โคจรา จ, เย ตสฺสา อารมฺมณา, เต ตสฺสา โคจรา, เย ตสฺสา โคจรา, เต ตสฺสา อารมฺมณา, เตน วุจฺจติ อตฺเถสุ ญาณํ อตฺถ\-ปฏิสมฺภิทา.}

\prose{ปญฺจ ธมฺเม สนฺทสฺเสตุํ พฺยญฺชน\-นิรุตฺตา\-ภิลาปา, ปญฺจ อตฺเถ สนฺทสฺเสตุํ พฺยญฺชน\-นิรุตฺตา\-ภิลาปา, อิมา ทส นิรุตฺติโย นิรุตฺติ\-ปฏิสมฺภิทาย อารมฺมณา เจว โหนฺติ โคจรา จ, เย ตสฺสา อารมฺมณา, เต ตสฺสา โคจรา. เย ตสฺสา โคจรา, เต ตสฺสา อารมฺมณา. เตน วุจฺจติ นิรุตฺตีสุ ญาณํ นิรุตฺติ\-ปฏิสมฺภิทา.}

\prose{ปญฺจสุ ธมฺเมสุ ญาณานิ, ปญฺจสุ อตฺเถสุ ญาณานิ, ทสสุ นิรุตฺตีสุ ญาณานิ, อิมานิ วีสติ ญาณานิ ปฏิภาน\-ปฏิสมฺภิทาย อารมฺมณา เจว โหนฺติ โคจรา จ, เย ตสฺสา อารมฺมณา, เต ตสฺสา โคจรา. เย ตสฺสา โคจรา, เต ตสฺสา อารมฺมณา. เตน วุจฺจติ ปฏิภาเนสุ ญาณํ ปฏิภาน\-ปฏิสมฺภิทา.}

\prose{กาเย กายานุปสฺสนา\-สติปฏฺฐาเน ปนฺนรส ธมฺมา, ปนฺนรส อตฺถา, ติํส นิรุตฺติโย, สฏฺฐิ ญาณานิ.}

\prose{(ข. ฆ) “อยํ เวทนาสุ ฯเปฯ. อยํ จิตฺเต ฯเปฯ. อยํ ธมฺเมสุ ธมฺมา\-นุปสฺสนา”ติ ปุพฺเพ อนนุสฺสุเตสุ ธมฺเมสุ จกฺขุํ อุทปาทิ ฯเปฯ อาโลโก อุทปาทิ ฯเปฯ “สา โข ปนายํ ธมฺเมสุ ธมฺมา\-นุปสฺสนา ภาเวตพฺพา”ติ ฯเปฯ “ภาวิตา”ติ ปุพฺเพ อนนุสฺสุเตสุ ธมฺเมสุ จกฺขุํ อุทปาทิ ฯเปฯ อาโลโก อุทปาทิ ฯเปฯ.}

\prose{\csromanfolio{337}ธมฺเมสุ ธมฺมา\-นุปสฺสนา สติปฏฺฐาเน ปนฺนรส ธมฺมา, ปนฺนรส อตฺถา, ติํส นิรุตฺติโย, สฏฺฐิ ญาณานิ.}

\prose{จตูสุ สติปฏฺฐาเนสุ สฏฺฐิ ธมฺมา, สฏฺฐิ อตฺถา, วีสติสต\-นิรุตฺติโย, จตฺตาลีสญฺจ เทฺว จ ญาณสตานิ.}

\csromanmatikaanchor{mtk.104}
\csromantocmarkref{subsection}{3. อิทฺธิปาทวาร}{mtk.104}
\bookmark[dest={mtk.104},level=2]{3. อิทฺธิปาทวาร}
\csromanheader{2}{3. \textbf{อิทฺธิปาทวาร}}
\proseitem{32}{“อยํ ฉนฺท\-สมาธิ\-ปธาน\-สงฺขาร\-สมนฺนาคโต อิทฺธิปาโท”ติ เม ภิกฺขเว ปุพฺเพ อนนุสฺสุเตสุ ธมฺเมสุ จกฺขุํ อุทปาทิ ฯเปฯ อาโลโก อุทปาทิ. “โส โข ปนายํ ฉนฺท\-สมาธิ\-ปธาน\-สงฺขาร\-สมนฺนาคโต อิทฺธิปาโท ภาเวตพฺโพ”ติ เม ภิกฺขเว ฯเปฯ “ภาวิโต”ติ เม ภิกฺขเว ปุพฺเพ อนนุสฺสุเตสุ ธมฺเมสุ จกฺขุํ อุทปาทิ ฯเปฯ อาโลโก อุทปาทิ.}

\prose{“อยํ วีริยสมาธิ ฯเปฯ. อยํ จิตฺตสมาธิ ฯเปฯ. อยํ วีมํสา\-สมาธิ\-ปธาน\-สงฺขาร\-สมนฺนาคโต อิทฺธิปาโท”ติ เม ภิกฺขเว ปุพฺเพ อนนุสฺสุเตสุ ธมฺเมสุ จกฺขุํ อุทปาทิ ฯเปฯ อาโลโก อุทปาทิ. “โส โข ปนายํ วีมํสา\-สมาธิ\-ปธาน\-สงฺขาร\-สมนฺนาคโต อิทฺธิปาโท ภาเวตพฺโพ”ติ เม ภิกฺขเว ปุพฺเพ อนนุสฺสุเตสุ ธมฺเมสุ จกฺขุํ อุทปาทิ ฯเปฯ อาโลโก อุทปาทิ.}

\prose{(ก) “อยํ ฉนฺท\-สมาธิ\-ปธาน\-สงฺขาร\-สมนฺนาคโต อิทฺธิปาโท”ติ ปุพฺเพ อนนุสฺสุเตสุ ธมฺเมสุ จกฺขุํ อุทปาทิ ฯเปฯ อาโลโก อุทปาทิ ฯเปฯ “โส โข ปนายํ ฉนฺท\-สมาธิ\-ปธาน\-สงฺขาร\-สมนฺนาคโต อิทฺธิปาโท ภาเวตพฺโพ”ติ ฯเปฯ “ภาวิโต”ติ ปุพฺเพ อนนุสฺสุเตสุ ธมฺเมสุ จกฺขุํ อุทปาทิ, ญาณํ อุทปาปาทิ, ปญฺญา อุทปาทิ, วิชฺชา อุทปาทิ, อาโลโก อุทปาทิ.}

\prose{จกฺขุํ อุทปาทีติ เกนฏฺเฐน, ญาณํ อุทปาทีติ เกนฏฺเฐน, ปญฺญา อุทปาทีติ เกนฏฺเฐน, วิชฺชา อุทปาทีติ เกนฏฺเฐน, อาโลโก อุทปาทีติ เกนฏฺเฐน? จกฺขุํ อุทปาทีติ ทสฺสนฏฺเฐน. ญาณํ อุทปาทีติ ญาตฏฺเฐน. ปญฺญา อุทปาทีติ ปชานนฏฺเฐน. วิชฺชา อุทปาทีติ ปฏิเวธ\-ฏฺเฐน. อาโลโก อุทปาทีติ โอภาสฏฺเฐน.}

\prose{จกฺขุํ ธมฺโม, ญาณํ ธมฺโม, ปญฺญา ธมฺโม, วิชฺชา ธมฺมา, อาโลโก ธมฺโม. อิเม ปญฺจ ธมฺมา ธมฺม\-ปฏิสมฺภิทาย อารมฺมณา เจว โหนฺติ โคจรา จ. \csromanfolio{338}เย ตสฺสา อารมฺมณา, เต ตสฺสา โคจรา. ย ตสฺสา โคจรา, เต ตสฺสา อารมฺมณา. เตน วุจฺจติ ธมฺเมสุ ญาณํ ธมฺม\-ปฏิสมฺภิทา.}

\prose{ทสฺสนฏฺโฐ อตฺโถ, ญาตฏฺโฐ อตฺโถ, ปชานนฏฺโฐ อตฺโถ, ปฏิเวธฏฺโฐ อตฺโถ, โอภาสฏฺโฐ อตฺโถ. อิเม ปญฺจ อตฺถา อตฺถ\-ปฏิสมฺภิทาย อารมฺมณา เจว โหนฺติ โคจรา จ. เย ตสฺสา อารมฺมณา, เต ตสฺสา โคจรา. เย ตสฺสา โคจรา, เต ตสฺสา อารมฺมณา. เตน วุจฺจติ อตฺเถสุ ญาณํ อตฺถ\-ปฏิสมฺภิทา.}

\prose{ปญฺจ ธมฺเม สนฺทสฺเสตุํ พฺยญฺชน\-นิรุตฺตา\-ภิลาปา, ปญฺจ อตฺเถ สนฺทสฺเสตุํ พฺยญฺชน\-นิรุตฺตา\-ภิลาปา. อิมา ทส นิรุตฺติโย นิรุตฺติ\-ปฏิสมฺภิทาย อารมฺมณา เจว โหนฺติ โคจรา จ. เย ตสฺสา อารมฺมณา, เต ตสฺสา โคจรา. เย ตสฺสา โคจรา, เต ตสฺสา อารมฺมณา. เตน วุจฺจติ นิรุตฺตีสุ ญาณํ นิรุตฺติ\-ปฏิสมฺภิทา.}

\prose{ปญฺจสุ ธมฺเมสุ ญาณานิ, ปญฺจสุ อตฺเถสุ ญาณานิ, ทสสุ นิรุตฺตีสุ ญาณานิ. อิมานิ วีสติ ญาณานิ ปฏิภาน\-ปฏิสมฺภิทาย อารมฺมณา เจว โหนฺติ โคจรา จ. เย ตสฺสา อารมฺมณา, เต ตสฺสา โคจรา. เย ตสฺสา โคจรา, เต ตสฺสา อารมฺมณา. เตน วุจฺจติ ปฏิภาเนสุ ญาณํ ปฏิภาน\-ปฏิสมฺภิทา.}

\prose{ฉนฺท\-สมาธิ\-ปธาน\-สงฺขาร\-สมนฺนาคเต อิทฺธิปาเท ปนฺนรส ธมฺมา, ปนฺนรส อตฺถา, ติํส นิรุตฺติโย, สฏฺฐิ ญาณานิ.}

\prose{(ข. ฆ) “อยํ วีริยสมาธิ ฯเปฯ. อยํ จิตฺตสมาธิ ฯเปฯ. อยํ วีมํสา\-สมาธิ\-ปธาน\-สงฺขาร\-สมนฺนาคโต อิทฺธิปาโท”ติ ปุพฺเพ อนนุสฺสุเตสุ ธมฺเมสุ จกฺขุํ อุทปาทิ ฯเปฯ อาโลโก อุทปาทิ ฯเปฯ “โส โข ปนายํ วีมํสา\-สมาธิ\-ปธาน\-สงฺขาร\-สมนฺนาคโต อิทฺธิปาโท ภาเวตพฺโพ”ติ ฯเปฯ “ภาวิโต”ติ ปุพฺเพ อนนุสฺสุเตสุ ธมฺเมสุ จกฺขุํ อุทปาทิ ฯเปฯ อาโลโก อุทปาทิ ฯเปฯ.}

\prose{วีมํสา\-สมาธิ\-ปธาน\-สงฺขาร\-สมนฺนาคเต อิทฺธิปาเท ปนฺนรส ธมฺมา, ปนฺนรส อตฺถา, ติํส นิรุตฺติโย, สฏฺฐิ ญาณานิ.}

\prose{จตูสุ อิทฺธิปาเทสุ สฏฺฐิ ธมฺมา, สฏฺฐิ อตฺถา, วีสติสต\-นิรุตฺติโย, จตฺตาลีสญฺจ เทฺว จ ญาณสตานิ.}

\csromanmatikaanchor{mtk.105}
\csromantocmarkref{subsection}{4. สตฺตโพธิสตฺตวาร}{mtk.105}
\bookmark[dest={mtk.105},level=2]{4. สตฺตโพธิสตฺตวาร}
\csromanheader{2}{4. สตฺต\-โพธิสตฺต\-วาร}
\proseitem{33}{\csromanfolio{339}“สมุทโย สมุทโย”ติ โข ภิกฺขเว วิปสฺสิสฺส โพธิสตฺตสฺส ปุพฺเพ อนนุสฺสุเตสุ ธมฺเมสุ จกฺขุํ อุทปาทิ ฯเปฯ อาโลโก อุทปาทิ. “นิโรโธ นิโรโธ”ติ โข ภิกฺขเว วิปสฺสิสฺส โพธิสตฺตสฺส ปุพฺเพ อนนุสฺสุเตสุ ธมฺเมสุ จกฺขุํ อุทปาทิ ฯเปฯ อาโลโก อุทปาทิ. วิปสฺสิสฺส โพธิสตฺตสฺส เวยฺยากรเณ ทส ธมฺมา, ทส อตฺถา, วีสติ นิรุตฺติโย, จตฺตาลีสํ\csromansharedfootnote{0b6b40f6629cb342}{จตฺตารีส (ก)} ญาณานิ.}

\prose{“สมุทโย สมุทโย”ติ โข ภิกฺขเว สิขิสฺส โพธิสตฺตสฺส ฯเปฯ เวสฺสภุสฺส โพธิสตฺตสฺส ฯเปฯ กกุสนฺธสฺส โพธิสตฺตสฺส ฯเปฯ โกณาคมนสฺส โพธิสตฺตสฺส ฯเปฯ กสฺสปสฺส โพธิสตฺตสฺส ปุพฺเพ อนนุสฺสุเตสุ ธมฺเมสุ จกฺขุํ อุทปาทิ ฯเปฯ อาโลโก อุทปาทิ. “นิโรโธ นิโรโธ”ติ โข ภิกฺขเว กสฺสปสฺส โพธิสตฺตสฺส ปุพฺเพ อนนุสฺสุเตสุ ธมฺเมสุ จกฺขุํ อุทปาทิ ฯเปฯ อาโลโก อุทปาทิ. กสฺสปสฺส โพธิสตฺตสฺส เวยฺยากรเณ ทส ธมฺมา, ทส อตฺถา, วีสติ นิรุตฺติโย จตฺตาลีสํ ญาณานิ.}

\prose{“สมุทโย สมุทโย”ติ โข ภิกฺขเว โคตมสฺส โพธิสตฺตสฺส ปุพฺเพ อนนุสฺสุเตสุ ธมฺเมสุ จกฺขุํ อุทปาทิ ฯเปฯ อาโลโก อุทปาทิ. “นิโรโธ นิโรโธ”ติ โข ภิกฺขเว โคตมสฺส โพธิสตฺตสฺส ปุพฺเพ อนนุสฺสุเตสุ ธมฺเมสุ จกฺขุํ อุทปาทิ ฯเปฯ อาโลโก อุทปาทิ. โคตมสฺส โพธิสตฺตสฺส เวยฺยากรเณ ทส ธมฺมา, ทส อตฺถา, วีสติ นิรุตฺติโย, จตฺตาลีสํ ญาณานิ.}

\prose{สตฺตนฺนํ โพธิสตฺตานํ สตฺตสุ เวยฺยากรเณสุ สตฺตติ ธมฺมา, สตฺตติ อตฺถา, จตฺตาลีสสตํ\csromansharedfootnote{cb18775890e60068}{จตฺตารีสสตา (สฺยา)} นิรุตฺติโย, อสีติ จ เทฺว จ ญาณสตานิ.}

\csromanmatikaanchor{mtk.106}
\csromantocmarkref{subsection}{5. อภิญฺญาทิวาร}{mtk.106}
\bookmark[dest={mtk.106},level=2]{5. อภิญฺญาทิวาร}
\csromanheader{2}{5. \textbf{อภิญฺญาทิวาร}}
\proseitem{34}{ยาวตา อภิญฺญาย อภิญฺญฏฺโฐ ญาโต ทิฏฺโฐ วิทิโต สจฺฉิกโต ผสฺสิโต ปญฺญาย. อผสฺสิโต ปญฺญาย อภิญฺญฏฺโฐ นตฺถีติ จกฺขุํ อุทปาทิ, ญาณํ อุทปาทิ, ปญฺญา อุทปาทิ, วิชฺชา อุทปาทิ, อาโลโก อุทปาทิ. อภิญฺญาย อภิญฺญฏฺเฐ ปญฺจวีสติ ธมฺมา, ปญฺจวีสติ อตฺถา, ปญฺญาสํ นิรุตฺติโย, สตํ ญาณานิ.}

\prose{\csromanfolio{340}ยาวตา ปริญฺญาย ปริญฺญฏฺโฐ ฯเปฯ. ยาวตา ปหานสฺส ปหานฏฺโฐ ฯเปฯ. ยาวตา ภาวนาย ภาวนฏฺโฐ ฯเปฯ. ยาวตา สจฺฉิกิริยาย สจฺฉิกิริยฏฺโฐ ญาโต ทิฏฺโฐ วิทิโต สจฺฉิกโต ผสฺสิโต ปญฺญาย. อผสฺสิโต ปญฺญาย สจฺฉิกิริยฏฺโฐ นตฺถีติ จกฺขุํ อุทปาทิ, ญาณํ อุทปาทิ, ปญฺญา อุทปาทิ, วิชฺชา อุทปาทิ, อาโลโก อุทปาทิ, สจฺฉิกิริยาย สจฺฉิกิริยฏฺเฐ ปญฺจวีสติ ธมฺมา, ปญฺจวีสติ อตฺถา, ปญฺญาสํ นิรุตฺติโย, สตํ ญาณานิ.}

\prose{อภิญฺญาย อภิญฺญฏฺเฐ ปริญฺญาย ปริญฺญฏฺเฐ ปหานาย ปหานฏฺเฐ. ภาวนาย ภาวนฏฺเฐ. สจฺฉิกิริยาย สจฺฉิกิริยฏฺเฐ ปญฺจวีสสตํ ธมฺมา ปญฺจวีสสตํ อตฺถา, อฑฺฒเตยฺยานิ นิรุตฺติสตานิ, ปญฺจ ญาณสตานิ.}

\csromanmatikaanchor{mtk.107}
\csromantocmarkref{subsection}{6. ขนฺธาทิวาร}{mtk.107}
\bookmark[dest={mtk.107},level=2]{6. ขนฺธาทิวาร}
\csromanheader{2}{6. \textbf{ขนฺธาทิวาร}}
\proseitem{35}{ยาวตา ขนฺธานํ ขนฺธฏฺโฐ ญาโต ทิฏฺโฐ วิทิโต สจฺฉิกโต ผสฺสิโต ปญฺญาย. อผสฺสิโต ปญฺญาย ขนฺธฏฺโฐ นตฺถีติ จกฺขุํ อุทปาทิ. ญาณํ อุทปาทิ, ปญฺญา อุทปาทิ, วิชฺชา อุทปาทิ, อาโลโก อุทปาทิ, ขนฺธานํ ขนฺธฏฺเฐ ปญฺจวีสติ ธมฺมา, ปญฺจวีสติ อตฺถา, ปญฺญาสํ นิรุตฺติโย, สตํ ญาณานิ.}

\prose{ยาวตา ธาตูนํ ธาตุฏฺโฐ ฯเปฯ. ยาวตา อายตนานํ อายตนฏฺโฐ ฯเปฯ. ยาวตา สงฺขตานํ สงฺขตฏฺโฐ ฯเปฯ. ยาวตา อสงฺขตสฺส อสงฺขตฏฺโฐ ญาโต ทิฏฺโฐ วิทิโต สจฺฉิกโต ผสฺสิโต ปญฺญาย. อผสฺสิโต ปญฺญาย อสงฺขตฏฺโฐ นตฺถีติ จกฺขุํ อุทปาทิ ฯเปฯ อาโลโก อุทปาทิ. อสงฺขตสฺส อสงฺขตฏฺเฐ ปญฺจวีสติ ธมฺมา, ปญฺจวีสติ อตฺถา, ปญฺญาสํ นิรุตฺติโย, สตํ ญาณานิ.}

\prose{ขนฺธานํ ขนฺธฏฺเฐ ธาตูนํ ธาตุฏฺเฐ อายตนานํ อายตนฏฺเฐ สงฺขตานํ สงฺขตฏฺเฐ อสงฺขตสฺส อสงฺขตฏฺเฐ ปญฺจวีส\-ติสตํ ธมฺมา, ปญฺจวีส\-ติสตํ อตฺถา, อฑฺฒเตยฺยานิ นิรุตฺติสตานิ, ปญฺจ ญาณสตานิ.}

\csromanmatikaanchor{mtk.108}
\csromantocmarkref{subsection}{7. สจฺจวาร}{mtk.108}
\bookmark[dest={mtk.108},level=2]{7. สจฺจวาร}
\csromanheader{2}{7. \textbf{สจฺจวาร}}
\proseitem{36}{ยาวตา ทุกฺขสฺส ทุกฺขฏฺโฐ ญาโต ทิฏฺโฐ วิทิโต สจฺฉิกโต ผสฺสิโต ปญฺญาย. อผสฺสิโต ปญฺญาย ทุกฺขฏฺโฐ นตฺถีติ จกฺขุํ อุทปาทิ, ญาณํ อุทปาทิ, ปญฺญา อุทปาทิ, วิชฺชา อุทปาทิ, อาโลโก อุทปาทิ. ทุกฺขสฺส}

\prose{\csromanfolio{341}ทุกฺขฏฺเฐ ปญฺจวีสติ ธมฺมา, ปญฺจวีสติ อตฺถา, ปญฺญาสํ นิรุตฺติโย, สตํ ญาณานิ.}

\prose{ยาวตา สมุทยสฺส สมุทยฏฺโฐ ฯเปฯ. ยาวตา นิโรธสฺส นิโรธฏฺโฐ ฯเปฯ. ยาวตา มคฺคสฺส มคฺคฏฺโฐ ญาโต ทิฏฺโฐ วิทิโต สจฺฉิกโต ผสฺสิโต ปญฺญาย. อผสฺสิโต ปญฺญาย มคฺคฏฺโฐ นตฺถีติ จกฺขุํ อุทปาทิ ฯเปฯ อาโลโก อุทปาทิ. มคฺคสฺส มคฺคฏฺเฐ ปญฺจวีสติ ธมฺมา, ปญฺจวีสติ อตฺถา. ปญฺญาสํ นิรุตฺติโย, สตํ ญาณานิ.}

\prose{จตูสุ อริยสจฺเจสุ สตํ ธมฺมา, สตํ อตฺถา, เทฺว นิรุตฺติสตานิ, จตฺตาริ ญาณสตานิ.}

\csromanmatikaanchor{mtk.109}
\csromantocmarkref{subsection}{8. ปฏิสมฺภิทาวาร}{mtk.109}
\bookmark[dest={mtk.109},level=2]{8. ปฏิสมฺภิทาวาร}
\csromanheader{2}{8. ปฏิสมฺภิทา\-วาร}
\proseitem{37}{ยาวตา อตฺถ\-ปฏิสมฺภิทาย อตฺถปฏิสมฺภิท\-ฏฺโฐ ญาโต ทิฏฺโฐ วิทิโต สจฺฉิกโต ผสฺสิโต ปญฺญาย. อผสฺสิโต ปญฺญาย อตฺถปฏิสมฺภิท\-ฏฺโฐ นตฺถีติ จกฺขุํ อุทปาทิ ฯเปฯ อาโลโก อุทปาทิ. อตฺถ\-ปฏิสมฺภิทาย อตฺถปฏิสมฺภิท\-ฏฺเฐ ปญฺจวีสติ ธมฺมา, ปญฺจวีสติ อตฺถา. ปญฺญาสํ นิรุตฺติโย, สตํ ญาณานิ.}

\prose{ยาวตา ธมฺม\-ปฏิสมฺภิทาย ธมฺมปฏิสมฺภิท\-ฏฺโฐ ฯเปฯ. ยาวตา นิรุตฺติ\-ปฏิสมฺภิทาย นิรุตฺติปฏิสมฺภิท\-ฏฺโฐ ฯเปฯ. ยาวตา ปฏิภาน\-ปฏิสมฺภิทาย ปฏิภานปฏิสมฺภิท\-ฏฺโฐ ญาโต ทิฏฺโฐ วิทิโต สจฺฉิกโต ผสฺสิโต ปญฺญาย. อผสฺสิโต ปญฺญาย ปฏิภานปฏิสมฺภิท\-ฏฺโฐ นตฺถีติ จกฺขุํ อุทปาทิ ฯเปฯ อาโลโก อุทปาทิ. ปฏิภานปฏิสมฺภิท\-ฏฺเฐ ปญฺจวีสติ ธมฺมา ปญฺจวีสติ อตฺถา, ปญฺญาสํ นิรุตฺติโย, สตํ ญาณานิ.}

\prose{จตูสุ ปฏิสมฺภิทาสุ สตํ ธมฺมา, สตํ อตฺถา, เทฺว นิรุตฺติสตานิ, จตฺตาริ ญาณสตานิ.}

\csromanmatikaanchor{mtk.110}
\csromantocmarkref{subsection}{9. ฉพุทฺธธมฺมวาร}{mtk.110}
\bookmark[dest={mtk.110},level=2]{9. ฉพุทฺธธมฺมวาร}
\csromanheader{2}{9. ฉพุทฺธธมฺม\-วาร}
\proseitem{38}{ยาวตา อินฺทฺริยปโรปริยตฺต\-ญาณํ ญาตํ ทิฏฺฐํ วิทิตํ สจฺฉิกตํ ผสฺสิตํ ปญฺญาย, อผสฺสิตํ ปญฺญาย อินฺทฺริยปโรปริยตฺต\-ญาณํ นตฺถีติ จกฺขุํ อุทปาทิ ฯเปฯ อาโลโก อุทปาทิ, อินฺทฺริยปโรปริยตฺต\-ญาเณ ปญฺจวีสติ ธมฺมา, ปญฺจวีสติ อตฺถา, ปญฺญาสํ นิรุตฺติโย, สตํ ญาณานิ.}

\prose{\csromanfolio{342}ยาวตา สตฺตานํ อาสยานุสเย ญาณํ ฯเปฯ. ยาวตา ยมกปาฏิหีเร ญาณํ ฯเปฯ. ยาวตา มหากรุณา\-สมาปตฺติยา ญาณํ ฯเปฯ. ยาวตา สพฺพญฺญุต\-ญฺญาณํ ฯเปฯ. ยาวตา อนาวรณํ ญาณํ ญาตํ ทิฏฺฐํ วิทิตํ สจฺฉิกตํ ผสฺสิตํ ปญฺญาย, อผสฺสิตํ ปญฺญาย อนาวรณํ ญาณํ นตฺถีติ จกฺขุํ อุทปาทิ, ญาณํ อุทปาทิ, ปญฺญา อุทปาทิ, วิชฺชา อุทปาทิ, อาโลโก อุทปาทิ. อนาวรเณ ญาเณ ปญฺจวีสติ ธมฺมา, ปญฺจวีสติ อตฺถา, ปญฺญาสํ นิรุตฺติโย, สตํ ญาณานิ.}

\prose{ฉสุ พุทฺธธมฺเมสุ ทิยฑฺฒสตํ ธมฺมา, ทิยฑฺฒสตํ อตฺถา, ตีณิ นิรุตฺติสตานิ, ฉ ญาณสตานิ.}

\prosewithcloserruled{ปฏิสมฺภิทา\-ธิกรเณ\csromansharedfootnote{131ce9afae816609}{ปฏิสมฺภิทาปกรเณ (สฺยา)} อฑฺฒ\-นว\-ธมฺมสตานิ\csromansharedfootnote{850e0f67835eda6b}{อฑฺฒนวมานิ ธมฺมสตานิ (สฺยา), อฑฺฒนวมธมฺมสตานิ (ก)}, อฑฺฒ\-นว\-อตฺถสตานิ, นิรุตฺติ\-สหสฺสญฺจ สตฺต จ นิรุตฺติสตานิ, ตีณิ จ ญาณสหสฺสานิ, จตฺตาริ จ ญาณสตานีติ.}{ปฏิสมฺภิทา\-กถา นิฏฺฐิตา.}
\csromanmatikaanchor{mtk.111}
\csromanmatikaanchor{mtk.112}
\csromantocmarknopage{section}{7. ธมฺมจกฺกกถา}{mtk.111}
\bookmark[dest={mtk.111},level=1]{7. ธมฺมจกฺกกถา}
\csromantocmarkref{subsection}{1. สจฺจวาร}{mtk.112}
\bookmark[dest={mtk.112},level=2]{1. สจฺจวาร}
\csromanmarkh{7. ธมฺมจกฺก\-กถา}\csromanmarkhii{1. สจฺจวาร}\csromanheadernomark{2}{7. ธมฺมจกฺก\-กถา 1. สจฺจวาร}
\proseitem{39}{เอกํ สมยํ ภควา พาราณสิยํ วิหรติ ฯเปฯ. อิติ หิทํ อายสฺมโต โกณฺฑญฺญสฺส “อญฺญาสิโกณฺฑญฺโญ”เตฺวว นามํ อโหสิ. (ก) “อิทํ ทุกฺขํ อริยสจฺจนฺ”ติ ปุพฺเพ อนนุสฺสุเตสุ ธมฺเมสุ จกฺขุํ อุทปาทิ, ญาณํ อุทปาทิ, ปญฺญา อุทปาทิ, วิชฺชา อุทปาทิ, อาโลโก อุทปาทิ.}

\prose{จกฺขุํ อุทปาทีติ เกนฏฺเฐน, ญาณํ อุทปาทีติ เกนฏฺเฐน, ปญฺญา อุทปาทีติ เกนฏฺเฐน, วิชฺชา อุทปาทีติ เกนฏฺเฐน, อาโลโก อุทาปาทีติ เกนฏฺเฐน? จกฺขุํ อุทปาทีติ ทสฺสนฏฺเฐน. ญาณํ อุทปาทีติ ญาตฏฺเฐน. ปญฺญา อุทปาทีติ ปชานนฏฺเฐน. วิชฺชา อุทปาทีติ ปฏิเวธ\-ฏฺเฐน. อาโลโก อุทปาทีติ โอภาสฏฺเฐน.}

\prose{จกฺขุํ ธมฺโม, ทสฺสนฏฺโฐ อตฺโถ. ญาณํ ธมฺโม, ญาตฏฺโฐ อตฺโถ. ปญฺญา ธมฺโม, ปชานนฏฺโฐ อตฺโถ. วิชฺชา ธมฺโม, ปฏิเวธฏฺโฐ อตฺโถ, อาโลโก ธมฺโม, โอภาสฏฺโฐ อตฺโถ. อิเม ปญฺจ ธมฺมา ปญฺจ อตฺถา}

\prose{\csromanfolio{343}ทุกฺขวตฺถุกา สจฺจวตฺถุกา สจฺจารมฺมณา สจฺจโคจรา สจฺจสงฺคหิตา สจฺจ\-ปริยาปนฺนา สจฺเจ สมุทาคตา สจฺเจ ฐิตา สจฺเจ ปติฏฺฐิตา.}

\proseitem{40}{ธมฺมจกฺกนฺติ เกนฏฺเฐน ธมฺมจกฺกํ? ธมฺมญฺจ ปวตฺเตติ จกฺกญฺจาติ ธมฺมจกฺกํ, จกฺกญฺจ ปวตฺเตติ ธมฺมญฺจาติ ธมฺมจกฺกํ, ธมฺเมน ปวตฺเตตีติ ธมฺมจกฺกํ, ธมฺมจริยาย ปวตฺเตตีติ ธมฺมจกฺกํ, ธมฺเม ฐิโต ปวตฺเตตีติ ธมฺมจกฺกํ, ธมฺเม ปติฏฺฐิโต ปวตฺเตตีติ ธมฺมจกฺกํ, ธมฺเม ปติฏฺฐาเปนฺโต ปวตฺเตตีติ ธมฺมจกฺกํ, ธมฺเม วสิปฺปตฺโต ปวตฺเตตีติ ธมฺมจกฺกํ, ธมฺเม วสิํ ปาเปนฺโต ปวตฺเตตีติ ธมฺมจกฺกํ, ธมฺเม ปารมิปฺปตฺโต ปวตฺเตตีติ ธมฺมจกฺกํ, ธมฺเม ปารมิํ ปาเปนฺโต ปวตฺเตตีติ ธมฺมจกฺกํ, ธมฺเม เวสารชฺช\-ปฺปตฺโต ปวตฺเตตีติ ธมฺมจกฺกํ, ธมฺเม เวสารชฺชํ ปาเปนฺโต ปวตฺเตตีติ ธมฺมจกฺกํ, ธมฺมํ สกฺกโรนฺโต ปวตฺเตตีติ ธมฺมจกฺกํ, ธมฺมํ ครุํ กโรนฺโต\csromansharedfootnote{bec34a25ac03e8d4}{ครุกโรนฺโต (สฺยา)} ปวตฺเตตีติ ธมฺมจกฺกํ, ธมฺมํ มาเนนฺโต ปวตฺเตตีติ ธมฺมจกฺกํ, ธมฺมํ ปูเชนฺโต ปวตฺเตตีติ ธมฺมจกฺกํ, ธมฺมํ อปจายมาโน ปวตฺเตตีติ ธมฺมจกฺกํ, ธมฺมทฺธโช ปวตฺเตตีติ ธมฺมจกฺกํ, ธมฺมเกตุ ปวตฺเตตีติ ธมฺมจกฺกํ, ธมฺมา\-ธิปเตยฺโย ปวตฺเตตีติ ธมฺมจกฺกํ. ตํ โข ปน ธมฺมจกฺกํ อปฺปฏิวตฺติยํ สมเณน วา พฺราหฺมเณน วา เทเวน วา มาเรน วา พฺรหฺมุนา วา เกนจิ วา โลกสฺมินฺติ ธมฺมจกฺกํ.}

\prose{สทฺธินฺทฺริยํ ธมฺโม, ตํ ธมฺมํ ปวตฺเตตีติ ธมฺมจกฺกํ, วีริยินฺทฺริยํ ธมฺโม, ตํ ธมฺมํ ปวตฺเตตีติ ธมฺมจกฺกํ, สตินฺทฺริยํ ธมฺโม, ตํ ธมฺมํ ปวตฺเตตีติ ธมฺมจกฺกํ, สมาธินฺทฺริยํ ธมฺโม, ตํ ธมฺมํ ปวตฺเตตีติ ธมฺมจกฺกํ, ปญฺญินฺทฺริยํ ธมฺโม ตํ ธมฺมํ ปวตฺเตตีติ ธมฺมจกฺกํ.  สทฺธาพลํ ธมฺโม, ตํ ธมฺมํ ปวตฺเตตีติ ธมฺมจกฺกํ, วีริยพลํ ธมฺโม, ตํ ธมฺมํ ปวตฺเตตีติ ธมฺมจกฺกํ, สติพลํ ธมฺโม, ตํ ธมฺมํ ปวตฺเตตีติ ธมฺมจกฺกํ, สมาธิพลํ ธมฺโม, ตํ ธมฺมํ ปวตฺเตตีติ ธมฺมจกฺกํ, ปญฺญาพลํ ธมฺโม, ตํ ธมฺมํ ปวตฺเตตีติ ธมฺมจกฺกํ.  สติสมฺโพชฺฌงฺโค ธมฺโม, ตํ ธมฺมํ ปวตฺเตตีติ ธมฺมจกฺกํ, ธมฺมวิจย\-สมฺโพชฺฌงฺโค ธมฺโม, ตํ ธมฺมํ ปวตฺเตตีติ ธมฺมจกฺกํ, วีริย\-สมฺโพชฺฌงฺโค ธมฺโม, ตํ ธมฺมํ ปวตฺเตตีติ ธมฺมจกฺกํ, ปีติสมฺโพชฺฌงฺโค ธมฺโม, ตํ ธมฺมํ ปวตฺเตตีติ ธมฺมจกฺกํ, ปสฺสทฺธิ\-สมฺโพชฺฌงฺโค ธมฺโม, ตํ ธมฺมํ ปวตฺเตตีติ ธมฺมจกฺกํ, สมาธิ\-สมฺโพชฺฌงฺโค ธมฺโม, ตํ ธมฺมํ ปวตฺเตตีติ ธมฺมจกฺกํ, อุเปกฺขา\-สมฺโพชฺฌงฺโค ธมฺโม, ตํ ธมฺมํ \csromanfolio{344}ปวตฺเตตีติ ธมฺมจกฺกํ.  สมฺมาทิฏฺฐิ ธมฺโม, ตํ ธมฺมํ ปวตฺเตตีติ ธมฺมจกฺกํ, สมฺมาสงฺกปฺโป ธมฺโม, ตํ ธมฺมํ ปวตฺเตตีติ ธมฺมจกฺกํ, สมฺมาวาจา ธมฺโม, ตํ ธมฺมํ ปวตฺเตตีติ ธมฺมจกฺกํ, สมฺมากมฺมนฺโต ธมฺโม, ตํ ธมฺมํ ปวตฺเตตีติ ธมฺมจกฺกํ, สมฺมาอาชีโว ธมฺโม, ตํ ธมฺมํ ปวตฺเตตีติ ธมฺมจกฺกํ, สมฺมาวายาโม ธมฺโม, ตํ ธมฺมํ ปวตฺเตตีติ ธมฺมจกฺกํ, สมฺมาสติ ธมฺโม, ตํ ธมฺมํ ปวตฺเตตีติ ธมฺมจกฺกํ, สมฺมาสมาธิ ธมฺโม, ตํ ธมฺมํ ปวตฺเตตีติ ธมฺมจกฺกํ.}

\prose{อาธิปเตยฺย\-ฏฺเฐน อินฺทฺริยํ ธมฺโม, ตํ ธมฺมํ ปวตฺเตตีติ ธมฺมจกฺกํ, อกมฺปิยฏฺเฐน พลํ ธมฺโม, ตํ ธมฺมํ ปวตฺเตตีติ ธมฺมจกฺกํ, นิยฺยา\-นิกฏฺเฐน โพชฺฌงฺโค ธมฺโม, ตํ ธมฺมํ ปวตฺเตตีติ ธมฺมจกฺกํ, เหตุฏฺเฐน มคฺโค ธมฺโม, ตํ ธมฺมํ ปวตฺเตตีติ ธมฺมจกฺกํ, อุปฏฺฐาน\-ฏฺเฐน สติปฏฺฐานา ธมฺโม, ตํ ธมฺมํ ปวตฺเตตีติ ธมฺมจกฺกํ, ปทหนฏฺเฐน สมฺมปฺปธานา ธมฺโม, ตํ ธมฺมํ ปวตฺเตตีติ ธมฺมจกฺกํ, อิชฺฌนฏฺเฐน อิทฺธิปาทา ธมฺโม, ตํ ธมฺมํ ปวตฺเตตีติ ธมฺมจกฺกํ, ตถฏฺเฐน สจฺจา ธมฺโม, ตํ ธมฺมํ ปวตฺเตตีติ ธมฺมจกฺกํ, อวิกฺเขป\-ฏฺเฐน สมโถ ธมฺโม, ตํ ธมฺมํ ปวตฺเตตีติ ธมฺมจกฺกํ, อนุปสฺสนฏฺเฐน วิปสฺสนา ธมฺโม, ตํ ธมฺมํ ปวตฺเตตีติ ธมฺมจกฺกํ, เอกรสฏฺเฐน สมถ\-วิปสฺสนา ธมฺโม, ตํ ธมฺมํ ปวตฺเตตีติ ธมฺมจกฺกํ, อนติวตฺตน\-ฏฺเฐน ยุคนทฺธํ ธมฺโม, ตํ ธมฺมํ ปวตฺเตตีติ ธมฺมจกฺกํ, สํวรฏฺเฐน สีลวิสุทฺธิ ธมฺโม, ตํ ธมฺมํ ปวตฺเตตีติ ธมฺมจกฺกํ อวิกฺเขป\-ฏฺเฐน จิตฺตวิสุทฺธิ ธมฺโม, ตํ ธมฺมํ ปวตฺเตตีติ ธมฺมจกฺกํ, ทสฺสนฏฺเฐน ทิฏฺฐิวิสุทฺธิ ธมฺโม, ตํ ธมฺมํ ปวตฺเตตีติ ธมฺมจกฺกํ, มุตฺตฏฺเฐน วิโมกฺโข ธมฺโม, ตํ ธมฺมํ ปวตฺเตตีติ ธมฺมจกฺกํ, ปฏิเวธ\-ฏฺเฐน วิชฺชา ธมฺโม, ตํ ธมฺมํ ปวตฺเตตีติ ธมฺมจกฺกํ, ปริจฺจาค\-ฏฺเฐน วิมุตฺติ ธมฺโม, ตํ ธมฺมํ ปวตฺเตตีติ ธมฺมจกฺกํ, สมุจฺเฉท\-ฏฺเฐน ขเย ญาณํ ธมฺโม, ตํ ธมฺมํ ปวตฺเตตีติ ธมฺมจกฺกํ, ปฏิปฺปสฺสทฺธ\-ฏฺเฐน อนุปฺปาเท ญาณํ ธมฺโม, ตํ ธมฺมํ ปวตฺเตตีติ ธมฺมจกฺกํ, ฉนฺโท มูลฏฺเฐน ธมฺโม, ตํ ธมฺมํ ปวตฺเตตีติ ธมฺมจกฺกํ, มนสิกาโร สมุฏฺฐาน\-ฏฺเฐน ธมฺโม, ตํ ธมฺมํ ปวตฺเตตีติ ธมฺมจกฺกํ, ผสฺโส สโมธาน\-ฏฺเฐน ธมฺโม, ตํ ธมฺมํ ปวตฺเตตีติ ธมฺมจกฺกํ, เวทนา สโมสรณ\-ฏฺเฐน ธมฺโม, ตํ ธมฺมํ ปวตฺเตตีติ ธมฺมจกฺกํ, สมาธิ ปมุขฏฺเฐน ธมฺโม, ตํ ธมฺมํ ปวตฺเตตีติ ธมฺมจกฺกํ, สติ อาธิปเตยฺย\-ฏฺเฐน ธมฺโม, ตํ ธมฺมํ ปวตฺเตตีติ ธมฺมจกฺกํ, ปญฺญา ตทุตฺต\-รฏฺเฐน ธมฺโม, ตํ ธมฺมํ ปวตฺเตตีติ ธมฺมจกฺกํ, วิมุตฺติ สารฏฺเฐน ธมฺโม, ตํ ธมฺมํ ปวตฺเตตีติ ธมฺมจกฺกํ, อมโตคธํ นิพฺพานํ ปริโยสาน\-ฏฺเฐน ธมฺโม, ตํ ธมฺมํ ปวตฺเตตีติ ธมฺมจกฺกํ.}

\prose{\csromanfolio{345}“ตํ โข ปนิทํ ทุกฺขํ อริยสจฺจํ ปริญฺเญยฺยนฺ”ติ ฯเปฯ “ปริญฺญาตนฺ”ติ ปุพฺเพ อนนุสฺสุเตสุ ธมฺเมสุ จกฺขุํ อุทปาทิ ฯเปฯ อาโลโก อุทปาทิ.}

\prose{จกฺขุํ อุทปาทีติ เกนฏฺเฐน ฯเปฯ อาโลโก อุทปาทีติ เกนฏฺเฐน? จกฺขุํ อุทปาทีติ ทสฺสนฏฺเฐน ฯเปฯ. อาโลโก อุทปาทีติ โอภาสฏฺเฐน. จกฺขุํ ธมฺโม, ทสฺสนฏฺโฐ อตฺโถ ฯเปฯ อาโลโก ธมฺโม โอภาสฏฺโฐ อตฺโถ. อิเม ปญฺจ ธมฺมา ปญฺจ อตฺถา ทุกฺขวตฺถุกา สจฺจวตฺถุกา สจฺจารมฺมณา สจฺจโคจรา สจฺจสงฺคหิตา สจฺจ\-ปริยาปนฺนา สจฺเจ สมุทาคตา สจฺเจ ฐิตา สจฺเจ ปติฏฺฐิตา.}

\prose{ธมฺมจกฺกนฺติ เกนฏฺเฐน ธมฺมจกฺกํ? ธมฺมญฺจ ปวตฺเตติ จกฺกญฺจาติ ธมฺมจกฺกํ, จกฺกญฺจ ปวตฺเตติ ธมฺมญฺจาติ ธมฺมจกฺกํ, ธมฺเมน ปวตฺเตตีติ ธมฺมจกฺกํ, ธมฺมจริยาย ปวตฺเตตีติ ธมฺมจกฺกํ, ธมฺเม ฐิโต ปวตฺเตตีติ ธมฺมจกฺกํ, ธมฺเม ปติฏฺฐิโต ปวตฺเตตีติ ธมฺมจกฺกํ ฯเปฯ อมโตคธํ นิพฺพานํ ปริโยสาน\-ฏฺเฐน ธมฺโม, ตํ ธมฺมํ ปวตฺเตตีติ ธมฺมจกฺกํ.}

\prose{(ข. ฆ) “อิทํ ทุกฺข\-สมุทยํ อริยสจฺจนฺ”ติ ปุพฺเพ อนนุสฺสุเตสุ ธมฺเมสุ จกฺขุํ อุทปาทิ ฯเปฯ อาโลโก อุทปาทิ ฯเปฯ “ตํ โข ปนิทํ ทุกฺข\-สมุทยํ อริยสจฺจํ ปหาตพฺพนฺ”ติ ฯเปฯ “ปหีนนฺ”ติ ปุพฺเพ อนนุสฺสุเตสุ ธมฺเมสุ จกฺขุํ อุทปาทิ ฯเปฯ อาโลโก อุทปาทิ.}

\prose{จกฺขุํ อุทปาทีติ เกนฏฺเฐน ฯเปฯ อาโลโก อุทปาทีติ เกนฏฺเฐน? \textbf{จกฺขุํ อุทปาที}ติ ทสฺสนฏฺเฐน ฯเปฯ. \textbf{อาโลโก อุทปาที}ติ โอภาสฏฺเฐน.}

\prose{จกฺขุํ ธมฺโม, ทสฺสนฏฺโฐ อตฺโถ ฯเปฯ อาโลโก ธมฺโม, โอภาสฏฺโฐ อตฺโถ. อิเม ปญฺจ ธมฺมา ปญฺจ อตฺถา สมุทย\-วตฺถุกา สจฺจวตฺถุกา ฯเปฯ นิโรธ\-วตฺถุกา สจฺจวตฺถุกา ฯเปฯ มคฺควตฺถุกา สจฺจวตฺถุกา สจฺจารมฺมณา สจฺจโคจรา สจฺจสงฺคหิตา สจฺจ\-ปริยาปนฺนา สจฺเจ สมุทาคตา สจฺเจ ฐิตา สจฺเจ ปติฏฺฐิตา.}

\prose{ธมฺมจกฺกนฺติ เกนฏฺเฐน ธมฺมจกฺกํ? ธมฺมญฺจ ปวตฺเตติ จกฺกญฺจาติ ธมฺมจกฺกํ, จกฺกญฺจ ปวตฺเตติ ธมฺมญฺจาติ ธมฺมจกฺกํ, ธมฺเมน ปวตฺเตตีติ ธมฺมจกฺกํ, ธมฺมจริยาย ปวตฺเตตีติ ธมฺมจกฺกํ, ธมฺเม ฐิโต ปวตฺเตตีติ ธมฺมจกฺกํ, ธมฺเม ปติฏฺฐิโต ปวตฺเตตีติ ธมฺมจกฺกํ ฯเปฯ อมโตคธํ นิพฺพานํ ปริโยสาน\-ฏฺเฐน ธมฺโม, ตํ ธมฺมํ ปวตฺเตตีติ ธมฺมจกฺกํ.}

\csromanmatikaanchor{mtk.113}
\csromantocmarkref{subsection}{2. สติปฏฺฐานวาร}{mtk.113}
\bookmark[dest={mtk.113},level=2]{2. สติปฏฺฐานวาร}
\csromanheader{2}{2. สติปฏฺฐาน\-วาร}
\proseitem{41}{\csromanfolio{346}“อยํ กาเย กายานุปสฺสนา”ติ เม ภิกฺขเว ปุพฺเพ อนนุสฺสุเตสุ ธมฺเมสุ จกฺขุํ อุทปาทิ ฯเปฯ อาโลโก อุทปาทิ. “สา โข ปนายํ กาเย กายานุปสฺสนา ภาเวตพฺพา”ติ เม ภิกฺขเว ฯเปฯ “ภาวิตา”ติ เม ภิกฺขเว ปุพฺเพ อนนุสฺสุเตสุ ธมฺเมสุ จกฺขุํ อุทปาทิ ฯเปฯ อาโลโก อุทปาทิ.}

\prose{“อยํ เวทนาสุ ฯเปฯ. อยํ จิตฺเต. อยํ ธมฺเมสุ ธมฺมา\-นุปสฺสนา”ติ เม ภิกฺขเว ปุพฺเพ อนนุสฺสุเตสุ ธมฺเมสุ จกฺขุํ อุทปาทิ ฯเปฯ อาโลโก อุทปาทิ. “สา โข ปนายํ ธมฺเมสุ ธมฺมา\-นุปสฺสนา ภาเวตพฺพา”ติ เม ภิกฺขเว ฯเปฯ “ภาวิตา”ติ เม ภิกฺขเว ปุพฺเพ อนนุสฺสุเตสุ ธมฺเมสุ จกฺขุํ อุทปาทิ ฯเปฯ อาโลโก อุทปาทิ.}

\prose{“อยํ กาเย กายานุปสฺสนา”ติ ปุพฺเพ อนนุสฺสุเตสุ ธมฺเมสุ จกฺขุํ อุทปาทิ ฯเปฯ อาโลโก อุทปาทิ ฯเปฯ “สา โข ปนายํ กาเย กายานุปสฺสนา ภาเวตพฺพา”ติ ฯเปฯ “ภาวิตา”ติ ปุพฺเพ อนนุสฺสุเตสุ ธมฺเมสุ จกฺขุํ อุทปาทิ ฯเปฯ อาโลโก อุทปาทิ.}

\prose{จกฺขุํ อุทปาทีติ เกนฏฺเฐน ฯเปฯ อาโลโก อุทปาทีติ เกนฏฺเฐน? \textbf{จกฺขุํ อุทปาที}ติ ทสฺสนฏฺเฐน ฯเปฯ. \textbf{อาโลโก} \textbf{อุทปาที}ติ โอภาสฏฺเฐน.}

\prose{จกฺขุํ ธมฺโม, ทสฺสนฏฺโฐ อตฺโถ ฯเปฯ อาโลโก ธมฺโม, โอภาสฏฺโฐ อตฺโถ. อิเม ปญฺจ ธมฺมา ปญฺจ อตฺถา กายวตฺถุกา สติปฏฺฐาน\-วตฺถุกา ฯเปฯ เวทนาวตฺถุกา สติปฏฺฐาน\-วตฺถุกา, จิตฺตวตฺถุกา สติปฏฺฐาน\-วตฺถุกา. ธมฺมวตฺถุกา สติปฏฺฐาน\-วตฺถุกา สติปฏฺฐา\-นารมฺมณา สติปฏฺฐาน\-โคจรา สติปฏฺฐาน\-สงฺคหิตา สติปฏฺฐาน\-ปริยาปนฺนา สติปฏฺฐาเน สมุทาคตา สติปฏฺฐาเน ฐิตา สติปฏฺฐาเน ปติฏฺฐิตา.}

\prose{ธมฺมจกฺกนฺติ เกนฏฺเฐน ธมฺมจกฺกํ? ธมฺมญฺจ ปวตฺเตติ จกฺกญฺจาติ ธมฺมจกฺกํ, จกฺกญฺจ ปวตฺเตติ ธมฺมญฺจาติ ธมฺมจกฺกํ, ธมฺเมน ปวตฺเตตีติ ธมฺมจกฺกํ, ธมฺมจริยาย ปวตฺเตตีติ ธมฺมจกฺกํ, ธมฺเม ฐิโต ปวตฺเตตีติ ธมฺมจกฺกํ, ธมฺเม ปติฏฺฐิโต ปวตฺเตตีติ ธมฺมจกฺกํ ฯเปฯ อมโตคธํ นิพฺพานํ ปริโยสาน\-ฏฺเฐน ธมฺโม, ตํ ธมฺมํ ปวตฺเตตีติ ธมฺมจกฺกํ.}

\csromanheader{2}{7. ธมฺมจกฺก\-กถา}
\csromanmatikaanchor{mtk.114}
\csromantocmarkref{subsection}{3. อิทฺธิปาทวาร}{mtk.114}
\bookmark[dest={mtk.114},level=2]{3. อิทฺธิปาทวาร}
\csromanheader{2}{3. \textbf{อิทฺธิปาทวาร}}
\proseitem{42}{\csromanfolio{347}“อยํ ฉนฺท\-สมาธิ\-ปธาน\-สงฺขาร\-สมนฺนาคโต อิทฺธิปาโท”ติ เม ภิกฺขเว ปุพฺเพ อนนุสฺสุเตสุ ธมฺเมสุ จกฺขุํ อุทปาทิ ฯเปฯ อาโลโก อุทปาทิ. “โส โข ปนายํ ฉนฺท\-สมาธิ\-ปธาน\-สงฺขาร\-สมนฺนาคโต อิทฺธิปาโท ภาเวตพฺโพ”ติ เม ภิกฺขเว ฯเปฯ “ภาวิโต”ติ เม ภิกฺขเว ปุพฺเพ อนนุสฺสุเตสุ ธมฺเมสุ จกฺขุํ อุทปาทิ ฯเปฯ อาโลโก อุทปาทิ.}

\prose{“อยํ วีริยสมาธิ ฯเปฯ. อยํ จิตฺตสมาธิ ฯเปฯ. อยํ วีมํสา\-สมาธิ\-ปธาน\-สงฺขาร\-สมนฺนาคโต อิทฺธิปาโท”ติ เม ภิกฺขเว ปุพฺเพ อนนุสฺสุเตสุ ธมฺเมสุ จกฺขุํ อุทปาทิ ฯเปฯ อาโลโก อุทปาทิ. “โส โข ปนายํ วีมํสา\-สมาธิ\-ปธาน\-สงฺขาร\-สมนฺนาคโต อิทฺธิปาโท ภาเวตพฺโพ”ติ เม ภิกฺขเว ฯเปฯ “ภาวิโต”ติ เม ภิกฺขเว ปุพฺเพ อนนุสฺสุเตสุ ธมฺเมสุ จกฺขุํ อุทปาทิ ฯเปฯ อาโลโก อุทปาทิ.}

\prose{“อยํ ฉนฺท\-สมาธิ\-ปธาน\-สงฺขาร\-สมนฺนาคโต อิทฺธิปาโท”ติ ปุพฺเพ อนนุสฺสุเตสุ ธมฺเมสุ จกฺขุํ อุทปาทิ ฯเปฯ อาโลโก อุทปาทิ ฯเปฯ “โส โข ปนายํ ฉนฺท\-สมาธิ\-ปธาน\-สงฺขาร\-สมนฺนาคโต อิทฺธิปาโท ภาเวตพฺโพ”ติ ฯเปฯ “ภาวิโต”ติ ปุพฺเพ อนนุสฺสุเตสุ ธมฺเมสุ จกฺขุํ อุทปาทิ, ญาณํ อุทปาทิ, ปญฺญา อุทปาทิ, วิชฺชา อุทปาทิ, อาโลโก อุทปาทิ.}

\prose{จกฺขุํ อุทปาทีติ เกนฏฺเฐน, ญาณํ อุทปาทีติ เกนฏฺเฐน, ปญฺญา อุทปาทีติ เกนฏฺเฐน วิชฺชา อุทปาทีติ เกนฏฺเฐน, อาโลโก อุทปาทีติ เกนฏฺเฐน? จกฺขุํ อุทปาทีติ ทสฺสนฏฺเฐน, ญาณํ อุทปาทีติ ญาตฏฺเฐน, ปญฺญา อุทปาทีติ ปชานนฏฺเฐน, วิชฺชา อุทปาทีติ ปฏิเวธ\-ฏฺเฐน, อาโลโก อุทปาทีติ โอภาสฏฺเฐน.}

\prose{จกฺขุํ ธมฺโม, ทสฺสนฏฺโฐ อตฺโถ. ญาณํ ธมฺโม, ญาตตฺโถ อตฺโถ. ปญฺญา ธมฺโม, ปชานนฏฺโฐ อตฺโถ. วิชฺชา ธมฺโม, ปฏิเวธฏฺโฐ อตฺโถ. อาโลโก ธมฺโม, โอภาสฏฺโฐ อตฺโถ. อิเม ปญฺจ ธมฺมา ปญฺจ อตฺถา ฉนฺทวตฺถุกา อิทฺธิปาท\-วตฺถุกา อิทฺธิปาทา\-รมฺมณา อิทฺธิปาท\-โคจรา อิทฺธิปาท\-สงฺคหิตา อิทฺธิปาท\-ปริยาปนฺนา อิทฺธิปาเท สมุทาคตา อิทฺธิปาเท ฐิตา อิทฺธิปาเท ปติฏฺฐิตา.}

\prose{\textbf{ธมฺมจกฺกนฺ}ติ เกนฏฺเฐน ธมฺมจกฺกํ? ธมฺมญฺจ ปวตฺเตติ จกฺกญฺจาติ ธมฺมจกฺกํ, จกฺกญฺจ ปวตฺเตติ ธมฺมญฺจาติ ธมฺมจกฺกํ, ธมฺเมน ปวตฺเตตีติ ธมฺมจกฺกํ, ธมฺมจริยาย \csromanfolio{348}ปวตฺเตตีติ ธมฺมจกฺกํ, ธมฺเม ฐิโต ปวตฺเตตีติ ธมฺมจกฺกํ, ธมฺเม ปติฏฺฐิโต ปวตฺเตตีติ ธมฺมจกฺกํ, ธมฺเม ปติฏฺฐาเปนฺโต ปวตฺเตตีติ ธมฺมจกฺกํ, ธมฺเม วสิปฺปตฺโต ปวตฺเตตีติ ธมฺมจกฺกํ, ธมฺเม วสิํ ปาเปนฺโต ปวตฺเตตีติ ธมฺมจกฺกํ, ธมฺเม ปารมิปฺปตฺโต ปวตฺเตตีติ ธมฺมจกฺกํ, ธมฺเม ปารมิํ ปาเปนฺโต ปวตฺเตตีติ ธมฺมจกฺกํ ฯเปฯ. ธมฺมํ อปจายมาโน ปวตฺเตตีติ ธมฺมจกฺกํ, ธมฺมทฺธโช ปวตฺเตตีติ ธมฺมจกฺกํ, ธมฺมเกตุ ปวตฺเตตีติ ธมฺมจกฺกํ, ธมฺมา\-ธิปเตยฺโย ปวตฺเตตีติ ธมฺมจกฺกํ. ตํ โข ปน ธมฺมจกฺกํ อปฺปฏิวตฺติยํ สมเณน วา พฺราหฺมเณน วา เทเวน วา มาเรน วา พฺรหฺมุนา วา เกนจิ วา โลกสฺมินฺติ ธมฺมจกฺกํ.}

\prose{สทฺธินฺทฺริยํ ธมฺโม, ตํ ธมฺมํ ปวตฺเตตีติ ธมฺมจกฺกํ ฯเปฯ อมโตคธํ นิพฺพานํ ปริโยสาน\-ฏฺเฐน ธมฺโม, ตํ ธมฺมํ ปวตฺเตตีติ ธมฺมจกฺกํ.}

\prose{“อยํ วีริย\-สมาธิ\-ปธาน\-สงฺขาร\-สมนฺนาคโต อิทฺธิปาโท”ติ ปุพฺเพ อนนุสฺสุเตสุ ธมฺเมสุ จกฺขุํ อุทปาทิ ฯเปฯ อาโลโก อุทปาทิ ฯเปฯ “โส โข ปนายํ วีริย\-สมาธิ\-ปธาน\-สงฺขาร\-สมนฺนาคโต อิทฺธิปาโท ภาเวตพฺโพ”ติ ฯเปฯ “ภาวิโต”ติ ปุพฺเพ อนนุสฺสุเตสุ ธมฺเมสุ จกฺขุํ อุทปาทิ ฯเปฯ อาโลโก อุทปาทิ.}

\prose{จกฺขุํ อุทปาทีติ เกนฏฺเฐน ฯเปฯ อาโลโก อุทปาทีติ เกนฏฺเฐน? \textbf{จกฺขุํ อุทปาที}ติ ทสฺสนฏฺเฐน. \textbf{อาโลโก อุทปาที}ติ โอภาสฏฺเฐน.}

\prose{จกฺขุํ ธมฺโม, ทสฺสนฏฺโฐ อตฺโถ ฯเปฯ อาโลโก ธมฺโม, โอภาสฏฺโฐ อตฺโถ. อิเม ปญฺจ ธมฺมา ปญฺจ อตฺถา วีริยวตฺถุกา อิทฺธิปาท\-วตฺถุกา ฯเปฯ จิตฺตวตฺถุกา อิทฺธิปาท\-วตฺถุกา. วีมํสา\-วตฺถุกา อิทฺธิปาท\-วตฺถุกา อิทฺธิปาทา\-รมฺมณา อิทฺธิปาท\-โคจรา อิทฺธิปาท\-สงฺคหิตา อิทฺธิปาท\-ปริยาปนฺนา อิทฺธิปาเท สมุทาคตา อิทฺธิปาเท ฐิตา อิทฺธิปาเท ปติฏฺฐิตา.}

\prosewithcloser{ธมฺมจกฺกนฺติ เกนฏฺเฐน ธมฺมจกฺกํ? ธมฺมญฺจ ปวตฺเตติ จกฺกญฺจาติ ธมฺมจกฺกํ, จกฺกญฺจ ปวตฺเตติ ธมฺมญฺจาติ ธมฺมจกฺกํ, ธมฺเมน ปวตฺเตตีติ ธมฺมจกฺกํ, ธมฺมจริยาย ปวตฺเตตีติ ธมฺมจกฺกํ, ธมฺเม ฐิโต ปวตฺเตตีติ ธมฺมจกฺกํ, ธมฺเม ปติฏฺฐิโต ปวตฺเตตีติ ธมฺมจกฺกํ ฯเปฯ อมโตคธํ นิพฺพานํ ปริโยสาน\-ฏฺเฐน ธมฺโม, ตํ ธมฺมํ ปวตฺเตตีติ ธมฺมจกฺกนฺติ.}{ธมฺมจกฺก\-กถา นิฏฺฐิตา.}
\csromanmatikaanchor{mtk.115}
\csromantocmarkref{section}{8. โลกุตฺตรกถา}{mtk.115}
\bookmark[dest={mtk.115},level=1]{8. โลกุตฺตรกถา}
\csromanheader{1}{8. โลกุตฺตรกถา}
\csromanheader{2}{8. โลกุตฺตรกถา}
\proseitem{43}{\csromanfolio{349}กตเม ธมฺมา โลกุตฺตรา? จตฺตาโร สติปฏฺฐานา จตฺตาโร สมฺมปฺปธานา จตฺตาโร อิทฺธิปาทา ปญฺจินฺทฺริยานิ ปญฺจ พลานิ สตฺต โพชฺฌงฺคา อริโย อฏฺฐงฺคิโก มคฺโค จตฺตาโร อริยมคฺคา จตฺตาริ จ สามญฺญผลานิ นิพฺพานญฺจ. อิเม ธมฺมา โลกุตฺตรา.}

\csromanmatikaanchor{mtk.116}
\csromantocmarkref{section}{9. พลกถา}{mtk.116}
\bookmark[dest={mtk.116},level=1]{9. พลกถา}
\prosewithcloserruled{โลกุตฺตราติ เกนฏฺเฐน โลกุตฺตรา? โลกํ ตรนฺตีติ โลกุตฺตรา, โลกา อุตฺตรนฺตีติ โลกุตฺตรา, โลกโต อุตฺตรนฺตีติ โลกุตฺตรา, โลกมฺหา อุตฺตรนฺตีติ โลกุตฺตรา, โลกํ อติกฺกมนฺตีติ โลกุตฺตรา, โลกํ สมติ\-กฺกม\-นฺตีติ โลกุตฺตรา. โลกํ สมติกฺกนฺตาติ โลกุตฺตรา, โลเกน อติเรกาหิ โลกุตฺตรา, โลกนฺตํ ตรนฺตีติ โลกุตฺตรา, โลกานิสฺสรนฺตีติ โลกุตฺตรา, โลกโต นิสฺสรนฺตีติ โลกุตฺตรา, โลกมฺหา นิสฺสรนฺตีติ โลกุตฺตรา, โลกา นิสฏาติ โลกุตฺตรา, โลเกน นิสฺสฏาติ โลกุตฺตรา, โลกมฺหา นิสฺสฏาติ โลกุตฺตรา.  โลเก น ติฏฺฐนฺตีติ โลกุตฺตรา, โลกสฺมิํ น ติฏฺฐนฺตีติ โลกุตฺตรา, โลเก น ลิมฺปนฺตีติ โลกุตฺตรา, โลเกน น ลิมฺปนฺตีติ โลกุตฺตรา, โลเก อสํลิตฺตาติ โลกุตฺตรา, โลเกน อสํลิตฺตาติ โลกุตฺตรา, โลเก อนุปลิตฺตาติ โลกุตฺตรา, โลเกน อนุปลิตฺตาติ โลกุตฺตรา. โลเก วิปฺปมุตฺตาติ โลกุตฺตรา โลเกน วิปฺปมุตฺตาติ โลกุตฺตรา, โลกา วิปฺปมุตฺตาติ โลกุตฺตรา, โลกโต วิปฺปมุตฺตาติ โลกุตฺตรา, โลกมฺหา วิปฺปมุตฺตาติ โลกุตฺตรา. โลเก วิสญฺญุตฺตาติ โลกุตฺตรา, โลเกน วิสญฺญุตฺตาติ โลกุตฺตรา, โลกา วิสญฺญุตฺตาติ โลกุตฺตรา, โลกสฺมิํ วิสญฺญุตฺตาติ โลกุตฺตรา, โลกโต วิสญฺญุตฺตาติ โลกุตฺตรา, โลกมฺหา วิสญฺญุตฺตาติ โลกุตฺตรา.  โลกา สุชฺฌนฺตีติ โลกุตฺตรา, โลกโต สุชฺฌนฺตีติ โลกุตฺตรา, โลกมฺหา สุชฺฌนฺตีติ โลกุตฺตรา, โลกา วิสุชฺฌนฺตีติ โลกุตฺตรา, โลกโต วิสุชฺฌนฺตีติ โลกุตฺตรา, โลกมฺหา วิสุชฺฌนฺตีติ โลกุตฺตรา.  โลกา วุฏฺฐหนฺตีติ\csromansharedfootnote{dcedded3d055d053}{อุทฺธรนฺตีติ (ก), อุฏฺฐหนฺตีติ (สี, ฏฺฐ)} โลกุตฺตรา, โลกโต วุฏฺฐหนฺตีติ โลกุตฺตรา, โลกมฺหา วุฏฺฐหนฺตีติ โลกุตฺตรา.  โลกา วิวฏฺฏนฺตีติ \csromanfolio{350}โลกุตฺตรา, โลกโต วิวฏฺฏนฺตีติ โลกุตฺตรา, โลกมฺหา วิวฏฺฏนฺตีติ โลกุตฺตรา.  โลเก น สชฺชนฺตีติ โลกุตฺตรา, โลเก น คยฺหนฺตีติ โลกุตฺตรา, โลเก น พชฺฌนฺตีติ โลกุตฺตรา.  โลกํ สมุจฺฉินฺทนฺตีติ โลกุตฺตรา, โลกํ สมุจฺฉินฺนตฺตาติ โลกุตฺตรา.  โลกํ ปฏิปฺปสฺสมฺเภนฺตีติ โลกุตฺตรา, โลกํ ปฏิปฺปสฺสมฺภิ\-ตตฺตาติ โลกุตฺตรา.  โลกสฺส อปถาติ โลกุตฺตรา, โลกสฺส อคตีติ โลกุตฺตรา, โลกสฺส อวิสยาติ โลกุตฺตรา, โลกสฺส อสาธารณาติ โลกุตฺตรา.  โลกํ วมนฺตีติ โลกุตฺตรา, โลกํ น ปจฺจาวมนฺตีติ โลกุตฺตรา. โลกํ ปชหนฺตีติ โลกุตฺตรา. โลกํ น อุปาทิยนฺตีติ โลกุตฺตรา.  โลกํ วิสีเนนฺตีติ โลกุตฺตรา, โลกํ น อุสฺสีเนนฺตีติ โลกุตฺตรา. โลกํ วิธูเปนฺตีติ โลกุตฺตรา, โลกํ น สํธูเปนฺตีติ โลกุตฺตรา. โลกํ สมติกฺกมฺม อภิภุยฺย ติฏฺฐนฺตีติ โลกุตฺตรา.}{โลกุตฺตรกถา นิฏฺฐิตา.}
\proseitem{44}{สาวตฺถิ\-นิทานํ\csromansharedfootnote{93a5b955b328e3e7}{สํ 3. 218 ปิฏฺเฐ ปสฺสิตพฺพา.}.  ปญฺจิมานิ ภิกฺขเว พลานิ. กตมานิ ปญฺจ? สทฺธาพลํ วีริยพลํ สติพลํ สมาธิพลํ ปญฺญาพลํ. อิมานิ โข ภิกฺขเว ปญฺจ พลานิ.}

\prose{อปิ จ อฏฺฐสฏฺฐิ พลานิ. สทฺธาพลํ, วีริยพลํ, สติพลํ, สมาธิพลํ, ปญฺญาพลํ, หิริพลํ, โอตฺตปฺปพลํ, ปฏิสงฺขานพลํ, ภาวนาพลํ, อนวชฺชพลํ, สงฺคหพลํ, ขนฺติพลํ, ปญฺญตฺติพลํ, นิชฺฌตฺติพลํ, อิสฺสริยพลํ, อธิฏฺฐานพลํ, สมถพลํ, วิปสฺสนาพลํ, ทส เสขพลานิ, ทส อเสขพลานิ, ทส ขีณาสว\-พลานิ, ทส อิทฺธิพลานิ, ทส ตถาคต\-พลานิ.}

\prose{กตมํ สทฺธาพลํ? อสฺสทฺธิเย น กมฺปตีติ สทฺธาพลํ, สหชาตานํ ธมฺมานํ อุปตฺถมฺภน\-ฏฺเฐน สทฺธาพลํ, กิเลสานํ ปริยาทาน\-ฏฺเฐน สทฺธาพลํ, ปฏิเวธา\-ทิ\-วิโสธ\-นฏฺเฐน สทฺธาพลํ, จิตฺตสฺส อธิฏฺฐาน\-ฏฺเฐน สทฺธาพลํ, จิตฺตสฺส โวทานฏฺเฐน สทฺธาพลํ, วิเสสาธิคม\-ฏฺเฐน สทฺธาพลํ, อุตฺตริ-}

\prose{\csromanfolio{351}ปฏิเวธ\-ฏฺเฐน สทฺธาพลํ, สจฺจาภิสมย\-ฏฺเฐน สทฺธาพลํ, นิโรเธ ปติฏฺฐา\-ปกฏฺเฐน สทฺธาพลํ, อิทํ สทฺธาพลํ. (1)}

\prose{กตมํ วีริยพลํ? โกสชฺเช น กมฺปตีติ วีริยพลํ, สหชาตานํ ธมฺมานํ อุปตฺถมฺภน\-ฏฺเฐน วีริยพลํ, กิเลสานํ ปริยาทาน\-ฏฺเฐน วีริยพลํ, ปฏิเวธา\-ทิ\-วิโสธ\-นฏฺเฐน วีริยพลํ, จิตฺตสฺส อธิฏฺฐาน\-ฏฺเฐน วีริยพลํ, จิตฺตสฺส โวทานฏฺเฐน วีริยพลํ, วิเสสาธิคม\-ฏฺเฐน วีริยพลํ, อุตฺตริ\-ปฏิเวธ\-ฏฺเฐน วีริยพลํ, สจฺจาภิสมย\-ฏฺเฐน วีริยพลํ, นิโรเธ ปติฏฺฐา\-ปกฏฺเฐน วีริยพลํ, อิทํ วีริยพลํ. (2)}

\prose{กตมํ สติพลํ? ปมาเท น กมฺปตีติ สติพลํ, สหชาตานํ ธมฺมานํ อุปตฺถมฺภน\-ฏฺเฐน สติพลํ ฯเปฯ นิโรเธ ปติฏฺฐา\-ปกฏฺเฐน สติพลํ, อิทํ สติพลํ. (3)}

\prose{กตมํ สมาธิพลํ? อุทฺธจฺเจ น กมฺปตีติ สมาธิพลํ, สหชาตานํ ธมฺมานํ อุปตฺถมฺภน\-ฏฺเฐน สมาธิพลํ ฯเปฯ นิโรเธ ปติฏฺฐา\-ปกฏฺเฐน สมาธิพลํ, อิทํ สมาธิพลํ. (4)}

\prose{กตมํ ปญฺญาพลํ? อวิชฺชาย น กมฺปตีติ ปญฺญาพลํ, สหชาตานํ ธมฺมานํ อุปตฺถมฺภน\-ฏฺเฐน ปญฺญาพลํ ฯเปฯ นิโรเธ ปติฏฺฐา\-ปกฏฺเฐน ปญฺญาพลํ, อิทํ ปญฺญาพลํ. (5)}

\prose{กตมํ หิริพลํ? เนกฺขมฺเมน กามจฺฉนฺทํ หิรียตีติ\csromansharedfootnote{b34774e9aa378b90}{หิริยตีติ (สฺยา, ก)} หิริพลํ, อพฺยาปาเทน พฺยาปาทํ หิรียตีติ หิริพลํ, อาโลกสญฺญาย ถินมิทฺธํ หิรียตีติ หิริพลํ, อวิกฺเขเปน อุทฺธจฺจํ หิรียตีติ หิริพลํ, ธมฺม\-ววตฺถาเนน วิจิกิจฺฉํ หิรียตีติ หิริพลํ, ญาเณน อวิชฺชํ หิรียตีติ หิริพลํ, ปาโมชฺเชน อรติํ หิรียตีติ หิริพลํ, ปฐเมน ฌาเนน นีวรเณ หิรียตีติ หิริพลํ ฯเปฯ อรหตฺต\-มคฺเคน สพฺพกิเลเส หิรียตีติ หิริพลํ, อิทํ หิริพลํ. (6)}

\prose{กตมํ โอตฺตปฺปพลํ? เนกฺขมฺเมน กามจฺฉนฺทํ โอตฺตปฺปตีติ โอตฺตปฺปพลํ, อพฺยาปาเทน พฺยาปาทํ โอตฺตปฺปตีติ โอตฺตปฺปพลํ, อาโลกสญฺญาย ถินมิทฺธํ โอตฺตปฺปตีติ โอตฺตปฺปพลํ, อวิกฺเขเปน อุทฺธจฺจํ โอตฺตปฺปตีติ โอตฺตปฺปพลํ, ธมฺม\-ววตฺถาเนน วิจิกิจฺฉํ โอตฺตปฺปตีติ โอตฺตปฺปพลํ, ญาเณน อวิชฺชํ โอตฺตปฺปตีติ}

\prose{\csromanfolio{352}โอตฺตปฺปพลํ, ปาโมชฺเชน อรติํ โอตฺตปฺปตีติ โอตฺตปฺปพลํ, ปฐเมน ฌาเนน นีวรเณ โอตฺตปฺปตีติ โอตฺตปฺปพลํ ฯเปฯ อรหตฺต\-มคฺเคน สพฺพกิเลเส โอตฺตปฺปตีติ โอตฺตปฺปพลํ, อิทํ โอตฺตปฺปพลํ. (7)}

\prose{กตมํ ปฏิสงฺขานพลํ? เนกฺขมฺเมน กามจฺฉนฺทํ ปฏิสงฺขาตีติ ปฏิสงฺขานพลํ. อพฺยาปาเทน พฺยาปาทํ ปฏิสงฺขาตีติ ปฏิสงฺขานพลํ, อาโลกสญฺญาย ถินมิทฺธํ ปฏิสงฺขาตีติ ปฏิสงฺขานพลํ, อวิกฺเขเปน อุทฺธจฺจํ ปฏิสงฺขาตีติ ปฏิสงฺขานพลํ, ธมฺม\-ววตฺถาเนน วิจิกิจฺฉํ ปฏิสงฺขาตีติ ปฏิสงฺขานพลํ, ญาเณน อวิชฺชํ ปฏิสงฺขาตีติ ปฏิสงฺขานพลํ, ปาโมชฺเชน อรติํ ปฏิสงฺขาตีติ ปฏิสงฺขานพลํ, ปฐเมน ฌาเนน นีวรเณ ปฏิสงฺขาตีติ ปฏิสงฺขานพลํ ฯเปฯ อรหตฺต\-มคฺเคน สพฺพกิเลเส ปฏิสงฺขาตีติ ปฏิสงฺขานพลํ, อิทํ ปฏิสงฺขานพลํ. (8)}

\prose{กตมํ ภาวนาพลํ? กามจฺฉนฺทํ ปชหนฺโต เนกฺขมฺมํ ภาเวตีติ ภาวนาพลํ, พฺยาปาทํ ปชหนฺโต อพฺยาปาทํ ภาเวตีติ ภาวนาพลํ, ถินมิทฺธํ ปชหนฺโต อาโลกสญฺญํ ภาเวตีติ ภาวนาพลํ, อุทฺธจฺจํ ปชหนฺโต อวิกฺเขปํ ภาเวตีติ ภาวนาพลํ, วิจิกิจฺฉํ ปชหนฺโต ธมฺม\-ววตฺถานํ ภาเวตีติ ภาวนาพลํ, อวิชฺชํ ปชหนฺโต ญาณํ ภาเวตีติ ภาวนาพลํ, อรติํ ปชหนฺโต ปาโมชฺชํ ภาเวตีติ ภาวนาพลํ, นีวรเณ ปชหนฺโต ปฐมํ ฌานํ ภาเวตีติ ภาวนาพลํ ฯเปฯ สพฺพกิเลเส ปชหนฺโต อรหตฺตมคฺคํ ภาเวตีติ ภาวนาพลํ, อิทํ ภาวนาพลํ. (9)}

\prose{กตมํ อนวชฺชพลํ? กามจฺฉนฺทสฺส ปหีนตฺตา เนกฺขมฺเม นตฺถิ กิญฺจิ วชฺชนฺติ อนวชฺชพลํ, พฺยาปาทสฺส ปหีนตฺตา อพฺยาปาเท นตฺถิ กิญฺจิ วชฺชนฺติ อนวชฺชพลํ, ถินมิทฺธสฺส ปหีนตฺตา อาโลกสญฺญาย นตฺถิ กิญฺจิ วชฺชนฺติ อนวชฺชพลํ, อุทฺธจฺจสฺส ปหีนตฺตา อวิกฺเขเป นตฺถิ กิญฺจิ วชฺชนฺติ อนวชฺชพลํ, วิจิกิจฺฉาย ปหีนตฺตา ธมฺม\-ววตฺถาเน นตฺถิ กิญฺจิ วชฺชนฺติ อนวชฺชพลํ, อวิชฺชาย ปหีนตฺตา ญาเณ นตฺถิ กิญฺจิ วชฺชนฺติ อนวชฺชพลํ, อรติยา ปหีนตฺตา ปาโมชฺเช นตฺถิ กิญฺจิ วชฺชนฺติ อนวชฺชพลํ, นีวรณานํ ปหีนตฺตา ปฐมชฺฌาเน นตฺถิ กิญฺจิ วชฺชนฺติ อนวชฺชพลํ ฯเปฯ สพฺพกิเลสานํ ปหีนตฺตา อรหตฺตมคฺเค นตฺถิ กิญฺจิ วชฺชนฺติ อนวชฺชพลํ, อิทํ อนวชฺชพลํ. (10)}

\prose{\csromanfolio{353}กตมํ สงฺคหพลํ? กามจฺฉนฺทํ ปชหนฺโต เนกฺขมฺม\-วเสน จิตฺตํ สงฺคณฺหาตีติ สงฺคหพลํ, พฺยาปาทํ ปชหนฺโต อพฺยาปาทวเสน จิตฺตํ สงฺคณฺหาตีติ สงฺคหพลํ, ถินมิทฺธํ ปชหนฺโต อาโลกสญฺญา\-วเสน จิตฺตํ สงฺคณฺหาตีติ สงฺคหพลํ ฯเปฯ สพฺพกิเลเส ปชหนฺโต อรหตฺตมคฺค\-วเสน จิตฺตํ สงฺคณฺหาตีติ สงฺคหพลํ, อิทํ สงฺคหพลํ. (11)}

\prose{กตมํ ขนฺติพลํ? กามจฺฉนฺทสฺส ปหีนตฺตา เนกฺขมฺมํ ขมตีติ ขนฺติพลํ, พฺยาปาทสฺส ปหีนตฺตา อพฺยาปาโท ขมตีติ ขนฺติพลํ, ถินมิทฺธสฺส ปหีนตฺตา อาโลกสญฺญา ขมตีติ ขนฺติพลํ, อุทฺธจฺจสฺส ปหีนตฺตา อวิกฺเขโป ขมตีติ ขนฺติพลํ, วิจิกิจฺฉาย ปหีนตฺตา ธมฺม\-ววตฺถานํ ขมตีติ ขนฺติพลํ, อวิชฺชาย ปหีนตฺตา ญาณํ ขมตีติ ขนฺติพลํ, อรติยา ปหีนตฺตา ปาโมชฺชํ ขมตีติ ขนฺติพลํ, นีวรณานํ ปหีนตฺตา ปฐมํ ฌานํ ขมตีติ ขนฺติพลํ ฯเปฯ สพฺพกิเลสานํ ปหีนตฺตา อรหตฺตมคฺโค ขมตีติ ขนฺติพลํ, อิทํ ขนฺติพลํ. (12)}

\prose{กตมํ ปญฺญตฺติพลํ? กามจฺฉนฺทํ ปชหนฺโต เนกฺขมฺม\-วเสน จิตฺตํ ปญฺญาเปตีติ ปญฺญตฺติพลํ, พฺยาปาทํ ปชหนฺโต อพฺยาปาทวเสน จิตฺตํ ปญฺญาเปตีติ ปญฺญตฺติพลํ, ถินมิทฺธํ ปชหนฺโต อาโลกสญฺญา\-วเสน จิตฺตํ ปญฺญาเปตีติ ปญฺญตฺติพลํ ฯเปฯ สพฺพกิเลเส ปชหนฺโต อรหตฺตมคฺค\-วเสน จิตฺตํ ปญฺญาเปตีติ ปญฺญตฺติพลํ, อิทํ ปญฺญตฺติพลํ. (13)}

\prose{กตมํ นิชฺฌตฺติพลํ? กามจฺฉนฺทํ ปชหนฺโต เนกฺขมฺม\-วเสน จิตฺตํ นิชฺฌาเปตีติ นิชฺฌตฺติพลํ, พฺยาปาทํ ปชหนฺโต อพฺยาปาทวเสน จิตฺตํ นิชฺฌาเปตีติ นิชฺฌตฺติพลํ, ถินมิทฺธํ ปชหนฺโต อาโลกสญฺญา\-วเสน จิตฺตํ นิชฺฌาเปตีติ นิชฺฌตฺติพลํ ฯเปฯ สพฺพกิเลเส ปชหนฺโต อรหตฺตมคฺค\-วเสน จิตฺตํ นิชฺฌาเปตีติ นิชฺฌตฺติพลํ, อิทํ นิชฺฌตฺติพลํ. (14)}

\prose{กตมํ อิสฺสริยพลํ? กามจฺฉนฺทํ ปชหนฺโต เนกฺขมฺม\-วเสน จิตฺตํ วสํ วตฺเตตีติ อิสฺสริยพลํ, พฺยาปาทํ ปชหนฺโต อพฺยาปาทวเสน จิตฺตํ วสํ วตฺเตตีติ อิสฺสริยพลํ, ถินมิทฺธํ ปชหนฺโต อาโลกสญฺญา\-วเสน จิตฺตํ วสํ วตฺเตตีติ อิสฺสริยพลํ ฯเปฯ สพฺพกิเลเส ปชหนฺโต อรหตฺตมคฺค\-วเสน จิตฺตํ วสํ วตฺเตตีติ อิสฺสริยพลํ, อิทํ อิสฺสริยพลํ. (15)}

\prose{กตมํ อธิฏฺฐานพลํ? กามจฺฉนฺทํ ปชหนฺโต เนกฺขมฺม\-วเสน จิตฺตํ อธิฏฺฐาตีติ อธิฏฺฐานพลํ, พฺยาปาทํ ปชหนฺโต อพฺยาปาทวเสน จิตฺตํ}

\prose{\csromanfolio{354}อธิฏฺฐาตีติ อธิฏฺฐานพลํ, ถินมิทฺธํ ปชหนฺโต อาโลกสญฺญา\-วเสน จิตฺตํ อธิฏฺฐาตีติ อธิฏฺฐานพลํ ฯเปฯ สพฺพกิเลเส ปชหนฺโต อรหตฺตมคฺค\-วเสน จิตฺตํ อธิฏฺฐาตีติ อธิฏฺฐานพลํ, อิทํ อธิฏฺฐานพลํ. (16)}

\prose{กตมํ สมถพลํ? เนกฺขมฺม\-วเสน จิตฺตสฺส เอกคฺคตา อวิกฺเขโป สมถพลํ, อพฺยาปาทวเสน จิตฺตสฺส เอกคฺคตา อวิกฺเขโป สมถพลํ. อาโลกสญฺญา\-วเสน จิตฺตสฺส เอกคฺคตา อวิกฺเขโป สมถพลํ ฯเปฯ ปฏินิสฺสคฺคา\-นุปสฺสี อสฺสาสวเสน จิตฺตสฺส เอกคฺคตา อวิกฺเขโป สมถพลํ, ปฏินิสฺสคฺคา\-นุปสฺสี ปสฺสาสวเสน จิตฺตสฺส เอกคฺคตา อวิกฺเขโป สมถพลํ.}

\prose{สมถพลนฺติ เกนฏฺเฐน สมถพลํ? ปฐเมน ฌาเนน นีวรเณ น กมฺปตีติ สมถพลํ, ทุติเยน ฌาเนน วิตกฺกวิจาเร น กมฺปตีติ สมถพลํ, ตติยน ฌาเนน ปีติยา น กมฺปตีติ สมถพลํ, จตุตฺเถน ฌาเนน สุขทุกฺเข น กมฺปตีติ สมถพลํ, อากาสา\-นญฺจา\-ยตนสมาปตฺติยา รูปสญฺญาย ปฏิฆสญฺญาย นานตฺตสญฺญาย น กมฺปตีติ สมถพลํ, วิญฺญาณญฺจายตน\-สมาปตฺติยา อากาสา\-นญฺจา\-ยตนสญฺญาย น กมฺปตีติ สมถพลํ, อากิญฺจญฺญายตน\-สมาปตฺติยา วิญฺญาณญฺจา\-ยตนสญฺญาย น กมฺปตีติ สมถพลํ, เนว\-สญฺญา\-นา\-สญฺญา\-ยตนสมาปตฺติยา อากิญฺจญฺญายตน\-สญฺญาย น กมฺปตีติ สมถพลํ. อุทฺธจฺเจ จ อุทฺธจฺจสหคต\-กิเลเส จ ขนฺเธ จ น กมฺปติ น จลติ น เวธตีติ สมถพลํ, อิทํ สมถพลํ. (17)}

\prose{กตมํ วิปสฺสนาพลํ? อนิจฺจา\-นุปสฺสนา วิปสฺสนาพลํ, ทุกฺขา\-นุปสฺสนา วิปสฺสนาพลํ ฯเปฯ ปฏินิสฺสคฺคา\-นุปสฺสนา วิปสฺสนาพลํ. รูเป อนิจฺจา\-นุปสฺสนา วิปสฺสนาพลํ, รูเป ทุกฺขา\-นุปสฺสนา วิปสฺสนาพลํ ฯเปฯ รูเป ปฏินิสฺสคฺคา\-นุปสฺสนา วิปสฺสนาพลํ, เวทนาย ฯเปฯ สญฺญาย. สงฺขาเรสุ. วิญฺญาเณ จกฺขุสฺมิํ ฯเปฯ ชรามรเณ อนิจฺจา\-นุปสฺสนา วิปสฺสนาพลํ, ชรามรเณ ทุกฺขา\-นุปสฺสนา วิปสฺสนาพลํ ฯเปฯ ชรามรเณ ปฏินิสฺสคฺคา\-นุปสฺสนา วิปสฺสนาพลํ.  วิปสฺสนา\-พลนฺติ เกนฏฺเฐน วิปสฺสนาพลํ? อนิจฺจา\-นุปสฺสนาย นิจฺจสญฺญาย น กมฺปตีติ วิปสฺสนาพลํ, ทุกฺขา\-นุปสฺสนาย สุขสญฺญาย น กมฺปตีติ วิปสฺสนาพลํ, อนตฺตา\-นุปสฺสนาย อตฺตสญฺญาย น กมฺปตีติ วิปสฺสนาพลํ, นิพฺพิทา\-นุปสฺสนาย นนฺทิยา น กมฺปตีติ วิปสฺสนาพลํ, วิราคา\-นุปสฺสนาย ราเค น กมฺปตีติ วิปสฺสนาพลํ,}

\prose{\csromanfolio{355}นิโรธา\-นุปสฺสนาย สมุทเย น กมฺปตีติ วิปสฺสนาพลํ, ปฏินิสฺสคฺคา\-นุปสฺสนาย อาทาเน น กมฺปตีติ วิปสฺสนาพลํ, อวิชฺชาย จ อวิชฺชา\-สหคต\-กิเลเส จ ขนฺเธ จ น กมฺปติ น จลติ น เวธตีติ วิปสฺสนาพลํ, อิทํ วิปสฺสนาพลํ. (18)}

\prose{กตมานิ ทส เสขพลานิ, ทส อเสขพลานิ? สมฺมาทิฏฺฐิํ\csromansharedfootnote{20682a42d835844b}{สมฺมาทิฏฺฐิ (ก) เอวมีทิเสสุ นวสุ ปเทสุ ทุติยนฺตวจเนน. ที 3. 225 ปิฏฺเฐ ปสฺสิตพฺพา.} สิกฺขตีติ เสขพลํ, ตตฺถ สิกฺขิตตฺตา อเสขพลํ. สมฺมาสงฺกปฺปํ สิกฺขตีติ เสขพลํ, ตตฺถ สิกฺขิตตฺตา อเสขพลํ. สมฺมาวาจํ ฯเปฯ. สมฺมากมฺมนฺตํ. สมฺมาอาชีวํ. สมฺมาวายามํ. สมฺมาสติํ. สมฺมาสมาธิํ. สมฺมาญาณํ ฯเปฯ. สมฺมาวิมุตฺติํ สสิกฺขตีติ เสขพลํ, ตตฺถ สิกฺขิตตฺตา อเสขพลํ. อิมานิ ทส เสขพลานิ, ทส อเสขพลานิ. \mbox{(20, 38)}}

\prose{กตมานิ ทส ขีณาสว\-พลานิ\csromansharedfootnote{d1f1451d3029f201}{อํ 3. 397 ปิฏฺเฐ ปสฺสิตพฺพา.}? อิธ ขีณาสวสฺส ภิกฺขุโน อนิจฺจโต สพฺเพ สงฺขารา ยถาภูตํ สมฺมปฺปญฺญาย สุทิฏฺฐา โหนฺติ. ยมฺปิ ขีณาสวสฺส ภิกฺขุโน อนิจฺจโต สพฺเพ สงฺขารา ยถาภูตํ สมฺมปฺปญฺญาย สุทิฏฺฐา โหนฺติ, อิทมฺปิ ขีณาสวสฺส ภิกฺขุโน พลํ โหติ. ยํ พลํ อาคมฺม ขีณาสโว ภิกฺขุ อาสวานํ ขยํ ปฏิชานาติ “ขีณา เม อาสวา”ติ. (1)}

\prose{ปุน จปรํ ขีณาสวสฺส ภิกฺขุโน องฺคารกาสูปมา กามา ยถาภูตํ สมฺมปฺปญฺญาย สุทิฏฺฐา โหนฺติ. ยมฺปิ ขีณาสวสฺส ภิกฺขุโน องฺคารกาสูปมา กามา ยถาภูตํ สมฺมปฺปญฺญาย สุทิฏฺฐา โหนฺติ, อิทมฺปิ ขีณาสวสฺส ภิกฺขุโน พลํ โหติ. ยํ พลํ อาคมฺม ขีณาสโว ภิกฺขุ อาสวานํ ขยํ ปฏิชานาติ “ขีณา เม อาสวา”ติ. (2)}

\prose{ปุน จปรํ ขีณาสวสฺส ภิกฺขุโน วิเวกนินฺนํ จิตฺตํ โหติ วิเวกโปณ วิเวก\-ปพฺภารํ วิเวกฏฺฐํ เนกฺขมฺมา\-ภิรตํ พฺยนฺตีภูตํ\csromansharedfootnote{d0a777e71703a2c9}{พฺยนฺติภูตํ (ก)} สพฺพโส อาสว\-ฏฺฐานิเยหิ ธมฺเมหิ. ยมฺปิ ขีณาสวสฺส ภิกฺขุโน วิเวกนินฺนํ จิตฺตํ โหติ วิเวกโปณํ วิเวก\-ปพฺภารํ วิเวกฏฺฐํ เนกฺขมฺมา\-ภิรตํ พฺยนฺตีภูตํ สพฺพโส อาสว\-ฏฺฐานิเยหิ ธมฺเมหิ, อิทมฺปิ ขีณาสวสฺส ภิกฺขุโน พลํ โหติ. ยํ พลํ อาคมฺม ขีณาสโว ภิกฺขุ อาสวานํ ขยํ ปฏิชานาติ “ขีณา เม อาสวา”ติ. (3)}

\prose{\csromanfolio{356}ปุน จปรํ ขีณาสวสฺส ภิกฺขุโน จตฺตาโร สติปฏฺฐานา ภาวิตา โหนฺติ สุภาวิตา. ยมฺปิ ขีณาสวสฺส ภิกฺขุโน จตฺตาโร สติปฏฺฐานา ภาวิตา โหนฺติ สุภาวิตา, อิทมฺปิ ขีณาสวสฺส ภิกฺขุโน พลํ โหติ. ยํ พลํ อาคมฺม ขีณาสโว ภิกฺขุ อาสวานํ ขยํ ปฏิชานาติ “ขีณา เม อาสวา”ติ. (4)}

\prose{ปุน จปรํ ขีณาสวสฺส ภิกฺขุโน จตฺตาโร สมฺมปฺปธานา ภาวิตา โหนฺติ สุภาวิตา ฯเปฯ. จตฺตาโร อิทฺธิปาทา ภาวิตา โหนฺติ สุภาวิตา. ปญฺจินฺทฺริยานิ ภาวิตานิ โหนฺติ สุภาวิตานิ. ปญฺจ พลานิ ภาวิตานิ โหนฺติ สุภาวิตานิ. สตฺต โพชฺฌงฺคา ภาวิตา โหนฺติ สุภาวิตา ฯเปฯ. อริโย อฏฺฐงฺคิโก มคฺโค ภาวิโต โหติ สุภาวิโต. ยมฺปิ ขีณาสวสฺส ภิกฺขุโน อฏฺฐงฺคิโก มคฺโค ภาวิโต โหติ สุภาวิโต, อิทมฺปิ ขีณาสวสฺส ภิกฺขุโน พลํ โหติ. ยํ พลํ อาคมฺม ขีณาสโว ภิกฺขุ อาสวานํ ขยํ ปฏิชานาติ “ขีณา เม อาสวา”ติ. อิมานิ ทส ขีณาสว\-พลานิ. (5-10, 48)}

\prose{กตมานิ ทส อิทฺธิพลานิ? อธิฏฺฐานา อิทฺธิ, วิกุพฺพนา อิทฺธิ, มโนมยา อิทฺธิ, ญาณวิปฺผารา อิทฺธิ, สมาธิ\-วิปฺผารา อิทฺธิ, อริยา อิทฺธิ, กมฺมวิปากชา อิทฺธิ, ปุญฺญวโต อิทฺธิ, วิชฺชามยา อิทฺธิ, ตตฺถ ตตฺถ สมฺมา ปโยคปฺปจฺจยา อิชฺฌนฏฺเฐน อิทฺธิ. อิมานิ ทส อิทฺธิพลานิ. \mbox{(10, 58)}}

\prose{กตมานิ ทส ตถาคต\-พลานิ? อิธ ตถาคโต ฐานญฺจ ฐานโต อฏฺฐานญฺจ อฏฺฐานโต ยถาภูตํ ปชานาติ, ยมฺปิ ตถาคโต ฐานญฺจ ฐานโต อฏฺฐานญฺจ อฏฺฐานโต ยถาภูตํ ปชานาติ, อิทมฺปิ ตถาคตสฺส ตถาคตพลํ โหติ, ยํ พลํ อาคมฺม ตถาคโต อาสภํ ฐานํ ปฏิชานาติ, ปริสาสุ สีหนาทํ นทติ, พฺรหฺมจกฺกํ ปวตฺเตติ. (1)}

\prose{ปุน จปรํ ตถาคโต อตีตา\-นาคต\-ปจฺจุปฺปนฺนานํ กมฺม\-สมาทานานํ ฐานโส เหตุโส วิปากํ ยถาภูตํ ปชานาติ, ยมฺปิ ตถาคโต อตีตา\-นาคต\-ปจฺจุปฺปนฺนานํ กมฺม\-สมาทานานํ ฐานโส เหตุโส วิปากํ ยถาภูตํ ปชานาติ, อิทมฺปิ ตถาคตสฺส ตถาคตพลํ โหติ. ยํ พลํ อาคมฺม ตถาคโต อาสภํ ฐานํ ปฏิชานาติ, ปริสาสุ สีหนาทํ นทติ, พฺรหฺมจกฺกํ ปวตฺเตติ. (2)}

\prose{\csromanfolio{357}ปุน จปรํ ตถาคโต สพฺพตฺถ\-คามินิํ ปฏิปทํ\csromansharedfootnote{bd5d260282307fd9}{สพฺพตฺถคามินีปฏิปทํ (สฺยา) ม 1. 100 ปิฏฺเฐ ปสฺสิตพฺพา.} ยถาภูตํ ปชานาติ, ยมฺปิ ตถาคโต สพฺพตฺถ\-คามินิํ ปฏิปทํ ยถาภูตํ ปชานาติ, อิทมฺปิ ตถาคตสฺส ตถาคตพลํ โหติ, ยํ พลํ อาคมฺม ตถาคโต อาสภํ ฐานํ ปฏิชานาติ, ปริสาสุ สีหนาทํ นทติ, พฺรหฺมจกฺกํ ปวตฺเตติ. (3)}

\prose{ปุน จปรํ ตถาคโต อเนกธาตุํ นานาธาตุํ โลกํ ยถาภูตํ ปชานาติ. ยมฺปิ ตถาคโต อเนกธาตุํ นานาธาตุํ โลกํ ยถาภูตํ ปชานาติ, อิทมฺปิ ตถาคตสฺส ฯเปฯ. (4)}

\prose{ปุน จปรํ ตถาคโต สตฺตานํ นานา\-ธิมุตฺติกตํ ยถาภูตํ ปชานาติ, ยมฺปิ ตถาคโต สตฺตานํ นานา\-ธิมุตฺติกตํ ยถาภูตํ ปชานาติ, อิทมฺปิ ตถาคตสฺส ฯเปฯ. (5)}

\prose{ปุน จปรํ ตถาคโต ปรสตฺตานํ ปรปุคฺคลานํ อินฺทฺริย\-ปโรปริยตฺตํ ยถาภูตํ ปชานาติ, ยมฺปิ ตถาคโต ปรสตฺตานํ ปรปุคฺคลานํ อินฺทฺริย\-ปโรปริยตฺตํ ยถาภูตํ ปชานาติ, อิทมฺปิ ตถาคตสฺส ฯเปฯ. (6)}

\prose{ปุน จปรํ ตถาคโต ฌาน\-วิโมกฺข\-สมาธิ\-สมาปตฺตีนํ สํกิเลสํ โวทานํ วุฏฺฐานํ ยถาภูตํ ปชานาติ, ยมฺปิ ตถาคโต ฌาน\-วิโมกฺข\-สมาธิ\-สมาปตฺตีนํ สํกิเลสํ โวทานํ วุฏฺฐานํ ยถาภูตํ ปชานาติ, อิทมฺปิ ตถาคตสฺส ฯเปฯ. (7)}

\prose{ปุน จปรํ ตถาคโต อเนกวิหิตํ ปุพฺเพนิวาสํ อนุสฺสรติ, เสยฺยถิทํ, เอกมฺปิ ชาติํ เทฺวปิ ชาติโย ฯเปฯ อิติ สาการํ สอุทฺเทสํ อเนกวิหิตํ ปุพฺเพนิวาสํ อนุสฺสรติ, ยมฺปิ ตถาคโต อเนกวิหิตํ ปุพฺเพนิวาสํ อนุสฺสรติ. เสยฺยถิทํ, เอกมฺปิ ชาติํ เทฺวปิ ชาติโย ฯเปฯ อิทมฺปิ ตถาคตสฺส ฯเปฯ. (8)}

\prose{ปุน จปรํ ตถาคโต ทิพฺเพน จกฺขุนา วิสุทฺเธน อติกฺกนฺต\-มานุสเกน สตฺเต ปสฺสติ จวมาเน อุปปชฺชมาเน ฯเปฯ ยมฺปิ ตถาคโต ทิพฺเพน จกฺขุนา วิสุทฺเธน อติกฺกนฺต\-มานุสเกน สตฺเต ปสฺสติ จวมาเน อุปปชฺชมาเน ฯเปฯ อิทมฺปิ ตถาคตสฺส ฯเปฯ. (9)}

\csromanmatikaanchor{mtk.117}
\csromantocmarkref{section}{10. สุญฺญกถา}{mtk.117}
\bookmark[dest={mtk.117},level=1]{10. สุญฺญกถา}
\prose{\csromanfolio{358}ปุน จปรํ ตถาคโต อาสวานํ ขยา อนาสวํ เจโตวิมุตฺติํ ปญฺญาวิมุตฺติํ ทิฏฺเฐว ธมฺเม สยํ อภิญฺญา สจฺฉิกตฺวา อุปสมฺปชฺช วิหรติ, ยมฺปิ ตถาคโต อาสวานํ ขยา อนาสวํ เจโตวิมุตฺติํ ปญฺญาวิมุตฺติํ ทิฏฺเฐว ธมฺเม สยํ อภิญฺญา สจฺฉิกตฺวา อุปสมฺปชฺช วิหรติ, อิทมฺปิ ตถาคตสฺส ตถาคตพลํ โหติ, ยํ พลํ อาคมฺม ตถาคโต อาสภํ ฐานํ ปฏิชานาติ, ปริสาสุ สีหนาทํ นทติ, พฺรหฺมจกฺกํ ปวตฺเตติ. อิมานิ ทส ตถาคต\-พลานิ. \mbox{(10, 68)}}

\proseitem{45}{เกนฏฺเฐน สทฺธาพลํ. เกนฏฺเฐน วีริยพลํ. เกนฏฺเฐน สติพลํ. เกนฏฺเฐน สมาธิพลํ. เกนฏฺเฐน ปญฺญาพลํ. เกนฏฺเฐน หิริพลํ. เกนฏฺเฐน โอตฺตปฺปพลํ. เกนฏฺเฐน ปฏิสงฺขานพลํ ฯเปฯ. เกนฏฺเฐน ตถาคตพลํ?}

\prosewithcloserruled{อสฺสทฺธิเย อกมฺปิยฏฺเฐน สทฺธาพลํ. โกสชฺเช อกมฺปิยฏฺเฐน วีริยพลํ. ปมาเท อกมฺปิยฏฺเฐน สติพลํ. อุทฺธจฺเจ อกมฺปิยฏฺเฐน สมาธิพลํ. อวิชฺชาย อกมฺปิยฏฺเฐน ปญฺญาพลํ. หิรียติ ปาปเก อกุสเล ธมฺเมติ หิริพลํ. โอตฺตปฺปติ ปาปเก อกุสเล ธมฺเมติ โอตฺตปฺปพลํ. ญาเณน กิเลเส ปฏิสงฺขาตีติ ปฏิสงฺขานพลํ. ตตฺถ ชาตา ธมฺมา เอกรสา โหนฺตีติ ภาวนาพลํ, ตตฺถ นตฺถิ กิญฺจิ วชฺชนฺติ อนวชฺชพลํ. เตน จิตฺตํ สงฺคณฺหาตีติ สงฺคหพลํ. ตํ ตสฺส ขมตีติ\csromansharedfootnote{bb685b3bb6e9eec6}{ตํ ขมตีติ (ก)} ขนฺติพลํ. เตน จิตฺตํ ปญฺญาเปตีติ ปญฺญตฺติพลํ, เตน จิตฺตํ นิชฺฌาเปตีติ นิชฺฌตฺติพลํ, เตน จิตฺตํ วสํ วตฺเตตีติ อิสฺสริยพลํ, เตน จิตฺตํ อธิฏฺฐาตีติ อธิฏฺฐานพลํ, เตน จิตฺตํ เอกคฺคนฺติ สมถพลํ, ตตฺถ ชาเต ธมฺเม อนุปสฺสตีติ วิปสฺสนาพลํ, ตตฺถ สิกฺขตีติ เสขพลํ, ตตฺถ สิกฺขิตตฺตา อเสขพลํ, เตน อาสวา ขีณาติ ขีณาสวพลํ. ตสฺส อิชฺฌตีติ อิทฺธิพลํ. อปฺปเมยฺย\-ฏฺเฐน ตถาคต\-พลนฺติ.}{พลกถา นิฏฺฐิตา.}
\proseitem{46}{เอวํ เม สุตํ\csromandash{}เอกํ สมยํ ภควา สาวตฺถิยํ วิหรติ เชตวเน อนาถ\-ปิณฺฑิกสฺส อาราเม. อถ โข อายสฺมา อานนฺโท เยน}

\prose{\csromanfolio{359}ภควา เตนุปสงฺกมิ, อุปสงฺกมิตฺวา ภควนฺตํ อภิวาเทตฺวา เอกมนฺตํ นิสีทิ, เอกมนฺตํ นิสินฺโน โข อายสฺมา อานนฺโท ภควนฺตํ เอตทโวจ\csromandash{}}

\prose{“สุญฺโญ โลโก สุญฺโญ โลโก”ติ ภนฺเต วุจฺจติ, กิตฺตาวตา นุ โข ภนฺเต “สุญฺโญ โลโก”ติ วุจฺจตีติ? ยสฺมา โข อานนฺท สุญฺญํ อตฺเตน วา อตฺตนิเยน วา, ตสฺมา “สุญฺโญ โลโก”ติ วุจฺจติ. กิญฺจานนฺท\csromansharedfootnote{3c4c1773843d6ad2}{กิญฺจ อานนฺท (สํ 2. 280 ปิฏฺเฐ)} สุญฺญํ อตฺเตน วา อตฺตนิเยน วา? จกฺขุ โข อานนฺท สุญฺญํ อตฺเตน วา อตฺตนิเยน วา, รูปา สุญฺญา อตฺเตน วา อตฺตนิเยน วา, จกฺขุวิญฺญาณํ สุญฺญํ อตฺเตน วา อตฺตนิเยน วา, จกฺขุ\-สมฺผสฺโส สุญฺโญ อตฺเตน วา อตฺตนิเยน วา, ยมฺปิทํ\csromansharedfootnote{d449944685e2e4c5}{ยทิทํ (ก) สํ 2. 280 ปิฏฺเฐ ปสฺสิตพฺพา. 3. สุญฺญํ สุญฺญํ (สฺยา)} จกฺขุ\-สมฺผสฺส\-ปจฺจยา อุปฺปชฺชติ เวทยิตํ สุขํ วา ทุกฺขํ วา อทุกฺขมสุขํ วา, ตมฺปิ สุญฺญํ อตฺเตน วา อตฺตนิเยน วา.}

\prose{โสตํ สุญฺญํ, สทฺทา สุญฺญา. ฆานํ สุญฺญํ, คนฺธา สุญฺญา. ชิวฺหา สุญฺญา, รสา สุญฺญา. กาโย สุญฺโญ, โผฏฺฐพฺพา สุญฺญา. มโน สุญฺโญ อตฺเตน วา อตฺตนิเยน วา. ธมฺมา สุญฺญา อตฺเตน วา อตฺตนิเยน วา. มโนวิญฺญาณํ สุญฺญํ อตฺเตน วา อตฺตนิเยน วา, มโนสมฺผสฺโส สุญฺโญ อตฺเตน วา อตฺตนิเยน วา, ยมฺปิทํ มโน\-สมฺผสฺส\-ปจฺจยา อุปฺปชฺชติ เวทยิตํ สุขํ วา ทุกฺขํ วา อทุกฺขมสุขํ วา, ตมฺปิ สุญฺญํ อตฺเตน วา อตฺตนิเยน วา. ยสฺมา โข อานนฺท สุญฺญํ อตฺเตน วา อตฺตนิเยน วา, ตสฺมา “สุญฺโญ โลโก”ติ วุจฺจตีติ.}

\csromanheader{2}{1. \textbf{มาติกา}}
\proseitem{47}{สุญฺญสุญฺญํ\csromansharedfootnote{7685312f99b9ff84}{ปริโยคาหนสุญฺญํ (สฺยา)}, สงฺขารสุญฺญํ, วิปริณาม\-สุญฺญํ, อคฺคสุญฺญํ, ลกฺขณสุญฺญํ, วิกฺขมฺภน\-สุญฺญํ, ตทงฺคสุญฺญํ, สมุจฺเฉท\-สุญฺญํ, ปฏิปฺปสฺสทฺธิ\-สุญฺญํ, นิสฺสรณ\-สุญฺญํ, อชฺฌตฺตสุญฺญํ, พหิทฺธาสุญฺญํ, ทุภโตสุญฺญํ, สภาคสุญฺญํ, วิสภาค\-สุญฺญํ, เอสนาสุญฺญํ, ปริคฺคห\-สุญฺญํ, ปฏิลาภ\-สุญฺญํ, ปฏิเวธ\-สุญฺญํ, เอกตฺตสุญฺญํ, นานตฺตสุญฺญํ, ขนฺติสุญฺญํ, อธิฏฺฐาน\-สุญฺญํ, ปริโยคาห\-ณสุญฺญํ\csromansharedfootnote{77e82cf7c384ec14}{สุญฺญํ สุญฺญํ (สฺยา)}, สมฺปชานสฺส ปวตฺต\-ปริยาทานํ สพฺพ\-สุญฺญตานํ ปรมตฺถ\-สุญฺญํ. (25)}

\csromanheader{2}{2. \textbf{นิทฺเทส}}
\proseitem{48}{\csromanfolio{360}กตมํ สุญฺญสุญฺญํ? จกฺขุ สุญฺญํ อตฺเตน วา อตฺตนิเยน วา นิจฺเจน วา ธุเวน วา สสฺสเตน วา อวิปริณาม\-ธมฺเมน วา.  โสตํ สุญฺญํ ฯเปฯ. ฆานํ สุญฺญํ. ชิวฺหา สุญฺญา. กาโย สุญฺโญ, มโน สุญฺโญ อตฺเตน วา อตฺตนิเยน วา นิจฺเจน วา ธุเวน วา สสฺสเตน วา อวิปริณาม\-ธมฺเมน วา. อิทํ สุญฺญสุญฺญํ. (1)}

\prose{กตมํ สงฺขารสุญฺญํ? ตโย สงฺขารา ปุญฺญา\-ภิสงฺขาโร อปุญฺญา\-ภิสงฺขาโร อาเนญฺชา\-ภิสงฺขาโร. ปุญฺญา\-ภิสงฺขาโร อปุญฺญา\-ภิสงฺขาเรน จ อาเนญฺชา\-ภิสงฺขาเรน จ สุญฺโญ, อปุญฺญา\-ภิสงฺขาโร ปุญฺญา\-ภิสงฺขาเรน จ อาเนญฺชา\-ภิสงฺขาเรน จ สุญฺโญ, อาเนญฺชา\-ภิสงฺขาโร ปุญฺญา\-ภิสงฺขาเรน จ อปุญฺญา\-ภิสงฺขาเรน จ สุญฺโญ. อิเม ตโย สงฺขารา.}

\prose{อปเรปิ ตโย สงฺขารา กายสงฺขาโร วจีสงฺขาโร จิตฺตสงฺขาโร. กายสงฺขาโร วจีสงฺขาเรน จ จิตฺต\-สงฺขาเรน จ สุญฺโญ, วจีสงฺขาโร กายสงฺขาเรน จ จิตฺต\-สงฺขาเรน จ สุญฺโญ, จิตฺตสงฺขาโร กายสงฺขาเรน จ วจีสงฺขาเรน จ สุญฺโญ. อิเม ตโย สงฺขารา.}

\prose{อปเรปิ ตโย สงฺขารา อตีตา สงฺขารา อนาคตา สงฺขารา ปจฺจุปฺปนฺนา สงฺขารา. อตีตา สงฺขารา อนาคเตหิ จ ปจฺจุปฺปนฺเนหิ จ สงฺขาเรหิ สุญฺญา, อนาคตา สงฺขารา อตีเตหิ จ ปจฺจุปฺปนฺเนหิ จ สงฺขาเรหิ สุญฺญา, ปจฺจุปฺปนฺนา สงฺขารา อตีเตหิ จ อนาคเตหิ จ สงฺขาเรหิ สุญฺญา. อิเม ตโย สงฺขารา. อิทํ สงฺขารสุญฺญํ. (2)}

\prose{กตมํ วิปริณาม\-สุญฺญํ? ชาตํ รูปํ สภาเวน สุญฺญํ, วิคตํ รูปํ วิปริณต\-ญฺเจว สุญฺญญฺจ, ชาตา เวทนา สภาเวน สุญฺญา, วิคตา เวทนา วิปริณตา เจว สุญฺญา จ ฯเปฯ ชาตา สญฺญา. ชาตา สงฺขารา. ชาตํ วิญฺญาณํ. ชาตํ จกฺขุ ฯเปฯ ชาโต ภโว สภาเวน สุญฺโญ, วิคโต ภโว วิปริณโต เจว สุญฺโญ จ. อิทํ วิปริณาม\-สุญฺญํ. (3)}

\prose{กตมํ อคฺคสุญฺญํ? อคฺคเมตํ ปทํ, เสฏฺฐเมตํ ปทํ, วิสิฏฺฐเมตํ\csromansharedfootnote{42634fc3a2e592e0}{วิเสฏฺฐเมตํ (ก)} ปทํ, ยทิทํ สพฺพ\-สงฺขาร\-สมโถ สพฺพู\-ปธิ\-ปฏินิสฺสคฺโค ตณฺหกฺขโย วิราโค นิโรโธ นิพฺพานํ. อิทํ อคฺคสุญฺญํ. (4)}

\prose{\csromanfolio{361}กตมํ ลกฺขณสุญฺญํ? เทฺว ลกฺขณานิ พาลลกฺขณญฺจ ปณฺฑิตลกฺขณญฺจ. พาลลกฺขณํ ปณฺฑิต\-ลกฺขเณน สุญฺญํ, ปณฺฑิต\-ลกฺขณํ พาลลกฺขเณน สุญฺญํ.  ตีณิ ลกฺขณานิ อุปฺปาท\-ลกฺขณํ วยลกฺขณํ ฐิตญฺญถตฺต\-ลกฺขณํ. อุปฺปาท\-ลกฺขณํ วยลกฺขเณน จ ฐิตญฺญถตฺต\-ลกฺขเณน จ สุญฺญํ, วยลกฺขณํ อุปฺปาท\-ลกฺขเณน จ ฐิตญฺญถตฺต\-ลกฺขเณน จ สุญฺญํ, ฐิตญฺญถตฺต\-ลกฺขณํ อุปฺปาท\-ลกฺขเณน จ วยลกฺขเณน จ สุญฺญํ.}

\prose{รูปสฺส อุปฺปาท\-ลกฺขณํ วยลกฺขเณน จ ฐิตญฺญถตฺต\-ลกฺขเณน จ สุญฺญํ, รูปสฺส วยลกฺขณํ อุปฺปาท\-ลกฺขเณน จ ฐิตญฺญถตฺต\-ลกฺขเณน จ สุญฺญํ, รูปสฺส ฐิตญฺญถตฺต\-ลกฺขณํ อุปฺปาท\-ลกฺขเณน จ วยลกฺขเณน จ สุญฺญํ.  เวทนาย ฯเปฯ. สญฺญาย. สงฺขารานํ. วิญฺญาณสฺส. จกฺขุสฺส ฯเปฯ. ชรามรณสฺส อุปฺปาท\-ลกฺขณํ วยลกฺขเณน จ ฐิตญฺญถตฺต\-ลกฺขเณน จ สุญฺญํ, ชรามรณสฺส วยลกฺขณํ อุปฺปาท\-ลกฺขเณน จ ฐิตญฺญถตฺต\-ลกฺขเณน จ สุญฺญํ, ชรามรณสฺส ฐิตญฺญถตฺต\-ลกฺขณํ อุปฺปาท\-ลกฺขเณน จ วยลกฺขเณน จ สุญฺญํ. อิทํ ลกฺขณสุญฺญํ. (5)}

\prose{กตมํ วิกฺขมฺภน\-สุญฺญํ? เนกฺขมฺเมน กามจฺฉนฺโท วิกฺขมฺภิโต เจว สุญฺโญ จ, อพฺยาปาเทน พฺยาปาโท วิกฺขมฺภิโต เจว สุญฺโญ จ, อาโลกสญฺญาย ถินมิทฺธํ วิกฺขมฺภิต\-ญฺเจว สุญฺญญฺจ, อวิกฺเขเปน อุทฺธจฺจํ วิกฺขมฺภิต\-ญฺเจว สุญฺญญฺจ, ธมฺม\-ววตฺถาเนน วิจิกิจฺฉา วิกฺขมฺภิตา เจว สุญฺญา จ, ญาเณน อวิชฺชา วิกฺขมฺภิตา เจว สุญฺญา จ, ปาโมชฺเชน อรติ วิกฺขมฺภิตา เจว สุญฺญา จ, ปฐเมน ฌาเนน นีวรณา วิกฺขมฺภิตา เจว สุญฺญา จ ฯเปฯ อรหตฺต\-มคฺเคน สพฺพกิเลสา วิกฺขมฺภิตา เจว สุญฺญา จ. อิทํ วิกฺขมฺภน\-สุญฺญํ. (6)}

\prose{กตมํ ตทงฺคสุญฺญํ? เนกฺขมฺเมน กามจฺฉนฺโท ตทงฺคสุญฺโญ, อพฺยาปาเทน พฺยาปาโท ตทงฺคสุญฺโญ, อาโลกสญฺญาย ถินมิทฺธํ ตทงฺคสุญฺญํ, อวิกฺเขเปน อุทฺธจฺจํ ตทงฺคสุญฺญํ, ธมฺม\-ววตฺถาเนน วิจิกิจฺฉา ตทงฺคสุญฺญา, ญาเณน อวิชฺชา ตทงฺคสุญฺญา, ปาโมชฺเชน อรติ ตทงฺคสุญฺญา, ปฐเมน ฌาเนน นีวรณา ตทงฺคสุญฺญา ฯเปฯ วิวฏฺฏนา\-นุปสฺสนาย สญฺโญคา\-ภินิเวโส ตทงฺคสุญฺโญ. อิทํ ตทงฺคสุญฺญํ. (7)}

\prose{กตมํ สมุจฺเฉท\-สุญฺญํ? เนกฺขมฺเมน กามจฺฉนฺโท สมุจฺฉินฺโน เจว สุญฺโญ จ. อพฺยาปาเทน พฺยาปาโท สมุจฺฉินฺโน เจว สุญฺโญ จ,}

\prose{\csromanfolio{362}อาโลกสญฺญาย ถินมิทฺธํ สมุจฺฉินฺน\-ญฺเจว สุญฺญญฺจ, อวิกฺเขเปน อุทฺธจฺจํ สมุจฺฉินฺน\-ญฺเจว สุญฺญญฺจ, ธมฺมวตฺถาเนน วิจิกิจฺฉา สมุจฺฉินฺนา เจว สุญฺญา จ, ญาเณน อวิชฺชา สมุจฺฉินฺนา เจว สุญฺญา จ, ปาโมชฺเชน อรติ สมุจฺฉินฺนา เจว สุญฺญา จ, ปฐเมน ฌาเนน นีวรณา สมุจฺฉินฺนา เจว สุญฺญา จ ฯเปฯ อรหตฺต\-มคฺเคน สพฺพกิเลสา สมุจฺฉินฺนา เจว สุญฺญา จ. อิทํ สมุจฺเฉท\-สุญฺญํ. (8)}

\prose{กตมํ ปฏิปฺปสฺสทฺธิ\-สุญฺญํ? เนกฺขมฺเมน กามจฺฉนฺโท ปฏิปฺปสฺสทฺโธ เจว สุญฺโญ จ, อพฺยาปาเทน พฺยาปาโท ปฏิปฺปสฺสทฺโธ เจว สุญฺโญ จ. อาโลกสญฺญาย ถินมิทฺธํ ปฏิปฺปสฺสทฺธ\-ญฺเจว สุญฺญญฺจ, อวิกฺเขเปน อุทฺธจฺจํ ปฏิปฺปสฺสทฺธ\-ญฺเจว สุญฺญญฺจ. ธมฺมวตฺถาเนน วิจิกิจฺฉา ปฏิปฺปสฺสทฺธา เจว สุญฺญา จ, ญาเณน อวิชฺชา ปฏิปฺปสฺสทฺธา เจว สุญฺญา จ. ปาโมชฺเชน อรติ ปฏิปฺปสฺสทฺธา เจว สุญฺญา จ, ปฐเมน ฌาเนน นีวรณา ปฏิปฺปสฺสทฺธา เจว สุญฺญา จ ฯเปฯ อรหตฺต\-มคฺเคน สพฺพกิเลสา ปฏิปฺปสฺสทฺธา เจว สุญฺญา จ. อิทํ ปฏิปฺปสฺสทฺธิ\-สุญฺญํ. (9)}

\prose{กตมํ นิสฺสรณ\-สุญฺญํ? เนกฺขมฺเมน กามจฺฉนฺโท นิสฺสโฏ เจว สุญฺโญ จ, อพฺยาปาเทน พฺยาปาโท นิสฺสโฏ เจว สุญฺโญ จ, อาโลกสญฺญาย ถินมิทฺธํ นิสฺสฏญฺเจว สุญฺญญฺจ, อวิกฺเขเปน อุทฺธจฺจํ นิสฺสฏญฺเจว สุญฺญญฺจ, ธมฺม\-ววตฺถาเนน วิจิกิจฺฉา นิสฺสฏา เจว สุญฺญา จ, ญาเณน อวิชฺชา นิสฺสฏา เจว สุญฺญา จ, ปาโมชฺเชน อรติ นิสฺสฏา เจว สุญฺญา จ, ปฐเมน ฌาเนน นีวรณา นิสฺสฏา เจว สุญฺญา จ ฯเปฯ อรหตฺต\-มคฺเคน สพฺพกิเลสา นิสฺสฏา เจว สุญฺญา จ. อิทํ นิสฺสรณ\-สุญฺญํ. (10)}

\prose{กตมํ อชฺฌตฺตสุญฺญํ? อชฺฌตฺตํ จกฺขุํ สุญฺญํ อตฺเตน วา อตฺตนิเยน วา นิจฺเจน วา ธุเวน วา สสฺสเตน วา อวิปริณาม\-ธมฺเมน วา. อชฺฌตฺตํ โสตํ สุญฺญํ. อชฺฌตฺตํ ฆานํ สุญฺญํ. อชฺฌตฺตํ ชิวฺหา สุญฺญา. อชฺฌตฺตํ กาโย สุญฺโญ. อชฺฌตฺตํ มโน สุญฺโญ อตฺเตน วา อตฺตนิเยน วา นิจฺเจน วา ธุเวน วา สสฺสเตน วา อวิปริณาม\-ธมฺเมน วา. อิทํ อชฺฌตฺตสุญฺญํ. (11)}

\prose{กตมํ พหิทฺธาสุญฺญํ? พหิทฺธา รูปา สุญฺญา ฯเปฯ. พหิทฺธา ธมฺมา สุญฺญา อตฺเตน วา อตฺตนิเยน วา นิจฺเจน วา ธุเวน วา สสฺสเตน วา อวิปริณาม\-ธมฺเมน วา. อิทํ พหิทฺธาสุญฺญํ. (12)}

\prose{\csromanfolio{363}กตมํ ทุภโตสุญฺญํ? ยญฺจ อชฺฌตฺตํ จกฺขุ เย จ พหิทฺธา รูปา, อุภยเมตํ สุญฺญํ\csromansharedfootnote{9275e5766d455eec}{อุภยโต ตํ สุญฺญา (สฺยา)} อตฺเตน วา อตฺตนิเยน วา นิจฺเจน วา ธุเวน วา สสฺสเตน วา อวิปริณาม\-ธมฺเมน วา. ยญฺจ อชฺฌตฺตํ โสตํ, เย จ พหิทฺธา สทฺทา ฯเปฯ. ยญฺจ อชฺฌตฺตํ ฆานํ, เย จ พหิทฺธา คนฺธา. ยา จ อชฺฌตฺตํ ชิวฺหา, เย จ พหิทฺธา รสา. โย จ อชฺฌตฺตํ กาโย, เย จ พหิทฺธา โผฏฺฐพฺพา. โย จ อชฺฌตฺตํ มโน, เย จ พหิทฺธา ธมฺมา, อุภยเมตํ สุญฺญํ อตฺเตน วา อตฺตนิเยน วา นิจฺเจน วา ธุเวน วา สสฺสเตน วา อวิปริณาม\-ธมฺเมน วา. อิทํ ทุภโตสุญฺญํ. (13)}

\prose{กตมํ \textbf{สภาคสุญฺญํ?} ฉ อชฺฌตฺติกานิ อายตนานิ สภาคานิ เจว สุญฺญานิ จ, ฉ พาหิรานิ อายตนานิ สภาคานิ เจว สุญฺญานิ จ, ฉ วิญฺญาณกายา สภาคา เจว สุญฺญา จ, ฉ ผสฺสกายา สภาคา เจว สุญฺญา จ, ฉ เวทนากายา สภาคา เจว สุญฺญา จ, ฉ สญฺญากายา สภาคา เจว สุญฺญา จ, ฉ เจตนากายา สภาคา เจว สุญฺญา จ. อิทํ สภาคสุญฺญํ. (14)}

\prose{กตมํ วิสภาค\-สุญฺญํ? ฉ อชฺฌตฺติกานิ อายตนานิ ฉหิ พาหิเรหิ อายตเนหิ วิสภาคานิ เจว สุญฺญานิ จ, ฉ พาหิรานิ อายตนานิ ฉหิ วิญฺญาณกาเยหิ วิสภาคานิ เจว สุญฺญานิ จ, ฉ วิญฺญาณกายา ฉหิ ผสฺสกาเยหิ วิสภาคา เจว สุญฺญา จ, ฉ ผสฺสกายา ฉหิ เวทนากาเยหิ วิสภาคา เจว สุญฺญา จ, ฉ เวทนากายา ฉหิ สญฺญากาเยหิ วิสภาคา เจว สุญฺญา จ, ฉ สญฺญากายา ฉหิ เจตนากาเยหิ วิสภาคา เจว สุญฺญา จ. อิทํ วิสภาค\-สุญฺญํ. (15)}

\prose{กตมํ เอสนาสุญฺญํ? เนกฺขมฺเมสนา กามจฺฉนฺเทน สุญฺญา, อพฺยาปาเทสนา พฺยาปาเทน สุญฺญา, อาโลกสญฺเญสนา ถินมิทฺเธน สุญฺญา, อวิกฺเขเปสนา อุทฺธจฺเจน สุญฺญา, ธมฺมววตฺถาเน\-สนา วิจิกิจฺฉาย สุญฺญา, ญาเณสนา อวิชฺชาย สุญฺญา, ปาโมชฺเชสนา อรติยา สุญฺญา, ปฐมชฺฌาเน\-สนา นีวรเณหิ สุญฺญา ฯเปฯ อรหตฺตมคฺเค\-สนา สพฺพกิเลเสหิ สุญฺญา. อิทํ เอสนาสุญฺญํ. (16)}

\prose{กตมํ ปริคฺคห\-สุญฺญํ? เนกฺขมฺม\-ปริคฺคโห กามจฺฉนฺเทน สุญฺโญ, อพฺยาปาท\-ปริคฺคโห พฺยาปาเทน สุญฺโญ, อาโลกสญฺญา\-ปริคฺคโห}

\prose{\csromanfolio{364}ถินมิทฺเธน สุญฺโญ, อวิกฺเขป\-ปริคฺคโห อุทฺธจฺเจน สุญฺโญ, ธมฺมววตฺถาน\-ปริคฺคโห วิจิกิจฺฉาย สุญฺโญ ญาณปริคฺคโห อวิชฺชาย สุญฺโญ, ปาโมชฺช\-ปริคฺคโห อรติยา สุญฺโญ, ปฐมชฺฌาน\-ปริคฺคโห นีวรเณหิ สุญฺโญ ฯเปฯ อรหตฺตมคฺคปริคฺคโต สพฺพกิเลเสหิ สุญฺโญ, อิทํ ปริคฺคห\-สุญฺญํ. (17)}

\prose{กตมํ ปฏิลาภ\-สุญฺญํ? เนกฺขมฺม\-ปฏิลาโภ กามจฺฉนฺเทน สุญฺโญ, อพฺยาปาท\-ปฏิลาโภ พฺยาปาเทน สุญฺโญ, อาโลก\-สญฺญาปฏิลาโภ ถินมิทฺเธน สุญฺโญ, อวิกฺเขป\-ปฏิลาโภ อุทฺธจฺเจน สุญฺโญ, ธมฺมววตฺถาน\-ปฏิลาโภ วิจิกิจฺฉาย สุญฺโญ, ญาณปฏิลาโภ อวิชฺชาย สุญฺโญ, ปาโมชฺช\-ปฏิลาโภ อรติยา สุญฺโญ, ปฐมชฺฌาน\-ปฏิลาโภ นีวรเณหิ สุญฺโญ ฯเปฯ อรหตฺตมคฺค\-ปฏิลาโภ สพฺพกิเลเสหิ สุญฺโญ. อิทํ ปฏิลาภ\-สุญฺญํ. (18)}

\prose{กตมํ ปฏิเวธ\-สุญฺญํ? เนกฺขมฺม\-ปฺปฏิเวโธ กามจฺฉนฺเทน สุญฺโญ, อพฺยาปาท\-ปฺปฏิเวโธ พฺยาปาเทน สุญฺโญ, อาโลกสญฺญา\-ปฺปฏิเวโธ ถินมิทฺเธน สุญฺโญ, อวิกฺเขป\-ปฺปฏิเวโธ อุทฺธจฺเจน สุญฺโญ, ธมฺมววตฺถาน\-ปฺปฏิเวโธ วิจิกิจฺฉาย สุญฺโญ, ญาณปฺปฏิเวโธ อวิชฺชาย สุญฺโญ, ปาโมชฺช\-ปฺปฏิเวโธ อรติยา สุญฺโญ, ปฐมชฺฌาน\-ปฺปฏิเวโธ นีวรเณหิ สุญฺโญ ฯเปฯ อรหตฺตมคฺค\-ปฺปฏิเวโธ สพฺพกิเลเสหิ สุญฺโญ. อิทํ ปฏิเวธ\-สุญฺญํ. (19)}

\prose{กตมํ เอกตฺตสุญฺญํ, นานตฺตสุญฺญํ? กามจฺฉนฺโท นานตฺตํ, เนกฺขมฺมํ เอกตฺตํ, เนกฺขมฺเม\-กตฺตํ เจตยโต กามจฺฉนฺเทน สุญฺญํ. พฺยาปาโท นานตฺตํ, อพฺยาปาโท เอกตฺตํ, อพฺยาปาเทกตฺตํ เจตยโต พฺยาปาเทน สุญฺญํ. ถินมิทฺธํ นานตฺตํ อาโลกสญฺญา เอกตฺตํ, อาโลกสญฺเญ\-กตฺตํ เจตยโต ถินมิทฺเธน สุญฺญํ. อุทฺธจฺจํ นานตฺตํ, อวิกฺเขโป เอกตฺตํ, อวิกฺเขเป\-กตฺตํ เจตยโต อุทฺธจฺเจน สุญฺญํ, วิจิกิจฺฉา นานตฺตํ, ธมฺม\-ววตฺถานํ เอกตฺตํ, ธมฺมววตฺถาเน\-กตฺตํ เจตยโต วิจิกิจฺฉาย สุญฺญํ. อวิชฺชา นานตฺตํ, ญาณํ เอกตฺตํ, ญาเณกตฺตํ, เจตยโต อวิชฺชาย สุญฺญํ. อรติ นานตฺตํ, ปโมชฺชํ เอกตฺตํ, ปาโมชฺเชกตฺตํ เจตยโต อรติยา สุญฺญํ. นีวรณา นานตฺตํ, ปฐมชฺฌานํ เอกตฺตํ, ปฐมชฺฌาเน\-กตฺตํ เจตยโต นีวรเณหิ สุญฺญํ ฯเปฯ สพฺพกิเลสา นานตฺตํ, อรหตฺตมคฺโค เอกตฺตํ,}

\prose{\csromanfolio{365}อรหตฺตมคฺเค\-กตฺตํ เจตยโต สพฺพกิเลเสหิ สุญฺญํ. อิทํ เอกตฺตสุญฺญํ, นานตฺตสุญฺญํ. \mbox{(20, 21)}}

\prose{กตมํ ขนฺติสุญฺญํ? เนกฺขมฺม\-ขนฺติ กามจฺฉนฺเทน สุญฺญา, อพฺยาปาทขนฺติ พฺยาปาเทน สุญฺญา, อาโลกสญฺญา\-ขนฺติ ถินมิทฺเธน สุญฺญา, อวิกฺเขป\-ขนฺติ อุทฺธจฺเจน สุญฺญา, ธมฺมววตฺถาน\-ขนฺติ วิจิกิจฺฉาย สุญฺญา, ญาณขนฺติ อวิชฺชาย สุญฺญา, ปาโมชฺชขนฺติ อรติยา สุญฺญา, ปฐมชฺฌาน\-ขนฺติ นีวรเณหิ สุญฺญา ฯเปฯ อรหตฺตมคฺค\-ขนฺติ สพฺพกิเลเสหิ สุญฺญา. อิทํ ขนฺติสุญฺญํ. (22)}

\prose{กตมํ อธิฏฺฐาน\-สุญฺญํ? เนกฺขมฺมา\-ธิฏฺฐานํ กามจฺฉนฺเทน สุญฺญํ, อพฺยาปาทา\-ธิฏฺฐานํ พฺยาปาเทน สุญฺญํ, อาโลกสญฺญา\-ธิฏฺฐานํ ถินมิทฺเธน สุญฺญํ, อวิกฺเขปา\-ธิฏฺฐานํ อุทฺธจฺเจน สุญฺญํ, ธมฺมววตฺถานา\-ธิฏฺฐานํ วิจิกิจฺฉาย สุญฺญํ, ญาณาธิฏฺฐานํ อวิชฺชาย สุญฺญํ, ปาโมชฺชา\-ธิฏฺฐานํ อรติยา สุญฺญํ, ปฐมชฺฌานา\-ธิฏฺฐานํ นีวรเณหิ สุญฺญํ ฯเปฯ อรหตฺตมคฺคา\-ธิฏฺฐานํ สพฺพกิเลเสหิ สุญฺญํ. อิทํ อธิฏฺฐาน\-สุญฺญํ. (23)}

\prose{กตมํ ปริโยคาห\-ณสุญฺญํ? เนกฺขมฺมปริโยคาหณํ กามจฺฉนฺเทน สุญฺญํ, อพฺยาปาทปริโยคาหณํ พฺยาปาเทน สุญฺญํ, อาโลกสญฺญาปริโยคาหณํ ถินมิทฺเธน สุญฺญํ, อวิกฺเขปปริโยคาหณํ อุทฺธจฺเจน สุญฺญํ, ธมฺม\-ววตฺถาน\-ปริโยคาหณํ วิจิกิจฺฉาย สุญฺญํ, ญาณปริโยคาหณํ อวิชฺชาย สุญฺญํ, ปาโมชฺชปริโยคาหณํ อรติยา สุญฺญํ, ปฐมชฺฌานปริโยคาหณํ นีวรเณหิ สุญฺญํ ฯเปฯ อรหตฺตมคฺคปริโยคาหณํ สพฺพกิเลเสหิ สุญฺญํ. อิทํ ปริโยคาห\-ณสุญฺญํ. (24)}

\prose{กตมํ สมฺปชานสฺส ปวตฺต\-ปริยาทานํ สพฺพ\-สุญฺญตานํ ปรมตฺถ\-สุญฺญํ? อิธ สมฺปชาโน เนกฺขมฺเมน กามจฺฉนฺทสฺส ปวตฺตํ ปริยาทิยติ, อพฺยาปาเทน พฺยาปาทสฺส ปวตฺตํ ปริยาทิยติ, อาโลกสญฺญาย ถินมิทฺธสฺส ปวตฺตํ ปริยาทิยติ, อวิกฺเขเปน อุทฺธจฺจสฺส ปวตฺตํ ปริยาทิยติ, ธมฺม\-ววตฺถาเนน วิจิกิจฺฉาย ปวตฺตํ ปริยาทิยติ, ญาเณน อวิชฺชาย ปวตฺตํ ปริยาทิยติ, ปาโมชฺเชน อรติยา ปวตฺตํ ปริยาทิยติ, ปฐเมน ฌาเนน นีวรณานํ ปวตฺตํ ปริยาทิยติ ฯเปฯ อรหตฺต\-มคฺเคน สพฺพกิเลสานํ ปวตฺตํ ปริยาทิยติ. อถ วา ปน สมฺปชานสฺส อนุปาทิเสสาย นิพฺพานธาตุยา ปรินิพฺพายนฺตสฺส อิทํ เจว จกฺขุปวตฺตํ ปริยาทิยติ, อญฺญญฺจ จกฺขุปวตฺตํ}

\prosewithcloser{\csromanfolio{366}น อุปฺปชฺชติ. อิทํ เจว โสตปวตฺตํ ฯเปฯ. ฆานปวตฺตํ. ชิวฺหาปวตฺตํ. กายปวตฺตํ. มโนปวตฺตํ ปริยาทิยติ, อญฺญญฺจ มโนปวตฺตํ น อุปฺปชฺชติ. อิทํ สมฺปชานสฺส ปวตฺต\-ปริยาทานํ สพฺพ\-สุญฺญตานํ ปรมตฺถสุญฺญนฺติ. (25)}{สุญฺญกถา นิฏฺฐิตา.}
\nitthitammidruled{ยุคนทฺธ\-วคฺโค ทุติโย.}
\csromancenter{\textbf{ตสฺสุทฺทานํ}}
\setlength{\csromangathapagewidth}{0pt}
\setlength{\csromangathawakwidth}{0pt}
\csromangathameasuregroup{ยุคนทฺธา สจฺจโพชฺฌงฺคา, \\ ปฏิสมฺภิทา ธมฺมจกฺกํ,}{\csromangathabat{ยุคนทฺธา สจฺจโพชฺฌงฺคา,}{เมตฺตา วิราคปญฺจมา.} \\ \csromangathabat{ปฏิสมฺภิทา ธมฺมจกฺกํ,}{โลกุตฺตรพลสุญฺญาติ.} \\ เอส นิกายธเรหิ ฐปิโต อสโม ทุติโย ปวโร วรวคฺโคติ.}
\csromangathasetleft{ยุคนทฺธา สจฺจโพชฺฌงฺคา, \\ ปฏิสมฺภิทา ธมฺมจกฺกํ,}
\csromangathagroup{\csromangathastanza{\csromangathabat{ยุคนทฺธา สจฺจโพชฺฌงฺคา,}{เมตฺตา วิราคปญฺจมา.} \\ \csromangathabat{ปฏิสมฺภิทา ธมฺมจกฺกํ,}{โลกุตฺตรพลสุญฺญาติ.} \\ เอส นิกายธเรหิ ฐปิโต อสโม ทุติโย ปวโร วรวคฺโคติ.}}

\csromanmatikaanchor{mtk.118}
\csromantocmarknopage{chapter}{3. ปญฺญาวคฺค}{mtk.118}
\bookmark[dest={mtk.118},level=0]{3. ปญฺญาวคฺค}
\chapterheadpageread{3. \textbf{ปญฺญาวคฺค}}
\proseitem{1}{\csromanfolio{367}อนิจฺจา\-นุปสฺสนา ภาวิตา พหุลีกตา กตมํ ปญฺญํ ปริปูเรติ, ทุกฺขา\-นุปสฺสนา ภาวิตา พหุลีกตา กตมํ ปญฺญํ ปริปูเรติ, อนตฺตา\-นุปสฺสนา ภาวิตา พหุลีกตา กตมํ ปญฺญํ ปริปูเรติ ฯเปฯ ปฏินิสฺสคฺคา\-นุปสฺสนา ภาวิตา พหุลีกตา กตมํ ปญฺญํ ปริปูเรติ?}

\prose{อนิจฺจา\-นุปสฺสนา ภาวิตา พหุลีกตา ชวนปญฺญํ ปริปูเรติ, ทุกฺขา\-นุปสฺสนา ภาวิตา พหุลีกตา นิพฺเพธิก\-ปญฺญํ ปริปูเรติ, อนตฺตา\-นุปสฺสนา ภาวิตา พหุลีกตา มหาปญฺญํ ปริปูเรติ, นิพฺพิทา\-นุปสฺสนา ภาวิตา พหุลีกตา ติกฺขปญฺญํ ปริปูเรติ, วิราคา\-นุปสฺสนา ภาวิตา พหุลีกตา วิปุลปญฺญํ ปริปูเรติ, นิโรธา\-นุปสฺสนา ภาวิตา พหุลีกตา คมฺภีรปญฺญํ ปริปูเรติ, ปฏินิสฺสคฺคา\-นุปสฺสนา ภาวิตา พหุลีกตา อสามนฺตปญฺญํ\csromansharedfootnote{f6f09e556e2217b5}{อสฺสามนฺตปญฺญํ (สฺยา)} ปริปูเรติ. อิมา สตฺต ปญฺญา ภาวิตา พหุลีกตา ปณฺฑิจฺจํ ปริปูเรนฺติ, อิมา อฏฺฐ ปญฺญา ภาวิตา พหุลีกตา ปุถุปญฺญํ ปริปูเรนฺติ, อิมา นว ปญฺญา ภาวิตา พหุลีกตา หาสปญฺญํ ปริปูเรนฺติ.}

\prose{หาสปญฺญา ปฏิภาน\-ปฏิสมฺภิทา, ตสฺสา อตฺถ\-ววตฺถานโต อตฺถ\-ปฏิสมฺภิทา อธิคตา โหติ สจฺฉิกตา ผสฺสิตา ปญฺญาย, ธมฺม\-ววตฺถานโต ธมฺมปฏิสมฺภีทา อธิคตา โหติ สจฺฉิกตา ผสฺสิตา ปญฺญาย, นิรุตฺติ\-ววตฺถานโต นิรุตฺติ\-ปฏิสมฺภิทา อธิคตา โหติ สจฺฉิกตา ผสฺสิตา ปญฺญาย, ปฏิภาน\-ววตฺถานโต ปฏิภาน\-ปฏิสมฺภิทา อธิคตา โหติ สจฺฉิกตา ผสฺสิตา ปญฺญาย. ตสฺสิมา จตสฺโส ปฏิสมฺภิทาโย อธิคตา โหนฺติ สจฺฉิกตา ผสฺสิตา ปญฺญาย.}

\prose{รูเป อนิจฺจา\-นุปสฺสนา ภาวิตา พหุลีกตา กตมํ ปญฺญํ ปริปูเรติ ฯเปฯ รูเป ปฏินิสฺสคฺคา\-นุปสฺสนา ภาวิตา พหุลีกตา กตมํ ปญฺญํ ปริปูเรติ? รูเป อนิจฺจา\-นุปสฺสนา ภาวิตา พหุลีกตา ชวนปญฺญํ ปริปูเรติ ฯเปฯ รูเป ปฏินิสฺสคฺคา\-นุปสฺสนา ภาวิตา พหุลีกตา อสามนฺตปญฺญํ ปริปูเรติ. อิมา สตฺต ปญฺญา ภาวิตา พหุลีกตา ปณฺฑิจฺจํ ปริปูเรนฺติ, อิมา อฏฺฐ \csromanfolio{368}ปญฺญา ภาวิตา พหุลีกตา ปุถุปญฺญํ ปริปูเรนฺติ, อิมา นว ปญฺญา ภาวิตา พหุลีกตา หาสปญฺญํ ปริปูเรนฺติ.}

\prose{หาสปญฺญา ปฏิภาน\-ปฏิสมฺภิทา, ตสฺสา อตฺถ\-ววตฺถานโต อตฺถ\-ปฏิสมฺภิทา อธิคตา โหติ สจฺฉิกตา ผสฺสิตา ปญฺญาย, ธมฺม\-ววตฺถานโต ธมฺม\-ปฏิสมฺภิทา อธิคตา โหติ สจฺฉิกตา ผสฺสิตา ปญฺญาย, นิรุตฺติ\-ววตฺถานโต นิรุตฺติ\-ปฏิสมฺภิทา อธิคตา โหติ สจฺฉิกตา ผสฺสิตา ปญฺญาย, ปฏิภาน\-ววตฺถานโต ปฏิภาน\-ปฏิสมฺภิทา อธิคตา โหติ สจฺฉิกตา ผสฺสิตา ปญฺญาย. ตสฺสิมา จตสฺโส ปฏิสมฺภิทาโย อธิคตา โหนฺติ สจฺฉิกตา ผสฺสิตา ปญฺญาย.}

\prose{เวทนาย ฯเปฯ. สญฺญาย. สงฺขาเรสุ. วิญฺญาเณ. จกฺขุสฺมิํ ฯเปฯ. ชรามรเณ อนิจฺจา\-นุปสฺสนา ภาวิตา พหุลีกตา กตมํ ปญฺญํ ปริปูเรติ ฯเปฯ ชรามรเณ ปฏินิสฺสคฺคา\-นุปสฺสนา ภาวิตา พหุลีกตา กตมํ ปญฺญํ ปริปูเรติ? ชรามรเณ อนิจฺจา\-นุปสฺสนา ภาวิตา พหุลีกตา ชวนปญฺญํ ปริปูเรติ ฯเปฯ ชรามรเณ ปฏินิสฺสคฺคา\-นุปสฺสนา ภาวิตา พหุลีกตา อสามนฺตปญฺญํ ปริปูเรติ. อิมา สตฺต ปญฺญา ภาวิตา พหุลีกตา ปณฺฑิจฺจํ ปริปูเรนฺติ, อิมา อฏฺฐ ปญฺญา ภาวิตา พหุลีกตา ปุถุปญฺญํ ปริปูเรนฺติ. อิมา นว ปญฺญา ภาวิตา พหุลีกตา หาสปญฺญํ ปริปูเรนฺติ.}

\prose{หาสปญฺญา ปฏิภาน\-ปฏิสมฺภิทา, ตสฺสา อตฺถ\-ววตฺถานโต อตฺถ\-ปฏิสมฺภิทา อธิคตา โหติ สจฺฉิกตา ผสฺสิตา ปญฺญาย. ธมฺม\-ววตฺถานโต ธมฺม\-ปฏิสมฺภิทา อธิคตา โหติ สจฺฉิกตา ผสฺสิตา ปญฺญาย. นิรุตฺติ\-ววตฺถานโต นิรุตฺติ\-ปฏิสมฺภิทา อธิคตา โหติ สจฺฉิกตา ผสฺสิตา ปญฺญาย, ปฏิภาน\-ววตฺถานโต ปฏิภาน\-ปฏิสมฺภิทา อธิคตา โหติ สจฺฉิกตา ผสฺสิตา ปญฺญาย, ตสฺสิมา จตสฺโส ปฏิสมฺภิทาโย อธิคตา โหนฺติ สจฺฉิกตา ผสฺสิตา ปญฺญาย.}

\prose{รูเป อนิจฺจา\-นุปสฺสนา ภาวิตา พหุลีกตา กตมํ ปญฺญํ ปริปูเรติ. อตีตา\-นาคต\-ปจฺจุปฺปนฺเน รูเป อนิจฺจา\-นุปสฺสนา ภาวิตา พหุลีกตา กตมํ ปญฺญํ ปริปูเรติ. รูเป ทุกฺขา\-นุปสฺสนา ภาวิตา พหุลีกตา กตมํ ปญฺญํ ปริปูเรติ. อตีตา\-นาคต\-ปจฺจุปฺปนฺเน รูเป ทุกฺขา\-นุปสฺสนา ภาวิตา พหุลีกตา กตมํ ปญฺญํ ปริปูเรติ. รูเป}

\prose{\csromanfolio{369}อนตฺตา\-นุปสฺสนา ภาวิตา พหุลีกตา กตมํ ปญฺญํ ปริปูเรติ. อตีตา\-นาคต\-ปจฺจุปฺปนฺเน รูเป อนตฺตา\-นุปสฺสนา ภาวิตา พหุลีกตา กตมํ ปญฺญํ ปริปูเรติ. รูเป นิพฺพิทา\-นุปสฺสนา ภาวิตา พหุลีกตา กตมํ ปญฺญํ ปริปูเรติ. อตีตา\-นาคต\-ปจฺจุปฺปนฺเน รูเป นิพฺพิทา\-นุปสฺสนา ภาวิตา พหุลีกตา กตมํ ปญฺญํ ปริปูเรติ. รูเป วิราคา\-นุปสฺสนา ภาวิตา พหุลีกตา กตมํ ปญฺญํ ปริปูเรติ. อตีตา\-นาคต\-ปจฺจุปฺปนฺเน รูเป วิราคา\-นุปสฺสนา ภาวิตา พหุลีกตา กตมํ ปญฺญํ ปริปูเรติ. รูเป นิโรธา\-นุปสฺสนา ภาวิตา พหุลีกตา กตมํ ปญฺญํ ปริปูเรติ. อตีตา\-นาคต\-ปจฺจุปฺปนฺเน รูเป นิโรธา\-นุปสฺสนา ภาวิตา พหุลีกตา กตมํ ปญฺญํ ปริปูเรติ. รูเป ปฏินิสฺสคฺคา\-นุปสฺสนา ภาวิตา พหุลีกตา กตมํ ปญฺญํ ปริปูเรติ. อตีตา\-นาคต\-ปจฺจุปฺปนฺเน รูเป ปฏินิสฺสคฺคา\-นุปสฺสนา ภาวิตา พหุลีกตา กตมํ ปญฺญํ ปริปูเรติ?}

\prose{รูเป อนิจฺจา\-นุปสฺสนา ภาวิตา พหุลีกตา ชวนปญฺญํ ปริปูเรติ. อตีตา\-นาคต\-ปจฺจุปฺปนฺเน รูเป อนิจฺจา\-นุปสฺสนา ภาวิตา พหุลีกตา ชวนปญฺญํ ปริปูเรติ. รูเป ทุกฺขา\-นุปสฺสนา ภาวิตา พหุลีกตา นิพฺเพธิก\-ปญฺญํ ปริปูเรติ. อตีตา\-นาคต\-ปจฺจุปฺปนฺเน รูเป ทุกฺขา\-นุปสฺสนา ภาวิตา พหุลีกตา ชวนปญฺญํ ปริปูเรติ. รูเป อนตฺตา\-นุปสฺสนา ภาวิตา พหุลีกตา มหาปญฺญํ ปริปูเรติ. อตีตา\-นาคต\-ปจฺจุปฺปนฺเน รูเป อนตฺตา\-นุปสฺสนา ภาวิตา พหุลีกตา ชวนปญฺญํ ปริปูเรติ. รูเป นิพฺพิทา\-นุปสฺสนา ภาวิตา พหุลีกตา ติกฺขปญฺญํ ปริปูเรติ. อตีตา\-นาคต\-ปจฺจุปฺปนฺเน รูเป นิพฺพิทา\-นุปสฺสนา ภาวิตา พหุลีกตา ชวนปญฺญํ ปริปูเรติ. รูเป วิราคา\-นุปสฺสนา ภาวิตา พหุลีกตา วิปุลปญฺญํ ปริปูเรติ. อตีตา\-นาคต\-ปจฺจุปฺปนฺเน รูเป วิราคา\-นุปสฺสนา ภาวิตา พหุลีกตา ชวนปญฺญํ ปริปูเรติ. รูเป นิโรธา\-นุปสฺสนา ภาวิตา พหุลีกตา คมฺภีรปญฺญํ ปริปูเรติ. อตีตา\-นาคต\-ปจฺจุปฺปนฺเน รูเป นิโรธา\-นุปสฺสนา ภาวิตา พหุลีกตา ชาวนปญฺญํ ปริปูเรติ. รูเป ปฏินิสฺสคฺคา\-นุปสฺสนา ภาวิตา พหุลีกตา อสามนฺตปญฺญํ ปริปูเรติ. อตีตา\-นาคต\-ปจฺจุปฺปนฺเน รูเป ปฏินิสฺสคฺคา\-นุปสฺสนา ภาวิตา พหุลีกตา ชวนปญฺญํ ปริปูเรติ. อิมา สตฺต ปญฺญา ภาวิตา พหุลีกตา ปณฺฑิจฺจํ ปริปูเรนฺติ. อิมา อฏฺฐ ปญฺญา ภาวิตา พหุลีกตา ปุถุปญฺญํ ปริปูเรนฺติ. อิมา นว ปญฺญา ภาวิตา พหุลีกตา หาสปญฺญํ ปริปูเรนฺติ.}

\prose{\csromanfolio{370}หาสปญฺญา ปฏิภาน\-ปฏิสมฺภิทา, ตสฺสา อตฺถ\-ววตฺถานโต อตฺถ\-ปฏิสมฺภิทา อธิคตา โหติ สจฺฉิกตา ผสฺสิตา ปญฺญาย, ธมฺม\-ววตฺถานโต ธมฺม\-ปฏิสมฺภิทา อธิคตา โหติ สจฺฉิกตา ผสฺสิตา ปญฺญาย, นิรุตฺติ\-ววตฺถานโต นิรุตฺติ\-ปฏิสมฺภิทา อธิคตา โหติ สจฺฉิกตา ผสฺสิตา ปญฺญาย, ปฏิภาน\-ววตฺถานโต ปฏิภาน\-ปฏิสมฺภิทา อธิคตา โหติ สจฺฉิกตา ผสฺสิตา ปญฺญาย. ตสฺสิมา จตสฺโส ปฏิสมฺภิทาโย อธิคตา โหนฺติ สจฺฉิกตา ผสฺสิตา ปญฺญาย.}

\prose{เวทนาย ฯเปฯ. สญฺญาย. สงฺขาเรสุ. วิญฺญาเณ. จกฺขุสฺมิํ ฯเปฯ. ชรามรเณ อนิจฺจา\-นุปสฺสนา ภาวิตา พหุลีกตา กตมํ ปญฺญํ ปริปูเรติ. อตีตา\-นาคต\-ปจฺจุปฺปนฺเน ชรามรเณ อนิจฺจา\-นุปสฺสนา ภาวิตา พหุลีกตา กตมํ ปญฺญํ ปริปูเรติ ฯเปฯ ชรามรเณ ปฏินิสฺสคฺคา\-นุปสฺสนา ภาวิตา พหุลีกตา กตมํ ปญฺญํ ปริปูเรติ. อตีตา\-นาคต\-ปจฺจุปฺปนฺเน ชรามรเณ ปฏินิสฺสคฺคา\-นุปสฺสนา ภาวิตา พหุลีกตา กตมํ ปญฺญํ ปริปูเรติ? ชรามรเณ อนิจฺจา\-นุปสฺสนา ภาวิตา พหุลีกตา ชวนปญฺญํ ปริปูเรติ. อตีตา\-นาคต\-ปจฺจุปฺปนฺเน ชรามรเณ อนิจฺจา\-นุปสฺสนา ภาวิตา พหุลีกตา ชวนปญฺญํ ปริปูเรติ ฯเปฯ. ตสฺสิมา จตสฺโส ปฏิสมฺภิทาโย อธิคตา โหนฺติ สจฺฉิกตา ผสฺสิตา ปญฺญาย.}

\proseitem{3}{จตฺตาโรเม ภิกฺขเว\csromansharedfootnote{ab5caa874af51fdd}{สํ 3. 359 ปิฏฺเฐปิ.} ธมฺมา ภาวิตา พหุลีกตา โสตาปตฺติผล\-สจฺฉิกิริยาย สํวตฺตนฺติ. กตเม จตฺตาโร? สปฺปุริส\-สํเสโว สทฺธมฺม\-สฺสวนํ โยนิโส\-มนสิกาโร ธมฺมานุธมฺม\-ปฏิปตฺติ. อิเม โข ภิกฺขเว จตฺตาโร ธมฺมา ภาวิตา พหุลีกตา โสตาปตฺติผล\-สจฺฉิกิริยาย สํวตฺตนฺติ.}

\prose{จตฺตาโรเม ภิกฺขเว ธมฺมา ภาวิตา พหุลีกตา สกทาคามิ\-ผล\-สจฺฉิกิริยาย สํวตฺตนฺติ ฯเปฯ อนาคามิ\-ผล\-สจฺฉิกิริยาย สํวตฺตนฺติ ฯเปฯ อรหตฺตผล\-สจฺฉิกิริยาย สํวตฺตนฺติ. กตเม จตฺตาโร, สปฺปุริส\-สํเสโว สทฺธมฺม\-สฺสวนํ โยนิโส\-มนสิกาโร ธมฺมานุธมฺม\-ปฏิปตฺติ. อิเม โข ภิกฺขเว จตฺตาโร ธมฺมา ภาวิตา พหุลีกตา อรหตฺตมคฺค\-ผลสจฺฉิกิริยาย สํวตฺตนฺติ.}

\prose{\csromanfolio{371}จตฺตาโรเม ภิกฺขเว ธมฺมา ภาวิตา พหุลีกตา ปญฺญา\-ปฏิลาภาย สํวตฺตนฺติ ฯเปฯ ปญฺญาพุทฺธิยา สํวตฺตนฺติ. ปญฺญาเวปุลฺลาย สํวตฺตนฺติ. มหาปญฺญตาย สํวตฺตนฺติ. ปุถุปญฺญตาย สํวตฺตนฺติ. วิปุล\-ปญฺญตาย สํวตฺตนฺติ. คมฺภีร\-ปญฺญตาย สํวตฺตนฺติ. อสามนฺต\-ปญฺญตาย\csromansharedfootnote{442078e3d1747a3d}{อสฺสามนฺตปญฺญตาย (สฺยา) สํ 3. 361 ปิฏฺเฐ ปสฺสิตพฺพา.} สํวตฺตนฺติ. ภูริปญฺญตาย สํวตฺตนฺติ. ปญฺญาพาหุลฺลาย สํวตฺตนฺติ. สีฆปญฺญตาย สํวตฺตนฺติ. ลหุปญฺญตาย สํวตฺตนฺติ. หาสปญฺญตาย สํวตฺตนฺติ. ชวน\-ปญฺญตาย สํวตฺตนฺติ. ติกฺข\-ปญฺญตาย สํวตฺตนฺติ. นิพฺเพธิก\-ปญฺญตาย สํวตฺตนฺติ. กตเม จตฺตาโร? สปฺปุริส\-สํเสโว สทฺธมฺม\-สฺสวนํ โยนิโส\-มนสิกาโร ธมฺมานุธมฺม\-ปฏิปตฺติ. อิเม โข ภิกฺขเว จตฺตาโร ธมฺมา ภาวิตา พหุลีกตา ปญฺญา\-ปฏิลาภาย สํวตฺตนฺติ. ปญฺญาพุทฺธิยา สํวตฺตนฺติ ฯเปฯ นิพฺเพธิก\-ปญฺญตาย สํวตฺตนฺติ.}

\csromanmatikaanchor{mtk.119}
\csromanmatikaanchor{mtk.120}
\csromantocmarknopage{section}{1. มหาปญฺญากถา}{mtk.119}
\bookmark[dest={mtk.119},level=1]{1. มหาปญฺญากถา}
\csromantocmarkref{subsection}{1. โสฬสปญฺญานิทฺเทส}{mtk.120}
\bookmark[dest={mtk.120},level=2]{1. โสฬสปญฺญานิทฺเทส}
\csromanheader{2}{1. โสฬส\-ปญฺญา\-นิทฺเทส}
\proseitem{4}{ปญฺญา\-ปฏิลาภาย สํวตฺตนฺตีติ กตโม ปญฺญาปฏิลาโภ? จตุนฺนํ มคฺคญาณานํ จตุนฺนํ ผลญาณานํ จตุนฺนํ ปฏิสมฺภิทา\-ญาณานํ ฉนฺนํ อภิญฺญาญาณานํ เตสตฺตตีนํ ญาณานํ สตฺต\-สตฺตตีนํ ญาณานํ ลาโภ ปฏิลาโภ ปตฺติ สมฺปตฺติ ผสฺสนา\csromansharedfootnote{926d0233078f5eb7}{ผุสนา (ก)} สจฺฉิกิริยา อุปสมฺปทา ปญฺญา\-ปฏิลาภาย สํวตฺตนฺตีติ อยํ ปญฺญา ปฏิลาโภ. (1)}

\prose{ปญฺญาพุทฺธิยา สํวตฺตนฺตีติ กตมา ปญฺญาพุทฺธิ? สตฺตนฺนํ จ เสกฺขานํ ปุถุชฺชน\-กลฺยาณกสฺส จ ปญฺญา วฑฺฒติ, อรหโต ปญฺญา วฑฺฒติ, วฑฺฒิภวฑฺฒนา ปญฺญาพุทฺธิยา สํวตฺตนฺติตี อยํ ปญฺญาพุทฺธิ. (2)}

\prose{ปญฺญาเวปุลฺลาย สํวตฺตนฺตีติ กตมํ ปญฺญาเวปุลฺลํ? สตฺตนฺนํ เสกฺขานํ ปุถุชฺชน\-กลฺยาณกสฺส จ ปญฺญาเวปุลฺลํ คจฺฉติ. อรหโต ปญฺญา เวปุลฺลคตา ปญฺญาเวปุลฺลาย สํวตฺตนฺตีติ อิทํ ปญฺญาเวปุลฺลํ. (3)}

\prose{มหาปญฺญตาย สํวตฺตนฺตีติ กตมา มหาปญฺญา? มหนฺเต อตฺเถ ปริคฺคณฺหาตีติ มหาปญฺญา, มหนฺเต ธมฺเม ปริคฺคณฺหาตีติ มหาปญฺญา, มหนฺตา นิรุตฺติโย ปริคฺคณฺหาตีติ มหาปญฺญา, มหนฺตานิ ปฏิภานานิ ปริคฺคณฺหาตีติ มหาปญฺญา, มหนฺเต สีลกฺขนฺเธ ปริคฺคณฺหาตีติ มหาปญฺญา, มหนฺเต สมาธิ\-กฺขนฺเธ ปริคฺคณฺหาตีติ มหาปญฺญา, มหนฺเต ปญฺญากฺขนฺเธ \csromanfolio{372}ปริคฺคณฺหาตีติ มหาปญฺญา, มหนฺเต วิมุตฺติ\-กฺขนฺเธ ปริคฺคณฺหาตีติ มหาปญฺญา, มหนฺเต วิมุตฺติ\-ญาณ\-ทสฺสนกฺขนฺเธ ปริคฺคณฺหาตีติ มหาปญฺญา, มหนฺตานิ ฐานาฏฺฐานานิ ปริคฺคณฺหาตีติ มหาปญฺญา, มหนฺตา วิหาร\-สมาปตฺติโย ปริคฺคณฺหาตีติ มหาปญฺญา, มหนฺตานิ อริยสจฺจานิ ปริคฺคณฺหาตีติ มหาปญฺญา, มหนฺเต สติปฏฺฐาเน ปริคฺคณฺหาตีติ มหาปญฺญา, มหนฺเต สมฺมปฺปธาเน ปริคฺคณฺหาตีติ มหาปญฺญา, มหนฺเต อิทฺธิปาเท ปริคฺคณฺหาตีติ มหาปญฺญา, มหนฺตานิ อินฺทฺริยานิ ปริคฺคณฺหาตีติ มหาปญฺญา, มหนฺตานิ พลานิ ปริคฺคณฺหาตีติ มหาปญฺญา, มหนฺเต โพชฺฌงฺเค ปริคฺคณฺหาตีติ มหาปญฺญา, มหนฺตํ อริยมคฺคํ\csromansharedfootnote{3e9bfc6da4a95482}{มหนฺเต อริยมคฺเค (ก)} ปริคฺคณฺหาตีติ มหาปญฺญา, มหนฺตานิ สามญฺญผลานิ ปริคฺคณฺหาตีติ มหาปญฺญา, มหนฺตา อภิญฺญาโย\csromansharedfootnote{aa53974adc778ce9}{มหาภิญฺญาโย (ก)} ปริคฺคณฺหาตีติ มหาปญฺญา, มหนฺตํ ปรมตฺถํ นิพฺพานํ ปริคฺคณฺหาตีติ มหาปญฺญา, มหาปญฺญตาย สํวตฺตนฺตีติ อยํ มหาปญฺญา. (4)}

\prose{ปุถุปญฺญตาย สํวตฺตนฺตีติ กตมา ปุถุปญฺญา? ปุถุนานา\-ขนฺเธสุ ญาณํ ปวตฺตตีติ ปุถุปญฺญา, ปุถุนานาธาตูสุ ญาณํ ปวตฺตตีติ ปุถุปญฺญา, ปุถุนานา\-อายตเนสุ ญาณํ ปวตฺตตีติ ปุถุปญฺญา, ปุถุนานา\-ปฏิจฺจสมุปฺปาเทสุ ญาณํ ปวตฺตตีติ ปุถุปญฺญา, ปุถุนานา\-สุญฺญตม\-นุป\-ลพฺเภสุ ญาณํ ปวตฺตตีติ ปุถุปญฺญา, ปุถุนานา\-อตฺเถสุ ญาณํ ปวตฺตตีติ ปุถุปญฺญา, ปุถุนา\-นาธมฺเมสุ ญาณํ ปวตฺตตีติ ปุถุปญฺญา, ปุถุนานา\-นิรุตฺตีสุ ญาณํ ปวตฺตตีติ ปุถุปญฺญา, ปุถุนานา\-ปฏิภาเนสุ ญาณํ ปวตฺตตีติ ปุถุปญฺญา, ปุถุนานา\-สีลกฺขนฺเธสุ ญาณํ ปวตฺตตีติ ปุถุปญฺญา, ปุถุนานา\-สมาธิกฺขนฺเธสุ ญาณํ ปวตฺตตีติ ปุถุปญฺญา, ปุถุนานา\-ปญฺญากฺขนฺเธสุ ญาณํ ปวตฺตตีติ ปุถุปญฺญา, ปุถุนานา\-วิมุตฺติกฺขนฺเธสุ ญาณํ ปวตฺตตีติ ปุถุปญฺญา. ปุถุนานา\-วิมุตฺติญาณทสฺสนกฺขนฺเธสุ ญาณํ ปวตฺตตีติ ปุถุปญฺญา, ปุถุนานา\-ฐานาฏฺฐาเนสุ ญาณํ ปวตฺตตีติ ปุถุปญฺญา, ปุถุนานา\-วิหารสมาปตฺตีสุ ญาณํ ปวตฺตตีติ ปุถุปญฺญา, ปุถุนานา\-อริยสจฺเจสุ ญาณํ ปวตฺตตีติ ปุถุปญฺญา. ปุถุนานา\-สติปฏฺฐาเนสุ ญาณํ ปวตฺตตีติ ปุถุปญฺญา, ปุถุนานา\-สมฺมปฺปธาเนสุ ญาณํ ปวตฺตตีติ ปุถุปญฺญา. ปุถุนานา\-อิทฺธิปาเทสุ ญาณํ ปวตฺตตีติ ปุถุปญฺญา. ปุถุนานา\-อินฺทฺริเยสุ ญาณํ ปวตฺตตีติ ปุถุปญฺญา. ปุถุนานา\-พเลสุ ญาณํ ปวตฺตตีติ ปุถุปญฺญา. ปุถุนานา\-โพชฺฌงฺเคสุ ญาณํ}

\prose{\csromanfolio{373}ปวตฺตตีติ ปุถุปญฺญา. ปุถุนานา\-อริยมคฺเค ญาณํ ปวตฺตตีติ ปุถุปญฺญา, ปุถุนานา\-สามญฺญผเลสุ ญาณํ ปวตฺตตีติ ปุถุปญฺญา. ปุถุนานา\-อภิญฺญาสุ ญาณํ ปวตฺตตีติ ปุถุปญฺญา. ปุถุชฺชน\-สาธารเณ ธมฺเม อติกฺกมฺม\csromansharedfootnote{a61c2441c7dd29b3}{สมติกฺกมฺม (สฺยา)} ปรมตฺเถ นิพฺพาเน ญาณํ ปวตฺตตีติ ปุถุปญฺญา. ปุถุปญฺญตาย สํวตฺตนฺตีติ อยํ ปุถุปญฺญา. (5)}

\prose{วิปุล\-ปญฺญตาย สํวตฺตนฺตีติ กตมา วิปุลปญฺญา? วิปุเล อตฺเถ ปริคฺคณฺหาตีติ วิปุลปญฺญา, วิปุเล ธมฺเม ปริคฺคณฺหาตีติ วิปุลปญฺญา, วิปุลา นิรุตฺติโย ปริคฺคณฺหาตีติ วิปุลปญฺญา, วิปุลานิ ปฏิภานานิ ปริคฺคณฺหาตีติ วิปุลปญฺญา, วิปุเล สีลกฺขนฺเธ ปริคฺคณฺหาตีติ วิปุลปญฺญา, วิปุเล สมาธิ\-กฺขนฺเธ ปริคฺคณฺหาตีติ วิปุลปญฺญา, วิปุเล ปญฺญากฺขนฺเธ ปริคฺคณฺหาตีติ วิปุลปญฺญา, วิปุเล วิมุตฺติ\-กฺขนฺเธ ปริคฺคณฺหาตีติ วิปุลปญฺญา, วิปุเล วิมุตฺติ\-ญาณ\-ทสฺสนกฺขนฺเธ ปริคฺคณฺหาตีติ วิปุลปญฺญา, วิปุลานิ ฐานาฏฺฐานานิ ปริคฺคณฺหาตีติ วิปุลปญฺญา, วิปุลา วิหาร\-สมาปตฺติโย ปริคฺคณฺหาตีติ วิปุลปญฺญา, วิปุลานิ อริยสจฺจานิ ปริคฺคณฺหาตีติ วิปุลปญฺญา, วิปุเล สติปฏฺฐาเน ปริคฺคณฺหาตีติ วิปุลปญฺญา, วิปุเล สมฺมปฺปธาเน ปริคฺคณฺหาตีติ วิปุลปญฺญา, วิปุเล อิทฺธิปาเท ปริคฺคณฺหาตีติ วิปุลปญฺญา, วิปุลานิ อินฺทฺริยานิ ปริคฺคณฺหาตีติ วิปุลปญฺญา, วิปุลานิ พลานิ ปริคฺคณฺหาตีติ วิปุลปญฺญา, วิปุเล โพชฺฌงฺเค ปริคฺคณฺหาตีติ วิปุลปญฺญา, วิปุลํ อริยมคฺคํ ปริคฺคณฺหาตีติ วิปุลปญฺญา, วิปุลานิ สามญฺญผลานิ ปริคฺคณฺหาตีติ วิปุลปญฺญา, วิปุลา อภิญฺญาโย ปริคฺคณฺหาตีติ วิปุลปญฺญา, วิปุลํ ปรมตฺถํ นิพฺพานํ ปริคฺคณฺหาตีติ วิปุลปญฺญา, วิปุล\-ปญฺญตาย สํวตฺตนฺตีติ อยํ วิปุลปญฺญา. (6)}

\prose{คมฺภีร\-ปญฺญตาย สํวตฺตนฺตีติ กตมา คมฺภีรปญฺญา? คมฺภีเรสุ ขนฺเธสุ ญาณํ ปวตฺตตีติ คมฺภีรปญฺญา, คมฺภีราสุ ธาตูสุ ญาณํ ปวตฺตตีติ คมฺภีรปญฺญา, คมฺภีเรสุ อายตเนสุ ญาณํ ปวตฺตตีติ คมฺภีรปญฺญา, คมฺภีเรสุ ปฏิจฺจ\-สมุปฺปาเทสุ ญาณํ ปวตฺตตีติ คมฺภีรปญฺญา. คมฺภีเรสุ สุญฺญตม\-นุป\-ลพฺเภสุ ญาณํ ปวตฺตตีติ คมฺภีรปญฺญา, คมฺภีเรสุ อตฺเถสุ ญาณํ ปวตฺตตีติ คมฺภีรปญฺญา, คมฺภีเรสุ ธมฺเมสุ ญาณํ ปวตฺตตีติ คมฺภีรปญฺญา, คมฺภีราสุ นิรุตฺตีสุ ญาณํ ปวตฺตตีติ คมฺภีรปญฺญา, คมฺภีเรสุ ปฏิภาเนสุ ญาณํ}

\prose{\csromanfolio{374}ปวตฺตตีติ คมฺภีรปญฺญา, คมฺภีเรสุ สีลกฺขนฺเธสุ ญาณํ ปวตฺตตีติ คมฺภีรปญฺญา, คมฺภีเรสุ สมาธิ\-กฺขนฺเธสุ ญาณํ ปวตฺตตีติ คมฺภีรปญฺญา, คมฺภีเรสุ ปญฺญา\-กฺขนฺเธสุ ญาณํ ปวตฺตตีติ คมฺภีรปญฺญา, คมฺภีเรสุ วิมุตฺติ\-กฺขนฺเธสุ ญาณํ ปวตฺตตีติ คมฺภีรปญฺญา, คมฺภีเรสุ วิมุตฺติ\-ญาณ\-ทสฺสนกฺขนฺเธสุ ญาณํ ปวตฺตตีติ คมฺภีรปญฺญา, คมฺภีเรสุ ฐานาฏฺฐาเนสุ ญาณํ ปวตฺทตีติ คมฺภีรปญฺญา, คมฺภีราสุ วิหาร\-สมาปตฺตีสุ ญาณํ ปวตฺตทีติ คมฺภีรปญฺญา, คมฺภีเรสุ อริยสจฺเจสุ ญาณํ ปวตฺตตีติ คมฺภีรปญฺญา, คมฺภีเรสุ สติปฏฺฐาเนสุ ญาณํ ปวตฺตตีติ คมฺภีรปญฺญา, คมฺภีเรสุ สมฺม\-ปฺปธาเนสุ ญาณํ ปวตฺตตีติ คมฺภีรปญฺญา, คมฺภีเรสุ อิทฺธิปาเทสุ ญาณํ ปวตฺตตีติ คมฺภีรปญฺญา, คมฺภีเรสุ อินฺทฺริเยสุ ญาณํ ปวตฺตตีติ คมฺภีรปญฺญา, คมฺภีเรสุ พเลสุ ญาณํ ปวตฺตตีติ คมฺภีรปญฺญา, คมฺภีเรสุ โพชฺฌงฺเคสุ ญาณํ ปวตฺตตีติ คมฺภีรปญฺญา, คมฺภีเร อริยมคฺเค ญาณํ ปวตฺตตีติ คมฺภีรปญฺญา, คมฺภีเรสุ สามญฺญผเลสุ ญาณํ ปวตฺตตีติ คมฺภีรปญฺญา, คมฺภีราสุ อภิญฺญาสุ ญาณํ ปวตฺตตีติ คมฺภีรปญฺญา, คมฺภีเร ปรมตฺเถ นิพฺพาเน ญาณํ ปวตฺตตีติ คมฺภีรปญฺญา, คมฺภีร\-ปญฺญตาย สํวตฺตนฺตีติ อยํ คมฺภีรปญฺญา. (7)}

\prose{อสามนฺต\-ปญฺญตาย สํวตฺตนฺตีติ กตมา อสามนฺตปญฺญา? ยสฺส ปุคฺคลสฺส อตฺถ\-ววตฺถานโต อตฺถ\-ปฏิสมฺภิทา อธิคตา โหติ สจฺฉิกตา ผสฺสิตา ปญฺญาย, ธมฺม\-ววตฺถานโต ธมฺม\-ปฏิสมฺภิทา อธิคตา โหติ สจฺฉิกตา ผสฺสิตา ปญฺญาย, นิรุตฺติ\-ววตฺถานโต นิรุตฺติ\-ปฏิสมฺภิทา อธิคตา โหติ สจฺฉิกตา ผสฺสิตา ปญฺญาย, ปฏิภาน\-ววตฺถานโต ปฏิภาน\-ปฏิสมฺภิทา อธิคตา โหติ สจฺฉิกตา ผสฺสิตา ปญฺญาย, ตสฺส อตฺเถ จ ธมฺเม จ นิรุตฺติยา จ ปฏิภาเน จ น อญฺโญ โกจิ สกฺโกติ อภิสมฺภวิตุํ, อนภิสมฺภวนีโย จ โส อญฺเญหีติ อสามนฺตปญฺโญ.}

\prose{ปุถุชฺชน\-กลฺยาณกสฺส ปญฺญา อฏฺฐมกสฺส ปญฺญาย ทูเร วิทูเร สุวิทูเร น สนฺติเก น สามนฺตา, ปุถุชฺชน\-กลฺยาณกํ อุปาทาย อฏฺฐมโก อสามนฺตปญฺโญ. อฏฺฐมกสฺส ปญฺญา โสตาปนฺนสฺส ปญฺญาย ทูเร วิทูเร สุวิทูเร น สนฺติเก น สามนฺตา, อฏฺฐมกํ อุปาทาย โสตาปนฺโน อสามนฺตปญฺโญ. โสตาปนฺนสฺส ปญฺญา สกทาคามิสฺส ปญฺญาย ทูเร วิทูเร สุวิทูเร น สนฺติเก น สามนฺตา, โสตาปนฺนํ อุปาทาย สกทาคามี อสามนฺตปญฺโญ. สกทาคามิสฺส ปญฺญา อนาคามิสฺส ปญฺญาย}

\prose{\csromanfolio{375}ทูเร วิทูเร สุวิทูเร น สนฺติเก น สามนฺตา, สกทาคามิํ อุปาทาย อนาคามี อสามนฺตปญฺโญ. อนาคามิสฺส ปญฺญา อรหโต ปญฺญาย ทูเร วิทูเร สุวิทูเร น สนฺติเก น สามนฺตา, อนาคามิํ อุปาทาย อรหา อสามนฺตปญฺโญ. อรหโต ปญฺญา ปจฺเจก\-สมฺพุทฺธสฺส\csromansharedfootnote{8eae175b86e866c6}{ปจฺเจกพุทฺธสฺส (สฺยา, ก) อฏฺฐกถา โอโลเกตพฺพา.} ปญฺญาย ทูเร วิทูเร สุวิทูเร น สนฺติเก น สามนฺตา, อรหนฺตํ อุปาทาย ปจฺเจกพุทฺโธ อสามนฺตปญฺโญ. ปจฺเจกพุทฺธญฺจ สเทวกญฺจ โลกํ อุปาทาย ตถาคโต อรหํ สมฺมาสมฺพุทฺโธ อคฺโค อสามนฺตปญฺโญ.}

\proseitem{5}{ปญฺญาปเภท\-กุสโล ปภินฺนญาโณ อธิคตปฺปฏิสมฺภิโท จตุเวสารชฺช\-ปฺปตฺโต ทสพลธารี ปุริสาสโภ ปุริสสีโห ปุริสนาโค ปุริสาชญฺโญ ปุริสโธรโยฺห อนนฺตญาโณ อนนฺตเตโช อนนฺตยโส อฑฺโฒ มหทฺธโน ธนวา เนตา วิเนตา อนุเนตา ปญฺญาเปตา นิชฺฌาเปตา เปกฺเขตา ปสาเทตา. โส หิ ภควา อนุปฺปนฺนสฺส มคฺคสฺส อุปฺปาเทตา, อสญฺชาตสฺส มคฺคสฺส สญฺชเนตา\csromansharedfootnote{5897d904e58ba5db}{สญฺชาเนตา (สฺยา)}, อนกฺขาตสฺส มคฺคสฺส อกฺขาตา, มคฺคญฺญู มคฺควิทู มคฺคโกวิโท มคฺคานุคามี\csromansharedfootnote{f182bf5222b432ac}{มคฺคานุคา (สฺยา) ขุ 7. 138 ปิฏฺเฐ ปสฺสิตพฺพา.} จ ปนสฺส เอตรหิ สาวกา วิหรนฺติ ปจฺฉา สมนฺนาคตา.}

\prose{โส หิ ภควา ชานํ ชานาติ, ปสฺสํ ปสฺสติ, จกฺขุภูโต ญาณภูโต ธมฺมภูโต พฺรหฺมภูโต วตฺตา ปวตฺตา อตฺถสฺส นินฺเนตา อมตสฺส ทาตา ธมฺมสฺสามี ตถาคโต, นตฺถิ ตสฺส ภควโต อญฺญาตํ อทิฏฺฐํ อวิทิตํ อสจฺฉิกตํ อผสฺสิตํ ปญฺญาย.  อตีตํ อนาคตํ ปจฺจุปฺปนฺนํ\csromansharedfootnote{c5699a5737c466cd}{อตีตานาคตปจฺจุปฺปนฺนํ (สฺยา) ขุ 7. 138 ปิฏฺเฐปิ.} อุปาทาย สพฺเพ ธมฺมา สพฺพากาเรน พุทฺธสฺส ภควโต ญาณมุเข อาปาถํ อาคจฺฉนฺติ, ยํ กิญฺจิ เนยฺยํ นาม อตฺถิ, ตํ สพฺพํ ชานิตพฺพํ\csromansharedfootnote{9af02245a4065f0a}{สพฺพํ ธมฺมํ ชานิตพฺพํ (ก)}, อตฺตตฺโถ วา ปรตฺโถ วา อุภยตฺโถ วา ทิฏฺฐธมฺมิโก วา อตฺโถ สมฺปรายิโก วา อตฺโถ อุตฺตาโน วา อตฺโถ คมฺภีโร วา อตฺโถ คูโฬฺห วา อตฺโถ ปฏิจฺฉนฺโน วา อตฺโถ เนยฺโย วา อตฺโถ นีโต วา อตฺโถ อนวชฺโช วา อตฺโถ นิกฺกิเลโส วา อตฺโถ โวทาโน วา อตฺโถ ปรมตฺโถ วา อตฺโถ, สพฺพํ ตํ อนฺโตพุทฺธญาเณ ปริวตฺตติ.}

\prose{\csromanfolio{376}สพฺพํ กายกมฺมํ พุทฺธสฺส ภควโต ญาณา\-นุปริวตฺติ, สพฺพํ วจีกมฺมํ พุทฺธสฺส ภควโต ญาณา\-นุปริวตฺติ, สพฺพํ มโนกมฺมํ พุทฺธสฺส ภควโต ญาณา\-นุปริวตฺติ. อตีเต พุทฺธสฺส ภควโต อปฺปฏิหตํ ญาณํ, อนาคเต พุทฺธสฺส ภควโต อปฺปฏิหตํ ญาณํ, ปจฺจุปฺปนฺเน พุทฺธสฺส ภควโต อปฺปฏิหตํ ญาณํ, ยาวตกํ เนยฺยํ, ตาวตกํ ญาณํ. ยาวตกํ ญาณํ, ตาวตกํ เนยฺยํ, เนยฺย\-ปริยนฺติกํ ญาณํ, ญาณ\-ปริยนฺติกํ เนยฺยํ, เนยฺยํ อติกฺกมิตฺวา ญาณํ นปฺปวตฺตติ, ญาณํ อติกฺกมิตฺวา เนยฺยปโถ นตฺถิ, อญฺญมญฺญ\-ปริยนฺต\-ฏฺฐายิโน เต ธมฺมา, ยถา ทฺวินฺนํ สมุคฺค\-ปฏลานํ สมฺมา ผุสิตานํ\csromansharedfootnote{48ab254ff482e7f1}{สุผุสฺสิตานํ (สฺยา, ก)} เหฏฺฐิมํ สมุคฺคปฏลํ อุปริมํ นาติวตฺตติ, อุปริมํ สมุคฺคปฏลํ เหฏฺฐิมํ นาติวตฺตติ, อญฺญมญฺญ\-ปริยนฺต\-ฏฺฐายิโน. เอวเมว พุทฺธสฺส ภควโต เนยฺยญฺจ ญาณญฺจ อญฺญมญฺญ\-ปริยนฺต\-ฏฺฐายิโน. ยาวตกํ เนยฺยํ, ตาวตกํ ญาณํ, ยาวตกํ ญาณํ, ตาวตกํ เนยฺยํ, เนยฺย\-ปริยนฺติกํ ญาณํ, ญาณ\-ปริยนฺติกํ เนยฺยํ, เนยฺยํ อติกฺกมิตฺวา ญาณํ นปฺปวตฺตติ, ญาณํ อติกฺกมิตฺวา เนยฺยปโถ นตฺถิ, อญฺญมญฺญ\-ปริยนฺต\-ฏฺฐายิโน เต ธมฺมา. สพฺพธมฺเมสุ พุทฺธสฺส ภควโต ญาณํ ปวตฺตติ.}

\prose{สพฺเพ ธมฺมา พุทฺธสฺส ภควโต อาวชฺชน\-ปฺปฏิพทฺธา อากงฺขปฺปฏิพทฺทา มนสิการ\-ปฺปฏิพทฺธา จิตฺตุปฺปาท\-ปฺปฏิพทฺธา, สพฺพสตฺเตสุ พุทฺธสฺส ภควโต ญาณํ ปวตฺตติ, สพฺเพสํ สตฺตานํ พุทฺโธ อาสยํ ชานาติ, อนุสยํ ชานาติ, จริตํ\csromansharedfootnote{bc024bcb431176df}{จริยํ (สฺยา) ขุ 7. 139 ปิฏฺเฐ ปสฺสิตพฺพา.} ชานาติ, อธิมุตฺติํ ชานาติ, อปฺปรชกฺเข มหารชกฺเข ติกฺขินฺทฺริเย มุทินฺทฺริเย สฺวากาเร ทฺวากาเร สุวิญฺญาปเย ทุวิญฺญาปเย ภพฺพาภพฺเพ สตฺเต ปชานาติ, สเทวโก โลโก สมารโก สพฺรหฺมโก สสฺสมณ\-พฺราหฺมณี ปชา สเทวมนุสฺสา อนฺโตพุทฺธญาเณ ปริวตฺตติ.}

\prose{ยถา เย เกจิ มจฺฉกจฺฉปา อนฺตมโส ติมิติมิงฺคลํ อุปาทาย อนฺโต\-มหาสมุทฺเท ปริวตฺตนฺติ, เอวเมวํ สเทวโก โลโก สมารโก สพฺรหฺมโก สสฺสมณ\-พฺราหฺมณี ปชา สเทวมนุสฺสา อนฺโตพุทฺธญาเณ ปริวตฺตติ.  ยถา เย เกจิ ปกฺขิโน อนฺตมโส ครุฬํ เวนเตยฺยํ อุปาทาย อากาสสฺส ปเทเส ปริวตฺตนฺติ. เอวเมว เยปิ เต สาริปุตฺตสมา ปญฺญาย, เตปิ พุทฺธญาณสฺส ปเทเส ปริวตฺตนฺติ, พุทฺธญาณํ}

\prose{\csromanfolio{377}เทวมนุสฺสานํ ปญฺญํ ผริตฺวา อติฆํสิตฺวา ติฏฺฐติ. เยปิ เต ขตฺติย\-ปณฺฑิตา พฺราหฺมณ\-ปณฺฑิตา คหปติ\-ปณฺฑิตา สมณปณฺฑิตา นิปุณา กตปรปฺปวาทา วาลเวธิรูปา, โวภินฺทนฺตา\csromansharedfootnote{8c0d862c67c1e117}{เต ภินฺทนฺตา (สฺยา, ก)} มญฺเญ จรนฺติ ปญฺญาคเตน ทิฏฺฐิคตานิ, เต ปญฺหํ อภิส\-งฺขริ\-ตฺวา อภิส\-งฺขริ\-ตฺวา ตถาคตํ อุปสงฺกมิตฺวา ปุจฺฉนฺติ คูฬฺหานิ จ ปฏิจฺฉนฺนานิ จ, กถิตา วิสชฺชิตา จ เต ปญฺหา ภควตา โหนฺติ นิทฺทิฏฺฐ\-การณา\csromansharedfootnote{cc18967b12849b08}{นิทฺทิฏฺฐิการณา (สฺยา)}, อุปกฺขิตฺตกา จ เต ภควโต สมฺปชฺชนฺติ. อถ โข ภควา ตตฺถ อติโรจติ ยทิทํ ปญฺญายาติ อคฺโค อสามนฺตปญฺโญ, อสามนฺต\-ปญฺญตาย สํวตฺตนฺตีติ อยํ อสามนฺตปญฺญา. (8)}

\proseitem{6}{\textbf{ภูริปญฺญตาย สํวตฺตนฺตี}ติ กตมา ภูริปญฺญา? ราคํ อภิภุยฺยตีติ ภูริปญฺญา, อภิภวิตาติ ภูริปญฺญา. โทสํ อภิภุยฺยตีติ ภูริปญฺญา, อภิภวิตาติ ภูริปญฺญา,. โมหํ อภิภุยฺยตีติ ภูริปญฺญา, อภิภวิตาติ ภูริปญฺญา. โกธํ ฯเปฯ. อุปนาหํ. มกฺขํ. ปฬาสํ. อิสฺสํ. มจฺฉริยํ. มายํ. สาเฐยฺยํ. ถมฺภํ. สารมฺภํ. มานํ. อติมานํ. มทํ. ปมาทํ. สพฺเพ กิเลเส. สพฺเพ ทุจฺจริเต. สพฺเพ อภิสงฺขาเร ฯเปฯ. สพฺเพ ภวคามิกมฺเม อภิภุยฺยตีติ ภูริปญฺญา, อภิภวิตาติ ภูริปญฺญา. ราโค อริ, ตํ อริํ มทฺทนิปญฺญาติ ภูริปญฺญา, โทโส อริ, ตํ อริํ มทฺทนิปญฺญาติ ภูริปญฺญา. โมโห อริ, ตํ อริํ มทฺทนิปญฺญาติ ภูริปญฺญา. โกโธ ฯเปฯ. อุปนาโห. มกฺโข. ปฬาโส. อิสฺสา. มจฺฉริยํ. มายา. สาเฐยฺยํ. ถมฺโภ. สารมฺโภ. มาโน. อติมาโน. มโท. ปมาโท. สพฺเพ กิเลสา. สพฺเพ ทุจฺจริตา. สพฺเพ อภิสงฺขารา ฯเปฯ. สพฺเพ ภวคามิกมฺมา อริ, ตํ อริํ มทฺทนิปญฺญาติ ภูริปญฺญา. ภูริ วุจฺจติ ปถวี\csromansharedfootnote{d775b7489019dc95}{ปฐวี (สฺยา)}, ตาย ปถวิสมาย วิตฺถตาย วิปุลาย ปญฺญาย สมนฺนาคโตติ ภูริปญฺญา. อปิ จ ปญฺญาย เมตํ อธิวจนํ “ภูริ เมธา ปริณายิกา”ติ ภูริปญฺญา, ภูริปญฺญตาย สํวตฺตนฺตีติ อยํ ภูริปญฺญา. (9)}

\prose{ปญฺญาพาหุลฺลาย สํวตฺตนฺตีติ กตมํ ปญฺญาพาหุลฺลํ? อิเธกจฺโจ ปญฺญาครุโก โหติ ปญฺญาจริโต ปญฺญาสโย ปญฺญาธิมุตฺโต ปญฺญาธโช ปญฺญาเกตุ ปญฺญา\-ธิปเตยฺโย วิจยพหุโล ปวิจยพหุโล โอกฺขายน\-พหุโล สโมกฺขายน\-พหุโล\csromansharedfootnote{1fe9d868ad771a22}{สมฺเปกฺขายนพหุโล (สฺยา, ก) ขุ 7. 393 ปิฏฺเฐ ปสฺสิตพฺพา.} สมฺเปกฺขายน\-ธมฺโม วิภูตวิหารี ตจฺจริโต ตคฺครุโก ตพฺพหุโล ตนฺนินฺโน ตปฺโปโณ \csromanfolio{378}ตปฺปพฺภาโร ตทธิมุตฺโต ตทธิปเตยฺโย, ยถา คณครุโก วุจฺจติ คณพาหุลิโกติ, จีวรครุโก วุจฺจติ “จีวรพาหุลิโก”ติ, ปตฺตครุโก วุจฺจติ “ปตฺตพาหุลิโก”ติ, เสนาสนครุโก วุจฺจติ “เสนาสนพาหุลิโก”ติ. เอวเมวํ อิเธกจฺโจ ปญฺญาครุโก โหติ ปญฺญาจริโต ปญฺญาสโย ปญฺญาธิมุตฺโต ปญฺญาธโช ปญฺญาเกตุ ปญฺญา\-ธิปเตยฺโย วิจยพหุโล ปวิจยพหุโล โอกฺขายน\-พหุโล สโมกฺขายน\-พหุโล สมฺเปกฺขายน\-ธมฺโม วิภูตวิหารี ตจฺจริโต ตคฺครุโก ตพฺพหุโล ตนฺนินฺโน ตปฺโปโณ ตปฺปพฺภาโร ตทธิมุตฺโต ตทธิปเตยฺโย, ปญฺญาพาหุลฺลาย สํวตฺตนฺตีติ อิทํ ปญฺญาพาหุลฺลํ. (10)}

\prose{สีฆปญฺญตาย สํวตฺตนฺตีติ กตมา สีฆปญฺญา? สีฆํ สีฆํ สีลานิ ปริปูเรตีติ สีฆปญฺญา, สีฆํ สีฆํ อินฺทฺริย\-สํวรํ ปริปูเรตีติ สีฆปญฺญา, สีฆํ สีฆํ โภชเน มตฺตญฺญุตํ ปริปูเรตีติ สีฆปญฺญา, สีฆํ สีฆํ ชาคริยานุโยคํ ปริปูเรตีติ สีฆปญฺญา, สีฆํ สีฆํ สีลกฺขนฺธํ ปริปูเรตีติ สีฆปญฺญา, สีฆํ สีฆํ สมาธิ\-กฺขนฺธํ ปริปูเรตีติ สีฆปญฺญา, สีฆํ สีฆํ ปญฺญากฺขนฺธํ ปริปูเรตีติ สีฆปญฺญา, สีฆํ สีฆํ วิมุตฺติ\-กฺขนฺธํ ปริปูเรตีติ สีฆปญฺญา, สีฆํ สีฆํ วิมุตฺติ\-ญาณ\-ทสฺสนกฺขนฺธํ ปริปูเรตีติ สีฆปญฺญา, สีฆํ สีฆํ ฐานาฏฺฐานานิ ปฏิวิชฺฌตีติ สีฆปญฺญา, สีฆํ สีฆํ วิหาร\-สมาปตฺติโย ปริปูเรตีติ สีฆปญฺญา, สีฆํ สีฆํ อริยสจฺจานิ ปฏิวิชฺฌตีติ สีฆปญฺญา, สีฆํ สีฆํ สติปฏฺฐาเน ภาเวตีติ สีฆปญฺญา, สีฆํ สีฆํ สมฺมปฺปธาเน ภาเวตีติ สีฆปญฺญา, สีฆํ สีฆํ อิทฺธิปาเท ภาเวตีติ สีฆปญฺญา, สีฆํ สีฆํ อินฺทฺริยานิ ภาเวตีติ สีฆปญฺญา, สีฆํ สีฆํ พลานิ ภาเวตีติ สีฆปญฺญา, สีฆํ สีฆํ โพชฺฌงฺเค ภาเวตีติ สีฆปญฺญา, สีฆํ สีฆํ อริยมคฺคํ ภาเวตีติ สีฆปญฺญา, สีฆํ สีฆํ สามญฺญผลานิ สจฺฉิกโรตีติ สีฆปญฺญา, สีฆํ สีฆํ อภิญฺญาโย ปฏิวิชฺฌตีติ สีฆปญฺญา, สีฆํ สีฆํ ปรมตฺถํ นิพฺพานํ สจฺฉิกโรตีติ สีฆปญฺญา, สีฆปญฺญตาย สํวตฺตนฺตีติ อยํ สีฆปญฺญา. (11)}

\prose{ลหุปญฺญตาย สํวตฺตนฺตีติ กตมา ลหุปญฺญา? ลหุํ ลหุํ สีลานิ ปริปูเรตีติ ลหุปญฺญา, ลหุํ ลหุํ อินฺทฺริย\-สํวรํ ปริปูเรตีติ ลหุปญฺญา, ลหุํ ลหุํ โภชเน มตฺตญฺญุตํ ปริปูเรตีติ ลหุปญฺญา, ลหุํ ลหุํ ชาคริยานุโยคํ ปริปูเรตีติ ลหุปญฺญา, ลหุํ ลหุํ สีลกฺขนฺธํ ฯเปฯ}

\prose{\csromanfolio{379}สมาธิ\-กฺขนฺธํ. ปญฺญากฺขนฺธํ. วิมุตฺติ\-กฺขนฺธํ. วิมุตฺติ\-ญาณ\-ทสฺสนกฺขนฺธํ ปริปูเรตีติ ลหุปญฺญา, ลหุํ ลหุํ ฐานาฏฺฐานานิ ปฏิวิชฺฌตีติ ลหุปญฺญา, ลหุํ ลหุํ วิหาร\-สมาปตฺติโย ปริปูเรตีติ ลหุปญฺญา, ลหุํ ลหุํ อริยสจฺจานิ ปฏิวิชฺฌตีติ ลหุปญฺญา, ลหุํ ลหุํ สติปฏฺฐาเน ภาเวตีติ ลหุปญฺญา, ลหุํ ลหุํ สมฺมปฺปธาเน ภาเวตีติ ลหุปญฺญา, ลหุํ ลหุํ อิทฺธิปาเท ภาเวตีติ ลหุปญฺญา, ลหุํ ลหุํ อินฺทฺริยานิ ภาเวตีติ ลหุปญฺญา, ลหุํ ลหุํ พลานิ ภาเวตีติ ลหุปญฺญา, ลหุํ ลหุํ โพชฺฌงฺเค ภาเวตีติ ลหุปญฺญา, ลหุํ ลหุํ อริยมคฺคํ ภาเวตีติ ลหุปญฺญา, ลหุํ ลหุํ สามญฺญผลานิ สจฺฉิกโรตีติ ลหุปญฺญา, ลหุํ ลหุํ อภิญฺญาโย ปฏิวิชฺฌตีติ ลหุปญฺญา, ลหุํ ลหุํ ปรมตฺถํ นิพฺพานํ สจฺฉิกโรตีติ ลหุปญฺญา, ลหุปญฺญตาย สํวตฺตนฺตีติ อยํ ลหุปญฺญา. (12)}

\prose{หาสปญฺญตาย สํวตฺตนฺตีติ กตมา หาสปญฺญา? อิเธกจฺโจ หาสพหุโล เวทพหุโล ตุฏฺฐิพหุโล ปาโมชฺชพหุโล สีลานิ ปริปูเรตีติ หาสปญฺญา, หาสพหุโล เวทพหุโล ตุฏฺฐิพหุโล ปาโมชฺชพหุโล อินฺทฺริย\-สํวรํ ปริปูเรตีติ หาสปญฺญา, หาสพหุโล เวทพหุโล ตุฏฺฐิพหุโล ปาโมชฺชพหุโล โภชเน มตฺตญฺญุตํ ปริปูเรตีติ หาสปญฺญา, หาสพหุโล เวทพหุโล ตุฏฺฐิพหุโล ปาโมชฺชพหุโล ชาคริยานุโยคํ ปริปูเรตีติ หาสปญฺญา, หาสพหุโล เวทพหุโล ตุฏฺฐิพหุโล ปาโมชฺชพหุโล สีลกฺขนฺธํ ฯเปฯ สมาธิ\-กฺขนฺธํ. ปญฺญากฺขนฺธํ. วิมุตฺติ\-กฺขนฺธํ. วิมุตฺติ\-ญาณ\-ทสฺสนกฺขนฺธํ ปริปูเรตีติ ฯเปฯ ฐานาฏฺฐานานิ ปฏิวิชฺฌตีติ. วิหาร\-สมาปตฺติโย ปริปูเรตีติ. อริยสจฺจานิ ปฏิวิชฺฌตีติ. สติปฏฺฐาเน ภาเวตีติ. สมฺมปฺปธาเน ภาเวตีติ. อิทฺธิปาเท ภาเวตีติ. อินฺทฺริยานิ ภาเวตีติ. พลานิ ภาเวตีติ. โพชฺฌงฺเค ภาเวตีติ. อริยมคฺคํ ภาเวตีติ ฯเปฯ สามญฺญผลานิ สจฺฉิกโรตีติ หาสปญฺญา, หาสพหุโล เวทพหุโล ตุฏฺฐิพหุโล ปาโมชฺชพหุโล อภิญฺญาโย ปฏิวิชฺฌตีติ หาสปญฺญา, หาสพหุโล เวทพหุโล ตุฏฺฐิพหุโล ปาโมชฺชพหุโล ปรมตฺถํ นิพฺพานํ สจฺฉิกโรตีติ หาสปญฺญา, หาสปญฺญตาย สํวตฺตนฺตีติ อยํ หาสปญฺญา. (13)}

\proseitem{7}{ชวน\-ปญฺญตาย สํวตฺตนฺตีติ กตมา ชวนปญฺญา? ยํ กิญฺจิ รูปํ อตีตา\-นาคต\-ปจฺจุปฺปนฺนํ อชฺฌตฺตํ วา พหิทฺธา วา โอฬาริกํ วา สุขุมํ วา หีนํ วา \csromanfolio{380}ปณีตํ วา ยํ ทูเร สนฺติเก วา, สพฺพํ รูปํ อนิจฺจโต ขิปฺปํ ชวตีติ ชวนปญฺญา, ทุกฺขโต ขปฺปํ ชวตีติ ชวนปญฺญา, อนตฺตโต ขิปฺปํ ชวตีติ ชวนปญฺญา. ยา กาจิ เวทนา ฯเปฯ. ยา กาจิ สญฺญา. เย เกจิ สงฺขารา. ยํ กิญฺจิ วิญฺญาณํ อตีตา\-นาคต\-ปจฺจุปฺปนฺนํ อชฺฌตฺตํ วา พหิทฺธา วา โอฬาริกํ วา สุขุมํ วา หีนํ วา ปณีตํ วา ยํ ทูเร สนฺติเก วา, สพฺพํ วิญฺญาณํ อนิจฺจโต ขิปฺปํ ชวตีติ ชวนปญฺญา, ทุกฺขโต ขิปฺปํ ชวตีติ ชวนปญฺญา, อนตฺตโต ขิปฺปํ ชวตีติ ชวนปญฺญา. จกฺขุ ฯเปฯ. ชรามรณํ อตีตา\-นาคต\-ปจฺจุปฺปนฺนํ อนิจฺจโต ขิปฺปํ ชวตีติ ชวนปญฺญา, ทุกฺขโต ขิปฺปํ ชวตีติ ชวนปญฺญา, อนตฺตโต ขิปฺปํ ชวตีติ ชวนปญฺญา.}

\prose{รูปํ อตีตา\-นาคต\-ปจฺจุปฺปนฺนํ อนิจฺจํ ขยฏฺเฐน ทุกฺขํ ภยฏฺเฐน อนตฺตา อสารกฏฺเฐนาติ ตุลยิตฺวา ตีรยิตฺวา วิภาวยิตฺวา วิภูตํ กตฺวา รูปนิโรเธ นิพฺพาเน ขิปฺปํ ชวตีติ ชวนปญฺญา.  เวทนา ฯเปฯ. สญฺญา. สงฺขารา. วิญฺญาณํ. จกฺขุ ฯเปฯ. ชรามรณํ อตีตา\-นาคต\-ปจฺจุปฺปนฺนํ อนิจฺจํ ขยฏฺเฐน ทุกฺขํ ภยฏฺเฐน อนตฺตา อสารกฏฺเฐนาติ ตุลยิตฺวา ตีรยิตฺวา วิภาวยิตฺวา วิภูตํ กตฺวา ชารามรณนิโรเธ\csromansharedfootnote{89971b43ef828230}{ชรามรณํ รูปนิโรเธ (สฺยา)} นิพฺพาเน ขิปฺปํ ชวตีติ ชวนปญฺญา.}

\prose{รูปํ อตีตา\-นาคต\-ปจฺจุปฺปนฺนํ อนิจฺจํ สงฺขตํ ปฏิจฺจ\-สมุปฺปนฺนํ ขยธมฺมํ วยธมฺมํ วิราคธมฺมํ นิโรธ\-ธมฺมนฺติ ตุลยิตฺวา ตีรยิตฺวา วิภาวยิตฺวา วิภูตํ กตฺวา รูปนิโรเธ นิพฺพาเน ขิปฺปํ ชวตีติ ชวนปญฺญา.  เวทนา ฯเปฯ. สญฺญา. สงฺขารา. วิญฺญาณํ. จกฺขุ ฯเปฯ. ชรามรณํ อตีตา\-นาคต\-ปจฺจุปฺปนฺนํ อนิจฺจํ สงฺขตํ ปฏิจฺจ\-สมุปฺปนฺนํ ขยธมฺมํ วยธมฺมํ วิราคธมฺมํ นิโรธ\-ธมฺมนฺติ ตุลยิตฺวา ตีรยิตฺวา วิภาวยิตฺวา วิภูตํ กตฺวา ชรามรณ\-นิโรเธ นิพฺพาเน ขิปฺปํ ชวตีติ ชวนปญฺญา, ชวน\-ปญฺญตาย สํวตฺตนฺตีติ อยํ ชวนปญฺญา. (14)}

\prose{ติกฺข\-ปญฺญตาย สํวตฺตนฺตีติ กตมา ติกฺขปญฺญา? ขิปฺปํ กิเลเส ฉินฺทตีติ ติกฺขปญฺญา, อุปฺปนฺนํ กามวิตกฺกํ นาธิวาเสติ ปชหติ วิโนเทติ พฺยนฺตีกโรติ\csromansharedfootnote{349b6ff3fa54a80e}{พฺยนฺติํกโรติ (ก)} อนภาวํ คเมตีติ ติกฺขปญฺญา, อุปฺปนฺนํ พฺยาปาท\-วิตกฺกํ นาธิวาเสติ ปชหติ วิโนเทติ พฺยนฺตีกโรติ อนภาวํ คเมตีติ ติกฺขปญฺญา, อุปฺปนฺนํ วิหิํสา\-วิตกฺกํ นาธิวาเสติ ฯเปฯ อุปฺปนฺนุปฺปนฺเน ปาปเก อกุสเล ธมฺเม นาธิวาเสติ ปชหติ วิโนเทติ พฺยนฺตีกโรติ}

\prose{\csromanfolio{381}อนภาวํ คเมตีติ ติกฺขปญฺญา, อุปฺปนฺนํ ราคํ นาธิวาเสติ ปชหติ วิโนเทติ พฺยนฺตีกโรติ อนภาวํ คเมตีติ ติกฺขปญฺญา. อุปฺปนฺนํ โทสํ ฯเปฯ. อุปฺปนฺนํ โมหํ. อุปฺปนฺนํ โกธํ. อุปฺปนฺนํ อุปนาหํ. มกฺขํ. ปฬาสํ. อิสฺสํ. มจฺฉริยํ. มายํ. สาเฐยฺยํ. ถมฺภํ. สารมฺภํ. มานํ. อติมานํ. มทํ. ปมาทํ. สพฺเพ กิเลเส. สพฺเพ ทุจฺจริเต. สพฺเพ อภิสงฺขาเร ฯเปฯ. สพฺเพ ภวคามิกมฺเม นาธิวาเสติ ปชหติ วิโนเทติ พฺยนฺตีกโรติ อนภาวํ คเมตีติ ติกฺขปญฺญา. เอกสฺมิํ อาสเน จตฺตาโร จ อริยมคฺคา จตฺตาริ จ สามญฺญผลานิ จตสฺโส ปฏิสมฺภิทาโย ฉ อภิญฺญาโย อธิคตา โหนฺติ สจฺฉิกตา ผสฺสิตา ปญฺญายาติ ติกฺขปญฺญา, ติกฺข\-ปญฺญตาย สํวตฺตนฺตีติ อยํ ติกฺขปญฺญา. (15)}

\prose{นิพฺเพธิก\-ปญฺญตาย สํวตฺตนฺตีติ กตมา นิพฺเพธิก\-ปญฺญา? อิเธกจฺโจ สพฺพ\-สงฺขาเรสุ อุพฺเพคพหุโล โหติ อุตฺตาสพหุโล อุกฺกณฺฐน\-พหุโล อรติพหุโล อนภิรติ\-พหุโล พหิมุโข น รมติ สพฺพ\-สงฺขาเรสุ, อนิพฺพิทฺธปุพฺพํ อปฺปทาลิตปุพฺพํ โลภกฺขนฺธํ นิพฺพิชฺฌติ ปทาเลตีติ นิพฺเพธิก\-ปญฺญา, อนิพฺพิทฺธปุพฺพํ อปฺปทาลิตปุพฺพํ โทสกฺขนฺธํ นิพฺพิชฺฌติ ปทาเลตีติ นิพฺเพธิก\-ปญฺญา, อนิพฺพิทฺธปุพฺพํ อปฺปทาลิตปุพฺพํ โมหกฺขนฺธํ นิพฺพิชฺฌติ ปทาเลตีติ นิพฺเพธิก\-ปญฺญา, อนิพฺพิทฺธปุพฺพํ อปฺปทาลิตปุพฺพํ โกธํ ฯเปฯ. อุปนาหํ. มกฺขํ. ปฬาสํ. อิสฺสํ. มจฺฉริยํ. มายํ. สาเฐยฺยํ. ถมฺภํ. สารมฺภํ. มานํ. อติมานํ. มทํ. ปมาทํ. สพฺเพ กิเลเส. สพฺเพ ทุจฺจริเต. สพฺเพ อภิสงฺขาเร ฯเปฯ. สพฺเพ ภวคามิกมฺเม นิพฺพิชฺฌติ ปทาเลตีติ นิพฺเพธิก\-ปญฺญา, นิพฺเพธิก\-ปญฺญตาย สํวตฺตนฺตีติ อยํ นิพฺเพธิก\-ปญฺญา. (16)}

\prose{อิมา โสฬส ปญฺญาโย. อิมาหิ โสฬสหิ ปญฺญาหิ สมนฺนาคโต ปุคฺคโล ปฏิสมฺภิท\-ปฺปตฺโต.}

\csromanmatikaanchor{mtk.121}
\csromantocmarkref{subsection}{2. ปุคฺคลวิเสสนิทฺเทส}{mtk.121}
\bookmark[dest={mtk.121},level=2]{2. ปุคฺคลวิเสสนิทฺเทส}
\csromanheader{2}{2. ปุคฺคล\-วิเส\-สนิทฺเทส}
\proseitem{8}{เทฺว ปุคฺคลา ปฏิสมฺภิท\-ปฺปตฺตา. เอโก ปุพฺพโยค\-สมฺปนฺโน, เอโก น ปุพฺพโยค\-สมฺปนฺโน. โย ปุพฺพโยค\-สมฺปนฺโน, โส เตน อติเรโก โหติ, อธิโก โหติ, วิเสโส โหติ, ตสฺส ญาณํ ปภิชฺชติ.}

\prose{เทฺว ปุคฺคลา ปฏิสมฺภิท\-ปฺปตฺตา. เทฺวปิ ปุพฺพโยค\-สมฺปนฺนา. เอโก พหุสฺสุโต, เอโก น พหุสฺสุโต. โย พหุสฺสุโต, โส เตน \csromanfolio{382}อติเรโก โหติ, อธิโก โหติ, วิเสโส โหติ, ตสฺส ญาณํ ปภิชฺชติ.}

\prose{เทฺว ปุคฺคลา ปฏิสมฺภิท\-ปฺปตฺตา. เทฺวปิ ปุพฺพโยค\-สมฺปนฺนา. เทฺวปิ พหุสฺสุตา. เอโก เทสนาพหุโล, เอโก น เทสนาพหุโล.โย เทสนาพหุโล, โส เตน อติเรโก โหติ, อธิโก โหติ, วิเสโส โหติ, ตสฺส ญาณํ ปภิชฺชติ.}

\prose{เทฺว ปุคฺคลา ปฏิสมฺภิท\-ปฺปตฺตา. เทฺวปิ ปุพฺพโยค\-สมฺปนฺนา. เทฺวปิ พหุสฺสุตา. เทฺวปิ เทสนาพหุลา. เอโก ครูปนิสฺสิโต, เอโก น ครูปนิสฺสิโต. โย ครูปนิสฺสิโต, โส เตน อติเรโก โหติ, อธิโก โหติ, วิเสโส โหติ, ตสฺส ญาณํ ปภิชฺชติ.}

\prose{เทฺว ปุคฺคลา ปฏิสมฺภิท\-ปฺปตฺตา. เทฺวปิ ปุพฺพโยค\-สมฺปนฺนา. เทฺวปิ พหุสฺสุตา. เทฺวปิ เทสนาพหุลา. เทฺวปิ ครูปนิสฺสิตา. เอโก วิหารพหุโล, เอโก น วิหารพหุโล. โย วิหารพหุโล, โส เตน อติเรโก โหติ, อธิโก โหติ, วิเสโส โหติ, ตสฺส ญาณํ ปภิชฺชติ.}

\prose{เทฺว ปุคฺคลา ปฏิสมฺภิท\-ปฺปตฺตา. เทฺวปิ ปุพฺพโยค\-สมฺปนฺนา. เทฺวปิ พหุสฺสุตา. เทฺวปิ เทสนาพหุลา. เทฺวปิ ครูปนิสฺสิตา. เทฺวปิ วิหารพหุลา. เอโก ปจฺจเวกฺขณา\-พหุโล, เอโก น ปจฺจเวกฺขณา\-พหุโล. โย ปจฺจเวกฺขณา\-พหุโล, โส เตน อติเรโก โหติ, อธิโก โหติ, วิเสโส โหติ, ตสฺส ญาณํ ปภิชฺชติ.}

\prose{เทฺว ปุคฺคลา ปฏิสมฺภิท\-ปฺปตฺตา. เทฺวปิ ปุพฺพโยค\-สมฺปนฺนา. เทฺวปิ พหุสฺสุตา. เทฺวปิ เทสนาพหุลา. เทฺวปิ ครูปนิสฺสิตา. เทฺวปิ วิหารพหุลา. เทฺวปิ ปจฺจเวกฺขณา\-พหุลา. เอโก เสข\-ปฏิสมฺภิท\-ปฺปตฺโต, เอโก อเสข\-ปฏิสมฺภิท\-ปฺปตฺโต. โย อเสข\-ปฏิสมฺภิท\-ปฺปตฺโต, โส เตน อติเรโก โหติ, อธิโก โหติ, วิเสโส โหติ, ตสฺส ญาณํ ปภิชฺชติ.}

\prose{เทฺว ปุคฺคลา ปฏิสมฺภิท\-ปฺปตฺตา. เทฺวปิ ปุพฺพโยค\-สมฺปนฺนา. เทฺวปิ พหุสฺสุตา. เทฺวปิ เทสนาพหุลา. เทฺวปิ ครูปนิสฺสิตา. เทฺวปิ วิหารพหุลา. เทฺวปิ ปจฺจเวกฺขณา\-พหุลา. เทฺวปิ อเสข\-ปฏิสมฺภิท\-ปฺปตฺตา. เอโก สาวก\-ปารมิปฺปตฺโต, เอโก น สาวก\-ปารมิปฺปตฺโต. โย สาวก\-ปารมิปฺปตฺโต,}

\csromanmatikaanchor{mtk.122}
\csromantocmarkref{section}{2. อิทฺธิกถา}{mtk.122}
\bookmark[dest={mtk.122},level=1]{2. อิทฺธิกถา}
\prose{\csromanfolio{383}โส เตน อติเรโก โหติ, อธิโก โหติ, วิเสโส โหติ, ตสฺส ญาณํ ปภิชฺชติ.}

\prose{เทฺว ปุคฺคลา ปฏิสมฺภิท\-ปฺปตฺตา. เทฺวปิ ปุพฺพโยค\-สมฺปนฺนา. เทฺวปิ พหุสฺสุตา. เทฺวปิ เทสนาพหุลา. เทฺวปิ ครูปนิสฺสิตา. เทฺวปิ วิหารพหุลา. เทฺวปิ ปจฺจเวกฺขณา\-พหุลา. เทฺวปิ อเสข\-ปฏิสมฺภิท\-ปฺปตฺตา. เอโก สาวก\-ปารมิปฺปตฺโต, เอโก ปจฺเจก\-สมฺพุทฺโธ. โย ปจฺเจก\-สมฺพุทฺโธ, โส เตน อติเรโก โหติ, อธิโก โหติ, วิเสโส โหติ, ตสฺส ญาณํ ปภิชฺชติ.}

\prosewithcloserruled{ปจฺเจกพุทฺธญฺจ สเทวกญฺจ โลกํ อุปาทาย ตถาคโต อรหํ สมฺมาสมฺพุทฺโธ อคฺโค ปฏิสมฺภิท\-ปฺปตฺโต ปญฺญาปเภท\-กุสโล ปภินฺนญาโณ อธิคตปฏิสมฺภิโท จตุเวสารชฺช\-ปฺปตฺโต ทสพลธารี ปุริสาสโภ ปุริสสีโห ฯเปฯ. เยปิ เต ขตฺติย\-ปณฺฑิตา พฺราหฺมณ\-ปณฺฑิตา คหปติ\-ปณฺฑิตา สมณปณฺฑิตา นิปุณา กตปรปฺปวาทา วาลเวธิรูปา โวภินฺทนฺตา มญฺเญ จรนฺติ ปญฺญาคเตน ทิฏฺฐิคตานิ. เต ปญฺหํ อภิส\-งฺขริ\-ตฺวา อภิสงฺขาริตฺวา ตถาคตํ อุปสงฺกมิตฺวา ปุจฺฉนฺติ คูฬฺหานิ จ ปฏิจฺฉนฺนานิ จ, กถิตา วิสชฺชิตา จ เต ปญฺหา ภควตา โหนฺติ นิทฺทิฏฺฐ\-การณา, อุปกฺขิตฺตกา จ เต ภควโต สมฺปชฺชนฺติ. อถ โข ภควา ตตฺถ อติโรจติ ยทิทํ ปญฺญายาติ อคฺโค ปฏิสมฺภิท\-ปฺปตฺโตติ.}{มหาปญฺญากถา นิฏฺฐิตา.}
\proseitem{9}{กา อิทฺธิ. กติ อิทฺธิโย. อิทฺธิยา กติ ภูมิโย. กติ ปาทา. กติ ปทานิ. กติ มูลานิ? \textbf{กา อิทฺธี}ติ อิชฺฌนฏฺเฐน อิทฺธิ.  กติ อิทฺธิโยติ ทส อิทฺธิโย.  \textbf{อิทฺธิยา กติ ภูมิโย}ติ อิทฺธิยา จตสฺโส ภูมิโย.  จตฺตาโร ปาทา.  อฏฺฐ ปทานิ.  โสฬส มูลานิ.}

\proseitem{10}{กตมา ทส อิทฺธิโย? อธิฏฺฐานา อิทฺธิ, วิกุพฺพนา อิทฺธิ, มโนมยา อิทฺธิ ญาณวิปฺผารา อิทฺธิ, สมาธิ\-วิปฺผารา อิทฺธิ, อริยา อิทฺธิ, กมฺมวิปากชา อิทฺธิ, \csromanfolio{384}ปุญฺญวโต อิทฺธิ, วิชฺชามยา อิทฺธิ, ตตฺถ ตตฺถ สมฺมา\-ปโยค\-ปจฺจยา\csromansharedfootnote{debbc82d538cb5cb}{สมฺมปฺปโยคปจฺจยา (สฺยา, ก)} อิชฺฌนฏฺเฐน อิทฺธิ.}

\prose{อิทฺธิยา กตมา จตสฺโส ภูมิโย? วิเวกชาภูมิ ปฐมํ ฌานํ, ปีติสุขภูมิ ทุติยํ ฌานํ, อุเปกฺขาสุขภูมิ ตติยํ ฌานํ, อทุกฺขมสุขาภูมิ จตุตฺถํ ฌานํ. อิทฺธิยา อิมา จตสฺโส ภูมิโย อิทฺธิลาภาย อิทฺธิ\-ปฏิลาภาย อิทฺธิ\-วิกุพฺพนตาย อิทฺธิ\-วิสวิตาย อิทฺธิ\-วสีภาวาย อิทฺธิ\-เวสารชฺชาย สํวตฺตนฺตีติ.}

\prose{อิทฺธิยา กตเม จตฺตาโร ปาทา? อิธ ภิกฺขุ ฉนฺท\-สมาธิ\-ปธาน\-สงฺขาร\-สมนฺนาคตํ อิทฺธิปาทํ ภาเวติ, จิตฺต\-สมาธิ\-ปธาน\-สงฺขาร\-สมนฺนาคตํ อิทฺธิปาทํ ภาเวติ, วีริย\-สมาธิ\-ปธาน\-สงฺขาร\-สมนฺนาคตํ อิทฺธิปาทํ ภาเวติ, วีมํสา\-สมาธิ\-ปธาน\-สงฺขาร\-สมนฺนาคตํ อิทฺธิปาทํ ภาเวติ. อิทฺธิยา อิเม จตฺตาโร ปาทา อิทฺธิลาภาย อิทฺธิ\-ปฏิลาภาย อิทฺธิ\-วิกุพฺพนตาย อิทฺธิ\-วิสวิตาย อิทฺธิ\-วสีภาวาย อิทฺธิ\-เวสารชฺชาย สํวตฺตนฺตีติ.}

\prose{อิทฺธิยา กตมานิ อฏฺฐ ปทานิ? ฉนฺทํ เจ ภิกฺขุ นิสฺสาย ลภติ สมาธิํ, ลภติ จิตฺตสฺส เอกคฺคตํ, ฉนฺโท น สมาธิ, สมาธิ น ฉนฺโท. อญฺโญ ฉนฺโท, อญฺโญ สมาธิ. วีริยํ เจ ภิกฺขุ นิสฺสาย ลภติ สมาธิํ, ลภติ จิตฺตสฺส เอกคฺคตํ. วีริยํ น สมาธิ, สมาธิ น วีริยํ. อญฺญํ วีริยํ, อญฺโญ สมาธิ. จิตฺตํ เจ ภิกฺขุ นิสฺสาย ลภติ สมาธิํ, ลภติ จิตฺตสฺส เอกคฺคตํ. จิตฺตํ น สมาธิ, สมาธิ น จิตฺตํ. อญฺญํ จิตฺตํ, อญฺโญ สมาธิ. วีมํสํ เจ ภิกฺขุ นิสฺสาย ลภติ สมาธิํ, ลภติ จิตฺตสฺส เอกคฺคตํ. วีมํสา น สมาธิ, สมาธิ น วีมํสา. อญฺญา วีมํสา, อญฺโญ สมาธิ. อิทฺธิยา อิมานิ อฏฺฐ ปทานิ อิทฺธิลาภาย อิทฺธิ\-ปฏิลาภาย อิทฺธิ\-วิกุพฺพนตาย อิทฺธิ\-วิสวิตาย อิทฺธิ\-วสีภาวาย อิทฺธิ\-เวสารชฺชาย สํวตฺตนฺตีติ.}

\prose{\textbf{อิทฺธิยา กตมานิ โสฬส มูลานิ?} อโนนตํ\csromansharedfootnote{c65c40770860ea06}{อโนณตํ (ก)} จิตฺตํ โกสชฺเช น อิญฺชตีติ อาเนญฺชํ, อนุนฺนตํ\csromansharedfootnote{b8b889d6dc4a4696}{อนุณฺณตํ (ก)} จิตฺตํ อุทฺธจฺเจ น อิญฺชตีติ อาเนญฺชํ, อนภินตํ จิตฺตํ ราเค น อิญฺชตีติ อาเนญฺชํ, อนปนตํ จิตฺตํ พฺยาปาเท น อิญฺชตีติ อาเนญฺชํ, อนิสฺสิตํ จิตฺตํ ทิฏฺฐิยา น อิญฺชตีติ อาเนญฺชํ, อปฺปฏิพทฺธํ จิตฺตํ ฉนฺทราเค น อิญฺชตีติ อาเนญฺชํ, วิปฺปมุตฺตํ จิตฺตํ กามราเค น อิญฺชตีติ อาเนญฺชํ, วิสญฺญุตฺตํ จิตฺตํ กิเลเส น อิญฺชตีติ อาเนญฺชํ, วิมริยาทิกตํ จิตฺตํ}

\csromanmatikaanchor{mtk.123}
\csromantocmarkref{section}{ทสอิทฺธินิทฺเทส}{mtk.123}
\bookmark[dest={mtk.123},level=1]{ทสอิทฺธินิทฺเทส}
\prose{\csromanfolio{385}กิเลสมริยาเท น อิญฺชตีติ อาเนญฺชํ, เอกตฺตคตํ\csromansharedfootnote{5166cc583c34c2b4}{เอกคฺคตํ (สฺยา)} จิตฺตํ นานตฺต\-กิเลเสหิ\csromansharedfootnote{0147577b812fb2b7}{นานตฺตกิเลเส (สฺยา)} น อิญฺชตีติ อาเนญฺชํ, สทฺธาย ปริคฺคหิตํ จิตฺตํ อสฺสทฺธิเย น อิญฺชตีติ อาเนญฺชํ, วีริเยน ปริคฺคหิตํ จิตฺตํ โกสชฺเช น อิญฺชตีติ อาเนญฺชํ, สติยา ปริคฺคหิตํ จิตฺตํ ปมาเท น อิญฺชตีติ อาเนญฺชํ, สมาธินา ปริคฺคหิตํ จิตฺตํ อุทฺธจฺเจ น อิญฺชตีติ อาเนญฺชํ, ปญฺญาย ปริคฺคหิตํ จิตฺตํ อวิชฺชาย น อิญฺชตีติ อาเนญฺชํ, โอภาสคตํ จิตฺตํ อวิชฺชนฺธกาเร น อิญฺชตีติ อาเนญฺชํ. อิทฺธิยา อิมานิ โสฬส มูลานิ อิทฺธิลาภาย อิทฺธิ\-ปฏิลาภาย อิทฺธิ\-วิกุพฺพนตาย อิทฺธิ\-วิสวิตาย อิทฺธิ\-วสีภาวาย อิทฺธิ\-เวสารชฺชาย สํวตฺตนฺตีติ.}

\prose{ทสอิทฺธินิทฺเทส}

\proseitem{10}{กตมา อธิฏฺฐานา อิทฺธิ? อิธ ภิกฺขุ อเนกวิหิตํ อิทฺธิวิธํ ปจฺจนุโภติ, เอโกปิ หุตฺวา พหุธา โหติ, พหุธาปิ หุตฺวา เอโก โหติ, อาวิภาวํ ติโรภาวํ ติโรกุฏฺฏํ ติโรปาการํ ติโรปพฺพตํ อสชฺชมาโน คจฺฉติ เสยฺยถาปิ อากาเส.ปถวิยาปิ อุมฺมุชฺช\-นิมุชฺชํ กโรติ เสยฺยถาปิ อุทเก. อุทเกปิ อภิชฺชมาเน\csromansharedfootnote{867d18ff3d461836}{อภิชฺชมาโน (ก) ที 1. 73 ปิฏฺเฐ ปสฺสิตพฺพา.} คจฺฉติ เสยฺยถาปิ ปถวิยํ. อากาเสปิ ปลฺลงฺเกน กมติ เสยฺยถาปิ ปกฺขี สกุโณ. อิเมปิ จนฺทิมสูริเย เอวํมหิทฺธิเก เอวํ\-มหา\-นุภาเว ปาณินา ปรามสติ ปริมชฺชติ, ยาว พฺรหฺมโลกาปิ กาเยน วสํ วตฺเตตีติ.}

\prose{อิธาติ อิมิสฺสา ทิฏฺฐิยา อิมิสฺสา ขนฺติยา อิมิสฺสา รุจิยา อิมสฺมิํ อาทาเย อิมสฺมิํ ธมฺเม อิมสฺมิํ วินเย อิมสฺมิํ ธมฺมวินเย อิมสฺมิํ ปาวจเน อิมสฺมิํ พฺรหฺมจริเย อิมสฺมิํ สตฺถุสาสเน, เตน วุจฺจติ อิธาติ. ภิกฺขูติ ปุถุชฺชน\-กลฺยาณโก วา โหติ ภิกฺขุ เสโข วา อรหา วา อกุปฺปธมฺโม. อเนกวิหิตํ อิทฺธิวิธํ ปจฺจนุโภตีติ นานปฺปการํ อิทฺธิวิธํ ปจฺจนุโภติ. เอโกปิ หุตฺวา พหุธา โหตีติ ปกติยา เอโก พหุกํ อาวชฺชติ, สตํ วา สหสฺสํ วา สตสหสฺสํ วา อาวชฺชติ, อาวชฺชิตฺวา ญาเณน อธิฏฺฐาติ “พหุโล โหมี”ติ\csromansharedfootnote{0cdcac8e1c49aa00}{โหตูติ (ก)}, พหุโล โหติ. ยถายสฺมา จูฬปนฺถโก\csromansharedfootnote{9b06c5f855b9abfe}{จุลฺลปนฺถโก (สฺยา)} เอโกปิ หุตฺวา พหุธา โหติ. เอวเมวํ \csromanfolio{386}โส อิทฺธิมา เจโตวสิปฺปตฺโต เอโกปิ หุตฺวา พหุธา โหติ. \textbf{พหุธาปิ หุตฺวา เอโก โหตี}ติ ปกติยา พหุโก เอกํ อาวชฺชติ, อาวชฺชิตฺวา ญาเณน อธิฏฺฐาติ “เอโก โหมี”ติ, เอโก โหติ. ยถายสฺมา จูฬปนฺถโก พหุธาปิ หุตฺวา เอโก โหติ. เอวเมวํ โส อิทฺธิมา เจโตวสิปฺปตฺโต พหุธาปิ หุตฺวา เอโก โหติ.}

\proseitem{11}{อาวิภาวนฺติ เกนจิ อนาวฏํ โหติ อปฺปฏิจฺฉนฺนํ วิวฏํ ปากฏํ. ติโรภาวนฺติ เกนจิ อาวฏํ โหติ ปฏิจฺฉนฺนํ ปิหิตํ ปฏิกุชฺชิตํ. ติโรกุฏฺฏํ ติโรปาการํ ติโรปพฺพตํ อสชฺชมาโน คจฺฉติเสยฺยถาปิ อากาเสติ ปกติยา อากาสกสิณ\-สมาปตฺติยา ลาภี โหติ, ติโรกุฏฺฏํ ติโรปาการํ ติโรปพฺพตํ อาวชฺชติ, อาวชฺชิตฺวา ญาเณน อธิฏฺฐาติ “อากาโส โหตู”ติ, อากาโส โหติ. ติโรกุฏฺฏํ ติโรปาการํ ติโรปพฺพตํ อสชฺชมาโน คจฺฉติ. ยถา มนุสฺสา ปกติยา อนิทฺธิมนฺโต เกนจิ อนาวเฏ อปริกฺขิตฺเต อสชฺชมานา คจฺฉนฺติ. เอวเมวํ โส อิทฺธิมา เจโตวสิปฺปตฺโต ติโรกุฏฺฏํ ติโรปาการํ ติโรปพฺพตํ อสชฺชมาโน คจฺฉติ เสยฺยถาปิ อากาเส.}

\prose{ปถวิยาปิ อุมฺมุชฺช\-นิมุชฺชํ กโรติ เสยฺยถาปิ อุทเกติ ปกติยา อาโปกสิณ\-สมาปตฺติยา ลาภี โหติ. ปถวิํ อาวชฺชติ, อาวชฺชิตฺวา ญาเณน อธิฏฺฐาติ “อุทกํ โหตู”ติ, อุทกํ โหติ. โส ปถวิยา อุมฺมุชฺช\-นิมุชฺชํ กโรติ. ยถา มนุสฺสา ปกติยา อนิทฺธิมนฺโต อุทเก อุมฺมุชฺช\-นิมุชฺชํ กโรนฺติ. เอวเมวํ โส อิทฺธิมา เจโตวสิปฺปตฺโต ปถวิยา อุมฺมุชฺช\-นิมุชฺชํ กโรติ เสยฺยถาปิ อุทเก.}

\prose{อุทเกปิ อภิชฺชมาเน คจฺฉติ เสยฺยถาปิ ปถวิยนฺติ ปกติยา ปถวีกสิณ\-สมาปตฺติยา ลาภี โหติ. อุทกํ อาวชฺชติ, อาวชฺชิตฺวา ญาเณน อธิฏฺฐาติ “ปถวี โหตู”ติ, ปถวี โหติ. โส อภิชฺชมาเน อุทเก คจฺฉติ. ยถา มนุสฺสา ปกติยา อนิทฺธิมนฺโต อภิชฺชมานาย ปถวิยา คจฺฉนฺติ. เอวเมวํ โส อิทฺธิมา เจโตวสิปฺปตฺโต อภิชฺชมาเน อุทเก คจฺฉติ เสยฺยถาปิ ปถวิยํ.}

\prose{อากาเสปิ ปลฺลงฺเกน กมติ เสยฺยถาปิ ปกฺขี สกุโณติ ปกติยา ปถวีกสิณ\-สมาปตฺติยา ลาภี โหติ. อากาสํ}

\prose{\csromanfolio{387}อาวชฺชติ, อาวชฺชิตฺวา ญาเณน อธิฏฺฐาติ “ปถวี โหตู”ติ, ปถวี โหติ. โส อากาเส อนฺตลิกฺเข จงฺกมติปิ ติฏฺฐติปิ นิสีทติปิ เสยฺยมฺปิ กปฺเปติ. ยถา มนุสฺสา ปกติยา อนิทฺธิมนฺโต ปถวิยา จงฺกมนฺติปิ ติฏฺฐนฺติปิ นิสีทนฺติปิ เสยฺยมฺปิ กปฺเปนฺติ. เอวเมวํ โส อิทฺธิมา เจโตวสิปฺปตฺโต อากาเส อนฺตลิกฺเข จงฺกมติปิ ติฏฺฐติปิ นิสีทติปิ เสยฺยมฺปิ กปฺเปติ เสยฺยถาปิ ปกฺขี สกุโณ.}

\proseitem{12}{อิเมปิ จนฺทิมสูริเย เอวํมหิทฺธิเก เอวํ\-มหา\-นุภาเว ปาณินา ปรามสติ ปริมชฺชตีติ อิธ โส อิทฺธิมา เจโตวสิปฺปตฺโต นิสินฺนโก วา นิปนฺนโก วา จนฺทิมสูริเย อาวชฺชติ, อาวชฺชิตฺวา ญาเณน อธิฏฺฐาติ “หตฺถปาเส โหตู”ติ, หตฺถปาเส โหติ. โส นิสินฺนโก วา นิปนฺนโก วา จนฺทิมสูริเย ปาณินา อามสติ ปรามสติ ปริมชฺชติ. ยถา มนุสฺสา ปกติยา อนิทฺธิมนฺโต กิญฺจิเทว รูปคตํ หตฺถปาเส อามสนฺติ ปรามสนฺติ ปริมชฺชนฺติ. เอวเมวํ โส อิทฺธิมา เจโตวสิปฺปตฺโต นิสินฺนโก วา นิปนฺนโก วา จนฺทิมสูริเย ปาณินา อามสติ ปรามสติ ปริมชฺชติ.}

\prose{ยาว พฺรหฺมโลกาปิ กาเยน วสํ วตฺเตตีติ สเจ โส อิทฺธิมา เจโตวสิปฺปตฺโต พฺรหฺมโลกํ คนฺตุกาโม โหติ, ทูเรปิ สนฺติเก อธิฏฺฐาติ “สนฺติเก โหตู”ติ, สนฺติเก โหติ. สนฺติเกปิ ทูเร อธิฏฺฐาติ “ทูเร โหตู”ติ, ทูเร โหติ. พหุกมฺปิ โถกํ อธิฏฺฐาติ “โถกํ โหตู”ติ, โถกํ โหติ. โถกมฺปิ พหุกํ อธิฏฺฐาติ “พหุกํ โหตู”ติ, พหุกํ โหติ. ทิพฺเพน จกฺขุนา ตสฺส พฺรหฺมุโน รูปํ ปสฺสติ, ทิพฺพาย โสตธาตุยา ตสฺส พฺรหฺมุโน สทฺทํ สุณาติ, เจโตปริย\-ญาเณน ตสฺส พฺรหฺมุโน จิตฺตํ ปชานาติ, สเจ โส อิทฺธิมา เจโตวสิปฺปตฺโต ทิสฺสมาเนน กาเยน พฺรหฺมโลกํ คนฺตุกาโม โหติ, กายวเสน จิตฺตํ ปริณาเมติ, กายวเสน จิตฺตํ อธิฏฺฐาติ. กายวเสน จิตฺตํ ปริณาเมตฺวา กายวเสน จิตฺตํ อธิฏฺฐหิตฺวา สุขสญฺญญฺจ ลหุสญฺญญฺจ โอกฺกมิตฺวา ทิสฺสมาเนน กาเยน พฺรหฺมโลกํ คจฺฉติ. สเจ โส อิทฺธิมา เจโตวสิปฺปตฺโต อทิสฺสมาเนน กาเยน พฺรหฺมโลกํ คนฺตุกาโม โหติ, จิตฺตวเสน กายํ ปริณาเมติ, จิตฺตวเสน กายํ อธิฏฺฐาติ, จิตฺตวเสน กายํ ปริณาเมตฺวา จิตฺตวเสน กายํ อธิฏฺฐหิตฺวา สุขสญฺญญฺจ ลหุสญฺญญฺจ \csromanfolio{388}โอกฺกมิตฺวา อทิสฺสมาเนน กาเยน พฺรหฺมโลกํ คจฺฉติ. โส ตสฺส พฺรหฺมุโน ปุรโต รูปิํ\csromansharedfootnote{837bef99d4a93bce}{รู (สฺยา, ก) ที 1. 73 ปิฏฺเฐ ปสฺสิตพฺพา.} อภินิมฺมินาติ มโนมยํ สพฺพงฺคปจฺจงฺคิํ\csromansharedfootnote{af4758f2baad75a0}{สพฺพงฺคปจฺจงฺคํ (สฺยา, ก)} อหีนินฺทฺริยํ, สเจ โส อิทฺธิมา จงฺกมติ, นิมฺมิโตปิ ตตฺถ จงฺกมติ. สเจ โส อิทฺธิมา ติฏฺฐติ, นิมฺมิโตปิ ตตฺถ ติฏฺฐติ. สเจ โส อิทฺธิมา นิสีทติ, นิมฺมิโตปิ ตตฺถ นิสีทติ. สเจ โส อิทฺธิมา เสยฺยํ กปฺเปติ, นิมฺมิโตปิ ตตฺถ เสยฺยํ กปฺเปติ. สเจ โส อิทฺธิมา ธูมายติ, นิมฺมิโตปิ ตตฺถ ธูมายติ. สเจ โส อิทฺธิมา ปชฺชลติ, นิมฺมิโตปิ ตตฺถ ปชฺชลติ. สเจ โส อิทฺธิมา ธมฺมํ ภาสติ, นิมฺมิโตปิ ตตฺถ ธมฺมํ ภาสติ. สเจ โส อิทฺธิมา ปญฺหํ ปุจฺฉติ, นิมฺมิโตปิ ตตฺถ ปญฺหํ ปุจฺฉติ. สเจ โส อิทฺธิมา ปญฺหํ ปุฏฺโฐ วิสชฺเชติ, นิมฺมิโตปิ ตตฺถ ปญฺหํ ปุฏฺโฐ วิสชฺเชติ. สเจ โส อิทฺธิมา เตน พฺรหฺมุนา สทฺธิํ สนฺติฏฺฐติ สลฺลปติ สากจฺฉํ สมาปชฺชติ, นิมฺมิโตปิ ตตฺถ เตน พฺรหฺมุนา สทฺธิํ สนฺติฏฺฐติ สลฺลปติ สากจฺฉํ สมาปชฺชติ. ยญฺญเทว โส อิทฺธิมา กโรติ, ตํตเทว หิ โส นิมฺมิโต กโรตีติ อยํ อธิฏฺฐานา อิทฺธิ. (1)}

\proseitem{13}{กตมา วิกุพฺพนา อิทฺธิ? สิขิสฺส ภควโต อรหโต สมฺมา\-สมฺพุทฺธสฺส อภิภู นาม สาวโก พฺรหฺมโลเก ฐิโต สหสฺสิ\-โลกธาตุํ\csromansharedfootnote{1d6879ff951d2cf1}{สหสฺสีโลกธาตุํ (สฺยา) วตฺถุ สํ 1. 157 ปิฏฺเฐ อาคตํ.} สเรน วิญฺญาเปสิ. โส ทิสฺสมาเนนปิ กาเยน ธมฺมํ เทเสสิ, อทิสฺสมาเนนปิ กาเยน ธมฺมํ เทเสสิ, ทิสฺสมาเนนปิ เหฏฺฐิเมน อุปฑฺฒกาเยน อทิสฺสมาเนนปิ อุปริเมน อุปฑฺฒกาเยน ธมฺมํ เทเสสิ. ทิสฺสมาเนนปิ อุปริเมน อุปฑฺฒกาเยน, อทิสฺสมาเนนปิ เหฏฺฐิเมน อุปฑฺฒกาเยน ธมฺมํ เทเสสิ. โส ปกติวณฺณํ วิชหิตฺวา กุมารกวณฺณํ วา ทสฺเสติ, นาควณฺณํ วา ทสฺเสติ, สุปณฺณวณฺณํ วา ทสฺเสติ, ยกฺขวณฺณํ วา ทสฺเสติ, อินฺทวณฺณํ วา\csromansharedfootnote{cb1ed66500bd003a}{อสุรวณฺณํ วา ทสฺสติ, อินฺทวณฺณํ วา (สฺยา)} ทสฺเสติ, เทววณฺณํ วา ทสฺเสติ, พฺรหฺมวณฺณํ วา ทสฺเสติ, สมุทฺทวณฺณํ วา ทสฺเสติ, ปพฺพตวณฺณํ วา ทสฺเสติ, วนวณฺณํ วา ทสฺเสติ, สีหวณฺณํ วา ทสฺเสติ, พฺยคฺฆวณฺณํ วา ทสฺเสติ, ทีปิวณฺณํ วา ทสฺเสติ, หตฺถิมฺปิ ทสฺเสติ, อสฺสมฺปิ ทสฺเสติ, รถมฺปิ ทสฺเสติ, ปตฺติมฺปิ ทสฺเสติ, วิวิธมฺปิ เสนาพฺยูหํ ทสฺเสตีติ อยํ วิกุพฺพนา อิทฺธิ. (2)}

\proseitem{14}{\csromanfolio{389}กตมา \textbf{มโนมยา อิทฺธิ?} อิธ ภิกฺขุ อิมมฺหา กายา อญฺญํ กายํ อภินิมฺมินาติ รูปิํ มโนมยํ สพฺพงฺคปจฺจงฺคิํ อหีนินฺทฺริยํ, เสยฺยถาปิ ปุริโส มุญฺชมฺหา อีสิกํ ปวาเหยฺย. ตสฺส เอวมสฺส “อยํ มุญฺโช, อยํ อีสิกา, อญฺโญ มุญฺโช, อญฺญา อีสิกา. มุญฺชมฺหาเตฺวว อีสิกา ปวาฬฺหา”ติ. เสยฺยถา วา ปน ปุริโส อสิํ โกสิยา ปวาเหยฺย. ตสฺส เอวมสฺส “อยํ อสิ, อยํ โกสิ, อญฺโญ อสิ, อญฺญา โกสิ. โกสิยาเตฺวว อสิ ปวาโฬฺห”ติ. เสยฺยถา วา ปน ปุริโส อหิํ กรณฺฑา อุทฺธเรยฺย. ตสฺส เอวมสฺส “อยํ อหิ, อยํ กรณฺโฑ. อญฺโญ อหิ, อญฺโญ กรณฺโฑ. กรณฺฑาเตฺวว อหิ อุพฺภโต”ติ. เอวเมวํ ภิกฺขุ อิมมฺหา กายา อญฺญํ กายํ อภินิมฺมินาติ รูปิํ มโนมยํ สพฺพงฺคปจฺจงฺคิํ อหีนินฺทฺริยํ. อยํ มโนมยา อิทฺธิ. (3)}

\proseitem{15}{กตมา ญาณวิปฺผารา อิทฺธิ? อนิจฺจา\-นุปสฺสนาย นิจฺจสญฺญาย ปหานฏฺโฐ อิชฺฌตีติ ญาณวิปฺผารา อิทฺธิ. ทุกฺขา\-นุปสฺสนาย สุขสญฺญาย. อนตฺตา\-นุปสฺสนาย อตฺตสญฺญาย. นิพฺพิทา\-นุปสฺสนาย นนฺทิยา. วิราคา\-นุปสฺสนาย ราคสฺส. นิโรธา\-นุปสฺสนาย สมุทยสฺส. ปฏินิสฺสคฺคา\-นุปสฺสนาย อาทานสฺส ปหานฏฺโฐ อิชฺฌตีติ ญาณวิปฺผารา อิทฺธิ. อายสฺมโต พากุลสฺส ญาณวิปฺผารา อิทฺธิ. อายสฺมโต สํกิจฺจสฺส ญาณวิปฺผารา อิทฺธิ. อายสฺมโต ภูตปาลสฺส ญาณวิปฺผารา อิทฺธิ. อยํ ญาณวิปฺผารา อิทฺธิ. (4)}

\proseitem{16}{กตมา สมาธิ\-วิปฺผารา อิทฺธิ? ปฐเมน ฌาเนน นีวรณานํ ปหานฏฺโฐ อิชฺฌตีติ สมาธิ\-วิปฺผารา อิทฺธิ. ทุติเยน ฌาเนน วิตกฺก\-วิจารานํ ปหานฏฺโฐ อิชฺฌตีติ สมาธิ\-วิปฺผารา อิทฺธิ. ตติเยน ฌาเนน ปีติยา ปหานฏฺโฐ อิชฺฌตีติ ฯเปฯ. จตุตฺเถน ฌาเนน สุขทุกฺขานํ ปหานฏฺโฐ อิชฺฌตีติ ฯเปฯ. อากาสา\-นญฺจา\-ยตนสมาปตฺติยา รูปสญฺญาย ปฏิฆสญฺญาย นานตฺตสญฺญาย ปหานฏฺโฐ อิชฺฌตีติ ฯเปฯ. วิญฺญาณญฺจายตน\-สมาปตฺติยา อากาสา\-นญฺจา\-ยตนสญฺญาย ปหานฏฺโฐ อิชฺฌตีติ ฯเปฯ. อากิญฺจญฺญายตน\-สมาปตฺติยา วิญฺญาณญฺจา\-ยตนสญฺญาย ปหานฏฺโฐ อิชฺฌตีติ ฯเปฯ. เนว\-สญฺญา\-นา\-สญฺญา\-ยตนสมาปตฺติยา อากิญฺจญฺญายตน\-สญฺญาย ปหานฏฺโฐ อิชฺฌตีติ สมาธิ\-วิปฺผารา อิทฺธิ, อายสฺมโต สาริปุตฺตสฺส สมาธิ\-วิปฺผารา อิทฺธิ, อายสฺมโต สญฺชีวสฺส สมาธิ\-วิปฺผารา \csromanfolio{390}อิทฺธิ, อายสฺมโต ขาณุ\-โกณฺฑญฺญสฺส สมาธิ\-วิปฺผารา อิทฺธิ, อุตฺตราย อุปาสิกาย สมาธิ\-วิปฺผารา อิทฺธิ, สามาวติยา อุปาสิกาย สมาธิ\-วิปฺผารา อิทฺธิ, อยํ สมาธิ\-วิปฺผารา อิทฺธิ. (5)}

\proseitem{117}{กตมา \textbf{อริยา อิทฺธิ?} อิธ ภิกฺขุ\csromansharedfootnote{936760ea94a16d61}{ม 3. 351 ปิฏฺเฐ ปสฺสิตพฺพา.} สเจ อากงฺขติ “ปฏิกูเล อปฺปฏิกูลสญฺญี วิหเรยฺยนฺ”ติ, อปฺปฏิกูลสญฺญี ตตฺถ วิหรติ. สเจ อากงฺขติ “อปฺปฏิกูเล ปฏิกูลสญฺญี วิหเรยฺยนฺ”ติ, ปฏิกูลสญฺญี ตตฺถ วิหรติ. สเจ อากงฺขติ “ปฏิกูเล จ อปฺปฏิกูเล จ อปฺปฏิกูลสญฺญี วิหเรยฺยนฺ”ติ, อปฺปฏิกูลสญฺญี ตตฺถ วิหรติ. สเจ อากงฺขติ “อปฺปฏิกูเล จ ปฏิกูเล จ ปฏิกูลสญฺญี วิหเรยฺยนฺ”ติ, ปฏิกูลสญฺญี ตตฺถ วิหรติ. สเจ อากงฺขติ “ปฏิกูเล จ อปฺปฏิกูเล จ ตทุภยํ อภินิวชฺเชตฺวา อุเปกฺขโก วิหเรยฺยํ สโต สมฺปชาโน”ติ อุเปกฺขโก ตตฺถ วิหรติ สโต สมฺปชาโน.}

\prose{กถํ \textbf{ปฏิกูเล อปฺปฏิกูลสญฺญี วิหรติ?} อนิฏฺฐสฺมิํ วตฺถุสฺมิํ เมตฺตาย วา ผรติ, ธาตุโต วา อุปสํหรติ, เอวํ ปฏิกูเล อปฺปฏิกูลสญฺญี วิหรติ.}

\prose{กถํ \textbf{อปฺปฏิกูเล ปฏิกูลสญฺญี วิหรติ?} อิฏฺฐสฺมิํ วตฺถุสฺมิํ อสุภาย วา ผรติ, อนิจฺจโต วา อุปสํหรติ, เอวํ อปฺปฏิกูเล ปฏิกูลสญฺญี วิหรติ.}

\prose{กถํ \textbf{ปฏิกูเล จ อปฺปฏิกูเล จ อปฺปฏิกูลสญฺญี วิหรติ?} อนิฏฺฐสฺมิํ จ อิฏฺฐสฺมิํ จ วตฺถุสฺมิํ เมตฺตาย วา ผรติ, ธาตุโต วา อุปสํหรติ. เอวํ ปฏิกูเล จ อปฺปฏิกูเล จ อปฺปฏิกูลญฺญี วิหรติ.}

\prose{กถํ \textbf{อปฺปฏิกูเล จ ปฏิกูเล จ ปฏิกูลสญฺญี วิหรติ?} อิฏฺฐสฺมิํ จ อนิฏฺฐสฺมิํ จ วตฺถุสฺมิํ อสุภาย วา ผรติ, อนิจฺจโต วา อุปสํหรติ, เอวํ อปฺปฏิกูเล จ ปฏิกูเล จ ปฏิกูลสญฺญี วิหรติ.}

\prose{\textbf{กถํ ปฏิกูเล จ อปฺป}ฏิกูเล จ ตทุภยํ อภินิวชฺเชตฺวา อุเปกฺขโก วิหรติ สโต สมฺปชาโน? อิธ ภิกฺขุ จกฺขุนา รูปํ ทิสฺวา เนว สุมโน โหติ น ทุมฺมโน, อุเปกฺขโก วิหรติ สโต สมฺปชาโน. โสเตน สทฺทํ สุตฺวา ฯเปฯ. ฆาเนน คนฺธํ ฆายิตฺวา. ชิวฺหาย รสํ สายิตฺวา.}

\vspace{\tipitakaparskip}
\csromancenter{อภิสมย\-กถา กาเยน โผฏฺฐพฺพํ ผุสิตฺวา. มนสา ธมฺมํ วิญฺญาย เนว สุมโน โหติ น ทุมฺมโน, อุเปกฺขโก วิหรติ สโต สมฺปชาโน. เอวํ ปฏิกูเล จ อปฺปฏิกูเล จ ตทุภยํ อภินิวชฺเชตฺวา อุเปกฺขโก วิหรติ สโต สมฺปชาโน, อยํ อริยา อิทฺธิ. (6)}
\proseitem{18}{\csromanfolio{391}กตมา \textbf{กมฺมวิปากชา อิทฺธิ?} สพฺเพสํ ปกฺขีนํ สพฺเพสํ เทวานํ เอกจฺจานํ มนุสฺสานํ เอกจฺจานํ วินิปาติกานํ, อยํ กมฺมวิปากชา อิทฺธิ. (7)}

\prose{กตมา ปุญฺญวโต อิทฺธิ? ราชา จกฺกวตฺตี\csromansharedfootnote{48449606da1ffa34}{จกฺกวตฺติ (สฺยา)} เวหาสํ คจฺฉติ สทฺธิํ จตุรงฺคินิยา เสนาย อนฺตมโส อสฺสพนฺธ\-โคปุริเส\csromansharedfootnote{d4e0b014e7f65b24}{อสฺสพนฺธโคปเก ปุริเส (สฺยา)} อุปาทาย. โชติกสฺส คหปติสฺส ปุญฺญวโต อิทฺธิ. ชฏิลสฺส คหปติสฺส ปุญฺญวโต อิทฺธิ. เมณฺฑกสฺส คหปติสฺส ปุญฺญวโต อิทฺธิ. โฆสิตสฺส คหปติสฺส ปุญฺญวโต อิทฺธิ. ปญฺจนฺนํ มหาปุญฺญานํ ปุญฺญวโต อิทฺธิ, อยํ ปุญฺญวโต อิทฺธิ. (8)}

\prose{กตมา \textbf{วิชฺชามยา อิทฺธิ?} วิชฺชาธรา วิชฺชํ ปริชปฺเปตฺวา เวหาสํ คจฺฉนฺติ, อากาเส อนฺตลิกฺเข หตฺถิมฺปิ ทสฺเสนฺติ, อสฺสมฺปิ ทสฺเสนฺติ, รถมฺปิ ทสฺเสนฺติ, ปตฺติมฺปิ ทสฺเสนฺติ, วิวิธมฺปิ เสนาพฺยูหํ ทสฺเสนฺติ, อยํ วิชฺชามยา อิทฺธิ. (9)}

\prosewithcloserruled{กถํ ตตฺถ ตตฺถ สมฺมา\-ปโยค\-ปจฺจยา อิชฺฌนฏฺเฐน อิทฺธิ? เนกฺขมฺเมน กามจฺฉนฺทสฺส ปหานฏฺโฐ อิชฺฌตีติ ตตฺถ ตตฺถ สมฺมา\-ปโยค\-ปจฺจยา อิชฺฌนฏฺเฐน อิทฺธิ, อพฺยาปาเทน พฺยาปาทสฺส ปหานฏฺโฐ อิชฺฌตีติ ตตฺถ ตตฺถ สมฺมา\-ปโยค\-ปจฺจยา อิชฺฌนฏฺเฐน อิทฺธิ ฯเปฯ อรหตฺต\-มคฺเคน สพฺพกิเลสานํ ปหานฏฺโฐ อิชฺฌตีติ ตตฺถ ตตฺถ สมฺมา\-ปโยค\-ปจฺจยา อิชฺฌนฏฺเฐน อิทฺธิ. เอวํ ตตฺถ ตตฺถ สมฺมา\-ปโยค\-ปจฺจยา อิชฺฌนฏฺเฐน อิทฺธิ. (10) อิมา ทส อิทฺธิโย.}{อิทฺธิกถา นิฏฺฐิตา.}
\csromanmatikaanchor{mtk.124}
\csromantocmarkref{section}{3. อภิสมยกถา}{mtk.124}
\bookmark[dest={mtk.124},level=1]{3. อภิสมยกถา}
\csromanheader{1}{3. อภิสมย\-กถา}
\proseitem{19}{อ\textbf{ภิสมโย}ติ เกน อภิสเมติ? จิตฺเตน อภิสเมติ.}

\prose{หญฺจิ จิตฺเตน อภิสเมติ, เตน หิ อญฺญาณี อภิสเมติ. น อญฺญาณี อภิสเมติ, ญาเณน อภิสเมติ.}

\prose{\csromanfolio{392}หญฺจิ ญาเณน อภิสเมติ, เตน หิ อจิตฺเตน จ ญาเณน จ อจิตฺตโก\csromansharedfootnote{dfb01bca7e7111ca}{เตน หิ อจิตฺตโก (สฺยา)} อภิสเมติ. น อจิตฺตโก อภิสเมติ, จิตฺเตน จ ญาเณน จ อภิสเมติ.}

\prose{หญฺจิ จิตฺเตน จ ญาเณน จ อภิสเมติ, เตน หิ กามาวจร\-จิตฺเตน จ ญาเณน จ อภิสเมติ, น กามาวจร\-จิตฺเตน จ ญาเณน จ อภิสเมติ.}

\prose{เตน หิ รูปาวจร\-จิตฺเตน จ ญาเณน จ อภิสเมติ, น รูปาวจร\-จิตฺเตน จ ญาเณน จ อภิสเมติ.}

\prose{เตน หิ อรูปาวจร\-จิตฺเตน จ ญาเณน จ อภิสเมติ. น อรูปาวจร\-จิตฺเตน จ ญาเณน จ อภิสเมติ.}

\prose{เตน หิ กมฺมสฺสกต\-จิตฺเตน จ ญาเณน จ อภิสเมติ, น กมฺมสฺสกต\-จิตฺเตน จ ญาเณน จ อภิสเมติ.}

\prose{เตน หิ สจฺจา\-นุโลมิก\-จิตฺเตน จ ญาเณน จ อภิสเมติ, น สจฺจา\-นุโลมิก\-จิตฺเตน จ ญาเณน จ อภิสเมติ.}

\prose{เตน หิ อตีตจิตฺเตน จ ญาเณน จ อภิสเมติ, น อตีตจิตฺเตน จ ญาเณน จ อภิสเมติ.}

\prose{เตน หิ อนาคตจิตฺเตน จ ญาเณน จ อภิสเมติ, น อนาคตจิตฺเตน จ ญาเณน จ อภิสเมติ.}

\prose{เตน หิ ปจฺจุปฺปนฺน\-โลกิยจิตฺเตน จ ญาเณน จ อภิสเมติ, น ปจฺจุปฺปนฺน\-โลกิยจิตฺเตน จ ญาเณน จ อภิสเมติ. โลกุตฺตรมคฺค\-กฺขเณ ปจฺจุปฺปนฺน\-จิตฺเตน จ ญาเณน จ อภิสเมติ.}

\prose{กถํ โลกุตฺตรมคฺค\-กฺขเณ ปจฺจุปฺปนฺน\-จิตฺเตน จ ญาเณน จ อภิสเมติ? โลกุตฺตรมคฺค\-กฺขเณ อุปฺปาทา\-ธิปเตยฺยํ จิตฺตํ ญาณสฺส เหตุ ปจฺจโย จ ตํสมฺปยุตฺตํ นิโรธโคจรํ ทสฺสนา\-ธิปเตยฺยํ ญาณํ จิตฺตสฺส เหตุ ปจฺจโย จ ตํสมฺปยุตฺตํ ญาณํ นิโรธโคจรํ, เอวํ โลกุตฺตรมคฺค\-กฺขเณ ปจฺจุปฺปนฺน\-จิตฺเตน จ ญาเณน จ อภิสเมติ.}

\csromanheader{2}{3. อภิสมย\-กถา}
\proseitem{20}{\csromanfolio{393}กิํ นุ เอตฺตโกเยว อภิสมโยติ? น หิ, โลกุตฺตรมคฺค\-กฺขเณ ทสฺสนาภิ\-สมโย สมฺมาทิฏฺฐิ, อภินิโรปนาภิ\-สมโย สมฺมาสงฺกปฺโป, ปริคฺคหาภิ\-สมโย สมฺมาวาจา, สมุฏฺฐานาภิ\-สมโย สมฺมากมฺมนฺโต, โวทานา\-ภิสมโย สมฺมาอาชีโว, ปคฺคหา\-ภิสมโย สมฺมาวายาโม, อุปฏฺฐานา\-ภิสมโย สมฺมาสติ, อวิกฺเขปา\-ภิสมโย สมฺมาสมาธิ.  อุปฏฺฐานา\-ภิสมโย สติสมฺโพชฺฌงฺโค, ปวิจยา\-ภิสมโย ธมฺมวิจย\-สมฺโพชฺฌงฺโค, ปคฺคหา\-ภิสมโย วีริย\-สมฺโพชฺฌงฺโค, ผรณา\-ภิสมโย ปีติสมฺโพชฺฌงฺโค, อุปสมาภิ\-สมโย ปสฺสทฺธิ\-สมฺโพชฺฌงฺโค, อวิกฺเขปา\-ภิสมโย สมาธิ\-สมฺโพชฺฌงฺโค, ปฏิสงฺขานา\-ภิสมโย อุเปกฺขา\-สมฺโพชฺฌงฺโค.  อสฺสทฺธิเย อกมฺปิยา\-ภิสมโย สทฺธาพลํ, โกสชฺเช อกมฺปิยา\-ภิสมโย วีริยพลํ, ปมาเท อกมฺปิยา\-ภิสมโย สติพลํ, อุทฺธจฺเจ อกมฺปิยา\-ภิสมโย สมาธิพลํ, อวิชฺชาย อกมฺปิยา\-ภิสมโย ปญฺญาพลํ.  อธิโมกฺขา\-ภิสมโย สทฺธินฺทฺริยํ, ปคฺคหา\-ภิสมโย วีริยินฺทฺริยํ, อุปฏฺฐานา\-ภิสมโย สตินฺทฺริยํ, อวิกฺเขปา\-ภิสมโย สมาธินฺทฺริยํ, ทสฺสนาภิ\-สมโย ปญฺญินฺทฺริยํ.  อาธิปเตยฺย\-ฏฺเฐน อินฺทฺริยาภิ\-สมโย, อกมฺปิยฏฺเฐน พลาภิสมโย, นิยฺยานฏฺเฐน โพชฺฌงฺคา\-ภิสมโย, เหตุฏฺเฐน มคฺคาภิสมโย, อุปฏฺฐาน\-ฏฺเฐน สติปฏฺฐานา\-ภิสมโย, ปทหฏฺเฐน สมฺมปฺปธานาภิ\-สมโย, อิชฺฌนฏฺเฐน อิทฺธิปาทา\-ภิสมโย, ตถฏฺเฐน สจฺจาภิสมโย, อวิกฺเขป\-ฏฺเฐน สมถา\-ภิสมโย, อนุปสฺสนฏฺเฐน วิปสฺสนาภิ\-สมโย, เอกรสฏฺเฐน สมถวิปสฺสนาภิ\-สมโย, อนติวตฺตน\-ฏฺเฐน ยุคนทฺธาภิ\-สมโย, สํวรฏฺเฐน สีลวิสุทฺธิ\-อภิสมโย, อวิกฺเขป\-ฏฺเฐน วิสุทฺธิ\-อภิสมโย, ทสฺสนฏฺเฐน ทิฏฺฐิวิสุทฺธิ\-อภิสมโย, มุตฺตฏฺเฐน วิโมกฺขา\-ภิสมโย, ปฏิเวธ\-ฏฺเฐน วิชฺชาภิสมโย, ปริจฺจาค\-ฏฺเฐน วิมุตฺติ\-อภิสมโย, สมุจฺเฉท\-ฏฺเฐน ขเย ญาณํ อภิสมโย.  ฉนฺโท มูลฏฺเฐน อภิสมโย, มนสิกาโร สมุฏฺฐาน\-ฏฺเฐน อภิสมโย, ผสฺโส สโมธาน\-ฏฺเฐน อภิสมโย, เวทนา สโมสรณ\-ฏฺเฐน อภิสมโย, สมาธิ ปมุขฏฺเฐน อภิสมโย, สติ อาธิปเตยฺย\-ฏฺเฐน อภิสมโย, ปญฺญา ตทุตฺต\-รฏฺเฐน อภิสมโย, วิมุตฺติ สารฏฺเฐน อภิสมโย, อมโตคธํ นิพฺพานํ ปริโยสาน\-ฏฺเฐน อภิสมโย.}

\proseitem{21}{\csromanfolio{394}กิํนุ เอตฺตโกเยว อภิสมโยติ? น หิ, โสตาปตฺติมคฺค\-กฺขเณ ทสฺสนาภิ\-สมโย สมฺมาทิฏฺฐิ ฯเปฯ อมโตคธํ นิพฺพานํ ปริโยสาน\-ฏฺเฐน อภิสมโย.}

\prose{กิํนุ เอตฺตโกเยว อภิสมโยติ? น หิ, โสตาปตฺติ\-ผลกฺขเณ ทสฺสนาภิ\-สมโย สมฺมาทิฏฺฐิ ฯเปฯ ปฏิปฺปสฺสทฺธ\-ฏฺเฐน อนุปฺปาเท ญาณํ อภิสมโย, ฉนฺโท มูลฏฺเฐน อภิสมโย ฯเปฯ อมโตคธํ นิพฺพานํ ปริโยสาน\-ฏฺเฐน อภิสมโย.}

\prose{กิํนุ เอตฺตโกเยว อภิสมโยติ? น หิ, สกทาคามิ\-มคฺคกฺขเณ ฯเปฯ สกทาคามิ\-ผลกฺขเณ. อนาคามิ\-มคฺคกฺขเณ. อนาคามิ\-ผลกฺขเณ. อรหตฺต\-มคฺคกฺขเณ ฯเปฯ. อรหตฺต\-ผลกฺขเณ ทสฺสนาภิ\-สมโย สมฺมาทิฏฺฐิ, อภินิโรปนาภิ\-สมโย สมฺมาสงฺกปฺโป ฯเปฯ ปฏิปฺปสฺสทฺธ\-ฏฺเฐน อนุปฺปาเท ญาณํ อภิสมโย. ฉนฺโท มูลฏฺเฐน อภิสมโย ฯเปฯ อมโตคธํ นิพฺพานํ ปริโยสาน\-ฏฺเฐน อภิสมโย.}

\prose{ยฺวายํ\csromansharedfootnote{97193d6a5b14314a}{สฺวายํ (สฺยา, อิ)} กิเลเส ปชหติ, อตีเต กิเลเส ปชหติ ฯเปฯ อนาคเต กิเลเส ปชหติ, ปจฺจุปฺปนฺเน กิเลเส ปชหติ? อตีเต กิเลเส ปชหตีติ. หญฺจิ อตีเต กิเลเส ปชหติ, เตน หิ ขีณํ เขเปติ, นิรุทฺธํ นิโรเธติ, วิคตํ วิคเมติ, อตฺถงฺคตํ อตฺถงฺคเมติ, อตีตํ ยํ น อตฺถิ, ตํ ปชหตีติ น อตีเต กิเลเส ปชหตีติ.  อนาคเต กิเลเส ปชหตีติ. หญฺจิ อนาคเต กิเลเส ปชหติ, เตน หิ อชาตํ ปชหติ, อนิพฺพตฺตํ ปชหติ, อนุปฺปนฺนํ ปชหติ, อปาตุภูตํ ปชหติ, อนาคตํ ยํ น อตฺถิ, ตํ ปชหตีติ น อนาคเต กิเลเส ปชหตีติ.  ปจฺจุปฺปนฺเน กิเลเส ปชหตีติ. หญฺจิ ปจฺจุปฺปนฺเน กิเลเส ปชหติ, เตน หิ รตฺโต ราคํ ปชหติ, ทุฏฺโฐ โทสํ ปชหติ, มูโฬฺห โมหํ ปชหติ, วินิพทฺโธ มานํ ปชหติ, ปรามฏฺโฐ ทิฏฺฐิํ ปชหติ, วิกฺเขปคโต อุทฺธจฺจํ ปชหติ, อนิฏฺฐงฺคโต วิจิกิจฺฉํ ปชหติ, ถามคโต อนุสยํ ปชหติ, กณฺห\-สุกฺกธมฺมา ยุคนทฺธา สมเมว วตฺตนฺติ. สํกิเลสิกา\csromansharedfootnote{612e14ef4c713d2e}{ตํสํกิเลสิกา (สฺยา, ก)} มคฺคภาวนา โหติ.}

\prose{น หิ อตีเต กิเลเส ปชหติ, น อนาคเต กิเลเส ปชหติ, น ปจฺจุปฺปนฺเน กิเลเส ปชหตีติ. หญฺจิ น อตีเต กิเลเส ปชหติ,}

\prosewithcloserruled{\csromanfolio{395}น อนาคเต ฯเปฯ. น ปจฺจุปฺปนฺเน กิเลเส ปชหติ, เตน หิ นตฺถิ มคฺคภาวนา, นตฺถิ ผลสจฺฉิกิริยา, นตฺถิ กิเลสปฺปหานํ, นตฺถิ ธมฺมา\-ภิสมโยติ. อตฺถิ มคฺคภาวนา, อตฺถิ ผลสจฺฉิกิริยา, อตฺถิ กิเลสปฺปหานํ, อตฺถิ ธมฺมา\-ภิสมโย. ยถา กถํ วิย, เสยฺยถาปิ ตรุโณ รุกฺโข อชาตผโล, ตเมนํ ปุริโส มูลํ ฉินฺเทยฺย, เย ตสฺส รุกฺขสฺส อชาตผลา, เต อชาตาเยว น ชายนฺติ. อนิพฺพตฺตาเยว น นิพฺพตฺตนฺติ, อนุปฺปนฺนาเยว น อุปฺปชฺชนฺติ, อปาตุภูตาเยว น ปาตุภวนฺติ. เอวเมวํ “อุปฺปาโท เหตุ อุปฺปาโท ปจฺจโย กิเลสานํ นิพฺพตฺติยา”ติ อุปฺปาเท อาทีนวํ ทิสฺวา อนุปฺปาเท จิตฺตํ ปกฺขนฺทติ, อนุปฺปาเท จิตฺตสฺส ปกฺขนฺทตฺตา เย อุปฺปาทปจฺจยา กิเลสา นิพฺพตฺเตยฺยุํ, เต อชาตาเยว น ชายนฺติ, อนิพฺพตฺตาเยว น นิพฺพตฺตนฺติ, อนุปฺปนฺนาเยว น อุปฺปชฺชนฺติ, อปาตุภูตาเยว น ปาตุภวนฺติ. เอวํ “เหตุนิโรธา ทุกฺขนิโรโธ, ปวตฺตํ เหตุ, นิมิตฺตํ เหตุ, อายูหนา เหตุ, อายูหนา ปจฺจโย กิเลสานํ นิพฺพตฺติยา”ติ อายูหเน อาทีนวํ ทิสฺวา อนายูหเน จิตฺตํ ปกฺขนฺทติ, อนายูหเน จิตฺตสฺส ปกฺขนฺทตฺตา เย อายูหนปจฺจยา กิเลสา นิพฺพตฺเตยฺยุํ, เต อชาตาเยว น ชายนฺติ, อนิพฺพตฺตาเยว น นิพฺพตฺตนฺติ, อนุปฺปนฺนาเยว น อุปฺปชฺชนฺติ, อปาตุภูตาเยว น ปาตุภวนฺติ. เอวํ เหตุนิโรธา ทุกฺขนิโรโธ. เอวํ อตฺถิ มคฺคภาวนา, อตฺถิ ผลสจฺฉิกิริยา, อตฺถิ กิเลสปฺปหานํ, อตฺถิ ธมฺมา\-ภิสมโยติ.}{อภิสมย\-กถา นิฏฺฐิตา.}
\proseitem{22}{สาวตฺถิ\-นิทานํ\csromansharedfootnote{a090ef5ead7c69e3}{สํ 3. 42 ปิฏฺเฐ ปสฺสิตพฺพา.}.  เสยฺยถาปิ ภิกฺขเว เย เกจิ พลกรณียา กมฺมนฺตา กรียนฺติ, สพฺเพ เต ปถวิํ นิสฺสาย ปถวิยํ ปติฏฺฐาย, เอวเมเต พลกรณียา กมฺมนฺตา กรียนฺติ. เอวเมวํ ภิกฺขเว ภิกฺขุ สีลํ นิสฺสาย สีเล ปติฏฺฐาย อริยํ อฏฺฐงฺคิกํ มคฺคํ ภาเวติ, อริยํ อฏฺฐงฺคิกํ มคฺคํ พหุลีกโรติ.}

\prose{กถญฺจ ภิกฺขเว ภิกฺขุ สีลํ นิสฺสาย สีเล ปติฏฺฐาย อริยํ อฏฺฐงฺคิกํ มคฺคํ ภาเวติ, อริยํ อฏฺฐงฺคิกํ มคฺคํ พหุลีกโรติ? อิธ ภิกฺขเว \csromanfolio{396}ภิกฺขุ สมฺมาทิฏฺฐิํ ภาเวติ วิเวกนิสฺสิตํ วิราคนิสฺสิตํ นิโรธ\-นิสฺสิตํ โวสคฺค\-ปริณามิํ. สมฺมาสงฺกปฺปํ ฯเปฯ. สมฺมาวาจํ. สมฺมากมฺมนฺตํ. สมฺมาอาชีวํ. สมฺมาวายามํ. สมฺมาสติํ. สมฺมาสมาธิํ ภาเวติ วิเวกนิสฺสิตํ วิราคนิสฺสิตํ นิโรธ\-นิสฺสิตํ โวสคฺค\-ปริณามิํ. เอวํ โข ภิกฺขเว ภิกฺขุ สีลํ นิสฺสาย สีเล ปติฏฺฐาย อริยํ อฏฺฐงฺคิกํ มคฺคํ ภาเวติ, อริยํ อฏฺฐงฺคิกํ มคฺคํ พหุลีกโรติ.}

\proseitem{23}{เสยฺยถาปิ ภิกฺขเว เย เกจิ\csromansharedfootnote{394edb582b577af2}{เยปิเม (สฺยา) สํ 3. 43 ปิฏฺเฐ ปสฺสิตพฺพา.} พีชคาม\-ภูตคามา วุทฺธิํ วิรูฬฺหิํ เวปุลฺลํ อาปชฺชนฺติ, สพฺเพ เต ปถวิํ นิสฺสาย ปถวิยํ ปติฏฺฐาย, เอวเมเต พีชคาม\-ภูตคามา วุทฺธิํ วิรุฬฺหิํ เวปุลฺลํ อาปชฺชนฺติ. เอวเมวํ โข ภิกฺขเว ภิกฺขุ สีลํ นิสฺสาย สีเล ปติฏฺฐาย อริยํ อฏฺฐงฺคิกํ มคฺคํ ภาเวนฺโต อริยํ อฏฺฐงฺคิกํ มคฺคํ พหุลีกโรนฺโต วุทฺธิํ วิรุฬฺหิํ เวปุลฺลํ ปาปุณาติ ธมฺเมสุ.}

\prose{กถญฺจ ภิกฺขเว ภิกฺขุ สีลํ นิสฺสาย สีเล ปติฏฺฐาย อริยํ อฏฺฐงฺคิกํ มคฺคํ ภาเวนฺโต อริยํ อฏฺฐงฺคิกํ มคฺคํ พหุลีกโรนฺโต วุทฺธิํ วิรุฬฺหิํ เวปุลฺลํ ปาปุณาติ ธมฺเมสุ? อิธ ภิกฺขเว ภิกฺขุ สมฺมาทิฏฺฐิํ ภาเวติ วิเวกนิสฺสิตํ วิราคนิสฺสิตํ นิโรธ\-นิสฺสิตํ โวสคฺค\-ปริณามิํ, สมฺมาสงฺกปฺปํ ภาเวติ ฯเปฯ สมฺมาวาจํ ภาเวติ. สมฺมากมฺมนฺตํ ภาเวติ. สมฺมาอาชีวํ ภาเวติ. สมฺมาวายามํ ภาเวติ. สมฺมาสติํ ภาเวติ. สมฺมาสมาธิํ ภาเวติ วิเวกนิสฺสิตํ วิราคนิสฺสิตํ นิโรธ\-นิสฺสิตํ โวสคฺค\-ปริณามิํ. เอวํ โข ภิกฺขเว ภิกฺขุ สีลํ นิสฺสาย สีเล ปติฏฺฐาย อริยํ อฏฺฐงฺคิกํ มคฺคํ ภาเวนฺโต อริยํ อฏฺฐงฺคิกํ มคฺคํ พหุลีกโรนฺโต วุฑฺฒิํ วิรุฬฺหิํ เวปุลฺลํ ปาปุณาติ ธมฺเมสูติ.}

\csromanmatikaanchor{mtk.125}
\csromanmatikaanchor{mtk.126}
\csromantocmarknopage{section}{4. วิเวกกถา}{mtk.125}
\bookmark[dest={mtk.125},level=1]{4. วิเวกกถา}
\csromantocmarkref{subsection}{1. มคฺคงฺคนิทฺเทส}{mtk.126}
\bookmark[dest={mtk.126},level=2]{1. มคฺคงฺคนิทฺเทส}
\csromanheader{2}{1. มคฺคงฺค\-นิทฺเทส}
\proseitem{24}{สมฺมาทิฏฺฐิยา ปญฺจ วิเวกา, ปญฺจ วิราคา, ปญฺจ นิโรธา, ปญฺจ โวสคฺคา, ทฺวาทส นิสฺสยา. สมฺมา\-สงฺกปฺปสฺส ฯเปฯ. สมฺมาวาจาย. สมฺมา\-กมฺมนฺตสฺส. สมฺมาอาชีวสฺส. สมฺมาวายามสฺส. สมฺมาสติยา. สมฺมา\-สมาธิสฺส ปญฺจ วิเวกา, ปญฺจ วิราคา, ปญฺจ นิโรธา, ปญฺจ โวสคฺคา, ทฺวาทส นิสฺสยา.}

\prose{\csromanfolio{397}สมฺมาทิฏฺฐิยา กตเม ปญฺจ วิเวกา? วิกฺขมฺภน\-วิเวโก ตทงฺควิเวโก สมุจฺเฉท\-วิเวโก ปฏิปฺปสฺสทฺธิ\-วิเวโก นิสฺสรณ\-วิเวโก. วิกฺขมฺภน\-วิเวโก จ นีวรณานํ ปฐมชฺฌานํ ภาวยโต, ตทงฺควิเวโก จ ทิฏฺฐิคตานํ นิพฺเพธ\-ภาคิยํ สมาธิํ ภาวยโต, สมุจฺเฉท\-วิเวโก จ โลกุตฺตรํ ขยคามิ\-มคฺคํ ภาวยโต, ปฏิปฺปสฺสทฺธิ\-วิเวโก จ ผลกฺขเณ, นิสฺสรณ\-วิเวโก จ นิโรโธ นิพฺพานํ. สมฺมาทิฏฺฐิยา อิเม ปญฺจ วิเวกา. อิเมสุ ปญฺจสุ วิเวเกสุ ฉนฺทชาโต โหติ สทฺธาธิมุตฺโต, จิตฺตํ จสฺส สฺวาธิฏฺฐิตํ.}

\prose{สมฺมาทิฏฺฐิยา กตเม ปญฺจ วิราคา? วิกฺขมฺภน\-วิราโค ตทงฺควิราโค สมุจฺเฉท\-วิราโค ปฏิปฺปสฺสทฺธิ\-วิราโค นิสฺสรณ\-วิราโค. วิกฺขมฺภน\-วิราโค จ นีวรณานํ ปฐมชฺฌานํ ภาวยโต, ตทงฺควิราโค จ ทิฏฺฐิคตานํ นิพฺเพธ\-ภาคิยํ สมาธิํ ภาวยโต, สมุจฺเฉท\-วิราโค จ โลกุตฺตรํ ขยคามิ\-มคฺคํ ภาวยโต, ปฏิปฺปสฺสทฺธิ\-วิราโค จ ผลกฺขเณ, นิสฺสรณ\-วิราโค จ นิโรโธ นิพฺพานํ. สมฺมาทิฏฺฐิยา อิเม ปญฺจ วิราคา. อิเมสุ ปญฺจสุ วิราเคสุ ฉนฺทชาโต โหติ สทฺธาธิมุตฺโต, จิตฺตํ จสฺส สฺวาธิฏฺฐิตํ.}

\prose{สมฺมาทิฏฺฐิยา กตเม ปญฺจ นิโรธา? วิกฺขมฺภน\-นิโรโธ ตทงฺคนิโรโธ สมุจฺเฉท\-นิโรโธ ปฏิปฺปสฺสทฺธิ\-นิโรโธ นิสฺสรณ\-นิโรโธ. วิกฺขมฺภน\-นิโรโธ จ นีวรณานํ ปฐมชฺฌานํ ภาวยโต, ตทงฺคนิโรโธ จ ทิฏฺฐิคตานํ นิพฺเพธ\-ภาคิยํ สมาธิํ ภาวยโต, สมุจฺเฉท\-นิโรโธ จ โลกุตฺตรํ ขยคามิ\-มคฺคํ ภาวยโต, ปฏิปฺปสฺสทฺธิ\-นิโรโธ จ ผลกฺขเณ, นิสฺสรณ\-นิโรโธ จ อมตา ธาตุ. สมฺมาทิฏฺฐิยา อิเม ปญฺจ นิโรธา. อิเมสุ ปญฺจสุ นิโรเธสุ ฉนฺทชาโต โหติ สทฺธาธิมุตฺโต, จิตฺตํ จสฺส สฺวาธิฏฺฐิตํ.}

\prose{สมฺมาทิฏฺฐิยา กตเม ปญฺจ โวสคฺคา? วิกฺขมฺภน\-โวสคฺโค ตทงฺคโวสคฺโค สมุจฺเฉท\-โวสคฺโค ปฏิปฺปสฺสทฺธิ\-โวสคฺโค นิสฺสรณ\-โวสคฺโค. วิกฺขมฺภน\-โวสคฺโค จ นีวรณานํ ปฐมชฺฌานํ ภาวยโต, ตทงฺคโวสคฺโค จ ทิฏฺฐิคตานํ นิพฺเพธ\-ภาคิยํ สมาธิํ ภาวยโต, สมุจฺเฉท\-โวสคฺโค จ โลกุตฺตรํ ขยคามิ\-มคฺคํ ภาวยโต, ปฏิปฺปสฺสทฺธิ\-โวสคฺโค จ ผลกฺขเณ, นิสฺสรณ\-โวสคฺโค จ นิโรโธ นิพฺพานํ. สมฺมาทิฏฺฐิยา อิเม ปญฺจ โวสคฺคา. อิเมสุ ปญฺจสุ โวสคฺเคสุ ฉนฺทชาโต}

\prose{\csromanfolio{398}โหติ สทฺธาธิมุตฺโต, จิตฺตํ จสฺส สฺวาธิฏฺฐิตํ. สมฺมาทิฏฺฐิยา อิเม ปญฺจ วิเวกา, ปญฺจ วิราคา, ปญฺจ นิโรธา, ปญฺจ โวสคฺคา, ทฺวาทส นิสฺสยา.}

\proseitem{25}{สมฺมา\-สงฺกปฺปสฺส ฯเปฯ. สมฺมาวาจาย. สมฺมา\-กมฺมนฺตสฺส. สมฺมาอาชีวสฺส. สมฺมาวายามสฺส. สมฺมาสติยา. สมฺมา\-สมาธิสฺส กตเม ปญฺจ วิเวกา? วิกฺขมฺภน\-วิเวโก ตทงฺควิเวโก สมุจฺเฉท\-วิเวโก ปฏิปฺปสฺสทฺธิ\-วิเวโก นิสฺสรณ\-วิเวโก. วิกฺขมฺภน\-วิเวโก จ นีวรณานํ ปฐมชฺฌานํ ภาวยโต, ตทงฺควิเวโก จ ทิฏฺฐิคตานํ นิพฺเพธ\-ภาคิยํ สมาธิํ ภาวยโต, สมุจฺเฉท\-วิเวโก จ โลกุตฺตรํ ขยคามิ\-มคฺคํ ภาวยโต, ปฏิปฺปสฺสทฺธิ\-วิเวโก จ ผลกฺขเณ, นิสฺสรณ\-วิเวโก จ นิโรโธ นิพฺพานํ. สมฺมา\-สมาธิสฺส อิเม ปญฺจ วิเวกา. อิเมสุ ปญฺจสุ วิเวเกสุ ฉนฺทชาโต โหติ สทฺธาธิมุตฺโต, จิตฺตํ จสฺส สฺวาธิฏฺฐิตํ.}

\prose{สมฺมา\-สมาธิสฺส กตเม ปญฺจ วิราคา? วิกฺขมฺภน\-วิราโค ตทงฺควิราโค สมุจฺเฉท\-วิราโค ปฏิปฺปสฺสทฺธิ\-วิราโค นิสฺสรณ\-วิราโค. วิกฺขมฺภน\-วิราโค จ นีวรณานํ ปฐมชฺฌานํ ภาวยโต, ตทงฺควิราโค จ ทิฏฺฐิคตานํ นิพฺเพธ\-ภาคิยํ สมาธิํ ภาวยโต, สมุจฺเฉท\-วิราโค จ โลกุตฺตรํ ขยคามิ\-มคฺคํ ภาวยโต, ปฏิปฺปสฺสทฺธิ\-วิราโค จ ผลกฺขเณ, นิสฺสรณ\-วิราโค จ นิโรโธ นิพฺพานํ. สมฺมา\-สมาธิสฺส อิเม ปญฺจ วิราคา. อิเมสุ ปญฺจสุ วิราเคสุ ฉนฺทชาโต โหติ สทฺธาธิมุตฺโต, จิตฺตํ จสฺส สฺวาธิฏฺฐิตํ.}

\prose{สมฺมา\-สมาธิสฺส กตเม ปญฺจ นิโรธา? วิกฺขมฺภน\-นิโรโธ ตทงฺคนิโรโธ สมุจฺเฉท\-นิโรโธ ปฏิปฺปสฺสทฺธิ\-นิโรโธ นิสฺสรณ\-นิโรโธ. วิกฺขมฺภน\-นิโรโธ จ นีวรณานํ ปฐมชฺฌานํ ภาวยโต, ตทงฺคนิโรโธ จ ทิฏฺฐิคตานํ นิพฺเพธ\-ภาคิยํ สมาธิํ ภาวยโต, สมุจฺเฉท\-นิโรโธ จ โลกุตฺตรํ ขยคามิ\-มคฺคํ ภาวยโต, ปฏิปฺปสฺสทฺธิ\-นิโรโธ จ ผลกฺขเณ, นิสฺสรณ\-นิโรโธ จ อมตา ธาตุ. สมฺมา\-สมาธิสฺส อิเม ปญฺจ นิโรธา. อิเมสุ ปญฺจสุ นิโรเธสุ ฉนฺทชาโต โหติ สทฺธาธิมุตฺโต, จิตฺตํ จสฺส สฺวาธิฏฺฐิตํ.}

\prose{สมฺมา\-สมาธิสฺส กตเม ปญฺจ โวสคฺคา? วิกฺขมฺภน\-โวสคฺโค ตทงฺคโวสคฺโค สมุจฺเฉท\-โวสคฺโค ปฏิปฺปสฺสทฺธิ\-โวสคฺโค นิสฺสรณ\-โวสคฺโค. วิกฺขมฺภน\-โวสคฺโค จ นีวรณานํ ปฐมชฺฌานํ ภาวยโต, ตทงฺคโวสคฺโค จ ทิฏฺฐิคตานํ นิพฺเพธ\-ภาคิยํ สมาธิํ ภาวยโต, สมุจฺเฉท\-โวสคฺโค จ}

\prose{\csromanfolio{399}โลกุตฺตรํ ขยคามิ\-มคฺคํ ภาวยโต, ปฏิปฺปสฺสทฺธิ\-โวสคฺโค จ ผลกฺขเณ, นิสฺสรณ\-โวสคฺโค จ นิโรโธ นิพฺพานํ. สมฺมา\-สมาธิสฺส อิเม ปญฺจ โวสคฺคา. อิเมสุ ปญฺจสุ โวสคฺเคสุ ฉนฺทชาโต โหติ สทฺธาธิมุตฺโต, จิตฺตํ จสฺส สฺวาธิฏฺฐิตํ. สมฺมา\-สมาธิสฺส อิเม ปญฺจ วิเวกา ปญฺจ วิราคา ปญฺจ นิโรธา ปญฺจ โวสคฺคา ทฺวาทส นิสฺสยา.}

\proseitem{26}{เสยฺยถาปิ ภิกฺขเว เย เกจิ พลกรณียา กมฺมนฺตา กรียนฺติ, สพฺเพ เต ปถวิํ นิสฺสาย ปถวิยํ ปติฏฺฐาย, เอวเมเต พลกรณียา กมฺมนฺตา กรียนฺติ. เอวเมวํ โข ภิกฺขเว ภิกฺขุ สีลํ นิสฺสาย สีเล ปติฏฺฐาย สตฺต โพชฺฌงฺเค ภาเวติ, สตฺต โพชฺฌงฺเค พหุลีกโรติ ฯเปฯ. สตฺต โพชฺฌงฺเค ภาเวนฺโต สตฺต โพชฺฌงฺเค พหุลีกโรนฺโต วุทฺธิํ วิรุฬฺหิํ เวปุลฺลํ ปาปุณาติ ธมฺเมสุ ฯเปฯ. ปญฺจ พลานิ ภาเวติ. ปญฺจ พลานิ พหุลีกโรติ ฯเปฯ. ปญฺจ พลานิ ภาเวนฺโต ปญฺจ พลานิ พหุลีกโรนฺโต วุทฺธิํ วิรุฬฺหิํ เวปุลฺลํ ปาปุณาติ ธมฺเมสุ ฯเปฯ. ปญฺจินฺทฺริยานิ ภาเวติ, ปญฺจินฺทฺริยานิ พหุลีกโรติ ฯเปฯ.}

\prose{เสยฺยถาปิ ภิกฺขเว เย เกจิ พีชคาม\-ภูตคามา วุทฺธิํ วิรุฬฺหิํ เวปุลฺลํ อาปชฺชนฺติ, สพฺเพ เต ปถวิํ นิสฺสาย ปถวิยํ ปติฏฺฐาย, เอวเมเต พีชคาม\-ภูตคามา วุทฺธิํ วิรุฬฺหิํ เวปุลฺลํ อาปชฺชนฺติ. เอวเมวํ โข ภิกฺขเว ภิกฺขุ สีลํ นิสฺสาย สีเล ปติฏฺฐาย ปญฺจินฺทฺริยานิ ภาเวนฺโต ปญฺจินฺทฺริยานิ พหุลีกโรนฺโต วุทฺธิํ วิรุฬฺหิํ เวปุลฺลํ ปาปุณาติ ธมฺเมสุ.}

\prose{กถญฺจ ภิกฺขเว ภิกฺขุ สีลํ นิสฺสาย สีเล ปติฏฺฐาย ปญฺจินฺทฺริยานิ ภาเวนฺโต ปญฺจินฺทฺริยานิ พหุลีกโรนฺโต วุทฺธิํ วิรุฬฺหิํ เวปุลฺลํ ปาปุณาติ ธมฺเมสุ? อิธ ภิกฺขเว ภิกฺขุ สทฺธินฺทฺริยํ ภาเวติ วิเวกนิสฺสิตํ วิราคนิสฺสิตํ นิโรธ\-นิสฺสิตํ โวสคฺค\-ปริณามิํ. วีริยินฺทฺริยํ ภาเวติ ฯเปฯ. สตินฺทฺริยํ ภาเวติ ฯเปฯ. สมาธินฺทฺริยํ ภาเวติ ฯเปฯ. ปญฺญินฺทฺริยํ ภาเวติ วิเวกนิสฺสิตํ วิราคนิสฺสิตํ นิโรธ\-นิสฺสิตํ โวสคฺค\-ปริณามิํ. เอวํ โข ภิกฺขเว ภิกฺขุ สีลํ นิสฺสาย ฯเปฯ. ปาปุณาติ ธมฺเมสูติ.}

\csromanmatikaanchor{mtk.127}
\csromantocmarkref{subsection}{2. อินฺทฺริยนิทฺเทส}{mtk.127}
\bookmark[dest={mtk.127},level=2]{2. อินฺทฺริยนิทฺเทส}
\csromanheader{2}{2. \textbf{อินฺทฺริยนิทฺเทส}}
\proseitem{27}{สทฺธิ\-นฺทฺริยสฺส ปญฺจ วิเวกา, ปญฺจ วิราคา, ปญฺจ นิโรธา, ปญฺจ โวสคฺคา, ทฺวาทส นิสฺสยา. วีริยิ\-นฺทฺริยสฺส ฯเปฯ. สตินฺทฺริยสฺส ฯเปฯ. สมาธิ\-นฺทฺริยสฺส ฯเปฯ. \csromanfolio{400}ปญฺญินฺทฺริยสฺส ปญฺจ วิเวกา, ปญฺจ วิราคา, ปญฺจ นิโรธา, ปญฺจ โวสคฺคา, ทฺวาทส นิสฺสยา.}

\prose{สทฺธิ\-นฺทฺริยสฺส กตเม ปญฺจ วิเวกา? วิกฺขมฺภน\-วิเวโก ตทงฺควิเวโก สมุจฺเฉท\-วิเวโก ปฏิปฺปสฺสทฺธิ\-วิเวโก นิสฺสรณ\-วิเวโก. วิกฺขมฺภน\-วิเวโก จ นีวรณานํ ปฐมชฺฌานํ ภาวยโต, ตทงฺควิเวโก จ ทิฏฺฐิคตานํ นิพฺเพธ\-ภาคิยํ สมาธิํ ภาวยโต, สมุจฺเฉท\-วิเวโก จ โลกุตฺตรํ ขยคามิ\-มคฺคํ ภาวยโต, ปฏิปฺปสฺสทฺธิ\-วิเวโก จ ผลกฺขเณ, นิสฺสรณ\-วิเวโก จ นิโรโธ นิพฺพานํ. สทฺธิ\-นฺทฺริยสฺส อิเม ปญฺจ วิเวกา. อิเมสุ ปญฺจสุ วิเวเกสุ ฉนฺทชาโต โหติ สทฺธาธิมุตฺโต, จิตฺตํ จสฺส สฺวาธิฏฺฐิตํ ฯเปฯ. สทฺธิ\-นฺทฺริยสฺส อิเม ปญฺจ วิเวกา, ปญฺจ วิราคา, ปญฺจ นิโรธา, ปญฺจ โวสคฺคา, ทฺวาทส นิสฺสยา.}

\prosewithcloserruled{วีริยิ\-นฺทฺริยสฺส ฯเปฯ. สตินฺทฺริยสฺส ฯเปฯ. สมาธิ\-นฺทฺริยสฺส ฯเปฯ. ปญฺญินฺทฺริยสฺส กตเม ปญฺจ วิเวกา. วิกฺขมฺภน\-วิเวโก ตทงฺควิเวโก สมุจฺเฉท\-วิเวโก ปฏิปฺปสฺสทฺธิ\-วิเวโก นิสฺสรณ\-วิเวโก ฯเปฯ. ปญฺญินฺทฺริยสฺส อิเม ปญฺจ วิเวกา, ปญฺจ วิราคา, ปญฺจ นิโรธา, ปญฺจ โวสคฺคา, ทฺวาทส นิสฺสยาติ.}{วิเวกกถา นิฏฺฐิตา.}
\csromanmatikaanchor{mtk.128}
\csromantocmarkref{section}{5. จริยากถา}{mtk.128}
\bookmark[dest={mtk.128},level=1]{5. จริยากถา}
\csromanheader{1}{5. \textbf{จริยากถา}}
\proseitem{28}{จริยาติ อฏฺฐ จฺอริยาโย อิริยาปถ\-จริยา อายตนจริยา สติจริยา สมาธิจริยา ญาณจริยา มคฺคจริยา ปตฺติจริยา โลกตฺถจริยาติ.}

\prose{อิริยาปถ\-จริยาติ จตูสุ อิริยาปเถสุ. อายตน\-จริยาติ ฉสุ อชฺฌตฺติก\-พาหิเรสุ อายตเนสุ. สติจริยาติ จตูสุ สติปฏฺฐาเนสุ. สมาธิจริยาติ จตูสุ ฌาเนสุ. ญาณจริยาติ จตูสุ อริยสจฺเจสุ. มฺอคฺคจริยาติ จตูสุ อริยมคฺเคสุ. ปตฺติจริยาติ จตูสุ สามญฺญผเลสุ. โลกตฺถจริยาติ ตถาคเตสุ อรหนฺเตสุ สมฺมา\-สมฺพุทฺเธสุ ปเทเส\csromansharedfootnote{4c57b1903e21cccd}{ปเทโส (สฺยา)} ปจฺเจก\-พุทฺเธสุ ปเทเส สาวเกสุ.}

\csromanmatikaanchor{mtk.129}
\csromantocmarkref{section}{6. ปาฏิหาริยกถา}{mtk.129}
\bookmark[dest={mtk.129},level=1]{6. ปาฏิหาริยกถา}
\csromanheader{1}{6. ปาฏิหาริย\-กถา}
\prose{\csromanfolio{401}อิริยาปถ\-จริยา จ ปณิธิ\-สมฺปนฺนานํ. อายตนจริยา จ อินฺทฺริเยสุ คุตฺตทฺวารานํ. สติจริยา จ อปฺปมาท\-วิหารีนํ. สมาธิจริยา จ อธิจิตฺตม\-นุยุตฺตานํ. ญาณจริยา จ พุทฺธิ\-สมฺปนฺนานํ. มคฺคจริยา จ สมฺมา\-ปฏิปนฺนานํ. ปตฺติจริยา จ อธิคต\-ผลานํ. โลกตฺถจริยา จ ตถาคตานํ อรหนฺตานํ สมฺมา\-สมฺพุทฺธานํ ปเทเส ปจฺเจก\-พุทฺธานํ ปเทเส สาวกานํ. อิมา อฏฺฐ จริยาโย.}

\proseitem{29}{อปราปิ อฏฺฐ จริยาโย, อธิมุจฺจนฺโต สทฺธาย จรติ, ปคฺคณฺหนฺโต วีริเยน จรติ, อุปฏฺฐาเปนฺโต สติยา จรติ, อวิกฺเขปํ กโรนฺโต สมาธินา จรติ, ปชานนฺโต ปญฺญาย จรติ, วิชานนฺโต วิญฺญาณ\-จริยาย จรติ, เอวํ ปฏิปนฺนสฺส กุสลา ธมฺมา อายาเปนฺตีติ อายตน\-จริยาย จรติ, เอวํ ปฏิปนฺโน วิเสสมธิคจฺฉตีติ วิเสสจริยาย จรติ. อิมา อฏฺฐ จริยาโย.}

\prosewithcloserruled{อปราปิ อฏฺฐ จริยาโย, ทสฺสนจริยา จ สมฺมาทิฏฺฐิยา, อภินิโรปน\-จริยา จ สมฺมา\-สงฺกปฺปสฺส, ปริคฺคห\-จริยา จ สมฺมาวาจาย, สมุฏฺฐาน\-จริยา จ สมฺมา\-กมฺมนฺตสฺส, โวทานจริยา จ สมฺมาอาชีวสฺส, ปคฺคหจริยา จ สมฺมาวายามสฺส, อุปฏฺฐาน\-จริยา จ สมฺมาสติยา, อวิกฺเขป\-จริยา จ สมฺมา\-สมาธิสฺส. อิมา อฏฺฐ จริยาโยติ.}{จริยากถา นิฏฺฐิตา.}
\csromanheader{2}{6. ปาฏิหาริย\-กถา}
\proseitem{30}{ตีณิมานิ ภิกฺขเว\csromansharedfootnote{7062f66464d47643}{อํ 1. 170 ปิฏฺเฐ ปสฺสิตพฺพา.} ปาฏิหาริยานิ. กตมานิ ตีณิ? อิทฺธิ\-ปาฏิหาริยํ อาเทสนา\-ปาฏิหาริยํ อนุสาสนี\-ปาฏิหาริยํ.}

\prose{กตมญฺจ ภิกฺขเว อิทฺธิ\-ปาฏิหาริยํ? อิธ ภิกฺขเว เอกจฺโจ อเนกวิหิตํ อิทฺธิวิธํ ปจฺจนุโภติ, เอโกปิ หุตฺวา พหุธา โหติ, พหุธาปิ หุตฺวา เอโก โหติ, อาวิภาวํ ติโรภาวํ ฯเปฯ ยาว พฺรหฺมโลกาปิ กาเยน วสํ วตฺเตติ. อิทํ วุจฺจติ ภิกฺขเว อิทฺธิ\-ปาฏิหาริยํ. (1)}

\prose{\csromanfolio{402}กตมญฺจ ภิกฺขเว อาเทสนา\-ปาฏิหาริยํ? อิธ ภิกฺขเว เอกจฺโจ นิมิตฺเตน อาทิสติ “เอวมฺปิ เต มโน, อิตฺถมฺปิ เต มโน, อิติปิ เต จิตฺตนฺ”ติ. โส พหุํ เจปิ อาทิสติ, ตเถว ตํ โหติ, โน อญฺญถา.  อิธ ปน ภิกฺขเว เอกจฺโจ น เหว โข นิมิตฺเตน อาทิสติ, อปิ จ โข มนุสฺสานํ วา อมนุสฺสานํ วา เทวตานํ วา สทฺทํ สุตฺวา อาทิสติ “เอวมฺปิ เต มโน, อิตฺถมฺปิ เต มโน, อิติปิ เต จิตฺตนฺ”ติ. โส พหุํ เจปิ อาทิสติ, ตเถว ตํ โหติ, โน อญฺญถา.  อิธ ปน ภิกฺขเว เอกจฺโจ น เหว โข นิมิตฺเตน อาทิสติ, นปิ มนุสฺสานํ วา อมนุสฺสานํ วา เทวตานํ วา สทฺทํ สุตฺวา อาทิสติ, อปิ จ โข วิตกฺกยโต วิจารยโต วิตกฺก\-วิปฺผาร\-สทฺทํ สุตฺวา อาทิสติ “เอวมฺปิ เต มโน, อิตฺถมฺปิ เต มโน, อิติปิ เต จิตฺตนฺ”ติ. โส พหุํ เจปิ อาทิสติ, ตเถว ตํ โหติ, โน อญฺญถา.  อิธ ปน ภิกฺขเว เอกจฺโจ นเหว โข นิมิตฺเตน อาทิสติ, นปิ มนุสฺสานํ วา อมนุสฺสานํ วา เทวตานํ วา สทฺทํ สุตฺวา อาทิสติ, นปิ วิตกฺกยโต วิจารยโต วิตกฺก\-วิปฺผาร\-สทฺทํ สุตฺวา อาทิสติ, อปิ จ โข อวิตกฺกํ อวิจารํ สมาธิํ สมาปนฺนสฺส เจตสา เจโต ปริจฺจ ปชานาติ “ยถา อิมสฺส โภโต มโนสงฺขารา ปณิหิตา อิมสฺส\csromansharedfootnote{e70df65b01fb13f2}{ตถา อิมสฺส (สฺยา), ยถา อิมสฺส (ก) อํ 1. 171 ปิฏฺเฐ ปสฺสิตพฺพา.} จิตฺตสฺส อนนฺตรา อมุกํ นาม วิตกฺกํ วิตกฺกยิสฺสตี”ติ\csromansharedfootnote{e7d7f450da7cc6bc}{วิตกฺเกสฺสตีติ (สฺยา) อํ 1. 171 ปิฏฺเฐปิ.}. โส พหุํ เจปิ อาทิสติ, ตเถว ตํ โหติ, โน อญฺญถา. อิทํ วุจฺจติ ภิกฺขเว อาเทสนา\-ปาฏิหาริยํ. (2)}

\prose{กตมญฺจ ภิกฺขเว อนุสาสนี\-ปาฏิหาริยํ? อิธ ภิกฺขเว เอกจฺโจ เอวมนุสาสติ “เอวํ วิตกฺเกถ, มา เอวํ วิตกฺกยิตฺถ. เอวํ มนสิ กโรถ, มา เอวํ มนสากตฺถ. อิทํ ปชหถ, อิทํ อุปสมฺปชฺช วิหรถา”ติ, อิทํ วุจฺจติ ภิกฺขเว อนุสาสนี\-ปาฏิหาริยํ. อิมานิ โข ภิกฺขเว ตีณิ ปาฏิหาริยานิ. (3)}

\proseitem{31}{เนกฺขมฺมํ อิชฺฌตีติ อิทฺธิ, กามจฺฉนฺทํ ปฏิหรตีติ ปาฏิหาริยํ. เย เตน เนกฺขมฺเมน สมนฺนาคตา, สพฺเพ เต วิสุทฺธจิตฺตา อนาวิลสงฺกปฺปาติ อาเทสนา\-ปาฏิหาริยํ. “ตํ โข ปน เนกฺขมฺมํ เอวํ อาเสวิตพฺพํ เอวํ}

\csromanheader{2}{6. ปาฏิหาริย\-กถา}
\prose{\csromanfolio{403}ภาเวตพฺพํ, เอวํ พหุลีกาตพฺพํ, เอวํ ตทนุธมฺมตา สติ อุปฏฺฐาเปตพฺพา”ติ อนุสาสนี\-ปาฏิหาริยํ. (1)}

\prose{อพฺยาปาโท อิชฺฌตีติ อิทฺธิ, พฺยาปาทํ ปฏิหรตีติ ปาฏิหาริยํ. เย เตน อพฺยาปาเทน สมนฺนาคตา, สพฺเพ เต วิสุทฺธจิตฺตา อนาวิลสงฺกปฺปาติ อาเทสนา\-ปาฏิหาริยํ. “โส โข ปน อพฺยาปาโท เอวํ อาเสวิตพฺโพ, เอวํ ภาเวตพฺโพ, เอวํ พหุลีกาตพฺโพ, เอวํ ตทนุธมฺมตา สติ อุปฏฺฐาเปตพฺพา”ติ อนุสาสนี\-ปาฏิหาริยํ. (2)}

\prose{อาโลกสญฺญา อิชฺฌตีติ อิทฺธิ, ถินมิทฺธํ ปฏิหรตีติ ปาฏิหาริยํ. เย ตาย อาโลกสญฺญาย สมนฺนาคตา, สพฺเพ เต วิสุทฺธจิตฺตา อนาวิลสงฺกปฺปาติ อาเทสนา\-ปาฏิหาริยํ. “สา โข ปน อาโลกสญฺญา เอวํ อาเสวิตพฺพา, เอวํ ภาเวตพฺพา, เอวํ พหุลีกาตพฺพา, เอวํ ตทนุธมฺมตา สติ อุปฏฺฐาเปตพฺพา”ติ อนุสาสนี\-ปาฏิหาริยํ. (3)}

\prose{อวิกฺเขโป อิชฺฌตีติ อิทฺธิ, อุทฺธจฺจํ ปฏิหรตีติ ปาฏิหาริยํ. เย เตน อวิกฺเขเปน สมนฺนาคตา, สพฺเพ เต วิสุทฺธจิตฺตา อนาวิลสงฺกปฺปาติ อาเทสนา\-ปาฏิหาริยํ. “โส โข ปน อวิกฺเขโป เอวํ อาเสวิตพฺโพ, เอวํ ภาเวตพฺโพ, เอวํ พหุลีกาตพฺโพ, เอวํ ตทนุธมฺมตา สติ อุปฏฺฐาเปตพฺพา”ติ อนุสาสนี\-ปาฏิหาริยํ. (4)}

\prose{ธมฺม\-ววตฺถานํ อิชฺฌตีติ อิทฺธิ, วิจิกิจฺฉํ ปฏิหรตีติ ปาฏิหาริยํ. เย เตน ธมฺม\-ววตฺถาเนน สมนฺนาคตา, สพฺเพ เต วิสุทฺธจิตฺตา อนาวิลสงฺกปฺปาติ อาเทสนา\-ปาฏิหาริยํ. “ตํ โข ปน ธมฺม\-ววตฺถานํ เอวํ อาเสวิตพฺพํ, เอวํ ภาเวตพฺพํ, เอวํ พหุลีกาตพฺพํ, เอวํ ตทนุธมฺมตา สติ อุปฏฺฐาเปตพฺพา”ติ อนุสาสนี\-ปาฏิหาริยํ. (5)}

\prose{ญาณํ อิชฺฌตีติ อิทฺธิ, อวิชฺชํ ปฏิหรตีติ ปาฏิหาริยํ. เย เตน ญาเณน สมนฺนาคตา, สพฺเพ เต วิสุทฺธจิตฺตา อนาวิลสงฺกปฺปาติ อาเทสนา\-ปาฏิหาริยํ. “ตํ โข ปน ญาณํ เอวํ อาเสวิตพฺพํ, เอวํ ภาเวตพฺพํ, เอวํ พหุลีกาตพฺพํ, เอวํ ตทนุธมฺมตา สติ อุปฏฺฐาเปตพฺพา”ติ อนุสาสนี\-ปาฏิหาริยํ. (6)}

\prose{ปาโมชฺชํ อิชฺฌตีติ อิทฺธิ, อรติํ ปฏิหรตีติ ปาฏิหาริยํ. เย เตน ปาโมชฺเชน สมนฺนาคตา, สพฺเพ เต วิสุทฺธจิตฺตา อนาวิลสงฺกปฺปาติ อาเทสนา\-ปาฏิหาริยํ. “ตํ โข ปน ปาโมชฺชํ เอวํ อาเสวิตพฺพํ, เอวํ}

\prose{\csromanfolio{404}ภาเวตพฺพํ, เอวํ พหุลีกาตพฺพํ, เอวํ ตทนุธมฺมตา สติ อุปฏฺฐาเปตพฺพา”ติ อนุสาสนี\-ปาฏิหาริยํ ฯเปฯ. (7)}

\prose{ปฐมํ ฌานํ อิชฺฌตีติ อิทฺธิ, นีวรเณ ปฏิหรตีติ ปาฏิหาริยํ. เย เตน ปฐเมน ฌาเนน สมนฺนาคตา, สพฺเพ เต วิสุทฺธจิตฺตา อนาวิลสงฺกปฺปาติ อาเทสนา\-ปาฏิหาริยํ. “ตํ โข ปน ปฐมํ ฌานํ เอวํ อาเสวิตพฺพํ, เอวํ ภาเวตพฺพํ, เอวํ พหุลีกาตพฺพํ, เอวํ ตทนุธมฺมตา สติ อุปฏฺฐาเปตพฺพา”ติ อนุสาสนี\-ปาฏิหาริยํ ฯเปฯ. \mbox{(26, 33)}}

\prose{อรหตฺตมคฺโค อิชฺฌตีติ อิทฺธิ, สพฺพกิเลเส ปฏิหรตีติ ปาฏิหาริยํ. เย เตน อรหตฺต\-มคฺเคน สมนฺนาคตา, สพฺเพ เต วิสุทฺธจิตฺตา อนาวิลสงฺกปฺปาติ อาเทสนา\-ปาฏิหาริยํ. “โส โข ปน อรหตฺตมคฺโค เอวํ อาเสวิตพฺโพ, เอวํ ภาเวตพฺโพ, เอวํ พหุลีกาตพฺโพ, เอวํ ตทนุธมฺมตา สติ อุปฏฺฐาเปตพฺพา”ติ อนุสาสนี\-ปาฏิหาริยํ. \mbox{(4, 37)}}

\proseitemwithcloserruled{32}{เนกฺขมฺมํ อิชฺฌตีติ อิทฺธิ, กามจฺฉนฺทํ ปฏิหรตีติ ปาฏิหาริยํ. ยา จ อิทฺธิ ยญฺจ ปาฏิหาริยํ, อิทํ วุจฺจติ อิทฺธิ\-ปาฏิหาริยํ.  อพฺยาปาโท อิชฺฌตีติ อิทฺธิ, พฺยาปาทํ ปฏิหรตีติ ปาฏิหาริยํ. ยา จ อิทฺธิ ยญฺจ ปาฏิหาริยํ, อิทํ วุจฺจติ อิทฺธิ\-ปาฏิหาริยํ.  อาโลกสญฺญา อิชฺฌตีติ อิทฺธิ, ถินมิทฺธํ ปฏิหรตีติ ปาฏิหาริยํ ฯเปฯ. อรหตฺตมคฺโค อิชฺฌตีติ อิทฺธิ, สพฺพกิเลเส ปฏิหรตีติ ปาฏิหาริยํ. ยา จ อิทฺธิ ยญฺจ ปาฏิหาริยํ, อิทํ วุจฺจติ อิทฺธิ\-ปาฏิหาริยนฺติ.}{ปาฏิหาริย\-กถา นิฏฺฐิตา.}
\csromanmatikaanchor{mtk.130}
\csromantocmarkref{section}{7. สมสีสกถา}{mtk.130}
\bookmark[dest={mtk.130},level=1]{7. สมสีสกถา}
\csromanheader{1}{7. \textbf{สมสีสกถา}}
\proseitem{33}{สพฺพธมฺมานํ\csromansharedfootnote{336545f304bfc39d}{สพฺเพสํ ธมฺมานํ (ก) 98 ปิฏฺเฐ ปสฺสิตพฺพา.} สมฺมา สมุจฺเฉเท นิโรเธ จ อนุปฏฺฐานตา \textbf{ปญฺญา สมสีสฏฺเฐ ญาณํ.}}

\prose{สพฺพ\-ธมฺมานนฺติ ปญฺจกฺขนฺธา ทฺวาทสา\-ยตนานิ อฏฺฐารส ธาตุโย กุสลา ธมฺมา อกุสลา ธมฺมา อพฺยากตา ธมฺมา กามาวจรา ธมฺมา รูปาวจรา ธมฺมา อรูปาวจรา ธมฺมา อปริยาปนฺนา ธมฺมา.  สมฺมา สมุจฺเฉเทติ}

\csromanheader{2}{7. สมสีสกถา}
\prose{\csromanfolio{405}เนกฺขมฺเมน กามจฺฉนฺทํ สมฺมา สมุจฺฉินฺทติ, อพฺยาปาเทน พฺยาปาทํ สมฺมา สมุจฺฉินฺทติ, อาโลกสญฺญาย ถินมิทฺธํ สมฺมา สมุจฺฉินฺทติ, อวิกฺเขเปน อุทฺธจฺจํ สมฺมา สมุจฺฉินฺทติ, ธมฺม\-ววตฺถาเนน วิจิกิจฺฉํ สมฺมา สมุจฺฉินฺทติ, ญาเณน อวิชฺชํ สมฺมา สมุจฺฉินฺทติ, ปาโมชฺเชน อรติํ สมฺมา สมุจฺฉินฺทติ, ปฐเมน ฌาเนน นีวรเณ สมฺมา สมุจฺฉินฺทติ ฯเปฯ อรหตฺต\-มคฺเคน สพฺเพกิเลเส สมฺมา สมุจฺฉินฺทติ.}

\prose{นิโรเธติ เนกฺขมฺเมน กามจฺฉนฺทํ นิโรเธติ, อพฺยาปาเทน พฺยาปาทํ นิโรเธติ, อาโลกสญฺญาย ถินมิทฺธํ นิโรเธติ, อวิกฺเขเปน อุทฺธจฺจํ นิโรเธติ, ธมฺม\-ววตฺถาเนน วิจิกิจฺฉํ นิโรเธติ, ญาเณน อวิชฺชํ นิโรเธติ ปาโมชฺเชน อรติํ นิโรเธติ, ปฐเมน ฌาเนน นีวรเณ นิโรเธติ ฯเปฯ อรหตฺต\-มคฺเคน สพฺพกิเลเส นิโรเธติ.}

\prose{อนุป\-ฏฺฐานตาติ เนกฺขมฺมํ ปฏิลทฺธสฺส กามจฺฉนฺโท น อุปฏฺฐาติ, อพฺยาปาทํ ปฏิลทฺธสฺส พฺยาปาโท น อุปฏฺฐาติ, อาโลกสญฺญํ ปฏิลทฺธสฺส ถินมิทฺธํ น อุปฏฺฐาติ, อวิกฺเขปํ ปฏิลทฺธสฺส อุทฺธจฺจํ น อุปฏฺฐาติ, ธมฺม\-ววตฺถานํ ปฏิลทฺธสฺส วิจิกิจฺฉา น อุปฏฺฐาติ, ญาณํ ปฏิลทฺธสฺส อวิชฺชา น อุปฏฺฐาติ, ปาโมชฺชํ ปฏิลทฺธสฺส อรติ น อุปฏฺฐาติ, ปฐมํ ฌานํ ปฏิลทฺธสฺส นีวรณา น อุปฏฺฐนฺติ ฯเปฯ อรหตฺตมคฺคํ ปฏิลทฺธสฺส สพฺพกิเลสา น อุปฏฺฐนฺติ.}

\prose{สมนฺติ กามจฺฉนฺทสฺส ปหีนตฺตา เนกฺขมฺมํ สมํ, พฺยาปาทสฺส ปหีนตฺตา อพฺยาปาโท สมํ, ถินมิทฺธสฺส ปหีนตฺตา อาโลกสญฺญา สมํ, อุทฺธจฺจสฺส ปหีนตฺตา อวิกฺเขโป สมํ, วิจิกิจฺฉาย ปหีนตฺตา ธมฺม\-ววตฺถานํ สมํ, อวิชฺชาย ปหีนตฺตา ญาณํ สมํ, อรติยา ปหีนตฺตา ปาโมชฺชํ สมํ, นีวรณานํ ปหีนตฺตา ปฐมํ ฌานํ สมํ ฯเปฯ สพฺพกิเลสานํ ปหีนตฺตา อรหตฺตมคฺโค สมํ.}

\prosewithcloser{สีสนฺติ เตรส สีสานิ, ปลิโพธ\-สีสญฺจ ตณฺหา, วินิพนฺธน\-สีสญฺจ มาโน, ปรามาส\-สีสญฺจ ทิฏฺฐิ, วิกฺเขป\-สีสญฺจ อุทฺธจฺจํ, สํกิเลส\-สีสญฺจ อวิชฺชา, อธิโมกฺข\-สีสญฺจ สทฺธา, ปคฺคหสีสญฺจ วีริยํ, อุปฏฺฐาน\-สีสญฺจ สติ, อวิกฺเขป\-สีสญฺจ สมาธิ, ทสฺสนสีสญฺจ ปญฺญา, ปวตฺตสีสญฺจ ชีวิตินฺทฺริยํ, โคจรสีสญฺจ วิโมกฺโข, สงฺขาร\-สีสญฺจ นิโรโธติ.}{สมสีสกถา นิฏฺฐิตา.}
\csromanmatikaanchor{mtk.131}
\csromantocmarkref{section}{8. สติปฏฺฐานกถา}{mtk.131}
\bookmark[dest={mtk.131},level=1]{8. สติปฏฺฐานกถา}
\csromanheader{1}{8. สติปฏฺฐาน\-กถา}
\proseitem{34}{\csromanfolio{406}สาวตฺถิ\-นิทานํ.  จตฺตาโรเม ภิกฺขเว สติปฏฺฐานา. กตเม จตฺตาโร? อิธ ภิกฺขเว ภิกฺขุ กาเย กายานุปสฺสี วิหรติ อาตาปี สมฺปชาโน สติมา วิเนยฺย โลเก อภิชฺฌา\-โทมนสฺสํ, เวทนาสุ. จิตฺเต. ธมฺเมสุ ธมฺมานุปสฺสี วิหรติ อาตาปี สมฺปชาโน สติมา วิเนยฺย โลเก อภิชฺฌา\-โทมนสฺสํ. อิเม โข ภิกฺขเว จตฺตาโร สติปฏฺฐานาติ.}

\proseitem{35}{(ก) กถํ กาเย กายานุปสฺสี วิหรติ? อิเธกจฺโจ ปถวีกายํ อนิจฺจโต อนุปสฺสติ, โน นิจฺจโต, ทุกฺขโต อนุปสฺสติ, โน สุขโต, อนตฺตโต อนุปสฺสติ, โน อตฺตโต, นิพฺพินฺทติ, โน นนฺทติ, วิรชฺชติ, โน รชฺชติ, นิโรเธติ, โน สมุเทติ, ปฏินิสฺสชฺชติ, โน อาทิยติ. อนิจฺจโต อนุปสฺสนฺโต นิจฺจสญฺญํ ปชหติ, ทุกฺขโต อนุปสฺสนฺโต สุขสญฺญํ ปชหติ, อนตฺตโต อนุปสฺสนฺโต อตฺตสญฺญํ ปชหติ, นิพฺพินฺทนฺโต นนฺทิํ ปชหติ, วิรชฺชนฺโต ราคํ ปชหติ, นิโรเธนฺโต สมุทยํ ปชหติ, ปฏินิสฺสชฺชนฺโต อาทานํ ปชหติ, อิเมหิ สตฺตหิ อากาเรหิ กายํ อนุปสฺสติ. กาโย อุปฏฺฐานํ, โน สติ, สติ อุปฏฺฐานญฺเจว สติ จ, ตาย สติยา เตน ญาเณน ตํ กายํ อนุปสฺสติ, เตน วุจฺจติ กาเย กายานุปสฺสนา\-สติปฏฺฐานา.}

\prose{ภาวนาติ จตสฺโส ภาวนา, ตตฺถ ชาตานํ ธมฺมานํ อนติวตฺตน\-ฏฺเฐน ภาวนา, อินฺทฺริยานํ เอกรสฏฺเฐน ภาวนา, ตทุปค\-วีริย\-วาห\-นฏฺเฐน ภาวนา, อาเสวนฏฺเฐน ภาวนา.}

\prose{อิเธกจฺโจ อาโปกายํ ฯเปฯ เตโชกายํ. วาโยกายํ. เกสกายํ. โลมกายํ. ฉวิกายํ. จมฺมกายํ. มํสกายํ. รุธิรกายํ. นฺหารุกายํ\csromansharedfootnote{d0ffc344b22307f1}{นหารุกายํ (สฺยา)}. อฏฺฐิกายํ. อฏฺฐิมิญฺช\-กายํ อนิจฺจโต อนุปสฺสติ, โน นิจฺจโต, ทุกฺขโต อนุปสฺสติ, โน สุขโต, อนตฺตโต อนุปสฺสติ, โน อตฺตโต, นิพฺพินฺทติ, โน นนฺทติ, วิรชฺชติ, โน รชฺชติ, นิโรเธติ, โน สมุเทติ, ปฏินิสฺสชฺชติ, โน อาทิยติ. อนิจฺจโต อนุปสฺสนฺโต นิจฺจสญฺญํ ปชหติ, ทุกฺขโต อนุปสฺสนฺโต สุขสญฺญํ ปชหติ, อนตฺตโต อนุปสฺสนฺโต อตฺตสญฺญํ ปชหติ,}

\csromanheader{2}{8. สติปฏฺฐาน\-กถา}
\prose{\csromanfolio{407}นิพฺพินฺทนฺโต นนฺทิํ ปชหติ, วิรชฺชนฺโต ราคํ ปชหติ, นิโรเธนฺโต สมุทยํ ปชหติ, ปฏินิสฺสชฺชนฺโต อาทานํ ปชหติ. อิเมหิ สตฺตหิ อากาเรหิ กายํ อนุปสฺสภิ. กาโย อุปฏฺฐานํ, โน สติ, สติ อุปฏฺฐานญฺเจว สติ จ, ตาย สติยา เตเน ญาเณน ตํ กายํ อนุปสฺสติ, เตน วุจฺจติ กาเย กายานุปสฺสนา\-สติปฏฺฐานา.}

\prose{ภาวนาติ จตสฺโส ภาวนา, ตตฺถ ชาตานํ ธมฺมานํ อนติวตฺตน\-ฏฺเฐน ภาวนา, อินฺทฺริยานํ เอกรสฏฺเฐน ภาวนา, ตทุปค\-วีริย\-วาห\-นฏฺเฐน ภาวนา, อาเสวนฏฺเฐน ภาวนา, เอวํ กาเย กายานุปสฺสี วิหรติ. (1)}

\prose{(ข) กถํ เวทนาสุ เวทนานุปสฺสี วิหรติ? อิเธกจฺโจ สุขํ เวทนํ อนิจฺจโต อนุปสฺสติ, โน นิจฺจโต ฯเปฯ ปฏินิสฺสชฺชติ, โน อาทิยติ. อนิจฺจโต อนุปสฺสนฺโต นิจฺจสญฺญํ ปชหติ ฯเปฯ ปฏินิสฺสชฺชนฺโต อาทานํ ปชหติ. อิเมหิ สตฺตหิ อากาเรหิ เวทนํ อนุปสฺสติ. เวทนา อุปฏฺฐานํ, โน สติ, สติ อุปฏฺฐานญฺเจว สติ จ, ตาย สติยา เตน ญาเณน ตํ เวทนํ อนุปสฺสติ, เตน วุจฺจติ เวทนาสุ เวทนานุปสฺสนา\-สติปฏฺฐานา.}

\prose{ภาวนาติ จตสฺโส ภาวนา ฯเปฯ อาเสวนฏฺเฐน ภาวนา ฯเปฯ. อิเธกจฺโจ ทุกฺขํ เวทนํ ฯเปฯ อทุกฺขมสุขํ เวทนํ. สามิสํ สุขํ เวทนํ. นิรามิสํ สุขํ เวทนํ. สามิสํ ทุกฺขํ เวทนํ. นิรามิสํ ทุกฺขํ เวทนํ. สามิสํ อทุกฺขมสุขํ เวทนํ. นิรามิสํ อทุกฺขมสุขํ เวทนํ. จกฺขุ\-สมฺผสฺสชํ เวทนํ. โสต\-สมฺผสฺสชํ เวทนํ.ฆาน\-สมฺผสฺสชํ เวทนํ. ชิวฺหา\-สมฺผสฺสชํ เวทนํ. กาย\-สมฺผสฺสชํ เวทนํ. มโน\-สมฺผสฺสชํ เวทนํ อนิจฺจโต อนุปสฺสติ, โน นิจฺจโต ฯเปฯ ปฏินิสฺสชฺชติ, โน อาทิยติ. อนิจฺจโต อนุปสฺสนฺโต นิจฺจสญฺญํ ปชหติ ฯเปฯ ปฏินิสฺสชฺชนฺโต อาทานํ ปชหติ. อิเมหิ สตฺตหิ อากาเรหิ เวทนํ อนุปสฺสติ. เวทนา อุปฏฺฐานํ, โน สติ, สติ อุปฏฺฐานญฺเจว สติ จ, ตาย สติยา เตน ญาเณน ตํ เวทนํ อนุปสฺสติ, เตน วุจฺจติ เวทนาสุ เวทนานุปสฺสนา\-สติปฏฺฐานา.}

\prose{\textbf{ภาวนา}ติ จตสฺโส ภาวนา ฯเปฯ เอวํ เวทนาสุ เวทนานุปสฺสี วิหรติ. (2)}

\prose{\csromanfolio{408}(ค) กถํ จิตฺเต จิตฺตานุปสฺสี วิหรติ? อิเธกจฺโจ สราคํ จิตฺตํ อนิจฺจโต อนุปสฺสติ, โน นิจฺจโต ฯเปฯ ปฏินิสฺสชฺชติ, โน อาทิยติ. อนิจฺจโต อนุปสฺสนฺโต นิจฺจสญฺญํ ปชหติ ฯเปฯ ปฏินิสฺสชฺชนฺโต อาทานํ ปชหติ. อิเมหิ สตฺตหิ อากาเรหิ จิตฺตํ อนุปสฺสติ, จิตฺตํ อุปฏฺฐานํ, โน สติ, สติ อุปฏฺฐานญฺเจว สติ จ, ตาย สติยา เตน ญาเณน ตํ จิตฺตํ อนุปสฺสติ, เตน วุจฺจติ จิตฺเต จิตฺตานุปสฺสนา\-สติปฏฺฐานา.}

\prose{\textbf{ภาวนา}ติ จตสฺโส ภาวนา ฯเปฯ อาเสวนฏฺเฐน ภาวนา.}

\prose{อิเธกจฺโจ วีตราคํ จิตฺตํ ฯเปฯ สโทสํ จิตฺตํ. วีตโทสํ จิตฺตํ. สโมหํ จิตฺตํ. วีตโมหํ จิตฺตํ. สํขิตฺตํ จิตฺตํ. วิกฺขิตฺตํ จิตฺตํ. มหคฺคตํ จิตฺตํ. อมหคฺคตํ จิตฺตํ. สอุตฺตรํ จิตฺตํ. อนุตฺตรํ จิตฺตํ. สมาหิตํ จิตฺตํ. อสมาหิตํ จิตฺตํ. วิมุตฺตํ จิตฺตํ. อวิมุตฺตํ จิตฺตํ. จกฺขุวิญฺญาณํ. โสตวิญฺญาณํ. ฆานวิญฺญาณํ. ชิวฺหาวิญฺญาณํ. กายวิญฺญาณํ. มโนวิญฺญาณํ อนิจฺจโต อนุปสฺสติ, โน นิจฺจโต ฯเปฯ ปฏินิสฺสชฺชติ, โน อาทิยติ. อนิจฺจโต อนุปสฺสนฺโต นิจฺจสญฺญํ ปชหติ ฯเปฯ ปฏินิสฺสชฺชนฺโต อาทานํ ปชหติ. อิเมหิ สตฺตหิ อากาเรหิ จิตฺตํ อนุปสฺสติ. จิตฺตํ อุปฏฺฐานํ โน สติ, สติ อุปฏฺฐานญฺเจว สติ จ, ตาย สติยา เตน ญาเณน ตํ จิตฺตํ อนุปสฺสติ, เตน วุจฺจติ จิตฺเต จิตฺตานุปสฺสนา\-สติปฏฺฐานา.}

\prose{\textbf{ภาวนา}ติ จตสฺโส ภาวนา ฯเปฯ อาเสวนฏฺเฐน ภาวนา. เอวํ จิตฺเต จิตฺตานุปสฺสี วิหรติ. (3)}

\prose{(ฆ) กถํ ธมฺมฺ\textbf{เอสุ ธมฺมานุปสฺสี วิหรติ?} อิเธกจฺโจ ฐเปตฺวา กายํ ฐเปตฺวา เวทนํ ฐปตฺวา จิตฺตํ ตทวเสเส ธมฺเม อนิจฺจโต อนุปสฺสติ, โน นิจฺจโต, ทุกฺขโต อนุปสฺสติ, โน สุขโต, อนตฺตโต อนุปสฺสติ, โน อตฺตโต, นิพฺพินฺทติ, โน นนฺทติ, วิรชฺชติ, โน รชฺชติ, นิโรเธติ, โน สมุเทติ, ปฏินิสฺสชฺชติ, โน อาทิยติ. อนิจฺจโต อนุปสฺสนฺโต นิจฺจสญฺญํ ปชหติ, ทุกฺขโต อนุปสฺสนฺโต สุขสญฺญํ ปชหติ, อนตฺตโต อนุปสฺสนฺโต อตฺตสญฺญํ ปชหติ, นิพฺพินฺทนฺโต นนฺทิํ ปชหติ, วิรชฺชนฺโต ราคํ ปชหติ, นิโรเธนฺโต สมุทยํ ปชหติ, ปฏินิสฺสชฺชนฺโต อาทานํ ปชหติ. อิเมหิ สตฺตหิ อากาเรหิ เต ธมฺเม อนุปสฺสติ. ธมฺมา อุปฏฺฐานํ, โน สติ, สติ อุปฏฺฐานญฺเจว สติ จ, ตาย สติยา เตน ญาเณน}

\csromanmatikaanchor{mtk.132}
\csromantocmarkref{section}{9. วิปสฺสนากถา}{mtk.132}
\bookmark[dest={mtk.132},level=1]{9. วิปสฺสนากถา}
\prose{\csromanfolio{409}เต ธมฺเม อนุปสฺสติ, เตน วุจฺจติ ธมฺเมสุ ธมฺมานุปสฺสนา\-สติปฏฺฐานา.}

\prosewithcloserruled{ภาวนาติ จตสฺโส ภาวนา, ตตฺถ ชาตานํ ธมฺมานํ อนติวตฺตน\-ฏฺเฐน ภาวนา, อินฺทฺริยานํ เอกรสฏฺเฐน ภาวนา, ตทุปค\-วีริย\-วาห\-นฏฺเฐน ภาวนา อาเสวนฏฺเฐน ภาวนา. เอวํ ธมฺเมสุ ธมฺมานุปสฺสี วิหรตีติ. (4)}{สติปฏฺฐาน\-กถา นิฏฺฐิตา.}
\proseitem{36}{เอวํ เม สุตํ\csromandash{}เอกํ สมยํ ภควา สาวตฺถิยํ วิหรติ เชตวเน อนาถ\-ปิณฺฑิกสฺส อาราเม. ตตฺร โข ภควา ภิกฺขู อามนฺเตสิ “ภิกฺขโว”ติ. “ภทนฺเต”ติ เต ภิกฺขู ภควโต ปจฺจสฺโสสุํ. ภควา เอตทโวจ\csromandash{}}

\prose{โส วต ภิกฺขเว ภิกฺขุ กญฺจิ\csromansharedfootnote{9d989739e822b11c}{กิญฺจิ (ก) อํ 2. 385 ปิฏฺเฐ ปสฺสิตพฺพา.} สงฺขารํ นิจฺจโต สมนุปสฺสนฺโต อนุโลมิกาย ขนฺติยา สมนฺนาคโต ภวิสฺสตีติ เนตํ ฐานํ วิชฺชติ, อนุโลมิกาย ขนฺติยา อสมนฺนาคโต สมฺมตฺต\-นิยามํ โอกฺกมิสฺสตีติ เนตํ ฐานํ วิชฺชติ, สมฺมตฺต\-นิยามํ อโนกฺกมมาโน โสตาปตฺติผลํ วา สกทาคามิผลํ วา อนาคามิผลํ วา อรหตฺตํ วา\csromansharedfootnote{5d135471398bc9f3}{อรหตฺตผลํ วา (สฺยา, ก) อํ 2. 384 ปิฏฺเฐ ปสฺสิตพฺพา.} สจฺฉิกริสฺสตีติ เนตํ ฐานํ วิชฺชติ. (1)}

\prose{โส วต ภิกฺขเว ภิกฺขุ สพฺพสงฺขาเร อนิจฺจโต สมนุปสฺสนฺโต อนุโลมิกาย ขนฺติยา สมนฺนาคโต ภวิสฺสตีติ ฐานเมตํ วิชฺชติ, อนุโลมิกาย ขนฺติยา สมนฺนาคโต สมฺมตฺต\-นิยามํ โอกฺกมิสฺสตีติ ฐานเมตํ วิชฺชติ, สมฺมตฺต\-นิยามํ โอกฺกมมาโน โสตาปตฺติผลํ วา สกทาคามิผลํ วา อนาคามิผลํ วา อรหตฺตํ วา สจฺฉิกริสฺสตีติ ฐานเมตํ วิชฺชติ. (2)}

\prose{โส วต ภิกฺขเว ภิกฺขุ กญฺจิ สงฺขารํ สุขโต สมนุปสฺสนฺโต อนุโลมิกาย ขนฺติยา สมนฺนาคโต ภวิสฺสตีติ เนตํ ฐานํ วิชฺชติ, อนุโลมิกาย ขนฺติยา อสมนฺนาคโต สมฺมตฺต\-นิยามํ โอกฺกมิสฺสตีติ \csromanfolio{410}เนตํ ฐานํ วิชฺชติ, สมฺมตฺต\-นิยามํ อโนกฺกมมาโน โสตาปตฺติผลํ วา สกทาคามิผลํ วา อนาคามิผลํ วา อรหตฺตํ วา สจฺฉิกริสฺสตีติ เนตํ ฐานํ วิชฺชติ. (3)}

\prose{โส วต ภิกฺขเว ภิกฺขุ สพฺพสงฺขาเร ทุกฺขโต สมนุปสฺสนฺโต อนุโลมิกาย ขนฺติยา สมนฺนาคโต ภวิสฺสตีติ ฐานเมตํ วิชฺชติ, อนุโลมิกาย ขนฺติยา สมนฺนาคโต สมฺมตฺต\-นิยามํ โอกฺกมิสฺสตีติ ฐานเมตํ วิชฺชติ, สมฺมตฺต\-นิยามํ โอกฺกมมาโน โสตาปตฺติผลํ วา สกทาคามิผลํ วา อนาคามิผลํ วา อรหตฺตํ วา สจฺฉิกริสฺสตีติ ฐานเมตํ วิชฺชติ. (4)}

\prose{โส วต ภิกฺขเว ภิกฺขุ กญฺจิ ธมฺมํ อตฺตโต สมนุปสฺสนฺโต อนุโลมิกาย ขนฺติยา สมนฺนาคโต ภวิสฺสตีติ เนตํ ฐานํ วิชฺชติ, อนุโลมิกาย ขนฺติยา อสมนฺนาคโต สมฺมตฺต\-นิยามํ โอกฺกมิสฺสตีติ เนตํ ฐานํ วิชฺชติ, สมฺมตฺต\-นิยามํ อโนกฺกมมาโน โสตาปตฺติผลํ วา สกทาคามิผลํ วา อนาคามิผลํ วา อรหตฺตํ วา สจฺฉิกริสฺสตีติ เนตํ ฐานํ วิชฺชติ. (5)}

\prose{โส วต ภิกฺขเว ภิกฺขุ สพฺพธมฺเม\csromansharedfootnote{6aff9ecab95a9e34}{ตญฺจิ ธมฺมํ (สฺยา), กิญฺจิ ธมฺมํ (ก) อํ 2. 385 ปิฏฺเฐ ปสฺสิตพฺพา.} อนตฺตโต สมนุปสฺสนฺโต อนุโลมิกาย ขนฺติยา สมนฺนาคโต ภวิสฺสตีติ ฐานเมตํ วิชฺชติ, อนุโลมิกาย ขนฺติยา สมนฺนาคโต สมฺมตฺต\-นิยามํ โอกฺกมิสฺสตีติ ฐานเมตํ วิชฺชติ, สมฺมตฺต\-นิยามํ โอกฺกมมาโน โสตาปตฺติผลํ วา สกทาคามิผลํ วา อนาคามิผลํ วา อรหตฺตํ วา สจฺฉิกริสฺสตีติ ฐานเมตํ วิชฺชติ. (6)}

\prose{โส วต ภิกฺขเว ภิกฺขุ นิพฺพานํ ทุกฺขโต สมนุปสฺสนฺโต อนุโลมิกาย ขนฺติยา สมนฺนาคโต ภวิสฺสตีติ เนตํ ฐานํ วิชฺชติ, อนุโลมิกาย ขนฺติยา อสมนฺนาคโต สมฺมตฺต\-นิยามํ โอกฺกมิสฺสตีติ เนตํ ฐานํ วิชฺชติ, สมฺมตฺต\-นิยามํ อโนกฺกมมาโน โสตาปตฺติผลํ วา สกทาคามิผลํ วา อนาคามิผลํ วา อรหตฺตํ วา สจฺฉิกริสฺสตีติ เนตํ ฐานํ วิชฺชติ. (7)}

\prose{โส วต ภิกฺขเว ภิกฺขุ นิพฺพานํ สุขโต สมนุปสฺสนฺโต อนุโลมิกาย ขนฺติยา สมนฺนาคโต ภวิสฺสตีติ ฐานเมตํ วิชฺชติ,}

\prose{\csromanfolio{411}อนุโลมิกาย ขนฺติยา สมนฺนาคโต สมฺมตฺต\-นิยามํ โอกฺกมิสฺสตีติ ฐานเมตํ วิชฺชติ, สมฺมตฺต\-นิยามํ โอกฺกมมาโน โสตาปตฺติผลํ วา สกทาคามิผลํ วา อนาคามิผลํ วา อรหตฺตํ วา สจฺฉิกริสฺสตีติ ฐานเมตํ วิชฺชติ. (8)}

\proseitem{37}{กติหากาเรหิ อนุโลมิกํ ขนฺติํ ปฏิลภติ, กติหากาเรหิ สมฺมตฺต\-นิยามํ โอกฺกมติ? จตฺตารีสาย อากาเรหิ อนุโลมิกํ ขนฺติํ ปฏิลภติ, จตฺตารีสาย อากาเรหิ สมฺมตฺต\-นิยามํ โอกฺกมติ.}

\prose{กตเมหิ จตฺตารีสาย อากาเรหิ อนุโลมิกํ ขนฺติํ ปฏิลภติ, กตเมหิ จตฺตารีสาย อากาเรหิ สมฺมตฺต\-นิยามํ โอกฺกมติ? ปญฺจกฺขนฺเธ อนิจฺจโต ทุกฺขโต โรคโต คณฺฑโต สลฺลโต อฆโต อาพาธโต ปรโต ปโลกโต อีติโต อุปทฺทวโต ภยโต อุปสคฺคโต จลโต ปภงฺคุโต\csromansharedfootnote{f6b826cbe920c19e}{ปภงฺคโต (สฺยา)} อทฺธุวโต อตาณโต\csromansharedfootnote{56276810f7ea44ce}{อตฺตาณโต (สฺยา)} อเลณโต อสรณโต ริตฺตโต ตุจฺฉโต สุญฺญโต อนตฺตโต อาทีนวโต วิปริณาม\-ธมฺมโต อสารกโต อฆมูลโต วธกโต วิภวโต สาสวโต สงฺขตโต มารามิสโต ชาติธมฺมโต ชราธมฺมโต พฺยาธิธมฺมโต มรณธมฺมโต โสกธมฺมโต ปริเทว\-ธมฺมโต อุปายา\-สธมฺมโต สํกิเลสิก\-ธมฺมโต.}

\proseitem{38}{ปญฺจกฺขนฺเธ อนิจฺจโต ปสฺสนฺโต อนุโลมิกํ ขนฺติํ ปฏิลภติ, “ปญฺจนฺนํ ขนฺธานํ นิโรโธ นิจฺจํ นิพฺพานนฺ”ติ ปสฺสนฺโต สมฺมตฺต\-นิยามํ โอกฺกมติ. ปญฺจกฺขนฺเธ ทุกฺขโต ปสฺสนฺโต อนุโลมิกํ ขนฺติํ ปฏิลภติ, “ปญฺจนฺนํ ขนฺธานํ นิโรโธ สุขํ นิพฺพานนฺ”ติ ปสฺสนฺโต สมฺมตฺต\-นิยามํ โอกฺกมติ. ปญฺจกฺขนฺเธ โรคโต ปสฺสนฺโต อนุโลมิกํ ขนฺติํ ปฏิลภติ, “ปญฺจนฺนํ ขนฺธานํ นิโรโธ อาโรคฺยํ นิพฺพานนฺ”ติ ปสฺสนฺโต สมฺมตฺต\-นิยามํ โอกฺกมติ. ปญฺจกฺขนฺเธ คณฺฑโต ปสฺสนฺโต อนุโลมิกํ ขนฺติํ ปฏิลภติ, “ปญฺจนฺนํ ขนฺธานํ นิโรโธ อคณฺฑํ\csromansharedfootnote{befabbaba5b0f20e}{นิคณฺโฑ (สฺยา)} นิพฺพานนฺ”ติ ปสฺสนฺโต สมฺมตฺต\-นิยามํ โอกฺกมติ. ปญฺจกฺขนฺเธ สลฺลโต ปสฺสนฺโต อนุโลมิกํ ขนฺติํ ปฏิลภติ, “ปญฺจนฺนํ ขนฺธานํ นิโรโธ วิสลฺลํ\csromansharedfootnote{29c6f795572be1f4}{นิสฺสลฺลํ (สฺยา)} นิพฺพานนฺ”ติ ปสฺสนฺโต สมฺมตฺต\-นิยามํ โอกฺกมติ. (5)}

\prose{\csromanfolio{412}ปญฺจกฺขนฺเธ อฆโต ปสฺสนฺโต อนุโลมิกํ ขนฺติํ ปฏิลภติ, “ปญฺจนฺนํ ขนฺธานํ นิโรโธ อนโฆ นิพฺพานนฺ”ติ ปสฺสนฺโต สมฺมตฺต\-นิยามํ โอกฺกมติ. ปญฺจกฺขนฺเธ อาพาธโต ปสฺสนฺโต อนุโลมิกํ ขนฺติํ ปฏิลภติ, “ปญฺจนฺนํ ขนฺธานํ นิโรโธ อนาพาธํ นิพฺพานนฺ”ติ ปสฺสนฺโต สมฺมตฺต\-นิยามํ โอกฺกมติ. ปญฺจกฺขนฺเธ ปรโต ปสฺสนฺโต อนุโลมิกํ ขนฺติํ ปฏิลภติ, “ปญฺจนฺนํ ขนฺธานํ นิโรโธ อปรปฺปจฺจยํ นิพฺพานนฺ”ติ ปสฺสนฺโต สมฺมตฺต\-นิยามํ โอกฺกมติ. ปญฺจกฺขนฺเธ ปโลกโต ปสฺสนฺโต อนุโลมิกํ ขนฺติํ ปฏิลภติ, “ปญฺจนฺนํ ขนฺธานํ นิโรโธ อปโลกธมฺโม นิพฺพานนฺ”ติ ปสฺสนฺโต สมฺมตฺต\-นิยามํ โอกฺกมติ. ปญฺจกฺขนฺเธ อีติโต ปสฺสนฺโต อนุโลมิกํ ขนฺติํ ปฏิลภติ, “ปญฺจนฺนํ ขนฺธานํ นิโรโธ อนีติกํ นิพฺพานนฺ”ติ ปสฺสนฺโต สมฺมตฺต\-นิยามํ โอกฺกมติ. (10)}

\prose{ปญฺจกฺขนฺเธ อุปทฺทวโต ปสฺสนฺโต อนุโลมิกํ ขนฺติํ ปฏิลภติ, “ปญฺจนฺนํ ขนฺธานํ นิโรโธ อนุปทฺทวํ นิพฺพานนฺ”ติ ปสฺสนฺโต สมฺมตฺต\-นิยามํ โอกฺกมติ. ปญฺจกฺขนฺเธ ภยโต ปสฺสนฺโต อนุโลมิกํ ขนฺติํ ปฏิลภติ, “ปญฺจนฺนํ ขนฺธานํ นิโรโธ อภยํ นิพฺพานนฺ”ติ ปสฺสนฺโต สมฺมตฺต\-นิยามํ โอกฺกมติ. ปญฺจกฺขนฺเธ อุปสคฺคโต ปสฺสนฺโต อนุโลมิกํ ขนฺติํ ปฏิลภติ, “ปญฺจนฺนํ ขนฺธานํ นิโรโธ อนุปสคฺคํ นิพฺพานนฺ”ติ ปสฺสนฺโต สมฺมตฺต\-นิยามํ โอกฺกมติ. ปญฺจกฺขนฺเธ จลโต ปสฺสนฺโต อนุโลมิกํ ขนฺติํ ปฏิลภติ, “ปญฺจนฺนํ ขนฺธานํ นิโรโธ อจลํ นิพฺพานนฺ”ติ ปสฺสนฺโต สมฺมตฺต\-นิยามํ โอกฺกมติ. ปญฺจกฺขนฺเธ ปภงฺคุโต ปสฺสนฺโต อนุโลมิกํ ขนฺติํ ปฏิลภติ, “ปญฺจนฺนํ ขนฺธานํ นิโรโธ อปภงฺคุ\csromansharedfootnote{95438e8c872b176d}{อปฺปภงฺคํ (สฺยา)} นิพฺพานนฺ”ติ ปสฺสนฺโต สมฺมตฺต\-นิยามํ โอกฺกมติ. (15)}

\prose{ปญฺจกฺขนฺเธ อทฺธุวโต ปสฺสนฺโต อนุโลมิกํ ขนฺติํ ปฏิลภติ, “ปญฺจนฺนํ ขนฺธานํ นิโรโธ ธุวํ นิพฺพานนฺ”ติ ปสฺสนฺโต สมฺมตฺต\-นิยามํ โอกฺกมติ. ปญฺจกฺขนฺเธ อตาณโต ปสฺสนฺโต อนุโลมิกํ ขนฺติํ ปฏิลภติ, “ปญฺจนฺนํ ขนฺธานํ นิโรโธ ตาณํ นิพฺพานนฺ”ติ ปสฺสนฺโต สมฺมตฺต\-นิยามํ โอกฺกมติ. ปญฺจกฺขนฺเธ อเลณโต ปสฺสนฺโต อนุโลมิกํ ขนฺติํ ปฏิลภติ, “ปญฺจนฺนํ ขนฺธานํ นิโรโธ เลณํ นิพฺพานนฺ”ติ ปสฺสนฺโต สมฺมตฺต\-นิยามํ โอกฺกมติ. ปญฺจกฺขนฺเธ อสรณโต ปสฺสนฺโต อนุโลมิกํ ขนฺติํ ปฏิลภติ, “ปญฺจนฺนํ ขนฺธานํ นิโรโธ สรณํ นิพฺพานนฺ”ติ ปสฺสนฺโต สมฺมตฺต\-นิยามํ โอกฺกมติ.}

\prose{\csromanfolio{413}ปญฺจกฺขนฺเธ ริตฺตโต ปสฺสนฺโต อนุโลมิกํ ขนฺติํ ปฏิลภติ, “ปญฺจนฺนํ ขนฺธานํ นิโรโธ อริตฺตํ นิพฺพานนฺ”ติ ปสฺสนฺโต สมฺมตฺต\-นิยามํ โอกฺกมติ. (20)}

\prose{ปญฺจกฺขนฺเธ ตุจฺฉโต ปสฺสนฺโต อนุโลมิกํ ขนฺติํ ปฏิลภติ, “ปญฺจนฺนํ ขนฺธานํ นิโรโธ อตุจฺฉํ นิพฺพานนฺ”ติ ปสฺสนฺโต สมฺมตฺต\-นิยามํ โอกฺกมติ. ปญฺจกฺขนฺเธ สุญฺญโต ปสฺสนฺโต อนุโลมิกํ ขนฺติํ ปฏิลภติ, “ปญฺจนฺนํ ขนฺธานํ นิโรโธ ปรมสุญฺญํ นิพฺพานนฺ”ติ ปสฺสนฺโต สมฺมตฺต\-นิยามํ โอกฺกมติ. ปญฺจกฺขนฺเธ อนตฺตโต ปสฺสนฺโต อนุโลมิกํ ขนฺติํ ปฏิลภติ, “ปญฺจนฺนํ ขนฺธานํ นิโรโธ ปรมตฺถํ นิพฺพานนฺ”ติ ปสฺสนฺโต สมฺมตฺต\-นิยามํ โอกฺกมติ. ปญฺจกฺขนฺเธ อาทีนวโต ปสฺสนฺโต อนุโลมิกํ ขนฺติํ ปฏิลภติ, “ปญฺจนฺนํ ขนฺธานํ นิโรโธ อนาทีนวํ นิพฺพานนฺ”ติ ปสฺสนฺโต สมฺมตฺต\-นิยามํ โอกฺกมติ. ปญฺจกฺขนฺเธ วิปริณาม\-ธมฺมโต ปสฺสนฺโต อนุโลมิกํ ขนฺติํ ปฏิลภติ, “ปญฺจนฺนํ ขนฺธานํ นิโรโธ อวิปริณาม\-ธมฺมํ นิพฺพานนฺ”ติ ปสฺสนฺโต สมฺมตฺต\-นิยามํ โอกฺกมติ. (25)}

\prose{ปญฺจกฺขนฺเธ อสารกโต ปสฺสนฺโต อนุโลมิกํ ขนฺติํ ปฏิลภติ, “ปญฺจนฺนํ ขนฺธานํ นิโรโธ สารํ นิพฺพานนฺ”ติ ปสฺสนฺโต สมฺมตฺต\-นิยามํ โอกฺกมติ. ปญฺจกฺขนฺเธ อฆมูลโต ปสฺสนฺโต อนุโลมิกํ ขนฺติํ ปฏิลภติ, “ปญฺจนฺนํ ขนฺธานํ นิโรโธ อนฆมูลํ นิพฺพานนฺ”ติ ปสฺสนฺโต สมฺมตฺต\-นิยามํ โอกฺกมติ. ปญฺจกฺขนฺเธ วธกโต ปสฺสนฺโต อนุโลมิกํ ขนฺติํ ปฏิลภติ, “ปญฺจนฺนํ ขนฺธานํ นิโรโธ อวธกํ นิพฺพานนฺ”ติ ปสฺสนฺโต สมฺมตฺต\-นิยามํ โอกฺกมติ. ปญฺจกฺขนฺเธ วิภวโต ปสฺสนฺโต อนุโลมิกํ ขนฺติํ ปฏิลภติ, “ปญฺจนฺนํ ขนฺธานํ นิโรโธ อวิภวํ นิพฺพานนฺ”ติ ปสฺสนฺโต สมฺมตฺต\-นิยามํ โอกฺกมติ. ปญฺจกฺขนฺเธ สาสวโต ปสฺสนฺโต อนุโลมิกํ ขนฺติํ ปฏิลภติ, “ปญฺจนฺนํ ขนฺธานํ นิโรโธ อนาสวํ นิพฺพานนฺ”ติ ปสฺสนฺโต สมฺมตฺต\-นิยามํ โอกฺกมติ. (30)}

\prose{ปญฺจกฺขนฺเธ สงฺขตโต ปสฺสนฺโต อนุโลมิกํ ขนฺติํ ปฏิลภติ, “ปญฺจนฺนํ ขนฺธานํ นิโรโธ อสงฺขตํ นิพฺพานนฺ”ติ ปสฺสนฺโต สมฺมตฺต\-นิยามํ โอกฺกมติ. ปญฺจกฺขนฺเธ มารามิสโต ปสฺสนฺโต อนุโลมิกํ ขนฺติํ ปฏิลภติ, “ปญฺจนฺนํ ขนฺธานํ นิโรโธ นิรามิสํ นิพฺพานนฺ”ติ ปสฺสนฺโต สมฺมตฺต\-นิยามํ โอกฺกมติ. ปญฺจกฺขนฺเธ ชาติธมฺมโต ปสฺสนฺโต อนุโลมิกํ ขนฺติํ ปฏิลภติ, “ปญฺจนฺนํ ขนฺธานํ นิโรโธ อชาตํ นิพฺพานนฺ”ติ ปสฺสนฺโต สมฺมตฺต\-นิยามํ โอกฺกมติ. ปญฺจกฺขนฺเธ ชราธมฺมโต ปสฺสนฺโต อนุโลมิกํ ขนฺติํ ปฏิลภติ, “ปญฺจนฺนํ}

\prose{\csromanfolio{414}ขนฺธานํ นิโรโธ อชรํ นิพฺพานนฺ”ติ ปสฺสนฺโต สมฺมตฺต\-นิยามํ โอกฺกมติ. ปญฺจกฺขนฺเธ พฺยาธิธมฺมโต ปสฺสนฺโต อนุโลมิกํ ขนฺติํ ปฏิลภติ, “ปญฺจนฺนํ ขนฺธานํ นิโรโธ อพฺยาธิ\csromansharedfootnote{c5effba3bb094d67}{อพฺยาธิธมฺมํ (สฺยา)} นิพฺพานนฺ”ติ ปสฺสนฺโต สมฺมตฺต\-นิยามํ โอกฺกมติ. (35)}

\prose{ปญฺจกฺขนฺเธ มรณธมฺมโต ปสฺสนฺโต อนุโลมิกํ ขนฺติํ ปฏิลภติ, “ปญฺจนฺนํ ขนฺธานํ นิโรโธ อมตํ นิพฺพานนฺ”ติ ปสฺสนฺโต สมฺมตฺต\-นิยามํ โอกฺกมติ. ปญฺจกฺขนฺเธ โสกธมฺมโต ปสฺสนฺโต อนุโลมิกํ ขนฺติํ ปฏิลภติ, “ปญฺจนฺนํ ขนฺธานํ นิโรโธ อโสกํ นิพฺพานนฺ”ติ ปสฺสนฺโต สมฺมตฺต\-นิยามํ โอกฺกมติ. ปญฺจกฺขนฺเธ ปริเทว\-ธมฺมโต ปสฺสนฺโต อนุโลมิกํ ขนฺติํ ปฏิลภติ, “ปญฺจนฺนํ ขนฺธานํ นิโรโธ อปริเทวํ นิพฺพานนฺ”ติ ปสฺสนฺโต สมฺมตฺต\-นิยามํ โอกฺกมติ. ปญฺจกฺขนฺเธ อุปายา\-สธมฺมโต ปสฺสนฺโต อนุโลมิกํ ขนฺติํ ปฏิลภติ, “ปญฺจนฺนํ ขนฺธานํ นิโรโธ อนุปายาสํ นิพฺพานนฺ”ติ ปสฺสนฺโต สมฺมตฺต\-นิยามํ โอกฺกมติ. ปญฺจกฺขนฺเธ สํกิเลสิก\-ธมฺมโต ปสฺสนฺโต อนุโลมิกํ ขนฺติํ ปฏิลภติ, “ปญฺจนฺนํ ขนฺธานํ นิโรโธ อสํกิลิฏฺฐํ นิพฺพานนฺ”ติ ปสฺสนฺโต สมฺมตฺต\-นิยามํ โอกฺกมติ. (40)}

\proseitem{39}{อนิจฺจโตติ อนิจฺจา\-นุปสฺสนา. ทุกฺขโตติ ทุกฺขา\-นุปสฺสนา. โรคโตติ ทุกฺขา\-นุปสฺสนา. คณฺฑโตติ ทุกฺขา\-นุปสฺสนา. สลฺลโตติ ทุกฺขา\-นุปสฺสนา. อฆโตติ ทุกฺขา\-นุปสฺสนา. อาพาธโตติ ทุกฺขา\-นุปสฺสนา. ปรโตติ อนตฺตา\-นุปสฺสนา. ปโลกโตติ อนิจฺจา\-นุปสฺสนา. ¢ติโตติ ทุกฺขา\-นุปสฺสนา. (10)}

\prose{อุปทฺทวโตติ ทุกฺขา\-นุปสฺสนา. ภยโตติ ทุกฺขา\-นุปสฺสนา. อุปสคฺคโตติ ทุกฺขา\-นุปสฺสนา. จลโตติ อนิจฺจา\-นุปสฺสนา. ปภงฺคุโตติ อนิจฺจา\-นุปสฺสนา. อทฺธุวโตติ อนิจฺจา\-นุปสฺสนา. อตาณโตติ ทุกฺขา\-นุปสฺสนา. อเลณโตติ ทุกฺขา\-นุปสฺสนา. อสรณโตติ ทุกฺขา\-นุปสฺสนา. ริตฺตโตติ อนตฺตา\-นุปสฺสนา. \mbox{(10, 20)}}

\prose{ตุจฺฉโตติ อนตฺตา\-นุปสฺสนา. สุญฺญโตติ อนตฺตา\-นุปสฺสนา. อนตฺตโตติ อนตฺตา\-นุปสฺสนา. อาทีนวโตติ ทุกฺขา\-นุปสฺสนา. วิปริณาม\-ธมฺมโตติ อนิจฺจา\-นุปสฺสนา. อสารกโตติ อนตฺตา\-นุปสฺสนา.}

\csromanmatikaanchor{mtk.133}
\csromantocmarkref{section}{10. มาติกากถา}{mtk.133}
\bookmark[dest={mtk.133},level=1]{10. มาติกากถา}
\csromanheader{1}{10. มาติกากถา}
\prose{\csromanfolio{415}อฆมูลโตติ ทุกฺขา\-นุปสฺสนา. วธกโตติ ทุกฺขา\-นุปสฺสนา. วิภวโตติ อนิจฺจา\-นุปสฺสนา. สาสวโตติ ทุกฺขา\-นุปสฺสนา. \mbox{(10, 30)}}

\prose{สงฺขตโตติ อนิจฺจา\-นุปสฺสนา. มารามิสโตติ ทุกฺขา\-นุปสฺสนา. ชาติธมฺมโตติ ทุกฺขา\-นุปสฺสนา. ชราธมฺมโตติ ทุกฺขา\-นุปสฺสนา. พฺยาธิ\-ธมฺมโตติ ทุกฺขา\-นุปสฺสนา. มรณ\-ธมฺมโตติ อนิจฺจา\-นุปสฺสนา. โสกธมฺมโตติ ทุกฺขา\-นุปสฺสนา. ปริเทว\-ธมฺมโตติ ทุกฺขา\-นุปสฺสนา. อุปายาสธมฺมโตติ ทุกฺขา\-นุปสฺสนา. สํกิเลสิก\-ธมฺมโตติ ทุกฺขา\-นุปสฺสนา. \mbox{(10, 40)}}

\prose{อิเมหิ จตฺตาลีสาย อากาเรหิ อนุโลมิกํ ขนฺติํ ปฏิลภติ. อิเมหิ จตฺตาลีสาย อากาเรหิ สมฺมตฺต\-นิยามํ โอกฺกมติ.}

\prose{อิเมหิ จตฺตาลีสาย อากาเรหิ อนุโลมิกํ ขนฺติํ ปฏิลภนฺตสฺส, อิเมหิ จตฺตาลีสาย อากาเรหิ สมฺมตฺต\-นิยามํ โอกฺกมนฺตสฺส กติ อนิจฺจา\-นุปสฺสนา, กติ ทุกฺขา\-นุปสฺสนา, กติ อนตฺตา\-นุปสฺสนา.}

\setlength{\csromangathapagewidth}{0pt}
\setlength{\csromangathawakwidth}{0pt}
\csromangathameasuregroup{}{ปญฺจวีสติ อนตฺตานุปสฺสนา, \\ ปญฺญาส อนิจฺจานุปสฺสนา. \\ สตํ ปญฺจวีสติ เจว, \\ ยานิ ทุกฺเข ปวุจฺจเรติ.}
\csromangathameasurewak{ปญฺจวีสติ อนตฺตานุปสฺสนา, \\ ปญฺญาส อนิจฺจานุปสฺสนา. \\ สตํ ปญฺจวีสติ เจว, \\ ยานิ ทุกฺเข ปวุจฺจเรติ.}
\setlength{\csromangathaleftcol}{0pt}
\csromangathagroupwithcloserruled{\csromangathastanza{\csromangathawak{ปญฺจวีสติ อนตฺตา\-นุปสฺสนา, \\ ปญฺญาส อนิจฺจา\-นุปสฺสนา. \\ สตํ ปญฺจวีสติ เจว, \\ ยานิ ทุกฺเข ปวุจฺจเรติ.}}}{วิปสฺสนากถา นิฏฺฐิตา.}
\csromanheader{2}{10. มาติกากถา}
\proseitem{40}{นิจฺฉาโต\csromansharedfootnote{115f9162e12c43ab}{นิจฺฉาโต มุจฺจตีติ (สฺยา)}, โมกฺโข วิโมกฺโข, วิชฺชาวิมุตฺติ, อธิสีลํ, อธิจิตฺตํ, อธิปญฺญา, ปสฺสทฺธิ, ญาณํ, ทสฺสนํ, วิสุทฺธิ\csromansharedfootnote{8896f1df55785188}{สุทฺธิ (สฺยา)}, เนกฺขมฺมํ, นิสฺสรณํ, ปวิเวโก, โวสคฺโค, จริยา, ฌานวิโมกฺโข, ภาวนา, อธิฏฺฐานํ ชีวิตํ. (19)}

\proseitem{41}{\csromanfolio{416}นิจฺฉาโตติ เนกฺขมฺเมน กามจฺฉนฺทโต นิจฺฉาโต, อพฺยาปาเทน พฺยาปาทโต นิจฺฉาโต ฯเปฯ ปฐเมน ฌาเนน นีวรเณหิ นิจฺฉาโต ฯเปฯ อรหตฺต\-มคฺเคน สพฺพกิเลเสหิ นิจฺฉาโต. (1)}

\prose{โมกฺโข วิโมกฺโขติ เนกฺขมฺเมน กามจฺฉนฺทโต มุจฺจตีติ โมกฺโข วิโมกฺโข, อพฺยาปาเทน พฺยาปาทโต มุจฺจตีติ โมกฺโข วิโมกฺโข ฯเปฯ ปฐเมน ฌาเนน นีวรเณหิ มุจฺจตีติ โมกฺโข วิโมกฺโข ฯเปฯ อรหตฺต\-มคฺเคน สพฺพกิเลเสหิ มุจฺจตีติ โมกฺโข วิโมกฺโข. (2)}

\prose{\textbf{วิชฺชาวิมุตฺตี}ติ เนกฺขมฺมํ วิชฺชตีติ วิชฺชา, กามจฺฉนฺทโต มุจฺจตีติ วิมุตฺติ. วิชฺชนฺโต มุจฺจติ, มุจฺจนฺโต วิชฺชตีติ วิชฺชาวิมุตฺติ. อพฺยาปาโท\csromansharedfootnote{4e96cf3421f4d135}{อพฺยาปาทํ (สฺยา)} วิชฺชตีติ วิชฺชา, พฺยาปาทโต วิมุจฺจตีติ วิมุตฺติ. วิชฺชนฺโต มุจฺจติ, มุจฺจนฺโต วิชฺชตีติ วิชฺชาวิมุตฺติ ฯเปฯ อรหตฺตมคฺโค วิชฺชตีติ วิชฺชา, สพฺพกิเลเสหิ มุจฺจตีติ วิมุตฺติ. วิชฺชนฺโต มุจฺจติ, มุจฺจนฺโต วิชฺชตีติ วิชฺชาวิมุตฺติ. (3)}

\prose{อธิสีลํ อธิจิตฺตํ อธิปญฺญาติ เนกฺขมฺเมน กามจฺฉนฺทํ สํวรฏฺเฐน สีลวิสุทฺธิ, อวิกฺเขป\-ฏฺเฐน จิตฺตวิสุทฺธิ, ทสฺสนฏฺเฐน ทิฏฺฐิวิสุทฺธิ. โย ตตฺถ สํวรฏฺโฐ, อยํ อธิสีลสิกฺขา. โย ตตฺถ อวิกฺเขปฏฺโฐ, อยํ อธิจิตฺต\-สิกฺขา. โย ตตฺถ ทสฺสนฏฺโฐ, อยํ อธิปญฺญา\-สิกฺขา. อพฺยาปาเทน พฺยาปาทํ สํวรฏฺเฐน สีลวิสุทฺธิ ฯเปฯ อรหตฺต\-มคฺเคน สพฺพกิเลเส สํวรฏฺเฐน สีลวิสุทฺธิ, อวิกฺเขป\-ฏฺเฐน จิตฺตวิสุทฺธิ. ทสฺสนฏฺเฐน ทิฏฺฐิวิสุทฺธิ. โย ตตฺถ สํวรฏฺโฐ, อยํ อธิสีลสิกฺขา. โย ตตฺถ อวิกฺเขปฏฺโฐ, อยํ อธิจิตฺต\-สิกฺขา. โย ตตฺถ ทสฺสนฏฺโฐ, อยํ อธิปญฺญา\-สิกฺขา. \mbox{(3, 6)}}

\prose{ปสฺสทฺธีติ เนกฺขมฺเมน กามจฺฉนฺทํ ปฏิปฺปสฺสมฺเภติ, อพฺยาปาเทน พฺยาปาทํ ปฏิปฺปสฺสมฺเภติ ฯเปฯ อรหตฺต\-มคฺเคน สพฺพกิเลเส ปฏิปฺปสฺสมฺเภติ. (7)}

\prose{\textbf{ญาณนฺ}ติ กามจฺฉนฺทสฺส ปหีนตฺตา เนกฺขมฺมํ ญาตฏฺเฐน ญาณํ, พฺยาปาทสฺส ปหีนตฺตา อพฺยาปาโท ญาตฏฺเฐน ญาณํ ฯเปฯ สพฺพกิเลสานํ ปหีนตฺตา อรหตฺตมคฺโค ญาตฏฺเฐน ญาณํ. (8)}

\prose{\textbf{ทสฺสนนฺ}ติ กามจฺฉนฺทสฺส ปหีนตฺตา เนกฺขมฺมํ ทิฏฺฐตฺตา ทสฺสนํ, พฺยาปาทสฺส ปหีนตฺตา อพฺยาปาโท ทิฏฺฐิตฺตา ทสฺสนํ ฯเปฯ สพฺพกิเลสานํ ปหีนตฺตา อรหตฺตมคฺโค ทิฏฺฐตฺตา ทสฺสนํ. (9)}

\csromanheader{2}{10. มาติกากถา}
\prose{\csromanfolio{417}วิสุทฺธีติ กามจฺฉนฺทํ ปชหนฺโต เนกฺขมฺเมน วิสุชฺฌติ, พฺยาปาทํ ปชหนฺโต อพฺยาปาเทน วิสุชฺฌติ ฯเปฯ สพฺพกิเลเส ปชหนฺโต อรหตฺต\-มคฺเคน วิสุชฺฌติ. (10)}

\prose{เนกฺขมฺมนฺติ กามานเมตํ นิสฺสรณํ, ยทิทํ เนกฺขมฺมํ. รูปานเมตํ นิสฺสรณํ, ยทิทํ อารุปฺปํ. ยํ โข ปน กิญฺจิ ภูตํ สงฺขตํ ปฏิจฺจ\-สมุปฺปนฺนํ, นิโรโธ ตสฺส เนกฺขมฺมํ. พฺยาปาทสฺส อพฺยาปาโท เนกฺขมฺมํ. ถินมิทฺธสฺส อาโลกสญฺญา เนกฺขมฺมํ ฯเปฯ สพฺพกิเลสานํ อรหตฺตมคฺโค เนกฺขมฺมํ. (11)}

\prose{นิสฺสรณนฺติ กามานเมตํ นิสฺสรณํ, ยทิทํ เนกฺขมฺมํ. รูปานเมตํ นิสฺสรณํ, ยทิทํ อารุปฺปํ. ยํ โข ปน กิญฺจิ ภูตํ สงฺขตํ ปฏิจฺจ\-สมุปฺปนฺนํ, นิโรโธ ตสฺส นิสฺสรณํ. กามจฺฉนฺทสฺส เนกฺขมฺมํ นิสฺสรณํ, พฺยาปาทสฺส อพฺยาปาโท นิสฺสรณํ ฯเปฯ สพฺพกิเลสานํ อรหตฺตมคฺโค นิสฺสรณํ. (12)}

\prose{\textbf{ปวิเวโก}ติ กามจฺฉนฺทสฺส เนกฺขมฺมํ ปวิเวโก ฯเปฯ สพฺพกิเลสานํ อรหตฺตมคฺโค ปวิเวโก. (13)}

\prose{โวสคฺโคติ เนกฺขมฺเมน กามจฺฉนฺทํ โวสชฺชตีติ โวสคฺโค, อพฺยาปาเทน พฺยาปาทํ โวสชฺชตีติ โวสคฺโค ฯเปฯ อรหตฺต\-มคฺเคน สพฺพกิเลเส โวสชฺชตีติ โวสคฺโค. (14)}

\prose{จริยาติ กามจฺฉนฺทํ ปชหนฺโต เนกฺขมฺเมน จรติ, พฺยาปาทํ ปชหนฺโต อพฺยาปาเทน จรติ ฯเปฯ สพฺพกิเลเส ปชหนฺโต อรหตฺต\-มคฺเคน จรติ. (15)}

\prose{\textbf{ฌานวิโมกฺโข}ติ เนกฺขมฺมํ ฌายตีติ ฌานํ, กามจฺฉนฺทํ ฌาเปตีติ ฌานํ, ฌายนฺโต มุจฺจตีติ ฌานวิโมกฺโข, ฌาเปนฺโต มุจฺจตีติ ฌานวิโมกฺโข, ฌายนฺตีติ ธมฺมา, ฌาเปนฺตีติ กิเลเส, ฌาเต จ ฌาเป จ ชานาตีติ ฌานฌายี\csromansharedfootnote{1bceed7e6b67cae0}{ฌานวิโมกฺโข (สฺยา) อฏฺฐกถา โอโลเกตพฺพา.}. อพฺยาปาโท ฌายตีติ ฌานํ, พฺยาปาทํ ฌาเปตีติ ฌานํ ฯเปฯ อาโลกสญฺญา ฌายตีติ ฌานํ, ถินมิทฺธํ ฌาเปตีติ ฌานํ ฯเปฯ}

\prose{\csromanfolio{418}อรหตฺตมคฺโค ฌายตีติ ฌานํ, สพฺพกิเลเส ฌาเปตีติ ฌานํ, ฌายนฺโต มุจฺจตีติ ฌานวิโมกฺโข, ฌาเปนฺโต มุจฺจตีติ ฌานวิโมกฺโข, ฌายนฺตีติ ธมฺมา, ฌาเปนฺตีติ กิเลเส, ฌาเต จ ฌาเป จ ชานาตีติ ฌานฌายี. (16)}

\proseitem{42}{ภาวนา อธิฏฺฐานํ ชีวิตนฺติ กามจฺฉนฺทํ ปชหนฺโต เนกฺขมฺมํ ภาเวตีติ ภาวนา\-สมฺปนฺโน, เนกฺขมฺม\-วเสน จิตฺตํ อธิฏฺฐาตีติ อธิฏฺฐาน\-สมฺปนฺโน, สฺวายํ เอวํ ภาวนา\-สมฺปนฺโน อธิฏฺฐาน\-สมฺปนฺโน สมํ ชีวติ, โน วิสมํ. สมฺมา ชีวติ, โน มิจฺฉา. วิสุทฺธํ ชีวติ, โน กิลิฏฺฐนฺติ อาชีวสมฺปนฺโน. สฺวายํ เอวํ ภาวนา\-สมฺปนฺโน อธิฏฺฐาน\-สมฺปนฺโน อาชีวสมฺปนฺโน ยญฺญเทว ปริสํ อุปสงฺกมติ ยทิ ขตฺติย\-ปริสํ ยทิ พฺราหฺมณปริยํ ยทิ คหปติ\-ปริสํ ยทิ สมณปริสํ วิสารโท อุปสงฺกมติ อมงฺกุภูโต, ตํ กิสฺส เหตุ, ตถา หิ โส ภาวนา\-สมฺปนฺโน อธิฏฺฐาน\-สมฺปนฺโน อาชีวสมฺปนฺโน.}

\prosewithcloser{พฺยาปาทํ ปชหนฺโต อพฺยาปาทํ ภาเวตีติ ภาวนา\-สมฺปนฺโน ฯเปฯ. ถินมิทฺธํ ปชหนฺโต อาโลกสญฺญํ ภาเวตีติ ภาวนา\-สมฺปนฺโน ฯเปฯ. อุทฺธจฺจํ ปชหนฺโต อวิกฺเขปํ ภาเวตีติ ภาวนา\-สมฺปนฺโน ฯเปฯ. วิจิกิจฺฉํ ปชหนฺโต ธมฺม\-ววตฺถานํ ภาเวตีติ ภาวนา\-สมฺปนฺโน ฯเปฯ. อวิชฺชํ ปชหนฺโต วิชฺชํ ภาเวตีติ ภาวนา\-สมฺปนฺโน ฯเปฯ. อรติํ ปชหนฺโต ปาโมชฺชํ ภาเวตีติ ภาวนา\-สมฺปนฺโน ฯเปฯ. นีวรเณ ปชหนฺโต ปฐมํ ฌานํ ภาเวตีติ ภาวนา\-สมฺปนฺโน ฯเปฯ. สพฺพกิเลเส ปชหนฺโต อรหตฺตมคฺคํ ภาเวตีติ ภาวนา\-สมฺปนฺโน. อรหตฺตมคฺค\-วเสน จิตฺตํ อธิฏฺฐาตีติ อธิฏฺฐาน\-สมฺปนฺโน. สฺวายํ เอวํ ภาวนา\-สมฺปนฺโน อธิฏฺฐาน\-สมฺปนฺโน สมํ ชีวติ, โน วิสมํ. สมฺมา ชีวติ, โน มิจฺฉา. วิสุทฺธํ ชีวติ, โน กิลิฏฺฐนฺติ อาชีวสมฺปนฺโน. สฺวายํ เอวํ ภาวนา\-สมฺปนฺโน อธิฏฺฐาน\-สมฺปนฺโน อาชีวสมฺปนฺโน ยญฺญเทว ปริสํ อุปสงฺกมติ ยทิ ขตฺติย\-ปริสํ ยทิ พฺราหฺมณ\-ปริสํ ยทิ คหปติ\-ปริสํ ยทิ สมณปริสํ วิสารโท อุปสงฺกมติ อมงฺกุภูโต, ตํ กิสฺส เหตุ, ตถา หิ โส ภาวนา\-สมฺปนฺโน อธิฏฺฐาน\-สมฺปนฺโน อาชีวสมฺปนฺโนติ. \mbox{(3, 19)}}{มาติกากถา นิฏฺฐิตา.}
\nitthitammidruled{ปญฺญาวคฺโค ตติโย.}
\csromanheader{2}{อุทฺทานคาถา}
\vspace{\tipitakaparskip}
\csromancenter{\textbf{ตสฺสุทฺทานํ}}
\setlength{\csromangathapagewidth}{0pt}
\setlength{\csromangathawakwidth}{0pt}
\csromangathameasuregroup{ปญฺญา อิทฺธิ อภิสมโย, \\ ปาฏิหาริ สมสีสิ, \\ ตติเย ปญฺญาวคฺคมฺหิ, \\ มหาวคฺโค ยุคนทฺโธ, \\ ตโยว วคฺคา อิมมฺหิ\textsuperscript{1}, \\ อนนฺตนยมคฺเคสุ, \\ นภํ จ ตารกากิณฺณํ, \\ กถิกานํ วิสาลาย,}{\csromangathabat{ปญฺญา อิทฺธิ อภิสมโย,}{วิเวโก จริยปญฺจโม.} \\ \csromangathabat{ปาฏิหาริ สมสีสิ,}{สติปฏฺฐานา วิปสฺสนา.} \\ \csromangathabat{ตติเย ปญฺญาวคฺคมฺหิ,}{มาติกาย จ เต ทสาติ.} \\ \csromangathabat{มหาวคฺโค ยุคนทฺโธ,}{ปญฺญาวคฺโค จ นามโต.} \\ \csromangathabat{ตโยว วคฺคา อิมมฺหิ\textsuperscript{1},}{ปฏิสมฺภิทาปกรเณ.} \\ \csromangathabat{อนนฺตนยมคฺเคสุ,}{คมฺภีโร สาครูปโม.} \\ \csromangathabat{นภํ จ ตารกากิณฺณํ,}{ถูโล ชาตสฺสโร ยถา.} \\ \csromangathabat{กถิกานํ วิสาลาย,}{โยคีนํ ญาณโชตนนฺติ.}}
\csromangathasetleft{ปญฺญา อิทฺธิ อภิสมโย, \\ ปาฏิหาริ สมสีสิ, \\ ตติเย ปญฺญาวคฺคมฺหิ, \\ มหาวคฺโค ยุคนทฺโธ, \\ ตโยว วคฺคา อิมมฺหิ\textsuperscript{1}, \\ อนนฺตนยมคฺเคสุ, \\ นภํ จ ตารกากิณฺณํ, \\ กถิกานํ วิสาลาย,}
\csromangathagroupwithclosermajor{\csromangathastanza{\csromanfolio{419}\csromangathabat{ปญฺญา อิทฺธิ อภิสมโย,}{วิเวโก จริยปญฺจโม.} \\ \csromangathabat{ปาฏิหาริ สมสีสิ,}{สติปฏฺฐานา วิปสฺสนา.}}\csromangathastanza{\csromangathabat{ตติเย ปญฺญาวคฺคมฺหิ,}{มาติกาย จ เต ทสาติ.} \\ \csromangathabat{มหาวคฺโค ยุคนทฺโธ,}{ปญฺญาวคฺโค จ นามโต.}}\csromangathastanza{\csromangathabat{ตโยว วคฺคา อิมมฺหิ\csromansharedfootnote{3a130e66abcae878}{ติวคฺโค ยสฺส วิกฺเขโป (สฺยา, ก)},}{ปฏิสมฺภิทา\-ปกรเณ.} \\ \csromangathabat{อนนฺตนยมคฺเคสุ,}{คมฺภีโร สาครูปโม.}}\csromangathastanza{\csromangathabat{นภํ จ ตารกากิณฺณํ,}{ถูโล ชาตสฺสโร ยถา.} \\ \csromangathabat{กถิกานํ วิสาลาย,}{โยคีนํ ญาณโชตนนฺติ.}}}{ปฏิสมฺภิทามคฺค\-ปาฬิ นิฏฺฐิตา.}
\orphannote{อาวชฺชน\-ปฏิพนฺธา (ก) เอวมีทิเสสุ ปเทสุ.}

